\documentclass[11pt,a4paper]{article}
\usepackage[T1]{fontenc}
\usepackage[utf8]{inputenc}
\usepackage{lmodern,microtype}
\usepackage[a4paper,margin=22mm,headheight=15pt]{geometry}
\usepackage{amsmath,amssymb,amsthm,mathtools}
\usepackage{booktabs,longtable,array,tabularx,enumitem,needspace}
\usepackage{xcolor,graphicx,fvextra,xurl,multicol,etoolbox}
\usepackage{hyperref,fancyhdr,lastpage}
\definecolor{ink}{HTML}{142D42}
\definecolor{accent}{HTML}{007A78}
\hypersetup{colorlinks=true,linkcolor=ink,urlcolor=accent,citecolor=ink,
 bookmarksnumbered=true,bookmarksdepth=3,
 pdftitle={OEIS Conjecture Hunt: Collected Proofs, Counterexamples, and Corrections, Rounds I--X},
 pdfauthor={KEIO}}
\newcommand{\oeis}[1]{\href{https://oeis.org/#1}{\texttt{#1}}}
\newcommand{\N}{\mathbb{N}}
\newtheorem{theorem}{Theorem}[section]
\newtheorem{lemma}[theorem]{Lemma}
\newtheorem{proposition}[theorem]{Proposition}
\newtheorem{corollary}[theorem]{Corollary}
\theoremstyle{definition}
\newtheorem{definition}[theorem]{Definition}
\newtheorem{remark}[theorem]{Remark}
\newtheorem{example}[theorem]{Example}
\setlength{\parindent}{0pt}
\setlength{\parskip}{5pt plus 1pt minus 1pt}
\setlength{\emergencystretch}{3em}
\setlist{nosep,leftmargin=1.5em}
\setcounter{tocdepth}{1}
\makeatletter
\patchcmd{\l@section}{1.5em}{3em}{}{}
\makeatother
\renewcommand{\thesection}{\arabic{part}.\arabic{section}}
\renewcommand{\theequation}{\arabic{part}.\arabic{equation}}
\renewcommand{\thetable}{\arabic{part}.\arabic{table}}
\renewcommand{\thefigure}{\arabic{part}.\arabic{figure}}
\renewcommand{\theHsection}{round.\arabic{part}.section.\arabic{section}}
\renewcommand{\theHequation}{round.\arabic{part}.equation.\arabic{equation}}
\renewcommand{\theHtheorem}{round.\arabic{part}.section.\arabic{section}.theorem.\arabic{theorem}}
\renewcommand{\theHtable}{round.\arabic{part}.table.\arabic{table}}
\renewcommand{\theHfigure}{round.\arabic{part}.figure.\arabic{figure}}
\pagestyle{fancy}\fancyhf{}
\lhead{\small\textcolor{ink}{OEIS Conjecture Hunt: Collected Proofs}}
\rhead{\small KEIO}
\cfoot{\small\thepage\ / \pageref{LastPage}}
\renewcommand{\headrulewidth}{0.3pt}
\allowdisplaybreaks
\begin{document}
\begin{center}
\vspace*{12mm}
{\Large\textcolor{accent}{CONSOLIDATED RESEARCH MANUSCRIPT}\par}
\vspace{10mm}
{\fontsize{29}{34}\selectfont\bfseries\textcolor{ink}{OEIS Conjecture Hunt}\par}
\vspace{5mm}
{\LARGE Collected Proofs, Counterexamples,\\and Corrections\par}
\vspace{5mm}
{\Large Rounds I--X\par}
\vspace{9mm}
{\Large KEIO\par}
\vspace{5mm}
{\large Consolidated edition: 21 October 2026\par}
{\normalsize Source manuscripts: 20 October 2026\par}
\end{center}
\vspace{5mm}
\begin{abstract}
This single English manuscript consolidates the complete mathematical
text of ten OEIS conjecture-hunt dossiers. Their historical selection
ledger contains 1,008 assertion packages: 107 in Round I, 101 in Round II,
and 100 in each of Rounds III--X. The subjects include triangular and
Hankel matrices, character sums, divisor identities, congruences,
formal power series, digital sequences, floor identities, asymptotics,
differential recurrences, and finite-state enumeration. The constituent
proofs, exact counterexamples, corrected statements, coefficient lists,
printed certificates, source attributions, and qualifications are retained.
One title, one contents list, continuous pagination, round-qualified
numbering, internal links, and an OEIS identifier locator provide a
common structure. The selection total is a project ledger, not a count
of worldwide-first discoveries or OEIS editorial closures. In particular,
the retrospective priority correction recorded in Round V for five
Round IV entries remains explicit. This is a consolidation of the
supplied manuscripts, not a new mathematical validation or priority audit.
\end{abstract}
\vfill
\begin{center}\small
A standalone pdfLaTeX source with all manuscript text embedded.\\
Original companion verification archives remain separate supporting records.
\end{center}
\clearpage
\section*{Editorial guide to the consolidated edition}
\addcontentsline{toc}{section}{Editorial guide to the consolidated edition}
\label{combined:guide}
The ten supplied manuscripts are reproduced in their original round order.
Their mathematical bodies have not been shortened: all written proofs,
recurrence signatures, tables, counterexamples, boundary qualifications,
and source discussions remain. The original separate cover pages and
contents lists have been replaced by the common front matter. Cover-page
outcome counts are collected below. Local definitions and reference lists
are retained with the arguments they support, even when they repeat
material from an earlier round.

\begin{center}\small
\renewcommand{\arraystretch}{1.18}
\begin{tabularx}{\textwidth}{@{}l r >{\raggedright\arraybackslash}X@{}}
\toprule
Round & Selected & Outcomes reported by the source manuscript\\
\midrule
I &107&104 proof-only; 1 mixed proof/refutation; 2 refutation-only.\\
II&101&100 proof-only; 1 mixed proof/refutation.\\
III&100&95 proof-only; 4 refutations; 1 mixed package.\\
IV&100&97 proof-only; 2 refutations; 1 boundary correction with a proved tail.\\
V&100&90 original claims proved; 2 refutations; 8 corrected statements proved.\\
VI&100&94 original claims proved; 5 recurrence-boundary corrections; 1 period-index correction.\\
VII&100&91 selected claims proved; 4 recurrence-boundary corrections; 4 other corrections; 1 asymptotic refutation.\\
VIII&100&98 selected claims proved; 1 support correction; 1 classification refutation.\\
IX&100&96 selected assertions proved; 4 source formulas corrected and proved.\\
X&100&97 selected assertions or questions proved; 3 source formulas corrected and proved.\\
\midrule
Total&\textbf{1,008}&Historical selection ledger; the categories retain each source's conventions.\\
\bottomrule
\end{tabularx}
\end{center}

\paragraph{Priority correction carried forward.}
Round V explicitly withdraws the earlier new-resolution assessment for
\oeis{A176605}, \oeis{A176607}, \oeis{A176610}, \oeis{A176611}, and
\oeis{A176612}: Georg Fischer's 26 January 2020 family derivation already
implies the selected recurrences. The Round IV arguments are retained,
together with the later correction. The 1,008 figure must therefore be
read as a selection total. Uncounted repairs, supplements, exclusions,
and previously known inputs retain the status assigned in the originals.

\paragraph{Numbering and navigation.}
Part $r$ is Round $r$. Numbered sections and automatic equation numbers
carry the round prefix; theorem numbers inherit the section prefix.
For example, Section 4.3 is original Round IV, Section 3. Descriptive
manual equation tags are kept as printed in the source and are local
to their round. All reference links point to the corresponding location
in this combined document, and all page references use continuous
pagination. The contents lists the main sections; PDF bookmarks retain
the finer hierarchy. The final identifier locator lists every mentioned
OEIS identifier, including uncounted and excluded entries, with links
to the rounds in which it occurs.

\paragraph{Printed mathematics and companion files.}
This \texttt{.tex} file compiles independently: it does not read any
external proof fragments, images, bibliography databases, or data files.
It includes all material printed in the ten supplied manuscripts.
Some original arguments additionally refer to expanded machine-readable
matrices, coefficient vectors, audit records, and verification programs
in the corresponding verification ZIP archives. Those payloads are not
newly transcribed into this manuscript. References to a round's original
filename, builder, README, or archive describe those preserved supporting
records. Standalone compilation should not be confused with embedding
every file from the original verification archives.

\paragraph{Editorial scope and provenance.}
Source dates, attributions, and statements about earlier audits remain
historical statements from the original dossiers. Consolidation does not
update OEIS status, certify novelty, supply missing external certificates,
or convert an internal check into external peer review. The edition was
prepared with AI assistance at KEIO's request. The source TeX files were
checked against the SHA-256 index in the supplied Rounds 01--10 archive
before assembly. The consolidation check compares the retained bodies,
mathematical environments, and references against those sources.
\clearpage
\tableofcontents
\clearpage

% ===== BEGIN ORIGINAL ROUND 1 =====
\clearpage
\setcounter{section}{0}\setcounter{equation}{0}\setcounter{theorem}{0}
\setcounter{table}{0}\setcounter{figure}{0}
\normalsize\renewcommand{\arraystretch}{1}\renewcommand{\baselinestretch}{1}\selectfont
\part[Round 1: Structural identities, arithmetic, floors, and generating functions]{Structural identities, arithmetic, floors, and generating functions}
\label{combined:round1}
\rhead{\small Round 1 / KEIO}
\section*{How to read and use this dossier}
\addcontentsline{toc}{section}{How to read and use this dossier}
The original goal was to settle about one hundred accessible OEIS problems.
The resulting packet contains 107 selected conjecture or empirical
assertion groups, each tied to a named entry and a specific mathematical
statement. The proofs resolve the specified statements for their full
stated ranges. Explicit starting indices and finite exceptions matter:
none are silently discarded.

The catalogue uses \textbf{P} for a proof, \textbf{R} for a refutation,
and \textbf{P/R} where one conjecture is proved and another is refuted in
the same entry. A394578 is the latter case: its recurrence is valid, while
its proposed ratio limit is the reciprocal of the actual value. Four
standalone index or definition repairs are presented separately and are
excluded from the headline total.

\paragraph{Counting convention.}
We use an entry-level assertion group, not an arbitrary number of corollaries.
All three A329649 conjectures count as one selection. Fixed-base or fixed-block
conjectures on different entries receive separate selections, even when one
parameterized theorem proves them together. The universal square assertion
copied across A113249--A113256 counts once; the A127751 copy of A127752,
A334110 copy of A334109, and A224511 consequence of A224515 add no count.
A284005 and A329369 share a transform identity, but A329369 also has a
separately proved Eulerian-sum assertion.

\paragraph{What the source audit establishes.}
The targeted statements were still conjectural or empirical in the inspected
OEIS text, subject to the source qualifications explicitly given below.
Known earlier solutions discovered in the entries or their linked documents
were excluded. An additional audit compared all 100 claims in Appendix A of
Tom Adamczewski's \emph{OEIS Open} with the selected entries. That appendix
contains the 100 costliest of 153 combined-model resolutions. Its only
overlapping entry was A224515, where the existing result concerns existence;
our selection is the separately stated uniqueness claim. The repository
metadata could not be inspected completely, so the audit does not establish
worldwide priority. See \url{https://arxiv.org/abs/2608.11941} and the audit
JSON accompanying this dossier.

\paragraph{What makes these proofs, rather than term fitting.}
The structural sections derive identities from definitions, with explicit
inductions or counting arguments. The floor sections bound remainders for
every index beyond a finite threshold using rational inequalities and
positive polynomial coefficients. The recurrence section applies exact
differential operators to known generating functions and reduces the result
to zero or to a finite boundary polynomial. The companion programs check
these identities with integer or rational arithmetic. Finite enumeration
is supplementary verification of indexing and construction; it is never
the reason an infinite assertion is accepted.

\paragraph{Scope of the supporting formulas.}
Where an established generating function or a previously proved recurrence
is used as a premise, the text identifies it. For example, A054120 uses
Hendel's already recorded proof of its first conjecture to prove its second.
The recurrence of A174193 is addressed from its defining generating function;
a separate tentative combinatorial interpretation is outside that selection.
The harmonic-block work covers the six specified entries, not every
polynomial in the broader A227012 discussion.

\paragraph{Companion files.}
The standalone \texttt{oeis\_hunt\_dossier.tex} contains the complete report.
The verification bundle contains the detailed ledger, per-entry mathematical
certificates, proof fragments, reproducible Python programs, and recorded
verification results. The JSON ledger preserves the exact selected statement
and its qualifications. OEIS page edits and submissions were not performed.

\clearpage
\section{Catalogue of the selected assertions}
\label{round1:sec:catalogue}
Within each group, entry identifiers link directly to the corresponding
OEIS source. The subsequent proof sections supply the complete arguments.
\Needspace{6\baselineskip}\subsection*{Structural and matrix (14)}
\begin{longtable}{@{}r l c p{0.62\textwidth}@{}}
\toprule No.&Entry&Result&Selected assertion\\\midrule\endhead
1&\oeis{A329649}&P&Nonvanishing, divisibility, and the exact quadratic L-value\\
2&\oeis{A355401}&P&The q-power triangle, first column, and inverse\\
3&\oeis{A355282}&P&The q-integer triangle and first column\\
4&\oeis{A350557}&P&Signed inverse identification\\
5&\oeis{A136097}&P&Gaussian central-entry divisibility for general q\\
6&\oeis{A146955}&P&Closed formula for the row sum\\
7&\oeis{A127752}&P&Inverse row sums have Feigenbaum parity\\
8&\oeis{A081199}&P&Collinear 4-subsets of a discrete cube\\
9&\oeis{A081200}&P&Collinear 5-subsets of a discrete cube\\
10&\oeis{A081201}&P&Collinear 6-subsets of a discrete cube\\
11&\oeis{A054120}&P&The second McGarvey identity and bivariate series\\
12&\oeis{A078121}&P&Binary-partition formula for matrix entries\\
13&\oeis{A078122}&P&Ternary-partition formulas for entries and row sums\\
14&\oeis{A113249}&P&Universal odd-index perfect-square assertion\\
\bottomrule\end{longtable}
\Needspace{6\baselineskip}\subsection*{Arithmetic and digital (19)}
\begin{longtable}{@{}r l c p{0.62\textwidth}@{}}
\toprule No.&Entry&Result&Selected assertion\\\midrule\endhead
15&\oeis{A049803}&P&Base-3 remainder recurrence and series\\
16&\oeis{A049804}&P&Base-4 remainder recurrence and series\\
17&\oeis{A330358}&P&Base-5 remainder recurrence and series\\
18&\oeis{A323689}&P&Integer coefficients; odd exactly at squares\\
19&\oeis{A324300}&P&Integer coefficients; odd exactly at k(k+2)\\
20&\oeis{A362560}&P&Divisibility is invariant under partition reversal\\
21&\oeis{A085765}&P&Growth exceeds every fixed power\\
22&\oeis{A100311}&P&The binary XOR recurrence\\
23&\oeis{A156834}&P&Prime-detecting inequality\\
24&\oeis{A156833}&P&Prime-detecting Mobius-sum inequality\\
25&\oeis{A008578}&P&The divisor-sum characterization\\
26&\oeis{A006068}&P&Inverse Gray-code recurrence\\
27&\oeis{A016789}&P&Negabinary XOR identity\\
28&\oeis{A284005}&P&Three binary-mask identities\\
29&\oeis{A329369}&P&Inverse mask transform and Eulerian sums\\
30&\oeis{A334109}&P&Exact first-occurrence minimization\\
31&\oeis{A216971}&P&Every entry in row n is divisible by n\\
32&\oeis{A334592}&P&Strict zero-count records in XOR triangles\\
33&\oeis{A224515}&P&Uniqueness for the XOR of consecutive squares\\
\bottomrule\end{longtable}
\Needspace{6\baselineskip}\subsection*{Floors and symmetric functions (23)}
\begin{longtable}{@{}r l c p{0.62\textwidth}@{}}
\toprule No.&Entry&Result&Selected assertion\\\midrule\endhead
34&\oeis{A203150}&P&Empirical periodic symmetric-function series\\
35&\oeis{A203164}&P&Periodic symmetric-function series and recurrence\\
36&\oeis{A203165}&P&Periodic symmetric-function series and recurrence\\
37&\oeis{A203166}&P&Periodic symmetric-function series and recurrence\\
38&\oeis{A203167}&P&Periodic symmetric-function series and recurrence\\
39&\oeis{A203231}&P&Periodic symmetric-function series and recurrence\\
40&\oeis{A203234}&P&Periodic symmetric-function series and recurrence\\
41&\oeis{A203235}&P&Periodic symmetric-function series and recurrence\\
42&\oeis{A203157}&P&Triangular-number symmetric-function recurrence\\
43&\oeis{A203159}&P&Even-integer symmetric-function recurrence\\
44&\oeis{A024220}&P&Arithmetic-progression floor series and recurrence\\
45&\oeis{A024385}&P&Arithmetic-progression floor series and recurrence\\
46&\oeis{A024399}&P&Arithmetic-progression floor series and recurrence\\
47&\oeis{A024389}&R&Empirical generating function fails at n=39\\
48&\oeis{A024390}&P&Empirical floor series with finite corrections\\
49&\oeis{A024403}&P&Empirical floor series with finite corrections\\
50&\oeis{A173693}&P&Rounded recurrence via a parity period\\
51&\oeis{A227012}&P&Exact floor formula and harmonic-block series\\
52&\oeis{A227013}&P&Exact floor formula and harmonic-block series\\
53&\oeis{A227014}&P&Exact floor formula and harmonic-block series\\
54&\oeis{A227015}&P&Exact floor formula and harmonic-block series\\
55&\oeis{A227016}&P&Exact floor formula and harmonic-block series\\
56&\oeis{A227017}&P&Exact floor formula and harmonic-block series\\
\bottomrule\end{longtable}
\Needspace{6\baselineskip}\subsection*{Generating functions and recurrences (51)}
\begin{longtable}{@{}r l c p{0.62\textwidth}@{}}
\toprule No.&Entry&Result&Selected assertion\\\midrule\endhead
57&\oeis{A174193}&P&Specified polynomial recurrence\\
58&\oeis{A346074}&P&Specified polynomial recurrence\\
59&\oeis{A247173}&P&Specified polynomial recurrence\\
60&\oeis{A247102}&P&Specified polynomial recurrence\\
61&\oeis{A234269}&P&Specified polynomial recurrence\\
62&\oeis{A268600}&P&Specified polynomial recurrence\\
63&\oeis{A268601}&P&Specified polynomial recurrence\\
64&\oeis{A108296}&P&Specified polynomial recurrence\\
65&\oeis{A244884}&P&Specified polynomial recurrence\\
66&\oeis{A244886}&P&Specified polynomial recurrence\\
67&\oeis{A257104}&P&Specified polynomial recurrence\\
68&\oeis{A099038}&P&Specified polynomial recurrence\\
69&\oeis{A337500}&P&Specified polynomial recurrence\\
70&\oeis{A127905}&P&Specified polynomial recurrence\\
71&\oeis{A392976}&P&Specified polynomial recurrence\\
72&\oeis{A186185}&P&Specified polynomial recurrence\\
73&\oeis{A220433}&P&Specified polynomial recurrence\\
74&\oeis{A071264}&P&Specified polynomial recurrence\\
75&\oeis{A392211}&P&Specified polynomial recurrence\\
76&\oeis{A091539}&P&Specified polynomial recurrence\\
77&\oeis{A027175}&P&Specified polynomial recurrence\\
78&\oeis{A027182}&P&Specified polynomial recurrence\\
79&\oeis{A177258}&P&Specified polynomial recurrence\\
80&\oeis{A177259}&P&Specified polynomial recurrence\\
81&\oeis{A177260}&P&Specified polynomial recurrence\\
82&\oeis{A351933}&P&Specified polynomial recurrence\\
83&\oeis{A001722}&P&Specified polynomial recurrence\\
84&\oeis{A345651}&P&Specified polynomial recurrence\\
85&\oeis{A111942}&P&Specified polynomial recurrence\\
86&\oeis{A394578}&P/R&Recurrence true; ratio limit is 4, not 1/4\\
87&\oeis{A396724}&P&Specified polynomial recurrence\\
88&\oeis{A121580}&P&Specified polynomial recurrence\\
89&\oeis{A180195}&P&Specified polynomial recurrence\\
90&\oeis{A121639}&P&Specified polynomial recurrence\\
91&\oeis{A121751}&P&Specified polynomial recurrence\\
92&\oeis{A121747}&P&Specified polynomial recurrence\\
93&\oeis{A121584}&P&Specified polynomial recurrence\\
94&\oeis{A121696}&P&Specified polynomial recurrence\\
95&\oeis{A121553}&P&Specified polynomial recurrence\\
96&\oeis{A187252}&P&Specified polynomial recurrence\\
97&\oeis{A001277}&P&Specified polynomial recurrence\\
98&\oeis{A001278}&P&Specified polynomial recurrence\\
99&\oeis{A001279}&P&Specified polynomial recurrence\\
100&\oeis{A001280}&P&Specified polynomial recurrence\\
101&\oeis{A045742}&P&Specified polynomial recurrence\\
102&\oeis{A121752}&P&Specified polynomial recurrence\\
103&\oeis{A027050}&P&Specified polynomial recurrence\\
104&\oeis{A027058}&P&Specified polynomial recurrence\\
105&\oeis{A027059}&P&Specified polynomial recurrence\\
106&\oeis{A027088}&R&Printed recurrence fails, including eventually\\
107&\oeis{A027089}&P&Specified polynomial recurrence\\
\bottomrule\end{longtable}
\subsection*{Four standalone repairs excluded from the count}
\begin{tabularx}{\textwidth}{@{}lX@{}}
\toprule Entry&Qualification\\\midrule
\oeis{A059837}&The printed recurrence uses a shifted sequence index; a corrected recurrence is proved.\\
\oeis{A142243}&The row-sum recurrence is shifted by one; the literal version is false and the corrected version is proved.\\
\oeis{A024388}&The title's block length disagrees with the data and cross-reference formula.\\
\oeis{A117436}&The diagonal formula requires a one-based column variable; a later formula already implies it.\\
\bottomrule
\end{tabularx}

\clearpage
% Proof packet: oeis_root_proofs.tex
\section{Quadratic character sums: all three assertions in A329649}
\label{round1:sec:characters}
The entry \oeis{A329649}, by Jianing Song (2019), defines
\[
 B(D)=\sum_{a=1}^{D}\chi_D(a)a^2,\qquad
 \chi_D(a)=\left(\frac{D}{a}\right),\quad D>0,\quad D\equiv0,1\pmod4.
\]
The symbol is the Kronecker symbol. The inspected entry still labels three
assertions as conjectures: nonvanishing; a divisibility statement with square,
$5$ times square, and odd powers of $2$ exceptions; and an exact $L$-value
formula for fundamental discriminants. We prove these assertions together.

\begin{theorem}\label{round1:thm:character}
Suppose $D=d f^2$ is not a square, where $d>1$ is the positive fundamental
discriminant of its quadratic field. Then
\begin{equation}\label{round1:eq:induced}
 B(D)=f^2 B(d)\prod_{p\mid f}(1-p\chi_d(p)),\qquad
 B(d)=\frac{d^{5/2}}{\pi^2}L(2,\chi_d)>0.
\end{equation}
In particular $B(D)\ne0$ for every allowed $D$, including squares.
If $D$ is neither a square nor $5$ times a square nor $2^t$ with odd
$t>1$, then $4D\mid B(D)$. If $D=5f^2$, then $4f^2\mid B(D)$.
For a fundamental discriminant $d>8$, the quotient $B(d)/(4d)$ is a
positive integer.
\end{theorem}

\begin{proof}
First let $\chi=\chi_d$ be nonprincipal, primitive, real and even. Character
orthogonality and the substitution $a\mapsto d-a$ give
\[
 \sum_{a=1}^d\chi(a)=0,\qquad \sum_{a=1}^d a\chi(a)=0.
\]
Consequently periodic lifting gives
\begin{equation}\label{round1:eq:lift}
 \sum_{a=1}^{Md}\chi(a)a^2=M B(d)\quad(M\ge1).
\end{equation}
Since $\chi_D(a)=\chi_d(a)\mathbf1_{\gcd(a,f)=1}$, inclusion--exclusion
over primes dividing $f$ but not $d$ gives
\begin{align*}
 B(D)&=\sum_{t\mid\operatorname{rad}(f),\ (t,d)=1}
  \mu(t)\chi_d(t)t^2\sum_{b=1}^{D/t}\chi_d(b)b^2\\
 &=f^2 B(d)\sum_{t\mid\operatorname{rad}(f)}\mu(t)\chi_d(t)t,
\end{align*}
which is the product in \eqref{round1:eq:induced}.

For completeness, the analytic step uses the classical quadratic Gauss-sum
identity
\[
 \sum_{a=1}^d\chi_d(a)e^{2\pi i ka/d}=\sqrt d\,\chi_d(k)
\]
and the absolutely convergent Fourier expansion
\[
 B_2(x)=x^2-x+\tfrac16
  =\frac1{\pi^2}\sum_{k\ge1}\frac{\cos(2\pi kx)}{k^2}
 \quad(0\le x\le1).
\]
Summing the latter against $\chi_d(a)$ at $x=a/d$ yields
$B(d)/d^2=\sqrt d\,L(2,\chi_d)/\pi^2$. The Euler product for
$L(2,\chi_d)$ is strictly positive. Every factor $1-p\chi_d(p)$ is
nonzero. If $D$ is a square, $\chi_D$ is a principal character and the
defining sum is positive. This proves nonvanishing.

We give the divisibility argument explicitly. For odd fundamental $d$,
pairing $a$ with $d-a$ gives
\[
 B(d)=2\sum_{1\le a<d/2}\chi_d(a)a(a-d),
\]
because the half-sum of the character is zero. Each $a(a-d)$ is even,
so $4\mid B(d)$. If $d$ has at least two prime factors, the Chinese
remainder theorem shows that $B(d)\equiv0\pmod p$ for every $p\mid d$:
the sum over the other prime factors contains a nonprincipal character
sum and vanishes. Since $d$ is squarefree, $d\mid B(d)$.
If $d=p$ is prime, then $p\equiv1\pmod4$ and Euler's criterion gives
\[
 B(p)\equiv\sum_{a=1}^{p-1}a^{(p-1)/2+2}\pmod p.
\]
This vanishes unless $p=5$, by the elementary power-sum identity over
$\mathbb F_p$. Thus $4d\mid B(d)$ for odd $d>5$, whereas $B(5)=4$.

Next let $d=4m$, where $m$ is odd, squarefree and $m\equiv3\pmod4$.
The primitive character satisfies $\chi_d(a+2m)=-\chi_d(a)$ for odd
$a$, and is even. Grouping the four residues
$a,2m-a,2m+a,4m-a$, for odd $0<a<m$ coprime to $m$, gives contribution
\[
 8m\chi_d(a)(m-a).
\]
It is divisible by $16m=4d$ because $m-a$ is even. These groups
partition all nonzero character terms.

Finally let $d=8m$, where $m$ is odd and squarefree. The corresponding
shift relation is $\chi_d(a+4m)=-\chi_d(a)$. The groups
$a,4m-a,4m+a,8m-a$, for odd $0<a<2m$ coprime to $m$, give
\[
 B(8m)=16m\sum_{\substack{1\le a<2m\\a\ {
m odd},\ (a,m)=1}}
       \chi_{8m}(a)(2m-a).
\]
For $m>1$ the sum has $\varphi(m)$ odd summands, an even number, so
$32m=4d$ divides $B(d)$. At $m=1$ direct evaluation gives $B(8)=16$.

Now apply \eqref{round1:eq:induced}. Fundamental $d>8$ immediately gives
$4D\mid B(D)$. At $d=5$ the product gives the asserted $4f^2$
divisibility. At $d=8$, any odd prime dividing $f$ supplies the even
factor $1-p\chi_8(p)$; if no such prime exists then $D=2^{3+2r}$.
These are precisely the exceptional nonsquare powers of $2$ in the
entry. This completes all three assertions.
\end{proof}

\begin{corollary}
For nonsquare $D=df^2$, the sign of $B(D)$ is
$(-1)^{s}$, where $s$ counts the distinct primes $p\mid f$ with
$\chi_d(p)=1$. In particular the negative terms are explained exactly
by the primes splitting in the primitive quadratic character.
\end{corollary}

\paragraph{Source correction.}
The entry's displayed $D=5$ numerical example omits a factor $4$ in one
expression: the correct value is $L(2,\chi_5)=4\pi^2/(25\sqrt5)$.
The general formula in its conjecture (c) has the correct normalization.
The classical Gauss-sum and Bernoulli Fourier formulas are imported
standard identities; the argument does not claim their novelty.

\section{A general cancellation theorem for triangular inverses}
\label{round1:sec:triangular}
Let $f(0)=0$ and let $f(1),f(2),\ldots$ be arbitrary elements of a
commutative ring. Define a unit lower triangular array $M$ by
$M(0,0)=1$, $M(n,0)=0$ for $n>0$, and
\[
 M(n,k)=M(n-1,k-1)+f(k)M(n-1,k)\quad(1\le k\le n).
\]
Write $I=M^{-1}$. Here $h_j$ and $e_j$ denote complete and elementary
symmetric polynomials, respectively, with $h_0=e_0=1$.

\begin{theorem}\label{round1:thm:trianglecancel}
For $1\le k\le n$, put $m=n-k$. Then
\begin{equation}\label{round1:eq:trianglecancel}
 \sum_{i=0}^{m}I(m,i)M(n-1+i,n-1)
   =h_m\bigl(f(m),f(m+1),\ldots,f(n-1)\bigr).
\end{equation}
At $m=0$ both sides equal $1$. In particular the first-column values
at row $n+1$ are $f(n)^n$, with the usual convention $0^0=1$.
\end{theorem}

\begin{proof}
The defining recurrence gives
$M(r+i,r)=h_i(f(1),\ldots,f(r))$. The coefficients of
$x\prod_{j=1}^{m-1}(x-f(j))$ form row $m$ of $I$, since these monic
polynomials and the monomials are inverse bases under the recurrence.
Thus $I(m,m-j)=(-1)^j e_j(f(1),\ldots,f(m-1))$.
The left side of \eqref{round1:eq:trianglecancel} is the coefficient of $z^m$ in
\[
 \frac{\prod_{j=1}^{m-1}(1-f(j)z)}
      {\prod_{j=1}^{n-1}(1-f(j)z)}
   =\prod_{j=m}^{n-1}(1-f(j)z)^{-1}.
\]
This is the asserted complete symmetric polynomial. The $m=0$ case
follows directly without using a product with a negative upper bound.
\end{proof}

\subsection{A355401: all its displayed q-triangle conjectures}
The construction in \oeis{A355401} takes $f(j)=q^{j-1}$ for $j\ge1$.
Theorem~\ref{round1:thm:trianglecancel} gives, for $1\le k\le n$,
\begin{equation}\label{round1:eq:qtriangle}
 T_q(n,k)=q^{(n-k)(n-k-1)}\genfrac{[}{]}{0pt}{}{n-1}{k-1}_q.
\end{equation}
Indeed, with $m=n-k$, factor $q^{m-1}$ from the $k$ variables in
$h_m$ and use
$h_m(1,q,\ldots,q^{k-1})=\genfrac{[}{]}{0pt}{}{m+k-1}{m}_q$.
This proves the full array formula, and hence the first-column formula
$T_q(n+1,1)=q^{n(n-1)}$. The identity is polynomial in $q$ and extends
to all specializations where the displayed definitions are interpreted
polynomially.

The inverse conjecture also follows. Set $c_m=q^{m(m-1)}$ and define
$g_0=1$ and
\[
 g_m=-\sum_{i=0}^{m-1}\genfrac{[}{]}{0pt}{}{m}{i}_q g_i c_{m-i}.
\]
Then the inverse has entries $R_q(n,k)=\genfrac{[}{]}{0pt}{}{n-1}{k-1}_q g_{n-k}$
for $k>0$, with $R_q(n,0)=\mathbf1_{n=0}$. To check this, use
\[
 \genfrac{[}{]}{0pt}{}{n-1}{j-1}_q\genfrac{[}{]}{0pt}{}{j-1}{k-1}_q
 =\genfrac{[}{]}{0pt}{}{n-1}{k-1}_q\genfrac{[}{]}{0pt}{}{n-k}{j-k}_q.
\]
The matrix product reduces to the displayed convolution defining $g$,
which is $0$ off the diagonal and $1$ on it. This also proves
$g_n=R_q(n+1,1)$. These formulas settle the special case $q=2$ and
the general version recorded in the comments of the same entry.

\subsection{A355282: the q-Stirling counterpart}
Here $f(j)=[j]_q=(q^j-1)/(q-1)$ and the construction is exactly that
of Theorem~\ref{round1:thm:trianglecancel}. It follows at once that
$T_q(n+1,1)=[n]_q^n$. For $m=n-k$ its full conjectured formula is
\begin{equation}\label{round1:eq:qstirling}
 T_q(n,k)=\frac{1}{(q-1)^m}
 \sum_{i=0}^m(-1)^i\binom{n-1}{i}
 \genfrac{[}{]}{0pt}{}{n-1-i}{k-1}_q q^{m(m-i)}.
\end{equation}
To prove it, use the translation identity in $k$ variables
\[
 h_m(x_1-1,\ldots,x_k-1)
 =\sum_{i=0}^m(-1)^i\binom{k+m-1}{i}h_{m-i}(x_1,\ldots,x_k).
\]
It follows by writing
$\prod_j(1-(x_j-1)z)^{-1}=(1+z)^{-k}
 \prod_j(1-x_j z/(1+z))^{-1}$ and extracting coefficients.
Apply it with $(x_1,\ldots,x_k)=(q^m,\ldots,q^{n-1})$ and use
$h_{m-i}(q^m,\ldots,q^{n-1})=q^{m(m-i)}\genfrac{[}{]}{0pt}{}{n-1-i}{k-1}_q$.
The resulting expression is \eqref{round1:eq:qstirling}. Although written
as a quotient, it has a removable specialization at $q=1$ by the
complete-symmetric-polynomial form. Both the first-column assertion
and the full formula in \oeis{A355282} are therefore proved.

\subsection{A350557: an explicit inverse}
Let $T$ be the array in \oeis{A350557}, and put $u_n=(2n-1)!!$,
including $u_0=1$. Its definition is
\[
 T(n,0)=u_n,\qquad
 T(n,k)=u_n\sum_{i=k}^n\binom{i-1}{k-1}\frac1{u_i}\quad(k\ge1).
\]
Let $L$ be lower bidiagonal with $L(n,n)=1$ and
$L(n,n-1)=-(2n-1)$, and let $P=1\oplus\mathrm{Pascal}$, so that
$P(n,k)=\binom{n-1}{k-1}$ for $n,k\ge1$ and its column $0$ is the
zeroth unit vector. The defining sum telescopes to $LT=P$.
Thus $T^{-1}=P^{-1}L$. Direct multiplication gives
\[
 (T^{-1})(n,0)=(-1)^n,
\]
and, for $1\le k\le n$,
\[
 (T^{-1})(n,k)=(-1)^{n-k}
 \left\{\binom{n-1}{k-1}+(2k+1)\binom{n-1}{k}\right\}.
\]
The bracket equals
$\binom{n-1}{k-1}(2k(n-k)+n)/k$, the defining expression for
\oeis{A350512}. If $J=\operatorname{diag}((-1)^n)$, this says
$T^{-1}=J A350512 J$ and hence $(A350512)^{-1}=JTJ$.
This is exactly the signed-inverse assertion of A350557.

\subsection{A136097: divisibility of a Gaussian central entry}
The inverse of the Gaussian array has entries
\[
 G^{-1}(n,k)=(-1)^{n-k}q^{\binom{n-k}{2}}\genfrac{[}{]}{0pt}{}{n}{k}_q.
\]
Multiplying by $G$ proves the formula using the finite $q$-binomial
identity $\sum_{j=0}^r(-1)^j q^{\binom j2}\genfrac{[}{]}{0pt}{}{r}{j}_q
=\prod_{j=0}^{r-1}(1-q^j)$, which is zero for $r>0$.
Moreover,
\[
 \frac{\genfrac{[}{]}{0pt}{}{2n}{n}_q}{[n+1]_q}
 =\genfrac{[}{]}{0pt}{}{2n}{n}_q-q\genfrac{[}{]}{0pt}{}{2n}{n-1}_q\in\mathbb Z[q].
\]
At $n=0$, interpret the second Gaussian coefficient as $0$.
Therefore the central inverse entry is divisible, at every integer
$q>1$, by $[n+1]_q q^{n(n-1)/2}$. This proves the general conjecture
in \oeis{A136097}. Its $q=2$ connection to the $q$-Catalan numbers
is already recorded in the entry; the proof uses classical identities
and is not a claim of a new $q$-Catalan theorem.

\section{A direct row sum and a diagonal indexing correction}
\label{round1:sec:directrows}
\subsection{A146955}
The row polynomial of \oeis{A146955}, for $n\ge1$, is
\[
 p_n(x)=(1+x)^n+2^{n-4}\sum_{m=1}^{n-1}(2^m+2m)(x^m+x^{n-m}).
\]
Its row sum is $p_n(1)$, and evaluation gives
\[
 p_n(1)=2^n+2^{n-3}\bigl((2^n-2)+n(n-1)\bigr)
 =\frac{2^n(2^n+6-n+n^2)}8.
\]
The empty sum at $n=1$ obeys the same calculation. This proves the
explicit row-sum conjecture without fitting any terms.

\subsection{A117436: state the column convention}
Write $c_d=d![z^d]\sec(2z)$. The established column exponential
generating function of \oeis{A117436} gives
\[
 T(k+d,k)=c_d\binom{k+d}{d}.
\]
With a \emph{one-based} column variable $x=k+1$, this is
$(-1)^d c_d\binom{-x}{d}$. Odd $d$ have $c_d=0$; the first nonzero
constants are $c_0=1,c_2=4,c_4=80,c_6=3904$.
Thus the proportionality conjecture has an exact formula under the
one-based convention. If its $x$ is instead the zero-based column
index used by $T(n,k)$, the literal statement is false: at $d=2,k=0$
it compares $4$ with a multiple of $\binom02=0$. We record this as
an indexing clarification, separately from unambiguous conjecture
proofs. The later explicit formula already in the entry implies the
corrected statement.

\section{A127752: the inverse row sums and their Feigenbaum parity}
\label{round1:sec:feigenbaum}
For an integer $c\ge1$, let
\[
 A_c(n,k)=\begin{cases}(cn+1)^{-1},&\lceil n/2\rceil\le k\le n,\\
 0,&\text{otherwise},\end{cases}
 \qquad n,k\ge0,
\]
and let $s_c(n)$ be the sum of row $n$ of $A_c^{-1}$. Then
\begin{equation}\label{round1:eq:rowvaluation}
 s_c(n)=c\bigl(\nu_2(n+1)+1\bigr)
       +(1-c)\mathbf1_{n+1\text{ is a power of }2}.
\end{equation}
Indeed $A_c s_c=\boldsymbol1$ gives
\[
 \sum_{k=\lceil n/2\rceil}^{n}s_c(k)=cn+1.
\]
Subtract the preceding equation. At even $n=2m\ge2$, the lower
endpoint is unchanged, so $s_c(2m)=c$. At odd $n=2m+1$, the
lower endpoint increases by one, so $s_c(2m+1)=c+s_c(m)$.
Together with $s_c(0)=1$, these relations prove
\eqref{round1:eq:rowvaluation} by induction: iterate through the trailing
binary $1$ digits, stopping at a positive even number or at $0$.
At $c=3$, the correction term in \eqref{round1:eq:rowvaluation} is even,
and hence
\[
 s_3(n)\equiv\nu_2(n+1)+1\pmod2.
\]
The right side is exactly the one-based sequence
$\oeis{A035263}(n+1)$, proving the conjecture in \oeis{A127752}.
The same assertion is repeated in \oeis{A127751}; that copy receives
no additional count. Related older inverse-matrix assertions with
$c=2$ are omitted from the claimed closures because a proof for their
underlying matrix is already linked in \oeis{A111967}.

\section{Full-length collinear subsets in a discrete cube}
\label{round1:sec:gridlines}
The entries \oeis{A081199}, \oeis{A081200}, and \oeis{A081201}
conjecture interpretations as collinear $q$-subsets of
$\{0,1,\ldots,q-1\}^d$, with $q=4,5,6$, respectively.
The following count proves all three.

\begin{theorem}
For all integers $q\ge2$ and $d\ge0$, the number of unordered
$q$-element collinear subsets of $\{0,\ldots,q-1\}^d$ is
\[
 \frac{(q+2)^d-q^d}{2}.
\]
\end{theorem}
\begin{proof}
For $d=0$ the count is zero. Otherwise a line containing two integer
points has a primitive integer direction vector $v$, and its integer
points have integer parameters in that direction. To see the latter
claim, choose integers $c_j$ with $\sum_j c_jv_j=1$; the parameter
of an integer displacement $t v$ is then $\sum_j c_j(t v_j)$.
Write the parameters of our $q$ points in increasing order as
$t_1<\cdots<t_q$. Their span is at least $q-1$. For any nonzero
coordinate $v_j$, containment in the cube gives
\[
 q-1\ge |v_j|(t_q-t_1)\ge q-1.
\]
Thus every nonzero $v_j$ equals $1$ or $-1$, the span is exactly
$q-1$, and the parameters are consecutive. Each subset is therefore
an entire $q$-point line. If the support of $v$ has size $k$, choose
the support in $\binom dk$ ways, its signs in $2^{k-1}$ ways after
identifying the two orientations, and the constant coordinates in
$q^{d-k}$ ways. Each varying coordinate is forced to run from one
boundary value to the other. The count is
\[
 \sum_{k=1}^d\binom dk2^{k-1}q^{d-k}
 =\frac{(q+2)^d-q^d}{2}.
\]
\end{proof}
Only the three explicitly conjectured instances are entered in the
closure ledger; the general theorem is not used to manufacture more
sequence entries.

\section{A054120: the second McGarvey assertion}
\label{round1:sec:mcgarvey}
Let $R(n,k)$ denote \oeis{A081577} in triangular coordinates, with
out-of-range entries zero. Its bivariate generating function, in
coordinates $(u,v)=(k,n-k)$, is
\[
 \sum_{n\ge0}\sum_{k=0}^n R(n,k)u^k v^{n-k}
 =\frac1{1-u-v-2uv}.
\]
Consequently $R$ obeys
$R(n,k)=R(n-1,k-1)+R(n-1,k)+2R(n-2,k-1)$ away from
the origin. Set $S(n,k)=R(n,k)-R(n-2,k-1)$. Then
\[
 \sum_{n,k} S(n,k)u^k v^{n-k}
 =\frac{1-uv}{1-u-v-2uv}.
\]
Its first three rows are $(1)$, $(1,1)$, and $(1,3,1)$;
thereafter it obeys the same recurrence. Hendel's proof already
recorded in \oeis{A054120} establishes that recurrence for its
geometrically defined array $T$, for $n\ge3$, with exactly these
initial rows. Induction therefore proves $S=T$. This proves the
second McGarvey assertion and the later tentative bivariate
generating function. The first assertion, explicitly proved in
the entry in 2015, is used as a cited premise and is not counted
as a new resolution.

\section{Matrix self-similarity and partitions into powers of a base}
\label{round1:sec:basepartitions}
Fix an integer $b\ge2$. Let $M$ be the lower unit triangular matrix
with first column $1$ and with
\[
 M^b(n,k)=M(n+1,k+1).
\]
These conditions recursively determine $M$ row by row. They are the
definitions of \oeis{A078121} and \oeis{A078122} at $b=2,3$.
The prose names of both entries reverse the indices in the
first-column condition; we use the row--column convention specified
by their examples and programs.
Let
\[
 p_r(N)=[x^N]\prod_{j=0}^r(1-x^{b^j})^{-1},\qquad p_r(N)=0\ (N<0).
\]

\begin{theorem}\label{round1:thm:basepartitions}
For $0\le k\le n$,
\[
 M(n,k)=p_{n-k}\bigl(b^n-b^{n-k}\bigr).
\]
Moreover $\sum_{k=0}^nM(n,k)$ counts partitions of $b^n$ into
powers of $b$.
\end{theorem}
\begin{proof}
For nonnegative integer $t$ define $P_0(t)=1$, and for $r\ge1$
define $P_r(0)=0$ and $P_r(t)=p_r(b^r(t-1))$ when $t\ge1$.
Classifying a partition according to the number of parts $b^r$
gives, for $t\ge1$,
\[
 P_r(t)=\sum_{u=0}^{t-1}p_{r-1}(b^r u).
\]
It follows, also at $t=0$, that
\begin{equation}\label{round1:eq:partitiondiff}
 P_r(t+1)-P_r(t)=P_{r-1}(bt+1).
\end{equation}
Induction on $r$ in this identity gives
\begin{equation}\label{round1:eq:partitionaddition}
 P_r(t+1)=\sum_{j=0}^rP_{r-j}(b^j t).
\end{equation}
For $r=0$ this just states $1=1$.

Deleting the first row and first column is an algebra homomorphism
$S$ on lower triangular matrices, because their entrywise matrix
products involve finite sums. Therefore
$S^k(M^t)=M^{b^k t}$. Put $F_r(t)=(M^t)(r,0)$. This gives
\[
 (M^t)(n,k)=F_{n-k}(b^k t).
\]
Since the first column of $M$ is $1$, multiplication by $M$ gives
\[
 F_r(t+1)=\sum_{j=0}^rF_{r-j}(b^j t),\qquad
 F_r(0)=\mathbf1_{r=0}.
\]
These are exactly \eqref{round1:eq:partitionaddition} and its initial
values. Induct first on $r$, proving the assertion for every $t$
at all smaller $r$, and then on $t$ for that fixed $r$. The term
$j=0$ uses the preceding $t$ at the same $r$, whereas every other
term uses a smaller $r$. This proves $F_r(t)=P_r(t)$ without
assuming a matrix formula in advance. At $t=1$ the desired formula
is $F_{n-k}(b^k)=p_{n-k}(b^n-b^{n-k})$.

Finally, classify partitions of $b^n$ by their largest part $b^r$,
$0\le r\le n$, and remove one such part. The resulting number for
that class is $p_r(b^n-b^r)$. Summing over $r=n-k$ gives the row sum.
\end{proof}
This proves the coefficient conjectures in both entries, and the
row-sum conjecture in A078122. The 2025 MathOverflow link from
A078121 discusses a different proposed closed form along fixed
diagonals; it is not used as a premise here.

\section{A113249: the universal perfect-square assertion}
\label{round1:sec:floretion}
For any integer $m$, define the sequence by
\[
 a_0=-1,\quad a_1=4,\quad a_2=3m^2-16,\quad
 a_3=(m^2-8)^2,
\]
\[
 a_n=-4a_{n-1}+4m^2a_{n-3}+m^4a_{n-4}\quad(n\ge4).
\]
The conjecture in \oeis{A113249}, repeated in
\oeis{A113250}--\oeis{A113256}, asks whether every odd-indexed
term is a square for every $m$. Define integers
\[
 c_0=1,\quad c_1=2,\qquad c_{r+2}=4c_{r+1}-m^2c_r.
\]
Then
\begin{equation}\label{round1:eq:squareodd}
 a_{2r+1}=c_{r+1}^{\,2}\quad(r\ge0).
\end{equation}
Here is an algebraic proof that includes degenerate parameters.
Temporarily take $m$ indeterminate and let $\alpha,\beta$ be the
roots of $z^2-4z+m^2=0$, so
$c_j=(\alpha^j+\beta^j)/2$. Set
\[
 A_n=\frac{m^{n+1}+(-m)^{n+1}
                  +(-\alpha)^{n+1}+(-\beta)^{n+1}}4.
\]
Each of the four geometric terms is annihilated by
\[
 (z^2-m^2)(z^2+4z+m^2)
 =z^4+4z^3-4m^2z-m^4.
\]
Direct evaluation of $A_0,A_1,A_2,A_3$ gives the four initial
values above, so $A_n=a_n$. At odd index $2r+1$ one obtains
\[
 A_{2r+1}=\frac{2m^{2r+2}+\alpha^{2r+2}+\beta^{2r+2}}4
 =\left(\frac{\alpha^{r+1}+\beta^{r+1}}2\right)^2,
\]
using $\alpha\beta=m^2$. Both sides of \eqref{round1:eq:squareodd}
are polynomials in $m$, so the identity specializes to every
integer $m$, including $m=0$ and $m=\pm2$. The $c_j$ recurrence
establishes their integrality directly.

The entry already attributes the $m=3$ case to Robert Israel
(2017), and subsequent formulas cover some other individual
parameters. We claim only a proof here of the remaining universal
statement. Its eight copied appearances receive one ledger item,
not eight; no priority claim is made for the elementary generalization.

\clearpage
% Proof packet: oeis_arithmetic_proofs.tex
\section{Arithmetic, digital, and parity conjectures}

The results in this section address precisely the assertions stated below.
Some of the cited entries contain further conjectures which are not addressed.
The first-occurrence assertion in A334109 and A334110 is the same conjecture
recorded twice.  Likewise the transform assertions in A284005 and A329369 are
equivalent formulations.  These overlaps should be retained in any count of
independent problems.  All source pages were inspected on October 20, 2026.

\subsection{Remainder sums: A049803, A049804, and A330358}

For an integer $b\ge2$ and $n\ge1$, define
\[
 S_b(n)=\sum_{j=1}^{\lfloor\log_b n\rfloor}(n\bmod b^j),
 \qquad L_b(n)=\lfloor\log_b n\rfloor+1.
\]
The three entries concern $b=3,4,5$, respectively.  Each records a recurrence
conjecture of Petros Hadjicostas from December 2019.
Sources: \url{https://oeis.org/A049803},
\url{https://oeis.org/A049804}, \url{https://oeis.org/A330358}.

\paragraph{Theorem.}
For $n\ge1$ and $0\le r<b$,
\[
 S_b(bn+r)=bS_b(n)+rL_b(n).
\]
Writing $A_b(x)=\sum_{n\ge1}S_b(n)x^n$, one also has the formal identity
\begin{align*}
 A_b(x)={}&b(1+x+\cdots+x^{b-1})A_b(x^b)\\
 &+\frac{x+2x^2+\cdots+(b-1)x^{b-1}}{1-x^b}
       \sum_{j\ge1}x^{b^j}.
\end{align*}

\paragraph{Proof.}
Put $L=L_b(n)$.  The number $bn+r$ has $L+1$ base-$b$ digits, and
\[
 (bn+r)\bmod b^j=b(n\bmod b^{j-1})+r
 \qquad(1\le j\le L).
\]
Sum this equality, using $n\bmod1=0$.  For the generating function, sum the
recurrence over residue classes and use
\[
 \sum_{n\ge1}L_b(n)z^n
 =\frac{\sum_{j\ge0}z^{b^j}}{1-z}.
\]
The omitted integers $1\le n<b$ have $S_b(n)=0$, so there is no additional
initial polynomial.  This proves both assertions.\hfill$\square$

\subsection{Delannoy divisor pairing: A323689 and A324300}

Sources: \url{https://oeis.org/A323689} and
\url{https://oeis.org/A324300}.  These entries conjecture restrictions on the
positions of odd coefficients in two series involving $i=\sqrt{-1}$.

\paragraph{Lemma.}
For $u,v\ge0$, put
\[
 D(u,v)=\sum_{j\ge0}2^j\binom uj\binom vj.
\]
Then $D(u,v)=D(v,u)$, every $D(u,v)$ is odd, and
\[
 [y^v]\frac{(y+i)^u}{(1+iy)^{u+1}}
 =i^{u-v}D(u,v).
\]

\paragraph{Proof.}
Symmetry and oddness follow from the displayed sum.  Summing first over $v$
and using $\sum_{v\ge0}\binom vjz^v=z^j/(1-z)^{j+1}$ gives
\[
 \sum_{v\ge0}D(u,v)z^v=\frac{(1+z)^u}{(1-z)^{u+1}}.
\]
Substitute $z=-iy$ and multiply by $i^u$.\hfill$\square$

\paragraph{Theorem.}
For A323689, $a(N)$ is odd if and only if $N$ is a square.
For A324300, $a(N)$ is odd if and only if $N=k(k+2)$ for some $k\ge0$.
Both series have integer coefficients under their specified constant-term
conventions.

\paragraph{Proof.}
For A323689 and $N\ge1$, coefficient extraction gives the finite sum
\[
 a(N)=\sum_{\substack{uv=N\\u,v\ge1}}i^{u-v}D(u,v).
\]
Pair the terms $(u,v)$ and $(v,u)$ when $u\ne v$.  Their sum is
\[
 \bigl(i^{u-v}+i^{v-u}\bigr)D(u,v),
\]
which is either zero or an even integer.  There is a diagonal term precisely
when $N$ is a square, and that term is odd.  The prescribed $a(0)=1$ is also
consistent.  This reasoning extracts only positive coefficients, so it does
not assume convergence of the original constant terms.

For A324300, the summand is multiplied by $x^u$ and $y$ is replaced by
$x^{u+1}$.  Consequently the relevant exponent is
$N=u+(u+1)v=(u+1)(v+1)-1$, and
\[
 a(N)=\sum_{\substack{(u+1)(v+1)=N+1\\u,v\ge0}}i^{u-v}D(u,v).
\]
The same pairing proves that the coefficient is odd exactly when $N+1$ is
a square.\hfill$\square$

\subsection{Reversing weighted partitions: A362560}

Source: \url{https://oeis.org/A362560}.
For a finite integer sequence $(y_1,\ldots,y_k)$ of positive total $n$, let
$W=\sum_{j=1}^k jy_j$.  The weighted sum of its reversal is
\[
 W_{\rm rev}=\sum_{j=1}^k(k+1-j)y_j=(k+1)n-W.
\]
Thus $n\mid W$ if and only if $n\mid W_{\rm rev}$, proving the stated
partition conjecture.  In fact
$\gcd(n,W)=\gcd(n,W_{\rm rev})$, and no ordering assumption on the sequence
is needed.\hfill$\square$

\subsection{Superpolynomial growth: A085765}

Sources: \url{https://oeis.org/A085765} and
\url{https://oeis.org/A086450}.  The parity conjecture on these pages was
already proved in 2007 and is not counted here.  We address the remaining
growth conjecture.

Let $b(0)=1$, $b(2n+1)=b(n)$, and
$b(2n)=\sum_{j=0}^n b(j)$.  Then A085765 is
$A(N)=\sum_{j=0}^N b(j)$.

\paragraph{Theorem.}
One has
\[
 \lim_{N\to\infty}\frac{\log A(N)}{\log N}=+\infty.
\]

\paragraph{Proof.}
The function $A$ is positive and increasing.  For $N\ge4$,
\begin{align*}
 A(N)&\ge\sum_{j=\lceil N/4\rceil}^{\lfloor N/2\rfloor}b(2j)
 =\sum_{j=\lceil N/4\rceil}^{\lfloor N/2\rfloor}A(j)\\
 &\ge\lfloor N/4\rfloor A(\lfloor N/4\rfloor).
\end{align*}
In particular, iteration at powers of $4$ gives
$A(4^r)\ge4^{r(r-1)/2}$.  If $4^r\le N<4^{r+1}$, then
\[
 \frac{\log A(N)}{\log N}
 \ge\frac{r(r-1)}{2(r+1)}\longrightarrow+\infty.
\]
The argument proves a limit, not just an unbounded subsequence.
\hfill$\square$

\subsection{A Pascal-parity XOR identity: A100311}

Source: \url{https://oeis.org/A100311}.
The selected 2012 conjecture is an immediate consequence of Pascal's
identity.  The page also cites Vladimir Shevelev's established polynomial
framework, so this proof does not support a claim of worldwide priority.

By definition,
\[
 a(n)=\sum_{k=0}^n\left(\binom nk\bmod2\right)8^k.
\]
The binary digits can be nonzero only at positions $3k$.  Multiplying by $8$
shifts them by three positions, while bitwise XOR adds their coefficients
modulo two.  Therefore Pascal's identity yields
\[
 (8a(n))\mathbin{\mathrm{XOR}}a(n)
 =\sum_k\left(\binom n{k-1}+\binom nk\bmod2\right)8^k=a(n+1).
\]
Together with $a(0)=1$, this proves the conjecture.  The same proof applies
with $8$ replaced by any $2^s$, $s\ge1$.\hfill$\square$

\subsection{A totient-weighted prime characterization: A156834}

Source: \url{https://oeis.org/A156834}.  The displayed formula is
\[
 a(n)=\sum_{d\mid n}\varphi(d)\binom{d+n/d-2}{d-1}.
\]

\paragraph{Theorem.}
For $n>1$, $a(n)=n$ if and only if $n$ is prime; otherwise $a(n)>n$.

\paragraph{Proof.}
Every binomial coefficient in the sum is at least one, hence
$a(n)\ge\sum_{d\mid n}\varphi(d)=n$.  For $d=1$ and $d=n$ the coefficient
equals one.  Any proper divisor $1<d<n$ has $d\ge2$ and $n/d\ge2$, making
its binomial coefficient strictly greater than one.  Since $\varphi(d)>0$,
the inequality is strict for every composite $n$, and is an equality for
every prime.\hfill$\square$

\subsection{Primitive compositions: A156833}

Source: \url{https://oeis.org/A156833}.  Its matrix-product definition gives
\[
 a(n)=\sum_{r\mid n}\mu(n/r)
       \sum_{d\mid r}d\binom{d+r/d-2}{d-1}.
\]

\paragraph{Theorem.}
For $n>1$, $a(n)=n$ if and only if $n$ is prime; otherwise $a(n)>n$.

\paragraph{Proof.}
Reverse the finite divisor sums to obtain
\[
 a(n)=\sum_{d\mid n}dP_d(n/d),\qquad
 P_d(m)=\sum_{q\mid m}\mu(m/q)\binom{d+q-2}{d-1}.
\]
The binomial coefficient $\binom{d+m-2}{d-1}$ counts integer tuples
$(x_1,\ldots,x_d)$ with $x_1\ge1$, all other coordinates nonnegative, and
sum $m$.  Sorting these tuples by the gcd of their coordinates shows, by
M\"obius inversion, that $P_d(m)$ counts the tuples with gcd one.
Thus all $P_d(m)$ are nonnegative, $P_d(1)=1$, and $P_1(m)=0$ for $m>1$.
The summand $d=n$ is exactly $n$.  If $n$ is prime it is the only nonzero
summand.  For composite $n$ choose $1<d<n$ and put $m=n/d\ge2$; the tuple
$(m-1,1,0,\ldots,0)$ has gcd one, so $P_d(m)\ge1$.  This gives the strict
inequality for composites.\hfill$\square$

\subsection{Abundancy and prime-factor multiplicity: A008578}

Source: \url{https://oeis.org/A008578}.  We prove Irina Gerasimova's 2013
characterization, recorded there together with partial observations.

\paragraph{Theorem.}
For a positive integer $n$,
\[
 \sigma(n)>n\Omega(n)
 \quad\Longleftrightarrow\quad n=1\text{ or }n\text{ is prime}.
\]
Among composites, equality $\sigma(n)=n\Omega(n)$ occurs only at $n=6$.

\paragraph{Proof.}
For a composite $n$, let $r=\omega(n)$ and $k=\Omega(n)$, and order its
distinct prime factors as $p_1<\cdots<p_r$.  Since $p_j\ge j+1$, if
$k\ge r+1$ then
\[
 \frac{\sigma(n)}n<\prod_{j=1}^r\frac{p_j}{p_j-1}
 \le\prod_{j=1}^r\frac{j+1}{j}=r+1\le k.
\]
If $k=r$, the integer is squarefree and $k\ge2$, so
\[
 \frac{\sigma(n)}n=\prod_{j=1}^k(1+p_j^{-1})
 \le\prod_{j=1}^k\frac{j+2}{j+1}=\frac{k+2}{2}\le k.
\]
Equality in the latter case requires $k=2$, $p_1=2$, $p_2=3$, giving $n=6$.
The values $n=1$ and prime $n$ satisfy the asserted strict inequality
directly.\hfill$\square$

\subsection{An inverse Gray-code recurrence: A006068}

Sources: \url{https://oeis.org/A006068},
\url{https://oeis.org/A063946}, \url{https://oeis.org/A053645}, and
\url{https://oeis.org/A053644}.
Let $G(n)=n\mathbin{\mathrm{XOR}}\lfloor n/2\rfloor
\mathbin{\mathrm{XOR}}\lfloor n/4\rfloor\mathbin{\mathrm{XOR}}\cdots$.
This is the inverse Gray-code permutation A006068.

\paragraph{Theorem.}
With the conventions of the cited entries,
\[
 G(n)=G(\mathrm{A053645}(\mathrm{A063946}(n)))+
       \mathrm{A053644}(n),\qquad n>0.
\]

\paragraph{Proof.}
Write $n=2^L+m$, $0\le m<2^L$.  The map $G$ is linear for XOR and
$G(2^L)=2^{L+1}-1$, whence
\[
 G(n)=2^L+(2^L-1-G(m)).
\]
For $L\ge1$, flipping the second most significant bit and deleting the
first bit produces $h=m\mathbin{\mathrm{XOR}}2^{L-1}$.
Therefore $G(h)=G(m)\mathbin{\mathrm{XOR}}(2^L-1)=2^L-1-G(m)$,
which proves the identity.  The case $n=1$ follows directly.
\hfill$\square$

\subsection{A negabinary XOR identity: A016789}

Sources: \url{https://oeis.org/A016789},
\url{https://oeis.org/A005351}, \url{https://oeis.org/A005352}.
Let $B(t)$ be the integer obtained by reading the negabinary digits of a
signed integer $t$ as ordinary binary digits.

\paragraph{Theorem.}
For every $n\ge0$,
\[
 n\mathbin{\mathrm{XOR}}B(n+1)\mathbin{\mathrm{XOR}}B(-n-1)=3n+2.
\]

\paragraph{Proof.}
Set $X(t)=B(t)\mathbin{\mathrm{XOR}}B(-t)$ for $t\ge0$.
Negabinary digit extraction gives, for every signed $u$,
$B(2u)=2B(-u)$ and $B(2u+1)=2B(-u)+1$.  Consequently
\[
 X(2u)=2X(u),\quad X(4u+1)=4X(u)+2,\quad
 X(4u+3)=4X(u+1)+2.
\]
Together with $X(0)=0$, these determine $X$ by induction.
The function $F(t)=t\mathbin{\mathrm{XOR}}3t$ obeys the same recurrences.
For the last recurrence use
\[
 u\mathbin{\mathrm{XOR}}(3u+2)
 =(u+1)\mathbin{\mathrm{XOR}}3(u+1):
\]
decrementing positive integers $t$ and $3t$ toggles identical bit masks,
because $v_2(t)=v_2(3t)$.
Hence $X(t)=t\mathbin{\mathrm{XOR}}3t$.  Applying the same decrement
observation with $t=n+1$ yields
\[
 (t-1)\mathbin{\mathrm{XOR}}X(t)=3t-1=3n+2,
\]
as required.\hfill$\square$

\subsection{Digital products and excedance sets: A284005 and A329369}

Sources: \url{https://oeis.org/A284005} and
\url{https://oeis.org/A329369}.
Write $\operatorname{wt}(n)$ for binary popcount and set
\[
 A(0)=1,\qquad A(n)=(1+\operatorname{wt}(n))A(\lfloor n/2\rfloor).
\]
This is A284005.  Equivalently,
$A(n)=\prod_{j\ge0}(1+\operatorname{wt}(\lfloor n/2^j\rfloor))$,
with all sufficiently late factors equal to one.

\paragraph{Theorem (trailing-zero identity).}
For $m,n\ge0$,
\[
 A(2^m(2n+1))=\sum_{k=0}^{m+1}\binom{m+1}{k}A(2^kn).
\]

\paragraph{Proof.}
One has $A(2^kn)=(\operatorname{wt}(n)+1)^kA(n)$ and
$A(2^m(2n+1))=(\operatorname{wt}(n)+2)^{m+1}A(n)$.
Both statements also hold at $n=0$.  The ordinary binomial theorem finishes
the proof.\hfill$\square$

\paragraph{Theorem (leading-bit identity).}
Let $n=2^d+m$, $0\le m<2^d$, $d\ge1$, and let $\epsilon_k$ be bit $k$ of $m$.
Then
\[
 A(n)=2A(m)+\sum_{k=0}^{d-1}A(m+2^k(1-\epsilon_k)).
\]

\paragraph{Proof.}
Put $c_k=1+\operatorname{wt}(\lfloor m/2^k\rfloor)$ for $0\le k\le d$.
Then $c_d=1$, $c_k-c_{k+1}=\epsilon_k$.  Define
\[
 P_{-1}=\prod_{j=0}^{d-1}c_j=A(m),\qquad
 P_k=\prod_{j=0}^k(c_j+1)\prod_{j=k+1}^{d-1}c_j.
\]
Now $A(n)=2P_{d-1}$, and $P_k=A(m+2^k)$ whenever $\epsilon_k=0$.
The identity
\[
 (c_{k+1}+1)P_k-(c_k+1)P_{k-1}=(1-\epsilon_k)P_k
\]
telescopes to
\[
 A(n)=(2+\operatorname{wt}(m))A(m)
       +\sum_{\epsilon_k=0}A(m+2^k),
\]
which is the claimed expression.\hfill$\square$

\paragraph{Theorem (submask transform).}
Let $C(n)=\mathrm{A329369}(n)$.  Then
\[
 A(n)=\sum_{k\subseteq_2 n}C(k),\qquad
 C(n)=\sum_{k\subseteq_2 n}(-1)^{\operatorname{wt}(n)-\operatorname{wt}(k)}A(k),
\]
where $k\subseteq_2n$ means each one-bit of $k$ is a one-bit of $n$.
In addition, for $L,r\ge0$,
\[
 \mathrm{A008292}(L+1,r+1)
 =\sum_{\substack{0\le k<2^L\\\operatorname{wt}(k)=r}}C(k).
\]

\paragraph{Proof.}
Choose $L$ with $n<2^L$ and put $M=L+1$.  Let the allowable excedance set
$S\subseteq\{1,\ldots,M-1\}$ contain $M-i$ exactly when bit $i-1$ of $n$
is one.  Write $r=|S|$ and list its complement as
$t_1<\cdots<t_{M-r}$.  A permutation with excedance set contained in $S$
must assign restricted row $t_j$ a column at most $t_j$.  Processing these
nested allowed sets in order gives exactly
\[
 r!\prod_{j=1}^{M-r}(t_j-j+1)
\]
permutations, the remaining $r$ rows being unrestricted.  Put
$b_i=|S\cap\{1,\ldots,i\}|$.  The last product equals
$\prod_{i=1}^{M-1}(1+b_i)=A(n)$: the factors at $i\in S$ are
$2,3,\ldots,r+1$, whereas the terminal restricted factor at $i=M$ is
$r+1$ and the unrestricted-row factor is $r!$.
Every exact excedance subset of $S$ corresponds to one submask of $n$,
giving the first identity.  The second is Boolean-lattice M\"obius
inversion.  The equivalent binomial-parity form follows by expanding
$(1+x)^n=\prod_{\text{bit }j\text{ of }n=1}(1+x^{2^j})$ over $\mathbb F_2$.
Finally, summing over all masks with $r$ one-bits counts permutations of
$L+1$ elements with $r$ excedances, namely the indicated Eulerian number.
This last statement is conjecture 7 on the A329369 page.\hfill$\square$

\subsection{First occurrences in a prime-exponent encoding: A334109 and A334110}

Sources: \url{https://oeis.org/A334109} and
\url{https://oeis.org/A334110}.  The same first-occurrence conjecture appears
in both entries.

Let $p_i$ denote the $i$th prime.  Define $f(2)=0$, $f(p)=1+f(p-1)$ for odd
primes, and extend $f$ completely additively, with $f(1)=0$.
This is A329697: if $p$ is the largest prime factor of a number which is
not a power of two, then
$f(n(p-1)/p)=f(n)-1$, exactly its printed stopping recursion.
Define
\[
 Q(e)=\prod_{\text{bit }j\text{ of }e=1}p_{j+1},\qquad F(e)=f(Q(e)).
\]
The displayed additive description of A334109 is
\[
 H(N)=\sum_i2^{i-1}F(e_i),\qquad N=\prod_i p_i^{e_i}.
\]
The sequence A334110 is $C(k)=Q(k)^2$.

\paragraph{Theorem.}
For every $k\ge0$, the least positive integer $N$ with $H(N)=k$ is $C(k)$.

\paragraph{Proof.}
For an integer $v\ge0$, put
$h(v)=\sum_{\text{bit }j\text{ of }v=1}(j+1)$.
Binary carrying never increases this weight, since two bits of weight
$j+1$ are replaced by one of weight $j+2\le2(j+1)$.  Thus $h$ is
subadditive for ordinary addition, and $h(v)\le v$.

We have $f(n)\le\log_2n$.  By complete additivity it suffices to check
primes inductively: for odd $p$,
\[
 f(p)=1+f((p-1)/2)\le\log_2(p-1)<\log_2p,
\]
while $p=2$ is immediate.  Bertrand's postulate implies
$p_{j+1}<2^{j+1}$, so $f(p_{j+1})\le j$ for $j\ge1$.  Consequently
\begin{align*}
 h(F(e))&\le\sum_{\text{bit }j\text{ of }e=1}h(f(p_{j+1}))\\
 &\le\sum_{\substack{j\ge1\\\text{bit }j\text{ of }e=1}}j
 \le\sum_{\substack{j\ge1\\\text{bit }j\text{ of }e=1}}2^{j-1}
 =\lfloor e/2\rfloor.
\end{align*}
The bit $j=0$ contributes zero because $f(2)=0$.

The function $C$ is submultiplicative under addition:
$C(a+b)\le C(a)C(b)$.  Indeed, a binary carry replaces $p_i^4$ by
$p_{i+1}^2$, which does not increase the product since
$p_{i+1}<2p_i\le p_i^2$.  Bertrand also gives, for $v\ge0$,
\[
 C(2^{i-1}v)=\prod_{\text{bit }j\text{ of }v=1}p_{i+j}^2
 \le p_i^{2h(v)},
\]
because $p_{i+j}\le p_i^{j+1}$ (use $p_{i+j}<2^jp_i$ for $j>0$).
It follows that
\[
 C(H(N))\le\prod_i C(2^{i-1}F(e_i))
 \le\prod_i p_i^{2h(F(e_i))}\le\prod_i p_i^{e_i}=N.
\]
Conversely $F(2)=f(3)=1$, so $H(C(k))=k$ for every $k\ge0$.
The lower bound is therefore attained, proving minimality, including
$k=0$ where $C(0)=1$.\hfill$\square$

\subsection{Row divisibility for functional digraphs: A216971}

Source: \url{https://oeis.org/A216971}.  Jon Perry's 2012 conjecture states
that $n$ divides every entry $T(n,k)$, where $T(n,k)$ counts functions on
$n$ labelled elements with exactly $k$ nonrecurrent vertices whose images
are recurrent.  Here $n\ge1$ and $0\le k\le n-1$.

\paragraph{Theorem.}
One has $n\mid T(n,k)$ throughout the stated range.

\paragraph{Proof.}
Let $U(x)$ be the rooted-tree series, characterized by $U=xe^U$.
The entry's bivariate exponential generating function is
\[
 \sum_{n\ge1}\sum_{k=0}^{n-1}T(n,k)\frac{x^n}{n!}y^k
 =\frac1{1-xe^{yU(x)}}-1
 =\sum_{r\ge1}x^r e^{ryU(x)}.
\]
This identity also follows by viewing a function as a permutation of its
$r$ cyclic vertices, with rooted trees attached to those vertices; $y$
marks each attached tree's root.  At $k=0$ the coefficient is $T(n,0)=n!$.
For $k\ge1$, put $m=n-r$.  Lagrange inversion gives
\[
 [x^m]U(x)^k
 =\frac{k}{m}[u^{m-k}]e^{mu}
 =\frac{k\,m^{m-k-1}}{(m-k)!}\qquad(m\ge k).
\]
Thus $T(n,k)$ is a sum, over $1\le r\le n-k$, of terms
\[
 \frac{n!\,r^k m^{m-k-1}}{(k-1)!(m-k)!},\qquad m=n-r.
\]
For $m>k$, dividing this term by $n$ leaves an integer: its factorial
factor is
\[
 \frac{(n-1)!}{(k-1)!(m-k)!}
 =\binom{n-1}{k-1}\frac{(n-k)!}{(m-k)!}\in\mathbb Z,
\]
and the remaining powers have nonnegative exponents.  At the endpoint
$m=k$, the coefficient $[x^k]U^k$ equals one, so the divided term is
$(n-1)!r^k/k!$, again an integer because $k\le n-1$.
Every summand is therefore divisible by $n$.\hfill$\square$

\subsection{Records of zeros in XOR triangles: A334592}

Source: \url{https://oeis.org/A334592}.  For a binary word of length $L$,
form the triangle by replacing each successive pair by its XOR until a
single entry remains.  The entry counts zeros when the first row is the
ordinary binary expansion of a positive integer $n$.

\paragraph{Theorem.}
The strict record indices of A334592 are exactly $n=2^{L-1}$, $L\ge1$;
the record value is $L(L-1)/2$.

\paragraph{Proof.}
Every such triangle with a nonzero first row of length $L$ contains at
least $L$ ones.  This follows by induction on $L$.  If its second row is
nonzero, the triangle beneath the first row has at least $L-1$ ones, and
the first row supplies another.  If the second row is zero, the first
row is constant; since it is nonzero, it already contains $L$ ones.
The assertion at $L=1$ is immediate.

There are $L(L+1)/2$ positions in total, so at most $L(L-1)/2$ zeros.
The word $10\cdots0$, representing $2^{L-1}$, has one one in every row
and attains this bound.  It is the first integer with $L$ binary digits.
For $L\ge2$ the bound exceeds $(L-1)(L-2)/2$, the maximum for shorter
words; no later word of the same length can exceed it.  This proves
the assertion, with $n=1$ as the initial record.\hfill$\square$

\subsection{Uniqueness for the XOR of consecutive squares: A224515 and A224511}

Sources: \url{https://oeis.org/A224515} and
\url{https://oeis.org/A224511}.  The former separately conjectures the
existence and uniqueness of $k\ge0$ with
$k^2\mathbin{\mathrm{XOR}}(k+1)^2=(2n+1)^2$.
Its existence assertion is explicitly acknowledged as proved in 2026;
see Tom Adamczewski, \emph{OEIS Open}, Appendix A,
\url{https://arxiv.org/abs/2608.11941}, and the AlphaProof Nexus file linked
from the entry.  We count only the still separately labelled uniqueness
assertion.  The resulting permutation assertion in A224511 is an
equivalent consequence and should not be counted as another independent
problem.  The existence argument is included for completeness and carries
no priority claim.

\paragraph{Theorem.}
For every positive integer $S\equiv1\pmod4$ there is exactly one $k\ge0$
such that $k^2\mathbin{\mathrm{XOR}}(k+1)^2=S$.
In particular A224515's least index is its only index, and A224511 is
a permutation of the positive odd integers.

\paragraph{Proof of uniqueness.}
Fix an odd $S=2M+1$.  Define, for every signed integer $k$,
\[
 F(k)=k+(k^2\mathbin{\mathrm{AND}}S).
\]
Only nonnegative squares enter the bitwise operation.  The identity
$u\mathbin{\mathrm{XOR}}S=u+S-2(u\mathbin{\mathrm{AND}}S)$ gives
\[
 k^2\mathbin{\mathrm{XOR}}(k+1)^2=S\quad\Longleftrightarrow\quad F(k)=M.
\]
Suppose $k\ne l$ have the same parity, and put $t=v_2(k-l)\ge1$.
Since $k+l$ is even, $k^2\equiv l^2\pmod{2^{t+1}}$.
Bitwise AND with a fixed mask preserves congruences modulo powers of two,
so
\[
 F(k)-F(l)\equiv k-l\not\equiv0\pmod{2^{t+1}}.
\]
Thus $F$ is injective on each parity class of signed integers.
If $F(l)=M$, then $-l-1$ is also a root, because it simply interchanges
the two squares in the XOR equation.  Two different nonnegative roots of
the same parity contradict injectivity directly.  If their parities
differ, reflect one: $k$ and $-l-1$ are distinct roots of the same parity,
also impossible.  This proves uniqueness for every odd target $S$.

\paragraph{Existence and the permutation consequence.}
Now assume $S\equiv1\pmod4$, so $M$ is even.  Choose $h$ with $2^h>S$.
There is an even $k$ with $0\le k<2^h$ and $F(k)\equiv M\pmod{2^h}$:
start with $k=0$ modulo two, and successively choose whether to add $2^i$
at stage $i\ge1$.  Indeed,
\[
 (k+2^i)^2\equiv k^2\pmod{2^{i+1}},\qquad
 F(k+2^i)\equiv F(k)+2^i\pmod{2^{i+1}},
\]
so exactly one choice matches the next target bit.
Since $0\le k^2\mathbin{\mathrm{AND}}S\le S<2^h$, this solution satisfies
$F(k)=M$ or $F(k)=M+2^h$.  In the latter case put $q=k-2^h<0$.
The squares $q^2$ and $k^2$ agree modulo $2^h$, and the mask $S$ has no
bits beyond this range, so $F(q)=F(k)-2^h=M$.  Reflecting gives the
nonnegative solution $-q-1$.

Every odd square is $1$ modulo eight and hence qualifies for this
existence theorem.  Conversely, the XOR of consecutive squares is always
odd.  Existence and uniqueness therefore say that each positive odd
square root is listed exactly once by A224511.\hfill$\square$

\clearpage
% Proof packet: oeis_floor_proofs.tex
\section{Arithmetic-progression floors and symmetric-function recurrences}
All sequence citations in this section refer to the primary OEIS entry at
\texttt{https://oeis.org/}\textit{A-number}, inspected on 20 October 2026.
The words ``conjecture'' and ``empirical'' were still present in the cited
statements. No priority beyond these inspected entries is asserted.

\subsection{One theorem for eight periodic symmetric-function entries}
Let $c_1,\ldots,c_L$ be nonzero numbers and repeat this block periodically.
Write $A_n=e_{n-1}(c_1,\ldots,c_n)$ for $n\geq1$ and put $A_0=0$.
Set $P=\prod_{j=1}^Lc_j$, $S=\sum_{j=1}^Lc_j^{-1}$, and
$p_r=\prod_{j=1}^rc_j$, $s_r=\sum_{j=1}^rc_j^{-1}$, with
$p_0=1,s_0=0$. Then, for $n=qL+r$, $0\leq r<L$,
\[
A_{qL+r}=P^qp_r(qS+s_r),\qquad
\sum_{n\geq0}A_nx^n=
\sum_{r=0}^{L-1}p_rx^r
\left(\frac{s_r}{1-Px^L}+\frac{SPx^L}{(1-Px^L)^2}\right).
\]
Consequently $A_n=2PA_{n-L}-P^2A_{n-2L}$ for $n\geq2L$.
\begin{proof}
Each product in $e_{n-1}$ omits exactly one factor, so
$e_{n-1}(c_1,\ldots,c_n)=(\prod c_j)\sum c_j^{-1}$.
Split the product and reciprocal sum into full periods and a prefix.
Summing $\sum qz^q=z/(1-z)^2$ proves the generating function.
Its numerator has degree at most $2L-1$, which proves the recurrence.
\end{proof}

This proves the conjectured recurrence and generating function in each of
the following eight distinct entries. The first has an empirical generating
function due to Colin Barker (2012) and a conjectured closed form due to
Bruno Berselli; the other conjectures are due to
R.~J. Mathar (2013). In the last column $N(x)$ means that the generating
function is $xN(x)/(1-Px^L)^2$.
\begin{longtable}{lll p{0.47\textwidth}}
\toprule
Entry&Block&$P$&$N(x)$\\\midrule
A203150&$(1,2)$&2&$1+3x+x^2$\\
A203164&$(4,1,2,3)$&24&$1+5x+14x^2+50x^3+176x^4+80x^5+64x^6$\\
A203165&$(3,4,1,2)$&24&$1+7x+19x^2+50x^3+126x^4+432x^5+144x^6$\\
A203166&$(1,2,3,4,5)$&120&$1+3x+11x^2+50x^3+274x^4+154x^5+188x^6+324x^7+576x^8$\\
A203167&$(2,2,1)$&4&$1+4x+8x^2+12x^3+16x^4$\\
A203231&$(3,1)$&3&$1+4x+9x^2$\\
A203234&$(1,1,1,2)$&2&$1+2x+3x^2+7x^3+5x^4+3x^5+x^6$\\
A203235&$(1,1,1,3)$&3&$1+2x+3x^2+10x^3+7x^4+4x^5+x^6$\\
\bottomrule
\end{longtable}
The table follows by finite substitution into the displayed formula,
not by fitting initial terms. For A203150 the same theorem gives
$A_{2q}=3q2^{q-1}$ and $A_{2q+1}=(3q+2)2^{q-1}$; substituting even
and odd $n$ into Berselli's printed expression with $r=\sqrt2$ gives
exactly these two formulas, so that conjecture is proved as well.
Related entries A203160--A203163 already
state their generating functions as facts, and A203233 explicitly records
Robert Israel's 2017 verification; these were excluded from the count.

\subsection{Two nonperiodic symmetric-function entries}
\paragraph{A203157.}
For the triangular numbers $T_j=j(j+1)/2$,
\[
A_n=e_{n-1}(T_1,\ldots,T_n)
=\frac{n!(n+1)!}{2^n}\sum_{j=1}^n\frac2{j(j+1)}
=\frac{n(n!)^2}{2^{n-1}}.
\]
The reciprocal sum telescopes to $2n/(n+1)$. Thus
$2(n-1)A_n=n^3A_{n-1}$ for $n\geq2$, proving the recurrence
conjectured by Mathar on 1 October 2016.

\paragraph{A203159.}
The same omitted-factor identity gives
$A_n=2^{n-1}n!H_n$ for the list $2,4,\ldots,2n$.
Since $A_n-2nA_{n-1}=2^{n-1}(n-1)!$, subtracting
$2(n-1)$ times the same equation at $n-1$ gives
\[
A_n-(4n-2)A_{n-1}+4(n-1)^2A_{n-2}=0\quad(n\geq3).
\]
This proves Mathar's 1 October 2016 conjecture.

\subsection{Three explicit floor recurrences}
For a finite list with power sums $p_j$, Newton's identities give
$e_2=(p_1^2-p_2)/2$ and $e_3=(p_1^3-3p_1p_2+2p_3)/6$.
Apply these identities to the indicated arithmetic progressions.

\paragraph{A024220.}
For the first $n+2$ positive integers congruent to $1\pmod3$,
\[
\frac{e_3}{e_1}=\frac{n(n+1)(9n^2+21n-2)}{24}=Q(n).
\]
The fractional part of $Q(n)$ is $1/3$ if $n\equiv1\pmod3$,
and zero otherwise. This is checked in the three residue classes;
alternatively the numerator modulo 24 has these values.
Thus $A_n=Q(n)-\frac13\mathbf1_{n\equiv1\ (3)}$ for every $n\geq1$.
Summation yields
\[
\sum_{n\geq1}A_nx^n=
\frac{x(2+11x+7x^2+8x^3-x^4)}{(1+x+x^2)(1-x)^5}.
\]
The denominator proves Mathar's conjectured signature
$(4,-6,5,-5,6,-4,1)$, valid for $n\geq8$.

\paragraph{A024399.}
For the first $n+2$ positive integers congruent to $2\pmod3$,
\[
\frac{e_3}{e_1}=\frac{n(n+1)(9n^2+33n+22)}{24}=Q(n).
\]
Its fractional part has the same three-class rule as A024220. Hence
\[
\sum_{n\geq1}A_nx^n=
\frac{x(5+11x+7x^2+5x^3-x^4)}{(1+x+x^2)(1-x)^5}.
\]
This proves the other conjectured recurrence with the same signature.

\paragraph{A024385.}
For the first $n+1$ positive integers congruent to $1\pmod4$,
\[
\frac{e_2}{e_1}=
\frac{n(12n^2+8n-5)}{6(2n+1)}
=n^2+\frac{n-3}{6}+\frac1{2(2n+1)}.
\]
For $n\geq2$ the positive last term is less than $1/6$, while
the preceding polynomial has fractional part in $\frac16\mathbb Z$.
The case $n=1$ is checked directly. Therefore
$A_n=n^2+\lfloor(n-3)/6\rfloor$ for all $n\geq1$.
Its generating function is
\[
\frac{x^2(3+3x+x^2+2x^3+2x^4+2x^5-x^6)}{(1-x)^2(1-x^6)}.
\]
This proves the conjectured signature $(2,-1,0,0,0,1,-2,1)$,
valid for $n\geq9$.

\subsection{Two empirical generating functions proved by finite corrections}
\paragraph{A024390.}
Use the first $n+3$ numbers congruent to $1\pmod4$ and take $e_4/e_3$.
Newton's identities through degree four give
\[
\frac{e_4}{e_3}=Q(n)+R(n),\quad
Q(n)=\frac{6n^2+9n-13}{12},\quad
R(n)=\frac{692n^2+3128n+2925}{60(2n+5)(4n^2+16n+9)}.
\]
The fractional parts of $Q(n)$ for $n\equiv1,2,3,0\pmod4$ are
$1/6,5/12,2/3,11/12$. For $n\geq16$, $0<R(n)<1/12$:
after clearing positive denominators the upper bound is equivalent to
$20n^3-216n^2-1319n-1350>0$, whose expansion at $n=16+u$ is
$20u^3+744u^2+7129u+4170>0$.
Exact substitution at $1\leq n\leq15$ gives
\[
A_n=\lfloor Q(n)\rfloor+\mathbf1_{n\in\{4,8,12\}}.
\]
Summing this identity proves Colin Barker's empirical formula (2014):
\[
\sum_{n\geq1}A_nx^n=
\frac{x^2(x^{16}-2x^{15}+x^{14}+x^4-2x^2-x-2)}
{(x-1)^3(x+1)(x^2+1)}.
\]

\paragraph{A024403.}
For the first $n+3$ numbers congruent to $2\pmod3$, the ratio $e_4/e_3$ is
\[
Q(n)+R(n),\quad Q(n)=\frac{3n^2+7n-4}{8},\quad
R(n)=\frac{153n^2+1002n+1600}{5(3n+10)(9n^2+51n+64)}.
\]
All fractional parts of $Q(n)$ lie in $\{0,1/4,1/2,3/4\}$.
For $n\geq3$, $0<R(n)<1/4$ follows by clearing denominators:
$135n^3+603n^2-498n-3200>0$ expands at $n=3+u$ as
$135u^3+1818u^2+6765u+4378$.
The two initial cases give
$A_n=\lfloor Q(n)\rfloor+\mathbf1_{n\in\{1,2\}}$.
Summation proves Barker's empirical generating function
\[
\frac{x(x^2-x+1)(x^8-x^7-x^5+x^4-x^3-x-1)}
{(x-1)^3(x^2+1)(x^4+1)}.
\]

\subsection{A024389: exact disproof of an empirical generating function}
For $S=\{1,5,\ldots,4n+9\}$, Newton's identities give
\[
\frac{e_4(S)}{e_2(S)}=
\frac{n(n+1)(240n^4+1920n^3+4760n^2+3072n-1217)}
{60(12n^2+56n+59)}.
\]
At $n=39$ this is $17588398646/20495$, with floor $858179$.
The coefficient of $x^{39}$ in Barker's displayed empirical function
is $858180$. The latter assertion is an exact rational coefficient
calculation: if $G=N/D$ is that function, generate $c_0,c_1,\ldots,c_{39}$
using $D_0c_n=N_n-\sum_{j=1}^{\deg D}D_jc_{n-j}$.
The accompanying verification script records the complete polynomial
$N,D$ and reproduces the contradiction with integer arithmetic.
Thus the empirical generating function, as printed, is false.

\subsection{A024388: index repair, counted separately}
The entry's title specifies the first $n+2$ numbers $1,5,9,\ldots$,
but its data and its cross-reference formula use the first $n+3$.
At $n=1$ the literal title gives $\lfloor45/59\rfloor=0$;
the displayed first term is $3$. We do not silently conflate these.
For the data-compatible interpretation one obtains
\[
\frac{e_3}{e_2}=
Q(n)+R(n),\quad Q(n)=\frac{18n^2+51n+11}{27},\quad
R(n)=\frac{2(118n+283)}{27(12n^2+56n+59)}.
\]
The fractional part of $Q(n)$ is periodic modulo 9, with numerators
$26,23,2,17,14,20,8,5,11$ over 27 in classes $1,\ldots,9$.
For $n\geq18$, $R(n)<1/27$ because
$4n^2-60n-169=4(n-18)^2+84(n-18)+47>0$.
Checking $1\leq n\leq17$ exactly yields
\[
A_n=\lfloor Q(n)\rfloor+\mathbf1_{n\in\{1,2,10\}}.
\]
This identity proves the printed empirical generating function for the
displayed data, once the title's index is repaired. This is listed as an
index-repair result, separate from unqualified confirmations.

\subsection{A173693: a ceiling preserves a finite recurrence}
The source specifies $b_0,b_1,b_2,b_3,b_4=0,0,0,0,1$ and
$b_n=b_{n-1}+b_{n-2}-b_{n-3}+b_{n-5}$, with $A_n=\lceil b_n/2\rceil$.
The parity sequence has the repeating block
\[
0000110010011111011100010101101.
\]
For a finite certificate, apply the recurrence modulo 2 through $n=35$:
$(b_{31},\ldots,b_{35})\bmod2=(0,0,0,0,1)$, exactly the initial
state. Determinism then proves period 31 for all subsequent terms.
Writing $r_n=b_n\bmod2$ gives $2A_n=b_n+r_n$.
If $E$ denotes backward shift, $1-E-E^2+E^3-E^5$ annihilates $b$,
and $1-E^{31}$ annihilates $r$. Therefore their product annihilates $A$,
proving Mathar's conjecture for $n\geq36$:
\begin{align*}
A_n={}&A_{n-1}+A_{n-2}-A_{n-3}+A_{n-5}+A_{n-31}\\
&-A_{n-32}-A_{n-33}+A_{n-34}-A_{n-36}.
\end{align*}

\subsection{A227012: the harmonic mean between consecutive cubes}
Let $M_n=(3n^2-3n+1)/(H_{n^3}-H_{(n-1)^3})$.
Then
\[
\lfloor M_n\rfloor=
\frac{2n^3-3n^2+n+2}{2}+
\left\lfloor\frac{\max(0,n-3)}4\right\rfloor\qquad(n\geq1).
\]
This proves Ramaharo's 2018 conjecture and, by summation, the entry's
conjectured ordinary generating function and seven-term recurrence.
Here is an exact inequality certificate, not an asymptotic argument.
Put $a=(n-1)^3,b=n^3,\Delta=b-a$ and
$Q=(4n^3-6n^2+3n+1)/4$.
We prove $Q<M_n<Q+1/4$ for $n\geq3$.

The standard signed remainder bound for the digamma expansion gives
\[
H_m=\log m+\gamma+\frac1{2m}-\frac1{12m^2}+\rho_m,
\qquad0<\rho_m<\frac1{120m^4}.
\]
See NIST DLMF, 5.11.2 and 5.11(ii),
\texttt{https://dlmf.nist.gov/5.11}.
Let $S=\sum_{j=1}^4( jn^j)^{-1}$ and
$C=(1/b-1/a)/2-(1/b^2-1/a^2)/12$. The elementary positive
logarithm series and its geometric remainder imply
\[
L=3S+C-\frac1{120a^4}<H_b-H_a<
U=3\left(S+\frac1{5n^4(n-1)}\right)+C+\frac1{120b^4}.
\]
For a reproducible polynomial positivity certificate define
\[
P_-(n)=480n^{12}(n-1)^6(\Delta-QU),\quad
P_+(n)=480n^6(n-1)^{12}((Q+1/4)L-\Delta).
\]
Both are polynomials of degree 17. The ascending coefficient lists of
$P_-(3+u)$ and $P_+(3+u)$ are given below; every coefficient is positive.
\begin{center}\small
\begin{tabular}{r|r|r}
$j$&$[u^j]P_-(3+u)$&$[u^j]P_+(3+u)$\\\hline
0&3559568768&21810535\\
1&22548059346&287558715\\
2&67082168268&1446297165\\
3&124564097140&4118075599\\
4&161774552436&7740002643\\
5&156012626130&10389087405\\
6&115733015648&10424077105\\
7&67475843205&8039537295\\
8&31310918742&4846668528\\
9&11632216336&2304171826\\
10&3460298664&865875390\\
11&819631074&256192530\\
12&152703988&59035670\\
13&21908256&10384770\\
14&2337384&1347990\\
15&174738&121770\\
16&8172&6840\\
17&180&180
\end{tabular}
\end{center}
Hence $\Delta/(Q+1/4)<L<H_b-H_a<U<\Delta/Q$ and the claimed
bounds for $M_n$ follow. Since $Q\in\frac14\mathbb Z$, they imply
$\lfloor M_n\rfloor=\lfloor Q\rfloor$ for $n\geq3$.
The cases $n=1,2$ give $1,4$ directly. Rearrangement of $\lfloor Q\rfloor$
gives the stated conjectured formula. Explicitly the generating function is
\[
\sum_{n\geq1}\lfloor M_n\rfloor x^n=
\frac{x(1+x+7x^2+6x^3+5x^4+5x^5-2x^7+x^8)}
{(1-x)^3(1-x^4)}.
\]
Only the cube-block conjectures are proved here; the much broader
arbitrary-polynomial assertion in the entry is not inferred from this case.

\subsection{Five further polynomial-block harmonic means}
\label{round1:sec:fiveharmonicmeans}
For each of the following five polynomials, put
\[
 A_n=\left\lfloor\frac{g(n)-g(n-1)}{H_{g(n)}-H_{g(n-1)}}\right\rfloor,
 \qquad n\geq1.
\]
The entries A227013--A227017 conjecture that these are linear recurrence
sequences and supply empirical generating functions. We prove their
generating functions exactly, with the indexing qualifications stated
below. A particularly compact form of the answers is
\begin{equation}\label{round1:eq:fiveharmonicanswers}
 A_n=\lfloor Q(n)\rfloor+\mathbf1_{n\in C}.
\end{equation}
Here are all the polynomials and exceptional sets:
\begin{gather*}
\mathrm{A227013}:\quad g(n)=n^4,\quad C=\{1,2\},\qquad
 Q(n)=n^4-2n^3+\frac53n^2-\frac23n+\frac12;\\
\mathrm{A227014}:\quad g(n)=n^5,\quad C=\{1,2\},\\
 Q(n)=n^5-\frac52n^4+\frac{35}{12}n^3-\frac{15}{8}n^2
             +\frac{85}{144}n+\frac{125}{288};\\
\mathrm{A227015}:\quad g(n)=n^3+n^2+n+1,\quad C=\{3\},\\
 Q(n)=n^3-\frac12n^2+\frac34n+1;\\
\mathrm{A227016}:\quad g(n)=\frac{n(n+1)(n+2)}6,\quad C=\{1\},\\
 Q(n)=\frac16n^3+\frac14n^2-\frac1{24}n+\frac7{16};\\
\mathrm{A227017}:\quad g(n)=\frac{n(3n-1)}2,\quad C=\{1\},\qquad
 Q(n)=\frac32n^2-2n+1.
\end{gather*}

\paragraph{Uniform rigorous bounds.}
Let $a=g(n-1)$, $b=g(n)$, $\Delta=b-a$, and
\[
 C_0=\frac12\left(\frac1b-\frac1a\right)
     -\frac1{12}\left(\frac1{b^2}-\frac1{a^2}\right).
\]
We only use this notation when $a>0$, handling the initial cases
separately. Choose a rational lower approximation $S$ to $\log(b/a)$
and a positive rational tail bound $T$ so that
$S<\log(b/a)<S+T$.
For $g(n)=n^d$ use
\[
 S=d\sum_{j=1}^{K}\frac1{jn^j},\qquad
 T=\frac{d}{(K+1)n^K(n-1)}.
\]
This follows by expanding $-d\log(1-1/n)$ and bounding each
remaining denominator $j$ below by $K+1$.
For the other three polynomials put $t=\Delta/(a+b)$ and use
\[
 S=2\sum_{j=0}^{K-1}\frac{t^{2j+1}}{2j+1},\qquad
 T=\frac{2t^{2K+1}}{(2K+1)(1-t^2)}.
\]
Here $0<t<1$; this is the positive series for
$2\operatorname{arctanh}t=\log(b/a)$ with the same geometric-tail
argument. The signed harmonic-number bound used in A227012 now gives
\[
 L=S+C_0-\frac1{120a^4}<H_b-H_a<
 U=S+T+C_0+\frac1{120b^4}.
\]

For each row in the next table, the exact polynomial certificates below
prove
\[
 E_-:=\Delta-(Q+\ell)U>0,\qquad
 E_+:=(Q+h)L-\Delta>0\qquad(n\geq N).
\]
The polynomials $a,b,\Delta,Q+\ell$ have positive constant terms and
nonnegative coefficients after substituting $n=N+u$. Thus they are
positive throughout the stated range. It follows in particular that
$L>0$, and division by positive quantities gives
\[
 Q(n)+\ell<\frac{\Delta}{H_b-H_a}<Q(n)+h.
\]
The last column gives the period of the fractional part of $Q$.
\begin{center}
\begin{tabular}{lrrrrr}
\toprule
Entry&$N$&$K$&$\ell$&$h$&fractional period $J$\\\midrule
A227013&4&6&$0$&$1/6$&3\\
A227014&30&7&$-1/288$&$1/288$&144\\
A227015&4&3&$0$&$1/4$&4\\
A227016&3&3&$-1/16$&$1/16$&8\\
A227017&2&3&$0$&$1/2$&2\\\bottomrule
\end{tabular}
\end{center}
For A227013 the fractional parts belong to $\{1/2,5/6\}$.
For A227014 they are odd multiples of $1/288$.
For A227015 they belong to $\{0,1/4,1/2,3/4\}$.
For A227016 they are odd multiples of $1/16$.
For A227017 they belong to $\{0,1/2\}$.
In each case $\lfloor Q\rfloor\leq Q+\ell$
and $Q+h\leq\lfloor Q\rfloor+1$. The strict inequalities for the
harmonic mean therefore place it strictly between these consecutive
integers, including when an endpoint of the bound is an integer. This proves
$A_n=\lfloor Q(n)\rfloor$ for $n\geq N$.
The stated periods can be checked globally by expanding
$Q(n+J)-Q(n)$ in the integer-valued basis
$\binom n0,\binom n1,\ldots$; all its coefficients are integers.
This is a polynomial identity, not an extrapolation from observed terms.

\paragraph{Explicit polynomial positivity certificates.}
For each entry let $P_-=D_-E_-$ and $P_+=D_+E_+$, using the following
positive denominators. In the A227015 row abbreviate
$B=2n^3-n^2+3n+1$; in the A227017 row put $V=3n^2-4n+2$.
\begin{align*}
\mathrm{A227013}:\quad
 D_-&=5040n^{16}(n-1)^8,\quad
 D_+=360n^8(n-1)^{16};\\
\mathrm{A227014}:\quad
 D_-&=120960n^{20}(n-1)^{10},\quad
 D_+=120960n^{10}(n-1)^{20};\\
\mathrm{A227015}:\quad
 D_-&=3360n^2(n+1)^4(n^2+1)^4(n^2-2n+2)^2B^5,\\
 D_+&=480n^4(n+1)^2(n^2+1)^2(n^2-2n+2)^4B^5;\\
\mathrm{A227016}:\quad
 D_-&=560n^4(n-1)^2(n+1)^4(n+2)^4(2n+1)^5,\\
 D_+&=40n^4(n-1)^4(n+1)^4(n+2)^2(2n+1)^5;\\
\mathrm{A227017}:\quad
 D_-&=105n^4(n-1)^2(3n-4)^2(3n-1)^4V^4,\\
 D_+&=15n^2(n-1)^4(3n-4)^4(3n-1)^2V^5.
\end{align*}
All factors are positive on the applicable ranges: in particular
$n^2-2n+2=(n-1)^2+1$, $B=n^2(2n-1)+3n+1$, and
$V=3(n-2/3)^2+2/3$ are positive there.
Exact expansion of $P_\pm(N+u)$ gives the following certificate.
Each ``minimum'' is over every coefficient, including the constant
coefficient; in particular no omitted coefficient is zero or negative.
\begin{center}
\begin{tabular}{lrrrr}
\toprule
Entry&$\deg P_-$&minimum coefficient&$\deg P_+$&minimum coefficient\\\midrule
A227013&23&1568&23&128\\
A227014&29&2100&29&2100\\
A227015&32&40320&32&5760\\
A227016&18&3360&18&240\\
A227017&17&2480058&21&2657205\\\bottomrule
\end{tabular}
\end{center}
The exact integer coefficient vectors, as well as the original numerator
and denominator polynomials, are supplied in
\texttt{oeis\_\allowbreak more\_\allowbreak means\_\allowbreak certificates.json}.
The companion \texttt{oeis\_\allowbreak more\_\allowbreak means\_\allowbreak verification.py} reconstructs
$S,T,L,U,E_\pm$ from the displayed definitions, checks the polynomial
identities and all coefficients with exact rational arithmetic, and
checks denominator positivity. Thus the table is a finite certificate
for every real $n\geq N$, rather than a finite test of the sequence.

\paragraph{Finite initial cases.}
For $n<N$, exact evaluation gives, respectively,
\[
(1,6,40),\qquad
(2,8,26),\qquad(1,2),\qquad(1)
\]
for A227013, A227015, A227016, A227017.
For A227014 the initial values through $n=29$ are
\begin{align*}
(&1,10,104,543,1883,5102,11717,23906,44626,77735,\\
 &128110,201769,305989,449428,642243,896212,1224852,\\
 &1643541,2169636,2822595,3624095,4598154,5771249,\\
 &7172438,8833478,10788947,13076362,15736301,18812521).
\end{align*}
These are certified by exact rational harmonic sums when $b\leq10000$.
For the larger finite endpoints the same $\operatorname{arctanh}$
bounds with $K=100$ enclose $\Delta/(H_b-H_a)$ between two rational
numbers having the same floor. This avoids floating-point assumptions.
Subtracting $\lfloor Q(n)\rfloor$ gives exactly the exceptional sets
in \eqref{round1:eq:fiveharmonicanswers} and finishes its proof for every $n$.

\paragraph{Generating functions and literal source qualifications.}
Write $Q(n)=\sum_{k=0}^{d}q_kn^k$, let
$\theta=x\,d/dx$, and set $f_r=Q(r)-\lfloor Q(r)\rfloor$.
The proved formulas imply the exact rational function
\begin{equation}\label{round1:eq:harmonicgf}
 G(x)=\sum_{n\geq1}A_nx^n
 =\sum_{k=0}^{d}q_k\theta^k\frac{x}{1-x}
 -\frac{\sum_{r=1}^{J}f_rx^r}{1-x^J}
 +\sum_{n\in C}x^n.
\end{equation}
All quantities in this expression are explicitly specified above.
Symbolic simplification gives exactly the empirical functions printed
in the five source entries. The accompanying verifier compares the
full rational functions, including the 151-coefficient numerator
listed by Moses for A227014, and records an exact zero difference.
Consequently a denominator for each is
$(1-x)^d(1-x^J)$, with $(d,J)=(4,3),(5,144),(3,4),(3,8),(2,2)$.
The respective numerator degrees are $9,151,10,12,5$, so the resulting
recurrences hold from $n=10,152,11,13,6$, respectively.

The source functions for A227013, A227014, A227016 and A227017
use $\sum_{n\geq1}A_nx^{n-1}$; one multiplies them by $x$ to obtain
our convention. A227015 already uses $G(x)$.
The prose recurrence signature in A227014 has sign errors;
the source's polynomial denominator and the theorem instead give
\begin{align*}
A_n={}&5A_{n-1}-10A_{n-2}+10A_{n-3}-5A_{n-4}+A_{n-5}\\
 &+A_{n-144}-5A_{n-145}+10A_{n-146}-10A_{n-147}
       +5A_{n-148}-A_{n-149}\qquad(n\geq152).
\end{align*}
Our result confirms the empirical generating function and establishes
this corrected recurrence; it does not endorse the mistyped signature.
The printed recurrence line in A227016 lacks an equality sign; the denominator above specifies its corrected recurrence. The starting ranges printed for several of these recurrences also precede the finite corrections, so the valid starting indices stated above must be used.
The printed trigonometric polynomial in A227013 is shifted by two
indices, and the printed closed form in A227017 is shifted by one
index (with the exceptional first term still separate). Equation
\eqref{round1:eq:fiveharmonicanswers} consistently uses the definitions'
offset $n=1$. We make no assertion here about additional geometric-mean
identities or the unrestricted polynomial conjecture in A227012.

\clearpage
% Proof packet: oeis_gf_proofs.tex

\section{Generating functions and exact recurrence certificates}
\label{round1:sec:gfproofs}

This section records 53 distinct OEIS entries checked on 20 October 2026.
There are 49 proved recurrence assertions, one entry whose recurrence is proved
and whose separate limit assertion is refuted (A394578), one false recurrence
(A027088), and two indexing repairs (A142243 and A059837).
The two indexing repairs should be tabulated separately from substantive closures.
The source pages still mark the targeted assertions as conjectural. That page
status does not establish first-publication priority: the results here are
independent certificates, and no submission or OEIS acceptance is claimed.
The formulas explicitly described as seed premises are existing source formulas
or definitions, not additional conjectures claimed to have been solved here.

For every entry, the certificate appendix prints the complete coefficient list
\(p_0(n),\ldots,p_r(n)\), meaning
\[
             \sum_{j=0}^{r}p_j(n)a_{n-j}+c(n)=0.
\]
Here \(c(n)=0\) except in A027175 and A027182. A stated lower bound on \(n\)
is part of the theorem. These are assertions about the actual sequence;
isolated zeros of \(p_0(n)\) do not create a problem because no division by
\(p_0(n)\) is used.

\subsection{Why the certificates prove infinitely many cases}

Let \(\theta=x\,d/dx\). For an ordinary generating function
\(F(x)=\sum_{n\geq0}a_nx^n\), extending any omitted initial indices by zero,
define
\[
 \mathcal L_p=\sum_{j=0}^{r}x^j p_j(\theta+j).
\]
Then coefficient extraction gives the exact identity
\[
 [x^n]\mathcal L_pF=\sum_{j=0}^{r}p_j(n)a_{n-j}.
 \tag{GF1}
\]
Consequently, an identity \(\mathcal L_pF=P(x)\), with \(P\) a polynomial
of degree less than \(N\), proves the recurrence for every \(n\geq N\).
For an exponential generating function \(F=\sum a_nx^n/n!\), first replace
\(p_j(n)\) by
\[
 q_j(n)=p_j(n)\prod_{h=0}^{r-j-1}(n-j-h).
 \tag{GF2}
\]
The coefficient in (GF1) is then the desired recurrence divided by
\((n-r)!\). We use this conversion only where indicated.

Every radical certificate writes
\[
  x=f(v),\qquad F=\sum_i R_i(v)\sqrt{D_i(v)}.
\]
The formal branch is the one specified by the source constant term.
The identity can be checked using only rational arithmetic, since
\[
 \theta\bigl(R_i\sqrt{D_i}\bigr)
 =\frac{f(v)}{f'(v)}
 \left(R_i'(v)+\frac{R_i(v)D_i'(v)}{2D_i(v)}\right)\sqrt{D_i(v)}.
 \tag{GF3}
\]
The appendix gives the exact resulting polynomial \(P\). All coefficients
of the nonconstant radical basis elements vanish. Thus the certificates are
explicit rational-function identities, not a test of a finite prefix of the
sequence. Clearing denominators converts each identity into a polynomial
identity over the rationals. The accompanying script
\texttt{oeis\_gf\_verify.py} performs precisely these operations using only
the accompanying exact certificate JSON, without source-page snapshots.

\subsection{Algebraic seed functions and elementary derivations}

For A174193, A247173, A247102, A234269, A268600, A268601 and A244886,
the ordinary generating function in the certificate is the entry's definition.
For A108296 and A244884 it is an existing FORMULA statement.
In particular, the A174193 result concerns its recurrence alone: the separately
conjectural interpretation involving almost-avoiding permutations is not proved.

For A346074, multiplying the defining convolution by \(x^n\) and summing gives
\(F=(1-x)^{-1}+x^5F^2\). Solving this quadratic with \(F(0)=1\) gives
the certified radical. For A268600 and A268601 use
\[
 \frac1{4-2\sqrt{1+4x}}=
 \frac1{3-4x}+\frac{\sqrt{1+4x}}{2(3-4x)}.
\]
This explains the two radical basis elements in their certificates.

Write
\[
 M(x)=\frac{1-x-\sqrt{1-2x-3x^2}}{2x^2}.
\]
For A244884, \(F=x+(1-x)M\). For A257104 the existing continued fraction is
\[
 F=\frac1{1-x-\displaystyle\frac{x^2}{1-x-
          \displaystyle\frac{x^2}{1-x+x^2(1-M)}}}.
\]
It also follows by the usual first-return decomposition at the restricted
Motzkin heights. Substitution in (GF3) proves the printed recurrences with
the bounds in the appendix. The bounds matter: the A244884 residual has
degree four and the A257104 residual has degree six.

For A099038, reindex the existing binomial sum by \(n=m+2k\).
The binomial coefficient extraction gives
\[
 F(x)=\operatorname{CT}_{z}
 \frac1{1-x(1+z)-x^2(1-z)/z}
 =\frac1{\sqrt{(1-x+x^2)^2-4x^3}}.
\]
Indeed \(\operatorname{CT}_z(A-Bz-C/z)^{-1}=(A^2-4BC)^{-1/2}\),
by expanding the geometric series and matching equal powers of \(z\).

For A337500 let \(C=C(x^2)=(1-\sqrt{1-4x^2})/(2x^2)\).
The generating function for unrestricted walks from 0 to displacement \(h\)
is \((xC)^{|h|}/\sqrt{1-4x^2}\). One proof is to solve
\(W_h=\mathbf1_{h=0}+x(W_{h-1}+W_{h+1})\), selecting the formal small
root for each half-line. Summing over \(-3\leq h\leq3\) gives
\[
 F=\frac{1+2xC+2x^2C^2+2x^3C^3}{\sqrt{1-4x^2}}.
\]
No restriction on intermediate positions is imposed.

\subsection{Rational parametrizations from Lagrange inversion}

For A127905 put \(w=x(1+w+w^3)\). Lagrange inversion gives
\(a_n=n[x^n]w\), hence \(F=xw'\). Setting \(v=w\) gives
\[
 x=\frac{v}{1+v+v^3},\qquad
 F=\frac{v(1+v+v^3)}{1-2v^3}.
\]
For A392976 the definition itself gives
\(x=v((1-v^2)^2-v)\) and \(F=v/x\).
For A186185 let \(v=T(z)\), where \(T=1+zT^3\) and \(z=x/(1-x)\).
Then
\[
 x=\frac{v-1}{v^3+v-1},\qquad F=\frac{v^2}{v^2-v+1},
\]
with \(v=1\) at \(x=0\). For A071264 the source gives
\(F=(1+xC^4)C^3\), so set \(v=C\) and \(x=(v-1)/v^2\), also about
\(v=1\). For A220433 use the existing exponential generating function:
\(x=v-v^2+v^3/6\), \(F=v\), and apply (GF2).

For A045742, the accepted binomial sum is, with \(m=n-1\),
\[
 a_n=\frac1m\sum_{k=1}^{m-1}
       k\binom{m+k-1}{k}\binom{3m}{m-1-k}.
\]
Using the derivative of \((1-u)^{-m}\), this becomes
\[
 a_n=[u^{m-1}]\frac{u}{1-u}
                   \left(\frac{(1+u)^3}{1-u}\right)^m.
\]
If \(w=x(1+w)^3/(1-w)\), Lagrange inversion applied to an antiderivative
of \(u/(1-u)\) gives
\[
 F=x^2 w'\frac{w}{1-w},\qquad
 x=f(v)=\frac{v(1-v)}{(1+v)^3},\qquad
 F=\frac{f(v)^2v}{(1-v)f'(v)}.
\]
The exact boundary is \(-2544x^3\); therefore the recurrence is proved
for \(n\geq4\). The source conjecture does not prescribe a conflicting
lower range.

For A392211 the accepted relation is
\[
 x=\frac{3+14F-20F^2-3\sqrt{1+4F-4F^2}}8.
\]
Set \(\sqrt{1+4F-4F^2}=1+vF\). Then
\[
 F=\frac{4-2v}{v^2+4},\qquad x=\frac{F(14-3v-20F)}8,
\]
about \(v=2\). Substitution into (GF3) gives the certificate in the appendix.

\subsection{Products, binomial sums, and factorial-series transforms}

A091539 has the existing product formula
\[
 a_n=\frac1{60}\prod_{j=0}^{n-1}(3j+2)
   \left(\prod_{j=0}^{n-1}3(j+1)-3\prod_{j=0}^{n-1}(3j+1)\right).
\]
The two product components have successive ratios
\(\rho(n)=3n(3n-1)\) and \(\rho(n)=(3n-1)(3n-2)\).
For each ratio the exact identity
\[
 \sum_{j=0}^{r}\frac{p_j(n)}{
                     \prod_{h=0}^{j-1}\rho(n-h)}=0
\]
holds. Multiplying by the corresponding product proves the recurrence.

The definition of A027170 and Pascal's identity give
\[
 T(n,k)=2\binom{n+2}{k+1}+\binom nk-4.
\]
Therefore, writing \(C_n=\binom{2n}n\),
\[
 \begin{split}
 \text{A027175: }a_n&=C_n\frac{9n+5}{2(n+1)}-4,\\
 \text{A027182: }a_n&=4\binom{2n+4}{n+2}
            +4\binom{2n+2}{n+1}+C_n+16n+40-72\,2^n.
 \end{split}
\]
For the second identity, symmetry of \(T(n,k)\) turns the central
self-convolution coefficient into \(\sum_{k=0}^nT(n,k)^2\).
Expanding the square and applying Vandermonde's identity yields the formula,
including the boundary corrections from \(k+1=1,\ldots,n+1\).
Now \(C_n/C_{n-1}=2(2n-1)/n\). Substitution in the respective recurrence
separates a central-binomial component, a polynomial component, and a
\(2^n\) component. Each vanishes identically; the polynomial component
includes the inhomogeneous term \(c(n)\) printed in the appendix.
The example value 120 on A027182 is inconsistent with its definition;
the actual row at \(n=2\) gives \(5^2+10^2+5^2=150\), as in its data.
This observation is not counted as a separate closure.

For A177258, A177259 and A177260 take \(q=2,3,4\), respectively, and use
the established series
\[
 Q=1+x+x^q,\qquad F_q=Q^{-1}f(x/Q),\qquad f(z)=\sum_{k\geq0}k!z^k.
\]
This is an identity of formal series, so no analytic convergence is needed.
Coefficient extraction proves \(z^2f'+(z-1)f+1=0\). The chain rule then gives
\[
 x^2QF_q'+Q\bigl(-1+x+(q-1)x^q\bigr)F_q+1-(q-1)x^q=0.
\]
Its coefficients are exactly
\[
 \begin{split}
 a_n={}&(n-1)a_{n-1}+(n-1)a_{n-2}+(q-2)a_{n-q}\\
       &+(n-1)a_{n-q-1}+(q-1)a_{n-2q},\qquad n\geq2q.
 \end{split}
\]
When two indices coincide their terms are added, as in the printed vectors.

For A351933 differentiate its defining exponential generating function
\(F=\exp(x/(1-x^2/2))\). It satisfies
\[
 (1-x^2/2)^2F'-(1+x^2/2)F=0.
\]
Taking \(4(n-1)![x^{n-1}]\) gives the displayed five-term recurrence.

For A001722 and A345651 the accepted Stirling sums give, respectively,
\[
 a_n=[u^2](5+u)^{\overline{n+2}},\qquad
 a_n=[u^4](4+u)^{\underline n}.
\]
Factoring the common rising product \((5+u)^{\overline{n-1}}\) out of
the first recurrence leaves exactly \(u^3\). Thus its \(u^2\) coefficient
is zero. Factoring \((4+u)^{\underline{n-4}}\) out of the second leaves
exactly \(u^4\). For \(n\geq9\) that common falling product contains a
factor \(u\), so its product with \(u^4\) has zero \(u^4\) coefficient.
The residual at \(n=8\) is 24, which explains the stated lower range.

\subsection{Exact elimination of factorial and parity recurrences}

Here is a reusable certificate rule. Fix \(n\) modulo 2, express each floor
at the indices \(m=n-j\) by its affine formula, and use the given short
recurrence to express all \(a_{n-j}\) in a fixed finite list of base
quantities. Normalize factorials by \((n-r)!\), using
\[
 \frac{(n-r+s)!}{(n-r)!}=\prod_{h=1}^{s}(n-r+h)
\]
and the reciprocal formula for negative \(s\). All coefficients are now
rational functions of \(n\). In each certificate the coefficient of every
base quantity in \(\sum p_j(n)a_{n-j}\) is identically zero, for both
parities. The explicit substitution rules and base lists below completely
specify the calculation; the appendix records its zero vectors.

For A121580 the existing factorial partial sum yields
\[
 a_m-a_{m-1}=\frac{(m-1)!(m^2-m+2)}2.
\]
For A180195 the defining alternating sum yields
\[
 a_m=-a_{m-1}+(m-1)!(m^2-m+1).
\]
For A121639 use the established identity
\(a_m=ma_{m-1}+(m-1)!-(m-3)!\), valid for \(m\geq4\).
Each of these three cases uses the base list
\(a_{n-r},(n-r)!\).

For A121751 put \(e_m=m\bmod2\). Its established recurrence is
\[
 a_m=2\left\lfloor\frac{m-1}2\right\rfloor a_{m-1}
 +\left(1-\left\lfloor\frac{m-1}2\right\rfloor
              \left\lfloor\frac{m-2}2\right\rfloor\right)a_{m-2}.
\]
Use \(\lfloor(m-1)/2\rfloor=(m-2+e_m)/2\) and
\(\lfloor(m-2)/2\rfloor=(m-2-e_m)/2\), with base list
\(a_{n-r},a_{n-r+1}\).

For A121752 the accepted recurrence is
\[
 \begin{split}
 a_m={}&(2m-3)a_{m-1}-(m-1)(m-3)a_{m-2}\\
 &+(m-2)!\left(m\left\lfloor\frac m2\right\rfloor
        -(m-2)\left\lfloor\frac{m-1}2\right\rfloor-1\right),
 \end{split}
\]
with \(a_1=1,a_2=2\). Define \(b_m=a_m-ma_{m-1}\). Then
\[
 b_m=(m-3)b_{m-1}+f_m,
 \quad f_m=\begin{cases}(2m-3)(m-2)!,&m\text{ even},\\
                         (m-2)(m-2)!,&m\text{ odd}.
          \end{cases}
\]
The base values are \(b_3=1,b_4=11\). Substituting each formula in the
opposite parity step proves by induction
\[
 b_m=\begin{cases}
 (m-2)!(2m^2-3m+2)/4,&m\geq4\text{ even},\\
 (m-3)!(m-1)^2(2m-5)/4,&m\geq3\text{ odd}.
 \end{cases}
\]
Use \(a_m=ma_{m-1}+b_m\) and base list \(a_{n-r},(n-r)!\).

For A121747 the accepted auxiliary sequence satisfies \(d_1=1,d_2=0\),
\(d_{2m}=3!+5!+\cdots+(2m-1)!\), and \(d_{2m+1}=-d_{2m}\).
Equivalently, for \(m\geq2\),
\[
 d_m=-d_{m-1}+\mathbf1_{2\mid m}(m-1)!,\qquad
 a_m=ma_{m-1}-d_{m-1}+(m-1)!\left\lfloor\frac m2\right\rfloor.
\]
The base list is \(a_{n-r},d_{n-r},(n-r)!\).

For A121584 the entry gives \(a=b+c\), with \(b=\text{A121580}\) and
\(c=\text{A121582}\). The latter source's established recurrence is
\[
 c_m=\frac{2m-3}{m-2}c_{m-1}-\frac{m-1}{m-2}c_{m-2}
                 +\frac{(m-1)!(m^2-3m+4)}2.
\]
Together with the rule for \(b\), it gives the base list
\(b_{n-r},c_{n-r},c_{n-r+1},(n-r)!\). Every coefficient vanishes.

For A121696 write \(f_m=m(m-2)!\lfloor(m-1)/2\rfloor\), so the source
recurrence is \(a_m-a_{m-2}=f_m\). Adding at adjacent indices and checking
the two parity classes proves, from \(a_1=0,a_2=1\),
\[
 a_m+a_{m-1}=\left\lceil\frac m2\right\rceil(m-1)!\quad(m\geq2).
\]
The conjectured residual reduces to
\[
 f_n-n(n-2)(a_{n-2}+a_{n-3})
 =n(n-2)!\left(\left\lfloor\frac{n-1}2\right\rfloor
                  -\left\lceil\frac{n-2}2\right\rceil\right)=0.
\]

For A121553 use \(a_n=n!(H_n+n(n-1)/4)\). For A187252 its existing
exponential generating function is
\[
 F=-\frac{2x-1+e^{2x}+4\log(1-x)}{4(1-x)}.
\]
Hence for \(n\geq1\), \(a_n=n!(H_n-1/4-E_n/4)\), where
\(E_n=\sum_{k=0}^n2^k/k!\). Divide the proposed recurrence by \(n!\),
replace \(H_{n-j}\) by \(H_n-\sum_{h=0}^{j-1}1/(n-h)\), and replace
\(E_{n-j}\) by \(E_n-\sum_{h=0}^{j-1}2^{n-h}/(n-h)!\).
The resulting identity is zero in
\(\mathbb Q(n)[H_n]\) or \(\mathbb Q(n)[H_n,E_n,2^n/n!]\), respectively.

For A111942 the source product formula gives
\(a_n/a_{n-1}=-\lfloor n/2\rfloor\) for \(n\geq2\).
Dividing the target by \(a_{n-2}\) and inserting the two ratios proves
\(4a_n+2a_{n-1}-(n-1)(n-2)a_{n-2}=0\) separately for even and odd \(n\).

\subsection{Four fixed columns of A136123}

The established bivariate ordinary generating function on
\url{https://oeis.org/A136123} is
\[
 G(x,t)=f(h(x,t)),\quad f(z)=\sum_{m\geq0}m!z^m,\quad
 h=\frac{(1-t)x^2-x}{(1-t)x^2-1}.
\]
The source identifies its columns \(k=1,2,3,4\) as A001277, A001278,
A001279 and A001280. Set \(F_0=f(x/(1+x))\). The factorial-series equation
already proved above gives
\[
 x^2(1+x)F_0'-F_0+(1+x)=0.
 \tag{GF4}
\]
Put \(K=h_t(1-h)/h^2\), \(J=-h_t/h^2\), and write
\(K_i=[t^i]K\), \(J_i=[t^i]J\). Differentiating in \(t\) gives
\(G_t=KG+J\). Therefore, if \([t^k]G=A_kF_0+B_k\),
\[
 \begin{split}
 A_0&=1,\quad B_0=0,\\
 (k+1)A_{k+1}&=\sum_{i=0}^kK_iA_{k-i},\\
 (k+1)B_{k+1}&=\sum_{i=0}^kK_iB_{k-i}+J_k.
 \end{split}
\]
This finite recursion derives, rather than guesses, the exact rational
functions printed in the appendix. Apply the conjectured operator to
\(A_kF_0+B_k\) and use (GF4) repeatedly. The coefficient of \(F_0\)
is zero, and the remaining polynomial is the displayed boundary. Thus
all four recurrence assertions follow from the source generating function.

\subsection{Three tribonacci triangles: a kernel proof}

Let \(T(n,k)\) be one of A027023 (type I), A027052 (type II), or
A027082 (type III). The common interior rule is
\[
 T(n,k)=T(n-1,k-3)+T(n-1,k-2)+T(n-1,k-1).
\]
Values outside a row are zero. The row at \(n=0\) is \((1)\).
Type I has left entries \((1,1,1)\) and rightmost entry 1; type II has
left entries \((1,0,1)\) and rightmost entry 1. Type III has left entries
\((1,1,1)\) and rightmost entry equal to the sum of the two rightmost
entries of the preceding row. The row at \(n=1\) is interpreted with
these boundary values, and each source's program supplies the same rule
for the first overlapping boundary row.
Explicitly, the rows at \(n=2\) are \((1,1,1,3,1)\),
\((1,0,1,2,1)\), and \((1,1,1,3,2)\) for types I, II, and III.

Reverse each row:
\[
 U_n(z)=\sum_{m=0}^{2n}T(n,2n-m)z^m,\quad
 B(x,z)=\sum_{n\geq0}U_n(z)x^n,\quad B_j(x)=[z^j]B(x,z).
\]
The kernel is \(K=1-x(z+1+z^{-1})\). Its small root is \(z=v\), with
\(x=v/(1+v+v^2)\), and \(v=x+O(x^2)\). For types I and II put
\[
 (\epsilon,L(z),\alpha)=
 \begin{cases}(1,z^2-1,1),&\text{type I},\\
              (0,z^2-z,0),&\text{type II}.
 \end{cases}
\]
Multiplication of the row polynomials, with the three left-edge terms
and the rightmost term corrected, gives
\[
 KB=1+\epsilon x+\frac{xL(z)}{1-xz^2}-xB_1-
                          \frac{x}{z(1-x)}.
 \tag{GF5}
\]
For completeness, the bulk operation is multiplication of \(U_{n-1}\)
by \(z+1+z^{-1}\). Its extraneous \(z^{-1}\) term contributes
\(-xB_0/z\); here \(B_0=(1-x)^{-1}\). Resetting the rightmost
coefficient to 1 contributes \(-xB_1\). The prescribed left boundary
contributes the geometric sum \(xL(z)/(1-xz^2)\). The exceptional row
\(n=1\) contributes \(\epsilon x\). This accounts for every term in (GF5).

Substitute the small root, a valid formal-series substitution, to obtain
\[
 B_1=\frac1x+\epsilon+\frac{L(v)}{1-xv^2}-\frac1{v(1-x)}.
\]
Taking the \(z^1\) and \(z^2\) coefficients of (GF5) gives
\[
 B_2=\frac{1-x}{x}B_1-B_0-(\alpha-1),\qquad
 B_3=\frac{1-x}{x}B_2-B_1-1-x(1-\alpha-1).
 \tag{GF6}
\]
Only type II uses \(B_3\) here, so its final term is zero.
The target series are \(B_1\) for A027050, \(B_2\) for A027058,
and \(B_3\) for A027059.

For type III, denote the column series by \(C_j\). The modified right
boundary instead gives
\[
 KB=1+x+\frac{x(z^2-1)}{1-xz^2}-\frac{xC_0}{z}.
 \tag{GF7}
\]
The same bulk calculation applies; the rightmost column is now propagated
by the rule instead of being reset to 1. Kernel substitution and extraction
of the first two coefficients yield
\[
 \begin{split}
 C_0&=\frac vx\left(1+x+\frac{x(v^2-1)}{1-xv^2}\right),\\
 C_1&=\frac{(1-x)C_0-1}{x},\\
 C_2&=\frac{1-x}{x}C_1-C_0.
 \end{split}
\]
Thus A027088 has generating function \(C_1\). A027089 starts at index 2,
so its zero-extended generating function is \(C_2-x\), since the full
column has an additional term at \(n=1\).
These expressions are rational in \(v\); the appendix prints their simplified
forms and the exact certificates. Four conjectures are proved. A027088 is
refuted as follows.

Its defining triangle gives \(a_1,\ldots,a_7=1,3,6,16,41,107,283\).
At \(n=7\), the printed coefficient vector is
\[
 (984,4841,-3247,-5426,-2686,-1223,213).
\]
Its dot product with \((283,107,41,16,6,3,1)\) is \(556944\), not zero.
There is also an eventual-recurrence obstruction: substituting
\(v=(1-x-\sqrt{1-2x-3x^2})/(2x)\) gives the radical form in the appendix.
The printed operator leaves a nonzero coefficient of that radical.
Since \(1-2x-3x^2=(1+x)(1-3x)\) is not a square in \(\mathbb Q(x)\),
the operator value is not rational and in particular is not a polynomial.
By (GF1) it cannot vanish at all sufficiently large coefficient indices.

\subsection{Interpolation and the two indexing repairs}

The polynomial interpolating \((1,0)\) and \((k,1)\),
\(2\leq k\leq n+1\), is uniquely
\[
 P_n(z)=1-\frac{\prod_{k=2}^{n+1}(k-z)}{n!}.
\]
For A394578 evaluation at \(z=n+1/2\), with its specified scale, gives
\[
 a_n=\frac{4^n}{2}+(-1)^n\operatorname{Cat}_{n-1}\quad(n\geq1),
 \qquad a_0=0,
\]
and consequently
\[
 F=\frac{2x}{1-4x}+\frac{1-\sqrt{1+4x}}2.
\]
Its recurrence follows from the exact certificate. Moreover,
\(\operatorname{Cat}_{n-1}\leq4^{n-1}/n\), so
\(a_n/(4^n/2)=1+O(1/n)\) and \(a_{n+1}/a_n\to4\).
This disproves the distinct printed limit \(1/4\).

For A396724 use the same interpolation polynomial with \(z=-3/2\),
and multiply by \(-4^n/2\). The identity
\[
 nP_n(z)=(n+1-z)P_{n-1}(z)+z-1
\]
follows directly from its product. Subtracting four times the scaled relation
one index earlier eliminates the geometric forcing and gives
\[
 na_n-(8n-4z)a_{n-1}+16(n-z)a_{n-2}=0.
\]
Putting \(z=-3/2\) proves the conjecture. A neighboring interpolation entry
already containing a proof was not counted.

For A142243, the row sums are
\(s_n=[z^n](1+z)^{2n}/\sqrt{1-4z}\). Lagrange coefficient summation gives
\[
 S(x)=\frac{1+v}{(1-v)\sqrt{1-4v}},\qquad x=\frac v{(1+v)^2}.
\]
The actual sums at indices 0 through 3 are \(1,4,20,106\). The printed
recurrence has residual \(-344\) at \(n=3\). Replacing \(n\) by \(n-1\)
in all its polynomial coefficients gives the certified homogeneous
operator identity \(\mathcal L S=0\) and proves the corrected recurrence
for \(n\geq3\). This is recorded as an index repair.

For A059837 the actual offset is 1 and
\[
 a_n=(n-1)!\binom{n-1}{\lfloor(n-1)/2\rfloor}.
\]
The printed recurrence at \(n=3\) evaluates to
\(4\cdot4-6\cdot1-48\cdot1=-38\). Its corrected version is
\[
 na_n-2(n-1)a_{n-1}-4(n-1)(n-2)^2a_{n-2}=0\quad(n\geq3).
\]
Divide by \(a_{n-2}\) and use the exact ratio
\(a_n/a_{n-1}=(n-1)^2/\lfloor n/2\rfloor\); both parity identities vanish.
The originally printed formula applies to the reindexed sequence
\(b_n=a_{n+1}\). The source also discusses offset zero nearby, so this
is conservatively classified as an index repair.
\subsection{Complete recurrence coefficients and exact certificates}
\subsubsection*{A174193 --- proved}
Source: \url{https://oeis.org/A174193}.
Coefficients of the algebraic ordinary generating function in the entry definition; the separate permutation interpretation is outside this result.
The certified recurrence range is $n\geq 4$.
Only the recurrence of the defined series is certified.
The original coefficient polynomials are:
\[
p_{0}(n)=- 404489 n^{2} + 1057390 n + 3732736
\]
\[
p_{1}(n)=3147815 n^{2} - 8680082 n - 14098754
\]
\[
p_{2}(n)=- 4672878 n^{2} + 1304044 n + 29029698
\]
\[
p_{3}(n)=- 8358160 n^{2} + 84680184 n - 162964424
\]
\[
p_{4}(n)=10287712 n^{2} - 78361536 n + 144300744
\]
For (GF3), the exact input is
\[
x=x
\]
with $F=\sum_i R_i\sqrt{D_i}$ and
\[
R_{0}=\frac{3 x^{2} - 1}{2 x^{3} - 2 x^{2}}
\]
\[
D_{0}=1
\]
\[
R_{1}=\frac{- 4 x^{4} - 30 x^{3} + 4 x^{2} + 5 x - 1}{32 x^{4} - 16 x^{3} + 2 x^{2}}
\]
\[
D_{1}=1 - 4 x
\]
The exact operator identity is $\mathcal L F=P(x)$, with every nonconstant radical coefficient zero, where
\[
P(x)=2708316 x^{3} + 2537736 x^{2} - 15245384 x + 3732736
\]

\subsubsection*{A346074 --- proved}
Source: \url{https://oeis.org/A346074}.
Sequence determined by the convolution whose ordinary generating function satisfies F=1/(1-x)+x\textasciicircum{}5 F\textasciicircum{}2, with F(0)=1.
The certified recurrence range is $n\geq 6$.
The original coefficient polynomials are:
\[
p_{0}(n)=n + 5
\]
\[
p_{1}(n)=- 2 n - 8
\]
\[
p_{2}(n)=n + 3
\]
\[
p_{3}(n)=0
\]
\[
p_{4}(n)=0
\]
\[
p_{5}(n)=10 - 4 n
\]
\[
p_{6}(n)=4 n - 12
\]
For (GF3), the exact input is
\[
x=x
\]
with $F=\sum_i R_i\sqrt{D_i}$ and
\[
R_{0}=\frac{1}{2 x^{5}}
\]
\[
D_{0}=1
\]
\[
R_{1}=- \frac{1}{2 x^{5}}
\]
\[
D_{1}=- \frac{4 x^{5}}{1 - x} + 1
\]
The exact operator identity is $\mathcal L F=P(x)$, with every nonconstant radical coefficient zero, where
\[
P(x)=5 - 4 x
\]

\subsubsection*{A247173 --- proved}
Source: \url{https://oeis.org/A247173}.
Coefficients of the radical ordinary generating function specified in the entry definition.
The certified recurrence range is $n\geq 5$.
The original coefficient polynomials are:
\[
p_{0}(n)=- 3 n^{2} + n + 4
\]
\[
p_{1}(n)=9 n^{2} - 3 n - 4
\]
\[
p_{2}(n)=18 n^{2} - 42 n + 16
\]
\[
p_{3}(n)=- 6 n^{2} + 2 n + 16
\]
\[
p_{4}(n)=- 15 n^{2} + 53 n - 12
\]
\[
p_{5}(n)=- 3 n^{2} + 13 n - 4
\]
For (GF3), the exact input is
\[
x=x
\]
with $F=\sum_i R_i\sqrt{D_i}$ and
\[
R_{0}=- \frac{1}{2 x^{2} + 2 x}
\]
\[
D_{0}=1
\]
\[
R_{1}=\frac{1 - 2 x}{2 x^{4} + 6 x^{3} - 10 x^{2} + 2 x}
\]
\[
D_{1}=x^{4} + 4 x^{3} - 2 x^{2} - 4 x + 1
\]
The exact operator identity is $\mathcal L F=P(x)$, with every nonconstant radical coefficient zero, where
\[
P(x)=- 4 x^{3} - 12 x^{2} + 4 x + 4
\]

\subsubsection*{A247102 --- proved}
Source: \url{https://oeis.org/A247102}.
Coefficients of the radical ordinary generating function specified in the entry definition.
The certified recurrence range is $n\geq 4$.
The original coefficient polynomials are:
\[
p_{0}(n)=n + 1
\]
\[
p_{1}(n)=- 2 n - 5
\]
\[
p_{2}(n)=21 - 24 n
\]
\[
p_{3}(n)=45 - 30 n
\]
\[
p_{4}(n)=18 - 9 n
\]
For (GF3), the exact input is
\[
x=x
\]
with $F=\sum_i R_i\sqrt{D_i}$ and
\[
R_{0}=\frac{- 2 x - 1}{2 x^{2} + 2 x}
\]
\[
D_{0}=1
\]
\[
R_{1}=\frac{- 3 x - 1}{6 x^{4} + 18 x^{3} + 10 x^{2} - 2 x}
\]
\[
D_{1}=- 3 x^{2} - 6 x + 1
\]
The exact operator identity is $\mathcal L F=P(x)$, with every nonconstant radical coefficient zero, where
\[
P(x)=9 x^{3} + 15 x^{2} + 6 x + 2
\]

\subsubsection*{A234269 --- proved}
Source: \url{https://oeis.org/A234269}.
Coefficients of the radical ordinary generating function specified in the entry definition.
The certified recurrence range is $n\geq 4$.
The original coefficient polynomials are:
\[
p_{0}(n)=- 2 n^{2} + 3 n + 5
\]
\[
p_{1}(n)=- 2 n^{2} + 8 n
\]
\[
p_{2}(n)=8 n^{2} - 20 n + 12
\]
\[
p_{3}(n)=16 n^{2} - 64 n + 58
\]
\[
p_{4}(n)=8 n^{2} - 44 n + 56
\]
For (GF3), the exact input is
\[
x=x
\]
with $F=\sum_i R_i\sqrt{D_i}$ and
\[
R_{0}=- \frac{1}{2 x}
\]
\[
D_{0}=1
\]
\[
R_{1}=\frac{2 x^{2} - 1}{8 x^{4} + 8 x^{3} - 2 x}
\]
\[
D_{1}=- 4 x^{3} - 4 x^{2} + 1
\]
The exact operator identity is $\mathcal L F=P(x)$, with every nonconstant radical coefficient zero, where
\[
P(x)=x^{2} \left(2 x + 3\right)
\]

\subsubsection*{A268600 --- proved}
Source: \url{https://oeis.org/A268600}.
Coefficients of 1/(2 sqrt(1-4x))+1/(4-2 sqrt(1+4x)).
The certified recurrence range is $n\geq 4$.
The original coefficient polynomials are:
\[
p_{0}(n)=- 3 n^{2} + 3 n
\]
\[
p_{1}(n)=40 n^{2} - 136 n + 96
\]
\[
p_{2}(n)=292 - 112 n
\]
\[
p_{3}(n)=- 640 n^{2} + 3840 n - 5600
\]
\[
p_{4}(n)=768 n^{2} - 4608 n + 6720
\]
For (GF3), the exact input is
\[
x=x
\]
with $F=\sum_i R_i\sqrt{D_i}$ and
\[
R_{0}=\frac{1}{3 - 4 x}
\]
\[
D_{0}=1
\]
\[
R_{1}=\frac{1}{2 - 8 x}
\]
\[
D_{1}=1 - 4 x
\]
\[
R_{2}=\frac{1}{6 - 8 x}
\]
\[
D_{2}=4 x + 1
\]
The exact operator identity is $\mathcal L F=P(x)$, with every nonconstant radical coefficient zero, where
\[
P(x)=12 x^{2} \left(4 x + 1\right)
\]

\subsubsection*{A268601 --- proved}
Source: \url{https://oeis.org/A268601}.
Coefficients of 1/(2 sqrt(1-4x))-1/(4-2 sqrt(1+4x)).
The certified recurrence range is $n\geq 4$.
The original coefficient polynomials are:
\[
p_{0}(n)=3 n^{2} - 3 n
\]
\[
p_{1}(n)=- 40 n^{2} + 136 n - 96
\]
\[
p_{2}(n)=112 n - 292
\]
\[
p_{3}(n)=640 n^{2} - 3840 n + 5600
\]
\[
p_{4}(n)=- 768 n^{2} + 4608 n - 6720
\]
For (GF3), the exact input is
\[
x=x
\]
with $F=\sum_i R_i\sqrt{D_i}$ and
\[
R_{0}=- \frac{1}{3 - 4 x}
\]
\[
D_{0}=1
\]
\[
R_{1}=\frac{1}{2 - 8 x}
\]
\[
D_{1}=1 - 4 x
\]
\[
R_{2}=- \frac{1}{6 - 8 x}
\]
\[
D_{2}=4 x + 1
\]
The exact operator identity is $\mathcal L F=P(x)$, with every nonconstant radical coefficient zero, where
\[
P(x)=12 x^{2} \left(4 x + 1\right)
\]

\subsubsection*{A108296 --- proved}
Source: \url{https://oeis.org/A108296}.
Sequence with the ordinary generating function displayed in the certificate, as listed in its FORMULA section.
The certified recurrence range is $n\geq 5$.
The original coefficient polynomials are:
\[
p_{0}(n)=n + 3
\]
\[
p_{1}(n)=- n - 2
\]
\[
p_{2}(n)=- 2 n - 2
\]
\[
p_{3}(n)=6 - 2 n
\]
\[
p_{4}(n)=n + 1
\]
\[
p_{5}(n)=3 n - 6
\]
For (GF3), the exact input is
\[
x=x
\]
with $F=\sum_i R_i\sqrt{D_i}$ and
\[
R_{0}=\frac{1}{2 x^{3}}
\]
\[
D_{0}=1
\]
\[
R_{1}=\frac{1}{2 x^{5} - 2 x^{3}}
\]
\[
D_{1}=- 3 x^{4} - 4 x^{3} - 2 x^{2} + 1
\]
The exact operator identity is $\mathcal L F=P(x)$, with every nonconstant radical coefficient zero, where
\[
P(x)=x + 3
\]

\subsubsection*{A244884 --- proved}
Source: \url{https://oeis.org/A244884}.
Coefficients of x+(1-x)M(x), where M is the Motzkin generating function.
The certified recurrence range is $n\geq 5$.
The original coefficient polynomials are:
\[
p_{0}(n)=n + 2
\]
\[
p_{1}(n)=- 3 n - 1
\]
\[
p_{2}(n)=- n
\]
\[
p_{3}(n)=3 n - 9
\]
For (GF3), the exact input is
\[
x=x
\]
with $F=\sum_i R_i\sqrt{D_i}$ and
\[
R_{0}=\frac{2 x^{3} + x^{2} - 2 x + 1}{2 x^{2}}
\]
\[
D_{0}=1
\]
\[
R_{1}=\frac{x - 1}{2 x^{2}}
\]
\[
D_{1}=- 3 x^{2} - 2 x + 1
\]
The exact operator identity is $\mathcal L F=P(x)$, with every nonconstant radical coefficient zero, where
\[
P(x)=3 x^{4} - 3 x^{3} - 5 x^{2} - x + 2
\]

\subsubsection*{A244886 --- proved}
Source: \url{https://oeis.org/A244886}.
Coefficients of the algebraic generating function in the definition.
The certified recurrence range is $n\geq 6$.
The original coefficient polynomials are:
\[
p_{0}(n)=n
\]
\[
p_{1}(n)=3 - 5 n
\]
\[
p_{2}(n)=2 n
\]
\[
p_{3}(n)=13 n - 30
\]
\[
p_{4}(n)=-3
\]
\[
p_{5}(n)=21 - 8 n
\]
\[
p_{6}(n)=9 - 3 n
\]
For (GF3), the exact input is
\[
x=x
\]
with $F=\sum_i R_i\sqrt{D_i}$ and
\[
R_{0}=\frac{- 2 x^{2} - 3 x + 1}{2 x^{4} + 4 x^{3} - 2 x^{2} - 6 x + 2}
\]
\[
D_{0}=1
\]
\[
R_{1}=\frac{1}{2 x^{4} + 4 x^{3} - 2 x^{2} - 6 x + 2}
\]
\[
D_{1}=- 3 x^{2} - 2 x + 1
\]
The exact operator identity is $\mathcal L F=P(x)$, with every nonconstant radical coefficient zero, where
\[
P(x)=x \left(3 x^{3} + 3 x^{2} + x - 1\right)
\]

\subsubsection*{A257104 --- proved}
Source: \url{https://oeis.org/A257104}.
Restricted Motzkin paths described by the continued fraction used in the proof.
The certified recurrence range is $n\geq 7$.
The original coefficient polynomials are:
\[
p_{0}(n)=2 - n
\]
\[
p_{1}(n)=7 n - 17
\]
\[
p_{2}(n)=34 - 14 n
\]
\[
p_{3}(n)=n + 22
\]
\[
p_{4}(n)=16 n - 89
\]
\[
p_{5}(n)=23 - 4 n
\]
\[
p_{6}(n)=3 n - 15
\]
For (GF3), the exact input is
\[
x=x
\]
with $F=\sum_i R_i\sqrt{D_i}$ and
\[
R_{0}=\frac{2 x^{4} - 3 x^{3} + 9 x^{2} - 8 x + 2}{2 x^{4} - 4 x^{3} + 14 x^{2} - 10 x + 2}
\]
\[
D_{0}=1
\]
\[
R_{1}=- \frac{x^{2}}{2 x^{4} - 4 x^{3} + 14 x^{2} - 10 x + 2}
\]
\[
D_{1}=- 3 x^{2} - 2 x + 1
\]
The exact operator identity is $\mathcal L F=P(x)$, with every nonconstant radical coefficient zero, where
\[
P(x)=\left(x^{2} + 3 x - 1\right) \left(3 x^{4} - 6 x^{3} + 4 x^{2} + 3 x - 2\right)
\]

\subsubsection*{A099038 --- proved}
Source: \url{https://oeis.org/A099038}.
Diagonal binomial/Krawtchouk sum from the entry; its constant-term representation is derived below.
The certified recurrence range is $n\geq 5$.
The original coefficient polynomials are:
\[
p_{0}(n)=n
\]
\[
p_{1}(n)=- n
\]
\[
p_{2}(n)=n
\]
\[
p_{3}(n)=3 - 3 n
\]
\[
p_{4}(n)=13 - 5 n
\]
\[
p_{5}(n)=n - 3
\]
For (GF3), the exact input is
\[
x=x
\]
with $F=\sum_i R_i\sqrt{D_i}$ and
\[
R_{0}=0
\]
\[
D_{0}=1
\]
\[
R_{1}=\frac{1}{x^{4} - 6 x^{3} + 3 x^{2} - 2 x + 1}
\]
\[
D_{1}=- 4 x^{3} + \left(x^{2} - x + 1\right)^{2}
\]
The exact operator identity is $\mathcal L F=P(x)$, with every nonconstant radical coefficient zero, where
\[
P(x)=0
\]

\subsubsection*{A337500 --- proved}
Source: \url{https://oeis.org/A337500}.
Walks with steps +1 and -1 whose final displacement has absolute value at most 3; intermediate positions are unrestricted.
The certified recurrence range is $n\geq 3$.
The original coefficient polynomials are:
\[
p_{0}(n)=- n^{2} + n + 12
\]
\[
p_{1}(n)=2 n^{2} - 4 n - 22
\]
\[
p_{2}(n)=4 n^{2} - 8 n + 4
\]
\[
p_{3}(n)=- 8 n^{2} + 24 n - 16
\]
For (GF3), the exact input is
\[
x=x
\]
with $F=\sum_i R_i\sqrt{D_i}$ and
\[
R_{0}=\frac{- x - 1}{x^{3}}
\]
\[
D_{0}=1
\]
\[
R_{1}=\frac{x^{3} + 2 x^{2} - x - 1}{4 x^{5} - x^{3}}
\]
\[
D_{1}=1 - 4 x^{2}
\]
The exact operator identity is $\mathcal L F=P(x)$, with every nonconstant radical coefficient zero, where
\[
P(x)=12
\]

\subsubsection*{A127905 --- proved}
Source: \url{https://oeis.org/A127905}.
a\_n is the coefficient of u\textasciicircum{}(n-1) in (1+u+u\textasciicircum{}3)\textasciicircum{}n.
The certified recurrence range is $n\geq 3$.
The original coefficient polynomials are:
\[
p_{0}(n)=4 n^{3} - 6 n^{2} + 2
\]
\[
p_{1}(n)=- 12 n^{3} + 24 n^{2} - 10 n
\]
\[
p_{2}(n)=12 n^{3} - 30 n^{2} + 18 n
\]
\[
p_{3}(n)=- 31 n^{3} + 93 n^{2} - 62 n
\]
For (GF3), the exact input is
\[
x=\frac{t}{t^{3} + t + 1}
\]
with $F=\sum_i R_i\sqrt{D_i}$ and
\[
R_{0}=\frac{t \left(t^{3} + t + 1\right)}{1 - 2 t^{3}}
\]
\[
D_{0}=1
\]
The exact operator identity is $\mathcal L F=P(x)$, with every nonconstant radical coefficient zero, where
\[
P(x)=0
\]

\subsubsection*{A392976 --- proved}
Source: \url{https://oeis.org/A392976}.
The compositional inverse w of t((1-t\textasciicircum{}2)\textasciicircum{}2-t), divided by x, supplies the ordinary generating function.
The certified recurrence range is $n\geq 5$.
The original coefficient polynomials are:
\[
\begin{aligned}
p_{0}(n)&=23495312375279 n^{5} - 97473596443030 n^{4}\\
&\quad + 77470631009665 n^{3} + 97473596443030 n^{2}\\
&\quad - 100965943384944 n
\end{aligned}
\]
\[
\begin{aligned}
p_{1}(n)&=- 73180721424774 n^{5} + 490203969890442 n^{4}\\
&\quad - 1274195207740074 n^{3} + 1488868655358738 n^{2}\\
&\quad - 631696696084332 n
\end{aligned}
\]
\[
\begin{aligned}
p_{2}(n)&=- 406745978528262 n^{5} + 2710034017555110 n^{4}\\
&\quad - 6497864886329538 n^{3} + 6606177703000266 n^{2}\\
&\quad - 2374863833644560 n - 36737022053016
\end{aligned}
\]
\[
\begin{aligned}
p_{3}(n)&=- 384334715888720 n^{5} + 1886690791154422 n^{4}\\
&\quad +2005293720367118 n^{3} - 26729399150503384 n^{2}\\
&\quad +53179415800700880 n - 33371926512356880
\end{aligned}
\]
\[
\begin{aligned}
p_{4}(n)&=78200611484375 n^{5} - 2829824426102500 n^{4}\\
&\quad +24949288577136625 n^{3} - 93338391900041900 n^{2}\\
&\quad +160425630781509240 n - 104626628147875680
\end{aligned}
\]
\[
\begin{aligned}
p_{5}(n)&=140685548031250 n^{5} - 2487209724443750 n^{4}\\
&\quad +17499620703856250 n^{3} - 61230529923966250 n^{2}\\
&\quad +106501850375442900 n - 73632743322117840
\end{aligned}
\]
For (GF3), the exact input is
\[
x=t \left(- t + \left(1 - t^{2}\right)^{2}\right)
\]
with $F=\sum_i R_i\sqrt{D_i}$ and
\[
R_{0}=\frac{1}{- t + \left(1 - t^{2}\right)^{2}}
\]
\[
D_{0}=1
\]
The exact operator identity is $\mathcal L F=P(x)$, with every nonconstant radical coefficient zero, where
\[
P(x)=0
\]

\subsubsection*{A186185 --- proved}
Source: \url{https://oeis.org/A186185}.
Coefficients of 1/(1-zT(z)), with z=x/(1-x) and T=1+zT\textasciicircum{}3.
The certified recurrence range is $n\geq 7$.
The original coefficient polynomials are:
\[
p_{0}(n)=4 n^{2} - 6 n + 2
\]
\[
p_{1}(n)=9 n^{2} - 127 n + 170
\]
\[
p_{2}(n)=- 570 n^{2} + 3680 n - 5590
\]
\[
p_{3}(n)=2579 n^{2} - 18451 n + 32892
\]
\[
p_{4}(n)=- 4899 n^{2} + 39585 n - 80634
\]
\[
p_{5}(n)=4272 n^{2} - 40026 n + 93330
\]
\[
p_{6}(n)=- 1395 n^{2} + 15345 n - 41850
\]
For (GF3), the exact input is
\[
x=\frac{t - 1}{t^{3} + t - 1}
\]
with $F=\sum_i R_i\sqrt{D_i}$ and
\[
R_{0}=\frac{t^{2}}{t^{2} - t + 1}
\]
\[
D_{0}=1
\]
The exact operator identity is $\mathcal L F=P(x)$, with every nonconstant radical coefficient zero, where
\[
P(x)=\left(2 - 2 x\right) \left(207 x^{3} - 243 x^{2} + 27 x + 1\right)
\]

\subsubsection*{A220433 --- proved}
Source: \url{https://oeis.org/A220433}.
The exponential generating function is the inverse under composition of t-t\textasciicircum{}2+t\textasciicircum{}3/6.
The certified recurrence range is $n\geq 3$.
The original coefficient polynomials are:
\[
p_{0}(n)=4
\]
\[
p_{1}(n)=18 - 12 n
\]
\[
p_{2}(n)=- 9 n^{2} + 36 n - 35
\]
For (GF3), the exact input is
\[
x=\frac{t^{3}}{6} - t^{2} + t
\]
with $F=\sum_i R_i\sqrt{D_i}$ and
\[
R_{0}=t
\]
\[
D_{0}=1
\]
Use the exponential conversion (GF2) before applying the operator.
The exact operator identity is $\mathcal L F=P(x)$, with every nonconstant radical coefficient zero, where
\[
P(x)=2 x^{2}
\]

\subsubsection*{A071264 --- proved}
Source: \url{https://oeis.org/A071264}.
Coefficients of (1+xC\textasciicircum{}4)C\textasciicircum{}3, where C=1+xC\textasciicircum{}2 is the Catalan series.
The certified recurrence range is $n\geq 3$.
The original coefficient polynomials are:
\[
p_{0}(n)=8587 n^{3} + 114480 n^{2} + 439415 n + 370002
\]
\[
p_{1}(n)=8587 n^{3} - 377960 n^{2} - 2268967 n - 1091596
\]
\[
p_{2}(n)=- 171740 n^{3} + 143858 n^{2} + 464342 n - 246668
\]
For (GF3), the exact input is
\[
x=\frac{t - 1}{t^{2}}
\]
with $F=\sum_i R_i\sqrt{D_i}$ and
\[
R_{0}=t^{3} \left(t^{2} \left(t - 1\right) + 1\right)
\]
\[
D_{0}=1
\]
The exact operator identity is $\mathcal L F=P(x)$, with every nonconstant radical coefficient zero, where
\[
P(x)=370002
\]

\subsubsection*{A392211 --- proved}
Source: \url{https://oeis.org/A392211}.
Algebraic series F specified by x=(3+14F-20F\textasciicircum{}2-3 sqrt(1+4F-4F\textasciicircum{}2))/8 and F(0)=0.
The certified recurrence range is $n\geq 6$.
The original coefficient polynomials are:
\[
p_{0}(n)=30632912640 n^{4} - 215110974720 n^{3} + 430902535680 n^{2} - 246424473600 n
\]
\[
\begin{aligned}
p_{1}(n)&=- 61010825216 n^{4} + 756124952576 n^{3}\\
&\quad - 3128904705856 n^{2} + 5009001827776 n\\
&\quad - 2575211249280
\end{aligned}
\]
\[
\begin{aligned}
p_{2}(n)&=- 1127075797485 n^{4} + 11295798763110 n^{3}\\
&\quad - 40811553201975 n^{2} + 64043018252310 n\\
&\quad - 37173001041840
\end{aligned}
\]
\[
\begin{aligned}
p_{3}(n)&=2244769385584 n^{4} - 34554389907976 n^{3}\\
&\quad + 195408589554464 n^{2} - 483915433738736 n\\
&\quad + 444700475929080
\end{aligned}
\]
\[
\begin{aligned}
p_{4}(n)&=765822816000 n^{4} - 9972711264000 n^{3}\\
&\quad + 48002376840000 n^{2} - 102102215016000 n\\
&\quad + 82139932008000
\end{aligned}
\]
\[
\begin{aligned}
p_{5}(n)&=- 1525270630400 n^{4} + 28054747596800 n^{3}\\
&\quad - 191649301734400 n^{2} + 575634929036800 n\\
&\quad - 640491407424000
\end{aligned}
\]
For (GF3), the exact input is
\[
x=\frac{\left(4 - 2 t\right) \left(- 3 t - \frac{80 - 40 t}{t^{2} + 4} + 14\right)}{8 t^{2} + 32}
\]
with $F=\sum_i R_i\sqrt{D_i}$ and
\[
R_{0}=\frac{4 - 2 t}{t^{2} + 4}
\]
\[
D_{0}=1
\]
The exact operator identity is $\mathcal L F=P(x)$, with every nonconstant radical coefficient zero, where
\[
P(x)=- 4200 x^{3} \left(36148218 x - 40772785\right)
\]

\subsubsection*{A142243 --- proved/corrected}
Source: \url{https://oeis.org/A142243}.
Row sums of T(n,k)=binom(2n,k) binom(2n-2k,n-k), with rows indexed from n=0.
The certified recurrence range is $n\geq 3$.
This entry is an indexing repair and is excluded from the substantive closure count.
The original coefficient polynomials are:
\[
p_{0}(n)=4 n^{2} + 6 n + 2
\]
\[
p_{1}(n)=- 81 n^{2} + 19 n - 8
\]
\[
p_{2}(n)=510 n^{2} - 770 n + 300
\]
\[
p_{3}(n)=- 1000 n^{2} + 2500 n - 1500
\]
The corrected coefficient vector is obtained by $n\mapsto n-1$ in every $p_j$.
For (GF3), the exact input is
\[
x=\frac{t}{\left(t + 1\right)^{2}}
\]
with $F=\sum_i R_i\sqrt{D_i}$ and
\[
R_{0}=0
\]
\[
D_{0}=1
\]
\[
R_{1}=\frac{t + 1}{\left(1 - 4 t\right) \left(1 - t\right)}
\]
\[
D_{1}=1 - 4 t
\]
The exact operator identity is $\mathcal L F=P(x)$, with every nonconstant radical coefficient zero, where
\[
P(x)=0
\]

\subsubsection*{A091539 --- proved}
Source: \url{https://oeis.org/A091539}.
The product formula in the entry, equivalently a fixed linear combination of two hypergeometric products.
The certified recurrence range is $n\geq 4$.
The original coefficient polynomials are:
\[
p_{0}(n)=1
\]
\[
p_{1}(n)=- 18 n^{2} + 21 n - 5
\]
\[
p_{2}(n)=81 n^{4} - 351 n^{3} + 531 n^{2} - 321 n + 60
\]
\[
\rho_{0}(n)=3 n \left(3 n - 1\right)
\]
\[
\rho_{1}(n)=\left(3 n - 2\right) \left(3 n - 1\right)
\]
The normalized recurrence residual is identically zero for each component.

\subsubsection*{A027175 --- proved}
Source: \url{https://oeis.org/A027175}, \url{https://oeis.org/A027170}.
The near-central entry T(2n-1,n-1) of the Pascal-derived array A027170.
The certified recurrence range is $n\geq 3$.
The original coefficient polynomials are:
\[
p_{0}(n)=2 n + 2
\]
\[
p_{1}(n)=- 7 n - 5
\]
\[
p_{2}(n)=6 - 4 n
\]
\[
c(n)=12 - 36 n
\]
In $a_n=\binom{2n}nR(n)+V(n)+E2^n$, the exact components are
\[
R(n)=\frac{9 n + 5}{2 n + 2}
\]
\[
V(n)=-4
\]
\[
E=0
\]
The central-binomial, polynomial (including $c(n)$), and exponential residuals are all exactly zero.

\subsubsection*{A027182 --- proved}
Source: \url{https://oeis.org/A027182}, \url{https://oeis.org/A027170}.
The central self-convolution coefficient of row n of A027170; symmetry makes it the sum of squares of that row.
The certified recurrence range is $n\geq 6$.
The original coefficient polynomials are:
\[
p_{0}(n)=2 n + 4
\]
\[
p_{1}(n)=- 17 n - 23
\]
\[
p_{2}(n)=47 n + 36
\]
\[
p_{3}(n)=- 44 n - 34
\]
\[
p_{4}(n)=80 - 4 n
\]
\[
p_{5}(n)=16 n - 72
\]
\[
c(n)=288 n - 1128
\]
In $a_n=\binom{2n}nR(n)+V(n)+E2^n$, the exact components are
\[
R(n)=\frac{81 n^{2} + 171 n + 66}{\left(n + 1\right) \left(n + 2\right)}
\]
\[
V(n)=16 n + 40
\]
\[
E=-72
\]
The central-binomial, polynomial (including $c(n)$), and exponential residuals are all exactly zero.

\subsubsection*{A177258 --- proved}
Source: \url{https://oeis.org/A177258}.
The adjacent-cycle derangement sequence with parameter q=2 and the listed double factorial sum.
The certified recurrence range is $n\geq 4$.
The original coefficient polynomials are:
\[
p_{0}(n)=1
\]
\[
p_{1}(n)=1 - n
\]
\[
p_{2}(n)=1 - n
\]
\[
p_{3}(n)=1 - n
\]
\[
p_{4}(n)=-1
\]
The exact identity is derived in the corresponding proof subsection above; all displayed residuals in the machine-readable certificate are zero.

\subsubsection*{A177259 --- proved}
Source: \url{https://oeis.org/A177259}.
The adjacent-cycle derangement sequence with parameter q=3 and the listed double factorial sum.
The certified recurrence range is $n\geq 6$.
The original coefficient polynomials are:
\[
p_{0}(n)=1
\]
\[
p_{1}(n)=1 - n
\]
\[
p_{2}(n)=1 - n
\]
\[
p_{3}(n)=-1
\]
\[
p_{4}(n)=1 - n
\]
\[
p_{5}(n)=0
\]
\[
p_{6}(n)=-2
\]
The exact identity is derived in the corresponding proof subsection above; all displayed residuals in the machine-readable certificate are zero.

\subsubsection*{A177260 --- proved}
Source: \url{https://oeis.org/A177260}.
The adjacent-cycle derangement sequence with parameter q=4 and the listed double factorial sum.
The certified recurrence range is $n\geq 8$.
The original coefficient polynomials are:
\[
p_{0}(n)=1
\]
\[
p_{1}(n)=1 - n
\]
\[
p_{2}(n)=1 - n
\]
\[
p_{3}(n)=0
\]
\[
p_{4}(n)=-2
\]
\[
p_{5}(n)=1 - n
\]
\[
p_{6}(n)=0
\]
\[
p_{7}(n)=0
\]
\[
p_{8}(n)=-3
\]
The exact identity is derived in the corresponding proof subsection above; all displayed residuals in the machine-readable certificate are zero.

\subsubsection*{A351933 --- proved}
Source: \url{https://oeis.org/A351933}.
Coefficients, in exponential generating function convention, of exp(x/(1-x\textasciicircum{}2/2)).
The certified recurrence range is $n\geq 4$.
The original coefficient polynomials are:
\[
p_{0}(n)=4
\]
\[
p_{1}(n)=-4
\]
\[
p_{2}(n)=- 4 n^{2} + 12 n - 8
\]
\[
p_{3}(n)=- 2 n^{2} + 6 n - 4
\]
\[
p_{4}(n)=n^{4} - 10 n^{3} + 35 n^{2} - 50 n + 24
\]
The exact identity is derived in the corresponding proof subsection above; all displayed residuals in the machine-readable certificate are zero.

\subsubsection*{A001722 --- proved}
Source: \url{https://oeis.org/A001722}.
The accepted Stirling-number sum, equal to [u\textasciicircum{}2](5+u) rising(n+2).
The certified recurrence range is $n\geq 3$.
The original coefficient polynomials are:
\[
p_{0}(n)=1
\]
\[
p_{1}(n)=- 3 n - 15
\]
\[
p_{2}(n)=3 n^{2} + 27 n + 61
\]
\[
p_{3}(n)=- n^{3} - 12 n^{2} - 48 n - 64
\]
After removing the common product specified in the proof, the exact remaining polynomial is
\[
Q(u)=u^{3}
\]

\subsubsection*{A345651 --- proved}
Source: \url{https://oeis.org/A345651}.
The signed-Stirling sum, equal to [u\textasciicircum{}4](4+u) falling n.
The certified recurrence range is $n\geq 9$.
The original coefficient polynomials are:
\[
p_{0}(n)=1
\]
\[
p_{1}(n)=4 n - 26
\]
\[
p_{2}(n)=6 n^{2} - 84 n + 295
\]
\[
p_{3}(n)=4 n^{3} - 90 n^{2} + 676 n - 1695
\]
\[
p_{4}(n)=n^{4} - 32 n^{3} + 384 n^{2} - 2048 n + 4096
\]
After removing the common product specified in the proof, the exact remaining polynomial is
\[
Q(u)=u^{4}
\]

\subsubsection*{A111942 --- proved}
Source: \url{https://oeis.org/A111942}.
The signed product (-1)\textasciicircum{}(n-1) floor((n-1)/2)! floor(n/2)! for n>=1.
The certified recurrence range is $n\geq 3$.
The original coefficient polynomials are:
\[
p_{0}(n)=4
\]
\[
p_{1}(n)=2
\]
\[
p_{2}(n)=- n^{2} + 3 n - 2
\]
The exact identity is derived in the corresponding proof subsection above; all displayed residuals in the machine-readable certificate are zero.

\subsubsection*{A059837 --- proved/corrected}
Source: \url{https://oeis.org/A059837}.
With the entry offset 1, a\_n=(n-1)! binom(n-1,floor((n-1)/2)).
The certified recurrence range is $n\geq 3$.
This entry is an indexing repair and is excluded from the substantive closure count.
The original coefficient polynomials are:
\[
p_{0}(n)=n + 1
\]
\[
p_{1}(n)=- 2 n
\]
\[
p_{2}(n)=- 4 n^{3} + 8 n^{2} - 4 n
\]
\[
\widetilde p_{0}(n)=n
\]
\[
\widetilde p_{1}(n)=2 - 2 n
\]
\[
\widetilde p_{2}(n)=\left(4 - 4 n\right) \left(n - 2\right)^{2}
\]
The original residual at $n=3$ is $-38$; both parity residuals of the corrected formula are zero.

\subsubsection*{A394578 --- proved/corrected}
Source: \url{https://oeis.org/A394578}.
Scaled evaluation at z=n+1/2 of the polynomial interpolating (1,0) and (k,1) for 2<=k<=n+1.
The certified recurrence range is $n\geq 3$.
The original coefficient polynomials are:
\[
p_{0}(n)=n
\]
\[
p_{1}(n)=2 - 4 n
\]
\[
p_{2}(n)=48 - 16 n
\]
\[
p_{3}(n)=64 n - 224
\]
For (GF3), the exact input is
\[
x=x
\]
with $F=\sum_i R_i\sqrt{D_i}$ and
\[
R_{0}=- \frac{1}{8 x - 2}
\]
\[
D_{0}=1
\]
\[
R_{1}=- \frac{1}{2}
\]
\[
D_{1}=4 x + 1
\]
The exact operator identity is $\mathcal L F=P(x)$, with every nonconstant radical coefficient zero, where
\[
P(x)=x \left(12 x + 1\right)
\]
The separate limit assertion is false: $a_{n+1}/a_n\to4$, as proved above.

\subsubsection*{A396724 --- proved}
Source: \url{https://oeis.org/A396724}.
The same interpolation polynomial evaluated at z=-3/2 and multiplied by -4\textasciicircum{}n/2.
The certified recurrence range is $n\geq 2$.
The original coefficient polynomials are:
\[
p_{0}(n)=n
\]
\[
p_{1}(n)=- 8 n - 6
\]
\[
p_{2}(n)=16 n + 24
\]
The exact identity is derived in the corresponding proof subsection above; all displayed residuals in the machine-readable certificate are zero.

\subsubsection*{A121580 --- proved}
Source: \url{https://oeis.org/A121580}.
First-column cell count for deco polyominoes, with accepted formula one half of sum\_\{j=0\}\textasciicircum{}\{n+1\} j!, minus n!.
The certified recurrence range is $n\geq 5$.
The original coefficient polynomials are:
\[
p_{0}(n)=1
\]
\[
p_{1}(n)=- n - 4
\]
\[
p_{2}(n)=3 n + 3
\]
\[
p_{3}(n)=6 - 4 n
\]
\[
p_{4}(n)=2 n - 6
\]
Using the exact substitution rules in the factorial/parity section, the coefficient residual vectors are
\[
\begin{array}{c|c}n\bmod2&\text{residual vector}\\\hline
0&(0,0)\\
1&(0,0)\\
\end{array}
\]

\subsubsection*{A180195 --- proved}
Source: \url{https://oeis.org/A180195}.
The alternating factorial partial sum in the definition.
The certified recurrence range is $n\geq 4$.
The original coefficient polynomials are:
\[
p_{0}(n)=1
\]
\[
p_{1}(n)=- n - 1
\]
\[
p_{2}(n)=-4
\]
\[
p_{3}(n)=n - 2
\]
Using the exact substitution rules in the factorial/parity section, the coefficient residual vectors are
\[
\begin{array}{c|c}n\bmod2&\text{residual vector}\\\hline
0&(0,0)\\
1&(0,0)\\
\end{array}
\]

\subsubsection*{A121639 --- proved}
Source: \url{https://oeis.org/A121639}.
The deco-polyomino statistic with the accepted recurrence a\_m=m a\_(m-1)+(m-1)!-(m-3)!.
The certified recurrence range is $n\geq 7$.
The original coefficient polynomials are:
\[
p_{0}(n)=1
\]
\[
p_{1}(n)=1 - 2 n
\]
\[
p_{2}(n)=n^{2} - 2 n
\]
\[
p_{3}(n)=2 n - 7
\]
\[
p_{4}(n)=- n^{2} + 8 n - 15
\]
Using the exact substitution rules in the factorial/parity section, the coefficient residual vectors are
\[
\begin{array}{c|c}n\bmod2&\text{residual vector}\\\hline
0&(0,0)\\
1&(0,0)\\
\end{array}
\]

\subsubsection*{A121751 --- proved}
Source: \url{https://oeis.org/A121751}.
The zero column of A121697, satisfying the established second-order floor recurrence.
The certified recurrence range is $n\geq 7$.
The original coefficient polynomials are:
\[
p_{0}(n)=256
\]
\[
p_{1}(n)=-384
\]
\[
p_{2}(n)=- 128 n^{2} + 640 n - 1072
\]
\[
p_{3}(n)=128 n^{2} - 864 n + 1392
\]
\[
p_{4}(n)=16 n^{4} - 224 n^{3} + 968 n^{2} - 1216 n - 84
\]
\[
p_{5}(n)=- 8 n^{4} + 152 n^{3} - 996 n^{2} + 2588 n - 2216
\]
\[
p_{6}(n)=n^{4} - 24 n^{3} + 207 n^{2} - 756 n + 992
\]
Using the exact substitution rules in the factorial/parity section, the coefficient residual vectors are
\[
\begin{array}{c|c}n\bmod2&\text{residual vector}\\\hline
0&(0,0)\\
1&(0,0)\\
\end{array}
\]

\subsubsection*{A121747 --- proved}
Source: \url{https://oeis.org/A121747}.
The first moment of A121745, with the established auxiliary sequence d\_m in its recurrence.
The certified recurrence range is $n\geq 6$.
The original coefficient polynomials are:
\[
p_{0}(n)=38 - 9 n
\]
\[
p_{1}(n)=9 n^{2} - 57 n + 81
\]
\[
p_{2}(n)=9 n^{3} - 37 n^{2} - 24 n + 127
\]
\[
p_{3}(n)=- 9 n^{4} + 102 n^{3} - 392 n^{2} + 525 n - 85
\]
\[
p_{4}(n)=- 28 n^{4} + 307 n^{3} - 1179 n^{2} + 1767 n - 711
\]
\[
p_{5}(n)=- 19 n^{4} + 224 n^{3} - 925 n^{2} + 1512 n - 720
\]
Using the exact substitution rules in the factorial/parity section, the coefficient residual vectors are
\[
\begin{array}{c|c}n\bmod2&\text{residual vector}\\\hline
0&(0,0,0)\\
1&(0,0,0)\\
\end{array}
\]

\subsubsection*{A121584 --- proved}
Source: \url{https://oeis.org/A121584}, \url{https://oeis.org/A121580}, \url{https://oeis.org/A121582}.
The sum of the deco-polyomino cell counts A121580 and A121582.
The certified recurrence range is $n\geq 7$.
The original coefficient polynomials are:
\[
p_{0}(n)=14 - 14 n
\]
\[
p_{1}(n)=14 n^{2} + 1731 n - 6995
\]
\[
p_{2}(n)=- 1731 n^{2} + 1440 n + 21729
\]
\[
p_{3}(n)=5562 n^{2} - 23904 n + 12008
\]
\[
p_{4}(n)=- 5987 n^{2} + 36181 n - 54220
\]
\[
p_{5}(n)=2142 n^{2} - 15434 n + 27464
\]
Using the exact substitution rules in the factorial/parity section, the coefficient residual vectors are
\[
\begin{array}{c|c}n\bmod2&\text{residual vector}\\\hline
0&(0,0,0,0)\\
1&(0,0,0,0)\\
\end{array}
\]

\subsubsection*{A121696 --- proved}
Source: \url{https://oeis.org/A121696}.
The sequence with a\_1=0,a\_2=1 and a\_m-a\_(m-2)=m(m-2)! floor((m-1)/2).
The certified recurrence range is $n\geq 4$.
The original coefficient polynomials are:
\[
p_{0}(n)=1
\]
\[
p_{1}(n)=0
\]
\[
p_{2}(n)=- n^{2} + 2 n - 1
\]
\[
p_{3}(n)=- n^{2} + 2 n
\]
The exact identity is derived in the corresponding proof subsection above; all displayed residuals in the machine-readable certificate are zero.

\subsubsection*{A121553 --- proved}
Source: \url{https://oeis.org/A121553}.
The accepted harmonic-number expression n! (H\_n+n(n-1)/4).
The certified recurrence range is $n\geq 5$.
The original coefficient polynomials are:
\[
p_{0}(n)=1
\]
\[
p_{1}(n)=- 2 n - 3
\]
\[
p_{2}(n)=n^{2} + 4 n - 3
\]
\[
p_{3}(n)=- 2 n^{2} + 2 n + 6
\]
\[
p_{4}(n)=2 n^{2} - 12 n + 18
\]
The exact identity is derived in the corresponding proof subsection above; all displayed residuals in the machine-readable certificate are zero.

\subsubsection*{A187252 --- proved}
Source: \url{https://oeis.org/A187252}.
The first moment of A187250, with the exponential generating function stated in the proof.
The certified recurrence range is $n\geq 5$.
The original coefficient polynomials are:
\[
p_{0}(n)=1
\]
\[
p_{1}(n)=- 2 n - 3
\]
\[
p_{2}(n)=n^{2} + 6 n - 7
\]
\[
p_{3}(n)=- 4 n^{2} + 8 n + 4
\]
\[
p_{4}(n)=4 n^{2} - 24 n + 36
\]
The exact identity is derived in the corresponding proof subsection above; all displayed residuals in the machine-readable certificate are zero.

\subsubsection*{A001277 --- proved}
Source: \url{https://oeis.org/A001277}, \url{https://oeis.org/A136123}.
Column k=1 of the permutation triangle A136123.
The certified recurrence range is $n\geq 5$.
The original coefficient polynomials are:
\[
p_{0}(n)=1
\]
\[
p_{1}(n)=- n
\]
\[
p_{2}(n)=0
\]
\[
p_{3}(n)=n - 1
\]
The derivation from A136123 gives $F=A F_0+B$.
\[
A=- \frac{1}{\left(x - 1\right) \left(x + 1\right)}
\]
\[
B=\frac{1}{x - 1}
\]
Reduction using (GF4) gives $\mathcal L F=P(x)$ and zero coefficient of $F_0$, with
\[
P(x)=x^{2}
\]

\subsubsection*{A001278 --- proved}
Source: \url{https://oeis.org/A001278}, \url{https://oeis.org/A136123}.
Column k=2 of the permutation triangle A136123.
The certified recurrence range is $n\geq 10$.
The original coefficient polynomials are:
\[
p_{0}(n)=1
\]
\[
p_{1}(n)=- n - 3
\]
\[
p_{2}(n)=3 n - 6
\]
\[
p_{3}(n)=2 n - 2
\]
\[
p_{4}(n)=17 - 6 n
\]
\[
p_{5}(n)=5 - n
\]
\[
p_{6}(n)=3 n - 12
\]
The derivation from A136123 gives $F=A F_0+B$.
\[
A=\frac{1 - 3 x}{2 \left(x - 1\right)^{2} \left(x + 1\right)}
\]
\[
B=\frac{2 x^{2} + 2 x - 1}{2 \left(x - 1\right)^{2} \left(x + 1\right)}
\]
Reduction using (GF4) gives $\mathcal L F=P(x)$ and zero coefficient of $F_0$, with
\[
P(x)=x^{4} \left(3 x + 1\right)
\]

\subsubsection*{A001279 --- proved}
Source: \url{https://oeis.org/A001279}, \url{https://oeis.org/A136123}.
Column k=3 of the permutation triangle A136123.
The certified recurrence range is $n\geq 12$.
The original coefficient polynomials are:
\[
p_{0}(n)=18 - 3 n
\]
\[
p_{1}(n)=3 n^{2} - 23 n + 67
\]
\[
p_{2}(n)=5 n^{2} - 58 n + 179
\]
\[
p_{3}(n)=n^{2} - 51 n + 128
\]
\[
p_{4}(n)=- 17 n^{2} + 183 n - 364
\]
\[
p_{5}(n)=- 4 n^{2} + 60 n - 176
\]
\[
p_{6}(n)=12 n^{2} - 108 n + 240
\]
The derivation from A136123 gives $F=A F_0+B$.
\[
A=\frac{\left(1 - 2 x\right) \left(6 x - 1\right)}{6 \left(x - 1\right)^{3} \left(x + 1\right)}
\]
\[
B=\frac{6 x^{4} + 12 x^{3} - 2 x^{2} - 6 x + 1}{6 \left(x - 1\right)^{3} \left(x + 1\right)^{2}}
\]
Reduction using (GF4) gives $\mathcal L F=P(x)$ and zero coefficient of $F_0$, with
\[
P(x)=0
\]

\subsubsection*{A001280 --- proved}
Source: \url{https://oeis.org/A001280}, \url{https://oeis.org/A136123}.
Column k=4 of the permutation triangle A136123.
The certified recurrence range is $n\geq 13$.
The original coefficient polynomials are:
\[
p_{0}(n)=n^{2} - 17 n + 72
\]
\[
p_{1}(n)=- n^{3} + 15 n^{2} - 52 n - 18
\]
\[
p_{2}(n)=2 n^{3} - 33 n^{2} + 143 n - 118
\]
\[
p_{3}(n)=5 n^{2} - 27 n + 16
\]
\[
p_{4}(n)=- 2 n^{3} + 24 n^{2} - 94 n + 120
\]
\[
p_{5}(n)=n^{3} - 12 n^{2} + 47 n - 60
\]
The derivation from A136123 gives $F=A F_0+B$.
\[
A=\frac{- 60 x^{3} + 60 x^{2} - 15 x + 1}{24 \left(x - 1\right)^{4} \left(x + 1\right)}
\]
\[
B=\frac{24 x^{6} + 72 x^{5} + 12 x^{4} - 66 x^{3} - 19 x^{2} + 12 x - 1}{24 \left(x - 1\right)^{4} \left(x + 1\right)^{3}}
\]
Reduction using (GF4) gives $\mathcal L F=P(x)$ and zero coefficient of $F_0$, with
\[
P(x)=0
\]

\subsubsection*{A045742 --- proved}
Source: \url{https://oeis.org/A045742}.
The established binomial sum for the face-count sequence: m=n-1 and the summation formula given below.
The certified recurrence range is $n\geq 4$.
The original coefficient polynomials are:
\[
p_{0}(n)=523 n^{3} - 2516 n^{2} + 1993 n
\]
\[
p_{1}(n)=- 4554 n^{3} + 40668 n^{2} - 112854 n + 76740
\]
\[
p_{2}(n)=- 56484 n^{3} + 328212 n^{2} - 505608 n + 150240
\]
\[
p_{3}(n)=491832 n^{3} - 4883976 n^{2} + 16099344 n - 17614080
\]
For (GF3), the exact input is
\[
x=\frac{t \left(1 - t\right)}{\left(t + 1\right)^{3}}
\]
with $F=\sum_i R_i\sqrt{D_i}$ and
\[
R_{0}=\frac{t^{3} \left(1 - t\right)}{\left(t + 1\right)^{6} \left(- \frac{3 t \left(1 - t\right)}{\left(t + 1\right)^{4}} - \frac{t}{\left(t + 1\right)^{3}} + \frac{1 - t}{\left(t + 1\right)^{3}}\right)}
\]
\[
D_{0}=1
\]
The exact operator identity is $\mathcal L F=P(x)$, with every nonconstant radical coefficient zero, where
\[
P(x)=- 2544 x^{3}
\]

\subsubsection*{A121752 --- proved}
Source: \url{https://oeis.org/A121752}.
The first moment of A121697, satisfying the established inhomogeneous second-order recurrence.
The certified recurrence range is $n\geq 7$.
The original coefficient polynomials are:
\[
p_{0}(n)=1223 - 460 n
\]
\[
p_{1}(n)=460 n^{2} + 460 n - 5569
\]
\[
p_{2}(n)=460 n^{3} - 3826 n^{2} + 7853 n + 1313
\]
\[
p_{3}(n)=- 460 n^{4} + 1840 n^{3} + 5501 n^{2} - 31045 n + 31726
\]
\[
p_{4}(n)=1223 n^{4} - 13205 n^{3} + 45787 n^{2} - 51389 n + 558
\]
\[
p_{5}(n)=- 852 n^{4} + 14150 n^{3} - 85636 n^{2} + 224000 n - 213312
\]
Using the exact substitution rules in the factorial/parity section, the coefficient residual vectors are
\[
\begin{array}{c|c}n\bmod2&\text{residual vector}\\\hline
0&(0,0)\\
1&(0,0)\\
\end{array}
\]

\subsubsection*{A027050 --- proved}
Source: \url{https://oeis.org/A027050}, \url{https://oeis.org/A027023}.
T(n,2n-1) in the tribonacci triangle A027023.
The certified recurrence range is $n\geq 12$.
The original coefficient polynomials are:
\[
p_{0}(n)=1 - n
\]
\[
p_{1}(n)=6 n - 9
\]
\[
p_{2}(n)=10 - 7 n
\]
\[
p_{3}(n)=38 - 8 n
\]
\[
p_{4}(n)=5 n - 23
\]
\[
p_{5}(n)=2 n - 5
\]
\[
p_{6}(n)=3 n - 12
\]
For (GF3), the exact input is
\[
x=\frac{t}{t^{2} + t + 1}
\]
with $F=\sum_i R_i\sqrt{D_i}$ and
\[
R_{0}=\frac{- 3 t^{4} - 3 t^{3} - 3 t^{2} - t}{t^{5} - t^{4} - 2 t^{2} - t - 1}
\]
\[
D_{0}=1
\]
The exact operator identity is $\mathcal L F=P(x)$, with every nonconstant radical coefficient zero, where
\[
P(x)=2 x^{3} \left(3 x^{3} + x^{2} - 3 x + 3\right)
\]

\subsubsection*{A027058 --- proved}
Source: \url{https://oeis.org/A027058}, \url{https://oeis.org/A027052}.
T(n,2n-2) in the tribonacci triangle A027052.
The certified recurrence range is $n\geq 12$.
The original coefficient polynomials are:
\[
p_{0}(n)=n
\]
\[
p_{1}(n)=4 - 7 n
\]
\[
p_{2}(n)=13 n - 10
\]
\[
p_{3}(n)=n - 34
\]
\[
p_{4}(n)=84 - 13 n
\]
\[
p_{5}(n)=3 n - 32
\]
\[
p_{6}(n)=6 - n
\]
\[
p_{7}(n)=3 n - 18
\]
For (GF3), the exact input is
\[
x=\frac{t}{t^{2} + t + 1}
\]
with $F=\sum_i R_i\sqrt{D_i}$ and
\[
R_{0}=\frac{- 2 t^{5} - 3 t^{4} - 3 t^{3} - t^{2} - t}{t^{5} - t^{4} - 2 t^{2} - t - 1}
\]
\[
D_{0}=1
\]
The exact operator identity is $\mathcal L F=P(x)$, with every nonconstant radical coefficient zero, where
\[
P(x)=x \left(3 x^{6} - 9 x^{4} - 24 x^{3} + 21 x^{2} - 8 x + 1\right)
\]

\subsubsection*{A027059 --- proved}
Source: \url{https://oeis.org/A027059}, \url{https://oeis.org/A027052}.
T(n,2n-3) in the tribonacci triangle A027052, with offset 2.
The certified recurrence range is $n\geq 12$.
The original coefficient polynomials are:
\[
p_{0}(n)=n + 1
\]
\[
p_{1}(n)=- 8 n - 2
\]
\[
p_{2}(n)=19 n - 5
\]
\[
p_{3}(n)=- 6 n - 30
\]
\[
p_{4}(n)=123 - 21 n
\]
\[
p_{5}(n)=8 n - 58
\]
\[
p_{6}(n)=n + 1
\]
\[
p_{7}(n)=6 n - 30
\]
For (GF3), the exact input is
\[
x=\frac{t}{t^{2} + t + 1}
\]
with $F=\sum_i R_i\sqrt{D_i}$ and
\[
R_{0}=\frac{- 2 t^{6} - 4 t^{5} - 2 t^{4} - 2 t^{3}}{t^{5} - t^{4} - 2 t^{2} - t - 1}
\]
\[
D_{0}=1
\]
The exact operator identity is $\mathcal L F=P(x)$, with every nonconstant radical coefficient zero, where
\[
P(x)=- 2 x^{3} \left(12 x^{4} + 7 x^{3} - 18 x^{2} + 19 x - 4\right)
\]

\subsubsection*{A027088 --- refuted}
Source: \url{https://oeis.org/A027088}, \url{https://oeis.org/A027082}.
T(n,2n-1) in the tribonacci triangle A027082, with offset 1.
The original coefficient polynomials are:
\[
p_{0}(n)=54 n^{2} - 201 n - 255
\]
\[
p_{1}(n)=292 n^{2} - 1316 n - 255
\]
\[
p_{2}(n)=- 218 n^{2} + 1213 n - 1056
\]
\[
p_{3}(n)=- 496 n^{2} + 3332 n - 4446
\]
\[
p_{4}(n)=- 50 n^{2} - 497 n + 3243
\]
\[
p_{5}(n)=- 52 n^{2} + 80 n + 765
\]
\[
p_{6}(n)=66 n^{2} - 645 n + 1494
\]
For (GF3), the exact input is
\[
x=x
\]
with $F=\sum_i R_i\sqrt{D_i}$ and
\[
R_{0}=\frac{x^{4} - 2 x^{3} - 4 x^{2} - 6 x + 3}{2 x \left(x^{3} + x^{2} + 3 x - 1\right)}
\]
\[
D_{0}=1
\]
\[
R_{1}=\frac{\left(x - 1\right) \left(x^{2} + 3\right)}{2 x \left(x^{3} + x^{2} + 3 x - 1\right)}
\]
\[
D_{1}=- 3 x^{2} - 2 x + 1
\]
At $n=7$ the exact recurrence residual is $556944$. Moreover $\mathcal L_pF=U_0+U_1\sqrt{1-2x-3x^2}$, with the following nonzero $U_1$; this proves the eventual obstruction.
Write $U_1=N_1/D_1^*$.
\[
\begin{aligned}
N_1&=- 2295 x^{12} - 9180 x^{11}\\
&\quad - 30030 x^{10} - 23112 x^{9}\\
&\quad - 73269 x^{8} - 77880 x^{7}\\
&\quad - 106692 x^{6} - 42480 x^{5}\\
&\quad +36015 x^{4} + 13428 x^{3}\\
&\quad +3474 x^{2} - 5160 x\\
&\quad +765
\end{aligned}
\]
\[
D_1^*=\left(x + 1\right)^{2} \left(3 x - 1\right)^{2} \left(x^{3} + x^{2} + 3 x - 1\right)^{3}
\]

\subsubsection*{A027089 --- proved}
Source: \url{https://oeis.org/A027089}, \url{https://oeis.org/A027082}.
T(n,2n-2) in the tribonacci triangle A027082, with offset 2.
The certified recurrence range is $n\geq 12$.
The original coefficient polynomials are:
\[
p_{0}(n)=- 244 n^{2} + 577 n + 2130
\]
\[
p_{1}(n)=1266 n^{2} - 3843 n - 5109
\]
\[
p_{2}(n)=- 718 n^{2} + 2422 n - 2840
\]
\[
p_{3}(n)=- 2348 n^{2} + 12638 n - 14820
\]
\[
p_{4}(n)=- 760 n^{2} + 2173 n + 1616
\]
\[
p_{5}(n)=- 502 n^{2} + 2245 n - 1123
\]
\[
p_{6}(n)=138 n^{2} - 1092 n + 2010
\]
For (GF3), the exact input is
\[
x=\frac{t}{t^{2} + t + 1}
\]
with $F=\sum_i R_i\sqrt{D_i}$ and
\[
R_{0}=- \frac{t}{t^{2} + t + 1} + \frac{- 3 t^{5} - 3 t^{4} - 4 t^{3} - t^{2} - t}{t^{3} - t^{2} - t - 1}
\]
\[
D_{0}=1
\]
The exact operator identity is $\mathcal L F=P(x)$, with every nonconstant radical coefficient zero, where
\[
P(x)=- x^{2} \left(1128 x^{5} - 5725 x^{4} - 6519 x^{3} - 1711 x^{2} - 3081 x - 2308\right)
\]


\clearpage
\section{Validation and prior-work exclusions}
\label{round1:sec:validation}
\paragraph{Structural and arithmetic validation.}
The quadratic-character proof was independently reviewed, including its
primitive reduction, congruences, and exceptional discriminants. Exact direct
evaluation checked the divisibility and nonvanishing normalizations on 300
admissible discriminants up to 600. Matrices constructed from the self-similarity
condition independently matched ordinary coin-change counts in nine binary
rows and seven ternary rows. The universal square formula was checked as a
polynomial identity in its parameter through index 23, in addition to the
characteristic-polynomial proof.

The arithmetic verifier independently enumerated partitions, permutations,
digital representations, and divisor sums within its documented bounds.
The prime-exponent first-occurrence theorem was checked against every index
up to 65537; its proof establishes the bound for all indices. For A216971,
every function on up to six labelled elements was enumerated. All binary
words through length 16 were checked for A334592. The consecutive-square XOR
map was checked for injectivity through index 200000, and its construction
was checked for every target congruent to 1 modulo 4 below 65536.

\paragraph{Exact certificate validation.}
The floor verifiers reconstruct their rational bounds and integer polynomial
coefficients, including finite exceptional indices. They compare complete
rational generating functions. The generating-function verifier applies
the recorded differential operators with exact arithmetic, retaining the
boundary polynomial and the correct ordinary or exponential normalization.
Each counterexample is evaluated with exact integers or rational numbers.
Program outputs are included so that the validation can be repeated.

\paragraph{Examples excluded because of earlier work.}
These exclusions illustrate why an old conjecture label alone was insufficient:
\begin{itemize}
\item A111967: a linked 2021 proof already settles the relevant matrix row-sum
assertion; related c=2 copies were not counted.
\item A054120: the first McGarvey assertion was proved by Hendel in 2015;
only the second assertion is selected here.
\item A224515: existence was already proved in 2026; uniqueness is selected.
\item A113249: Israel's 2017 proof covers m=3; only the universal parameter
assertion is selected, with all its copies counted once.
\item A024219, A024401, A075002, A105696, and the inspected A220119 divisibility
assertion had recorded prior resolutions, including updates in 2026.
\item A203233, A269027, A155828, A001969, A008364, A008365, and A008486
also had prior proofs or refutations and were not included as new selections.
\end{itemize}
The complete audit scope and its access limitation are recorded in
\texttt{oeis\_\allowbreak open\_\allowbreak prior\_\allowbreak work\_\allowbreak audit.json}. This dossier makes no
claim that all relevant literature has been searched exhaustively.

\paragraph{Standard identities used.}
The quadratic character proof uses the classical quadratic Gauss sum,
Dirichlet Euler product, and Fourier expansion of the second Bernoulli
polynomial. Reference sections are NIST DLMF
\url{https://dlmf.nist.gov/27.10},
\url{https://dlmf.nist.gov/25.15}, and
\url{https://dlmf.nist.gov/24.8}.
Other standard tools used explicitly in the proofs are Lagrange inversion,
Newton's symmetric-function identities, the finite q-binomial theorem,
and Bertrand's postulate. The report claims no novelty for these tools.
% ===== END ORIGINAL ROUND 1 =====

% ===== BEGIN ORIGINAL ROUND 2 =====
\clearpage
\setcounter{section}{0}\setcounter{equation}{0}\setcounter{theorem}{0}
\setcounter{table}{0}\setcounter{figure}{0}
\normalsize\renewcommand{\arraystretch}{1}\renewcommand{\baselinestretch}{1}\selectfont
\part[Round 2: Exact sums, prime-power matrices, connected arrays, and recurrences]{Exact sums, prime-power matrices, connected arrays, and recurrences}
\label{combined:round2}
\rhead{\small Round 2 / KEIO}
\section*{Scope, counting, and how to use the packet}
\addcontentsline{toc}{section}{Scope, counting, and how to use the packet}
This second round provides 101 selected conjecture or empirical
assertion groups, with no counted entry repeated from the first round.
The headline total consists of 100 proof-only groups and one mixed
group, \oeis{A004737}: its proposed matrix property has a small exact
counterexample, while its normalized Gram-matrix pattern is proved for
every size. The cumulative total under the same counting convention is
208 selected groups.

\paragraph{The unit being counted.}
Each row of the catalogue is one entry-level selection. Multiple related
conjectures in one entry count once. Fixed-width or fixed-color assertions
on different entries count separately even when one general theorem
proves them together. Copied universal assertions and aliases receive
no extra count. Thus 22 connected-array selections share three methods;
they are 22 source assertions, not 22 unrelated mathematical discoveries.

\paragraph{Selected scope is explicit.}
Several entries contain further questions that are outside the selected
result. For example, \oeis{A376336} is resolved only for existence of a
tau-zero modulo each modulus, not for its proposed quadratic bound.
\oeis{A290815} is resolved for the odd-modulus reciprocal-square
equivalence, not for the general-exponent extension. The source wording,
formal statement, applicable range, and proof are tied together in the
ledger and the following sections.

\paragraph{Qualifications do not add to the total.}
The normalized transform theorem for \oeis{A153197} is included as an
uncounted note: its normalization is essential, and the source also
records established fixed-point formulas. Index-only or definition-only
repairs, such as \oeis{A395454} and \oeis{A178113}, are excluded.
The independent divisibility proof found for \oeis{A098275} was excluded
after identifying an earlier correct proof by Roberto Blanco G\'omez,
submitted on 14 October 2026:
\url{https://aixiv.online/abs/2610.01676}.

\paragraph{What is verified.}
The infinite assertions are proved by algebra, induction, bijections,
global inequalities, or complete finite-state certificates. An initial
segment agreeing with a guessed recurrence is never the proof. In the
floor family, the finite state space is exhaustive, and exact matrix or
rational-function identities imply every subsequent coefficient. In the
generating-function family, exact differential identities imply the
stated recurrence beyond an explicitly checked boundary polynomial.
Finite numerical or rational comparisons are labeled supplementary
regression checks.

\paragraph{Source audit.}
The inspected entries still mark the selected statements conjectural or
empirical, with the qualifications stated here. Later confirmations in
the source text and identified linked literature were removed from the
count. A second check compared all selected IDs with the 97 distinct
entry IDs appearing in the 100-claim Appendix A of Tom Adamczewski's
\emph{OEIS Open}; there is no overlap. That appendix covers the 100
costliest of 153 combined-model resolutions, not every result in its
repositories or every published OEIS proof. The audit therefore does
not establish worldwide novelty. Its exact scope is retained in
\texttt{r2\_prior\_work\_audit.json}.

\paragraph{Files.}
The supplied \texttt{oeis\_hunt\_round2.tex} is standalone and contains
this complete report. The ZIP includes the ledger, proof fragments,
full certificates, the programs that check them, and the validation
outputs. The report can be recompiled with pdfLaTeX; verification uses
Python and SymPy. See \texttt{README\_ROUND2.txt}. No OEIS edit or
submission is represented by this packet.
\clearpage
\section{Catalogue of the selected assertions}
\label{round2:sec:r2catalogue}
\textbf{P}: proved. \textbf{P/R}: some selected assertions proved and
another refuted in the same entry. Each entry ID links to its primary
source; the following sections contain the proof and exact scope.

\Needspace{6\baselineskip}\subsection*{Structural identities and exact sums (20)}
\begin{longtable}{@{}r l c p{0.62\textwidth}@{}}
\toprule No.&Entry&Result&Selected assertion\\\midrule\endhead
1&\oeis{A350900}&P&Double-gcd convolution for arbitrary arithmetic functions\\
2&\oeis{A347293}&P&Residue-product gcd convolution and multiplicativity\\
3&\oeis{A130540}&P&Even-column and alternating row sums\\
4&\oeis{A001222}&P&Complex two-parameter convolution identity\\
5&\oeis{A265649}&P&Both finite sums for the generalized Bessel triangle\\
6&\oeis{A097724}&P&Chebyshev translation coefficients\\
7&\oeis{A102086}&P&The row-polynomial recurrence\\
8&\oeis{A379907}&P&General Riordan convolution and column recurrence\\
9&\oeis{A380191}&P&Central-binomial weighted columns\\
10&\oeis{A004737}&P/R&E-matrix refutation and exact Gram-matrix formula\\
11&\oeis{A005773}&P&Order-comparable coprime divisors and Motzkin meanders\\
12&\oeis{A376182}&P&The triangular enumeration is a permutation\\
13&\oeis{A378277}&P&Alternating harmonic Fibonacci row sums\\
14&\oeis{A349802}&P&Exact Euler transforms of binary Lyndon columns\\
15&\oeis{A163774}&P&A single binomial sum for the defining double sum\\
16&\oeis{A152920}&P&Binomially weighted rows and A193844\\
17&\oeis{A105794}&P&Factorial and Bernoulli weighted sums\\
18&\oeis{A132393}&P&The full Vandermonde-squared nested-sum identity\\
19&\oeis{A397048}&P&The first signed-Stirling convolution conjecture\\
20&\oeis{A112302}&P&Limit of the Lerch-transcendent recurrence\\
\bottomrule\end{longtable}

\Needspace{6\baselineskip}\subsection*{Arithmetic and congruences (11)}
\begin{longtable}{@{}r l c p{0.62\textwidth}@{}}
\toprule No.&Entry&Result&Selected assertion\\\midrule\endhead
21&\oeis{A334595}&P&XOR-triangle right-boundary records and value one\\
22&\oeis{A133910}&P&Divisibility by the index\\
23&\oeis{A051283}&P&The A346418 exponent characterization\\
24&\oeis{A354995}&P&The nonzero positions of a unitary-divisor difference\\
25&\oeis{A376336}&P&Existence of a zero of tau modulo every modulus\\
26&\oeis{A087726}&P&All prime-power counts of square-zero matrices\\
27&\oeis{A322327}&P&The general parameter multiplicative function\\
28&\oeis{A033955}&P&Strict bound for the sum of prime remainders\\
29&\oeis{A166306}&P&Nondivisibility by 3 of the Bernoulli-prime sums\\
30&\oeis{A234307}&P&Alternating gcd-sum identity\\
31&\oeis{A290815}&P&Odd-modulus reciprocal-square characterization\\
\bottomrule\end{longtable}

\Needspace{6\baselineskip}\subsection*{Connected arrays and floor identities (28)}
\begin{longtable}{@{}r l c p{0.62\textwidth}@{}}
\toprule No.&Entry&Result&Selected assertion\\\midrule\endhead
32&\oeis{A166756}&P&Empirical connected-coloring generating function\\
33&\oeis{A166757}&P&Empirical connected-coloring generating function\\
34&\oeis{A166758}&P&Empirical connected-coloring generating function\\
35&\oeis{A166759}&P&Empirical connected-coloring generating function\\
36&\oeis{A166760}&P&Empirical connected-coloring generating function\\
37&\oeis{A166762}&P&The ninth-order connected-coloring recurrence\\
38&\oeis{A166776}&P&Empirical connected-coloring generating function\\
39&\oeis{A166777}&P&Empirical connected-coloring generating function\\
40&\oeis{A166778}&P&Empirical connected-coloring generating function\\
41&\oeis{A166779}&P&Empirical connected-coloring generating function\\
42&\oeis{A166780}&P&Empirical connected-coloring generating function\\
43&\oeis{A166781}&P&Empirical connected-coloring generating function\\
44&\oeis{A166782}&P&Empirical connected-coloring generating function\\
45&\oeis{A166783}&P&Empirical connected-coloring generating function\\
46&\oeis{A166786}&P&Empirical connected-coloring generating function\\
47&\oeis{A166796}&P&Empirical connected-coloring generating function\\
48&\oeis{A166797}&P&Empirical connected-coloring generating function\\
49&\oeis{A166798}&P&Empirical connected-coloring generating function\\
50&\oeis{A166799}&P&Empirical connected-coloring generating function\\
51&\oeis{A166800}&P&Empirical connected-coloring generating function\\
52&\oeis{A166801}&P&Empirical connected-coloring generating function\\
53&\oeis{A166802}&P&Empirical connected-coloring generating function\\
54&\oeis{A024172}&P&The empirical symmetric-function floor series\\
55&\oeis{A024180}&P&The empirical shifted symmetric-function floor series\\
56&\oeis{A184534}&P&Cube-root fractional-part reciprocal series\\
57&\oeis{A002620}&P&Exponential floor identity with parameter 4\\
58&\oeis{A001840}&P&Exponential floor identity with parameter 6\\
59&\oeis{A227018}&P&Quartic-binomial harmonic means and their series\\
\bottomrule\end{longtable}

\Needspace{6\baselineskip}\subsection*{Generating functions and recurrences (42)}
\begin{longtable}{@{}r l c p{0.62\textwidth}@{}}
\toprule No.&Entry&Result&Selected assertion\\\midrule\endhead
60&\oeis{A089354}&P&The specified polynomial-coefficient recurrence\\
61&\oeis{A247471}&P&The specified polynomial-coefficient recurrence\\
62&\oeis{A135925}&P&The specified polynomial-coefficient recurrence\\
63&\oeis{A218045}&P&The specified polynomial-coefficient recurrence\\
64&\oeis{A006343}&P&The specified polynomial-coefficient recurrence\\
65&\oeis{A158826}&P&The specified polynomial-coefficient recurrence\\
66&\oeis{A075045}&P&The specified polynomial-coefficient recurrence\\
67&\oeis{A082397}&P&The specified polynomial-coefficient recurrence\\
68&\oeis{A190160}&P&The specified polynomial-coefficient recurrence\\
69&\oeis{A001723}&P&The specified polynomial-coefficient recurrence\\
70&\oeis{A001724}&P&The specified polynomial-coefficient recurrence\\
71&\oeis{A027051}&P&The specified polynomial-coefficient recurrence\\
72&\oeis{A027027}&P&The specified polynomial-coefficient recurrence\\
73&\oeis{A027028}&P&The specified polynomial-coefficient recurrence\\
74&\oeis{A027056}&P&The specified polynomial-coefficient recurrence\\
75&\oeis{A027087}&P&The specified polynomial-coefficient recurrence\\
76&\oeis{A025258}&P&The specified polynomial-coefficient recurrence\\
77&\oeis{A157125}&P&The specified polynomial-coefficient recurrence\\
78&\oeis{A186335}&P&The specified polynomial-coefficient recurrence\\
79&\oeis{A133158}&P&The specified polynomial-coefficient recurrence\\
80&\oeis{A157021}&P&The specified polynomial-coefficient recurrence\\
81&\oeis{A173993}&P&The specified polynomial-coefficient recurrence\\
82&\oeis{A141353}&P&The specified polynomial-coefficient recurrence\\
83&\oeis{A078999}&P&The specified polynomial-coefficient recurrence\\
84&\oeis{A000159}&P&The specified polynomial-coefficient recurrence\\
85&\oeis{A000181}&P&The specified polynomial-coefficient recurrence\\
86&\oeis{A000185}&P&The specified polynomial-coefficient recurrence\\
87&\oeis{A080894}&P&The specified polynomial-coefficient recurrence\\
88&\oeis{A272261}&P&The specified polynomial-coefficient recurrence\\
89&\oeis{A144575}&P&The specified polynomial-coefficient recurrence\\
90&\oeis{A053532}&P&The specified polynomial-coefficient recurrence\\
91&\oeis{A250917}&P&The specified polynomial-coefficient recurrence\\
92&\oeis{A292461}&P&The specified polynomial-coefficient recurrence\\
93&\oeis{A126089}&P&The specified polynomial-coefficient recurrence\\
94&\oeis{A108629}&P&The specified polynomial-coefficient recurrence\\
95&\oeis{A239230}&P&The specified polynomial-coefficient recurrence\\
96&\oeis{A032115}&P&The specified polynomial-coefficient recurrence\\
97&\oeis{A168452}&P&The specified polynomial-coefficient recurrence\\
98&\oeis{A005573}&P&Terminating hypergeometric formula\\
99&\oeis{A368633}&P&Odd coefficients occur exactly at indices 2\textasciicircum{}k-1\\
100&\oeis{A027116}&P&Rational series and the constant-coefficient recurrence\\
101&\oeis{A027117}&P&Rational series and the constant-coefficient recurrence\\
\bottomrule\end{longtable}

\clearpage
% Source fragment: r2_root_proofs.tex
\section{Structural identities, transforms, and exact sums}
\label{round2:sec:r2root}
This section resolves 20 counted entry-level selections and includes one
uncounted normalization note. Unless a different range
is stated, triangular arrays have indices $0\le k\le n$. Every source is
the corresponding linked OEIS entry, inspected on 20 October 2026.
The statements below specify the selected assertions; unrelated conjectures
in the same entry are not implicitly included.

\subsection{\texorpdfstring{\oeis{A350900}: a double-gcd convolution}{A350900: a double-gcd convolution}}
For an arbitrary arithmetic function $f$, put
\[
 A_f(n)=\sum_{i,k=1}^{n}
 f\left(\frac{\gcd(i,n)}{\gcd(i,k,n)}\right),\qquad
 C(n)=\sum_{d\mid n}d\varphi(d).
\]
The two selected conjectures follow from the identity
\[
 \boxed{A_f=(f\varphi)*C.}
\]
Here multiplication inside $f\varphi$ is pointwise and $*$ is Dirichlet
convolution. To prove it, fix $d=\gcd(i,n)$; there are $\varphi(n/d)$ choices
of $i$. For $e\mid d$, exactly $(n/d)\varphi(e)$ choices of $k\in[1,n]$
satisfy $d/\gcd(d,k)=e$. Consequently
\[
 A_f(n)=\sum_{d\mid n}\frac nd\varphi(n/d)
                 \sum_{e\mid d}f(e)\varphi(e).
\]
Associativity gives $(\operatorname{id}\,\varphi)*1*(f\varphi)$, which is
the displayed result. If $f$ is multiplicative, all three factors are
multiplicative. Taking $f(n)=n$ gives the stated row-sum convolution
\oeis{A002618}$*$\oeis{A057660}. The auxiliary identity defining $C$ is
already recorded in \oeis{A057660}; it is an input, not a new selection.

\subsection{\texorpdfstring{\oeis{A347293}: a residue-product convolution}{A347293: a residue-product convolution}}
For any arithmetic function $f$, including nonmultiplicative ones, the
selected formula is
\[
 \boxed{\sum_{a,b\bmod n}f(\gcd(1+ab,n))
 =\sum_{d\mid n}(n/d)^2\varphi(d)(\mu*f)(d).}
\]
Indeed $f(m)=\sum_{d\mid m}(\mu*f)(d)$. For each $d\mid n$, the congruence
$ab\equiv-1\pmod d$ has $\varphi(d)$ solutions modulo $d$: $a$ must be a
unit and then $b$ is determined. Each solution has $(n/d)^2$ lifts modulo
$n$. Interchanging finite sums proves the formula. Thus the function is
$\operatorname{id}^{,2}*(\varphi(\mu*f))$, and is multiplicative when
$f$ is. For $f=\operatorname{id}$, this becomes
$\operatorname{id}^{,2}*\varphi^2$, giving the proposed identification
with \oeis{A127473}. The argument includes $n=1$.

\subsection{\texorpdfstring{\oeis{A130540}: even and alternating row sums}{A130540: even and alternating row sums}}
The array is $T(n,k)=\sigma(n/k)$ if $k\mid n$, and zero otherwise, for
$1\le k\le n$. Let $R(n)=\sum_{d\mid n}\sigma(d)$, and set $R(t)=0$ when
$t$ is not a positive integer. The even-column sum is exactly
\[
 E(n)=\sum_{2j\mid n}\sigma(n/(2j))=R(n/2).
\]
This proves the second selected conjecture, including the zero values for
odd $n$. The alternating sum, with a positive sign at odd $k$, is
\[
 R(n)-2R(n/2)=\sum_{d\mid n}\sigma_{\rm odd}(d),
\]
because $\sigma_{\rm odd}(d)=\sigma(d)-2\sigma(d/2)$. The last expression
is the defining divisor sum for \oeis{A288417}, proving the third
conjecture. The separate identification of the odd-column sum with
\oeis{A327096} is not counted or asserted here; only its elementary
formula $R(n)-R(n/2)$ follows from this proof.

\subsection{\texorpdfstring{\oeis{A001222}: the two-parameter convolution}{A001222: the two-parameter convolution}}
Let $\Omega$ count prime factors with multiplicity and $\omega$ count
distinct prime factors. For all complex $x,y$ and positive integers $n$,
\[
 (x+y)^{\Omega(n)}
 =\sum_{d\mid n}x^{\Omega(d)}
  (x+y)^{\Omega(n/d)-\omega(n/d)}y^{\omega(n/d)}.
\]
Use $0^0=1$. Both sides define multiplicative functions of $n$, so it
suffices to take $n=p^e$. The right side is
\[
 x^e+y\sum_{i=0}^{e-1}x^i(x+y)^{e-i-1}=(x+y)^e.
\]
This is a polynomial telescoping identity, so no division by $x$, $y$,
or $x+y$ is required. It therefore includes all degenerate parameter
values. The source already credits the normalized case $x+y=1$ to
earlier work; the selected assertion is the displayed general-parameter
version, with that antecedent explicitly retained.

\subsection{\texorpdfstring{\oeis{A265649}: two finite sums}{A265649: two finite sums}}
Write $p_n(x)=\sum_k T(n,k)x^k$, where $T(0,0)=1$ and
$T(n,k)=T(n-1,k-1)+(2n-1+k)T(n-1,k)$. Its exponential generating
function is
\[
 E(z,x)=\sum_{n\ge0}p_n(x)\frac{z^n}{n!}
       =g\exp(x(g-1)),\qquad g=(1-2z)^{-1/2}.
\]
For example, the recurrence gives
$p_{n+1}=(2n+1+x)p_n+xp_n'$, and direct differentiation verifies the
corresponding differential equation and the initial value $E(0,x)=1$.
Hence $T(n,k)=\frac{n!}{k!}[z^n]g(g-1)^k$.

\paragraph{First sum.}
For $0\le k\le n$,
\[
 T(n,k)=2^{n-k}(n-k)!\binom nk
 \sum_{j=0}^{n-k}(-1/4)^j\binom{2j+k}{j}\binom n{j+k}.
\]
After division by $n!/k!$ and summation in $n$, the proposed right side
becomes
\[
 \frac{z^k}{(1-2z)^{k+1}}
 \sum_{j\ge0}\binom{2j+k}{j}
       \left(\frac{-z}{2(1-2z)}\right)^j.
\]
The Catalan identity
$\sum_{j\ge0}\binom{2j+k}{j}w^j=C(w)^k/\sqrt{1-4w}$, with
$C(w)=2/(1+\sqrt{1-4w})$, follows from the binomial theorem (or Lagrange
inversion). Here $\sqrt{1-4w}=g$. Substitution, using
$2zg^2=g^2-1$, reduces the last display to $g(g-1)^k$, as required.

\paragraph{Second sum.}
Let $B(m,j)$ be the reverse Bessel coefficients \oeis{A001497}, and let
$L(m,k)$ be the signed Laguerre coefficients \oeis{A021009}. For $n\ge1$,
\[
 T(n,k)=(-1)^k\sum_{j=0}^{n-1}B(n-1,j)L(j+1,k).
\]
Put $R_m(x)=m!L_m(-x)=\sum_k m!\binom mk x^k/k!$. Then
\[
 \sum_{m\ge0}R_m(x)\frac{y^m}{m!}
 =\frac{1}{1-y}\exp\left(\frac{xy}{1-y}\right),
 \qquad
 \sum_{m\ge0}\sum_j B(m,j)y^j\frac{z^m}{m!}
 =g\exp(y(1-1/g)).
\]
Consequently the derivative, with respect to $z$, of the exponential
generating function of the proposed rows is
\[
 g\sum_{j\ge0}R_{j+1}(x)\frac{(1-1/g)^j}{j!}
 =g^3(1+xg)e^{x(g-1)}=E_z(z,x).
\]
This proves all rows $n\ge1$. The empty sum at $n=0$ is not claimed to
equal $T(0,0)$; its initial row is stated separately.

\subsection{\texorpdfstring{\oeis{A097724}: Chebyshev translation}{A097724: Chebyshev translation}}
Let $u_n(x)=U_n(x/2)$, so that $u_0=1$, $u_1=x$, and
$u_n=xu_{n-1}-u_{n-2}$. The selected connection-coefficient conjecture is
\[
 \boxed{u_n(x+1)=\sum_{k=0}^n T(n,k)u_k(x).}
\]
In this orientation it agrees with the examples in the entry, whose
description reverses a row when listing the basis in descending degree.
The known column generating functions are $G(z)(zG(z))^k$, where
\[
 z^2G^2-(1-z+z^2)G+1=0,\qquad G(0)=1.
\]
Therefore
\[
 \sum_{n\ge0}\sum_k T(n,k)u_k(x)z^n
 =\frac{G}{1-xzG+z^2G^2}
 =\frac{1}{1-(x+1)z+z^2}.
\]
The final expression is the generating function of $u_n(x+1)$.
The related entry \oeis{A191579} adds no separate count.

\subsection{\texorpdfstring{\oeis{A102086}: rows of a self-squaring triangle}{A102086: rows of a self-squaring triangle}}
The lower triangular matrix $T$ has diagonal $T(n,n)=n+1$ and
$T(n,k)=(T^2)(n-1,k)$ for $k<n$. If $L$ is the one-step downward shift
and $D=\operatorname{diag}(1,2,3,\ldots)$, this is precisely
\[
 T=D+LT^2.
\]
The conjectured polynomial recurrence is $P(n,n)=x^n$, $P(n,0)=1$, and
$P(n,k)=P(n,k-1)+(k+1)P(n-1,k)$ for $0<k<n$. We prove that
$P(n+1,n)=\sum_jT(n,j)x^j$.

For $n\ge k$, define
\[
 Q(n,k)=\sum_{j=0}^k (T^{\,n-k})(k,j)x^j.
\]
Clearly $Q(k,k)=x^k$ and $Q(n,0)=1$. Multiplying the matrix equation on
the right by $T^{m-1}$ gives $T^m=DT^{m-1}+LT^{m+1}$. Its row $k$,
with $m=n-k\ge1$, says
\[
 Q(n,k)=(k+1)Q(n-1,k)+Q(n,k-1).
\]
Uniqueness in the triangular recurrence gives $Q=P$, and the result
follows by taking $n-k=1$. The same proof works for any diagonal
$D=\operatorname{diag}(1,d_1,d_2,\ldots)$, with $k+1$ replaced by $d_k$.
The source links an earlier proof of a different column-vector algorithm;
the selected target here is its later row-polynomial assertion.

\subsection{\texorpdfstring{\oeis{A379907}: a Riordan convolution and column relation}{A379907: a Riordan convolution and column relation}}
For ordinary Riordan arrays $A=(g,f)$ and $B=(u,v)$, meaning
$A(n,k)=[t^n]g(t)f(t)^k$ and similarly for $B$, the general selected
identity is
\[
 \sum_{i=0}^{n-k}A(n-k,i)B(k+i,k)
 =[t^n]g(t)u(f(t))
       \left(\frac{t\,v(f(t))}{f(t)}\right)^k.
\]
Since $B(k+i,k)=[s^i]u(s)(v(s)/s)^k$, substitution in $g(t)f(t)^i$
proves the identity immediately. The quotients are well-defined formal
series because $f$ and $v$ have order one.

For the particular array in this entry, put
$G=C(t/(1+t))/(1+t)$ and $F=tC(t/(1+t))$, so $T=(G,F)$.
The Catalan equation yields $F^2=(1+t)F-t(1+t)$. Multiplication by
$GF^{k-2}$ and coefficient extraction prove, for $k\ge2$,
\[
 T(n,k)=T(n,k-1)+T(n-1,k-1)
        -T(n-1,k-2)-T(n-2,k-2),
\]
where entries outside the triangle are zero. Both assertions belong to
one counted selection.

\subsection{\texorpdfstring{\oeis{A380191}: central-binomial weights}{A380191: central-binomial weights}}
Let $D=1+xD^3$ with $D(0)=1$. The established Riordan form of the source
array is $T=(2-D,xD)$. For each $k\ge0$, the generating function of the
selected weighted column sum is
\[
 \sum_{n\ge k}\sum_{i\ge0}\binom{2i}{i}T(n,k+i)x^n
 =\frac{(2-D)(xD)^k}{\sqrt{1-4xD}}.
\]
Since $xD=(D-1)/D^2$ and the square root has constant term one,
$\sqrt{1-4xD}=(2-D)/D$. The result is $D(x)(xD(x))^k$, exactly the
column generating function of \oeis{A110616}. Thus the identification
holds for every $n\ge k\ge0$.

\subsection{\texorpdfstring{\oeis{A004737}: a false matrix claim and a true Gram formula}{A004737: a false matrix claim and a true Gram formula}}
Let $q\ge1$, $\theta_a=(2a-1)\pi/(2q)$, and
\[
 E_{a,j}=\frac{\sin(j\theta_a)}{\sin\theta_a}
        =U_{j-1}(\cos\theta_a),\qquad 1\le a,j\le q.
\]
The source's proposed integrality and symmetry of $E$ fail already at
$q=2$:
\[
 E=\begin{pmatrix}1&\sqrt2\\1&-\sqrt2\end{pmatrix}.
\]
However, the normalized Gram matrix has the exact formula
\[
 \boxed{\frac1q(E^TE)_{ij}=
 \begin{cases}\min(i,j),&i\equiv j\pmod2,\\0,&i\not\equiv j\pmod2.
 \end{cases}}
\]
To prove it, the average over $a=1,\ldots,q$ of $\cos(\ell\theta_a)$ is
zero for $1\le\ell\le2q-2$: odd $\ell$ cancel in symmetric pairs, and
even $\ell$ give a geometric sum. Hence the average of $U_r(\cos\theta_a)$
is one for even $r$ and zero for odd $r$, for $0\le r\le2q-2$.
Apply the product identity
\[
 U_{i-1}(t)U_{j-1}(t)
 =\sum_{r=0}^{\min(i,j)-1}U_{i+j-2-2r}(t).
\]
All summands have the same parity, proving the formula. In particular
the Gram matrix is integral and symmetric, and its even antidiagonals
give the stated increasing-then-decreasing integer pattern. This entry
is one mixed proof/refutation selection.

\subsection{\texorpdfstring{\oeis{A153197}: formal convergence of alternating transforms}{A153197: formal convergence of alternating transforms}}
\textbf{Uncounted qualification.} The source also records an explicit
iteration program and established fixed-point formulas. The precise
normalization below is mathematically useful, but the relationship between
the original wording and its broader conjectural extension is not sharp
enough to count this as another settled open assertion.

We use the prepending convention described with the source conjecture.
For a formal power series $F$, define
\[
 (\mathcal BF)(x)=\frac{1}{1-x}F\left(\frac{x}{1-x}\right),
 \qquad
 (\mathcal HG)(x)=\frac{1}{1-x-x^2G(x)}.
\]
Here $\mathcal B$ is the binomial transform. To see the normalization of
$\mathcal H$, prepend $1$ to $G$ to get $A=1+xG$, apply the invert
transform $A/(1-xA)$, and prepend $1$ once more; the result is
$1+xA/(1-xA)=\mathcal HG$.

The $x$-adic order of $\mathcal BF-\mathcal B\widetilde F$ equals that
of $F-\widetilde F$. Direct subtraction shows that $\mathcal H$ raises
the order of a difference by two. Therefore $\mathcal H\mathcal B$ is
a strict contraction in the complete ring of formal power series. Its
iterates from any initial series converge coefficient by coefficient to
a unique fixed point $F_*$. The interleaved iterates converge to
$G_*=\mathcal BF_*$. This proves the selected convergence assertion,
and in fact removes the need for the initial Hankel-transform condition.

For identification, the Jacobi continued fraction
\[
 F_* = \cfrac{1}{1-x-\cfrac{x^2}{1-2x-
                  \cfrac{x^2}{1-3x-\cdots}}}
\]
has diagonal coefficients $1,2,3,\ldots$ and all off-diagonal squared
coefficients equal to one. The binomial transform adds one to every
diagonal coefficient; this follows by substituting $x/(1-x)$ in each
level of the fraction and canceling $1-x$. Applying $\mathcal H$ then
puts back the first level with coefficient one, so this fraction is the
fixed point just proved unique. It is \oeis{A006789}; its binomial
transform is the defining series for \oeis{A153197}. The source's less
precise extension to arbitrary Hankel transforms is outside this selection.

\subsection{\texorpdfstring{\oeis{A005773}: universally ordered coprime divisors}{A005773: universally ordered coprime divisors}}
We give a precise version of the source's empirical comparison assertion:
an inequality is deducible from prime order alone when it holds for every
increasing list of primes $p_1<\cdots<p_n$. For disjoint index sets $I,J$,
the coprime divisors $d_I=\prod_{i\in I}p_i$ and $d_J=\prod_{j\in J}p_j$
satisfy $d_I\le d_J$ for every such list if and only if
\[
 |J\cap\{r,\ldots,n\}|\ge |I\cap\{r,\ldots,n\}|
 \quad(1\le r\le n).
\]
Sufficiency follows by matching each index of $I$, starting with the
largest, to a distinct index of $J$ that is at least as large. Extra
factors in $J$ are positive integers greater than one.

For necessity, suppose a suffix has excess $\delta>0$ indices from $I$.
Fix the earlier primes. By repeated application of Bertrand's theorem,
the remaining $m$ primes can all be chosen between $X$ and $2^mX$ for
arbitrarily large $X$. The ratio $d_I/d_J$ is then at least a fixed
positive constant times $X^\delta/2^{m|J\cap\{r,\ldots,n\}|}$, and
eventually exceeds one, a contradiction.

Read indices in order $n,n-1,\ldots,1$, writing an up-step for $J$, a
down-step for $I$, and a horizontal step for neither. The displayed
condition is precisely that every initial height is nonnegative. These
are Motzkin meanders of length $n$, counted by \oeis{A005773}$(n+1)$
with the source's stated length-to-index shift. The empty pair gives the initial case.
Other assertions attached to this entry are not included.

\subsection{\texorpdfstring{\oeis{A376182}: the triangular enumeration is a permutation}{A376182: the triangular enumeration is a permutation}}
For $n\ge1$ set
\[
 C_n=\frac{2n^2+4n+1-(-1)^n}{4},\qquad
 T(n,k)=\begin{cases}C_n-k,&k\text{ odd},\\C_n-2n+k,&k\text{ even}.
 \end{cases}
\]
Every $C_n$ is even. In row $n$, the odd entries form the block of odd
integers from $C_n-2\lceil n/2\rceil+1$ to $C_n-1$. The even entries
form the block from $C_n-2n+2$ to $C_n-2\lceil n/2\rceil$, when this
block is nonempty. Since
\[
 C_n-C_{n-1}=2\lceil n/2\rceil,\qquad C_n-C_{n-2}=2n,
\]
each odd block starts two after the preceding odd block ends, and the
same is true of the nonempty even blocks. The first odd value is one
and the first even value is two; the block endpoints tend to infinity.
Thus every positive integer occurs exactly once, proving the conjecture.

\subsection{\texorpdfstring{\oeis{A378277}: alternating harmonic Fibonacci rows}{A378277: alternating harmonic Fibonacci rows}}
Take $F_0=0,F_1=1$ and $F_{-1}=1$. The denominator array has
$D(n,k)=F_kF_{k+2}$ for $k<n$, and $D(n,n)=F_nF_{n+1}$. The selected
identity is
\[
 H_n:=\sum_{k=1}^n\frac{(-1)^{k-1}}{D(n,k)}
 =\frac{F_{n-2}}{F_{n+1}},\qquad n\ge1.
\]
The initial case is $H_1=1$. For $n\ge2$, subtraction of consecutive
rows gives
\[
 H_n-H_{n-1}=\frac{2(-1)^{n-1}}{F_nF_{n+1}}.
\]
The proposed right side has the same difference, because
$F_nF_{n-2}-F_{n+1}F_{n-3}=2(-1)^{n-1}$. This last identity follows
from the Fibonacci recurrence and Cassini's identity, including $n=2$
with the stated negative-index convention. Induction proves the claim.

\subsection{\texorpdfstring{\oeis{A349802}: exact Euler transforms of Lyndon columns}{A349802: exact Euler transforms of Lyndon columns}}
Let $T(n,k)$ count binary Lyndon words of length $n$ beginning with
exactly $k$ zeros, and let
$E_k(x)=\prod_{n\ge1}(1-x^n)^{-T(n,k)}$ for $k\ge1$. The precise
form of the source's Fibonacci and higher-bonacci assertion is
\[
 \boxed{E_1(x)=\frac1{1-x-x^2},\qquad
 E_k(x)=\frac{1-x-\cdots-x^k}{1-x-\cdots-x^{k+1}}\quad(k\ge2).}
\]
In particular, for $k\ge2$ the series is
$1+x^{k+1}/(1-x-\cdots-x^{k+1})$; this states its initial terms as well
as its higher-bonacci tail, avoiding an unspecified shift.

Every nonconstant binary Lyndon word starts at a longest cyclic run of
zeros. Decompose its cyclic word into blocks $0^j1$. Primitivity is
preserved by this block encoding, and the blocks are a weighted alphabet
with weights $1,2,\ldots,k+1$. Beginning with exactly $k$ zeros is
equivalent to having maximum block weight $k+1$.

For any weighted alphabet of inventory $S(x)$, unique factorization of
words into a nonincreasing product of Lyndon words gives
\[
 \prod_{\text{primitive necklaces }w}(1-x^{|w|})^{-1}
 =\frac1{1-S(x)}.
\]
This identity can alternatively be obtained by taking logarithms and
partitioning words by their primitive period. Dividing the identity for
weights $1,\ldots,k+1$ by the one for weights $1,\ldots,k$ proves the
second formula. For $k=1$, the one-letter word $0$ contributes the
additional factor $(1-x)^{-1}$, giving the first formula. The identity
and its normalization are the counted result; the constant term is not
silently identified with a different convention for higher-bonacci numbers.

\subsection{\texorpdfstring{\oeis{A163774}: collapsing the binomial sum}{A163774: collapsing the binomial sum}}
Using the established even swing-number identity
$\operatorname{swing}(2i)=\binom{2i}{i}$, the definition is
\[
 a(n)=\sum_{k=0}^n\sum_{i=k}^n
 (-1)^{n-i}\binom{n-k}{n-i}\binom{2i}{i}.
\]
The finite sums can be interchanged. The hockey-stick identity gives
$\sum_{k=0}^i\binom{n-k}{n-i}=\binom{n+1}{i}$. Therefore
\[
 \boxed{a(n)=\sum_{i=0}^n(-1)^{n-i}\binom{n+1}{i}\binom{2i}{i},}
\]
which is the selected conjecture for every $n\ge0$.

\subsection{\texorpdfstring{\oeis{A152920}: weighted rows}{A152920: weighted rows}}
The array is $T(N,k)=(2N-k)2^{k-1}$, with its $k=0$ value understood
as $N$. Define
\[
 B(N,p)=\frac1{N+p+1}\sum_{k=0}^{N}\binom{p+k}{p}T(N,k).
\]
The source conjectures $B(n+1,p)=A(n,p)$, where
$A(n,p)=\oeis{A193844}(n+p,n)$. The known bivariate generating function
of the latter triangle, $1/((1-xy)(1-x-2xy))$, gives for fixed $p\ge0$
\[
 A_p(x)=\sum_{n\ge0}A(n,p)x^n
       =\frac1{(1-x)(1-2x)^{p+1}}.
\]
Indeed substitute $x=u,y=v/u$ and extract $[u^p]$.
Put $N=k+r$ in the unnormalized sum, and use the negative-binomial
series and its derivative. The result is
\begin{align*}
 \sum_{N\ge0}\sum_{k=0}^N\binom{p+k}{p}(2N-k)2^{k-1}x^N
 &=\frac{x((p+2)-(p+3)x)}{(1-x)^2(1-2x)^{p+2}}\\
 &=x\big((p+2)A_p(x)+xA_p'(x)\big).
\end{align*}
The coefficient of $x^N$ is $(N+p+1)A(N-1,p)$ for $N\ge1$.
Division by the nonzero factor proves the result.

\subsection{\texorpdfstring{\oeis{A105794}: factorial and Bernoulli sums}{A105794: factorial and Bernoulli sums}}
The defining connection formula is
$(x-1)^n=\sum_kT(n,k)x_{\underline{k}}$, equivalently the column
exponential generating function is $e^{-z}(e^z-1)^k/k!$.
The two selected 2025 conjectures follow directly from this representation.

First, set $x=-m-1$ for any integer $m\ge0$. Since
$(-m-1)_{\underline{k}}=(-1)^k(k+m)!/m!$, one obtains
\[
 \sum_{k=0}^n(-1)^{n-k}T(n,k)\frac{(k+m)!}{m!}=(m+2)^n.
\]
Second, with Bernoulli convention $B_1=-1/2$,
\[
 \sum_{k=1}^n(-1)^kT(n,k)\frac{(k-1)!}{k+1}=(-1)^n-B_n.
\]
For the proof, the exponential generating function of the left side is
\[
 e^{-z}\sum_{k\ge1}\frac{(-1)^k(e^z-1)^k}{k(k+1)}
 =e^{-z}-\frac{z}{e^z-1}.
\]
Here we used
$\sum_{k\ge1}(-1)^kw^k/(k(k+1))=1-(1+w)\log(1+w)/w$.
Coefficient extraction proves the result, including $n=0$.

\subsection{\texorpdfstring{\oeis{A132393}: a Vandermonde-squared multisum}{A132393: a Vandermonde-squared multisum}}
Let $c(a,b)$ denote unsigned Stirling numbers of the first kind. The
selected May 2026 conjecture is, for integers $n,m\ge0$,
\begin{align*}
 c(n+m+1,m+1)
 ={}&(n+m)!\prod_{i=1}^{m}(n+i)^i\\
 &\times\sum_{n\ge j_1\ge\cdots\ge j_m\ge0}
 (-1)^{nm+\sum_tj_t}
 \prod_{1\le q<p\le m}(j_q-j_p+p-q)^2\\
 &\hspace{22mm}\times
 \prod_{t=1}^{m}
 \frac{(n+t-1)_{\underline{j_t}}}
 {(n-j_t+t)^{m+1}(j_t+m-t)!}.
\end{align*}
All products and the nested sum for $m=0$ have value one, giving the
correct $n!$; below assume $m\ge1$ and put $N=n+m$.

Set $r_t=j_t+m-t$ and $s_t=N-r_t$. This is a bijection from the
summation region to $1\le s_1<\cdots<s_m\le N$. The squared product
becomes the squared Vandermonde $V(s)^2$, and the sign becomes
$(-1)^{m(m+1)/2+\sum s_t}$. The factorial identity
\[
 \prod_{i=1}^{m}(n+i)^i\prod_{t=1}^{m}(n+t-1)!=(N!)^m
\]
and $\binom{N-1}{s-1}=\frac{s}{N}\binom Ns$ reduce the proposed right
side exactly to
\[
 N!(-1)^{m(m+1)/2}
 \sum_{1\le s_1<\cdots<s_m\le N}V(s)^2
 \prod_{t=1}^{m}\frac{(-1)^{s_t}\binom N{s_t}}{s_t^m}.
\]
By Cauchy--Binet, the sum is $\det H$, where
\[
 H_{a,b}=\sum_{s=1}^{N}(-1)^s\binom Ns s^{a+b-m},
 \qquad 0\le a,b<m.
\]
Let $h_j$ be the complete homogeneous symmetric function of degree $j$
in $1,1/2,\ldots,1/N$, with $h_0=1$ and $h_j=0$ for $j<0$.
Partial fractions give
\[
 \prod_{r=1}^{N}(1-z/r)^{-1}
 =\sum_{s=1}^{N}\frac{(-1)^{s-1}\binom Ns}{1-z/s}.
\]
Thus $H_{a,b}=-h_{m-a-b}$ whenever $m-a-b\ge0$. When the exponent
$a+b-m$ is positive, it is at most $m-2<N$; the corresponding binomial
finite difference is zero, giving the same formula with $h_j=0$.
Reverse the $m$ columns. The resulting entries are $-h_{1-a+b}$, hence
\[
 \det H=(-1)^{m(m+1)/2}\det[h_{1-a+b}]_{a,b=0}^{m-1}
       =(-1)^{m(m+1)/2}e_m(1,1/2,\ldots,1/N).
\]
The last determinant identity is the single-column Jacobi--Trudi
identity; it also follows recursively by multiplying
$\sum h_jz^j$ by $\sum(-1)^je_jz^j$ to obtain one.
Finally $N!e_m(1,1/2,\ldots,1/N)=c(N+1,m+1)$, from
$t\prod_{r=1}^{N}(t+r)$. This proves the full multisum, including all
boundary cases stated above.

\subsection{\texorpdfstring{\oeis{A397048}: a Stirling convolution}{A397048: a Stirling convolution}}
The defining triangle is
\[
 T(n,k)=\frac1{k!}\sum_{i=0}^{k}(-1)^{k-i}\binom ki(n+i)^n.
\]
We prove the first conjecture only:
\[
 T(n,k)=\sum_{i=0}^{n-k}s(n,k+i)S(n+k+i,n+k),
\]
where $s$ is the signed first-kind Stirling triangle \oeis{A048994} and
$S$ is the second-kind triangle \oeis{A048993}. The elementary and
complete symmetric representations of these numbers turn the right
side into
\[
 [z^{n-k}]
 \frac{\prod_{j=1}^{n-1}(1-jz)}{\prod_{j=1}^{n+k}(1-jz)}
 =h_{n-k}(n,n+1,\ldots,n+k).
\]
For distinct $x_0,\ldots,x_k$, partial fractions of
$\prod_i(1-x_i z)^{-1}$ give
\[
 h_{n-k}(x_0,\ldots,x_k)
 =\sum_{i=0}^k\frac{x_i^n}{\prod_{j\ne i}(x_i-x_j)}.
\]
At $x_i=n+i$, the denominator is $(-1)^{k-i}i!(k-i)!$, exactly the
defining finite difference. The case $n=k=0$ is checked separately as
$1=1$. The source's other conjectures are not needed for this selection.

\subsection{\texorpdfstring{\oeis{A112302}: a Lerch-transcendent limit}{A112302: a Lerch-transcendent limit}}
Define $L_0=1$ and
\[
 nL_n=\sum_{k=0}^{n-1}\Phi(1/2,n-k,1)L_k,\qquad n\ge1,
\]
where $\Phi(1/2,r,1)=\sum_{j\ge0}2^{-j}(j+1)^{-r}$. The selected
2016 conjecture states that $L_n$ tends to Somos's quadratic recurrence
constant $\prod_{j\ge1}j^{2^{-j}}$.

Logarithmic differentiation of $F(x)=\sum_{n\ge0}L_nx^n$ gives, for
$|x|<1$,
\[
 F(x)=\prod_{j\ge0}\left(1-\frac{x}{j+1}\right)^{-2^{-j}}
     =\frac{H(x)}{1-x}.
\]
The product $H$, obtained by restricting to $j\ge1$, converges
analytically on $|x|<2$ and has nonnegative coefficients $h_r$.
Thus $L_n=\sum_{r=0}^nh_r$ and $L_n\uparrow H(1)$. Finally
\begin{align*}
 \log H(1)
 &=\sum_{j\ge1}2^{-j}\log\frac{j+1}{j}\\
 &=\sum_{j\ge1}2^{-j}\log j,
\end{align*}
by shifting the absolutely convergent series. This is the logarithm of
the proposed constant. The proof also supplies the explicit error bound
$0\le H(1)-L_n\le\rho^{-n-1}H(\rho)$ for every $1<\rho<2$.
The product representation of the constant was already known; the
selected result is convergence of the particular Lerch recurrence.

\clearpage
% Source fragment: r2_arith_body.tex
\section{Arithmetic and congruences}
\subsection{A334595: right boundaries of XOR triangles}
\noindent Source: \url{https://oeis.org/A334595}. Peter Kagey's entry lists the two assertions below as conjectures.

\begin{theorem}
For the binary value $a(n)$ obtained by reading the right boundary of the XOR triangle of $n$, $a(n)=1$ exactly when $n$ is a power of two. The strict record positions are $1$ and $2^k+1$ for $k\geq1$.
\end{theorem}
\begin{proof}
Let the input word have length $L$ and leading digit $1$. Reverse it, obtaining digits $x_0,\ldots,x_{L-1}$. The successive right-boundary digits, read from top to bottom, are
\[
y_j=\sum_{i=0}^j\binom ji x_i\pmod2,\qquad 0\leq j<L.
\]
The lower Pascal matrix $P=(\binom ji)$ is an involution over $\mathbb F_2$, since its square has $(j,i)$ entry $\binom ji2^{j-i}$, equal to $1$ for $j=i$ and $0$ otherwise. Thus the boundary map is one-to-one on words of length $L$, including words with leading zeros.

The input $10\cdots0$ has boundary $00\cdots01$, hence value $1$. Invertibility proves that no other length-$L$ input gives this value. The all-ones boundary has inverse $x=(1,0,\ldots,0)$: this follows from $\sum_{i=0}^j\binom ji=2^j$. For $L\geq2$ the corresponding unreversed input has leading digit zero and is inadmissible. Therefore every admissible boundary has value at most $2^L-2$.

The input $10\cdots01$ is the sum over $\mathbb F_2$ of $10\cdots0$ and $00\cdots01$. Their boundaries are respectively $00\cdots01$ and $11\cdots11$, so its boundary is $11\cdots10$, of value $2^L-2$. The first input of length $L$ is $2^{L-1}$ and has value $1$; the second, $2^{L-1}+1$, attains this length's maximum. All inputs of shorter length have value at most $2^{L-1}-2$, except the single length-one value $1$, which is still smaller when $L=2$. Thus each second input is a strict record and no subsequent input of the same length is. The initial input $1$ is the remaining record.
\end{proof}

\subsection{A133910: divisibility by the index}
\noindent Source: \url{https://oeis.org/A133910}. The target is Ridouane Oudra's January 2026 conjecture. We use Hari Rajesh's October 2025 product formula explicitly recorded there.

For $n>1$ write $n=\prod_p p^{e_p}$ and put
\[
P(n)=\max_{p\mid n}p^{e_p}.
\]
The set A051283 consists of $n>1$ such that $n>P(n)^2$.
\begin{theorem}
For A133910, $n\mid a(n)$ if and only if $n=1$ or $n\in\mathrm{A051283}$.
\end{theorem}
\begin{proof}
The recorded formula is
\[
a(n)=\prod_{p\mid n}p^{\lfloor\log_p n\rfloor-e_p}.
\]
Consequently $n\mid a(n)$ is equivalent to $\lfloor\log_p n\rfloor\geq2e_p$ for every $p\mid n$, or to $n\geq P(n)^2$. Equality is impossible for $n>1$: if $P(n)=p^{e_p}$ and $n=P(n)^2=p^{2e_p}$, the exact exponent of $p$ in $n$ would be both $e_p$ and $2e_p$. Therefore the condition is $n>P(n)^2$. At $n=1$, $a(1)=1$.
\end{proof}

\subsection{A051283: an exponent in an initial least common multiple}
\noindent Sources: \url{https://oeis.org/A051283} and \url{https://oeis.org/A346418}. Only the characterization using A346418 is selected here. The other binomial-gcd conjecture has an acknowledged 2026 proof and is excluded.

\begin{theorem}
The members of A051283 are exactly the integers $n$ for which $\mathrm{A346418}(n)>1$.
\end{theorem}
\begin{proof}
For $n>1$, A346418 is the greatest exponent $r$ such that $n^r$ divides $L_n=\operatorname{lcm}(1,\ldots,n)$. Hence its value is greater than one exactly when $n^2\mid L_n$. Since $v_p(L_n)=\lfloor\log_p n\rfloor$, this is equivalent to $2e_p\leq\lfloor\log_p n\rfloor$ for every exact prime-power factor $p^{e_p}$ of $n$. As in the preceding proof, the condition is $n>P(n)^2$. The defined initial value of A346418 is $1$, so $n=1$ belongs to neither set.
\end{proof}

\subsection{A354995: when the unitary-divisor difference vanishes}
\noindent Sources: \url{https://oeis.org/A354995} and \url{https://oeis.org/A354933}. The latter entry identifies its value $U(n)$ as the least unitary divisor of $n$ at least $\sqrt n$. Put $P(1)=U(1)=1$.

\begin{theorem}
The difference $U(n)-P(n)$ is positive exactly at the members of A051283, and is zero elsewhere. Thus its nonzero positions are exactly the positions conjectured in A354995.
\end{theorem}
\begin{proof}
Every unitary divisor is a product of whole exact prime-power blocks of $n$. If $n>P(n)^2$, then every divisor at least $\sqrt n$ is greater than $P(n)$, giving $U(n)>P(n)$. If $n<P(n)^2$, the block $P(n)$ itself is a unitary divisor larger than $\sqrt n$, so $U(n)\leq P(n)$. A unitary divisor $d<P(n)$ cannot contain that block, and therefore divides $n/P(n)<\sqrt n$. Thus none of these smaller divisors is eligible, and $U(n)=P(n)$. Equality $n=P(n)^2$ is impossible for $n>1$ by the exact-exponent argument above. The initial difference is $0$.
\end{proof}

\subsection{A376336: existence of a zero of tau modulo every modulus}
\noindent Sources: \url{https://oeis.org/A376336} and \url{https://oeis.org/A000594}.

Only existence is proved; the separate proposed bound $a(n)\leq n^2$ remains outside this result. The standard Hecke recurrence used below follows from Mordell's theorem on Ramanujan's tau function, as referenced in A000594.

\begin{theorem}
For every positive integer $N$ there is a positive integer $m$ for which $N\mid\tau(m)$.
\end{theorem}
\begin{proof}
The case $N=1$ is immediate. Choose a prime $p$ not dividing $N$. Set $u_r=\tau(p^r)$, so $u_0=1$, $u_1=\tau(p)$, and
\[
u_r=\tau(p)u_{r-1}-p^{11}u_{r-2}\quad(r\geq2).
\]
The matrix $M=\left(\begin{smallmatrix}\tau(p)&-p^{11}\\1&0\end{smallmatrix}\right)$ has determinant $p^{11}$, a unit modulo $N$. It therefore belongs to the finite group $\mathrm{GL}_2(\mathbb Z/N\mathbb Z)$ and has some finite order $h$. Since its lower-left entry is $1$, $M\ne I$ modulo $N$, so $h\geq2$.

For every $r\geq1$, recurrence and induction give
\[
M^r\binom10=\binom{u_r}{u_{r-1}}\pmod N.
\]
Taking $r=h$ gives $u_{h-1}\equiv0\pmod N$. The integer $m=p^{h-1}$ has the required property, proving the existence of a least such positive $m$ by well-ordering.
\end{proof}

\subsection{A087726: square-zero matrices over prime-power rings}
\noindent Source: \url{https://oeis.org/A087726}. The prime-modulus case was already proved by Manfred Scheucher in 2015; the proof below covers every exponent and the full squarefree characterization.

\begin{theorem}
Let $A(N)$ count matrices $X\in M_2(\mathbb Z/N\mathbb Z)$ satisfying $X^2=0$. For every prime $p$ and $k\geq1$,
\[
A(p^k)=p^{2k}\left(1+\left\lfloor\frac{k}{2}\right\rfloor(1-p^{-2})\right).
\]
Equivalently, this is Robert Israel's displayed formula
\[
\frac{p^{2k-2}}4\bigl(1+3p^2+2k(p^2-1)+(-1)^k(p^2-1)\bigr).
\]
Moreover $A(N)=N^2$ exactly when $N$ is squarefree, including $N=1$.
\end{theorem}
\begin{proof}
Fix $R=\mathbb Z/p^k\mathbb Z$. Call a matrix primitive if its reduction modulo $p$ is nonzero. A primitive square-zero matrix reduces to a nonzero nilpotent map on the two-dimensional vector space over $\mathbb F_p$. There is a vector $\bar v$ for which $(X\bar v,\bar v)$ is a basis. Indeed, choose $\bar v$ outside the one-dimensional kernel; the nonzero vector $X\bar v$ lies in that kernel and is independent of $\bar v$.

Lift $\bar v$ to $v\in R^2$. The columns $(Xv,v)$ have unit determinant, hence form an $R$-basis. In this basis $X$ has matrix $J=\left(\begin{smallmatrix}0&1\\0&0\end{smallmatrix}\right)$. Therefore all primitive square-zero matrices form exactly the conjugacy orbit of $J$.

The invertible matrices commuting with $J$ are precisely $\left(\begin{smallmatrix}a&b\\0&a\end{smallmatrix}\right)$ with $a$ a unit and $b$ arbitrary. Their number is $p^{2k-1}(p-1)$. Since
\[
|\mathrm{GL}_2(R)|=p^{4(k-1)}(p^2-1)(p^2-p),
\]
the number of primitive solutions is $p^{2k-2}(p^2-1)$.

For $k\geq2$, a nonprimitive matrix is uniquely $X=pY$ with $Y$ taken modulo $p^{k-1}$. Its square vanishes modulo $p^k$ exactly when $Y^2$ vanishes modulo $p^{k-2}$. Each matrix modulo $p^{k-2}$ has $p^4$ lifts modulo $p^{k-1}$. With $A(p^0)=1$, we obtain
\[
A(p^k)=p^{2k-2}(p^2-1)+p^4A(p^{k-2}).
\]
For $k=1$ there are $p^2-1$ primitive solutions and the zero matrix, so $A(p)=p^2$. Division by $p^{2k}$ and induction on each parity of $k$ gives the stated formula.

The Chinese remainder theorem identifies the matrix ring modulo a product of coprime moduli with the product of their matrix rings, preserving $X^2=0$. Hence $A$ is multiplicative. Each ratio $A(p^k)/p^{2k}$ is $1$ for $k=1$ and strictly greater than $1$ for $k\geq2$. Their product equals $1$ exactly for squarefree $N$.
\end{proof}

\subsection{A322327: the general parameter statement}
\noindent Source: \url{https://oeis.org/A322327}. The general-$k$ assertion is explicitly conjectured; the $k=2$ instance is already recorded as a formula.

\begin{theorem}
Let $k\in\mathbb Z$ be fixed and let $E(n)$ be the product of the positive exponents in the prime factorization of $n$, with $E(1)=1$. Then $f_k(n)=E(n)k^{\omega(n)}$, using $0^0=1$, is multiplicative and $f_k(p^e)=ke$ for $e>0$.
\end{theorem}
\begin{proof}
For $n=\prod_{j=1}^r p_j^{e_j}$ one has
\[
f_k(n)=\left(\prod_{j=1}^r e_j\right)k^r
=\prod_{j=1}^r(ke_j).
\]
The formula holds for the empty product at $n=1$, including $k=0$. For coprime integers the lists of prime powers are disjoint, so the displayed product for their product is the product of their individual values. A single prime-power factor gives $ke$.
\end{proof}

\subsection{A033955: a strict bound for the sum of prime remainders}
\noindent Source: \url{https://oeis.org/A033955}. The target is Alonso del Arte's July 2026 conjecture. We use Jitsuro Nagura, ``On the interval containing at least one prime number,'' \emph{Proceedings of the Japan Academy} 28 (1952), 177--181, \url{https://doi.org/10.3792/pja/1195570997}. That theorem provides a prime in $(x,6x/5)$ whenever $x\geq25$.

\begin{theorem}
For the $n$th prime $q=p_n$, set $S(q)=\sum_{p<q,\ p\text{ prime}}(q\bmod p)$. If $n>6$, then $S(q)>q$. The equality $S(p_n)=p_n$ occurs only for $n=6$.
\end{theorem}
\begin{proof}
Assume $q\geq79$. Apply Nagura's theorem at $x=q/3,q/2,3q/5,18q/25$. This produces four distinct primes in the disjoint intervals
\[
\left(\frac q3,\frac{2q}5\right),\quad
\left(\frac q2,\frac{3q}5\right),\quad
\left(\frac{3q}5,\frac{18q}{25}\right),\quad
\left(\frac{18q}{25},\frac{108q}{125}\right).
\]
The first corresponding quotient is $2$ and the other three are $1$. Their remainders are respectively greater than $q/5,2q/5,7q/25,17q/125$. All other remainders are nonnegative. Consequently
\[
S(q)>\left(\frac15+\frac25+\frac7{25}+\frac{17}{125}\right)q
=\frac{127}{125}q>q.
\]
The primes $17\leq q<79$ and their exact sums are
\[
\begin{array}{c|rrrrrrrr}
q&17&19&23&29&31&37&41&43\\ \hline
S(q)&18&27&29&46&56&70&74&88
\end{array}
\]
\[
\begin{array}{c|rrrrrrr}
q&47&53&59&61&67&71&73\\ \hline
S(q)&98&134&147&171&200&217&252.
\end{array}
\]
Each exceeds $q$. For the first six primes the sums are $0,1,3,4,8,13$, establishing the stated unique equality. These finite sums can be checked by the supplied exact verifier; the table exhausts the range not covered by the theorem.
\end{proof}

\subsection{A166306: the empirical residue restriction}
\noindent Source: \url{https://oeis.org/A166306}. Only the assertion about residues modulo $9$ is selected.

\begin{theorem}
No term of A166306 is divisible by $3$. Equivalently its residue modulo $9$ always lies in $\{1,2,4,5,7,8\}$.
\end{theorem}
\begin{proof}
The first term corresponds to $B_1$ and is $1$. For an even Bernoulli index $2m$, $m\geq1$, let
\[
\mathcal P=\{p\text{ prime}:p-1\mid2m\},\qquad D=\prod_{p\in\mathcal P}p.
\]
By the definition using the von Staudt--Clausen primes, the term is $\sum_{p\in\mathcal P}D/p$. The prime $3$ belongs to $\mathcal P$, and $D$ is squarefree. Modulo $3$ every summand except $D/3$ vanishes; $D/3$ is not divisible by $3$. Thus the sum is not divisible by $3$.
\end{proof}

\subsection{A234307: an alternating gcd identity}
\noindent Source: \url{https://oeis.org/A234307}. Peter Bala's January 2024 assertion is proved below.

\begin{theorem}
For every $n\geq1$,
\[
\sum_{i=1}^n\gcd(2n-i,i)
=\frac14\sum_{k=1}^{4n}(-1)^k\gcd(k,8n).
\]
\end{theorem}
\begin{proof}
Write $G(m)=\sum_{k=1}^m\gcd(k,m)$. Pairing $k$ with $2n-k$ gives
\[
\sum_{i=1}^n\gcd(2n-i,i)=\frac{G(2n)-n}{2}.
\]
The full alternating sum at modulus $8n$ is
\[
H=\sum_{k=1}^{8n}(-1)^k\gcd(k,8n)=4G(4n)-G(8n),
\]
because the even-index contribution is $2G(4n)$. If $T$ is the half sum from $1$ through $4n$, reflection about $4n$ shows $H=2T+4n$: the two endpoint terms are $4n$ and $8n$.

The formula $G(m)=\sum_{d\mid m}d\varphi(m/d)$ proves that $G$ is multiplicative. Directly,
\[
G(2^r)=(r+2)2^{r-1}\qquad(r\geq1).
\]
Writing $2n=2^r u$ with $u$ odd and $r\geq1$, this yields
\[
4G(4n)-G(8n)=4G(2n).
\]
Therefore $T=2G(2n)-2n$, and $T/4=(G(2n)-n)/2$, as required.
\end{proof}

\subsection{A290815: the odd-modulus reciprocal-square characterization}
\noindent Source: \url{https://oeis.org/A290815}. This proves the empirical odd-term characterization of August 2019. It does not resolve the separate conjecture for arbitrary exponents.

\begin{theorem}
For every odd integer $m\geq1$, let
\[
S_j(m)=\sum_{\substack{1\leq k\leq m\\(k,m)=1}}\frac1{k^j}.
\]
Then $m^3\mid\operatorname{num}(S_1(m))$ if and only if $m^2\mid\operatorname{num}(S_2(m))$, where the numerator is in lowest terms.
\end{theorem}
\begin{proof}
For $m=1$ both divisibilities are automatic. Assume $m>1$ is odd. All denominators are coprime to $m$, so congruences can be computed in the localization of $\mathbb Z$ obtained by inverting integers coprime to $m$. Divisibility of a reduced numerator by $m^r$ is exactly congruence to zero modulo $m^r$ in this ring.

The substitution $k\mapsto m-k$ permutes the reduced residues. Modulo $m$, it gives $S_3\equiv-S_3$, so $S_3\equiv0\pmod m$ because $2$ is a unit. Expansion of $(m-k)^{-1}$ modulo $m^3$ gives
\[
\frac1{m-k}\equiv-\frac1k-\frac m{k^2}-\frac{m^2}{k^3}\pmod{m^3}.
\]
Summing and again using the permutation gives
\[
2S_1\equiv-mS_2-m^2S_3\equiv-mS_2\pmod{m^3}.
\]
Since $2$ is invertible, $S_1\equiv0\pmod{m^3}$ is equivalent to $S_2\equiv0\pmod{m^2}$, proving both directions.
\end{proof}

\clearpage
% Source fragment: r2_floor_proofs.tex
\section{Second round: connected arrays and exact floor identities}
This section resolves 28 distinct entry-level assertions, all disjoint from
the first-round ledger. The primary sources are the current OEIS entries
at \texttt{https://oeis.org/}\textit{A-number}, inspected on 20 October 2026.
The selected formulas remain explicitly conjectural or empirical in those
entries. The source status is not a claim of global publication priority.
We exclude A166761, which records Howroyd's 2024 confirmation, and nearby
entries whose formulas are already stated as established or which contain
no conjectural assertion.

\subsection{Five connected-coloring formulas from a two-state transfer matrix}
For $q\geq1$, let $a_q(n)$ count $n\times2$ arrays over $q$ labeled colors,
using every color, in which the cells of each color form one connected
region under horizontal and vertical adjacency. Define
\[
B(x,t)=\frac{xt\bigl(1+t+(t-1)x\bigr)}
{(1-x)^2-xt(x+3)-xt^2(x+1)}.
\]
Then the complete ordinary generating function is
\begin{equation}\label{round2:eq:r2widthtwo}
 \sum_{n\geq1}a_q(n)x^n=q![t^q]B(x,t).
\end{equation}
Taking $q=3,4,5,6,7$ proves the empirical generating functions
in the five entries A166756--A166760.

\begin{proof}
First forget the names of the colors, and give each monochromatic region
a new name at its first appearance. Mark its introduction by $t$.
Once a color has disappeared from the last processed row, connectivity
prevents it from reappearing. The last row has one of two types:
$M$ (both cells the same color) or $D$ (two different colors).
In width two, each active color has a single connected component.

For an $M$ row, the next $M$ row either retains its color or starts a new
one, giving $1+t$. A next $D$ row can retain the old color in either
position and introduce one new color, or introduce two new colors,
giving $2t+t^2$. For a $D$ row, a next $M$ row can retain either of its
two colors or introduce a new one, giving $2+t$. A next $D$ row can
retain both colors in their original positions, retain exactly one
in its original position and introduce a new one, or introduce two
new colors, giving $1+2t+t^2$. A color moved only diagonally would be
disconnected and is excluded. Thus the exact transfer matrix is
\[
 T(t)=\begin{pmatrix}1+t&2t+t^2\\2+t&1+2t+t^2\end{pmatrix}.
\]
The first row contributes $(t,t^2)$, and both frontier types are accepting.
Therefore the generating function for the unnamed region partitions is
$x(t,t^2)(I-xT(t))^{-1}(1,1)^{\mathsf T}=B(x,t)$.
A partition into exactly $q$ nonempty regions has exactly $q!$ labelings
by the required colors, proving \eqref{round2:eq:r2widthtwo}.
\end{proof}

For a directly reproducible coefficient extraction, write
$b_q(x)=[t^q]B(x,t)$. Then
\[
b_1=\frac{x}{1-x},\qquad b_2=\frac{x(1+3x)}{(1-x)^3},\qquad
b_q=\frac{x(x+3)b_{q-1}+x(x+1)b_{q-2}}{(1-x)^2}\quad(q\geq3).
\]
The following table gives the resulting source functions in the form
\[
G_q(x)=\frac{q!x^h\sum_{j=0}^m v_jx^j}{(1-x)^{2q-1}}.
\]
All coefficient lists are ascending, starting with $v_0$.
\begin{center}\small
\begin{tabular}{lrrp{0.58\textwidth}}
\toprule Entry&$q$&$h$&$(v_0,\ldots,v_m)$\\\midrule
A166756&3&2&\(4,\allowbreak 9,\allowbreak 2,\allowbreak 1\)\\
A166757&4&2&\(1,\allowbreak 14,\allowbreak 27,\allowbreak 13,\allowbreak 8,\allowbreak 1\)\\
A166758&5&3&\(7,\allowbreak 48,\allowbreak 84,\allowbreak 60,\allowbreak 43,\allowbreak 12,\allowbreak 2\)\\
A166759&6&3&\(1,\allowbreak 34,\allowbreak 163,\allowbreak 273,\allowbreak 246,\allowbreak 196,\allowbreak 83,\allowbreak 25,\allowbreak 3\)\\
A166760&7&4&\(10,\allowbreak 144,\allowbreak 552,\allowbreak 917,\allowbreak 958,\allowbreak 827,\allowbreak 452,\allowbreak 187,\allowbreak 44,\allowbreak 5\)\\
\bottomrule\end{tabular}\end{center}

\subsection{Sixteen formulas from deleting repeated rows}
Consider arrays whose rows are lexicographically nondecreasing. Some of
the entries additionally require the columns to be lexicographically
nondecreasing. In each case all prescribed colors must occur and each
color must be connected.

Let $c_j$ be the number of valid arrays with exactly $j$ distinct rows,
with those rows in strictly increasing order, and no repeated row.
The following identity holds for every positive integer $n$:
\begin{equation}\label{round2:eq:r2rowcompression}
 a(n)=\sum_{j\geq1}c_j\binom{n-1}{j-1},\qquad
 \sum_{n\geq1}a(n)x^n=\sum_{j\geq1}c_j\left(\frac{x}{1-x}\right)^j.
\end{equation}
The sums are finite because there are only $q^w$ possible rows for width
$w$ and $q$ colors.

\begin{proof}
Equal rows occur consecutively. Contracting vertical edges between two
identical consecutive rows deletes one copy and preserves the connected
components of every color. It also preserves which colors occur, the
row ordering, and every lexicographic comparison between columns.
Conversely, duplicating a row preserves the same properties. Thus a
valid array with $j$ distinct rows is uniquely a valid strict-row array
and a choice of $j$ positive multiplicities summing to $n$.
There are $\binom{n-1}{j-1}$ such choices. Summing their generating
functions proves both identities.
\end{proof}

\paragraph{Exact finite computation of $c_j$.}
For completeness, the finite enumeration used in the coefficient table
is specified here. A state consists of the last row $r$, a partition
$\pi$ of its positions into connected components in the processed grid,
the set $U$ of colors already used, and (when needed) one bit for each
adjacent pair of columns indicating whether the lexicographic comparison
has already become strict. A partition block contains positions of only
one color. Colors in $U$ absent from $r$ have been completed.

To append a proposed row $r'$, perform these exact steps:
\begin{enumerate}
\item Require $r'>r$ for the strict-row enumeration. For each unsettled
column comparison, reject if the new left entry exceeds the right
entry, and mark it settled if the left entry is smaller.
\item Reject a color that was completed earlier. Form the disjoint union
of the old frontier components and the new cells. Join vertically
adjacent equal colors and horizontally adjacent equal colors in $r'$.
\item If an old component has no representative in the new row, its
color must have no other component and must be absent from the new
row. Otherwise reject. Such a color is now completed.
\item Restrict the resulting component partition to $r'$ and update $U$.
At the end accept exactly when every color has occurred and every color
still in the frontier has a single component.
\end{enumerate}
These rules are equivalent to extendible monochromatic connectivity:
a component that has left the frontier can never join a future cell,
while all components still meeting the frontier remain available to join.
Consequently the state paths are in bijection with the arrays being counted.

Every transition strictly increases the row, so the state graph is
acyclic. Starting with polynomial $1$ at the empty state, add $z$ times
each state's polynomial to all its successors, in increasing row order.
The sum over accepting states is exactly $\sum c_jz^j$. This algorithm
terminates after at most $q^w$ rows and establishes that every unlisted
coefficient is zero. It does not infer a formula from an initial segment
of the infinite sequence.

The table gives the complete nonzero coefficient vectors, with $h$ the
first listed index. All $c_j$ for $j<h$ or beyond the list are zero.
``Columns'' means that column ordering is required in addition to row
ordering. Substitution into \eqref{round2:eq:r2rowcompression} gives exactly
the empirical generating function in each of these sixteen entries.
\small
\begin{longtable}{lrrcrp{0.55\textwidth}}
\toprule Entry&$w$&$q$&Columns&$h$&$(c_h,c_{h+1},\ldots)$\\\midrule
\endhead
A166776&2&3&no&2&\(12,\allowbreak 27,\allowbreak 18,\allowbreak 4\)\\
A166777&2&4&no&2&\(12,\allowbreak 92,\allowbreak 175,\allowbreak 138,\allowbreak 52,\allowbreak 8\)\\
A166778&2&5&no&3&\(160,\allowbreak 765,\allowbreak 1306,\allowbreak 1096,\allowbreak 516,\allowbreak 136,\allowbreak 16\)\\
A166779&2&6&no&3&\(120,\allowbreak 1920,\allowbreak 6866,\allowbreak 10823,\allowbreak 9332,\allowbreak 4928,\allowbreak 1656,\allowbreak 336,\allowbreak 32\)\\
A166780&2&7&no&4&\(2730,\allowbreak 22666,\allowbreak 66234,\allowbreak 97749,\allowbreak 85222,\allowbreak 48204,\allowbreak 18544,\allowbreak 4848,\allowbreak 800,\allowbreak 64\)\\
A166781&3&2&no&1&\(4,\allowbreak 13,\allowbreak 14,\allowbreak 6\)\\
A166782&3&3&no&1&\(6,\allowbreak 84,\allowbreak 280,\allowbreak 443,\allowbreak 392,\allowbreak 192,\allowbreak 44\)\\
A166783&3&4&no&2&\(252,\allowbreak 2120,\allowbreak 7143,\allowbreak 13448,\allowbreak 16006,\allowbreak 12584,\allowbreak 6492,\allowbreak 2056,\allowbreak 320\)\\
A166786&4&2&no&1&\(6,\allowbreak 25,\allowbreak 56,\allowbreak 68,\allowbreak 36,\allowbreak 8\)\\
A166796&2&3&yes&2&\(6,\allowbreak 14,\allowbreak 9,\allowbreak 2\)\\
A166797&2&4&yes&2&\(8,\allowbreak 51,\allowbreak 92,\allowbreak 70,\allowbreak 26,\allowbreak 4\)\\
A166798&2&5&yes&3&\(105,\allowbreak 452,\allowbreak 722,\allowbreak 576,\allowbreak 262,\allowbreak 68,\allowbreak 8\)\\
A166799&2&6&yes&3&\(90,\allowbreak 1282,\allowbreak 4234,\allowbreak 6262,\allowbreak 5109,\allowbreak 2578,\allowbreak 840,\allowbreak 168,\allowbreak 16\)\\
A166800&2&7&yes&4&\(2016,\allowbreak 15421,\allowbreak 42152,\allowbreak 58628,\allowbreak 48404,\allowbreak 26088,\allowbreak 9652,\allowbreak 2456,\allowbreak 400,\allowbreak 32\)\\
A166801&3&3&yes&1&\(1,\allowbreak 25,\allowbreak 83,\allowbreak 118,\allowbreak 92,\allowbreak 40,\allowbreak 8\)\\
A166802&3&4&yes&2&\(67,\allowbreak 586,\allowbreak 1950,\allowbreak 3528,\allowbreak 3973,\allowbreak 2924,\allowbreak 1401,\allowbreak 410,\allowbreak 58\)\\
\bottomrule\end{longtable}
\normalsize
The finite graphs have between 11 and 1437 reachable states. The full
vectors and graph sizes are in
\texttt{r2\_\allowbreak floor\_\allowbreak hardin\_\allowbreak certificates.json};
the accompanying verifier reconstructs the graphs from the preceding
rules and checks every rational-function identity. It also compares the
frontier implementation to independent whole-grid connectivity searches
on small arrays, as an implementation audit.

\subsection{A166762: an exact 67-state recurrence certificate}
Here the width is three, there are three colors, and no row or column
ordering is imposed. Use the same frontier states and connectivity
rules, but permit every next row. There are exactly 67 reachable states
and 627 transitions, including the empty starting state indexed by 0.
Let $T$ be the integer adjacency matrix with rows as source states, and
let $b$ be the column vector whose entry is 1 precisely at an accepting
state. Each path specifies its rows uniquely, so
\[
 a(n)=e_0^{\mathsf T}T^n b.
\]
The exact finite matrix certificate is
\begin{equation}\label{round2:eq:r2matrixcertificate}
 (T-I)^7(T-2I)^2b=24e_0,\qquad Te_0=0.
\end{equation}
All 67 states, their successor lists, the accepting vector, and the
residual vector are supplied in
\texttt{r2\_\allowbreak floor\_\allowbreak hardin\_\allowbreak automaton.json}.
The verifier regenerates every state and transition, and verifies
\eqref{round2:eq:r2matrixcertificate} using integer vector operations.
The second identity also follows directly because no transition returns
to the empty state.

Expanding $(z-1)^7(z-2)^2$ and applying
$e_0^{\mathsf T}T^{n-9}$ to the first identity proves, for every $n\geq10$,
\begin{align*}
a(n)={}&11a(n-1)-53a(n-2)+147a(n-3)-259a(n-4)\\
 &+301a(n-5)-231a(n-6)+113a(n-7)-32a(n-8)+4a(n-9).
\end{align*}
This is the entry's empirical recurrence. The same construction gives
the initial values
\[
(a(1),\ldots,a(9))=(6,174,1548,8100,32208,108456,327042,912330,2403570).
\]

\subsection{Two arithmetic-progression floor generating functions}
\paragraph{A024172.}
For $S=\{1,2,\ldots,n\}$ and $n\geq2$, Newton's identities yield
\[
\frac{e_3(S)}{e_2(S)}=
\frac{n(n-2)(n+1)}{2(3n+2)}=Q(n)+\frac8{27(3n+2)},\qquad
Q(n)=\frac{n^2}6-\frac{5n}{18}-\frac4{27}.
\]
The denominator is positive. Since $Q(n+9)-Q(n)=3n+11$, its fractional
part has period nine. For $n\equiv1,\ldots,9\pmod9$ its numerators over
27 are $20,26,14,11,17,5,2,8,23$. For $n\geq3$ the positive remainder
is less than $1/27$. The case $n=2$ gives a correction of one. Thus
\[
a(n)=\lfloor Q(n)\rfloor+\mathbf1_{n=2}\qquad(n\geq2).
\]
Summing the polynomial part and the nine fractional residues proves
Barker's 2014 empirical generating function:
\[
\sum_{n\geq2}a(n)x^n=
\frac{x^4(x^4-x^3+x^2-x+1)(x^5-x^3-x^2-x-1)}
{(x-1)^3(x^2+x+1)(x^6+x^3+1)}.
\]

\paragraph{A024180.}
For $S=\{2,3,\ldots,n+3\}$, Newton's identities give
\[
\frac{e_3(S)}{e_2(S)}=
\frac{n(n+5)(n^2+9n+22)}{2(3n^2+29n+72)}=Q(n)+R(n),
\]
\[
Q(n)=\frac{n^2}6+\frac{13n}{18}+\frac5{27},\qquad
R(n)=-\frac{8(8n+45)}{27(3n^2+29n+72)}.
\]
Again all denominators are positive on $n\geq1$.
The identity $Q(n+9)-Q(n)=3n+20$ proves period nine for the fractional
part, whose numerators over 27 are $2,8,23,20,26,14,11,17,5$.
For $n\geq7$, $-2/27<R(n)<0$: the nontrivial inequality is equivalent to
\[
n^2-n-36=(n-7)^2+13(n-7)+6>0.
\]
Checking $n=1,\ldots,6$ gives only the correction $-1$ at $n=1$.
Consequently $a(n)=\lfloor Q(n)\rfloor-\mathbf1_{n=1}$ for every $n\geq1$,
and exact summation proves Barker's empirical function
\[
\sum_{n\geq1}a(n)x^n=
-\frac{x^2(x^{10}-2x^9+x^7+x^4+x^2-x+2)}
{(x-1)^3(x^2+x+1)(x^6+x^3+1)}.
\]

\subsection{A184534: reciprocal fractional parts of cube roots}
Let $a(n)=\lfloor1/\{(n^3+4)^{1/3}\}\rfloor$ for $n\geq1$.
The exact formula is
\begin{equation}\label{round2:eq:r2cuberoot}
a(n)=\left\lfloor\frac{3n^2}{4}\right\rfloor+
\mathbf1_{n=1}+\mathbf1_{n=3}.
\end{equation}
Put $\delta=(n^3+4)^{1/3}-n$. Since $n^3<n^3+4<(n+1)^3$,
this is the fractional part in the definition. Expanding the cube gives
$0<\delta<4/(3n^2)$, and rationalization gives
\[
\frac1\delta=\frac{3n^2}{4}+\frac{3n\delta+\delta^2}{4}.
\]
For $n\geq5$, the positive residual is smaller than
\[
\frac1n+\frac4{9n^4}\leq\frac15+\frac4{9\cdot5^4}<\frac14.
\]
The fractional part of $3n^2/4$ is either 0 or $3/4$, so its floor
is unchanged. For $n=1,2,3,4$, comparing $n^3+4$ to the rational cubes
$(n+1/(v+1))^3$ and $(n+1/v)^3$, with $v=1,3,7,12$, respectively,
gives the exact remaining values. This proves \eqref{round2:eq:r2cuberoot}.
Its generating function is exactly Barker's empirical one:
\[
\frac{3x(1+x)}{4(1-x)^3}-\frac{3x}{4(1-x^2)}+x+x^3
=\frac{x(x^6-2x^5+x^4-x^2-x-1)}{(x-1)^3(x+1)}.
\]
The separate empirical closed form printed in the entry uses a shifted
index; the count here is for the generating-function assertion.

\subsection{A002620 and A001840: two exponential floor identities}
For every real $t>0$ define
\[
L(t)=\frac{t\coth t-1}{t^2}.
\]
We first prove the global inequalities
\begin{equation}\label{round2:eq:r2coth}
\frac1{3+t^2}<L(t)<\frac13.
\end{equation}
The coefficients of $t^{2k+1}$, $k\geq1$, in the two entire functions
\[
(1+t^2/3)\sinh t-t\cosh t,\qquad
(3+t^2)t\cosh t-(3+2t^2)\sinh t
\]
are respectively
\[
\frac{4k(k-1)}{3(2k+1)!},\qquad
\frac{8k^2(k-1)}{(2k+1)!}.
\]
They vanish for $k=1$ and are strictly positive for $k\geq2$; the
lower coefficients also vanish. The convergent series are therefore
positive for $t>0$. Dividing by positive $\sinh t$ and rearranging
proves \eqref{round2:eq:r2coth}.

\paragraph{A002620.}
The selected empirical identity of Forberg (2013) is
\[
\left\lfloor\frac{n}{e^{4/n}-1}\right\rfloor
=\left\lfloor\frac{(n-1)^2}{4}\right\rfloor\qquad(n\geq1).
\]
For $n\geq2$ put $t=2/n\leq1$. Then
\[
\frac{n}{e^{4/n}-1}=\frac{n^2}{4}-\frac n2+L(t),\qquad
\frac14<L(t)<\frac13.
\]
The base polynomial has fractional part 0 for even $n$ and $3/4$
for odd $n$. In either case adding $L(t)$ gives the same floor as
adding $1/4$, which is the claimed right-hand side. For $n=1$,
$0<1/(e^4-1)<1$, so the identity also holds.

\paragraph{A001840.}
Here Forberg's selected empirical identity is
\[
\left\lfloor\frac{n+3}{e^{6/(n+3)}-1}\right\rfloor
=\left\lfloor\frac{(n+1)(n+2)}6\right\rfloor\qquad(n\geq0).
\]
Set $m=n+3\geq3$ and $t=3/m\leq1$. The left-hand expression before
taking its floor equals
\[
\frac{m^2}{6}-\frac m2+\frac32L(t),\qquad
\frac38<\frac32L(t)<\frac12.
\]
The fractional part of $m^2/6-m/2$ is either 0 or $2/3$.
Adding any number in $(3/8,1/2)$ consequently gives the same floor
as adding $1/3$. This is exactly the right-hand side, proving the
identity for every permitted index. No other conjecture appearing
in these two entries is included in this result.

\subsection{A227018: harmonic means between consecutive binomial values}
Put $g(n)=n(n+1)(n+2)(n+3)/24$ and
\[
M_n=\frac{g(n)-g(n-1)}{H_{g(n)}-H_{g(n-1)}},\qquad
Q(n)=\frac{n^4}{24}+\frac{n^3}6+\frac{11n^2}{72}-\frac n{36}
       +\frac{119}{270}.
\]
We prove, for every $n\geq1$,
\begin{equation}\label{round2:eq:r2harmonicfloor}
\lfloor M_n\rfloor=\lfloor Q(n)\rfloor+
\mathbf1_{n\in\{1,2,5,8\}}.
\end{equation}
This proves the entry's conjectured generating function.

Let $a=g(n-1)$, $b=g(n)$, $\Delta=b-a$, and $t=2/(n+1)$.
For $n\geq2$ all these denominators are nonzero and $0<t<1$.
Since $b/a=(n+3)/(n-1)$,
$\log(b/a)=2\operatorname{arctanh}t$.
Define the explicit rational quantities
\[
S=2\sum_{j=0}^{3}\frac{t^{2j+1}}{2j+1},\qquad
R=\frac{2t^9}{9(1-t^2)},\qquad
C=\frac12\left(\frac1b-\frac1a\right)
 -\frac1{12}\left(\frac1{b^2}-\frac1{a^2}\right).
\]
The positive logarithm series gives $S<\log(b/a)<S+R$.
The signed harmonic-number remainder bound is
\[
H_m=\log m+\gamma+\frac1{2m}-\frac1{12m^2}+\rho_m,\qquad
0<\rho_m<\frac1{120m^4}.
\]
This is the positive-real signed remainder in the digamma expansion;
see NIST DLMF 5.11.2 and 5.11(ii),
\texttt{https://dlmf.nist.gov/5.11}.
It follows that
\[
L=S+C-\frac1{120a^4}<H_b-H_a<
U=S+R+C+\frac1{120b^4}.
\]

The following exact polynomial certificate proves the required bounds
throughout the infinite range $n\geq13$:
\begin{align*}
P_-(n)&=42525n^4(n-1)^2(n+1)^7(n+2)^4(n+3)^4(\Delta-QU),\\
P_+(n)&=945n^4(n-1)^4(n+1)^7(n+2)^4(n+3)^2
             ((Q+1/270)L-\Delta).
\end{align*}
The two expressions are polynomials of degrees 18 and 20. Their complete
ascending coefficients after substituting $n=13+u$ are in the following
table. All coefficients within each polynomial's degree are positive.
\begin{center}\small
\begin{tabular}{r|r|r}
$j$&$[u^j]P_-(13+u)$&$[u^j]P_+(13+u)$\\\hline
0&55835947384574015510839296&109658760095545747439616\\
1&70973316689756907301797888&265101895223921156751360\\
2&42589713597153350782777344&252456303046151176335360\\
3&16031633945158486688025600&138768983076755766398976\\
4&4241979336413711433369600&51220726939926411314688\\
5&837862165053501355509760&13737347950186928090880\\
6&128026064500098106824832&2805670927599765213312\\
7&15474053555433829401408&449601046378547440320\\
8&1499750431545219474368&57646663206762721184\\
9&117430034768907000512&5989088391494318208\\
10&7445872019338020336&507962140667685096\\
11&381415588809560224&35285154142911672\\
12&15667255218382072&2006166896915838\\
13&509062269622464&92930287845888\\
14&12796039020312&3475232824328\\
15&240104119968&103354758768\\
16&3167005416&2388148548\\
17&26204160&41330304\\
18&102360&504096\\
19&--&3864\\
20&--&14\\
\end{tabular}\end{center}
All displayed denominator factors are positive for $n\geq13$, as are
$Q$ and $\Delta$. Hence
$\Delta-QU>0$ and $(Q+1/270)L-\Delta>0$. In particular $L>0$, and
division by positive quantities gives the rigorous enclosure
\[
Q(n)<M_n<Q(n)+\frac1{270}\qquad(n\geq13).
\]
The identity
\[
Q(n+36)-Q(n)=6n^3+342n^2+8435n+77957
\]
proves period 36 for the fractional part. Its numerators over 270,
in residue classes $1,\ldots,36$, are
\begin{gather*}
209,269,164,29,269,209,254,269,119,209,134,29,\\
29,269,74,119,269,119,74,269,29,29,134,209,\\
119,269,254,209,269,29,164,269,209,119,134,119.
\end{gather*}
They are at most 269, so the strict enclosure proves equality of floors.
Exact rational harmonic sums for $n=1,\ldots,12$ give
\[
1,3,9,24,51,95,164,266,407,598,850,1174,
\]
which differ from $\lfloor Q(n)\rfloor$ precisely by the four corrections
in \eqref{round2:eq:r2harmonicfloor}.

Finally, if $Q(n)=\sum_{k=0}^4q_kn^k$, $f_r=\{Q(r)\}$, and
$\theta=x\,d/dx$, the exact generating function is
\[
G(x)=\sum_{k=0}^4q_k\theta^k\frac{x}{1-x}
-\frac{\sum_{r=1}^{36}f_rx^r}{1-x^{36}}+x+x^2+x^5+x^8.
\]
This supplies an explicit reconstructible rational expression for every
coefficient. Its denominator is $(1-x)^4(1-x^{36})$, and exact
simplification yields the entry's conjectured function multiplied by
$x$, because the source uses $\sum_{n\geq1}a(n)x^{n-1}$.
The full source numerator and the exact identity check are in
\texttt{r2\_\allowbreak floor\_\allowbreak harmonic\_\allowbreak certificate.json}
and its companion verifier. This result concerns the specific
binomial-block assertion, not an unrestricted polynomial generalization.

\clearpage
% Source fragment: r2_gf_proofs.tex

\section{Round 2: generating functions and exact recurrence proofs}
\label{round2:sec:r2gfproofs}

This section proves explicitly conjectural assertions attached to 42 distinct
OEIS entries, using pages checked on 20 October 2026. All entries are disjoint
from the first-round ledger. Forty entries concern a recurrence (in three cases
a generating-function assertion is also proved); A005573 concerns a
hypergeometric evaluation, and A368633 concerns parity. Each entry counts once.
Existing generating functions and finite sums are identified as source premises;
their presence is not presented as a new result. No first-publication priority,
external submission, or OEIS acceptance is claimed.

The on-page audit removed A001712: its 2020 continuation explicitly derives the
recurrence despite the surviving conjecture label. The normalization correction
in A065919 is also excluded. Unrelated conjectures on a retained page are outside
the asserted scope unless expressly named below.

\subsection{Exact certificates and their logical meaning}

Write \(\theta=x\,d/dx\). For \(F=\sum_{n\geq0}a_nx^n\), with omitted
initial indices extended by zero, define
\[
 \mathcal L_p=\sum_{j=0}^{r}x^j p_j(\theta+j).
\]
Then, solely by coefficient extraction,
\[
 [x^n]\mathcal L_pF=\sum_{j=0}^{r}p_j(n)a_{n-j}.
 \tag{R2-G1}
\]
An exact identity \(\mathcal L_pF=P(x)\), where \(P\) is a polynomial of
degree below \(N\), proves the printed recurrence for all \(n\geq N\).
The bounds in the individual records are part of the conclusions. No division
by a leading recurrence coefficient is used. If the entry gives no lower
bound, we state an explicit sufficient bound; in particular, A135925 requires
\(n\geq14\), A239230 requires \(n\geq8\), and A032115 requires \(n\geq6\).

For an exponential generating function \(F=\sum a_nx^n/n!\), set
\[
 q_j(n)=p_j(n)\prod_{h=0}^{r-j-1}(n-j-h).
 \tag{R2-G2}
\]
For \(n\geq r\), the coefficient of \(x^n\) in \(\mathcal L_qF\) is the
desired recurrence divided by \((n-r)!\). Thus a polynomial boundary works
in the exponential convention as well.

For the rational and radical certificates write
\[
 x=f(t),\qquad F=\sum_iR_i(t)\sqrt{D_i(t)}.
\]
The specified formal branch selects the series at \(x=0\). Ordinary
differentiation gives the exact reduction rule
\[
 \theta(R_i\sqrt{D_i})=
 \frac{f(t)}{f'(t)}
 \left(R_i'(t)+\frac{R_i(t)D_i'(t)}{2D_i(t)}\right)\sqrt{D_i(t)}.
 \tag{R2-G3}
\]
The per-entry certificate prints all \(p_j\), all seed functions, and the
resulting boundary polynomial. The nonrational basis coefficients are zero.
These identities can be checked by clearing denominators in the rational
function field. They do not rely on fitting or testing finitely many terms.
The standalone script \texttt{r2\_gf\_verify.py} performs those exact operations
from \texttt{r2\_gf\_certificates.json}, without source-page snapshots or
network access.

\subsection{Elementary generating-function derivations}

For A025258, let \(F=\sum_{n\geq1}a_nx^n\). The defining convolution gives
\[
 F-x+x^2-x^3=F(F-x+x^2),
 \qquad F=x\left(1-x+\frac{x^2}{1-F}\right).
\]
The formal root with \(F(0)=0\) is the radical printed in its certificate.
This also agrees with the existing December 2025 formula of Seiichi Manyama.

For A075045, use Manyama's existing finite sum
\[
 a_n=\sum_{k=0}^n2^{n-k}\binom{3n+4}{k}
     =[u^n]\frac{(1+u)^{3n+4}}{1-2u}.
\]
The formal change of variable \(t=x\phi(t)\) in coefficient residues gives
the diagonal Lagrange identity
\[
 \sum_{n\geq0}[u^n]\psi(u)\phi(u)^n x^n
 =\frac{\psi(t)}{1-x\phi'(t)}.
\]
Indeed, substitute \(x=u/\phi(u)\) in the residue computing each coefficient;
its Jacobian is precisely \(1-x\phi'(u)\), which cancels the displayed
denominator. Here \(\phi(u)=(1+u)^3\), \(\psi(u)=(1+u)^4/(1-2u)\), so
\[
 x=\frac{t}{(1+t)^3},\qquad F=\frac{(1+t)^5}{(1-2t)^2}.
\]
Putting \(g=t/(1+t)\) gives \(g(1-g)^2=x\) and
\(F=(3g-1)^{-2}(1-g)^{-3}\), proving the separately empirical generating
function as well as the recurrence. That empirical function is not a premise.

For A082397 put \(h_k=\binom{k}{\lfloor(k-1)/2\rfloor}\), taking \(h_0=0\),
and \(C(z)=(1-\sqrt{1-4z})/(2z)\). Splitting even and odd indices gives
\[
 h_{2m}=\binom{2m}{m}-C_m,\qquad
 h_{2m+1}=2\binom{2m}{m}-C_m,
\]
and hence
\[
 H(z)=\sum_{k\geq0}h_kz^k
 =\frac{1+2z}{\sqrt{1-4z^2}}-(1+z)C(z^2).
\]
In the accepted alternating binomial sum, summing first over \(n\geq k\)
uses \(\sum_{n\geq k}\binom{n+1}{k+1}x^n=x^k/(1-x)^{k+2}\).
Consequently \(F(x)=-H(-x/(1-x))/(1-x)^2\), the certified radical.

For A186335 use its existing double sum. Let \(m=n-k-j\) and
\(\ell=k-j-m\); the nonzero terms have \(j,m,\ell\geq0\), with
\(n=2j+2m+\ell\). Thus
\[
 F=\sum_{j,m,\ell\geq0}
 \binom{m+\ell}{m}\binom{j+m+\ell}{j}\binom{2m}{m}
 x^{2j+2m+\ell}.
\]
Summing in \(j\), then \(\ell\), with the binomial series gives
\[
 F=\frac1{1-x-x^2}\sum_{m\geq0}\binom{2m}{m}
       \left(\frac{x^2}{1-x-x^2}\right)^m
 =\frac1{\sqrt{(1-x-x^2)(1-x-5x^2)}}.
\]

For A133158, the defining twice-binomial transform implies
\(F(x)=(1-2x)^{-1}B(x/(1-2x))\). The existing formula in A026641 is
\[
 B(z)=\frac2{3\sqrt{1-4z}-1+4z}.
\]
Substitution and rationalization yield the printed certificate.
For A157021 the defining binomial sum gives
\[
 F(x)=\frac{1-2x+2x^2}{(1-x)^3}
 C\left(\frac{x^4}{(1-x)^4}\right):
\]
indeed \(\sum_{k\geq0}C_{\lfloor k/2\rfloor}z^{2k}=(1+z^2)C(z^4)\),
followed by the usual binomial transform \(z=x/(1-x)\).
For A141353, the definition immediately gives \(F=C(x)+2x/(1-2x)\).

\subsection{Rational parameters for nested and inverse series}

The remaining substitutions in this group are consequences of existing
source formulas. For A089354 set \(t=g\) in
\(g=1+xg^3\), \(F=g/(1+xg^2)\); this gives
\(x=(t-1)/t^3\), \(F=t^2/(2t-1)\), about \(t=1\).
For A218045 use Kruchinin's November 2014 equation
\[
 F=xG,\qquad G=x\left(1+\frac{G(G+1)}{1-G}\right)^2
 =x\left(\frac{1+G^2}{1-G}\right)^2.
\]
Putting \(t=G\) gives \(x=t(1-t)^2/(1+t^2)^2\), \(F=xt\), about \(t=0\).
The later exponential formula on that page is inconsistent with its zero
constant coefficient; it is not used.

For A006343, Bala's June 2015 inverse-series formula is \(F=w/x\),
\(x=w/(1+w^2C(w))\). Setting \(t=C(w)\), so \(w=(t-1)/t^2\), gives
\[
 x=\frac{t(t-1)}{t^3+(t-1)^2},\qquad
 F=1+\frac{(t-1)^2}{t^3},\qquad t(0)=1.
\]
For A158826, composition of inverses shows that the third iterate of
\(xC(x)\) is the inverse of the third iterate of \(t(1-t)\). Therefore
\[
 x=t-3t^2+6t^3-9t^4+10t^5-8t^6+4t^7-t^8,
 \qquad F=t,
\]
with \(t(0)=0\), exactly as the existing source formula states.

For A078999 use Irvine's July 2025 relation
\[
 a_n=\tfrac12\bigl(A078995(n+1)-(5n+6)A002293(n+1)\bigr).
\]
The accepted generating functions are \(t=1+xt^4\) for A002293 and
\(H=t^2/(3t-4)^2\) for A078995 (van Hoeij, May 2013). Shifting the indices
in the relation gives
\[
 F=\frac{H-t-5xt'}{2x},\qquad x=\frac{t-1}{t^4},\qquad t(0)=1.
\]

For A239230 define
\[
 v=\sqrt{1-4x},\qquad w=\sqrt{5+2v-3v^2},\qquad
 B=\frac{-w+v+5}{4}.
\]
Its defining generating function is \(F=-xB'/B\). The conic is rationally
parametrized by the line \(w=2+t(v-1)\):
\[
 v=1-\frac{4(t+1)}{t^2+3},\qquad
 w=2-\frac{4t(t+1)}{t^2+3},\qquad x=\frac{1-v^2}{4}.
\]
The branch at \(t=-1\) has \(v=1,w=2,x=0\), and \(dx/dt\ne0\).
Consequently \(F=-xB_t/(x_tB)\) is the certified rational function.

\subsection{Stirling sums as coefficients of rising products}

Write \((z)^{\overline m}=\prod_{i=0}^{m-1}(z+i)\).
For A001723 and A001724 the source formula is
\[
 a_n=\sum_{k=0}^n(-1)^{n+k}\binom{k+d}{d}5^k s(n+d,k+d),
 \qquad d=3,4,
\]
where \(s\) is the signed Stirling number of the first kind. Expanding the
rising product, the two signs cancel and give
\[
 a_n=[u^d](5+u)^{\overline{n+d}}.
\]
For either printed coefficient list, with \(r=d+1\), direct polynomial
multiplication gives the finite identity
\[
 \sum_{j=0}^{r}p_j(n)
     \prod_{h=j}^{r-1}(u+5+n+d-1-h)=u^{d+1}.
 \tag{R2-G4}
\]
For \(n\geq r\), factor the common product
\((5+u)^{\overline{n-1}}\) from the recurrence. Its remaining factor is
\(u^{d+1}\), so its coefficient of \(u^d\) is zero. This proves the two
recurrences for every stated index.

\subsection{The three reversed tribonacci triangles}

Let the triangle be A027023 (type I), A027052 (type II), or A027082
(type III), with rows indexed from \(n=0\). All have the interior rule
\[
 T(n,k)=T(n-1,k-3)+T(n-1,k-2)+T(n-1,k-1),
\]
and values outside a row are zero. The initial row is \((1)\).
Types I and II have left triples \((1,1,1)\), \((1,0,1)\), respectively,
and rightmost entry 1. Type III has left triple \((1,1,1)\), while its
rightmost entry sums the two rightmost entries of the preceding row.
The row at \(n=1\) is respectively \((1,1,1)\), \((1,0,1)\), \((1,1,1)\).
At \(n=2\) the rows are \((1,1,1,3,1)\), \((1,0,1,2,1)\), and
\((1,1,1,3,2)\), fixing the overlapping boundary conventions.

Reverse each row and put
\[
 U_n(z)=\sum_{m=0}^{2n}T(n,2n-m)z^m,\qquad
 B(x,z)=\sum_{n\geq0}U_n(z)x^n,\qquad B_j=[z^j]B.
\]
The bulk multiplication is \(z+1+z^{-1}\). Correcting its extraneous
\(z^{-1}\) coefficient, its reset rightmost coefficient, and its three
leftmost coefficients gives, for types I and II,
\[
 (1-x(z+1+z^{-1}))B
 =1+\epsilon x+\frac{xL(z)}{1-xz^2}-xB_1-\frac{x}{z(1-x)},
 \tag{R2-G5}
\]
where \((\epsilon,L)=(1,z^2-1)\) or \((0,z^2-z)\), respectively.
Here \(B_0=1/(1-x)\); its unwanted shift produces the last term.
The right-edge reset produces \(-xB_1\); the left-edge corrections form
the geometric series \(xL(z)/(1-xz^2)\); the exceptional first row produces
\(\epsilon x\). Thus every correction follows from the definition.

Let \(t=x+O(x^2)\) be the small root of the kernel, so
\(x=t/(1+t+t^2)\). Substitution into (R2-G5) gives
\[
 B_1=\frac1x+\epsilon+\frac{L(t)}{1-xt^2}
                   -\frac1{t(1-x)}.
\]
For type I, coefficient extraction at \(z^1,z^2,z^3\) gives
\[
 B_2=\frac{1-x}{x}B_1-B_0,\quad
 B_3=\frac{1-x}{x}B_2-B_1-1+x,\quad
 B_4=\frac{1-x}{x}B_3-B_2.
\]
Therefore the target functions are \(B_2-x\) for A027051, \(B_3\) for
A027027, and \(B_4\) for A027028. The subtraction in A027051 removes the
column's extra \(n=1\) term, absent from its specified offset. Type II's
\(B_1\) gives A027056.

For type III the rightmost coefficient propagates instead of being reset,
and the same row calculation gives
\[
 (1-x(z+1+z^{-1}))B
 =1+x+\frac{x(z^2-1)}{1-xz^2}-\frac{xB_0}{z}.
\]
Kernel substitution therefore gives the target A027087 as
\[
 B_0=\frac{t}{x}\left(1+x+\frac{x(t^2-1)}{1-xt^2}\right).
\]
All five expressions are rational in \(t\); (R2-G3) proves their recurrences.

\subsection{Two central diagonals in A027113}

This fourth triangle has left boundary \((1,1,2)\), the same interior rule,
and rightmost entry equal to the sum of the two rightmost entries of the
preceding row. Let \(c_k=[x^k](1-x-x^2-x^3)^{-1}\). Induction using the
row rule proves \(T(n,k)=c_k\) whenever \(0\leq k\leq n+1\), with the
initial row treated separately. Put
\[
 D_j(x)=\sum_{n\geq j}T(n,n+j)x^n\quad(j\geq0).
\]
The first two functions are \(D_0=C\), \(D_1=(C-1-x)/x\), where
\(C=(1-x-x^2-x^3)^{-1}\). For \(n\geq j+1\), the defining rule says
\(d_j(n)=d_j(n-1)+d_{j-1}(n-1)+d_{j-2}(n-1)\).
Using \(d_2(2)=3,d_3(3)=7,d_4(4)=16\), from the initial rows, yields
\[
 \begin{aligned}
 (1-x)D_2&=x(D_0+D_1)-x,\\
 (1-x)D_3&=x(D_2+D_1)-2x^2,\\
 (1-x)D_4&=x(D_3+D_2)-3x^3.
 \end{aligned}
\]
Solving gives precisely Barker's two conjectural rational generating functions:
\[
 D_3=\frac{x^3(7+2x+x^2-2x^3)}{(1-x)^2(1-x-x^2-x^3)},
\]
\[
 D_4=\frac{x^4(16-5x-x^2-5x^3+3x^4)}{(1-x)^3(1-x-x^2-x^3)}.
\]
Multiplying by the denominators proves the paired constant-coefficient
recurrences, with the source's bounds \(n>7\), \(n>9\), respectively.

\subsection{Fixed diagonals of the menage hit polynomials}

The existing Jovovic bivariate generating function in A058087 is
\[
 P(x,y)=(1-x(y-1))\sum_{m\geq0}
       \frac{m!(xy)^m}{(1+x(y-1))^{2m+1}}.
\]
Its coefficient is \([x^ny^k]P=T(n,k)\), with row coefficients in
descending-power order. Replacing \((x,y)\) by \((xt,1/t)\) reverses
that order. Write
\[
 H(x,t)=\frac{1-x(1-t)}{1+x(1-t)}f\left(\frac{x}{(1+x(1-t))^2}\right),
 \qquad f(z)=\sum_{m\geq0}m!z^m.
\]
Then \([t^d]H\) gives A000159, A000181, and A000185 for \(d=3,4,5\).
Everything here is formal; no convergence of the factorial series is asserted.
The factorial identity gives
\[
 z^2f'(z)+(z-1)f(z)+1=0.
 \tag{R2-G6}
\]
For \(h=x/(1+x(1-t))^2\), put
\[
 K=\frac{h_t(1-h)}{h^2},\qquad J=-\frac{h_t}{h^2},\qquad
 f(h)=\sum_{k\geq0}(\alpha_kF_0+\beta_k)t^k,
\]
where \(F_0=f(x/(1+x)^2)\), \(\alpha_0=1,\beta_0=0\).
Equation (R2-G6) gives the finite rational recursion
\[
 \begin{aligned}
 (k+1)\alpha_{k+1}&=\sum_{i=0}^k[t^i]K\,\alpha_{k-i},\\
 (k+1)\beta_{k+1}&=\sum_{i=0}^k[t^i]K\,\beta_{k-i}+[t^k]J.
 \end{aligned}
\]
Multiply by the rational prefactor of \(H\). This derives the printed
\(F=A_dF_0+B_d\). For \(h_0=x/(1+x)^2\), differentiating (R2-G6) also gives
\[
 F_0'=UF_0+V,\qquad
 U=\frac{h_0'(1-h_0)}{h_0^2},\qquad V=-\frac{h_0'}{h_0^2}.
\]
The exact operator calculation on the basis \((F_0,1)\) gives zero coefficient
of \(F_0\) and the stated polynomial boundary. The three recurrence assertions
follow. Their separate prime-divisibility claims are not included.

\subsection{Exponential and symmetric-square differential spaces}

For A080894, A272261, A144575, A053532 and A250917, the EGF \(F\) is
specified in its individual record. Its logarithmic derivative has the form
\(F'/F=\ell_0+\ell_1y\), \(y^2=D(x)\). Writing
\(\theta^kF=F(A_k+B_ky)\), with \((A_0,B_0)=(1,0)\), gives
\[
 \begin{aligned}
 A_{k+1}&=x(A_k'+\ell_0A_k+D\ell_1B_k),\\
 B_{k+1}&=x\left(B_k'+\frac{D'}{2D}B_k+\ell_1A_k+\ell_0B_k\right).
 \end{aligned}
 \tag{R2-G7}
\]
Applying the coefficients \(q_j\) from (R2-G2) gives exactly zero in both
components. This proves all five recurrences. In A053532 there is an especially
short direct calculation: \(F'/F=x^3/(2(1-x))\), so
\(2(1-x)F'=x^3F\); equating coefficients yields the printed assertion.

For A032115, Howroyd's existing EGF is
\[
 F=1+\frac{e^{-x}}{2(1-2x)}-\frac{1-2x}{2}e^x.
\]
On \(R(x)e^{cx}\), Euler differentiation replaces \(R\) by
\(x(R'+cR)\). The two exponential components vanish under the adjusted
operator, and the rational component gives \(2x^5\). This proves the
recurrence for \(n\geq6\); a universal assertion beginning at \(n=5\)
would fail there.

For A168452 let \(G=\sum_{n\geq0}C_nC_{n+1}x^n\); its defining
self-convolution gives \(F=G^2\). The coefficient ratio
\[
 \frac{C_nC_{n+1}}{C_{n-1}C_n}
 =\frac{4(2n-1)(2n+1)}{(n+1)(n+2)}
\]
implies
\[
 \bigl((\theta+1)(\theta+2)-4x(2\theta+1)(2\theta+3)\bigr)G=2.
\]
Thus \(G''=AG+BG'+C\), with the three rational coefficients printed below.
The six functions \(G^2,GG',(G')^2,G,G',1\) form a differential space.
For a polynomial expression \(R(x,G,H)\), \(H=G'\), its Euler derivative is
\[
 x\left(\partial_xR+H\partial_GR+(AG+BH+C)\partial_HR\right).
 \tag{R2-G8}
\]
The conjectured operator applied to \(G^2\) reduces to \(72-1152x\), with
all other basis coefficients zero. Hence the recurrence holds for \(n\geq4\).

\subsection{The hypergeometric evaluation and the parity conjecture}

For A005573, the proposed evaluation is
\[
 a_n=6^n\,{}_2F_1\left(\frac12,-n;2;\frac23\right),\qquad n\geq0.
\]
The series terminates at \(k=n\). The identities
\[
 \frac{(-n)_k}{k!}=(-1)^k\binom nk,
 \qquad \frac{(1/2)_k}{(2)_k}=\frac{C_k}{4^k}
\]
turn its \(k\)-th summand into
\((-1)^k\binom nk C_k6^{n-k}\), exactly the existing finite-sum formula.
This proves the evaluation term by term.

For A368633, the defining equation is
\[
 F(x)=1+2xF(x)^2-xF(-x)^2.
\]
It uniquely determines integral coefficients:
the coefficient of degree \(n>0\) only uses earlier coefficients on its
right-hand side. Modulo 2 it becomes
\[
 F(x)=1+xF(x)^2=1+xF(x^2).
\]
Here \(-x=x\) in characteristic 2 and the Frobenius identity applies
coefficientwise. Therefore \(a_0\) is odd, \(a_{2m}\) is even for \(m\geq1\),
and \(a_{2m+1}\equiv a_m\pmod2\). Repeatedly applying this recursion shows
\[
 F(x)\equiv\sum_{k\geq0}x^{2^k-1}\pmod2.
\]
Thus precisely the indices \(n=2^k-1\), including \(n=0\), have odd
coefficients. The page's separate D-finite conjecture is not used or counted.

\subsection{Per-entry statements and exact certificates}

In the following records, every displayed coefficient vector means
\(\sum_{j=0}^{r}p_j(n)a_{n-j}=0\) in the stated range. All coefficients and
boundaries are exact integers or rational functions. For a radical certificate,
the ordinary convention is used unless the exponential convention is stated;
\(D_0=1\) denotes its rational component. Each source URL identifies the
current conjecture and the seed premise, and additional seed URLs are listed
where required.

\subsection{A089354}
Source: \url{https://oeis.org/A089354}. Offset 0.

\textbf{Conclusion.} The selected printed recurrence holds for every integer \(n\geq3\). Its complete coefficient list is:
\[
\begin{aligned}
p_{0}(n)&=128 n^{2} + 64 n
\end{aligned}
\]
\[
\begin{aligned}
p_{1}(n)&=- 1136 n^{2} + 1640 n \\
&\quad -528
\end{aligned}
\]
\[
\begin{aligned}
p_{2}(n)&=1760 n^{2} - 7802 n \\
&\quad +7848
\end{aligned}
\]
\[
\begin{aligned}
p_{3}(n)&=513 n^{2} - 2052 n \\
&\quad +1995
\end{aligned}
\]
The source premise is the existing implicit GF, with the rational substitution derived above. The branch is \(t(0)=1\).
The ordinary generating function has certificate
\[
F=\sum_{i=0}^{0}R_i\sqrt{D_i}.
\]
\[
x=\frac{t - 1}{t^{3}}.
\]
\[
R_{0}=\frac{t^{2}}{2 t - 1}.
\]
\[
D_{0}=1.
\]
Applying (R2-G3) gives the exact identity
\[
\mathcal L_pF=- 4 x \left(19 x + 6\right).
\]
Every remaining radical coefficient is zero. Equations (R2-G1) prove the stated conclusion.

\subsection{A247471}
Source: \url{https://oeis.org/A247471}. Offset 0.

\textbf{Conclusion.} The selected printed recurrence holds for every integer \(n\geq8\). Its complete coefficient list is:
\[
\begin{aligned}
p_{0}(n)&=n + 3
\end{aligned}
\]
\[
\begin{aligned}
p_{1}(n)&=- 3 n - 3
\end{aligned}
\]
\[
\begin{aligned}
p_{2}(n)&=0
\end{aligned}
\]
\[
\begin{aligned}
p_{3}(n)&=n - 9
\end{aligned}
\]
\[
\begin{aligned}
p_{4}(n)&=4 n - 18
\end{aligned}
\]
\[
\begin{aligned}
p_{5}(n)&=n - 9
\end{aligned}
\]
\[
\begin{aligned}
p_{6}(n)&=9 - n
\end{aligned}
\]
Use the existing Deutsch formula \(F=x^3(1-x-x^2)g^2\), \(g=1+xg+x^2g+x^3g^2\). The formal solution is \(g=(1-x-x^2-\sqrt{(1-x-x^2)^2-4x^3})/(2x^3)\), with \(g(0)=1\). Substitution gives the displayed coefficients.
The ordinary generating function has certificate
\[
F=\sum_{i=0}^{1}R_i\sqrt{D_i}.
\]
\[
x=x.
\]
\[
R_{0}=- \frac{\left(x^{2} + x - 1\right) \left(x^{4} - x^{2} - 2 x + 1\right)}{2 x^{3}}.
\]
\[
D_{0}=1.
\]
\[
R_{1}=- \frac{\left(x^{2} + x - 1\right)^{2}}{2 x^{3}}.
\]
\[
D_{1}=\left(x^{2} - 3 x + 1\right) \left(x^{2} + x + 1\right).
\]
Applying (R2-G3) gives the exact identity
\[
\mathcal L_pF=- 2 x^{3} \left(x^{2} - x + 3\right) \left(x^{2} + x - 1\right).
\]
Every remaining radical coefficient is zero. Equations (R2-G1) prove the stated conclusion.

\subsection{A135925}
Source: \url{https://oeis.org/A135925}. Offset 0.

\textbf{Conclusion.} The selected printed recurrence holds for every integer \(n\geq14\). Its complete coefficient list is:
\[
\begin{aligned}
p_{0}(n)&=30715031990396 n^{2} - 433554237731400 n \\
&\quad +402839205741004
\end{aligned}
\]
\[
\begin{aligned}
p_{1}(n)&=- 291262905844111 n^{2} + 4745710679447853 n \\
&\quad -12935393149282124
\end{aligned}
\]
\[
\begin{aligned}
p_{2}(n)&=805570866653979 n^{2} - 14936203714114791 n \\
&\quad +60755127774344160
\end{aligned}
\]
\[
\begin{aligned}
p_{3}(n)&=- 393057731655406 n^{2} + 8294532445122504 n \\
&\quad -43776275319524372
\end{aligned}
\]
\[
\begin{aligned}
p_{4}(n)&=- 539125155360312 n^{2} + 12669158483525292 n \\
&\quad -74126175519017400
\end{aligned}
\]
The seed is the entry definition. Put \(s=\sqrt{1-4x}\); the defining function is the product \((x^2+x-1)(1+2x+s)(1-5x+2x^2+(3x-1)s)^2(-1+3x+2x^2+(1-x)s)\). Reducing modulo \(s^2=1-4x\) gives the two coefficients below. Its polynomial boundary explains the lower bound 14.
The ordinary generating function has certificate
\[
F=\sum_{i=0}^{1}R_i\sqrt{D_i}.
\]
\[
x=x.
\]
\[
\begin{aligned}
R_{0}&=16 x^{9} - 160 x^{8} \\
&\quad - 824 x^{7} + 1024 x^{6} \\
&\quad +1240 x^{5} - 2152 x^{4} \\
&\quad +1056 x^{3} - 216 x^{2} \\
&\quad +16 x
\end{aligned}
\]
\[
D_{0}=1.
\]
\[
\begin{aligned}
R_{1}&=48 x^{8} + 72 x^{7} \\
&\quad - 640 x^{6} - 72 x^{5} \\
&\quad +1016 x^{4} - 720 x^{3} \\
&\quad +184 x^{2} - 16 x
\end{aligned}
\]
\[
D_{1}=1 - 4 x.
\]
Applying (R2-G3) gives the exact identity
\[
\begin{aligned}
\mathcal L_pF&=- 8628263825301312 x^{13} - 56627815207112768 x^{12} \\
&\quad - 82029355854002112 x^{11} + 107522236010828544 x^{10} \\
&\quad +267411638401515712 x^{9} + 26860851665524032 x^{8} \\
&\quad - 96203187527862528 x^{7} + 17483921103890240 x^{6}
\end{aligned}
\]
Every remaining radical coefficient is zero. Equations (R2-G1) prove the stated conclusion.

\subsection{A218045}
Source: \url{https://oeis.org/A218045}. Offset 0.

\textbf{Conclusion.} The selected printed recurrence holds for every integer \(n\geq5\). Its complete coefficient list is:
\[
\begin{aligned}
p_{0}(n)&=2 n^{3} - 5 n^{2} \\
&\quad +3 n
\end{aligned}
\]
\[
\begin{aligned}
p_{1}(n)&=- 22 n^{3} + 100 n^{2} \\
&\quad +70 - 148 n
\end{aligned}
\]
\[
\begin{aligned}
p_{2}(n)&=- 24 n^{3} + 336 n^{2} \\
&\quad +1400 - 1272 n
\end{aligned}
\]
\[
\begin{aligned}
p_{3}(n)&=512 n^{3} - 5696 n^{2} \\
&\quad +21024 n - 25760
\end{aligned}
\]
\[
\begin{aligned}
p_{4}(n)&=512 n^{3} - 6656 n^{2} \\
&\quad +28544 n - 40320
\end{aligned}
\]
Use only Kruchinin's accepted implicit GF described above. The parameter is \(t=G\), with \(t(0)=0\).
The ordinary generating function has certificate
\[
F=\sum_{i=0}^{0}R_i\sqrt{D_i}.
\]
\[
x=\frac{t \left(t - 1\right)^{2}}{\left(t^{2} + 1\right)^{2}}.
\]
\[
R_{0}=\frac{t^{2} \left(t - 1\right)^{2}}{\left(t^{2} + 1\right)^{2}}.
\]
\[
D_{0}=1.
\]
Applying (R2-G3) gives the exact identity
\[
\mathcal L_pF=2 x^{2} \left(16 x^{2} - 16 x + 1\right).
\]
Every remaining radical coefficient is zero. Equations (R2-G1) prove the stated conclusion.

\subsection{A006343}
Source: \url{https://oeis.org/A006343}. Offset 0.

\textbf{Conclusion.} The selected printed recurrence holds for every integer \(n\geq5\). Its complete coefficient list is:
\[
\begin{aligned}
p_{0}(n)&=- n^{2} - 2 n \\
&\quad +3
\end{aligned}
\]
\[
\begin{aligned}
p_{1}(n)&=11 n^{2} - 2 n \\
&\quad -45
\end{aligned}
\]
\[
\begin{aligned}
p_{2}(n)&=- 37 n^{2} + 82 n \\
&\quad +87
\end{aligned}
\]
\[
\begin{aligned}
p_{3}(n)&=29 n^{2} - 125 n \\
&\quad +78
\end{aligned}
\]
\[
\begin{aligned}
p_{4}(n)&=61 n^{2} - 289 n \\
&\quad +318
\end{aligned}
\]
\[
\begin{aligned}
p_{5}(n)&=- 155 n^{2} + 1085 n \\
&\quad -1860
\end{aligned}
\]
Use Bala's established inverse-series formula and the substitution derived above, with \(t(0)=1\).
The ordinary generating function has certificate
\[
F=\sum_{i=0}^{0}R_i\sqrt{D_i}.
\]
\[
x=\frac{t \left(t - 1\right)}{t^{3} + t^{2} - 2 t + 1}.
\]
\[
R_{0}=\frac{t^{3} + t^{2} - 2 t + 1}{t^{3}}.
\]
\[
D_{0}=1.
\]
Applying (R2-G3) gives the exact identity
\[
\mathcal L_pF=98 x^{2} - 36 x + 3.
\]
Every remaining radical coefficient is zero. Equations (R2-G1) prove the stated conclusion.

\subsection{A158826}
Source: \url{https://oeis.org/A158826}. Offset 1.

\textbf{Conclusion.} The selected printed recurrence holds for every integer \(n\geq6\). Its complete coefficient list is:
\[
\begin{aligned}
p_{0}(n)&=1053 n^{4} - 6318 n^{3} \\
&\quad +11583 n^{2} - 6318 n
\end{aligned}
\]
\[
\begin{aligned}
p_{1}(n)&=- 22824 n^{4} + 186156 n^{3} \\
&\quad - 546336 n^{2} + 678276 n \\
&\quad -295272
\end{aligned}
\]
\[
\begin{aligned}
p_{2}(n)&=191184 n^{4} - 1999920 n^{3} \\
&\quad +7835448 n^{2} - 13605720 n \\
&\quad +8810064
\end{aligned}
\]
\[
\begin{aligned}
p_{3}(n)&=- 769024 n^{4} + 9966464 n^{3} \\
&\quad - 48753200 n^{2} + 106646584 n \\
&\quad -87964536
\end{aligned}
\]
\[
\begin{aligned}
p_{4}(n)&=1470464 n^{4} - 23144448 n^{3} \\
&\quad +137212672 n^{2} - 363259008 n \\
&\quad +362432208
\end{aligned}
\]
\[
\begin{aligned}
p_{5}(n)&=- 1048576 n^{4} + 19922944 n^{3} \\
&\quad - 141787136 n^{2} + 447954944 n \\
&\quad -530100480
\end{aligned}
\]
The seed is the inverse of the third iterate of \(t(1-t)\), as derived above. Its local branch has \(t(0)=0\).
The ordinary generating function has certificate
\[
F=\sum_{i=0}^{0}R_i\sqrt{D_i}.
\]
\[
x=- t \left(t - 1\right) \left(t^{2} - t + 1\right) \left(t^{4} - 2 t^{3} + 2 t^{2} - t + 1\right).
\]
\[
R_{0}=t.
\]
\[
D_{0}=1.
\]
Applying (R2-G3) gives the exact identity
\[
\mathcal L_pF=120 x^{4} \left(16 x - 29\right).
\]
Every remaining radical coefficient is zero. Equations (R2-G1) prove the stated conclusion.

\subsection{A075045}
Source: \url{https://oeis.org/A075045}. Offset 0.

\textbf{Conclusion.} The selected printed recurrence holds for every integer \(n\geq2\). Its complete coefficient list is:
\[
\begin{aligned}
p_{0}(n)&=112 n^{3} + 296 n^{2} \\
&\quad +208 n + 24
\end{aligned}
\]
\[
\begin{aligned}
p_{1}(n)&=- 1512 n^{3} - 2862 n^{2} \\
&\quad - 1320 n - 66
\end{aligned}
\]
\[
\begin{aligned}
p_{2}(n)&=5103 n^{3} + 5832 n^{2} \\
&\quad - 567 n - 648
\end{aligned}
\]
The diagonal Lagrange argument above proves the generating function from the existing finite sum. The branch has \(t(0)=0\). This also proves the empirical GF separately printed in the entry.
The ordinary generating function has certificate
\[
F=\sum_{i=0}^{0}R_i\sqrt{D_i}.
\]
\[
x=\frac{t}{\left(t + 1\right)^{3}}.
\]
\[
R_{0}=\frac{\left(t + 1\right)^{5}}{\left(2 t - 1\right)^{2}}.
\]
\[
D_{0}=1.
\]
Applying (R2-G3) gives the exact identity
\[
\mathcal L_pF=24.
\]
Every remaining radical coefficient is zero. Equations (R2-G1) prove the stated conclusion.

\subsection{A082397}
Source: \url{https://oeis.org/A082397}. Offset 1.

\textbf{Conclusion.} The selected printed recurrence holds for every integer \(n\geq4\). Its complete coefficient list is:
\[
\begin{aligned}
p_{0}(n)&=n + 2
\end{aligned}
\]
\[
\begin{aligned}
p_{1}(n)&=- 3 n - 2
\end{aligned}
\]
\[
\begin{aligned}
p_{2}(n)&=- n
\end{aligned}
\]
\[
\begin{aligned}
p_{3}(n)&=3 n
\end{aligned}
\]
The alternating binomial-sum derivation above identifies the GF. Only the recurrence is targeted; the separate partial-sum claim involving A342912 remains outside this packet.
The ordinary generating function has certificate
\[
F=\sum_{i=0}^{1}R_i\sqrt{D_i}.
\]
\[
x=x.
\]
\[
R_{0}=\frac{2 x - 1}{2 x^{2} \left(x - 1\right)}.
\]
\[
D_{0}=1.
\]
\[
R_{1}=\frac{1}{2 x^{2} \left(x - 1\right) \left(x + 1\right)}.
\]
\[
D_{1}=- \left(x + 1\right) \left(3 x - 1\right).
\]
Applying (R2-G3) gives the exact identity
\[
\mathcal L_pF=3 x.
\]
Every remaining radical coefficient is zero. Equations (R2-G1) prove the stated conclusion.

\subsection{A190160}
Source: \url{https://oeis.org/A190160}. Offset 0.

\textbf{Conclusion.} The selected printed recurrence holds for every integer \(n\geq9\). Its complete coefficient list is:
\[
\begin{aligned}
p_{0}(n)&=n + 2
\end{aligned}
\]
\[
\begin{aligned}
p_{1}(n)&=- 5 n - 6
\end{aligned}
\]
\[
\begin{aligned}
p_{2}(n)&=8 n + 6
\end{aligned}
\]
\[
\begin{aligned}
p_{3}(n)&=- 2 n - 9
\end{aligned}
\]
\[
\begin{aligned}
p_{4}(n)&=25 - 8 n
\end{aligned}
\]
\[
\begin{aligned}
p_{5}(n)&=6 n - 12
\end{aligned}
\]
\[
\begin{aligned}
p_{6}(n)&=9 n - 63
\end{aligned}
\]
\[
\begin{aligned}
p_{7}(n)&=96 - 15 n
\end{aligned}
\]
\[
\begin{aligned}
p_{8}(n)&=7 n - 49
\end{aligned}
\]
\[
\begin{aligned}
p_{9}(n)&=8 - n
\end{aligned}
\]
The existing implicit equation is \(G=1+xG+x^2G(x+(1-x)^2(G-xG-1))/(1-x)\). Multiplication by \(1-x\) gives \(AG^2+BG+C=0\), where \(A=x^2(1-x)^3\), \(B=x^3-x^2(1-x)^2-(1-x)^2\), and \(C=1-x\). The root \(G=(-B-\sqrt{B^2-4AC})/(2A)\), with \(G(0)=1\), is the certificate.
The ordinary generating function has certificate
\[
F=\sum_{i=0}^{1}R_i\sqrt{D_i}.
\]
\[
x=x.
\]
\[
R_{0}=- \frac{x^{4} - 3 x^{3} + 2 x^{2} - 2 x + 1}{2 x^{2} \left(x - 1\right)^{3}}.
\]
\[
D_{0}=1.
\]
\[
R_{1}=\frac{1}{2 x^{2} \left(x - 1\right)^{3}}.
\]
\[
D_{1}=\left(x^{4} - x^{3} - 2 x^{2} + 1\right) \left(x^{4} - 5 x^{3} + 6 x^{2} - 4 x + 1\right).
\]
Applying (R2-G3) gives the exact identity
\[
\mathcal L_pF=- 2 \left(x - 1\right) \left(x^{5} - 3 x^{4} + 4 x^{3} + 2 x^{2} - 3 x + 1\right).
\]
Every remaining radical coefficient is zero. Equations (R2-G1) prove the stated conclusion.

\subsection{A001723}
Source: \url{https://oeis.org/A001723}. Offset 0.

\textbf{Conclusion.} The selected printed recurrence holds for every integer \(n\geq4\). Its complete coefficient list is:
\[
\begin{aligned}
p_{0}(n)&=1
\end{aligned}
\]
\[
\begin{aligned}
p_{1}(n)&=- 4 n - 22
\end{aligned}
\]
\[
\begin{aligned}
p_{2}(n)&=6 n^{2} + 60 n \\
&\quad +151
\end{aligned}
\]
\[
\begin{aligned}
p_{3}(n)&=- 4 n^{3} - 54 n^{2} \\
&\quad - 244 n - 369
\end{aligned}
\]
\[
\begin{aligned}
p_{4}(n)&=n^{4} + 16 n^{3} \\
&\quad +96 n^{2} + 256 n \\
&\quad +256
\end{aligned}
\]
Use the rising-product proof (R2-G4), with \(d=3\), \(r=4\). The finite polynomial identity has residual \(u^{4}\); multiplying by \((5+u)^{\overline{n-1}}\) and extracting \([u^{3}]\) proves the recurrence.

\subsection{A001724}
Source: \url{https://oeis.org/A001724}. Offset 0.

\textbf{Conclusion.} The selected printed recurrence holds for every integer \(n\geq5\). Its complete coefficient list is:
\[
\begin{aligned}
p_{0}(n)&=1
\end{aligned}
\]
\[
\begin{aligned}
p_{1}(n)&=- 5 n - 30
\end{aligned}
\]
\[
\begin{aligned}
p_{2}(n)&=10 n^{2} + 110 n \\
&\quad +305
\end{aligned}
\]
\[
\begin{aligned}
p_{3}(n)&=- 10 n^{3} - 150 n^{2} \\
&\quad - 755 n - 1275
\end{aligned}
\]
\[
\begin{aligned}
p_{4}(n)&=5 n^{4} + 90 n^{3} \\
&\quad +610 n^{2} + 1845 n \\
&\quad +2101
\end{aligned}
\]
\[
\begin{aligned}
p_{5}(n)&=- n^{5} - 20 n^{4} \\
&\quad - 160 n^{3} - 640 n^{2} \\
&\quad - 1280 n - 1024
\end{aligned}
\]
Use the rising-product proof (R2-G4), with \(d=4\), \(r=5\). The finite polynomial identity has residual \(u^{5}\); multiplying by \((5+u)^{\overline{n-1}}\) and extracting \([u^{4}]\) proves the recurrence.

\subsection{A027051}
Source: \url{https://oeis.org/A027051}. Additional established seed sources: \url{https://oeis.org/A027023}. Offset 2.

\textbf{Conclusion.} The selected printed recurrence holds for every integer \(n\geq9\). Its complete coefficient list is:
\[
\begin{aligned}
p_{0}(n)&=n
\end{aligned}
\]
\[
\begin{aligned}
p_{1}(n)&=5 - 7 n
\end{aligned}
\]
\[
\begin{aligned}
p_{2}(n)&=13 n - 18
\end{aligned}
\]
\[
\begin{aligned}
p_{3}(n)&=n - 13
\end{aligned}
\]
\[
\begin{aligned}
p_{4}(n)&=64 - 13 n
\end{aligned}
\]
\[
\begin{aligned}
p_{5}(n)&=3 n - 25
\end{aligned}
\]
\[
\begin{aligned}
p_{6}(n)&=2 - n
\end{aligned}
\]
\[
\begin{aligned}
p_{7}(n)&=3 n - 15
\end{aligned}
\]
The reversed-row kernel derivation above identifies this column GF from its defining triangle.
The ordinary generating function has certificate
\[
F=\sum_{i=0}^{0}R_i\sqrt{D_i}.
\]
\[
x=\frac{t}{t^{2} + t + 1}.
\]
\[
R_{0}=\frac{\mathcal N}{\mathcal D}.
\]
\[
\begin{aligned}
\mathcal N&=- 4 t^{7} - 8 t^{6} \\
&\quad - 11 t^{5} - 9 t^{4} \\
&\quad - 5 t^{3} - t^{2}
\end{aligned}
\]
\[
\begin{aligned}
\mathcal D&=t^{7} - 3 t^{4} \\
&\quad - 3 t^{3} - 4 t^{2} \\
&\quad - 2 t - 1
\end{aligned}
\]
\[
D_{0}=1.
\]
Applying (R2-G3) gives the exact identity
\[
\mathcal L_pF=- x^{2} \left(9 x^{6} + x^{5} - 20 x^{4} - 2 x^{3} + 29 x^{2} + x - 2\right).
\]
Every remaining radical coefficient is zero. Equations (R2-G1) prove the stated conclusion.

\subsection{A027027}
Source: \url{https://oeis.org/A027027}. Additional established seed sources: \url{https://oeis.org/A027023}. Offset 2.

\textbf{Conclusion.} The selected printed recurrence holds for every integer \(n\geq9\). Its complete coefficient list is:
\[
\begin{aligned}
p_{0}(n)&=n + 1
\end{aligned}
\]
\[
\begin{aligned}
p_{1}(n)&=- 8 n - 1
\end{aligned}
\]
\[
\begin{aligned}
p_{2}(n)&=19 n - 14
\end{aligned}
\]
\[
\begin{aligned}
p_{3}(n)&=- 6 n - 2
\end{aligned}
\]
\[
\begin{aligned}
p_{4}(n)&=89 - 21 n
\end{aligned}
\]
\[
\begin{aligned}
p_{5}(n)&=8 n - 45
\end{aligned}
\]
\[
\begin{aligned}
p_{6}(n)&=n - 4
\end{aligned}
\]
\[
\begin{aligned}
p_{7}(n)&=6 n - 24
\end{aligned}
\]
The reversed-row kernel derivation above identifies this column GF from its defining triangle.
The ordinary generating function has certificate
\[
F=\sum_{i=0}^{0}R_i\sqrt{D_i}.
\]
\[
x=\frac{t}{t^{2} + t + 1}.
\]
\[
R_{0}=\frac{\mathcal N}{\mathcal D}.
\]
\[
\begin{aligned}
\mathcal N&=- 4 t^{8} - 8 t^{7} \\
&\quad - 12 t^{6} - 11 t^{5} \\
&\quad - 7 t^{4} - 3 t^{3} \\
&\quad - t^{2}
\end{aligned}
\]
\[
\begin{aligned}
\mathcal D&=t^{7} - 3 t^{4} \\
&\quad - 3 t^{3} - 4 t^{2} \\
&\quad - 2 t - 1
\end{aligned}
\]
\[
D_{0}=1.
\]
Applying (R2-G3) gives the exact identity
\[
\mathcal L_pF=x^{2} \left(24 x^{6} - 15 x^{5} - 35 x^{4} + 4 x^{3} + 8 x^{2} - 13 x + 3\right).
\]
Every remaining radical coefficient is zero. Equations (R2-G1) prove the stated conclusion.

\subsection{A027028}
Source: \url{https://oeis.org/A027028}. Additional established seed sources: \url{https://oeis.org/A027023}. Offset 2.

\textbf{Conclusion.} The selected printed recurrence holds for every integer \(n\geq8\). Its complete coefficient list is:
\[
\begin{aligned}
p_{0}(n)&=- n^{2} + 4 n \\
&\quad +12
\end{aligned}
\]
\[
\begin{aligned}
p_{1}(n)&=6 n^{2} - 27 n \\
&\quad -33
\end{aligned}
\]
\[
\begin{aligned}
p_{2}(n)&=- 7 n^{2} + 31 n
\end{aligned}
\]
\[
\begin{aligned}
p_{3}(n)&=- 8 n^{2} + 62 n \\
&\quad -114
\end{aligned}
\]
\[
\begin{aligned}
p_{4}(n)&=5 n^{2} - 38 n \\
&\quad +84
\end{aligned}
\]
\[
\begin{aligned}
p_{5}(n)&=2 n^{2} - 11 n \\
&\quad +15
\end{aligned}
\]
\[
\begin{aligned}
p_{6}(n)&=3 n^{2} - 21 n \\
&\quad +36
\end{aligned}
\]
The reversed-row kernel derivation above identifies this column GF from its defining triangle.
The ordinary generating function has certificate
\[
F=\sum_{i=0}^{0}R_i\sqrt{D_i}.
\]
\[
x=\frac{t}{t^{2} + t + 1}.
\]
\[
R_{0}=\frac{\mathcal N}{\mathcal D}.
\]
\[
\begin{aligned}
\mathcal N&=- 4 t^{9} - 8 t^{8} \\
&\quad - 12 t^{7} - 12 t^{6} \\
&\quad - 7 t^{5} - 5 t^{4} \\
&\quad - t^{3} - t^{2}
\end{aligned}
\]
\[
\begin{aligned}
\mathcal D&=t^{7} - 3 t^{4} \\
&\quad - 3 t^{3} - 4 t^{2} \\
&\quad - 2 t - 1
\end{aligned}
\]
\[
D_{0}=1.
\]
Applying (R2-G3) gives the exact identity
\[
\mathcal L_pF=- x^{2} \left(x - 1\right) \left(36 x^{4} + 3 x^{3} - 2 x^{2} - 29 x + 16\right).
\]
Every remaining radical coefficient is zero. Equations (R2-G1) prove the stated conclusion.

\subsection{A027056}
Source: \url{https://oeis.org/A027056}. Additional established seed sources: \url{https://oeis.org/A027052}. Offset 1.

\textbf{Conclusion.} The selected printed recurrence holds for every integer \(n\geq7\). Its complete coefficient list is:
\[
\begin{aligned}
p_{0}(n)&=1 - n
\end{aligned}
\]
\[
\begin{aligned}
p_{1}(n)&=6 n - 8
\end{aligned}
\]
\[
\begin{aligned}
p_{2}(n)&=3 - 7 n
\end{aligned}
\]
\[
\begin{aligned}
p_{3}(n)&=52 - 8 n
\end{aligned}
\]
\[
\begin{aligned}
p_{4}(n)&=5 n - 29
\end{aligned}
\]
\[
\begin{aligned}
p_{5}(n)&=2 n - 4
\end{aligned}
\]
\[
\begin{aligned}
p_{6}(n)&=3 n - 15
\end{aligned}
\]
The reversed-row kernel derivation above identifies this column GF from its defining triangle.
The ordinary generating function has certificate
\[
F=\sum_{i=0}^{0}R_i\sqrt{D_i}.
\]
\[
x=\frac{t}{t^{2} + t + 1}.
\]
\[
R_{0}=- \frac{2 t^{2} \left(t^{2} + t + 1\right)}{\left(t^{2} + 1\right) \left(t^{3} - t^{2} - t - 1\right)}.
\]
\[
D_{0}=1.
\]
Applying (R2-G3) gives the exact identity
\[
\mathcal L_pF=- 2 x^{2} \left(x - 1\right) \left(5 x - 1\right).
\]
Every remaining radical coefficient is zero. Equations (R2-G1) prove the stated conclusion.

\subsection{A027087}
Source: \url{https://oeis.org/A027087}. Additional established seed sources: \url{https://oeis.org/A027082}. Offset 0.

\textbf{Conclusion.} The selected printed recurrence holds for every integer \(n\geq7\). Its complete coefficient list is:
\[
\begin{aligned}
p_{0}(n)&=3 n
\end{aligned}
\]
\[
\begin{aligned}
p_{1}(n)&=9 - 15 n
\end{aligned}
\]
\[
\begin{aligned}
p_{2}(n)&=7 n - 4
\end{aligned}
\]
\[
\begin{aligned}
p_{3}(n)&=25 n - 67
\end{aligned}
\]
\[
\begin{aligned}
p_{4}(n)&=17 n - 44
\end{aligned}
\]
\[
\begin{aligned}
p_{5}(n)&=19 n - 97
\end{aligned}
\]
\[
\begin{aligned}
p_{6}(n)&=5 n - 32
\end{aligned}
\]
\[
\begin{aligned}
p_{7}(n)&=3 n - 21
\end{aligned}
\]
The reversed-row kernel derivation above identifies this column GF from its defining triangle.
The ordinary generating function has certificate
\[
F=\sum_{i=0}^{0}R_i\sqrt{D_i}.
\]
\[
x=\frac{t}{t^{2} + t + 1}.
\]
\[
R_{0}=- \frac{\left(t + 1\right) \left(2 t^{2} + t + 1\right)}{t^{3} - t^{2} - t - 1}.
\]
\[
D_{0}=1.
\]
Applying (R2-G3) gives the exact identity
\[
\mathcal L_pF=x \left(x + 1\right)^{2} \left(x^{3} - 13 x^{2} + 7 x - 3\right).
\]
Every remaining radical coefficient is zero. Equations (R2-G1) prove the stated conclusion.

\subsection{A025258}
Source: \url{https://oeis.org/A025258}. Offset 1.

\textbf{Conclusion.} The selected printed recurrence holds for every integer \(n\geq6\). Its complete coefficient list is:
\[
\begin{aligned}
p_{0}(n)&=- n
\end{aligned}
\]
\[
\begin{aligned}
p_{1}(n)&=2 n - 3
\end{aligned}
\]
\[
\begin{aligned}
p_{2}(n)&=9 - 3 n
\end{aligned}
\]
\[
\begin{aligned}
p_{3}(n)&=6 n - 27
\end{aligned}
\]
\[
\begin{aligned}
p_{4}(n)&=6 - n
\end{aligned}
\]
The convolution argument above derives the radical from the definition; the same expression is an established Manyama formula on the page.
The ordinary generating function has certificate
\[
F=\sum_{i=0}^{1}R_i\sqrt{D_i}.
\]
\[
x=x.
\]
\[
R_{0}=- \frac{x^{2} - x - 1}{2}.
\]
\[
D_{0}=1.
\]
\[
R_{1}=- \frac{1}{2}.
\]
\[
D_{1}=x^{4} - 6 x^{3} + 3 x^{2} - 2 x + 1.
\]
Applying (R2-G3) gives the exact identity
\[
\mathcal L_pF=- x \left(x^{4} - x^{3} + 6 x^{2} - 3 x + 1\right).
\]
Every remaining radical coefficient is zero. Equations (R2-G1) prove the stated conclusion.

\subsection{A157125}
Source: \url{https://oeis.org/A157125}. Offset 0.

\textbf{Conclusion.} The selected printed recurrence holds for every integer \(n\geq5\). Its complete coefficient list is:
\[
\begin{aligned}
p_{0}(n)&=n^{2} + n \\
&\quad -6
\end{aligned}
\]
\[
\begin{aligned}
p_{1}(n)&=-4
\end{aligned}
\]
\[
\begin{aligned}
p_{2}(n)&=2 n^{2} - 2 n \\
&\quad -8
\end{aligned}
\]
\[
\begin{aligned}
p_{3}(n)&=4 n^{2} - 14 n \\
&\quad +4
\end{aligned}
\]
\[
\begin{aligned}
p_{4}(n)&=n^{2} - 3 n \\
&\quad -4
\end{aligned}
\]
\[
\begin{aligned}
p_{5}(n)&=4 n^{2} - 18 n \\
&\quad +14
\end{aligned}
\]
The accepted GF is \((1-x)(\sqrt{1+x^2+4x^3}-\sqrt{1+x^2})/(2x^3\sqrt{1+x^2})\). Multiplying the quotient of radicals by \(\sqrt{1+x^2}\) gives the single-radical certificate.
The ordinary generating function has certificate
\[
F=\sum_{i=0}^{1}R_i\sqrt{D_i}.
\]
\[
x=x.
\]
\[
R_{0}=\frac{x - 1}{2 x^{3}}.
\]
\[
D_{0}=1.
\]
\[
R_{1}=- \frac{x - 1}{2 x^{3} \left(x^{2} + 1\right)}.
\]
\[
D_{1}=\left(x^{2} + 1\right) \left(4 x^{3} + x^{2} + 1\right).
\]
Applying (R2-G3) gives the exact identity
\[
\mathcal L_pF=- 2 \left(x^{3} + 3\right).
\]
Every remaining radical coefficient is zero. Equations (R2-G1) prove the stated conclusion.

\subsection{A186335}
Source: \url{https://oeis.org/A186335}. Offset 0.

\textbf{Conclusion.} The selected printed recurrence holds for every integer \(n\geq4\). Its complete coefficient list is:
\[
\begin{aligned}
p_{0}(n)&=n
\end{aligned}
\]
\[
\begin{aligned}
p_{1}(n)&=1 - 2 n
\end{aligned}
\]
\[
\begin{aligned}
p_{2}(n)&=5 - 5 n
\end{aligned}
\]
\[
\begin{aligned}
p_{3}(n)&=6 n - 9
\end{aligned}
\]
\[
\begin{aligned}
p_{4}(n)&=5 n - 10
\end{aligned}
\]
The triple reindexing and binomial-series summation above derive this GF directly from the established double sum.
The ordinary generating function has certificate
\[
F=\sum_{i=0}^{1}R_i\sqrt{D_i}.
\]
\[
x=x.
\]
\[
R_{0}=0.
\]
\[
D_{0}=1.
\]
\[
R_{1}=\frac{1}{\left(x^{2} + x - 1\right) \left(5 x^{2} + x - 1\right)}.
\]
\[
D_{1}=\left(x^{2} + x - 1\right) \left(5 x^{2} + x - 1\right).
\]
Applying (R2-G3) gives the exact identity
\[
\mathcal L_pF=0.
\]
Every remaining radical coefficient is zero. Equations (R2-G1) prove the stated conclusion.

\subsection{A133158}
Source: \url{https://oeis.org/A133158}. Additional established seed sources: \url{https://oeis.org/A026641}. Offset 0.

\textbf{Conclusion.} The selected printed recurrence holds for every integer \(n\geq3\). Its complete coefficient list is:
\[
\begin{aligned}
p_{0}(n)&=2 n
\end{aligned}
\]
\[
\begin{aligned}
p_{1}(n)&=12 - 19 n
\end{aligned}
\]
\[
\begin{aligned}
p_{2}(n)&=48 n - 66
\end{aligned}
\]
\[
\begin{aligned}
p_{3}(n)&=72 - 36 n
\end{aligned}
\]
The second binomial-transform calculation above derives the function from the established GF of A026641.
The ordinary generating function has certificate
\[
F=\sum_{i=0}^{1}R_i\sqrt{D_i}.
\]
\[
x=x.
\]
\[
R_{0}=- \frac{1}{2 \left(3 x - 2\right)}.
\]
\[
D_{0}=1.
\]
\[
R_{1}=\frac{3}{2 \left(3 x - 2\right) \left(6 x - 1\right)}.
\]
\[
D_{1}=\left(2 x - 1\right) \left(6 x - 1\right).
\]
Applying (R2-G3) gives the exact identity
\[
\mathcal L_pF=- x.
\]
Every remaining radical coefficient is zero. Equations (R2-G1) prove the stated conclusion.

\subsection{A157021}
Source: \url{https://oeis.org/A157021}. Offset 0.

\textbf{Conclusion.} The selected printed recurrence holds for every integer \(n\geq5\). Its complete coefficient list is:
\[
\begin{aligned}
p_{0}(n)&=n^{3} - 7 n^{2} \\
&\quad +128 - 12 n
\end{aligned}
\]
\[
\begin{aligned}
p_{1}(n)&=- 5 n^{3} + 41 n^{2} \\
&\quad - 26 n - 360
\end{aligned}
\]
\[
\begin{aligned}
p_{2}(n)&=10 n^{3} - 94 n^{2} \\
&\quad +224 n + 128
\end{aligned}
\]
\[
\begin{aligned}
p_{3}(n)&=- 10 n^{3} + 106 n^{2} \\
&\quad +336 - 364 n
\end{aligned}
\]
\[
\begin{aligned}
p_{4}(n)&=n^{3} - 7 n^{2} \\
&\quad +4 n + 32
\end{aligned}
\]
\[
\begin{aligned}
p_{5}(n)&=3 n^{3} - 39 n^{2} \\
&\quad +174 n - 264
\end{aligned}
\]
The target is the second, shorter conjectural recurrence on the page (the one with five backward shifts). Its seed GF follows from the defining binomial sum as shown above.
The ordinary generating function has certificate
\[
F=\sum_{i=0}^{1}R_i\sqrt{D_i}.
\]
\[
x=x.
\]
\[
R_{0}=- \frac{\left(x - 1\right) \left(2 x^{2} - 2 x + 1\right)}{2 x^{4}}.
\]
\[
D_{0}=1.
\]
\[
R_{1}=\frac{2 x^{2} - 2 x + 1}{2 x^{4} \left(x - 1\right)}.
\]
\[
D_{1}=- \left(x^{2} + 2 x - 1\right) \left(3 x^{2} - 2 x + 1\right).
\]
Applying (R2-G3) gives the exact identity
\[
\mathcal L_pF=- 16 \left(2 x^{3} - 10 x^{2} + 15 x - 8\right).
\]
Every remaining radical coefficient is zero. Equations (R2-G1) prove the stated conclusion.

\subsection{A173993}
Source: \url{https://oeis.org/A173993}. Offset 0.

\textbf{Conclusion.} The selected printed recurrence holds for every integer \(n\geq6\). Its complete coefficient list is:
\[
\begin{aligned}
p_{0}(n)&=n + 1
\end{aligned}
\]
\[
\begin{aligned}
p_{1}(n)&=- 8 n - 2
\end{aligned}
\]
\[
\begin{aligned}
p_{2}(n)&=19 n - 5
\end{aligned}
\]
\[
\begin{aligned}
p_{3}(n)&=-36
\end{aligned}
\]
\[
\begin{aligned}
p_{4}(n)&=208 - 56 n
\end{aligned}
\]
\[
\begin{aligned}
p_{5}(n)&=68 n - 286
\end{aligned}
\]
\[
\begin{aligned}
p_{6}(n)&=120 - 24 n
\end{aligned}
\]
The seed is the existing GF \((\sqrt{(1-2x)(1-2x-4x^2+4x^3)}-2x^2+4x-1)/(2x(1-4x+3x^2))\). The target is its recurrence only.
The ordinary generating function has certificate
\[
F=\sum_{i=0}^{1}R_i\sqrt{D_i}.
\]
\[
x=x.
\]
\[
R_{0}=- \frac{2 x^{2} - 4 x + 1}{2 x \left(x - 1\right) \left(3 x - 1\right)}.
\]
\[
D_{0}=1.
\]
\[
R_{1}=\frac{1}{2 x \left(x - 1\right) \left(3 x - 1\right)}.
\]
\[
D_{1}=- \left(2 x - 1\right) \left(4 x^{3} - 4 x^{2} - 2 x + 1\right).
\]
Applying (R2-G3) gives the exact identity
\[
\mathcal L_pF=\left(x - 1\right) \left(10 x^{3} - 10 x^{2} + 5 x - 1\right).
\]
Every remaining radical coefficient is zero. Equations (R2-G1) prove the stated conclusion.

\subsection{A141353}
Source: \url{https://oeis.org/A141353}. Offset 0.

\textbf{Conclusion.} The selected printed recurrence holds for every integer \(n\geq4\). Its complete coefficient list is:
\[
\begin{aligned}
p_{0}(n)&=n + 1
\end{aligned}
\]
\[
\begin{aligned}
p_{1}(n)&=2 - 8 n
\end{aligned}
\]
\[
\begin{aligned}
p_{2}(n)&=20 n - 28
\end{aligned}
\]
\[
\begin{aligned}
p_{3}(n)&=40 - 16 n
\end{aligned}
\]
The definition gives \(F=C(x)+2x/(1-2x)\), with \(C\) the Catalan series and the constant term fixed at 1.
The ordinary generating function has certificate
\[
F=\sum_{i=0}^{1}R_i\sqrt{D_i}.
\]
\[
x=x.
\]
\[
R_{0}=- \frac{4 x^{2} - 2 x + 1}{2 x \left(2 x - 1\right)}.
\]
\[
D_{0}=1.
\]
\[
R_{1}=- \frac{1}{2 x}.
\]
\[
D_{1}=1 - 4 x.
\]
Applying (R2-G3) gives the exact identity
\[
\mathcal L_pF=8 x^{3} - 12 x^{2} + 1.
\]
Every remaining radical coefficient is zero. Equations (R2-G1) prove the stated conclusion.

\subsection{A078999}
Source: \url{https://oeis.org/A078999}. Additional established seed sources: \url{https://oeis.org/A078995}, \url{https://oeis.org/A002293}. Offset 0.

\textbf{Conclusion.} The selected printed recurrence holds for every integer \(n\geq3\). Its complete coefficient list is:
\[
\begin{aligned}
p_{0}(n)&=- 2937742094763 n^{4} - 6138004010652 n^{3} \\
&\quad - 461143675071 n^{2} + 5117093595162 n \\
&\quad +2377975354344
\end{aligned}
\]
\[
\begin{aligned}
p_{1}(n)&=39322791505152 n^{4} + 130827262124736 n^{3} \\
&\quad - 77390801861328 n^{2} - 57276849171240 n \\
&\quad -8278877159568
\end{aligned}
\]
\[
\begin{aligned}
p_{2}(n)&=46619105624064 n^{4} - 1454666767368192 n^{3} \\
&\quad +2115306373226496 n^{2} - 874201548149760 n \\
&\quad +84277928590848
\end{aligned}
\]
\[
\begin{aligned}
p_{3}(n)&=- 1473031998275584 n^{4} + 7263078429229056 n^{3} \\
&\quad - 12705856454918144 n^{2} + 9076166541705216 n \\
&\quad -2081737996369920
\end{aligned}
\]
The source cross-sequence formula and the two established GFs give the rational substitution derived above, about \(t=1\).
The ordinary generating function has certificate
\[
F=\sum_{i=0}^{0}R_i\sqrt{D_i}.
\]
\[
x=\frac{t - 1}{t^{4}}.
\]
\[
R_{0}=\frac{t^{5} \left(3 t - 2\right)}{\left(3 t - 4\right)^{2}}.
\]
\[
D_{0}=1.
\]
Applying (R2-G3) gives the exact identity
\[
\mathcal L_pF=- 216 \left(6397991648 x - 11009145159\right).
\]
Every remaining radical coefficient is zero. Equations (R2-G1) prove the stated conclusion.

\subsection{A000159}
Source: \url{https://oeis.org/A000159}. Additional established seed sources: \url{https://oeis.org/A058087}. Offset 3.

\textbf{Conclusion.} The selected printed recurrence holds for every integer \(n\geq9\). Its complete coefficient list is:
\[
\begin{aligned}
p_{0}(n)&=2082154 - 504614 n
\end{aligned}
\]
\[
\begin{aligned}
p_{1}(n)&=504614 n^{2} - 3362985 n \\
&\quad +5118150
\end{aligned}
\]
\[
\begin{aligned}
p_{2}(n)&=1280831 n^{2} - 7397886 n \\
&\quad +6461565
\end{aligned}
\]
\[
\begin{aligned}
p_{3}(n)&=746598 n^{2} - 2913543 n \\
&\quad -1336090
\end{aligned}
\]
\[
\begin{aligned}
p_{4}(n)&=- 405481 n^{2} + 6175011 n \\
&\quad -15469320
\end{aligned}
\]
\[
\begin{aligned}
p_{5}(n)&=- 375862 n^{2} + 4098537 n \\
&\quad -8846430
\end{aligned}
\]
\[
\begin{aligned}
p_{6}(n)&=1121260 - 375862 n
\end{aligned}
\]
Use the fixed-diagonal construction above with \(d=3\). In its notation,
\[
F=A F_0+B,\qquad F_0=f\left(\frac{x}{(1+x)^2}\right).
\]
\[
A=- \frac{2 \left(2 x^{5} - 3 x^{4} - 5 x^{3} - 4 x^{2} - 3 x - 2\right)}{3 \left(x + 1\right)^{2}}.
\]
\[
B=\frac{2 \left(2 x^{5} - x^{4} - 8 x^{3} - 7 x^{2} - 5 x - 2\right)}{3 \left(x + 1\right)^{2}}.
\]
The coefficient of \(F_0\) in \(\mathcal L_pF\) is zero, while its rational part is
\[
\mathcal L_pF=4 x^{3} \left(379025 x^{2} - 2587 x + 284156\right).
\]
The factorial-series differential equation therefore proves the recurrence.

\subsection{A000181}
Source: \url{https://oeis.org/A000181}. Additional established seed sources: \url{https://oeis.org/A058087}. Offset 4.

\textbf{Conclusion.} The selected printed recurrence holds for every integer \(n\geq8\). Its complete coefficient list is:
\[
\begin{aligned}
p_{0}(n)&=25332 n^{4} - 428952 n^{3} \\
&\quad +2720193 n^{2} - 7656417 n \\
&\quad +8070516
\end{aligned}
\]
\[
\begin{aligned}
p_{1}(n)&=- 25332 n^{5} + 453232 n^{4} \\
&\quad - 3030389 n^{3} + 9016781 n^{2} \\
&\quad +1452000 - 10446214 n
\end{aligned}
\]
\[
\begin{aligned}
p_{2}(n)&=- 24280 n^{5} + 333424 n^{4} \\
&\quad - 1619342 n^{3} + 3451376 n^{2} \\
&\quad +1805760 - 3470928 n
\end{aligned}
\]
\[
\begin{aligned}
p_{3}(n)&=1052 n^{5} - 40656 n^{4} \\
&\quad +266063 n^{3} - 549153 n^{2} \\
&\quad +49850 n + 655200
\end{aligned}
\]
\[
\begin{aligned}
p_{4}(n)&=1052 n^{4} - 7960 n^{3} \\
&\quad +9379 n^{2} + 46459 n \\
&\quad -94080
\end{aligned}
\]
Use the fixed-diagonal construction above with \(d=4\). In its notation,
\[
F=A F_0+B,\qquad F_0=f\left(\frac{x}{(1+x)^2}\right).
\]
\[
A=- \frac{4 x^{6} - 12 x^{5} - 13 x^{4} - 6 x^{3} - 3 x^{2} - 4 x - 4}{6 \left(x + 1\right)^{2}}.
\]
\[
B=\frac{4 x^{7} - 4 x^{6} - 33 x^{5} - 40 x^{4} - 20 x^{3} - 15 x^{2} - 12 x - 4}{6 \left(x + 1\right)^{3}}.
\]
The coefficient of \(F_0\) in \(\mathcal L_pF\) is zero, while its rational part is
\[
\mathcal L_pF=0.
\]
The factorial-series differential equation therefore proves the recurrence.

\subsection{A000185}
Source: \url{https://oeis.org/A000185}. Additional established seed sources: \url{https://oeis.org/A058087}. Offset 5.

\textbf{Conclusion.} The selected printed recurrence holds for every integer \(n\geq9\). Its complete coefficient list is:
\[
\begin{aligned}
p_{0}(n)&=482 n^{4} - 9603 n^{3} \\
&\quad +71606 n^{2} - 236835 n \\
&\quad +293150
\end{aligned}
\]
\[
\begin{aligned}
p_{1}(n)&=- 482 n^{5} + 10099 n^{4} \\
&\quad - 79756 n^{3} + 285961 n^{2} \\
&\quad +149292 - 426904 n
\end{aligned}
\]
\[
\begin{aligned}
p_{2}(n)&=- 496 n^{5} + 8178 n^{4} \\
&\quad - 50261 n^{3} + 150426 n^{2} \\
&\quad +192618 - 243765 n
\end{aligned}
\]
\[
\begin{aligned}
p_{3}(n)&=- 14 n^{5} - 49 n^{4} \\
&\quad - 619 n^{3} + 13174 n^{2} \\
&\quad +61248 - 51690 n
\end{aligned}
\]
\[
\begin{aligned}
p_{4}(n)&=- 14 n^{4} + 321 n^{3} \\
&\quad - 2338 n^{2} + 6939 n \\
&\quad -7308
\end{aligned}
\]
Use the fixed-diagonal construction above with \(d=5\). In its notation,
\[
F=A F_0+B,\qquad F_0=f\left(\frac{x}{(1+x)^2}\right).
\]
\[
A=- \frac{4 x^{7} - 20 x^{6} - 9 x^{5} + 10 x^{4} + 15 x^{3} + 9 x^{2} - 4}{15 \left(x + 1\right)^{2}}.
\]
\[
B=\frac{4 x^{9} - 8 x^{8} - 57 x^{7} - 69 x^{6} + 2 x^{5} + 58 x^{4} + 34 x^{3} - 3 x^{2} - 12 x - 4}{15 \left(x + 1\right)^{4}}.
\]
The coefficient of \(F_0\) in \(\mathcal L_pF\) is zero, while its rational part is
\[
\mathcal L_pF=0.
\]
The factorial-series differential equation therefore proves the recurrence.

\subsection{A080894}
Source: \url{https://oeis.org/A080894}. Offset 0.

\textbf{Conclusion.} The selected printed recurrence holds for every integer \(n\geq3\). Its complete coefficient list is:
\[
\begin{aligned}
p_{0}(n)&=3 - 2 n
\end{aligned}
\]
\[
\begin{aligned}
p_{1}(n)&=- 2 n^{3} + 9 n^{2} \\
&\quad +1 - 9 n
\end{aligned}
\]
\[
\begin{aligned}
p_{2}(n)&=4 n^{4} - 14 n^{3} \\
&\quad +11 n^{2} + 5 n \\
&\quad -6
\end{aligned}
\]
\[
\begin{aligned}
p_{3}(n)&=6 n^{5} - 51 n^{4} \\
&\quad +162 n^{3} - 237 n^{2} \\
&\quad +156 n - 36
\end{aligned}
\]
The defining or established EGF is
\[
F=\frac{e^{\frac{1}{2 x}} e^{- \frac{\sqrt{- 3 x^{2} - 2 x + 1}}{2 x}}}{e^{\frac{1}{2}}}.
\]
For the reduction (R2-G7), its exact logarithmic derivative is \(F'/F=\ell_0+\ell_1\sqrt D\), with
\[
D=- \left(x + 1\right) \left(3 x - 1\right).
\]
\[
\ell_0=- \frac{1}{2 x^{2}}.
\]
\[
\ell_1=\frac{x - 1}{2 x^{2} \left(x + 1\right) \left(3 x - 1\right)}.
\]
The adjusted operator has zero residual in both components: \(\mathcal L_qF=F(0+0\sqrt D)=0\). Thus (R2-G2) proves the recurrence.

\subsection{A272261}
Source: \url{https://oeis.org/A272261}. Offset 0.

\textbf{Conclusion.} The selected printed recurrence holds for every integer \(n\geq2\). Its complete coefficient list is:
\[
\begin{aligned}
p_{0}(n)&=1
\end{aligned}
\]
\[
\begin{aligned}
p_{1}(n)&=2 - 4 n
\end{aligned}
\]
\[
\begin{aligned}
p_{2}(n)&=-4
\end{aligned}
\]
The defining or established EGF is
\[
F=\frac{e^{\sqrt{1 - 4 x}}}{e \sqrt{1 - 4 x}}.
\]
For the reduction (R2-G7), its exact logarithmic derivative is \(F'/F=\ell_0+\ell_1\sqrt D\), with
\[
D=1 - 4 x.
\]
\[
\ell_0=- \frac{2}{4 x - 1}.
\]
\[
\ell_1=\frac{2}{4 x - 1}.
\]
The adjusted operator has zero residual in both components: \(\mathcal L_qF=F(0+0\sqrt D)=0\). Thus (R2-G2) proves the recurrence.

\subsection{A144575}
Source: \url{https://oeis.org/A144575}. Offset 0.

\textbf{Conclusion.} The selected printed recurrence holds for every integer \(n\geq4\). Its complete coefficient list is:
\[
\begin{aligned}
p_{0}(n)&=1
\end{aligned}
\]
\[
\begin{aligned}
p_{1}(n)&=3 - 2 n
\end{aligned}
\]
\[
\begin{aligned}
p_{2}(n)&=- 3 n^{2} + 12 n \\
&\quad -10
\end{aligned}
\]
\[
\begin{aligned}
p_{3}(n)&=9 - 6 n
\end{aligned}
\]
\[
\begin{aligned}
p_{4}(n)&=- 9 n^{2} + 36 n \\
&\quad -27
\end{aligned}
\]
The defining or established EGF is
\[
F=e e^{- \sqrt{- 3 x^{2} - 2 x + 1}}.
\]
For the reduction (R2-G7), its exact logarithmic derivative is \(F'/F=\ell_0+\ell_1\sqrt D\), with
\[
D=- \left(x + 1\right) \left(3 x - 1\right).
\]
\[
\ell_0=0.
\]
\[
\ell_1=- \frac{3 x + 1}{\left(x + 1\right) \left(3 x - 1\right)}.
\]
The adjusted operator has zero residual in both components: \(\mathcal L_qF=F(0+0\sqrt D)=0\). Thus (R2-G2) proves the recurrence.

\subsection{A053532}
Source: \url{https://oeis.org/A053532}. Offset 0.

\textbf{Conclusion.} The selected printed recurrence holds for every integer \(n\geq4\). Its complete coefficient list is:
\[
\begin{aligned}
p_{0}(n)&=2
\end{aligned}
\]
\[
\begin{aligned}
p_{1}(n)&=2 - 2 n
\end{aligned}
\]
\[
\begin{aligned}
p_{2}(n)&=0
\end{aligned}
\]
\[
\begin{aligned}
p_{3}(n)&=0
\end{aligned}
\]
\[
\begin{aligned}
p_{4}(n)&=- n^{3} + 6 n^{2} \\
&\quad +6 - 11 n
\end{aligned}
\]
The defining or established EGF is
\[
F=\frac{e^{- \frac{x}{2}} e^{- \frac{x^{2}}{4}} e^{- \frac{x^{3}}{6}}}{\sqrt{1 - x}}.
\]
For the reduction (R2-G7), its exact logarithmic derivative is \(F'/F=\ell_0+\ell_1\sqrt D\), with
\[
D=1.
\]
\[
\ell_0=- \frac{x^{3}}{2 \left(x - 1\right)}.
\]
\[
\ell_1=0.
\]
The adjusted operator has zero residual in both components: \(\mathcal L_qF=F(0+0\sqrt D)=0\). Thus (R2-G2) proves the recurrence.

\subsection{A250917}
Source: \url{https://oeis.org/A250917}. Offset 0.

\textbf{Conclusion.} The selected printed recurrence holds for every integer \(n\geq5\). Its complete coefficient list is:
\[
\begin{aligned}
p_{0}(n)&=2
\end{aligned}
\]
\[
\begin{aligned}
p_{1}(n)&=20 - 11 n
\end{aligned}
\]
\[
\begin{aligned}
p_{2}(n)&=n^{3} + 9 n^{2} \\
&\quad +164 - 116 n
\end{aligned}
\]
\[
\begin{aligned}
p_{3}(n)&=- 4 n^{4} + 35 n^{3} \\
&\quad +n^{2} - 317 n \\
&\quad +342
\end{aligned}
\]
\[
\begin{aligned}
p_{4}(n)&=- 36 n^{4} + 408 n^{3} \\
&\quad - 1782 n^{2} + 3702 n \\
&\quad -3168
\end{aligned}
\]
\[
\begin{aligned}
p_{5}(n)&=24 n^{4} - 396 n^{3} \\
&\quad +2424 n^{2} - 6516 n \\
&\quad +6480
\end{aligned}
\]
The defining or established EGF is
\[
F=e^{\frac{1}{2 x^{2}}} e^{- \frac{3}{2 x}} e^{- \frac{\sqrt{1 - 4 x}}{2 x^{2}}} e^{\frac{\sqrt{1 - 4 x}}{2 x}}.
\]
For the reduction (R2-G7), its exact logarithmic derivative is \(F'/F=\ell_0+\ell_1\sqrt D\), with
\[
D=1 - 4 x.
\]
\[
\ell_0=\frac{3 x - 2}{2 x^{3}}.
\]
\[
\ell_1=- \frac{2 x^{2} - 7 x + 2}{2 x^{3} \left(4 x - 1\right)}.
\]
The adjusted operator has zero residual in both components: \(\mathcal L_qF=F(0+0\sqrt D)=0\). Thus (R2-G2) proves the recurrence.

\subsection{A292461}
Source: \url{https://oeis.org/A292461}. Offset 0.

\textbf{Conclusion.} The selected printed recurrence holds for every integer \(n\geq6\). Its complete coefficient list is:
\[
\begin{aligned}
p_{0}(n)&=n
\end{aligned}
\]
\[
\begin{aligned}
p_{1}(n)&=3 - 2 n
\end{aligned}
\]
\[
\begin{aligned}
p_{2}(n)&=3 - n
\end{aligned}
\]
\[
\begin{aligned}
p_{3}(n)&=9 - 2 n
\end{aligned}
\]
\[
\begin{aligned}
p_{4}(n)&=n - 6
\end{aligned}
\]
The seed is exactly the ordinary GF in the entry definition.
The ordinary generating function has certificate
\[
F=\sum_{i=0}^{1}R_i\sqrt{D_i}.
\]
\[
x=x.
\]
\[
R_{0}=- \frac{x^{2} + x - 1}{2}.
\]
\[
D_{0}=1.
\]
\[
R_{1}=\frac{1}{2}.
\]
\[
D_{1}=\left(x^{2} - 3 x + 1\right) \left(x^{2} + x + 1\right).
\]
Applying (R2-G3) gives the exact identity
\[
\mathcal L_pF=x^{3} \left(x^{2} - x + 3\right).
\]
Every remaining radical coefficient is zero. Equations (R2-G1) prove the stated conclusion.

\subsection{A126089}
Source: \url{https://oeis.org/A126089}. Offset 0.

\textbf{Conclusion.} The selected printed recurrence holds for every integer \(n\geq1\). Its complete coefficient list is:
\[
\begin{aligned}
p_{0}(n)&=4 - n
\end{aligned}
\]
\[
\begin{aligned}
p_{1}(n)&=4 n^{2} - 22 n \\
&\quad +30
\end{aligned}
\]
The seed is exactly the exponential GF in the entry definition. The target is its conjectural first-order recurrence, rather than its other already unqualified recurrence.
The exponential generating function has certificate
\[
F=\sum_{i=0}^{1}R_i\sqrt{D_i}.
\]
\[
x=x.
\]
\[
R_{0}=0.
\]
\[
D_{0}=1.
\]
\[
R_{1}=1 - 2 x.
\]
\[
D_{1}=1 - 4 x.
\]
Applying (R2-G3) gives the exact identity
\[
\mathcal L_qF=0.
\]
Every remaining radical coefficient is zero. Equations (R2-G1) and (R2-G2) prove the stated conclusion.

\subsection{A108629}
Source: \url{https://oeis.org/A108629}. Offset 1.

\textbf{Conclusion.} The selected printed recurrence holds for every integer \(n\geq6\). Its complete coefficient list is:
\[
\begin{aligned}
p_{0}(n)&=n
\end{aligned}
\]
\[
\begin{aligned}
p_{1}(n)&=3 - n
\end{aligned}
\]
\[
\begin{aligned}
p_{2}(n)&=12 - 4 n
\end{aligned}
\]
\[
\begin{aligned}
p_{3}(n)&=30 - 9 n
\end{aligned}
\]
\[
\begin{aligned}
p_{4}(n)&=27 - 7 n
\end{aligned}
\]
\[
\begin{aligned}
p_{5}(n)&=18 - 4 n
\end{aligned}
\]
The seed is the ordinary GF in the entry definition. Its zero constant coefficient agrees with the index-1 offset.
The ordinary generating function has certificate
\[
F=\sum_{i=0}^{1}R_i\sqrt{D_i}.
\]
\[
x=x.
\]
\[
R_{0}=\frac{x + 1}{2 \left(x^{2} + x + 1\right)}.
\]
\[
D_{0}=1.
\]
\[
R_{1}=- \frac{1}{2 \left(x^{2} + x + 1\right)}.
\]
\[
D_{1}=- 4 x^{3} - 3 x^{2} - 2 x + 1.
\]
Applying (R2-G3) gives the exact identity
\[
\mathcal L_pF=x \left(x^{3} + 3 x^{2} + x + 1\right).
\]
Every remaining radical coefficient is zero. Equations (R2-G1) prove the stated conclusion.

\subsection{A239230}
Source: \url{https://oeis.org/A239230}. Offset 0.

\textbf{Conclusion.} The selected printed recurrence holds for every integer \(n\geq8\). Its complete coefficient list is:
\[
\begin{aligned}
p_{0}(n)&=17521226928 n^{4} - 141607619832 n^{3} \\
&\quad +342315933792 n^{2} - 298512866472 n \\
&\quad +80283325584
\end{aligned}
\]
\[
\begin{aligned}
p_{1}(n)&=- 108516679788 n^{4} + 1055343845440 n^{3} \\
&\quad - 3529733053572 n^{2} + 4962562394720 n \\
&\quad -2512676462064
\end{aligned}
\]
\[
\begin{aligned}
p_{2}(n)&=- 57647707646 n^{4} + 974519385068 n^{3} \\
&\quad - 6210429875914 n^{2} + 17628548676196 n \\
&\quad -18287840662872
\end{aligned}
\]
\[
\begin{aligned}
p_{3}(n)&=1104332263320 n^{4} - 16688474493280 n^{3} \\
&\quad +99299523282780 n^{2} - 277054887180164 n \\
&\quad +303328616888592
\end{aligned}
\]
\[
\begin{aligned}
p_{4}(n)&=- 587491214125 n^{4} + 9941738070620 n^{3} \\
&\quad - 75070680472775 n^{2} + 281912285021344 n \\
&\quad -413197788157152
\end{aligned}
\]
\[
\begin{aligned}
p_{5}(n)&=- 2373849695822 n^{4} + 52217689535398 n^{3} \\
&\quad - 423872944767808 n^{2} + 1515169459097264 n \\
&\quad -2019443387466624
\end{aligned}
\]
\[
\begin{aligned}
p_{6}(n)&=2628244359392 n^{4} - 56698793715560 n^{3} \\
&\quad +459016033763548 n^{2} - 1650515137948516 n \\
&\quad +2220852187679664
\end{aligned}
\]
\[
\begin{aligned}
p_{7}(n)&=- 1193212957536 n^{4} + 26793378151488 n^{3} \\
&\quad - 224248504659624 n^{2} + 828694144073232 n \\
&\quad -1140184384447680
\end{aligned}
\]
The conic parametrization above follows from the defining nested radical. The local parameter is \(t=-1\) at \(x=0\); both square roots have their positive constant terms. The target is the recurrence, not the separate double-sum conjecture.
The ordinary generating function has certificate
\[
F=\sum_{i=0}^{0}R_i\sqrt{D_i}.
\]
\[
x=\frac{2 \left(t - 1\right)^{2} \left(t + 1\right)}{\left(t^{2} + 3\right)^{2}}.
\]
\[
R_{0}=\frac{4 t \left(t - 1\right) \left(t + 1\right)}{\left(t + 3\right) \left(t^{2} + 1\right) \left(t^{2} - 4 t - 1\right)}.
\]
\[
D_{0}=1.
\]
Applying (R2-G3) gives the exact identity
\[
\begin{aligned}
\mathcal L_pF&=1138854670092 x^{7} - 801911343684 x^{6} \\
&\quad - 1058136530640 x^{5} - 414143562744 x^{4} \\
&\quad +103948668000 x^{3}
\end{aligned}
\]
Every remaining radical coefficient is zero. Equations (R2-G1) prove the stated conclusion.

\subsection{A032115}
Source: \url{https://oeis.org/A032115}. Offset 0.

\textbf{Conclusion.} The selected printed recurrence holds for every integer \(n\geq6\). Its complete coefficient list is:
\[
\begin{aligned}
p_{0}(n)&=1
\end{aligned}
\]
\[
\begin{aligned}
p_{1}(n)&=- 2 n - 2
\end{aligned}
\]
\[
\begin{aligned}
p_{2}(n)&=4 n - 4
\end{aligned}
\]
\[
\begin{aligned}
p_{3}(n)&=2
\end{aligned}
\]
\[
\begin{aligned}
p_{4}(n)&=11 - 4 n
\end{aligned}
\]
\[
\begin{aligned}
p_{5}(n)&=2 n - 8
\end{aligned}
\]
The established EGF and the componentwise derivative rule are derived above. The exact certificate is
\[
\mathcal L_qF=2x^5.
\]
Its boundary is why the statement begins at 6.

\subsection{A168452}
Source: \url{https://oeis.org/A168452}. Offset 0.

\textbf{Conclusion.} The selected printed recurrence holds for every integer \(n\geq4\). Its complete coefficient list is:
\[
\begin{aligned}
p_{0}(n)&=3 n^{5} + 33 n^{4} \\
&\quad +135 n^{3} + 255 n^{2} \\
&\quad +222 n + 72
\end{aligned}
\]
\[
\begin{aligned}
p_{1}(n)&=- 128 n^{5} - 784 n^{4} \\
&\quad - 1588 n^{3} - 1232 n^{2} \\
&\quad - 300 n
\end{aligned}
\]
\[
\begin{aligned}
p_{2}(n)&=1536 n^{5} + 2304 n^{4} \\
&\quad - 1344 n^{3} - 2688 n^{2} \\
&\quad +288 - 1536 n
\end{aligned}
\]
\[
\begin{aligned}
p_{3}(n)&=- 12288 n^{4} + 33792 n^{3} \\
&\quad +23040 - 49152 n
\end{aligned}
\]
\[
\begin{aligned}
p_{4}(n)&=- 65536 n^{5} + 655360 n^{4} \\
&\quad - 2539520 n^{3} + 4751360 n^{2} \\
&\quad +1474560 - 4276224 n
\end{aligned}
\]
Use the Catalan-product ODE and the six-dimensional calculation (R2-G8) above. Its coefficients in \(G''=AG+BG'+C\) are
\[
A=- \frac{2 \left(6 x - 1\right)}{x^{2} \left(16 x - 1\right)}.
\]
\[
B=- \frac{4 \left(12 x - 1\right)}{x \left(16 x - 1\right)}.
\]
\[
C=- \frac{2}{x^{2} \left(16 x - 1\right)}.
\]
The exact six-component residual is
\[
(0,0,0,0,0,72-1152x).
\]
Hence \(\mathcal L_p(G^2)=72-1152x\), proving the recurrence.

\subsection{A005573}
Source: \url{https://oeis.org/A005573}. Offset 0.

The assertion is the terminating hypergeometric evaluation proved term by term above, for every integer \(n\geq0\). Its proof needs no recurrence certificate.

\subsection{A368633}
Source: \url{https://oeis.org/A368633}. Offset 0.

The assertion is the complete parity characterization proved above, for every integer \(n\geq0\). Its proof uses the defining functional equation over \(\mathbb F_2[[x]]\).

\subsection{A027116}
Source: \url{https://oeis.org/A027116}. Additional established seed sources: \url{https://oeis.org/A027113}. Offset 3.

\textbf{Conclusion.} The selected printed recurrence holds for every integer \(n\geq8\). Its complete coefficient list is:
\[
\begin{aligned}
p_{0}(n)&=1
\end{aligned}
\]
\[
\begin{aligned}
p_{1}(n)&=-3
\end{aligned}
\]
\[
\begin{aligned}
p_{2}(n)&=2
\end{aligned}
\]
\[
\begin{aligned}
p_{3}(n)&=0
\end{aligned}
\]
\[
\begin{aligned}
p_{4}(n)&=1
\end{aligned}
\]
\[
\begin{aligned}
p_{5}(n)&=-1
\end{aligned}
\]
The central-diagonal derivation above proves both the conjectural GF and this recurrence.
The ordinary generating function has certificate
\[
F=\sum_{i=0}^{1}R_i\sqrt{D_i}.
\]
\[
x=x.
\]
\[
R_{0}=\frac{x^{3} \left(2 x^{3} - x^{2} - 2 x - 7\right)}{\left(x - 1\right)^{2} \left(x^{3} + x^{2} + x - 1\right)}.
\]
\[
D_{0}=1.
\]
\[
R_{1}=0.
\]
\[
D_{1}=1.
\]
Applying (R2-G3) gives the exact identity
\[
\mathcal L_pF=- x^{3} \left(2 x^{3} - x^{2} - 2 x - 7\right).
\]
Every remaining radical coefficient is zero. Equations (R2-G1) prove the stated conclusion.

\subsection{A027117}
Source: \url{https://oeis.org/A027117}. Additional established seed sources: \url{https://oeis.org/A027113}. Offset 4.

\textbf{Conclusion.} The selected printed recurrence holds for every integer \(n\geq10\). Its complete coefficient list is:
\[
\begin{aligned}
p_{0}(n)&=1
\end{aligned}
\]
\[
\begin{aligned}
p_{1}(n)&=-4
\end{aligned}
\]
\[
\begin{aligned}
p_{2}(n)&=5
\end{aligned}
\]
\[
\begin{aligned}
p_{3}(n)&=-2
\end{aligned}
\]
\[
\begin{aligned}
p_{4}(n)&=1
\end{aligned}
\]
\[
\begin{aligned}
p_{5}(n)&=-2
\end{aligned}
\]
\[
\begin{aligned}
p_{6}(n)&=1
\end{aligned}
\]
The central-diagonal derivation above proves both the conjectural GF and this recurrence.
The ordinary generating function has certificate
\[
F=\sum_{i=0}^{1}R_i\sqrt{D_i}.
\]
\[
x=x.
\]
\[
R_{0}=\frac{x^{4} \left(3 x^{4} - 5 x^{3} - x^{2} - 5 x + 16\right)}{\left(x - 1\right)^{3} \left(x^{3} + x^{2} + x - 1\right)}.
\]
\[
D_{0}=1.
\]
\[
R_{1}=0.
\]
\[
D_{1}=1.
\]
Applying (R2-G3) gives the exact identity
\[
\mathcal L_pF=x^{4} \left(3 x^{4} - 5 x^{3} - x^{2} - 5 x + 16\right).
\]
Every remaining radical coefficient is zero. Equations (R2-G1) prove the stated conclusion.

\clearpage
\section{Validation, exclusions, and provenance}
\paragraph{Independent mathematical review.}
The full \oeis{A132393} proof was independently checked through its
change of variables, factorial and sign cancellations, moment determinant,
and final symmetric-function identity. The \oeis{A349802}, \oeis{A005773},
and \oeis{A152920} derivations received separate checks of their
normalization and indexing. The \oeis{A153197} contraction proof is sound
under the displayed normalization and remains uncounted.

The connected-array transition rules were independently reviewed against
the whole-grid definition. A second reconstruction checked all 21
rational generating-function identities, the 67-state/627-edge matrix
certificate for \oeis{A166762}, and independent small-grid enumerations.
The arithmetic arguments were reviewed directly, including the
prime-power square-zero conjugacy count and the odd-modulus localization
argument. The certificate engine retains exact rational arithmetic and
the ordinary/exponential normalization explicitly.

\paragraph{Reproducible checks.}
The structural packet includes exact finite checks for all 21 written
arguments, of which 20 are counted. These checks supplement the proofs.
The arithmetic verifier covers all 11 counted results and labels the
additional \oeis{A098275} rediscovery separately. The three floor programs
reconstruct all 28 all-index certificates. The generating-function program
checks all retained operator, product, hypergeometric, and parity
certificates directly from JSON. The verification bundle records the
actual outputs and includes a one-command runner.

\paragraph{Some prior-resolution exclusions.}
The following are examples of candidates removed after reading beyond
the surviving conjecture label:
\begin{itemize}
\item \oeis{A001712}: the later Hadjicostas continuation already proves
the selected recurrence.
\item \oeis{A025273}: the entry credits Israel's 2016 resolution.
\item \oeis{A032123}, \oeis{A348410}, and \oeis{A001711}: linked 2026
proofs already address the selected targets.
\item \oeis{A166761}: Howroyd's 2024 confirmation is recorded in the source.
\item \oeis{A059297}, \oeis{A059259}, and \oeis{A100631}: later source
material already establishes the contemplated claims.
\item \oeis{A037223}: the proposed universal statement already has a
recorded counterexample.
\item \oeis{A098275}: the prior 14 October 2026 paper proves the same
divisibility statement, so our different proof adds no counted target.
\end{itemize}
The \oeis{A065919} exponential/ordinary generating-function normalization
issue was likewise kept outside the count. The same policy applies to
the index and definition discrepancies noted earlier.

\paragraph{Standard references and existing inputs.}
Each proof identifies its entry-specific established formulas. Classical
tools such as Cauchy--Binet, the symmetric-function determinant identity,
Lagrange inversion, the Chinese remainder theorem, and Bertrand's theorem
are used as standard results. The harmonic remainder enclosure uses NIST
DLMF \url{https://dlmf.nist.gov/5.11}. The prime-interval estimate is
Nagura's original 1952 theorem, cited in the relevant proof. The extra
published-results comparison is documented against
\url{https://arxiv.org/abs/2608.11941}.

The resulting dossier is a reviewable mathematical and computational
packet. Public submission, editorial acceptance, exhaustive literature
priority, and proof-assistant formalization are separate steps and are
not claimed here.
% ===== END ORIGINAL ROUND 2 =====

% ===== BEGIN ORIGINAL ROUND 3 =====
\clearpage
\setcounter{section}{0}\setcounter{equation}{0}\setcounter{theorem}{0}
\setcounter{table}{0}\setcounter{figure}{0}
\normalsize\renewcommand{\arraystretch}{1}\renewcommand{\baselinestretch}{1}\selectfont
\part[Round 3: Pell zeros, modular periods, formal congruences, and transfer certificates]{Pell zeros, modular periods, formal congruences, and transfer certificates}
\label{combined:round3}
\rhead{\small Round 3 / KEIO}
\section*{Round overview}
\addcontentsline{toc}{section}{Round 3 overview and scope}
We assemble a final third-round dossier containing 100 selected
OEIS assertion packages. Eighteen are retained and audited from the
supplied provisional manuscript, and 82 are added here. The outcomes
are 95 positive proofs, four refutations, and one entry with both a
refutation and a complete replacement congruence. The arguments include
an explicit infinite Pell family of zero partial sums, prime divisibility
in finite algebraic rings, a general modular-period theorem for integral
Stirling transforms, finite differential quotients, exact formal-series
congruences, arithmetic restrictions at twin-prime centers, and
definition-derived transfer certificates for 41 empirical array
recurrences. Three universal prime-divisibility conjectures are refuted
by complete zero-free residue cycles. A fourth pure-period conjecture
fails modulo 27. Every counted entry is assigned a precise target,
source, proof location, and reproducible check. The word ``additional''
describes this project's provenance and does not assert historical
priority.

\section*{Scope and counting}
\paragraph{What the total means.}
One catalogue row is one specified assertion package associated with a
distinct OEIS identifier. Several assertions in the same entry count once;
an alias, a copied universal assertion, or an uncounted corollary adds
nothing. Different width or color parameters on different source entries
count separately even when one theorem proves them together. Thus the
41 array rows share a transfer method; they do not represent 41 unrelated
mathematical methods. The 100 selections comprise 18 inherited and 82
additional packages. They have no identifier overlap with the previous
208 selections, giving a project total of 308 under this convention.

\paragraph{Relation to the provisional manuscript.}
The inherited part consists of six elementary/formal-series packages and
twelve algebraic-power coefficient packages. Their full arguments are
included below, so this edition is readable without the earlier file.
The provisional manuscript's proposed 100-entry candidate queue is not
used as proof evidence. The three constant-coefficient existence questions
left unresolved there are now refuted, at the primes 71, 241, and 167.
The transferred congruence for \oeis{A396845} remains an uncounted
consequence of the general theorem.

\paragraph{Positive, negative, and mixed outcomes.}
The four refuted selections are \oeis{A188570}, \oeis{A377109},
\oeis{A377113}, and \oeis{A000364}. The mixed entry is
\oeis{A396844}: its proposed mod-16 pattern fails at index 17, and a
complete replacement law is proved. Accordingly, ``100 resolved''
does not mean that 100 proposed statements are true.

\paragraph{The selected scope is part of each result.}
For example, \oeis{A348295} is selected for infinitely many zeros,
not its separately treated unboundedness question.
\oeis{A396803}, \oeis{A397241}, \oeis{A396808}, and \oeis{A397356}
are selected for the stated parity assertions; other congruences or
asymptotic questions on those pages are not included.
\oeis{A378580} is selected for its mod-4 assertion, not the question
about convergence radius. The \oeis{A000364} counterexample concerns
a pure period dividing $\varphi(m)$ from index 1, not eventual
periodicity. Ancillary erroneous small-prime examples in the
algebraic-power entries are explicitly separated from the universal
infinitude and eventual-gap statements proved here.

\paragraph{What is a proof and what is a check.}
All-index conclusions follow from the written algebraic, combinatorial,
or arithmetic arguments, or from a proved reduction to a complete finite
certificate. In the transfer family every state and every transition is
enumerated from the definition, and a complete integer vector identity
proves the recurrence for all subsequent indices. The three zero-free
prime certificates cover an entire proved period. Other finite-prefix
calculations are supplementary checks of transcription, boundaries,
and implementations. No claim of proof-assistant formalization or
external refereeing is made.

\paragraph{Source and prior-work audit.}
Primary entries and relevant linked material were checked on 20 October
2026. The separate audit file records comparison with all 100 costliest
claims in Appendix A of
\href{https://arxiv.org/html/2608.11941v1}{the OEIS Open paper}.
The only identifier overlap is \oeis{A348295}; that paper's result is
unboundedness, whereas the selected assertion here is the infinite
zero set. This is a claim-level distinction, not an assertion that the
identifier is absent from the paper. The audit is limited to the stated
appendix, inspected linked work, and targeted searches, rather than an
exhaustive worldwide literature review.
Known results used as inputs are credited. For example, the Catwalk
binomial enumeration is due to Guy, and the earlier recurrences of
Sela Fried and Robert Israel are not recounted as new results here.
\oeis{A141117} was excluded after locating its linked 2026 proof;
\oeis{A002050} is an uncounted known consequence of Barsky's congruences.
This dossier does not report an OEIS edit, submission, or editorial
acceptance.

\paragraph{Files.}
The standalone source is \nolinkurl{oeis_hunt_round3.tex}.
The portable bundle contains the complete 100-row ledger,
\nolinkurl{oeis_hunt_round3_ledger.json}, all proof fragments,
four verifiers, 90 complete transfer graphs, exact counterexample
certificates, audit notes, and \nolinkurl{run_all.py}.

\clearpage
\section*{Catalogue of the 100 selected packages}
\addcontentsline{toc}{section}{Catalogue of 100 packages}
P = proved; R = refuted; M = mixed proof and refutation.
I = inherited from the supplied manuscript; N = additional this round.
The page column links to the proof. Each identifier links to its primary
OEIS record. Full target qualifications are given in the proof and ledger.
\begingroup\small
\setlength{\tabcolsep}{3pt}
\renewcommand{\arraystretch}{1.13}
\begin{longtable}{@{}r l c c p{0.57\textwidth} r@{}}
\toprule No.&Entry&Result&Origin&Selected target&Page\\\midrule
\endfirsthead
\toprule No.&Entry&Result&Origin&Selected target&Page\\\midrule
\endhead
\midrule\multicolumn{6}{r}{Continued on the next page}\\\endfoot
\bottomrule\endlastfoot
1&\oeis{A348295}&P&I&Infinite Pell family of zero partial sums&\pageref{round3:r3:thm:pell}\\[2pt]
2&\oeis{A014865}&P&I&Oddness, restricted prime factors, and compositeness&\pageref{round3:r3:thm:b15}\\[2pt]
3&\oeis{A048868}&P&I&Prime-majority reduced residues force an even modulus&\pageref{round3:r3:thm:phi}\\[2pt]
4&\oeis{A131418}&P&I&Digit-sum equality forces divisibility by 9&\pageref{round3:r3:thm:digits}\\[2pt]
5&\oeis{A396843}&P&I&Every coefficient is odd; full mod-9 law included&\pageref{round3:r3:thm:congruence}\\[2pt]
6&\oeis{A396844}&M&I&Mod-16 counterexample and complete replacement law&\pageref{round3:r3:thm:congruence}\\[2pt]
7&\oeis{A188571}&P&I&Infinite prime divisibility and eventual zero-gap periods&\pageref{round3:thm:r3mod-biquadratic}\\[2pt]
8&\oeis{A188572}&P&I&Infinite prime divisibility and eventual zero-gap periods&\pageref{round3:thm:r3mod-biquadratic}\\[2pt]
9&\oeis{A188573}&P&I&Infinite prime divisibility and eventual zero-gap periods&\pageref{round3:thm:r3mod-biquadratic}\\[2pt]
10&\oeis{A377110}&P&I&Infinite prime divisibility and eventual zero-gap periods&\pageref{round3:thm:r3mod-biquadratic}\\[2pt]
11&\oeis{A377111}&P&I&Infinite prime divisibility and eventual zero-gap periods&\pageref{round3:thm:r3mod-biquadratic}\\[2pt]
12&\oeis{A377112}&P&I&Infinite prime divisibility and eventual zero-gap periods&\pageref{round3:thm:r3mod-biquadratic}\\[2pt]
13&\oeis{A377114}&P&I&Infinite prime divisibility and eventual zero-gap periods&\pageref{round3:thm:r3mod-biquadratic}\\[2pt]
14&\oeis{A377115}&P&I&Infinite prime divisibility and eventual zero-gap periods&\pageref{round3:thm:r3mod-biquadratic}\\[2pt]
15&\oeis{A377116}&P&I&Infinite prime divisibility and eventual zero-gap periods&\pageref{round3:thm:r3mod-biquadratic}\\[2pt]
16&\oeis{A377117}&P&I&Infinite prime divisibility and eventual zero-gap periods&\pageref{round3:thm:r3mod-cubic}\\[2pt]
17&\oeis{A377118}&P&I&Infinite prime divisibility and eventual zero-gap periods&\pageref{round3:thm:r3mod-cubic}\\[2pt]
18&\oeis{A377119}&P&I&Infinite prime divisibility and eventual zero-gap periods&\pageref{round3:thm:r3mod-cubic}\\[2pt]
19&\oeis{A195194}&P&N&Nested-reversion parity: power-of-two indices&\pageref{round3:r3entry:A195194}\\[2pt]
20&\oeis{A397241}&P&N&All coefficients of the diagonal relation are odd&\pageref{round3:r3entry:A397241}\\[2pt]
21&\oeis{A396808}&P&N&Parity support at zero and powers of two&\pageref{round3:r3entry:A396808}\\[2pt]
22&\oeis{A397356}&P&N&Parity support where n+1 is a power of two&\pageref{round3:r3entry:A397356}\\[2pt]
23&\oeis{A396803}&P&N&Iterated exponential composition: a(n) = n mod 2&\pageref{round3:r3entry:A396803}\\[2pt]
24&\oeis{A368634}&P&N&Odd coefficients exactly at n = 2\textasciicircum{}j-1&\pageref{round3:r3entry:A368634}\\[2pt]
25&\oeis{A393172}&P&N&Complete mod-4 support at powers of two&\pageref{round3:r3entry:A393172}\\[2pt]
26&\oeis{A389536}&P&N&Complete mod-4 support at n = 2\textasciicircum{}j+1&\pageref{round3:r3entry:A389536}\\[2pt]
27&\oeis{A378580}&P&N&Complete mod-4 support at square indices&\pageref{round3:r3entry:A378580}\\[2pt]
28&\oeis{A188570}&R&N&No term divisible by 71: universal-prime claim false&\pageref{round3:thm:r3mod-constant-counterexamples}\\[2pt]
29&\oeis{A377109}&R&N&No term divisible by 241: universal-prime claim false&\pageref{round3:thm:r3mod-constant-counterexamples}\\[2pt]
30&\oeis{A377113}&R&N&No term divisible by 167: universal-prime claim false&\pageref{round3:thm:r3mod-constant-counterexamples}\\[2pt]
31&\oeis{A004123}&P&N&Stirling-transform period dividing Carmichael lambda(m)&\pageref{round3:thm:r3mod-stirling}\\[2pt]
32&\oeis{A306082}&P&N&Stirling-transform period dividing Carmichael lambda(m)&\pageref{round3:thm:r3mod-stirling}\\[2pt]
33&\oeis{A006531}&P&N&Stirling-transform period dividing Carmichael lambda(m)&\pageref{round3:thm:r3mod-stirling}\\[2pt]
34&\oeis{A167137}&P&N&Stirling-transform period dividing Carmichael lambda(m)&\pageref{round3:thm:r3mod-stirling}\\[2pt]
35&\oeis{A080253}&P&N&Stirling-transform period dividing Carmichael lambda(m)&\pageref{round3:thm:r3mod-stirling}\\[2pt]
36&\oeis{A305550}&P&N&Stirling-transform period dividing Carmichael lambda(m)&\pageref{round3:thm:r3mod-stirling}\\[2pt]
37&\oeis{A354242}&P&N&Stirling-transform period dividing Carmichael lambda(m)&\pageref{round3:thm:r3mod-stirling}\\[2pt]
38&\oeis{A301921}&P&N&Stirling-transform period dividing Carmichael lambda(m)&\pageref{round3:thm:r3mod-stirling}\\[2pt]
39&\oeis{A320352}&P&N&Stirling-transform period dividing Carmichael lambda(m)&\pageref{round3:thm:r3mod-stirling}\\[2pt]
40&\oeis{A354253}&P&N&Stirling-transform period dividing Carmichael lambda(m)&\pageref{round3:thm:r3mod-stirling}\\[2pt]
41&\oeis{A005046}&P&N&Eventual modular periodicity for even-block partitions&\pageref{round3:thm:r3mod-even-blocks}\\[2pt]
42&\oeis{A006154}&P&N&Eventual modular periodicity for ordered odd blocks&\pageref{round3:thm:r3mod-odd-ordered}\\[2pt]
43&\oeis{A121860}&P&N&Factorial divisor sum is eventually 2 modulo m&\pageref{round3:thm:r3mod-factorial}\\[2pt]
44&\oeis{A258899}&P&N&Explicit modular tail 2+2\textasciicircum{}n and eventual period&\pageref{round3:thm:r3mod-factorial}\\[2pt]
45&\oeis{A258903}&P&N&Explicit modular tail 3*2\textasciicircum{}n+2*3\textasciicircum{}n and eventual period&\pageref{round3:thm:r3mod-factorial}\\[2pt]
46&\oeis{A000364}&R&N&Pure period dividing phi(m) fails at m = 27&\pageref{round3:thm:r3mod-euler}\\[2pt]
47&\oeis{A394757}&P&N&Twin-prime centers: gcd is 2 or 3; every term is 18-divisible&\pageref{round3:r3arith:thm:sigma}\\[2pt]
48&\oeis{A394399}&P&N&Gcd-3 centers: twice-square form and set relations&\pageref{round3:r3arith:thm:sigma}\\[2pt]
49&\oeis{A394815}&P&N&Gcd-2 centers are 36-divisible; disjoint-union identity&\pageref{round3:r3arith:thm:sigma}\\[2pt]
50&\oeis{A048702}&P&N&Both empirical formulas for even binary palindromes&\pageref{round3:r3arith:thm:binary}\\[2pt]
51&\oeis{A230089}&P&N&Complex-power formula and rational generating function&\pageref{round3:r3arith:thm:fourclasses}\\[2pt]
52&\oeis{A389425}&P&N&Unit gaps have density one in the sequence index&\pageref{round3:r3arith:thm:density}\\[2pt]
53&\oeis{A092278}&P&N&Empirical generating function for floor((3n+4)/16)&\pageref{round3:r3arith:thm:floor}\\[2pt]
54&\oeis{A005558}&P&N&The displayed lattice-walk recurrence(s), column y = 0&\pageref{round3:r3arith:thm:walkrec}\\[2pt]
55&\oeis{A005559}&P&N&The displayed lattice-walk recurrence(s), column y = 1&\pageref{round3:r3arith:thm:walkrec}\\[2pt]
56&\oeis{A005560}&P&N&The displayed lattice-walk recurrence(s), column y = 2&\pageref{round3:r3arith:thm:walkrec}\\[2pt]
57&\oeis{A005561}&P&N&The displayed lattice-walk recurrence(s), column y = 3&\pageref{round3:r3arith:thm:walkrec}\\[2pt]
58&\oeis{A005562}&P&N&The displayed lattice-walk recurrence(s), column y = 4&\pageref{round3:r3arith:thm:walkrec}\\[2pt]
59&\oeis{A205355}&P&N&Empirical array recurrence of order 4&\pageref{round3:r3entry:A205355}\\[2pt]
60&\oeis{A205356}&P&N&Empirical array recurrence of order 5&\pageref{round3:r3entry:A205356}\\[2pt]
61&\oeis{A205357}&P&N&Empirical array recurrence of order 12&\pageref{round3:r3entry:A205357}\\[2pt]
62&\oeis{A205358}&P&N&Empirical array recurrence of order 19&\pageref{round3:r3entry:A205358}\\[2pt]
63&\oeis{A206377}&P&N&Empirical array recurrence of order 5&\pageref{round3:r3entry:A206377}\\[2pt]
64&\oeis{A206378}&P&N&Empirical array recurrence of order 6&\pageref{round3:r3entry:A206378}\\[2pt]
65&\oeis{A206379}&P&N&Empirical array recurrence of order 11&\pageref{round3:r3entry:A206379}\\[2pt]
66&\oeis{A206380}&P&N&Empirical array recurrence of order 20&\pageref{round3:r3entry:A206380}\\[2pt]
67&\oeis{A209709}&P&N&Empirical array recurrence of order 6&\pageref{round3:r3entry:A209709}\\[2pt]
68&\oeis{A209729}&P&N&Empirical array recurrence of order 3&\pageref{round3:r3entry:A209729}\\[2pt]
69&\oeis{A209807}&P&N&Empirical array recurrence of order 4&\pageref{round3:r3entry:A209807}\\[2pt]
70&\oeis{A209808}&P&N&Empirical array recurrence of order 7&\pageref{round3:r3entry:A209808}\\[2pt]
71&\oeis{A209809}&P&N&Empirical array recurrence of order 14&\pageref{round3:r3entry:A209809}\\[2pt]
72&\oeis{A209810}&P&N&Empirical array recurrence of order 33&\pageref{round3:r3entry:A209810}\\[2pt]
73&\oeis{A209811}&P&N&Empirical array recurrence of order 91&\pageref{round3:r3entry:A209811}\\[2pt]
74&\oeis{A209822}&P&N&Empirical array recurrence of order 3&\pageref{round3:r3entry:A209822}\\[2pt]
75&\oeis{A209823}&P&N&Empirical array recurrence of order 5&\pageref{round3:r3entry:A209823}\\[2pt]
76&\oeis{A209824}&P&N&Empirical array recurrence of order 9&\pageref{round3:r3entry:A209824}\\[2pt]
77&\oeis{A209825}&P&N&Empirical array recurrence of order 21&\pageref{round3:r3entry:A209825}\\[2pt]
78&\oeis{A209826}&P&N&Empirical array recurrence of order 51&\pageref{round3:r3entry:A209826}\\[2pt]
79&\oeis{A209841}&P&N&Empirical array recurrence of order 4&\pageref{round3:r3entry:A209841}\\[2pt]
80&\oeis{A209842}&P&N&Empirical array recurrence of order 6&\pageref{round3:r3entry:A209842}\\[2pt]
81&\oeis{A209843}&P&N&Empirical array recurrence of order 10&\pageref{round3:r3entry:A209843}\\[2pt]
82&\oeis{A209844}&P&N&Empirical array recurrence of order 24&\pageref{round3:r3entry:A209844}\\[2pt]
83&\oeis{A209845}&P&N&Empirical array recurrence of order 59&\pageref{round3:r3entry:A209845}\\[2pt]
84&\oeis{A209890}&P&N&Empirical array recurrence of order 2&\pageref{round3:r3entry:A209890}\\[2pt]
85&\oeis{A209891}&P&N&Empirical array recurrence of order 3&\pageref{round3:r3entry:A209891}\\[2pt]
86&\oeis{A209892}&P&N&Empirical array recurrence of order 9&\pageref{round3:r3entry:A209892}\\[2pt]
87&\oeis{A209893}&P&N&Empirical array recurrence of order 23&\pageref{round3:r3entry:A209893}\\[2pt]
88&\oeis{A209894}&P&N&Empirical array recurrence of order 67&\pageref{round3:r3entry:A209894}\\[2pt]
89&\oeis{A209962}&P&N&Empirical array recurrence of order 6&\pageref{round3:r3entry:A209962}\\[2pt]
90&\oeis{A209963}&P&N&Empirical array recurrence of order 10&\pageref{round3:r3entry:A209963}\\[2pt]
91&\oeis{A209964}&P&N&Empirical array recurrence of order 18&\pageref{round3:r3entry:A209964}\\[2pt]
92&\oeis{A209965}&P&N&Empirical array recurrence of order 44&\pageref{round3:r3entry:A209965}\\[2pt]
93&\oeis{A210045}&P&N&Empirical array recurrence of order 4&\pageref{round3:r3entry:A210045}\\[2pt]
94&\oeis{A210046}&P&N&Empirical array recurrence of order 6&\pageref{round3:r3entry:A210046}\\[2pt]
95&\oeis{A210047}&P&N&Empirical array recurrence of order 14&\pageref{round3:r3entry:A210047}\\[2pt]
96&\oeis{A210048}&P&N&Empirical array recurrence of order 32&\pageref{round3:r3entry:A210048}\\[2pt]
97&\oeis{A210088}&P&N&Empirical array recurrence of order 4&\pageref{round3:r3entry:A210088}\\[2pt]
98&\oeis{A210089}&P&N&Empirical array recurrence of order 6&\pageref{round3:r3entry:A210089}\\[2pt]
99&\oeis{A210091}&P&N&Empirical array recurrence of order 22&\pageref{round3:r3entry:A210091}\\[2pt]
100&\oeis{A368628}&P&N&Odd coefficients exactly at n = (4\textasciicircum{}j-1)/3&\pageref{round3:sec:r3gf-parity}\\[2pt]
\end{longtable}\endgroup\clearpage
\section{Inherited arguments: Pell zeros and elementary obstructions}\label{round3:sec:r3-inherited}
\subsection{A348295: an infinite Pell family of zeros}\label{round3:r3entry:A348295}
Put $\alpha=\sqrt2-1$ and
\[
S(N)=\sum_{k=1}^{N}(-1)^{\lfloor k\alpha\rfloor},\qquad S(0)=0.
\]
Let Pell numbers satisfy $P_0=0$, $P_1=1$, $P_{j+2}=2P_{j+1}+P_j$.

\begin{theorem}\label{round3:r3:thm:pell}
For each integer $m\ge1$,
\[
\boxed{S(2P_{2m})=0.}
\]
In particular, the OEIS zero-index sequence \oeis{A348299} is infinite; the ``infinitely many zeros'' conjecture in \oeis{A348295} holds.
\end{theorem}
\begin{proof}
Set $p=P_{2m}$, $q=P_{2m-1}$, and $\delta=(\sqrt2-1)^{2m}$. Standard Pell identities give
\[
p\alpha=q-\delta,\quad \gcd(p,q)=1,\quad p\text{ even},\quad q\text{ odd}.
\]
Using $P_{2m}=((1+\sqrt2)^{2m}-(\sqrt2-1)^{2m})/(2\sqrt2)$,
\[
2p\delta=\frac{1-(\sqrt2-1)^{4m}}{\sqrt2}<1,
\quad\text{hence}\quad 2\delta<\frac1p.
\]
For every $1\le k<2p$, $k\ne p$, the fractional part of $kq/p$ is $r/p$ for some $1\le r<p$. Since
$0<k\delta/p<2\delta<1/p$, subtraction produces no integer carry, and
\[
\{k\alpha\}=\frac{r}{p}-\frac{k\delta}{p}
\le 1-\frac1p <1-2\delta.
\]
Therefore, using $2p\alpha=2q-2\delta$,
\[
\lfloor (2p-k)\alpha\rfloor=2q-1-\lfloor k\alpha\rfloor.
\]
Thus the terms indexed by $k$ and $2p-k$ have opposite signs and cancel, apart from the central index $k=p$. At $k=p$ the sign is $(-1)^{q-1}=+1$, while the final term at $k=2p$ has sign $(-1)^{2q-1}=-1$. The total is zero.
\end{proof}
\begin{remark}
The first positive witnesses are $4,24,140,816,4756,27720,161564,\ldots$. A separate 2026 OEIS Open result addresses the unboundedness conjecture of \oeis{A348295}; we do not count that published result as ours. Combining unboundedness with the infinite zero family implies every fixed nonnegative value is visited infinitely often.
\end{remark}

\subsection{A014865: a general least-prime-factor obstruction}\label{round3:r3entry:A014865}
Let $b\ge2$ be an integer, $s_1=1$, and $s_n=b s_{n-1}+n$ for $n\ge2$. The identity
\begin{equation}\label{round3:r3:eq:closed}
(b-1)^2s_n=b^{n+1}-(n+1)b+n=b(b^n-1)-n(b-1)
\end{equation}
holds for every $n\ge1$.

\begin{theorem}\label{round3:r3:thm:factor}
If $n>1$ and $n\mid s_n$, then the least prime factor $p$ of $n$ divides $b(b-1)$.
\end{theorem}
\begin{proof}
Let $p$ be the least prime factor of $n$. Reducing \eqref{round3:r3:eq:closed} modulo $p$ and using $p\mid n,s_n$ gives
$b(b^n-1)\equiv0\pmod p$. If $p\mid b$, the conclusion holds. Otherwise $b^n\equiv1\pmod p$, so the order of $b$ modulo $p$ divides both $n$ and $p-1$. As $p$ is the least prime factor of $n$, one has $\gcd(n,p-1)=1$. The order is one, yielding $b\equiv1\pmod p$, hence $p\mid b-1$.
\end{proof}
\begin{theorem}\label{round3:r3:thm:b15}
For $b=15$, every $n$ satisfying $n\mid s_n$ is odd. Moreover, either $n=1$ or $n$ has a factor in $\{3,5,7\}$; thus every such $n>7$ is composite.
\end{theorem}
\begin{proof}
If $n$ is even and $a=v_2(n)\ge1$, the lifting-the-exponent formula gives
\[
v_2(15^n-1)=v_2(14)+v_2(16)+v_2(n)-1=a+4.
\]
Using \eqref{round3:r3:eq:closed} with $(b-1)^2=196$, the numerator
$15(15^n-1)-14n$ has valuation exactly $a+1$, because the first term has valuation $a+4$ and the second valuation $a+1$. Therefore
\[
\boxed{v_2(s_n)=a-1<v_2(n)},
\]
which precludes $n\mid s_n$. So $n$ is odd. By Theorem~\ref{round3:r3:thm:factor} its smallest prime factor must divide $15\cdot14=2\cdot3\cdot5\cdot7$; oddness eliminates $2$. Any prime term must consequently be $3$, $5$, or $7$.
\end{proof}
\begin{remark}
This proves all three assertions displayed under \oeis{A014865}. The general obstruction also applies to the related divisor-recursion sequences \oeis{A014866} ($b=16$) and \oeis{A014868} ($b=18$), but these latter applications are not counted as settled conjectures.
\end{remark}

\subsection{A048868 and A131418: two short structural proofs}
\begin{theorem}\label{round3:r3:thm:phi}
If the reduced residue system modulo $n$ contains more primes than nonprimes, then $n$ is even (\oeis{A048868}).
\end{theorem}
\begin{proof}
For odd $n\ge3$, the totatives in $\{1,\dots,n-1\}$ consist of exactly $\varphi(n)/2$ odd and $\varphi(n)/2$ even integers, under the involution $r\mapsto n-r$. Among even totatives, only $2$ can be prime, and among odd totatives, $1$ is certainly nonprime. Thus the reduced residue system has at least $(\varphi(n)/2-1)+1=\varphi(n)/2$ nonprimes. Its number of primes is at most $\varphi(n)/2$, contrary to a strict prime majority. The case $n=1$ is immediate.
\end{proof}
\begin{theorem}\label{round3:r3:thm:digits}
Let $s_b(n)$ denote the base-$b$ digit sum for $b\ge2$. If
$s_b(n)=s_b(n+s_b(n))$, then $(b-1)\mid n$. In particular every term of \oeis{A131418} is divisible by $9$.
\end{theorem}
\begin{proof}
Since $s_b(t)\equiv t\pmod{b-1}$, the hypothesis implies
$n\equiv n+s_b(n)\pmod{b-1}$. Hence $s_b(n)\equiv0$, whence $n\equiv0\pmod{b-1}$.
\end{proof}

\subsection{A396843 and A396844: self-composed generating functions}
For an integer $c\ge2$, let $A_c(x)=\sum_{n\ge1}a_c(n)x^n\in\mathbb{Z}[[x]]$ be the solution with $a_c(1)=1$ of
\begin{equation}\label{round3:r3:eq:functional}
A_c\bigl(xA_c(x)-cxA_c(x)^2\bigr)=x^2.
\end{equation}
Such a solution exists uniquely: for $n\ge2$, the coefficient of $x^{n+1}$ on the left is $a_c(n)$ plus an integer polynomial in $a_c(1),\ldots,a_c(n-1)$. Thus its vanishing determines $a_c(n)$ successively.

\begin{theorem}\label{round3:r3:thm:congruence}
For each integer $c\ge2$, one has in $\mathbb{Z}[[x]]$,
\begin{equation}\label{round3:r3:eq:universal}
\boxed{A_c(x)\equiv x+c\sum_{j\ge0}(-1)^j x^{2^j+1}\pmod{c^2}.}
\end{equation}
Moreover, when $c$ is odd, every coefficient $a_c(n)$ is odd.
\end{theorem}
\begin{proof}
Modulo $c$, Equation~\eqref{round3:r3:eq:functional} reduces to $A(xA)=x^2$. The series $A=x$ solves this equation and uniqueness gives $A_c\equiv x\pmod c$. Write $A_c=x+cH$ with $H\in x^2\mathbb{Z}[[x]]$. Working modulo $c^2$,
\[
xA_c-cxA_c^2\equiv x^2+c\bigl(xH-x^3\bigr),
\]
so composition in \eqref{round3:r3:eq:functional} yields
\[
xH-x^3+H(x^2)\equiv0\pmod c.
\]
Equivalently $H(x)\equiv x^2-x^{-1}H(x^2)\pmod c$. Iteration, or coefficient induction, gives the unique solution
$H(x)\equiv\sum_{j\ge0}(-1)^jx^{2^j+1}\pmod c$. This proves \eqref{round3:r3:eq:universal}.

If $c$ is odd, reduce \eqref{round3:r3:eq:functional} modulo $2$. There $c\equiv1$. The all-ones series $U(x)=x/(1+x)$ satisfies
\[
U\bigl(xU+xU^2\bigr)=U\!\left(\frac{x^2}{1+x^2}\right)=x^2
\quad\text{in }\mathbb F_2[[x]].
\]
Uniqueness gives $A_c\equiv U\pmod2$, meaning every coefficient is odd.
\end{proof}
\begin{corollary}[\oeis{A396843}, $c=3$]
All terms of \oeis{A396843} are odd. In addition, for $n\ge2$,
\[
a_3(n)\equiv\begin{cases}
3(-1)^j\pmod9,&n=2^j+1,\\
0\pmod9,&\text{otherwise}.
\end{cases}
\]
This supplies a full modular classification, including $a_3(5)\equiv3\pmod9$. The source already records a July 2026 counterexample to its older mod-$9$ conjecture; that refutation is not counted again. The counted target for this entry is the all-odd conjecture.
\end{corollary}
\begin{corollary}[\oeis{A396844}, $c=4$: a refutation and replacement]
For $n\ge2$,
\[
a_4(n)\equiv\begin{cases}
4(-1)^j\pmod{16},&n=2^j+1,\\
0\pmod{16},&\text{otherwise}.
\end{cases}
\]
The OEIS conjecture that indices $n>8$ not of the form $2\cdot4^k+1$
give multiples of $16$ is false:
\[
\boxed{n=17=2^4+1,\quad a_4(17)=113101684573952804\equiv4\pmod{16}.}
\]
In fact every index $n=2^{2j}+1$ with $j\ge2$ is such a counterexample; the conjecture omits the residue-$4$ branch.
\end{corollary}
\begin{corollary}[\oeis{A396845}, $c=5$]
Every coefficient is odd, and for $n\ge2$ the complete residue modulo $25$ is $5(-1)^j$ if $n=2^j+1$, and $0$ otherwise. This consequence is retained from the supplied manuscript but is not counted as a selected open conjecture.
\end{corollary}


\paragraph{Provenance and scope.} These six counted entry-level selections are retained from the supplied exploratory manuscript, after rechecking the proofs and the current OEIS statements. A348299 is an alias for the zero-index question and adds no count. The published OEIS Open result for A348295 concerns unboundedness (Conjecture 1), whereas the selected Pell construction proves infinitely many zeros (Conjecture 2). The auxiliary applications A014866, A014868, and A396845 add no count.

\section{New coefficient congruences}\label{round3:sec:r3-root-new}
All congruences between ordinary generating functions in this section are
coefficientwise. An exponential generating function is explicitly identified
as such; its integer coefficients are never reduced modulo a prime by
illegally inverting factorials.

\subsection{A195194: nested reversions}\label{round3:r3entry:A195194}
Write \(R_k(x)=\operatorname{Rev}(x-kxR_{k+1}(x))\), where
\(\operatorname{Rev}\) denotes compositional inversion and
\(R_k=x+O(x^2)\). The coefficients of any fixed degree stabilize as the
nesting depth increases: multiplication by \(x\) raises the order of a
change by one, and inversion preserves the order of a change.
The sequence \oeis{A195194} is \([x^n]R_1(x)\), \(n\ge1\).
\begin{theorem}
The term of \oeis{A195194} at index \(n\) is odd if and only if
\(n\) is a power of \(2\).
\end{theorem}
\begin{proof}
Reduction of compositional inverses commutes with reduction modulo \(2\).
The factor \(2\) gives \(R_2(x)\equiv x\), and consequently
\[
 R_1(x)\equiv\operatorname{Rev}(x+x^2)
       =\sum_{j\ge0}x^{2^j}\pmod2.
\]
The last identity follows by squaring: if \(S=\sum_{j\ge0}x^{2^j}\),
then \(S+S^2=x\). Its linear coefficient is \(1\), so it is the unique
compositional inverse.
\end{proof}

\subsection{A397241: a diagonal-power relation}\label{round3:r3entry:A397241}
Here \(A(x)=\sum_{n\ge0}a(n)x^n\), \(a(0)=a(1)=1\), and the defining
relation is
\[
 n[x^n]A(x)^n=(n-1)[x^n]A(x)^{n+1}\qquad(n>1).
\]
\begin{theorem}
Every coefficient of \oeis{A397241} is odd.
\end{theorem}
\begin{proof}
For a series with constant term \(1\), its coefficient \(a(n)\) occurs
in \([x^n]A^k\) with multiplier \(k\). Thus the difference of the two
sides has multiplier \(n^2-(n-1)(n+1)=1\) on \(a(n)\).
The initial coefficients therefore determine a unique series over
every commutative ring.

We show that \(U(x)=(1-x)^{-1}\) satisfies the relation over
\(\mathbb F_2\). If \(n\) is even, the required vanishing is
\(\binom{2n}{n}\equiv0\pmod2\). If \(n>1\) is odd, it is
\(\binom{2n-1}{n}\equiv0\pmod2\).
Indeed, with \(s_2(n)\) denoting the number of \(1\)'s in the binary
expansion, the factorial valuation formula gives
\[
 v_2\binom{2n}{n}=s_2(n),\qquad
 v_2\binom{2n-1}{n}=s_2(n)-1.
\]
For an odd \(n>1\), \(s_2(n)\ge2\). Uniqueness now gives \(A\equiv U\).
Only the parity conjecture is selected here; the source's mod-\(3\),
mod-\(4\), and asymptotic assertions are not included in this count.
\end{proof}

\subsection{A396808: another diagonal-power relation}\label{round3:r3entry:A396808}
The definition, again with \(a(0)=a(1)=1\), is
\[
 (n+1)[x^n]A(x)^{n+1}
       =n[x^n]A(x)^{n+2}\qquad(n>1).
\]
\begin{theorem}
For \(n>0\), the term of \oeis{A396808} at index \(n\) is odd if
and only if \(n\) is a power of \(2\).
\end{theorem}
\begin{proof}
The coefficient multiplier is \((n+1)^2-n(n+2)=1\), so reduction
modulo \(2\) still gives a unique solution. Let
\[
 R(x)=\frac{1+\sqrt{1-4x}}2=1-\sum_{n\ge1}C_{n-1}x^n,
\]
where \(C_j\) are the Catalan numbers. Lagrange inversion, applied to
\(1-R=x/(1-(1-R))\), gives, for integers \(k>n\ge1\),
\[
 [x^n]R(x)^k=(-1)^n\frac{k}{k-n}\binom{k-n}{n}.
\]
For \(n>2\), both coefficients with \(k=n+1,n+2\) are zero.
For \(n=2\), they are \(0\) and \(2\), respectively; the defining
relation therefore holds modulo \(2\) also at this boundary.
Moreover \(R\equiv1+x\pmod{x^2,2}\), so uniqueness gives \(A\equiv R\).
Writing \(S=1-R\), the equation \(R^2-R+x=0\) reduces to
\(S^2+S=x\); hence \(S=\sum_{j\ge0}x^{2^j}\).
This proves the exact parity pattern. The separate mod-\(3\)
conjectures and asymptotic limit are not claimed here.
\end{proof}

\subsection{A397356: a reciprocal diagonal relation}\label{round3:r3entry:A397356}
For \oeis{A397356}, \(A(0)=A'(0)=1\) and
\[
 [x^n]A(x)^{-n^2}=[x^n]A(x)^{-(n^2-1)}\qquad(n>1).
\]
\begin{theorem}
The coefficient at index \(n\ge0\) is odd if and only if \(n+1\)
is a power of \(2\).
\end{theorem}
\begin{proof}
The coefficient \(a(n)\) has multiplier \(-1\) in the difference,
so this system is triangular with a unique solution modulo \(2\).
Let \(C=1+xC^2\) be the usual Catalan generating function. Substituting
it into that difference gives
\[
 [x^n](1-C)C^{-n^2}=-[x^{n-1}]C^{\,2-n^2}.
\]
If \(n\) is even, the exponent is even; modulo \(2\) its series has
only even powers, while \(n-1\) is odd.

If \(n>1\) is odd, the Lagrange coefficient formula yields
\[
 [x^{n-1}]C^{\,2-n^2}
 =\frac{2-n^2}{n(2-n)}
       \binom{n(2-n)}{n-1}
 =\frac{2-n^2}{n(2-n)}
       \binom{n^2-n-2}{n-1}.
\]
The prefactor has odd numerator and odd denominator, so it is enough
to show that the last binomial coefficient is even.
Put \(t=v_2(n-1)\ge1\). The bit at position \(t\) of \(n-1\) is \(1\),
whereas
\[
 n^2-n-2=n(n-1)-2\equiv2^t-2\pmod{2^{t+1}}
\]
has bit \(0\) there. The parity criterion for binomial coefficients
therefore gives the required evenness. Thus \(A\equiv C\pmod2\).
Finally \(C=1+xC^2\) gives
\[
 C(x)\equiv\sum_{j\ge0}x^{2^j-1}\pmod2,
\]
which proves the assertion. The source's additional mod-\(3\)
conjecture remains outside the selected scope.
\end{proof}

\subsection{A396803: exponential generating functions and parity}\label{round3:r3entry:A396803}
Let \(\ell\ge0\) and write \(A^{\circ\ell}\) for the \(\ell\)-fold
compositional iterate, with \(A^{\circ0}(x)=x\). Consider
\[
 A(x)=x\exp(A^{\circ\ell}(x))
     =\sum_{n\ge1}a_\ell(n)\frac{x^n}{n!}.
\]
Its coefficients are integers: composition and exponential of
integer exponential series use integer Bell polynomials, and
multiplication by \(x\) multiplies the \(n\)-th coefficient by \(n\).
The same operations give coefficientwise existence and uniqueness,
since the coefficient of degree \(n\) on the right depends only on
coefficients of \(A\) of degrees less than \(n\).

\begin{lemma}\label{round3:r3:bell-parity}
If the integer sequence \(b_1,b_2,\ldots\) satisfies \(b_1=1\)
and \(b_{2j}\) is even for every \(j\ge1\), then the complete
Bell polynomial \(B_{2m}(b_1,\ldots,b_{2m})\) is odd.
\end{lemma}
\begin{proof}
Reduce the block weights modulo \(2\). It suffices to count set
partitions of \(2m\) labels into allowed odd-sized blocks; the
allowed sizes are those \(j\) with \(b_j\) odd.
Act on the labels by a fixed-point-free involution consisting of
\(m\) prescribed transpositions. Partitions not fixed by this
involution pair off.

In a fixed partition, no odd-sized block can be fixed, since a fixed
block would be a union of transposition pairs. Blocks therefore come
in exchanged pairs. A pair of blocks of size \(r\) occupies exactly
\(r\) of the prescribed label pairs, one label from each in each block.
For a prescribed group of \(r\) pairs there are \(2^{r-1}\) ways to
choose these two unordered blocks. If \(r>1\), that factor is even.
The sole odd contribution has \(r=1\) for every block pair, namely
the all-singleton partition, which is allowed because \(b_1=1\).
\end{proof}

\begin{theorem}
For every integer \(\ell\ge0\) and every \(n\ge1\),
\[
 a_\ell(n)\equiv n\pmod2.
\]
In particular, this proves the parity conjecture in
\oeis{A396803}, where \(\ell=3\).
\end{theorem}
\begin{proof}
Write \(B=A^{\circ\ell}\), with exponential coefficients \(b_j\).
The equation immediately gives
\[
 a_\ell(n)=n\,B_{n-1}(b_1,\ldots,b_{n-1}),
\]
so every even-indexed coefficient of \(A\) is even.
Composition preserves this property for exponential series with
zero constant term: in a Bell-polynomial expansion at an even index,
an even-sized block contributes an even inner coefficient, while if
all blocks have odd size, their number is even and the outer
coefficient is even. Consequently all \(b_{2j}\) are even and
\(b_1=1\). For odd \(n\), Lemma~\ref{round3:r3:bell-parity} makes
the displayed Bell polynomial odd. This proves both directions.
The selected claim is parity; other congruences printed in the
OEIS entry are not silently included.
\end{proof}

\subsection{A368634: a parity-sensitive quadratic equation}\label{round3:r3entry:A368634}
Let
\[
 A(x)=1+\frac{x}{2}\bigl(3A(x)^2-A(-x)^2\bigr)
       =\sum_{n\ge0}a(n)x^n.
\]
\begin{theorem}
For \oeis{A368634}, \(a(n)\) is odd exactly when \(n=2^j-1\),
where \(j\ge0\).
\end{theorem}
\begin{proof}
For \(n\ge1\), comparing coefficients gives the integer recurrence
\[
 a(n)=
 \begin{cases}1,&n\text{ odd},\\2,&n\text{ even}\end{cases}
 \ \sum_{i=0}^{n-1}a(i)a(n-1-i).
\]
Thus positive even indices give even values. For odd \(n\), the
off-diagonal convolution terms pair off modulo \(2\), leaving
\(a(n)\equiv a((n-1)/2)\). Starting from \(a(0)=1\), repeated
application gives precisely the stated indices.
\end{proof}

\subsection{A393172: a complete congruence modulo four}\label{round3:r3entry:A393172}
The defining equation is \(A(x-x^2)=x+A(x)^2\), with
\(A(x)=x+O(x^2)\).
\begin{theorem}
For \oeis{A393172},
\[
 A(x)\equiv x+2\sum_{j\ge1}x^{2^j}\pmod4.
\]
Thus, for \(n>1\), \(a(n)\equiv2\pmod4\) exactly at powers of \(2\),
and \(a(n)\equiv0\pmod4\) otherwise.
\end{theorem}
\begin{proof}
The equation is triangular with coefficient \(1\) on each new
coefficient of \(A\). Modulo \(2\), its solution is \(A=x\).
Write \(A=x+2H\). Reduction modulo \(4\) gives
\[
 x-x^2+2H(x-x^2)\equiv x+x^2\pmod4,
\]
hence \(H(x+x^2)=x^2\) in \(\mathbb F_2[[x]]\).
The inverse of \(x+x^2\) is \(S=\sum_{j\ge0}x^{2^j}\),
so \(H=S^2=\sum_{j\ge1}x^{2^j}\) modulo \(2\).
\end{proof}

\subsection{A389536: a shifted dyadic support}\label{round3:r3entry:A389536}
For \oeis{A389536}, \(A=x+O(x^2)\) satisfies
\(xA(x)=A(x^2+2x^3)\).
\begin{theorem}
\[
 A(x)\equiv x+2\sum_{j\ge0}x^{2^j+1}\pmod4.
\]
In particular, for \(n>1\), \(a(n)\equiv2\pmod4\) exactly when
\(n=2^j+1\), and is divisible by \(4\) at every other index.
\end{theorem}
\begin{proof}
For \(n>1\), the coefficient \(a(n)\) is determined by coefficients
of \(A\) with indices at most \(\lfloor(n+1)/2\rfloor<n\).
Modulo \(2\), \(xA=A(x^2)\) forces \(A=x\).
Put \(A=x+2H\). Modulo \(4\),
\[
 x^2+2xH(x)=x^2+2x^3+2H(x^2),
\]
so \(H=x^2+x^{-1}H(x^2)\) modulo \(2\).
Iterating this order-increasing relation gives the asserted sum.
\end{proof}

\subsection{A378580: a theta-series congruence}\label{round3:r3entry:A378580}
The definition is
\[
 A(x/A(x))=1+2\sum_{j\ge1}x^{j^2},\qquad A(0)=1.
\]
\begin{theorem}
For \oeis{A378580},
\[
 A(x)\equiv1+2\sum_{j\ge1}x^{j^2}\pmod4.
\]
Thus a positive-indexed term is \(2\) modulo \(4\) exactly at square
indices and is \(0\) modulo \(4\) elsewhere.
\end{theorem}
\begin{proof}
The coefficient of degree \(n\) in the left side depends on the
new coefficient \(a(n)\) with multiplier \(1\); hence the equation
defines an integer series recursively. Since \(x/A(x)\) has linear
coefficient \(1\), composition with it is invertible modulo \(2\).
The defining equation therefore first gives \(A\equiv1\pmod2\).
Write \(A=1+2H\). Then \(x/A(x)\equiv x\pmod2\), and hence
\[
 A(x/A(x))=1+2H(x/A(x))
          \equiv1+2H(x)\pmod4.
\]
Comparing with the defining theta series proves the result.
The separate question about the radius of convergence is not selected.
\end{proof}

\clearpage
% Requires amsmath, amssymb, amsthm, hyperref and longtable.
% The enclosing document supplies theorem/proof/remark environments and \oeis.
\section{Prime divisibility and modular periodicity}
\label{round3:sec:r3mod}

This section settles 31 selected entry-level assertions: 12 inherited from
the supplied draft and 19 additional assertions.  There are 27 proofs and
four refutations.  Each OEIS identifier below links to the primary record
checked on 20 October 2026.  ``Additional'' describes this project's
provenance; it is not a claim of worldwide mathematical priority.  The
repeated general conjecture about $G(e^x-1)$ is proved once and contributes
no count beyond the ten distinct sequence-specific assertions to which it
is applied.  The established consequence for \oeis{A002050} is explicitly
excluded from the count.

\subsection{Twelve inherited assertions about algebraic powers}
\label{round3:subsec:r3mod-inherited}

The selected assertion in each of the following twelve records is that
every prime divides infinitely many terms and that the successive gaps
between all such indices are eventually periodic:
\[
\begin{array}{c|ccc}
\text{power}&[\sqrt2]&[\sqrt3]&[\sqrt6]\\\hline
(1+\sqrt2+\sqrt3)^n&\oeis{A188571}&\oeis{A188572}&\oeis{A188573}\\
(2+\sqrt2+\sqrt3)^n&\oeis{A377110}&\oeis{A377111}&\oeis{A377112}\\
(3+\sqrt2+\sqrt3)^n&\oeis{A377114}&\oeis{A377115}&\oeis{A377116}
\end{array}
\]
The remaining three entries are the coefficients of $1,t,t^2$ in
$(t+t^2)^n$, where $t^3=2$:
\[
\begin{array}{c|ccc}
\text{coefficient}&1&t&t^2\\\hline
\text{entry}&\oeis{A377117}&\oeis{A377118}&\oeis{A377119}.
\end{array}
\]
The assertion is about the infinite tail, including the full zero-index
set, not merely one infinite subsequence of zeros.  The finite-state
argument below supplies the required statement about the full set.

\begin{lemma}\label{round3:lem:r3mod-gaps}
Let $a_n$ be an integer sequence that is eventually periodic modulo a
prime $p$.  If $p\mid a_n$ for infinitely many $n$, then the increasing
enumeration of all these indices has eventually periodic successive gaps.
\end{lemma}
\begin{proof}
Choose $N,T$ such that $a_{n+T}\equiv a_n\pmod p$ for $n\ge N$.
Let $r_1<\cdots<r_s$ be the zero indices in $[N,N+T)$.
Infinitely many zeros imply $s\ge1$.  After the finite initial portion,
the zero set is exactly
\[
\{r_i+kT:1\le i\le s,\ k\ge0\}.
\]
Its repeating gap block is
$(r_2-r_1,\ldots,r_s-r_{s-1},T+r_1-r_s)$, interpreted as $(T)$ if
$s=1$.
\end{proof}

\begin{theorem}\label{round3:thm:r3mod-biquadratic}
For $c=1,2,3$, every prime divides infinitely many terms of each of the
three nonconstant coefficient sequences in $(c+\sqrt2+\sqrt3)^n$.
Their full zero-index sets have eventually periodic gaps.
\end{theorem}
\begin{proof}
Work in the free rank-four ring
\[
R=\mathbb Z[u,v]/(u^2-2,v^2-3),\qquad \beta_c=c+u+v,
\]
with basis $1,u,v,uv$.  These coordinates agree with the indicated real
algebraic coefficients.  The determinant of multiplication by $\beta_c$
is its norm
\[
N(\beta_c)=\prod_{\epsilon,\delta\in\{1,-1\}}
(c+\epsilon\sqrt2+\delta\sqrt3)=c^4-10c^2+1.
\]
Its values for $c=1,2,3$ are respectively $-8,-23,-8$.  Whenever this
determinant is nonzero modulo $p$, multiplication is invertible on
$R/pR$; in particular $\beta_c$ is a unit there.

For $p\ge5$, Fermat's theorem gives $u^{p^2}=u$ and $v^{p^2}=v$:
for example
$u^{p^2}=2^{(p^2-1)/2}u=u$, since $(p^2-1)/2$ is divisible by $p-1$.
Frobenius therefore gives $\beta_c^{p^2}=\beta_c$.  If also
$p\nmid N(\beta_c)$, cancellation yields
\begin{equation}\label{round3:eq:r3mod-quadratic-period}
\beta_c^{p^2-1}=1\quad\text{in }R/pR.
\end{equation}
Each nonconstant coefficient is consequently zero at every positive
multiple of $p^2-1$.

Here are all exceptional primes, including the nonunit cases.
Modulo $2$,
\[
\beta_1^2=\beta_3^2=0,\qquad \beta_2^2=1.
\]
Thus the first and third families are identically zero after $n=1$,
and the second family has zeros at every positive even index.
Modulo $3$ one has $\beta_c^3=c-u$.  For $c=1,2$ this is a nonzero
element of the field $\mathbb F_3[u]/(u^2+1)$, so
\[
\beta_1^{24}=\beta_2^{24}=1.
\]
For $c=3$, $(\beta_3^3)^4=(-u)^4=1$, giving $\beta_3^{12}=1$.
These identities again supply infinitely many simultaneous zeros of
the three nonconstant coefficients.

The only remaining case is $c=2$, $p=23$.  Since
$2^{11}=3^{11}=1\pmod{23}$, Frobenius now gives
\[
\beta_2^{23}=\beta_2,\qquad
\beta_2^{n+22}=\beta_2^n\quad(n\ge1).
\]
Direct integer expansions give
\[
[u]\beta_2^3=23,\qquad [v]\beta_2^5=529=23^2,
\qquad [uv]\beta_2^1=0.
\]
These yield zeros at the positive indices congruent to $3,5,1$ modulo
$22$, respectively.  No cancellation of the nonunit $\beta_2$ has been
used.

Finally, for any fixed $p$ the sequence $\beta_c^n$ is an orbit under
multiplication by $\beta_c$ in the finite ring $R/pR$.  Every such
orbit is eventually periodic.  Each coordinate is therefore eventually
periodic; the infinite zero families already proved and
Lemma~\ref{round3:lem:r3mod-gaps} prove the assertion about all successive gaps.
\end{proof}

\begin{theorem}\label{round3:thm:r3mod-cubic}
Every prime divides infinitely many terms of each coefficient sequence
of $(t+t^2)^n$ in the basis $1,t,t^2$, with $t^3=2$.  The associated
zero-index gaps are eventually periodic.
\end{theorem}
\begin{proof}
In $T=\mathbb Z[t]/(t^3-2)$ put $\eta=t+t^2$.  Exact multiplication gives
\[
\eta^3=6(1+t+t^2)=6(1+\eta).
\]
For $p=2$ or $3$ we have $\eta^3=0$ in $T/pT$, so every coefficient
vanishes for all $n\ge3$.

Let $p\ge5$.  The identity
$\eta(\eta^2-6)=6$ shows that $\eta$ is a unit.  The polynomial
$t^3-2$ is squarefree modulo $p$, since it has no common root with
$3t^2$.  The algebra $T/pT$ is therefore a product of finite fields
of degrees among $1,2,3$.  Each such degree divides $6$, so the
$p^6$-power map is the identity on every factor.  Thus
\[
\eta^{p^6}=\eta,\qquad \eta^{p^6-1}=1.
\]
Writing $L=p^6-1$, the two nonconstant coefficients vanish at every
positive multiple of $L$.  The constant coefficient of $\eta$ is zero,
so the constant coefficient vanishes at every $1+kL$, $k\ge0$.
All three zero sets are infinite.  Periodicity of the powers in the
finite ring, followed by Lemma~\ref{round3:lem:r3mod-gaps}, proves the gap claim.
\end{proof}

\begin{remark}[Source scope and a finite initial discrepancy]
The displayed data of \oeis{A377117} have value $1$ at $n=1$, whereas
the coefficient definition gives $[1]\eta=0$.  This is precisely a
one-term difference, as can be checked from the printed generating
function.  The true constant-coefficient generating function is
\[
\frac{1-2x^2}{1-6x^2-6x^3},
\]
whereas the printed one is
\[
\frac{1+x-2x^2-6x^3-6x^4}{1-6x^2-6x^3}
=\frac{1-2x^2}{1-6x^2-6x^3}+x.
\]
Both tails therefore satisfy the selected assertion.  This finite
correction is not a separate counted result.

Some ancillary short-period examples printed in the source comments
do not agree with the definitions.  For example, modulo $2$ the
nonconstant coefficients for $c=1$ vanish at every $n\ge2$, so the
eventual zero-index gap is $1$.  The twelve results above settle the
universal infinitude and eventual-gap assertions; they do not validate
the source's separate numerical period lists.
The two-page
\href{https://oeis.org/A188571/a188571_1.pdf}{Sela Fried note linked from
A188571, dated 1 January 2025} proves the already known recurrence and
generating function, not the modular assertion selected here.  Those
earlier recurrence and generating-function results are not counted.
\end{remark}

\subsection{Three constant-coefficient conjectures refuted by complete cycles}
\label{round3:subsec:r3mod-countercycles}

\begin{theorem}\label{round3:thm:r3mod-constant-counterexamples}
Let $C_c(n)=[1](c+u+v)^n$ in the ring above.  Then
\[
\begin{array}{c|c|c|c}
\text{entry}&c&p&\text{conclusion for every }n\ge0\\\hline
\oeis{A188570}&1&71&C_1(n)\not\equiv0\pmod{71}\\
\oeis{A377109}&2&241&C_2(n)\not\equiv0\pmod{241}\\
\oeis{A377113}&3&167&C_3(n)\not\equiv0\pmod{167}.
\end{array}
\]
Consequently each entry's conjecture that every prime divides infinitely
many terms is false.
\end{theorem}
\begin{proof}
For the three respective primes, choose square roots $r^2=2$ and
$s^2=3$ as follows.  All four entries in the last column are nonzero.
\[
\begin{array}{c|c|c|c|c}
c&p&r&s&(c+r+s,c+r-s,c-r+s,c-r-s)\bmod p\\\hline
1&71&12&28&(41,56,17,32)\\
2&241&22&56&(80,209,36,165)\\
3&167&13&62&(78,121,52,95)
\end{array}
\]
Adding the four evaluations under $(u,v)=(\pm r,\pm s)$ cancels
all nonconstant coordinates.  Hence, if the last column is
$(w_1,w_2,w_3,w_4)$,
\begin{equation}\label{round3:eq:r3mod-four-powers}
C_c(n)\equiv 4^{-1}\sum_{i=1}^4 w_i^n\pmod p.
\end{equation}
By Fermat's theorem the right-hand side has period $p-1$ for all
$n\ge0$.  The complete residue lists in the following certificate,
covering $0\le n<p-1$, contain no zero.  This proves the claim for
every $n$, not merely for a tested finite initial segment.

The lists can be checked directly from
\eqref{round3:eq:r3mod-four-powers}, or from the integer coefficient update
\begin{equation}\label{round3:eq:r3mod-vector-update}
(a,b,d,e)\longmapsto
(ca+2b+3d,\ a+cb+3e,\ a+cd+2e,\ b+d+ce)
\end{equation}
starting at $(1,0,0,0)$ and reducing modulo $p$ at each step.
The machine-readable certificate includes all four coordinates at
every step and verifies return to the initial vector after $p-1$
steps.  Its independent four-power check uses
\eqref{round3:eq:r3mod-four-powers}.
\end{proof}

\paragraph{Complete zero-free residue certificates.}
Each line is labeled by the index of its first value; values continue
in increasing order of $n$.  The three blocks contain respectively
$70$, $240$, and $166$ residues.  Every listed value is nonzero.

\noindent\textbf{A188570 modulo 71 (period 70).}
\begingroup\small
\begin{verbatim}
  0:   1,   1,   6,  16,   9,  12,  18,  33,  13,  63,  62,  59
 12:   6,  56,  13,  13,  37,  14,  29,  39,  60,  44,  24,  50
 24:   1,  30,  13,  59,  29,  29,  31,  35,  32,   4,  45,  35
 36:  15,   9,  38,  15,  46,  63,   3,   3,  23,  63,  36,  52
 48:  25,  23,   3,  49,  40,  66,  19,  21,  63,  63,  36,  53
 60:  65,  45,  22,   7,  58,  55,  19,  63,  54,  36
\end{verbatim}
\endgroup

\noindent\textbf{A377109 modulo 241 (period 240).}
\begingroup\small
\begin{verbatim}
  0:   1,   2,   9,  38, 185, 199, 110, 139,  66, 110, 169,  71
 12:  45,  62,  52,  98, 113,  60,  33,  51, 137,  52,  54,  22
 24: 227, 103, 158, 199, 162,  97,  68,  57,  44,  36,  57, 188
 36: 225,  22, 207,  16,  60,  70,  15, 233,  64, 186,  37, 129
 48:  16, 135, 193,  83, 142,  89,  89, 239,  87, 131, 206,  36
 60:  44,   8,  42, 218, 176, 132, 224, 176,  59, 216,  67, 124
 72: 165, 160,   1,  23,  34,   7, 142,  90, 181, 177, 223, 169
 84:  13,  25, 180,  53, 172,  10, 183, 203,  46, 148, 233, 237
 96: 195,  23, 194,  60, 137, 197,  25, 136, 144, 205,  75, 188
108: 198, 175,  54, 240, 221,  73, 186, 121, 212, 193, 199, 218
120: 121,  40,  13, 216, 153, 221, 125, 204, 187,  72, 165, 134
132:  77,  40,  37,  33, 233,  98,  37, 229, 105,  74,  69,  87
144: 107,  65, 154,  21, 194,  75,  53, 233, 164,  74,  61, 125
156: 193,  44, 222,  81, 181,  32,  11,  55,  96, 164,  22,  64
168: 136, 173, 197,  20, 110, 111, 104,  63, 208,  93, 202,  99
180:  76, 227,  27, 153,  55, 170, 228, 113,  27, 238,  82, 189
192:  45, 122, 238,  86,  66, 226, 127,  25,  60, 215, 227, 106
204: 222,  47, 195, 118,  52, 213, 179,  25,  78, 126, 218, 172
216:  74,  61, 198, 238, 105, 219,  40, 201,  24, 167,  71,  10
228: 230, 153,  39, 175, 100, 111, 190,  58, 180, 215, 214,  42
\end{verbatim}
\endgroup

\noindent\textbf{A377113 modulo 167 (period 166).}
\begingroup\small
\begin{verbatim}
  0:   1,   3,  14,  72,  66, 157,  43,  24, 114, 125, 151, 139
 12:  99, 147, 111,  60,  10, 143,  34,  86, 134,  99,  26,  70
 24:   9,  70,   4, 131, 152, 152, 120,  90, 137, 151,  62,  60
 36: 157,  88,  29, 150, 159,  76,  12, 121, 166,  23,  45, 114
 48: 150, 130, 124,  39, 114, 131,  88,  74,  41,   3, 150,  53
 60:  19, 110,  53, 138,  80, 149, 139,  56,  10,   9, 128,  64
 72: 157, 107, 142,  34,  51, 108, 150, 102,  49, 156, 134,  83
 84: 141,  42,  24,  27,  74,  19, 130, 150,  89,  25, 115,  74
 96:  78,  60,  90, 104, 124,  42,  92,  28, 131,  82, 139,  63
108: 125,  44, 166,  47,  46, 124,  42, 137, 104,  65, 110,  39
120:  81,  46,  74,  58, 129,  77, 127, 116,  31, 143,  89,  31
132:  61, 108, 144, 152, 158, 145,  63,   3, 144, 102,  45, 149
144:  11,  59,  54, 106, 146, 152,  85,  17,  85,  57, 110,  22
156:   9,  33,  99,  10,  92,   2,  87, 108, 137,  82
\end{verbatim}
\endgroup


\subsection{One general Stirling theorem and ten new applications}
\label{round3:subsec:r3mod-stirling}

Write $S(n,k)$ for Stirling numbers of the second kind.  For $m\ge2$
define
\[
e(m)=\max_{p\mid m}v_p(m),\qquad
\lambda(m)=\exp((\mathbb Z/m\mathbb Z)^\times),
\]
where $\lambda$ is the Carmichael function.  It divides $\varphi(m)$.

\begin{theorem}\label{round3:thm:r3mod-stirling}
For any integer sequence $(b_k)_{k\ge0}$, put
\[
A_n=\sum_{k=0}^{n}b_k\,k!S(n,k).
\]
For every $m\ge2$ and every $n\ge e(m)$,
\begin{equation}\label{round3:eq:r3mod-stirling-period}
A_{n+\lambda(m)}\equiv A_n\pmod m.
\end{equation}
Equivalently, every exponential generating function of the form
$G(e^x-1)$ with $G(t)\in\mathbb Z[[t]]$ has this eventual period.
In particular its least eventual period divides $\varphi(m)$.
\end{theorem}
\begin{proof}
Since $m\mid k!$ for $k\ge m$, and $S(n,k)=0$ for $k>n$,
\[
A_n\equiv\sum_{k=0}^{m-1}b_k\,k!S(n,k)
=\sum_{k=0}^{m-1}b_k\sum_{j=0}^{k}
(-1)^{k-j}\binom{k}{j}j^n\pmod m.
\]
Let $p^a\Vert m$.  If $p\mid j$ and $n\ge e(m)\ge a$, then
$j^n$ and $j^{n+\lambda(m)}$ both vanish modulo $p^a$.
If $p\nmid j$, then
$j^{\lambda(m)}\equiv1\pmod{p^a}$ because
$\lambda(p^a)\mid\lambda(m)$.  Thus every summand has the stated
period modulo every prime power dividing $m$.  The Chinese remainder
theorem proves \eqref{round3:eq:r3mod-stirling-period}.

The familiar identity
$(e^x-1)^k=\sum_{n\ge0}k!S(n,k)x^n/n!$ proves the equivalent
generating-function formulation by coefficient comparison.  A least
eventual period divides any eventual period, and
$\lambda(m)\mid\varphi(m)$.  Modulo $1$ the assertion is immediate.
\end{proof}

Each of the following ten source records explicitly proposes eventual
periodicity with least period dividing $\varphi(m)$.  Their known
definitions or formulas give the indicated integral weights, so
Theorem~\ref{round3:thm:r3mod-stirling} proves each selected assertion with
the stronger period $\lambda(m)$.
\begin{longtable}{p{0.16\textwidth}p{0.75\textwidth}}
\hline
Entry & Weights $b_k$, or their ordinary generating function $G(t)$\\\hline
\endfirsthead
\hline Entry & Weights or ordinary generating function\\\hline
\endhead
\oeis{A004123}&$b_k=2^k$, so $G(t)=(1-2t)^{-1}$.  The displayed
OEIS offset is $1$: its $a(n)$ is $A_{n-1}$.\\[3pt]
\oeis{A306082}&$b_k$ counts partitions of $k$ into squares;
$G(t)=\prod_{r\ge1}(1-t^{r^2})^{-1}$.\\[3pt]
\oeis{A006531}&$b_0=1$ and $b_k=M_{k-1}$ for $k\ge1$, where
$M_j$ is the $j$th Motzkin number.\\[3pt]
\oeis{A167137}&$b_k$ is the ordinary partition number;
$G(t)=\prod_{r\ge1}(1-t^r)^{-1}$.\\[3pt]
\oeis{A080253}&$G(t)=(1+t)/(1-2t-t^2)$, so
$b_0=1$, $b_1=3$, $b_k=2b_{k-1}+b_{k-2}$.\\[3pt]
\oeis{A305550}&$b_k$ counts partitions into distinct parts;
$G(t)=\prod_{r\ge1}(1+t^r)$.\\[3pt]
\oeis{A354242}&$b_k=\binom{2k}{k}$, so
$G(t)=(1-4t)^{-1/2}$.\\[3pt]
\oeis{A301921}&$G(t)=\displaystyle
\frac{1}{1-\frac{t}{1-\frac{t^2}{1-\frac{t^3}{1-\cdots}}}}$.\\[7pt]
\oeis{A320352}&$b_k=F_k$ with $F_0=0$, $F_1=1$;
$G(t)=t/(1-t-t^2)$.\\[3pt]
\oeis{A354253}&$b_k=2^k\binom{2k}{k}$, so
$G(t)=(1-8t)^{-1/2}$.\\\hline
\end{longtable}
For the continued fraction, the coefficient of any fixed degree is
unchanged after sufficiently deep truncation, because successive
numerators have strictly increasing positive degrees.  Every finite
truncation is obtained by reciprocating series with constant term $1$
over $\mathbb Z[[t]]$; hence every coefficient of its formal limit is
an integer.  This justifies the hypothesis for \oeis{A301921} without
appealing to a numerically fitted series.  The infinite products likewise
have only finitely many relevant factors at any fixed degree.

The index threshold is $n\ge e(m)$ for the nine zero-offset sequences
and $n\ge e(m)+1$ in the displayed indexing of \oeis{A004123}.
For \oeis{A080253}, for example, the source exponential generating
function $e^x/(2-e^{2x})$ becomes
$(1+t)/(1-2t-t^2)$ upon setting $t=e^x-1$, as required.
The formulas used for the weights are existing identities on the cited
pages; the new selected targets here are their modular conjectures.

\begin{remark}[A deliberately uncounted known consequence]
The same proof applies to \oeis{A002050}, with
$G(t)=(t+t^2)/(1-t)$.  However, for $n>0$ this sequence is
$2\,\oeis{A000670}(n)-1$, and the primary record for
\oeis{A000670} already cites Barsky's congruence with period
$(p-1)p^{h-1}$ modulo $p^h$ for $n\ge h$.  Combining these known
congruences by the Chinese remainder theorem immediately gives a period
dividing $\varphi(m)$ for \oeis{A002050}.  It contributes zero to
the 31-entry count.
\end{remark}

\subsection{Two eventual-periodicity proofs from polynomial differential equations}
\label{round3:subsec:r3mod-ode}

\begin{lemma}\label{round3:lem:r3mod-finite-quotient}
Let $D=\sum_{i=1}^{d}Q_i(X_1,\ldots,X_d)\partial_{X_i}$ be a
polynomial derivation over $\mathbb Z$, let $H$ be any integer
polynomial, and let $\mathcal L=D+H$ act on polynomials, with $H$
acting by multiplication.  Fix an integer point
$\boldsymbol c=(c_1,\ldots,c_d)$ and an integer polynomial $P$.
For every positive modulus $m$, the sequence
\[
\bigl(\mathcal L^nP\bigr)(\boldsymbol c)\pmod m
\]
is eventually periodic.
\end{lemma}
\begin{proof}
For $m\ge2$, work in
\[
Q=(\mathbb Z/m\mathbb Z)[X_1,\ldots,X_d] /
(X_1^m-c_1^m,\ldots,X_d^m-c_d^m).
\]
The ideal is stable under $D$, since the derivative of each generator
is $mX_i^{m-1}Q_i=0$ modulo $m$.  It is also stable under multiplication
by $H$.  Consequently $\mathcal L$ is a well-defined endomorphism of
the finite additive module $Q$.  The quotient has basis the $m^d$
monomials whose individual exponents are less than $m$, and therefore
has $m^{m^d}$ elements.  The orbit of $P$ is eventually periodic by
the pigeonhole principle and deterministic iteration.  Evaluation at
$\boldsymbol c$ is well-defined on the quotient, and yields the stated
eventual period.  Modulo $1$ there is nothing to prove.
\end{proof}

\begin{theorem}\label{round3:thm:r3mod-even-blocks}
The sequence \oeis{A005046} is eventually periodic modulo every
positive integer.
\end{theorem}
\begin{proof}
This sequence counts partitions of a $2n$-element set into blocks of
even size and has
\[
F(x)=\exp(\cosh x-1)
=\sum_{n\ge0}a(n)\frac{x^{2n}}{(2n)!}.
\]
Put $X=\sinh x$ and $Y=\cosh x$.  Since $X'=Y$, $Y'=X$ and
$F'=XF$, induction gives
\[
F^{(j)}(x)=P_j(X,Y)F(x),\qquad P_0=1,
\qquad P_{j+1}=(Y\partial_X+X\partial_Y+X)P_j.
\]
All $P_j$ are integer polynomials, and $F(0)=1$,
$(X(0),Y(0))=(0,1)$.  Apply
Lemma~\ref{round3:lem:r3mod-finite-quotient} with the ideal
$(X^m,Y^m-1)$.  It shows that the full derivative sequence
$F^{(j)}(0)$ is eventually periodic modulo $m$.  Every fixed-step
subsequence of an eventually periodic sequence is eventually periodic:
if a full-sequence period is $T$, then advancing $n$ by $T$ advances
the derivative index $2n$ by $2T$.  The even derivative subsequence
$a(n)=F^{(2n)}(0)$ therefore has the required property.
\end{proof}

\begin{theorem}\label{round3:thm:r3mod-odd-ordered}
The sequence \oeis{A006154} is eventually periodic modulo every
positive integer.
\end{theorem}
\begin{proof}
The exponential generating function of ordered set partitions into
odd-sized blocks is $F(x)=1/(1-\sinh x)$.  Put
$X=\sinh x$, $Y=\cosh x$ and $Z=F(x)$.  Then
\[
X'=Y,\qquad Y'=X,\qquad Z'=YZ^2,
\qquad (X(0),Y(0),Z(0))=(0,1,1).
\]
For the integer polynomial derivation
\[
D=Y\partial_X+X\partial_Y+YZ^2\partial_Z,
\]
the chain rule gives $a(n)=(D^nZ)(0,1,1)$.
Lemma~\ref{round3:lem:r3mod-finite-quotient}, with $H=0$ and ideal
$(X^m,Y^m-1,Z^m-1)$, proves eventual periodicity.
\end{proof}

The selected source conjectures for these two entries ask for existence
of an eventual period.  The proof gives existence for every modulus,
without asserting a period dividing $\varphi(m)$ or a sharp bound on
its length.

\subsection{Three factorial divisor sums with explicit modular tails}
\label{round3:subsec:r3mod-factorial}

\begin{theorem}\label{round3:thm:r3mod-factorial}
For positive integers $r,s$ define, for $n\ge1$,
\[
B_{r,s}(n)=\sum_{d\mid n}r^d s^{n/d}\frac{n!}{d!(n/d)!}.
\]
For every positive integer $m$ and all $n\ge2m+4$,
\begin{equation}\label{round3:eq:r3mod-factorial-tail}
B_{r,s}(n)\equiv r s^n+s r^n\pmod m.
\end{equation}
In particular its least eventual period divides $\lambda(m)$, with
the convention $\lambda(1)=1$.
\end{theorem}
\begin{proof}
The case $m=1$ is immediate, so assume $m\ge2$.
The terms $d=1$ and $d=n$ give $rs^n$ and $sr^n$, respectively.
For any other divisor put $e=n/d$; then $d,e\ge2$ and $de=n$.
The inequality $(d-2)(e-2)\ge0$ implies
\[
d+e\le\frac n2+2,\qquad N=n-d-e\ge\frac n2-2\ge m.
\]
Thus the corresponding factorial quotient can be written as
\[
\frac{n!}{d!e!}=\binom{n}{d,e,N}N!,
\]
which is an integer multiple of $m$.  All nonendpoint terms vanish
modulo $m$, proving \eqref{round3:eq:r3mod-factorial-tail}.

For any integer $q$, the proof of
Theorem~\ref{round3:thm:r3mod-stirling} also showed directly that
$q^{n+\lambda(m)}\equiv q^n\pmod m$ once $n\ge e(m)$.
Since $2m+4\ge e(m)$, the displayed tail has period $\lambda(m)$.
\end{proof}

This settles exactly the following three printed modular conjectures:
\[
\begin{array}{c|c|c|l}
\text{entry}&r&s& a(n)\pmod m\text{ for }n\ge2m+4\\\hline
\oeis{A121860}&1&1&2\\
\oeis{A258899}&2&1&2+2^n\\
\oeis{A258903}&2&3&3\cdot2^n+2\cdot3^n.
\end{array}
\]
The first source asks for eventual constancy equal to $2$ for $m\ge3$;
the proof supplies it for every modulus.  The latter two ask for an
eventual period dividing $\varphi(m)$, which follows from the stronger
period $\lambda(m)$.  The divisor-sum identities are the existing
source formulas, not conjectures inferred from an initial numerical list.

\subsection{A pure-period conjecture for Euler numbers fails modulo 27}
\label{round3:subsec:r3mod-euler}

The selected May 2023 conjecture in \oeis{A000364} concerns the sequence
starting at index $1$: it asks for a \emph{pure} period dividing
$\varphi(m)$ for every modulus $m$.  This starting-index requirement
is essential to the following refutation.

\begin{theorem}\label{round3:thm:r3mod-euler}
For the secant numbers $a(n)=|E_{2n}|$ of \oeis{A000364}, the
sequence $(a(n))_{n\ge1}$ has no period dividing $\varphi(27)=18$
modulo $27$.  Hence the stated all-moduli pure-period conjecture is false.
\end{theorem}
\begin{proof}
From
\[
\left(\sum_{n\ge0}a(n)\frac{x^{2n}}{(2n)!}\right)\cos x=1
\]
one obtains $a(0)=1$ and the exact recurrence
\begin{equation}\label{round3:eq:r3mod-euler-recurrence}
a(n)=\sum_{k=0}^{n-1}(-1)^{n-k+1}\binom{2n}{2k}a(k)
\qquad(n\ge1).
\end{equation}
Its residues modulo $27$ at indices $0,1,\ldots,19$ are
\[
1,1,5,7,8,4,11,1,14,25,17,22,20,19,23,16,26,13,2,10.
\]
In particular,
\[
a(1)=1,\qquad
a(19)=23489580527043108252017828576198947741
\equiv10\pmod{27}.
\]
Any pure period $T$ dividing $18$ from index $1$ would imply
$a(19)\equiv a(1)\pmod{27}$, by iterating the period $18/T$ times.
The two residues contradict that implication.
\end{proof}

This refutes only the stated pure-period claim.  It does not refute
eventual periodicity or the established odd-prime theorem of Knuth and
Buckholtz, which the source separately credits.  The recurrence and
the 20-value residue certificate above make the counterexample exact;
no approximation to the Euler numbers is used.

\subsection{Reproducible checks and what they certify}
The standard-library program \nolinkurl{r3_modular_verify.py} writes
\nolinkurl{r3_modular_certificates.json}.  Its complete cycle certificates
prove the three zero-free-prime claims by the finite reduction already
proved above, and its exact secant-number recurrence certifies the
fourth counterexample.  For the universal positive statements, the
written arguments supply the quantifiers; numerical checks merely
guard the definitions, exceptional primes, offsets, and algebraic
identities.  The checks include the biquadratic exceptional cases,
cubic powers for primes below $300$, Stirling representations and
congruences through modulus $40$, both polynomial operators through
derivative $50$ and modulus $12$, and factorial-tail identities through
modulus $30$.  All five verifier components pass.  The accompanying
entry ledger records all 31 counted identifiers and the uncounted
\oeis{A002050}; it verifies disjointness from the 208 prior counted
identifiers supplied to this round.

\clearpage
\section{Twelve arithmetic and lattice-walk assertion groups}
\label{round3:sec:r3arith}
This section settles the specified conjectural formulas or comments in
twelve distinct OEIS entries.  Several assertions in the same entry count
once.  All source pages were checked on 20 October 2026.  The lattice-walk
binomial enumeration is an established input, not a claimed new result;
the selected targets there are the recurrences still labeled conjectural.

\subsection{A394757, A394399, A394815: twin-prime centers and divisor sums}
Let $\sigma(n)=\sum_{d\mid n}d$.  Call $n$ a twin-prime center when both
$n-1$ and $n+1$ are prime.  The three entries respectively impose
$\gcd(n,\sigma(n))$ prime, equal to $3$, and equal to $2$.

\begin{theorem}\label{round3:r3arith:thm:sigma}
If $n$ is a twin-prime center and $g=\gcd(n,\sigma(n))$ is prime, then
$g\in\{2,3\}$ and $18\mid n$.  More precisely:
\begin{enumerate}
\item if $g=3$, then $n=2m^2$ for an integer $m$ divisible by $3$;
\item if $g=2$, then both $v_2(n)$ and $v_3(n)$ are positive even
integers, and in particular $36\mid n$.
\end{enumerate}
Consequently \oeis{A394757} is the disjoint union of \oeis{A394399} and
\oeis{A394815}.  Also, \oeis{A394399} is precisely the intersection of
the twin-prime centers, the numbers having prime $\gcd(n,\sigma(n))$,
and the numbers of the form $2m^2$.
\end{theorem}
\begin{proof}
The center $n=4$ gives $\gcd(4,7)=1$, so it does not occur here.  Every
other twin-prime center is a multiple of $6$: both surrounding primes
are odd, and one of $n-1,n+1$ would be divisible by $3$ unless $3\mid n$.
In what follows $n$ is therefore even and divisible by $3$.

The standard parity criterion for the divisor sum follows directly from
its prime factorization: $\sigma(n)$ is odd if and only if the exponent
of every odd prime in $n$ is even.  Equivalently, $n$ is a square or
twice a square.  If $g$ is odd, the evenness of $n$ forces $\sigma(n)$
to be odd.  A square center greater than $4$ is impossible because
$n-1=(\sqrt n-1)(\sqrt n+1)$ would be composite.  Hence $n=2m^2$.
Now $a=v_2(n)$ is odd, so
\[
 \sigma(2^a)=2^{a+1}-1\equiv0\pmod3.
\]
Since $3\mid n$, the prime $g$ is divisible by $3$, and thus $g=3$.
Moreover $3\mid m$, giving $18\mid n$.

It remains to consider $g=2$.  Put $a=v_2(n)\ge1$ and $b=v_3(n)\ge1$.
If $a$ were odd, the same factor $\sigma(2^a)$ would make
$3\mid\gcd(n,\sigma(n))$, a contradiction.  Thus $a$ is positive
and even, and $4\mid n$.  If $b$ were odd, then
\[
 \sigma(3^b)=1+3+\cdots+3^b\equiv0\pmod4,
\]
by grouping successive pairs.  Multiplicativity of $\sigma$ would imply
$4\mid\sigma(n)$ and hence $4\mid g$, again a contradiction.
Therefore $b$ is positive and even, proving $36\mid n$.

The disjoint-union statement now follows from the only possible prime
values $2$ and $3$.  For the stated intersection, the forward inclusion
has just been proved.  Conversely, an even number $n=2m^2$ has odd
$\sigma(n)$; its gcd cannot equal $2$, so a prime gcd at a twin-prime
center must be $3$.
\end{proof}
\paragraph{Exact source scope.}
This proves both displayed conjectures of \oeis{A394757}, the
twice-a-square conjecture and conjectural intersection of
\oeis{A394399}, and both the divisibility and disjoint-union assertions
of \oeis{A394815}.  The checked entries link only numerical tables;
they do not contain an earlier proof of these selected claims.

\subsection{A048702: two formulas for even binary palindromes}
Let $a(0)=0$, and let $a(n)$ be the $n$th positive even-length binary
palindrome divided by $3$, in increasing order.
\begin{theorem}\label{round3:r3arith:thm:binary}
For $n\ge1$, put $\ell=\lfloor\log_2 n\rfloor$.  Then
\begin{align}
 a(n)&=2^\ell\sum_{i=1}^{n}2^{-v_2(i)},\label{round3:r3arith:binarysum}\\
 &=n2^\ell-\sum_{j=1}^{\ell}2^{\ell-j}
       \left\lfloor\frac n{2^j}\right\rfloor.
       \label{round3:r3arith:binaryfloor}
\end{align}
\end{theorem}
\begin{proof}
Write $L=\ell+1$ and $n=\sum_{t=0}^{L-1}b_t2^t$ with
$b_t\in\{0,1\}$.  The palindrome having $n$ as its first half is
\[
 P(n)=2^Ln+\sum_{t=0}^{L-1}b_t2^{L-1-t}.
\]
These are all positive even-length binary palindromes, in increasing
order: comparison first by length and then by the first half gives the
same order as comparison of $n$.  Thus $a(n)=P(n)/3$.

For every integer $v\ge0$,
$2^{-v}=1-\sum_{j=1}^{v}2^{-j}$.  Summing this identity over $i\le n$
shows that the right sides of \eqref{round3:r3arith:binarysum} and
\eqref{round3:r3arith:binaryfloor} agree.  In the latter expression, the
contribution of a single bit $b_t=1$ is
\[
 2^{L-1}\left(2^t-\sum_{j=1}^{t}2^{t-2j}\right)
   =\frac{2^{L+t}+2^{L-1-t}}3.
\]
Summing over the bits therefore gives exactly $P(n)/3$.
\end{proof}
\paragraph{Exact source scope.}
The two displayed conjectures in \oeis{A048702} are proved for their
positive-index domain.  The existing 2019 comment proving divisibility
of even binary palindromes by $3$ is acknowledged; it does not establish
these two summation identities.  The separate unproved Maple algorithms
are not included in the selected assertion group.

\subsection{A230089: a residue-class formula}
The sequence is defined for $n\ge1$ by $a(n)=n$ for odd $n$,
$a(n)=2$ for $n\equiv2\pmod4$, and $a(n)=4$ for $n\equiv0\pmod4$.
\begin{theorem}\label{round3:r3arith:thm:fourclasses}
For $n\ge1$,
\[
 a(n)=\frac{3+3(-1)^n+(-i)^n+i^n+n-(-1)^n n}{2}.
\]
Its ordinary generating function, starting at index $1$, is
\[
 \sum_{n\ge1}a(n)x^n
 =-\frac{x(4x^5-x^4-2x^3-2x^2-2x-1)}
 {(x-1)^2(x+1)^2(x^2+1)}.
\]
\end{theorem}
\begin{proof}
For odd $n$, the two imaginary powers cancel, the constant terms cancel,
and the numerator is $2n$.  For $n\equiv0\pmod4$ the numerator is
$6+2=8$; for $n\equiv2\pmod4$ it is $6-2=4$.
This proves the asserted expression in every residue class.  Summing
the defining classes gives
\[
 \sum_{n\ge1}a(n)x^n
 =\frac{x(1+x^2)}{(1-x^2)^2}
   +\frac{2x^2+4x^4}{1-x^4},
\]
which simplifies to the displayed rational function.
\end{proof}
The selected source assertion is the conjectural closed form and its
accompanying generating function in \oeis{A230089}; both are counted
together as one entry.

\subsection{A389425: density one of unit gaps}
Let $s_1<s_2<\cdots$ enumerate the positive integers $k$ for which
$6k-5$ and $6k-1$ are both composite (\oeis{A199716}).
The sequence \oeis{A389425} is $a(j)=s_{j+1}-s_j$.
\begin{theorem}\label{round3:r3arith:thm:density}
The set $\{j\ge1:a(j)=1\}$ has natural density $1$.
\end{theorem}
\begin{proof}
We first recall, with an elementary justification, that the primes have
natural density zero.  For any fixed finite set $\mathcal P$ of primes,
all primes outside $\mathcal P$ avoid multiples of its elements.  The
density of integers avoiding these multiples is
$\prod_{p\in\mathcal P}(1-1/p)$.  Taking $\mathcal P$ to be the primes
up to $Y$, these products tend to zero: their reciprocals satisfy
\[
 \prod_{p\le Y}(1-1/p)^{-1}
 =\sum_{\substack{m\ge1\\p\mid m\Rightarrow p\le Y}}\frac1m
 \ge\sum_{m\le Y}\frac1m\longrightarrow\infty.
\]
Thus $\pi(X)=o(X)$, without requiring the prime number theorem.

Let $S=\{s_1,s_2,\ldots\}$ and
$E(X)=\#(\{1,\ldots,\lfloor X\rfloor\}\setminus S)$.
Every excluded $k>1$ has at least one of $6k-5,6k-1$ prime; the
exceptional possibility that one is $1$ contributes at most one.
Consequently
\[
 E(X)\le2\pi(6X)+1=o(X).
\]
It follows that $S$ has density one and $s_N\sim N$.
For every $j<N$ with $s_{j+1}-s_j>1$, the integer $s_j+1$ is excluded
from $S$.  These witnesses are distinct and at most $s_N$.  Hence
\[
 \#\{j<N:a(j)>1\}\le E(s_N)=o(s_N)=o(N),
\]
which is exactly the asserted natural-density statement.
\end{proof}
This supplies a proof of the explicitly conjectural frequency assertion
in \oeis{A389425}.  It is a statement about density in the sequence
index, which is why the passage from $s_N$ to $N$ is included.

\subsection{A092278: the empirical floor generating function}
\begin{theorem}\label{round3:r3arith:thm:floor}
For $a(n)=\lfloor(3n+4)/16\rfloor$, $n\ge0$,
\[
 \sum_{n\ge0}a(n)x^n
 =\frac{x^4+x^{10}+x^{15}}{(1-x)(1-x^{16})}
 =\frac{x^4(x^{11}+x^6+1)}{x^{17}-x^{16}-x+1}.
\]
\end{theorem}
\begin{proof}
One has $a(0)=0$ and $a(n+16)=a(n)+3$.
The increment $a(n)-a(n-1)$, for $n\ge1$, is $1$ exactly in
the classes $n\equiv4,10,15\pmod{16}$ and is $0$ in all other
classes.  Therefore, if $G(x)=\sum_{n\ge0}a(n)x^n$, then
\[
 (1-x)G(x)=\frac{x^4+x^{10}+x^{15}}{1-x^{16}}.
\]
This identity in formal power series proves the claimed expression.
\end{proof}
The target is the generating function labeled empirical in
\oeis{A092278}; this is one assertion group.

\subsection{A005558--A005562: five lattice-walk recurrence claims}
Let $w_y(n)$ be column $y$ of \oeis{A052174}, with $y\ge0$ and
$n\ge y$.  The established enumeration in R.~K.~Guy,
\emph{Catwalks, Sandsteps and Pascal Pyramids}, J.~Integer Sequences
3 (2000), Article 00.1.6, gives the following two formulas, also
recorded in \oeis{A052174}.  With $\binom n{-1}=0$,
\begin{align*}
 w_y(y+2k)&=\binom{y+2k+1}{k}
   \left(\binom{y+2k}{k}-\binom{y+2k}{k-1}\right),\\
 w_y(y+2k+1)&=\binom{y+2k+2}{k}\binom{y+2k+1}{k+1}
   -\binom{y+2k+2}{k+1}\binom{y+2k+1}{k-1}.
\end{align*}
Both formulas hold for $k\ge0$.  The relevant OEIS offsets are
$n=y$, not $n=0$ for every column.  In particular
$w_y(y)=1$ and $w_y(y+1)=y+1$.

\begin{theorem}\label{round3:r3arith:thm:walkrec}
For all integers $y\ge0$ and $n\ge y+2$,
\begin{align}
 &(n-y+1)(n-y)(2n+1)(n+y+3)(n+y+2)w_y(n)\notag\\
 &\quad-4n(n+1)(2n^2+4n+2y^2+4y+3)w_y(n-1)\notag\\
 &\quad-16n^2(n-1)(2n+3)(n+1)w_y(n-2)=0.
 \label{round3:r3arith:walkrec}
\end{align}
Its instances $y=0,1,2,3,4$ prove the selected recurrence conjectures
in \oeis{A005558}, \oeis{A005559}, \oeis{A005560},
\oeis{A005561}, and \oeis{A005562}, respectively.
\end{theorem}
\begin{proof}
Writing the known binomial expressions as factorials gives
\begin{align*}
 E_y(k):=w_y(y+2k)
 &=\frac{(y+1)(y+2k+1)!(y+2k)!}{k!^2(y+k+1)!^2},\\
 O_y(k):=w_y(y+2k+1)
 &=\frac{(y+1)(y+2k+2)!(y+2k+1)!}
 {k!(k+1)!(y+k+2)!(y+k+1)!}.
\end{align*}
These are positive.  The adjacent ratios reduce to
\[
 \frac{w_y(n-1)}{w_y(n)}=\frac{u}{n(n+1)},\qquad
 \frac{w_y(n-2)}{w_y(n)}=\frac{uv}{n^2(n^2-1)},
\]
where the exact values are
\[
\begin{array}{c|c|c}
 n&u&v\\\hline
 y+2k&k(y+k+1)&k(y+k+1)\\
 y+2k+1&(k+1)(y+k+2)&k(y+k+1).
\end{array}
\]
In the stated range $n\ge y+2$, each denominator is nonzero.
After division by $w_y(n)$, the left side of
\eqref{round3:r3arith:walkrec} becomes
\[
 (n-y+1)(n-y)(2n+1)(n+y+3)(n+y+2)
 -4(2n^2+4n+2y^2+4y+3)u-16(2n+3)uv.
\]
Substituting either row of the table and $n=y+2k$ or $n=y+2k+1$
gives the zero polynomial in $k,y$.  This is a polynomial identity,
not a finite-term inference.  The verifier independently expands both
polynomials over $\mathbb Z[k,y]$ and records zero residuals.
\end{proof}

\paragraph{Correspondence with the five printed recurrences.}
For $y=0$, divide \eqref{round3:r3arith:walkrec} by $n(n+1)$ to obtain
\[
 (2n+1)(n+3)(n+2)w_0(n)
 -4(2n^2+4n+3)w_0(n-1)
 -16n(2n+3)(n-1)w_0(n-2)=0\quad(n\ge2).
\]
For $y=1$, divide by $n$ and multiply by $-1$; the result is the
recurrence printed in \oeis{A005559}, for $n\ge3$.
For $y=2,3,4$ the formula already has the printed normalization;
the constants $2y^2+4y+3$ are $19,33,51$, and the valid ranges are
$n\ge4,5,6$, respectively.  Thus the source offsets and boundary
conditions agree, rather than merely matching a shifted numerical list.

\begin{theorem}\label{round3:r3arith:thm:walksecond}
The second recurrence conjectured in \oeis{A005558} also holds:
\begin{align*}
 &(n+3)(n+2)w_0(n)-4(n^2+3n+1)w_0(n-1)\\
 &\qquad+16(-n^2+n+1)w_0(n-2)
 +64(n-1)(n-2)w_0(n-3)=0\qquad(n\ge3).
\end{align*}
\end{theorem}
\begin{proof}
Use the ratio table above with $y=0$.  To obtain the third ratio,
multiply $w_0(n-2)/w_0(n)$ by $w/((n-2)(n-1))$, where
$w=k(k-1)$ for $n=2k$ and $w=k(k+1)$ for $n=2k+1$.
Substitution into the displayed recurrence, followed by multiplication
by $n^2(n-1)(n+1)$, yields in both parity cases the polynomial identity
\begin{align*}
0={}&(n+3)(n+2)n^2(n-1)(n+1)\\
&-4(n^2+3n+1)un(n-1)
 +16(-n^2+n+1)uv+64uvw.
\end{align*}
The identity is verified by inserting the stated values of $n,u,v,w$;
all denominators canceled above are nonzero for $n\ge3$.
\end{proof}

\paragraph{Prior-work and counting boundary.}
The binomial enumeration and existing special-function generating
functions are credited to their earlier sources.  This section provides
explicit algebraic proofs of Mathar's recurrence conjectures dated
2 April 2017, which remain labeled conjectural on the five checked
pages.  It makes no claim of a new lattice-walk enumeration or of
bibliographic priority for consequences of the known formulas.
The two recurrences in \oeis{A005558} count as one entry; its other
conjecturally worded formulas are outside the selected target.

\paragraph{Reproducible certificates.}
The files \nolinkurl{r3_arith_verify.py} and
\nolinkurl{r3_arith_certificates.json} contain exact symbolic certificates
for the two general parity identities and the two third-order parity
identities.  They also check the complete residue periods for
\oeis{A230089} and \oeis{A092278}.  Finite arithmetic and walk checks
are recorded separately as implementation checks, not as substitutes
for the proofs above.  Source URLs, original assertion scopes, offsets,
and count decisions are preserved in \nolinkurl{r3_arith_results.json}.

\clearpage
\section{Exact transfer certificates for 41 array recurrences}
\label{round3:sec:r3gf}
The 41 claims in this section are empirical recurrences explicitly printed
on the cited OEIS pages, checked on 20 October 2026.  They concern different
widths, color bounds, or local constraints.  They are counted by distinct
entry, not as 41 unrelated general methods.  The proof below applies to all
indices and is followed by the complete recurrence signatures.  Finite
prefix agreement is used only to check definitions and indexing.

\subsection{A definition-derived transfer theorem}
Let $w\ge2$.  A row is a word of length $w$ over a fixed alphabet of $j$
colors.  Its canonical form replaces its first distinct color by $0$, its
next new color by $1$, and so on.  Let $\mathcal R_{j,w}$ be the set of
canonical rows with at most $j$ distinct colors.  A canonical row $r$ with
$k(r)$ distinct colors has
\[
 m_j(r)=(j)_{k(r)}=j(j-1)\cdots(j-k(r)+1)
\]
actual labelings.  Suppose the allowed relation $C(r,v)$ between consecutive
rows is invariant under simultaneous color permutations.  Define
\[
 T_{j,C}(r,s)=\#\{v\in\{0,\ldots,j-1\}^{w}:
       \operatorname{can}(v)=s,\ C(r,v)\}.
\]
This is a complete finite matrix: each possible next row is examined once.
If $L_j(n)$ is the number of labeled arrays with $n+1$ rows, then
\begin{equation}\label{round3:eq:r3gf-transfer}
 L_j(n)=m_j^{\mathsf T}T_{j,C}^{n}\mathbf1.
\end{equation}
Indeed, the number of legal completions of an actual row depends only on its
canonical form.  Induction on the number of remaining rows gives the right
side; the weights $m_j$ count all choices of the first row.

For an array using exactly $q$ colors, renaming colors in their order of
first appearance gives one canonical array per $q!$ labelings.  Inclusion--
exclusion therefore gives
\begin{equation}\label{round3:eq:r3gf-exactcolors}
 U_q^{=}(n)=\sum_{j=1}^{q}
 \frac{(-1)^{q-j}}{j!(q-j)!}L_j(n).
\end{equation}
An array with at most $q$ colors has a unique actual number of colors, so
\begin{equation}\label{round3:eq:r3gf-atmostcolors}
 U_q^{\le}(n)=\sum_{j=1}^{q}\left(
 \frac1{j!}\sum_{h=0}^{q-j}\frac{(-1)^h}{h!}\right)L_j(n).
\end{equation}
These are precisely the ``new values introduced in row major order''
conditions in the definitions.  No assumption that all permitted colors
occur is made unless the entry explicitly requires it.

For a $2\times2$ block $\left(\begin{smallmatrix}u&v\\w&t\end{smallmatrix}\right)$,
the local rules used below are:
\begin{itemize}
\item $D_2$: $\#\{u,v,w,t\}=2$;
\item $D_{12}$: $\#\{u,v,w,t\}\le2$;
\item $H_h$: $\mathbf1_{u=t}+\mathbf1_{v=w}=h$, for a fixed $h\in\{0,1,2\}$.
\end{itemize}
For $D_2$ and $D_{12}$ the condition holds at every adjacent column pair.
For the diagonal-equality family the same value $h$ is used throughout the
array, so the count is the sum of the three $H_h$ counts.  Since $n\ge1$
and $w\ge2$, every array has a $2\times2$ block and these three cases are
disjoint.  This explains the source wording ``the same number of equal
diagonal or antidiagonal elements.''

The two edge-statistic entries use the actual numeric labels, so color
permutations are not taken.  Their states are all $4^2=16$ rows, initial
weights are all one, and the final count is divided by $4$ as specified in
their definitions.  For edge sums the four values
$u+v,v+t,t+w,w+u$ must be distinct; for clockwise edge differences the four
values $v-u,t-v,w-t,u-w$ must be distinct.

\subsection{The finite identities that certify every index}
For a displayed signature $(c_1,\ldots,c_d)$ put
\[
 R(z)=z^d-\sum_{i=1}^{d}c_i z^{d-i},\qquad
 Q(x)=1-\sum_{i=1}^{d}c_i x^i.
\]
For every transfer component of every entry listed below, the supplied
integer certificate verifies the vector identity
\begin{equation}\label{round3:eq:r3gf-kernel}
 \boxed{\ T R(T)\mathbf1=0.\ }
\end{equation}
There is no floating point arithmetic and no inferred recurrence in this
check.  The state set and every matrix entry are obtained from the explicit
finite definitions above.  For $n\ge d+1$, multiplying
\eqref{round3:eq:r3gf-kernel} by $m^{\mathsf T}T^{n-d-1}$ gives
\[
 m^{\mathsf T}T^n\mathbf1
   =\sum_{i=1}^{d}c_i m^{\mathsf T}T^{n-i}\mathbf1.
\]
The linear combinations \eqref{round3:eq:r3gf-exactcolors} and
\eqref{round3:eq:r3gf-atmostcolors}, and the disjoint sum over $h$, preserve this
identity.  Thus the precise all-index conclusion for each entry is
\begin{equation}\label{round3:eq:r3gf-recurrence}
 \boxed{a(n)=\sum_{i=1}^{d}c_i a(n-i)\quad(n\ge d+1).}
\end{equation}
The source index begins at $n=1$; this range uses only defined terms.

The complete integer matrices, their initial weights and canonical rows are
in \nolinkurl{r3_gf_transition_matrices.json}.  The verifier reconstructs
every matrix from its combinatorial predicate, compares it with that file,
and recomputes \eqref{round3:eq:r3gf-kernel}.  This is a reproducible finite proof
certificate for the universal recurrence, rather than a test of an
arbitrarily long initial segment.  The files
\nolinkurl{r3_gf_array_certificates.json} and
\nolinkurl{r3_gf_verify.py} record the signatures and checks.

For $F(x)=\sum_{n\ge1}a(n)x^n$, this also gives $QF=P$ with $\deg P\le d$.
The polynomial $P$ is computed from the first $d$ values obtained directly
from \eqref{round3:eq:r3gf-transfer}.  Seventeen correctly indexed empirical
generating functions in the source entries agree exactly with this
construction.  The separately printed generating function in A209962
needs an additional factor $x$ under this indexing; its empirical recurrence
is proved and counted, while the indexing correction adds no count.

\subsection{Individual recurrence signatures}
In the following entries, $D$ is the sum of the dimensions of the transfer
components with nonzero coefficients.  All of their component identities
\eqref{round3:eq:r3gf-kernel} hold.  The coefficient list gives $c_1,c_2,\ldots,c_d$
in that order.  Each entry constitutes one counted assertion group.

\subsubsection{A205355}\label{round3:r3entry:A205355}
Arrays with $n+1$ rows, width $w=3$, at most $3$ colors up to order of first appearance, with rule $H_0,H_1,H_2$.

The recurrence \eqref{round3:eq:r3gf-recurrence} holds with $d=4$, for $n\ge 5$. The certificate has $D=18$ states. Source: \url{https://oeis.org/A205355}.

\[ (c_1,\ldots,c_d)=(15,-76,146,-84).\]

The empirical generating function on the cited page is also obtained from $QF=P$ with the exact boundary polynomial stored in the certificate.

\subsubsection{A205356}\label{round3:r3entry:A205356}
Arrays with $n+1$ rows, width $w=4$, at most $3$ colors up to order of first appearance, with rule $H_0,H_1,H_2$.

The recurrence \eqref{round3:eq:r3gf-recurrence} holds with $d=5$, for $n\ge 6$. The certificate has $D=45$ states. Source: \url{https://oeis.org/A205356}.

\[ (c_1,\ldots,c_d)=(23,-187,645,-912,432).\]

The empirical generating function on the cited page is also obtained from $QF=P$ with the exact boundary polynomial stored in the certificate.

\subsubsection{A205357}\label{round3:r3entry:A205357}
Arrays with $n+1$ rows, width $w=5$, at most $3$ colors up to order of first appearance, with rule $H_0,H_1,H_2$.

The recurrence \eqref{round3:eq:r3gf-recurrence} holds with $d=12$, for $n\ge 13$. The certificate has $D=126$ states. Source: \url{https://oeis.org/A205357}.

\begin{Verbatim}[fontsize=\scriptsize,breaklines=true]
41, -634, 4322, -6905, -79071, 489544, -808128, -1807684, 9900708, -17129232, 13584240,
-4147200
\end{Verbatim}

\subsubsection{A205358}\label{round3:r3entry:A205358}
Arrays with $n+1$ rows, width $w=6$, at most $3$ colors up to order of first appearance, with rule $H_0,H_1,H_2$.

The recurrence \eqref{round3:eq:r3gf-recurrence} holds with $d=19$, for $n\ge 20$. The certificate has $D=369$ states. Source: \url{https://oeis.org/A205358}.

\begin{Verbatim}[fontsize=\scriptsize,breaklines=true]
72, -2111, 31963, -244721, 393930, 9129842, -78749065, 200421324, 702551851,
-6318391361, 15148890613, 7356872888, -116533400840, 228856999872, -38635874496,
-523003405824, 885155081472, -617973276672, 165112971264
\end{Verbatim}

\subsubsection{A206377}\label{round3:r3entry:A206377}
Arrays with $n+1$ rows, width $w=3$, at most $4$ colors up to order of first appearance, with rule $H_0,H_1,H_2$.

The recurrence \eqref{round3:eq:r3gf-recurrence} holds with $d=5$, for $n\ge 6$. The certificate has $D=30$ states. Source: \url{https://oeis.org/A206377}.

\[ (c_1,\ldots,c_d)=(36,-377,1398,-2064,1008).\]

The empirical generating function on the cited page is also obtained from $QF=P$ with the exact boundary polynomial stored in the certificate.

\subsubsection{A206378}\label{round3:r3entry:A206378}
Arrays with $n+1$ rows, width $w=4$, at most $4$ colors up to order of first appearance, with rule $H_0,H_1,H_2$.

The recurrence \eqref{round3:eq:r3gf-recurrence} holds with $d=6$, for $n\ge 7$. The certificate has $D=72$ states. Source: \url{https://oeis.org/A206378}.

\[ (c_1,\ldots,c_d)=(74,-1401,9064,-21992,21312,-7056).\]

\subsubsection{A206379}\label{round3:r3entry:A206379}
Arrays with $n+1$ rows, width $w=5$, at most $4$ colors up to order of first appearance, with rule $H_0,H_1,H_2$.

The recurrence \eqref{round3:eq:r3gf-recurrence} holds with $d=11$, for $n\ge 12$. The certificate has $D=204$ states. Source: \url{https://oeis.org/A206379}.

\begin{Verbatim}[fontsize=\scriptsize,breaklines=true]
160, -5537, 25556, 1495902, -23370552, 81849864, 544237824, -5119960320, 15199301376,
-19180578816, 8497004544
\end{Verbatim}

\subsubsection{A206380}\label{round3:r3entry:A206380}
Arrays with $n+1$ rows, width $w=6$, at most $4$ colors up to order of first appearance, with rule $H_0,H_1,H_2$.

The recurrence \eqref{round3:eq:r3gf-recurrence} holds with $d=20$, for $n\ge 21$. The certificate has $D=660$ states. Source: \url{https://oeis.org/A206380}.

\begin{Verbatim}[fontsize=\scriptsize,breaklines=true]
394, -38562, 1246392, 6334675, -1215511426, 25096715040, -111433495776, -2633554263696,
38447016473440, -124529929487104, -973109346072576, 8528195948079360,
-13461810857735680, -67310021223342080, 238037998037753856, 44544280591749120,
-947327372288188416, 643157082054918144, 987246769084563456, -892353207989698560
\end{Verbatim}

\subsubsection{A209709}\label{round3:r3entry:A209709}
One quarter of the labeled arrays on colors $0,1,2,3$, width $w=2$, with distinct clockwise edge differences.

The recurrence \eqref{round3:eq:r3gf-recurrence} holds with $d=6$, for $n\ge 7$. The certificate has $D=16$ states. Source: \url{https://oeis.org/A209709}.

\[ (c_1,\ldots,c_d)=(2,34,30,-76,-24,16).\]

The empirical generating function on the cited page is also obtained from $QF=P$ with the exact boundary polynomial stored in the certificate.

\subsubsection{A209729}\label{round3:r3entry:A209729}
One quarter of the labeled arrays on colors $0,1,2,3$, width $w=2$, with distinct edge sums.

The recurrence \eqref{round3:eq:r3gf-recurrence} holds with $d=3$, for $n\ge 4$. The certificate has $D=16$ states. Source: \url{https://oeis.org/A209729}.

\[ (c_1,\ldots,c_d)=(3,14,4).\]

The empirical generating function on the cited page is also obtained from $QF=P$ with the exact boundary polynomial stored in the certificate.

\subsubsection{A209807}\label{round3:r3entry:A209807}
Arrays with $n+1$ rows, width $w=2$, at most $4$ colors up to order of first appearance, with rule $D_2$.

The recurrence \eqref{round3:eq:r3gf-recurrence} holds with $d=4$, for $n\ge 5$. The certificate has $D=5$ states. Source: \url{https://oeis.org/A209807}.

\[ (c_1,\ldots,c_d)=(8,-7,-28,-12).\]

The empirical generating function on the cited page is also obtained from $QF=P$ with the exact boundary polynomial stored in the certificate.

\subsubsection{A209808}\label{round3:r3entry:A209808}
Arrays with $n+1$ rows, width $w=3$, at most $4$ colors up to order of first appearance, with rule $D_2$.

The recurrence \eqref{round3:eq:r3gf-recurrence} holds with $d=7$, for $n\ge 8$. The certificate has $D=10$ states. Source: \url{https://oeis.org/A209808}.

\[ (c_1,\ldots,c_d)=(17,-68,-56,397,-181,-264,108).\]

The empirical generating function on the cited page is also obtained from $QF=P$ with the exact boundary polynomial stored in the certificate.

\subsubsection{A209809}\label{round3:r3entry:A209809}
Arrays with $n+1$ rows, width $w=4$, at most $4$ colors up to order of first appearance, with rule $D_2$.

The recurrence \eqref{round3:eq:r3gf-recurrence} holds with $d=14$, for $n\ge 15$. The certificate has $D=24$ states. Source: \url{https://oeis.org/A209809}.

\begin{Verbatim}[fontsize=\scriptsize,breaklines=true]
26, -68, -1905, 6356, 40342, -117881, -302966, 604526, 940543, -1163096, -1137756,
834912, 402624, -124416
\end{Verbatim}

\subsubsection{A209810}\label{round3:r3entry:A209810}
Arrays with $n+1$ rows, width $w=5$, at most $4$ colors up to order of first appearance, with rule $D_2$.

The recurrence \eqref{round3:eq:r3gf-recurrence} holds with $d=33$, for $n\ge 34$. The certificate has $D=68$ states. Source: \url{https://oeis.org/A209810}.

\begin{Verbatim}[fontsize=\scriptsize,breaklines=true]
66, -1211, -3327, 287711, -1650227, -18911756, 200787144, 249811709, -8470618231,
12934099929, 169407260904, -526263544629, -1680919496686, 8328865779706, 6270909877563,
-70313845544213, 27029214503869, 333056045453117, -377395472742367, -825772616226370,
1585680300164445, 703085339056680, -3069417122224407, 900452612086569, 2425264270015320,
-1808654163691437, -347830308417366, 652709084826024, -96050625683220, -59518267835508,
14239281350520, 1128398046480, -309936391200
\end{Verbatim}

\subsubsection{A209811}\label{round3:r3entry:A209811}
Arrays with $n+1$ rows, width $w=6$, at most $4$ colors up to order of first appearance, with rule $D_2$.

The recurrence \eqref{round3:eq:r3gf-recurrence} holds with $d=91$, for $n\ge 92$. The certificate has $D=220$ states. Source: \url{https://oeis.org/A209811}.

\begin{Verbatim}[fontsize=\scriptsize,breaklines=true]
107, -1878, -137789, 4495053, 60979281, -3382378658, -5378727169, 1356204907007,
-5005172801248, -338085898941018, 2348610841442425, 56536290959562784,
-541259839742806307, -6608801096459681129, 81395189596556559983, 548394351778255652595,
-8770207024501159865529, -31666973586868041812990, 710376576458807557547216,
1128924193484746151881180, -44551314523662218100563425, -7263266321979453932360270,
2206376031045169589972758106, -2026353407686910357730744032,
-87454472438177627226033574345, 157141357413725898086880662843,
2799017239399371977020592788297, -7178294228784280740564781901801,
-72691433711711173410574983496660, 239443617254126262209848548396745,
1533257431497069166591149239829307, -6221922488596066027751549217359670,
-26159175017183962402068162814758501, 129627573400367716297066655136898331,
356671191170640145323147937051920800, -2198351761876285273302384497226420197,
-3777310666553059903322898621430359569, 30602786767129664524288240851030652355,
28834501346077162987374441766139145544, -351203971753422209686115577891709785746,
-116701428260928979532142632668200005284, 3327353660864163071142600026522029460808,
-560017644613994507027394206455325197980, -25995896123168019862614781805551174366491,
15575168550871264099563278351155746483043, 166872533834247680226333166810473007305778,
-165583631919347715240453270733787800300967,
-874071453968014666788967372998585172431989,
1206235128308924857940819522655801838918602,
3692161451201415426488274381201249258574504,
-6665858797205359780634341310843509685136667,
-12319864530826648929075899837324201902878036,
28901545602486379348422349423810863556824093,
31170391469749070397533351890857301477542909,
-99502674405964027780060847135394733515338115,
-53857770398272002266481193998251821786322739,
272482776349046561779905216334225226036875695,
37540447794350619447986912229887791121752431,
-590415367694103245739906836722138721711566806,
110686266599666467047553539311810235014535318,
1000112928372531572212817768326632301093822884,
-473363054517674986586028935552949302192543622,
-1296049222794051073104369660002451905377315569,
983069991779873415748977900612905249562346419,
1234590759437986624743595568834744319557269623,
-1367316646006905669194495899803555096031034679,
-790329430414292201826669496840547383132415040,
1352458505939904189893228637165571481776553934,
242440486825297403697854903558758025775196236,
-955596064883055418717334290995333836848088796,
91237272508018656852555859600497039867416040,
470364980986922495221701297389250520645203036,
-149414611561722834497002967906083481125218504,
-151612194930909389895225913796358009798829200,
83459176699483481099731422805986592523856192,
27203058188450643926088113201093410799384736,
-26266525941365495565985041894404510362335360,
-835328187570826089554010524357221294549056,
4821949963807386895787459003961042725066880,
-679112903590477098800215366046658653702208,
-478646169765785125170403007240661020667264, 134388350157539529383186764724725440221184,
19845613671413486485308605949798443768832, -10375428852018050542986066625716287592960,
118385596283169479151035951638780557312, 342248518750406303246673137377100599296,
-26440425803133463527797481814195814400, -4234051645813953234637647944641339392,
454953491082351367104806724121657344, 13225918069327774097709415191674880,
-1738900086789955366482634541629440
\end{Verbatim}

\subsubsection{A209822}\label{round3:r3entry:A209822}
Arrays with $n+1$ rows, width $w=2$, at most $3$ colors up to order of first appearance, with rule $D_{12}$.

The recurrence \eqref{round3:eq:r3gf-recurrence} holds with $d=3$, for $n\ge 4$. The certificate has $D=3$ states. Source: \url{https://oeis.org/A209822}.

\[ (c_1,\ldots,c_d)=(6,-3,-2).\]

The empirical generating function on the cited page is also obtained from $QF=P$ with the exact boundary polynomial stored in the certificate.

\subsubsection{A209823}\label{round3:r3entry:A209823}
Arrays with $n+1$ rows, width $w=3$, at most $3$ colors up to order of first appearance, with rule $D_{12}$.

The recurrence \eqref{round3:eq:r3gf-recurrence} holds with $d=5$, for $n\ge 6$. The certificate has $D=6$ states. Source: \url{https://oeis.org/A209823}.

\[ (c_1,\ldots,c_d)=(12,-16,-6,13,-2).\]

The empirical generating function on the cited page is also obtained from $QF=P$ with the exact boundary polynomial stored in the certificate.

\subsubsection{A209824}\label{round3:r3entry:A209824}
Arrays with $n+1$ rows, width $w=4$, at most $3$ colors up to order of first appearance, with rule $D_{12}$.

The recurrence \eqref{round3:eq:r3gf-recurrence} holds with $d=9$, for $n\ge 10$. The certificate has $D=15$ states. Source: \url{https://oeis.org/A209824}.

\[ (c_1,\ldots,c_d)=(24,-62,-226,445,475,-811,29,147,-20).\]

\subsubsection{A209825}\label{round3:r3entry:A209825}
Arrays with $n+1$ rows, width $w=5$, at most $3$ colors up to order of first appearance, with rule $D_{12}$.

The recurrence \eqref{round3:eq:r3gf-recurrence} holds with $d=21$, for $n\ge 22$. The certificate has $D=42$ states. Source: \url{https://oeis.org/A209825}.

\begin{Verbatim}[fontsize=\scriptsize,breaklines=true]
60, -839, 1399, 31874, -122506, -318771, 1769870, 389301, -8009315, 3097168, 13833529,
-8760492, -9461275, 7615892, 2191722, -2581040, 92008, 286200, -54000, -1264, 480
\end{Verbatim}

\subsubsection{A209826}\label{round3:r3entry:A209826}
Arrays with $n+1$ rows, width $w=6$, at most $3$ colors up to order of first appearance, with rule $D_{12}$.

The recurrence \eqref{round3:eq:r3gf-recurrence} holds with $d=51$, for $n\ge 52$. The certificate has $D=123$ states. Source: \url{https://oeis.org/A209826}.

\begin{Verbatim}[fontsize=\scriptsize,breaklines=true]
138, -5291, 41948, 1246056, -19527125, -94011205, 2626583096, 1511468746, -177763031682,
138302103902, 7019042194083, -7969961832191, -174075869696471, 186648296877104,
2832030056843165, -2415279667981283, -31006473277373149, 19111371400840450,
232526299989981431, -99505929279539438, -1213858407562141252, 370715377868919185,
4475559340570646909, -1096935236489998102, -11786638113405890713, 2781294795279031273,
22309407784237294666, -5920732121027419916, -30324789854047144258, 9761270047165567745,
29307097858909263706, -11646482334059398589, -19688084504632376578, 9649239854975506279,
8805727983122157768, -5380377784965746096, -2398306480372906352, 1943991227752730806,
304993982412468654, -428499039605489966, 11979961616427444, 51791333004986360,
-7840293455999272, -2807693417208336, 702780093944768, 47980140467616, -22670306024704,
380884158080, 247814909440, -6850234368, -945684480
\end{Verbatim}

\subsubsection{A209841}\label{round3:r3entry:A209841}
Arrays with $n+1$ rows, width $w=2$, at most $4$ colors up to order of first appearance, with rule $D_{12}$.

The recurrence \eqref{round3:eq:r3gf-recurrence} holds with $d=4$, for $n\ge 5$. The certificate has $D=5$ states. Source: \url{https://oeis.org/A209841}.

\[ (c_1,\ldots,c_d)=(11,-30,4,16).\]

The empirical generating function on the cited page is also obtained from $QF=P$ with the exact boundary polynomial stored in the certificate.

\subsubsection{A209842}\label{round3:r3entry:A209842}
Arrays with $n+1$ rows, width $w=3$, at most $4$ colors up to order of first appearance, with rule $D_{12}$.

The recurrence \eqref{round3:eq:r3gf-recurrence} holds with $d=6$, for $n\ge 7$. The certificate has $D=10$ states. Source: \url{https://oeis.org/A209842}.

\[ (c_1,\ldots,c_d)=(22,-127,90,288,-400,128).\]

The empirical generating function on the cited page is also obtained from $QF=P$ with the exact boundary polynomial stored in the certificate.

\subsubsection{A209843}\label{round3:r3entry:A209843}
Arrays with $n+1$ rows, width $w=4$, at most $4$ colors up to order of first appearance, with rule $D_{12}$.

The recurrence \eqref{round3:eq:r3gf-recurrence} holds with $d=10$, for $n\ge 11$. The certificate has $D=24$ states. Source: \url{https://oeis.org/A209843}.

\begin{Verbatim}[fontsize=\scriptsize,breaklines=true]
43, -459, -321, 13084, -11850, -84560, 143072, 6528, -98304, 32768
\end{Verbatim}

\subsubsection{A209844}\label{round3:r3entry:A209844}
Arrays with $n+1$ rows, width $w=5$, at most $4$ colors up to order of first appearance, with rule $D_{12}$.

The recurrence \eqref{round3:eq:r3gf-recurrence} holds with $d=24$, for $n\ge 25$. The certificate has $D=68$ states. Source: \url{https://oeis.org/A209844}.

\begin{Verbatim}[fontsize=\scriptsize,breaklines=true]
101, -3119, 23301, 332146, -4841510, 189458, 228888364, -753564284, -2816048344,
16949078144, -1771169504, -117752841744, 156873217312, 244896903072, -582512363712,
-47046663680, 765499169280, -294034673664, -383848898560, 236348588032, 55918985216,
-46495170560, -1491075072, 1811939328
\end{Verbatim}

\subsubsection{A209845}\label{round3:r3entry:A209845}
Arrays with $n+1$ rows, width $w=6$, at most $4$ colors up to order of first appearance, with rule $D_{12}$.

The recurrence \eqref{round3:eq:r3gf-recurrence} holds with $d=59$, for $n\ge 60$. The certificate has $D=220$ states. Source: \url{https://oeis.org/A209845}.

\begin{Verbatim}[fontsize=\scriptsize,breaklines=true]
223, -16175, 333859, 9171561, -425242479, 693710971, 161486863949, -1475551022652,
-27865453094464, 411423026735736, 2219874285551176, -55658326952391880,
-36672282912577248, 4395265143027791120, -7555217898067289496, -220424722751437565016,
708740451169041114912, 7370084221594588613424, -32339860483036642598592,
-168530209053002435663904, 936812779773450938868224, 2642855316304854396118720,
-18715488567630894450654592, -27472172546451820502458112, 268022858676750241561673216,
162950499399519898498751232, -2808765207968050863945458688, -49471709803971865187176448,
21774673081608046584375029760, -8998227402678417097333478400,
-125431619657959886331737444352, 93687752847328791095550070784,
536479435921838898006754967552, -549614260424646029255096614912,
-1693398692148090451777893171200, 2166425771093637945125752602624,
3893649542776241227342317355008, -6029664436570554686166090645504,
-6361043711060989783077333499904, 12040057186303971555099777433600,
7016919022564616606584512970752, -17259291576535507977088989659136,
-4570925557827729604857032605696, 17587140735069580470774635954176,
764365413232018874812935962624, -12491242705025707087031177314304,
1408210287098130144028902031360, 5997551919805635046752510803968,
-1329515975696725414894169489408, -1863857265638179678020068966400,
547108429772846613377035796480, 353189996395406454860564922368,
-117801254695507003051819925504, -37321328417622966875887501312,
12758180005724856234398449664, 1870234589327458736214638592,
-593761192111663779818766336, -25884281547576941708574720, 7740585278261102326579200
\end{Verbatim}

\subsubsection{A209890}\label{round3:r3entry:A209890}
Arrays with $n+1$ rows, width $w=2$, at most $3$ colors up to order of first appearance, with rule $D_2$.

The recurrence \eqref{round3:eq:r3gf-recurrence} holds with $d=2$, for $n\ge 3$. The certificate has $D=3$ states. Source: \url{https://oeis.org/A209890}.

\[ (c_1,\ldots,c_d)=(4,4).\]

The empirical generating function on the cited page is also obtained from $QF=P$ with the exact boundary polynomial stored in the certificate.

\subsubsection{A209891}\label{round3:r3entry:A209891}
Arrays with $n+1$ rows, width $w=3$, at most $3$ colors up to order of first appearance, with rule $D_2$.

The recurrence \eqref{round3:eq:r3gf-recurrence} holds with $d=3$, for $n\ge 4$. The certificate has $D=6$ states. Source: \url{https://oeis.org/A209891}.

\[ (c_1,\ldots,c_d)=(8,4,-16).\]

The empirical generating function on the cited page is also obtained from $QF=P$ with the exact boundary polynomial stored in the certificate.

\subsubsection{A209892}\label{round3:r3entry:A209892}
Arrays with $n+1$ rows, width $w=4$, at most $3$ colors up to order of first appearance, with rule $D_2$.

The recurrence \eqref{round3:eq:r3gf-recurrence} holds with $d=9$, for $n\ge 10$. The certificate has $D=15$ states. Source: \url{https://oeis.org/A209892}.

\[ (c_1,\ldots,c_d)=(14,31,-366,-210,2564,40,-4096,-96,1152).\]

\subsubsection{A209893}\label{round3:r3entry:A209893}
Arrays with $n+1$ rows, width $w=5$, at most $3$ colors up to order of first appearance, with rule $D_2$.

The recurrence \eqref{round3:eq:r3gf-recurrence} holds with $d=23$, for $n\ge 24$. The certificate has $D=42$ states. Source: \url{https://oeis.org/A209893}.

\begin{Verbatim}[fontsize=\scriptsize,breaklines=true]
37, -224, -2979, 29449, 46557, -1073860, 1135122, 15431403, -33807620, -96498020,
305903160, 230102448, -1252267328, 63674496, 2499686528, -1105423872, -2340302336,
1583680512, 822790144, -753893376, -7176192, 78643200, -10485760
\end{Verbatim}

\subsubsection{A209894}\label{round3:r3entry:A209894}
Arrays with $n+1$ rows, width $w=6$, at most $3$ colors up to order of first appearance, with rule $D_2$.

The recurrence \eqref{round3:eq:r3gf-recurrence} holds with $d=67$, for $n\ge 68$. The certificate has $D=123$ states. Source: \url{https://oeis.org/A209894}.

\begin{Verbatim}[fontsize=\scriptsize,breaklines=true]
59, 104, -42972, 225150, 13058670, -99533244, -2197387878, 19875267460, 229851257770,
-2366141610633, -15894063929471, 186716781277733, 749911518956427, -10357596907410544,
-24242602552629694, 419545423349562323, 513419955614790263, -12743636649243029327,
-5588724654850929951, 295948232414959255193, -41030128867263773257,
-5330465080579574799097, 3096535616962170688023, 75245900033979090263854,
-70343205102961136452140, -838515904703767168682000, 1033530787470818142569456,
7408191483521721497325888, -11102771711988258151207232, -51956877639237774786347264,
91197930041539107221974784, 288637005364861044268135424, -584918449332204408531975168,
-1261391919782270562967646208, 2960105091371020869876793344,
4274172957537436605506322432, -11873183283623499865326174208,
-10900689087152926975005589504, 37756515336414977978883178496,
19498582274969148853931212800, -94868879756767967625610788864,
-18957098993963844384982564864, 187014088298245520981590278144,
-10916938203208559688348073984, -285860793239614379638877323264,
79644827524571502161338302464, 332665796495918740601567182848,
-159698020481282454813413474304, -286182176633624129437528227840,
196804266468189564590637449216, 172676696186043581657438486528,
-164414144679340498216040792064, -64875360508321568620864864256,
93966732830765490489645858816, 8984417088496218246521815040,
-35640178748572809806303199232, 4005503615190852715862294528,
8337517133060640365135003648, -2350211162187270322185568256,
-1012217318863117434857979904, 494800755740146443961237504, 26276649640154457035505664,
-44578672501261648598663168, 5284342715267323252441088, 1083745547797691306606592,
-284446307629606977404928, 17414014635957968437248
\end{Verbatim}

\subsubsection{A209962}\label{round3:r3entry:A209962}
Arrays with $n+1$ rows, width $w=2$, exactly $4$ colors up to order of first appearance, with rule $D_{12}$.

The recurrence \eqref{round3:eq:r3gf-recurrence} holds with $d=6$, for $n\ge 7$. The certificate has $D=7$ states. Source: \url{https://oeis.org/A209962}.

\[ (c_1,\ldots,c_d)=(16,-83,132,56,-88,-32).\]

The source's separate generating function omits a factor $x$.
With the stated offset the correct version is
\[
 F(x)=\frac{9x^3(1+2x)}{(1-x)(1-4x)(1-5x-2x^2)(1-6x-4x^2)}.
\]
The initial values $a(1)=a(2)=0$, $a(3)=9$ explain the indexing issue.
The count for this entry is the proved empirical recurrence, not a separate
claim for the correction.

\subsubsection{A209963}\label{round3:r3entry:A209963}
Arrays with $n+1$ rows, width $w=3$, exactly $4$ colors up to order of first appearance, with rule $D_{12}$.

The recurrence \eqref{round3:eq:r3gf-recurrence} holds with $d=10$, for $n\ge 11$. The certificate has $D=15$ states. Source: \url{https://oeis.org/A209963}.

\begin{Verbatim}[fontsize=\scriptsize,breaklines=true]
33, -374, 1586, -1093, -4559, 7212, -420, -4336, 2208, -256
\end{Verbatim}

\subsubsection{A209964}\label{round3:r3entry:A209964}
Arrays with $n+1$ rows, width $w=4$, exactly $4$ colors up to order of first appearance, with rule $D_{12}$.

The recurrence \eqref{round3:eq:r3gf-recurrence} holds with $d=18$, for $n\ge 19$. The certificate has $D=38$ states. Source: \url{https://oeis.org/A209964}.

\begin{Verbatim}[fontsize=\scriptsize,breaklines=true]
66, -1487, 11648, 14141, -454021, 667500, 5458068, -11933206, -23623289, 65454499,
26595728, -134617798, 33919512, 79370592, -38830848, -7503360, 6127616, -655360
\end{Verbatim}

\subsubsection{A209965}\label{round3:r3entry:A209965}
Arrays with $n+1$ rows, width $w=5$, exactly $4$ colors up to order of first appearance, with rule $D_{12}$.

The recurrence \eqref{round3:eq:r3gf-recurrence} holds with $d=44$, for $n\ge 45$. The certificate has $D=109$ states. Source: \url{https://oeis.org/A209965}.

\begin{Verbatim}[fontsize=\scriptsize,breaklines=true]
160, -9858, 286721, -3505459, -7704489, 640517305, -4759487502, -21219049490,
420859791361, -878171254221, -12643232573808, 66259330277918, 122870759491994,
-1553376671736484, 1219539182999473, 17359779478680752, -39146051553201926,
-91438263542603406, 373898684344261096, 116854653298320024, -1799083158352232088,
998160450833133144, 4801216485839890224, -5451035382131524304, -7054956173126826080,
12732134852398029184, 4672289972159061760, -16728156971141573696, 1059265242914046976,
13059429141341396224, -4207097515631301632, -5978493921676298752, 3137966625793618432,
1468906693869291520, -1139164366776567808, -134802883891748864, 213153759748931584,
-11638547304415232, -18765649522720768, 2857764298162176, 564292134895616,
-120413963157504, -2136276467712, 869730877440
\end{Verbatim}

\subsubsection{A210045}\label{round3:r3entry:A210045}
Arrays with $n+1$ rows, width $w=2$, exactly $3$ colors up to order of first appearance, with rule $D_2$.

The recurrence \eqref{round3:eq:r3gf-recurrence} holds with $d=4$, for $n\ge 5$. The certificate has $D=5$ states. Source: \url{https://oeis.org/A210045}.

\[ (c_1,\ldots,c_d)=(7,-6,-20,-8).\]

The empirical generating function on the cited page is also obtained from $QF=P$ with the exact boundary polynomial stored in the certificate.

\subsubsection{A210046}\label{round3:r3entry:A210046}
Arrays with $n+1$ rows, width $w=3$, exactly $3$ colors up to order of first appearance, with rule $D_2$.

The recurrence \eqref{round3:eq:r3gf-recurrence} holds with $d=6$, for $n\ge 7$. The certificate has $D=10$ states. Source: \url{https://oeis.org/A210046}.

\[ (c_1,\ldots,c_d)=(14,-41,-66,100,56,-32).\]

The empirical generating function on the cited page is also obtained from $QF=P$ with the exact boundary polynomial stored in the certificate.

\subsubsection{A210047}\label{round3:r3entry:A210047}
Arrays with $n+1$ rows, width $w=4$, exactly $3$ colors up to order of first appearance, with rule $D_2$.

The recurrence \eqref{round3:eq:r3gf-recurrence} holds with $d=14$, for $n\ge 15$. The certificate has $D=23$ states. Source: \url{https://oeis.org/A210047}.

\begin{Verbatim}[fontsize=\scriptsize,breaklines=true]
24, -89, -977, 3094, 13063, -28268, -71414, 90536, 163472, -114928, -148256, 54080,
35328, -9216
\end{Verbatim}

\subsubsection{A210048}\label{round3:r3entry:A210048}
Arrays with $n+1$ rows, width $w=5$, exactly $3$ colors up to order of first appearance, with rule $D_2$.

The recurrence \eqref{round3:eq:r3gf-recurrence} holds with $d=32$, for $n\ge 33$. The certificate has $D=58$ states. Source: \url{https://oeis.org/A210048}.

\begin{Verbatim}[fontsize=\scriptsize,breaklines=true]
58, -992, 1114, 104383, -609244, -3158300, 31872941, 14255768, -695247797, 851352494,
7722215712, -17461446135, -44607174085, 153341052250, 107386696583, -724717887559,
133505885446, 1890011566228, -1409205824752, -2511295843328, 3234462976416,
1147782485632, -3211415530368, 620083812608, 1267243773440, -658602110976, -85186113536,
120800202752, -16266002432, -5194055680, 1509949440, -104857600
\end{Verbatim}

\subsubsection{A210088}\label{round3:r3entry:A210088}
Arrays with $n+1$ rows, width $w=2$, exactly $3$ colors up to order of first appearance, with rule $D_{12}$.

The recurrence \eqref{round3:eq:r3gf-recurrence} holds with $d=4$, for $n\ge 5$. The certificate has $D=5$ states. Source: \url{https://oeis.org/A210088}.

\[ (c_1,\ldots,c_d)=(10,-27,10,8).\]

The empirical generating function on the cited page is also obtained from $QF=P$ with the exact boundary polynomial stored in the certificate.

\subsubsection{A210089}\label{round3:r3entry:A210089}
Arrays with $n+1$ rows, width $w=3$, exactly $3$ colors up to order of first appearance, with rule $D_{12}$.

The recurrence \eqref{round3:eq:r3gf-recurrence} holds with $d=6$, for $n\ge 7$. The certificate has $D=10$ states. Source: \url{https://oeis.org/A210089}.

\[ (c_1,\ldots,c_d)=(20,-112,122,61,-106,16).\]

The empirical generating function on the cited page is also obtained from $QF=P$ with the exact boundary polynomial stored in the certificate.

\subsubsection{A210091}\label{round3:r3entry:A210091}
Arrays with $n+1$ rows, width $w=5$, exactly $3$ colors up to order of first appearance, with rule $D_{12}$.

The recurrence \eqref{round3:eq:r3gf-recurrence} holds with $d=22$, for $n\ge 23$. The certificate has $D=58$ states. Source: \url{https://oeis.org/A210091}.

\begin{Verbatim}[fontsize=\scriptsize,breaklines=true]
92, -2759, 28247, -12894, -1142474, 3601421, 11970542, -56246539, -20466947, 259395248,
-85275847, -451433420, 270874469, 310376692, -241516822, -72716144, 82685288, -2658056,
-9212400, 1726736, 40928, -15360
\end{Verbatim}

\Needspace{0.65\textheight}
\section{A368628: the exact parity support}
\label{round3:sec:r3gf-parity}
Let $A(x)=\sum_{n\ge0}a(n)x^n$ be the unique formal series with $A(0)=1$
satisfying
\[
 A(x)=1+\frac{x}{2}\bigl(A(x)^2-A(-x)^2\bigr)
        +\frac{x}{2}\bigl(A(x)^4+A(-x)^4\bigr).
\]
The right side determines successive coefficients from earlier ones, so
existence and uniqueness are coefficientwise.  The divisions by $2$
extract odd and even coefficients and preserve integrality.
The explicitly conjectured claim on \url{https://oeis.org/A368628} is
\[
 \boxed{a(n)\text{ is odd}\quad\Longleftrightarrow\quad
 n=\frac{4^k-1}{3}\text{ for some integer }k\ge0.}
\]
To prove it, write $A=E+O$ for its even and odd parts.  Taking these parts in
the defining equation gives integral equations
\[
 E=1+2xEO,\qquad O=x(E^4+6E^2O^2+O^4).
\]
In $\mathbb F_2[[x]]$, therefore, $E=1$ and $O=x(1+O^4)$.
Since $A=1+O$, this is $A=1+xA^4=1+xA(x^4)$.  Iterating the last equation
in the $x$-adic topology gives the unique solution
\[
 A(x)\equiv\sum_{k\ge0}x^{1+4+\cdots+4^{k-1}}
       =\sum_{k\ge0}x^{(4^k-1)/3}\pmod2.
\]
The exponents are distinct, proving both directions of the claim, including
$n=0$.  The algebraic equation and order-68 recurrence credited to Robert
Israel in July 2026 are prior results and are not counted here; this entry
counts only the parity conjecture.  The accompanying check computes exact
coefficients from the original defining equation, independently of this
modulo-two simplification.

\clearpage
\section{Reproduction and certificate inventory}
\label{round3:sec:r3-reproduce}
The accompanying archive is self-contained for checking the mathematical
certificates: after extraction, run the following command in its directory.
\begin{Verbatim}[fontsize=\small]
python3 run_all.py
\end{Verbatim}
Python 3.10 or later and SymPy are required; the recorded environment used
Python 3.12.14 and SymPy 1.14.0. The root and modular verifiers use only the
Python standard library. The arithmetic and transfer verifiers use SymPy
for exact polynomial or rational-function identities. Verification reads
no websites and does not need the earlier rounds or the provisional file.
The requirements file records the tested SymPy version.

\begin{longtable}{p{0.23\textwidth}r p{0.62\textwidth}}
\toprule Component&Count&Certificate and role\\\midrule
\endhead
Root arguments&15&Formal-series calculations and exact arithmetic
checks supplement the written all-index proofs.\\[3pt]
Modular arguments&31&Complete residue cycles and an exact Euler-number
recurrence certify four refutations; the general positive results have
written proofs, with checks for representations and exceptional primes.\\[3pt]
Arithmetic arguments&12&Exact zero polynomial identities certify the
walk recurrences. Complete residue checks certify the short periodic
formulas; bounded arithmetic computations are supplementary.\\[3pt]
Transfer and parity&42&41 complete matrix identities certify the
empirical recurrences. The separate parity theorem has a formal proof;
its exact initial coefficients check the implementation.\\\midrule
Total&100&95 proved, 4 refuted, 1 mixed.\\\bottomrule
\end{longtable}

\paragraph{Transfer data.}
The 90 distinct complete graphs in
\nolinkurl{r3_gf_transition_matrices.json} are reconstructed from their
predicates and compared entry by entry. Each graph has at most 187 states;
the largest sum of component dimensions for an entry is 660.
The recurrence of \oeis{A209811} has order 91.
The certificate test is $T R(T)\mathbf1=0$ in every component;
it is not a fitted scalar recurrence. Initial values and the explicitly
specified starting index complete the statement for each sequence.

\paragraph{Review and reproducibility limits.}
The packet includes independent mathematical cross-checks of the
inherited elementary arguments, the new root congruences, the arithmetic
core, the transfer theorem, and the modular-period theorems.
One off-by-one error in an intermediate binomial transformation for
\oeis{A397356} was corrected during review; the final expression and
its exact regression check agree. These reviews improve reliability
but are not external peer review. The archive contains the actual
programs and finite data so that a reader can inspect or reproduce the
computational portions of the proofs.

\paragraph{Rebuilding this document.}
The supplied \texttt{.tex} contains the entire document; it does not
depend on external proof-fragment inputs. Compile it twice with pdfLaTeX
after installing the standard packages listed in its preamble. The
optional builder \nolinkurl{r3_build_dossier.py} regenerates the standalone
source and ledger from the included component records and fragments.
% ===== END ORIGINAL ROUND 3 =====

% ===== BEGIN ORIGINAL ROUND 4 =====
\clearpage
\setcounter{section}{0}\setcounter{equation}{0}\setcounter{theorem}{0}
\setcounter{table}{0}\setcounter{figure}{0}
\normalsize\renewcommand{\arraystretch}{1}\renewcommand{\baselinestretch}{1}\selectfont
\part[Round 4: Modular obstructions, arithmetic, determinants, and recurrence certificates]{Modular obstructions, arithmetic, determinants, and recurrence certificates}
\label{combined:round4}
\rhead{\small Round 4 / KEIO}
\section*{Round overview}
\addcontentsline{toc}{section}{Round 4 overview and scope}
This fourth-round dossier resolves 100 selected assertion packages on
100 distinct OEIS entries, disjoint from the 308 entries counted in the
preceding three rounds. There are 97 proof-only packages, two refutations,
and one boundary correction. The last concerns \oeis{A243642}: its
unqualified order-six recurrence has residual $24$ at the first admissible
index $n=6$, and holds for every $n\ge7$. Since the source specifies no
starting index, this is recorded as a boundary correction rather than
an unqualified refutation of an explicitly stated lower bound.

The work comprises 14 packages on congruences and formal series, nine
on arithmetic and error estimates, 35 on generating functions, and 42 on
locally constrained arrays. The two refutations are particularly explicit:
\oeis{A126156} satisfies
$a(n)\equiv1+3n(n-1)\pmod4$, which rules out the proposed eventual period
dividing $\varphi(4)$; the recurrence proposed for \oeis{A190363} fails at
$n=161$ and is also incompatible with its irrational limiting slope.

The 42 array recurrences follow from complete integer transfer-matrix
identities. The generating-function recurrences follow from exact
differential identities with all boundary terms retained. Other statements
have written all-index arguments, accompanied by supplementary checks.
The verification archive includes the matrices, symbolic certificates,
component proofs, source ledger and code needed to repeat these checks
without network access.

\subsection*{What the count means}
One catalogue row is one selected assertion package on one A-number.
Several claims proved together on the same entry add one count. Different
entries in a parameter family may use the same theorem; in particular,
the 42 array results use a common transfer construction. The catalogue
therefore does not represent 100 unrelated proof methods.

The selected primary-page assertions are explicitly conjectural,
empirical, or described as probable. An entry can also contain established
formulas and other unresolved questions. Only the assertion specified
in the catalogue and proof is included. Established inputs are credited;
proving a selected comment does not resolve every question on its page.
Examples of deliberately limited scope are the reciprocal-square comment
in \oeis{A320192}, part (ii) in \oeis{A228884}, the 2019 divisor-gap
assertion in \oeis{A258409}, and the mod-3 statements in
\oeis{A368630} and \oeis{A368631}.

The source audit used primary pages consulted on 20 October 2026,
relevant linked work and targeted searches. Previously proved candidates
were excluded when found. The separate comparison with the OEIS Open
paper covers only its stated Appendix A of 100 costliest claims, not
every result in that project. Some retrieved pages may be cached; a
consultation date does not certify the site's instantaneous editorial
state. This dossier asserts the mathematical results given here, and
does not certify worldwide priority or OEIS editorial acceptance.
No OEIS submission or edit is part of this dossier.

The reported reviews are independent calculations and readings by other
agents in the same assistant workflow. They are not external peer review
or proof-assistant formalization. The exact finite certificates are part
of the proofs; bounded sequence-prefix checks are identified separately.

\clearpage
\section{Catalogue of the 100 selected entries}
\label{round4:sec:r4catalogue}
\textbf{P}: proved. \textbf{R}: refuted. \textbf{B}: boundary corrected
with the corrected tail proved. The page reference leads to the individual
argument or recurrence signature. Full definitions, attribution, formulas
and scope notes are retained in \texttt{oeis\_hunt\_round4\_ledger.json}.

\small
\begin{longtable}{@{}r p{21mm} p{103mm} c r@{}}
\toprule
No.&Entry&Selected assertion package&Result&Page\\
\midrule\endfirsthead
\toprule No.&Entry&Selected assertion package&Result&Page\\
\midrule\endhead
\midrule\multicolumn{5}{r}{\emph{Continued on the next page}}\\\endfoot
\bottomrule\endlastfoot
\multicolumn{5}{@{}l}{\rule{0pt}{4ex}\textbf{Congruences and formal series}}\\
1&\oeis{A000657}&Eventual modular periods dividing Euler phi&P&\pageref{round4:r4entry:A000657}\\
2&\oeis{A064618}&Eventual modular periods for the squared-factorial transform&P&\pageref{round4:r4entry:A064618}\\
3&\oeis{A126156}&All-index mod-4 obstruction to the proposed period bound&R&\pageref{round4:r4entry:A126156}\\
4&\oeis{A396099}&Four mod-4 assertions for iterated composition&P&\pageref{round4:r4entry:A396099}\\
5&\oeis{A378575}&Every coefficient is 1 modulo 4&P&\pageref{round4:r4entry:A378575}\\
6&\oeis{A378578}&Odd-valued terms at binary indices with no adjacent ones&P&\pageref{round4:r4entry:A378578}\\
7&\oeis{A380709}&Binomial parity for a coefficientwise absolute-value equation&P&\pageref{round4:r4entry:A380709}\\
8&\oeis{A380711}&Binomial parity and its complete binary support&P&\pageref{round4:r4entry:A380711}\\
9&\oeis{A369083}&Parity of the signed quadratic functional equation&P&\pageref{round4:r4entry:A369083}\\
10&\oeis{A361552}&Half-coefficient parity equals the shifted partition sequence&P&\pageref{round4:r4entry:A361552}\\
11&\oeis{A359914}&Half-coefficient parity equals the partition sequence&P&\pageref{round4:r4entry:A359914}\\
12&\oeis{A380056}&All coefficients are 1 mod 4; multiples of four give 0 mod 5&P&\pageref{round4:r4entry:A380056}\\
13&\oeis{A368630}&Every exponential coefficient is 1 modulo 3&P&\pageref{round4:r4entry:A368630}\\
14&\oeis{A368631}&Every exponential coefficient is 1 modulo 3&P&\pageref{round4:r4entry:A368631}\\
\multicolumn{5}{@{}l}{\rule{0pt}{4ex}\textbf{Arithmetic and exact bounds}}\\
15&\oeis{A190363}&Proposed recurrence fails at 161 and cannot hold eventually&R&\pageref{round4:r4entry:A190363}\\
16&\oeis{A336838}&Divisibility survives reduction of a shifted divisor mean&P&\pageref{round4:r4entry:A336838}\\
17&\oeis{A320192}&The uniquely closest reciprocal-square partial sum (comment only)&P&\pageref{round4:r4entry:A320192}\\
18&\oeis{A111205}&An explicit infinite family of the required composite numbers&P&\pageref{round4:r4entry:A111205}\\
19&\oeis{A025060}&Every integer has the signed distinct nonzero representation&P&\pageref{round4:r4entry:A025060}\\
20&\oeis{A228884}&Prime factors of the gcd determinant are below composite n&P&\pageref{round4:r4entry:A228884}\\
21&\oeis{A228885}&Two totient-product determinant formulas and the squarefree sign&P&\pageref{round4:r4entry:A228885}\\
22&\oeis{A258409}&Gcd of consecutive divisor gaps (2019 assertion only)&P&\pageref{round4:r4entry:A258409}\\
23&\oeis{A176615}&The square-root error term, with an explicit constant 3&P&\pageref{round4:r4entry:A176615}\\
\multicolumn{5}{@{}l}{\rule{0pt}{4ex}\textbf{Generating-function certificates}}\\
24&\oeis{A210496}&Conjectured polynomial recurrence of order 8&P&\pageref{round4:r4entry:A210496}\\
25&\oeis{A102318}&Conjectured polynomial recurrence of order 5&P&\pageref{round4:r4entry:A102318}\\
26&\oeis{A162481}&Conjectured polynomial recurrence of order 4&P&\pageref{round4:r4entry:A162481}\\
27&\oeis{A116391}&Conjectured polynomial recurrence of order 5&P&\pageref{round4:r4entry:A116391}\\
28&\oeis{A114464}&Conjectured polynomial recurrence of order 5&P&\pageref{round4:r4entry:A114464}\\
29&\oeis{A097332}&Conjectured polynomial recurrence of order 3&P&\pageref{round4:r4entry:A097332}\\
30&\oeis{A128057}&Conjectured polynomial recurrence of order 2&P&\pageref{round4:r4entry:A128057}\\
31&\oeis{A100095}&Conjectured polynomial recurrence of order 5&P&\pageref{round4:r4entry:A100095}\\
32&\oeis{A171851}&Conjectured polynomial recurrence of order 9&P&\pageref{round4:r4entry:A171851}\\
33&\oeis{A171847}&Conjectured polynomial recurrence of order 10&P&\pageref{round4:r4entry:A171847}\\
34&\oeis{A171849}&Conjectured polynomial recurrence of order 10&P&\pageref{round4:r4entry:A171849}\\
35&\oeis{A191792}&Conjectured polynomial recurrence of order 8&P&\pageref{round4:r4entry:A191792}\\
36&\oeis{A186338}&Conjectured polynomial recurrence of order 4&P&\pageref{round4:r4entry:A186338}\\
37&\oeis{A182899}&Conjectured polynomial recurrence of order 8&P&\pageref{round4:r4entry:A182899}\\
38&\oeis{A187256}&Conjectured polynomial recurrence of order 4&P&\pageref{round4:r4entry:A187256}\\
39&\oeis{A184120}&Conjectured polynomial recurrence of order 4&P&\pageref{round4:r4entry:A184120}\\
40&\oeis{A186940}&Conjectured polynomial recurrence of order 5&P&\pageref{round4:r4entry:A186940}\\
41&\oeis{A176610}&Conjectured polynomial recurrence of order 5&P&\pageref{round4:r4entry:A176610}\\
42&\oeis{A179648}&Conjectured polynomial recurrence of order 4&P&\pageref{round4:r4entry:A179648}\\
43&\oeis{A166228}&Conjectured polynomial recurrence of order 3&P&\pageref{round4:r4entry:A166228}\\
44&\oeis{A050168}&Conjectured polynomial recurrence of order 3&P&\pageref{round4:r4entry:A050168}\\
45&\oeis{A072100}&Conjectured polynomial recurrence of order 3&P&\pageref{round4:r4entry:A072100}\\
46&\oeis{A176611}&Conjectured polynomial recurrence of order 5&P&\pageref{round4:r4entry:A176611}\\
47&\oeis{A176612}&Conjectured polynomial recurrence of order 5&P&\pageref{round4:r4entry:A176612}\\
48&\oeis{A176605}&Conjectured polynomial recurrence of order 4&P&\pageref{round4:r4entry:A176605}\\
49&\oeis{A176607}&Conjectured polynomial recurrence of order 4&P&\pageref{round4:r4entry:A176607}\\
50&\oeis{A045743}&Conjectured polynomial recurrence of order 4&P&\pageref{round4:r4entry:A045743}\\
51&\oeis{A308726}&Conjectured polynomial recurrence of order 5&P&\pageref{round4:r4entry:A308726}\\
52&\oeis{A285195}&Conjectured polynomial recurrence of order 6&P&\pageref{round4:r4entry:A285195}\\
53&\oeis{A243642}&First-index exception and the corrected recurrence from n = 7&B&\pageref{round4:r4entry:A243642}\\
54&\oeis{A243816}&Conjectured polynomial recurrence of order 6&P&\pageref{round4:r4entry:A243816}\\
55&\oeis{A000986}&Conjectured polynomial recurrence of order 5&P&\pageref{round4:r4entry:A000986}\\
56&\oeis{A121754}&Conjectured polynomial recurrence of order 5&P&\pageref{round4:r4entry:A121754}\\
57&\oeis{A004395}&Conjectured polynomial recurrence of order 4&P&\pageref{round4:r4entry:A004395}\\
58&\oeis{A052731}&The factorial-weighted shift of A002026&P&\pageref{round4:r4entry:A052731}\\
\multicolumn{5}{@{}l}{\rule{0pt}{4ex}\textbf{Finite-state array certificates}}\\
59&\oeis{A210005}&Empirical array recurrence of order 4&P&\pageref{round4:r4entry:A210005}\\
60&\oeis{A210006}&Empirical array recurrence of order 6&P&\pageref{round4:r4entry:A210006}\\
61&\oeis{A210007}&Empirical array recurrence of order 10&P&\pageref{round4:r4entry:A210007}\\
62&\oeis{A210008}&Empirical array recurrence of order 19&P&\pageref{round4:r4entry:A210008}\\
63&\oeis{A210054}&Empirical array recurrence of order 4&P&\pageref{round4:r4entry:A210054}\\
64&\oeis{A210055}&Empirical array recurrence of order 6&P&\pageref{round4:r4entry:A210055}\\
65&\oeis{A210056}&Empirical array recurrence of order 11&P&\pageref{round4:r4entry:A210056}\\
66&\oeis{A210057}&Empirical array recurrence of order 28&P&\pageref{round4:r4entry:A210057}\\
67&\oeis{A210079}&Empirical array recurrence of order 5&P&\pageref{round4:r4entry:A210079}\\
68&\oeis{A210080}&Empirical array recurrence of order 7&P&\pageref{round4:r4entry:A210080}\\
69&\oeis{A210081}&Empirical array recurrence of order 19&P&\pageref{round4:r4entry:A210081}\\
70&\oeis{A210092}&Empirical array recurrence of order 52&P&\pageref{round4:r4entry:A210092}\\
71&\oeis{A210100}&Empirical array recurrence of order 3&P&\pageref{round4:r4entry:A210100}\\
72&\oeis{A210103}&Empirical array recurrence of order 22&P&\pageref{round4:r4entry:A210103}\\
73&\oeis{A210118}&Empirical array recurrence of order 3&P&\pageref{round4:r4entry:A210118}\\
74&\oeis{A210119}&Empirical array recurrence of order 7&P&\pageref{round4:r4entry:A210119}\\
75&\oeis{A210129}&Empirical array recurrence of order 11&P&\pageref{round4:r4entry:A210129}\\
76&\oeis{A210130}&Empirical array recurrence of order 32&P&\pageref{round4:r4entry:A210130}\\
77&\oeis{A210136}&Empirical array recurrence of order 6&P&\pageref{round4:r4entry:A210136}\\
78&\oeis{A210137}&Empirical array recurrence of order 10&P&\pageref{round4:r4entry:A210137}\\
79&\oeis{A210138}&Empirical array recurrence of order 19&P&\pageref{round4:r4entry:A210138}\\
80&\oeis{A210139}&Empirical array recurrence of order 48&P&\pageref{round4:r4entry:A210139}\\
81&\oeis{A210158}&Empirical array recurrence of order 4&P&\pageref{round4:r4entry:A210158}\\
82&\oeis{A210159}&Empirical array recurrence of order 7&P&\pageref{round4:r4entry:A210159}\\
83&\oeis{A210160}&Empirical array recurrence of order 16&P&\pageref{round4:r4entry:A210160}\\
84&\oeis{A210161}&Empirical array recurrence of order 40&P&\pageref{round4:r4entry:A210161}\\
85&\oeis{A210167}&Empirical array recurrence of order 5&P&\pageref{round4:r4entry:A210167}\\
86&\oeis{A210168}&Empirical array recurrence of order 10&P&\pageref{round4:r4entry:A210168}\\
87&\oeis{A210169}&Empirical array recurrence of order 24&P&\pageref{round4:r4entry:A210169}\\
88&\oeis{A210260}&Empirical array recurrence of order 4&P&\pageref{round4:r4entry:A210260}\\
89&\oeis{A210261}&Empirical array recurrence of order 8&P&\pageref{round4:r4entry:A210261}\\
90&\oeis{A210262}&Empirical array recurrence of order 18&P&\pageref{round4:r4entry:A210262}\\
91&\oeis{A210302}&Empirical array recurrence of order 5&P&\pageref{round4:r4entry:A210302}\\
92&\oeis{A210303}&Empirical array recurrence of order 7&P&\pageref{round4:r4entry:A210303}\\
93&\oeis{A210304}&Empirical array recurrence of order 20&P&\pageref{round4:r4entry:A210304}\\
94&\oeis{A210313}&Empirical array recurrence of order 6&P&\pageref{round4:r4entry:A210313}\\
95&\oeis{A210314}&Empirical array recurrence of order 10&P&\pageref{round4:r4entry:A210314}\\
96&\oeis{A210315}&Empirical array recurrence of order 21&P&\pageref{round4:r4entry:A210315}\\
97&\oeis{A210316}&Empirical array recurrence of order 50&P&\pageref{round4:r4entry:A210316}\\
98&\oeis{A210397}&Empirical array recurrence of order 3&P&\pageref{round4:r4entry:A210397}\\
99&\oeis{A210398}&Empirical array recurrence of order 5&P&\pageref{round4:r4entry:A210398}\\
100&\oeis{A210400}&Empirical array recurrence of order 32&P&\pageref{round4:r4entry:A210400}\\
\end{longtable}\normalsize

\clearpage
\section{Congruences, formal series, and a persistent counterexample}
\label{round4:r4:root}

All ordinary power series in this section are formal. Congruences of exponential
generating functions refer to their \emph{integer derivative sequences}, not to
ordinary coefficients with factorial denominators. A conjecture is counted only
at its explicitly named OEIS entry; auxiliary consequences add no count.

\subsection{Finite Euler moments and the median Euler numbers}
\label{round4:r4entry:A000657}
Write
\[
 \operatorname{sech}t=\sum_{r\geq0}E_r\frac{t^r}{r!};
 \qquad E_0=1,\quad E_2=-1,\quad E_4=5.
\]
These are the signed Euler numbers, and $E_{2j+1}=0$.
\begin{lemma}\label{round4:r4:Euler-moments}
For every odd positive integer $M$ and every $r\geq0$,
\[
 E_r\equiv\sum_{j=0}^{M-1}(-1)^j(2j+1)^r\pmod M.
\]
\end{lemma}
\begin{proof}
The finite geometric sum gives the exact identity
\[
 2\sum_{j=0}^{M-1}(-1)^je^{(2j+1)t}
       =(1+e^{2Mt})\operatorname{sech}t.
\]
After taking $r$ derivatives at zero, the right side is $2E_r$ plus
$\sum_{h=1}^r\binom rh(2M)^hE_{r-h}$. Every summand in this last expression
is divisible by $2M$. Divide the integer identity by $2$.
\end{proof}

\begin{theorem}[The selected periodicity conjecture in \oeis{A000657}]
The median Euler numbers are eventually periodic modulo every positive integer
$m$, with a period dividing $\varphi(m)$.
\end{theorem}
\begin{proof}
We use two existing descriptions recorded in the entry: Kruchinin's Euler-number
formula and the Jacobi continued fraction,
\[
 a_n=(-1)^n\sum_{k=0}^n\binom nk E_{n+k},\qquad
 \sum_{n\geq0}a_nx^n
 =\cfrac1{1-b_0x-\cfrac{\lambda_1x^2}
 {1-b_1x-\cfrac{\lambda_2x^2}{1-b_2x-\cdots}}},
\]
where
\[
 b_i=8i^2+4i+1,\qquad \lambda_i=i^2(4i-3)(4i-1)\quad(i\geq1).
\]
Their use does not count as a new enumeration formula.
For odd $M$, Lemma~\ref{round4:r4:Euler-moments} and the binomial theorem imply
\[
 a_n\equiv\sum_{j=0}^{M-1}(-1)^j
       \bigl(-(2j+1)(2j+2)\bigr)^n\pmod M.                 \tag{R1}
\]
Let $e(M)=\max_{p\mid M}v_p(M)$ for $M>1$. Every scalar power sequence on
the right is periodic with period $\lambda(M)$ after index $e(M)$:
modulo $p^\alpha$, a nonunit has become zero by index $\alpha$, while
a unit has period dividing the Carmichael exponent $\lambda(M)$.
Chinese remaindering proves this assertion modulo $M$.

It remains to handle powers of $2$, for which we prove eventual zero.
Fix $a\geq1$, and put $K=2^{\lceil a/2\rceil}$. Let $T$ be the $K$ by $K$
integer tridiagonal matrix indexed by $0,\ldots,K-1$, with diagonal $b_i$,
superdiagonal $1$, and subdiagonal $T_{i,i-1}=\lambda_i$.
The continued-fraction path expansion gives
\[
                         a_n\equiv (T^n)_{0,0}\pmod{2^a}
                         \quad(n\geq0).                 \tag{R2}
\]
Indeed, every omitted excursion reaching height $K$ must return across the
edge of weight $\lambda_K$, and $2^a\mid\lambda_K$.
Modulo $2$, all diagonal entries are $1$ and $\lambda_i\equiv i$.
In consecutive pairs of indices, $T$ is block upper triangular, with
diagonal blocks $\left(\begin{smallmatrix}1&1\\1&1\end{smallmatrix}\right)$.
Thus its characteristic polynomial is $X^K$ over $\mathbb F_2$.
Cayley--Hamilton gives $T^K=2U$ for an integer matrix $U$, so
$T^{aK}=2^aU^a$. Equation (R2) yields $a_n\equiv0\pmod{2^a}$ for
$n\geq aK$.

Write $m=2^aM$ with $M$ odd. Combine this eventual zero with (R1).
For $M>1$ one may take the period $\lambda(M)$ and threshold
$\max(e(M),a2^{\lceil a/2\rceil})$, omitting the second argument when $a=0$.
For $M=1$ take period $1$. These periods divide $\varphi(m)$.
The modulus $1$ is immediate. This proves the stated conjecture, but makes
no assertion about the separate exponential continued-fraction conjecture.
\end{proof}

\subsection{A factorial Stirling transform}
\label{round4:r4entry:A064618}
\begin{theorem}[The selected conjecture in \oeis{A064618}]
For
\[
 a_n=\sum_{k=0}^n(k!)^2S(n,k),
\]
the reduction modulo $m\geq2$ has period dividing $\lambda(m)$, hence dividing
$\varphi(m)$, for every $n\geq e(m)=\max_{p\mid m}v_p(m)$.
\end{theorem}
\begin{proof}
This is a new entry-level application of the factorial Stirling-transform
argument used in Round III. We reproduce the entire argument.
More generally let $c_k$ be any integer sequence and put
$u_n=\sum_k c_k k!S(n,k)$. Terms with $k\geq m$ vanish modulo $m$. For
each remaining $k$,
\[
 k!S(n,k)=\sum_{j=0}^k(-1)^{k-j}\binom kj j^n.
\]
Each scalar $j^n$ has the claimed eventual period by the prime-power argument
in (R1). Sum this finite collection of congruences. Set $c_k=k!$.
The modulus $1$ is trivial.
\end{proof}

\subsection{A period bound which fails forever}
\label{round4:r4entry:A126156}
\begin{theorem}[Refutation for \oeis{A126156}]
Define $a_n$ by
\[
 \sqrt{\sec(\sqrt2\,t)}=\sum_{n\geq0}a_n\frac{t^{2n}}{(2n)!}.
\]
Then
\[
                  a_n\equiv 1+3n(n-1)\pmod4 \quad(n\geq0).       \tag{R3}
\]
Consequently $a_{n+2}-a_n\equiv2\pmod4$ for every $n\geq0$.
The entry's conjecture of eventual period dividing $\varphi(m)$ is false:
for $m=4$ no tail has period $1$ or $2$.
\end{theorem}
\begin{proof}
Set $y(t)=\tan(\sqrt2\,t)/\sqrt2$ and
$f(t)=\sqrt{\sec(\sqrt2\,t)}$. Then $y'=1+2y^2$ and $f'=yf$.
Repeated differentiation therefore gives
\[
 a_n=(L^{2n}1)(0),\qquad L=(1+2y^2)\frac d{dy}+y.
\]
In particular the coefficients being evaluated are integers. Direct expansion
yields
\[
 L^2=(1+2y^2)^2\frac{d^2}{dy^2}
       +6y(1+2y^2)\frac d{dy}+1+3y^2.
\]
This preserves the even polynomials. On writing $z=y^2$, its reduction modulo
$4$ is the operator
\[
                  H:Q(z)\longmapsto (1+3z)Q(z)+2Q'(z).
\]
Put $v=1+3z$. In expanding $H^n1$, terms with two or more derivative
operators have a factor $4$. For $n\geq2$ the remaining terms are
\[
 H^n1\equiv v^n+
    2\sum_{i=0}^{n-1}v^i\frac d{dz}v^{n-1-i}
   =v^n+3n(n-1)v^{n-2}\pmod4.
\]
Evaluate at $z=0$. The cases $n=0,1$ give $1$ separately.
Replacing $n$ by $n+2$ in (R3) gives the persistent difference $2$.
This is an all-index obstruction, not a finite-prefix disproof of
eventual periodicity. The sequence does have the exact period $4$ modulo $4$.
\end{proof}

\subsection{Iterated composition modulo four}
\label{round4:r4entry:A396099}
\begin{theorem}[All four selected congruences in \oeis{A396099}]
Let $A=x+A^{\circ2}A^{\circ3}$, with $A=x+O(x^2)$.
Its coefficients are all odd, and for $n=3,4,5,6,\ldots$ their residues
modulo $4$ repeat $1,3,3,1$. Moreover,
\[
 A^{\circ2}\equiv x+2x^2\pmod4,\qquad
 A-xA^{\circ3}\equiv x+\frac{2x^3}{1-x}\pmod4.
\]
\end{theorem}
\begin{proof}
The coefficient of degree $n$ on the right of the defining equation depends
only on coefficients of $A$ of smaller degree: both factors in the product
have positive order. Thus it defines a unique integral series, and the same
uniqueness holds over $\mathbb Z/4\mathbb Z$.
Consider the rational series
\[
 F=x+x^2+\frac{x^3(1+x)^3}{1-x^4}
   =\frac{x+x^2+x^3+3x^4+2x^5}{1-x^4},\qquad B=x+2x^2.
\]
Clearing denominators, or simply substituting these rational expressions,
gives the following identities over $\mathbb Z/4\mathbb Z$:
\[
 F(F(x))=B,\qquad F=x+B\,F(B),\qquad
 F-xF(B)=x+\frac{2x^3}{1-x}.                              \tag{R4}
\]
Every denominator used here has constant term $1$. Hence these are valid
formal-series identities, not evaluations away from poles.
The first two identities show that $F$ satisfies the defining equation
modulo $4$, because $F^{\circ3}=F(B)$; uniqueness gives $A\equiv F$.
Its displayed expansion and the third identity prove all the assertions.
The accompanying exact verifier independently clears the three denominators
in (R4) and checks every coefficient of the resulting polynomial numerators.
\end{proof}

\label{round4:r4entry:A378575}
\begin{theorem}[The selected conjecture in \oeis{A378575}]
For the unique series $A=x+xA^{\circ5}$, every coefficient of positive degree
is $1$ modulo $4$.
\end{theorem}
\begin{proof}
The leading factor $x$ makes the equation triangular in its coefficients.
For $U=x/(1-x)$ one has $U^{\circ j}=x/(1-jx)$ by induction.
Therefore $U^{\circ5}=U$ modulo $4$, and
$x+xU^{\circ5}=x+xU=U$ modulo $4$.
Uniqueness over $\mathbb Z/4\mathbb Z$ gives $A\equiv U$.
The same proof works for iteration counts $j\equiv1\pmod4$, but these
unlisted specializations add no entry to the count.
\end{proof}

\subsection{Removing signs and odd factors}
\label{round4:r4entry:A378578}
\begin{theorem}[The selected conjecture in \oeis{A378578}]
Starting with $a_0=1$, obtain $a_n$ by removing all factors of $3$ from
$[x^{n-1}]A(x)^3$. Then, for every $n\geq0$,
\[
                    a_n\equiv\frac1{2n+1}\binom{3n}{n}\pmod2.
\]
\end{theorem}
\begin{proof}
The coefficient used to define $a_n$ is positive and involves only preceding
terms, so the definition gives unique positive integers.
Removing a power of $3$ preserves parity. Thus the reduced series satisfies
$A=1+xA^3$ over $\mathbb F_2$.
The unique integral solution $T=1+xT^3$ has coefficients
$[x^n]T=\binom{3n}{n}/(2n+1)$ by Lagrange inversion.
Reduction and uniqueness prove the assertion. By Lucas's theorem this
coefficient is odd precisely when the binary expansion of $n$ contains no
adjacent ones: this is the no-carry condition for $n+2n$.
The additional empirical run-length descriptions are not separately counted.
\end{proof}

\label{round4:r4entry:A380709}
\begin{theorem}[The selected conjecture in \oeis{A380709}]
For $A=1+x\,|1/A|^3$, where absolute value is taken coefficient by coefficient,
\[
                         a_n\equiv\binom{4n-1}{n}\pmod2
                         \quad(n\geq0).
\]
\end{theorem}
\begin{proof}
The equation is triangular and integral. Taking an absolute value preserves
parity, so $A=1+xA^{-3}$ modulo $2$.
The integral solution $B=1+xB^{-3}$ has, for $n\geq1$,
\[
 [x^n]B=\frac1n\binom{-3n}{n-1}
       =\frac{(-1)^{n-1}}{4n-1}\binom{4n-1}{n}.
\]
The denominator $4n-1$ is odd, giving the required parity. Both constant
terms are $1$, with $\binom{-1}{0}=1$.
\end{proof}

\label{round4:r4entry:A380711}
\begin{theorem}[Both parity descriptions in \oeis{A380711}]
For $A=1+xA\,|A^{-3}|$, the coefficients satisfy
\[
 a_n\equiv\frac1{3n-1}\binom{3n-1}{n}\pmod2\quad(n\geq0).
\]
For $n>0$, $a_n$ is odd exactly when $n=2j+1$ and the binary expansion of
$j$ has no adjacent ones.
\end{theorem}
\begin{proof}
Reduction of the defining equation gives $A=1+xA^{-2}$ over $\mathbb F_2$.
For the integral solution $B=1+xB^{-2}$, Lagrange inversion gives
\[
 [x^n]B=\frac{(-1)^{n-1}}{3n-1}\binom{3n-1}{n}\quad(n\geq1).
\]
This is an identity of integers, so reducing it modulo $2$ does not require
inverting the possibly even denominator $3n-1$. At $n=0$ the displayed
quotient is $-1$, congruent to the constant coefficient $1$.

For the second description let $T=1+tT^3$ over $\mathbb F_2$.
Then $F=1+xT(x^2)$ satisfies $F=1+xF^{-2}$: indeed
$T(x^2)(1+x^2T(x^2)^2)=1$.
By uniqueness $A=F$ modulo $2$. The ternary Catalan parity condition proved
above identifies precisely the indicated indices. The entry's separate
analytic radius and asymptotic observations are outside this target.
\end{proof}

\subsection{An algebraic parity identity}
\label{round4:r4entry:A369083}
\begin{theorem}[The selected conjecture in \oeis{A369083}]
For $A=1+x(5A(x)^2-A(-x)^2)/4$ one has
\[
                         a_n\equiv\binom{4n+3}{n}\pmod2
                         \quad(n\geq0).
\]
\end{theorem}
\begin{proof}
Write $A=E+O$, separating the even and odd degrees.
The defining equation is equivalent to the integral triangular system
\[
                E=1+3xEO,\qquad O=x(E^2+O^2).
\]
It proves integrality without dividing by $4$ in a finite field.
Modulo $2$ we have $O=xA^2$ and hence
\[
                 A=1+xA^2+x^2A^3+x^3A^4.                \tag{R5}
\]
The constant term $1$ solution of (R5) is unique.
Let $T=1+xT^4$ and put $B=T^3$. In characteristic $2$,
$1+xB=T^{-1}$, so $B(1+xB)^3=1$, which is (R5).
Lagrange inversion gives
\[
                  [x^n]T^3=\frac3{4n+3}\binom{4n+3}{n}.
\]
Both $3$ and $4n+3$ are odd. Therefore these coefficients have exactly the
parity asserted.
\end{proof}

\subsection{Two implicit pentagonal series}
\label{round4:r4entry:A361552}\label{round4:r4entry:A359914}
Let $g_j=j(3j+1)/2$, and let
\[
              P(x)=\prod_{k\geq1}(1-x^k)
                   =\sum_{j\in\mathbb Z}(-1)^j x^{g_j}.
\]
The last equality is Euler's pentagonal identity; $1/P(x)=\sum p(n)x^n$
is the partition generating function.
\begin{theorem}[The selected parity claims in \oeis{A361552} and \oeis{A359914}]
Let $A$ and $B$, respectively, be the constant-term-$1$ solutions
\[
 2x=\sum_{j\in\mathbb Z}x^{g_j}
       (A^{3j}-A^{-3j-1}),                               \tag{R6}
\]
\[
 2=\sum_{j\in\mathbb Z}(-1)^jx^{j(3j-1)/2}
       B^{3j}(1+x^jB).                                  \tag{R7}
\]
Then
\[
 A\equiv1+\frac{2x}{P(x)}\pmod4,\qquad
 B\equiv1+2\left(\frac1{P(x)}-1\right)\pmod4.              \tag{R8}
\]
In particular $[x^n]A/2\equiv p(n-1)\pmod2$ for $n\geq1$, and
$[x^n]B/2\equiv p(n)\pmod2$ for $n\geq1$.
\end{theorem}
\begin{proof}
All sums are well defined degree by degree. At index $j=0$ the right sides
are $1-A^{-1}$ and $1+B$, respectively; every other term has a positive
power of $x$ after the two summands in (R7) are separated. Consequently
the equations determine unique integral coefficients, with leading
coefficient $1$ on each unknown coefficient. The constant series $1$
solves both reduced equations modulo $2$, proving that $A=1+2H$ and
$B=1+2J$ with $H,J\in x\mathbb Z[[x]]$.

For every integer exponent $r$, $(1+2H)^r\equiv1+2rH\pmod4$.
Thus (R6) reduces to
$2x\equiv2H\sum_jx^{g_j}\pmod4$.
After dividing this integer coefficient congruence by $2$ and reducing
modulo $2$, the signs in $P$ disappear and $H=x/P$.

In (R7), replace $j$ by $-j$ in its first summand to write the entire sum as
\[
          \sum_{j\in\mathbb Z}(-1)^jx^{g_j}
                         (B^{-3j}+B^{3j+1}).
\]
The parenthesis is $2+2J$ modulo $4$. Thus
$1=(1+J)P$ modulo $2$, giving $J=1/P-1$.
This proves (R8) and both parity claims.
\end{proof}
The cross-entry relation $[x^n]A\equiv[x^{n-1}]B\pmod4$ follows for
$n>1$ and adds no count. The reverse comparison printed in A359914
requires $n\geq1$: its stated $n=0$ endpoint compares $1$ with $2$.
This endpoint correction is recorded but is not the assertion selected
for counting at A359914.

\subsection{An elementary exponential congruence}
\label{round4:r4entry:A380056}
\begin{theorem}[Both selected congruences in \oeis{A380056}]
If
\[
            \frac{e^t-1}{\cos(2t)}=\sum_{n\geq1}a_n\frac{t^n}{n!},
\]
then $a_n\equiv1\pmod4$ for all $n\geq1$, and $5\mid a_{4r}$ for all
$r\geq1$.
\end{theorem}
\begin{proof}
Let $h_j$ be the $j$th derivative of $\sec(2t)$ at zero.
Then $h_0=1$, all odd $h_j$ vanish, and
$h_{2k}=(-4)^kE_{2k}$.
Consequently
\[
                   a_n=\sum_{j=0}^{n-1}\binom nj h_j.     \tag{R9}
\]
All terms of this sum except $j=0$ vanish modulo $4$.
Modulo $5$, $h_j\equiv E_j$, because $-4\equiv1$ and the odd derivatives
vanish. Lemma~\ref{round4:r4:Euler-moments} with $M=5$ changes (R9) into
\[
 a_n\equiv\sum_{j=0}^4(-1)^j
           \bigl((2j+2)^n-(2j+1)^n\bigr)\pmod5.
\]
For positive $n$ divisible by $4$, Fermat's theorem makes each nonzero
base contribute $1$. Only the bases $2j+1\equiv0,-1$ matter: they occur
at $j=2,4$, both with positive sign, and contribute $1$ and $-1$.
Their sum is zero.
\end{proof}

\subsection{Cubing an integral exponential series}
\label{round4:r4entry:A368630}\label{round4:r4entry:A368631}
\begin{lemma}\label{round4:r4:Hurwitz}
Let $B=\sum_{n\geq0}b_nx^n/n!$ have integral derivatives and $b_0=1$.
For a prime $p$, every positive-degree derivative of $B^p$ is divisible
by $p$. Consequently every derivative of $\exp(xB^p)$ is $1$ modulo $p$.
\end{lemma}
\begin{proof}
The $n$th derivative of $B^p$ is the multinomial sum over ordered
$p$-tuples $(n_1,\ldots,n_p)$ with total $n$. Rotation permutes these tuples.
Nonconstant tuples have orbits of size $p$ and equal summands. A fixed tuple
has $n_1=\cdots=n_p=k$; if $n>0$, then $k>0$ and its multinomial factor
$(pk)!/(k!)^p$ has valuation $s_p(k)\geq1$, by Legendre's formula.
This proves the first assertion.
The derivatives $c_n$ of $xB^p$ satisfy $c_1=1$ and $p\mid c_n$ for
$n>1$. The derivatives of its exponential are the complete Bell polynomials,
which have integer coefficients. Their reduction is
$B_n(1,0,0,\ldots)=1$, proving the second assertion.
\end{proof}

\begin{theorem}[The selected mod-$3$ claims in \oeis{A368630}, \oeis{A368631}]
The derivatives $a_n$ of either of the series
\[
 \begin{array}{ll}
 \text{A368630:}& A(x/A(x)^2)=\exp(xA(x)),\\
 \text{A368631:}& A(x/A(x))=\exp(xA(x)^2)
 \end{array}
\]
are $1$ modulo $3$ for every $n\geq0$.
\end{theorem}
\begin{proof}
In the two cases set $r=2$ and $r=1$, respectively. Reversion of
$x/A(x)^r$ gives the equivalent pair of formal identities
\[
                         A=\exp(xB^3),\qquad
                         B=A(xB^r),\qquad A(0)=B(0)=1.   \tag{R10}
\]
For example, putting $z=xB(x)^r$ makes
$z/A(z)^r=x$ and transforms the original equation into the first identity;
the converse substitution is identical.
The equations determine the derivatives recursively: $a_n$ in the first
identity uses only $b_0,\ldots,b_{n-1}$, and $b_n$ in the second uses
$a_0,\ldots,a_n$ and preceding $b$'s. Products, integral powers and
composition of integral exponential series with zero constant term in
the inner series preserve integral derivatives, by the integer Bell
polynomial formulas. Thus (R10) proves integrality inductively.
Apply Lemma~\ref{round4:r4:Hurwitz} with $p=3$.
The two separate parity conjectures on each page have not been selected.
\end{proof}


\clearpage
\section{Elementary arithmetic, exact counterexamples, and a square-sum graph}
\label{round4:r4:arithmetic}

This section addresses nine explicitly selected assertions on nine distinct
OEIS entries.  A result about a comment is counted only at the entry containing
that comment.  In particular, the result for \oeis{A320192} concerns its
reciprocal-square approximation comment, and the result for \oeis{A228884}
concerns part (ii) of its conjecture.  We do not count already established
claims at those entries as new results.  Computations accompanying this section
are exact checks of the arguments; they are not substitutes for the proofs.

\subsection{A counterexample to the proposed recurrence in A190363}
\label{round4:r4entry:A190363}

The definition and its equivalent formula on \oeis{A190363} give
\[
 a(n)=2n+\left\lfloor\frac{n\sqrt5}{2}\right\rfloor
          +\left\lfloor\frac n4\right\rfloor,\qquad n\geq1.
\]
The recurrence proposed by Harvey P. Dale on January 28, 2025 has its only
nonzero coefficients at positions $4,17,21$, with values $1,1,-1$.

\begin{theorem}
The proposed recurrence
\[
                 a(n)=a(n-4)+a(n-17)-a(n-21)
\]
fails at $n=161$.  It cannot hold for all sufficiently large $n$ either.
\end{theorem}
\begin{proof}
For each of the following pairs $(n,b)$ the inequalities
$4b^2\leq5n^2<4(b+1)^2$ hold:
\[
 (161,180),\quad(157,175),\quad(144,160),\quad(140,156).
\]
These are integer inequalities and certify all four irrational floors exactly.
Consequently
\[
 a(161)=542,\quad a(157)=528,\quad a(144)=484,\quad a(140)=471.
\]
The recurrence predicts $528+484-471=541$, not $542$.

To rule out an eventual version, suppose the recurrence holds after some
index and put $b_n=a(n)-a(n-4)$.  Then $b_n=b_{n-17}$ eventually.
Since $4$ and $17$ are coprime, the seventeen indices $n-4j$
($0\leq j\leq16$) run through all residue classes modulo $17$.
Their $b$-values therefore have a constant integer sum $B$ for all
sufficiently large $n$.  Telescoping gives
\[
                    a(n)-a(n-68)=B.
\]
Consequently $\lim_{n\to\infty}a(n)/n=B/68$ is rational.
The defining floor formula instead gives the irrational limit
$2+\sqrt5/2+1/4$, a contradiction.
\end{proof}

This is a counterexample to the asserted universal recurrence.  It is not a
claim based on a numerical fit.  The accompanying verifier additionally checks
that $161$ is the first failure among the permissible indices $n\geq22$.

\subsection{Cancellation in a divisor mean: A336838}
\label{round4:r4entry:A336838}

Let $p$ be prime and let $q$ be the next prime after $p$.
The prime shift defining \oeis{A003961} sends $p^e$ to $q^e$.
Thus \oeis{A336838}, the reduced numerator of the mean of the divisors of
that shifted number, satisfies
\[
 a(p^e)=\frac{S}{\gcd(S,e+1)},\qquad S=1+q+\cdots+q^e.
\]

\begin{theorem}[The selected conjecture in \oeis{A336838}]
If $e\geq1$ is odd, then $(q+1)/2$ divides $a(p^e)$.
\end{theorem}
\begin{proof}
Write $e+1=2h$.  Pairing successive terms gives
\[
                   S=(q+1)\sum_{j=0}^{h-1}(q^2)^j.
\]
We use the elementary valuation fact
\[
 v_\ell(1+u+\cdots+u^{h-1})=v_\ell(h)                 \tag{A1}
\]
when $\ell$ is odd and $u\equiv1\pmod\ell$, or when $\ell=2$ and
$u\equiv1\pmod4$.  For completeness, factor $h=\ell^t r$ with
$\ell\nmid r$.  The geometric factor of length $r$ is a unit modulo
$\ell$.  In successively enlarging a length by a factor $\ell$, the new
factor is $1+v+\cdots+v^{\ell-1}$ with $v$ satisfying the same congruence.
The binomial expansion at $v=1$ shows this factor has valuation exactly one:
for odd $\ell$ the constant term is $\ell$ and all remaining terms are
divisible by $\ell^2$; for $\ell=2$ it is $1+v\equiv2\pmod4$.
This proves (A1).

If $\ell\mid q+1$ is odd, (A1) with $u=q^2$ gives
\[
 v_\ell(S)=v_\ell(q+1)+v_\ell(e+1).
\]
Since $q$ is odd, $q^2\equiv1\pmod8$, so it also gives
\[
 v_2(S)=v_2(q+1)+v_2(e+1)-1.
\]
Cancelling $\gcd(S,e+1)$ can remove at most $v_\ell(e+1)$ factors of
each prime $\ell$.  The reduced numerator therefore retains at least
$v_\ell(q+1)$ for odd $\ell\mid q+1$, and at least $v_2(q+1)-1$ for
$\ell=2$.  These are exactly the exponents required by $(q+1)/2$.
\end{proof}

\subsection{The closest reciprocal-square partial sum: A320192}
\label{round4:r4entry:A320192}

Write $H_m^{(2)}=\sum_{k=1}^m k^{-2}$, with $H_0^{(2)}=0$.
The assertion selected here is Andrew Howroyd's November 24, 2018
\emph{comment} in \oeis{A320192}; the separate formula for its main sequence,
which involves finite Euler products, is outside this result's scope.

\begin{theorem}
For every integer $n\geq1$, the unique partial sum $H_m^{(2)}$ closest to
$\zeta(2)-1/n$ is $H_n^{(2)}$.
\end{theorem}
\begin{proof}
Put $R_n=\sum_{k=n+1}^{\infty}k^{-2}$.
The strict integral comparison for the decreasing function $x^{-2}$ gives
$R_n<\int_n^\infty x^{-2}\,dx=1/n$.
On each unit interval this function is strictly convex, so its integral is
strictly smaller than the trapezoidal area.  Summing these inequalities gives
\[
 \frac1n=\int_n^\infty x^{-2}\,dx
       <\frac1{2n^2}+R_n.
\]
Hence
\[
 0<H_n^{(2)}-\left(\zeta(2)-\frac1n\right)
       =\frac1n-R_n<\frac1{2n^2}.                     \tag{A2}
\]
Since $H_n^{(2)}-H_{n-1}^{(2)}=1/n^2$, the target lies strictly between
the midpoint of these two sums and $H_n^{(2)}$.  Thus $H_n^{(2)}$ is
closer than $H_{n-1}^{(2)}$.  Every earlier sum is still farther below,
and every later sum is still farther above, proving uniqueness.
\end{proof}

\subsection{An infinite explicit family for A111205}
\label{round4:r4entry:A111205}

The entries of \oeis{A111205} are nonsquare composite integers $k$ for which
$\gcd(k,2\lfloor\sqrt{k}\rfloor+1)>1$.  The second argument is the
difference between the two consecutive squares surrounding $k$.

\begin{theorem}
There are infinitely many such integers.  In fact every
\[
                  k_t=(3t+1)^2+2=3(3t^2+2t+1),\qquad t\geq1,
\]
belongs to \oeis{A111205}.
\end{theorem}
\begin{proof}
Set $r=3t+1\geq4$.  Then
$r^2<r^2+2<(r+1)^2$, so $k_t$ is nonsquare and its integer square root is
$r$.  Its displayed factorization has both factors greater than one, so it is
composite.  Both $k_t$ and $2r+1=6t+3$ are divisible by $3$.
Finally, $k_t$ is strictly increasing with $t$.
\end{proof}

\subsection{The signed representation conjecture in A025060}
\label{round4:r4entry:A025060}

The original sequence uses positive $i<j<k$.  We address Jon Perry's
April 21, 2013 conjecture in \oeis{A025060} concerning the variant that
allows negative integers while retaining distinctness and excluding zero.

\begin{theorem}
Every integer $N$ has a representation
$N=ij+jk+ki$ by three distinct nonzero integers $i<j<k$.
\end{theorem}
\begin{proof}
Choose $a=|N|+2$ and put
\[
              i=-a,\qquad j=a+1,\qquad k=N+a(a+1).
\]
Then $i+j=1$ and $ij=-a(a+1)$, giving
$ij+(i+j)k=N$.
Also $i<0<j$, and
\[
             k-j=N+a^2-1\geq-|N|+(|N|+2)^2-1>0.
\]
The three integers have all the required properties.
\end{proof}

\subsection{Prime divisors of a gcd determinant: A228884(ii)}
\label{round4:r4entry:A228884}

Let $D_n=\det(\gcd(i-j,n))_{0\leq i,j<n}$.
We prove only conjecture (ii) in \oeis{A228884}.  Fedor Petrov's February
2021 answer to Zhi-Wei Sun's MathOverflow question proves positivity and
gives an eigenvalue formula that also implies the divisibility assertion
in part (i); these are not counted here.  We recall that Fourier formula
with attribution and give its short derivation so that this proof is complete.
The prior answer is at
\url{https://mathoverflow.net/questions/384719}.

\begin{theorem}
If $n$ is composite, every prime divisor of $D_n$ is smaller than $n$.
\end{theorem}
\begin{proof}
The identity $\gcd(t,n)=\sum_{d\mid\gcd(t,n)}\varphi(d)$ and the finite
geometric sum show that the eigenvalue at the character indexed by
$k\in\{0,\ldots,n-1\}$ is
\[
 \lambda_n(k)=\sum_{d\mid\gcd(k,n)}d\varphi(n/d).     \tag{A3}
\]
Indeed, the sum of the character over multiples of a divisor $e$ of $n$
is $n/e$ if $(n/e)\mid k$, and zero otherwise; substitute $d=n/e$.
The Fourier characters are a basis, so $D_n=\prod_{k=0}^{n-1}\lambda_n(k)$.
All these eigenvalues are positive integers.

For $n=\prod p^a$, the divisor sum in (A3) factors over prime powers.
At one such power, put $j=\min(v_p(k),a)$, interpreting $j=a$ when
$k\equiv0\pmod{p^a}$.  Direct summation gives
\[
 \lambda_{p^a}(k)=
 \begin{cases}
  (j+1)p^{a-1}(p-1),&j<a,\\
  p^{a-1}\bigl((a+1)p-a\bigr),&j=a.
 \end{cases}                                                   \tag{A4}
\]
Every prime factor of a value in the first line divides one of $p$, $p-1$,
or $j+1\leq a$; each is smaller than the composite number $n$.

It remains to consider $B=(a+1)p-a$.
If $a=1$, the fact that $n$ is composite and $p$ occurs to exponent one
implies $n\geq2p>2p-1=B$.
If $a\geq2$, then $B\leq p^a$.  At $a=2$ the difference is
$p^2-(3p-2)=(p-1)(p-2)\geq0$; increasing $a$ increases $p^a$ by
$(p-1)p^a$ and increases $B$ by only $p-1$.
Equality is possible only for $p=2,a=2$.
Thus $B<n$ except possibly when $n=4$ and $B=4$, whose only prime
factor is $2<4$.
All prime factors of all eigenvalues in (A3) are therefore smaller than
$n$, and the same holds for their product.
\end{proof}

\subsection{A coprimality determinant: A228885}
\label{round4:r4entry:A228885}

Let $M_n(i,j)$ be $1$ if $\gcd(i+j,n)=1$ and $0$ otherwise, with
rows and columns indexed by residues modulo $n$.  Let $a(n)=\det M_n$.
This is the definition of \oeis{A228885}.

\begin{theorem}
For every $n\geq1$,
\[
 a(n)=(-1)^{\lfloor n/2\rfloor}\mu(n)^2
                   \prod_{k=1}^n\varphi(\gcd(n,k))
     =(-1)^{\lfloor n/2\rfloor}\mu(n)^2
                   \prod_{d\mid n}\varphi(d)^{\varphi(n/d)}.     \tag{A5}
\]
In particular, the conjectured sign
$(-1)^{n(n-1)/2}a(n)>0$ holds whenever $n$ is squarefree.
\end{theorem}
\begin{proof}
If $p^2\mid n$, adding $n/p$ to a row index does not change whether
$i+j$ is divisible by any prime divisor of $n$.  Two different rows are
equal and the determinant is zero.  This observation was already recorded
in the entry; it agrees with the factor $\mu(n)^2$ in (A5).

Now assume $n$ is squarefree and set
$C_n(i,j)=1_{\gcd(i-j,n)=1}$.  The Chinese remainder theorem simultaneously
reindexes rows and columns and identifies
\[
                      C_n\cong\bigotimes_{p\mid n}(J_p-I_p).
\]
The $p\times p$ matrix $J_p-I_p$ has eigenvalue $p-1$ once and eigenvalue
$-1$ with multiplicity $p-1$.  Consequently
\[
       |\det C_n|=\prod_{p\mid n}(p-1)^{n/p}.          \tag{A6}
\]
The sign is positive for odd $n$.  For even squarefree $n$, the factor
with $p=2$ contributes sign $(-1)^{n/2}=-1$, and all odd-prime factors
contribute positive sign.

The column permutation $j\mapsto-j$ changes $C_n$ into $M_n$.
It has $\lfloor(n-1)/2\rfloor$ transpositions.  Combining its sign with
the preceding sign of $C_n$ gives $(-1)^{\lfloor n/2\rfloor}$ in
both the odd and the even squarefree cases.
For squarefree $n$, the factor $p-1$ occurs in
$\varphi(\gcd(n,k))$ precisely when $p\mid k$, which happens $n/p$ times
for $1\leq k\leq n$.  Thus its product over $k$ is exactly (A6).
Grouping the same product according to $d=\gcd(n,k)$ gives its divisor
form in (A5), because there are $\varphi(n/d)$ such values of $k$.
The case $n=1$ is immediate.

Finally $n(n-1)/2$ and $\lfloor n/2\rfloor$ have the same parity, and
the product in (A6) is positive.  This proves the sign assertion.
\end{proof}

The two product formulas in (A5), listed as conjectures by Ridouane Oudra
on March 13, 2025, and the older sign assertion are treated as a single
entry-level result.  They do not contribute three counts.

\subsection{Consecutive divisor gaps: A258409}
\label{round4:r4entry:A258409}

For $n\geq2$, let $1=d_1<d_2<\cdots<d_t=n$ be its positive divisors.
The definition of \oeis{A258409} is $a(n)=\gcd(d_1-1,\ldots,d_t-1)$.

\begin{theorem}[The October 16, 2019 conjecture in \oeis{A258409}]
\[
                 a(n)=\gcd(d_2-d_1,d_3-d_2,\ldots,d_t-d_{t-1}).
\]
\end{theorem}
\begin{proof}
Any integer dividing every $d_i-1$ also divides their successive
differences $d_{i+1}-d_i$.  Conversely, if it divides all those
differences, it divides every telescoping sum
\[
               d_j-1=d_j-d_1=\sum_{i=1}^{j-1}(d_{i+1}-d_i).
\]
The two lists have exactly the same common divisors, hence the same
positive greatest common divisor.
\end{proof}

The earlier prime-divisor characterization in this entry was proved by
Robert Israel on September 14, 2016.  It is not part of the present count.
The displayed gap statement is the numerical content of the later
composition-of-sequences conjecture; no new count is assigned to its
auxiliary sequence identifiers.

\subsection{A square-root error term for A176615}
\label{round4:r4entry:A176615}

Let $a(n)$ count pairs $1\leq i<j\leq n$ for which $i+j$ is a square.
The entry \oeis{A176615} gives the leading asymptotic and describes an
$O(\sqrt n)$ error as probable.  We prove that stronger error estimate.

\begin{theorem}
For every integer $n\geq1$,
\[
 \left|a(n)-\frac{2\sqrt2-2}{3}n^{3/2}\right|<3\sqrt n.
\]
In particular the error term is $O(\sqrt n)$.
\end{theorem}
\begin{proof}
For a fixed upper endpoint $j$, the possible squares lie strictly above
$j$ and at most $2j-1$.  Therefore
\[
 a(n)=T(n)-U(n),\quad
 T(n)=\sum_{j=1}^n\lfloor\sqrt{2j-1}\rfloor,\quad
 U(n)=\sum_{j=1}^n\lfloor\sqrt j\rfloor.
\]
Put $L=\lfloor\sqrt n\rfloor$ and $M=\lfloor\sqrt{2n}\rfloor$.
Counting each threshold $r^2$ gives the exact identities
\[
 U(n)=L(n+1)-\frac{L(L+1)(2L+1)}6
     =Ln-\frac{L^3}3-\frac{L^2}2+\frac{5L}6,
\]
\[
 \begin{split}
 T(n)&=Mn-\sum_{r=1}^M\left\lfloor\frac{r^2}2\right\rfloor\\
     &=Mn-\frac{M^3}6-\frac{M^2}4+\frac M6+\varepsilon_M,
 \qquad \varepsilon_M=\frac{1-(-1)^M}{8}\in\{0,1/4\}.
 \end{split}                                                  \tag{A7}
\]
If $M^2=2n$, its threshold contributes zero, so the formula remains
valid with this choice of $M$.

Write $x=\sqrt n$, $y=\sqrt{2n}$, $L=x-\alpha$ and $M=y-\beta$, where
$0\leq\alpha,\beta<1$.  Expanding (A7) gives
\[
 U(n)=\frac23 n^{3/2}-\frac n2
       +x\left(\frac56+\alpha-\alpha^2\right)+u(\alpha),
 \quad u(t)=\frac{t^3}3-\frac{t^2}2-\frac{5t}6,
\]
\[
 T(n)=\frac{2\sqrt2}{3}n^{3/2}-\frac n2
       +y\left(\frac16+\frac{\beta-\beta^2}2\right)
       +v(\beta)+\varepsilon_M,
 \quad v(t)=\frac{t^3}6-\frac{t^2}4-\frac t6.
\]
On $0\leq t\leq1$, we have $0\leq t-t^2\leq1/4$,
$-1\leq u(t)\leq0$ and $-1/4\leq v(t)\leq0$.
The latter two bounds follow, for example, by differentiating the two
polynomials, which are decreasing on this interval.
The two terms $-n/2$ cancel.  The triangle inequality now bounds the
remaining error by
\[
 \left(\frac{13}{12}+\frac{7\sqrt2}{24}\right)\sqrt n+\frac54
 \leq\left(\frac{13}{12}+\frac{7\sqrt2}{24}+\frac54\right)\sqrt n
 <3\sqrt n.
\]
This proves the bound uniformly for all $n\geq1$.
\end{proof}


\clearpage
\section{Thirty-five generating-function and recurrence packages}
\label{round4:sec:r4gf}
This section treats 35 distinct OEIS entries absent from the 308 selected
entries in rounds 1--3. The selected assertions were explicitly labeled
\emph{Conjecture} in the primary pages consulted on 20 October 2026.
Thirty-four packages prove the selected assertion with an explicit index
range. A243642 additionally requires a boundary correction: the source
does not specify a starting index; the first natural index gives residual
24, while the recurrence holds at every subsequent index. That correction
and its proof count as one package.

The defining formulas, established generating functions and existing
shorter recurrences are prior input. In particular, the existing
Kotesovec recurrences for A114464, A100095 and A000986 and the other Mathar
recurrence for A128057 are not presented as newly discovered results.
A102318 contributes only its first conjectured recurrence. A002026 is
used as a supporting reference, not a further counted problem. These are
entry-level selected assertions, often sharing the same differentiation
method; no claim of worldwide priority or external acceptance is made.

\subsection{Exact coefficient-extraction theorem}
For an OGF $F(x)=\sum_{n\ge0}a_nx^n$, extending undefined initial indices
by zero only for this formal-series notation, put $\theta=x\,d/dx$.
For polynomials $p_0(n),\ldots,p_r(n)$ define
\[
 \mathcal L_p=\sum_{j=0}^{r}x^j p_j(\theta+j).
\]
Directly on monomials,
\begin{equation}\label{round4:eq:r4gf-extraction}
 [x^n]\mathcal L_pF=\sum_{j=0}^{r}p_j(n)a_{n-j}.
\end{equation}
Consequently the exact identity $\mathcal L_pF=P(x)$, with $P$ a
polynomial, proves the displayed homogeneous recurrence at every
$n>\deg P$ for which all source indices are defined. No agreement of a
finite numerical prefix is used to infer the identity.

All 31 algebraic cases below are certified over a rational function
field. Write $x=\xi(t)$ and
\[
 F(\xi(t))=\sum_{i=0}^{s}R_i(t)Y_i(t),\qquad
 Y_i(t)^2=D_i(t),\quad D_0=Y_0=1.
\]
For cases without a parameter, $t=x$ and $\xi(t)=t$. All fractions are
interpreted in the formal Laurent field; any apparent poles cancel in
the stated analytic branch of $F$. Define
\[
 \mathcal T_i R=\frac{\xi(t)}{\xi'(t)}
       \left(R'(t)+\frac{D_i'(t)}{2D_i(t)}R(t)\right),\qquad
 c_m(t)=\sum_{j=0}^{r}\xi(t)^j[n^m]p_j(n+j).
\]
The chain rule gives $\theta(RY_i)=(\mathcal T_iR)Y_i$, and hence the
finite certificate consists of the rational identities
\begin{equation}\label{round4:eq:r4gf-rational-cert}
 \sum_m c_m\mathcal T_0^mR_0=P(\xi(t)),\qquad
 \sum_m c_m\mathcal T_i^mR_i=0\quad(1\le i\le s).
\end{equation}
Every coefficient $p_j$, every $R_i,D_i,\xi$ and the complete polynomial
$P$ are supplied below. Substitution in \eqref{round4:eq:r4gf-rational-cert}
and clearing denominators yields zero polynomials over $\mathbb Q$.
Thus this is an exact finite certificate with a specified verification
rule, not a numerical guess. The accompanying Python verifier repeats
these rational operations, compares all recurrence coefficients with the
retained source wording, and independently checks initial coefficients
and branch choices.

For clarity, A162481 gives a small fully expanded instance. Its Euler
operator is $c_1\theta+c_0$, where
\[
 c_1=1-8x+10x^2-4x^3+x^4,\quad
 c_0=1-6x+14x^2-4x^3+x^4.
\]
Here $R_0=1/(2x)$, $R_1=-1/(2x)$ and
$D_1=1-4x/(1-x)^3$. The two residuals reduce respectively to
\[
 \frac{c_0-c_1}{2x}=1+2x,\qquad
 -\frac1{2x}\left(c_0-c_1+
                  \frac{c_1xD_1'}{2D_1}\right)=0.
\]
The second equality follows on inserting $D_1$, or from
$c_1xD_1'+2D_1(c_0-c_1)=0$. The remaining algebraic certificates use
exactly the same finite rule, with the supplied data.

\subsection{Generating functions derived from the definitions}
The entries defined by a Catalan transform use
$C(z)=(1-\sqrt{1-4z})/(2z)$. For example,
$(1-x)^{-3}C(x/(1-x)^3)$ immediately gives the A162481 seed.
The three peakless-Motzkin statistics A171851, A171847 and A171849 use
the established seed $g=1+xg+x^2g(g-1)$, so
\[
 g=\frac{1-x+x^2-\sqrt{1-2x-x^2-2x^3+x^4}}{2x^2}.
\]
Substitution in their respective established statistic OGFs gives the
$R_i,D_i$ below. A191792 uses the branch $C(0)=1$ of
$x^2(1+x^2)C^2-(1+x^2+x^4)C+1+x^2=0$, followed by $F=C/(1-xC)$.
A182899 similarly uses $c=1+xc+x^2c+x^3c^2$.

For A176610, A176611, A176612, A176605 and A176607 a direct derivation
from the defining convolution avoids any inference from data. If
$a_0=1$, $a_1=m$ and
\[
 a_{n+1}=\sum_{p=0}^{n}a_pa_{n-p}+k(n+1)+\ell\qquad(n\ge1),
\]
then summing after multiplication by $x^{n+1}$ gives
\[
 F=xF^2+1+(m-1)x+
       \frac{(k+\ell)x^2}{1-x}+\frac{kx^2}{(1-x)^2}.
\]
The parameters $(m,k,\ell)$ are respectively
$(0,1,1),(1,1,1),(2,1,1),(1,0,1),(4,0,1)$.
The solution with constant term one is $(1-\sqrt D)/(2x)$, with
$D=1-4x\{1+(m-1)x+(k+\ell)x^2/(1-x)+kx^2/(1-x)^2\}$.

For A072100 let
$B(x)=\sum_{k\ge0}\binom{k}{\lfloor k/2\rfloor}x^k$.
Splitting even and odd indices in the central-binomial expansion gives
\[
 B(x)=\frac{(1+2x)/\sqrt{1-4x^2}-1}{2x},\qquad
 F(x)=\frac{x}{1-x}+\frac{x^2B(x)}{1-x}.
\]
The second identity is exactly its defining partial sum, with offset one.
For A050168 its displayed binomial formula similarly gives
$F=(1+x)(\sqrt{(1+2x)/(1-2x)}-1)/(2x)$, with $a_0=1$.
For A186940 the entry's printed ``offset one'' OGF must be divided by
$x$ to obtain the zero-offset OGF used here. This adjustment is part
of the proof, and is checked independently against the source terms.

For A045743 use $t=T(x)$, $T(0)=1$, and
$xT^4=(1-T)(1-3T+T^2)$; thus
$\xi(t)=(1-t)(1-3t+t^2)/t^4$ and $F=x(T-1)$.
For A308726, A285195 and A243642 take $t=\sqrt{1-4x}$,
so $t(0)=1$ and $\xi(t)=(1-t^2)/4$.
The second radical has constant term one in the first two cases and
constant term two in A243642.
For A243816 set $u=\sqrt{1+4x^2}$,
$t=(u-1)/(2x)$; then $t(0)=0$, $x=t/(1-t^2)$,
$u=(1+t^2)/(1-t^2)$ and the outer radical is $\sqrt{1-4t}$
with constant term one. These local branches have nonzero $\xi'(t_0)$,
so the formal chain rule is legitimate.

\subsection{The thirty-one algebraic certificates}
In each item $D_0=1$ is implicit. Unless a parameter and base point
are stated, the variable is $x$, $\xi=x$, and each radical has constant
term one. The selected recurrence is
$\sum_{j=0}^r p_j(n)a_{n-j}=0$ in the stated range.

\subsubsection{A210496: order 8, Oct 08 2016}
\label{round4:r4gf:A210496}\label{round4:r4entry:A210496}
Primary source: \url{https://oeis.org/A210496}. The generating function is established input, credited to the entry and its link to Jelinek, Mansour and Shattuck (2013), Theorem 4.27. It is not claimed as a new generating function.
\begin{align*}
p_{0}(n)&=n + 2\\
p_{1}(n)&=- 2 \left(3 n + 4\right)\\
p_{2}(n)&=4 \left(2 n + 3\right)\\
p_{3}(n)&=13 n - 40\\
p_{4}(n)&=- 5 \left(7 n - 18\right)\\
p_{5}(n)&=6 \left(2 n - 1\right)\\
p_{6}(n)&=22 \left(n - 7\right)\\
p_{7}(n)&=134 - 19 n\\
p_{8}(n)&=2 \left(2 n - 15\right).
\end{align*}
The exact generating-function representation is $F=\sum_i R_i\sqrt{D_i}$, with
\begin{align*}
R_{0}&=- \frac{3 x^{5} - 15 x^{4} + 14 x^{3} - 2 x^{2} - 3 x + 1}{2 x^{2} \left(x - 1\right)^{2} \left(x^{2} - 3 x + 1\right)},\\
R_{1}&=- \frac{\left(x - 1\right) \left(x + 1\right)}{2 x^{2} \left(x^{2} - 3 x + 1\right)},\\
D_{1}&=\left(x - 1\right) \left(4 x^{2} + x - 1\right).
\end{align*}
For these data the complete rational certificate \eqref{round4:eq:r4gf-rational-cert} gives
\[ P(x)=3 x^{7} - 30 x^{6} + 16 x^{5} - 32 x^{4} + 8 x^{3} + 16 x^{2} - 11 x + 2.\]
Every nonconstant radical component is identically zero. By \eqref{round4:eq:r4gf-extraction}, the recurrence therefore holds for every $n\ge8$. The source offset is 0.
\subsubsection{A102318: order 5, Jun 08 2016}
\label{round4:r4gf:A102318}\label{round4:r4entry:A102318}
Primary source: \url{https://oeis.org/A102318}. The source contains two conjectured recurrences. Only the first, of order five, is selected here; the second is outside this package.
\begin{align*}
p_{0}(n)&=- \left(n - 1\right) \left(n + 1\right)\\
p_{1}(n)&=2 \left(5 n^{2} - 9 n + 1\right)\\
p_{2}(n)&=- 2 \left(15 n^{2} - 58 n + 49\right)\\
p_{3}(n)&=2 \left(10 n^{2} - 76 n + 123\right)\\
p_{4}(n)&=\left(n - 3\right) \left(31 n - 55\right)\\
p_{5}(n)&=- 30 \left(n - 4\right) \left(n - 3\right).
\end{align*}
The exact generating-function representation is $F=\sum_i R_i\sqrt{D_i}$, with
\begin{align*}
R_{0}&=\frac{1}{2 x},\\
R_{1}&=- \frac{1}{4 x},\\
R_{2}&=- \frac{1}{4 x},\\
D_{1}&=- \frac{3 x - 1}{x + 1}\\
D_{2}&=\frac{5 x - 1}{x - 1}.
\end{align*}
For these data the complete rational certificate \eqref{round4:eq:r4gf-rational-cert} gives
\[ P(x)=11 x^{2} - 6 x + 1.\]
Every nonconstant radical component is identically zero. By \eqref{round4:eq:r4gf-extraction}, the recurrence therefore holds for every $n\ge5$. The source offset is 0.
\subsubsection{A162481: order 4, Dec 11 2011}
\label{round4:r4gf:A162481}\label{round4:r4entry:A162481}
Primary source: \url{https://oeis.org/A162481}. The defining formula or generating function is established input from the cited primary entry. Only the selected, explicitly conjectural identity is counted.
\begin{align*}
p_{0}(n)&=n + 1\\
p_{1}(n)&=- 2 \left(4 n - 1\right)\\
p_{2}(n)&=2 \left(5 n - 3\right)\\
p_{3}(n)&=- 4 \left(n - 2\right)\\
p_{4}(n)&=n - 3.
\end{align*}
The exact generating-function representation is $F=\sum_i R_i\sqrt{D_i}$, with
\begin{align*}
R_{0}&=\frac{1}{2 x},\\
R_{1}&=- \frac{1}{2 x},\\
D_{1}&=\frac{x^{3} - 3 x^{2} + 7 x - 1}{\left(x - 1\right)^{3}}.
\end{align*}
For these data the complete rational certificate \eqref{round4:eq:r4gf-rational-cert} gives
\[ P(x)=2 x + 1.\]
Every nonconstant radical component is identically zero. By \eqref{round4:eq:r4gf-extraction}, the recurrence therefore holds for every $n\ge4$. The source offset is 0.
\subsubsection{A116391: order 5, Jan 23 2020}
\label{round4:r4gf:A116391}\label{round4:r4entry:A116391}
Primary source: \url{https://oeis.org/A116391}. The defining formula or generating function is established input from the cited primary entry. Only the selected, explicitly conjectural identity is counted.
\begin{align*}
p_{0}(n)&=n\\
p_{1}(n)&=n\\
p_{2}(n)&=- 3 \left(3 n - 4\right)\\
p_{3}(n)&=- 3 \left(3 n - 4\right)\\
p_{4}(n)&=20 \left(n - 3\right)\\
p_{5}(n)&=20 \left(n - 3\right).
\end{align*}
The exact generating-function representation is $F=\sum_i R_i\sqrt{D_i}$, with
\begin{align*}
R_{0}&=- \frac{x}{\left(x + 1\right) \left(5 x^{2} - 1\right)},\\
R_{1}&=- \frac{1}{\left(x + 1\right) \left(5 x^{2} - 1\right)},\\
D_{1}&=- \left(2 x - 1\right) \left(2 x + 1\right).
\end{align*}
For these data the complete rational certificate \eqref{round4:eq:r4gf-rational-cert} gives
\[ P(x)=x.\]
Every nonconstant radical component is identically zero. By \eqref{round4:eq:r4gf-extraction}, the recurrence therefore holds for every $n\ge5$. The source offset is 0.
\subsubsection{A114464: order 5, Nov 26 2012}
\label{round4:r4gf:A114464}\label{round4:r4entry:A114464}
Primary source: \url{https://oeis.org/A114464}. Kotesovec already recorded a different, shorter recurrence on 13 February 2014. That existing result and the established generating function are credited. The selected target is Mathar's still-labeled order-five recurrence of 26 November 2012.
\begin{align*}
p_{0}(n)&=n + 1\\
p_{1}(n)&=3 - 5 n\\
p_{2}(n)&=2 \left(3 n - 7\right)\\
p_{3}(n)&=- 2 \left(3 n - 11\right)\\
p_{4}(n)&=5 n - 27\\
p_{5}(n)&=7 - n.
\end{align*}
The exact generating-function representation is $F=\sum_i R_i\sqrt{D_i}$, with
\begin{align*}
R_{0}&=- \frac{x^{3} - 3 x^{2} + x - 1}{2 x},\\
R_{1}&=\frac{x - 1}{2 x},\\
D_{1}&=\left(x^{2} + 1\right) \left(x^{2} - 4 x + 1\right).
\end{align*}
For these data the complete rational certificate \eqref{round4:eq:r4gf-rational-cert} gives
\[ P(x)=- 3 x^{4} + 4 x^{3} - 6 x^{2} + 1.\]
Every nonconstant radical component is identically zero. By \eqref{round4:eq:r4gf-extraction}, the recurrence therefore holds for every $n\ge5$. The source offset is 0.
\subsubsection{A097332: order 3, Nov 26 2012}
\label{round4:r4gf:A097332}\label{round4:r4entry:A097332}
Primary source: \url{https://oeis.org/A097332}. The defining formula or generating function is established input from the cited primary entry. Only the selected, explicitly conjectural identity is counted.
\begin{align*}
p_{0}(n)&=n + 1\\
p_{1}(n)&=- 3 n\\
p_{2}(n)&=5 - n\\
p_{3}(n)&=3 \left(n - 2\right).
\end{align*}
The exact generating-function representation is $F=\sum_i R_i\sqrt{D_i}$, with
\begin{align*}
R_{0}&=- \frac{x + 1}{2 x \left(x - 1\right)},\\
R_{1}&=\frac{1}{2 x \left(x - 1\right)},\\
D_{1}&=- \left(x + 1\right) \left(3 x - 1\right).
\end{align*}
For these data the complete rational certificate \eqref{round4:eq:r4gf-rational-cert} gives
\[ P(x)=x + 1.\]
Every nonconstant radical component is identically zero. By \eqref{round4:eq:r4gf-extraction}, the recurrence therefore holds for every $n\ge3$. The source offset is 0.
\subsubsection{A128057: order 2, Dec 02 2014}
\label{round4:r4gf:A128057}\label{round4:r4entry:A128057}
Primary source: \url{https://oeis.org/A128057}. A different order-three D-finite recurrence is already recorded by Mathar on the same date. The selected target is only the still-labeled order-two recurrence.
\begin{align*}
p_{0}(n)&=n \left(5 n - 9\right)\\
p_{1}(n)&=4\\
p_{2}(n)&=4 \left(n - 2\right) \left(5 n - 4\right).
\end{align*}
The exact generating-function representation is $F=\sum_i R_i\sqrt{D_i}$, with
\begin{align*}
R_{0}&=0,\\
R_{1}&=\frac{x + 1}{4 x^{2} + 1},\\
D_{1}&=4 x^{2} + 1.
\end{align*}
For these data the complete rational certificate \eqref{round4:eq:r4gf-rational-cert} gives
\[ P(x)=0.\]
Every nonconstant radical component is identically zero. By \eqref{round4:eq:r4gf-extraction}, the recurrence therefore holds for every $n\ge2$. The source offset is 0.
\subsubsection{A100095: order 5, Nov 24 2012}
\label{round4:r4gf:A100095}\label{round4:r4entry:A100095}
Primary source: \url{https://oeis.org/A100095}. Kotesovec already recorded a shorter recurrence on 12 February 2014. The selected target is only Mathar's still-labeled order-five recurrence of 24 November 2012.
\begin{align*}
p_{0}(n)&=2 - n\\
p_{1}(n)&=3 - n\\
p_{2}(n)&=9 n - 22\\
p_{3}(n)&=9 n - 31\\
p_{4}(n)&=- 20 \left(n - 3\right)\\
p_{5}(n)&=- 20 \left(n - 4\right).
\end{align*}
The exact generating-function representation is $F=\sum_i R_i\sqrt{D_i}$, with
\begin{align*}
R_{0}&=- \frac{x}{5 x^{2} - 1},\\
R_{1}&=\frac{x^{2}}{\left(2 x - 1\right) \left(2 x + 1\right) \left(5 x^{2} - 1\right)},\\
D_{1}&=- \left(2 x - 1\right) \left(2 x + 1\right).
\end{align*}
For these data the complete rational certificate \eqref{round4:eq:r4gf-rational-cert} gives
\[ P(x)=x \left(x + 1\right).\]
Every nonconstant radical component is identically zero. By \eqref{round4:eq:r4gf-extraction}, the recurrence therefore holds for every $n\ge5$. The source offset is 0.
\subsubsection{A171851: order 9, Jul 22 2022}
\label{round4:r4gf:A171851}\label{round4:r4entry:A171851}
Primary source: \url{https://oeis.org/A171851}. The defining formula or generating function is established input from the cited primary entry. Only the selected, explicitly conjectural identity is counted.
\begin{align*}
p_{0}(n)&=- 2 \left(n + 2\right) \left(1088 n - 6417\right)\\
p_{1}(n)&=10968 n^{2} - 55789 n - 59252\\
p_{2}(n)&=- 2 \left(6704 n^{2} - 45808 n + 27647\right)\\
p_{3}(n)&=2 \left(2264 n^{2} - 18729 n + 42256\right)\\
p_{4}(n)&=- 2 \left(11968 n^{2} - 82555 n + 88362\right)\\
p_{5}(n)&=24904 n^{2} - 243375 n + 511724\\
p_{6}(n)&=- 4 \left(1088 n^{2} - 9137 n + 25680\right)\\
p_{7}(n)&=2 \left(6616 n^{2} - 65333 n + 164938\right)\\
p_{8}(n)&=- 2 \left(n - 7\right) \left(5616 n - 25895\right)\\
p_{9}(n)&=\left(n - 8\right) \left(2264 n - 10805\right).
\end{align*}
The exact generating-function representation is $F=\sum_i R_i\sqrt{D_i}$, with
\begin{align*}
R_{0}&=- \frac{x^{5} - 2 x^{4} + 3 x^{3} - 2 x^{2} + 5 x - 2}{2 x^{2} \left(x^{2} - 3 x + 1\right) \left(x^{2} + x + 1\right)},\\
R_{1}&=\frac{x^{3} - x^{2} + 3 x - 2}{2 x^{2} \left(x^{2} - 3 x + 1\right) \left(x^{2} + x + 1\right)},\\
D_{1}&=\left(x^{2} - 3 x + 1\right) \left(x^{2} + x + 1\right).
\end{align*}
For these data the complete rational certificate \eqref{round4:eq:r4gf-rational-cert} gives
\[ P(x)=30 x^{3} \left(103 x^{2} - 260 x + 1051\right).\]
Every nonconstant radical component is identically zero. By \eqref{round4:eq:r4gf-extraction}, the recurrence therefore holds for every $n\ge9$. The source offset is 0.
\subsubsection{A171847: order 10, Jul 22 2022}
\label{round4:r4gf:A171847}\label{round4:r4entry:A171847}
Primary source: \url{https://oeis.org/A171847}. The defining formula or generating function is established input from the cited primary entry. Only the selected, explicitly conjectural identity is counted.
\begin{align*}
p_{0}(n)&=n + 2\\
p_{1}(n)&=- 5 \left(n + 1\right)\\
p_{2}(n)&=7 n - 4\\
p_{3}(n)&=- 2 \left(3 n - 7\right)\\
p_{4}(n)&=13 n - 18\\
p_{5}(n)&=- 11 \left(n - 3\right)\\
p_{6}(n)&=13 n - 60\\
p_{7}(n)&=- 2 \left(3 n - 11\right)\\
p_{8}(n)&=7 n - 38\\
p_{9}(n)&=- 5 \left(n - 7\right)\\
p_{10}(n)&=n - 8.
\end{align*}
The exact generating-function representation is $F=\sum_i R_i\sqrt{D_i}$, with
\begin{align*}
R_{0}&=- \frac{x^{4} + 2 x^{3} - x^{2} + 2 x - 1}{2 x^{2} \left(x^{2} - 3 x + 1\right) \left(x^{2} + x + 1\right)},\\
R_{1}&=- \frac{x^{2} - x + 1}{2 x^{2} \left(x^{2} - 3 x + 1\right) \left(x^{2} + x + 1\right)},\\
D_{1}&=\left(x^{2} - 3 x + 1\right) \left(x^{2} + x + 1\right).
\end{align*}
For these data the complete rational certificate \eqref{round4:eq:r4gf-rational-cert} gives
\[ P(x)=2 x^{3} \left(x^{4} - 2 x^{3} + 3 x^{2} - 4 x + 5\right).\]
Every nonconstant radical component is identically zero. By \eqref{round4:eq:r4gf-extraction}, the recurrence therefore holds for every $n\ge10$. The source offset is 0.
\subsubsection{A171849: order 10, Jul 22 2022}
\label{round4:r4gf:A171849}\label{round4:r4entry:A171849}
Primary source: \url{https://oeis.org/A171849}. The defining formula or generating function is established input from the cited primary entry. Only the selected, explicitly conjectural identity is counted.
\begin{align*}
p_{0}(n)&=n + 1\\
p_{1}(n)&=- 5 n\\
p_{2}(n)&=7 n - 11\\
p_{3}(n)&=- 2 \left(3 n - 10\right)\\
p_{4}(n)&=13 n - 31\\
p_{5}(n)&=- 11 \left(n - 4\right)\\
p_{6}(n)&=13 n - 73\\
p_{7}(n)&=- 2 \left(3 n - 14\right)\\
p_{8}(n)&=7 n - 45\\
p_{9}(n)&=- 5 \left(n - 8\right)\\
p_{10}(n)&=n - 9.
\end{align*}
The exact generating-function representation is $F=\sum_i R_i\sqrt{D_i}$, with
\begin{align*}
R_{0}&=\frac{x^{4} - 2 x^{3} + x^{2} - 2 x + 1}{2 x \left(x^{2} - 3 x + 1\right) \left(x^{2} + x + 1\right)},\\
R_{1}&=- \frac{x^{2} - x + 1}{2 x \left(x^{2} - 3 x + 1\right) \left(x^{2} + x + 1\right)},\\
D_{1}&=\left(x^{2} - 3 x + 1\right) \left(x^{2} + x + 1\right).
\end{align*}
For these data the complete rational certificate \eqref{round4:eq:r4gf-rational-cert} gives
\[ P(x)=- 4 x^{3} \left(x - 1\right) \left(x + 1\right).\]
Every nonconstant radical component is identically zero. By \eqref{round4:eq:r4gf-extraction}, the recurrence therefore holds for every $n\ge10$. The source offset is 0.
\subsubsection{A191792: order 8, Jul 22 2022}
\label{round4:r4gf:A191792}\label{round4:r4entry:A191792}
Primary source: \url{https://oeis.org/A191792}. The defining formula or generating function is established input from the cited primary entry. Only the selected, explicitly conjectural identity is counted.
\begin{align*}
p_{0}(n)&=n + 1\\
p_{1}(n)&=-2\\
p_{2}(n)&=- 2 \left(n - 1\right)\\
p_{3}(n)&=-4\\
p_{4}(n)&=- 5 \left(n - 3\right)\\
p_{5}(n)&=4\\
p_{6}(n)&=- 2 \left(n - 5\right)\\
p_{7}(n)&=2\\
p_{8}(n)&=n - 7.
\end{align*}
The exact generating-function representation is $F=\sum_i R_i\sqrt{D_i}$, with
\begin{align*}
R_{0}&=- \frac{1}{2 x},\\
R_{1}&=\frac{1}{2 x \left(x^{4} - 2 x^{3} + x^{2} - 2 x + 1\right)},\\
D_{1}&=\left(x^{4} - 2 x^{3} + x^{2} - 2 x + 1\right) \left(x^{4} + 2 x^{3} + x^{2} + 2 x + 1\right).
\end{align*}
For these data the complete rational certificate \eqref{round4:eq:r4gf-rational-cert} gives
\[ P(x)=- \left(x - 1\right) \left(x + 1\right) \left(x^{4} + 3 x^{2} + 1\right).\]
Every nonconstant radical component is identically zero. By \eqref{round4:eq:r4gf-extraction}, the recurrence therefore holds for every $n\ge8$. The source offset is 0.
\subsubsection{A186338: order 4, Nov 17 2011}
\label{round4:r4gf:A186338}\label{round4:r4entry:A186338}
Primary source: \url{https://oeis.org/A186338}. The defining formula or generating function is established input from the cited primary entry. Only the selected, explicitly conjectural identity is counted.
\begin{align*}
p_{0}(n)&=n + 1\\
p_{1}(n)&=- 3 \left(4 n - 1\right)\\
p_{2}(n)&=15 \left(3 n - 4\right)\\
p_{3}(n)&=- 6 \left(11 n - 26\right)\\
p_{4}(n)&=16 \left(2 n - 7\right).
\end{align*}
The exact generating-function representation is $F=\sum_i R_i\sqrt{D_i}$, with
\begin{align*}
R_{0}&=\frac{3 x - 1}{2 x \left(2 x - 1\right)},\\
R_{1}&=\frac{1}{2 x \left(2 x - 1\right)},\\
D_{1}&=- \left(x - 1\right) \left(16 x^{2} - 9 x + 1\right).
\end{align*}
For these data the complete rational certificate \eqref{round4:eq:r4gf-rational-cert} gives
\[ P(x)=- \left(2 x - 1\right) \left(6 x^{2} - 3 x + 1\right).\]
Every nonconstant radical component is identically zero. By \eqref{round4:eq:r4gf-extraction}, the recurrence therefore holds for every $n\ge4$. The source offset is 0.
\subsubsection{A182899: order 8, Jul 22 2022}
\label{round4:r4gf:A182899}\label{round4:r4entry:A182899}
Primary source: \url{https://oeis.org/A182899}. The defining formula or generating function is established input from the cited primary entry. Only the selected, explicitly conjectural identity is counted.
\begin{align*}
p_{0}(n)&=n\\
p_{1}(n)&=3 - 4 n\\
p_{2}(n)&=2 n - 3\\
p_{3}(n)&=0\\
p_{4}(n)&=11 \left(n - 3\right)\\
p_{5}(n)&=0\\
p_{6}(n)&=2 n - 9\\
p_{7}(n)&=21 - 4 n\\
p_{8}(n)&=n - 6.
\end{align*}
The exact generating-function representation is $F=\sum_i R_i\sqrt{D_i}$, with
\begin{align*}
R_{0}&=- \frac{x^{2} + x - 1}{\left(x^{2} - 3 x + 1\right) \left(x^{2} + x + 1\right)},\\
R_{1}&=- \frac{1}{\left(x^{2} - 3 x + 1\right) \left(x^{2} + x + 1\right)},\\
D_{1}&=\left(x^{2} - 3 x + 1\right) \left(x^{2} + x + 1\right).
\end{align*}
For these data the complete rational certificate \eqref{round4:eq:r4gf-rational-cert} gives
\[ P(x)=2 x^{3} \left(x^{2} - x + 3\right).\]
Every nonconstant radical component is identically zero. By \eqref{round4:eq:r4gf-extraction}, the recurrence therefore holds for every $n\ge8$. The source offset is 0.
\subsubsection{A187256: order 4, Nov 16 2011}
\label{round4:r4gf:A187256}\label{round4:r4entry:A187256}
Primary source: \url{https://oeis.org/A187256}. The entry links Barry's 2019 and 2021 papers. Their relevant discussions give existing generating-function, convolution and Hankel-transform results. These are credited as prior work; the selected linear recurrence is the separately labeled conjecture.
\begin{align*}
p_{0}(n)&=n + 2\\
p_{1}(n)&=- 2 \left(2 n + 1\right)\\
p_{2}(n)&=2 \left(n - 1\right)\\
p_{3}(n)&=- 2 \left(2 n - 5\right)\\
p_{4}(n)&=n - 4.
\end{align*}
The exact generating-function representation is $F=\sum_i R_i\sqrt{D_i}$, with
\begin{align*}
R_{0}&=\frac{\left(x - 1\right)^{2}}{2 x^{2}},\\
R_{1}&=- \frac{1}{2 x^{2}},\\
D_{1}&=\left(x^{2} + 1\right) \left(x^{2} - 4 x + 1\right).
\end{align*}
For these data the complete rational certificate \eqref{round4:eq:r4gf-rational-cert} gives
\[ P(x)=- 2 \left(x - 1\right) \left(x + 1\right).\]
Every nonconstant radical component is identically zero. By \eqref{round4:eq:r4gf-extraction}, the recurrence therefore holds for every $n\ge4$. The source offset is 0.
\subsubsection{A184120: order 4, Nov 17 2011}
\label{round4:r4gf:A184120}\label{round4:r4entry:A184120}
Primary source: \url{https://oeis.org/A184120}. The defining formula or generating function is established input from the cited primary entry. Only the selected, explicitly conjectural identity is counted.
\begin{align*}
p_{0}(n)&=n + 1\\
p_{1}(n)&=2 \left(2 n + 1\right)\\
p_{2}(n)&=4 \left(n - 1\right)\\
p_{3}(n)&=4 \left(2 n - 5\right)\\
p_{4}(n)&=4 \left(n - 3\right).
\end{align*}
The exact generating-function representation is $F=\sum_i R_i\sqrt{D_i}$, with
\begin{align*}
R_{0}&=\frac{1}{2 x},\\
R_{1}&=- \frac{1}{2 x},\\
D_{1}&=\frac{2 x^{2} + 1}{2 x^{2} + 4 x + 1}.
\end{align*}
For these data the complete rational certificate \eqref{round4:eq:r4gf-rational-cert} gives
\[ P(x)=1 - 2 x^{2}.\]
Every nonconstant radical component is identically zero. By \eqref{round4:eq:r4gf-extraction}, the recurrence therefore holds for every $n\ge4$. The source offset is 0.
\subsubsection{A186940: order 5, Nov 17 2011}
\label{round4:r4gf:A186940}\label{round4:r4entry:A186940}
Primary source: \url{https://oeis.org/A186940}. The defining formula or generating function is established input from the cited primary entry. Only the selected, explicitly conjectural identity is counted.
\begin{align*}
p_{0}(n)&=3 \left(n + 1\right)\\
p_{1}(n)&=10 - 17 n\\
p_{2}(n)&=- 6 \left(n + 1\right)\\
p_{3}(n)&=10 \left(2 n - 1\right)\\
p_{4}(n)&=- 9 \left(n - 1\right)\\
p_{5}(n)&=n - 2.
\end{align*}
The exact generating-function representation is $F=\sum_i R_i\sqrt{D_i}$, with
\begin{align*}
R_{0}&=- \frac{2 x^{2} - 3 x - 1}{2 x \left(x^{3} - 3 x^{2} + x + 3\right)},\\
R_{1}&=- \frac{1}{2 x \left(x^{3} - 3 x^{2} + x + 3\right)},\\
D_{1}&=x^{2} - 6 x + 1.
\end{align*}
For these data the complete rational certificate \eqref{round4:eq:r4gf-rational-cert} gives
\[ P(x)=- x^{3} + 9 x^{2} - 7 x + 3.\]
Every nonconstant radical component is identically zero. By \eqref{round4:eq:r4gf-extraction}, the recurrence therefore holds for every $n\ge5$. The source offset is 0.
\subsubsection{A176610: order 5, Nov 17 2011}
\label{round4:r4gf:A176610}\label{round4:r4entry:A176610}
Primary source: \url{https://oeis.org/A176610}. The defining formula or generating function is established input from the cited primary entry. Only the selected, explicitly conjectural identity is counted.
\begin{align*}
p_{0}(n)&=n + 1\\
p_{1}(n)&=2 - 7 n\\
p_{2}(n)&=19 n - 29\\
p_{3}(n)&=110 - 37 n\\
p_{4}(n)&=36 \left(n - 4\right)\\
p_{5}(n)&=- 12 \left(n - 5\right).
\end{align*}
The exact generating-function representation is $F=\sum_i R_i\sqrt{D_i}$, with
\begin{align*}
R_{0}&=\frac{1}{2 x},\\
R_{1}&=- \frac{1}{2 x},\\
D_{1}&=\frac{12 x^{4} - 24 x^{3} + 13 x^{2} - 6 x + 1}{\left(x - 1\right)^{2}}.
\end{align*}
For these data the complete rational certificate \eqref{round4:eq:r4gf-rational-cert} gives
\[ P(x)=6 x^{4} - 18 x^{3} + 18 x^{2} - 5 x + 1.\]
Every nonconstant radical component is identically zero. By \eqref{round4:eq:r4gf-extraction}, the recurrence therefore holds for every $n\ge5$. The source offset is 0.
\subsubsection{A179648: order 4, Nov 17 2011}
\label{round4:r4gf:A179648}\label{round4:r4entry:A179648}
Primary source: \url{https://oeis.org/A179648}. The defining formula or generating function is established input from the cited primary entry. Only the selected, explicitly conjectural identity is counted.
\begin{align*}
p_{0}(n)&=n + 1\\
p_{1}(n)&=2 \left(2 n + 1\right)\\
p_{2}(n)&=- 4 \left(n - 1\right)\\
p_{3}(n)&=- 4 \left(2 n - 5\right)\\
p_{4}(n)&=4 \left(n - 3\right).
\end{align*}
The exact generating-function representation is $F=\sum_i R_i\sqrt{D_i}$, with
\begin{align*}
R_{0}&=\frac{1}{2 x},\\
R_{1}&=- \frac{1}{2 x},\\
D_{1}&=\frac{2 x^{2} - 1}{2 x^{2} - 4 x - 1}.
\end{align*}
For these data the complete rational certificate \eqref{round4:eq:r4gf-rational-cert} gives
\[ P(x)=2 x^{2} + 1.\]
Every nonconstant radical component is identically zero. By \eqref{round4:eq:r4gf-extraction}, the recurrence therefore holds for every $n\ge4$. The source offset is 0.
\subsubsection{A166228: order 3, Nov 17 2011}
\label{round4:r4gf:A166228}\label{round4:r4entry:A166228}
Primary source: \url{https://oeis.org/A166228}. The defining formula or generating function is established input from the cited primary entry. Only the selected, explicitly conjectural identity is counted.
\begin{align*}
p_{0}(n)&=n + 1\\
p_{1}(n)&=4 - 5 n\\
p_{2}(n)&=1 - 5 n\\
p_{3}(n)&=n - 2.
\end{align*}
The exact generating-function representation is $F=\sum_i R_i\sqrt{D_i}$, with
\begin{align*}
R_{0}&=- \frac{x - 1}{2 x \left(x + 1\right)},\\
R_{1}&=- \frac{1}{2 x \left(x + 1\right)},\\
D_{1}&=x^{2} - 6 x + 1.
\end{align*}
For these data the complete rational certificate \eqref{round4:eq:r4gf-rational-cert} gives
\[ P(x)=x + 1.\]
Every nonconstant radical component is identically zero. By \eqref{round4:eq:r4gf-extraction}, the recurrence therefore holds for every $n\ge3$. The source offset is 0.
\subsubsection{A050168: order 3, Nov 26 2012}
\label{round4:r4gf:A050168}\label{round4:r4entry:A050168}
Primary source: \url{https://oeis.org/A050168}. The defining formula or generating function is established input from the cited primary entry. Only the selected, explicitly conjectural identity is counted.
\begin{align*}
p_{0}(n)&=n + 1\\
p_{1}(n)&=n - 3\\
p_{2}(n)&=- 2 \left(2 n - 1\right)\\
p_{3}(n)&=- 4 \left(n - 3\right).
\end{align*}
The exact generating-function representation is $F=\sum_i R_i\sqrt{D_i}$, with
\begin{align*}
R_{0}&=- \frac{x + 1}{2 x},\\
R_{1}&=\frac{x + 1}{2 x},\\
D_{1}&=- \frac{2 x + 1}{2 x - 1}.
\end{align*}
For these data the complete rational certificate \eqref{round4:eq:r4gf-rational-cert} gives
\[ P(x)=\left(x + 1\right)^{2}.\]
Every nonconstant radical component is identically zero. By \eqref{round4:eq:r4gf-extraction}, the recurrence therefore holds for every $n\ge3$. The source offset is 0.
\subsubsection{A072100: order 3, Nov 26 2012}
\label{round4:r4gf:A072100}\label{round4:r4entry:A072100}
Primary source: \url{https://oeis.org/A072100}. The defining formula or generating function is established input from the cited primary entry. Only the selected, explicitly conjectural identity is counted.
\begin{align*}
p_{0}(n)&=1 - n\\
p_{1}(n)&=n + 1\\
p_{2}(n)&=2 \left(2 n - 7\right)\\
p_{3}(n)&=- 4 \left(n - 3\right).
\end{align*}
The exact generating-function representation is $F=\sum_i R_i\sqrt{D_i}$, with
\begin{align*}
R_{0}&=- \frac{x}{2 \left(x - 1\right)},\\
R_{1}&=\frac{x}{2 \left(x - 1\right) \left(2 x - 1\right)},\\
D_{1}&=- \left(2 x - 1\right) \left(2 x + 1\right).
\end{align*}
For these data the complete rational certificate \eqref{round4:eq:r4gf-rational-cert} gives
\[ P(x)=x^{2}.\]
Every nonconstant radical component is identically zero. By \eqref{round4:eq:r4gf-extraction}, the recurrence therefore holds for every $n\ge4$. The source offset is 1.
\subsubsection{A176611: order 5, Mar 01 2016}
\label{round4:r4gf:A176611}\label{round4:r4entry:A176611}
Primary source: \url{https://oeis.org/A176611}. The defining formula or generating function is established input from the cited primary entry. Only the selected, explicitly conjectural identity is counted.
\begin{align*}
p_{0}(n)&=n + 1\\
p_{1}(n)&=2 - 7 n\\
p_{2}(n)&=3 \left(5 n - 7\right)\\
p_{3}(n)&=74 - 25 n\\
p_{4}(n)&=24 \left(n - 4\right)\\
p_{5}(n)&=- 8 \left(n - 5\right).
\end{align*}
The exact generating-function representation is $F=\sum_i R_i\sqrt{D_i}$, with
\begin{align*}
R_{0}&=\frac{1}{2 x},\\
R_{1}&=- \frac{1}{2 x},\\
D_{1}&=\frac{8 x^{4} - 16 x^{3} + 9 x^{2} - 6 x + 1}{\left(x - 1\right)^{2}}.
\end{align*}
For these data the complete rational certificate \eqref{round4:eq:r4gf-rational-cert} gives
\[ P(x)=4 x^{4} - 12 x^{3} + 12 x^{2} - 3 x + 1.\]
Every nonconstant radical component is identically zero. By \eqref{round4:eq:r4gf-extraction}, the recurrence therefore holds for every $n\ge5$. The source offset is 0.
\subsubsection{A176612: order 5, Feb 29 2016}
\label{round4:r4gf:A176612}\label{round4:r4entry:A176612}
Primary source: \url{https://oeis.org/A176612}. The defining formula or generating function is established input from the cited primary entry. Only the selected, explicitly conjectural identity is counted.
\begin{align*}
p_{0}(n)&=n + 1\\
p_{1}(n)&=2 - 7 n\\
p_{2}(n)&=11 n - 13\\
p_{3}(n)&=38 - 13 n\\
p_{4}(n)&=12 \left(n - 4\right)\\
p_{5}(n)&=- 4 \left(n - 5\right).
\end{align*}
The exact generating-function representation is $F=\sum_i R_i\sqrt{D_i}$, with
\begin{align*}
R_{0}&=\frac{1}{2 x},\\
R_{1}&=- \frac{1}{2 x},\\
D_{1}&=\frac{4 x^{4} - 8 x^{3} + 5 x^{2} - 6 x + 1}{\left(x - 1\right)^{2}}.
\end{align*}
For these data the complete rational certificate \eqref{round4:eq:r4gf-rational-cert} gives
\[ P(x)=2 x^{4} - 6 x^{3} + 6 x^{2} - x + 1.\]
Every nonconstant radical component is identically zero. By \eqref{round4:eq:r4gf-extraction}, the recurrence therefore holds for every $n\ge5$. The source offset is 0.
\subsubsection{A176605: order 4, Feb 29 2016}
\label{round4:r4gf:A176605}\label{round4:r4entry:A176605}
Primary source: \url{https://oeis.org/A176605}. The defining formula or generating function is established input from the cited primary entry. Only the selected, explicitly conjectural identity is counted.
\begin{align*}
p_{0}(n)&=n + 1\\
p_{1}(n)&=- 2 \left(3 n - 1\right)\\
p_{2}(n)&=9 n - 13\\
p_{3}(n)&=- 8 \left(n - 3\right)\\
p_{4}(n)&=4 \left(n - 4\right).
\end{align*}
The exact generating-function representation is $F=\sum_i R_i\sqrt{D_i}$, with
\begin{align*}
R_{0}&=\frac{1}{2 x},\\
R_{1}&=- \frac{1}{2 x},\\
D_{1}&=\frac{4 x^{3} - 4 x^{2} + 5 x - 1}{x - 1}.
\end{align*}
For these data the complete rational certificate \eqref{round4:eq:r4gf-rational-cert} gives
\[ P(x)=- 2 x^{3} + 4 x^{2} - 2 x + 1.\]
Every nonconstant radical component is identically zero. By \eqref{round4:eq:r4gf-extraction}, the recurrence therefore holds for every $n\ge4$. The source offset is 0.
\subsubsection{A176607: order 4, Feb 29 2016}
\label{round4:r4gf:A176607}\label{round4:r4entry:A176607}
Primary source: \url{https://oeis.org/A176607}. The defining formula or generating function is established input from the cited primary entry. Only the selected, explicitly conjectural identity is counted.
\begin{align*}
p_{0}(n)&=n + 1\\
p_{1}(n)&=- 2 \left(3 n - 1\right)\\
p_{2}(n)&=11 - 3 n\\
p_{3}(n)&=16 \left(n - 3\right)\\
p_{4}(n)&=- 8 \left(n - 4\right).
\end{align*}
The exact generating-function representation is $F=\sum_i R_i\sqrt{D_i}$, with
\begin{align*}
R_{0}&=\frac{1}{2 x},\\
R_{1}&=- \frac{1}{2 x},\\
D_{1}&=- \frac{8 x^{3} - 8 x^{2} - 5 x + 1}{x - 1}.
\end{align*}
For these data the complete rational certificate \eqref{round4:eq:r4gf-rational-cert} gives
\[ P(x)=4 x^{3} - 8 x^{2} + 4 x + 1.\]
Every nonconstant radical component is identically zero. By \eqref{round4:eq:r4gf-extraction}, the recurrence therefore holds for every $n\ge4$. The source offset is 0.
\subsubsection{A045743: order 4, Jul 26 2022}
\label{round4:r4gf:A045743}\label{round4:r4entry:A045743}
Primary source: \url{https://oeis.org/A045743}. The defining formula or generating function is established input from the cited primary entry. Only the selected, explicitly conjectural identity is counted.
\begin{align*}
p_{0}(n)&=5 n \left(n - 2\right) \left(n - 1\right) \left(1353818 n - 3651663\right)\\
p_{1}(n)&=16 \left(n - 2\right) \left(n - 1\right) \left(2236406 n^{2} - 6864025 n + 2014500\right)\\
p_{2}(n)&=- 16 \left(n - 2\right) \left(46338048 n^{3} - 326486432 n^{2} + 734160494 n - 506411475\right)\\
p_{3}(n)&=8 \left(2 n - 7\right) \left(26592864 n^{3} - 436441232 n^{2} + 2071128458 n - 3030735075\right)\\
p_{4}(n)&=128 \left(2 n - 9\right) \left(4 n - 19\right) \left(4 n - 17\right) \left(94246 n - 334675\right).
\end{align*}
Use $x=\xi(t)=- \frac{\left(t - 1\right) \left(t^{2} - 3 t + 1\right)}{t^{4}}$, $t(0)=1$.
The exact generating-function representation is $F=\sum_i R_i\sqrt{D_i}$, with
\begin{align*}
R_{0}&=- \frac{\left(t - 1\right)^{2} \left(t^{2} - 3 t + 1\right)}{t^{4}},\\
R_{1}&=0,\\
D_{1}&=1.
\end{align*}
For these data the complete rational certificate \eqref{round4:eq:r4gf-rational-cert} gives
\[ P(x)=- 6 x^{3} \left(8739520 x^{2} - 37127156 x - 14413953\right).\]
Every nonconstant radical component is identically zero. By \eqref{round4:eq:r4gf-extraction}, the recurrence therefore holds for every $n\ge6$. The source offset is 2.
\subsubsection{A308726: order 5, Jan 27 2020}
\label{round4:r4gf:A308726}\label{round4:r4entry:A308726}
Primary source: \url{https://oeis.org/A308726}. The stack-sorting interpretation and established generating function are prior input, with the entry citing Mansour, Skogman and Smith (2017--2019). The selected target is the 2020 polynomial recurrence.
\begin{align*}
p_{0}(n)&=3 n \left(n - 1\right) \left(n + 1\right)\\
p_{1}(n)&=- n \left(n - 1\right) \left(67 n - 101\right)\\
p_{2}(n)&=2 \left(n - 1\right) \left(286 n^{2} - 1112 n + 1089\right)\\
p_{3}(n)&=- 4 \left(580 n^{3} - 4200 n^{2} + 10106 n - 8049\right)\\
p_{4}(n)&=24 \left(184 n^{3} - 1784 n^{2} + 5770 n - 6221\right)\\
p_{5}(n)&=- 96 \left(2 n - 9\right) \left(4 n - 17\right) \left(4 n - 15\right).
\end{align*}
Use $x=\xi(t)=- \frac{\left(t - 1\right) \left(t + 1\right)}{4}$, $t(0)=1$.
The exact generating-function representation is $F=\sum_i R_i\sqrt{D_i}$, with
\begin{align*}
R_{0}&=\frac{t^{2} - 4 t + 1}{t \left(t - 1\right) \left(t + 1\right)},\\
R_{1}&=\frac{2}{\left(t - 1\right) \left(t + 1\right)},\\
D_{1}&=2 t - 1.
\end{align*}
For these data the complete rational certificate \eqref{round4:eq:r4gf-rational-cert} gives
\[ P(x)=- 12 x^{2} \left(8 x^{2} - 2 x + 1\right).\]
Every nonconstant radical component is identically zero. By \eqref{round4:eq:r4gf-extraction}, the recurrence therefore holds for every $n\ge5$. The source offset is 0.
\subsubsection{A285195: order 6, Jan 25 2020}
\label{round4:r4gf:A285195}\label{round4:r4entry:A285195}
Primary source: \url{https://oeis.org/A285195}. The established generating function and the He--Shapiro Riordan-array reference are prior input. Only the displayed conjectural recurrence is selected.
\begin{align*}
p_{0}(n)&=126 n \left(n - 1\right) \left(n + 1\right)\\
p_{1}(n)&=- 3 n \left(n - 1\right) \left(691 n - 1037\right)\\
p_{2}(n)&=\left(n - 1\right) \left(12023 n^{2} - 48049 n + 47736\right)\\
p_{3}(n)&=- 2 \left(12889 n^{3} - 93970 n^{2} + 221427 n - 169506\right)\\
p_{4}(n)&=- 4 \left(248 n^{3} - 12718 n^{2} + 88433 n - 160113\right)\\
p_{5}(n)&=16 \left(2918 n^{3} - 41574 n^{2} + 198532 n - 317121\right)\\
p_{6}(n)&=240 \left(2 n - 9\right) \left(4 n - 21\right) \left(4 n - 19\right).
\end{align*}
Use $x=\xi(t)=- \frac{\left(t - 1\right) \left(t + 1\right)}{4}$, $t(0)=1$.
The exact generating-function representation is $F=\sum_i R_i\sqrt{D_i}$, with
\begin{align*}
R_{0}&=- \frac{6}{\left(t - 5\right) \left(t + 1\right)},\\
R_{1}&=- \frac{2}{\left(t - 5\right) \left(t + 1\right)},\\
D_{1}&=2 t - 1.
\end{align*}
For these data the complete rational certificate \eqref{round4:eq:r4gf-rational-cert} gives
\[ P(x)=- 90 x^{2} \left(4 x^{3} + 80 x^{2} - 32 x + 3\right).\]
Every nonconstant radical component is identically zero. By \eqref{round4:eq:r4gf-extraction}, the recurrence therefore holds for every $n\ge6$. The source offset is 0.
\subsubsection{A243642: order 6, Jan 25 2020}
\label{round4:r4gf:A243642}\label{round4:r4entry:A243642}
Primary source: \url{https://oeis.org/A243642}. The defining formula or generating function is established input from the cited primary entry. Only the selected, explicitly conjectural identity is counted.
\textbf{Boundary correction.} The source gives no starting index.
At its first natural admissible index $n=6$, the displayed recurrence
has residual $24$, not zero. Thus its interpretation as holding at all
admissible indices is refuted. The corrected tail assertion is proved
for every $n\ge7$. We do not attribute an explicit $n=6$ claim to the
source, and the counterexample and corrected theorem count together once.
\begin{align*}
p_{0}(n)&=16 n \left(n - 2\right) \left(n - 1\right)\\
p_{1}(n)&=- 4 \left(n - 2\right) \left(n - 1\right) \left(43 n - 54\right)\\
p_{2}(n)&=2 \left(n - 2\right) \left(52 n^{2} + 695 n - 2439\right)\\
p_{3}(n)&=3479 n^{3} - 42714 n^{2} + 169789 n - 219234\\
p_{4}(n)&=- 2 \left(4633 n^{3} - 64020 n^{2} + 293825 n - 447594\right)\\
p_{5}(n)&=4 \left(2 n - 11\right) \left(296 n^{2} - 3643 n + 11187\right)\\
p_{6}(n)&=24 \left(n - 7\right) \left(2 n - 13\right) \left(2 n - 11\right).
\end{align*}
Use $x=\xi(t)=- \frac{\left(t - 1\right) \left(t + 1\right)}{4}$, $t(0)=1$.
The exact generating-function representation is $F=\sum_i R_i\sqrt{D_i}$, with
\begin{align*}
R_{0}&=\frac{\left(t - 3\right) \left(t + 1\right)}{8},\\
R_{1}&=- \frac{t + 1}{8},\\
D_{1}&=t^{2} + 10 t - 7.
\end{align*}
The chosen square root of $D_1$ has value $2$ at $t=1$.
For these data the complete rational certificate \eqref{round4:eq:r4gf-rational-cert} gives
\[ P(x)=12 x^{3} \left(2 x^{3} + 44 x^{2} + 31 x - 2\right).\]
Every nonconstant radical component is identically zero. By \eqref{round4:eq:r4gf-extraction}, the recurrence therefore holds for every $n\ge7$. The source offset is 0.
\subsubsection{A243816: order 6, Jan 25 2020}
\label{round4:r4gf:A243816}\label{round4:r4entry:A243816}
Primary source: \url{https://oeis.org/A243816}. The defining formula or generating function is established input from the cited primary entry. Only the selected, explicitly conjectural identity is counted.
\begin{align*}
p_{0}(n)&=2 n \left(n - 1\right) \left(2 n - 3\right)\\
p_{1}(n)&=- \left(n - 1\right) \left(47 n^{2} - 221 n + 260\right)\\
p_{2}(n)&=4 \left(38 n^{3} - 300 n^{2} + 777 n - 660\right)\\
p_{3}(n)&=- 4 \left(94 n^{3} - 1006 n^{2} + 3617 n - 4355\right)\\
p_{4}(n)&=16 \left(n - 5\right) \left(64 n^{2} - 510 n + 1041\right)\\
p_{5}(n)&=- 16 \left(n - 6\right) \left(n - 5\right) \left(47 n - 221\right)\\
p_{6}(n)&=1920 \left(n - 7\right) \left(n - 6\right) \left(n - 5\right).
\end{align*}
Use $x=\xi(t)=- \frac{t}{\left(t - 1\right) \left(t + 1\right)}$, $t(0)=0$.
The exact generating-function representation is $F=\sum_i R_i\sqrt{D_i}$, with
\begin{align*}
R_{0}&=- \frac{2 t - 1}{2 \left(t - 1\right) \left(t + 1\right)},\\
R_{1}&=\frac{1}{2 \left(t - 1\right) \left(t + 1\right)},\\
D_{1}&=1 - 4 t.
\end{align*}
For these data the complete rational certificate \eqref{round4:eq:r4gf-rational-cert} gives
\[ P(x)=- 4 x^{2} \left(8 x^{2} + 8 x + 1\right).\]
Every nonconstant radical component is identically zero. By \eqref{round4:eq:r4gf-extraction}, the recurrence therefore holds for every $n\ge6$. The source offset is 0.
\subsection{A000986: the exponential generating function}
\label{round4:r4gf:A000986}\label{round4:r4entry:A000986}
Primary source: \url{https://oeis.org/A000986}. The selected conjecture
is Mathar's order-five relation dated 4 August 2013. Kotesovec's
different order-four recurrence of 13 February 2014 is already known
and is explicitly credited as such. The established EGF is
\[
 E(x)=(1-x)^{-1/2}\exp\left(-x-\frac{x^2}{4}+
                             \frac{x}{2(1-x)}\right).
\]
Its logarithmic derivative is
\[
 L(x)=\frac{E'}E=\frac1{2(1-x)}-1-\frac x2+
                                  \frac1{2(1-x)^2}.
\]
The conjectured recurrence uses
\begin{align*}
p_{0}(n)&=2\\
p_{1}(n)&=- 2 \left(2 n - 1\right)\\
p_{2}(n)&=2 \left(n^{2} - 2 n - 1\right)\\
p_{3}(n)&=- 2 \left(n - 4\right) \left(n - 2\right)\\
p_{4}(n)&=\left(n - 3\right) \left(n - 2\right) \left(n - 1\right)\\
p_{5}(n)&=- \left(n - 4\right) \left(n - 3\right) \left(n - 2\right).
\end{align*}
Put $b_n=a_n/n!$ and $r=5$. For $n\ge r$, multiplication of
the recurrence by $1/(n-r)!$ yields
\[
 \sum_{j=0}^{r}q_j(n)b_{n-j}=0,\qquad
 q_j(n)=p_j(n)(n-j)_{r-j},
\]
where $(z)_h=z(z-1)\cdots(z-h+1)$.
Let $U_0=1$, $U_{m+1}=x(U_m'+LU_m)$, and
$c_m(x)=\sum_jx^j[n^m]q_j(n+j)$. Then
$\theta^mE=U_mE$ by induction. Exact rational simplification gives
\[
 \sum_m c_m(x)U_m(x)=0.
\]
All its inputs have been specified explicitly above, and the verifier
reconstructs every $q_j,c_m,U_m$ and clears denominators to obtain the
zero polynomial. Thus $\mathcal L_qE=0$, proving the selected recurrence
for every $n\ge5$. The factorial conversion is essential: no OGF
operator is incorrectly applied directly to the integer sequence.

\subsection{A121754: a harmonic-number certificate}
\label{round4:r4gf:A121754}\label{round4:r4entry:A121754}
Primary source: \url{https://oeis.org/A121754}; selected conjecture
dated 26 July 2022. The established recurrence immediately preceding
it in the entry is
\begin{align*}
 a_n-(2n-3)a_{n-1}+(n-1)(n-3)a_{n-2}
 &=(n-2)!\left(n-2+\frac{n-1}{2}(1+(-1)^{n-1})\right),\\
 a_1&=0,\qquad a_2=1.
\end{align*}
It gives $a_3=6$. Define
\[
 R(n)=\frac{8n^3-2n^2-14n+7}{16},\quad
 V(n)=\frac{2n-1}{16},\quad W(n)=-\frac{n(n-1)}2.
\]
Then the exact closed form is
\begin{equation}\label{round4:eq:r4gf-harmonic}
 a_n=(n-2)!\{R(n)+V(n)(-1)^{n-1}+W(n)H_{n-1}\}
 \qquad(n\ge2).
\end{equation}
One derivation sets $b_n=a_n/(n-2)!$ for $n\ge2$, defines the auxiliary
value $b_1=0$, and puts $d_n=b_n-b_{n-1}$ for $n\ge2$.
The given recurrence becomes, for $n\ge4$,
\[
 (n-2)d_n-(n-1)d_{n-1}
 =(n-2)\left\{n-2+\frac{n-1}{2}(1+(-1)^{n-1})\right\}.
\]
Consequently $e_n=d_n/(n-1)$ satisfies
\[
 e_n-e_{n-1}=\frac32-\frac1{n-1}+\frac{(-1)^{n-1}}2,
\qquad e_2=1.
\]
The same difference equation at $n=3$ follows directly from $b_2=1$
and $b_3=6$, so summation may start with $e_2$.
Summation gives
$e_n=(6n-3+(-1)^{n-1})/4-H_{n-1}$. A second summation uses
\[
 \sum_{j=1}^NjH_j=\frac{N(N+1)}2H_N-\frac{N(N-1)}4,\quad
 \sum_{j=1}^Nj(-1)^j=\frac{(2N+1)(-1)^N-1}{4},
\]
and yields \eqref{round4:eq:r4gf-harmonic}. Equivalently, substitute that
formula into the known recurrence, divide by $(n-4)!$ for $n\ge4$,
and use $H_{n-2}=H_{n-1}-1/(n-1)$: the rational residual is identically
zero. The explicit values $a_2=1,a_3=6$ cover the indices omitted by
that division, so induction proves the formula for every $n\ge2$.

The conjectured homogeneous recurrence has the following coefficients.
\begin{align*}
p_{0}(n)&=16 \left(n + 1\right)\\
p_{1}(n)&=- 16 n^{2} - 178 n + 531\\
p_{2}(n)&=- 16 n^{3} + 178 n^{2} - 393 n - 510\\
p_{3}(n)&=16 n^{4} + 98 n^{3} - 1439 n^{2} + 4222 n - 3623\\
p_{4}(n)&=- 146 n^{4} + 1479 n^{3} - 4483 n^{2} + 3054 n + 2841\\
p_{5}(n)&=\left(n - 6\right) \left(n - 4\right)^{2} \left(130 n - 311\right).
\end{align*}
For independent symbols $H,\varepsilon$, put
\[
 B_j=R(n-j)+(-1)^jV(n-j)\varepsilon+
 W(n-j)\left(H-\sum_{h=1}^{j}\frac1{n-h}\right).
\]
The complete rational identity certificate is
\[
 \sum_{j=0}^{5}p_j(n)
       \left(\prod_{h=1}^{5-j}(n-7+h)\right)B_j=0.
\]
Its coefficients of $1,H,\varepsilon$ all vanish identically in
$\mathbb Q(n)$; this is directly checked by clearing denominators.
Set $H=H_{n-1}$, $\varepsilon=(-1)^{n-1}$ and multiply by $(n-7)!$.
The result is the desired recurrence for every $n\ge7$.
At $n=6$, the required values
$(a_1,\ldots,a_6)=(0,1,6,31,211,1530)$ give residual zero by direct
substitution. Hence the selected recurrence holds for every $n\ge6$.

\subsection{A004395: three polynomial certificates}
\label{round4:r4gf:A004395}\label{round4:r4entry:A004395}
Primary source: \url{https://oeis.org/A004395}; selected conjecture
dated 14 September 2025. The definition by successive ratios is
\[
 a_0=1,\qquad a_n=\left\lfloor\frac{2n+1}{3}\right\rfloor a_{n-1}.
\]
The conjectured recurrence has coefficients
\begin{align*}
p_{0}(n)&=243\\
p_{1}(n)&=-216\\
p_{2}(n)&=- 54 \left(2 n - 3\right)\\
p_{3}(n)&=- 6 \left(12 n^{3} - 60 n^{2} + 87 n - 28\right)\\
p_{4}(n)&=4 \left(n - 3\right) \left(2 n - 7\right) \left(2 n - 5\right).
\end{align*}
For $n=3k+s$, $s\in\{0,1,2\}$, divide the recurrence residual
by the positive integer $a_{n-4}$. The result is
\[
 \sum_{j=0}^{4}p_j(3k+s)
       \prod_{h=j}^{3}\left(2k+
                   \left\lfloor\frac{2(s-h)+1}{3}\right\rfloor\right).
\]
For each of the three stated values of $s$, this is a polynomial in
$k$ and its complete expansion is zero. These three finite identities
are exactly the retained certificates; they prove the recurrence for
all $n\ge4$. The proof does not assume convergence of the source's
formal hypergeometric generating function.

\subsection{A052731: the factorial-weighted A002026 identity}
\label{round4:r4gf:A052731}\label{round4:r4entry:A052731}
Primary sources: \url{https://oeis.org/A052731} and
\url{https://oeis.org/A002026}. The selected conjecture, dated
16 October 2013, is $a_n=n!\,\mathrm{A002026}(n-2)$.
Let $D=1-2x-3x^2$ and let
\[
 M=\frac{1-x-\sqrt D}{2x^2}.
\]
The quadratic identity $M=1+xM+x^2M^2$ follows either by squaring
or by substitution. The established OGF of A002026 is $B=xM^2$;
the established EGF of A052731 is
\[
 E=\frac{1-x-\sqrt D}{2x}-\frac{1+x-\sqrt D}{2}.
\]
Since $x^2M=(1-x-\sqrt D)/2$, we have
\[
 E=xM-x-x^2M=x(M-1-xM)=x^3M^2=x^2B.
\]
Equating coefficients proves
$a_n=n!\,\mathrm{A002026}(n-2)$ for every $n\ge2$.
The source has offset zero, with $a_0=a_1=a_2=0$. The supporting A002026 entry is
not counted separately. This is an exact formal-series identity.

\subsection{Reproducibility and prior-work scope}
The files \texttt{r4\_gf\_certificates.json},
\texttt{r4\_gf\_engine.py} and \texttt{r4\_gf\_verify.py} retain the
complete exact inputs and the finite rational checks. The verifier
uses only rational arithmetic, polynomial identities and elementary
formal-series coefficient arithmetic. The compact primary-source records
retain the selected mathematical statements, attributions, offsets,
numeric prefixes and source URLs; raw webpage text is not required.

The Ekhad--Yang 2017 output proving 33 Mathar recurrence conjectures
was checked by A-number; none of these 35 retained targets occurs in
that output. A185010, which does occur there, was excluded. A032123
and A045406 were excluded after their 2026 proof references were
identified, and A162477 was excluded because its entry explicitly
credits a prior proof. Established generating functions, shorter
recurrences and linked literature are credited above. This is a
targeted source and prior-work audit, not an exhaustive worldwide
priority search. No OEIS entry was edited or externally declared closed.



\clearpage
\section{Forty-two empirical array recurrences}
\label{round4:sec:r4arrays}
The claims in this section are constant-coefficient recurrences marked
\emph{Empirical} in their OEIS entries, consulted on 20 October 2026.
They concern different widths, color requirements, or local constraints.
All 42 A-numbers are absent from the 308 selected entries in rounds 1--3.
They are 42 entry-level assertions proved by one general transfer method;
they are not presented as 42 unrelated methods.  The cited entries link
numerical data tables and record no proof of the selected recurrences.
This source check is not an exhaustive claim of priority over all literature.

\subsection{The arrays and their exact finite state spaces}
For integers $q\ge1$, $w\ge2$ and a set $S\subseteq\{1,2,3,4\}$,
consider arrays with $n+1$ rows and $w$ columns.  Every adjacent
$2\times2$ block must contain a number of distinct colors belonging to $S$.
Colors are renamed $0,1,\ldots$ in their order of first appearance when
the array is read row by row.  The superscript $=$ requires all $q$ colors
to occur; $\le$ permits any positive number at most $q$.
These are the precise conventions of the entries below.

For a row $r$, write $\operatorname{can}(r)$ for its renaming in order of
first appearance.  Let $\mathcal R_{j,w}$ contain all such canonical rows
with at most $j$ colors.  For $r\in\mathcal R_{j,w}$ let $k(r)$ be its
number of colors and put $m_j(r)=(j)_{k(r)}$.  Define the integer matrix
\begin{equation}\label{round4:eq:r4arr-matrix}
 T_{j,w,S}(r,s)=\#\left\{v\in\{0,\ldots,j-1\}^{w}:
 \begin{array}{l}
 \operatorname{can}(v)=s,\\
 \#\{r_i,r_{i+1},v_i,v_{i+1}\}\in S\quad(1\le i<w)
 \end{array}\right\}.
\end{equation}
All possible successor rows are included, including rows using colors
absent from the preceding row.  Simultaneous color permutations preserve
the local condition, so the number of legal completions of a row depends
only on its canonical form.  Consequently the labeled count is
\begin{equation}\label{round4:eq:r4arr-labeled}
 L_j(n)=m_j^{\mathsf T}T_{j,w,S}^{\,n}\mathbf1.
\end{equation}
The factor $m_j(r)$ counts the actual labelings of the first row, and each
subsequent multiplication adds one row.  This proves
\eqref{round4:eq:r4arr-labeled} by induction on the number of rows.

An array that uses exactly $q$ colors has exactly $q!$ labelings on a fixed
$q$-element alphabet.  Inclusion--exclusion for omitted colors gives
\begin{align}
 U_q^{=}(n)&=\sum_{j=1}^{q}
       \frac{(-1)^{q-j}}{j!(q-j)!}L_j(n),\label{round4:eq:r4arr-exact}\\
 U_q^{\le}(n)&=\sum_{j=1}^{q}
       \left(\frac1{j!}\sum_{h=0}^{q-j}\frac{(-1)^h}{h!}\right)L_j(n).
       \label{round4:eq:r4arr-atmost}
\end{align}
The second identity follows by summing the first over the actual number
of colors.  Thus there is no assumption that every allowed color occurs
unless the superscript is $=$.  Empty-alphabet counts vanish because every
array has at least one cell.

\subsection{An all-index certificate theorem}
For a displayed recurrence coefficient list $c=(c_1,\ldots,c_d)$ set
\[
 R(z)=z^d-\sum_{i=1}^{d}c_i z^{d-i},\qquad
 Q(x)=1-\sum_{i=1}^{d}c_i x^i.
\]
For every component with nonzero coefficient in
\eqref{round4:eq:r4arr-exact} or \eqref{round4:eq:r4arr-atmost}, the retained exact
integer certificate establishes
\begin{equation}\label{round4:eq:r4arr-cert}
 \boxed{\ T R(T)\mathbf1=\mathbf0.\ }
\end{equation}
For $n\ge d+1$, multiply this identity on the left by
$m_j^{\mathsf T}T^{n-d-1}$, and then take the linear combination in
\eqref{round4:eq:r4arr-exact} or \eqref{round4:eq:r4arr-atmost}.  This proves
\begin{equation}\label{round4:eq:r4arr-conclusion}
 \boxed{\ a(n)=\sum_{i=1}^{d}c_i a(n-i)\quad(n\ge d+1).\ }
\end{equation}
Every term on the right is defined, since all selected OEIS entries start
at $n=1$.  This is a universal proof from a finite matrix identity; it does
not infer a recurrence from agreement with a finite sequence prefix.

The complete certificate contains 90 distinct matrices, each with at most
122 states; the largest sum of component dimensions is 155.  The matrices
are in \texttt{r4\_arrays\_matrices.json}.  The verifier rebuilds every
state and edge directly from \eqref{round4:eq:r4arr-matrix}, compares the result
with the stored matrix, and checks every coordinate in
\eqref{round4:eq:r4arr-cert} by integer arithmetic.  No recurrence-fitting step
or numerical approximation is used.  The following pseudocode specifies
the finite identity check completely:
\begin{Verbatim}[fontsize=\small]
v[0] = [1 for each canonical row]
for k = 1,...,d+1:
    for each canonical row r:
        v[k][r] = sum(T[r,s]*v[k-1][s] for all s)
for each canonical row r:
    assert v[d+1][r] == sum(c[i]*v[d+1-i][r] for i=1,...,d)
\end{Verbatim}
The coefficient signatures below, the finite definition
\eqref{round4:eq:r4arr-matrix}, and this algorithm specify every certificate
without relying on empirical data.  Stored matrices make the arithmetic
easy to inspect and reproduce.  The source prefixes are compared only as
a separate check of indexing and interpretation; all 42 computed prefixes
are different, so there are no aliases within this section.

For $F(x)=\sum_{n\ge1}a(n)x^n$, the proved recurrence gives $Q(x)F(x)=P(x)$,
where $\deg P\le d$ and
\[
 [x^k]P(x)=\sum_{\substack{0\le i\le\min(k,d)\\k-i\ge1}}q_i a(k-i),
 \quad q_0=1,\quad q_i=-c_i\ (i\ge1).
\]
The values used here are calculated from the defining transfer matrices.
All 16 generating functions explicitly given as empirical or conjectural
on the retained source pages agree with this rational function exactly.
Their numerator and denominator coefficient lists are included in
\texttt{r4\_arrays\_certificates.json}.  They add no extra entry count.
Other assertions on an entry, such as a separately printed closed form,
are outside the selected claim unless explicitly addressed here.

\subsection{A fully displayed small example: A210100}
For A210100, $(q,w,S)=(3,2,\{1,3\})$ and at most three colors are allowed.
The only nonzero weights in \eqref{round4:eq:r4arr-atmost} are $1/2$ for $j=1$
and $1/6$ for $j=3$.  For $j=1$, $T=(1)$ and $m=(1)$.  For $j=3$ the
canonical rows are $00,01$, and direct enumeration gives
\[
 T=\begin{pmatrix}1&2\\1&4\end{pmatrix},\qquad
 m=\begin{pmatrix}3\\6\end{pmatrix}.
\]
For example, after $00$ there is one constant successor (the same color)
and two nonconstant successors using the other two colors; after $01$
there is one constant successor using the third color and four
nonconstant successors containing that color and one of $0,1$.
The characteristic polynomial of the latter matrix is $z^2-5z+2$.
Hence both components are annihilated by
\[
 (z-1)(z^2-5z+2)=z^3-6z^2+7z-2.
\]
Therefore $a(n)=6a(n-1)-7a(n-2)+2a(n-3)$ for $n\ge4$.
The transfer formula gives $a(1),a(2),a(3)=7,30,135$, yielding
\[
 F(x)=\frac{x(7-12x+4x^2)}{(1-x)(1-5x+2x^2)}.
\]
This proves both the empirical recurrence and the conjectural generating
function recorded at \url{https://oeis.org/A210100}.  It illustrates the
same exact certificate mechanism used for the larger state spaces.

\subsection{Individual recurrence signatures and parameters}
For each entry below, $D$ is the total dimension of the components with
nonzero weight.  The list gives $c_1,\ldots,c_d$ in order, including zero
coefficients.  Each signature satisfies \eqref{round4:eq:r4arr-cert} in every
component and therefore has the all-index range stated in
\eqref{round4:eq:r4arr-conclusion}.
\subsubsection{A210005}
\label{round4:r4entry:A210005}
The arrays have $n+1$ rows, width $w=2$, exactly $3$ colors up to order of first appearance, and $S=\{2,3\}$.
The recurrence \eqref{round4:eq:r4arr-conclusion} holds for $n\ge 5$ with order $d=4$ and component dimension $D=5$. Source: \url{https://oeis.org/A210005}.
\[ (c_1,\ldots,c_d)=(11,-16,-34,-12).\]
The source generating function also agrees exactly with the proved rational function $P/Q$.
\subsubsection{A210006}
\label{round4:r4entry:A210006}
The arrays have $n+1$ rows, width $w=3$, exactly $3$ colors up to order of first appearance, and $S=\{2,3\}$.
The recurrence \eqref{round4:eq:r4arr-conclusion} holds for $n\ge 7$ with order $d=6$ and component dimension $D=10$. Source: \url{https://oeis.org/A210006}.
\[ (c_1,\ldots,c_d)=(30,-109,-290,96,136,-48).\]
The source generating function also agrees exactly with the proved rational function $P/Q$.
\subsubsection{A210007}
\label{round4:r4entry:A210007}
The arrays have $n+1$ rows, width $w=4$, exactly $3$ colors up to order of first appearance, and $S=\{2,3\}$.
The recurrence \eqref{round4:eq:r4arr-conclusion} holds for $n\ge 11$ with order $d=10$ and component dimension $D=23$. Source: \url{https://oeis.org/A210007}.
\begin{Verbatim}[fontsize=\scriptsize,breaklines=true]
75, -30, -7001, -18089, 52126, 91400, -118464, -92992, 85632, -13824
\end{Verbatim}
\subsubsection{A210008}
\label{round4:r4entry:A210008}
The arrays have $n+1$ rows, width $w=5$, exactly $3$ colors up to order of first appearance, and $S=\{2,3\}$.
The recurrence \eqref{round4:eq:r4arr-conclusion} holds for $n\ge 20$ with order $d=19$ and component dimension $D=58$. Source: \url{https://oeis.org/A210008}.
\begin{Verbatim}[fontsize=\scriptsize,breaklines=true]
218, -672, -102513, 336965, 9441327, -64650448, 29276745, 733248704, -1733314135,
-814273096, 6961302814, -7307156292, -92994368, 3740262880, -1180127616, -466935296,
173254656, 18587648, -5406720
\end{Verbatim}
\subsubsection{A210054}
\label{round4:r4entry:A210054}
The arrays have $n+1$ rows, width $w=2$, at most $4$ colors up to order of first appearance, and $S=\{1,2,4\}$.
The recurrence \eqref{round4:eq:r4arr-conclusion} holds for $n\ge 5$ with order $d=4$ and component dimension $D=5$. Source: \url{https://oeis.org/A210054}.
\[ (c_1,\ldots,c_d)=(13,-48,52,-16).\]
The source generating function also agrees exactly with the proved rational function $P/Q$.
\subsubsection{A210055}
\label{round4:r4entry:A210055}
The arrays have $n+1$ rows, width $w=3$, at most $4$ colors up to order of first appearance, and $S=\{1,2,4\}$.
The recurrence \eqref{round4:eq:r4arr-conclusion} holds for $n\ge 7$ with order $d=6$ and component dimension $D=10$. Source: \url{https://oeis.org/A210055}.
\[ (c_1,\ldots,c_d)=(29,-266,902,-968,176,128).\]
The source generating function also agrees exactly with the proved rational function $P/Q$.
\subsubsection{A210056}
\label{round4:r4entry:A210056}
The arrays have $n+1$ rows, width $w=4$, at most $4$ colors up to order of first appearance, and $S=\{1,2,4\}$.
The recurrence \eqref{round4:eq:r4arr-conclusion} holds for $n\ge 12$ with order $d=11$ and component dimension $D=24$. Source: \url{https://oeis.org/A210056}.
\begin{Verbatim}[fontsize=\scriptsize,breaklines=true]
66, -1511, 14904, -58950, 21360, 330508, -503712, -93160, 369120, -9504, -69120
\end{Verbatim}
\subsubsection{A210057}
\label{round4:r4entry:A210057}
The arrays have $n+1$ rows, width $w=5$, at most $4$ colors up to order of first appearance, and $S=\{1,2,4\}$.
The recurrence \eqref{round4:eq:r4arr-conclusion} holds for $n\ge 29$ with order $d=28$ and component dimension $D=68$. Source: \url{https://oeis.org/A210057}.
\begin{Verbatim}[fontsize=\scriptsize,breaklines=true]
176, -12383, 457521, -9583668, 109074306, -432638439, -4419650349, 64230480793,
-252540170918, -595255299234, 8433913817808, -21427208559976, -38069222537892,
289460472761408, -336643109299632, -883702002797616, 2460214504431920, -361513807179872,
-4377499678974656, 3253938371960768, 2790969020407040, -3141539845016064,
-669537234834176, 1122282161785600, 46793169024000, -143518636118016, -1001036759040,
3148485230592
\end{Verbatim}
\subsubsection{A210079}
\label{round4:r4entry:A210079}
The arrays have $n+1$ rows, width $w=2$, exactly $4$ colors up to order of first appearance, and $S=\{1,3,4\}$.
The recurrence \eqref{round4:eq:r4arr-conclusion} holds for $n\ge 6$ with order $d=5$ and component dimension $D=7$. Source: \url{https://oeis.org/A210079}.
\[ (c_1,\ldots,c_d)=(17,-71,67,-8,-4).\]
The source generating function also agrees exactly with the proved rational function $P/Q$.
\subsubsection{A210080}
\label{round4:r4entry:A210080}
The arrays have $n+1$ rows, width $w=3$, exactly $4$ colors up to order of first appearance, and $S=\{1,3,4\}$.
The recurrence \eqref{round4:eq:r4arr-conclusion} holds for $n\ge 8$ with order $d=7$ and component dimension $D=15$. Source: \url{https://oeis.org/A210080}.
\[ (c_1,\ldots,c_d)=(42,-322,274,1887,-3124,76,1168).\]
\subsubsection{A210081}
\label{round4:r4entry:A210081}
The arrays have $n+1$ rows, width $w=4$, exactly $4$ colors up to order of first appearance, and $S=\{1,3,4\}$.
The recurrence \eqref{round4:eq:r4arr-conclusion} holds for $n\ge 20$ with order $d=19$ and component dimension $D=38$. Source: \url{https://oeis.org/A210081}.
\begin{Verbatim}[fontsize=\scriptsize,breaklines=true]
111, -1551, -11620, 269796, -346276, -10657824, 37514734, 105399645, -595931403,
25518343, 3128667030, -3607436824, -4027147296, 9243102400, -2810113344, -3425981120,
1936238208, 61521408, -60604416
\end{Verbatim}
\subsubsection{A210092}
\label{round4:r4entry:A210092}
The arrays have $n+1$ rows, width $w=6$, exactly $3$ colors up to order of first appearance, and $S=\{1,2\}$.
The recurrence \eqref{round4:eq:r4arr-conclusion} holds for $n\ge 53$ with order $d=52$ and component dimension $D=155$. Source: \url{https://oeis.org/A210092}.
\begin{Verbatim}[fontsize=\scriptsize,breaklines=true]
202, -14123, 380572, -1438616, -99274709, 1155724795, 8643300216, -166589849398,
-274497031426, 11515136131550, -1832292455645, -457188662253503, 336001687563753,
11327503957451248, -9113460943291491, -183665203305943843, 123571425473428963,
2003525661152721986, -990601469663807369, -14981189128638351022, 5154521066328382780,
78057653461845959313, -19250224843040180931, -287532733033011400278,
58417217021953987815, 757126134053256036905, -155693459113620706806,
-1433722830312214278540, 348602065891707730366, 1950547820706182800257,
-595414185159687071974, -1887300745304252275773, 725686784875169133118,
1269686648151447607271, -608745622735310244088, -568946968704783843248,
341945871757434843792, 155435605971618737334, -124110444593762302930,
-19948113914003483822, 27435918496367785268, -714926210446370056, -3322485605775126312,
498971087766745072, 180395158795278272, -44929945871997536, -3093399295952128,
1451280469739136, -24128771207680, -15867004438528, 437469315072, 60523806720
\end{Verbatim}
\subsubsection{A210100}
\label{round4:r4entry:A210100}
The arrays have $n+1$ rows, width $w=2$, at most $3$ colors up to order of first appearance, and $S=\{1,3\}$.
The recurrence \eqref{round4:eq:r4arr-conclusion} holds for $n\ge 4$ with order $d=3$ and component dimension $D=3$. Source: \url{https://oeis.org/A210100}.
\[ (c_1,\ldots,c_d)=(6,-7,2).\]
The source generating function also agrees exactly with the proved rational function $P/Q$.
\subsubsection{A210103}
\label{round4:r4entry:A210103}
The arrays have $n+1$ rows, width $w=5$, at most $3$ colors up to order of first appearance, and $S=\{1,3\}$.
The recurrence \eqref{round4:eq:r4arr-conclusion} holds for $n\ge 23$ with order $d=22$ and component dimension $D=42$. Source: \url{https://oeis.org/A210103}.
\begin{Verbatim}[fontsize=\scriptsize,breaklines=true]
39, -424, -100, 25342, -98910, -353224, 2683456, -679949, -24899773, 39419608, 84219756,
-230024712, -56603460, 492890768, -143861456, -452015392, 208017984, 175455296,
-65216128, -32317440, 2433024, 925696
\end{Verbatim}
\subsubsection{A210118}
\label{round4:r4entry:A210118}
The arrays have $n+1$ rows, width $w=2$, at most $4$ colors up to order of first appearance, and $S=\{2,4\}$.
The recurrence \eqref{round4:eq:r4arr-conclusion} holds for $n\ge 4$ with order $d=3$ and component dimension $D=5$. Source: \url{https://oeis.org/A210118}.
\[ (c_1,\ldots,c_d)=(10,-19,-14).\]
The source generating function also agrees exactly with the proved rational function $P/Q$.
\subsubsection{A210119}
\label{round4:r4entry:A210119}
The arrays have $n+1$ rows, width $w=3$, at most $4$ colors up to order of first appearance, and $S=\{2,4\}$.
The recurrence \eqref{round4:eq:r4arr-conclusion} holds for $n\ge 8$ with order $d=7$ and component dimension $D=10$. Source: \url{https://oeis.org/A210119}.
\[ (c_1,\ldots,c_d)=(24,-174,368,303,-888,-340,240).\]
The source generating function also agrees exactly with the proved rational function $P/Q$.
\subsubsection{A210129}
\label{round4:r4entry:A210129}
The arrays have $n+1$ rows, width $w=4$, at most $4$ colors up to order of first appearance, and $S=\{1,3,4\}$.
The recurrence \eqref{round4:eq:r4arr-conclusion} holds for $n\ge 12$ with order $d=11$ and component dimension $D=24$. Source: \url{https://oeis.org/A210129}.
\begin{Verbatim}[fontsize=\scriptsize,breaklines=true]
95, -1, -14266, 20367, 435937, -1093869, -1923590, 8479272, -8472184, 1936944, 631296
\end{Verbatim}
\subsubsection{A210130}
\label{round4:r4entry:A210130}
The arrays have $n+1$ rows, width $w=5$, at most $4$ colors up to order of first appearance, and $S=\{1,3,4\}$.
The recurrence \eqref{round4:eq:r4arr-conclusion} holds for $n\ge 33$ with order $d=32$ and component dimension $D=68$. Source: \url{https://oeis.org/A210130}.
\begin{Verbatim}[fontsize=\scriptsize,breaklines=true]
319, -12448, -307770, 19003857, -47940781, -8062797435, 95630724963, 949633978540,
-23587913597875, 72618564306781, 1459146419334837, -13658086185029599,
17603864407218421, 283502331705919989, -1303129591973246857, -279014374825669271,
13815024039103815024, -22992880145677918265, -44182446460842415541,
152128697073903743642, -14505606850313450236, -344933214811274872678,
261141694209971459720, 296779187097877685464, -389773691540956876144,
-44632888312447155024, 204261362669091285120, -46639179100779322176,
-30753710235878999040, 11357654705402941440, 637302851625194496, -414118546177001472
\end{Verbatim}
\subsubsection{A210136}
\label{round4:r4entry:A210136}
The arrays have $n+1$ rows, width $w=2$, exactly $4$ colors up to order of first appearance, and $S=\{1,2,4\}$.
The recurrence \eqref{round4:eq:r4arr-conclusion} holds for $n\ge 7$ with order $d=6$ and component dimension $D=7$. Source: \url{https://oeis.org/A210136}.
\[ (c_1,\ldots,c_d)=(18,-111,266,-180,-24,32).\]
The source generating function also agrees exactly with the proved rational function $P/Q$.
\subsubsection{A210137}
\label{round4:r4entry:A210137}
The arrays have $n+1$ rows, width $w=3$, exactly $4$ colors up to order of first appearance, and $S=\{1,2,4\}$.
The recurrence \eqref{round4:eq:r4arr-conclusion} holds for $n\ge 11$ with order $d=10$ and component dimension $D=15$. Source: \url{https://oeis.org/A210137}.
\begin{Verbatim}[fontsize=\scriptsize,breaklines=true]
40, -590, 3962, -11899, 12350, 3806, -12980, 4512, 1056, -256
\end{Verbatim}
\subsubsection{A210138}
\label{round4:r4entry:A210138}
The arrays have $n+1$ rows, width $w=4$, exactly $4$ colors up to order of first appearance, and $S=\{1,2,4\}$.
The recurrence \eqref{round4:eq:r4arr-conclusion} holds for $n\ge 20$ with order $d=19$ and component dimension $D=38$. Source: \url{https://oeis.org/A210138}.
\begin{Verbatim}[fontsize=\scriptsize,breaklines=true]
89, -3068, 51966, -443001, 1546826, 1718332, -24576952, 30663994, 107350909, -226371596,
-140409398, 501187032, -47667796, -343494224, 74492360, 92532416, -19372128, -8588160,
1382400
\end{Verbatim}
\subsubsection{A210139}
\label{round4:r4entry:A210139}
The arrays have $n+1$ rows, width $w=5$, exactly $4$ colors up to order of first appearance, and $S=\{1,2,4\}$.
The recurrence \eqref{round4:eq:r4arr-conclusion} holds for $n\ge 49$ with order $d=48$ and component dimension $D=109$. Source: \url{https://oeis.org/A210139}.
\begin{Verbatim}[fontsize=\scriptsize,breaklines=true]
235, -23547, 1326017, -46312598, 1033233394, -14208694709, 96208712962, 267300062626,
-11424901483364, 86498377558175, -31767158001113, -3617589382708047, 20693266346928210,
10888192487943923, -546652409630123972, 1681515865137829063, 4013358879516874873,
-33983772523000820266, 33111863198202404690, 257960652876911394845,
-727947325958421224552, -568221315698344240576, 4807060185026767170712,
-3133719554772897792592, -15172922164215860043716, 23613973519613807140264,
21941571653572275412952, -66462530816578512662688, -2423064632835660970480,
101519954129757123683008, -38482562973699671210464, -91174177912735792628544,
59493374679631303605248, 48485959363726885547520, -44332290747363931594240,
-14202007929251348447232, 18965402349674525710848, 1540656279076359243264,
-4753117347999007374336, 269751803727187415040, 663177631600048187392,
-96002320931723030528, -44409102700361863168, 9151174318820646912, 918429987777970176,
-240590747390312448, -2948910065123328, 1511272910684160
\end{Verbatim}
\subsubsection{A210158}
\label{round4:r4entry:A210158}
The arrays have $n+1$ rows, width $w=2$, at most $4$ colors up to order of first appearance, and $S=\{2,3\}$.
The recurrence \eqref{round4:eq:r4arr-conclusion} holds for $n\ge 5$ with order $d=4$ and component dimension $D=5$. Source: \url{https://oeis.org/A210158}.
\[ (c_1,\ldots,c_d)=(16,-19,-80,-36).\]
The source generating function also agrees exactly with the proved rational function $P/Q$.
\subsubsection{A210159}
\label{round4:r4entry:A210159}
The arrays have $n+1$ rows, width $w=3$, at most $4$ colors up to order of first appearance, and $S=\{2,3\}$.
The recurrence \eqref{round4:eq:r4arr-conclusion} holds for $n\ge 8$ with order $d=7$ and component dimension $D=10$. Source: \url{https://oeis.org/A210159}.
\[ (c_1,\ldots,c_d)=(55,-186,-1230,2723,159,-1604,468).\]
\subsubsection{A210160}
\label{round4:r4entry:A210160}
The arrays have $n+1$ rows, width $w=4$, at most $4$ colors up to order of first appearance, and $S=\{2,3\}$.
The recurrence \eqref{round4:eq:r4arr-conclusion} holds for $n\ge 17$ with order $d=16$ and component dimension $D=24$. Source: \url{https://oeis.org/A210160}.
\begin{Verbatim}[fontsize=\scriptsize,breaklines=true]
172, 2086, -61767, -221678, 6709264, -21628195, -28904236, 171575402, 51603, -460002202,
137911776, 512783040, -179902696, -200966160, 54584192, -885504
\end{Verbatim}
\subsubsection{A210161}
\label{round4:r4entry:A210161}
The arrays have $n+1$ rows, width $w=5$, at most $4$ colors up to order of first appearance, and $S=\{2,3\}$.
The recurrence \eqref{round4:eq:r4arr-conclusion} holds for $n\ge 41$ with order $d=40$ and component dimension $D=68$. Source: \url{https://oeis.org/A210161}.
\begin{Verbatim}[fontsize=\scriptsize,breaklines=true]
651, 2491, -3033274, 75213828, 2871779009, -168547329017, 3025868330721,
-14913244152146, -227860409342090, 3284704577622354, -4647468768394855,
-151912308978989829, 746051218249104195, 2759862594209092375, -24875375814765882871,
-12826708391148092367, 412811120432989292070, -271333439668914274397,
-3949195606845988432506, 4878572324314521633479, 22621015632021698182303,
-35706585467435371532954, -76058853813329983983298, 141387729457080507256936,
137247708151497848656089, -312567964628714713985554, -99662573283231077847331,
367361272953844857693772, -21661696257415105631489, -209997903963825729024970,
56676111775642707907746, 53600108892623612777764, -22156986755762524106826,
-5080148629481596023904, 3198197489711999573012, 14089783129953149268,
-171981655847047928312, 15964369340316597080, 1991492927343812928, -247274191930612320
\end{Verbatim}
\subsubsection{A210167}
\label{round4:r4entry:A210167}
The arrays have $n+1$ rows, width $w=2$, exactly $4$ colors up to order of first appearance, and $S=\{2,4\}$.
The recurrence \eqref{round4:eq:r4arr-conclusion} holds for $n\ge 6$ with order $d=5$ and component dimension $D=7$. Source: \url{https://oeis.org/A210167}.
\[ (c_1,\ldots,c_d)=(14,-55,22,132,56).\]
The source generating function also agrees exactly with the proved rational function $P/Q$.
\subsubsection{A210168}
\label{round4:r4entry:A210168}
The arrays have $n+1$ rows, width $w=3$, exactly $4$ colors up to order of first appearance, and $S=\{2,4\}$.
The recurrence \eqref{round4:eq:r4arr-conclusion} holds for $n\ge 11$ with order $d=10$ and component dimension $D=15$. Source: \url{https://oeis.org/A210168}.
\begin{Verbatim}[fontsize=\scriptsize,breaklines=true]
32, -362, 1648, -1561, -7568, 11440, 11360, -14768, -6400, 3840
\end{Verbatim}
\subsubsection{A210169}
\label{round4:r4entry:A210169}
The arrays have $n+1$ rows, width $w=4$, exactly $4$ colors up to order of first appearance, and $S=\{2,4\}$.
The recurrence \eqref{round4:eq:r4arr-conclusion} holds for $n\ge 25$ with order $d=24$ and component dimension $D=38$. Source: \url{https://oeis.org/A210169}.
\begin{Verbatim}[fontsize=\scriptsize,breaklines=true]
63, -1431, 12729, 3336, -726316, 2597781, 14573977, -80969613, -143169465, 1105985892,
755457488, -8134284268, -2199693420, 33945204080, 3581733232, -80566448928, -3198195264,
104642408704, 1573518080, -66575366656, -375523328, 15376932864, -161906688, -629932032
\end{Verbatim}
\subsubsection{A210260}
\label{round4:r4entry:A210260}
The arrays have $n+1$ rows, width $w=2$, exactly $4$ colors up to order of first appearance, and $S=\{3\}$.
The recurrence \eqref{round4:eq:r4arr-conclusion} holds for $n\ge 5$ with order $d=4$ and component dimension $D=7$. Source: \url{https://oeis.org/A210260}.
\[ (c_1,\ldots,c_d)=(12,-18,-64,-24).\]
The source generating function also agrees exactly with the proved rational function $P/Q$.
\subsubsection{A210261}
\label{round4:r4entry:A210261}
The arrays have $n+1$ rows, width $w=3$, exactly $4$ colors up to order of first appearance, and $S=\{3\}$.
The recurrence \eqref{round4:eq:r4arr-conclusion} holds for $n\ge 9$ with order $d=8$ and component dimension $D=15$. Source: \url{https://oeis.org/A210261}.
\begin{Verbatim}[fontsize=\scriptsize,breaklines=true]
32, -225, 0, 2856, -2952, -4100, 3664, -384
\end{Verbatim}
\subsubsection{A210262}
\label{round4:r4entry:A210262}
The arrays have $n+1$ rows, width $w=4$, exactly $4$ colors up to order of first appearance, and $S=\{3\}$.
The recurrence \eqref{round4:eq:r4arr-conclusion} holds for $n\ge 19$ with order $d=18$ and component dimension $D=38$. Source: \url{https://oeis.org/A210262}.
\begin{Verbatim}[fontsize=\scriptsize,breaklines=true]
63, -318, -15234, 103150, 1267432, -7283556, -49152712, 160312712, 877545008,
-931302192, -5537877952, 2212548864, 14304785408, -4034892800, -15800012800, 5312643072,
6186336256, -2686451712
\end{Verbatim}
\subsubsection{A210302}
\label{round4:r4entry:A210302}
The arrays have $n+1$ rows, width $w=2$, exactly $4$ colors up to order of first appearance, and $S=\{1,3\}$.
The recurrence \eqref{round4:eq:r4arr-conclusion} holds for $n\ge 6$ with order $d=5$ and component dimension $D=7$. Source: \url{https://oeis.org/A210302}.
\[ (c_1,\ldots,c_d)=(15,-57,41,10,-8).\]
The source generating function also agrees exactly with the proved rational function $P/Q$.
\subsubsection{A210303}
\label{round4:r4entry:A210303}
The arrays have $n+1$ rows, width $w=3$, exactly $4$ colors up to order of first appearance, and $S=\{1,3\}$.
The recurrence \eqref{round4:eq:r4arr-conclusion} holds for $n\ge 8$ with order $d=7$ and component dimension $D=15$. Source: \url{https://oeis.org/A210303}.
\[ (c_1,\ldots,c_d)=(35,-315,729,1132,-3804,688,1536).\]
\subsubsection{A210304}
\label{round4:r4entry:A210304}
The arrays have $n+1$ rows, width $w=4$, exactly $4$ colors up to order of first appearance, and $S=\{1,3\}$.
The recurrence \eqref{round4:eq:r4arr-conclusion} holds for $n\ge 21$ with order $d=20$ and component dimension $D=38$. Source: \url{https://oeis.org/A210304}.
\begin{Verbatim}[fontsize=\scriptsize,breaklines=true]
81, -1553, 2061, 172685, -899169, -6513363, 44110455, 105719944, -862490004, -709417008,
7693929440, 1058374912, -32657532864, 6969586944, 63976482560, -23949775872,
-51839258624, 18903613440, 14501412864, -3227516928
\end{Verbatim}
\subsubsection{A210313}
\label{round4:r4entry:A210313}
The arrays have $n+1$ rows, width $w=2$, exactly $4$ colors up to order of first appearance, and $S=\{2,3\}$.
The recurrence \eqref{round4:eq:r4arr-conclusion} holds for $n\ge 7$ with order $d=6$ and component dimension $D=7$. Source: \url{https://oeis.org/A210313}.
\[ (c_1,\ldots,c_d)=(24,-141,-24,718,768,216).\]
The source generating function also agrees exactly with the proved rational function $P/Q$.
\subsubsection{A210314}
\label{round4:r4entry:A210314}
The arrays have $n+1$ rows, width $w=3$, exactly $4$ colors up to order of first appearance, and $S=\{2,3\}$.
The recurrence \eqref{round4:eq:r4arr-conclusion} holds for $n\ge 11$ with order $d=10$ and component dimension $D=15$. Source: \url{https://oeis.org/A210314}.
\begin{Verbatim}[fontsize=\scriptsize,breaklines=true]
79, -1474, 1450, 39515, -30297, -122076, 99228, 43912, -53472, 11232
\end{Verbatim}
\subsubsection{A210315}
\label{round4:r4entry:A210315}
The arrays have $n+1$ rows, width $w=4$, exactly $4$ colors up to order of first appearance, and $S=\{2,3\}$.
The recurrence \eqref{round4:eq:r4arr-conclusion} holds for $n\ge 22$ with order $d=21$ and component dimension $D=38$. Source: \url{https://oeis.org/A210315}.
\begin{Verbatim}[fontsize=\scriptsize,breaklines=true]
237, -8494, -300237, 2482313, 58239270, -296444067, -2842201417, 12050668766,
41834205193, -197110345949, -109615007438, 1042214343696, -265321484840, -2283447354472,
1334391576272, 2102283712768, -1614381867264, -554944840192, 578116512768, -98061852672,
1530150912
\end{Verbatim}
\subsubsection{A210316}
\label{round4:r4entry:A210316}
The arrays have $n+1$ rows, width $w=5$, exactly $4$ colors up to order of first appearance, and $S=\{2,3\}$.
The recurrence \eqref{round4:eq:r4arr-conclusion} holds for $n\ge 51$ with order $d=50$ and component dimension $D=109$. Source: \url{https://oeis.org/A210316}.
\begin{Verbatim}[fontsize=\scriptsize,breaklines=true]
848, -122300, -5801743, 682041320, -1213046081, -1079782780850, 27573638241302,
79725630193847, -11873109205022518, 131395804871513578, 986201712895823243,
-29652374531996895310, 132452849847042036360, 1622629993578919995350,
-18253509106659216838968, 6849069634272644477350, 714587571710513391501643,
-2760444329694926932104549, -10516539662994187322703315, 86912495830872109675164069,
-9563542017692602626829178, -1248365123628345697575639065, 2249181879462811987679557346,
9086435408361391418501120428, -30736207455851587697978075075,
-25923589316759701283979550795, 202500201415812975536465857367,
-72715906863721870790944600743, -720781514789318502643514797441,
801786923731748401090054422513, 1286931802588307290422443408052,
-2522061032842347970441473677732, -674119438641704924434781252930,
3733616049201706012714214101006, -1013770420222518200606193339716,
-2707380660785061103024598191772, 1579508813514184554109584102836,
942689284708178387720861782888, -859718979462997536027199561344,
-125446929953546119395609975056, 229541474235873867982539795760,
-8019298577799431458900026560, -31686621384563077338679511744,
3853494757916521671974788352, 2204406306231678899205504512,
-367673092093538607145367552, -67950978187980533942738944, 13061961692650600938110976,
613878827815294878744576, -133694231899508024279040
\end{Verbatim}
\subsubsection{A210397}
\label{round4:r4entry:A210397}
The arrays have $n+1$ rows, width $w=2$, at most $4$ colors up to order of first appearance, and $S=\{1,3\}$.
The recurrence \eqref{round4:eq:r4arr-conclusion} holds for $n\ge 4$ with order $d=3$ and component dimension $D=5$. Source: \url{https://oeis.org/A210397}.
\[ (c_1,\ldots,c_d)=(10,-5,-4).\]
The source generating function also agrees exactly with the proved rational function $P/Q$.
\subsubsection{A210398}
\label{round4:r4entry:A210398}
The arrays have $n+1$ rows, width $w=3$, at most $4$ colors up to order of first appearance, and $S=\{1,3\}$.
The recurrence \eqref{round4:eq:r4arr-conclusion} holds for $n\ge 6$ with order $d=5$ and component dimension $D=10$. Source: \url{https://oeis.org/A210398}.
\[ (c_1,\ldots,c_d)=(28,-123,-20,500,-384).\]
The source generating function also agrees exactly with the proved rational function $P/Q$.
\subsubsection{A210400}
\label{round4:r4entry:A210400}
The arrays have $n+1$ rows, width $w=5$, at most $4$ colors up to order of first appearance, and $S=\{1,3\}$.
The recurrence \eqref{round4:eq:r4arr-conclusion} holds for $n\ge 33$ with order $d=32$ and component dimension $D=68$. Source: \url{https://oeis.org/A210400}.
\begin{Verbatim}[fontsize=\scriptsize,breaklines=true]
222, -14466, 294924, 4656995, -258980402, 1519478244, 65765309856, -899946571108,
-6066410464024, 162105406266400, -88886529769536, -13947758522827904, 56368494845199616,
595808254992314112, -4168954910728206336, -10590192881040593920, 140772489264317126656,
-38165648086702518272, -2466112169355411496960, 4370610701672862203904,
22131193152696389959680, -66520641034825424371712, -84130431929983157141504,
444218891805352251097088, -24524753263238930497536, -1357521916200660102021120,
906383329842391263739904, 1679955754896158917394432, -1768503113779591826636800,
-469035266983925668380672, 970910637009369175162880, -255356668034757049712640
\end{Verbatim}


\clearpage
\section{Reproduction and verification}
\label{round4:sec:r4reproduce}
The accompanying archive is self-contained. It contains the complete
standalone manuscript, the 100-entry ledger, four proof fragments,
four main verifier programs, exact certificates, complete array matrices,
and the count and prior-work audits. The proofs state their hypotheses,
indexing conventions and selected scopes explicitly.

\subsection{Running the verification}
Use Python 3.10 or later with SymPy. The delivered run was checked with
Python 3.12 and SymPy 1.14.0. From the extracted directory run:
\begin{Verbatim}[fontsize=\small]
python3 r4_run_all.py
\end{Verbatim}
This runs the component verifiers and independent array enumeration,
checks all 100 unique A-numbers against the preceding 308, and writes
a combined JSON summary. All mathematical checks use exact integers or
rational expressions. No network connection is needed.

\begin{center}\small
\begin{tabularx}{\textwidth}{@{}lX@{}}
\toprule
Program&What is checked\\
\midrule
\texttt{r4\_root\_verify.py}&Fourteen congruence packages, including
exact rational identities and supplementary sequence checks.\\
\texttt{r4\_arith\_verify.py}&Nine arithmetic packages, exact floor
certificates, determinant formulas, tail enclosures and floor sums.\\
\texttt{r4\_gf\_verify.py}&Thirty-five symbolic certificates, recurrence
coefficients, boundary ranges and independently expanded initial terms.\\
\texttt{r4\_arrays\_verify.py}&Forty-two recurrences, all 90 complete
matrices reconstructed from their definitions, and 16 source generating
functions.\\
\texttt{r4\_gf\_review\_arrays.py}&A separate direct enumeration of
small global color patterns, checking the transfer interpretation.\\
\bottomrule
\end{tabularx}
\end{center}

\subsection{Why the finite certificates prove infinite statements}
For an array recurrence, the defining count is a fixed linear combination
of $m^{\mathsf T}T^n\mathbf1$. The complete identity
$TR(T)\mathbf1=0$ remains true after multiplication by every nonnegative
power of $T$, proving the recurrence for all admissible indices.
Checking a long initial prefix alone would not give this implication.

For the algebraic generating functions, the recurrence is translated
to a differential operator. The verifier computes its exact action in
a rational radical basis. All nonconstant basis components vanish and
the remaining boundary polynomial is displayed. Its degree gives the
precise range in which every recurrence coefficient is zero. This also
detects the initial exception in \oeis{A243642}. The three special
recurrence proofs and the factorial-shift identity have
separate exact derivations in their sections.

For the remaining statements, the written arguments establish the
quantified result. The accompanying finite checks guard against
transcription and indexing errors and are not substituted for those
arguments. For example, the persistent mod-4 obstruction in
\oeis{A126156} follows from an operator identity for every index, whereas
the extra computed initial terms are a consistency check.

\subsection{Rebuilding and integrity}
The delivered \texttt{oeis\_hunt\_round4.tex} contains every proof directly
and uses no external \verb|\input| files. Compile it twice with pdfLaTeX
and the packages listed in its preamble. To reconstruct the TEX and the
ledger from the included component files, run:
\begin{Verbatim}[fontsize=\small]
python3 r4_build_dossier.py
\end{Verbatim}
The archive's README describes optional certificate rebuilds. Its
\texttt{manifest.json} lists SHA-256 digests. Rerunning verifiers can
regenerate reports and timing fields; those report digests may then change.

\subsection{Source and priority records}
Each entry in the catalogue links to its primary OEIS page. The portable
records retain mathematical statements, source attribution and consulted
URLs; full scraped webpages are omitted. The component audit files record
known prior results, excluded candidates and limits of the source review.
Three successful array candidates were held out to keep the selected
catalogue at exactly 100; they contribute no count or unprovided proof
to the final packet.

The count after this round is $107+101+100+100=408$ selected entry-level
packages under the project's counting convention. This is a project
ledger total, not a claim that 408 entries have been edited or formally
closed by the OEIS editors.
% ===== END ORIGINAL ROUND 4 =====

% ===== BEGIN ORIGINAL ROUND 5 =====
\clearpage
\setcounter{section}{0}\setcounter{equation}{0}\setcounter{theorem}{0}
\setcounter{table}{0}\setcounter{figure}{0}
\normalsize\renewcommand{\arraystretch}{1}\renewcommand{\baselinestretch}{1}\selectfont
\part[Round 5: Universal congruences, arithmetic, theta values, and priority corrections]{Universal congruences, arithmetic, theta values, and priority corrections}
\label{combined:round5}
\rhead{\small Round 5 / KEIO}
\section*{Round overview}
\addcontentsline{toc}{section}{Round 5 overview and scope}
This fifth-round dossier treats 100 selected assertion packages on
100 distinct OEIS entries, disjoint from the 408 A-numbers previously
selected by this project. The outcomes are 90 proofs of the original
selected claims, two refutations, and eight corrected statements proved
in full. The corrections comprise six recurrence starting-index
qualifications, one malformed recurrence factor, and one theta-sum
index normalization. Several assertions on one entry count only once.

The four parts contain 13 packages on formal series, periodicity and
theta evaluations; 12 on arithmetic, averages and bounds; 30 on
generating functions; and 45 on constrained arrays. The array results
are proved by complete integer transfer matrices and vector identities.
The generating-function recurrences are proved by exact differential
identities, with every boundary term retained. The remaining results
have written quantified proofs, accompanied by supplementary calculations.

The two refutations are explicit. For \oeis{A051699}, the nearest-prime
distance sum at $N=2^{20}$ is $5542944$, exceeding the stated upper
bound $\tfrac38N\log N$. For \oeis{A129366}, the proposed recurrence
has residual $-73968$ at $n=6$ and cannot hold on any eventual tail;
a corrected recurrence is proved. The eight correction cases are
separated from these genuine counterexamples in the catalogue.

\subsection*{Scope and counting convention}
Each catalogue row specifies one assertion package on one A-number.
Different members of a parameter family can use a common theorem;
the count does not mean 100 unrelated proof methods. The primary
statements were labeled conjectural, empirical, or expressed as
unproved bounds in the consulted sources. Established generating
functions and classical identities are credited as inputs. A selected
comment can be settled while other questions on the same entry remain
outside scope. For example, this dossier does not settle the asymptotic
constant in \oeis{A051699}, the mod-3 or asymptotic assertions in
\oeis{A397242}, or every modular statement on \oeis{A396804}.

The six boundary cases have no starting index specified in the source.
We give the precise valid tail and exhibit the earlier admissible
exceptions; we do not attribute an explicit false starting-index claim
to the author. For \oeis{A279013}, an undefined printed token is replaced
by a specified factor and the repaired formula is proved. For
\oeis{A243763}, both the original-index and zero-based evaluations are
proved, and their factor $e^{\pi}$ is explained.

\subsection*{Retrospective correction: five Round IV entries}
\textbf{The earlier new-resolution assessment is withdrawn for
\oeis{A176605}, \oeis{A176607}, \oeis{A176610}, \oeis{A176611},
and \oeis{A176612}.} Georg Fischer's general derivation dated
26 January 2020 already implies their selected recurrences.
The Round IV proofs remain mathematically valid, but these five
claims should not have been presented as newly resolved. The exact
parameter substitutions and operator comparisons are documented in
Section~\ref{round5:sec:r5retrospective} and the accompanying audit certificate.
None of these five entries contributes to the present 100.
The preceding total of 408 is therefore used only as an exclusion
ledger of selected A-numbers, not as a certified total of new discoveries.

\subsection*{Evidence and limits}
Primary OEIS pages, relevant linked work and targeted literature searches
were consulted on 20 October 2026; retrieved pages may be cached.
The OEIS Open comparison covers only Appendix A: 100 costliest claims
on 97 distinct A-numbers, with no present overlap. It does not cover
the full project or constitute an exhaustive priority search.

The results are mathematical claims supported by the written arguments
and certificates here. No worldwide-priority or OEIS editorial-acceptance
claim is made, and no OEIS edit or submission was performed. The
cross-reviews were performed by other agents in the same AI-assisted
workflow; they are not external peer review or proof-assistant formalization.

\clearpage
\section{Catalogue of the 100 selected entries}
\label{round5:sec:r5catalogue}
\textbf{P}: original selected claim proved. \textbf{R}: refuted.
\textbf{B}: boundary corrected. \textbf{F}: malformed formula corrected.
\textbf{I}: index normalization corrected. Every corrected statement
is proved. Page references lead to the individual argument or recurrence
signature; the JSON ledger retains exact scopes and attribution.

\small
\begin{longtable}{@{}r p{21mm} p{103mm} c r@{}}
\toprule
No.&Entry&Selected assertion package&Result&Page\\
\midrule\endfirsthead
\toprule No.&Entry&Selected assertion package&Result&Page\\
\midrule\endhead
\midrule\multicolumn{5}{r}{\emph{Continued on the next page}}\\\endfoot
\bottomrule\endlastfoot
\multicolumn{5}{@{}l}{\rule{0pt}{4ex}\textbf{Congruences, periodicity and theta values}}\\
1&\oeis{A338632}&All three continued-fraction congruences, modulo 4 and 3&P&\pageref{round5:r5entry:A338632}\\
2&\oeis{A338636}&All three continued-fraction congruences, modulo 8 and 3&P&\pageref{round5:r5entry:A338636}\\
3&\oeis{A397242}&Odd precisely when the index plus one is a power of two&P&\pageref{round5:r5entry:A397242}\\
4&\oeis{A396804}&The coefficient at index n is congruent to n modulo 4&P&\pageref{round5:r5entry:A396804}\\
5&\oeis{A383165}&Eventual periodicity modulo every positive integer&P&\pageref{round5:r5entry:A383165}\\
6&\oeis{A383166}&Eventual periodicity modulo every positive integer&P&\pageref{round5:r5entry:A383166}\\
7&\oeis{A179929}&An eventual modular period dividing Euler phi&P&\pageref{round5:r5entry:A179929}\\
8&\oeis{A210063}&The empirical exponential-weighted theta sum&P&\pageref{round5:r5entry:A210063}\\
9&\oeis{A243763}&Two theta sums with original and shifted indices distinguished&I&\pageref{round5:r5entry:A243763}\\
10&\oeis{A107635}&The empirical exponential-weighted theta sum&P&\pageref{round5:r5entry:A107635}\\
11&\oeis{A002512}&The empirical exponential-weighted theta sum&P&\pageref{round5:r5entry:A002512}\\
12&\oeis{A134416}&The empirical exponential-weighted theta sum&P&\pageref{round5:r5entry:A134416}\\
13&\oeis{A030204}&The empirical exponential-weighted theta sum&P&\pageref{round5:r5entry:A030204}\\
\multicolumn{5}{@{}l}{\rule{0pt}{4ex}\textbf{Arithmetic, averages and exact bounds}}\\
14&\oeis{A010786}&The floor-factorial logarithmic constant and a uniform error bound&P&\pageref{round5:r5entry:A010786}\\
15&\oeis{A034444}&The 2022 Hultman-polynomial Dirichlet-series family&P&\pageref{round5:r5entry:A034444}\\
16&\oeis{A051699}&Exact counterexample to the stated uniform nearest-prime bound&R&\pageref{round5:r5entry:A051699}\\
17&\oeis{A129455}&The cubic gcd triangle is a binomial coefficient times a sixth power&P&\pageref{round5:r5entry:A129455}\\
18&\oeis{A161789}&Existence of the empirical mean and a certified interval&P&\pageref{round5:r5entry:A161789}\\
19&\oeis{A187011}&Prime-distance determinant sign and nth-root divergence&P&\pageref{round5:r5entry:A187011}\\
20&\oeis{A276646}&The decimal-prefix average has sharp upper bound 0.55&P&\pageref{round5:r5entry:A276646}\\
21&\oeis{A278083}&Primitive-quadruple multiplicativity and the three-square limit&P&\pageref{round5:r5entry:A278083}\\
22&\oeis{A280098}&Complete classification of the divisor-sum moduli&P&\pageref{round5:r5entry:A280098}\\
23&\oeis{A300984}&Prime support and exponent restrictions for the divisor-sum condition&P&\pageref{round5:r5entry:A300984}\\
24&\oeis{A339387}&The remainder-sum inequality, strict at nonsquarefree indices&P&\pageref{round5:r5entry:A339387}\\
25&\oeis{A394454}&Every prime divides infinitely many central binomial coefficients&P&\pageref{round5:r5entry:A394454}\\
\multicolumn{5}{@{}l}{\rule{0pt}{4ex}\textbf{Generating-function certificates}}\\
26&\oeis{A025259}&Conjectured polynomial recurrence of order 4; exact starting index supplied&B&\pageref{round5:r5entry:A025259}\\
27&\oeis{A025272}&Conjectured polynomial recurrence of order 5&P&\pageref{round5:r5entry:A025272}\\
28&\oeis{A034267}&Conjectured polynomial recurrence of order 2&P&\pageref{round5:r5entry:A034267}\\
29&\oeis{A089324}&Conjectured polynomial recurrence of order 6&P&\pageref{round5:r5entry:A089324}\\
30&\oeis{A097331}&Conjectured polynomial recurrence of order 3&P&\pageref{round5:r5entry:A097331}\\
31&\oeis{A102319}&Conjectured polynomial recurrence of order 6&P&\pageref{round5:r5entry:A102319}\\
32&\oeis{A104624}&Conjectured polynomial recurrence of order 3&P&\pageref{round5:r5entry:A104624}\\
33&\oeis{A107587}&Conjectured polynomial recurrence of order 5&P&\pageref{round5:r5entry:A107587}\\
34&\oeis{A108600}&Conjectured polynomial recurrence of order 6&P&\pageref{round5:r5entry:A108600}\\
35&\oeis{A116388}&Conjectured polynomial recurrence of order 7&P&\pageref{round5:r5entry:A116388}\\
36&\oeis{A129366}&The proposed recurrence fails on every tail; a replacement is proved&R&\pageref{round5:r5entry:A129366}\\
37&\oeis{A134424}&Conjectured polynomial recurrence of order 10&P&\pageref{round5:r5entry:A134424}\\
38&\oeis{A162482}&Conjectured polynomial recurrence of order 6&P&\pageref{round5:r5entry:A162482}\\
39&\oeis{A164586}&Conjectured polynomial recurrence of order 7&P&\pageref{round5:r5entry:A164586}\\
40&\oeis{A174013}&Conjectured polynomial recurrence of order 7&P&\pageref{round5:r5entry:A174013}\\
41&\oeis{A182879}&Conjectured polynomial recurrence of order 5&P&\pageref{round5:r5entry:A182879}\\
42&\oeis{A182894}&Conjectured polynomial recurrence of order 7&P&\pageref{round5:r5entry:A182894}\\
43&\oeis{A182904}&Conjectured polynomial recurrence of order 7&P&\pageref{round5:r5entry:A182904}\\
44&\oeis{A189176}&Conjectured polynomial recurrence of order 2; exact starting index supplied&B&\pageref{round5:r5entry:A189176}\\
45&\oeis{A190166}&Conjectured polynomial recurrence of order 6&P&\pageref{round5:r5entry:A190166}\\
46&\oeis{A218185}&Conjectured polynomial recurrence of order 4; exact starting index supplied&B&\pageref{round5:r5entry:A218185}\\
47&\oeis{A239201}&Conjectured polynomial recurrence of order 3&P&\pageref{round5:r5entry:A239201}\\
48&\oeis{A243764}&Conjectured polynomial recurrence of order 6; exact starting index supplied&B&\pageref{round5:r5entry:A243764}\\
49&\oeis{A262768}&Conjectured polynomial recurrence of order 2; exact starting index supplied&B&\pageref{round5:r5entry:A262768}\\
50&\oeis{A270490}&Conjectured polynomial recurrence of order 4&P&\pageref{round5:r5entry:A270490}\\
51&\oeis{A270530}&Conjectured polynomial recurrence of order 4&P&\pageref{round5:r5entry:A270530}\\
52&\oeis{A273351}&Conjectured polynomial recurrence of order 5&P&\pageref{round5:r5entry:A273351}\\
53&\oeis{A279013}&Repair and proof of the malformed recurrence factor&F&\pageref{round5:r5entry:A279013}\\
54&\oeis{A279014}&Conjectured polynomial recurrence of order 4; exact starting index supplied&B&\pageref{round5:r5entry:A279014}\\
55&\oeis{A387084}&Conjectured polynomial recurrence of order 4&P&\pageref{round5:r5entry:A387084}\\
\multicolumn{5}{@{}l}{\rule{0pt}{4ex}\textbf{Finite-state array certificates}}\\
56&\oeis{A209505}&Empirical array recurrence of order 2&P&\pageref{round5:r5entry:A209505}\\
57&\oeis{A209506}&Empirical array recurrence of order 3&P&\pageref{round5:r5entry:A209506}\\
58&\oeis{A209507}&Empirical array recurrence of order 5&P&\pageref{round5:r5entry:A209507}\\
59&\oeis{A209508}&Empirical array recurrence of order 7&P&\pageref{round5:r5entry:A209508}\\
60&\oeis{A209509}&Empirical array recurrence of order 13&P&\pageref{round5:r5entry:A209509}\\
61&\oeis{A209530}&Empirical array recurrence of order 3&P&\pageref{round5:r5entry:A209530}\\
62&\oeis{A209531}&Empirical array recurrence of order 3&P&\pageref{round5:r5entry:A209531}\\
63&\oeis{A209546}&Empirical array recurrence of order 4&P&\pageref{round5:r5entry:A209546}\\
64&\oeis{A209547}&Empirical array recurrence of order 6&P&\pageref{round5:r5entry:A209547}\\
65&\oeis{A209548}&Empirical array recurrence of order 11&P&\pageref{round5:r5entry:A209548}\\
66&\oeis{A209549}&Empirical array recurrence of order 20&P&\pageref{round5:r5entry:A209549}\\
67&\oeis{A209710}&Empirical array recurrence of order 19&P&\pageref{round5:r5entry:A209710}\\
68&\oeis{A209711}&Empirical array recurrence of order 64&P&\pageref{round5:r5entry:A209711}\\
69&\oeis{A209851}&Empirical array recurrence of order 6&P&\pageref{round5:r5entry:A209851}\\
70&\oeis{A209852}&Empirical array recurrence of order 20&P&\pageref{round5:r5entry:A209852}\\
71&\oeis{A209853}&Empirical array recurrence of order 68&P&\pageref{round5:r5entry:A209853}\\
72&\oeis{A209906}&Empirical array recurrence of order 6&P&\pageref{round5:r5entry:A209906}\\
73&\oeis{A209907}&Empirical array recurrence of order 19&P&\pageref{round5:r5entry:A209907}\\
74&\oeis{A209908}&Empirical array recurrence of order 65&P&\pageref{round5:r5entry:A209908}\\
75&\oeis{A209953}&Empirical array recurrence of order 3&P&\pageref{round5:r5entry:A209953}\\
76&\oeis{A209955}&Empirical array recurrence of order 14&P&\pageref{round5:r5entry:A209955}\\
77&\oeis{A209956}&Empirical array recurrence of order 36&P&\pageref{round5:r5entry:A209956}\\
78&\oeis{A210070}&Empirical array recurrence of order 6&P&\pageref{round5:r5entry:A210070}\\
79&\oeis{A210071}&Empirical array recurrence of order 20&P&\pageref{round5:r5entry:A210071}\\
80&\oeis{A210072}&Empirical array recurrence of order 71&P&\pageref{round5:r5entry:A210072}\\
81&\oeis{A210082}&Empirical array recurrence of order 53&P&\pageref{round5:r5entry:A210082}\\
82&\oeis{A210128}&Empirical array recurrence of order 5&P&\pageref{round5:r5entry:A210128}\\
83&\oeis{A210149}&Empirical array recurrence of order 5&P&\pageref{round5:r5entry:A210149}\\
84&\oeis{A210150}&Empirical array recurrence of order 13&P&\pageref{round5:r5entry:A210150}\\
85&\oeis{A210151}&Empirical array recurrence of order 48&P&\pageref{round5:r5entry:A210151}\\
86&\oeis{A210263}&Empirical array recurrence of order 50&P&\pageref{round5:r5entry:A210263}\\
87&\oeis{A210269}&Empirical array recurrence of order 3&P&\pageref{round5:r5entry:A210269}\\
88&\oeis{A210270}&Empirical array recurrence of order 9&P&\pageref{round5:r5entry:A210270}\\
89&\oeis{A210271}&Empirical array recurrence of order 25&P&\pageref{round5:r5entry:A210271}\\
90&\oeis{A210272}&Empirical array recurrence of order 69&P&\pageref{round5:r5entry:A210272}\\
91&\oeis{A210305}&Empirical array recurrence of order 53&P&\pageref{round5:r5entry:A210305}\\
92&\oeis{A210328}&Empirical array recurrence of order 6&P&\pageref{round5:r5entry:A210328}\\
93&\oeis{A210329}&Empirical array recurrence of order 20&P&\pageref{round5:r5entry:A210329}\\
94&\oeis{A210330}&Empirical array recurrence of order 72&P&\pageref{round5:r5entry:A210330}\\
95&\oeis{A210348}&Empirical array recurrence of order 5&P&\pageref{round5:r5entry:A210348}\\
96&\oeis{A210349}&Empirical array recurrence of order 20&P&\pageref{round5:r5entry:A210349}\\
97&\oeis{A210350}&Empirical array recurrence of order 72&P&\pageref{round5:r5entry:A210350}\\
98&\oeis{A210383}&Empirical array recurrence of order 3&P&\pageref{round5:r5entry:A210383}\\
99&\oeis{A210384}&Empirical array recurrence of order 6&P&\pageref{round5:r5entry:A210384}\\
100&\oeis{A210386}&Empirical array recurrence of order 42&P&\pageref{round5:r5entry:A210386}\\
\end{longtable}\normalsize

\Needspace{6\baselineskip}
\section{Formal-series congruences, periodicity, and theta values}
\label{round5:r5:root}

The thirteen entry-level packages in this section use exact formal identities.
The classical theta evaluations and the finite-ring argument are credited
tools, not claims of new general theorems. Every selected assertion was still
marked conjectural or empirical in the retrieved OEIS entry. Other assertions
in the same entry are outside scope unless explicitly included below.

\subsection{Two implicit continued fractions}
Let $A_r(x)$, $r=1,2$, be the series with constant term $1$ determined by
\[
 1=A_r-\cfrac{x}{A_r-\cfrac{3^r x}{A_r-
       \cfrac{5^r x}{A_r-\cfrac{7^r x}{\ddots}}}}.
\]
The cases $r=1,2$ define \oeis{A338632} and \oeis{A338636}, respectively.
Every denominator has constant term $1$. Changing a tail beyond depth $N$
changes the fraction only in degrees at least $N$. Thus the fraction exists
formally. The equation is $A_r=1+xH(A_r,x)$, so its coefficient of degree $n$
determines $[x^n]A_r$ from earlier coefficients, over any commutative ring.
This proves existence, uniqueness, and integrality without division by $n$.

\begin{theorem}[A338632]\label{round5:r5entry:A338632}
For $n>1$, $a(n)\equiv2\pmod4$. For $n>0$, its residue modulo $3$ is
$1$ exactly when the ternary expansion of $n$ ends in $1$ and all other digits
are $0$ or $1$. The residue is $2$ exactly when the final digit is $0$ or $2$
and all other digits are $0$ or $1$; otherwise it is $0$.
These are all three congruence conjectures selected from the entry.
\end{theorem}
\begin{proof}
Modulo $4$, the successive numerators have weights $1,3,1,3,\ldots$.
Let the two tail values be $U,V$. Then $U=A_1-x/V$,
$V=A_1-3x/U$, and the defining equation says $U=1$. Therefore
\[
 A_1^2-(1+3x)A_1+2x=0\pmod4.
\]
This polynomial equation has a unique root with constant term $1$: its
derivative with respect to $A_1$ has constant term $1$ modulo $4$.
For $R=1+x+2x^2/(1-x)$, direct multiplication gives
\[
 R^2-(1+3x)R+2x=\frac{4x^4}{(1-x)^2}.
\]
Hence $A_1\equiv R\pmod4$, proving the first statement.

Modulo $3$, the second numerator vanishes, so $1=A_1-x/A_1$ and
$A_1^2-A_1-x=0$. Put
\[
 H(x)=(1+x)^2\prod_{j\ge1}(1+x^{3^j})\quad\text{in }\mathbb F_3[[x]].
\]
For the product ending at $j=M$, squaring and using
$(1+z)^3=1+z^3$ gives $(1+x)(1+x^{3^{M+1}})$.
Consequently $H^2=1+x$. The unique root is $A_1=(1+H)/2$.
In positive degree, the coefficient of $H$ is $2$ for final ternary digit
$1$, is $1$ for final digit $0$ or $2$, and is $0$ if a higher digit is $2$.
Multiplication by $1/2=2$ gives the stated residues.
The digit descriptions are the established descriptions of
\oeis{A191107} and \oeis{A186776}; their own assertions add no count here.
\end{proof}

\begin{theorem}[A338636]\label{round5:r5entry:A338636}
For $n>1$, $a(n)\equiv0\pmod8$. The two mod-$3$ support descriptions
in the preceding theorem hold for this sequence as well.
\end{theorem}
\begin{proof}
Every odd square is $1$ modulo $8$. Thus the whole continued fraction
has one repeated tail $U=A_2-x/U$. Since $U=1$, it follows that
$A_2=1+x\pmod8$. Modulo $3$, the numerator $3^2x$ vanishes.
The same equation $A_2^2-A_2-x=0$ and the same uniqueness and digit
argument apply. This proves all three displayed congruence conjectures.
\end{proof}

\subsection{A397242: parity from diagonal powers}
\begin{theorem}\label{round5:r5entry:A397242}
The coefficient $a(n)$ of \oeis{A397242}, including $a(0)=1$, is odd
if and only if $n+1$ is a power of $2$.
\end{theorem}
\begin{proof}
We first reproduce, with attribution to Round III, the integer-series lemma
used there for \oeis{A397241}; that earlier entry is not counted again.
Let $B(t)=1+t+\cdots$ satisfy
\[
 n[t^n]B(t)^n=(n-1)[t^n]B(t)^{n+1}\qquad(n>1).
\]
The coefficient of $b(n)$ in the difference is
$n^2-(n-1)(n+1)=1$. Therefore this defines a unique integer series.
Over $\mathbb F_2$, $B=(1-t)^{-1}$ satisfies it. Indeed the required
vanishing is $\binom{2n}{n}$ when $n$ is even and
$\binom{2n-1}{n}$ when $n>1$ is odd. Their $2$-adic valuations are
$s_2(n)$ and $s_2(n)-1$, respectively; an odd $n>1$ has at least two
binary $1$'s. Thus every coefficient of $B$ is odd.

For completeness we derive the change of variables relating $B$ to the
present entry. Define the unique integer series $A=1+x+\cdots$ by
$A(x)=B(xA(x))$. Write $t=xA(x)$, so its inverse is $x=t/B(t)$.
Lagrange inversion gives, for $n>1$,
\[
 [t^n]B(t)^n=n[x^n]\log A(x),\qquad
 [t^n]B(t)^{n+1}=(n+1)[x^n]A(x).
\]
The diagonal relation therefore gives
$[x^n]\log A=(n^2-1)a(n)/n^2$.
Together with $a(0)=a(1)=1$, this is exactly the defining equation
\[
 A(x)=\exp\left(x+\sum_{n\ge2}\frac{n^2-1}{n^2}a(n)x^n\right).
\]
It determines $a(n)$ uniquely over $\mathbb Q$, since the coefficient
left after moving its two occurrences to one side is $1/n^2$.
Hence the constructed integral series is the source sequence.

Reducing $A=B(xA)$ modulo $2$ gives $A=(1-xA)^{-1}$, or
$A=1+xA^2$. The unique solution is
\[
 A(x)=\sum_{j\ge0}x^{2^j-1}\pmod2,
\]
as follows immediately by squaring. The other mod-$3$ assertions and
the asymptotic observation in the entry are not selected.
\end{proof}

\subsection{A396804: a Hurwitz-series congruence}
An integer Hurwitz series means $F(x)=\sum f_nx^n/n!$ with $f_n\in\mathbb Z$.
Composition of zero-constant such series, and the exponential of one,
are given by integer Bell polynomials. Congruence therefore respects
these operations; no reduction of $1/n!$ in a finite field is intended.

\begin{lemma}\label{round5:r5:hurwitz}
Let $F(x)=xe^x$. Then $F^{\circ2}\equiv x\pmod2$ and
$F^{\circ4}\equiv x\pmod4$, coefficientwise in the Hurwitz sense.
\end{lemma}
\begin{proof}
Lagrange inversion gives $n![x^n]F^{-1}(x)=(-n)^{n-1}$.
For every $n\ge1$ this is congruent to $n$, the coefficient of $F$,
modulo $2$. Thus $F^{-1}\equiv F$ and $F^{\circ2}\equiv x$.
Write $F^{\circ2}=x+2G$, with integral Hurwitz coefficients in $G$.
Since $G\circ(x+2G)\equiv G\pmod2$,
\[
 (x+2G)\circ(x+2G)=x+2G+2G\circ(x+2G)\equiv x\pmod4.
\]
\end{proof}

\begin{theorem}\label{round5:r5entry:A396804}
For the exponential series $A=x\exp(A^{\circ4})$ defining
\oeis{A396804}, one has $a(n)\equiv n\pmod4$ for every $n\ge1$.
This proves both its mod-$4$ and its parity conjectures as one package.
\end{theorem}
\begin{proof}
The equation determines its Hurwitz coefficients by induction: its
degree-$n$ coefficient depends only on those of $A$ of degree below $n$,
and all defining coefficient polynomials are integral. Hence existence,
integrality, and uniqueness hold modulo $4$ as well as over $\mathbb Z$.
By the lemma, $F=xe^x$ satisfies
$x\exp(F^{\circ4})\equiv xe^x=F\pmod4$. Uniqueness gives $A\equiv F$.
The source's obsolete mod-$3$ claim already has a published counterexample
in the entry; it is expressly excluded. Neither the replacement mod-$3$
claim nor the mod-$5$ claim is being asserted here.
\end{proof}

\subsection{A383165 and A383166: eventual periodicity}
\begin{lemma}[Finite differential algebra]\label{round5:r5:finite-ode}
Suppose $y_i'(x)=P_i(y_1(x),\ldots,y_d(x))$, where each $P_i$ is an
integer polynomial and each $y_i(0)$ is an integer. For any integer
polynomial $Q$, the derivatives $Q(y(x))^{(n)}|_{x=0}$ are integers
and are eventually periodic modulo every positive integer $m$.
\end{lemma}
\begin{proof}
Translate the initial values to zero, and set
$D=\sum_i P_i\partial/\partial y_i$ in the translated coordinates.
The $n$-th derivative is the constant term of $D^nQ$, proving integrality.
For $m\ge2$, consider the finite ring
\[
 R_m=(\mathbb Z/m\mathbb Z)[y_1,\ldots,y_d]/(y_1^m,\ldots,y_d^m).
\]
The derivation descends to this ring since $D(y_i^m)=m y_i^{m-1}P_i=0$.
Evaluation at zero also descends. The orbit $Q,DQ,D^2Q,\ldots$ in
this finite ring eventually repeats, and iteration of the same map
then repeats forever. Taking its constant term proves the assertion.
For $m=1$ the assertion is immediate. This argument gives existence of
a period; it does not claim a bound involving Euler's totient.
\end{proof}

Set
\[
 L(x)=\log\frac{e^{2x}+1}{2},\quad U(x)=e^{-L(x)},\quad
 H_2(x)=L(x)^2/2,\quad H_3(x)=L(x)^3/6.
\]
These functions satisfy the integer polynomial system
\[
 L'=2-U,\qquad U'=U^2-2U,\qquad
 H_2'=L(2-U),\qquad H_3'=H_2(2-U),
\]
with initial values $(L,U,H_2,H_3)=(0,1,0,0)$.

\begin{theorem}[A383165]\label{round5:r5entry:A383165}
The exponential coefficients of $H_2$ are eventually periodic modulo
every positive integer, as conjectured in \oeis{A383165}.
\end{theorem}
\begin{proof}
Apply Lemma~\ref{round5:r5:finite-ode} with $Q=H_2$ to the displayed system.
The polynomial system also supplies integrality at primes dividing $2$;
there is no division by $2$ in the finite ring.
\end{proof}

\begin{theorem}[A383166]\label{round5:r5entry:A383166}
The exponential coefficients of $H_3$ are eventually periodic modulo
every positive integer, as conjectured in \oeis{A383166}.
\end{theorem}
\begin{proof}
Apply the same lemma with $Q=H_3$. This handles primes $2$ and $3$
without inverting $6$. The source's sample modulus-$5$ period is a
supplementary regression in the code; the proof is valid for all moduli.
The linked Graves-method note supplies differentiation identities, not
a proof of these modular-periodicity claims.
\end{proof}

\subsection{A179929: a totient period with an exceptional prime}
\begin{theorem}\label{round5:r5entry:A179929}
For the exponential coefficients $a(n)$ of $3/(1+2e^{-3x})$,
the sequence modulo any $m\ge1$ is eventually periodic with period
dividing $\varphi(m)$, proving the selected conjecture in \oeis{A179929}.
\end{theorem}
\begin{proof}
First take a prime power $p^r$ with $p\ne3$. Put $t=e^{3x}-1$.
The established exponential generating function gives
\[
 \frac{3}{1+2e^{-3x}}=\frac{1+t}{1+t/3}
   =1+\sum_{k\ge1}2(-1)^{k-1}3^{-k}t^k.
\]
Consequently, for $n>0$,
\[
 a(n)=\sum_{k\ge1}2(-1)^{k-1}3^{-k}
       \sum_{j=0}^k(-1)^{k-j}\binom{k}{j}(3j)^n.
\]
This is interpreted via $k!S(n,k)$, so terms with $k>n$ vanish.
Modulo $p^r$, terms with $k\ge pr$ vanish as well, since $p^r\mid k!$.
The remaining finite sum has fixed coefficients and bases $3j$.
For $n\ge r$, nonunit bases vanish modulo $p^r$; unit bases have a
period dividing $\lambda(p^r)$ by Euler's theorem (or the exponent of
the unit group). Thus this part has that eventual period.

For $p=3$, use $u'=3u-u^2$, $u(0)=1$. Its derivatives are
$a(n)=(D^n u)(1)$ for $D=(3u-u^2)d/du$.
Modulo $3$, $D=-u^2d/du$ and
\[
 D^3u^j=-j(j+1)(j+2)u^{j+3}=0\pmod3.
\]
Therefore $D^3$ sends every integer polynomial into $3\mathbb Z[u]$,
and iteration gives $D^{3r}\mathbb Z[u]\subseteq3^r\mathbb Z[u]$.
So $a(n)=0\pmod{3^r}$ for $n\ge3r$.
Write $m=3^rM$ with $3\nmid M$. The Chinese remainder theorem gives
an eventual period $\lambda(M)$, with $\lambda(1)=1$, starting no later
than $\max(\{1,3r\}\cup\{v_p(M):p\mid M\})$.
Since $\lambda(M)\mid\varphi(M)\mid\varphi(m)$, the desired bound follows.
\end{proof}

\subsection{Classical theta evaluations used below}
Here the nome is $q=e^{-\pi}$, and
\[
 T=\frac{\pi^{1/4}}{\Gamma(3/4)},\quad
 \phi(z)=\sum_{n\in\mathbb Z}z^{n^2},\quad
 \psi(z)=\sum_{n\ge0}z^{n(n+1)/2},\quad
 \chi(z)=\prod_{n\ge0}(1+z^{2n+1}).
\]
All these products and series converge in $|z|<1$; the nonzero theta
denominators used below follow from their convergent product formulas.
We use the following classical special values, giving their derivation
to make the substitutions independently reviewable:
\begin{align}
 \phi(q)&=T,&\phi(-q)&=2^{-1/4}T,\label{round5:r5:theta-base}\\
 \psi(q)&=2^{-5/8}q^{-1/8}T,&
 \psi(q^2)&=2^{-5/4}q^{-1/4}T,\label{round5:r5:psi-base}\\
 \psi(q^4)&=\tfrac14q^{-1/2}T\sqrt{2-\sqrt2},&
 \chi(q)&=2^{1/4}q^{1/24}.\label{round5:r5:chi-base}
\end{align}

At lattice parameter $i$, the theta transformation gives
$\theta_2(q)=\theta_4(q)$, and the quartic identity gives
$\theta_2(q)=\theta_4(q)=2^{-1/4}\theta_3(q)$.
The elliptic-integral relation
$K(1/\sqrt2)=\pi\theta_3(q)^2/2$, together with
$K(1/\sqrt2)=\Gamma(1/4)^2/(4\sqrt\pi)$ and the gamma reflection
formula, gives $\theta_3(q)=T$.
These are standard identities; see NIST DLMF
\href{https://dlmf.nist.gov/20.7}{\S20.7},
\href{https://dlmf.nist.gov/20.9}{\S20.9}, and
\href{https://dlmf.nist.gov/19.20}{\S19.20}.

The elementary series relation
$\theta_2(z)=2z^{1/4}\psi(z^2)$ gives $\psi(q^2)$.
The duplication formulas
\[
 \theta_2(z)^2=2\theta_2(z^2)\theta_3(z^2),\qquad
 \theta_2(z^2)^2=\tfrac12(\theta_3(z)^2-\theta_4(z)^2)
\]
give $\psi(q)$ and $\psi(q^4)$, respectively.
Finally let $E(q)=\prod_{n\ge1}(1-q^n)$ and
$\eta(i)=q^{1/12}E(q^2)$. The product identities give
$\phi(q)=E(q^2)\chi(q)^2$ and
$2\eta(i)^3=\theta_2(q)\theta_3(q)\theta_4(q)$.
Hence $E(q^2)=q^{-1/12}T/\sqrt2$ and
$\chi(q)^2=\sqrt2 q^{1/12}$, proving the last value.
All roots are the positive real roots.

\paragraph{Source and priority scope.}
The six following statements were labelled empirical by Simon Plouffe
in September 2025. His
\href{https://arxiv.org/abs/2509.15609}{\emph{Numbers in the base $e^\pi$}}
describes the numerical search and explicitly discusses shifting the
initial index. The general theta identities above are classical; these
packages are exact, source-specific deductions, not claims of six new
special-function identities or of worldwide priority. Decimal-constant
entries cross-referenced by the source are not separately counted.

\begin{theorem}[A210063]\label{round5:r5entry:A210063}
For the coefficients of $\psi(x^4)/\phi(x)$,
\[
 \sum_{n\ge0}a(n)e^{-\pi n}
   =\frac14 e^{\pi/2}\sqrt{2-\sqrt2}.
\]
\end{theorem}
\begin{proof}
The source's established generating function is analytic at $q$.
By \eqref{round5:r5:theta-base} and \eqref{round5:r5:chi-base}, its value is
$(q^{-1/2}T\sqrt{2-\sqrt2}/4)/T$, exactly the asserted expression.
\end{proof}

\begin{theorem}[A243763: index correction]\label{round5:r5entry:A243763}
The sequence has offset $1$, and its generating function is
$\sum_{n\ge1}a(n)x^n=x\phi(x)^3\psi(x^2)^4$. Its exact evaluations are
\[
 \sum_{n\ge1}a(n)e^{-\pi n}=\frac{T^7}{32},\qquad
 \sum_{j\ge0}a(j+1)e^{-\pi j}=\frac{e^\pi T^7}{32}.
\]
Thus the empirical right-hand side is correct for the zero-based list
of terms; the displayed sum must use $a(j+1)$, not the original $a(j)$.
\end{theorem}
\begin{proof}
At $x=q$, substitute $\phi(q)=T$ and
$\psi(q^2)=2^{-5/4}q^{-1/4}T$ to get
$qT^3(2^{-5/4}q^{-1/4}T)^4=T^7/32$.
Shifting the exponent down by one multiplies this by $q^{-1}=e^\pi$.
This is classified as an index correction with a proof of the corrected
formula, not as a counterexample to the intended normalized value.
\end{proof}

\begin{theorem}[A107635]\label{round5:r5entry:A107635}
For the coefficients of $\chi(x)^3$,
\[
 \sum_{n\ge0}a(n)e^{-\pi n}=2^{3/4}e^{-\pi/8}.
\]
\end{theorem}
\begin{proof}
The generating function $\chi(x)^3$ is an established formula in the
entry. Cubing $\chi(q)=2^{1/4}q^{1/24}$ gives the result.
\end{proof}

\begin{theorem}[A002512]\label{round5:r5entry:A002512}
For the coefficients of $\chi(x)^{10}/\phi(x)^4$,
\[
 \sum_{n\ge0}a(n)e^{-\pi n}
   =\frac{4\sqrt2\,e^{-5\pi/12}\Gamma(3/4)^4}{\pi}.
\]
\end{theorem}
\begin{proof}
The defining generating function at $q$ is
$(2^{1/4}q^{1/24})^{10}/T^4=2^{5/2}q^{5/12}T^{-4}$.
Replacing $T$ and $q$ by their definitions gives precisely the formula.
\end{proof}

\begin{theorem}[A134416]\label{round5:r5entry:A134416}
For the source generating function $1/(\phi(x)\phi(-x)^4)$,
\[
 \sum_{n\ge0}a(n)e^{-\pi n}
   =\frac{2\Gamma(3/4)^5}{\pi^{5/4}}.
\]
\end{theorem}
\begin{proof}
Michael Somos's established generating-function identity in the entry
gives $1/(T(2^{-1/4}T)^4)=2T^{-5}$ at $q$, as required.
\end{proof}

\begin{theorem}[A030204]\label{round5:r5entry:A030204}
For the coefficients of $\psi(x)\phi(-x)$,
\[
 \sum_{n\ge0}a(n)e^{-\pi n}
   =\frac{2^{-7/8}e^{\pi/8}\sqrt\pi}{\Gamma(3/4)^2}.
\]
\end{theorem}
\begin{proof}
The source's established generating function evaluates to
$(2^{-5/8}q^{-1/8}T)(2^{-1/4}T)=2^{-7/8}q^{-1/8}T^2$.
This is the empirical expression, since $2^{-7/8}=2^{1/8}/2$.
\end{proof}


\Needspace{6\baselineskip}
\section{Arithmetic, averages, and exact counterexamples}
\label{round5:r5:arithmetic}

This section resolves twelve selected packages at twelve different OEIS
entries.  Nine packages prove conjectures, two prove explicitly empirical
observations, and one refutes a stated uniform bound.  The last
classification matters: the counterexample in \oeis{A051699} does not resolve
the separate conjecture about the asymptotic constant.  Known theorems used
below are credited explicitly.  The accompanying standard-library verifier
\texttt{r5\_arith\_verify.py} produces
\texttt{r5\_arith\_certificates.json}; finite checks supplement the universal
arguments and do not establish any of their limits by extrapolation.

\subsection{The floor-factorial constant in A010786}
\label{round5:r5entry:A010786}

For $n\geq1$, put $F(n)=\prod_{j=1}^n\lfloor n/j\rfloor$.
The October 3, 2018 comment in \oeis{A010786} conjectures that the constant
in its leading logarithmic asymptotic is \oeis{A085361}.

\begin{theorem}
With
\[
 C=\sum_{r\geq1}\frac{\zeta(r+1)-1}{r}
   =0.7885305659\ldots,
\]
one has, for every $n\geq1$,
\[
                 0\leq Cn-\log F(n)\leq\log n+1.
\]
In particular $\log F(n)\sim Cn$ and also
$\log F(n)\sim C(n-\tfrac12\log(2\pi n))$.
\end{theorem}
\begin{proof}
For any positive integer $h$,
$\log h=\sum_{m=2}^h\log(m/(m-1))$.  Interchanging two finite sums gives
\[
 \log F(n)=\sum_{m=2}^n
          \left\lfloor\frac nm\right\rfloor\log\frac m{m-1}.
\]
The convergent positive series
$C=\sum_{m\geq2}m^{-1}\log(m/(m-1))$ has the displayed zeta form because
$-\log(1-1/m)=\sum_{r\geq1}(r m^r)^{-1}$ and all terms are nonnegative.
Subtracting the preceding finite identity from $Cn$ yields
\[
 Cn-\log F(n)=
 \sum_{m=2}^n\left\{\frac nm\right\}\log\frac m{m-1}
       +n\sum_{m>n}\frac1m\log\frac m{m-1}.
\]
Both terms are nonnegative.  The first is at most $\log n$, by telescoping.
The inequality $\log(1+u)\leq u$ bounds the second by
$n\sum_{m>n}(m(m-1))^{-1}=1$.  The conclusions follow.
\end{proof}

\paragraph{Scope and audit.}
The constant and its alternative series are already known; the selected
claim identifies this constant in the floor-factorial asymptotic.
The linked Alladi--Grinstead discussion was inspected and supplies the known
constant, rather than a new claim of its discovery.  Here $\sim$ means ratio
asymptotics.  The theorem does not assert a second-order expansion with
coefficient $-C/2$ and bounded remainder.

\subsection{The Hultman Dirichlet-series family in A034444}
\label{round5:r5entry:A034444}

The selected target is Werner Schulte's February 17, 2022 conjecture, not
the 2018 unitary-divisor divisibility conjecture, which the current entry
already marks proved.  Write $H(i,k)=\oeis{A164652}(i,k)$ for the Hultman
numbers, and define the multiplicative function $b_i$ by $b_i(1)=1$ and
\[
 b_i(p^e)=\frac{i+2}{i!}\sum_{k=1}^{i+1}H(i,k)e^{k-1}
       \qquad (i\geq0,\ e\geq1).
\]

\begin{theorem}
For every integer $i\geq0$ and every complex $s$ with $\Re(s)>1$,
\[
                   \sum_{n\geq1}\frac{b_i(n)}{n^s}
                   =\frac{\zeta(s)^{i+2}}{\zeta((i+2)s)}.
\]
\end{theorem}
\begin{proof}
We use the established Hultman formula
\[
 H(i,k)=\begin{cases}
 c(i+2,k)/\binom{i+2}{2},&i-k\text{ odd},\\
 0,&i-k\text{ even},
 \end{cases}
\]
where $c$ denotes unsigned Stirling numbers of the first kind.
This is the known input recorded in \oeis{A164652}, and appears through
the cycle polynomial in Section 3 of B\'ona--Flynn,
\href{https://arxiv.org/pdf/0811.0740}{arXiv:0811.0740}.
Let $P_i(x)=\prod_{j=1}^{i+1}(x+j)$.  Separating the two parities in
$xP_i(x)=\sum_k c(i+2,k)x^k$ gives
\[
 \sum_kH(i,k)x^k
 =\frac{x\{P_i(x)+(-1)^iP_i(-x)\}}{(i+1)(i+2)}.
\]
Consequently, for $e\geq1$,
\[
 b_i(p^e)=\frac{\prod_{j=1}^{i+1}(e+j)-
                         \prod_{j=1}^{i+1}(e-j)}{(i+1)!}
 =\binom{e+i+1}{i+1}-\binom{e-1}{i+1},
\]
where the second binomial is zero when $1\leq e\leq i+1$.
With $m=i+2$, the local power series therefore is
\[
                 1+\sum_{e\geq1}b_i(p^e)t^e
                       =\frac{1-t^m}{(1-t)^m}.
\]
Euler multiplication proves the assertion.  Absolute convergence in
$\Re(s)>1$ follows, for example, by bounding each nonnegative local
coefficient by $\binom{e+m-1}{m-1}$, the coefficient of $(1-t)^{-m}$.
\end{proof}

The known Hultman formula is not counted again.  The special cases
\oeis{A226602} and \oeis{A286779}, and the parameter choices $i$, do not add
to the one-entry count.

\subsection{An exact counterexample to the uniform bound in A051699}
\label{round5:r5entry:A051699}

Let $a(k)$ be the distance from $k$ to its nearest prime and
$S(N)=\sum_{k=1}^N a(k)$.  The selected September 18, 2005 assertion in
\oeis{A051699} includes $S(N)<\tfrac38N\log N$ for every $N>130$.

\begin{theorem}
This upper bound is false.  At $N=2^{20}=1048576$,
\[
 S(N)=5542944>5505024>\frac38N\log N.
\]
\end{theorem}
\begin{proof}
The exact certificate uses the sieve of Eratosthenes through $N+100$.
It gives $\pi(N)=82025$, last prime $1048573$, and next prime $1048583$.
For consecutive primes $p<q$, the total nearest-prime distance in the
completed interval from $p$ to $q$ is $\lfloor(q-p)^2/4\rfloor$.
Summing this expression over all completed prime gaps up to $1048573$
gives $5542937$.  The three remaining distances are $1,2,3$; the distance
at $1$ is $1$.  Thus $S(N)=5542937+6+1=5542944$.
The certificate includes the entire gap histogram and the verifier also
recomputes the sum directly, one integer at a time.

The logarithm comparison is certified without floating point:
\[
 e^{7/10}>1+\frac7{10}+\frac{(7/10)^2}{2}
                 +\frac{(7/10)^3}{6}
            =\frac{12013}{6000}>2.
\]
Hence $\log2<7/10$, and
$\frac38N\log N<\frac38N\cdot20\cdot\frac7{10}=5505024$.
The difference between the exact sum and this larger rational upper bound
is $37920>0$.
\end{proof}

\paragraph{Scope and audit.}
This is a complete refutation of the explicit uniform bound.  It does not
prove or refute the separate proposed asymptotic $S(N)\sim C N\log N$,
nor determine its constant.  The current entry contains later numerical
estimates inconsistent with the old bound, but no exact counterexample
to that bound is supplied there.

\subsection{Sixth powers in the cubic gcd triangle A129455}
\label{round5:r5entry:A129455}

Put $G(0)=1$ and
$G(n)=\prod_{1\leq r,s,t\leq n}\gcd(r,s,t)$ for $n\geq1$.
Thus $G(n)=\oeis{A129454}(n+1)$.  The triangle in \oeis{A129455} is
\[
                   T(n,k)=\frac{G(n)}{G(k)G(n-k)},\qquad 0\leq k\leq n.
\]
Peter Bala's April 16, 2007 observations propose integrality, divisibility
by $\binom nk$, and a perfect-sixth-power quotient.

\begin{theorem}
For all integers $n\geq0$ and $0\leq k\leq n$, there is a positive integer
$Y(n,k)$ such that $T(n,k)=\binom nkY(n,k)^6$.
\end{theorem}
\begin{proof}
For every prime $p$, threshold counting gives
\[
 v_p(G(n))=\sum_{r\geq1}\left\lfloor\frac n{p^r}\right\rfloor^3.
\]
Indeed $v_p(\gcd(u,v,w))=\sum_{r\geq1}
\mathbf1_{p^r\mid u}\mathbf1_{p^r\mid v}\mathbf1_{p^r\mid w}$.
This valuation identity is already a known result: Ma's
\href{https://arxiv.org/html/2509.13354v1}{2025 generalized-factorials paper},
Section 2.2, includes the general power-$q$ version.  It is used here as
an input, and is not counted as a new solution for \oeis{A129454}.

For a fixed $p^r$, set
$u=\lfloor k/p^r\rfloor$, $v=\lfloor(n-k)/p^r\rfloor$, and
$\epsilon=\lfloor n/p^r\rfloor-u-v\in\{0,1\}$.
Legendre's formula shows that the contribution to
$v_p(T(n,k)/\binom nk)$ is
\[
 (u+v+\epsilon)^3-u^3-v^3-\epsilon
 =\begin{cases}
 3uv(u+v),&\epsilon=0,\\
 3(u+v)(u+1)(v+1),&\epsilon=1.
 \end{cases}
\]
Both expressions are nonnegative multiples of six.  In the first case,
$uv(u+v)$ is even; in the second, either $u+v$ is even or one of $u+1,v+1$
is even.  Every prime valuation of the positive rational quotient is
therefore a nonnegative multiple of six.  Unique factorization proves
the assertion, including the empty-product boundary cases.
\end{proof}

Only the sixth-power package is selected.  The incorrectly shifted
$\lfloor n/p^r\rfloor^3$ statement attached to the $n-1$ convention in
\oeis{A129454} is not used to inflate the count.

\subsection{The empirical mean in A161789}
\label{round5:r5entry:A161789}

For $n\geq1$, let $a(n)=\max\{j\geq1:2^j-1\mid n\}$.
The selected target is Amiram Eldar's June 30, 2025 empirical asymptotic
mean in \oeis{A161789}.  Its separate gcd conjecture was proved by
Robert Israel on August 30, 2026; that one-page proof was inspected and
is not claimed again here.

\begin{theorem}
The mean $\lim_{N\to\infty}N^{-1}\sum_{n=1}^Na(n)=C$ exists, and
\[
                       1.9249135089<C<1.9249135166.
\]
More generally it has the convergent rational approximation scheme below,
with an explicit error bound for every truncation.
\end{theorem}
\begin{proof}
Put $M_j=2^j-1$ and
$a_L(n)=\max\{1\leq j\leq L:M_j\mid n\}$.
This function is periodic with period
$P_L=\operatorname{lcm}(M_1,\ldots,M_L)$, so it has a rational mean $C_L$.
Pointwise,
\[
 0\leq a(n)-a_L(n)
       \leq\sum_{j>L}(j-1)\mathbf1_{M_j\mid n}.
\]
Averaging over any finite initial interval gives the uniform bound
\[
 0\leq\frac1N\sum_{n\leq N}(a(n)-a_L(n))
 \leq\sum_{j>L}\frac{j-1}{2^j-1}
 \leq E_L:=\frac{2(L+1)}{2^{L+1}-1}.
\]
For the last inequality, use
$(2^j-1)^{-1}\leq 2^{-j}/(1-2^{-L-1})$ and
$\sum_{j>L}(j-1)2^{-j}=(L+1)2^{-L}$.
The $C_L$ are increasing and bounded; their limit $C$ satisfies
$0\leq C-C_L\leq E_L$.  Taking the lower and upper limits of the finite
means, then letting $L\to\infty$, proves their convergence to $C$.

For reproducible exact evaluation, represent
$a_L(n)=\sum_{d\mid n}c_{L,d}$, with finite support.
Start with $c_{1,1}=1$.  The identity
\[
 a_j(n)=a_{j-1}(n)+
                 (j-a_{j-1}(n))\mathbf1_{M_j\mid n}
\]
gives an integer update: retain the old coefficients, add $j$ at $M_j$,
and subtract each old $c_{j-1,d}$ at $\operatorname{lcm}(d,M_j)$.
Then
$C_L=\sum_d c_{L,d}/d=P_L^{-1}\sum_d c_{L,d}(P_L/d)$.
The exact $L=32$ certificate has $350208$ nonzero coefficients and gives
\[
\begin{split}
 C_{32}={}&
 \frac{207177202723641066584134331391358901063030104744958452260684812233899814}
      {107629356729483771508846897326023675486658403727012548610030251254582045},\\
 E_{32}={}&\frac{66}{8589934591}.
\end{split}
\]
Exact rational comparisons yield
$1.9249135089<C_{32}\leq C\leq C_{32}+E_{32}<1.9249135166$.
The displayed enclosure therefore certifies the initially observed digits;
it is not a large-$N$ numerical fit.
\end{proof}

\subsection{Prime distance determinants in A187011}
\label{round5:r5entry:A187011}

Let $p_1<p_2<\cdots$ be the primes and
$D_n=\det(|p_i-p_j|)_{1\leq i,j\leq n}$, with $D_0=1$.
Both parts of Zhi-Wei Sun's August 30, 2013 conjecture in \oeis{A187011}
are selected.

\begin{theorem}
For every $n\geq2$,
\[
 D_n=(-1)^{n-1}2^{n-2}(p_n-2)
                  \prod_{j=1}^{n-1}(p_{j+1}-p_j).
\]
In particular $(-1)^{n-1}D_n>0$ and $|D_n|^{1/n}\to\infty$.
\end{theorem}
\begin{proof}
First let $x_1<\cdots<x_n$ be arbitrary real numbers and set
$d_j=x_{j+1}-x_j>0$.  Apply successive-difference operations to both rows
and columns of $A=(|x_i-x_j|)$. Equivalently, apply the congruence
transformation $BAB^{\mathsf T}$, where $B$ is the lower bidiagonal
matrix with diagonal entries $1$ and subdiagonal entries $-1$,
whose determinant is $1$.  The resulting matrix is
\[
 \begin{pmatrix}0&d^{\mathsf T}\\d&-2\operatorname{diag}(d_1,\ldots,d_{n-1})\end{pmatrix}.
\]
The Schur complement gives its determinant as
\[
 (-2)^{n-1}\Bigl(\prod_jd_j\Bigr)
                           \frac12\sum_jd_j,
\]
which proves the formula and its sign.

For the growth claim we use the classical, unconditional prime number
theorem in each fixed arithmetic progression:
$\pi(X;q,a)/\pi(X)\to1/\varphi(q)$ when $(a,q)=1$.
For a primary modern reference stating and refining this theorem, see
Thorner--Zaman,
\href{https://arxiv.org/pdf/2108.10878}{arXiv:2108.10878}, Section 1.1.
No uniformity in growing moduli and no conjecture on prime tuples is used.

Fix $h\geq1$.  For any fixed product $Q$ of odd primes $\ell$ not dividing
$h$, a pair of primes $p,p+h$ larger than the prime factors of $Q$ must
lie in one of $\prod_{\ell\mid Q}(\ell-2)$ admissible classes modulo $Q$.
The fixed-modulus theorem therefore gives
\[
 \limsup_{X\to\infty}
 \frac{\#\{p\leq X:p,p+h\text{ prime}\}}{\pi(X)}
 \leq \prod_{\ell\mid Q}\frac{\ell-2}{\ell-1}.
\]
Let $Q$ include successively all the eligible primes.  The product tends
to zero, since the sum of prime reciprocals diverges and only finitely
many primes were omitted.  Thus each fixed difference $h$ occurs in
$o(\pi(X))$ prime pairs up to $X$.

For every fixed integer $K\geq1$, it follows that among the first $n-1$
consecutive prime gaps only $o(n)$ have size at most $K$.
All gaps are at least $1$, so
\[
 \prod_{j=1}^{n-1}(p_{j+1}-p_j)
                          \geq(K+1)^{n-1-o(n)}.
\]
The determinant formula implies
$\liminf_{n\to\infty}|D_n|^{1/n}\geq2(K+1)$.
Since $K$ is arbitrary, the desired divergence follows.
\end{proof}

The general distance-matrix calculation and the classical prime-distribution
input are not presented as new theorems.  Their application settles the two
specific statements still labelled conjectural in this entry.

\subsection{The sharp decimal-average bound in A276646}
\label{round5:r5entry:A276646}

Write $d(k)$ for the number of decimal digits in $k$ and
$S(n)=\sum_{k=1}^n k/10^{d(k)}$.  Thus the notation $0.k$ in
\oeis{A276646} means $k/10^{d(k)}$.

\begin{theorem}
For all $n\geq1$, $S(n)/n<11/20=0.55$, and
$\sup_{n\geq1}S(n)/n=11/20$.
\end{theorem}
\begin{proof}
The sum of the $d$-digit block is
\[
 \sum_{k=10^{d-1}}^{10^d-1}\frac{k}{10^d}
                     =\frac{99\cdot10^{d-1}-9}{20}.
\]
Summing the blocks gives the exact endpoint formula
\[
 S(10^d-1)=\frac{11(10^d-1)-9d}{20}.
\]
Let $A=10^{d-1}-1$ and $B=10^d-1$.
Within the block, $20S(n)$ is a quadratic polynomial in $n$ with positive
leading coefficient, because
\[
 S(n)=S(A)+\frac{n(n+1)-A(A+1)}{2\cdot10^d}.
\]
Hence $D(n)=11n-20S(n)$ is concave on the real interval $[A,B]$.
Its endpoint values are $9(d-1)$ and $9d$, so it is nonnegative
throughout.  If $d\geq2$ it is strictly positive.  For $d=1$, the left
endpoint is $n=0$ and strict concavity (or the formula) gives positivity
at each integer $1\leq n\leq9$ as well.
Finally the endpoint ratios tend to $11/20$, proving sharpness.
\end{proof}

This proves the explicit bound, and does not count the numerator and
denominator companion entries separately.

\subsection{Primitive quadruples in A278083}
\label{round5:r5entry:A278083}

For odd $m\geq1$, let $6b(m)$ count primitive ordered integral quadruples
$(x_1,x_2,x_3,x_4)$ with
\[
             \sum x_i=2m,\qquad\sum x_i^2=2m^2.
\]
Set $b(m)=0$ for even $m$.  Then $a(n)=b(2n-1)$ is \oeis{A278083}.

\begin{theorem}
The function $b$ is multiplicative, with $b(1)=1$, $b(2^e)=0$ for $e\geq1$,
and, for every odd prime $p$,
\[
 b(p^e)=p^{e-1}(p-\chi(p)),\qquad
 \chi(p)=(-1)^{(p-1)/2}.
\]
For odd $m\geq1$ the additional October 14, 2025 limit conjecture holds:
\[
 b(m)=\frac{\varphi(m)}6
                   \lim_{k\to\infty}\frac{r_3(m^{2k})}{m^k},
\]
where $r_3(t)$ counts all ordered integral representations by three squares.
\end{theorem}
\begin{proof}
Use the change of variables
\[
 a=x_1+x_2-m,\quad b=x_1+x_3-m,\quad c=x_1+x_4-m.
\]
The two constraints are equivalent to $a^2+b^2+c^2=m^2$, with inverse
\[
 (x_1,x_2,x_3,x_4)
 =\tfrac12(m+a+b+c,\ m+a-b-c,\ m-a+b-c,\ m-a-b+c).
\]
For odd $m$, the three-square equation forces $a+b+c$ odd, so all four
inverse coordinates are integers.  The transformation is a bijection.
The four coordinates cannot all be even, since their square sum is
$2m^2\equiv2\pmod4$.  For an odd prime $p$, divisibility of all four
coordinates implies $p\mid m,a,b,c$; conversely these divisibilities imply
$p\mid x_i$ for every $i$.  Also any prime dividing $a,b,c$ divides $m$,
by the three-square equation.  Thus the bijection preserves primitivity.

The required known input is Gauss's primitive three-square count
\[
 R_3(m^2)=6\prod_{p^e\parallel m}p^{e-1}(p-\chi(p))
                    \qquad(m\text{ odd}).
\]
For a primary research source giving precisely this formula, see
H\"urlimann, \emph{Exact and Asymptotic Evaluation of the Number of Distinct
Primitive Cuboids}, J. Integer Sequences 18 (2015), Article 15.2.5,
\href{https://emis.de/ft/3840}{equation (2)}; it credits Gauss and
Cooper--Hirschhorn.  The bijection proves $6b(m)=R_3(m^2)$ and therefore
the first assertion, including multiplicativity after adjoining the
zero value on even integers.

The same paper's Theorem 2 records the known total-count formula
\[
 r_3(m^2)=6\prod_{p^e\parallel m}
 \left(\sum_{j=0}^{e}p^j-\chi(p)\sum_{j=0}^{e-1}p^j\right).
\]
Apply it to $m^k$, divide by $m^k$, and let $k\to\infty$.
Each factor tends to $(p-\chi(p))/(p-1)$, so multiplication by
$\varphi(m)/6$ gives $m\prod_{p\mid m}(1-\chi(p)/p)=b(m)$.
For $m=1$ the same assertion follows from $r_3(1)=6$.
\end{proof}

The classical representation formulas are credited inputs.  The selected
OEIS resolution is the explicit bijection and application to both stated
conjectures, counted together as one package.

\subsection{All admissible divisor-sum moduli in A280098}
\label{round5:r5entry:A280098}

The selected conjecture is the classification, proposed on December 25,
2016, of positive $k$ for which $\sigma(kn-1)/k$ is always integral.
The identical assertion at \oeis{A295012} is an uncounted duplicate.
Since $\sigma$ here is defined for positive arguments, the domain is
$n\geq1$ with $kn-1\geq1$; thus $k=1$ starts at $n=2$.

\begin{theorem}
The admissible moduli are exactly $1,3,4,6,8,12,24$.
\end{theorem}
\begin{proof}
Suppose $k>1$ is admissible.  Choose any unit $r$ modulo $k$.
Dirichlet's theorem provides distinct primes
$p\equiv r\pmod k$ and $q\equiv-r^{-1}\pmod k$.
Their product is $kn-1$ for a positive integer $n$.  Admissibility gives
\[
 0\equiv\sigma(pq)=(p+1)(q+1)
       \equiv(r+1)(1-r^{-1})=r-r^{-1}\pmod k.
\]
Thus every unit modulo $k$ satisfies $r^2\equiv1\pmod k$.
Dirichlet's theorem is a classical input; for a formalized primary account
see Carneiro, \href{https://arxiv.org/abs/1608.02029}{arXiv:1608.02029}.

If an odd prime $\ell$ divides $k$, CRT lifts each nonzero residue modulo
$\ell$ to a unit modulo $k$.  The polynomial $X^2-1$ has at most two roots
over $\mathbb F_\ell$, forcing $\ell=3$.
The factor $9$ is impossible because the unit $4$ is not a square root of
$1$ modulo $9$; the factor $16$ is impossible because $3^2\not\equiv1$
modulo $16$.  CRT again lifts these local obstructions.
Therefore $k\mid24$.  The candidate $k=2$ fails at $n=1$, since
$\sigma(1)=1$.

Conversely, let $k\in\{3,4,6,8,12,24\}$ and $m=kn-1>0$.
Every unit modulo such $k$ has square $1$.  A divisor $d$ of $m$ is a unit,
and its complementary divisor satisfies
$m/d\equiv-d^{-1}\equiv-d\pmod k$.  Their sum is divisible by $k$.
There is no unpaired divisor: $m$ cannot be a square because $m\equiv-1$
modulo $3$ or modulo $4$, and at least one of these moduli divides $k$.
Pairing all divisors proves $k\mid\sigma(m)$.  The case $k=1$ is immediate.
\end{proof}

\subsection{The prime support restriction in A300984}
\label{round5:r5entry:A300984}

An integer $N$ belongs to \oeis{A300984} when the sum $U(N)$ of its
squarefree divisors and the sum $V(N)$ of its nonsquarefree divisors
are both positive squarefree integers.

\begin{theorem}
Every such $N$ has the conjectured form $2^i p^j$ for an odd prime $p$.
In fact $i,j\geq2$ and $p\equiv1\pmod{12}$ are necessary.
\end{theorem}
\begin{proof}
Unique factorization gives
$U(N)=\prod_{p\mid N}(1+p)$.
Two distinct odd primes dividing $N$ would make $4\mid U(N)$,
contradicting squarefreeness.  Thus $N$ has at most one odd prime factor,
which proves the proposed form (allowing initially zero exponents).

For the stronger necessary conditions, a pure power of $2$ has either
$V=0$ or $4\mid V$, and an odd prime power has either $V=0$ or $p^2\mid V$.
The value $N=1$ also has $V=0$.  Hence both primes occur.
If $N=2p^j$, then $V(N)=3p^2\sigma(p^{j-2})$ for $j\geq2$, and $V=0$
for $j=1$.  If $N=2^i p$ with $i\geq2$, then
$V(N)=(2^{i+1}-4)(p+1)$ is divisible by $4$.
This proves $i,j\geq2$.
Now $U(N)=3(p+1)$ squarefree implies $4\nmid p+1$ and $3\nmid p+1$.
It also excludes $p=3$.  Thus $p\equiv1\pmod4$ and $p\equiv1\pmod3$,
giving $p\equiv1\pmod{12}$.
\end{proof}

These are necessary conditions.  No converse or density claim about the
allowed primes is asserted; the entry's separate majority-of-primes
observation is not counted as resolved.

\subsection{Remainder pairing in A339387}
\label{round5:r5entry:A339387}

For $n\geq1$ write
\[
 A(n)=\sum_{k=1}^n\left(\frac{\operatorname{lcm}(n,k)}{\gcd(n,k)}\bmod k\right),
 \qquad
 B(n)=\sum_{k=1}^n\left(\frac{\operatorname{lcm}(n,k)}{\gcd(n,k)}\bmod n\right).
\]
These are respectively \oeis{A339387} and \oeis{A339384}.

\begin{theorem}
$B(n)\geq A(n)$ always.  If $n$ is not squarefree then $B(n)>A(n)$.
Consequently the conjecture that $A(n)=B(n)$ implies $n$ squarefree is true.
\end{theorem}
\begin{proof}
The case $n=1$ is immediate.  For $n>1$, group the terms by
$d=\gcd(n,k)$, and put $g=n/d$, $k=dt$.  For $g>1$ the variable $t$
runs over $1\leq t<g$ coprime to $g$, and
$\operatorname{lcm}(n,k)/\gcd(n,k)=gt$.
The respective remainders are
\[
                     t(g\bmod d),\qquad g(t\bmod d).
\]
Set $r=g\bmod d$ and let
$N_d=\#\{1\leq t<g:(t,g)=1,\ t\bmod d>r\}$.
The involution $t\mapsto g-t$ preserves this set of totatives, and
\[
 (t\bmod d)+((g-t)\bmod d)
       =r+d\,\mathbf1_{t\bmod d>r}.
\]
Also $\sum_{(t,g)=1,\ 1\leq t<g}t=g\varphi(g)/2$.
It follows that the difference of the two grouped contributions is
\[
 g\sum(t\bmod d)-r\sum t=\frac{gd}{2}N_d=\frac n2N_d\geq0.
\]
The remaining group $g=1$ consists of $k=n$ and contributes $1$ to both
sums.  Therefore
$B(n)-A(n)=\frac n2\sum_{d\mid n,\ n/d>1}N_d$.

If $p^2\mid n$, choose $d=p$.  Then $g=n/p$ is divisible by $p$, so
$r=0$, while the totative $t=1$ has $t\bmod p=1>0$.
Thus $N_p\geq1$, making the difference strictly positive.
\end{proof}

The existing prime-power evaluations and divisibility equivalences in the
two entries were inspected; they are not counted again.  Only the displayed
implication and its strengthening by a nonnegative difference are selected.

\subsection{Every prime occurs infinitely often in A394454's divisor sums}
\label{round5:r5entry:A394454}

The entry \oeis{A394454} sums the distinct prime divisors of $\binom{2n}{n}$.
Its selected conjecture states that, for each prime $p$, divisibility by
$p$ occurs at infinitely many indices.

\begin{theorem}
For every prime $p$ and integer $r\geq1$,
\[
 v_p\binom{2(p^r-1)}{p^r-1}=r.
\]
Thus each prime contributes to the defining sum infinitely often, and
the sequence has infinite upper limit.
\end{theorem}
\begin{proof}
Put $n=p^r-1$.  Legendre's formula gives
\[
 v_p\binom{2n}{n}
   =\sum_{j\geq1}\left(\left\lfloor\frac{2n}{p^j}\right\rfloor
                  -2\left\lfloor\frac n{p^j}\right\rfloor\right).
\]
For $1\leq j\leq r$ the two floors are respectively
$2p^{r-j}-1$ and $p^{r-j}-1$, so each summand is $1$.
For $j>r$ both floors vanish, since $2(p^r-1)<p^{r+1}$.
This proves the exact valuation, including $p=2$.
The indices $p^r-1$ are distinct as $r$ grows.  Taking $r=1$ and
arbitrarily large primes gives $a(p-1)\geq p$, hence the final conclusion.
\end{proof}

The conclusion concerns occurrence as a prime divisor of the central
binomial coefficient, hence as a summand in $a(n)$.  It does not assert
that each prime divides the integer $a(n)$ itself.  This is an elementary
application of classical factorial valuations, not a new theorem about
the distribution of primes.


\Needspace{6\baselineskip}
\section{Thirty generating-function and recurrence packages}
\label{round5:sec:r5gf}
This section treats thirty distinct OEIS entries absent from the 408
previously counted entries. The count is 22 direct proofs, six boundary
qualifications, one repair of a malformed printed factor, and one
refutation with a proved replacement. Each entry counts once. The
six boundary cases have no starting index specified in their source;
we supply that index and retain all earlier admissible exceptions.
The separate A129366 refutation is stronger: its original recurrence
cannot become valid after any finite index. The source wording of
A279013 contains an undefined token, which is repaired explicitly.

All selected statements were labeled \emph{Conjecture} on their consulted
primary OEIS entries. Known generating functions, combinatorial
enumerations, closed forms and other recurrences are credited inputs.
This is a selected-statement proof collection, with no assertion that
the underlying generating functions or proof methods are new.
The limited prior-work audit and the historical correction at the end
of this section prevent source labels from being treated as guarantees
of worldwide novelty. No external submission or acceptance is claimed.

\subsection{A finite exact certificate proves an infinite recurrence}
Let $F(x)=\sum_{n\ge0}a_nx^n$, filling indices before the source offset
with zero solely for formal-series notation. Put $\theta=x\,d/dx$ and
\[
\mathcal L_p=\sum_{j=0}^{r}x^j p_j(\theta+j).
\]
Since $\theta x^k=kx^k$, coefficient extraction gives
\begin{equation}\label{round5:eq:r5gf-extract}
[x^n]\mathcal L_pF=\sum_{j=0}^{r}p_j(n)a_{n-j}.
\end{equation}
An exact identity $\mathcal L_pF=P(x)$ with polynomial $P$ therefore
proves the recurrence for every $n>\deg P$ for which the source indices
are defined. Coefficients of $P$ give every exceptional initial residual.

For the 29 algebraic cases, write $x=\xi(t)$ and
\[
F(\xi(t))=\sum_{i=0}^{s}R_i(t)\sqrt{D_i(t)},\qquad D_0=1.
\]
The case without a change of variable means $t=x$ and $\xi=t$.
All radicals below have constant term one at their indicated base point.
Define rational operators and coefficients by
\[
T_iR=\frac{\xi}{\xi'}\left(R'+\frac{D_i'}{2D_i}R\right),
\qquad c_m(t)=\sum_{j=0}^{r}\xi(t)^j[n^m]p_j(n+j).
\]
The chain rule proves $\theta(R\sqrt{D_i})=(T_iR)\sqrt{D_i}$.
Thus each supplied certificate consists of the exact rational identities
\begin{equation}\label{round5:eq:r5gf-cert}
\sum_m c_mT_0^mR_0=P(\xi(t)),\qquad
\sum_m c_mT_i^mR_i=0\quad(i\ge1).
\end{equation}
Every $p_j,R_i,D_i,\xi$ and the complete boundary polynomial $P$ is
printed below. Inserting these finite expressions and clearing
denominators gives zero polynomials over $\mathbb Q$; this is a
checkable algebraic certificate, not inference from a numerical prefix.
Equation \eqref{round5:eq:r5gf-extract} then completes each all-index proof.
The separate A034267 case uses a rational identity between successive
hypergeometric terms.

For a fully expanded small illustration, A104624 has
$F=\sqrt{1+4x^2}-2x$ and
\[
\mathcal L_p=(1+x)\bigl((1+4x^2)\theta-4x^2\bigr).
\]
Its action on $\sqrt{1+4x^2}$ is zero, because
$\theta\sqrt{1+4x^2}=4x^2/\sqrt{1+4x^2}$.
Its action on $-2x$ is $-2x(1+x)$.
This gives the complete boundary polynomial and proves the recurrence
at every $n\ge3$. The other algebraic entries apply precisely the same
finite rule, with possibly more derivatives and radicals.

\subsection{Indexing, seeds and parameter choices}
The Catalan and Motzkin generating functions are
\[
C(x)=\frac{1-\sqrt{1-4x}}{2x},\qquad
M(x)=\frac{1-x-\sqrt{1-2x-3x^2}}{2x^2}.
\]
For A162482, substituting $F=(1-x)^{-3}M(x/(1-x)^3)$ into
$M=1+xM+x^2M^2$ gives
$x^2F^2-((1-x)^3-x)F+1=0$, selecting the branch $F(0)=1$.
The A129366 definition gives
$F=(C(x)-xC(x^2))/(1-x)$ by summing its Catalan partial sums.
A134424 uses the existing statistic formula with
$g=1+xg+x^2g+x^2g^2$, $g(0)=1$; solving this quadratic gives
the displayed radical data. The other ordinary algebraic seeds are
the established generating functions identified in their individual
source notes. A025259 requires multiplication of the source's
offset-zero displayed GF by $x$, because its sequence begins at index 1.

For A243764, A270490, A270530, A279013 and A279014, set
\[
t=xC(x)=\frac{1-\sqrt{1-4x}}2,
\qquad x=\xi(t)=t(1-t),\qquad t(0)=0.
\]
Then $C(x)=1/(1-t)$, so the nested square roots reduce to the rational
radicands below. For example, A243764 becomes
$F=-(1-t)(1+\sqrt{1-8t^2})/2$.
The usual Catalan binomial kernel is
\[
\sum_{n\ge0}\binom{2n-1}{n-k}x^n
=\frac{1-t}{1-2t}\left(\frac{t}{1-t}\right)^k\quad(k\ge0),
\]
using $\binom{-1}{0}=1$ at the single initial case.
Substituting the Catalan and Fibonacci generating functions for
the $k$-weights gives the A279013 and A279014 representations.
For A270530 the analogous kernel is
\[
\sum_{n\ge0}\binom{2n+2}{n-k}x^n
=\frac1{(1-t)^2(1-2t)}\left(\frac{t}{1-t}\right)^k.
\]
Both identities follow by the binomial expansion after $x=t(1-t)$
or by formal coefficient extraction; thus the sum definitions select
the same analytic branches as the listed generating functions.

For A387084 let $g=1+xg^5$, $g(0)=1$, and $t=g-1$.
Then $\xi(t)=t/(1+t)^5$, with $\xi'(0)=1$.
The central-binomial series $B=\sum_n\binom{5n}{n}x^n$ is
$B=g/(5-4g)$, by Lagrange inversion. The defining expression
$B/\sqrt{1+4(B-1)/5}$ therefore equals $g/\sqrt{5-4g}$.
This proves the parameterized seed $F=(1+t)/\sqrt{1-4t}$.
For every parameter used here $\xi'(0)=1$, so the formal inverse
exists. Apparent poles in individual $R_i$ are interpreted in the
Laurent field and cancel in $F$.

\subsection{The individual statements and complete certificates}
Each displayed recurrence means $\sum_jp_j(n)a_{n-j}=0$ in the stated
range. In the one refuted case the original coefficients are denoted
$q_j$ and the proved replacement coefficients are denoted $p_j$.


\subsubsection{A025259: Sep 29 2012}
\label{round5:r5entry:A025259}

Source: \url{https://oeis.org/A025259}. Selected conjecture: R. J. Mathar, Sep 29 2012. The defining formula or established generating function is credited input from the primary entry; only the selected conjectural recurrence is treated.

\textbf{Boundary qualification.} The source does not give a starting index. $n=5$ gives residual $1$. The exact valid tail is $n\ge6$. No explicit smaller starting index is attributed to the source.

\begin{align*}
p_{0}(n)&=n\\
p_{1}(n)&=- 2 \left(2 n - 3\right)\\
p_{2}(n)&=6 \left(n - 3\right)\\
p_{3}(n)&=- 4 \left(2 n - 9\right)\\
p_{4}(n)&=n - 6.
\end{align*}

Use $t=x$ and $\xi=x$.

The exact representation is $F=\sum_iR_i\sqrt{D_i}$, with $D_0=1$ and

\begin{align*}
R_{0}&=- \frac{x^{2} - 2 x - 1}{2},\\
R_{1}&=- \frac{1}{2},\\
D_{1}&=x^{4} - 8 x^{3} + 6 x^{2} - 4 x + 1.
\end{align*}

Inserting these data in \eqref{round5:eq:r5gf-cert} gives exactly

\[\mathcal L_pF=P(x)=x \left(x^{4} + 9 x^{2} - 6 x + 2\right).\]

All radical residuals in \eqref{round5:eq:r5gf-cert} vanish as rational identities. Therefore \eqref{round5:eq:r5gf-extract} proves the displayed recurrence for every integer $n\ge6$. The source offset is 1.

\subsubsection{A025272: Nov 21 2014}
\label{round5:r5entry:A025272}

Source: \url{https://oeis.org/A025272}. Selected conjecture: R. J. Mathar, Nov 21 2014. Kotesovec already supplied the displayed generating function and a different order-four recurrence in 2015. Only the separately labeled Mathar order-five recurrence is selected.

\begin{align*}
p_{0}(n)&=n\\
p_{1}(n)&=n + 1\\
p_{2}(n)&=- 10 \left(2 n - 5\right)\\
p_{3}(n)&=2 \left(2 n - 9\right)\\
p_{4}(n)&=2 \left(12 n - 67\right)\\
p_{5}(n)&=20 \left(n - 7\right).
\end{align*}

Use $t=x$ and $\xi=x$.

The exact representation is $F=\sum_iR_i\sqrt{D_i}$, with $D_0=1$ and

\begin{align*}
R_{0}&=\frac{1}{2},\\
R_{1}&=- \frac{1}{2},\\
D_{1}&=4 x^{4} + 4 x^{3} - 4 x + 1.
\end{align*}

Inserting these data in \eqref{round5:eq:r5gf-cert} gives exactly

\[\mathcal L_pF=P(x)=- x \left(5 x + 1\right) \left(4 x^{3} + 3 x^{2} - 1\right).\]

All radical residuals in \eqref{round5:eq:r5gf-cert} vanish as rational identities. Therefore \eqref{round5:eq:r5gf-extract} proves the displayed recurrence for every integer $n\ge6$. The source offset is 1.

\subsubsection{A034267: Feb 10 2025}
\label{round5:r5entry:A034267}

Source: \url{https://oeis.org/A034267}. Selected conjecture: R. J. Mathar, Feb 10 2025. The diagonal formula follows directly from A034261, whose formula is credited to Michael Somos. The same diagonal formula was added by Amiram Eldar in 2025. These formulas are prior input.

\begin{align*}
p_{0}(n)&=- \left(n + 2\right) \left(11 n - 7\right)\\
p_{1}(n)&=2 \left(23 n^{2} + 44 n + 30\right)\\
p_{2}(n)&=- 4 \left(n + 5\right) \left(2 n - 3\right).
\end{align*}

\textbf{Proof.} For $n\ge1$ the known formula is
$a_n=C_n n(n^2+n+1)/(n+2)$, where $C_n$ is the Catalan number.
Cancellation of factorials gives, for $n\ge2$,

\[\frac{a_{n-1}}{a_n}=r(n)=\frac{\left(n - 1\right) \left(n + 2\right) \left(n^{2} - n + 1\right)}{2 n \left(2 n - 1\right) \left(n^{2} + n + 1\right)}.\]

The same quotient gives $a_0/a_1=r(1)=0$ directly.
The finite rational identity
\[
p_0(n)+p_1(n)r(n)+p_2(n)r(n)r(n-1)=0
\]
follows by inserting the displayed polynomials and clearing
denominators. Its denominators are nonzero for integer $n\ge2$.
Multiplication by $a_n$ proves the recurrence for every $n\ge2$.
The source values $a_0=0,a_1=1,a_2=7$ also verify the first case.


\subsubsection{A089324: Nov 24 2012}
\label{round5:r5entry:A089324}

Source: \url{https://oeis.org/A089324}. Selected conjecture: R. J. Mathar, Nov 24 2012. The defining formula or established generating function is credited input from the primary entry; only the selected conjectural recurrence is treated.

\begin{align*}
p_{0}(n)&=n + 1\\
p_{1}(n)&=- 3 \left(n + 1\right)\\
p_{2}(n)&=13 - 5 n\\
p_{3}(n)&=18 \left(n - 2\right)\\
p_{4}(n)&=7 - 5 n\\
p_{5}(n)&=- 3 \left(n - 5\right)\\
p_{6}(n)&=n - 5.
\end{align*}

Use $t=x$ and $\xi=x$.

The exact representation is $F=\sum_iR_i\sqrt{D_i}$, with $D_0=1$ and

\begin{align*}
R_{0}&=- \frac{\left(x - 1\right)^{2}}{2 x \left(x^{2} - 3 x + 1\right)},\\
R_{1}&=\frac{1}{2 x \left(x^{2} - 3 x + 1\right)},\\
D_{1}&=\left(x^{2} - 2 x - 1\right) \left(x^{2} + 2 x - 1\right).
\end{align*}

Inserting these data in \eqref{round5:eq:r5gf-cert} gives exactly

\[\mathcal L_pF=P(x)=- \left(x - 1\right) \left(x + 1\right) \left(x^{2} - 4 x + 1\right).\]

All radical residuals in \eqref{round5:eq:r5gf-cert} vanish as rational identities. Therefore \eqref{round5:eq:r5gf-extract} proves the displayed recurrence for every integer $n\ge6$. The source offset is 0.

\subsubsection{A097331: Dec 02 2012}
\label{round5:r5entry:A097331}

Source: \url{https://oeis.org/A097331}. Selected conjecture: R. J. Mathar, Dec 02 2012. Fung Lam already supplied the shorter step-two recurrence in 2014. We prove only the separately conjectured Mathar order-three recurrence from the defining generating function.

\begin{align*}
p_{0}(n)&=n + 1\\
p_{1}(n)&=n\\
p_{2}(n)&=- 4 \left(n - 2\right)\\
p_{3}(n)&=- 4 \left(n - 3\right).
\end{align*}

Use $t=x$ and $\xi=x$.

The exact representation is $F=\sum_iR_i\sqrt{D_i}$, with $D_0=1$ and

\begin{align*}
R_{0}&=\frac{2 x + 1}{2 x},\\
R_{1}&=- \frac{1}{2 x},\\
D_{1}&=- \left(2 x - 1\right) \left(2 x + 1\right).
\end{align*}

Inserting these data in \eqref{round5:eq:r5gf-cert} gives exactly

\[\mathcal L_pF=P(x)=\left(x + 1\right) \left(2 x + 1\right).\]

All radical residuals in \eqref{round5:eq:r5gf-cert} vanish as rational identities. Therefore \eqref{round5:eq:r5gf-extract} proves the displayed recurrence for every integer $n\ge3$. The source offset is 0.

\subsubsection{A102319: Feb 20 2015}
\label{round5:r5entry:A102319}

Source: \url{https://oeis.org/A102319}. Selected conjecture: R. J. Mathar, Feb 20 2015. Only the first of the two listed conjectured recurrences, of order six, is counted. The second is outside the selected claim.

\begin{align*}
p_{0}(n)&=n \left(n - 1\right)\\
p_{1}(n)&=- 4 \left(n - 1\right) \left(3 n - 4\right)\\
p_{2}(n)&=53 n^{2} - 221 n + 232\\
p_{3}(n)&=- 8 \left(13 n^{2} - 85 n + 134\right)\\
p_{4}(n)&=51 n^{2} - 563 n + 1308\\
p_{5}(n)&=4 \left(n - 4\right) \left(29 n - 93\right)\\
p_{6}(n)&=- 105 \left(n - 5\right) \left(n - 4\right).
\end{align*}

Use $t=x$ and $\xi=x$.

The exact representation is $F=\sum_iR_i\sqrt{D_i}$, with $D_0=1$ and

\begin{align*}
R_{0}&=0,\\
R_{1}&=\frac{1}{2 \left(x - 1\right) \left(5 x - 1\right)},\\
R_{2}&=- \frac{1}{2 \left(x + 1\right) \left(3 x - 1\right)},\\
D_{1}&=\left(x - 1\right) \left(5 x - 1\right),\\
D_{2}&=- \left(x + 1\right) \left(3 x - 1\right).
\end{align*}

Inserting these data in \eqref{round5:eq:r5gf-cert} gives exactly

\[\mathcal L_pF=P(x)=0.\]

All radical residuals in \eqref{round5:eq:r5gf-cert} vanish as rational identities. Therefore \eqref{round5:eq:r5gf-extract} proves the displayed recurrence for every integer $n\ge6$. The source offset is 0.

\subsubsection{A104624: Nov 15 2012}
\label{round5:r5entry:A104624}

Source: \url{https://oeis.org/A104624}. Selected conjecture: R. J. Mathar, Nov 15 2012. The defining formula or established generating function is credited input from the primary entry; only the selected conjectural recurrence is treated.

\begin{align*}
p_{0}(n)&=n\\
p_{1}(n)&=n - 1\\
p_{2}(n)&=4 \left(n - 3\right)\\
p_{3}(n)&=4 \left(n - 4\right).
\end{align*}

Use $t=x$ and $\xi=x$.

The exact representation is $F=\sum_iR_i\sqrt{D_i}$, with $D_0=1$ and

\begin{align*}
R_{0}&=- 2 x,\\
R_{1}&=1,\\
D_{1}&=4 x^{2} + 1.
\end{align*}

Inserting these data in \eqref{round5:eq:r5gf-cert} gives exactly

\[\mathcal L_pF=P(x)=- 2 x \left(x + 1\right).\]

All radical residuals in \eqref{round5:eq:r5gf-cert} vanish as rational identities. Therefore \eqref{round5:eq:r5gf-extract} proves the displayed recurrence for every integer $n\ge3$. The source offset is 0.

\subsubsection{A107587: Feb 20 2015}
\label{round5:r5entry:A107587}

Source: \url{https://oeis.org/A107587}. Selected conjecture: R. J. Mathar, Feb 20 2015. Eremin 2021, equation 2.10, proves the generating function used here. His OEIS check of the selected recurrence through 800 is finite verification, not used as an all-index proof. Only the first listed recurrence is selected.

\begin{align*}
p_{0}(n)&=n \left(n + 2\right)\\
p_{1}(n)&=- 5 n^{2} - n + 3\\
p_{2}(n)&=10 n^{2} - 16 n + 3\\
p_{3}(n)&=- 10 n^{2} + 34 n - 27\\
p_{4}(n)&=- \left(n - 3\right) \left(11 n - 5\right)\\
p_{5}(n)&=15 \left(n - 4\right) \left(n - 3\right).
\end{align*}

Use $t=x$ and $\xi=x$.

The exact representation is $F=\sum_iR_i\sqrt{D_i}$, with $D_0=1$ and

\begin{align*}
R_{0}&=0,\\
R_{1}&=\frac{1}{4 x^{2}},\\
R_{2}&=- \frac{1}{4 x^{2}},\\
D_{1}&=5 x^{2} - 2 x + 1,\\
D_{2}&=- \left(x + 1\right) \left(3 x - 1\right).
\end{align*}

Inserting these data in \eqref{round5:eq:r5gf-cert} gives exactly

\[\mathcal L_pF=P(x)=0.\]

All radical residuals in \eqref{round5:eq:r5gf-cert} vanish as rational identities. Therefore \eqref{round5:eq:r5gf-extract} proves the displayed recurrence for every integer $n\ge5$. The source offset is 0.

\subsubsection{A108600: Aug 24 2013}
\label{round5:r5entry:A108600}

Source: \url{https://oeis.org/A108600}. Selected conjecture: R. J. Mathar, Aug 24 2013. The freely braided permutation enumeration and generating function are prior work of Green--Losonczy and Mansour. The current entry includes the initial value at zero, added by Heinz in January 2025. Only the later Mathar recurrence is targeted.

\begin{align*}
p_{0}(n)&=1 - n\\
p_{1}(n)&=7 n - 10\\
p_{2}(n)&=- 2 \left(4 n - 1\right)\\
p_{3}(n)&=- 8 \left(2 n - 11\right)\\
p_{4}(n)&=n - 1\\
p_{5}(n)&=- 3 \left(n - 2\right)\\
p_{6}(n)&=- 2 \left(2 n - 11\right).
\end{align*}

Use $t=x$ and $\xi=x$.

The exact representation is $F=\sum_iR_i\sqrt{D_i}$, with $D_0=1$ and

\begin{align*}
R_{0}&=- \frac{x^{3} - 4 x^{2} - 7 x + 2}{2 \left(x^{4} + 4 x - 1\right)},\\
R_{1}&=- \frac{x \left(x + 1\right)^{2}}{2 \left(x^{4} + 4 x - 1\right)},\\
D_{1}&=1 - 4 x.
\end{align*}

Inserting these data in \eqref{round5:eq:r5gf-cert} gives exactly

\[\mathcal L_pF=P(x)=- x^{5} + 15 x^{4} + 28 x^{3} - 12 x^{2} - 3 x + 1.\]

All radical residuals in \eqref{round5:eq:r5gf-cert} vanish as rational identities. Therefore \eqref{round5:eq:r5gf-extract} proves the displayed recurrence for every integer $n\ge6$. The source offset is 0.

\subsubsection{A116388: Jun 22 2016}
\label{round5:r5entry:A116388}

Source: \url{https://oeis.org/A116388}. Selected conjecture: R. J. Mathar, Jun 22 2016. The defining formula or established generating function is credited input from the primary entry; only the selected conjectural recurrence is treated.

\begin{align*}
p_{0}(n)&=n\\
p_{1}(n)&=- 7 \left(n - 1\right)\\
p_{2}(n)&=2 \left(4 n - 9\right)\\
p_{3}(n)&=25 n - 58\\
p_{4}(n)&=65 - 18 n\\
p_{5}(n)&=199 - 52 n\\
p_{6}(n)&=135 - 31 n\\
p_{7}(n)&=- 6 \left(n - 5\right).
\end{align*}

Use $t=x$ and $\xi=x$.

The exact representation is $F=\sum_iR_i\sqrt{D_i}$, with $D_0=1$ and

\begin{align*}
R_{0}&=- \frac{1}{2 \left(x^{3} + 3 x^{2} + 2 x - 1\right)},\\
R_{1}&=- \frac{2 x^{2} + 3 x - 1}{2 \left(x + 1\right) \left(3 x - 1\right) \left(x^{3} + 3 x^{2} + 2 x - 1\right)},\\
D_{1}&=- \left(x + 1\right) \left(3 x - 1\right).
\end{align*}

Inserting these data in \eqref{round5:eq:r5gf-cert} gives exactly

\[\mathcal L_pF=P(x)=- x \left(3 x^{3} + 3 x^{2} + x - 1\right).\]

All radical residuals in \eqref{round5:eq:r5gf-cert} vanish as rational identities. Therefore \eqref{round5:eq:r5gf-extract} proves the displayed recurrence for every integer $n\ge7$. The source offset is 0.

\subsubsection{A129366: Feb 05 2015}
\label{round5:r5entry:A129366}

Source: \url{https://oeis.org/A129366}. Selected conjecture: R. J. Mathar, Feb 05 2015. The defining formula or established generating function is credited input from the primary entry; only the selected conjectural recurrence is treated.

\textbf{Original conjecture refuted.} Its original coefficients are:

\begin{align*}
q_{0}(n)&=n \left(n - 1\right) \left(12 n + 35\right)\\
q_{1}(n)&=\left(n - 1\right) \left(12 n^{2} - 701 n + 1236\right)\\
q_{2}(n)&=2 \left(6 n^{3} - 385 n^{2} + 2285 n - 3432\right)\\
q_{3}(n)&=- 4 \left(405 n^{3} - 5313 n^{2} + 19970 n - 23175\right)\\
q_{4}(n)&=8 \left(156 n^{3} - 1724 n^{2} + 5498 n - 5175\right)\\
q_{5}(n)&=16 \left(393 n^{3} - 4981 n^{2} + 20393 n - 26820\right)\\
q_{6}(n)&=- 32 \left(n - 5\right) \left(2 n - 9\right) \left(93 n - 268\right).
\end{align*}

The defining Catalan sum gives $(a_0,\ldots,a_6)=(1,1,3,7,21,61,193)$. Exact substitution gives

\[\sum_{j=0}^{6}q_j(6)a_{6-j}=-73968\ne0.\]

This cannot be repaired by choosing a later starting index. Applying $\mathcal L_q$ to the same generating function gives

\[
Q_1(x)=\frac{N(x)}{x(4x-1)^2},\qquad Q_2(x)\in\mathbb Q(x),
\]
\begin{align*}
N(x)&=11040x^6-3200x^5+26064x^4\\
&\quad-8152x^3-2320x^2+756x-23.
\end{align*}

Here $\mathcal L_qF=Q_1\sqrt{1-4x}+Q_2\sqrt{1-4x^2}$.
The square classes of $1-4x$ and $1-4x^2$ are independent in
$\mathbb Q(x)$: each factorization has simple roots, and their product
also has a simple root. The biquadratic field therefore has an
automorphism changing the sign of $\sqrt{1-4x}$ while fixing the other
radical. If this residual were rational, that automorphism would force
$2Q_1\sqrt{1-4x}=0$, contrary to $Q_1\ne0$.
It is consequently not a polynomial, so infinitely many recurrence
residuals are nonzero. The original formula is not eventually valid.

\textbf{Proved replacement.} Write $F=U+V$ with the two radical
terms specified below, and let $\ell_1=U'/U$, $\ell_2=V'/V$.
Both are rational functions. Defining
\[
A=-\frac{\ell_1'+\ell_1^2-\ell_2'-\ell_2^2}{\ell_1-\ell_2},
\qquad B=-\ell_1'-\ell_1^2-A\ell_1
\]
gives $U''+AU'+BU=V''+AV'+BV=0$ directly.
Clearing denominators, writing derivatives in terms of $\theta$ and
extracting coefficients yields the following replacement. Its exact
certificate is also verified independently by \eqref{round5:eq:r5gf-cert}.

\begin{align*}
p_{0}(n)&=n \left(n + 1\right)\\
p_{1}(n)&=- n \left(7 n - 5\right)\\
p_{2}(n)&=2 \left(n - 2\right) \left(7 n - 9\right)\\
p_{3}(n)&=12 \left(4 n - 9\right)\\
p_{4}(n)&=- 8 \left(7 n^{2} - 41 n + 63\right)\\
p_{5}(n)&=16 \left(7 n^{2} - 59 n + 126\right)\\
p_{6}(n)&=- 32 \left(n - 5\right) \left(2 n - 9\right).
\end{align*}

Use $t=x$ and $\xi=x$.

The exact representation is $F=\sum_iR_i\sqrt{D_i}$, with $D_0=1$ and

\begin{align*}
R_{0}&=0,\\
R_{1}&=\frac{1}{2 x \left(x - 1\right)},\\
R_{2}&=- \frac{1}{2 x \left(x - 1\right)},\\
D_{1}&=1 - 4 x,\\
D_{2}&=- \left(2 x - 1\right) \left(2 x + 1\right).
\end{align*}

Inserting these data in \eqref{round5:eq:r5gf-cert} gives exactly

\[\mathcal L_pF=P(x)=0.\]

All radical residuals in \eqref{round5:eq:r5gf-cert} vanish as rational identities. Therefore \eqref{round5:eq:r5gf-extract} proves the displayed recurrence for every integer $n\ge6$. The source offset is 0.

\subsubsection{A134424: Jul 24 2022}
\label{round5:r5entry:A134424}

Source: \url{https://oeis.org/A134424}. Selected conjecture: R. J. Mathar, Jul 24 2022. The defining formula or established generating function is credited input from the primary entry; only the selected conjectural recurrence is treated.

\begin{align*}
p_{0}(n)&=- \left(n + 2\right) \left(5 n - 7\right)\\
p_{1}(n)&=- \left(n + 1\right) \left(5 n - 127\right)\\
p_{2}(n)&=135 n^{2} - 655 n - 42\\
p_{3}(n)&=2 \left(5 n^{2} - 275 n - 108\right)\\
p_{4}(n)&=- 725 n^{2} + 4941 n - 5734\\
p_{5}(n)&=- 235 n^{2} + 1880 n - 1173\\
p_{6}(n)&=725 n^{2} - 6659 n + 12606\\
p_{7}(n)&=2 \left(5 n^{2} + 195 n - 1988\right)\\
p_{8}(n)&=- 135 n^{2} + 1505 n - 3358\\
p_{9}(n)&=- \left(n - 9\right) \left(5 n + 87\right)\\
p_{10}(n)&=\left(n - 10\right) \left(5 n - 33\right).
\end{align*}

Use $t=x$ and $\xi=x$.

The exact representation is $F=\sum_iR_i\sqrt{D_i}$, with $D_0=1$ and

\begin{align*}
R_{0}&=\frac{\left(x^{2} + 1\right) \left(x^{4} + 2 x^{3} - 3 x^{2} - 2 x + 1\right)}{2 x^{2} \left(x^{2} - x - 1\right) \left(x^{2} + 3 x - 1\right)},\\
R_{1}&=\frac{\left(x^{2} + 1\right) \left(x^{2} + x - 1\right)}{2 x^{2} \left(x^{2} - x - 1\right) \left(x^{2} + 3 x - 1\right)},\\
D_{1}&=\left(x^{2} - x - 1\right) \left(x^{2} + 3 x - 1\right).
\end{align*}

Inserting these data in \eqref{round5:eq:r5gf-cert} gives exactly

\[\mathcal L_pF=P(x)=12 x^{2} \left(x^{2} + 1\right) \left(x^{2} + 24 x - 1\right).\]

All radical residuals in \eqref{round5:eq:r5gf-cert} vanish as rational identities. Therefore \eqref{round5:eq:r5gf-extract} proves the displayed recurrence for every integer $n\ge10$. The source offset is 0.

\subsubsection{A162482: Feb 10 2015}
\label{round5:r5entry:A162482}

Source: \url{https://oeis.org/A162482}. Selected conjecture: R. J. Mathar, Feb 10 2015. The defining formula or established generating function is credited input from the primary entry; only the selected conjectural recurrence is treated.

\begin{align*}
p_{0}(n)&=n + 2\\
p_{1}(n)&=- 4 \left(2 n + 1\right)\\
p_{2}(n)&=18 \left(n - 1\right)\\
p_{3}(n)&=- 13 \left(2 n - 5\right)\\
p_{4}(n)&=17 \left(n - 4\right)\\
p_{5}(n)&=- 3 \left(2 n - 11\right)\\
p_{6}(n)&=n - 7.
\end{align*}

Use $t=x$ and $\xi=x$.

The exact representation is $F=\sum_iR_i\sqrt{D_i}$, with $D_0=1$ and

\begin{align*}
R_{0}&=- \frac{x^{3} - 3 x^{2} + 4 x - 1}{2 x^{2}},\\
R_{1}&=- \frac{1}{2 x^{2}},\\
D_{1}&=\left(x^{3} - 3 x^{2} + 2 x - 1\right) \left(x^{3} - 3 x^{2} + 6 x - 1\right).
\end{align*}

Inserting these data in \eqref{round5:eq:r5gf-cert} gives exactly

\[\mathcal L_pF=P(x)=2 \left(x - 1\right)^{2} \left(2 x + 1\right).\]

All radical residuals in \eqref{round5:eq:r5gf-cert} vanish as rational identities. Therefore \eqref{round5:eq:r5gf-extract} proves the displayed recurrence for every integer $n\ge6$. The source offset is 0.

\subsubsection{A164586: Feb 10 2015}
\label{round5:r5entry:A164586}

Source: \url{https://oeis.org/A164586}. Selected conjecture: R. J. Mathar, Feb 10 2015. The defining formula or established generating function is credited input from the primary entry; only the selected conjectural recurrence is treated.

\begin{align*}
p_{0}(n)&=n + 1\\
p_{1}(n)&=n - 1\\
p_{2}(n)&=- 2 \left(2 n + 1\right)\\
p_{3}(n)&=- 2 \left(7 n - 10\right)\\
p_{4}(n)&=- 14 \left(n - 2\right)\\
p_{5}(n)&=- 2 \left(2 n - 3\right)\\
p_{6}(n)&=n - 7\\
p_{7}(n)&=n - 5.
\end{align*}

Use $t=x$ and $\xi=x$.

The exact representation is $F=\sum_iR_i\sqrt{D_i}$, with $D_0=1$ and

\begin{align*}
R_{0}&=- \frac{1}{2 x \left(x + 1\right)},\\
R_{1}&=- \frac{1}{2 x \left(x + 1\right) \left(x^{3} + 2 x^{2} + 2 x - 1\right)},\\
D_{1}&=\left(x^{3} - 2 x^{2} - 2 x - 1\right) \left(x^{3} + 2 x^{2} + 2 x - 1\right).
\end{align*}

Inserting these data in \eqref{round5:eq:r5gf-cert} gives exactly

\[\mathcal L_pF=P(x)=x^{4} + 2 x^{3} + 2 x + 1.\]

All radical residuals in \eqref{round5:eq:r5gf-cert} vanish as rational identities. Therefore \eqref{round5:eq:r5gf-extract} proves the displayed recurrence for every integer $n\ge7$. The source offset is 0.

\subsubsection{A174013: Jan 12 2013}
\label{round5:r5entry:A174013}

Source: \url{https://oeis.org/A174013}. Selected conjecture: R. J. Mathar, Jan 12 2013. The defining formula or established generating function is credited input from the primary entry; only the selected conjectural recurrence is treated.

\begin{align*}
p_{0}(n)&=n + 1\\
p_{1}(n)&=n\\
p_{2}(n)&=- 12 \left(n - 1\right)\\
p_{3}(n)&=3 \left(3 n - 8\right)\\
p_{4}(n)&=6 \left(6 n - 23\right)\\
p_{5}(n)&=- 11 \left(n - 2\right)\\
p_{6}(n)&=253 - 49 n\\
p_{7}(n)&=- 20 \left(n - 6\right).
\end{align*}

Use $t=x$ and $\xi=x$.

The exact representation is $F=\sum_iR_i\sqrt{D_i}$, with $D_0=1$ and

\begin{align*}
R_{0}&=- \frac{x^{2} + 3 x - 1}{2 x \left(x^{2} + x - 1\right)},\\
R_{1}&=- \frac{1}{2 x \left(x^{2} + x - 1\right)},\\
D_{1}&=5 x^{4} + 6 x^{3} - x^{2} - 2 x + 1.
\end{align*}

Inserting these data in \eqref{round5:eq:r5gf-cert} gives exactly

\[\mathcal L_pF=P(x)=- \left(4 x + 1\right) \left(6 x^{4} - 3 x^{2} + x - 1\right).\]

All radical residuals in \eqref{round5:eq:r5gf-cert} vanish as rational identities. Therefore \eqref{round5:eq:r5gf-extract} proves the displayed recurrence for every integer $n\ge7$. The source offset is 0.

\subsubsection{A182879: Jun 14 2016}
\label{round5:r5entry:A182879}

Source: \url{https://oeis.org/A182879}. Selected conjecture: R. J. Mathar, Jun 14 2016. Kotesovec 2017 already gives a different order-four recurrence. Its existence is credited; the selected target is Mathar's order-five formula.

\begin{align*}
p_{0}(n)&=\left(n - 1\right) \left(310 n - 781\right)\\
p_{1}(n)&=- 882 n^{2} + 3845 n - 3679\\
p_{2}(n)&=214 n^{2} - 2903 n + 3751\\
p_{3}(n)&=- 358 n^{2} - 459 n + 1955\\
p_{4}(n)&=834 n^{2} - 3631 n + 4065\\
p_{5}(n)&=- \left(n - 3\right) \left(262 n - 663\right).
\end{align*}

Use $t=x$ and $\xi=x$.

The exact representation is $F=\sum_iR_i\sqrt{D_i}$, with $D_0=1$ and

\begin{align*}
R_{0}&=0,\\
R_{1}&=- \frac{x \left(x^{3} - 2 x^{2} - 1\right)}{\left(x^{2} - 3 x + 1\right)^{2} \left(x^{2} + x + 1\right)^{2}},\\
D_{1}&=\left(x^{2} - 3 x + 1\right) \left(x^{2} + x + 1\right).
\end{align*}

Inserting these data in \eqref{round5:eq:r5gf-cert} gives exactly

\[\mathcal L_pF=P(x)=0.\]

All radical residuals in \eqref{round5:eq:r5gf-cert} vanish as rational identities. Therefore \eqref{round5:eq:r5gf-extract} proves the displayed recurrence for every integer $n\ge5$. The source offset is 0.

\subsubsection{A182894: Jun 14 2016}
\label{round5:r5entry:A182894}

Source: \url{https://oeis.org/A182894}. Selected conjecture: R. J. Mathar, Jun 14 2016. The defining formula or established generating function is credited input from the primary entry; only the selected conjectural recurrence is treated.

\begin{align*}
p_{0}(n)&=n\\
p_{1}(n)&=3 - 4 n\\
p_{2}(n)&=n - 3\\
p_{3}(n)&=-3\\
p_{4}(n)&=3 \left(5 n - 14\right)\\
p_{5}(n)&=6 \left(n - 3\right)\\
p_{6}(n)&=6 \left(n - 4\right)\\
p_{7}(n)&=- 4 \left(n - 6\right).
\end{align*}

Use $t=x$ and $\xi=x$.

The exact representation is $F=\sum_iR_i\sqrt{D_i}$, with $D_0=1$ and

\begin{align*}
R_{0}&=\frac{x \left(x + 1\right)}{4 x^{3} + 2 x^{2} + 2 x - 1},\\
R_{1}&=- \frac{1}{4 x^{3} + 2 x^{2} + 2 x - 1},\\
D_{1}&=\left(x^{2} - 3 x + 1\right) \left(x^{2} + x + 1\right).
\end{align*}

Inserting these data in \eqref{round5:eq:r5gf-cert} gives exactly

\[\mathcal L_pF=P(x)=x \left(2 x^{4} + 3 x^{2} - x - 1\right).\]

All radical residuals in \eqref{round5:eq:r5gf-cert} vanish as rational identities. Therefore \eqref{round5:eq:r5gf-extract} proves the displayed recurrence for every integer $n\ge7$. The source offset is 0.

\subsubsection{A182904: Jun 14 2016}
\label{round5:r5entry:A182904}

Source: \url{https://oeis.org/A182904}. Selected conjecture: R. J. Mathar, Jun 14 2016. The defining formula or established generating function is credited input from the primary entry; only the selected conjectural recurrence is treated.

\begin{align*}
p_{0}(n)&=n\\
p_{1}(n)&=n - 2\\
p_{2}(n)&=10 - 7 n\\
p_{3}(n)&=- 3 \left(n - 2\right)\\
p_{4}(n)&=2 - n\\
p_{5}(n)&=14 - 5 n\\
p_{6}(n)&=18 - 5 n\\
p_{7}(n)&=3 \left(n - 4\right).
\end{align*}

Use $t=x$ and $\xi=x$.

The exact representation is $F=\sum_iR_i\sqrt{D_i}$, with $D_0=1$ and

\begin{align*}
R_{0}&=0,\\
R_{1}&=\frac{1}{x^{6} - 2 x^{5} - x^{4} - x^{2} - 2 x + 1},\\
D_{1}&=x^{6} - 2 x^{5} - x^{4} - x^{2} - 2 x + 1.
\end{align*}

Inserting these data in \eqref{round5:eq:r5gf-cert} gives exactly

\[\mathcal L_pF=P(x)=0.\]

All radical residuals in \eqref{round5:eq:r5gf-cert} vanish as rational identities. Therefore \eqref{round5:eq:r5gf-extract} proves the displayed recurrence for every integer $n\ge7$. The source offset is 0.

\subsubsection{A189176: Nov 15 2011}
\label{round5:r5entry:A189176}

Source: \url{https://oeis.org/A189176}. Selected conjecture: R. J. Mathar, Nov 15 2011. Eldar 2025 gives the Catalan closed form for positive indices. The exceptional initial value at zero is essential to the boundary qualification below.

\textbf{Boundary qualification.} The source does not give a starting index. $n=2$ gives residual $-5$. The exact valid tail is $n\ge3$. No explicit smaller starting index is attributed to the source.

\begin{align*}
p_{0}(n)&=n + 1\\
p_{1}(n)&=- 3 \left(3 n - 1\right)\\
p_{2}(n)&=10 \left(2 n - 3\right).
\end{align*}

Use $t=x$ and $\xi=x$.

The exact representation is $F=\sum_iR_i\sqrt{D_i}$, with $D_0=1$ and

\begin{align*}
R_{0}&=- \frac{x - 1}{2 x},\\
R_{1}&=- \frac{5 x - 1}{2 x \left(4 x - 1\right)},\\
D_{1}&=1 - 4 x.
\end{align*}

Inserting these data in \eqref{round5:eq:r5gf-cert} gives exactly

\[\mathcal L_pF=P(x)=- 5 x^{2} - 2 x + 1.\]

All radical residuals in \eqref{round5:eq:r5gf-cert} vanish as rational identities. Therefore \eqref{round5:eq:r5gf-extract} proves the displayed recurrence for every integer $n\ge3$. The source offset is 0.

\subsubsection{A190166: Apr 09 2019}
\label{round5:r5entry:A190166}

Source: \url{https://oeis.org/A190166}. Selected conjecture: R. J. Mathar, Apr 09 2019. The defining formula or established generating function is credited input from the primary entry; only the selected conjectural recurrence is treated.

\begin{align*}
p_{0}(n)&=1 - n\\
p_{1}(n)&=3 n - 4\\
p_{2}(n)&=- 2 \left(n - 1\right)\\
p_{3}(n)&=3 \left(n - 2\right)\\
p_{4}(n)&=- 2 \left(n - 3\right)\\
p_{5}(n)&=3 n - 8\\
p_{6}(n)&=3 - n.
\end{align*}

Use $t=x$ and $\xi=x$.

The exact representation is $F=\sum_iR_i\sqrt{D_i}$, with $D_0=1$ and

\begin{align*}
R_{0}&=0,\\
R_{1}&=\frac{x}{\left(x^{2} - 3 x + 1\right) \left(x^{2} - x + 1\right) \left(x^{2} + x + 1\right)},\\
D_{1}&=\left(x^{2} - 3 x + 1\right) \left(x^{2} + x + 1\right).
\end{align*}

Inserting these data in \eqref{round5:eq:r5gf-cert} gives exactly

\[\mathcal L_pF=P(x)=0.\]

All radical residuals in \eqref{round5:eq:r5gf-cert} vanish as rational identities. Therefore \eqref{round5:eq:r5gf-extract} proves the displayed recurrence for every integer $n\ge6$. The source offset is 0.

\subsubsection{A218185: Oct 08 2016}
\label{round5:r5entry:A218185}

Source: \url{https://oeis.org/A218185}. Selected conjecture: R. J. Mathar, Oct 08 2016. Yildiz 2012, Proposition 6.13, supplies the generating function. Mathar's2016 recurrence is the selected target, with its initial boundary stated explicitly.

\textbf{Boundary qualification.} The source does not give a starting index. $n=4$ gives residual $-16$. The exact valid tail is $n\ge5$. No explicit smaller starting index is attributed to the source.

\begin{align*}
p_{0}(n)&=n \left(n - 2\right) \left(n - 1\right)\\
p_{1}(n)&=- 12 \left(n - 2\right) \left(n - 1\right) \left(2 n - 5\right)\\
p_{2}(n)&=4 \left(n - 2\right) \left(52 n^{2} - 312 n + 465\right)\\
p_{3}(n)&=- 48 \left(2 n - 7\right) \left(8 n^{2} - 56 n + 95\right)\\
p_{4}(n)&=256 \left(n - 5\right) \left(2 n - 9\right) \left(2 n - 7\right).
\end{align*}

Use $t=x$ and $\xi=x$.

The exact representation is $F=\sum_iR_i\sqrt{D_i}$, with $D_0=1$ and

\begin{align*}
R_{0}&=\frac{4 x - 1}{4},\\
R_{1}&=- \frac{1}{4},\\
R_{2}&=\frac{1}{4},\\
R_{3}&=\frac{1}{4},\\
D_{1}&=1 - 8 x,\\
D_{2}&=\left(4 x - 1\right) \left(8 x - 1\right),\\
D_{3}&=1 - 4 x.
\end{align*}

Inserting these data in \eqref{round5:eq:r5gf-cert} gives exactly

\[\mathcal L_pF=P(x)=- 16 x^{4}.\]

All radical residuals in \eqref{round5:eq:r5gf-cert} vanish as rational identities. Therefore \eqref{round5:eq:r5gf-extract} proves the displayed recurrence for every integer $n\ge5$. The source offset is 0.

\subsubsection{A239201: Oct 07 2016}
\label{round5:r5entry:A239201}

Source: \url{https://oeis.org/A239201}. Selected conjecture: R. J. Mathar, Oct 07 2016. The defining formula or established generating function is credited input from the primary entry; only the selected conjectural recurrence is treated.

\begin{align*}
p_{0}(n)&=1 - n\\
p_{1}(n)&=7 n - 11\\
p_{2}(n)&=25 - 11 n\\
p_{3}(n)&=5 \left(n - 3\right).
\end{align*}

Use $t=x$ and $\xi=x$.

The exact representation is $F=\sum_iR_i\sqrt{D_i}$, with $D_0=1$ and

\begin{align*}
R_{0}&=- \frac{x}{2 \left(x - 1\right)},\\
R_{1}&=\frac{3 x}{2 \left(x - 1\right) \left(5 x - 1\right)},\\
D_{1}&=\left(x - 1\right) \left(5 x - 1\right).
\end{align*}

Inserting these data in \eqref{round5:eq:r5gf-cert} gives exactly

\[\mathcal L_pF=P(x)=x^{2}.\]

All radical residuals in \eqref{round5:eq:r5gf-cert} vanish as rational identities. Therefore \eqref{round5:eq:r5gf-extract} proves the displayed recurrence for every integer $n\ge4$. The source offset is 1.

\subsubsection{A243764: Jun 07 2016}
\label{round5:r5entry:A243764}

Source: \url{https://oeis.org/A243764}. Selected conjecture: R. J. Mathar, Jun 07 2016. The defining formula or established generating function is credited input from the primary entry; only the selected conjectural recurrence is treated.

\textbf{Boundary qualification.} The source does not give a starting index. $n=6$ gives residual $24576$. The exact valid tail is $n\ge7$. No explicit smaller starting index is attributed to the source.

\begin{align*}
p_{0}(n)&=245 n \left(n - 2\right) \left(n - 1\right)\\
p_{1}(n)&=- 140 \left(n - 2\right) \left(n - 1\right) \left(28 n - 81\right)\\
p_{2}(n)&=4 \left(n - 2\right) \left(5812 n^{2} - 41488 n + 72261\right)\\
p_{3}(n)&=- 16 \left(3368 n^{3} - 45144 n^{2} + 191968 n - 260859\right)\\
p_{4}(n)&=- 128 \left(152 n^{3} + 768 n^{2} - 16748 n + 45483\right)\\
p_{5}(n)&=2048 \left(2 n - 11\right) \left(74 n^{2} - 715 n + 1629\right)\\
p_{6}(n)&=- 98304 \left(n - 7\right) \left(2 n - 13\right) \left(2 n - 11\right).
\end{align*}

Use $x=\xi(t)=- t \left(t - 1\right)$, $t(0)=0$.

The exact representation is $F=\sum_iR_i\sqrt{D_i}$, with $D_0=1$ and

\begin{align*}
R_{0}&=\frac{t - 1}{2},\\
R_{1}&=\frac{t - 1}{2},\\
D_{1}&=1 - 8 t^{2}.
\end{align*}

Inserting these data in \eqref{round5:eq:r5gf-cert} gives exactly

\[\mathcal L_pF=P(x)=12 x^{3} \left(2048 x^{3} - 4096 x^{2} + 1352 x - 105\right).\]

All radical residuals in \eqref{round5:eq:r5gf-cert} vanish as rational identities. Therefore \eqref{round5:eq:r5gf-extract} proves the displayed recurrence for every integer $n\ge7$. The source offset is 0.

\subsubsection{A262768: Jun 07 2016}
\label{round5:r5entry:A262768}

Source: \url{https://oeis.org/A262768}. Selected conjecture: R. J. Mathar, Jun 07 2016. The defining formula or established generating function is credited input from the primary entry; only the selected conjectural recurrence is treated.

\textbf{Boundary qualification.} The source does not give a starting index. $n=2$ gives residual $-96$. The exact valid tail is $n\ge3$. No explicit smaller starting index is attributed to the source.

\begin{align*}
p_{0}(n)&=- 2 \left(n + 2\right) \left(49 n - 95\right)\\
p_{1}(n)&=325 \left(n - 1\right) \left(n + 1\right)\\
p_{2}(n)&=2 \left(2 n - 5\right) \left(67 n - 38\right).
\end{align*}

Use $t=x$ and $\xi=x$.

The exact representation is $F=\sum_iR_i\sqrt{D_i}$, with $D_0=1$ and

\begin{align*}
R_{0}&=\frac{\left(x - 1\right) \left(x + 1\right)}{2 x^{2}},\\
R_{1}&=\frac{2 x^{3} + 3 x^{2} + 2 x - 1}{2 x^{2} \left(4 x - 1\right)},\\
D_{1}&=1 - 4 x.
\end{align*}

Inserting these data in \eqref{round5:eq:r5gf-cert} gives exactly

\[\mathcal L_pF=P(x)=- 96 x^{2}.\]

All radical residuals in \eqref{round5:eq:r5gf-cert} vanish as rational identities. Therefore \eqref{round5:eq:r5gf-extract} proves the displayed recurrence for every integer $n\ge3$. The source offset is 0.

\subsubsection{A270490: Jun 07 2016}
\label{round5:r5entry:A270490}

Source: \url{https://oeis.org/A270490}. Selected conjecture: R. J. Mathar, Jun 07 2016. The defining formula or established generating function is credited input from the primary entry; only the selected conjectural recurrence is treated.

\begin{align*}
p_{0}(n)&=3 n \left(n - 2\right) \left(n + 2\right)\\
p_{1}(n)&=- 4 \left(n + 1\right) \left(8 n^{2} - 23 n + 12\right)\\
p_{2}(n)&=16 n \left(2 n - 5\right) \left(3 n - 4\right)\\
p_{3}(n)&=8 \left(2 n - 3\right) \left(4 n - 7\right)\\
p_{4}(n)&=- 64 \left(n - 3\right) \left(2 n - 5\right) \left(2 n - 3\right).
\end{align*}

Use $x=\xi(t)=- t \left(t - 1\right)$, $t(0)=0$.

The exact representation is $F=\sum_iR_i\sqrt{D_i}$, with $D_0=1$ and

\begin{align*}
R_{0}&=- \frac{1}{2 t^{2} \left(2 t - 1\right)},\\
R_{1}&=\frac{1}{2 t^{2} \left(2 t - 1\right)},\\
D_{1}&=- \left(2 t - 1\right) \left(2 t + 1\right).
\end{align*}

Inserting these data in \eqref{round5:eq:r5gf-cert} gives exactly

\[\mathcal L_pF=P(x)=- 2 x \left(8 x - 3\right).\]

All radical residuals in \eqref{round5:eq:r5gf-cert} vanish as rational identities. Therefore \eqref{round5:eq:r5gf-extract} proves the displayed recurrence for every integer $n\ge4$. The source offset is 0.

\subsubsection{A270530: Jun 07 2016}
\label{round5:r5entry:A270530}

Source: \url{https://oeis.org/A270530}. Selected conjecture: R. J. Mathar, Jun 07 2016. Only the first listed recurrence is selected. The known Catalan-binomial sum and radical generating function are inputs.

\begin{align*}
p_{0}(n)&=2 n \left(n + 1\right) \left(2 n + 3\right)\\
p_{1}(n)&=- n \left(77 n^{2} + 27 n - 4\right)\\
p_{2}(n)&=549 n^{3} - 987 n^{2} + 686 n - 168\\
p_{3}(n)&=- 20 \left(2 n - 3\right) \left(43 n^{2} - 104 n + 70\right)\\
p_{4}(n)&=500 \left(n - 2\right) \left(2 n - 5\right) \left(2 n - 3\right).
\end{align*}

Use $x=\xi(t)=- t \left(t - 1\right)$, $t(0)=0$.

The exact representation is $F=\sum_iR_i\sqrt{D_i}$, with $D_0=1$ and

\begin{align*}
R_{0}&=\frac{1}{2 t \left(t - 1\right) \left(2 t - 1\right)},\\
R_{1}&=\frac{1}{2 t \left(t - 1\right)^{2} \left(2 t - 1\right)},\\
D_{1}&=\left(t - 1\right) \left(5 t - 1\right).
\end{align*}

Inserting these data in \eqref{round5:eq:r5gf-cert} gives exactly

\[\mathcal L_pF=P(x)=0.\]

All radical residuals in \eqref{round5:eq:r5gf-cert} vanish as rational identities. Therefore \eqref{round5:eq:r5gf-extract} proves the displayed recurrence for every integer $n\ge4$. The source offset is 0.

\subsubsection{A273351: Jun 06 2016}
\label{round5:r5entry:A273351}

Source: \url{https://oeis.org/A273351}. Selected conjecture: R. J. Mathar, Jun 06 2016. The known bargraph generating function is an input. Only the first of the two separately conjectured recurrences is selected.

\begin{align*}
p_{0}(n)&=n \left(13 n - 40\right)\\
p_{1}(n)&=- 55 n^{2} + 211 n - 129\\
p_{2}(n)&=38 n^{2} - 212 n + 255\\
p_{3}(n)&=- 6 n^{2} + 56 n - 75\\
p_{4}(n)&=13 n^{2} - 66 n + 3\\
p_{5}(n)&=- \left(n - 6\right) \left(3 n - 13\right).
\end{align*}

Use $t=x$ and $\xi=x$.

The exact representation is $F=\sum_iR_i\sqrt{D_i}$, with $D_0=1$ and

\begin{align*}
R_{0}&=- \frac{x + 1}{2},\\
R_{1}&=- \frac{x^{3} + x^{2} + x - 1}{2 \left(x - 1\right) \left(x^{3} + x^{2} + 3 x - 1\right)},\\
D_{1}&=\left(x - 1\right) \left(x^{3} + x^{2} + 3 x - 1\right).
\end{align*}

Inserting these data in \eqref{round5:eq:r5gf-cert} gives exactly

\[\mathcal L_pF=P(x)=- 28 x^{2}.\]

All radical residuals in \eqref{round5:eq:r5gf-cert} vanish as rational identities. Therefore \eqref{round5:eq:r5gf-extract} proves the displayed recurrence for every integer $n\ge7$. The source offset is 2.

\subsubsection{A279013: Jan 27 2020}
\label{round5:r5entry:A279013}

Source: \url{https://oeis.org/A279013}. Selected conjecture: R. J. Mathar, Jan 27 2020. The defining formula or established generating function is credited input from the primary entry; only the selected conjectural recurrence is treated.

\textbf{Printed-formula repair.} The source has the undefined token \texttt{(2*n7)}. The coefficient of $a(n-4)$ below uses \texttt{(2*n-7)}. This repaired, well-defined formula is what is proved; the printed token is not silently normalized.

\begin{align*}
p_{0}(n)&=2 n \left(n + 1\right) \left(2 n - 3\right)\\
p_{1}(n)&=- n \left(101 n^{2} - 312 n + 203\right)\\
p_{2}(n)&=995 n^{3} - 5570 n^{2} + 9567 n - 4984\\
p_{3}(n)&=- 2 \left(2393 n^{3} - 19300 n^{2} + 50494 n - 42835\right)\\
p_{4}(n)&=4 \left(2 n - 7\right) \left(1408 n^{2} - 9889 n + 16690\right)\\
p_{5}(n)&=- 2600 \left(n - 5\right) \left(2 n - 9\right) \left(2 n - 7\right).
\end{align*}

Use $x=\xi(t)=- t \left(t - 1\right)$, $t(0)=0$.

The exact representation is $F=\sum_iR_i\sqrt{D_i}$, with $D_0=1$ and

\begin{align*}
R_{0}&=- \frac{\left(t - 1\right)^{2}}{2 t \left(2 t - 1\right)},\\
R_{1}&=- \frac{t - 1}{2 t \left(2 t - 1\right)},\\
D_{1}&=\left(t - 1\right) \left(5 t - 1\right).
\end{align*}

Inserting these data in \eqref{round5:eq:r5gf-cert} gives exactly

\[\mathcal L_pF=P(x)=- 6 x^{2} \left(52 x^{2} - 14 x + 1\right).\]

All radical residuals in \eqref{round5:eq:r5gf-cert} vanish as rational identities. Therefore \eqref{round5:eq:r5gf-extract} proves the displayed recurrence for every integer $n\ge5$. The source offset is 0.

\subsubsection{A279014: Mar 12 2017}
\label{round5:r5entry:A279014}

Source: \url{https://oeis.org/A279014}. Selected conjecture: R. J. Mathar, Mar 12 2017. Only the first listed conjecture is selected. The Fibonacci-binomial definition and listed generating function are prior input.

\textbf{Boundary qualification.} The source does not give a starting index. $n=4$ gives residual $49$. The exact valid tail is $n\ge5$. No explicit smaller starting index is attributed to the source.

\begin{align*}
p_{0}(n)&=n \left(3 n - 62\right)\\
p_{1}(n)&=3 n^{2} + 362 n - 247\\
p_{2}(n)&=- 171 n^{2} + 220 n + 162\\
p_{3}(n)&=417 n^{2} - 2570 n + 3551\\
p_{4}(n)&=2 \left(2 n - 7\right) \left(27 n - 59\right).
\end{align*}

Use $x=\xi(t)=- t \left(t - 1\right)$, $t(0)=0$.

The exact representation is $F=\sum_iR_i\sqrt{D_i}$, with $D_0=1$ and

\begin{align*}
R_{0}&=\frac{\left(t - 1\right)^{3}}{\left(2 t - 1\right) \left(t^{2} - 3 t + 1\right)},\\
R_{1}&=0,\\
D_{1}&=1.
\end{align*}

Inserting these data in \eqref{round5:eq:r5gf-cert} gives exactly

\[\mathcal L_pF=P(x)=x^{3} \left(49 x - 159\right).\]

All radical residuals in \eqref{round5:eq:r5gf-cert} vanish as rational identities. Therefore \eqref{round5:eq:r5gf-extract} proves the displayed recurrence for every integer $n\ge5$. The source offset is 0.

\subsubsection{A387084: Aug 19 2025}
\label{round5:r5entry:A387084}

Source: \url{https://oeis.org/A387084}. Selected conjecture: R. J. Mathar, Aug 19 2025. The already listed generating functions involving B and the quintic Fuss--Catalan series g are credited input. The selected claim is Mathar's August 2025 order-four recurrence.

\begin{align*}
p_{0}(n)&=3902464 n \left(8 n - 5\right) \left(8 n - 3\right) \left(8 n - 1\right) \left(8 n + 1\right)\\
p_{1}(n)&=- 1005260800000 n^{5} + 5475840000000 n^{4}\\
&\quad - 13076601280000 n^{3} + 16038028320000 n^{2}\\
&\quad +2384377435440 - 9827447160880 n\\
p_{2}(n)&=16756875000000 n^{5} - 138724375000000 n^{4}\\
&\quad +465871453125000 n^{3} - 788390515625000 n^{2}\\
&\quad +665727720750000 n - 221182599750000\\
p_{3}(n)&=- 76019287109375 n^{5} + 773712158203125 n^{4}\\
&\quad - 2961907958984375 n^{3} + 5069256591796875 n^{2}\\
&\quad +67940625000000 - 3308797753906250 n\\
p_{4}(n)&=- 95367431640625 n^{5} + 1621246337890625 n^{4}\\
&\quad - 11005401611328125 n^{3} + 37288665771484375 n^{2}\\
&\quad +42583007812500000 - 63060607910156250 n.
\end{align*}

Use $x=\xi(t)=\frac{t}{\left(t + 1\right)^{5}}$, $t(0)=0$.

The exact representation is $F=\sum_iR_i\sqrt{D_i}$, with $D_0=1$ and

\begin{align*}
R_{0}&=0,\\
R_{1}&=- \frac{t + 1}{4 t - 1},\\
D_{1}&=1 - 4 t.
\end{align*}

Inserting these data in \eqref{round5:eq:r5gf-cert} gives exactly

\[\mathcal L_pF=P(x)=0.\]

All radical residuals in \eqref{round5:eq:r5gf-cert} vanish as rational identities. Therefore \eqref{round5:eq:r5gf-extract} proves the displayed recurrence for every integer $n\ge4$. The source offset is 0.

\subsection{Verification and scope of prior-work audit}
The portable verifier uses exact rational arithmetic to check all29
algebraic certificates and the one hypergeometric identity. It compares
each selected polynomial with the retained primary-source mathematical
statement, treating the two explicit formula changes separately.
It independently expands the generating functions through sixteen
source terms for each entry, for 480 matched source terms, and checks
ten of the sum definitions independently as well. These finite checks
guard indexing and branch choices; the rational identities provide the
proofs for all indices.

The thirty retained IDs are disjoint from the 408 earlier counted IDs,
from the 33 explicit Ekhad--Yang 2017 proof IDs checked in the preceding
round, and from the retained OEIS Open Appendix A comparison list.
The last comparison covers only that stated appendix baseline, not
every problem or result in the OEIS Open project. Current primary-entry
proof notices were screened, and known covered entries were removed.
The audit does not assert absence of a proof anywhere in the literature.

\subsection{Retrospective correction to five Round IV priority assessments}
\label{round5:sec:r5retrospective}
The current A176604 page links Georg Fischer's dated 26 January 2020
derivation for a parameterized convolution family:
\url{https://oeis.org/A176604/a176604_2.txt}.
Its general differential equation and explicit parameter table cover
the previously counted A176605, A176607, A176610, A176611 and A176612.
For the last three, substituting $(m,k,j)=(0,1,1),(1,1,1),(2,1,1)$
gives exactly the Round IV recurrence coefficients.
For the first two, substituting $(1,0,1)$ and $(4,0,1)$ gives
\[
\mathcal L_{\rm Fischer}=(1-x)\mathcal L_{\rm Round IV}.
\]
Cancelling the common factor in Fischer's homogeneous equation for
the radical part gives the order-four equations used in Round IV.
All coefficient identities are retained in the accompanying historical
audit and checked exactly. Therefore those five Round IV proofs remain
mathematically valid reproductions, but their assessment as previously
unresolved claims must be withdrawn. They are not counted in Round V.
The family is excluded from the current hunt; its source retaining
a conjecture label on neighboring entries does not establish novelty.



\Needspace{6\baselineskip}
\section{Forty-five empirical array recurrences}
\label{round5:sec:r5arrays}
This section proves 45 different entry-level recurrence assertions marked
\emph{Empirical} on the cited OEIS pages, consulted on 20 October 2026.
The underlying arrays differ in width, alphabet, or local condition.
All 45 A-numbers are absent from the 408 entries counted in rounds I--IV.
The argument uses one general transfer principle with two precisely
specified state constructions.  The count is a count of selected OEIS
assertions, not a count of unrelated mathematical techniques.
The consulted pages contain data links and no proof of the selected
recurrence.  This limited source audit does not establish global priority.

\subsection{Signed clockwise differences: the new transfer construction}
Forty-one entries count, or take a stated scalar multiple of the count of,
arrays on $\{0,\ldots,q-1\}$ with $n+1$ rows and fixed width $w$.
For a block $\left(\begin{smallmatrix}a&b\\c&d\end{smallmatrix}\right)$,
define
\[
 e(a,b,c,d)=\#\{b-a,d-b,c-d,a-c\}.
\]
These are ordinary signed integer differences; no reduction modulo $q$
or absolute-value operation is used.  The local requirement is
$e(a,b,c,d)\in S$, where $S\subseteq\{1,2,3,4\}$ is specified below.
The source normalization $\lambda$ is one of $1$, $1/2$, $1/4$ and is
part of the definition, not an assumption about color orbits.

Let $H(r_1,\ldots,r_w)=(r_w,\ldots,r_1)$ and
$C(r_1,\ldots,r_w)=(q-1-r_1,\ldots,q-1-r_w)$.
Let $\mathcal O_{q,w}$ be the lexicographically least representatives of
orbits under $\{I,H,C,HC\}$, and write $\gamma(r)$ for a representative.
Horizontal reflection permutes and negates the clockwise differences;
simultaneous complementation negates them.  Hence each operation preserves
their number, so it preserves the local condition.  No other color
permutation is used in this construction.
For $r,s\in\mathcal O_{q,w}$ define
\begin{align}
 m(r)&=\#\{v\in\{0,\ldots,q-1\}^w:\gamma(v)=r\},\label{round5:eq:r5arr-mass}\\
 T(r,s)&=\#\left\{v\in\{0,\ldots,q-1\}^{w}:
 \begin{array}{l}
 \gamma(v)=s,\\
 e(r_i,r_{i+1},v_i,v_{i+1})\in S\quad(1\le i<w)
 \end{array}\right\}.\label{round5:eq:r5arr-edgeT}
\end{align}
Every row in an orbit has the same number of legal continuations to
each orbit: apply the same symmetry to both rows.  Thus quotienting loses
no information needed for counting.  Summing the actual first-row
labelings and adding each subsequent row gives, by induction,
\begin{equation}\label{round5:eq:r5arr-edgecount}
 a(n)=\lambda m^{\mathsf T}T^n\mathbf1\qquad(n\ge1).
\end{equation}
The $n$ transitions correspond to $n+1$ rows.  The vector $m$ is not
replaced by a uniform orbit size; fixed rows and two-element orbits are
counted with their true sizes.

\subsection{Four entries with colors renamed in first-occurrence order}
The remaining four entries restrict the number of distinct \emph{cell
values} in each block and identify global color renamings by reading the
array row by row.  Their state construction adapts the method developed
for the Round IV array packet.
For $j$ available labeled colors let $\mathcal R_{j,w}$ be the row
patterns obtained by renaming colors in order of first appearance, and
write $\operatorname{can}(v)$ for this row pattern.
For $r\in\mathcal R_{j,w}$ with $k(r)$ distinct colors set
$m_j(r)=(j)_{k(r)}=j!/(j-k(r))!$, and define
\begin{equation}\label{round5:eq:r5arr-colorT}
 T_{j,w,S}(r,s)=\#\left\{v\in\{0,\ldots,j-1\}^{w}:
 \begin{array}{l}
 \operatorname{can}(v)=s,\\
 \#\{r_i,r_{i+1},v_i,v_{i+1}\}\in S\quad(1\le i<w)
 \end{array}\right\}.
\end{equation}
All successor rows are enumerated, including rows using colors absent
from the previous row.  Simultaneous color permutations preserve this
local condition.  The same induction therefore gives the labeled count
$L_j(n)=m_j^{\mathsf T}T_{j,w,S}^{n}\mathbf1$.
Dividing a surjective labeled count by $q!$, and applying
inclusion--exclusion to omitted colors, gives
\begin{align}
 U_q^{=}(n)&=\sum_{j=1}^{q}
       \frac{(-1)^{q-j}}{j!(q-j)!}L_j(n),\label{round5:eq:r5arr-exact}\\
 U_q^{\le}(n)&=\sum_{j=1}^{q}
       \left(\frac1{j!}\sum_{h=0}^{q-j}\frac{(-1)^h}{h!}\right)L_j(n).
       \label{round5:eq:r5arr-atmost}
\end{align}
The second identity is the sum of the exact-color identity over the
actual number of colors used.  In particular, the requirement to use
all $q$ colors is imposed only when the source says so.  The empty
alphabet contributes zero because the arrays are nonempty.

\subsection{The all-index proof from exact finite certificates}
For the coefficient vector $c=(c_1,\ldots,c_d)$ printed for an entry, set
\[
 R(z)=z^d-\sum_{i=1}^d c_i z^{d-i},\qquad
 Q(x)=1-\sum_{i=1}^d c_i x^i.
\]
For its edge-difference matrix, or for every color matrix carrying a
nonzero weight in \eqref{round5:eq:r5arr-exact}--\eqref{round5:eq:r5arr-atmost}, the
complete exact integer certificate establishes
\begin{equation}\label{round5:eq:r5arr-cert}
 \boxed{\ T R(T)\mathbf1=\mathbf0.\ }
\end{equation}
For $n\ge d+1$, multiply \eqref{round5:eq:r5arr-cert} by
$m^{\mathsf T}T^{n-d-1}$, then by the appropriate scalar weight, and
sum the color components when necessary.  Equations
\eqref{round5:eq:r5arr-edgecount}--\eqref{round5:eq:r5arr-atmost} give
\begin{equation}\label{round5:eq:r5arr-recurrence}
 \boxed{\ a(n)=\sum_{i=1}^{d}c_i a(n-i)\quad(n\ge d+1).\ }
\end{equation}
All indices on the right are at least 1, as required by every selected
entry's offset.  We make no use of a source value $a(0)$.  The transfer
formula can have an auxiliary nonintegral value at $n=0$ for a source
normalized by $1/2$; this does not affect any asserted sequence value.

The proof package contains all 56 integer matrices (42,001 nonzero
weighted edges), all Krylov vectors $v_k=T^k\mathbf1$ for
$0\le k\le d+1$, and every residual coordinate in
\eqref{round5:eq:r5arr-cert}.  The largest single matrix has 574 states and the
largest recurrence has order 72.  The complete matrices are in
\texttt{r5\_arrays\_matrices.json}; the vectors, coefficient lists,
weights and residuals are in \texttt{r5\_arrays\_certificates.json}.
These are finite proof certificates, supplied as part of this proof.
The verifier rebuilds every state and every transition directly from
\eqref{round5:eq:r5arr-edgeT} or \eqref{round5:eq:r5arr-colorT}, then performs the
following exact calculation:
\begin{Verbatim}[fontsize=\small]
v[0] = [1 for each state]
for k = 1,...,d+1:
    v[k][r] = sum(T[r,s] * v[k-1][s] for all states s)
for every state r:
    assert v[d+1][r] == sum(c[i] * v[d+1-i][r] for i=1,...,d)
\end{Verbatim}
The matrix definitions, complete coefficient signatures below and this
finite algorithm specify the entire certificate check without any fitted
terms or numerical approximations.  Multiplication by further powers of
$T$ is the reason the conclusion holds for every indicated index.
The separate comparison with all 687 printed source-prefix terms checks
the source interpretation and indexing; it is not used to infer the
recurrences.  The 20-term computed prefixes are pairwise distinct, so
none of these 45 selected sequences are identical to another in this
section.

For $F(x)=\sum_{n\ge1}a(n)x^n$, the proved recurrence yields $QF=P$,
where $\deg P\le d$ and
\[
 [x^k]P(x)=\sum_{\substack{0\le i\le\min(k,d)\\k-i\ge1}}
               q_i a(k-i),\qquad q_0=1,\quad q_i=-c_i\ (i\ge1).
\]
The initial values in this formula are calculated from the defining
matrices.  All 17 rational generating functions marked empirical or
conjectural on the retained pages agree exactly with this $P/Q$.
The full numerator, denominator and initial-value lists are in the
certificate file.  These generating functions add no extra entry count.
Other assertions on a page, such as a separately displayed closed form,
are outside the selected claim unless expressly listed as proved here.

\subsection{A fully displayed example: A209953}
Here $q=3$, $w=2$, $S=\{1,3,4\}$ and $\lambda=1$.
The four representatives are $00,01,02,11$, with orbit sizes
$2,4,2,1$.  The exact construction gives
\[
 T=\begin{pmatrix}
 2&4&2&1\\2&2&2&1\\2&4&1&1\\2&4&2&1
 \end{pmatrix},\qquad
 m=\begin{pmatrix}2\\4\\2\\1\end{pmatrix}.
\]
Direct multiplication gives
\[
 v_1=\begin{pmatrix}9\\7\\8\\9\end{pmatrix},\quad
 v_2=\begin{pmatrix}71\\57\\63\\71\end{pmatrix},\quad
 v_3=\begin{pmatrix}567\\453\\504\\567\end{pmatrix}
       =6v_2+15v_1+6v_0.
\]
Therefore $R(T)\mathbf1=0$ for $R(z)=z^3-6z^2-15z-6$, which
in particular implies \eqref{round5:eq:r5arr-cert}.  The initial sequence
values are $71,567,4521$, and
\[
 F(x)=\frac{x(71+141x+54x^2)}{1-6x-15x^2-6x^3}.
\]
This proves the source recurrence for $n\ge4$ and its empirical
generating function.  The matrices for larger widths are checked by
exactly the same coordinate identity.

\subsection{Source audit, exclusions and earlier work}
The source status is the explicit word \emph{Empirical} attached to the
particular recurrence, not the OEIS editorial status ``approved.''
The audit compares all IDs against the 408 previously counted entries.
It also finds no intersection with the stored list of 97 distinct A-numbers
appearing in the 100 costly claims in Appendix A of \emph{OEIS Open};
that is a limited appendix audit,
not an audit of that project's entire repository or all mathematical
literature.  A209722 was excluded because its current page explicitly
identifies it as an elementary shifted form of a named parent sequence.

A210082, A210263 and A210305 were investigated provisionally during
Round IV but were not in its final counted 100.  Their current source
claims and complete certificates were rebuilt here.  Their inclusion
means first inclusion in our counted packets, not first observation
within this project.  The color-orbit method and part of its software
are credited to the earlier Round IV packet.  No OEIS entry was edited
or submitted by this work, and no official OEIS closure is claimed.

\subsection{Individual assertions and complete recurrence signatures}
For each entry, the coefficient list is $c_1,\ldots,c_d$ in that order,
including zero coefficients.  The given matrix dimensions refer to all
nonzero weighted components.  Every displayed signature satisfies
\eqref{round5:eq:r5arr-cert}, and the source's $a(1),a(2),a(3)$ are also printed
to fix the indexing unambiguously.
\subsubsection{A209505}
\label{round5:r5entry:A209505}
Arrays on $0,\ldots,2$ have $n+1$ rows and width $w=3$. The number of distinct signed clockwise edge differences in each block lies in $S=\{2,4\}$, and the count is multiplied by $\lambda=1/2$.
The recurrence \eqref{round5:eq:r5arr-recurrence} holds for every $n\ge 3$, with order $d=2$ and component dimensions $(10)$ (total $D=10$). Source: \url{https://oeis.org/A209505}.
The first three values are $(27,123,561)$.
\[ (c_1,\ldots,c_d)=(5,-2).\]
The source's empirical/conjectural rational generating function is also proved by the exact identity $QF=P$.
\subsubsection{A209506}
\label{round5:r5entry:A209506}
Arrays on $0,\ldots,2$ have $n+1$ rows and width $w=4$. The number of distinct signed clockwise edge differences in each block lies in $S=\{2,4\}$, and the count is multiplied by $\lambda=1/2$.
The recurrence \eqref{round5:eq:r5arr-recurrence} holds for every $n\ge 4$, with order $d=3$ and component dimensions $(25)$ (total $D=25$). Source: \url{https://oeis.org/A209506}.
The first three values are $(81,561,3906)$.
\[ (c_1,\ldots,c_d)=(9,-15,6).\]
The source's empirical/conjectural rational generating function is also proved by the exact identity $QF=P$.
\subsubsection{A209507}
\label{round5:r5entry:A209507}
Arrays on $0,\ldots,2$ have $n+1$ rows and width $w=5$. The number of distinct signed clockwise edge differences in each block lies in $S=\{2,4\}$, and the count is multiplied by $\lambda=1/2$.
The recurrence \eqref{round5:eq:r5arr-recurrence} holds for every $n\ge 6$, with order $d=5$ and component dimensions $(70)$ (total $D=70$). Source: \url{https://oeis.org/A209507}.
The first three values are $(243,2559,27225)$.
\[ (c_1,\ldots,c_d)=(16,-65,92,-48,8).\]
The source's empirical/conjectural rational generating function is also proved by the exact identity $QF=P$.
\subsubsection{A209508}
\label{round5:r5entry:A209508}
Arrays on $0,\ldots,2$ have $n+1$ rows and width $w=6$. The number of distinct signed clockwise edge differences in each block lies in $S=\{2,4\}$, and the count is multiplied by $\lambda=1/2$.
The recurrence \eqref{round5:eq:r5arr-recurrence} holds for every $n\ge 8$, with order $d=7$ and component dimensions $(196)$ (total $D=196$). Source: \url{https://oeis.org/A209508}.
The first three values are $(729,11673,189801)$.
\[ (c_1,\ldots,c_d)=(30,-291,1278,-2901,3519,-2152,516).\]
\subsubsection{A209509}
\label{round5:r5entry:A209509}
Arrays on $0,\ldots,2$ have $n+1$ rows and width $w=7$. The number of distinct signed clockwise edge differences in each block lies in $S=\{2,4\}$, and the count is multiplied by $\lambda=1/2$.
The recurrence \eqref{round5:eq:r5arr-recurrence} holds for every $n\ge 14$, with order $d=13$ and component dimensions $(574)$ (total $D=574$). Source: \url{https://oeis.org/A209509}.
The first three values are $(2187,53247,1323270)$.
\begin{Verbatim}[fontsize=\scriptsize,breaklines=true]
55, -1109, 11330, -67206, 247404, -582440, 881876, -846764, 499200, -172400, 33152,
-3264, 128
\end{Verbatim}
\subsubsection{A209530}
\label{round5:r5entry:A209530}
Arrays on $0,\ldots,2$ have $n+1$ rows and width $w=3$. The number of distinct signed clockwise edge differences in each block lies in $S=\{2\}$, and the count is multiplied by $\lambda=1/2$.
The recurrence \eqref{round5:eq:r5arr-recurrence} holds for every $n\ge 4$, with order $d=3$ and component dimensions $(10)$ (total $D=10$). Source: \url{https://oeis.org/A209530}.
The first three values are $(9,25,65)$.
\[ (c_1,\ldots,c_d)=(1,8,-8).\]
The source's empirical/conjectural rational generating function is also proved by the exact identity $QF=P$.
\subsubsection{A209531}
\label{round5:r5entry:A209531}
Arrays on $0,\ldots,2$ have $n+1$ rows and width $w=5$. The number of distinct signed clockwise edge differences in each block lies in $S=\{2\}$, and the count is multiplied by $\lambda=1/2$.
The recurrence \eqref{round5:eq:r5arr-recurrence} holds for every $n\ge 4$, with order $d=3$ and component dimensions $(70)$ (total $D=70$). Source: \url{https://oeis.org/A209531}.
The first three values are $(33,193,1025)$.
\[ (c_1,\ldots,c_d)=(1,32,-32).\]
The source's empirical/conjectural rational generating function is also proved by the exact identity $QF=P$.
\subsubsection{A209546}
\label{round5:r5entry:A209546}
Arrays on $0,\ldots,3$ have $n+1$ rows and width $w=2$. The number of distinct signed clockwise edge differences in each block lies in $S=\{2\}$, and the count is multiplied by $\lambda=1/4$.
The recurrence \eqref{round5:eq:r5arr-recurrence} holds for every $n\ge 5$, with order $d=4$ and component dimensions $(6)$ (total $D=6$). Source: \url{https://oeis.org/A209546}.
The first three values are $(7,20,61)$.
\[ (c_1,\ldots,c_d)=(5,-6,0,2).\]
The source's empirical/conjectural rational generating function is also proved by the exact identity $QF=P$.
\subsubsection{A209547}
\label{round5:r5entry:A209547}
Arrays on $0,\ldots,3$ have $n+1$ rows and width $w=3$. The number of distinct signed clockwise edge differences in each block lies in $S=\{2\}$, and the count is multiplied by $\lambda=1/4$.
The recurrence \eqref{round5:eq:r5arr-recurrence} holds for every $n\ge 7$, with order $d=6$ and component dimensions $(20)$ (total $D=20$). Source: \url{https://oeis.org/A209547}.
The first three values are $(20,105,562)$.
\[ (c_1,\ldots,c_d)=(8,-14,-3,20,-8,-2).\]
\subsubsection{A209548}
\label{round5:r5entry:A209548}
Arrays on $0,\ldots,3$ have $n+1$ rows and width $w=4$. The number of distinct signed clockwise edge differences in each block lies in $S=\{2\}$, and the count is multiplied by $\lambda=1/4$.
The recurrence \eqref{round5:eq:r5arr-recurrence} holds for every $n\ge 12$, with order $d=11$ and component dimensions $(72)$ (total $D=72$). Source: \url{https://oeis.org/A209548}.
The first three values are $(61,562,5075)$.
\begin{Verbatim}[fontsize=\scriptsize,breaklines=true]
15, -59, 27, 267, -488, 111, 312, -174, -30, 18, 2
\end{Verbatim}
\subsubsection{A209549}
\label{round5:r5entry:A209549}
Arrays on $0,\ldots,3$ have $n+1$ rows and width $w=5$. The number of distinct signed clockwise edge differences in each block lies in $S=\{2\}$, and the count is multiplied by $\lambda=1/4$.
The recurrence \eqref{round5:eq:r5arr-recurrence} holds for every $n\ge 21$, with order $d=20$ and component dimensions $(272)$ (total $D=272$). Source: \url{https://oeis.org/A209549}.
The first three values are $(191,3051,46515)$.
\begin{Verbatim}[fontsize=\scriptsize,breaklines=true]
26, -177, 43, 3477, -10179, -6160, 62922, -59442, -82557, 141262, 24872, -112175, 5660,
41275, -1936, -7049, -319, 400, 56, 2
\end{Verbatim}
\subsubsection{A209710}
\label{round5:r5entry:A209710}
Arrays on $0,\ldots,3$ have $n+1$ rows and width $w=3$. The number of distinct signed clockwise edge differences in each block lies in $S=\{4\}$, and the count is multiplied by $\lambda=1/4$.
The recurrence \eqref{round5:eq:r5arr-recurrence} holds for every $n\ge 20$, with order $d=19$ and component dimensions $(20)$ (total $D=20$). Source: \url{https://oeis.org/A209710}.
The first three values are $(175,2710,40600)$.
\begin{Verbatim}[fontsize=\scriptsize,breaklines=true]
9, 155, -718, -5017, 21111, 51148, -251141, -85963, 1092718, -438881, -1923975, 1461077,
1300250, -1385402, -184926, 434126, -73408, -12888, 1440
\end{Verbatim}
\subsubsection{A209711}
\label{round5:r5entry:A209711}
Arrays on $0,\ldots,3$ have $n+1$ rows and width $w=4$. The number of distinct signed clockwise edge differences in each block lies in $S=\{4\}$, and the count is multiplied by $\lambda=1/4$.
The recurrence \eqref{round5:eq:r5arr-recurrence} holds for every $n\ge 65$, with order $d=64$ and component dimensions $(72)$ (total $D=72$). Source: \url{https://oeis.org/A209711}.
The first three values are $(1241,40600,1277340)$.
\begin{Verbatim}[fontsize=\scriptsize,breaklines=true]
15, 1008, -7876, -328904, 1926537, 52684864, -273951748, -4866393539, 24240653273,
284467521073, -1401913368926, -11199860935130, 55383623152305, 310221810205502,
-1550263153851120, -6235493002893665, 31634306714134207, 92998581142818874,
-480833347610958673, -1045915789998724699, 5532194062321947358, 8975769923693909839,
-48756799709225176849, -59299615298869844599, 332041426928348290121,
303700687565429488003, -1758256324787498237829, -1212811908997938780012,
7270465412741619864800, 3797566692114290775489, -23537814548092748856931,
-9379511582307506472416, 59726661084682197089538, 18401973290976709236198,
-118737454342588779056334, -28914958687246461124391, 184561198236942732186820,
36699265124914024312999, -223429523113044656251384, -37861805098615603112860,
209387912188279839779389, 31751887956162117388848, -150583340090115952826545,
-21431798119336812426970, 82105995000110413699734, 11401867924374467330034,
-33394113408712074250948, -4638404024514496142694, 9913173513982954582346,
1389760315203109342656, -2086440149254353119856, -293279980685597445408,
299468489710238488568, 41250131799294016000, -27806351608888814224,
-3591128372719351552, 1556040144706074432, 173746695509329664, -48068935504144192,
-3984150632877312, 742241025193984, 34721251272704, -4259930669056, -57036898304
\end{Verbatim}
\subsubsection{A209851}
\label{round5:r5entry:A209851}
Arrays on $0,\ldots,3$ have $n+1$ rows and width $w=2$. The number of distinct signed clockwise edge differences in each block lies in $S=\{1,4\}$, and the count is multiplied by $\lambda=1/4$.
The recurrence \eqref{round5:eq:r5arr-recurrence} holds for every $n\ge 7$, with order $d=6$ and component dimensions $(6)$ (total $D=6$). Source: \url{https://oeis.org/A209851}.
The first three values are $(25,182,1308)$.
\[ (c_1,\ldots,c_d)=(4,29,-28,-110,82,-8).\]
The source's empirical/conjectural rational generating function is also proved by the exact identity $QF=P$.
\subsubsection{A209852}
\label{round5:r5entry:A209852}
Arrays on $0,\ldots,3$ have $n+1$ rows and width $w=3$. The number of distinct signed clockwise edge differences in each block lies in $S=\{1,4\}$, and the count is multiplied by $\lambda=1/4$.
The recurrence \eqref{round5:eq:r5arr-recurrence} holds for every $n\ge 21$, with order $d=20$ and component dimensions $(20)$ (total $D=20$). Source: \url{https://oeis.org/A209852}.
The first three values are $(182,2829,43091)$.
\begin{Verbatim}[fontsize=\scriptsize,breaklines=true]
11, 139, -972, -4071, 28751, 27819, -332527, 170449, 1380897, -1533070, -2378942,
3871609, 1392648, -4192842, 445370, 1893206, -633950, -236300, 99664, 2096
\end{Verbatim}
\subsubsection{A209853}
\label{round5:r5entry:A209853}
Arrays on $0,\ldots,3$ have $n+1$ rows and width $w=4$. The number of distinct signed clockwise edge differences in each block lies in $S=\{1,4\}$, and the count is multiplied by $\lambda=1/4$.
The recurrence \eqref{round5:eq:r5arr-recurrence} holds for every $n\ge 69$, with order $d=68$ and component dimensions $(72)$ (total $D=72$). Source: \url{https://oeis.org/A209853}.
The first three values are $(1308,43091,1386131)$.
\begin{Verbatim}[fontsize=\scriptsize,breaklines=true]
20, 963, -12793, -312156, 3571443, 49605457, -553758368, -4457972689, 52612881372,
243417930360, -3275949406092, -8240107440663, 140666492498369, 160116943902217,
-4326535316148209, -817021352845345, 98006065095893985, -49203318829660909,
-1668148030882549725, 1723999493198171149, 21637008392085651838, -31773088576337392703,
-215843167651091106385, 399719679271900244237, 1664018938648323383901,
-3697329744805424755681, -9917310558011937270066, 25982170466786070131093,
45447978346529331965091, -141117691880001505315599, -157856587989686977889953,
598069670716001721149653, 401968717089020637279250, -1987776050324662023234413,
-686660571916972008322547, 5191091867684820767344742, 522733262632635990820808,
-10646069432263823711177799, 915809068137700417207761, 17098571599529269338686083,
-3943750897982196417975435, -21398685175106614262349856, 7368362612864864601239602,
20709879068969481472447106, -9015916934517168198203128, -15335247867307455505750474,
7824208837288323351242620, 8561662593050307213736807, -4920492104705864543228925,
-3532336472390984912564596, 2242642650064493959464073, 1047216070797663839803065,
-731775898101400159982255, -214107873494035007863503, 167132111313457973560540,
28238564794974400827975, -25825554003258952415666, -2098980376078516736024,
2571400424053876584470, 52829982196209769868, -153859888970921746848,
2954600261893327248, 4977837135363185688, -211003278535383648, -72160266595102240,
3672896320734208, 335148663461120, -16066999343104
\end{Verbatim}
\subsubsection{A209906}
\label{round5:r5entry:A209906}
Arrays on $0,\ldots,3$ have $n+1$ rows and width $w=2$. The number of distinct signed clockwise edge differences in each block lies in $S=\{2,4\}$, and the count is multiplied by $\lambda=1/2$.
The recurrence \eqref{round5:eq:r5arr-recurrence} holds for every $n\ge 7$, with order $d=6$ and component dimensions $(6)$ (total $D=6$). Source: \url{https://oeis.org/A209906}.
The first three values are $(62,562,5008)$.
\[ (c_1,\ldots,c_d)=(7,21,-30,-30,12,8).\]
The source's empirical/conjectural rational generating function is also proved by the exact identity $QF=P$.
\subsubsection{A209907}
\label{round5:r5entry:A209907}
Arrays on $0,\ldots,3$ have $n+1$ rows and width $w=3$. The number of distinct signed clockwise edge differences in each block lies in $S=\{2,4\}$, and the count is multiplied by $\lambda=1/2$.
The recurrence \eqref{round5:eq:r5arr-recurrence} holds for every $n\ge 20$, with order $d=19$ and component dimensions $(20)$ (total $D=20$). Source: \url{https://oeis.org/A209907}.
The first three values are $(562,13098,295448)$.
\begin{Verbatim}[fontsize=\scriptsize,breaklines=true]
16, 227, -1440, -8244, 46409, 53926, -406100, -117231, 1484559, 166049, -2587430,
-386782, 2096907, 472125, -660956, -166954, 44616, 4488, -192
\end{Verbatim}
\subsubsection{A209908}
\label{round5:r5entry:A209908}
Arrays on $0,\ldots,3$ have $n+1$ rows and width $w=4$. The number of distinct signed clockwise edge differences in each block lies in $S=\{2,4\}$, and the count is multiplied by $\lambda=1/2$.
The recurrence \eqref{round5:eq:r5arr-recurrence} holds for every $n\ge 66$, with order $d=65$ and component dimensions $(72)$ (total $D=72$). Source: \url{https://oeis.org/A209908}.
The first three values are $(5008,295448,16664508)$.
\begin{Verbatim}[fontsize=\scriptsize,breaklines=true]
55, 923, -47071, -231407, 15527254, 7880250, -2573669497, 3681229288, 249190408738,
-579450712511, -15545777999632, 42210038348626, 666118006004289, -1886332941998678,
-20450809195283436, 57107262824986158, 462837772656461540, -1233589210852233412,
-7872410572016992734, 19630990804163291243, 101986218315556639540,
-235100512508203944296, -1015682793390488125566, 2150002789746681905996,
7825721744951506498234, -15163164422139297144181, -46841616343312521746791,
82995111044360045734710, 218301161066990302616716, -353790955288161781138915,
-792601249983629087376567, 1175946145985688378644800, 2239982343619860900022721,
-3045446389006222439252829, -4916701577362858573256478, 6130589140664657413790376,
8353918309455200803527969, -9556132008935803121684739, -10939049915929290500808672,
11474688682455975032157676, 10978124927194304878629570, -10543831886760113161536049,
-8383988700731830638859763, 7352253368826483234978608, 4826725129284052501397921,
-3849463056922623709810553, -2067854150675962492097465, 1492952670444499632838573,
647468497254091793151276, -421469451029786639157580, -144455861921375175788532,
84661307211707006760020, 22156825180654649239287, -11740439782116785076012,
-2219666541575216243128, 1077870975515912576792, 134607253301714192944,
-61608386055612989120, -4358301391935808576, 1995845043302066048, 56288530795526656,
-31330477096792576, 102010241398784, 162473872068608, -2652426567680
\end{Verbatim}
\subsubsection{A209953}
\label{round5:r5entry:A209953}
Arrays on $0,\ldots,2$ have $n+1$ rows and width $w=2$. The number of distinct signed clockwise edge differences in each block lies in $S=\{1,3,4\}$, and the count is multiplied by $\lambda=1$.
The recurrence \eqref{round5:eq:r5arr-recurrence} holds for every $n\ge 4$, with order $d=3$ and component dimensions $(4)$ (total $D=4$). Source: \url{https://oeis.org/A209953}.
The first three values are $(71,567,4521)$.
\[ (c_1,\ldots,c_d)=(6,15,6).\]
The source's empirical/conjectural rational generating function is also proved by the exact identity $QF=P$.
\subsubsection{A209955}
\label{round5:r5entry:A209955}
Arrays on $0,\ldots,2$ have $n+1$ rows and width $w=4$. The number of distinct signed clockwise edge differences in each block lies in $S=\{1,3,4\}$, and the count is multiplied by $\lambda=1$.
The recurrence \eqref{round5:eq:r5arr-recurrence} holds for every $n\ge 15$, with order $d=14$ and component dimensions $(25)$ (total $D=25$). Source: \url{https://oeis.org/A209955}.
The first three values are $(4521,260201,14922177)$.
\begin{Verbatim}[fontsize=\scriptsize,breaklines=true]
48, 720, -8224, -127687, -63722, 3510704, 8214978, -17816696, -46577704, 37523312,
70734912, -30102144, -26296320, 10088448
\end{Verbatim}
\subsubsection{A209956}
\label{round5:r5entry:A209956}
Arrays on $0,\ldots,2$ have $n+1$ rows and width $w=5$. The number of distinct signed clockwise edge differences in each block lies in $S=\{1,3,4\}$, and the count is multiplied by $\lambda=1$.
The recurrence \eqref{round5:eq:r5arr-recurrence} holds for every $n\ge 37$, with order $d=36$ and component dimensions $(70)$ (total $D=70$). Source: \url{https://oeis.org/A209956}.
The first three values are $(36057,5569307,856216289)$.
\begin{Verbatim}[fontsize=\scriptsize,breaklines=true]
188, -2606, -449432, 4278439, 344696470, -1695340318, -106831277262, 260315668742,
14665133182628, -32541316090784, -1037044125842224, 3010738321041520, 39283976226587808,
-157080544206380960, -734266820477791744, 4221150804291177792, 4462953823466462848,
-55438390445568782720, 33810272263704976640, 347725953909047112192,
-562768188874634440704, -933791709840933227520, 2653805933866673033216,
418755395615456186368, -5521527001945241739264, 2418407310629868339200,
5278645092513311916032, -4226822592850201149440, -2085348951772738617344,
2615796490878429691904, 190791165053952327680, -696828367550619320320,
49927158958133346304, 76557678286106263552, -6934923112613937152, -2692851001743900672
\end{Verbatim}
\subsubsection{A210070}
\label{round5:r5entry:A210070}
Arrays on $0,\ldots,3$ have $n+1$ rows and width $w=2$. The number of distinct signed clockwise edge differences in each block lies in $S=\{1,3\}$, and the count is multiplied by $\lambda=1/4$.
The recurrence \eqref{round5:eq:r5arr-recurrence} holds for every $n\ge 7$, with order $d=6$ and component dimensions $(6)$ (total $D=6$). Source: \url{https://oeis.org/A210070}.
The first three values are $(33,313,2908)$.
\[ (c_1,\ldots,c_d)=(9,9,-50,-54,12,16).\]
The source's empirical/conjectural rational generating function is also proved by the exact identity $QF=P$.
\subsubsection{A210071}
\label{round5:r5entry:A210071}
Arrays on $0,\ldots,3$ have $n+1$ rows and width $w=3$. The number of distinct signed clockwise edge differences in each block lies in $S=\{1,3\}$, and the count is multiplied by $\lambda=1/4$.
The recurrence \eqref{round5:eq:r5arr-recurrence} holds for every $n\ge 21$, with order $d=20$ and component dimensions $(20)$ (total $D=20$). Source: \url{https://oeis.org/A210071}.
The first three values are $(313,7653,183628)$.
\begin{Verbatim}[fontsize=\scriptsize,breaklines=true]
40, -410, 162, 13887, -37312, -114069, 441850, 251010, -1865698, 353991, 3283497,
-1672828, -2287997, 1529824, 444317, -334882, -35990, 21286, 648, -288
\end{Verbatim}
\subsubsection{A210072}
\label{round5:r5entry:A210072}
Arrays on $0,\ldots,3$ have $n+1$ rows and width $w=4$. The number of distinct signed clockwise edge differences in each block lies in $S=\{1,3\}$, and the count is multiplied by $\lambda=1/4$.
The recurrence \eqref{round5:eq:r5arr-recurrence} holds for every $n\ge 72$, with order $d=71$ and component dimensions $(72)$ (total $D=72$). Source: \url{https://oeis.org/A210072}.
The first three values are $(2908,183628,11333893)$.
\begin{Verbatim}[fontsize=\scriptsize,breaklines=true]
105, -2719, -15451, 1390943, -7924538, -266634929, 2783228062, 26184293659,
-399963589595, -1390685894189, 33792993132607, 30902632289869, -1890381555307832,
822619990486604, 74406664819656249, -89835422690076184, -2138656368143245036,
3544878797500349305, 46008317343363322887, -87903800402488590784,
-753473801980270591181, 1535436320348069965003, 9505348644817473245840,
-19817177440215843837829, -93126842954944481973745, 194008472781630232252845,
712417629951787189697006, -1464351302320119255078001, -4269039597123343663401254,
8611866713944156678864641, 20062337369297929637737396, -39721870514776593763639623,
-73883618013286155451079859, 144186462856893414613730955, 212594575285455273050900740,
-412162806244228323032035765, -475362727557671730453508543, 926067583198793638831571069,
818943669820992562798006291, -1628407201463016755238157088,
-1073377817747961549292055653, 2225907870727063061101510830,
1050655346169817900693995129, -2343867586300105570589047976,
-746446528111441994989443963, 1879795853747926271266149480, 366426635285816505070213167,
-1132874535089578641767501432, -111323955733240541695109281,
505194638908821997108469349, 12927626382477733740240979, -163887799475213923454810367,
4251111577638305441346089, 37966681387926744734306902, -2249309141628534316063828,
-6153601732703873799844070, 476254714828361557941333, 681177457210232099839986,
-56897891948249201993222, -49910567815506758541560, 4016897111990219118375,
2316946754025069893095, -163269106694177521346, -63879794902532341859,
3536992747961553187, 943996081685851752, -35450092561754998, -6289331583619410,
135657995058528, 14840999586360, -115072392784
\end{Verbatim}
\subsubsection{A210082}
\label{round5:r5entry:A210082}
Arrays have $n+1$ rows, width $w=5$, exactly $4$ colors renamed in row-major first-occurrence order, and $S=\{1,3,4\}$ distinct cell values per block.
The recurrence \eqref{round5:eq:r5arr-recurrence} holds for every $n\ge 54$, with order $d=53$ and component dimensions $(1,16,41,51)$ (total $D=109$). Source: \url{https://oeis.org/A210082}.
The first three values are $(9405,2718363,735082380)$.
\begin{Verbatim}[fontsize=\scriptsize,breaklines=true]
357, -24956, 287902, 26074079, -902939413, 1277468261, 399614425317, -6321060969799,
-24184081408184, 1587694606642173, -12981021635371611, -72434732281423272,
1851336712434285028, -8298795143833957130, -55370910495748475918, 699885760777699195037,
-1271589380489406898905, -16850531904836117165732, 97437529757163175050052,
48421575192093992037435, -1974487710666877857698347, 4494801106206276016732911,
15972527052167522580745893, -85997411174773544395221229, 8070664493704334395329420,
710134757867676419060179075, -1161429886184152718241829591,
-2761558002058205806660337818, 9607464701171282796386424036,
1672650823286895899011692366, -38888570744662356012685589228,
29594770811521649020052740416, 85957722600879503060086236632,
-132835061319571827563642760504, -91415452353353247701405254720,
282171566706458815056468521216, -6031025209778057302947810368,
-346443281245790759531269100096, 139125028729243971214443853952,
252488578502803235293338745600, -172347974003740353484283547648,
-104528726263890830851423499264, 103392729149084501549131203584,
21426153350278051580162213888, -33786660119662982226892468224,
-1224295150346030527568994304, 6066487242636218074798030848,
-214584923867248084083081216, -569047585790917406986076160, 30222671766403263944196096,
24646600389530110604083200, -800959542877582339866624, -383347881721865554624512
\end{Verbatim}
\subsubsection{A210128}
\label{round5:r5entry:A210128}
Arrays have $n+1$ rows, width $w=3$, at most $4$ colors renamed in row-major first-occurrence order, and $S=\{1,3,4\}$ distinct cell values per block.
The recurrence \eqref{round5:eq:r5arr-recurrence} holds for every $n\ge 6$, with order $d=5$ and component dimensions $(1,4,5)$ (total $D=10$). Source: \url{https://oeis.org/A210128}.
The first three values are $(81,2609,84381)$.
\[ (c_1,\ldots,c_d)=(35,-81,-153,492,-292).\]
The source's empirical/conjectural rational generating function is also proved by the exact identity $QF=P$.
\subsubsection{A210149}
\label{round5:r5entry:A210149}
Arrays on $0,\ldots,3$ have $n+1$ rows and width $w=2$. The number of distinct signed clockwise edge differences in each block lies in $S=\{1,3,4\}$, and the count is multiplied by $\lambda=1/4$.
The recurrence \eqref{round5:eq:r5arr-recurrence} holds for every $n\ge 6$, with order $d=5$ and component dimensions $(6)$ (total $D=6$). Source: \url{https://oeis.org/A210149}.
The first three values are $(57,820,11783)$.
\[ (c_1,\ldots,c_d)=(11,46,36,-10,-8).\]
The source's empirical/conjectural rational generating function is also proved by the exact identity $QF=P$.
\subsubsection{A210150}
\label{round5:r5entry:A210150}
Arrays on $0,\ldots,3$ have $n+1$ rows and width $w=3$. The number of distinct signed clockwise edge differences in each block lies in $S=\{1,3,4\}$, and the count is multiplied by $\lambda=1/4$.
The recurrence \eqref{round5:eq:r5arr-recurrence} holds for every $n\ge 14$, with order $d=13$ and component dimensions $(20)$ (total $D=20$). Source: \url{https://oeis.org/A210150}.
The first three values are $(820,42689,2219034)$.
\begin{Verbatim}[fontsize=\scriptsize,breaklines=true]
49, 200, -2256, -3750, 29143, -5802, -84658, 57585, 56482, -40584, -14146, 6952, 1280
\end{Verbatim}
\subsubsection{A210151}
\label{round5:r5entry:A210151}
Arrays on $0,\ldots,3$ have $n+1$ rows and width $w=4$. The number of distinct signed clockwise edge differences in each block lies in $S=\{1,3,4\}$, and the count is multiplied by $\lambda=1/4$.
The recurrence \eqref{round5:eq:r5arr-recurrence} holds for every $n\ge 49$, with order $d=48$ and component dimensions $(72)$ (total $D=72$). Source: \url{https://oeis.org/A210151}.
The first three values are $(11783,2219034,417014920)$.
\begin{Verbatim}[fontsize=\scriptsize,breaklines=true]
145, 8371, -28223, -5099034, -6848727, 1239650435, 930448375, -145737260866,
174912243963, 8264631635949, -22596313142751, -230989466281390, 920787848497955,
3145294614986535, -17388014234115316, -20446231184431219, 183336119039290408,
35747712256243592, -1196660533703535191, 358715517055014311, 5151447312682008123,
-2881286940145422064, -15254072639096425931, 10457263543037282604, 31951721178336409991,
-23277237121895822264, -48167665472082343099, 34598067257613352063,
52718006780308703710, -35619157696064520092, -41957052053167847864,
25829709988559823846, 24165370410041239916, -13259152815213783132, -9956951932632964712,
4797896668438751400, 2878768836324951744, -1208406171336576976, -567099023575525152,
207214883425141632, 72862743480025024, -23352420721108352, -5706693362437376,
1628318996755456, 241536697018368, -62481141760000, -4133487042560, 987100119040
\end{Verbatim}
\subsubsection{A210263}
\label{round5:r5entry:A210263}
Arrays have $n+1$ rows, width $w=5$, exactly $4$ colors renamed in row-major first-occurrence order, and $S=\{3\}$ distinct cell values per block.
The recurrence \eqref{round5:eq:r5arr-recurrence} holds for every $n\ge 51$, with order $d=50$ and component dimensions $(1,16,41,51)$ (total $D=109$). Source: \url{https://oeis.org/A210263}.
The first three values are $(4298,592969,74289956)$.
\begin{Verbatim}[fontsize=\scriptsize,breaklines=true]
208, -11777, 94953, 11282777, -300230858, -2030020415, 154714365717, -835409506883,
-32817201385328, 387387151685882, 3044940024025706, -64835114300366009,
-42786707245576098, 5751830532275106073, -16103656230733623821, -295460333548848141092,
1596987697112378159646, 8716400092711373757801, -76829522478341353618635,
-117025527365549564447848, 2255386434235503948211582, -1090527640994479807979772,
-42924560123670593643529800, 79873007729548579232819824, 530452740825373237454249528,
-1678996015416401623433215328, -3974597159842265998807443728,
20442608956556896189560069024, 12572243416557611196531929600,
-158760447373287871211600674496, 62371746018858818643273251840,
790240531776664873390830228736, -895213241848646797338489708544,
-2390400685449497424892238006272, 4529311964034293254493883850752,
3661856275780121034325608316928, -12436379430730778229284273520640,
-157232863171166993110527508480, 19218418014847258443838877335552,
-8524407983882505977655434674176, -15750453996892525372114665996288,
12632938151949082801824296075264, 5555135058401540564501416378368,
-7628049341872581906805584560128, -8980990683676708801829928960,
1936146568314526563464720154624, -327642732580228949614402732032,
-172830348605676587877088100352, 30197447438230952553989799936,
5737493600858405860068556800
\end{Verbatim}
\subsubsection{A210269}
\label{round5:r5entry:A210269}
Arrays on $0,\ldots,2$ have $n+1$ rows and width $w=2$. The number of distinct signed clockwise edge differences in each block lies in $S=\{2,3\}$, and the count is multiplied by $\lambda=1/2$.
The recurrence \eqref{round5:eq:r5arr-recurrence} holds for every $n\ge 4$, with order $d=3$ and component dimensions $(4)$ (total $D=4$). Source: \url{https://oeis.org/A210269}.
The first three values are $(35,273,2130)$.
\[ (c_1,\ldots,c_d)=(8,0,-12).\]
The source's empirical/conjectural rational generating function is also proved by the exact identity $QF=P$.
\subsubsection{A210270}
\label{round5:r5entry:A210270}
Arrays on $0,\ldots,2$ have $n+1$ rows and width $w=3$. The number of distinct signed clockwise edge differences in each block lies in $S=\{2,3\}$, and the count is multiplied by $\lambda=1/2$.
The recurrence \eqref{round5:eq:r5arr-recurrence} holds for every $n\ge 10$, with order $d=9$ and component dimensions $(10)$ (total $D=10$). Source: \url{https://oeis.org/A210270}.
The first three values are $(273,5558,113248)$.
\begin{Verbatim}[fontsize=\scriptsize,breaklines=true]
20, 26, -364, -218, 1203, -418, -394, 216, -20
\end{Verbatim}
\subsubsection{A210271}
\label{round5:r5entry:A210271}
Arrays on $0,\ldots,2$ have $n+1$ rows and width $w=4$. The number of distinct signed clockwise edge differences in each block lies in $S=\{2,3\}$, and the count is multiplied by $\lambda=1/2$.
The recurrence \eqref{round5:eq:r5arr-recurrence} holds for every $n\ge 26$, with order $d=25$ and component dimensions $(25)$ (total $D=25$). Source: \url{https://oeis.org/A210271}.
The first three values are $(2130,113248,6028871)$.
\begin{Verbatim}[fontsize=\scriptsize,breaklines=true]
48, 480, -9807, -58253, 657082, 2138431, -18778839, -19795420, 235514744, -62769938,
-1259927930, 1276469157, 2637995603, -4298139072, -1484220465, 5231313990, -949603112,
-2426962356, 931229528, 400609008, -188316850, -12522776, 5412968, -237208, -448
\end{Verbatim}
\subsubsection{A210272}
\label{round5:r5entry:A210272}
Arrays on $0,\ldots,2$ have $n+1$ rows and width $w=5$. The number of distinct signed clockwise edge differences in each block lies in $S=\{2,3\}$, and the count is multiplied by $\lambda=1/2$.
The recurrence \eqref{round5:eq:r5arr-recurrence} holds for every $n\ge 70$, with order $d=69$ and component dimensions $(70)$ (total $D=70$). Source: \url{https://oeis.org/A210272}.
The first three values are $(16620,2307956,321050788)$.
\begin{Verbatim}[fontsize=\scriptsize,breaklines=true]
142, 1390, -255743, -110496, 169565270, -665183207, -52878902879, 395903765675,
8133200677492, -92576230139190, -542663156454254, 10497292789475974, 2631885281757262,
-646106145808214387, 1826441384938919701, 22362684125873346283, -123476318341203413496,
-395337011643732047682, 4142092565872768583256, 605782653333712666179,
-83518787705085031195627, 136634888376530178952716, 1040134849078032981074015,
-3402074914964904960822357, -7239465059314247652144549, 44377737415311976751838594,
9840741933860801981979443, -361590798269938925507784670, 321818684009046992762532501,
1883576717146322164819296451, -3479890138790632205298850558,
-5878393363654009470130124838, 19078667220689449803911834552,
7392854341921657854446418742, -65742784478773744652167403362,
19627223671361316822942496555, 149839903328239191224944472173,
-117630629870435594695245179193, -226767728498105405443631444233,
284875215327416958030542435899, 219032259251618705695478104114,
-424657211702672231000865863606, -114298568153964486240525208337,
426463674134772011008406591220, -1465111296208263054958667449,
-298361702108034534055790446739, 50042360589416264853656878477,
147499557244192463582697992372, -39645514435604456524417564586,
-51787606560690998053903082159, 16872384336675821704842993456,
12907214080006987000746559571, -4471894138938904280092501035,
-2270649243797947159981070488, 757433256288002048852712547, 277898940359626427914691228,
-80741549607760164976554158, -22888243320669979845951670, 5181436657383846894790730,
1185029940722673827637793, -185251468573962367963557, -33904216065585670936623,
3275032115968437164460, 413089312994017424066, -25348341919215966472,
-1274153174398735370, 12493913728073620, 443471547890804, -2001148574736
\end{Verbatim}
\subsubsection{A210305}
\label{round5:r5entry:A210305}
Arrays have $n+1$ rows, width $w=5$, exactly $4$ colors renamed in row-major first-occurrence order, and $S=\{1,3\}$ distinct cell values per block.
The recurrence \eqref{round5:eq:r5arr-recurrence} holds for every $n\ge 54$, with order $d=53$ and component dimensions $(1,16,41,51)$ (total $D=109$). Source: \url{https://oeis.org/A210305}.
The first three values are $(4553,655743,85724987)$.
\begin{Verbatim}[fontsize=\scriptsize,breaklines=true]
260, -23288, 929838, -12001245, -334728110, 13677246260, -97982783232, -3038968483483,
61195800410360, 21617308239752, -10672008917494402, 77526891307988953,
740183820173087714, -11993209513343105860, 1757788347100408132, 831215814617340342320,
-3626826186455581358000, -28284540559276107423616, 255592495497158190025072,
251472921681558655747376, -8953225998033769175695072, 16464543868595186192082752,
176604280245762880053747200, -734118266281625015141049088,
-1732590776838801409725204224, 14878214282327021473824928768,
-1632048876867875141848579072, -176456313876275191138951528448,
275325919515911232333562363904, 1236150931517543249393523490816,
-3718761489967820580413373988864, -4219444988326081346462792581120,
26532685864866441003684492673024, -4031028009268447870327055187968,
-113394942239733332464880215654400, 106471669450351021801461776384000,
288206551250917417355764897415168, -481709443052870836912613364858880,
-384039040102328706184157404856320, 1142268907564109813860532685373440,
91027236998423324239131310030848, -1580464812314641136289277739532288,
473969325227938496307658542284800, 1268980715592062152644700409954304,
-689711550621859287710203188871168, -548172194555268379337156183719936,
403997167828185646769666181300224, 106908215762795739559038382243840,
-102992983528709537221967712616448, -7280608274553819303010390507520,
10221635153981636051140175986688, 41096544975325809960739143680,
-236382646173102461890791997440
\end{Verbatim}
\subsubsection{A210328}
\label{round5:r5entry:A210328}
Arrays on $0,\ldots,3$ have $n+1$ rows and width $w=2$. The number of distinct signed clockwise edge differences in each block lies in $S=\{2,3\}$, and the count is multiplied by $\lambda=1/4$.
The recurrence \eqref{round5:eq:r5arr-recurrence} holds for every $n\ge 7$, with order $d=6$ and component dimensions $(6)$ (total $D=6$). Source: \url{https://oeis.org/A210328}.
The first three values are $(39,406,4200)$.
\[ (c_1,\ldots,c_d)=(12,-7,-108,22,258,124).\]
The source's empirical/conjectural rational generating function is also proved by the exact identity $QF=P$.
\subsubsection{A210329}
\label{round5:r5entry:A210329}
Arrays on $0,\ldots,3$ have $n+1$ rows and width $w=3$. The number of distinct signed clockwise edge differences in each block lies in $S=\{2,3\}$, and the count is multiplied by $\lambda=1/4$.
The recurrence \eqref{round5:eq:r5arr-recurrence} holds for every $n\ge 21$, with order $d=20$ and component dimensions $(20)$ (total $D=20$). Source: \url{https://oeis.org/A210329}.
The first three values are $(406,11877,344659)$.
\begin{Verbatim}[fontsize=\scriptsize,breaklines=true]
47, -547, -255, 35216, -129394, -465060, 2920628, 320982, -20531472, 17232541, 55329403,
-75077657, -49764339, 111287472, -12747232, -52250238, 26800488, -2501204, -451078,
-5036
\end{Verbatim}
\subsubsection{A210330}
\label{round5:r5entry:A210330}
Arrays on $0,\ldots,3$ have $n+1$ rows and width $w=4$. The number of distinct signed clockwise edge differences in each block lies in $S=\{2,3\}$, and the count is multiplied by $\lambda=1/4$.
The recurrence \eqref{round5:eq:r5arr-recurrence} holds for every $n\ge 73$, with order $d=72$ and component dimensions $(72)$ (total $D=72$). Source: \url{https://oeis.org/A210330}.
The first three values are $(4200,344659,27838382)$.
\begin{Verbatim}[fontsize=\scriptsize,breaklines=true]
140, -5048, -23365, 4584810, -58220704, -1208071283, 30838348016, 46781002450,
-6610272669799, 33388421310781, 724931448827996, -7330140469767465, -38529835529169361,
750770095780529564, 104547760010327276, -45822745228789445629, 124886909490991773867,
1772584705154196067963, -9188006643225676802179, -42487551119663686790240,
369441396567930016447963, 493046341331884330242325, -9840007964400110991140626,
4577246401662065605851681, 184119139789121757095000706, -324697880502964050967877701,
-2457999961706566468164419681, 7357259533993676894189678070,
22913914256580375682856411562, -105938294940954623488202172345,
-133797398187861506534787512518, 1088287205984201969751840912702,
214149173009648038478719960791, -8295244297697700252477102952711,
4551271399893045426167401874520, 47492148653867277925166063188222,
-55467328086449248198445554895418, -203044727836439247258180379708115,
364953161964814879802809348845513, 629000600003594042909200245131002,
-1656973614766749241452655210247331, -1286419122570055663502048783109264,
5521435181568923625231432851494067, 1085213406458212188709424799285359,
-13763133278895135251639489780210047, 2932860737587723748361214175457686,
25672766856385812407325333095611470, -13928357412518894891097071827999092,
-35342311049626777492661042089369341, 30149006440809583122467985587701935,
34731032869049446788157006292902188, -42409921614806508539702333650807347,
-22523081462687078837390075946641420, 41474885258750816250324684644171213,
7300575902610538123306992258616268, -28518904133762111651213915708675288,
1573935469653866497475056486719028, 13609710325408504350514230640881256,
-3000959034483091478767274705019179, -4366242148634467525612510604850896,
1533561432404703029191784717830830, 891878386621906486208161657560730,
-421653290626187252410928504168962, -105776373439489265946108677516538,
66067934140142616805443885636242, 6007572598581250574435738268266,
-5638660751066341761268282571708, -71565691223184653218166474032,
231250260842273331695302457980, -3237179075004335463427133468,
-3230599808041743511640862256, -41722134296876122805175792
\end{Verbatim}
\subsubsection{A210348}
\label{round5:r5entry:A210348}
Arrays on $0,\ldots,3$ have $n+1$ rows and width $w=2$. The number of distinct signed clockwise edge differences in each block lies in $S=\{1,2,3\}$, and the count is multiplied by $\lambda=1/4$.
The recurrence \eqref{round5:eq:r5arr-recurrence} holds for every $n\ge 6$, with order $d=5$ and component dimensions $(6)$ (total $D=6$). Source: \url{https://oeis.org/A210348}.
The first three values are $(40,431,4615)$.
\[ (c_1,\ldots,c_d)=(14,-30,-66,104,136).\]
The source's empirical/conjectural rational generating function is also proved by the exact identity $QF=P$.
\subsubsection{A210349}
\label{round5:r5entry:A210349}
Arrays on $0,\ldots,3$ have $n+1$ rows and width $w=3$. The number of distinct signed clockwise edge differences in each block lies in $S=\{1,2,3\}$, and the count is multiplied by $\lambda=1/4$.
The recurrence \eqref{round5:eq:r5arr-recurrence} holds for every $n\ge 21$, with order $d=20$ and component dimensions $(20)$ (total $D=20$). Source: \url{https://oeis.org/A210349}.
The first three values are $(431,13566,423694)$.
\begin{Verbatim}[fontsize=\scriptsize,breaklines=true]
55, -890, 4374, 14813, -179200, 198626, 1718220, -3858631, -4143905, 14290154, 1274736,
-19132266, 4460692, 9184880, -3036040, -1419448, 517280, 1792, -4096, -128
\end{Verbatim}
\subsubsection{A210350}
\label{round5:r5entry:A210350}
Arrays on $0,\ldots,3$ have $n+1$ rows and width $w=4$. The number of distinct signed clockwise edge differences in each block lies in $S=\{1,2,3\}$, and the count is multiplied by $\lambda=1/4$.
The recurrence \eqref{round5:eq:r5arr-recurrence} holds for every $n\ge 73$, with order $d=72$ and component dimensions $(72)$ (total $D=72$). Source: \url{https://oeis.org/A210350}.
The first three values are $(4615,423694,38415623)$.
\begin{Verbatim}[fontsize=\scriptsize,breaklines=true]
191, -12666, 362418, -2471771, -110401217, 2621836067, -5950956298, -457658156996,
4931495900625, 22595991223312, -660371000230224, 1576600613573351, 40726756485462035,
-261581043509847118, -1180848375506028374, 15214038374148596375, 2740239137031201481,
-500647751869384591185, 1003334958924814164002, 10265837091234582479490,
-38885756005297761240461, -130612308191256893882321, 812041422755207796165679,
855729427631059422140114, -11175171560989023064509960, 1889581061866764605910891,
107928399415545111392724496, -107784289315410029448855934, -748128702571716925965136011,
1270446999046412386430183438, 3712144693238748687188701435,
-9124381122510874786617773257, -12697180476385310458871841925,
45745095998148294054172238533, 25879826918923108106955540193,
-167256385704846228147053848315, -5522734439691358725661862744,
452571003883272632598446053394, -162676959471579995234071136543,
-904948271559219280422768304046, 626035293669715175604860994661,
1314278201887400107774175032607, -1355020035697843080612169545698,
-1325873311191946166966833531231, 1973302357536165503278014656818,
819196798824108568727701573827, -2014556097985763744645141103153,
-141292825105011889583185641794, 1444488136170597993772231897282,
-247655714150765676981442117027, -712507042523512207592090491261,
261206533429284400319401980540, 230922165671337271495130560074,
-130296618009089433855843407234, -44578783099713406658548437066,
38556121946149222868083061076, 3691991189510051391638768180,
-6995991014442444122356572748, 257789834532082975263071992, 772349630438652610156030230,
-86947033325326193209399996, -50550945071670308168118284, 7804233987414798680212112,
1880536771253356175158040, -321393299282915797600248, -37501945544789391775192,
6220742886607802957056, 344840850367579506624, -53336079334554500352,
-987610693682276352, 164858395676078080, -1278984298302976
\end{Verbatim}
\subsubsection{A210383}
\label{round5:r5entry:A210383}
Arrays on $0,\ldots,2$ have $n+1$ rows and width $w=2$. The number of distinct signed clockwise edge differences in each block lies in $S=\{1,2,3\}$, and the count is multiplied by $\lambda=1$.
The recurrence \eqref{round5:eq:r5arr-recurrence} holds for every $n\ge 4$, with order $d=3$ and component dimensions $(4)$ (total $D=4$). Source: \url{https://oeis.org/A210383}.
The first three values are $(73,597,4881)$.
\[ (c_1,\ldots,c_d)=(9,-6,-6).\]
The source's empirical/conjectural rational generating function is also proved by the exact identity $QF=P$.
\subsubsection{A210384}
\label{round5:r5entry:A210384}
Arrays on $0,\ldots,2$ have $n+1$ rows and width $w=3$. The number of distinct signed clockwise edge differences in each block lies in $S=\{1,2,3\}$, and the count is multiplied by $\lambda=1$.
The recurrence \eqref{round5:eq:r5arr-recurrence} holds for every $n\ge 7$, with order $d=6$ and component dimensions $(10)$ (total $D=10$). Source: \url{https://oeis.org/A210384}.
The first three values are $(597,13399,300653)$.
\[ (c_1,\ldots,c_d)=(26,-76,-93,144,66,-32).\]
The source's empirical/conjectural rational generating function is also proved by the exact identity $QF=P$.
\subsubsection{A210386}
\label{round5:r5entry:A210386}
Arrays on $0,\ldots,2$ have $n+1$ rows and width $w=5$. The number of distinct signed clockwise edge differences in each block lies in $S=\{1,2,3\}$, and the count is multiplied by $\lambda=1$.
The recurrence \eqref{round5:eq:r5arr-recurrence} holds for every $n\ge 43$, with order $d=42$ and component dimensions $(70)$ (total $D=70$). Source: \url{https://oeis.org/A210386}.
The first three values are $(39909,6747079,1140445889)$.
\begin{Verbatim}[fontsize=\scriptsize,breaklines=true]
223, -9461, 38458, 3549663, -32588685, -508261230, 5706283402, 35105920095,
-470735135093, -1154545872838, 21185965166884, 11225843446480, -544537281641561,
302744827612151, 8161346038514766, -9528641858555588, -73104114337370148,
109184395751663568, 406815111946067415, -670442664892782833, -1455165639539505801,
2462053942187126418, 3420041227155372378, -5685667280406325350, -5317279202551076247,
8464088502302311604, 5417981379313428776, -8206507857103512021, -3524350348222019545,
5179718539525222811, 1395392552338466122, -2099353043062207344, -304790403886902836,
529249988183090028, 26637678526607096, -78134322897382024, 1358756843570944,
6043523752505920, -329711911635072, -200633167870464, 9138104982528, 2395048366080
\end{Verbatim}


\clearpage
\section{Reproduction and verification}
\label{round5:sec:r5reproduce}
\begingroup\small
The archive includes this standalone manuscript, the numbered ledger,
four proof fragments, four main verifiers, an independent array checker,
complete transfer matrices, exact algebraic certificates, source-scope
audits and cross-review reports. Full scraped webpages are omitted.

\subsection{Running the offline checks}
Use Python 3.10 or later with SymPy and mpmath. From the extracted
directory run:
\begin{Verbatim}[fontsize=\small]
python3 r5_run_all.py
\end{Verbatim}
The runner checks all 100 unique A-numbers against the preceding 408,
checks the outcome counts and proof labels, then runs all five programs.
It writes \texttt{oeis\_hunt\_round5\_verification\_summary.json}.
No network connection is required.

\begin{center}\small
\begin{tabularx}{\textwidth}{@{}lX@{}}
\toprule
Program&Verification scope\\
\midrule
\texttt{r5\_root\_verify.py}&Thirteen formal-series and theta packages:
exact identities and supplementary coefficient and numerical checks.\\
\texttt{r5\_arith\_verify.py}&Twelve arithmetic packages: exact integer
counterexamples, rational enclosures and bounded consistency checks
supporting the separate written proofs.\\
\texttt{r5\_gf\_verify.py}&Thirty exact all-index certificates, all
recurrence boundary terms, 480 independently expanded initial terms,
and the five retrospective Fischer operator comparisons.\\
\texttt{r5\_arrays\_verify.py}&Forty-five recurrences, all 56 complete
matrices rebuilt from their definitions, and 17 source generating
functions proved without additional counts.\\
\texttt{r5\_arrays\_root\_verify.py}&Ninety independent initial values
across all 45 arrays, using full rows and a separate column dynamic
program without importing the quotient-matrix engine.\\
\bottomrule
\end{tabularx}
\end{center}

\subsection{Finite proof certificates and bounded checks}
For an array recurrence, the defining count is a linear combination of
$m^{\mathsf T}T^n\mathbf1$. The complete identity
$TR(T)\mathbf1=0$ remains true after multiplication by every
nonnegative power of $T$, proving the recurrence on the stated infinite
range. The largest component has 574 states, and the largest recurrence
order is 72. Every coordinate is checked with exact integer arithmetic.
An initial prefix by itself would not imply the recurrence.

For the generating functions, the selected recurrence is translated
to a differential operator. Its exact action in a rational radical
basis gives the displayed polynomial boundary term, whose degree
determines the valid starting index. A nonzero irrational component
in the \oeis{A129366} residual rules out every eventual tail of the
original recurrence. The repaired recurrence has an exact certificate.

The remaining universal conclusions follow from the written arguments.
Their bounded computations check transcription and indexing. In
particular, the 80-digit theta checks are supplementary; the exact
evaluations follow from the cited classical identities and stated
generating functions. No numerical match is used as the proof of an
infinite identity or limit.

\subsection{Rebuilding, source records and integrity}
The delivered \texttt{oeis\_hunt\_round5.tex} contains all arguments
directly and needs no external \verb|\input| files. Compile it twice
with pdfLaTeX and the standard packages listed in its preamble.
To rebuild the standalone TEX, catalogue and audit index, run:
\begin{Verbatim}[fontsize=\small]
python3 r5_build_dossier.py
\end{Verbatim}
Optional certificate reconstruction commands and dependencies are
listed in the README. \texttt{manifest.json} gives SHA-256 digests of
the delivered files. Running the verifiers regenerates reports and
timing fields, so those report digests can then change.

The metadata retain source attribution, indexing, exclusions and audit
limits. The five retrospective corrections remain explicit. Project
selection counts do not certify first discovery or OEIS editorial closure.
\endgroup
% ===== END ORIGINAL ROUND 5 =====

% ===== BEGIN ORIGINAL ROUND 6 =====
\clearpage
\setcounter{section}{0}\setcounter{equation}{0}\setcounter{theorem}{0}
\setcounter{table}{0}\setcounter{figure}{0}
\normalsize\renewcommand{\arraystretch}{1}\renewcommand{\baselinestretch}{1}\selectfont
\part[Round 6: Formal series, Bernoulli arithmetic, valuation formulas, and finite states]{Formal series, Bernoulli arithmetic, valuation formulas, and finite states}
\label{combined:round6}
\rhead{\small Round 6 / KEIO}
\begingroup\small
\section*{Round overview}
\addcontentsline{toc}{section}{Round 6 overview and scope}
This sixth-round dossier treats 100 selected assertion packages on
100 distinct OEIS entries, disjoint from the 508 A-numbers selected in
the preceding five rounds. It proves 94 original selected claims,
supplies and proves the valid starting range of five conjectural
recurrences, and corrects and proves one periodicity starting index.
Multiple assertions or generating-function consequences at one entry
count only once.

The four parts contain 18 packages on formal series and theta values,
11 on arithmetic and exact congruences, 30 on generating functions,
and 41 on constrained arrays. The formal-series proofs work over
integer coefficient rings, including divided-power series where
appropriate. The arithmetic arguments retain their precise source
scope. The generating-function proofs use exact differential identities
with boundary polynomials. The array proofs use complete integer
transfer matrices and vector annihilator certificates.

\subsection*{What the count means}
Every catalogue row names one selected assertion package on one
distinct A-number. Different members of a family may follow from
one general lemma; the count does not mean 100 unrelated methods.
The consulted primary statements were marked conjectural or empirical.
Established generating functions, different previously known recurrences,
and classical identities are credited inputs.

The selected scope can be narrower than all questions on an entry.
For example, the packages for \oeis{A396805} and \oeis{A396806}
prove parity only, and leave the other modular conjectures on their
pages outside scope. The mod-4 package for \oeis{A389537} does not
claim its separate sign-run observation. Such limits are recorded
in the individual arguments and the source ledger.

The number 508 is a ledger of previously selected A-numbers, not a
certified count of worldwide-new discoveries. The five Round IV
priority corrections disclosed in Round V remain in force and all
those IDs remain excluded. Known prior results found during this
round's audit were excluded or explicitly credited as inputs.

\subsection*{Six qualifications and corrections}
For five recurrences, the source gives no starting index. We show an
earlier admissible exception and prove the exact valid tail; we do
not attribute an explicitly false starting-index assertion to the
source author. The residual is the left side of the displayed source
recurrence, in its stated normalization.
\begin{center}
\begin{tabular}{@{}lrrl@{}}
\toprule
Entry&Exceptional index&Residual&Proved range\\
\midrule
\oeis{A102880}&4&3&$n\ge5$\\
\oeis{A121320}&3&1&$n\ge4$\\
\oeis{A141351}&3&2&$n\ge4$\\
\oeis{A278023}&4&156&$n\ge5$\\
\oeis{A385252}&5&$-6912$&$n\ge6$\\
\bottomrule
\end{tabular}
\end{center}
For \oeis{A396798}, the stated mod-8 pattern $(1,1,5,5)$ starts at
$n=2$, with $a(1)=1$. Starting it at $n=1$, as printed, would predict
$a(3)\equiv5$, whereas $a(3)=9\equiv1\pmod8$. The corrected pattern
and all seven other source iterate-congruence statements are proved
together in its single package.

\subsection*{Evidence and limits}
Primary OEIS entries, relevant linked work and targeted literature
searches were consulted on 20 October 2026; retrieved pages may be
cached. The retained OEIS Open comparison covers Appendix A only:
100 costliest claims on 97 distinct A-numbers, with no current overlap.
It is not a comparison with every result in that project or a complete
historical priority search.

The written arguments and exact certificates support the stated
mathematical conclusions. No worldwide-first-proof or OEIS editorial
acceptance claim is made, and no OEIS edit or submission was performed.
Cross-reviews are by separate agents within the same AI-assisted
workflow; they are not external peer review or proof-assistant
formalization.
\endgroup

\clearpage
\section{Catalogue of the 100 selected entries}
\label{round6:sec:r6catalogue}
\textbf{P}: original selected claim proved.
\textbf{B}: recurrence boundary qualified and proved.
\textbf{I}: period starting index corrected and proved.
Page references lead to the individual argument or complete recurrence
signature. The JSON ledger preserves exact scope, indexing and attribution.

\small
\begin{longtable}{@{}r p{21mm} p{103mm} c r@{}}
\toprule
No.&Entry&Selected assertion package&Result&Page\\
\midrule\endfirsthead
\toprule No.&Entry&Selected assertion package&Result&Page\\
\midrule\endhead
\midrule\multicolumn{5}{r}{\emph{Continued on the next page}}\\\endfoot
\bottomrule\endlastfoot
\multicolumn{5}{@{}l}{\rule{0pt}{4ex}\textbf{Formal series and theta values}}\\
1&\oeis{A393856}&Every term is 1 modulo 5; odd iff the positive index is a power of 2&P&\pageref{round6:r6entry:A393856}\\
2&\oeis{A393857}&Every term is 1 modulo 6&P&\pageref{round6:r6entry:A393857}\\
3&\oeis{A393858}&Every term is 1 modulo 7; odd iff the positive index is a power of 2&P&\pageref{round6:r6entry:A393858}\\
4&\oeis{A393859}&Every term is 1 modulo 8&P&\pageref{round6:r6entry:A393859}\\
5&\oeis{A396793}&Every coefficient after the first is divisible by 3&P&\pageref{round6:r6entry:A396793}\\
6&\oeis{A396794}&Every coefficient after the first is divisible by 4&P&\pageref{round6:r6entry:A396794}\\
7&\oeis{A396795}&Every coefficient after the first is divisible by 5&P&\pageref{round6:r6entry:A396795}\\
8&\oeis{A393170}&Complete parity and residue-two classification modulo four&P&\pageref{round6:r6entry:A393170}\\
9&\oeis{A393173}&Complete parity and residue-two classification modulo four&P&\pageref{round6:r6entry:A393173}\\
10&\oeis{A393171}&Complete residue-four support modulo eight&P&\pageref{round6:r6entry:A393171}\\
11&\oeis{A389537}&Residue-two coefficients are exactly at three times powers of two&P&\pageref{round6:r6entry:A389537}\\
12&\oeis{A396798}&Correct the first period start and prove all eight iterate congruences&I&\pageref{round6:r6entry:A396798}\\
13&\oeis{A396807}&Every integer iterate has coefficients k\textasciicircum{}(n-1) modulo ten&P&\pageref{round6:r6entry:A396807}\\
14&\oeis{A396805}&The exponential coefficient is odd exactly at odd positive indices&P&\pageref{round6:r6entry:A396805}\\
15&\oeis{A396806}&The exponential coefficient is odd exactly at odd positive indices&P&\pageref{round6:r6entry:A396806}\\
16&\oeis{A096920}&The empirical exponential-weighted theta sum&P&\pageref{round6:r6entry:A096920}\\
17&\oeis{A083365}&The empirical exponential-weighted theta sum&P&\pageref{round6:r6entry:A083365}\\
18&\oeis{A273228}&The empirical exponential-weighted theta sum&P&\pageref{round6:r6entry:A273228}\\
\multicolumn{5}{@{}l}{\rule{0pt}{4ex}\textbf{Arithmetic and exact congruences}}\\
19&\oeis{A061002}&Harmonic-number numerator and quotient-denominator identity&P&\pageref{round6:r6entry:A061002}\\
20&\oeis{A063956}&The conjectured product identity for simple prime factors&P&\pageref{round6:r6entry:A063956}\\
21&\oeis{A068991}&Squarefree members divide the indicated Bernoulli denominator&P&\pageref{round6:r6entry:A068991}\\
22&\oeis{A069208}&The general multiplicative-function convolution identity&P&\pageref{round6:r6entry:A069208}\\
23&\oeis{A116184}&The harmonic-numerator progression and prime-residue restriction&P&\pageref{round6:r6entry:A116184}\\
24&\oeis{A133907}&Equivalence of the binomial and power-sum least-prime rules&P&\pageref{round6:r6entry:A133907}\\
25&\oeis{A262735}&The selected trinomial-triangle diagonal-sum identity&P&\pageref{round6:r6entry:A262735}\\
26&\oeis{A299406}&The Liouville product identity for the Dirichlet-series coefficients&P&\pageref{round6:r6entry:A299406}\\
27&\oeis{A376484}&Integrality and divisibility of the trigonometric power sum&P&\pageref{round6:r6entry:A376484}\\
28&\oeis{A377278}&The reciprocal-row-sum determinant quotient&P&\pageref{round6:r6entry:A377278}\\
29&\oeis{A381355}&Divisibility by the prime at the coefficient index&P&\pageref{round6:r6entry:A381355}\\
\multicolumn{5}{@{}l}{\rule{0pt}{4ex}\textbf{Generating-function certificates}}\\
30&\oeis{A025268}&Conjectured polynomial recurrence of order 5&P&\pageref{round6:r6entry:A025268}\\
31&\oeis{A025269}&Conjectured polynomial recurrence of order 4&P&\pageref{round6:r6entry:A025269}\\
32&\oeis{A025270}&Conjectured polynomial recurrence of order 4&P&\pageref{round6:r6entry:A025270}\\
33&\oeis{A025275}&Conjectured polynomial recurrence of order 4&P&\pageref{round6:r6entry:A025275}\\
34&\oeis{A026030}&Conjectured polynomial recurrence of order 4&P&\pageref{round6:r6entry:A026030}\\
35&\oeis{A026571}&Conjectured polynomial recurrence of order 5&P&\pageref{round6:r6entry:A026571}\\
36&\oeis{A026743}&Conjectured polynomial recurrence of order 8&P&\pageref{round6:r6entry:A026743}\\
37&\oeis{A052183}&Conjectured polynomial recurrence of order 2&P&\pageref{round6:r6entry:A052183}\\
38&\oeis{A095981}&Conjectured polynomial recurrence of order 7&P&\pageref{round6:r6entry:A095981}\\
39&\oeis{A102880}&Conjectured polynomial recurrence of order 4; exact starting index supplied&B&\pageref{round6:r6entry:A102880}\\
40&\oeis{A105864}&Conjectured polynomial recurrence of order 4&P&\pageref{round6:r6entry:A105864}\\
41&\oeis{A114584}&Conjectured polynomial recurrence of order 7&P&\pageref{round6:r6entry:A114584}\\
42&\oeis{A118093}&Conjectured polynomial recurrence of order 3&P&\pageref{round6:r6entry:A118093}\\
43&\oeis{A121320}&Conjectured polynomial recurrence of order 3; exact starting index supplied&B&\pageref{round6:r6entry:A121320}\\
44&\oeis{A141351}&Conjectured polynomial recurrence of order 3; exact starting index supplied&B&\pageref{round6:r6entry:A141351}\\
45&\oeis{A157003}&Conjectured polynomial recurrence of order 6&P&\pageref{round6:r6entry:A157003}\\
46&\oeis{A159772}&Conjectured polynomial recurrence of order 6&P&\pageref{round6:r6entry:A159772}\\
47&\oeis{A160823}&Conjectured polynomial recurrence of order 5&P&\pageref{round6:r6entry:A160823}\\
48&\oeis{A162533}&Conjectured polynomial recurrence of order 5&P&\pageref{round6:r6entry:A162533}\\
49&\oeis{A174107}&Conjectured polynomial recurrence of order 4&P&\pageref{round6:r6entry:A174107}\\
50&\oeis{A174808}&Conjectured polynomial recurrence of order 5&P&\pageref{round6:r6entry:A174808}\\
51&\oeis{A190169}&Conjectured polynomial recurrence of order 8&P&\pageref{round6:r6entry:A190169}\\
52&\oeis{A208355}&Conjectured polynomial recurrence of order 2&P&\pageref{round6:r6entry:A208355}\\
53&\oeis{A228178}&Conjectured polynomial recurrence of order 4&P&\pageref{round6:r6entry:A228178}\\
54&\oeis{A243022}&Conjectured polynomial recurrence of order 5&P&\pageref{round6:r6entry:A243022}\\
55&\oeis{A243255}&Conjectured polynomial recurrence of order 4&P&\pageref{round6:r6entry:A243255}\\
56&\oeis{A254314}&Conjectured polynomial recurrence of order 5&P&\pageref{round6:r6entry:A254314}\\
57&\oeis{A257388}&Conjectured polynomial recurrence of order 3&P&\pageref{round6:r6entry:A257388}\\
58&\oeis{A278023}&Conjectured polynomial recurrence of order 4; exact starting index supplied&B&\pageref{round6:r6entry:A278023}\\
59&\oeis{A385252}&Conjectured polynomial recurrence of order 5; exact starting index supplied&B&\pageref{round6:r6entry:A385252}\\
\multicolumn{5}{@{}l}{\rule{0pt}{4ex}\textbf{Finite-state array certificates}}\\
60&\oeis{A183586}&Empirical array recurrence of order 2&P&\pageref{round6:r6entry:A183586}\\
61&\oeis{A183587}&Empirical array recurrence of order 4&P&\pageref{round6:r6entry:A183587}\\
62&\oeis{A183589}&Empirical array recurrence of order 31&P&\pageref{round6:r6entry:A183589}\\
63&\oeis{A183596}&Empirical array recurrence of order 4&P&\pageref{round6:r6entry:A183596}\\
64&\oeis{A210176}&Empirical array recurrence of order 4&P&\pageref{round6:r6entry:A210176}\\
65&\oeis{A210177}&Empirical array recurrence of order 8&P&\pageref{round6:r6entry:A210177}\\
66&\oeis{A210178}&Empirical array recurrence of order 16&P&\pageref{round6:r6entry:A210178}\\
67&\oeis{A210406}&Empirical array recurrence of order 2&P&\pageref{round6:r6entry:A210406}\\
68&\oeis{A210407}&Empirical array recurrence of order 4&P&\pageref{round6:r6entry:A210407}\\
69&\oeis{A210408}&Empirical array recurrence of order 9&P&\pageref{round6:r6entry:A210408}\\
70&\oeis{A210409}&Empirical array recurrence of order 31&P&\pageref{round6:r6entry:A210409}\\
71&\oeis{A233921}&Empirical array recurrence of order 6&P&\pageref{round6:r6entry:A233921}\\
72&\oeis{A233982}&Empirical array recurrence of order 3&P&\pageref{round6:r6entry:A233982}\\
73&\oeis{A233983}&Empirical array recurrence of order 13&P&\pageref{round6:r6entry:A233983}\\
74&\oeis{A233984}&Empirical array recurrence of order 25&P&\pageref{round6:r6entry:A233984}\\
75&\oeis{A234162}&Empirical array recurrence of order 4&P&\pageref{round6:r6entry:A234162}\\
76&\oeis{A234163}&Empirical array recurrence of order 17&P&\pageref{round6:r6entry:A234163}\\
77&\oeis{A234220}&Empirical array recurrence of order 3&P&\pageref{round6:r6entry:A234220}\\
78&\oeis{A234221}&Empirical array recurrence of order 6&P&\pageref{round6:r6entry:A234221}\\
79&\oeis{A234222}&Empirical array recurrence of order 11&P&\pageref{round6:r6entry:A234222}\\
80&\oeis{A234557}&Empirical array recurrence of order 5&P&\pageref{round6:r6entry:A234557}\\
81&\oeis{A234558}&Empirical array recurrence of order 14&P&\pageref{round6:r6entry:A234558}\\
82&\oeis{A234559}&Empirical array recurrence of order 14&P&\pageref{round6:r6entry:A234559}\\
83&\oeis{A234762}&Empirical array recurrence of order 5&P&\pageref{round6:r6entry:A234762}\\
84&\oeis{A234763}&Empirical array recurrence of order 13&P&\pageref{round6:r6entry:A234763}\\
85&\oeis{A234764}&Empirical array recurrence of order 27&P&\pageref{round6:r6entry:A234764}\\
86&\oeis{A234875}&Empirical array recurrence of order 4&P&\pageref{round6:r6entry:A234875}\\
87&\oeis{A234876}&Empirical array recurrence of order 9&P&\pageref{round6:r6entry:A234876}\\
88&\oeis{A234877}&Empirical array recurrence of order 11&P&\pageref{round6:r6entry:A234877}\\
89&\oeis{A235011}&Empirical array recurrence of order 5&P&\pageref{round6:r6entry:A235011}\\
90&\oeis{A235012}&Empirical array recurrence of order 8&P&\pageref{round6:r6entry:A235012}\\
91&\oeis{A235080}&Empirical array recurrence of order 4&P&\pageref{round6:r6entry:A235080}\\
92&\oeis{A235081}&Empirical array recurrence of order 10&P&\pageref{round6:r6entry:A235081}\\
93&\oeis{A235271}&Empirical array recurrence of order 4&P&\pageref{round6:r6entry:A235271}\\
94&\oeis{A235272}&Empirical array recurrence of order 9&P&\pageref{round6:r6entry:A235272}\\
95&\oeis{A235273}&Empirical array recurrence of order 9&P&\pageref{round6:r6entry:A235273}\\
96&\oeis{A235303}&Empirical array recurrence of order 5&P&\pageref{round6:r6entry:A235303}\\
97&\oeis{A235304}&Empirical array recurrence of order 14&P&\pageref{round6:r6entry:A235304}\\
98&\oeis{A235305}&Empirical array recurrence of order 14&P&\pageref{round6:r6entry:A235305}\\
99&\oeis{A235312}&Empirical array recurrence of order 6&P&\pageref{round6:r6entry:A235312}\\
100&\oeis{A235313}&Empirical array recurrence of order 18&P&\pageref{round6:r6entry:A235313}\\
\end{longtable}\normalsize

\Needspace{6\baselineskip}
\section{Formal-series congruences and three theta values}
\label{round6:r6:root}

This section treats eighteen distinct entries. Every counted target is an
explicit conjecture or empirical equality in its cited OEIS entry. The two
exponential-series entries below are selected for parity only. The other
congruence packages specify exactly which assertions they settle. Classical
theta identities, earlier project lemmas, and Sundaram's comparison method
are credited inputs. They are not presented as new general mathematics.

\subsection{Four nested-reversion sequences}
For an integer $k\ge1$, let $A_k=x+O(x^2)$ satisfy
\[
 A_k^{-1}(x)=x-x\frac{A_k(kx)}{k}.
\]
The superscript $-1$ here means compositional inversion. The apparent
division by $k$ is harmless: $A_k(kx)/k=\sum_{n\ge1}a_k(n)k^{n-1}x^n$.
Since the coefficient of degree $n$ of a compositional inverse is
$-a_k(n)$ plus an integer polynomial in earlier coefficients, this equation
determines a unique integer series, over any coefficient ring after reducing
the displayed integer polynomials.

\begin{lemma}\label{round6:r6:reversion}
For all $n\ge1$, $a_k(n)\equiv1\pmod{k+1}$. If $k$ is even, $a_k(n)$
is odd exactly when $n$ is a power of $2$.
\end{lemma}
\begin{proof}
Modulo $k+1$ we have $A_k(kx)/k\equiv-A_k(-x)$. Set
$F=x/(1-x)$. Its inverse is $x/(1+x)$, and
\[
 x+xF(-x)=x-\frac{x^2}{1+x}=\frac{x}{1+x}=F^{-1}(x).
\]
Thus $F$ satisfies the reduced equation; triangular uniqueness proves
$A_k\equiv F\pmod{k+1}$.

For even $k$, the integer series $A_k(kx)/k$ reduces to $x$ modulo $2$.
Consequently $A_k^{-1}\equiv x+x^2$. The series
$S(x)=\sum_{j\ge0}x^{2^j}$ satisfies $S+S^2=x$ in $\mathbb F_2[[x]]$,
so it is the unique inverse of $x+x^2$. Hence $A_k\equiv S\pmod2$.
\end{proof}

\begin{theorem}[A393856]\label{round6:r6entry:A393856}
For $k=4$, every coefficient is $1$ modulo $5$, and the odd coefficients
occur exactly at positive indices that are powers of $2$.
\end{theorem}
\begin{proof}Apply Lemma~\ref{round6:r6:reversion} with $k=4$. These are both
congruence conjectures in \oeis{A393856} (Paul D. Hanna, March 11, 2026).
\end{proof}
\begin{theorem}[A393857]\label{round6:r6entry:A393857}
For $k=5$, every coefficient is $1$ modulo $6$.
\end{theorem}
\begin{proof}Apply Lemma~\ref{round6:r6:reversion} with $k=5$ to the defining
equation in \oeis{A393857}. This proves its displayed congruence conjecture.
\end{proof}
\begin{theorem}[A393858]\label{round6:r6entry:A393858}
For $k=6$, every coefficient is $1$ modulo $7$, and the odd coefficients
occur exactly at positive indices that are powers of $2$.
\end{theorem}
\begin{proof}Apply Lemma~\ref{round6:r6:reversion} with $k=6$ to
\oeis{A393858}; both displayed congruence conjectures follow.
\end{proof}
\begin{theorem}[A393859]\label{round6:r6entry:A393859}
For $k=7$, every coefficient is $1$ modulo $8$.
\end{theorem}
\begin{proof}Apply Lemma~\ref{round6:r6:reversion} with $k=7$ to
\oeis{A393859}. All four applications have the source range $n\ge1$.
\end{proof}

\subsection{Divisibility in three products of iterates}
Here $A^{\circ j}$ always denotes composition, not an ordinary power.
\begin{lemma}\label{round6:r6:product}
Let $r,s$ be positive integers and $m=r+s$. The unique rational series
$A=x+O(x^2)$ satisfying
\[
 A^{\circ r}(x)A^{\circ s}(x)=x^2+m^2x^3
\]
belongs to $x+m x^2\mathbb Z[[x]]$.
\end{lemma}
\begin{proof}
In the original equation, the coefficient of $a(n)$ at degree $n+1$ is
$r+s=m$; all remaining terms depend on earlier coefficients. This gives
existence and uniqueness over $\mathbb Q$.

To prove the stronger integrality statement without dividing a congruence
by $m$, put $A=x+mB$, with $B\in x^2\mathbb Z[[x]]$ to be constructed.
For an indeterminate $t$, induction under composition gives
\[
 (x+tB)^{\circ j}=x+jtB+t^2H_j(t,B),
\]
where the coefficients of $H_j$ are integer polynomials in the coefficients
of $B$. Every term of $H_j$ involves at least two occurrences of $B$ or its
formal derivatives. More precisely, changing $B$ by $c x^n+O(x^{n+1})$
changes $H_j$ only in degrees at least $n+1$; this follows by expanding
$B(x+tB)-B(x)$ with the binomial theorem, since $\operatorname{ord} B\ge2$.
After multiplication and substitution $t=m$, the equation becomes
\[
 B+H_r+H_s+rs\frac{B^2}{x}
   +m\frac{rBH_s+sBH_r}{x}+m^2\frac{H_rH_s}{x}=x^2,
\]
where $H_j$ now means $H_j(m,B)$. Each displayed quotient has integer
coefficients. All terms other than the first $B$ depend at degree $n$
only on $b(2),\ldots,b(n-1)$. Thus this is a monic triangular recursion
for integer $b(n)$, starting with $b(2)=1$. The resulting $A=x+mB$ solves
the original equation and hence is its unique rational solution.
\end{proof}

\begin{theorem}[A396793]\label{round6:r6entry:A396793}
For the sequence determined by $A\,A^{\circ2}=x^2+9x^3$, every $a(n)$
with $n>1$ is divisible by $3$.
\end{theorem}
\begin{proof}Use $(r,s,m)=(1,2,3)$ in Lemma~\ref{round6:r6:product}.
This is the selected conjecture in \oeis{A396793}.
\end{proof}
\begin{theorem}[A396794]\label{round6:r6entry:A396794}
For $A\,A^{\circ3}=x^2+16x^3$, every $a(n)$ with $n>1$ is divisible by $4$.
\end{theorem}
\begin{proof}Use $(r,s,m)=(1,3,4)$; this is the conjecture in
\oeis{A396794}.
\end{proof}
\begin{theorem}[A396795]\label{round6:r6entry:A396795}
For $A^{\circ2}A^{\circ3}=x^2+25x^3$, every $a(n)$ with $n>1$ is divisible by $5$.
\end{theorem}
\begin{proof}Use $(r,s,m)=(2,3,5)$; this is the conjecture in
\oeis{A396795}. All three source assertions are due to Paul D. Hanna,
June 2026.
\end{proof}

\subsection{Two diagonal-zero series and Catalan valuations}
For $d=1,2$, let $A_d(0)=1$ be determined by
\[
 [x^n]\frac1{A_d(x)-n^d x}=0\qquad(n\ge1).
\]
The coefficient of $a_d(n)$ is $-1$, so these equations define unique
integer series, with uniqueness also after reduction modulo $4$.
Let $C_j=\binom{2j}{j}/(j+1)$ and put
\[
 C(x)=\sum_{j\ge0}(-1)^j C_j x^{2j},\qquad
 U(x)=1+x+x^2C(x).
\]
The Catalan equation gives $C=1-x^2C^2$. Consequently the following are
exact identities in integer formal series:
\begin{equation}\label{round6:r6:catalan-inverses}
 U-x=C^{-1},\qquad U^{-1}=1-xC.
\end{equation}

\begin{lemma}\label{round6:r6:diagonal}
Both $A_1$ and $A_2$ are congruent to $U$ modulo $4$.
\end{lemma}
\begin{proof}
For $d=2$, an even $n$ has $n^2\equiv0\pmod4$ and
$[x^n]U^{-1}=0$; an odd $n$ has $n^2\equiv1\pmod4$ and
$[x^n](U-x)^{-1}=[x^n]C=0$. Thus $U$ satisfies every required diagonal
equation modulo $4$.

For $d=1$, the classes $n\equiv0,1\pmod4$ follow in the same way.
For $n\equiv2\pmod4$,
\[
 (U-2x)^{-1}\equiv U^{-1}+2xU^{-2}\pmod4.
\]
The first summand has no positive even coefficients. Modulo $2$ the
square $U^{-2}$ has only even degrees, so the second summand also has
zero coefficient at this even $n$. Finally, for $n\equiv3\pmod4$,
\[
 (U-3x)^{-1}\equiv C+2xC^2\pmod4.
\]
The first summand has only even degrees, and $C^2$ modulo $2$ has degrees
divisible by $4$. Since $n-1\equiv2\pmod4$, the required coefficient
vanishes. Triangular uniqueness proves the lemma in both cases.
\end{proof}

\begin{theorem}[A393170]\label{round6:r6entry:A393170}
For $d=1$ and $n>0$, $a_1(n)$ is odd exactly when $n$ is a power of $2$.
Moreover, $a_1(n)\equiv2\pmod4$ exactly when $n/2$ is the sum of two
distinct powers of $2$.
\end{theorem}
\begin{theorem}[A393173]\label{round6:r6entry:A393173}
The same two assertions hold for $d=2$.
\end{theorem}
\begin{proof}[Proof of both theorems]
Lemma~\ref{round6:r6:diagonal} gives $a_d(1)\equiv1$, $a_d(n)\equiv0$ at odd
$n>1$, and
\[
 a_d(2k)\equiv(-1)^{k-1}C_{k-1}\pmod4\qquad(k\ge1).
\]
Legendre's factorial formula gives
$v_2(C_{k-1})=s_2(k)-1$, where $s_2(k)$ is the number of $1$'s in the
binary expansion of $k$. Thus oddness is equivalent to $s_2(k)=1$,
and residue $2$ modulo $4$ is equivalent to $s_2(k)=2$; a sign does not
change the latter residue. Including $n=1$ proves both statements.
These are the two congruence conjectures in each of \oeis{A393170} and
\oeis{A393173} (Paul D. Hanna, March 5, 2026). Their asymptotic assertions
are outside this selection.
\end{proof}

\subsection{A393171: a mod-eight consequence of a Catalan substitution}
\begin{theorem}\label{round6:r6entry:A393171}
Let $A=x+O(x^2)$ satisfy $A(x)=x+A(C_*(x))^2$, where
$C_*(x)=x+C_*(x)^2$. For $n>2$, its coefficient is $4$ modulo $8$
exactly when $n=2^j$ or $2^j+1$ for some $j\ge1$; all other such
coefficients are divisible by $8$.
\end{theorem}
\begin{proof}
Put $f=x-x^2$, so $C_*=f^{-1}$, and set $B=A\circ C_*$. Then
$B(f(x))=A(x)=x+B(x)^2$. This is the defining series of
\oeis{A393172}, used in Round III and not counted again here. We reproduce
the required short lemma. The equation for $B$ is monic triangular.
Modulo $2$ its unique solution is $B=x$, so write $B=x+2H$.
Reduction modulo $4$ gives $H(x+x^2)=x^2$ in $\mathbb F_2[[x]]$.
Since the inverse of $x+x^2$ is $S=\sum_{j\ge0}x^{2^j}$,
\[
 H\equiv S^2=\sum_{j\ge1}x^{2^j}\pmod2.
\]
Now square $B$ modulo $8$:
\[
 A=x+B^2\equiv x+x^2+4xH+4H^2\pmod8.
\]
The two supports in positive degrees are $2^j+1$ for $j\ge1$ and
$2^j$ for $j\ge2$, respectively, and do not overlap. This is exactly
the asserted set after imposing $n>2$. It proves the displayed conjecture
in \oeis{A393171}.
\end{proof}

\subsection{A389537: lifting the inverse of a dyadic series}
\begin{theorem}\label{round6:r6entry:A389537}
Let $R(x)=\sum_{j\ge0}(-1)^j x^{2^j}$ and $A=R^{-1}$, the series in
\oeis{A389537}. Then
\[
 A(x)\equiv x+x^2+2\sum_{j\ge0}x^{3\cdot2^j}\pmod4.
\]
Thus for $n>2$ its coefficient is $2$ modulo $4$ exactly at
$n=3\cdot2^j$, and is otherwise divisible by $4$.
\end{theorem}
\begin{proof}
Modulo $2$, $R=\sum x^{2^j}$ has inverse $B=x+x^2$.
For $j\ge1$, the binomial theorem gives
\[
 B^{2^j}\equiv x^{2^j}+2x^{3\cdot2^{j-1}}+x^{2^{j+1}}\pmod4.
\]
Indeed, among the interior binomial coefficients in row $2^j$, only
the central one has $2$-adic valuation one: for $0<i<2^j$ the valuation
is $j-v_2(i)$. Substitution in $R(B)$ cancels consecutive endpoint
terms and yields
\[
 R(B)\equiv x+2\sum_{j\ge0}x^{3\cdot2^j}\pmod4.
\]
Also $R'(x)\equiv1\pmod2$. Therefore for any integer series $H$,
$R(B+2H)\equiv R(B)+2H\pmod4$ by the binomial theorem.
Choosing $H=\sum x^{3\cdot2^j}$ makes the composition equal to $x$
modulo $4$, so uniqueness of the compositional inverse proves the result.
The separate observation about runs of signs in the source is not selected.
\end{proof}

\subsection{A396798: an index correction and all eight iterate formulas}
Let $A=x+A^{\circ4}A^{\circ5}$. This equation has a unique integer
solution, and a unique solution over any coefficient ring: a change in
degree $n$ in $A$ changes the product of its zero-constant iterates only
in degrees at least $n+1$.

\begin{theorem}\label{round6:r6entry:A396798}
Put $F_j=x/(1-jx)$ and $h=x^4/(1-x)^3$. For every integer $k$,
\[
 A^{\circ k}(x)\equiv
 \begin{cases}
 F_k(x),&k\ \text{even},\\
 F_k(x)+4h(x),&k\ \text{odd}
 \end{cases}\pmod8.
\]
In particular $a(1)=1$, and the coefficients for $n\ge2$ repeat
$[1,1,5,5]$ modulo $8$. The first conjecture on the consulted page
incorrectly starts that repeating list at $n=1$; already $a(3)=9\equiv1$,
whereas that start would predict $5$. Its seven other listed iterate
conjectures, with their stated starting indices, are true.
\end{theorem}
\begin{proof}
Set $R=F_1+4h$. Since $F_r\circ F_s=F_{r+s}$, the first-order
composition expansion modulo $8$ gives
\[
 R^{\circ2}\equiv F_2+
 4\bigl(F_1'(F_1)h+h(F_1)\bigr)=F_2+\frac{8x^4}{(1-2x)^3}
 \equiv F_2\pmod8.
\]
It follows that $R^{\circ 2j}\equiv F_{2j}$. For odd iterates,
$h(F_{2j})\equiv h(x)\pmod2$, so
$R^{\circ(2j+1)}\equiv F_{2j+1}+4h$. This holds for negative iterates
as well, because composition and inversion preserve coefficientwise
congruences for series with linear coefficient $1$.

To check the defining equation for $R$, the same formulas give
$R^{\circ4}\equiv F_4$ and $R^{\circ5}\equiv F_5+4h$. Direct rational
simplification gives
\[
 R(x)-x-F_4(x)\cdot(F_5(x)+4h(x))=
 \frac{8x^3(1-3x)(5x^2-2x+1)}{(x-1)^3(4x-1)(5x-1)}.
\]
Every denominator here has unit constant term, so this is zero modulo
$8$. By uniqueness, $A\equiv R$, and the iterate formulas follow.

For odd $k$, coefficient extraction gives
$[x^n]A^{\circ k}\equiv k^{n-1}+4\binom{n-2}{2}$ for $n\ge2$,
with the binomial term interpreted as zero at $n=2,3$. The second term
repeats $0,0,4,4$ from $n=2$. The four odd iterates therefore have
patterns $[1,1,5,5]$, $[3,1,7,5]$, $[5,1,1,5]$, and $[7,1,3,5]$,
respectively. For even $k=2,4,6,8$, $k^{n-1}$ vanishes modulo $8$
after $n=3,2,3,1$, respectively. These are exactly the source ranges.
We count the entire eight-assertion package once, as an index correction.
\end{proof}

\subsection{A396807: comparison with a rational iterate}
\begin{theorem}\label{round6:r6entry:A396807}
For $A=x+A^{\circ5}A^{\circ6}$ and every integer $k$,
\[
 A^{\circ k}(x)\equiv\frac{x}{1-kx}\pmod{10}.
\]
Thus $a(n)\equiv1\pmod{10}$ for $n\ge1$, and more generally
$[x^n]A^{\circ k}\equiv k^{n-1}\pmod{10}$, with $0^0=1$.
\end{theorem}
\begin{proof}
We use the rational comparison method in Shiva Priyan Shanmuga Sundaram's
September 2026 proof for \oeis{A396797}; that already proved entry is
excluded. His note treats iterates $3,4$ modulo $6$. For the present
iterates $5,6$ the required separate computation is
\[
 (1-5x)(1-6x)=1-11x+30x^2\equiv1-x\pmod{10}.
\]
Consequently $F=x/(1-x)$ satisfies $F\equiv x+F^{\circ5}F^{\circ6}$
modulo $10$. The same degree argument as in the previous subsection
proves uniqueness, so $A\equiv F$. Composition gives the assertion for
positive $k$; inversion extends it to negative $k$, and $k=0$ gives $x$.
This proves both conjectures in \oeis{A396807}. The method is an explicit
extension of a credited proof, not a claim of a newly invented technique.
\end{proof}

\subsection{A396805 and A396806: parity of exponential coefficients}
An integer Hurwitz series is written $G=\sum g_n x^n/n!$ with
$g_n\in\mathbb Z$. Its coefficient congruences refer to $g_n$.
Products, composition of zero-constant series, and exponentiation of a
zero-constant series use integer binomial or Bell polynomials. We never
reduce a factorial denominator by inverting it modulo a prime.

\begin{lemma}\label{round6:r6:hurwitz}
For every positive integer $\ell$, the unique Hurwitz solution
\[
 A_\ell=x\exp(A_\ell^{\circ\ell})
\]
has $a_\ell(n)\equiv n\pmod2$ for all $n\ge1$.
\end{lemma}
\begin{proof}
The degree-$n$ Hurwitz coefficient on the right is $n$ times a Bell
polynomial depending only on coefficients of $A_\ell$ below degree $n$.
Thus the equation has a unique integral solution, and is uniquely solvable
after reduction modulo $2$.

Let $F=xe^x$, whose $n$-th Hurwitz coefficient is $n$. By Lagrange
inversion its compositional inverse has coefficient $(-n)^{n-1}$.
This is congruent to $n$ modulo $2$, including $n=1$, so
$F^{\circ2}\equiv x\pmod2$ (the lemma used in Round V is reproduced here).
Also the tree series $T=xe^T$ has Hurwitz coefficient $n^{n-1}$, hence
$T\equiv F\pmod2$. Reduction of its defining equation gives
$F\equiv xe^F\pmod2$.

If $\ell$ is even then $F^{\circ\ell}\equiv x$, and
$x\exp(F^{\circ\ell})\equiv xe^x=F$. If $\ell$ is odd then
$F^{\circ\ell}\equiv F$, and the tree-series identity gives the same
conclusion. Uniqueness proves $A_\ell\equiv F$ in both cases.
\end{proof}

\begin{theorem}[A396805]\label{round6:r6entry:A396805}
For $\ell=5$, $a(n)$ is odd if and only if $n$ is odd.
\end{theorem}
\begin{proof}Apply Lemma~\ref{round6:r6:hurwitz} with $\ell=5$.
Only this explicit parity conjecture is selected from \oeis{A396805};
its mod-$3$ and mod-$5$ conjectures are not claimed here.
\end{proof}
\begin{theorem}[A396806]\label{round6:r6entry:A396806}
For $\ell=6$, $a(n)$ is odd if and only if $n$ is odd.
\end{theorem}
\begin{proof}Apply Lemma~\ref{round6:r6:hurwitz} with $\ell=6$.
Only this parity conjecture is selected from \oeis{A396806}; its mod-$3$
and mod-$6$ assertions are outside scope. These two entries are distinct
sequences, each counted once, despite sharing a general proof.
\end{proof}

\subsection{Three empirical special values}
Put
\[
 q=e^{-\pi},\qquad T=\frac{\pi^{1/4}}{\Gamma(3/4)},\qquad
 \phi(z)=\sum_{j\in\mathbb Z}z^{j^2},\quad
 \psi(z)=\sum_{j\ge0}z^{j(j+1)/2},\quad
 \chi(z)=\prod_{j\ge0}(1+z^{2j+1}).
\]
We use the following classical theta evaluations:
\begin{equation}\label{round6:r6:theta-values}
 \phi(q)=T,\quad \phi(-q)=2^{-1/4}T,\quad
 \psi(q)=2^{-5/8}q^{-1/8}T,\quad
 \chi(q)=2^{1/4}q^{1/24},\quad
 \psi(-q)^4=2^{-3}q^{-1/2}T^4.
\end{equation}
For clarity, these are established inputs, not numerical discoveries.
Here is their derivation from standard identities. At nome $e^{-\pi}$,
the elliptic modulus is $1/\sqrt2$; the theta--elliptic relation and
$K(1/\sqrt2)=\Gamma(1/4)^2/(4\sqrt\pi)$ give
$\theta_3(0,q)=T$ by the gamma reflection formula. The theta quartic
identity gives $\theta_4(0,q)=2^{-1/4}T$.
Duplication gives $\psi(q)^2=\phi(q)\psi(q^2)$ and
$\psi(q^2)=\theta_2(0,q)/(2q^{1/4})=2^{-5/4}q^{-1/4}T$,
which proves the stated $\psi(q)$ value.
For the remaining product, put $E(z)=\prod_{n\ge1}(1-z^n)$ and
$\eta(i)=q^{1/12}E(q^2)$. The Jacobi products give
$\phi(q)=E(q^2)\chi(q)^2$ and
$2\eta(i)^3=\theta_2(0,q)\theta_3(0,q)\theta_4(0,q)$.
Thus $E(q^2)=q^{-1/12}T/\sqrt2$ and
$\chi(q)^2=\sqrt2 q^{1/12}$, proving its stated value.
Finally, applying the same duplication identity
$\psi(z)^2=\phi(z)\psi(z^2)$ at $z=-q$ gives
$\psi(-q)^2=\phi(-q)\psi(q^2)$, and hence the fourth-power value
in \eqref{round6:r6:theta-values}.
The standard identities and elliptic special value are documented in
DLMF \S\S19.20, 20.7, 20.9 and the theta reference linked in the entries.
All signs are fixed by the positive real products at $0<q<1$.

\begin{theorem}[A096920]\label{round6:r6entry:A096920}
For the sequence with generating function $\psi(x)\chi(x)$,
\[
 \sum_{n\ge0}a(n)e^{-n\pi}
 =2^{-3/8}e^{\pi/12}\frac{\pi^{1/4}}{\Gamma(3/4)}.
\]
\end{theorem}
\begin{proof}
The known theta-product generating function is recorded in
\oeis{A096920}. Multiplying the $\psi(q)$ and $\chi(q)$ values gives
$2^{-3/8}q^{-1/12}T$, exactly the empirical equality attributed there
to Simon Plouffe, September 17, 2025. The factors are analytic in
$|x|<1$, so the absolutely convergent evaluation is legitimate.
\end{proof}

\begin{theorem}[A083365]\label{round6:r6entry:A083365}
For the sequence with generating function $\psi(x)/\phi(x)$,
\[
 \sum_{n\ge0}a(n)e^{-n\pi}=2^{-5/8}e^{\pi/8}.
\]
\end{theorem}
\begin{proof}
Divide the corresponding values in \eqref{round6:r6:theta-values}.
This is the empirical equality in \oeis{A083365} (Simon Plouffe,
September 17, 2025). The Jacobi product for $\phi$ has no zeros in
$|x|<1$, so the quotient's Taylor series converges at $q$.
\end{proof}

\begin{theorem}[A273228]\label{round6:r6entry:A273228}
For the sequence with generating function $1/\psi(-x)^4$,
\[
 \sum_{n\ge0}a(n)e^{-n\pi}
 =8e^{-\pi/2}\frac{\Gamma(3/4)^4}{\pi}.
\]
\end{theorem}
\begin{proof}
The source generating function is a previously known partition product;
its derivation is not counted as new. Substituting the last value in
\eqref{round6:r6:theta-values} gives $8q^{1/2}T^{-4}$, which is the empirical
equality in \oeis{A273228} (Simon Plouffe, September 21, 2025).
Its product has no zeros in the unit disk, so inversion and evaluation
are valid. The three special values use classical identities and share
one method; they are three distinct source assertions, not three new
theta theorems. The linked decimal-constant entries add no counts.
\end{proof}


\Needspace{6\baselineskip}
\section{Arithmetic identities and exact modular certificates}
\label{round6:r6:arithmetic}

This section resolves eleven selected packages at eleven distinct OEIS
entries.  Each entry still explicitly labels the selected statement a
conjecture in the primary page inspected on October 20, 2026.  Two packages
also settle an accompanying empirical observation.  The count is one per
entry, including when a theorem proves more than one statement there.
The standard-library program \texttt{r6\_arith\_verify.py} recomputes
\texttt{r6\_arith\_certificates.json} using exact integer and rational
arithmetic.  The modular coefficient table for \oeis{A116184} is a finite
part of its proof; the other finite checks supplement universal arguments.
Known formulas used as inputs are credited below.  The audit concerns the
inspected entries, relevant cross-references and targeted searches; it does
not establish worldwide priority.

\subsection{The harmonic denominator ratio in A061002}
\label{round6:r6entry:A061002}

Let $H_m=\sum_{j=1}^m1/j$, and write $\operatorname{num}$ and
$\operatorname{den}$ for numerator and positive denominator in lowest
terms.  The selected statement is Gary Detlefs's February 20, 2013
conjecture in \oeis{A061002}.

\begin{theorem}
For every prime $p>3$,
\[
 \frac{\operatorname{den}(H_p/H_{p-1})}
 {\operatorname{num}(H_{p-1}/p^2)}=p^3.
\]
\end{theorem}
\begin{proof}
Write $H_{p-1}=A/B$ in lowest terms.  None of its summands has a denominator
divisible by $p$, so $p\nmid B$.  We first recall the elementary harmonic
form of Wolstenholme's theorem, with a proof to specify the input precisely.
In the localization of the integers at $p$,
\[
 \frac1{p-j}\equiv-\frac1j-\frac p{j^2}\pmod{p^2}.
\]
Summing for $1\leq j<p$ gives
\[
 2H_{p-1}\equiv-p\sum_{j=1}^{p-1}j^{-2}\pmod{p^2}.
\]
Inversion permutes the nonzero residues modulo $p$, and hence
\[
 \sum_{j=1}^{p-1}j^{-2}\equiv\sum_{j=1}^{p-1}j^2
 =\frac{p(p-1)(2p-1)}6\equiv0\pmod p.
\]
Here $6$ and $2$ are invertible because $p>3$.  Thus $p^2\mid A$.

Now
\[
 \frac{H_p}{H_{p-1}}=\frac{pA+B}{pA},\qquad
 \gcd(pA+B,pA)=\gcd(B,pA)=1.
\]
Its denominator is therefore $pA$.  Also
\[
 \frac{H_{p-1}}{p^2}=\frac{A/p^2}{B}
\]
is already reduced, with numerator $A/p^2$.  The asserted quotient is
$pA/(A/p^2)=p^3$.  This reasoning also covers primes for which $p^3$ or a
higher power divides $A$.
\end{proof}

\paragraph{Scope and verification.}
The classical integrality underlying the definition of \oeis{A061002}
is a known input, not an additional resolution.  The selected target is
the denominator ratio.  The verifier checks every prime $5\leq p\leq200$
using reduced rational harmonic numbers and matches the displayed prefix.

\subsection{The two-variable gcd identity in A063956}
\label{round6:r6entry:A063956}

Define $f(n)=\sum_{p:\,v_p(n)=1}p$, the sum of prime divisors occurring
to exponent exactly one.  Velin Yanev's February 17, 2019 formula in
\oeis{A063956} is marked conjectural.  The additive prime-power
description, also stated by Amiram Eldar in July 2024, is credited as a
known description of the sequence.

\begin{theorem}
For all positive integers $n,m$,
\[
 f(nm)=f(n)+f(m)-f(\gcd(n^2,m))-f(\gcd(n,m^2)).
\]
\end{theorem}
\begin{proof}
Fix a prime $p$ and put $a=v_p(n)$, $b=v_p(m)$.  The coefficient of $p$
on each side is governed by
\[
 [a+b=1]=[a=1]+[b=1]
       -[\min(2a,b)=1]-[\min(a,2b)=1],
\]
where brackets denote indicators.  If $a=0$ or $b=0$, both subtracted
indicators are zero and the identity follows immediately.  If $a,b>0$,
then $a+b\ne1$, while
$\min(2a,b)=1$ precisely when $b=1$, and
$\min(a,2b)=1$ precisely when $a=1$.  The two positive indicators cancel
the two negative ones.  Multiplying by $p$ and summing over primes proves
the result.
\end{proof}

\paragraph{Verification.}
Independent factorization checks cover all $1\leq n,m\leq200$.
The proof's case distinction applies to every exponent; the finite table
in the certificate is a supplementary check, not a bound on the exponents.

\subsection{Squarefree Bernoulli-denominator divisibility in A068991}
\label{round6:r6entry:A068991}

Let
\[
 S(k)=\sum_{d\mid k}\frac{\sigma(d)}{\varphi(d)}.
\]
The selected conjecture on \oeis{A068991} says that a squarefree integer
$k$ for which $S(k)$ is an integer divides the denominator of $B_{2k}$.
Here $B_j$ denotes the Bernoulli numbers.

\begin{theorem}
If $k$ is squarefree and $S(k)\in\mathbb Z$, then
\[
 k\mid\operatorname{den}(B_{2k}).
\]
\end{theorem}
\begin{proof}
Write $r=\omega(k)$.  Multiplicativity and squarefreeness give
\[
 S(k)=\prod_{p\mid k}\left(1+\frac{p+1}{p-1}\right)
     =\frac{2^r k}{\varphi(k)}.
\]
Consequently $\varphi(k)\mid2^r k$.  For each odd prime $p\mid k$, write
\[
 p-1=2^{\alpha_p}u_p,\qquad u_p\ \hbox{odd},\quad\alpha_p\geq1.
\]
The odd part of $\varphi(k)$ divides the odd part of the squarefree
integer $k$.  Thus every $u_p$ is squarefree, its prime factors divide
$k$, and each odd prime can occur among all the $u_p$ at most once.

If $k$ is odd, there are $r$ values of $\alpha_p$ and
$\sum_{p\mid k}\alpha_p\leq r$.  Each is at least one, so all equal one.
It follows that $p-1\mid2k$ for every $p\mid k$.

Suppose next that $k$ is even, and let $t=r-1$ be the number of its odd
prime divisors.  We have
\[
 \sum_{p\mid k,\ p>2}\alpha_p\leq t+2,
 \qquad
 \sum_{p\mid k,\ p>2}(\alpha_p-1)\leq2.
\]
It suffices to show that all $\alpha_p\leq2$.  The only alternative
allowed by the displayed inequality is one prime with $\alpha_p=3$
and all the other odd prime divisors having $\alpha_p=1$.
We show this alternative is impossible.

On the set of the $t$ odd prime divisors of $k$, draw an arrow $q\to p$
whenever $q\mid p-1$.  Every arrow goes from a smaller prime to a larger
one.  Every vertex has outdegree at most one, because an odd prime occurs
in the product of all the $u_p$ at most once.  A vertex with no incoming
arrow has $u_p=1$ and hence $p=2^{\alpha_p}+1$.
Under the alternative being considered, $\alpha_p$ is either $1$ or $3$;
the only prime of those two forms is $3$, since $2^3+1=9$ is composite.
In particular $3$ is present, is the unique vertex of indegree zero,
and all other vertices have indegree at least one.

The largest vertex has outdegree zero.  There are therefore at most
$t-1$ arrows and, by the preceding indegree count, at least $t-1$ arrows.
It follows that every vertex except $3$ has indegree one and every vertex
except the largest has outdegree one.  Since arrows strictly increase,
the graph is a single chain starting at $3$.  Along that chain, a
successor is $2^{\alpha_p}$ times its predecessor plus one, because
$u_p$ is squarefree and has exactly one prime divisor.
Before the unique exponent $3$, a successor uses exponent $1$.
The chain can then begin only $3,7$: a further such successor would be
$15$, which is composite.  If exponent $3$ follows $3$, its value would
be $25$; if it follows $7$, its value would be $57$.  Both are composite.
The alternative is impossible.

Thus, for even $k$, each $\alpha_p\leq2$ and $u_p$ divides the odd part
of $k$, so again $p-1\mid2k$ for every $p\mid k$.  The prime $2$ causes
no restriction; $k=1$ and $k=2$ are also covered directly.

Finally, the classical von Staudt--Clausen theorem states
\[
 \operatorname{den}(B_{2k})
 =\prod_{q\ \text{prime},\ q-1\mid2k}q.
\]
Every prime divisor of $k$ occurs in this product.  Since $k$ is
squarefree, the desired divisibility follows.
\end{proof}

\paragraph{Known input and verification.}
The Bernoulli denominator theorem is used in its classical form; see
\url{https://dlmf.nist.gov/24.10.i}.  It is not counted as a new result.
The verifier checks the derived condition $p-1\mid2k$ for every
qualifying squarefree $k\leq100000$.  It also checks the initial cases
against Bernoulli numbers computed by their defining rational recurrence,
independently of the denominator theorem.

\subsection{The powerful-divisor transform in A069208}
\label{round6:r6entry:A069208}

Call $d$ powerful if every nonzero prime exponent in $d$ is at least two;
$1$ is powerful.  Let $\chi(d)$ be its indicator, namely
\oeis{A112526}.  Werner Schulte's January 23, 2025 conjecture on
\oeis{A069208} concerns the transform
\[
 P_f(n)=\sum_{\substack{d\mid n\\d\ \text{powerful}}}f(d).
\]

\begin{theorem}
For every multiplicative arithmetic function $f$ with $f(1)=1$, the
function $P_f$ is multiplicative.  For any arithmetic function $f$,
$P_f$ is the inverse M\"obius transform of the pointwise product $\chi f$;
equivalently,
\[
 P_f=\mathbf1*(\chi f),\qquad \mu*P_f=\chi f.
\]
Here $*$ denotes Dirichlet convolution and $\mathbf1(n)=1$.
\end{theorem}
\begin{proof}
For coprime $u,v$, unique factorization shows
$\chi(uv)=\chi(u)\chi(v)$: the prime exponents of $uv$ are precisely
those of $u$ and $v$.  Thus $\chi f$ is multiplicative whenever $f$ is.
Every divisor of $uv$ has a unique form $d_1d_2$ with $d_1\mid u$ and
$d_2\mid v$, so
\[
 P_f(uv)=\sum_{d_1\mid u}\sum_{d_2\mid v}
     \chi(d_1)f(d_1)\chi(d_2)f(d_2)=P_f(u)P_f(v).
\]
Independently of multiplicativity, the defining sum is exactly
$\mathbf1*(\chi f)$.  Convolving with $\mu$, and using
$\mu*\mathbf1=\varepsilon$, gives the second formula.
\end{proof}

\paragraph{Scope and verification.}
This is one package for the explicitly conjectured general transform.
The entry's older formula for its base sequence and the already known
powerful-number indicator are not additional counts.  The verifier tests
$f(n)=1,n,n^2,(-1)^{\Omega(n)},\tau(n)$ through $300$, checking both
M\"obius inversion and every coprime product that stays within the bound.

\subsection{A complete harmonic congruence certificate for A116184}
\label{round6:r6entry:A116184}

For $n\geq1$, put
\[
 H(36,n)=\sum_{j=1}^{36}j^{-n}.
\]
The entry \oeis{A116184}, authored by Alexander Adamchuk on April 8,
2007, conjectures that every $3+36k$ belongs to the sequence defined by
$37^3\mid\operatorname{num}H(36,n)$.  It also records the empirical
observation that any prime term greater than $37$ is $-1$ modulo $37$.

\begin{theorem}
For every $k\geq0$,
\[
 37^3\mid\operatorname{num}H(36,3+36k).
\]
If a prime $\ell>37$ satisfies
$37^3\mid\operatorname{num}H(36,\ell)$, then
$\ell\equiv-1\pmod{37}$.
\end{theorem}
\begin{proof}
All denominators are relatively prime to $37$.  Therefore the numerator
condition is equivalent to $S(n)=0$ modulo $M=37^3=50653$, where
\[
 S(n)=\sum_{j=1}^{36}u_j^n\pmod M,
 \qquad u_j=j^{-1}\pmod M.
\]
Choose least nonnegative residues for the $u_j$.  Fermat's theorem gives
$u_j^{36}\equiv1\pmod{37}$, so define $h_j$ by
\[
 u_j^{36}\equiv1+37h_j\pmod M,\qquad0\leq h_j<37^2.
\]
For $0\leq r<36$ define the following least nonnegative residues:
\[
 \begin{split}
 C_0(r)&=\sum_j u_j^r\pmod{37^3},\\
 C_1(r)&=\sum_j u_j^r h_j\pmod{37^2},\\
 C_2(r)&=\sum_j u_j^r h_j^2\pmod{37}.
 \end{split}
\]
The binomial theorem proves, for every $k\geq0$,
\begin{equation}
 S(r+36k)\equiv C_0(r)+37kC_1(r)
            +37^2\binom{k}{2}C_2(r)\pmod{37^3}.
 \label{round6:r6eq:harmonic-quadratic}
\end{equation}
Terms of binomial degree at least three vanish modulo $37^3$.
Direct exact evaluation of these 36-term sums gives
\[
 (C_0(3),C_1(3),C_2(3))=(0,0,0).
\]
Equation~\eqref{round6:r6eq:harmonic-quadratic} proves the first assertion for
all $k$, without a bound on $k$.

For the prime assertion, $\ell>37$ has $\gcd(\ell,6)=1$.
Its residue $r$ modulo $36$ therefore belongs to the twelve units modulo
$36$.  For the eleven residues other than $35$, the exact coefficient
values are given below, with $C_0=37^2 A$, $C_1=37B$, and $C_2=D$.
All entries and polynomial equalities in the table are independently
recomputed by the accompanying verifier.
\[
\begin{array}{c|rrrr}
 r&A&B&D&c\\\hline
 1&24&13&24&12\\
 5&1&12&5&21\\
 7&29&2&5&21\\
 11&19&3&30&15\\
 13&20&13&36&18\\
 17&23&18&12&6\\
 19&30&34&6&3\\
 23&21&26&23&30\\
 25&19&27&30&15\\
 29&33&20&21&29\\
 31&23&24&4&2
\end{array}
\]
For every row,
\[
 A+Bk+D\binom{k}{2}
       \equiv c(k-r)(k-r-1)\pmod{37},\qquad c\ne0.
\]
Thus $S(r+36k)=0$ modulo $37^3$ forces
$k\equiv r$ or $r+1$ modulo $37$.  Since $36\equiv-1\pmod{37}$,
this means $\ell=r+36k\equiv0$ or $-1\pmod{37}$.
The zero residue is impossible for a prime $\ell>37$.

For the remaining residue $r=35$, exact evaluation gives
\[
 (C_0(35),C_1(35),C_2(35))=(17760,203,30).
\]
Already modulo $37^2$, the necessary condition $S(35+36k)=0$ reduces,
after division by $37$, to
\[
 36+18k=18(k-35)\equiv0\pmod{37}.
\]
This forces $k\equiv35\pmod{37}$ and hence
$\ell\equiv35-k\equiv0\pmod{37}$, again impossible.  The second
assertion follows.
\end{proof}

\paragraph{Reproducibility of the finite part.}
The following standard-library calculation specifies the coefficient
certificate completely; it uses neither approximate harmonic sums nor
large reduced rational powers.
\begin{verbatim}
p = 37
M = p**3
u = [pow(j, -1, M) for j in range(1, p)]
h = [(pow(v, 36, M) - 1)//p for v in u]
for r in range(36):
    a = [pow(v, r, M) for v in u]
    C0 = sum(a) % M
    C1 = sum(v*w for v, w in zip(a, h)) % (p*p)
    C2 = sum(v*w*w for v, w in zip(a, h)) % p
    print(r, C0, C1, C2)
\end{verbatim}
The included verifier also proves the table's polynomial factorizations
by comparing all three coefficients.  As an independent check,
$S(n)$ is periodic with period $36\cdot37^2=49284$ by Euler's theorem,
and the verifier enumerates every residue in this full period directly.
There are $2644$ zero residues, of which $407$ are relatively prime to
$222$; all of those $407$ are $36$ modulo $37$.
It also matches all 53 terms displayed on the source page.
The theorem is therefore supported both by the small quadratic
certificate and by a complete-period check, not a sampled prefix.

\subsection{The power-sum criterion in A133907}
\label{round6:r6entry:A133907}

Let $a(n)$ be the least prime $p$ for which
$\binom{n+p}{p}\equiv1\pmod p$.  The selected conjecture in
\oeis{A133907}, due to Thomas Ordowski on July 29, 2018, replaces the
binomial test by a power-sum test.  Robert Israel's May 11, 2017 floor
criterion on the same entry is a known input, and is also derived below.

\begin{theorem}
For every integer $n\geq1$ and prime $p$,
\[
 \binom{n+p}{p}\equiv1\pmod p
 \quad\Longleftrightarrow\quad
 \sum_{k=1}^n k^{p-1}\equiv n\pmod p.
\]
Consequently the least primes satisfying the two tests are equal.
\end{theorem}
\begin{proof}
Write $n=qp+r$ with $0\leq r<p$.  In $\mathbb F_p[x]$,
\[
 (1+x)^{n+p}\equiv(1+x^p)^{q+1}(1+x)^r.
\]
Extracting the coefficient of $x^p$ gives
$\binom{n+p}{p}\equiv q+1\pmod p$.  On the other hand, Fermat's
theorem makes $k^{p-1}$ equal to $1$ modulo $p$ for nonmultiples of $p$
and to $0$ for multiples.  There are exactly $q$ multiples among
$1,\ldots,n$, and hence
\[
 \sum_{k=1}^n k^{p-1}\equiv n-q\pmod p.
\]
Each test is therefore equivalent to $q\equiv0\pmod p$.
The set of qualifying primes is nonempty because every prime $p>n$
has $q=0$, so equality of the least primes follows as well.
\end{proof}

\paragraph{Verification.}
The verifier compares the actual binomial residue and the independently
summed powers for all $1\leq n\leq200$ and all primes $p\leq223$.
The chosen prime bound also guarantees a qualifying prime for each of
these $n$.

\subsection{The trinomial diagonal sum in A262735}
\label{round6:r6entry:A262735}

Let $a(n)$ be defined by
\[
 A(x)=\sum_{n\geq0}a(n)x^n
 =\frac{x(1-x-x^2)}{(1-x)(1-x-2x^2-x^3)}.
\]
Write $T(m,j)=[y^j](1+y+y^2)^m$, the ordinary trinomial coefficients
of \oeis{A027907}.  The selected statement in \oeis{A262735} is
the conjectural diagonal sum on the entry authored by Werner Schulte on
September 29, 2015; the formula line itself is undated.

\begin{theorem}
For every $n\geq0$,
\[
 a(n+1)=\sum_{j=0}^{\lfloor n/2\rfloor}T(n-j,2j).
\]
\end{theorem}
\begin{proof}
Let $b(n)$ denote the sum on the right.  Put $m=n-j$, so that
$0\leq j\leq m$.  All subsequent operations are valid as formal power
series: each coefficient receives only finitely many contributions.
Taking the even part of the trinomial polynomial gives
\[
 \begin{split}
 \sum_{n\geq0}b(n)x^n
 &=\sum_{m\geq0}x^m\sum_{j=0}^mT(m,2j)x^j\\
 &=\frac12\sum_{m\geq0}x^m
   \left((1+\sqrt{x}+x)^m+(1-\sqrt{x}+x)^m\right)\\
 &=\frac12\left(
 \frac1{1-x-x^2-x\sqrt{x}}
 +\frac1{1-x-x^2+x\sqrt{x}}\right)\\
 &=\frac{1-x-x^2}{(1-x-x^2)^2-x^3}\\
 &=\frac{1-x-x^2}{(1-x)(1-x-2x^2-x^3)}=\frac{A(x)}x.
 \end{split}
\]
The square-root notation denotes an even part in the formal variable
$\sqrt{x}$; the resulting expression contains only integer powers of
$x$.  Coefficient comparison proves the formula, including $n=0$.
\end{proof}

\paragraph{Verification.}
The verifier builds the trinomial rows by polynomial multiplication and
compares their diagonal sums with the independently supplied fourth-order
recurrence for $a(n)$ through $n=101$.  It also checks the polynomial
denominator factorization exactly.  The rational generating function and
the known trinomial definition are inputs, not extra resolutions.

\subsection{The Liouville twist in A299406}
\label{round6:r6entry:A299406}

The defining Dirichlet series of \oeis{A299406} is
\[
 \sum_{n\geq1}\frac{a(n)}{n^s}
 =\frac{\zeta(s)\zeta(6s)}{\zeta(2s)\zeta(3s)}.
\]
Let $b(n)=\oeis{A210826}(n)$ be defined by its Lambert series
\[
 \sum_{n\geq1}\frac{b(n)x^n}{1-x^n}=\sum_{m\geq1}x^{m^3},
\]
and let $\lambda(n)=(-1)^{\Omega(n)}$.  The conjecture on Werner
Schulte's 2018 entry is the following identity; its formula line is undated.

\begin{theorem}
For all $n\geq1$, $a(n)=\lambda(n)b(n)$.
\end{theorem}
\begin{proof}
Extracting coefficients from the Lambert series gives
\[
 \sum_{d\mid n}b(d)=[n\text{ is a cube}].
\]
M\"obius inversion yields $b=\mu*\chi_3$, where $\chi_3$ is the
indicator of cubes.  In particular $b$ is multiplicative and its
prime-power series is
\[
 B_p(t)=\sum_{e\geq0}b(p^e)t^e=\frac{1-t}{1-t^3}.
\]
These standard properties of $b$ already appear on \oeis{A210826}
in comments by R.~J. Mathar and Benoit Cloitre from 2012.
Since $\lambda$ is completely multiplicative, its pointwise twist
replaces $t$ by $-t$ in each prime-power series:
\[
 \sum_{e\geq0}\lambda(p^e)b(p^e)t^e
   =B_p(-t)=\frac{1+t}{1+t^3}.
\]
The defining zeta quotient for $a$ has Euler factor
\[
 \frac{(1-t^2)(1-t^3)}{(1-t)(1-t^6)}
 =\frac{1+t}{1+t^3}.
\]
Equality of every local coefficient, followed by unique factorization,
proves the result.  This can be read entirely as an identity of formal
Dirichlet series, so no analytic convergence assumption is needed.
\end{proof}

\paragraph{Verification.}
For $n\leq1000$, the verifier constructs $a$ by Dirichlet convolution
from its zeta quotient and $b$ independently from its Lambert divisor
recursion, then compares $a(n)$ with $\lambda(n)b(n)$.  The six-periodic
local coefficient formula is checked as a third representation.

\subsection{Integrality of the sine-power array A376484}
\label{round6:r6entry:A376484}

For $n,k\geq0$, define
\[
 A(n,k)=4^k\sum_{j=1}^n
             \sin^{2k}\left(\frac{2\pi j}{2n+1}\right).
\]
The source \oeis{A376484}, by Cheng-Jun Li, September 24, 2024,
conjectures integrality and observes divisibility by $2n+1$ when $n,k>0$.
The empty row has value zero; at $k=0$ the sum has $n$ terms of value one.

\begin{theorem}
For $k\geq1$ and $m=2n+1$,
\[
 A(n,k)=m\left(
   \frac12\binom{2k}{k}
   +\sum_{\ell=1}^{\lfloor k/m\rfloor}
             (-1)^\ell\binom{2k}{k+\ell m}\right).
\]
Thus all entries are integers, and $m\mid A(n,k)$ for $n,k>0$.
\end{theorem}
\begin{proof}
Let $\zeta=e^{2\pi i/m}$.  The elementary Laurent identity
\[
 (2-z-z^{-1})^k
   =\sum_{r=-k}^k(-1)^r\binom{2k}{k+r}z^r
\]
follows by expanding $(-1)^kz^{-k}(1-z)^{2k}$.
Since $m$ is odd, multiplication by $2$ permutes residues modulo $m$.
The terms indexed by $j$ and $m-j$ in a full sine-power sum are equal,
and the term at $j=0$ vanishes because $k\geq1$.  Therefore
\[
 \begin{split}
 2A(n,k)
 &=\sum_{j=0}^{m-1}(2-\zeta^{2j}-\zeta^{-2j})^k\\
 &=\sum_{j=0}^{m-1}(2-\zeta^j-\zeta^{-j})^k\\
 &=m\sum_{\substack{-k\leq r\leq k\\m\mid r}}
                     (-1)^r\binom{2k}{k+r}.
 \end{split}
\]
The last step is the root-of-unity filter
$\sum_{j=0}^{m-1}\zeta^{jr}=m[m\mid r]$.
Pair $r=\ell m$ with $r=-\ell m$.  Their binomial coefficients agree,
and $(-1)^{\ell m}=(-1)^\ell$ because $m$ is odd.  Dividing by two
gives the displayed formula.
Finally,
$\binom{2k}{k}=2\binom{2k-1}{k-1}$ is even for $k\geq1$,
so the expression in parentheses is an integer.
The boundary values $A(0,k)=0$ and $A(n,0)=n$ settle the remaining cases.
\end{proof}

\paragraph{Scope and verification.}
Integrality and the accompanying divisibility observation form one
package.  The quotient array \oeis{A376499} is not an additional count.
The source's other row and diagonal identifications are not needed here.
The verifier compares the formula against exact cyclic convolution by
$2-z-z^{-1}$ for $0\leq n\leq12$, $0\leq k\leq30$, and matches all
66 displayed antidiagonal terms.  No floating-point sine evaluations
are used.

\subsection{The harmonic row sums in A377278}
\label{round6:r6entry:A377278}

Put $m_j=2^j-1$ and define the denominator triangle
\[
 T(n,k)=m_km_{k+1}\quad(1\leq k<n),\qquad T(n,n)=m_n.
\]
The conjecture in \oeis{A377278}, an entry authored by Werner Schulte
on October 22, 2024, identifies the sum of reciprocals in a row with a
determinant ratio.  The conjecture comment itself is undated.
Let $P_n=\prod_{j=1}^n m_j=\oeis{A005329}(n)$ and let
$D_n=\oeis{A204243}(n)$, the determinant of the $n\times n$ matrix
whose first row and first column are all $1$, whose remaining diagonal
entries are $m_2,\ldots,m_n$, and whose other entries are zero.

\begin{theorem}
For every $n\geq1$,
\[
 \sum_{k=1}^n\frac1{T(n,k)}=\frac{D_n}{P_n}.
\]
\end{theorem}
\begin{proof}
Since $m_{k+1}=2m_k+1$,
\[
 \frac1{m_km_{k+1}}=\frac1{m_k}-\frac2{m_{k+1}}.
\]
Summing for $1\leq k<n$ and adding $1/m_n$ gives
\[
 \sum_{k=1}^n\frac1{T(n,k)}
 =1-\sum_{j=2}^n\frac1{m_j}.
\]
For the determinant, subtract $1/m_j$ times row $j$ from row $1$ for
each $j=2,\ldots,n$.  This leaves all first-row entries except the first
equal to zero, and the first entry becomes
$1-\sum_{j=2}^n1/m_j$.  Expansion along that row yields
\[
 D_n=\left(1-\sum_{j=2}^n\frac1{m_j}\right)
       \prod_{j=2}^n m_j.
\]
Because $m_1=1$, the product is $P_n$.  The empty sums and products
give the same conclusion for $n=1$.
\end{proof}

\paragraph{Known input and verification.}
The determinant formula itself was already given by Robert Israel on
\oeis{A204243} on November 30, 2015; the proof above recalls it for
completeness.  The selected conjecture is its connection to the later
harmonic triangle, not discovery of the determinant evaluation.
For $1\leq n\leq20$, the verifier compares reduced rational row sums
with determinants obtained by integer Bareiss elimination.

\subsection{Prime divisibility of the logarithmic derivative in A381355}
\label{round6:r6entry:A381355}

Let $p_n$ be the $n$th prime.  Define the formal series
$F(x)=1+x+\sum_{n\geq2}f_nx^n$ by
\[
 [x^n]F(x)^{p_n}=0\qquad(n>1),
\]
and set
\[
 \frac{xF'(x)}{F(x)}=\sum_{n\geq1}a_nx^n.
\]
The conjecture in \oeis{A381355}, an entry authored by Paul D. Hanna
on March 11, 2025, is $p_n\mid a_n$ for $n>1$.  The comment itself is
undated.  The related entry \oeis{A381353} contains
the defining coefficients $f_n$ and repeats the divisibility observation;
it is not counted separately.

\begin{theorem}
The defining equations uniquely determine $F\in\mathbb Z[[x]]$.
All $a_n$ are integers, and $p_n\mid a_n$ for every $n>1$.
\end{theorem}
\begin{proof}
We first prove integrality of $F$ by induction.  Suppose that the
coefficients below degree $n$ are integers, put $p=p_n$, and write
$F_{<n}=\sum_{j=0}^{n-1}f_jx^j$.  The defining equation is
\[
 pf_n+[x^n]F_{<n}^{p}=0.
\]
The freshman's dream gives $F_{<n}(x)^p\equiv F_{<n}(x^p)\pmod p$.
Since $p_n>n$, its coefficient of $x^n$ is zero modulo $p$.
Thus $f_n=-[x^n]F_{<n}^p/p$ is an integer, uniquely determined.
Starting from $f_0=f_1=1$ proves the assertion about $F$.
As a formal series with constant term $1$, $F$ has an inverse in
$\mathbb Z[[x]]$, so all coefficients of $xF'/F$ are integers.

Fix $n>1$ again and put $p=p_n$.  Differentiating $F^p$ and using
the defining equation gives the exact integer identity
\[
 p[x^{n-1}]F^{p-1}F'=n[x^n]F^p=0,
\]
so $[x^{n-1}]F^{p-1}F'=0$.  On the other hand,
$F(x)^p\equiv F(x^p)\equiv1\pmod{(p,x^p)}$, and hence
\[
 F^{p-1}\equiv F^{-1}\pmod{(p,x^p)}.
\]
The degree $n-1$ is less than $p$.  Multiplication by $F'$ and
extraction of that coefficient therefore yields
\[
 a_n=[x^{n-1}]\frac{F'}F
 \equiv[x^{n-1}]F^{p-1}F'=0\pmod p.
\]
This proves the conjectured divisibility for every $n>1$.
\end{proof}

\paragraph{Verification.}
The verifier constructs $F$ through degree $50$ by exact truncated
polynomial powers, checks all defining equations, and computes $xF'/F$
both by its coefficient recurrence and by an independent reciprocal
series product.  It checks every prime divisibility in that range and
matches the 25 displayed source terms.  The general proof does not
depend on a bound for the coefficients or on numerical testing.

\subsection*{Excluded stale statements and counting boundary}

Two initially promising entries were removed from the counted set after
following their cross-references.  The limit conjecture in
\oeis{A048651} follows from the known cumulative-sum relation in
\oeis{A177510} and Vaclav Kotesovec's 2014 asymptotic already stated on
\oeis{A008930}; its separate Schulte series was also explicitly derived
by Peter Bala in 2020.  Neither is counted.  The closure conjecture in
\oeis{A348713} already has a posted 2026 proof attributed to Bernard
Schott and is excluded as well.  Easy secondary recurrences in
\oeis{A349890} and \oeis{A385590} are left uncounted while their principal
permutation conjectures remain outside this dossier's conclusions.
The detailed rejection ledger is \texttt{r6\_arith\_audit\_notes.json}.


\Needspace{6\baselineskip}
\section{Thirty generating-function and recurrence packages}
\label{round6:sec:r6gf}
This section treats thirty distinct OEIS entries absent from the 508
previously counted A-IDs. The count is twenty-five direct proofs and
five boundary qualifications. Every selected formula is explicitly
labeled \emph{Conjecture} in the consulted primary entry. In each
boundary case the source does not specify a starting index; the exact
valid tail and all earlier admissible exceptions are stated below.
These qualifications do not attribute an unstated starting index to
the conjecture's author.

Known generating functions, enumerations, transforms, closed forms and
other recurrences are credited inputs. Thus these are proofs of selected
source assertions, not claims that the sequences, generating functions,
methods, or mathematical consequences are newly discovered. In
particular three entries already contain different polynomial
recurrences; those are identified individually. The prior-work audit
has a stated finite scope and does not establish worldwide priority.

\subsection{Exact finite certificates and all-index consequences}
Let $F(x)=\sum_{n\ge0}a_nx^n$, filling indices before the source offset
with zero only to write a formal power series. Put $\theta=x\,d/dx$ and
\[
\mathcal L_p=\sum_{j=0}^{r}x^jp_j(\theta+j).
\]
Since $\theta x^k=kx^k$, for $n$ such that all source indices are defined,
\begin{equation}\label{round6:eq:r6gf-extract}
[x^n]\mathcal L_pF=\sum_{j=0}^{r}p_j(n)a_{n-j}.
\end{equation}
Consequently an exact identity $\mathcal L_pF=P(x)$, with $P$ a
polynomial, proves the homogeneous recurrence for every $n>\deg P$.
The coefficient $[x^n]P$ gives the exact initial residual. Taking
$n\ge r+o$, where $o$ is the original offset, ensures that no extension
of the source's sequence is used in a claimed recurrence instance.

For the twenty-nine algebraic cases below, write $x=\xi(t)$ and
\[
F(\xi(t))=R_0(t)+R_1(t)\sqrt{D(t)},\qquad D_0=1,\quad D_1=D.
\]
Without a change of variable this means $t=x$ and $\xi=x$.
Every radical takes the formal branch with constant term one.
For rational functions define
\[
T_iR=\frac{\xi}{\xi'}\left(R'+\frac{D_i'}{2D_i}R\right),
\qquad c_m(t)=\sum_{j=0}^{r}\xi(t)^j[n^m]p_j(n+j).
\]
The chain rule gives
$\theta(R\sqrt{D_i})=(T_iR)\sqrt{D_i}$.
Thus the exact certificate for each entry is the following pair of
rational identities:
\begin{equation}\label{round6:eq:r6gf-cert}
\sum_m c_mT_0^mR_0=P(\xi(t)),\qquad
\sum_m c_mT_1^mR_1=0.
\end{equation}
All inputs $p_j,R_0,R_1,D,\xi$ and the full polynomial $P$ are printed
below. Substitution followed by clearing denominators gives two zero
polynomials over $\mathbb Q$. This is a finite algebraic certificate
of an infinite identity. No linear independence of radical terms is
needed: vanishing of each displayed residual suffices. Equation
\eqref{round6:eq:r6gf-extract} completes the proof for every stated index.

As an expanded illustration, for A025269 let
$D=1-8x+12x^2+12x^3+16x^4$ and $F=(1-\sqrt D)/2$.
Its displayed coefficients give
\[
\mathcal L_p=D\theta-\frac{xD'}2.
\]
Using $\theta\sqrt D=xD'/(2\sqrt D)$, one obtains
\[
\mathcal L_pF=-\frac{xD'\sqrt D}{4}
-\frac{xD'}4+\frac{xD'\sqrt D}{4}
=2x-6x^2-9x^3-16x^4.
\]
Hence the recurrence holds for every $n\ge5$. All other algebraic
certificates use the same chain rule, with at most four derivatives.
The remaining A052183 certificate is a rational identity between
successive factorial terms, proved separately below.

\subsection{Definitions, indexing and the three rational parameters}
Write
\[
C(x)=\frac{1-\sqrt{1-4x}}{2x},\qquad
M(x)=\frac{1-x-\sqrt{1-2x-3x^2}}{2x^2}.
\]
These satisfy $C=1+xC^2$ and $M=1+xM+x^2M^2$, with constant term one.
The large Schroeder series $S$ satisfies $S=1+xS+xS^2$.
For the four convolution entries A025268, A025269, A025270 and A025275,
the definition is $a_n=\sum_{k=1}^{n-1}a_ka_{n-k}$ for $n\ge5$.
It makes $F-F^2$ a polynomial of degree at most four, fixed by their
four initial values. Solving the quadratic with $F(0)=0$ yields
exactly the radicands listed below. This also derives their credited
generating functions directly from the definitions.

For A026030 the reflection principle gives
$a_n=\binom{2n}{n-1}-\binom{2n}{n-5}$ for $n\ge1$.
The Somos generating function used below starts with $x$ and preserves
this offset. A026571 starts at index two, and its given generating
function already has valuation two. In A228178, the source displays
$G=(xC+2x^2C^4)/(1-x)$, while its data and stated offset index the
coefficient of $x$ as $a_0$. Its supplied coefficient-list program
uses this same convention. The indexed series used here is therefore
$F=G/x=(C+2xC^4)/(1-x)$. This shift is explicit and is not a separately
counted result. All other ordinary seeds are the established formulas
identified in the individual source notes.

For A159772 let $t$ count an ordered rooted tree whose nodes all have
zero through three children. The ordered list of children gives
$t=x(1+t+t^2+t^3)$. Allowing unrestricted root degree and retaining the
root factor gives $F=x/(1-t)$, hence
\[
\xi(t)=\frac{t}{1+t+t^2+t^3},\qquad
F(\xi(t))=\frac{t}{1-t^4}.
\]
This is Howroyd's already listed interpretation, used as input; the
original binary-tree enumeration is not asserted anew.

For A243022 put $t=xC(x^2)$. The Catalan equation gives
$t=x(1+t^2)$ and $C(x^2)=1+t^2$. Thus its defining nested series is
\[
\xi(t)=\frac{t}{1+t^2},\qquad
-\frac1{C(x^2)C(xC(x^2))}
=-\frac{1+\sqrt{1-4t}}{2(1+t^2)}.
\]
For A243255 put $t=xM(x)$. Then $t=x(1+t+t^2)$ and
\[
\xi(t)=\frac{t}{1+t+t^2},\qquad
F=xM(x)C(xM(x))=tC(t)=\frac{1-\sqrt{1-4t}}2.
\]
In all three cases $\xi(0)=0$ and $\xi'(0)=1$, so the formal inverse
$t(x)$ exists uniquely. Apparent poles in individual $R_i$ are
interpreted in the Laurent field; their sum is the power series
specified by the input definition.

\subsection{The thirty statements and their complete certificates}
Each recurrence below means $\sum_jp_j(n)a_{n-j}=0$ in its stated
range. Sources were consulted on 20 October 2026. This date records
retrieval, not a claim that every search cache was refreshed that day.


\subsubsection{A025268: Jan 25 2015}
\label{round6:r6entry:A025268}

Source: \url{https://oeis.org/A025268}. Selected conjecture: R. J. Mathar, Jan 25 2015. Kotesovec supplied the generating function and a different order-four recurrence on January 25, 2015. The selected target is the separately conjectured Mathar order-five relation. Both are consequences of the known convolution definition; no novelty is claimed for that definition or generating function.

\textbf{Statement.} With the source offset $o=1$, the following recurrence holds for every integer $n\ge6$.

\begin{align*}
p_{0}(n)&=n\\
p_{1}(n)&=n + 1\\
p_{2}(n)&=- 10 \left(2 n - 5\right)\\
p_{3}(n)&=2 \left(2 n - 9\right)\\
p_{4}(n)&=2 \left(14 n - 79\right)\\
p_{5}(n)&=40 \left(n - 7\right).
\end{align*}

Use $t=x$ and $\xi=x$.

The credited input has the exact form $F=R_0+R_1\sqrt D$, with

\begin{align*}
R_0&=\frac{1}{2},\\
R_1&=- \frac{1}{2},\\
D&=\left(2 x - 1\right) \left(4 x^{3} + 4 x^{2} + 2 x - 1\right).
\end{align*}

\textbf{Proof certificate.} Inserting the displayed data into \eqref{round6:eq:r6gf-cert} gives exactly

\[\mathcal L_pF=P(x)=- x \left(5 x + 1\right) \left(8 x^{3} + 3 x^{2} - 1\right).\]

The radical residual is identically zero as a rational function. The rational residual equals the displayed $P$; both assertions are verified by clearing denominators in \eqref{round6:eq:r6gf-cert}. Taking coefficients in \eqref{round6:eq:r6gf-extract} proves the recurrence for every $n\ge6$, with the initial residuals exactly as stated.

\subsubsection{A025269: Nov 21 2014}
\label{round6:r6entry:A025269}

Source: \url{https://oeis.org/A025269}. Selected conjecture: R. J. Mathar, Nov 21 2014. Kotesovec's January 25, 2015 generating function is credited input. The selected assertion is Mathar's separately labeled recurrence.

\textbf{Statement.} With the source offset $o=1$, the following recurrence holds for every integer $n\ge5$.

\begin{align*}
p_{0}(n)&=n\\
p_{1}(n)&=- 4 \left(2 n - 3\right)\\
p_{2}(n)&=12 \left(n - 3\right)\\
p_{3}(n)&=6 \left(2 n - 9\right)\\
p_{4}(n)&=16 \left(n - 6\right).
\end{align*}

Use $t=x$ and $\xi=x$.

The credited input has the exact form $F=R_0+R_1\sqrt D$, with

\begin{align*}
R_0&=\frac{1}{2},\\
R_1&=- \frac{1}{2},\\
D&=\left(4 x - 1\right) \left(4 x^{3} + 4 x^{2} + 4 x - 1\right).
\end{align*}

\textbf{Proof certificate.} Inserting the displayed data into \eqref{round6:eq:r6gf-cert} gives exactly

\[\mathcal L_pF=P(x)=- x \left(16 x^{3} + 9 x^{2} + 6 x - 2\right).\]

The radical residual is identically zero as a rational function. The rational residual equals the displayed $P$; both assertions are verified by clearing denominators in \eqref{round6:eq:r6gf-cert}. Taking coefficients in \eqref{round6:eq:r6gf-extract} proves the recurrence for every $n\ge5$, with the initial residuals exactly as stated.

\subsubsection{A025270: Nov 21 2014}
\label{round6:r6entry:A025270}

Source: \url{https://oeis.org/A025270}. Selected conjecture: R. J. Mathar, Nov 21 2014. Kotesovec's January 25, 2015 generating function is credited input. The selected assertion is Mathar's separately labeled recurrence.

\textbf{Statement.} With the source offset $o=1$, the following recurrence holds for every integer $n\ge5$.

\begin{align*}
p_{0}(n)&=n\\
p_{1}(n)&=- 2 \left(2 n - 3\right)\\
p_{2}(n)&=- 4 \left(n - 3\right)\\
p_{3}(n)&=6 \left(2 n - 9\right)\\
p_{4}(n)&=20 \left(n - 6\right).
\end{align*}

Use $t=x$ and $\xi=x$.

The credited input has the exact form $F=R_0+R_1\sqrt D$, with

\begin{align*}
R_0&=\frac{1}{2},\\
R_1&=- \frac{1}{2},\\
D&=20 x^{4} + 12 x^{3} - 4 x^{2} - 4 x + 1.
\end{align*}

\textbf{Proof certificate.} Inserting the displayed data into \eqref{round6:eq:r6gf-cert} gives exactly

\[\mathcal L_pF=P(x)=- x \left(20 x^{3} + 9 x^{2} - 2 x - 1\right).\]

The radical residual is identically zero as a rational function. The rational residual equals the displayed $P$; both assertions are verified by clearing denominators in \eqref{round6:eq:r6gf-cert}. Taking coefficients in \eqref{round6:eq:r6gf-extract} proves the recurrence for every $n\ge5$, with the initial residuals exactly as stated.

\subsubsection{A025275: May 30 2014}
\label{round6:r6entry:A025275}

Source: \url{https://oeis.org/A025275}. Selected conjecture: R. J. Mathar, May 30 2014. Kotesovec's January 25, 2015 generating function is credited input. The selected assertion is Mathar's separately labeled recurrence.

\textbf{Statement.} With the source offset $o=1$, the following recurrence holds for every integer $n\ge5$.

\begin{align*}
p_{0}(n)&=n\\
p_{1}(n)&=- 2 \left(2 n - 3\right)\\
p_{2}(n)&=4 \left(n - 3\right)\\
p_{3}(n)&=- 2 \left(2 n - 9\right)\\
p_{4}(n)&=8 \left(n - 6\right).
\end{align*}

Use $t=x$ and $\xi=x$.

The credited input has the exact form $F=R_0+R_1\sqrt D$, with

\begin{align*}
R_0&=\frac{1}{2},\\
R_1&=- \frac{1}{2},\\
D&=\left(2 x - 1\right) \left(4 x^{3} + 2 x - 1\right).
\end{align*}

\textbf{Proof certificate.} Inserting the displayed data into \eqref{round6:eq:r6gf-cert} gives exactly

\[\mathcal L_pF=P(x)=- x \left(8 x^{3} - 3 x^{2} + 2 x - 1\right).\]

The radical residual is identically zero as a rational function. The rational residual equals the displayed $P$; both assertions are verified by clearing denominators in \eqref{round6:eq:r6gf-cert}. Taking coefficients in \eqref{round6:eq:r6gf-extract} proves the recurrence for every $n\ge5$, with the initial residuals exactly as stated.

\subsubsection{A026030: Jun 15 2014}
\label{round6:r6entry:A026030}

Source: \url{https://oeis.org/A026030}. Selected conjecture: R. J. Mathar, Jun 15 2014. Stephan's 2005 binomial formula and Somos's 2012 generating function are prior input. Drube's 2022 raised-path enumeration is also credited. The displayed first generating function uses a shifted indexing convention; the Somos formula used here starts at the source index one.

\textbf{Statement.} With the source offset $o=1$, the following recurrence holds for every integer $n\ge5$.

\begin{align*}
p_{0}(n)&=n + 5\\
p_{1}(n)&=- 2 \left(5 n + 16\right)\\
p_{2}(n)&=35 n + 47\\
p_{3}(n)&=- 2 \left(25 n - 14\right)\\
p_{4}(n)&=12 \left(2 n - 5\right).
\end{align*}

Use $t=x$ and $\xi=x$.

The credited input has the exact form $F=R_0+R_1\sqrt D$, with

\begin{align*}
R_0&=\frac{\left(2 x - 1\right)^{2} \left(x^{2} - 4 x + 1\right)}{2 x^{5}},\\
R_1&=\frac{\left(x - 1\right) \left(2 x - 1\right) \left(3 x - 1\right)}{2 x^{5}},\\
D&=1 - 4 x.
\end{align*}

\textbf{Proof certificate.} Inserting the displayed data into \eqref{round6:eq:r6gf-cert} gives exactly

\[\mathcal L_pF=P(x)=6 x \left(2 x - 1\right)^{2}.\]

The radical residual is identically zero as a rational function. The rational residual equals the displayed $P$; both assertions are verified by clearing denominators in \eqref{round6:eq:r6gf-cert}. Taking coefficients in \eqref{round6:eq:r6gf-extract} proves the recurrence for every $n\ge5$, with the initial residuals exactly as stated.

\subsubsection{A026571: Jun 23 2013}
\label{round6:r6entry:A026571}

Source: \url{https://oeis.org/A026571}. Selected conjecture: R. J. Mathar, Jun 23 2013. Gomez Aiza's 2024 Trees with flowers paper, pp. 18--19, supplies the generating function and related asymptotics. The selected target is the 2013 recurrence. The generating function already starts at degree two; no extra shift is made.

\textbf{Statement.} With the source offset $o=2$, the following recurrence holds for every integer $n\ge7$.

\begin{align*}
p_{0}(n)&=n + 2\\
p_{1}(n)&=- 3 \left(n + 1\right)\\
p_{2}(n)&=- 3 n\\
p_{3}(n)&=11 \left(n - 1\right)\\
p_{4}(n)&=2 \left(n - 6\right)\\
p_{5}(n)&=- 4 \left(2 n - 7\right).
\end{align*}

Use $t=x$ and $\xi=x$.

The credited input has the exact form $F=R_0+R_1\sqrt D$, with

\begin{align*}
R_0&=- \frac{1}{2 x^{2}},\\
R_1&=\frac{\left(x + 1\right) \left(2 x - 1\right)}{2 x^{2} \left(4 x^{2} + x - 1\right)},\\
D&=\frac{4 x^{2} + x - 1}{x - 1}.
\end{align*}

\textbf{Proof certificate.} Inserting the displayed data into \eqref{round6:eq:r6gf-cert} gives exactly

\[\mathcal L_pF=P(x)=- 2 x^{2} \left(x - 2\right).\]

The radical residual is identically zero as a rational function. The rational residual equals the displayed $P$; both assertions are verified by clearing denominators in \eqref{round6:eq:r6gf-cert}. Taking coefficients in \eqref{round6:eq:r6gf-extract} proves the recurrence for every $n\ge7$, with the initial residuals exactly as stated.

\subsubsection{A026743: Jan 13 2023}
\label{round6:r6entry:A026743}

Source: \url{https://oeis.org/A026743}. Selected conjecture: R. J. Mathar, Jan 13 2023. Callan's January 17, 2016 generating function is the prior input. The selected target is the 2023 recurrence.

\textbf{Statement.} With the source offset $o=0$, the following recurrence holds for every integer $n\ge8$.

\begin{align*}
p_{0}(n)&=n\\
p_{1}(n)&=-2\\
p_{2}(n)&=20 - 11 n\\
p_{3}(n)&=14\\
p_{4}(n)&=39 n - 152\\
p_{5}(n)&=-22\\
p_{6}(n)&=268 - 41 n\\
p_{7}(n)&=-6\\
p_{8}(n)&=- 12 \left(n - 6\right).
\end{align*}

Use $t=x$ and $\xi=x$.

The credited input has the exact form $F=R_0+R_1\sqrt D$, with

\begin{align*}
R_0&=- \frac{\left(2 x + 1\right) \left(x^{2} + 1\right)}{2 \left(x^{4} + 4 x^{2} - 1\right)},\\
R_1&=\frac{3 x^{2} - 1}{2 \left(x^{4} + 4 x^{2} - 1\right)},\\
D&=- \frac{2 x + 1}{2 x - 1}.
\end{align*}

\textbf{Proof certificate.} Inserting the displayed data into \eqref{round6:eq:r6gf-cert} gives exactly

\[\mathcal L_pF=P(x)=2 x^{2} \left(2 x + 1\right) \left(3 x^{4} - 6 x^{2} + 1\right).\]

The radical residual is identically zero as a rational function. The rational residual equals the displayed $P$; both assertions are verified by clearing denominators in \eqref{round6:eq:r6gf-cert}. Taking coefficients in \eqref{round6:eq:r6gf-extract} proves the recurrence for every $n\ge8$, with the initial residuals exactly as stated.

\subsubsection{A052183: Sep 27 2012}
\label{round6:r6entry:A052183}

Source: \url{https://oeis.org/A052183}. Selected conjecture: R. J. Mathar, Sep 27 2012. The factorial formula in the sequence definition is the prior input; it is related to Carlitz's two-line array enumeration and the Fuss--Catalan numbers. Only the explicitly conjectured three-term formula is selected.

\textbf{Statement.} With the source offset $o=0$, the following recurrence holds for every integer $n\ge2$.

\begin{align*}
p_{0}(n)&=2 n \left(2 n + 1\right)\\
p_{1}(n)&=- 47 n^{2} + 50 n - 12\\
p_{2}(n)&=15 \left(3 n - 5\right) \left(3 n - 4\right).
\end{align*}

\textbf{Proof.} The defining factorial formula is
\[
a_n=\frac{n+2}{2n+1}\binom{3n}{n},\qquad n\ge0.
\]
Cancellation of factorials gives
\[
r(n)=\frac{a_{n-1}}{a_n}
=\frac{2n(n+1)(2n+1)}{3(n+2)(3n-1)(3n-2)},\qquad n\ge1.
\]
The identity $p_0+p_1r(n)+p_2r(n)r(n-1)=0$ reduces, after division
by $2n(2n+1)$ and clearing denominators, to
\begin{align*}
&3(n+2)(3n-1)(3n-2)\\
&\quad +(n+1)(-47n^2+50n-12)
+10n(n-1)(2n-1)=0.
\end{align*}
Expansion verifies this polynomial identity term by term. The cleared
denominators are nonzero for every integer $n\ge2$. Multiplication
by $a_n$ proves the recurrence throughout that range. Its first case
also reads $20\cdot12-100\cdot3+30\cdot2=0$.


\subsubsection{A095981: Nov 23 2014}
\label{round6:r6entry:A095981}

Source: \url{https://oeis.org/A095981}. Selected conjecture: R. J. Mathar, Nov 23 2014. Callan's known first-return decomposition and its generating function are the prior input. The empty path is excluded by this entry's plateau convention, so its initial values are zero, zero, one.

\textbf{Statement.} With the source offset $o=0$, the following recurrence holds for every integer $n\ge7$.

\begin{align*}
p_{0}(n)&=n + 2\\
p_{1}(n)&=- 5 n - 4\\
p_{2}(n)&=6 n\\
p_{3}(n)&=4 n - 13\\
p_{4}(n)&=43 - 13 n\\
p_{5}(n)&=3 \left(3 n - 13\right)\\
p_{6}(n)&=4 - n\\
p_{7}(n)&=7 - n.
\end{align*}

Use $t=x$ and $\xi=x$.

The credited input has the exact form $F=R_0+R_1\sqrt D$, with

\begin{align*}
R_0&=- \frac{x^{3} - x^{2} - 2 x + 1}{2 x^{2} \left(x - 1\right)},\\
R_1&=\frac{1}{2 x^{2} \left(x - 1\right)},\\
D&=\left(x^{3} - x^{2} + 1\right) \left(x^{3} + 3 x^{2} - 4 x + 1\right).
\end{align*}

\textbf{Proof certificate.} Inserting the displayed data into \eqref{round6:eq:r6gf-cert} gives exactly

\[\mathcal L_pF=P(x)=x^{2} \left(x^{4} - x^{3} + 6 x^{2} - 9 x + 4\right).\]

The radical residual is identically zero as a rational function. The rational residual equals the displayed $P$; both assertions are verified by clearing denominators in \eqref{round6:eq:r6gf-cert}. Taking coefficients in \eqref{round6:eq:r6gf-extract} proves the recurrence for every $n\ge7$, with the initial residuals exactly as stated.

\subsubsection{A102880: Nov 09 2012}
\label{round6:r6entry:A102880}

Source: \url{https://oeis.org/A102880}. Selected conjecture: R. J. Mathar, Nov 09 2012. The defining Chebyshev transform of the Catalan series is prior input. The initial residual below is retained explicitly.

\textbf{Statement.} With the source offset $o=0$, the following recurrence holds for every integer $n\ge5$.

\textbf{Boundary qualification.} The source specifies no starting index. $n=4$ gives residual $3$. These are all failures between the first index with all source terms defined and the proved tail. The coefficient formula itself is unchanged.

\begin{align*}
p_{0}(n)&=\left(n - 3\right) \left(n + 1\right)\\
p_{1}(n)&=- 2 \left(n - 3\right) \left(2 n - 1\right)\\
p_{2}(n)&=2 \left(n^{2} - 4 n + 1\right)\\
p_{3}(n)&=- 2 \left(n - 1\right) \left(2 n - 7\right)\\
p_{4}(n)&=\left(n - 5\right) \left(n - 1\right).
\end{align*}

Use $t=x$ and $\xi=x$.

The credited input has the exact form $F=R_0+R_1\sqrt D$, with

\begin{align*}
R_0&=- \frac{\left(x - 1\right) \left(x + 1\right)}{2 x},\\
R_1&=\frac{\left(x - 1\right) \left(x + 1\right)}{2 x},\\
D&=\frac{x^{2} - 4 x + 1}{x^{2} + 1}.
\end{align*}

\textbf{Proof certificate.} Inserting the displayed data into \eqref{round6:eq:r6gf-cert} gives exactly

\[\mathcal L_pF=P(x)=3 \left(x - 1\right) \left(x + 1\right) \left(x^{2} + 1\right).\]

The radical residual is identically zero as a rational function. The rational residual equals the displayed $P$; both assertions are verified by clearing denominators in \eqref{round6:eq:r6gf-cert}. Taking coefficients in \eqref{round6:eq:r6gf-extract} proves the recurrence for every $n\ge5$, with the initial residuals exactly as stated.

\subsubsection{A105864: Nov 15 2011}
\label{round6:r6entry:A105864}

Source: \url{https://oeis.org/A105864}. Selected conjecture: R. J. Mathar, Nov 15 2011. Barry's 2009 transformation paper and the already listed Catalan transform are prior input. A 2017 FindSequenceFunction program on the entry is computational guessing; it is not used as an all-index proof.

\textbf{Statement.} With the source offset $o=0$, the following recurrence holds for every integer $n\ge4$.

\begin{align*}
p_{0}(n)&=n + 1\\
p_{1}(n)&=- 2 \left(2 n - 1\right)\\
p_{2}(n)&=- 2 \left(n - 1\right)\\
p_{3}(n)&=2 \left(2 n - 3\right)\\
p_{4}(n)&=n - 3.
\end{align*}

Use $t=x$ and $\xi=x$.

The credited input has the exact form $F=R_0+R_1\sqrt D$, with

\begin{align*}
R_0&=\frac{1}{2 x},\\
R_1&=- \frac{1}{2 x},\\
D&=\frac{x^{2} + 4 x - 1}{\left(x - 1\right) \left(x + 1\right)}.
\end{align*}

\textbf{Proof certificate.} Inserting the displayed data into \eqref{round6:eq:r6gf-cert} gives exactly

\[\mathcal L_pF=P(x)=x^{2} + 1.\]

The radical residual is identically zero as a rational function. The rational residual equals the displayed $P$; both assertions are verified by clearing denominators in \eqref{round6:eq:r6gf-cert}. Taking coefficients in \eqref{round6:eq:r6gf-extract} proves the recurrence for every $n\ge4$, with the initial residuals exactly as stated.

\subsubsection{A114584: Feb 16 2018}
\label{round6:r6entry:A114584}

Source: \url{https://oeis.org/A114584}. Selected conjecture: R. J. Mathar, Feb 16 2018. The established Motzkin-path generating function is prior input; the general forbidden-pattern treatment of Asinowski, Bacher, Banderier and Gittenberger (2019) is credited. Only the separately labeled recurrence is selected.

\textbf{Statement.} With the source offset $o=0$, the following recurrence holds for every integer $n\ge7$.

\begin{align*}
p_{0}(n)&=n + 2\\
p_{1}(n)&=n + 2\\
p_{2}(n)&=- 3 \left(3 n - 2\right)\\
p_{3}(n)&=13 - 7 n\\
p_{4}(n)&=4 n - 13\\
p_{5}(n)&=- 6 \left(n - 5\right)\\
p_{6}(n)&=n - 7\\
p_{7}(n)&=3 \left(n - 8\right).
\end{align*}

Use $t=x$ and $\xi=x$.

The credited input has the exact form $F=R_0+R_1\sqrt D$, with

\begin{align*}
R_0&=\frac{x^{3} - x + 1}{2 x^{2}},\\
R_1&=- \frac{1}{2 x^{2}},\\
D&=\left(x^{3} - 3 x + 1\right) \left(x^{3} + x + 1\right).
\end{align*}

\textbf{Proof certificate.} Inserting the displayed data into \eqref{round6:eq:r6gf-cert} gives exactly

\[\mathcal L_pF=P(x)=- 2 \left(3 x + 1\right) \left(2 x^{3} - 1\right).\]

The radical residual is identically zero as a rational function. The rational residual equals the displayed $P$; both assertions are verified by clearing denominators in \eqref{round6:eq:r6gf-cert}. Taking coefficients in \eqref{round6:eq:r6gf-extract} proves the recurrence for every $n\ge7$, with the initial residuals exactly as stated.

\subsubsection{A118093: Apr 05 2018}
\label{round6:r6entry:A118093}

Source: \url{https://oeis.org/A118093}. Selected conjecture: R. J. Mathar, Apr 05 2018. The rooted toroidal hypermap enumeration is prior work. Coserea's November 6, 2018 generating function and the already listed finite sum are credited input.

\textbf{Statement.} With the source offset $o=3$, the following recurrence holds for every integer $n\ge6$.

\begin{align*}
p_{0}(n)&=n \left(5 n - 17\right)\\
p_{1}(n)&=- 15 \left(n - 1\right) \left(5 n - 16\right)\\
p_{2}(n)&=12 \left(20 n^{2} - 103 n + 140\right)\\
p_{3}(n)&=32 \left(2 n - 5\right) \left(5 n - 12\right).
\end{align*}

Use $t=x$ and $\xi=x$.

The credited input has the exact form $F=R_0+R_1\sqrt D$, with

\begin{align*}
R_0&=- \frac{4 x^{2} - 7 x + 1}{8 \left(x + 1\right) \left(8 x - 1\right)},\\
R_1&=- \frac{3 x - 1}{8 \left(x + 1\right) \left(8 x - 1\right)},\\
D&=1 - 8 x.
\end{align*}

\textbf{Proof certificate.} Inserting the displayed data into \eqref{round6:eq:r6gf-cert} gives exactly

\[\mathcal L_pF=P(x)=- 6 x^{3}.\]

The radical residual is identically zero as a rational function. The rational residual equals the displayed $P$; both assertions are verified by clearing denominators in \eqref{round6:eq:r6gf-cert}. Taking coefficients in \eqref{round6:eq:r6gf-extract} proves the recurrence for every $n\ge6$, with the initial residuals exactly as stated.

\subsubsection{A121320: Jun 22 2016}
\label{round6:r6entry:A121320}

Source: \url{https://oeis.org/A121320/internal}. Selected conjecture: R. J. Mathar, Jun 22 2016. The plane-tree statistic generating function, reformulated by Georg Fischer on April 6, 2020, is prior input. This does not belong to Fischer's excluded A176604 parameter family.

\textbf{Statement.} With the source offset $o=0$, the following recurrence holds for every integer $n\ge4$.

\textbf{Boundary qualification.} The source specifies no starting index. $n=3$ gives residual $1$. These are all failures between the first index with all source terms defined and the proved tail. The coefficient formula itself is unchanged.

\begin{align*}
p_{0}(n)&=2 - n\\
p_{1}(n)&=5 n - 12\\
p_{2}(n)&=8 - 3 n\\
p_{3}(n)&=- 2 \left(2 n - 5\right).
\end{align*}

Use $t=x$ and $\xi=x$.

The credited input has the exact form $F=R_0+R_1\sqrt D$, with

\begin{align*}
R_0&=- \frac{x^{2}}{2 \left(x^{2} + x - 1\right)},\\
R_1&=\frac{x^{2}}{2 \left(4 x - 1\right) \left(x^{2} + x - 1\right)},\\
D&=1 - 4 x.
\end{align*}

\textbf{Proof certificate.} Inserting the displayed data into \eqref{round6:eq:r6gf-cert} gives exactly

\[\mathcal L_pF=P(x)=x^{3}.\]

The radical residual is identically zero as a rational function. The rational residual equals the displayed $P$; both assertions are verified by clearing denominators in \eqref{round6:eq:r6gf-cert}. Taking coefficients in \eqref{round6:eq:r6gf-extract} proves the recurrence for every $n\ge4$, with the initial residuals exactly as stated.

\subsubsection{A141351: Oct 15 2014}
\label{round6:r6entry:A141351}

Source: \url{https://oeis.org/A141351}. Selected conjecture: R. J. Mathar, Oct 15 2014. The Catalan-plus-one definition and its other combinatorial interpretations are prior input. Only the separately conjectured recurrence is selected.

\textbf{Statement.} With the source offset $o=0$, the following recurrence holds for every integer $n\ge4$.

\textbf{Boundary qualification.} The source specifies no starting index. $n=3$ gives residual $2$. These are all failures between the first index with all source terms defined and the proved tail. The coefficient formula itself is unchanged.

\begin{align*}
p_{0}(n)&=n + 1\\
p_{1}(n)&=- 2 \left(3 n - 1\right)\\
p_{2}(n)&=9 n - 13\\
p_{3}(n)&=- 2 \left(2 n - 5\right).
\end{align*}

Use $t=x$ and $\xi=x$.

The credited input has the exact form $F=R_0+R_1\sqrt D$, with

\begin{align*}
R_0&=- \frac{2 x^{2} - x + 1}{2 x \left(x - 1\right)},\\
R_1&=- \frac{1}{2 x},\\
D&=1 - 4 x.
\end{align*}

\textbf{Proof certificate.} Inserting the displayed data into \eqref{round6:eq:r6gf-cert} gives exactly

\[\mathcal L_pF=P(x)=2 x^{3} - 6 x^{2} + 1.\]

The radical residual is identically zero as a rational function. The rational residual equals the displayed $P$; both assertions are verified by clearing denominators in \eqref{round6:eq:r6gf-cert}. Taking coefficients in \eqref{round6:eq:r6gf-extract} proves the recurrence for every $n\ge4$, with the initial residuals exactly as stated.

\subsubsection{A157003: Nov 19 2014}
\label{round6:r6entry:A157003}

Source: \url{https://oeis.org/A157003}. Selected conjecture: R. J. Mathar, Nov 19 2014. Kotesovec supplied a different order-five recurrence on February 1, 2015. Barry's 2023 elliptic-curve paper also gives the known generating function. The selected target is the separately labeled Mathar order-six relation; no Hankel-transform or pattern-avoidance conjecture is claimed here.

\textbf{Statement.} With the source offset $o=0$, the following recurrence holds for every integer $n\ge6$.

\begin{align*}
p_{0}(n)&=n + 1\\
p_{1}(n)&=n + 2\\
p_{2}(n)&=29 - 21 n\\
p_{3}(n)&=3 n - 16\\
p_{4}(n)&=40 \left(n - 3\right)\\
p_{5}(n)&=- 2 \left(2 n - 7\right)\\
p_{6}(n)&=- 10 \left(2 n - 9\right).
\end{align*}

Use $t=x$ and $\xi=x$.

The credited input has the exact form $F=R_0+R_1\sqrt D$, with

\begin{align*}
R_0&=- \frac{1}{2 x \left(x - 1\right) \left(x + 1\right)},\\
R_1&=\frac{1}{2 x \left(x - 1\right) \left(x + 1\right)},\\
D&=4 x^{3} - 4 x + 1.
\end{align*}

\textbf{Proof certificate.} Inserting the displayed data into \eqref{round6:eq:r6gf-cert} gives exactly

\[\mathcal L_pF=P(x)=- \left(5 x + 1\right) \left(3 x^{2} - 1\right).\]

The radical residual is identically zero as a rational function. The rational residual equals the displayed $P$; both assertions are verified by clearing denominators in \eqref{round6:eq:r6gf-cert}. Taking coefficients in \eqref{round6:eq:r6gf-extract} proves the recurrence for every $n\ge6$, with the initial residuals exactly as stated.

\subsubsection{A159772: May 30 2014}
\label{round6:r6entry:A159772}

Source: \url{https://oeis.org/A159772}. Selected conjecture: R. J. Mathar, May 30 2014. The binary-tree enumeration and Howroyd's 2017 rooted-tree generating function are prior input. Kotesovec's December 5, 2017 order-four recurrence is explicitly credited. The selected target is Mathar's separately conjectured order-six relation.

\textbf{Statement.} With the source offset $o=1$, the following recurrence holds for every integer $n\ge7$.

\begin{align*}
p_{0}(n)&=2 \left(n - 1\right) \left(2 n - 3\right)\\
p_{1}(n)&=- 43 n^{2} + 172 n - 177\\
p_{2}(n)&=4 \left(44 n^{2} - 266 n + 411\right)\\
p_{3}(n)&=- 8 \left(43 n^{2} - 358 n + 741\right)\\
p_{4}(n)&=96 \left(3 n^{2} - 29 n + 69\right)\\
p_{5}(n)&=- 128 \left(n - 6\right) \left(n - 4\right)\\
p_{6}(n)&=512 \left(n - 7\right) \left(n - 6\right).
\end{align*}

Use $x=\xi(t)=\frac{t}{\left(t + 1\right) \left(t^{2} + 1\right)}$, with $t(0)=0$.

The credited input has the exact form $F=R_0+R_1\sqrt D$, with

\begin{align*}
R_0&=- \frac{t}{\left(t - 1\right) \left(t + 1\right) \left(t^{2} + 1\right)},\\
R_1&=0,\\
D&=1.
\end{align*}

\textbf{Proof certificate.} Inserting the displayed data into \eqref{round6:eq:r6gf-cert} gives exactly

\[\mathcal L_pF=P(x)=- x^{2} \left(32 x^{3} - 24 x^{2} - 12 x + 3\right).\]

The radical residual is identically zero as a rational function. The rational residual equals the displayed $P$; both assertions are verified by clearing denominators in \eqref{round6:eq:r6gf-cert}. Taking coefficients in \eqref{round6:eq:r6gf-extract} proves the recurrence for every $n\ge7$, with the initial residuals exactly as stated.

\subsubsection{A160823: Nov 16 2011}
\label{round6:r6entry:A160823}

Source: \url{https://oeis.org/A160823}. Selected conjecture: R. J. Mathar, Nov 16 2011. Barry's established large-Schroeder transform and its finite sum are prior input.

\textbf{Statement.} With the source offset $o=0$, the following recurrence holds for every integer $n\ge5$.

\begin{align*}
p_{0}(n)&=n + 2\\
p_{1}(n)&=- 3 \left(n + 1\right)\\
p_{2}(n)&=- 3 \left(n - 2\right)\\
p_{3}(n)&=11 n - 20\\
p_{4}(n)&=11 - 5 n\\
p_{5}(n)&=4 - n.
\end{align*}

Use $t=x$ and $\xi=x$.

The credited input has the exact form $F=R_0+R_1\sqrt D$, with

\begin{align*}
R_0&=\frac{x^{2} + x - 1}{2 x^{2} \left(x - 1\right)},\\
R_1&=\frac{1}{2 x^{2} \left(x - 1\right)},\\
D&=x^{4} + 6 x^{3} - 5 x^{2} - 2 x + 1.
\end{align*}

\textbf{Proof certificate.} Inserting the displayed data into \eqref{round6:eq:r6gf-cert} gives exactly

\[\mathcal L_pF=P(x)=- \left(x - 2\right) \left(x^{2} - x + 1\right).\]

The radical residual is identically zero as a rational function. The rational residual equals the displayed $P$; both assertions are verified by clearing denominators in \eqref{round6:eq:r6gf-cert}. Taking coefficients in \eqref{round6:eq:r6gf-extract} proves the recurrence for every $n\ge5$, with the initial residuals exactly as stated.

\subsubsection{A162533: Nov 16 2011}
\label{round6:r6entry:A162533}

Source: \url{https://oeis.org/A162533}. Selected conjecture: R. J. Mathar, Nov 16 2011. Barry's October 13, 2009 radical generating function is prior input. Its first two printed forms agree. The final factorization on the entry has a separate typo: the factors must contain 1-2x, rather than 1-x. That typo is not silently used and is not counted as a second result.

\textbf{Statement.} With the source offset $o=0$, the following recurrence holds for every integer $n\ge5$.

\begin{align*}
p_{0}(n)&=n\\
p_{1}(n)&=4 - 5 n\\
p_{2}(n)&=2 \left(4 n - 7\right)\\
p_{3}(n)&=- 4 \left(n - 3\right)\\
p_{4}(n)&=- 4 \left(n - 2\right)\\
p_{5}(n)&=4 \left(n - 4\right).
\end{align*}

Use $t=x$ and $\xi=x$.

The credited input has the exact form $F=R_0+R_1\sqrt D$, with

\begin{align*}
R_0&=0,\\
R_1&=\frac{x - 1}{\left(2 x^{2} - 2 x + 1\right) \left(2 x^{2} + 2 x - 1\right)},\\
D&=- \left(2 x^{2} - 2 x + 1\right) \left(2 x^{2} + 2 x - 1\right).
\end{align*}

\textbf{Proof certificate.} Inserting the displayed data into \eqref{round6:eq:r6gf-cert} gives exactly

\[\mathcal L_pF=P(x)=0.\]

The radical residual is identically zero as a rational function. The rational residual equals the displayed $P$; both assertions are verified by clearing denominators in \eqref{round6:eq:r6gf-cert}. Taking coefficients in \eqref{round6:eq:r6gf-extract} proves the recurrence for every $n\ge5$, with the initial residuals exactly as stated.

\subsubsection{A174107: Nov 15 2011}
\label{round6:r6entry:A174107}

Source: \url{https://oeis.org/A174107}. Selected conjecture: R. J. Mathar, Nov 15 2011. The defining Catalan transform and Gladkovskii's June 27, 2013 radical form are prior input.

\textbf{Statement.} With the source offset $o=0$, the following recurrence holds for every integer $n\ge4$.

\begin{align*}
p_{0}(n)&=n + 1\\
p_{1}(n)&=- 2 \left(3 n - 1\right)\\
p_{2}(n)&=7 \left(n - 1\right)\\
p_{3}(n)&=- 2 \left(3 n - 5\right)\\
p_{4}(n)&=n - 3.
\end{align*}

Use $t=x$ and $\xi=x$.

The credited input has the exact form $F=R_0+R_1\sqrt D$, with

\begin{align*}
R_0&=\frac{1}{2 x},\\
R_1&=- \frac{1}{2 x},\\
D&=\frac{x^{2} - 5 x + 1}{x^{2} - x + 1}.
\end{align*}

\textbf{Proof certificate.} Inserting the displayed data into \eqref{round6:eq:r6gf-cert} gives exactly

\[\mathcal L_pF=P(x)=- \left(x - 1\right) \left(x + 1\right).\]

The radical residual is identically zero as a rational function. The rational residual equals the displayed $P$; both assertions are verified by clearing denominators in \eqref{round6:eq:r6gf-cert}. Taking coefficients in \eqref{round6:eq:r6gf-extract} proves the recurrence for every $n\ge4$, with the initial residuals exactly as stated.

\subsubsection{A174808: Feb 10 2015}
\label{round6:r6entry:A174808}

Source: \url{https://oeis.org/A174808}. Selected conjecture: R. J. Mathar, Feb 10 2015. The established large-Schroeder transform, including Somos's 2014 continued fraction, is prior input.

\textbf{Statement.} With the source offset $o=0$, the following recurrence holds for every integer $n\ge5$.

\begin{align*}
p_{0}(n)&=n + 1\\
p_{1}(n)&=4 - 5 n\\
p_{2}(n)&=13 - 11 n\\
p_{3}(n)&=- 3 \left(n - 1\right)\\
p_{4}(n)&=3 \left(n - 4\right)\\
p_{5}(n)&=n - 5.
\end{align*}

Use $t=x$ and $\xi=x$.

The credited input has the exact form $F=R_0+R_1\sqrt D$, with

\begin{align*}
R_0&=- \frac{x^{2} + x - 1}{2 x \left(x + 1\right)},\\
R_1&=- \frac{1}{2 x \left(x + 1\right)},\\
D&=x^{4} + 2 x^{3} - 5 x^{2} - 6 x + 1.
\end{align*}

\textbf{Proof certificate.} Inserting the displayed data into \eqref{round6:eq:r6gf-cert} gives exactly

\[\mathcal L_pF=P(x)=\left(2 x + 1\right) \left(x^{2} + x + 1\right).\]

The radical residual is identically zero as a rational function. The rational residual equals the displayed $P$; both assertions are verified by clearing denominators in \eqref{round6:eq:r6gf-cert}. Taking coefficients in \eqref{round6:eq:r6gf-extract} proves the recurrence for every $n\ge5$, with the initial residuals exactly as stated.

\subsubsection{A190169: Apr 09 2019}
\label{round6:r6entry:A190169}

Source: \url{https://oeis.org/A190169}. Selected conjecture: R. J. Mathar, Apr 09 2019. Deutsch's odd-level horizontal-step statistic generating function is prior input. This is a distinct statistic and A-ID from the previously treated even-level entry A190166.

\textbf{Statement.} With the source offset $o=0$, the following recurrence holds for every integer $n\ge8$.

\begin{align*}
p_{0}(n)&=- \left(n - 1\right) \left(n + 1\right)\\
p_{1}(n)&=- n \left(n - 19\right)\\
p_{2}(n)&=2 \left(n - 1\right) \left(7 n - 40\right)\\
p_{3}(n)&=- \left(n - 2\right) \left(17 n - 97\right)\\
p_{4}(n)&=2 \left(9 n^{2} - 64 n + 119\right)\\
p_{5}(n)&=- 17 \left(n - 5\right) \left(n - 4\right)\\
p_{6}(n)&=\left(n - 5\right) \left(19 n - 59\right)\\
p_{7}(n)&=- 2 \left(n - 6\right) \left(8 n - 21\right)\\
p_{8}(n)&=2 \left(n - 7\right) \left(2 n - 5\right).
\end{align*}

Use $t=x$ and $\xi=x$.

The credited input has the exact form $F=R_0+R_1\sqrt D$, with

\begin{align*}
R_0&=- \frac{1}{2 x},\\
R_1&=\frac{x^{4} - 2 x^{3} + x^{2} - 2 x + 1}{2 x \left(x^{2} - 3 x + 1\right) \left(x^{2} - x + 1\right) \left(x^{2} + x + 1\right)},\\
D&=\left(x^{2} - 3 x + 1\right) \left(x^{2} + x + 1\right).
\end{align*}

\textbf{Proof certificate.} Inserting the displayed data into \eqref{round6:eq:r6gf-cert} gives exactly

\[\mathcal L_pF=P(x)=- 8 x^{3}.\]

The radical residual is identically zero as a rational function. The rational residual equals the displayed $P$; both assertions are verified by clearing denominators in \eqref{round6:eq:r6gf-cert}. Taking coefficients in \eqref{round6:eq:r6gf-extract} proves the recurrence for every $n\ge8$, with the initial residuals exactly as stated.

\subsubsection{A208355: Aug 04 2015}
\label{round6:r6entry:A208355}

Source: \url{https://oeis.org/A208355}. Selected conjecture: R. J. Mathar, Aug 04 2015. The known duplicated-Catalan formula and Russell's January 2024 generating function are prior input.

\textbf{Statement.} With the source offset $o=0$, the following recurrence holds for every integer $n\ge2$.

\begin{align*}
p_{0}(n)&=- \left(n - 2\right) \left(n + 3\right)\\
p_{1}(n)&=-4\\
p_{2}(n)&=4 \left(n - 1\right)^{2}.
\end{align*}

Use $t=x$ and $\xi=x$.

The credited input has the exact form $F=R_0+R_1\sqrt D$, with

\begin{align*}
R_0&=- \frac{\left(x - 1\right) \left(2 x + 1\right)}{2 x^{3}},\\
R_1&=- \frac{x + 1}{2 x^{3}},\\
D&=- \left(2 x - 1\right) \left(2 x + 1\right).
\end{align*}

\textbf{Proof certificate.} Inserting the displayed data into \eqref{round6:eq:r6gf-cert} gives exactly

\[\mathcal L_pF=P(x)=6.\]

The radical residual is identically zero as a rational function. The rational residual equals the displayed $P$; both assertions are verified by clearing denominators in \eqref{round6:eq:r6gf-cert}. Taking coefficients in \eqref{round6:eq:r6gf-extract} proves the recurrence for every $n\ge2$, with the initial residuals exactly as stated.

\subsubsection{A228178: Aug 25 2013}
\label{round6:r6entry:A228178}

Source: \url{https://oeis.org/A228178}. Selected conjecture: R. J. Mathar, Aug 25 2013. The source offset is zero, but its displayed generating function begins with x. Its supplied PARI Vec program extracts the list beginning with that coefficient. Accordingly the proof identifies a(n) with the coefficient of degree n+1 in G, and divides G by x. This explicitly resolves the source's indexing convention; the separate 'Apparently partial sums' claim is not counted.

\textbf{Statement.} With the source offset $o=0$, the following recurrence holds for every integer $n\ge4$.

\begin{align*}
p_{0}(n)&=2 \left(n + 3\right)\\
p_{1}(n)&=- 2 \left(7 n + 11\right)\\
p_{2}(n)&=29 n + 7\\
p_{3}(n)&=19 - 21 n\\
p_{4}(n)&=2 \left(2 n - 5\right).
\end{align*}

Use $t=x$ and $\xi=x$.

The credited input has the exact form $F=R_0+R_1\sqrt D$, with

\begin{align*}
R_0&=- \frac{5 x^{2} - 8 x + 2}{2 x^{3} \left(x - 1\right)},\\
R_1&=\frac{x^{2} - 4 x + 2}{2 x^{3} \left(x - 1\right)},\\
D&=1 - 4 x.
\end{align*}

\textbf{Proof certificate.} Inserting the displayed data into \eqref{round6:eq:r6gf-cert} gives exactly

\[\mathcal L_pF=P(x)=5 x^{2} - 4 x + 6.\]

The radical residual is identically zero as a rational function. The rational residual equals the displayed $P$; both assertions are verified by clearing denominators in \eqref{round6:eq:r6gf-cert}. Taking coefficients in \eqref{round6:eq:r6gf-extract} proves the recurrence for every $n\ge4$, with the initial residuals exactly as stated.

\subsubsection{A243022: Jul 15 2017}
\label{round6:r6entry:A243022}

Source: \url{https://oeis.org/A243022}. Selected conjecture: R. J. Mathar, Jul 15 2017. The nested Catalan generating function in the definition is prior input. Rational parametrization below is used solely to certify the selected polynomial recurrence.

\textbf{Statement.} With the source offset $o=0$, the following recurrence holds for every integer $n\ge5$.

\begin{align*}
p_{0}(n)&=2 n \left(n - 1\right) \left(2 n - 3\right) \left(5 n - 12\right)\\
p_{1}(n)&=- \left(n - 1\right) \left(85 n^{3} - 459 n^{2} + 776 n - 400\right)\\
p_{2}(n)&=- 4 \left(40 n^{4} - 396 n^{3} + 1455 n^{2} - 2324 n + 1360\right)\\
p_{3}(n)&=4 \left(n - 4\right) \left(170 n^{3} - 1258 n^{2} + 3003 n - 2185\right)\\
p_{4}(n)&=16 \left(n - 5\right) \left(n - 4\right) \left(20 n^{2} - 118 n + 143\right)\\
p_{5}(n)&=- 272 \left(n - 6\right) \left(n - 5\right) \left(n - 4\right) \left(5 n - 7\right).
\end{align*}

Use $x=\xi(t)=\frac{t}{t^{2} + 1}$, with $t(0)=0$.

The credited input has the exact form $F=R_0+R_1\sqrt D$, with

\begin{align*}
R_0&=- \frac{1}{2 \left(t^{2} + 1\right)},\\
R_1&=- \frac{1}{2 \left(t^{2} + 1\right)},\\
D&=1 - 4 t.
\end{align*}

\textbf{Proof certificate.} Inserting the displayed data into \eqref{round6:eq:r6gf-cert} gives exactly

\[\mathcal L_pF=P(x)=4 x^{2} \left(23 x + 1\right).\]

The radical residual is identically zero as a rational function. The rational residual equals the displayed $P$; both assertions are verified by clearing denominators in \eqref{round6:eq:r6gf-cert}. Taking coefficients in \eqref{round6:eq:r6gf-extract} proves the recurrence for every $n\ge5$, with the initial residuals exactly as stated.

\subsubsection{A243255: Jul 15 2017}
\label{round6:r6entry:A243255}

Source: \url{https://oeis.org/A243255}. Selected conjecture: R. J. Mathar, Jul 15 2017. The nested Motzkin--Catalan generating function in the definition is prior input.

\textbf{Statement.} With the source offset $o=0$, the following recurrence holds for every integer $n\ge4$.

\begin{align*}
p_{0}(n)&=2 n \left(2 n + 1\right)\\
p_{1}(n)&=- 65 n^{2} + 119 n - 42\\
p_{2}(n)&=3 \left(97 n^{2} - 373 n + 342\right)\\
p_{3}(n)&=- 9 \left(n - 3\right) \left(23 n - 68\right)\\
p_{4}(n)&=- 567 \left(n - 4\right) \left(n - 3\right).
\end{align*}

Use $x=\xi(t)=\frac{t}{t^{2} + t + 1}$, with $t(0)=0$.

The credited input has the exact form $F=R_0+R_1\sqrt D$, with

\begin{align*}
R_0&=\frac{1}{2},\\
R_1&=- \frac{1}{2},\\
D&=1 - 4 t.
\end{align*}

\textbf{Proof certificate.} Inserting the displayed data into \eqref{round6:eq:r6gf-cert} gives exactly

\[\mathcal L_pF=P(x)=- 6 x \left(4 x - 1\right).\]

The radical residual is identically zero as a rational function. The rational residual equals the displayed $P$; both assertions are verified by clearing denominators in \eqref{round6:eq:r6gf-cert}. Taking coefficients in \eqref{round6:eq:r6gf-extract} proves the recurrence for every $n\ge4$, with the initial residuals exactly as stated.

\subsubsection{A254314: Jun 07 2016}
\label{round6:r6entry:A254314}

Source: \url{https://oeis.org/A254314}. Selected conjecture: R. J. Mathar, Jun 07 2016. The established quadratic generating function, tree interpretation and Gomez Aiza's 2024 asymptotic analysis are prior input. Only the separately conjectured recurrence is selected.

\textbf{Statement.} With the source offset $o=0$, the following recurrence holds for every integer $n\ge5$.

\begin{align*}
p_{0}(n)&=n\\
p_{1}(n)&=- 2 \left(5 n - 6\right)\\
p_{2}(n)&=2 \left(17 n - 39\right)\\
p_{3}(n)&=- 6 \left(8 n - 27\right)\\
p_{4}(n)&=25 n - 114\\
p_{5}(n)&=- 2 \left(n - 6\right).
\end{align*}

Use $t=x$ and $\xi=x$.

The credited input has the exact form $F=R_0+R_1\sqrt D$, with

\begin{align*}
R_0&=- \frac{x^{2} - 6 x + 3}{2 \left(2 x - 1\right)},\\
R_1&=\frac{1}{2 \left(2 x - 1\right)},\\
D&=\left(x - 1\right) \left(x^{3} - 11 x^{2} + 7 x - 1\right).
\end{align*}

\textbf{Proof certificate.} Inserting the displayed data into \eqref{round6:eq:r6gf-cert} gives exactly

\[\mathcal L_pF=P(x)=- x \left(12 x^{3} - 21 x^{2} + 14 x - 3\right).\]

The radical residual is identically zero as a rational function. The rational residual equals the displayed $P$; both assertions are verified by clearing denominators in \eqref{round6:eq:r6gf-cert}. Taking coefficients in \eqref{round6:eq:r6gf-extract} proves the recurrence for every $n\ge5$, with the initial residuals exactly as stated.

\subsubsection{A257388: Sep 24 2016}
\label{round6:r6entry:A257388}

Source: \url{https://oeis.org/A257388}. Selected conjecture: R. J. Mathar, Sep 24 2016. The known colored Motzkin-path generating function and its quadratic equation, also listed by Gutkovskiy in 2020, are prior input.

\textbf{Statement.} With the source offset $o=0$, the following recurrence holds for every integer $n\ge3$.

\begin{align*}
p_{0}(n)&=n + 2\\
p_{1}(n)&=- 8 \left(n + 1\right)\\
p_{2}(n)&=4 \left(3 n + 1\right)\\
p_{3}(n)&=8 \left(2 n - 3\right).
\end{align*}

Use $t=x$ and $\xi=x$.

The credited input has the exact form $F=R_0+R_1\sqrt D$, with

\begin{align*}
R_0&=\frac{1}{2 x^{2}},\\
R_1&=\frac{1}{2 x^{2} \left(4 x - 1\right)},\\
D&=\left(4 x - 1\right) \left(4 x^{2} + 4 x - 1\right).
\end{align*}

\textbf{Proof certificate.} Inserting the displayed data into \eqref{round6:eq:r6gf-cert} gives exactly

\[\mathcal L_pF=P(x)=- 2 \left(2 x - 1\right).\]

The radical residual is identically zero as a rational function. The rational residual equals the displayed $P$; both assertions are verified by clearing denominators in \eqref{round6:eq:r6gf-cert}. Taking coefficients in \eqref{round6:eq:r6gf-extract} proves the recurrence for every $n\ge3$, with the initial residuals exactly as stated.

\subsubsection{A278023: Jun 24 2018}
\label{round6:r6entry:A278023}

Source: \url{https://oeis.org/A278023}. Selected conjecture: R. J. Mathar, Jun 24 2018. The generating function printed in the sequence definition is prior input. Baril's irreducible-permutation interpretation is credited; the initial residual below qualifies the selected recurrence.

\textbf{Statement.} With the source offset $o=0$, the following recurrence holds for every integer $n\ge5$.

\textbf{Boundary qualification.} The source specifies no starting index. $n=4$ gives residual $156$. These are all failures between the first index with all source terms defined and the proved tail. The coefficient formula itself is unchanged.

\begin{align*}
p_{0}(n)&=n \left(3 n^{2} - 12 n + 11\right)\\
p_{1}(n)&=- \left(3 n - 5\right) \left(3 n^{2} - 9 n + 4\right)\\
p_{2}(n)&=- 2 \left(2 n - 5\right) \left(3 n^{2} - 6 n + 2\right)\\
p_{3}(n)&=n \left(3 n^{2} - 12 n + 11\right)\\
p_{4}(n)&=- 2 \left(2 n - 5\right) \left(3 n^{2} - 6 n + 2\right).
\end{align*}

Use $t=x$ and $\xi=x$.

The credited input has the exact form $F=R_0+R_1\sqrt D$, with

\begin{align*}
R_0&=- \frac{2 x^{3} + x + 1}{2 \left(x^{3} + x + 1\right)},\\
R_1&=\frac{\left(x - 1\right) \left(2 x + 1\right)}{2 \left(4 x - 1\right) \left(x^{3} + x + 1\right)},\\
D&=1 - 4 x.
\end{align*}

\textbf{Proof certificate.} Inserting the displayed data into \eqref{round6:eq:r6gf-cert} gives exactly

\[\mathcal L_pF=P(x)=2 x \left(78 x^{3} - 3 x^{2} - x + 1\right).\]

The radical residual is identically zero as a rational function. The rational residual equals the displayed $P$; both assertions are verified by clearing denominators in \eqref{round6:eq:r6gf-cert}. Taking coefficients in \eqref{round6:eq:r6gf-extract} proves the recurrence for every $n\ge5$, with the initial residuals exactly as stated.

\subsubsection{A385252: Jul 31 2025}
\label{round6:r6entry:A385252}

Source: \url{https://oeis.org/A385252}. Selected conjecture: R. J. Mathar, Jul 31 2025. The known ternary-string generating function and finite-sum formula are prior input. The definition's exceptional value a(0)=0 is retained and is essential to the boundary residual.

\textbf{Statement.} With the source offset $o=0$, the following recurrence holds for every integer $n\ge6$.

\textbf{Boundary qualification.} The source specifies no starting index. $n=5$ gives residual $-6912$. These are all failures between the first index with all source terms defined and the proved tail. The coefficient formula itself is unchanged.

\begin{align*}
p_{0}(n)&=n\\
p_{1}(n)&=36 - 37 n\\
p_{2}(n)&=4 \left(131 n - 245\right)\\
p_{3}(n)&=- 16 \left(221 n - 605\right)\\
p_{4}(n)&=192 \left(59 n - 213\right)\\
p_{5}(n)&=- 6912 \left(2 n - 9\right).
\end{align*}

Use $t=x$ and $\xi=x$.

The credited input has the exact form $F=R_0+R_1\sqrt D$, with

\begin{align*}
R_0&=\frac{72 x^{2} - 12 x + 1}{2 \left(4 x - 1\right) \left(9 x - 1\right)},\\
R_1&=\frac{12 x - 1}{2 \left(8 x - 1\right) \left(9 x - 1\right)},\\
D&=1 - 8 x.
\end{align*}

\textbf{Proof certificate.} Inserting the displayed data into \eqref{round6:eq:r6gf-cert} gives exactly

\[\mathcal L_pF=P(x)=- 16 x^{2} \left(432 x^{3} - 360 x^{2} + 69 x - 4\right).\]

The radical residual is identically zero as a rational function. The rational residual equals the displayed $P$; both assertions are verified by clearing denominators in \eqref{round6:eq:r6gf-cert}. Taking coefficients in \eqref{round6:eq:r6gf-extract} proves the recurrence for every $n\ge6$, with the initial residuals exactly as stated.

\subsection{Verification and the precise prior-work audit}
The portable verifier checks twenty-nine exact differential identities
and one exact hypergeometric identity. It parses the retained
primary-source mathematical statements and verifies every recurrence
coefficient. It independently expands the input generating functions
without using the target recurrence, matching twenty source terms per
entry, or 600 terms in total. Additional finite-sum and convolution
checks guard several defining inputs. For the three parameters it
uses Lagrange inversion,
$[x^n]t(x)=\frac1n[u^{n-1}]\phi(u)^n$ when $t=x\phi(t)$.
These numerical checks guard branch and indexing choices; the rational
identities are the all-index proofs.

All thirty IDs are disjoint from the previous 508 selected IDs,
from the 85-entry parameter table in Georg Fischer's January 26, 2020
general proof for A176604, from the 33 explicit Ekhad--Yang 2017 proof
IDs, and from the retained 97-ID OEIS Open Appendix A baseline.
The latter represents the 100 costliest Appendix A claims among 153
combined-model resolutions, not all results of the OEIS Open project.
The same finite baseline limitation applies to the accompanying JSON
audit. No completeness of a global literature search is claimed.

Primary-entry proof notices were inspected. A396722 was excluded
because its entry contains Rodriguez's September 28, 2026 proof;
A000781 was excluded because its entry contains Blanco Gomez's
October 1, 2026 proof. A217333 was rejected as a shifted alias of the
already proved A025273 recurrence, and A080243 was conservatively
excluded as a signed little-Schroeder alias. None contributes to the
thirty. The previously disclosed Round IV Fischer-family priority
correction remains in force; no member of that 85-entry family is
counted here.

The sources linked from the selected entries were checked selectively:
Gomez Aiza's 2024 \emph{Trees with flowers}, Drube's 2022
\emph{Raised k-Dyck paths}, Barry's 2009 transformation paper and 2023
elliptic-curve paper, the 2019 forbidden-pattern paper of Asinowski
and coauthors, Rowland's binary-tree paper, Baril's irreducible
permutation paper, and the 2023 parking-function paper of Adeniran
and Pudwell. Their known enumerations or generating functions are
credited above. Targeted searches for the selected IDs or Mathar's
name in those documents did not identify an explicit proof of these
selected recurrence formulas. This restricted observation is not an
absence-of-prior-proof theorem. Several broad web queries returned
irrelevant results and are not treated as evidence of novelty.
In particular A025268, A157003 and A159772 already have different
polynomial recurrences on their OEIS pages. This collection proves
the separately labeled formulas from their established inputs; it
does not portray those sequences as previously unknown.



\Needspace{6\baselineskip}
\section{Forty-one empirical array recurrences}
\label{round6:sec:r6arrays}
This section proves forty-one different entry-level recurrence assertions
explicitly marked \emph{Empirical} on their primary OEIS pages, consulted
on 20 October 2026.  Twenty-five displayed empirical or conjectural
rational generating functions are proved within the same packages.
The source definitions and conjectures are due to R.~H.~Hardin; the
additional generating-function conjectures are credited on the relevant
pages to Colin Barker.  The forty-one A-numbers are absent from the 508
entries selected in rounds I--V.  Several entries share a transfer
principle: forty-one is the number of selected OEIS assertions, not the
number of unrelated proof techniques.

\subsection{Local conditions and the counting model}
All arrays have $n+1$ rows, fixed width $w\ge2$, and entries in
$\{0,\ldots,q-1\}$. For a block
$B=\left(\begin{smallmatrix}a&b\\c&d\end{smallmatrix}\right)$ put
\[
 \Delta(B)=(b-a,d-b,c-d,a-c).
\]
These differences are ordinary signed integers.  The selected local
conditions are the following; $S$ and $s$ are specified for each entry.
\begin{description}
\item[Cell values (C).] $\#\{a,b,c,d\}\in S$.
\item[Stress (D).] $|a+d-b-c|=s$.  When the source prohibits equal
adjacent elements, additionally every entry of $\Delta(B)$ is nonzero.
\item[Absolute edge sum (U).] $\sum_{e\in\Delta(B)}|e|=s$ and every
edge difference is nonzero. The sum includes all four edges with
multiplicity.
\item[Squared edge sum (V).] $\sum_{e\in\Delta(B)}e^2=s$ and every
edge difference is nonzero, again including all four edges.
\item[One permitted orientation (J).] Exactly one of
$\max\Delta(B)\le1$ and $\min\Delta(B)\ge-1$ holds. These are the
clockwise and counterclockwise bounds in the source definition.
\end{description}
Every horizontal or vertical neighboring pair of a selected array
belongs to a $2\times2$ block, since $n+1\ge2$ and $w\ge2$.
Thus the blockwise nonzero condition enforces precisely the source's
prohibition of equal adjacent elements, including at the outer boundary.
Diagonal neighbors are not subject to that prohibition.

\subsection{Definition-derived transfer matrices}
For C, canonicalize a row by renaming its colors in order of first
appearance.  For D, U, V and J, canonicalize under the four operations
consisting of horizontal reflection, simultaneous complementation
$x\mapsto q-1-x$, their product, and the identity.
Write $\gamma$ for the appropriate canonicalization and
$\mathcal O_{q,w}$ for the lexicographically ordered representative rows.
For each representative $r$ define its actual orbit mass
\begin{equation}\label{round6:eq:r6arr-mass}
 m(r)=\#\{v\in\{0,\ldots,q-1\}^w:\gamma(v)=r\}.
\end{equation}
For the appropriate local predicate $P$, define the nonnegative integer
matrix
\begin{equation}\label{round6:eq:r6arr-T}
 T(r,s)=\#\left\{v\in\{0,\ldots,q-1\}^{w}:
 \begin{array}{l}
 \gamma(v)=s,\\
 P(r_i,r_{i+1},v_i,v_{i+1})\quad(1\le i<w)
 \end{array}\right\}.
\end{equation}
For C, any simultaneous color permutation preserves $P$.
For the other predicates, horizontal reflection and complementation
permute and negate the edge differences and preserve the absolute
stress. They preserve U and V and interchange the two bounds in J,
preserving their exclusive disjunction. Consequently every row in an
orbit has the same number of successors in each orbit. No arbitrary
color permutation is used for D, U, V or J.

Adding rows, starting with the actual $m(r)$ labeled first rows, therefore
gives the exact labeled count
\begin{equation}\label{round6:eq:r6arr-labeled}
 L_q(n)=m^{\mathsf T}T^n\mathbf1\qquad(n\ge1).
\end{equation}
The exponent is $n$ because an array has $n+1$ rows. Fixed representatives
and smaller orbits retain their correct masses. In particular, a row
that cannot extend to a legal array simply has no outgoing edges.
For a scalar-normalized source definition, $a(n)=\lambda L_q(n)$,
with precisely the printed scalar $\lambda$; no orbit interpretation of
that scalar is assumed.

Seven C entries additionally identify global color renamings by row-major
first-occurrence order. Let $U_q^{=}$ count such arrays using exactly
$q$ colors, and $U_q^{\le}$ those using at most $q$ colors.
Surjective labelings have $q!$ different color renamings. Inclusion--exclusion
on omitted colors and then summation over the number actually used give
\begin{align}
 U_q^{=}(n)&=\sum_{j=1}^{q}\frac{(-1)^{q-j}}{j!(q-j)!}L_j(n),
 \label{round6:eq:r6arr-exact}\\
 U_q^{\le}(n)&=\sum_{j=1}^{q}
 \left(\frac1{j!}\sum_{h=0}^{q-j}\frac{(-1)^h}{h!}\right)L_j(n).
 \label{round6:eq:r6arr-atmost}
\end{align}
The empty alphabet contributes zero. A requirement to use every color
is imposed only when the source explicitly includes it. In the other
four C entries, the source normalization is simply $L_4/24$ or $L_3/6$.
These constructions adapt the project's earlier color-orbit and
geometric-orbit transfer methods; the present local predicates and
selected source assertions are specified afresh here.

\subsection{Exact certificates and the all-index conclusion}
For each printed coefficient signature $c=(c_1,\ldots,c_d)$ set
\[
 R(z)=z^d-\sum_{i=1}^d c_i z^{d-i},\qquad
 Q(x)=1-\sum_{i=1}^d c_i x^i.
\]
For every matrix in that entry with a nonzero scalar weight, the complete
exact integer certificate establishes
\begin{equation}\label{round6:eq:r6arr-cert}
 \boxed{T R(T)\mathbf1=\mathbf0.}
\end{equation}
For $n\ge d+1$, multiply this identity by $m^{\mathsf T}T^{n-d-1}$
and by the scalar weight, then add the weighted components.
Equations \eqref{round6:eq:r6arr-labeled}--\eqref{round6:eq:r6arr-atmost} prove
\begin{equation}\label{round6:eq:r6arr-recurrence}
 \boxed{a(n)=\sum_{i=1}^{d}c_i a(n-i)\qquad(n\ge d+1).}
\end{equation}
All source indices on the right are at least 1. The auxiliary value
obtained by setting $n=0$ in the transfer formula is not an asserted
OEIS term and may be nonintegral under a scalar normalization.

The certificate file contains the forty-one signatures, every matrix
component's exact weight, all vectors $v_k=T^k\mathbf1$ for
$0\le k\le d+1$, and every residual coordinate in
\eqref{round6:eq:r6arr-cert}. There are 49 distinct complete matrices and
6,105 nonzero weighted edges; the largest matrix has 342 states and the
largest selected recurrence has order 31. The reproducible exact check is
\begin{Verbatim}[fontsize=\small]
v[0] = [1 for every state]
for k = 1,...,d+1:
    v[k][r] = sum(T[r,s] * v[k-1][s] for every state s)
for every state r:
    assert v[d+1][r] == sum(c[i] * v[d+1-i][r] for i=1,...,d)
\end{Verbatim}
The verifier rebuilds every state and transition directly from the
finite definition \eqref{round6:eq:r6arr-T}, compares the stored matrices and
all certificate coordinates, and uses only exact integers and fractions.
This is a computer-assisted finite proof: further powers of $T$, not
agreement with many numerical terms, establish the universal conclusion.
The matrix definitions and all coefficient signatures below specify
the entire finite calculation. The complete expanded certificates are
supplied in \texttt{r6\_arrays\_matrices.json} and
\texttt{r6\_arrays\_certificates.json}.

For $F(x)=\sum_{n\ge1}a(n)x^n$ the proved recurrence yields $QF=P$,
where $\deg P\le d$ and
\begin{equation}\label{round6:eq:r6arr-numerator}
 [x^k]P=\sum_{\substack{0\le i\le\min(k,d)\\k-i\ge1}}
 q_i a(k-i),\qquad q_0=1,\quad q_i=-c_i\ (i\ge1).
\end{equation}
Every initial value in this expression is calculated from the defining
transfer matrix. For the 25 source-stated rational generating functions,
exact polynomial cross-multiplication proves equality with this $P/Q$;
the relevant numerator $P$ is also displayed below. The source's original
rational expression is retained in the statement and certificate files.
No generating function adds a second count to its A-number.
Separately stated radical closed forms are not additional selected claims.
All 832 printed source-prefix terms agree with the independent matrix
counts; that comparison checks indexing and interpretation, not the
universal recurrence. The forty-one computed 20-term prefixes are distinct.

\subsection{A short hand-checkable example: A210406}
Here $q=4$, $w=2$ and $S=\{3,4\}$. A legal array uses at least three
colors, so each orbit under four-color permutations has size $24$.
For the two row representatives $00,01$, direct enumeration gives
\[
 T=\begin{pmatrix}0&6\\2&10\end{pmatrix},\qquad
 m=\begin{pmatrix}4\\12\end{pmatrix},\qquad
 a(n)=\frac1{24}m^{\mathsf T}T^n\mathbf1\quad(n\ge1).
\]
The identity $T^2=10T+12I$ proves $a(n)=10a(n-1)+12a(n-2)$
for $n\ge3$. Also $a(1)=7$ and $a(2)=78$, so
\[
 F(x)=\frac{x(7+8x)}{1-10x-12x^2}.
\]
This proves precisely the empirical recurrence and the conjectural
rational generating function on \url{https://oeis.org/A210406}.
The larger selected matrices use the same exact finite argument.

\subsection{Source audit and counting limits}
The status relevant to selection is the word \emph{Empirical} attached
to each recurrence, not OEIS's editorial label ``approved.'' The
consulted primary pages link numerical data tables and do not record
a proof of their selected recurrence. Representative exact-A-number
and family proof searches found no such proof among the inspected
results. This is a limited primary-source audit, not a global-priority
claim or a proof that no prior proof exists.

All forty-one IDs are disjoint from the 508 previously selected project
entries, and from the retained 97 distinct A-numbers in the 100 costly
claims in Appendix A of \emph{OEIS Open}. That appendix check does not
cover all repository results or all mathematical literature. Four larger
certificate-verified candidates and two further sourced candidates were
left outside this forty-one-entry selection and add no count. No OEIS entry
has been edited or submitted by this work.

\subsection{Individual definitions and complete recurrence signatures}
Every list below is $c_1,\ldots,c_d$ in order, including zeros, in
\eqref{round6:eq:r6arr-recurrence}. The displayed component dimensions include
all nonzero weighted components. Initial values begin at $n=1$.
\Needspace{7\baselineskip}
\subsubsection{A183586}
\label{round6:r6entry:A183586}
Arrays have $n+1$ rows and width $w=2$ on the alphabet $0,\ldots,3$. Predicate C requires $S=\{3\}$ distinct values per block. The labeled count is multiplied by $\lambda=1/24$.
The recurrence holds for every $n\ge 3$, with order $d=2$ and component dimensions $(2)$ (total $D=2$). Source: \url{https://oeis.org/A183586}.
The first three values are $(6,56,520)$.
\[(c_1,\ldots,c_d)=(8,12).\]
The source rational generating function is also proved: $F=P/Q$, with $Q=1-\sum_{i=1}^{d}c_i x^i$ and
\[P(x)=6x+8x^{2}.\]
\Needspace{7\baselineskip}
\subsubsection{A183587}
\label{round6:r6entry:A183587}
Arrays have $n+1$ rows and width $w=3$ on the alphabet $0,\ldots,3$. Predicate C requires $S=\{3\}$ distinct values per block. The labeled count is multiplied by $\lambda=1/24$.
The recurrence holds for every $n\ge 5$, with order $d=4$ and component dimensions $(5)$ (total $D=5$). Source: \url{https://oeis.org/A183587}.
The first three values are $(56,1243,27380)$.
\[(c_1,\ldots,c_d)=(24,-30,-296,192).\]
The source rational generating function is also proved: $F=P/Q$, with $Q=1-\sum_{i=1}^{d}c_i x^i$ and
\[P(x)=56x-101x^{2}-772x^{3}+512x^{4}.\]
\Needspace{7\baselineskip}
\subsubsection{A183589}
\label{round6:r6entry:A183589}
Arrays have $n+1$ rows and width $w=5$ on the alphabet $0,\ldots,3$. Predicate C requires $S=\{3\}$ distinct values per block. The labeled count is multiplied by $\lambda=1/24$.
The recurrence holds for every $n\ge 32$, with order $d=31$ and component dimensions $(51)$ (total $D=51$). Source: \url{https://oeis.org/A183589}.
The first three values are $(4832,603766,74491726)$.
\begin{Verbatim}[fontsize=\scriptsize,breaklines=true]
179, -6470, -111366, 8421556, -27797881, -3641734984, 33282463706, 695892591667,
-8771266871058, -59332612133336, 1069043025480468, 1401139251128208,
-67977787226742392, 103924226454772488, 2296282350133287328, -7835601773053827936,
-39951799309240069504, 217732921395304638464, 281480768985183547392,
-3088710469400605992960, 1060118307218454011904, 23235106821777545527296,
-29813502870904911822848, -86205910856152217812992, 176340455396913407066112,
120352469518364633464832, -413035836961378449489920, 23412756898836805320704,
336580017414471294124032, -87482121395567755001856, -52610526709749173452800
\end{Verbatim}
\Needspace{7\baselineskip}
\subsubsection{A183596}
\label{round6:r6entry:A183596}
Arrays have $n+1$ rows and width $w=3$ on the alphabet $0,\ldots,2$. Predicate C requires $S=\{3\}$ distinct values per block. The labeled count is multiplied by $\lambda=1/6$.
The recurrence holds for every $n\ge 5$, with order $d=4$ and component dimensions $(5)$ (total $D=5$). Source: \url{https://oeis.org/A183596}.
The first three values are $(27,204,1479)$.
\[(c_1,\ldots,c_d)=(8,-3,-16,2).\]
The source rational generating function is also proved: $F=P/Q$, with $Q=1-\sum_{i=1}^{d}c_i x^i$ and
\[P(x)=27x-12x^{2}-72x^{3}+9x^{4}.\]
\Needspace{7\baselineskip}
\subsubsection{A210176}
\label{round6:r6entry:A210176}
Arrays have $n+1$ rows and width $w=2$ on the alphabet $0,\ldots,3$. Predicate C requires $S=\{3,4\}$ distinct values per block. The array uses all $4$ colors and is renamed in row-major first-occurrence order.
The recurrence holds for every $n\ge 5$, with order $d=4$ and component dimensions $(1,2,2,2)$ (total $D=7$). Source: \url{https://oeis.org/A210176}.
The first three values are $(1,51,744)$.
\[(c_1,\ldots,c_d)=(14,-26,-68,-24).\]
The source rational generating function is also proved: $F=P/Q$, with $Q=1-\sum_{i=1}^{d}c_i x^i$ and
\[P(x)=x+37x^{2}+56x^{3}+20x^{4}.\]
\Needspace{7\baselineskip}
\subsubsection{A210177}
\label{round6:r6entry:A210177}
Arrays have $n+1$ rows and width $w=3$ on the alphabet $0,\ldots,3$. Predicate C requires $S=\{3,4\}$ distinct values per block. The array uses all $4$ colors and is renamed in row-major first-occurrence order.
The recurrence holds for every $n\ge 9$, with order $d=8$ and component dimensions $(1,4,5,5)$ (total $D=15$). Source: \url{https://oeis.org/A210177}.
The first three values are $(51,2316,79017)$.
\begin{Verbatim}[fontsize=\scriptsize,breaklines=true]
39, -211, -523, 2954, -1030, -4272, 2096, -192
\end{Verbatim}
\Needspace{7\baselineskip}
\subsubsection{A210178}
\label{round6:r6entry:A210178}
Arrays have $n+1$ rows and width $w=4$ on the alphabet $0,\ldots,3$. Predicate C requires $S=\{3,4\}$ distinct values per block. The array uses all $4$ colors and is renamed in row-major first-occurrence order.
The recurrence holds for every $n\ge 17$, with order $d=16$ and component dimensions $(1,8,14,15)$ (total $D=38$). Source: \url{https://oeis.org/A210178}.
The first three values are $(744,79017,7363642)$.
\begin{Verbatim}[fontsize=\scriptsize,breaklines=true]
93, 199, -28711, 35176, 2453690, -5602318, -51561340, 96505796, 378763976, -653899024,
-988302720, 1830649984, 381532160, -1548998656, 621346816, -44040192
\end{Verbatim}
\Needspace{7\baselineskip}
\subsubsection{A210406}
\label{round6:r6entry:A210406}
Arrays have $n+1$ rows and width $w=2$ on the alphabet $0,\ldots,3$. Predicate C requires $S=\{3,4\}$ distinct values per block. The array uses at most $4$ colors and is renamed in row-major first-occurrence order.
The recurrence holds for every $n\ge 3$, with order $d=2$ and component dimensions $(1,2,2)$ (total $D=5$). Source: \url{https://oeis.org/A210406}.
The first three values are $(7,78,864)$.
\[(c_1,\ldots,c_d)=(10,12).\]
The source rational generating function is also proved: $F=P/Q$, with $Q=1-\sum_{i=1}^{d}c_i x^i$ and
\[P(x)=7x+8x^{2}.\]
\Needspace{7\baselineskip}
\subsubsection{A210407}
\label{round6:r6entry:A210407}
Arrays have $n+1$ rows and width $w=3$ on the alphabet $0,\ldots,3$. Predicate C requires $S=\{3,4\}$ distinct values per block. The array uses at most $4$ colors and is renamed in row-major first-occurrence order.
The recurrence holds for every $n\ge 5$, with order $d=4$ and component dimensions $(1,4,5)$ (total $D=10$). Source: \url{https://oeis.org/A210407}.
The first three values are $(78,2520,80496)$.
\[(c_1,\ldots,c_d)=(31,40,-280,96).\]
The source rational generating function is also proved: $F=P/Q$, with $Q=1-\sum_{i=1}^{d}c_i x^i$ and
\[P(x)=78x+102x^{2}-744x^{3}+256x^{4}.\]
\Needspace{7\baselineskip}
\subsubsection{A210408}
\label{round6:r6entry:A210408}
Arrays have $n+1$ rows and width $w=4$ on the alphabet $0,\ldots,3$. Predicate C requires $S=\{3,4\}$ distinct values per block. The array uses at most $4$ colors and is renamed in row-major first-occurrence order.
The recurrence holds for every $n\ge 10$, with order $d=9$ and component dimensions $(1,8,15)$ (total $D=24$). Source: \url{https://oeis.org/A210408}.
The first three values are $(864,80496,7380816)$.
\begin{Verbatim}[fontsize=\scriptsize,breaklines=true]
81, 1151, -13041, -117154, 478740, 445664, -2873152, 2422784, -196608
\end{Verbatim}
\Needspace{7\baselineskip}
\subsubsection{A210409}
\label{round6:r6entry:A210409}
Arrays have $n+1$ rows and width $w=5$ on the alphabet $0,\ldots,3$. Predicate C requires $S=\{3,4\}$ distinct values per block. The array uses at most $4$ colors and is renamed in row-major first-occurrence order.
The recurrence holds for every $n\ge 32$, with order $d=31$ and component dimensions $(1,16,51)$ (total $D=68$). Source: \url{https://oeis.org/A210409}.
The first three values are $(9576,2574592,677993552)$.
\begin{Verbatim}[fontsize=\scriptsize,breaklines=true]
276, -844, -647137, 5856940, 402216721, -5862695617, -54643074114, 1382343039112,
-3102790608576, -89911455972644, 660007977690996, 263588914646272, -18055429128783448,
51374341139205968, 113931146918069408, -752286799961769696, 423125539061919552,
3650072001026978048, -6040761986311450112, -5385080244386082816, 17461611564608397312,
-2457130524259205120, -16604319853477298176, 6207057768744222720, 6168586199070932992,
-2183260302370406400, -667625816443584512, 191550881093648384, 17950332629811200,
-3348386568732672, 49529562857472
\end{Verbatim}
\Needspace{7\baselineskip}
\subsubsection{A233921}
\label{round6:r6entry:A233921}
Arrays have $n+1$ rows and width $w=2$ on the alphabet $0,\ldots,4$. Predicate U has target $s=8$, with every horizontal and vertical neighboring pair unequal.
The recurrence holds for every $n\ge 7$, with order $d=6$ and component dimensions $(9)$ (total $D=9$). Source: \url{https://oeis.org/A233921}.
The first three values are $(98,518,2792)$.
\[(c_1,\ldots,c_d)=(6,5,-46,5,60,20).\]
The source rational generating function is also proved: $F=P/Q$, with $Q=1-\sum_{i=1}^{d}c_i x^i$ and
\[P(x)=98x-70x^{2}-806x^{3}+396x^{4}+1288x^{5}+400x^{6}.\]
\Needspace{7\baselineskip}
\subsubsection{A233982}
\label{round6:r6entry:A233982}
Arrays have $n+1$ rows and width $w=2$ on the alphabet $0,\ldots,4$. Predicate V has target $s=10$, with every horizontal and vertical neighboring pair unequal.
The recurrence holds for every $n\ge 4$, with order $d=3$ and component dimensions $(9)$ (total $D=9$). Source: \url{https://oeis.org/A233982}.
The first three values are $(56,236,976)$.
\[(c_1,\ldots,c_d)=(2,8,4).\]
The source rational generating function is also proved: $F=P/Q$, with $Q=1-\sum_{i=1}^{d}c_i x^i$ and
\[P(x)=56x+124x^{2}+56x^{3}.\]
\Needspace{7\baselineskip}
\subsubsection{A233983}
\label{round6:r6entry:A233983}
Arrays have $n+1$ rows and width $w=3$ on the alphabet $0,\ldots,4$. Predicate V has target $s=10$, with every horizontal and vertical neighboring pair unequal.
The recurrence holds for every $n\ge 14$, with order $d=13$ and component dimensions $(39)$ (total $D=39$). Source: \url{https://oeis.org/A233983}.
The first three values are $(236,1540,9632)$.
\begin{Verbatim}[fontsize=\scriptsize,breaklines=true]
3, 53, -116, -892, 1620, 5804, -9766, -13812, 23148, 8224, -16848, 2560, 256
\end{Verbatim}
\Needspace{7\baselineskip}
\subsubsection{A233984}
\label{round6:r6entry:A233984}
Arrays have $n+1$ rows and width $w=4$ on the alphabet $0,\ldots,4$. Predicate V has target $s=10$, with every horizontal and vertical neighboring pair unequal.
The recurrence holds for every $n\ge 26$, with order $d=25$ and component dimensions $(169)$ (total $D=169$). Source: \url{https://oeis.org/A233984}.
The first three values are $(976,9632,92320)$.
\begin{Verbatim}[fontsize=\scriptsize,breaklines=true]
0, 152, 150, -8071, -8102, 220582, 154134, -3549162, -957558, 35319910, -6392950,
-218035822, 131758738, 804014815, -803217908, -1602211348, 2304441408, 1290558096,
-3076728768, 201005888, 1628959232, -560044032, -245846016, 128188416, -7864320
\end{Verbatim}
\Needspace{7\baselineskip}
\subsubsection{A234162}
\label{round6:r6entry:A234162}
Arrays have $n+1$ rows and width $w=2$ on the alphabet $0,\ldots,4$. Predicate D has stress $s=2$.
The recurrence holds for every $n\ge 5$, with order $d=4$ and component dimensions $(9)$ (total $D=9$). Source: \url{https://oeis.org/A234162}.
The first three values are $(136,766,4296)$.
\[(c_1,\ldots,c_d)=(4,41,-132,-264).\]
The source rational generating function is also proved: $F=P/Q$, with $Q=1-\sum_{i=1}^{d}c_i x^i$ and
\[P(x)=136x+222x^{2}-4344x^{3}-6480x^{4}.\]
\Needspace{7\baselineskip}
\subsubsection{A234163}
\label{round6:r6entry:A234163}
Arrays have $n+1$ rows and width $w=3$ on the alphabet $0,\ldots,4$. Predicate D has stress $s=2$.
The recurrence holds for every $n\ge 18$, with order $d=17$ and component dimensions $(39)$ (total $D=39$). Source: \url{https://oeis.org/A234163}.
The first three values are $(766,4970,32070)$.
\begin{Verbatim}[fontsize=\scriptsize,breaklines=true]
17, 16, -1651, 5750, 45373, -252572, -278797, 2737923, 518738, -11300724, 169584,
18786600, -2246544, -13289184, 2960064, 3359232, -1119744
\end{Verbatim}
\Needspace{7\baselineskip}
\subsubsection{A234220}
\label{round6:r6entry:A234220}
Arrays have $n+1$ rows and width $w=2$ on the alphabet $0,\ldots,4$. Predicate D has stress $s=3$.
The recurrence holds for every $n\ge 4$, with order $d=3$ and component dimensions $(9)$ (total $D=9$). Source: \url{https://oeis.org/A234220}.
The first three values are $(104,436,1824)$.
\[(c_1,\ldots,c_d)=(4,20,-80).\]
The source rational generating function is also proved: $F=P/Q$, with $Q=1-\sum_{i=1}^{d}c_i x^i$ and
\[P(x)=104x+20x^{2}-2000x^{3}.\]
\Needspace{7\baselineskip}
\subsubsection{A234221}
\label{round6:r6entry:A234221}
Arrays have $n+1$ rows and width $w=3$ on the alphabet $0,\ldots,4$. Predicate D has stress $s=3$.
The recurrence holds for every $n\ge 7$, with order $d=6$ and component dimensions $(39)$ (total $D=39$). Source: \url{https://oeis.org/A234221}.
The first three values are $(436,1560,5612)$.
\[(c_1,\ldots,c_d)=(7,20,-224,144,1680,-2880).\]
The source rational generating function is also proved: $F=P/Q$, with $Q=1-\sum_{i=1}^{d}c_i x^i$ and
\[P(x)=436x-1492x^{2}-14028x^{3}+47904x^{4}+106320x^{5}-360000x^{6}.\]
\Needspace{7\baselineskip}
\subsubsection{A234222}
\label{round6:r6entry:A234222}
Arrays have $n+1$ rows and width $w=4$ on the alphabet $0,\ldots,4$. Predicate D has stress $s=3$.
The recurrence holds for every $n\ge 12$, with order $d=11$ and component dimensions $(169)$ (total $D=169$). Source: \url{https://oeis.org/A234222}.
The first three values are $(1824,5612,17448)$.
\begin{Verbatim}[fontsize=\scriptsize,breaklines=true]
9, 20, -390, 460, 5520, -14240, -26400, 115776, -13824, -299520, 276480
\end{Verbatim}
\Needspace{7\baselineskip}
\subsubsection{A234557}
\label{round6:r6entry:A234557}
Arrays have $n+1$ rows and width $w=2$ on the alphabet $0,\ldots,4$. Predicate D has stress $s=4$.
The recurrence holds for every $n\ge 6$, with order $d=5$ and component dimensions $(9)$ (total $D=9$). Source: \url{https://oeis.org/A234557}.
The first three values are $(70,220,618)$.
\[(c_1,\ldots,c_d)=(3,18,-54,-80,240).\]
The source rational generating function is also proved: $F=P/Q$, with $Q=1-\sum_{i=1}^{d}c_i x^i$ and
\[P(x)=70x+10x^{2}-1302x^{3}-80x^{4}+6000x^{5}.\]
\Needspace{7\baselineskip}
\subsubsection{A234558}
\label{round6:r6entry:A234558}
Arrays have $n+1$ rows and width $w=3$ on the alphabet $0,\ldots,4$. Predicate D has stress $s=4$.
The recurrence holds for every $n\ge 15$, with order $d=14$ and component dimensions $(39)$ (total $D=39$). Source: \url{https://oeis.org/A234558}.
The first three values are $(220,546,1240)$.
\begin{Verbatim}[fontsize=\scriptsize,breaklines=true]
4, 30, -132, -331, 1720, 1530, -11280, -1324, 39136, -12360, -67968, 39456, 46080,
-34560
\end{Verbatim}
\Needspace{7\baselineskip}
\subsubsection{A234559}
\label{round6:r6entry:A234559}
Arrays have $n+1$ rows and width $w=4$ on the alphabet $0,\ldots,4$. Predicate D has stress $s=4$.
The recurrence holds for every $n\ge 15$, with order $d=14$ and component dimensions $(169)$ (total $D=169$). Source: \url{https://oeis.org/A234559}.
The first three values are $(618,1240,2390)$.
\begin{Verbatim}[fontsize=\scriptsize,breaklines=true]
4, 30, -132, -331, 1720, 1530, -11280, -1324, 39136, -12360, -67968, 39456, 46080,
-34560
\end{Verbatim}
\Needspace{7\baselineskip}
\subsubsection{A234762}
\label{round6:r6entry:A234762}
Arrays have $n+1$ rows and width $w=2$ on the alphabet $0,\ldots,3$. Predicate J permits exactly one orientation whose four edge jumps are all at most $+1$.
The recurrence holds for every $n\ge 6$, with order $d=5$ and component dimensions $(6)$ (total $D=6$). Source: \url{https://oeis.org/A234762}.
The first three values are $(56,240,916)$.
\[(c_1,\ldots,c_d)=(1,14,-1,-25,-10).\]
The source rational generating function is also proved: $F=P/Q$, with $Q=1-\sum_{i=1}^{d}c_i x^i$ and
\[P(x)=56x+184x^{2}-108x^{3}-420x^{4}-156x^{5}.\]
\Needspace{7\baselineskip}
\subsubsection{A234763}
\label{round6:r6entry:A234763}
Arrays have $n+1$ rows and width $w=3$ on the alphabet $0,\ldots,3$. Predicate J permits exactly one orientation whose four edge jumps are all at most $+1$.
The recurrence holds for every $n\ge 14$, with order $d=13$ and component dimensions $(20)$ (total $D=20$). Source: \url{https://oeis.org/A234763}.
The first three values are $(240,1624,9404)$.
\begin{Verbatim}[fontsize=\scriptsize,breaklines=true]
8, 21, -252, 125, 1818, -2050, -4364, 6237, 2860, -5787, 336, 1311, -262
\end{Verbatim}
\Needspace{7\baselineskip}
\subsubsection{A234764}
\label{round6:r6entry:A234764}
Arrays have $n+1$ rows and width $w=4$ on the alphabet $0,\ldots,3$. Predicate J permits exactly one orientation whose four edge jumps are all at most $+1$.
The recurrence holds for every $n\ge 28$, with order $d=27$ and component dimensions $(72)$ (total $D=72$). Source: \url{https://oeis.org/A234764}.
The first three values are $(916,9404,71280)$.
\begin{Verbatim}[fontsize=\scriptsize,breaklines=true]
3, 137, -284, -6386, 9604, 145993, -160399, -1876827, 1488326, 14432970, -8040903,
-68366990, 25519880, 201876831, -47552365, -371183764, 52834992, 418743519, -37830792,
-280238238, 20541150, 104060106, -8730936, -18816300, 2167560, 1263600, -194400
\end{Verbatim}
\Needspace{7\baselineskip}
\subsubsection{A234875}
\label{round6:r6entry:A234875}
Arrays have $n+1$ rows and width $w=2$ on the alphabet $0,\ldots,4$. Predicate D has stress $s=1$, with every horizontal and vertical neighboring pair unequal.
The recurrence holds for every $n\ge 5$, with order $d=4$ and component dimensions $(9)$ (total $D=9$). Source: \url{https://oeis.org/A234875}.
The first three values are $(32,68,128)$.
\[(c_1,\ldots,c_d)=(0,5,0,-4).\]
The source rational generating function is also proved: $F=P/Q$, with $Q=1-\sum_{i=1}^{d}c_i x^i$ and
\[P(x)=32x+68x^{2}-32x^{3}-72x^{4}.\]
\Needspace{7\baselineskip}
\subsubsection{A234876}
\label{round6:r6entry:A234876}
Arrays have $n+1$ rows and width $w=3$ on the alphabet $0,\ldots,4$. Predicate D has stress $s=1$, with every horizontal and vertical neighboring pair unequal.
The recurrence holds for every $n\ge 10$, with order $d=9$ and component dimensions $(39)$ (total $D=39$). Source: \url{https://oeis.org/A234876}.
The first three values are $(68,148,268)$.
\begin{Verbatim}[fontsize=\scriptsize,breaklines=true]
2, 10, -21, -29, 68, 22, -73, 0, 22
\end{Verbatim}
The source rational generating function is also proved: $F=P/Q$, with $Q=1-\sum_{i=1}^{d}c_i x^i$ and
\[P(x)=68x+12x^{2}-708x^{3}+40x^{4}+2332x^{5}-416x^{6}-2684x^{7}+440x^{8}+924x^{9}.\]
\Needspace{7\baselineskip}
\subsubsection{A234877}
\label{round6:r6entry:A234877}
Arrays have $n+1$ rows and width $w=4$ on the alphabet $0,\ldots,4$. Predicate D has stress $s=1$, with every horizontal and vertical neighboring pair unequal.
The recurrence holds for every $n\ge 12$, with order $d=11$ and component dimensions $(169)$ (total $D=169$). Source: \url{https://oeis.org/A234877}.
The first three values are $(128,268,472)$.
\begin{Verbatim}[fontsize=\scriptsize,breaklines=true]
2, 11, -23, -38, 87, 46, -130, -16, 78, 0, -16
\end{Verbatim}
The source rational generating function is also proved: $F=P/Q$, with $Q=1-\sum_{i=1}^{d}c_i x^i$ and
\[\begin{aligned}P(x)&=128x+12x^{2}-1472x^{3}+168x^{4}+5700x^{5}-1248x^{6}-9048x^{7}+1968x^{8}\\ &\quad +5984x^{9}-768x^{10}-1376x^{11}.\end{aligned}\]
\Needspace{7\baselineskip}
\subsubsection{A235011}
\label{round6:r6entry:A235011}
Arrays have $n+1$ rows and width $w=3$ on the alphabet $0,\ldots,4$. Predicate D has stress $s=3$, with every horizontal and vertical neighboring pair unequal.
The recurrence holds for every $n\ge 6$, with order $d=5$ and component dimensions $(39)$ (total $D=39$). Source: \url{https://oeis.org/A235011}.
The first three values are $(104,224,488)$.
\[(c_1,\ldots,c_d)=(3,4,-14,-4,16).\]
The source rational generating function is also proved: $F=P/Q$, with $Q=1-\sum_{i=1}^{d}c_i x^i$ and
\[P(x)=104x-88x^{2}-600x^{3}+208x^{4}+832x^{5}.\]
\Needspace{7\baselineskip}
\subsubsection{A235012}
\label{round6:r6entry:A235012}
Arrays have $n+1$ rows and width $w=4$ on the alphabet $0,\ldots,4$. Predicate D has stress $s=3$, with every horizontal and vertical neighboring pair unequal.
The recurrence holds for every $n\ge 9$, with order $d=8$ and component dimensions $(169)$ (total $D=169$). Source: \url{https://oeis.org/A235012}.
The first three values are $(264,488,936)$.
\begin{Verbatim}[fontsize=\scriptsize,breaklines=true]
4, 4, -30, 7, 74, -46, -60, 48
\end{Verbatim}
The source rational generating function is also proved: $F=P/Q$, with $Q=1-\sum_{i=1}^{d}c_i x^i$ and
\[P(x)=264x-568x^{2}-2072x^{3}+4136x^{4}+5408x^{5}-9792x^{6}-4672x^{7}+7552x^{8}.\]
\Needspace{7\baselineskip}
\subsubsection{A235080}
\label{round6:r6entry:A235080}
Arrays have $n+1$ rows and width $w=2$ on the alphabet $0,\ldots,4$. Predicate D has stress $s=4$, with every horizontal and vertical neighboring pair unequal.
The recurrence holds for every $n\ge 5$, with order $d=4$ and component dimensions $(9)$ (total $D=9$). Source: \url{https://oeis.org/A235080}.
The first three values are $(38,90,202)$.
\[(c_1,\ldots,c_d)=(2,6,-10,-5).\]
The source rational generating function is also proved: $F=P/Q$, with $Q=1-\sum_{i=1}^{d}c_i x^i$ and
\[P(x)=38x+14x^{2}-206x^{3}-86x^{4}.\]
\Needspace{7\baselineskip}
\subsubsection{A235081}
\label{round6:r6entry:A235081}
Arrays have $n+1$ rows and width $w=3$ on the alphabet $0,\ldots,4$. Predicate D has stress $s=4$, with every horizontal and vertical neighboring pair unequal.
The recurrence holds for every $n\ge 11$, with order $d=10$ and component dimensions $(39)$ (total $D=39$). Source: \url{https://oeis.org/A235081}.
The first three values are $(90,198,418)$.
\begin{Verbatim}[fontsize=\scriptsize,breaklines=true]
3, 10, -33, -28, 111, 20, -111, -27, 30, 10
\end{Verbatim}
The source rational generating function is also proved: $F=P/Q$, with $Q=1-\sum_{i=1}^{d}c_i x^i$ and
\[\begin{aligned}P(x)&=90x-72x^{2}-1076x^{3}+666x^{4}+4086x^{5}-1444x^{6}-5108x^{7}-514x^{8}\\ &\quad +1544x^{9}+452x^{10}.\end{aligned}\]
\Needspace{7\baselineskip}
\subsubsection{A235271}
\label{round6:r6entry:A235271}
Arrays have $n+1$ rows and width $w=2$ on the alphabet $0,\ldots,4$. Predicate D has stress $s=5$.
The recurrence holds for every $n\ge 5$, with order $d=4$ and component dimensions $(9)$ (total $D=9$). Source: \url{https://oeis.org/A235271}.
The first three values are $(40,100,208)$.
\[(c_1,\ldots,c_d)=(0,10,0,-24).\]
The source rational generating function is also proved: $F=P/Q$, with $Q=1-\sum_{i=1}^{d}c_i x^i$ and
\[P(x)=40x+100x^{2}-192x^{3}-480x^{4}.\]
\Needspace{7\baselineskip}
\subsubsection{A235272}
\label{round6:r6entry:A235272}
Arrays have $n+1$ rows and width $w=3$ on the alphabet $0,\ldots,4$. Predicate D has stress $s=5$.
The recurrence holds for every $n\ge 10$, with order $d=9$ and component dimensions $(39)$ (total $D=39$). Source: \url{https://oeis.org/A235272}.
The first three values are $(100,208,380)$.
\begin{Verbatim}[fontsize=\scriptsize,breaklines=true]
1, 15, -15, -80, 80, 180, -180, -144, 144
\end{Verbatim}
The source rational generating function is also proved: $F=P/Q$, with $Q=1-\sum_{i=1}^{d}c_i x^i$ and
\[\begin{aligned}P(x)&=100x+108x^{2}-1328x^{3}-1156x^{4}+6236x^{5}+3864x^{6}-12312x^{7}-4032x^{8}\\ &\quad +8640x^{9}.\end{aligned}\]
\Needspace{7\baselineskip}
\subsubsection{A235273}
\label{round6:r6entry:A235273}
Arrays have $n+1$ rows and width $w=4$ on the alphabet $0,\ldots,4$. Predicate D has stress $s=5$.
The recurrence holds for every $n\ge 10$, with order $d=9$ and component dimensions $(169)$ (total $D=169$). Source: \url{https://oeis.org/A235273}.
The first three values are $(208,380,632)$.
\begin{Verbatim}[fontsize=\scriptsize,breaklines=true]
1, 15, -15, -80, 80, 180, -180, -144, 144
\end{Verbatim}
The source rational generating function is also proved: $F=P/Q$, with $Q=1-\sum_{i=1}^{d}c_i x^i$ and
\[\begin{aligned}P(x)&=208x+172x^{2}-2868x^{3}-1932x^{4}+13924x^{5}+6792x^{6}-28248x^{7}-7488x^{8}\\ &\quad +20160x^{9}.\end{aligned}\]
\Needspace{7\baselineskip}
\subsubsection{A235303}
\label{round6:r6entry:A235303}
Arrays have $n+1$ rows and width $w=2$ on the alphabet $0,\ldots,5$. Predicate D has stress $s=6$.
The recurrence holds for every $n\ge 6$, with order $d=5$ and component dimensions $(12)$ (total $D=12$). Source: \url{https://oeis.org/A235303}.
The first three values are $(70,210,518)$.
\[(c_1,\ldots,c_d)=(3,13,-39,-40,120).\]
The source rational generating function is also proved: $F=P/Q$, with $Q=1-\sum_{i=1}^{d}c_i x^i$ and
\[P(x)=70x-1022x^{3}+3600x^{5}.\]
\Needspace{7\baselineskip}
\subsubsection{A235304}
\label{round6:r6entry:A235304}
Arrays have $n+1$ rows and width $w=3$ on the alphabet $0,\ldots,5$. Predicate D has stress $s=6$.
The recurrence holds for every $n\ge 15$, with order $d=14$ and component dimensions $(63)$ (total $D=63$). Source: \url{https://oeis.org/A235304}.
The first three values are $(210,518,1106)$.
\begin{Verbatim}[fontsize=\scriptsize,breaklines=true]
4, 25, -112, -231, 1260, 875, -7280, -224, 22736, -7980, -36288, 21456, 23040, -17280
\end{Verbatim}
\Needspace{7\baselineskip}
\subsubsection{A235305}
\label{round6:r6entry:A235305}
Arrays have $n+1$ rows and width $w=4$ on the alphabet $0,\ldots,5$. Predicate D has stress $s=6$.
The recurrence holds for every $n\ge 15$, with order $d=14$ and component dimensions $(342)$ (total $D=342$). Source: \url{https://oeis.org/A235305}.
The first three values are $(518,1106,2118)$.
\begin{Verbatim}[fontsize=\scriptsize,breaklines=true]
4, 25, -112, -231, 1260, 875, -7280, -224, 22736, -7980, -36288, 21456, 23040, -17280
\end{Verbatim}
\Needspace{7\baselineskip}
\subsubsection{A235312}
\label{round6:r6entry:A235312}
Arrays have $n+1$ rows and width $w=2$ on the alphabet $0,\ldots,6$. Predicate D has stress $s=7$.
The recurrence holds for every $n\ge 7$, with order $d=6$ and component dimensions $(16)$ (total $D=16$). Source: \url{https://oeis.org/A235312}.
The first three values are $(112,392,1120)$.
\[(c_1,\ldots,c_d)=(0,28,0,-252,0,720).\]
The source rational generating function is also proved: $F=P/Q$, with $Q=1-\sum_{i=1}^{d}c_i x^i$ and
\[P(x)=112x+392x^{2}-2016x^{3}-7056x^{4}+8640x^{5}+30240x^{6}.\]
\Needspace{7\baselineskip}
\subsubsection{A235313}
\label{round6:r6entry:A235313}
Arrays have $n+1$ rows and width $w=3$ on the alphabet $0,\ldots,6$. Predicate D has stress $s=7$.
The recurrence holds for every $n\ge 19$, with order $d=18$ and component dimensions $(100)$ (total $D=100$). Source: \url{https://oeis.org/A235313}.
The first three values are $(392,1120,2744)$.
\begin{Verbatim}[fontsize=\scriptsize,breaklines=true]
4, 47, -200, -901, 4204, 8957, -48440, -47194, 334096, 101948, -1410080, 170136,
3549696, -1446912, -4861440, 2954880, 2764800, -2073600
\end{Verbatim}


\clearpage
\section{Reproduction and verification}
\label{round6:sec:r6reproduce}
\begingroup\small
The archive includes this standalone manuscript, the numbered ledger,
four proof fragments, four main verifiers, an independent array checker,
complete transfer matrices, exact algebraic certificates, source-scope
audits and four cross-review reports. Full scraped webpages are omitted.

\subsection{Running the offline checks}
Use Python 3.10 or later with SymPy and mpmath. From the extracted
directory run:
\begin{Verbatim}[fontsize=\small]
python3 r6_run_all.py
\end{Verbatim}
The runner checks the 100 distinct IDs, the 508-ID exclusion ledger,
the proof labels and the reviewed proof/code hashes. It executes the
following independent components without network access.
\begin{center}
\begin{tabularx}{\linewidth}{@{}lX@{}}
\toprule
Script&Purpose\\
\midrule
\texttt{r6\_root\_verify.py}&Exact algebraic and modular certificates,
with bounded series and numerical checks supplementing 18 written proofs.\\
\texttt{r6\_arith\_verify.py}&Exact arithmetic checks for 11 packages,
including the complete modular certificate for \oeis{A116184}.\\
\texttt{r6\_arith\_independent\_check.py}&Independent common-denominator
enumeration of its complete 49,284-residue period.\\
\texttt{r6\_gf\_verify.py}&30 exact all-index differential certificates
and 600 source-prefix checks, retaining all boundary terms.\\
\texttt{r6\_arrays\_verify.py}&41 recurrences from 49 complete rebuilt
integer matrices and their vector identities; 25 source GFs.\\
\texttt{r6\_arrays\_root\_verify.py}&82 independent initial counts
using columns and global used-color sets, without importing the row
engine or using its inclusion--exclusion normalization.\\
\bottomrule
\end{tabularx}
\end{center}

\subsection{Why the finite certificates prove every index}
For an array recurrence the defining count is a linear combination of
$m^{\mathsf T}T^n\mathbf1$. The complete identity
$TR(T)\mathbf1=0$ remains true after multiplication by every
nonnegative power of $T$, proving the recurrence on the stated infinite
range. All 6,105 nonzero weighted edges are rebuilt from the local
definitions; the largest matrix has 342 states and the largest
recurrence has order 31. Every coordinate is checked with integers.
Agreement of 832 source-prefix values is an interpretation check,
not the proof of the universal recurrence.

For a generating-function recurrence, the exact action of its
differential operator is evaluated in a rational algebraic basis.
The residual is the displayed boundary polynomial; coefficients
beyond its degree vanish identically. This proves the stated
infinite tail, including each qualified starting range.

The other universal conclusions follow from the written algebraic,
valuation, modular, or analytic arguments. Their bounded calculations
check transcription and indexing. In particular, the theta evaluations
follow from exact classical identities, with analytic convergence
justified; their high-precision numerical checks are supplementary.

\subsection{Rebuilding and integrity}
The delivered \texttt{oeis\_hunt\_round6.tex} contains all proof text
directly and needs no external \verb|\input| files. Compile it twice
with pdfLaTeX and the standard packages in its preamble. To rebuild
the manuscript, catalogue and combined audit index, run:
\begin{Verbatim}[fontsize=\small]
python3 r6_build_dossier.py
\end{Verbatim}
The README records optional certificate reconstruction commands.
\texttt{manifest.json} gives SHA-256 digests of the delivered files.
Running the verifiers regenerates reports and timing fields, so those
report digests can change after a successful rerun. Review bindings
for immutable proof and code files are checked separately.

Source records retain authorship, the selected assertion, indexing,
excluded candidates and search limits. The 100-entry selection count
does not certify historical first discovery or official OEIS closure.
\endgroup
% ===== END ORIGINAL ROUND 6 =====

% ===== BEGIN ORIGINAL ROUND 7 =====
\clearpage
\setcounter{section}{0}\setcounter{equation}{0}\setcounter{theorem}{0}
\setcounter{table}{0}\setcounter{figure}{0}
\normalsize\renewcommand{\arraystretch}{1}\renewcommand{\baselinestretch}{1}\selectfont
\part[Round 7: Divisibility, arithmetic asymptotics, and exact array recurrences]{Divisibility, arithmetic asymptotics, and exact array recurrences}
\label{combined:round7}
\rhead{\small Round 7 / KEIO}
\begingroup\small
\section*{Round overview}
\addcontentsline{toc}{section}{Round 7 overview and scope}
This seventh-round dossier resolves 100 selected assertion packages on
100 distinct OEIS entries, disjoint from the 608 A-numbers selected in
Rounds I--VI. It proves 91 selected claims, specifies and proves the
valid tail of four conjectural recurrences, corrects and proves four
formulas or indexing/support statements, and refutes one asymptotic
conjecture. Multiple assertions on one entry count once.

The four parts contain 14 packages on divisibility and formal congruences
(including one theta-index correction), 10 on arithmetic and asymptotics,
31 on generating functions, and 45 on constrained arrays. The written
arguments use integer coefficient rings, exact divisor and lattice
formulas, differential identities with boundary polynomials, and finite
transfer operators with integer vector certificates.

\subsection*{What the count means}
Each catalogue row is one precisely scoped assertion package attached to
one distinct A-number. Several rows can share a general lemma. The count
does not assert 100 unrelated methods or 100 new general theorems.
The consulted source statements explicitly mark the selected claims as
conjectural or empirical. Established formulas, classical identities and
previous project lemmas are credited inputs.

Some entries contain further questions outside the selected scope.
In particular, the exponential-series packages do not claim their
separate asymptotic observations. The chosen scope and every correction
are recorded in the proofs and the machine-readable ledger. Two
additional certified array candidates were held out when the total
reached 100; they are not counted in this dossier.

The previous 608 is an exclusion ledger, not a certified discovery total.
The five Round IV priority withdrawals disclosed in Round V remain in
force; all those A-IDs remain excluded. New consequences at previously
selected A397242 and equivalent quotient/decimal entries add no counts.

\subsection*{Qualifications and corrections}
Four recurrence sources omit a starting index. The following earlier
admissible residuals are nonzero; the proved tails start one index later.
We do not attribute an unstated false starting index to their authors.
\begin{center}
\begin{tabular}{@{}lrrl@{}}
\toprule Entry&Exceptional index&Residual&Proved range\\\midrule
\oeis{A163824}&3&$-6$&$n\ge4$\\
\oeis{A165201}&3&2&$n\ge4$\\
\oeis{A190788}&7&16&$n\ge8$\\
\oeis{A228180}&3&12&$n\ge4$\\\bottomrule
\end{tabular}
\end{center}
The four corrected packages are \oeis{A030212} (the special-value sum
requires the actual offset), \oeis{A108630} (replace coefficient
arguments $n$ by $n+1$), \oeis{A168494} (one duplicated $a(n-2)$ should
be $a(n-1)$), and \oeis{A366561} (the support rule requires the correct
2-adic restriction). The proposed bound in \oeis{A133205} is refuted by
a divergent normalized summatory lower bound.

Two further input issues are disclosed within otherwise proved packages:
the defining range in the title of \oeis{A300732} contradicts its initial
value, and a displayed generating function in \oeis{A026270} is
inconsistent with its defining triangle. Each proof explicitly derives
and uses the consistent definition; these do not add extra counts.

\subsection*{Evidence and limits}
Primary entries, relevant linked material and targeted searches were
consulted on 20 October 2026; retrieval may use cached pages. The retained
OEIS Open comparison covers the 100 costliest claims in Appendix A on
97 distinct A-numbers, with no current overlap. It does not cover all
results in that project. This is not an exhaustive historical priority
audit and no worldwide-first-proof claim is made.

No OEIS edit or submission was performed. Mathematical resolution of a
selected assertion here is distinct from editorial acceptance by OEIS.
Cross-review is by separate agents within the same AI-assisted workflow,
not external peer review or proof-assistant formalization.
\endgroup

\clearpage
\section{Catalogue of the 100 selected entries}
\label{round7:sec:r7catalogue}
\textbf{P}: selected claim proved.
\textbf{B}: recurrence boundary qualified and proved.
\textbf{C}: formula or support rule corrected and proved.
\textbf{I}: indexing corrected and proved.
\textbf{R}: selected asymptotic conjecture refuted.
Page references lead to the argument or full recurrence signature.

\small
\begin{longtable}{@{}r p{21mm} p{103mm} c r@{}}
\toprule No.&Entry&Selected assertion package&Result&Page\\
\midrule\endfirsthead
\toprule No.&Entry&Selected assertion package&Result&Page\\
\midrule\endhead
\midrule\multicolumn{5}{r}{\emph{Continued on the next page}}\\\endfoot
\bottomrule\endlastfoot

\multicolumn{5}{@{}l}{\textbf{Divisibility and formal congruences}}\\[2pt]

1&\oeis{A266489}&Index divisibility and even coefficients after index 1&P&\pageref{round7:r7entry:A266489}\\

2&\oeis{A300732}&Index divisibility, with the defining range clarified&P&\pageref{round7:r7entry:A300732}\\

3&\oeis{A292394}&Square of the index divides every positive-index term&P&\pageref{round7:r7entry:A292394}\\

4&\oeis{A300734}&Square-index divisibility with exponent 2n\textasciicircum{}2&P&\pageref{round7:r7entry:A300734}\\

5&\oeis{A395833}&Divisibility by 2n-1 and by 3&P&\pageref{round7:r7entry:A395833}\\

6&\oeis{A397348}&Odd coefficients exactly at n=2\textasciicircum{}j-1&P&\pageref{round7:r7entry:A397348}\\

7&\oeis{A397345}&Odd coefficients exactly at n=2\textasciicircum{}j-1&P&\pageref{round7:r7entry:A397345}\\

8&\oeis{A397245}&Complete ternary description of coefficients modulo 3&P&\pageref{round7:r7entry:A397245}\\

9&\oeis{A397349}&Powers-of-two parity and alternating residues modulo 3&P&\pageref{round7:r7entry:A397349}\\

10&\oeis{A397346}&Period-four residue pattern modulo 8&P&\pageref{round7:r7entry:A397346}\\

11&\oeis{A397347}&Period-four logarithmic-derivative pattern modulo 4&P&\pageref{round7:r7entry:A397347}\\

12&\oeis{A396846}&Powers of three modulo 3 and period four modulo 8&P&\pageref{round7:r7entry:A396846}\\

13&\oeis{A397244}&Fibbinary characterization of odd coefficients&P&\pageref{round7:r7entry:A397244}\\

14&\oeis{A030212}&Correct the one-index shift in the empirical gamma value&I&\pageref{round7:r7entry:A030212}\\

\multicolumn{5}{@{}l}{\textbf{Arithmetic and asymptotic assertions}}\\[2pt]

15&\oeis{A092843}&The shifted-totient summatory asymptotic&P&\pageref{round7:r7entry:A092843}\\

16&\oeis{A108625}&The general alternating binomial identity&P&\pageref{round7:r7entry:A108625}\\

17&\oeis{A133205}&Refute the proposed summatory big-O bound&R&\pageref{round7:r7entry:A133205}\\

18&\oeis{A140110}&The gcd-sum characterization&P&\pageref{round7:r7entry:A140110}\\

19&\oeis{A179873}&Monotonicity of the odd antiharmonic means&P&\pageref{round7:r7entry:A179873}\\

20&\oeis{A309263}&Membership of 4 times every nonnegative power of 5&P&\pageref{round7:r7entry:A309263}\\

21&\oeis{A356459}&The exact leading constant in the asymptotic&P&\pageref{round7:r7entry:A356459}\\

22&\oeis{A366561}&Local formulas and the corrected zero-support condition&C&\pageref{round7:r7entry:A366561}\\

23&\oeis{A375838}&The conjectured triangle recurrence&P&\pageref{round7:r7entry:A375838}\\

24&\oeis{A385001}&The fixed-parameter partition summatory asymptotics&P&\pageref{round7:r7entry:A385001}\\

\multicolumn{5}{@{}l}{\textbf{Generating-function certificates}}\\[2pt]

25&\oeis{A026026}&Conjectured recurrence of order 2&P&\pageref{round7:r7entry:A026026}\\

26&\oeis{A026270}&Conjectured recurrence of order 3&P&\pageref{round7:r7entry:A026270}\\

27&\oeis{A059279}&Conjectured recurrence of order 3&P&\pageref{round7:r7entry:A059279}\\

28&\oeis{A096471}&Conjectured recurrence of order 5&P&\pageref{round7:r7entry:A096471}\\

29&\oeis{A106391}&Conjectured recurrence of order 6&P&\pageref{round7:r7entry:A106391}\\

30&\oeis{A107026}&Conjectured recurrence of order 2&P&\pageref{round7:r7entry:A107026}\\

31&\oeis{A108630}&Conjectured recurrence of order 6; coefficient argument corrected&I&\pageref{round7:r7entry:A108630}\\

32&\oeis{A110198}&Conjectured recurrence of order 5&P&\pageref{round7:r7entry:A110198}\\

33&\oeis{A119975}&Conjectured recurrence of order 3&P&\pageref{round7:r7entry:A119975}\\

34&\oeis{A134389}&Conjectured recurrence of order 3&P&\pageref{round7:r7entry:A134389}\\

35&\oeis{A143020}&Conjectured recurrence of order 2&P&\pageref{round7:r7entry:A143020}\\

36&\oeis{A155051}&Conjectured recurrence of order 4&P&\pageref{round7:r7entry:A155051}\\

37&\oeis{A157418}&Conjectured recurrence of order 3&P&\pageref{round7:r7entry:A157418}\\

38&\oeis{A163824}&Conjectured recurrence of order 3; starting range specified&B&\pageref{round7:r7entry:A163824}\\

39&\oeis{A165201}&Conjectured recurrence of order 3; starting range specified&B&\pageref{round7:r7entry:A165201}\\

40&\oeis{A168494}&Conjectured recurrence of order 4; term index corrected&C&\pageref{round7:r7entry:A168494}\\

41&\oeis{A168503}&Conjectured recurrence of order 4&P&\pageref{round7:r7entry:A168503}\\

42&\oeis{A174016}&Conjectured recurrence of order 4&P&\pageref{round7:r7entry:A174016}\\

43&\oeis{A184018}&Conjectured recurrence of order 4&P&\pageref{round7:r7entry:A184018}\\

44&\oeis{A185020}&Conjectured recurrence of order 2&P&\pageref{round7:r7entry:A185020}\\

45&\oeis{A190725}&Conjectured recurrence of order 4&P&\pageref{round7:r7entry:A190725}\\

46&\oeis{A190736}&Conjectured recurrence of order 5&P&\pageref{round7:r7entry:A190736}\\

47&\oeis{A190788}&Conjectured recurrence of order 7; starting range specified&B&\pageref{round7:r7entry:A190788}\\

48&\oeis{A191319}&Conjectured recurrence of order 4&P&\pageref{round7:r7entry:A191319}\\

49&\oeis{A191398}&Conjectured recurrence of order 6&P&\pageref{round7:r7entry:A191398}\\

50&\oeis{A220899}&Conjectured recurrence of order 5&P&\pageref{round7:r7entry:A220899}\\

51&\oeis{A228180}&Conjectured recurrence of order 3; starting range specified&B&\pageref{round7:r7entry:A228180}\\

52&\oeis{A278472}&Conjectured recurrence of order 3&P&\pageref{round7:r7entry:A278472}\\

53&\oeis{A372991}&Conjectured recurrence of order 3&P&\pageref{round7:r7entry:A372991}\\

54&\oeis{A383702}&Conjectured recurrence of order 3&P&\pageref{round7:r7entry:A383702}\\

55&\oeis{A386830}&Conjectured recurrence of order 2&P&\pageref{round7:r7entry:A386830}\\

\multicolumn{5}{@{}l}{\textbf{Finite-state array certificates}}\\[2pt]

56&\oeis{A209550}&Empirical array recurrence of order 37&P&\pageref{round7:r7entry:A209550}\\

57&\oeis{A210179}&Empirical array recurrence of order 50&P&\pageref{round7:r7entry:A210179}\\

58&\oeis{A233637}&Empirical array recurrence of order 2&P&\pageref{round7:r7entry:A233637}\\

59&\oeis{A233638}&Empirical array recurrence of order 7&P&\pageref{round7:r7entry:A233638}\\

60&\oeis{A233639}&Empirical array recurrence of order 13&P&\pageref{round7:r7entry:A233639}\\

61&\oeis{A233640}&Empirical array recurrence of order 31&P&\pageref{round7:r7entry:A233640}\\

62&\oeis{A233684}&Empirical array recurrence of order 3&P&\pageref{round7:r7entry:A233684}\\

63&\oeis{A233685}&Empirical array recurrence of order 9&P&\pageref{round7:r7entry:A233685}\\

64&\oeis{A233686}&Empirical array recurrence of order 15&P&\pageref{round7:r7entry:A233686}\\

65&\oeis{A233710}&Empirical array recurrence of order 8&P&\pageref{round7:r7entry:A233710}\\

66&\oeis{A233711}&Empirical array recurrence of order 27&P&\pageref{round7:r7entry:A233711}\\

67&\oeis{A233749}&Empirical array recurrence of order 9&P&\pageref{round7:r7entry:A233749}\\

68&\oeis{A233750}&Empirical array recurrence of order 44&P&\pageref{round7:r7entry:A233750}\\

69&\oeis{A233885}&Empirical array recurrence of order 10&P&\pageref{round7:r7entry:A233885}\\

70&\oeis{A233886}&Empirical array recurrence of order 32&P&\pageref{round7:r7entry:A233886}\\

71&\oeis{A233910}&Empirical array recurrence of order 11&P&\pageref{round7:r7entry:A233910}\\

72&\oeis{A233911}&Empirical array recurrence of order 38&P&\pageref{round7:r7entry:A233911}\\

73&\oeis{A233951}&Empirical array recurrence of order 4&P&\pageref{round7:r7entry:A233951}\\

74&\oeis{A233952}&Empirical array recurrence of order 19&P&\pageref{round7:r7entry:A233952}\\

75&\oeis{A234146}&Empirical array recurrence of order 8&P&\pageref{round7:r7entry:A234146}\\

76&\oeis{A234147}&Empirical array recurrence of order 37&P&\pageref{round7:r7entry:A234147}\\

77&\oeis{A234651}&Empirical array recurrence of order 5&P&\pageref{round7:r7entry:A234651}\\

78&\oeis{A234652}&Empirical array recurrence of order 15&P&\pageref{round7:r7entry:A234652}\\

79&\oeis{A234653}&Empirical array recurrence of order 21&P&\pageref{round7:r7entry:A234653}\\

80&\oeis{A234731}&Empirical array recurrence of order 7&P&\pageref{round7:r7entry:A234731}\\

81&\oeis{A234732}&Empirical array recurrence of order 23&P&\pageref{round7:r7entry:A234732}\\

82&\oeis{A234733}&Empirical array recurrence of order 23&P&\pageref{round7:r7entry:A234733}\\

83&\oeis{A234753}&Empirical array recurrence of order 6&P&\pageref{round7:r7entry:A234753}\\

84&\oeis{A234754}&Empirical array recurrence of order 17&P&\pageref{round7:r7entry:A234754}\\

85&\oeis{A234770}&Empirical array recurrence of order 6&P&\pageref{round7:r7entry:A234770}\\

86&\oeis{A234771}&Empirical array recurrence of order 23&P&\pageref{round7:r7entry:A234771}\\

87&\oeis{A235091}&Empirical array recurrence of order 5&P&\pageref{round7:r7entry:A235091}\\

88&\oeis{A235092}&Empirical array recurrence of order 16&P&\pageref{round7:r7entry:A235092}\\

89&\oeis{A235093}&Empirical array recurrence of order 22&P&\pageref{round7:r7entry:A235093}\\

90&\oeis{A235216}&Empirical array recurrence of order 4&P&\pageref{round7:r7entry:A235216}\\

91&\oeis{A235217}&Empirical array recurrence of order 8&P&\pageref{round7:r7entry:A235217}\\

92&\oeis{A235218}&Empirical array recurrence of order 21&P&\pageref{round7:r7entry:A235218}\\

93&\oeis{A235282}&Empirical array recurrence of order 3&P&\pageref{round7:r7entry:A235282}\\

94&\oeis{A235283}&Empirical array recurrence of order 6&P&\pageref{round7:r7entry:A235283}\\

95&\oeis{A235284}&Empirical array recurrence of order 6&P&\pageref{round7:r7entry:A235284}\\

96&\oeis{A235291}&Empirical array recurrence of order 4&P&\pageref{round7:r7entry:A235291}\\

97&\oeis{A235292}&Empirical array recurrence of order 10&P&\pageref{round7:r7entry:A235292}\\

98&\oeis{A235293}&Empirical array recurrence of order 10&P&\pageref{round7:r7entry:A235293}\\

99&\oeis{A235294}&Empirical array recurrence of order 10&P&\pageref{round7:r7entry:A235294}\\

100&\oeis{A235314}&Empirical array recurrence of order 18&P&\pageref{round7:r7entry:A235314}\\

\end{longtable}
\normalsize
\clearpage

\section{Divisibility, formal congruences, and an indexing correction}
\label{round7:r7:root}

This section contains fourteen selected A-IDs. Five concern zero diagonals
under substitution, eight concern implicitly defined power series, and one
corrects the index in an empirical theta evaluation. Different parameter
specializations share lemmas; the count is a count of source assertion
packages, not of independent general theorems. Asymptotic observations in
the same entries are outside the selected scope.

\subsection{Divisibility forced by a zero diagonal}
Throughout this subsection $A(x)=1+x+\sum_{n\ge2}a(n)x^n$.
For positive integers $m_n$, impose
\begin{equation}\label{round7:r7:zero-diagonal}
 [x^n]A\!\left(\frac{x}{A(x)^{m_n}}\right)=0\qquad(n>1).
\end{equation}
The coefficient of $a(n)$ in this equation is $1$; all other terms depend
only on $a(1),\ldots,a(n-1)$. Negative integer powers of a series with
constant term $1$ have integer coefficients. Thus the defining equations
give a unique integer series, and the same triangular uniqueness holds
over every quotient of $\mathbb Z$.

\begin{lemma}\label{round7:r7:diagonal-divisibility}
If $d\ge1$ and $n^d\mid m_n$ for every $n>1$, the solution of
\eqref{round7:r7:zero-diagonal} satisfies $n^d\mid a(n)$ for $n\ge1$.
If instead $m_n=2n-1$, then $(2n-1)\mid a(n)$ for $n\ge1$.
\end{lemma}
\begin{proof}
For $1\le k<n$ put $r=n-k$ and
$c_{n,k}=[x^r]A(x)^{-m_n k}$. The defining recursion is
\[
 a(n)=-\sum_{k=1}^{n-1}a(k)c_{n,k}.
\]
Formal differentiation gives the integer identity
\begin{equation}\label{round7:r7:derivative-divisibility}
 r c_{n,k}=-m_n k[x^{r-1}]A'(x)A(x)^{-m_nk-1}.
\end{equation}
Fix a prime $p$ with $e=v_p(n)>0$. Assume inductively that $k^d\mid a(k)$
for all $k<n$. If $v_p(k)\ge e$, the factor $a(k)$ is divisible by
$p^{de}$. Otherwise $v_p(n-k)=v_p(k)$, and
\eqref{round7:r7:derivative-divisibility} shows $p^{de}\mid c_{n,k}$.
Each summand is therefore divisible by $p^{de}$, proving the first claim
for every prime dividing $n$. The initial value $a(1)=1$ starts the induction.

For the second claim put $m=2n-1$. Equation
\eqref{round7:r7:derivative-divisibility} gives $m\mid r c_{n,k}$.
Since $2k-1=m-2r$, it also gives $m\mid (2k-1)c_{n,k}$.
The induction hypothesis $(2k-1)\mid a(k)$ now makes every summand a
multiple of $m$. Again the base case is $a(1)=1$.
\end{proof}

We will also use the elementary fact
\begin{equation}\label{round7:r7:binomial-lucas}
 p\mid {pr\choose r}\qquad(r\ge1, p\text{ prime}).
\end{equation}
Indeed, write the binomial expansion in base $p$ using
$(1+t)^p=1+t^p$ over $\mathbb F_p$. At the least position with a
nonzero digit of $r$, the corresponding digit of $pr$ is zero. The
coefficient in question therefore vanishes. This is the needed special
case of Lucas's congruence.

\begin{theorem}[A266489]\label{round7:r7entry:A266489}
For $m_n=n$, one has $n\mid a(n)$ for every $n\ge1$, and $a(n)$ is even
for every $n\ge2$.
\end{theorem}
\begin{proof}
The divisibility is Lemma~\ref{round7:r7:diagonal-divisibility}, with $d=1$.
For parity, test $U(x)=1+x$ in the defining equation modulo $2$:
\[
 [x^n]U\!\left(\frac{x}{U(x)^n}\right)
 =[x^{n-1}](1+x)^{-n}
 =(-1)^{n-1}{2n-2\choose n-1}\equiv0\pmod2
 \quad(n>1).
\]
Equation~\eqref{round7:r7:binomial-lucas} and triangular uniqueness give
$A\equiv1+x\pmod2$. These are both conjectures C1 and C2 in
\oeis{A266489} (Paul D. Hanna, February 7, 2016).
The equivalent integrality assertion in \oeis{A268293} is a consequence
and is not counted a second time.
\end{proof}

\begin{theorem}[A300732]\label{round7:r7entry:A300732}
For $m_n=2n$, the sequence with initial values $a(0)=a(1)=1$ and
\eqref{round7:r7:zero-diagonal} for $n>1$ satisfies $n\mid a(n)$ for $n\ge1$.
\end{theorem}
\begin{proof}
Apply Lemma~\ref{round7:r7:diagonal-divisibility} with $d=1$.
The conjectured divisibility in \oeis{A300732} is thereby proved.
There is a separate defining-range typo in its title: the title says
$n\ge1$, although its listed $a(1)=1$, example, and program use $n>1$.
At $n=1$ the displayed coefficient is $a(1)=1$, not zero. Our theorem
uses the explicit initial values and the consistent range $n>1$; it
does not assert the existence of a series satisfying that contradictory
literal title and $a(1)=1$ simultaneously.
\end{proof}

\begin{theorem}[A292394]\label{round7:r7entry:A292394}
For $m_n=n^2$, one has $n^2\mid a(n)$ for every $n\ge1$.
\end{theorem}
\begin{proof}
Apply Lemma~\ref{round7:r7:diagonal-divisibility} with $d=2$ to the conjecture in
\oeis{A292394} (Paul D. Hanna, September 15, 2017).
The integrality of \oeis{A292395} is the same assertion after division by
$n^2$ and adds no count.
\end{proof}

\begin{theorem}[A300734]\label{round7:r7entry:A300734}
For $m_n=2n^2$, one has $n^2\mid a(n)$ for every $n\ge1$.
\end{theorem}
\begin{proof}
Use $d=2$ in Lemma~\ref{round7:r7:diagonal-divisibility}. This proves the
displayed divisibility conjecture in \oeis{A300734} (March 11, 2018).
\end{proof}

\begin{theorem}[A395833]\label{round7:r7entry:A395833}
For $m_n=2n-1$, one has $(2n-1)\mid a(n)$ for every $n\ge1$, and
$3\mid a(n)$ for every $n\ge2$.
\end{theorem}
\begin{proof}
The first assertion is the second part of
Lemma~\ref{round7:r7:diagonal-divisibility}. To prove the second, the test series
$U=1+x$ gives, for $n>1$,
\[
 [x^n]U\!\left(\frac{x}{U(x)^{2n-1}}\right)
 =(-1)^{n-1}{3n-3\choose n-1}\equiv0\pmod3.
\]
Use \eqref{round7:r7:binomial-lucas} with $p=3$ and triangular uniqueness.
This settles both displayed congruence conjectures in \oeis{A395833}
(Paul D. Hanna, May 7, 2026).
\end{proof}

\subsection{A uniform implicit exponential family}
For a positive integer $c$, define $F_c(x)=\sum_{n\ge0}f_c(n)x^n$ by
\begin{equation}\label{round7:r7:exp-definition}
 F_c=\exp\left(x+\sum_{n\ge2}\frac{cn^2-1}{cn^2}f_c(n)x^n\right),
 \qquad f_c(0)=f_c(1)=1.
\end{equation}
Also write $B_c=xF_c'/F_c=\sum_{n\ge1}b_c(n)x^n$.
The selected entries are $F_2$ (A397346), $F_3$ (A397348),
$F_4$ (A397245), $F_5$ (A397345), and $B_2,B_3,B_4$
(A397347, A397349, A396846). The $c=1$ parity case was already selected
as A397242 in Round V; it is credited background and is not counted again.

\begin{lemma}\label{round7:r7:integer-differential}
Equation~\eqref{round7:r7:exp-definition} defines a unique integer series.
It is equivalent over $\mathbb Q[[x]]$ to
\begin{equation}\label{round7:r7:monic-ode}
 F=1+x+c\left\{x(F-1)\frac{F'+xF''}{F}
                    +x^2\frac{(F')^2}{F^2}\right\}.
\end{equation}
With $F(0)=1,F'(0)=1$, the latter equation has a unique solution over
every commutative coefficient ring. In particular reductions modulo
integers can be checked in \eqref{round7:r7:monic-ode} without dividing by $n$.
\end{lemma}
\begin{proof}
Let $\theta=x\,d/dx$. Applying $\theta^2$ to the logarithm of
\eqref{round7:r7:exp-definition} gives
\[
 \theta^2\log F=\theta^2F-\frac{F-1-x}{c}.
\]
Using $\theta^2\log F=\theta^2F/F-(\theta F)^2/F^2$ yields
\eqref{round7:r7:monic-ode}. Conversely the constant and linear coefficients
of the logarithmic equation are fixed, so reversing this calculation
over $\mathbb Q$ recovers \eqref{round7:r7:exp-definition}.
At degree $n\ge2$ the right side of \eqref{round7:r7:monic-ode} involves only
earlier coefficients: its two nonlinear terms have an extra factor of
$x$ beyond the possible appearance of $f_c(n)$. Inversion of a unit
series and formal differentiation use integer coefficient polynomials.
The recursion is therefore monic, proving all the assertions.

For reproducibility, an alternative exact recursion follows by comparing
coefficients in $F'=F(\log F)'$:
\begin{equation}\label{round7:r7:fast-recursion}
 f_c(n)=cn f_c(n-1)+n\sum_{j=2}^{n-1}(cj^2-1)
                  \frac{f_c(j)}{j}f_c(n-j),\qquad n\ge2.
\end{equation}
Induction proves $n\mid f_c(n)$, so all its divisions are exact.
The series $B_c$ is integral because $F_c$ is an integer unit series.
\end{proof}

\begin{lemma}\label{round7:r7:small-characteristics}
Put $S_p(x)=\sum_{j\ge0}x^{p^j}$ in $\mathbb F_p[[x]]$.
For odd $c$,
\[
 F_c\equiv \frac{S_2(x)}x\pmod2,
 \qquad B_c\equiv S_2(x)\pmod2.
\]
For $c\equiv1\pmod3$,
\begin{equation}\label{round7:r7:mod-three}
 F_c\equiv\frac{S_3+S_3^2-x}{x^2}\pmod3,
 \qquad B_c\equiv S_3\pmod3.
\end{equation}
For $c\equiv0\pmod3$, $F_c\equiv1+x$ and
$B_c\equiv x/(1+x)\pmod3$.
\end{lemma}
\begin{proof}
The last claim follows immediately from \eqref{round7:r7:monic-ode}.
For the first two, $S_p'=1$ and $S_p=x+S_p^p$.
In characteristic $2$ take $R=S_2/x=1/(1-S_2)$.
In characteristic $3$ take
\[
 R=\frac1{1-x-S_3^2}=\frac{S_3+S_3^2-x}{x^2}.
\]
The second equality follows on substituting $x=S_3-S_3^3$; its right
side is a genuine power series with constant term $1$.
In either characteristic direct differentiation gives
$\theta\log R=S_p$ and hence $\theta^2\log R=x$.
For characteristic $3$, for example,
\[
 \frac{xR'}R=\frac{x(1-S_3)}{1-x-S_3^2}=S_3,
\]
because $1-x-S_3^2=(1-S_3)^2(1+S_3)$.
It follows that $\theta^2R=R(x+S_p^2)$.
In both cases $R(1-x-S_p^2)=1$, so
\[
 R-1-x=\theta^2R-\theta^2\log R.
\]
This is \eqref{round7:r7:monic-ode} with $c=1$. Its uniqueness proves the
claims for the specified residues of $c$.
\end{proof}

\begin{theorem}[A397348]\label{round7:r7entry:A397348}
For $F_3$, $f_3(n)$ is odd if and only if $n+1$ is a power of $2$.
\end{theorem}
\begin{proof}
Apply Lemma~\ref{round7:r7:small-characteristics} with $c=3$. It proves the
parity conjecture in \oeis{A397348} for every $n\ge0$.
The asymptotic observation in that entry is not selected.
\end{proof}

\begin{theorem}[A397345]\label{round7:r7entry:A397345}
For $F_5$, $f_5(n)$ is odd if and only if $n+1$ is a power of $2$.
\end{theorem}
\begin{proof}
Use the same lemma with $c=5$. This is the parity conjecture in
\oeis{A397345}. These two applications extend the earlier $c=1$
project result; the common family argument is not counted separately.
\end{proof}

\begin{theorem}[A397245]\label{round7:r7entry:A397245}
For $F_4$ and every $n\ge0$, $f_4(n)\equiv1\pmod3$ precisely when
$n+2=3^j$ or $2\cdot3^j$ for a nonnegative integer $j$.
One has $f_4(n)\equiv2\pmod3$ precisely when
$n+2=3^i+3^j$ with $0\le i<j$. All other residues are zero.
\end{theorem}
\begin{proof}
Use \eqref{round7:r7:mod-three}. The three types of exponents in $S_3+S_3^2$
have, respectively, one ternary digit $1$, one ternary digit $2$, or
two ternary digits $1$. They are disjoint. The corresponding coefficients
are $1,1,2$. Subtracting $x$ removes the sole negative-power term before
division by $x^2$. This proves both congruence conjectures in
\oeis{A397245}. The same lemma additionally settles the previously
unselected mod-$3$ assertions for A397242; that already selected A-ID
receives no new count here.
\end{proof}

\begin{theorem}[A397349]\label{round7:r7entry:A397349}
For $B_3$, $b_3(n)$ is odd exactly when $n$ is a power of $2$.
For $n\ge1$, $b_3(n)\equiv(-1)^{n-1}\pmod3$.
\end{theorem}
\begin{proof}
The identity $B_3=xF_3'/F_3$ follows by differentiating the defining
logarithmic series in \oeis{A397349}; the coefficient relation is
$b_3(n)=(3n^2-1)f_3(n)/(3n)$ for $n>1$, with $b_3(1)=1$.
The two reductions in Lemma~\ref{round7:r7:small-characteristics} give
$B_3\equiv S_2\pmod2$ and $B_3\equiv x/(1+x)\pmod3$.
These settle both displayed congruence conjectures.
\end{proof}

\subsection{Exact rational certificates modulo eight}
For a unit series $R$ set
\[
 H_c(R)=R-1-x-cx(R-1)(R'+xR'')/R-cx^2(R')^2/R^2.
\]
Define the rational functions
\begin{align*}
 R_2&=1+x+\frac{4x^2+2x^3+2x^5}{1-x^4},\\
 R_4&=1+x+\frac{4x^3}{(1+x)^2},\qquad
 J=x^4-2x^3+4x^2+1,\quad K=5x^3+3x^2+3x+1.
\end{align*}
The following finite polynomial identity is an all-coefficient certificate:
\begin{equation}\label{round7:r7:rational-certificate}
 H_2(R_2)=\frac{8x^3 N(x)}{(x-1)^3(x^2+1)^3J^2},
\end{equation}
where
\begin{align*}
 N(x)={}&5x^{12}-18x^{11}+54x^{10}-95x^9+118x^8-28x^7-68x^6\\
        &+174x^5-83x^4+62x^3-2x^2+x+8.
\end{align*}
It is verified by ordinary rational differentiation and clearing
denominators; each denominator has constant term $1$ or $-1$.
Consequently $F_2\equiv R_2\pmod8$ by monic uniqueness.

For $c=4$, equation~\eqref{round7:r7:monic-ode} first gives $F_4\equiv1+x\pmod4$.
Substitute $1+x$ into the expression in braces and reduce it modulo $2$:
\[
 \frac{x^2}{1+x}+\frac{x^2}{(1+x)^2}
 =\frac{x^2(x+2)}{(1+x)^2}\equiv\frac{x^3}{(1+x)^2}\pmod2.
\]
Thus $F_4\equiv R_4\pmod8$.
Two further exact identities certify the logarithmic derivatives:
\begin{align}
 \frac{xR_2'}{R_2}-\frac{x+3x^2+3x^3+3x^4}{1-x^4}
 &=\frac{4x^2(x^6-x^5+2x^4+3x^2+2x-1)}
 {(x-1)(x+1)(x^2+1)J},\label{round7:r7:log-two}\\
 \frac{xR_4'}{R_4}-\frac{x+7x^2+5x^3+7x^4}{1-x^4}
 &=\frac{8x^2(5x^4+2x^3+6x^2+x+1)}
 {(x-1)(x^2+1)K}.\label{round7:r7:log-four}
\end{align}

\begin{theorem}[A397346]\label{round7:r7entry:A397346}
Starting at $n=2$, the residues of $f_2(n)$ modulo $8$ repeat
$4,2,0,2$.
\end{theorem}
\begin{proof}
The coefficients of $R_2-1-x$ have precisely this pattern. The identity
\eqref{round7:r7:rational-certificate} therefore proves the conjecture in
\oeis{A397346} for its stated range.
\end{proof}

\begin{theorem}[A397347]\label{round7:r7entry:A397347}
Starting at $n=1$, the residues of $b_2(n)$ modulo $4$ repeat
$1,3,3,3$.
\end{theorem}
\begin{proof}
The definition in \oeis{A397347} gives $B_2=xF_2'/F_2$.
Formal differentiation and inversion preserve coefficientwise congruences
of integer unit series. Reduce \eqref{round7:r7:log-two} modulo $4$ and use
$F_2\equiv R_2\pmod8$. This proves the selected residue pattern.
\end{proof}

\begin{theorem}[A396846]\label{round7:r7entry:A396846}
For $n\ge1$, $b_4(n)\equiv1\pmod3$ if $n$ is a power of $3$, and
$b_4(n)\equiv0\pmod3$ otherwise. Its residues modulo $8$, starting at
$n=1$, repeat $1,7,5,7$.
\end{theorem}
\begin{proof}
The logarithmic definition in \oeis{A396846} gives $B_4=xF_4'/F_4$;
the correct coefficient relation is
$b_4(n)=(4n^2-1)f_4(n)/(4n)$ for $n>1$.
Lemma~\ref{round7:r7:small-characteristics} gives $B_4\equiv S_3\pmod3$.
For the second assertion use $F_4\equiv R_4\pmod8$ and
\eqref{round7:r7:log-four}. Both source congruence conjectures are settled.
\end{proof}

\subsection{A binary restriction from a cubic equation}
\begin{theorem}[A397244]\label{round7:r7entry:A397244}
Let $V(x)=\sum_{n\ge0}v(n)x^n$, $v(0)=v(1)=1$, be defined by
\[
 (2n)[x^n]V^{2n}=(2n-1)[x^n]V^{2n+1}\qquad(n>1).
\]
For $n>0$, $v(n)$ is odd exactly when $n=2k+1$ and the binary expansion
of $k$ has no adjacent $1$ digits. Equivalently $k$ is a Fibbinary number.
\end{theorem}
\begin{proof}
The coefficient of $v(n)$ in the difference of the two sides is
$(2n)^2-(2n-1)(2n+1)=1$. Hence this is again an integer monic
triangular definition, valid after reduction modulo $2$.
Let $U=1+x/U^2$, an integer unit series. Lagrange inversion, with
$W=U-1=x(1+W)^{-2}$, gives for $n>1$
\[
 [x^n]U^{2n+1}
 =\frac{2n+1}{n}[t^{n-1}](1+t)^{2n}(1+t)^{-2n}=0.
\]
Thus $U$ satisfies the defining relation for $V$ modulo $2$, and
$V\equiv U\pmod2$ by uniqueness.
In characteristic $2$, $D=U-1$ satisfies $D=x+D^3$.
The integer solution $E=x+E^3$ has
\[
 [x^{2k+1}]E=\frac1{2k+1}{3k\choose k},\qquad [x^{2k}]E=0.
\]
These follow from Lagrange inversion applied to $E=x/(1-E^2)$.
Its reduction is $D$. Since $2k+1$ is odd, the displayed odd-index
coefficient has the parity of ${3k\choose k}$.
The base-$2$ binomial rule says this coefficient is odd precisely when
adding $k$ and $2k$ produces no carry, equivalently $k\mathbin{\&}(2k)=0$.
That is exactly the absence of adjacent binary $1$ digits. This proves
the parity conjecture in \oeis{A397244}. Its asymptotic assertion is not
selected, and no potentially misindexed transform formula from that
entry is needed.
\end{proof}

\subsection{An offset correction in an empirical theta value}
\begin{theorem}[A030212: corrected index]\label{round7:r7entry:A030212}
Use the actual OEIS indexing $a(1)=1,a(2)=-4,a(3)=0,\ldots$ for
Glaisher's sequence A030212. Then
\begin{align*}
 \sum_{n\ge1}a(n)e^{-n\pi}
  &=\frac{\pi^{5/2}}{64\Gamma(3/4)^{10}},\\
 \sum_{n\ge1}a(n)e^{-(n-1)\pi}
  &=\frac{e^\pi\pi^{5/2}}{64\Gamma(3/4)^{10}}.
\end{align*}
The second value is the constant printed in the empirical formula;
it requires a shift of the exponent by one relative to the actual index.
\end{theorem}
\begin{proof}
Put $q=e^{-\pi}$, $T=\pi^{1/4}/\Gamma(3/4)$,
$\phi(z)=\sum_{j\in\mathbb Z}z^{j^2}$, and
$\psi(z)=\sum_{j\ge0}z^{j(j+1)/2}$.
The previously recorded generating function in \oeis{A030212} is
\[
 \sum_{n\ge1}a(n)z^n
 =z\prod_{j\ge1}(1-z^j)^4(1-z^{2j})^2(1-z^{4j})^4
 =z\phi(z)^2\psi(-z)^8.
\]
The product/theta identity is also an immediate application of the Jacobi
products; it is credited classical input, not a new identity.
For completeness the required special values are
\[
 \phi(q)=T,\qquad \phi(-q)=2^{-1/4}T,\qquad
 \psi(q^2)=2^{-5/4}q^{-1/4}T.
\]
At nome $e^{-\pi}$ the elliptic modulus is $1/\sqrt2$.
The theta--elliptic relation, the classical value
$K(1/\sqrt2)=\Gamma(1/4)^2/(4\sqrt\pi)$, and the gamma reflection
formula yield $\theta_3(0,q)=T$.
Symmetry gives $\theta_2(0,q)=\theta_4(0,q)$, and
$\theta_3^4=\theta_2^4+\theta_4^4$ gives the factor $2^{-1/4}$.
Finally $\theta_2(0,q)=2q^{1/4}\psi(q^2)$.
The standard duplication identity
$\psi(-q)^2=\phi(-q)\psi(q^2)$ therefore implies
$\psi(-q)^4=2^{-3}q^{-1/2}T^4$.
Substitution in $q\phi(q)^2\psi(-q)^8$ gives $T^{10}/64$, proving
both equalities. All products are analytic for $|z|<1$, so these sums
converge absolutely at $q$.

The consulted source's empirical expression starts at $n=0$ although
the entry has offset $1$. Even if one naturally extends $a(0)=0$, its
unshifted left side is smaller than the printed constant by the factor
$e^{-\pi}$. Reindexing $a(n)$ as $a(n+1)$, or replacing the exponent
$n$ by $n-1$ and starting at $1$, repairs the formula. The linked decimal
entry A388390 is not counted separately. The theta evaluations used here
were already reproduced in Round VI and are independently included above;
the selected contribution is this source-index correction.
\end{proof}

The classical identities in the last proof are documented in
\href{https://dlmf.nist.gov/20.7}{DLMF \S20.7},
\href{https://dlmf.nist.gov/20.9}{\S20.9}, and
\href{https://dlmf.nist.gov/19.20}{\S19.20}. Source records, boundary
notes, and exact rational certificates accompany the reproducibility
packet. Finite coefficient checks supplement the universal arguments;
they are not their logical basis.

\section{Arithmetic identities, asymptotics, and two corrections}
\label{round7:r7:arith}

The ten packages in this section concern explicitly uncertain statements on
the cited OEIS entries, inspected on 20 October 2026.  Eight selected
statements are proved.  The proposed bound in A133205 is false, and the
zero-support assertion in A366561 needs a correction; its accompanying
general formulas are proved.  Known identities used as inputs are credited.
Finite computations in the accompanying script check transcription and
independent definitions; the all-index arguments are the proofs below.

\subsection{A092843: a shifted-totient divisor sum}
\label{round7:r7entry:A092843}
\paragraph{Source and target.}
\href{https://oeis.org/A092843}{A092843} defines
$a(n)=\sum_{d\mid n,\ d>1}\varphi(d-1)$.
Kotesovec's 25 June 2024 conjecture is
$\sum_{n\leq N}a(n)\sim N^2/2$.

\paragraph{Proof.}
Put $\Phi(y)=\sum_{1\leq m\leq y}\varphi(m)$, with $\Phi(0)=0$.
The classical elementary estimate
\[
 \Phi(y)=\frac{y^2}{2\zeta(2)}+O(y\log(2y))
\]
follows directly from
$2\Phi(y)-1=\sum_{d\leq y}\mu(d)\lfloor y/d\rfloor^2$:
replace each square by $(y/d)^2+O(y/d)$, and bound the tail of
$\sum\mu(d)/d^2=1/\zeta(2)$.  Reversing two finite sums gives
\[
 \sum_{n\leq N}a(n)
 =\sum_{j\leq N}\Phi(\lfloor N/j\rfloor-1)
 =\frac{N^2}{2\zeta(2)}\sum_{j\leq N}\frac1{j^2}
    +O\!\left(N\sum_{j\leq N}\frac{\log(2N/j)}j\right).
\]
The error is $O(N\log^2(2N))$, and the omitted tail of the zeta sum
contributes only $O(N)$.  Consequently
\[
 \boxed{\displaystyle \sum_{n\leq N}a(n)
       =\frac{N^2}{2}+O(N\log^2(2N)).}
\]
This proves the conjecture with an error term.

\subsection{A108625: an arbitrary-parameter binomial identity}
\label{round7:r7entry:A108625}
\paragraph{Source and scope.}
The 12 October 2024 conjecture in
\href{https://oeis.org/A108625}{A108625} asks, for every positive integer
$m$ and every $n\geq0$, whether
\[
 \sum_{k=0}^n(-1)^{n+k}\binom nk T(mn,k)
 =\frac{((m+1)n)!}{((m-1)n)!(n!)^2}.
\]
The entry reports verification for $m\leq10$.  We use and credit
Peter Bala's 26 February 2024 identity
$T(r,k)=\sum_i\binom ri\binom{r+i}i\binom ki$;
the selected claim is the universal $m$ assertion, not a new formula for
the crystal-ball array or any of its unrelated congruences.

\paragraph{Proof.}
For $0\leq i\leq n$, the finite binomial theorem gives
\[
 \sum_{k=i}^n(-1)^{n-k}\binom nk\binom ki
 =\binom ni\sum_{\ell=0}^{n-i}(-1)^{n-i-\ell}
       \binom{n-i}\ell
 =\begin{cases}1&i=n,\\0&i<n.\end{cases}
\]
Substitute the credited formula for $T(r,k)$ and interchange finite sums.
The result is $\binom rn\binom{r+n}n$ for every $r\geq n$.
Putting $r=mn$ gives the stated factorial quotient.  This also includes
$n=0$, when both sides are $1$.

\subsection{A133205: the proposed logarithmic bound is false}
\label{round7:r7entry:A133205}
\paragraph{Source and result.}
\href{https://oeis.org/A133205}{A133205} is completely multiplicative with
$a(p)=p(p+1)/2$.  Kotesovec's 28 January 2025 conjecture asserts
$S(N):=\sum_{n\leq N}a(n)=O(N^3/\log N)$.
In fact,
\[
 \boxed{\displaystyle \frac{S(N)\log N}{N^3}\longrightarrow+\infty.}
\]

\paragraph{Proof.}
We use the classical prime number theorem, as recorded in
\href{https://dlmf.nist.gov/27.12}{DLMF 27.12}.  Partial summation gives
\[
 \sum_{q\leq X\atop q\ {\rm prime}} a(q)
 =\frac12\sum_{q\leq X}(q^2+q)
 \sim\frac{X^3}{6\log X}.
\]
Fix $M\geq2$.  The integers $pq$, with $p\leq M<q$ both prime, are
distinct for different pairs.  Positivity and multiplicativity therefore
give
\[
 S(N)\geq\sum_{p\leq M\atop p\ {\rm prime}}
 a(p)\sum_{M<q\leq N/p\atop q\ {\rm prime}}a(q).
\]
Here $M$ is fixed before $N$ tends to infinity.  Hence
\[
 \liminf_{N\to\infty}\frac{S(N)\log N}{N^3}
 \geq \frac1{12}\sum_{p\leq M}\frac{p+1}{p^2}.
\]
The right side tends to infinity with $M$, because $\sum_p1/p$ diverges.
For completeness, convergence of that sum would make
$\prod_p(1-1/p)^{-1}$ bounded, whereas its finite products dominate
arbitrarily long harmonic sums.  For any prescribed real $C$, choose
$M$ so the displayed finite lower bound exceeds $C$; the definition of
the lower limit then proves the boxed assertion.  Thus no constant in
the proposed big-$O$ bound can exist.  No numerical threshold or
unproved assertion about prime pairs is needed.

\subsection{A140110: the gcd-sum characterization}
\label{round7:r7entry:A140110}
\paragraph{Source and target.}
\href{https://oeis.org/A140110}{A140110} contains those positive integers
whose successive divisor ratios reduce to ratios of consecutive integers.
Ratajczak's 9 January 2017 empirical assertion identifies precisely these
integers by
\[
 n=1+\sum_{i=1}^{\tau(n)-1}\gcd(d_i,d_{i+1}),
 \qquad 1=d_1<\cdots<d_{\tau(n)}=n.
\]
The entry's separate, already proved divisibility-by-$4$-or-$6$ statement
is not part of this selection.

\paragraph{Proof.}
Let $g_i=\gcd(d_i,d_{i+1})$.  Since $g_i$ divides the positive difference
$d_{i+1}-d_i$, one has $g_i\leq d_{i+1}-d_i$.  Summing gives
$\sum_i g_i\leq n-1$.  Equality holds exactly when every individual
inequality is equality.  Write $d_i=g_i u$ and $d_{i+1}=g_i v$ with
$\gcd(u,v)=1$.  The equality $g_i=d_{i+1}-d_i$ is exactly $v-u=1$.
This is the required reduced-ratio condition, proving both directions.
For $n=1$ the empty sum also gives the correct result.

\subsection{A179873: monotonicity and oddness of integral means}
\label{round7:r7entry:A179873}
\paragraph{Source and target.}
\href{https://oeis.org/A179873}{A179873} lists the integral values, in
increasing order of $h$, of
\[
 B(h)=\frac{\sum_{1\leq k<h,\ (k,h)=1}k^2}
             {\sum_{1\leq k<h,\ (k,h)=1}k},
 \qquad B(1)=1.
\]
Its comment conjectures that the list is nondecreasing and consists of
odd integers.  The general arithmetic formula for $B(h)$ already appears
in Hasler's program on \href{https://oeis.org/A175505}{A175505} and in
Eldar's later programs; we derive it here and use it as a credited input.

\paragraph{Proof.}
For $h>1$, pairing $k$ with $h-k$ gives
$\sum_{(k,h)=1}k=h\varphi(h)/2$.  Inclusion-exclusion gives
\[
 \sum_{(k,h)=1}k^2
 =\sum_{d\mid h}\mu(d)d^2\sum_{j=1}^{h/d}j^2
 =\frac{h^2\varphi(h)}3+\frac h6\prod_{p\mid h}(1-p).
\]
Put $t=h/\operatorname{rad}(h)$ and $\varepsilon=(-1)^{\omega(h)}$.
Then
\[
 B(h)=\frac{2h+\varepsilon/t}{3}.
\]
If $B(h)$ is integral, $3tB(h)=2ht+\varepsilon$ forces $t=1$.
Thus every selected $h>1$ is squarefree and
$B(h)=(2h+\mu(h))/3$.  Its numerator is odd, so an integral $B(h)$ is
odd.  If $h'<h$ are two selected indices greater than $1$, then
\[
 B(h)-B(h')
 =\frac{2(h-h')+\mu(h)-\mu(h')}{3}\geq0.
\]
The separately prescribed $B(1)=1$ is no larger than any selected value
with $h>1$.  This proves the complete conjecture, including its initial
repeated value $1$.

\subsection{A309263: an infinite family of admissible divisors}
\label{round7:r7entry:A309263}
\paragraph{Source and target.}
\href{https://oeis.org/A309263}{A309263} selects the members of A140110
that are not divisible by $6$.  Its conjecture says that $4\cdot5^r$
is present for every $r\geq0$.

\paragraph{Proof.}
The complete ordered divisor list of $4\cdot5^r$ consists of the blocks
\[
 (5^j,\ 2\cdot5^j,\ 4\cdot5^j),\qquad 0\leq j\leq r.
\]
They occur in precisely this order since $4\cdot5^j<5^{j+1}$.
Within a block each successive ratio is $2/1$; between blocks it is
$5/4$.  Both are ratios of consecutive coprime integers.  Finally
$3\nmid4\cdot5^r$, so the exclusion of multiples of $6$ never applies.

\subsection{A356459: the exact asymptotic constant}
\label{round7:r7entry:A356459}
\paragraph{Source and target.}
\href{https://oeis.org/A356459}{A356459} defines
$a(n)=n!\sum_{k\leq n}\sum_{d\mid k}d/(k/d)!$.
Kotesovec's 12 August 2025 conjecture has the form
$a(n)\sim c\,n!n^2$ with $c=0.5732\ldots$.

\paragraph{Proof.}
Writing $k=rd$ gives the exact identity
\[
 \frac{a(n)}{n!}
 =\frac12\sum_{r=1}^n\frac{\lfloor n/r\rfloor
              (\lfloor n/r\rfloor+1)}{r!}.
\]
For $u\geq1$,
$|\lfloor u\rfloor(\lfloor u\rfloor+1)-u^2|\leq u$.
The convergent sums of the resulting errors give
\[
 \boxed{\displaystyle
 \frac{a(n)}{n!}=c n^2+O(n),\qquad
 c=\frac12\sum_{r=1}^{\infty}\frac1{r!r^2}.}
\]
More explicitly, the absolute error is at most
\[
 \frac n2\sum_{r\leq n}\frac1{r!r}
 +\frac{n^2}{2}\sum_{r>n}\frac1{r!r^2},
\]
which is $O(n)$.  The verification file encloses $c$ in an exact rational
interval: after $M$ terms the omitted tail is at most
$\frac{1}{2(M+1)!(M+1)^2}\frac{M+2}{M+1}$.
This identifies the conjectured decimal constant without treating a
decimal fit as a proof.

\subsection{A366561: general counting formulas and the missing zeros}
\label{round7:r7entry:A366561}
\paragraph{Source and selected scope.}
Let
\[
 T(n,k)=\#\{(x,y)\in(\mathbb Z/n\mathbb Z)^2:
                    \gcd(x^2-y^2,n)=k\}.
\]
The 17 June 2025 conjectures of Oudra on
\href{https://oeis.org/A366561}{A366561} propose a support criterion and
several equivalent divisor-sum formulas.  The formulas are correct.
The proposed criterion $T(n,k)=0$ exactly when $k\nmid n$ is false:
\[
 \boxed{T(4,2)=0\quad\hbox{although }2\mid4.}
\]
Squares modulo $4$ are $0$ and $1$, so their difference cannot be $2$.
The corrected criterion, for positive $n,k$, is
\[
 \boxed{T(n,k)>0\iff k\mid n\ \hbox{and}\
       \neg\bigl(v_2(n)\geq2\ \hbox{and}\ v_2(k)=1\bigr).}
\]
The diagonal and first-column identifications with A062803 and A082953
were already recorded on those entries; they are credited consequences,
not additional selected problems.

\paragraph{Local enumeration.}
The Chinese remainder theorem makes the count a product over prime
powers.  Let $C_p(a,b)=T(p^a,p^b)$ for $a\geq1$ and $0\leq b\leq a$.
For odd $p$, the substitution $(u,v)=(x-y,x+y)$ is bijective.  Counting
the possible valuations of $u$ and $v$ gives
\[
 C_p(a,b)=
 \begin{cases}
 (b+1)(p-1)^2p^{2a-b-2},&b<a,\\
 ((a+1)p-a)p^{a-1},&b=a.
 \end{cases}
\]
For $b<a$, there are $b+1$ possible pairs of valuations, each contributing
$(p-1)^2p^{2a-b-2}$.  For $b=a$, sum over $v_p(u)=i<a$ and include $u=0$;
the result is $a(p-1)p^{a-1}+p^a$.

For $p=2$, the same map has a kernel of size $2$ and image the pairs of
equal parity.  Odd $u,v$ give $2^{2a-1}$ original pairs.  Even $u,v$ give
$uv=4u'v'$, with $u',v'$ modulo $2^{a-1}$.  Thus
\[
 C_2(a,b)=
 \begin{cases}
 2^{2a-1},&b=0,\\
 0,&b=1<a,\\
 (b-1)2^{2a-b-1},&2\leq b<a,\\
 a2^a,&b=a.
 \end{cases}
\]
For the last case and $a\geq2$, reduce $u',v'$ modulo $2^{a-2}$.
There are four lifts, and the original map contributes another factor
$2$; the product-zero count modulo $2^e$ is $(e+2)2^{e-1}$ for $e\geq1$.
The cases $a=1,2$ are respectively $2$ and $8$ by direct parity counting.
This proves the local table, including all boundary cases.  Every odd
prime local count is positive; the displayed $2$-adic zero is the only
possible additional obstruction.  The corrected support follows.

\paragraph{The source formulas.}
For $k\mid n$ define
\[
 F(n,k)=\gcd(n,2)(-1)^k\varphi(n/k)
       \sum_{d\mid k}(-1)^{k/d}\frac{k}{d}\varphi(dn/k).
\]
For odd $n$ the two signs cancel, giving the source's unsigned formula.
The expression factors into local factors: signs depend only on the
$2$-part of $k/d$, and every other factor is multiplicative.  For odd
$p$, its local sum gives the first table.  At $2$, if $1\leq b<a$, the
inner signed sum is $(b-1)2^{a-1}$; if $b=a$, it is $a2^{a-1}$.
For $b=0$ it is $-\varphi(2^a)$ and the outer sign is negative.
Multiplication by the outer factors gives exactly the second table.
Therefore $T(n,k)=F(n,k)$ for every $k\mid n$.

Finally let $H(n,k)=\varphi(n/k)$ if $k\mid n$, and $0$ otherwise, as in
\href{https://oeis.org/A054523}{A054523}.  Extending the count by zero for
$k>n$, the source's all-positive-integer version is also proved:
\[
 T(n,k)=\gcd(n,2)(-1)^k H(n,k)
       \sum_{d\mid k}(-1)^{k/d}\frac{k}{d}H(dn,k).
\]
When $k\nmid n$ both sides vanish; when $k\mid n$ it is the identity
already proved.  The empirical row-sum assertion also follows, since
the gcd values partition the $n^2$ ordered pairs.

\subsection{A375838: the triangular recursion for rooted chains}
\label{round7:r7entry:A375838}
\paragraph{Source and target.}
Kurkov's 21 June 2025 conjecture in
\href{https://oeis.org/A375838}{A375838} defines
\[
 R(n,n)=1,\qquad
 R(n,k)=(k+1)\left(\sum_{i=0}^{n-1}R(n-1,i)
                   +\sum_{j=0}^{k-1}R(n-1,j)\right),\quad k<n,
\]
and asserts that $R(n,0)=a(n)$.  The known rooted-chain recurrence,
credited to Mohapatra on 1 July 2025, is
\[
 a(0)=1,\qquad
 a(n)=1+\sum_{r=1}^{n-1}c(n,r)a(r)\quad(n\geq1),
\]
where $c(n,r)$ denotes the unsigned Stirling number of the first kind.

\paragraph{Proof.}
Write $b(n)=R(n,0)$, and for $k<n$ put $Q(n,k)=R(n,k)/(k+1)$.
Subtracting consecutive defining equations gives
\[
 Q(n,k)=Q(n,k-1)+kQ(n-1,k-1),\qquad Q(n,0)=b(n).
\]
Induction on $k$, using the Stirling recurrence, yields
\[
 Q(n,k)=\sum_{j=0}^k c(k+1,k+1-j)b(n-j),\qquad k<n.
\]
On the other hand, the definition and the boundary value
$R(n-1,n-1)=1$ give
$Q(n,n-1)=2b(n)-1$.  Substitution of $k=n-1$ in the previous identity
therefore gives
\[
 2b(n)-1=\sum_{r=1}^n c(n,r)b(r)
         =b(n)+\sum_{r=1}^{n-1}c(n,r)b(r).
\]
Thus $b$ satisfies exactly the credited recurrence for $a$, including
$b(0)=1$.  Inductive uniqueness proves $b(n)=a(n)$ for all $n\geq0$.
The related \href{https://mathoverflow.net/questions/496310/recursion-for-a375838}
{MathOverflow question} was also inspected; its retrieved page contains
the conjecture and comments but no proof of it.

\subsection{A385001: every fixed column has the proposed mean}
\label{round7:r7entry:A385001}
\paragraph{Source and target.}
In \href{https://oeis.org/A385001}{A385001}, a partition with distinct
part sizes $m_1<\cdots<m_k$ and positive multiplicities $r_1,\ldots,r_k$
has $\prod_i r_i$ choices of designated summands.
Kotesovec's 1 August 2025 conjecture states, for each fixed $k\geq1$,
\[
 \sum_{j=1}^N T(j,k)\sim
 \frac{\pi^{2k}}{(2k)!(2k+1)!}\,N^{2k}.
\]
The underlying MacMahon generating functions are classical and appear
in \href{https://oeis.org/A060043}{A060043} and
\href{https://arxiv.org/abs/1010.5769}{Andrews--Rose (2010)}.
Only the marked summatory limit is selected here.

\paragraph{Proof.}
For a fixed tuple $\mathbf m=(m_1,\ldots,m_k)$, let
\[
 S_{\mathbf m}(N)=
 \sum_{r_1,\ldots,r_k\geq1\atop \sum_i m_i r_i\leq N}\prod_i r_i.
\]
The Riemann sum with coordinates $u_i=m_i r_i/N$ gives
\[
 \lim_{N\to\infty}\frac{S_{\mathbf m}(N)}{N^{2k}}
 =\frac1{\prod_i m_i^2}
   \int_{u_i\geq0,\ \sum u_i\leq1}\prod_i u_i\,d\mathbf u
 =\frac1{(2k)!\prod_i m_i^2}.
\]
The last simplex integral follows by iterated beta integration.
There is a summable bound independent of $N$: if every $m_i\leq N$,
\[
 S_{\mathbf m}(N)
 \leq\prod_i\sum_{r\leq N/m_i}r
 \leq\frac{N^{2k}}{\prod_i m_i^2};
\]
otherwise the sum is zero.  Dominated convergence for the countable
sum over increasing tuples is therefore applicable.  It gives
\[
 \lim_{N\to\infty}N^{-2k}\sum_{j\leq N}T(j,k)
 =\frac1{(2k)!}\sum_{m_1<\cdots<m_k}\frac1{m_1^2\cdots m_k^2}.
\]
The classical product
\[
 \prod_{m\geq1}\left(1+\frac{z^2}{m^2}\right)
 =\frac{\sinh(\pi z)}{\pi z}
\]
(\href{https://dlmf.nist.gov/4.36.E1}{DLMF 4.36.1}) identifies the final
sum as $\pi^{2k}/(2k+1)!$.  This proves the conjecture for every fixed
$k$, with no Tauberian assumption and no numerical extrapolation.

\section{Generating functions and exact recurrence proofs}
\label{round7:r7:gf}

This section treats 31 distinct selected assertions.  Twenty-five recurrences
are proved as stated, four require an explicit starting-index qualification,
and two printed recurrences require a correction.  Each target was explicitly
marked conjectural in the consulted OEIS entry.  Existing generating functions,
finite sums, and product definitions are credited inputs.  The selected scope
does not claim that the sequence, its generating function, or the method of
holonomic elimination is new.

The retained list excludes all 608 previously selected IDs and the inspected
Fischer 2020 and Ekhad--Yang 2017 families.  It also avoids the 97 IDs in the
inspected OEIS Open Appendix A baseline; that appendix is only a bounded part
of that project.  In particular, A065942 and A226302 were removed because their
marked recurrences are exact left multiples of already recorded stronger
recurrences.  A359643 was removed after the official Alkahest guide identified
an earlier lower-order recurrence result.  The audit file records these
decisions and their limited scope.  A remaining conjecture label does not by
itself establish historical novelty or official closure.

\subsection*{The coefficient certificate}
Write $a_n$ for the sequence with its actual OEIS offset $s$ and set
$F(x)=\sum_{n\geq s}a_nx^n$, with zero coefficients below $s$.
Every ordinary recurrence below uses the convention
\begin{equation}
 R_n=\sum_{j=0}^r p_j(n)a_{n-j},\qquad
 L=\sum_{j=0}^r x^j p_j(\theta+j),\qquad
 \theta=x\frac{d}{dx}.
 \label{round7:r7:gf-operator}
\end{equation}
Indeed, $\theta x^{n-j}=(n-j)x^{n-j}$, so $[x^n]LF=R_n$.
Consequently the exact identity $LF=P(x)$ proves $R_n=0$ for all
$n>\deg P$.  We additionally require $n\geq s+r$, so every sequence
index in the displayed recurrence belongs to its stated domain.
The coefficients of $P$ give the exceptional initial residuals exactly.
No inference from agreement of finitely many terms is involved.

Here is a finite rational-arithmetic recipe for checking every displayed
identity.  If $x=\xi(t)$, put $M=\xi/\xi'$, and write
$F=\sum_i R_i(t)\sqrt{D_i(t)}$, where $D_0=1$ and every square root is
the formal branch with constant term one.  Define
\begin{equation}
 H_{i,0}=R_i,\qquad
 H_{i,m+1}=M\left(H_{i,m}'+\frac{D_i'}{2D_i}H_{i,m}\right),\qquad
 c_m=\sum_{j=0}^r\xi^j[z^m]p_j(z+j).
 \label{round7:r7:gf-rational-check}
\end{equation}
Then $LF=\sum_i\sqrt{D_i}\sum_m c_mH_{i,m}$.
Thus it suffices to check the rational identities
$\sum_m c_mH_{0,m}=P(\xi)$ and
$\sum_m c_mH_{i,m}=0$ for $i>0$.
For an unparametrized series take $t=x$, $\xi=x$, $M=x$.
These identities use only differentiation, multiplication, and cancellation
of rational functions; their exact numerators vanish.  The manuscript gives
all $p_j$, all generating functions, and every boundary polynomial, so this
recipe is an independently reproducible finite certificate.  The portable
verification script carries it out over exact rational numbers and also
checks offsets and independently generated initial coefficients.

\subsection*{Two formal-series facts used below}
Let $t=x\phi(t)$ with $\phi(0)=1$.  The unique solution has $t(0)=0$.
For every formal series $B$,
\begin{equation}
 \sum_{n\geq0}x^n[u^n]B(u)\phi(u)^n
 =\frac{B(t)}{1-t\phi'(t)/\phi(t)}.
 \label{round7:r7:gf-diagonal}
\end{equation}
To prove this, extract $[x^n]$ on the right by a formal residue and
substitute $x=u/\phi(u)$.  The differential
$dx=(1-u\phi'(u)/\phi(u))\,du/\phi(u)$ cancels the displayed denominator,
leaving $\operatorname{Res}_u B(u)\phi(u)^n u^{-n-1}du$.
This is exactly the coefficient on the left, including $n=0$.

For an exponential generating function $F(x)=\sum a_nx^n/n!$ and a
recurrence of order $r$, replace $p_j(n)$ in
\eqref{round7:r7:gf-operator} by
\begin{equation}
 q_j(n)=p_j(n)\prod_{h=0}^{r-j-1}(n-j-h).
 \label{round7:r7:gf-egf}
\end{equation}
The resulting coefficient of $x^n$ is $R_n/(n-r)!$ for $n\geq r$,
because $q_j(n)/(n-j)!=p_j(n)/(n-r)!$.  An empty product is one.
This explains the conversion used for A096471.



\subsection{A026026: the selected recurrence}

\label{round7:r7entry:A026026}

\paragraph{Source and selected assertion.}

\href{https://oeis.org/A026026}{A026026}, consulted on 20 October 2026, records the conjecture attributed to R. J. Mathar, Jun 22 2013.

The source offset is $s=2$.  In the convention of \eqref{round7:r7:gf-operator}, the selected recurrence is $R_n=0$ for every integer $n\geq 4$, with $r=2$ and

\[\begin{aligned}
p_{0}(n)&=- \left(n + 3\right) \left(11 n - 38\right)\\
p_{1}(n)&=2 \left(35 n^{2} - 86 n - 102\right)\\
p_{2}(n)&=- 4 \left(2 n - 5\right) \left(13 n - 30\right)\\
\end{aligned}\]

\paragraph{Known inputs and scope.}

The Catalan generating function is an existing input. Krause (May 9, 2026) also records a binomial difference, and Drube (2022) gives the path setting. Only Mathar's marked order-two recurrence is selected; neither the closed form nor the path interpretation is claimed as new.

\paragraph{Proof.}

With $C=(1-\sqrt{1-4x})/(2x)$, use the recorded series $x^2(1+x^2C^4)C^3$, whose lowest degree is the source offset two.

\[y^2=1 - 4 x,\qquad y(0)=1.\]

\[C=\frac{1-y}{2x},\qquad F=x^2(1+x^2C^4)C^3.\]

Exact rational differentiation and cancellation in \eqref{round7:r7:gf-rational-check} give

\[LF=P(x),\qquad P(x)=80 x^{2}.\]

Coefficient extraction in \eqref{round7:r7:gf-operator} therefore proves the displayed recurrence for all $n\geq 4$.

The stated lower bound also ensures $n\geq s+r=4$, so no term below the source offset is used.


\subsection{A026270: the selected recurrence}

\label{round7:r7entry:A026270}

\paragraph{Source and selected assertion.}

\href{https://oeis.org/A026270}{A026270}, consulted on 20 October 2026, records the conjecture attributed to R. J. Mathar, Jun 23 2013.

The source offset is $s=2$.  In the convention of \eqref{round7:r7:gf-operator}, the selected recurrence is $R_n=0$ for every integer $n\geq 5$, with $r=3$ and

\[\begin{aligned}
p_{0}(n)&=n + 3\\
p_{1}(n)&=- 3 \left(n + 1\right)\\
p_{2}(n)&=- n - 1\\
p_{3}(n)&=3 \left(n - 5\right)\\
\end{aligned}\]

\paragraph{Known inputs and scope.}

The generating function printed on A026270 has inconsistent constant and initial terms. We derive the needed generating function from the defining boundary of A026268 and the credited diagonal generating function on A026269. The selected Mathar recurrence is true from n=5; correction of the auxiliary input is disclosed separately and is not an extra count.

\paragraph{Proof.}

Let $d_n=T(n,n)$ and $a_n=T(n,n-1)$ for the defining triangle
\href{https://oeis.org/A026268}{A026268}.  Its boundary rule gives
$d_n=d_{n-1}+a_{n-1}$ for $n\geq3$, with $d_2=a_2=1$.
The credited diagonal series in
\href{https://oeis.org/A026269}{A026269} is
\[
 D(x)=\sum_{n\geq2}d_nx^n=
 \frac{4x^2(1-x^2)}{(1-x+\sqrt{1-2x-3x^2})^2}.
\]
Summing the boundary rule gives $D-x^2=xD+xF$, hence
\[
 F(x)=\frac{(1-x)D(x)-x^2}{x}.
\]
The expression printed directly as a GF on A026270 is
$-1+(1-x)D(x)$; it has constant term $-1$ and fails to reproduce the
listed offset-two sequence.  It is not a premise of this proof.
The boundary-derived series below begins $x^2+2x^3+6x^4+15x^5$.
The selected recurrence itself needs no coefficient correction.

\[y^2=- 3 x^{2} - 2 x + 1,\qquad y(0)=1.\]

\[F(x)=\frac{4 x \left(1 - x\right) \left(1 - x^{2}\right)}{\left(- x + y + 1\right)^{2}} - x.\]

Exact rational differentiation and cancellation in \eqref{round7:r7:gf-rational-check} give

\[LF=P(x),\qquad P(x)=7 x^{4} + 5 x^{2}.\]

Coefficient extraction in \eqref{round7:r7:gf-operator} therefore proves the displayed recurrence for all $n\geq 5$.

The stated lower bound also ensures $n\geq s+r=5$, so no term below the source offset is used.


\subsection{A059279: the selected recurrence}

\label{round7:r7entry:A059279}

\paragraph{Source and selected assertion.}

\href{https://oeis.org/A059279}{A059279}, consulted on 20 October 2026, records the conjecture attributed to R. J. Mathar, Nov 15 2011.

The source offset is $s=0$.  In the convention of \eqref{round7:r7:gf-operator}, the selected recurrence is $R_n=0$ for every integer $n\geq 3$, with $r=3$ and

\[\begin{aligned}
p_{0}(n)&=n + 1\\
p_{1}(n)&=- 2 \left(4 n - 1\right)\\
p_{2}(n)&=4 \left(4 n - 5\right)\\
p_{3}(n)&=- 4 \left(2 n - 5\right)\\
\end{aligned}\]

\paragraph{Known inputs and scope.}

The Catalan substitution is the definition. Greubel (January 4, 2017) and Gutkovskiy (June 30, 2020) record equivalent algebraic formulas. The linked Gomez Aiza material is background on the same enumerative family; the selected assertion is only the still-marked recurrence.

\paragraph{Proof.}

The defining Catalan substitution is $((1-x)/(1-2x))C(x(1-x)/(1-2x))$.  Its radical form is the following expression.

\[y^2=- \frac{4 x \left(1 - x\right)}{1 - 2 x} + 1,\qquad y(0)=1.\]

\[F(x)=\frac{1 - y}{2 x}.\]

Exact rational differentiation and cancellation in \eqref{round7:r7:gf-rational-check} give

\[LF=P(x),\qquad P(x)=2 x^{2} - 2 x + 1.\]

Coefficient extraction in \eqref{round7:r7:gf-operator} therefore proves the displayed recurrence for all $n\geq 3$.

The stated lower bound also ensures $n\geq s+r=3$, so no term below the source offset is used.


\subsection{A096471: the selected recurrence}

\label{round7:r7entry:A096471}

\paragraph{Source and selected assertion.}

\href{https://oeis.org/A096471}{A096471}, consulted on 20 October 2026, records the conjecture attributed to R. J. Mathar, Jul 06 2023.

The source offset is $s=2$.  In the convention of \eqref{round7:r7:gf-operator}, the selected recurrence is $R_n=0$ for every integer $n\geq 7$, with $r=5$ and

\[\begin{aligned}
p_{0}(n)&=1\\
p_{1}(n)&=-3\\
p_{2}(n)&=- 2 n^{2} + 8 n - 3\\
p_{3}(n)&=\left(n - 2\right) \left(4 n - 15\right)\\
p_{4}(n)&=\left(n - 3\right) \left(n^{3} - 9 n^{2} + 22 n - 7\right)\\
p_{5}(n)&=- \left(n - 5\right)^{2} \left(n - 4\right) \left(n - 3\right)\\
\end{aligned}\]

\paragraph{Known inputs and scope.}

The cycle decomposition and its logarithmic EGF are known inputs recorded in the entry. The selected assertion is Mathar's July 2023 recurrence, with the ordinary-to-exponential coefficient conversion and source offset made explicit.

\paragraph{Proof.}

A permutation with exactly one even cycle consists of one
even cycle and an arbitrary set of odd cycles.  The cycle EGF therefore is
\[
 F(x)=\left(\sum_{m\geq1}\frac{x^{2m}}{2m}\right)
       \exp\left(\sum_{m\geq0}\frac{x^{2m+1}}{2m+1}\right)
     =-\frac12H(x)\ell(x),
\]
where $H=\sqrt{(1+x)/(1-x)}$ and $\ell=\log(1-x^2)$.
This also derives the known EGF recorded in the entry.
Its derivatives satisfy
\[
 \frac{H'}H=\frac1{1-x^2},\qquad \ell'=-\frac{2x}{1-x^2}.
\]
Use \eqref{round7:r7:gf-egf} with $r=5$.  On the basis $(H,H\ell)$,
the Euler derivative acts on the rational coefficients by
\[
 (R_0,R_1)\longmapsto
 \left(x\left(R_0'+\frac{R_0-2xR_1}{1-x^2}\right),
       x\left(R_1'+\frac{R_1}{1-x^2}\right)\right).
\]
Starting from $(R_0,R_1)=(0,-1/2)$ and applying
$\sum_{j=0}^5x^jq_j(\theta+j)$ gives exactly $(0,0)$.
This is a rational identity, checked by the same finite derivative recipe
as \eqref{round7:r7:gf-rational-check}.  Thus $R_n=0$ for every $n\geq7$,
the first index for which all six source terms have index at least two.


\subsection{A106391: the selected recurrence}

\label{round7:r7entry:A106391}

\paragraph{Source and selected assertion.}

\href{https://oeis.org/A106391}{A106391}, consulted on 20 October 2026, records the conjecture attributed to R. J. Mathar, Feb 20 2015.

The source offset is $s=0$.  In the convention of \eqref{round7:r7:gf-operator}, the selected recurrence is $R_n=0$ for every integer $n\geq 6$, with $r=6$ and

\[\begin{aligned}
p_{0}(n)&=- \left(n - 2\right) \left(n + 1\right)\\
p_{1}(n)&=4 n \left(3 n - 7\right)\\
p_{2}(n)&=- 49 n^{2} + 193 n - 148\\
p_{3}(n)&=8 \left(n - 4\right) \left(9 n - 19\right)\\
p_{4}(n)&=5 n^{2} + 115 n - 418\\
p_{5}(n)&=- 12 \left(n - 4\right) \left(7 n - 19\right)\\
p_{6}(n)&=45 \left(n - 5\right) \left(n - 4\right)\\
\end{aligned}\]

\paragraph{Known inputs and scope.}

The finite binomial sum and Bessel EGF are existing inputs. The order-six recurrence is the selected source-marked assertion. The entry's second, order-four conjecture is not counted separately. The ordinary GF used here is derived directly by two binomial transforms.

\paragraph{Proof.}

Put $b_n=\binom{2n}{n-1}$, with $b_0=0$.
The recorded defining sum is
$a_n=\sum_k\binom n{2k}b_{n-2k}$.
Since $b_n=\binom{2n}n-C_n$, its generating function is
\[
 B(x)=\frac{1-2x}{2x\sqrt{1-4x}}-\frac1{2x}.
\]
The defining even-binomial transform equals the mean of the transforms
with weights $1$ and $-1$.  Therefore
\[
 F(x)=\frac12\left(\frac{B(x/(1-x))}{1-x}
                    +\frac{B(x/(1+x))}{1+x}\right),
\]
which is the three-component radical expression displayed below.
This derives the needed ordinary GF from the existing finite sum.

\[
 F(x)=-\frac1{2x}
 +\frac{1-3x}{4x\sqrt{1-6x+5x^2}}
 +\frac{1-x}{4x\sqrt{1-2x-3x^2}}.
\]
For \eqref{round7:r7:gf-rational-check}, take
\[
 (D_0,D_1,D_2)=(1,1-6x+5x^2,1-2x-3x^2),
\]
\[
 (R_0,R_1,R_2)=\left(-\frac1{2x},\frac{1-3x}{4xD_1},
                           \frac{1-x}{4xD_2}\right).
\]
The apparent negative powers cancel.

Exact rational differentiation and cancellation in \eqref{round7:r7:gf-rational-check} give

\[LF=P(x),\qquad P(x)=14 x^{3} - 8 x^{2} + 2 x.\]

Coefficient extraction in \eqref{round7:r7:gf-operator} therefore proves the displayed recurrence for all $n\geq 6$.

The stated lower bound also ensures $n\geq s+r=6$, so no term below the source offset is used.


\subsection{A107026: the selected recurrence}

\label{round7:r7entry:A107026}

\paragraph{Source and selected assertion.}

\href{https://oeis.org/A107026}{A107026}, consulted on 20 October 2026, records the conjecture attributed to R. J. Mathar, Feb 20 2015.

The source offset is $s=0$.  In the convention of \eqref{round7:r7:gf-operator}, the selected recurrence is $R_n=0$ for every integer $n\geq 2$, with $r=2$ and

\[\begin{aligned}
p_{0}(n)&=3 n \left(3 n - 2\right) \left(3 n - 1\right) \left(22 n^{2} - 62 n + 43\right)\\
p_{1}(n)&=- 8 \left(1892 n^{5} - 8280 n^{4} + 13330 n^{3} - 9660 n^{2} + 3048 n - 315\right)\\
p_{2}(n)&=128 \left(2 n - 3\right) \left(4 n - 7\right) \left(4 n - 5\right) \left(22 n^{2} - 18 n + 3\right)\\
\end{aligned}\]

\paragraph{Known inputs and scope.}

The sequence is a Riordan inverse row sum. Manyama (November 8, 2025) records the quartic-tree generating function used as input. Only the second marked recurrence (order two) is selected; the first, order-four conjecture is outside the count.

\paragraph{Proof.}

The inverse of the defining Riordan array
$(1/(1+x),x/(1+x)^4)$ has row-sum series $(1+t)/(1-t)$,
where $t=x(1+t)^4$: inversion gives first component $1+t$ and
second component $t$, and summing its columns divides by $1-t$.
Equivalently use Manyama's November 8, 2025 formula
$F=g/(2-g)$ with $g=1+xg^4$, and put $g=1+t$.
This proves the rational parametrization used below.

\[x=\xi(t)=\frac{t}{\left(t + 1\right)^{4}},\qquad t(0)=0,\quad \xi'(0)=1.\]

\[F=- \frac{t + 1}{t - 1}.\]

Use $D_0=1$, $R_0=F(\xi(t))$, and $M=\xi/\xi'$ in \eqref{round7:r7:gf-rational-check}.  All expressions are rational in $t$.

Exact rational differentiation and cancellation in \eqref{round7:r7:gf-rational-check} give

\[LF=P(x),\qquad P(x)=- 84 x.\]

Coefficient extraction in \eqref{round7:r7:gf-operator} therefore proves the displayed recurrence for all $n\geq 2$.

The stated lower bound also ensures $n\geq s+r=2$, so no term below the source offset is used.


\subsection{A108630: correction of the polynomial index}

\label{round7:r7entry:A108630}

\paragraph{Source and selected assertion.}

\href{https://oeis.org/A108630}{A108630}, consulted on 20 October 2026, records the conjecture attributed to R. J. Mathar, Sep 27 2014.

The source offset is $s=1$.  In the convention of \eqref{round7:r7:gf-operator}, the selected recurrence is $R_n=0$ for every integer $n\geq 7$, with $r=6$ and

\[\begin{aligned}
p_{0}(n)&=n\\
p_{1}(n)&=3 - n\\
p_{2}(n)&=- 8 \left(n - 3\right)\\
p_{3}(n)&=60 - 17 n\\
p_{4}(n)&=81 - 19 n\\
p_{5}(n)&=- 12 \left(n - 5\right)\\
p_{6}(n)&=- 4 \left(n - 6\right)\\
\end{aligned}\]

\paragraph{Known inputs and scope.}

The generating function is the sequence definition at offset one. The printed recurrence fails at n=7 with residual 168. Replacing each polynomial coefficient p\_j(n) by p\_j(n+1), with sequence indices unchanged, yields the proved recurrence for n>=7. This is a correction, not a proof of the literal printed statement.

\paragraph{The literal source statement and its correction.}

The relevant source values are \[ (a_1,\ldots,a_{7})=(1,1,2,6,17,54,178). \]

Substituting them into the printed recurrence gives residual $168\ne0$ at $n=7$.

The corrected coefficients above are $p_j^{\rm source}(n+1)$, while the sequence terms remain $a_{n-j}$.  The coefficient shift is essential; the literal recurrence is refuted.

\paragraph{Proof.}

Use the credited generating function displayed below, with the formal branch $y(0)=1$.

\[y^2=- 4 x^{4} - 8 x^{3} - 7 x^{2} - 2 x + 1,\qquad y(0)=1.\]

\[F(x)=\frac{x - y + 1}{2 x^{2} + 2 x + 2}.\]

Exact rational differentiation and cancellation in \eqref{round7:r7:gf-rational-check} give

\[LF=P(x),\qquad P(x)=2 x^{5} + 6 x^{4} + 6 x^{3} + 3 x^{2} + x.\]

Coefficient extraction in \eqref{round7:r7:gf-operator} therefore proves the displayed recurrence for all $n\geq 7$.

The stated lower bound also ensures $n\geq s+r=7$, so no term below the source offset is used.


\subsection{A110198: the selected recurrence}

\label{round7:r7entry:A110198}

\paragraph{Source and selected assertion.}

\href{https://oeis.org/A110198}{A110198}, consulted on 20 October 2026, records the conjecture attributed to R. J. Mathar, Nov 15 2012.

The source offset is $s=0$.  In the convention of \eqref{round7:r7:gf-operator}, the selected recurrence is $R_n=0$ for every integer $n\geq 5$, with $r=5$ and

\[\begin{aligned}
p_{0}(n)&=n\\
p_{1}(n)&=1 - 3 n\\
p_{2}(n)&=n\\
p_{3}(n)&=2 - n\\
p_{4}(n)&=3 n - 5\\
p_{5}(n)&=2 - n\\
\end{aligned}\]

\paragraph{Known inputs and scope.}

The antidiagonal generating function and a double binomial sum are already recorded. They are input formulas, not discoveries. Only the marked order-five recurrence is selected.

\paragraph{Proof.}

Use the credited generating function displayed below, with the formal branch $y(0)=1$.

\[y^2=\left(x^{2} - 3 x + 1\right) \left(x^{2} + x + 1\right),\qquad y(0)=1.\]

\[F(x)=\frac{1}{y \left(1 - x\right)}.\]

Exact rational differentiation and cancellation in \eqref{round7:r7:gf-rational-check} give

\[LF=P(x),\qquad P(x)=0.\]

Coefficient extraction in \eqref{round7:r7:gf-operator} therefore proves the displayed recurrence for all $n\geq 5$.

The stated lower bound also ensures $n\geq s+r=5$, so no term below the source offset is used.


\subsection{A119975: the selected recurrence}

\label{round7:r7entry:A119975}

\paragraph{Source and selected assertion.}

\href{https://oeis.org/A119975}{A119975}, consulted on 20 October 2026, records the conjecture attributed to R. J. Mathar, Dec 10 2011.

The source offset is $s=0$.  In the convention of \eqref{round7:r7:gf-operator}, the selected recurrence is $R_n=0$ for every integer $n\geq 3$, with $r=3$ and

\[\begin{aligned}
p_{0}(n)&=n + 1\\
p_{1}(n)&=n - 4\\
p_{2}(n)&=9 - 13 n\\
p_{3}(n)&=- 21 \left(n - 2\right)\\
\end{aligned}\]

\paragraph{Known inputs and scope.}

The ordinary GF is explicitly credited to Kotesovec's June 26, 2013 correction. The Bessel EGF is also known. Only Mathar's December 2011 recurrence is selected.

\paragraph{Proof.}

Use the credited generating function displayed below, with the formal branch $y(0)=1$.

\[y^2=- 7 x^{2} - 2 x + 1,\qquad y(0)=1.\]

\[F(x)=\frac{-1 + \frac{3 x + 1}{y}}{4 x}.\]

Exact rational differentiation and cancellation in \eqref{round7:r7:gf-rational-check} give

\[LF=P(x),\qquad P(x)=x + 1.\]

Coefficient extraction in \eqref{round7:r7:gf-operator} therefore proves the displayed recurrence for all $n\geq 3$.

The stated lower bound also ensures $n\geq s+r=3$, so no term below the source offset is used.


\subsection{A134389: the selected recurrence}

\label{round7:r7entry:A134389}

\paragraph{Source and selected assertion.}

\href{https://oeis.org/A134389}{A134389}, consulted on 20 October 2026, records the conjecture attributed to R. J. Mathar, Oct 25 2012.

The source offset is $s=0$.  In the convention of \eqref{round7:r7:gf-operator}, the selected recurrence is $R_n=0$ for every integer $n\geq 3$, with $r=3$ and

\[\begin{aligned}
p_{0}(n)&=n - 1\\
p_{1}(n)&=- 2 n\\
p_{2}(n)&=- 4 \left(n - 4\right)\\
p_{3}(n)&=8 \left(n - 3\right)\\
\end{aligned}\]

\paragraph{Known inputs and scope.}

The quadratic GF is already recorded for the transform defining this sequence. Only the conjectured polynomial recurrence is selected.

\paragraph{Proof.}

Use the credited generating function displayed below, with the formal branch $y(0)=1$.

\[y^2=1 - 4 x^{2},\qquad y(0)=1.\]

\[F(x)=\frac{x y}{2 \left(1 - 2 x\right)^{2}} + \frac{2 - x}{2 - 4 x}.\]

Exact rational differentiation and cancellation in \eqref{round7:r7:gf-rational-check} give

\[LF=P(x),\qquad P(x)=5 x^{2} - 2 x - 1.\]

Coefficient extraction in \eqref{round7:r7:gf-operator} therefore proves the displayed recurrence for all $n\geq 3$.

The stated lower bound also ensures $n\geq s+r=3$, so no term below the source offset is used.


\subsection{A143020: the selected recurrence}

\label{round7:r7entry:A143020}

\paragraph{Source and selected assertion.}

\href{https://oeis.org/A143020}{A143020}, consulted on 20 October 2026, records the conjecture attributed to R. J. Mathar, May 10 2018.

The source offset is $s=2$.  In the convention of \eqref{round7:r7:gf-operator}, the selected recurrence is $R_n=0$ for every integer $n\geq 4$, with $r=2$ and

\[\begin{aligned}
p_{0}(n)&=- \left(n - 1\right) \left(n + 1\right) \left(39 n^{2} - 136 n + 60\right)\\
p_{1}(n)&=- 60 \left(18 n^{2} - 57 n + 7\right)\\
p_{2}(n)&=12 \left(3 n - 8\right) \left(3 n - 7\right) \left(39 n^{2} - 58 n - 37\right)\\
\end{aligned}\]

\paragraph{Known inputs and scope.}

Deutsch's enumerative formula, the cubic GF for connected noncrossing graphs, and Howroyd's 2017 reversion program are known inputs. Only Mathar's May 2018 order-two recurrence is selected; the Flajolet--Noy setting is credited background.

\paragraph{Proof.}

The credited GF is $F=g(g-x)/x$, where
$g^3+g^2-3xg+2x^2=0$ is the connected-noncrossing-graph series.
Its branch starts $g=x+O(x^2)$.  Substitute $g=x(1+t)$.
After cancellation the cubic becomes
$x=t(1-t)/(1+t)^3$, and $F=xt(1+t)$.
Since $\xi'(0)=1$, this specifies the branch uniquely.

\[x=\xi(t)=- \frac{t \left(t - 1\right)}{\left(t + 1\right)^{3}},\qquad t(0)=0,\quad \xi'(0)=1.\]

\[F=xt(1+t).\]

Use $D_0=1$, $R_0=F(\xi(t))$, and $M=\xi/\xi'$ in \eqref{round7:r7:gf-rational-check}.  All expressions are rational in $t$.

Exact rational differentiation and cancellation in \eqref{round7:r7:gf-rational-check} give

\[LF=P(x),\qquad P(x)=168 x^{2}.\]

Coefficient extraction in \eqref{round7:r7:gf-operator} therefore proves the displayed recurrence for all $n\geq 4$.

The stated lower bound also ensures $n\geq s+r=4$, so no term below the source offset is used.


\subsection{A155051: the selected recurrence}

\label{round7:r7entry:A155051}

\paragraph{Source and selected assertion.}

\href{https://oeis.org/A155051}{A155051}, consulted on 20 October 2026, records the conjecture attributed to R. J. Mathar, Feb 05 2015.

The source offset is $s=0$.  In the convention of \eqref{round7:r7:gf-operator}, the selected recurrence is $R_n=0$ for every integer $n\geq 4$, with $r=4$ and

\[\begin{aligned}
p_{0}(n)&=n + 2\\
p_{1}(n)&=-2\\
p_{2}(n)&=4 - 5 n\\
p_{3}(n)&=8\\
p_{4}(n)&=4 \left(n - 3\right)\\
\end{aligned}\]

\paragraph{Known inputs and scope.}

The Catalan-substitution GF is the definition. Select the first Mathar recurrence, of order four. The second Mathar recurrence and Zumkeller's unrelated composition conjecture are outside the selected assertion.

\paragraph{Proof.}

The defining series is $C(x^2)(1+x)/(1-x)$.  Substituting the Catalan radical gives the following expression.

\[y^2=1 - 4 x^{2},\qquad y(0)=1.\]

\[F(x)=\frac{\left(1 - y\right) \left(x + 1\right)}{2 x^{2} \left(1 - x\right)}.\]

Exact rational differentiation and cancellation in \eqref{round7:r7:gf-rational-check} give

\[LF=P(x),\qquad P(x)=2 x^{2} + 4 x + 2.\]

Coefficient extraction in \eqref{round7:r7:gf-operator} therefore proves the displayed recurrence for all $n\geq 4$.

The stated lower bound also ensures $n\geq s+r=4$, so no term below the source offset is used.


\subsection{A157418: the selected recurrence}

\label{round7:r7entry:A157418}

\paragraph{Source and selected assertion.}

\href{https://oeis.org/A157418}{A157418}, consulted on 20 October 2026, records the conjecture attributed to R. J. Mathar, Jul 24 2012.

The source offset is $s=1$.  In the convention of \eqref{round7:r7:gf-operator}, the selected recurrence is $R_n=0$ for every integer $n\geq 4$, with $r=3$ and

\[\begin{aligned}
p_{0}(n)&=n \left(7 n - 29\right)\\
p_{1}(n)&=- 21 n^{2} + 136 n - 192\\
p_{2}(n)&=- 2 \left(28 n^{2} - 165 n + 210\right)\\
p_{3}(n)&=- 4 \left(n - 3\right) \left(7 n - 22\right)\\
\end{aligned}\]

\paragraph{Known inputs and scope.}

The GF belongs to the Guo--Sit Rota--Baxter enumeration (2006/2010). Spiegel (June 5, 2026) records another Catalan expression. These are credited inputs; only Mathar's marked recurrence is selected.

\paragraph{Proof.}

Use the credited generating function displayed below, with the formal branch $y(0)=1$.

\[y^2=- 4 x^{2} - 4 x + 1,\qquad y(0)=1.\]

\[F(x)=\frac{- 2 x^{2} - x - y \left(x + 2\right) + 2}{2 \left(x + 1\right)^{2}}.\]

Exact rational differentiation and cancellation in \eqref{round7:r7:gf-rational-check} give

\[LF=P(x),\qquad P(x)=- 64 x^{2} - 22 x.\]

Coefficient extraction in \eqref{round7:r7:gf-operator} therefore proves the displayed recurrence for all $n\geq 4$.

The stated lower bound also ensures $n\geq s+r=4$, so no term below the source offset is used.


\subsection{A163824: the recurrence with its initial boundary}

\label{round7:r7entry:A163824}

\paragraph{Source and selected assertion.}

\href{https://oeis.org/A163824}{A163824}, consulted on 20 October 2026, records the conjecture attributed to R. J. Mathar, Jul 24 2012.

The source offset is $s=0$.  In the convention of \eqref{round7:r7:gf-operator}, the selected recurrence is $R_n=0$ for every integer $n\geq 4$, with $r=3$ and

\[\begin{aligned}
p_{0}(n)&=n \left(n^{2} - 6 n + 11\right)\\
p_{1}(n)&=- 9 n^{3} + 56 n^{2} - 119 n + 60\\
p_{2}(n)&=2 \left(13 n^{3} - 83 n^{2} + 193 n - 150\right)\\
p_{3}(n)&=- 12 \left(2 n - 5\right) \left(n^{2} - 4 n + 6\right)\\
\end{aligned}\]

\paragraph{Known inputs and scope.}

Bevan's recorded central-binomial GF with a rational correction is used as an input. The recurrence needs n>=4; its first fully indexed value n=3 has residual -6. No all-n version is asserted.

\paragraph{Proof.}

Use the credited generating function displayed below, with the formal branch $y(0)=1$.

\[y^2=1 - 4 x,\qquad y(0)=1.\]

\[F(x)=- \frac{x \left(1 - x\right)}{\left(1 - 3 x\right) \left(1 - 2 x\right)} + \frac{1}{y}.\]

Exact rational differentiation and cancellation in \eqref{round7:r7:gf-rational-check} give

\[LF=P(x),\qquad P(x)=- 6 x^{3} + 2 x^{2} - 6 x.\]

Coefficient extraction in \eqref{round7:r7:gf-operator} therefore proves the displayed recurrence for all $n\geq 4$.

The first fully indexed value is $n=3$, where $R_{3}=-6$.  This is exactly the coefficient of $x^{3}$ in $P$ and explains the stated starting index.


\subsection{A165201: the recurrence with its initial boundary}

\label{round7:r7entry:A165201}

\paragraph{Source and selected assertion.}

\href{https://oeis.org/A165201}{A165201}, consulted on 20 October 2026, records the conjecture attributed to R. J. Mathar, Nov 15 2011.

The source offset is $s=0$.  In the convention of \eqref{round7:r7:gf-operator}, the selected recurrence is $R_n=0$ for every integer $n\geq 4$, with $r=3$ and

\[\begin{aligned}
p_{0}(n)&=n \left(n - 3\right)\\
p_{1}(n)&=- 2 \left(4 n^{2} - 15 n + 10\right)\\
p_{2}(n)&=15 n^{2} - 69 n + 80\\
p_{3}(n)&=2 \left(n - 2\right) \left(2 n - 5\right)\\
\end{aligned}\]

\paragraph{Known inputs and scope.}

The Catalan composition is the definition. The source recurrence needs n>=4; n=3 has residual 2. This is a boundary qualification of the selected recurrence.

\paragraph{Proof.}

Let $C=(1-\sqrt{1-4x})/(2x)$ be the Catalan series.  The defining expression $1/(1-xC^3)$ rationalizes to the following $F$.

\[y^2=1 - 4 x,\qquad y(0)=1.\]

\[F(x)=\frac{- 2 x^{2} - 3 x + y \left(1 - x\right) + 1}{- 2 x^{2} - 8 x + 2}.\]

Exact rational differentiation and cancellation in \eqref{round7:r7:gf-rational-check} give

\[LF=P(x),\qquad P(x)=2 x^{3} + 2 x^{2}.\]

Coefficient extraction in \eqref{round7:r7:gf-operator} therefore proves the displayed recurrence for all $n\geq 4$.

The first fully indexed value is $n=3$, where $R_{3}=2$.  This is exactly the coefficient of $x^{3}$ in $P$ and explains the stated starting index.


\subsection{A168494: correction of a duplicated shift}

\label{round7:r7entry:A168494}

\paragraph{Source and selected assertion.}

\href{https://oeis.org/A168494}{A168494}, consulted on 20 October 2026, records the conjecture attributed to R. J. Mathar, Sep 30 2012.

The source offset is $s=0$.  In the convention of \eqref{round7:r7:gf-operator}, the selected recurrence is $R_n=0$ for every integer $n\geq 4$, with $r=4$ and

\[\begin{aligned}
p_{0}(n)&=n + 1\\
p_{1}(n)&=4 - 11 n\\
p_{2}(n)&=5 \left(7 n - 11\right)\\
p_{3}(n)&=92 - 37 n\\
p_{4}(n)&=6 \left(2 n - 7\right)\\
\end{aligned}\]

\paragraph{Known inputs and scope.}

The period-three continued fraction and its algebraic GF are recorded inputs. The printed recurrence uses a(n-2) twice and fails at n=4 with residual 200. Replacing its first a(n-2), the term with coefficient 4-11n, by a(n-1) proves the corrected relation.

\paragraph{The literal source statement and its correction.}

The relevant source values are \[ (a_0,\ldots,a_{4})=(1,1,2,7,32). \]

Substituting them into the printed recurrence gives residual $200\ne0$ at $n=4$.

The printed term $(4-11n)a_{n-2}$ must instead be $(4-11n)a_{n-1}$.  The separate term $5(7n-11)a_{n-2}$ is retained.  The literal recurrence is refuted.

\paragraph{Proof.}

Use the credited generating function displayed below, with the formal branch $y(0)=1$.

\[y^2=- 12 x^{3} + 25 x^{2} - 10 x + 1,\qquad y(0)=1.\]

\[F(x)=\frac{x - y + 1}{6 x \left(1 - x\right)}.\]

Exact rational differentiation and cancellation in \eqref{round7:r7:gf-rational-check} give

\[LF=P(x),\qquad P(x)=x^{3} + 3 x^{2} - 5 x + 1.\]

Coefficient extraction in \eqref{round7:r7:gf-operator} therefore proves the displayed recurrence for all $n\geq 4$.

The stated lower bound also ensures $n\geq s+r=4$, so no term below the source offset is used.


\subsection{A168503: the selected recurrence}

\label{round7:r7entry:A168503}

\paragraph{Source and selected assertion.}

\href{https://oeis.org/A168503}{A168503}, consulted on 20 October 2026, records the conjecture attributed to R. J. Mathar, Feb 10 2015.

The source offset is $s=0$.  In the convention of \eqref{round7:r7:gf-operator}, the selected recurrence is $R_n=0$ for every integer $n\geq 4$, with $r=4$ and

\[\begin{aligned}
p_{0}(n)&=n + 1\\
p_{1}(n)&=5 - 13 n\\
p_{2}(n)&=6 \left(8 n - 13\right)\\
p_{3}(n)&=- 12 \left(5 n - 13\right)\\
p_{4}(n)&=12 \left(2 n - 7\right)\\
\end{aligned}\]

\paragraph{Known inputs and scope.}

Barry's period-three continued fraction and its quadratic GF are known inputs. The selected claim is the marked order-four recurrence. This is not an alias of A168494: its defining GF and recurrence coefficients differ.

\paragraph{Proof.}

Use the credited generating function displayed below, with the formal branch $y(0)=1$.

\[y^2=- 24 x^{3} + 36 x^{2} - 12 x + 1,\qquad y(0)=1.\]

\[F(x)=\frac{1 - y}{6 x \left(1 - x\right)}.\]

Exact rational differentiation and cancellation in \eqref{round7:r7:gf-rational-check} give

\[LF=P(x),\qquad P(x)=6 x^{2} - 6 x + 1.\]

Coefficient extraction in \eqref{round7:r7:gf-operator} therefore proves the displayed recurrence for all $n\geq 4$.

The stated lower bound also ensures $n\geq s+r=4$, so no term below the source offset is used.


\subsection{A174016: the selected recurrence}

\label{round7:r7entry:A174016}

\paragraph{Source and selected assertion.}

\href{https://oeis.org/A174016}{A174016}, consulted on 20 October 2026, records the conjecture attributed to R. J. Mathar, Nov 13 2012.

The source offset is $s=0$.  In the convention of \eqref{round7:r7:gf-operator}, the selected recurrence is $R_n=0$ for every integer $n\geq 4$, with $r=4$ and

\[\begin{aligned}
p_{0}(n)&=n + 1\\
p_{1}(n)&=- 6 n\\
p_{2}(n)&=12 \left(n - 1\right)\\
p_{3}(n)&=-12\\
p_{4}(n)&=- 8 \left(2 n - 7\right)\\
\end{aligned}\]

\paragraph{Known inputs and scope.}

The recorded algebraic GF for the row sums of A174014 is an input. Only Mathar's marked order-four recurrence is selected.

\paragraph{Proof.}

Use the credited generating function displayed below, with the formal branch $y(0)=1$.

\[y^2=8 x^{3} + 4 x^{2} - 4 x + 1,\qquad y(0)=1.\]

\[F(x)=\frac{4 x + y - 1}{2 x \left(1 - 2 x\right)}.\]

Exact rational differentiation and cancellation in \eqref{round7:r7:gf-rational-check} give

\[LF=P(x),\qquad P(x)=- 8 x^{3} + 6 x^{2} - 2 x + 1.\]

Coefficient extraction in \eqref{round7:r7:gf-operator} therefore proves the displayed recurrence for all $n\geq 4$.

The stated lower bound also ensures $n\geq s+r=4$, so no term below the source offset is used.


\subsection{A184018: the selected recurrence}

\label{round7:r7entry:A184018}

\paragraph{Source and selected assertion.}

\href{https://oeis.org/A184018}{A184018}, consulted on 20 October 2026, records the conjecture attributed to R. J. Mathar, Nov 15 2011.

The source offset is $s=0$.  In the convention of \eqref{round7:r7:gf-operator}, the selected recurrence is $R_n=0$ for every integer $n\geq 4$, with $r=4$ and

\[\begin{aligned}
p_{0}(n)&=n + 1\\
p_{1}(n)&=- 2 \left(3 n - 1\right)\\
p_{2}(n)&=3 \left(n - 1\right)\\
p_{3}(n)&=2 \left(3 n - 5\right)\\
p_{4}(n)&=n - 3\\
\end{aligned}\]

\paragraph{Known inputs and scope.}

The Catalan substitution defines the sequence; Hanna's September 6, 2011 algebraic identity is credited. Only the marked recurrence is selected, not its known generating function or growth constant.

\paragraph{Proof.}

The defining Catalan substitution is $C(x/(1-x-x^2))/(1-x-x^2)$.  Substituting the elementary expression for $C$ gives the following $F$.

\[y^2=\left(- x^{2} - 5 x + 1\right) \left(- x^{2} - x + 1\right),\qquad y(0)=1.\]

\[F(x)=\frac{- x^{2} - x - y + 1}{2 x \left(- x^{2} - x + 1\right)}.\]

Exact rational differentiation and cancellation in \eqref{round7:r7:gf-rational-check} give

\[LF=P(x),\qquad P(x)=x^{2} + 1.\]

Coefficient extraction in \eqref{round7:r7:gf-operator} therefore proves the displayed recurrence for all $n\geq 4$.

The stated lower bound also ensures $n\geq s+r=4$, so no term below the source offset is used.


\subsection{A185020: the selected recurrence}

\label{round7:r7entry:A185020}

\paragraph{Source and selected assertion.}

\href{https://oeis.org/A185020}{A185020}, consulted on 20 October 2026, records the conjecture attributed to R. J. Mathar, Oct 08 2016.

The source offset is $s=0$.  In the convention of \eqref{round7:r7:gf-operator}, the selected recurrence is $R_n=0$ for every integer $n\geq 2$, with $r=2$ and

\[\begin{aligned}
p_{0}(n)&=n \left(n + 1\right)\\
p_{1}(n)&=- 4 n \left(2 n - 1\right)\\
p_{2}(n)&=- 8 \left(2 n - 3\right) \left(2 n - 1\right)\\
\end{aligned}\]

\paragraph{Known inputs and scope.}

The product C\_n*A002605(n+1) is the definition. The second-order constant-coefficient recurrence of A002605 is known and credited. Dividing out Catalan factors proves the selected polynomial recurrence. The generic gauge method is elementary and not claimed as new.

\paragraph{Proof.}

The defining product is $a_n=C_nS_n$, where
$S_n=A002605(n+1)$, $S_0=1$, $S_1=2$, and
$S_n=2S_{n-1}+2S_{n-2}$.  This constant-coefficient recurrence is a
known input, recorded in the source programs as well as the defining
sequence.  The Catalan backward ratio is
\[
 \frac{C_{n-1}}{C_n}=\frac{n+1}{2(2n-1)}.
\]
Dividing the displayed $R_n$ by $C_n$ and using this ratio twice gives
\[
 \frac{R_n}{C_n}=n(n+1)(S_n-2S_{n-1}-2S_{n-2})=0.
\]
All divisions are legitimate for $n\geq2$.  This proves the selected
recurrence for every such $n$.


\subsection{A190725: the selected recurrence}

\label{round7:r7entry:A190725}

\paragraph{Source and selected assertion.}

\href{https://oeis.org/A190725}{A190725}, consulted on 20 October 2026, records the conjecture attributed to R. J. Mathar, Jul 24 2012.

The source offset is $s=0$.  In the convention of \eqref{round7:r7:gf-operator}, the selected recurrence is $R_n=0$ for every integer $n\geq 4$, with $r=4$ and

\[\begin{aligned}
p_{0}(n)&=n \left(2 n - 3\right)\\
p_{1}(n)&=- 6 n^{2} + 15 n - 8\\
p_{2}(n)&=- 2 \left(16 n^{2} - 32 n + 11\right)\\
p_{3}(n)&=- 6 n^{2} + 9 n - 2\\
p_{4}(n)&=\left(n - 2\right) \left(2 n - 1\right)\\
\end{aligned}\]

\paragraph{Known inputs and scope.}

Munarini's entry already records a nonconjectural recurrence of order six. The selected Mathar assertion is a shorter order-four recurrence, not a reindexing or scalar multiple of that known formula. Its lower order and boundary are proved from the credited GF; the existence of a recurrence is not claimed as new.

\paragraph{Proof.}

Use the credited generating function displayed below, with the formal branch $y(0)=1$.

\[y^2=x^{2} - 6 x + 1,\qquad y(0)=1.\]

\[F(x)=\frac{3 x - y + 3}{y \left(2 x^{2} + 6 x + 2\right)}.\]

Exact rational differentiation and cancellation in \eqref{round7:r7:gf-rational-check} give

\[LF=P(x),\qquad P(x)=- 2 x.\]

Coefficient extraction in \eqref{round7:r7:gf-operator} therefore proves the displayed recurrence for all $n\geq 4$.

The stated lower bound also ensures $n\geq s+r=4$, so no term below the source offset is used.


\subsection{A190736: the selected recurrence}

\label{round7:r7entry:A190736}

\paragraph{Source and selected assertion.}

\href{https://oeis.org/A190736}{A190736}, consulted on 20 October 2026, records the conjecture attributed to R. J. Mathar, Jul 24 2012.

The source offset is $s=0$.  In the convention of \eqref{round7:r7:gf-operator}, the selected recurrence is $R_n=0$ for every integer $n\geq 5$, with $r=5$ and

\[\begin{aligned}
p_{0}(n)&=n \left(11 n - 35\right)\\
p_{1}(n)&=- 3 \left(33 n^{2} - 149 n + 136\right)\\
p_{2}(n)&=2 \left(77 n^{2} - 377 n + 396\right)\\
p_{3}(n)&=- 209 n^{2} + 1061 n - 1200\\
p_{4}(n)&=- 12 \left(11 n - 30\right)\\
p_{5}(n)&=4 \left(n - 4\right) \left(11 n - 24\right)\\
\end{aligned}\]

\paragraph{Known inputs and scope.}

Munarini's entry already records a nonconjectural order-nine recurrence. The selected Mathar assertion is a shorter order-five recurrence, not a reindexing or scalar multiple of that known formula. It is proved from the credited GF; the already-known holonomicity is not counted as a result.

\paragraph{Proof.}

Use the credited generating function displayed below, with the formal branch $y(0)=1$.

\[y^2=4 x^{2} - 8 x + 1,\qquad y(0)=1.\]

\[F(x)=\frac{- 2 x^{2} - 5 x - y \left(x + 2\right) + 4}{2 x^{3} + 4 x^{2} - 2 x + 2}.\]

Exact rational differentiation and cancellation in \eqref{round7:r7:gf-rational-check} give

\[LF=P(x),\qquad P(x)=54 x^{3} - 102 x^{2} - 108 x.\]

Coefficient extraction in \eqref{round7:r7:gf-operator} therefore proves the displayed recurrence for all $n\geq 5$.

The stated lower bound also ensures $n\geq s+r=5$, so no term below the source offset is used.


\subsection{A190788: the recurrence with its initial boundary}

\label{round7:r7entry:A190788}

\paragraph{Source and selected assertion.}

\href{https://oeis.org/A190788}{A190788}, consulted on 20 October 2026, records the conjecture attributed to R. J. Mathar, Nov 28 2013.

The source offset is $s=0$.  In the convention of \eqref{round7:r7:gf-operator}, the selected recurrence is $R_n=0$ for every integer $n\geq 8$, with $r=7$ and

\[\begin{aligned}
p_{0}(n)&=1 - n\\
p_{1}(n)&=3 n - 5\\
p_{2}(n)&=2 \left(3 n - 7\right)\\
p_{3}(n)&=- 4 \left(6 n - 19\right)\\
p_{4}(n)&=4 \left(n - 3\right)\\
p_{5}(n)&=4 \left(11 n - 53\right)\\
p_{6}(n)&=- 16 \left(3 n - 16\right)\\
p_{7}(n)&=16 \left(n - 6\right)\\
\end{aligned}\]

\paragraph{Known inputs and scope.}

The algebraic GF is the definition. The selected recurrence is valid from n=8; its first fully indexed value n=7 has residual 16. No claim at that exceptional index is retained.

\paragraph{Proof.}

Use the credited generating function displayed below, with the formal branch $y(0)=1$.

\[y^2=1 - 4 x^{2},\qquad y(0)=1.\]

\[F(x)=\frac{y \left(x - 1\right)}{x + y \left(x - 1\right)}.\]

Exact rational differentiation and cancellation in \eqref{round7:r7:gf-rational-check} give

\[LF=P(x),\qquad P(x)=16 x^{7} - 32 x^{6} + 12 x^{5} + 4 x^{4} + 4 x^{3} - 3 x^{2} - 2 x + 1.\]

Coefficient extraction in \eqref{round7:r7:gf-operator} therefore proves the displayed recurrence for all $n\geq 8$.

The first fully indexed value is $n=7$, where $R_{7}=16$.  This is exactly the coefficient of $x^{7}$ in $P$ and explains the stated starting index.


\subsection{A191319: the selected recurrence}

\label{round7:r7entry:A191319}

\paragraph{Source and selected assertion.}

\href{https://oeis.org/A191319}{A191319}, consulted on 20 October 2026, records the conjecture attributed to R. J. Mathar, Dec 07 2017.

The source offset is $s=0$.  In the convention of \eqref{round7:r7:gf-operator}, the selected recurrence is $R_n=0$ for every integer $n\geq 4$, with $r=4$ and

\[\begin{aligned}
p_{0}(n)&=2 - n\\
p_{1}(n)&=2\\
p_{2}(n)&=5 \left(n - 2\right)\\
p_{3}(n)&=-2\\
p_{4}(n)&=- 4 \left(n - 2\right)\\
\end{aligned}\]

\paragraph{Known inputs and scope.}

The recorded GF for pyramid weights in dispersed Dyck paths is an input. Only the marked order-four recurrence is selected.

\paragraph{Proof.}

Use the credited generating function displayed below, with the formal branch $y(0)=1$.

\[y^2=1 - 4 x^{2},\qquad y(0)=1.\]

\[F(x)=\frac{x^{2}}{y \left(1 - 2 x\right) \left(1 - x^{2}\right)}.\]

Exact rational differentiation and cancellation in \eqref{round7:r7:gf-rational-check} give

\[LF=P(x),\qquad P(x)=0.\]

Coefficient extraction in \eqref{round7:r7:gf-operator} therefore proves the displayed recurrence for all $n\geq 4$.

The stated lower bound also ensures $n\geq s+r=4$, so no term below the source offset is used.


\subsection{A191398: the selected recurrence}

\label{round7:r7entry:A191398}

\paragraph{Source and selected assertion.}

\href{https://oeis.org/A191398}{A191398}, consulted on 20 October 2026, records the conjecture attributed to R. J. Mathar, Jun 14 2016.

The source offset is $s=0$.  In the convention of \eqref{round7:r7:gf-operator}, the selected recurrence is $R_n=0$ for every integer $n\geq 6$, with $r=6$ and

\[\begin{aligned}
p_{0}(n)&=- \left(n - 3\right) \left(n + 2\right)\\
p_{1}(n)&=3 n^{2} - 3 n - 14\\
p_{2}(n)&=2 \left(n^{2} - 7 n + 8\right)\\
p_{3}(n)&=- 4 \left(3 n^{2} - 12 n + 10\right)\\
p_{4}(n)&=7 n^{2} - 31 n + 38\\
p_{5}(n)&=4\\
p_{6}(n)&=4 \left(n - 2\right)^{2}\\
\end{aligned}\]

\paragraph{Known inputs and scope.}

The recorded dispersed-Dyck-path GF is an input. Only the marked order-six recurrence is selected.

\paragraph{Proof.}

Use the credited generating function displayed below, with the formal branch $y(0)=1$.

\[y^2=1 - 4 x^{2},\qquad y(0)=1.\]

\[F(x)=\frac{2}{- 2 x^{3} - x + y \left(1 - x\right) + 1}.\]

Exact rational differentiation and cancellation in \eqref{round7:r7:gf-rational-check} give

\[LF=P(x),\qquad P(x)=- 4 x^{3} - 4 x^{2} - 8 x + 6.\]

Coefficient extraction in \eqref{round7:r7:gf-operator} therefore proves the displayed recurrence for all $n\geq 6$.

The stated lower bound also ensures $n\geq s+r=6$, so no term below the source offset is used.


\subsection{A220899: the selected recurrence}

\label{round7:r7entry:A220899}

\paragraph{Source and selected assertion.}

\href{https://oeis.org/A220899}{A220899}, consulted on 20 October 2026, records the conjecture attributed to R. J. Mathar, Jul 17 2014.

The source offset is $s=2$.  In the convention of \eqref{round7:r7:gf-operator}, the selected recurrence is $R_n=0$ for every integer $n\geq 7$, with $r=5$ and

\[\begin{aligned}
p_{0}(n)&=2 \left(n + 1\right) \left(2 n + 1\right)\\
p_{1}(n)&=- 23 n^{2} + 31 n - 6\\
p_{2}(n)&=- 2 \left(4 n^{2} + 9 n - 25\right)\\
p_{3}(n)&=2 \left(23 n^{2} - 142 n + 174\right)\\
p_{4}(n)&=4 \left(n - 5\right) \left(n + 8\right)\\
p_{5}(n)&=- 23 \left(n - 6\right) \left(n - 5\right)\\
\end{aligned}\]

\paragraph{Known inputs and scope.}

The map enumeration of Kitaev, Salimov, Severs and Ulfarsson and the composition using A000139 and A220898 are credited inputs. The offset-two GF removes the extension 2-x before the exact certificate is applied. Only the marked order-five recurrence is selected.

\paragraph{Proof.}

The credited map composition is
$(1-x)G(x/(1+x))$, where
\href{https://oeis.org/A000139}{A000139} has coefficients
$2(3n)!/((2n+1)!(n+1)!)$.
Its parameterization is $z=t(1-t)^2$ and
$G(z)=(2-3t)/(1-t)^2$.  It can also be checked directly by Lagrange
inversion: for $n\geq1$ the coefficient is
\[
 \frac1n\left(\binom{3n+1}{n-1}
                    -3\binom{3n}{n-2}\right)
 =\frac{2(3n)!}{(2n+1)!(n+1)!},
\]
and its constant term is $2$.
Thus, with $u=t(1-t)^2$, set $x=u/(1-u)$.
The composition has initial extension $2-x$; subtracting it gives
the actual offset-two series $F$.  This boundary adjustment is included
before applying the recurrence operator.

\[x=\xi(t)=- \frac{t \left(t - 1\right)^{2}}{t^{3} - 2 t^{2} + t - 1},\qquad t(0)=0,\quad \xi'(0)=1.\]

\[F=(1-x)\frac{2-3t}{(1-t)^2}-2+x.\]

Use $D_0=1$, $R_0=F(\xi(t))$, and $M=\xi/\xi'$ in \eqref{round7:r7:gf-rational-check}.  All expressions are rational in $t$.

Exact rational differentiation and cancellation in \eqref{round7:r7:gf-rational-check} give

\[LF=P(x),\qquad P(x)=40 x^{4} + 112 x^{3}.\]

Coefficient extraction in \eqref{round7:r7:gf-operator} therefore proves the displayed recurrence for all $n\geq 7$.

The stated lower bound also ensures $n\geq s+r=7$, so no term below the source offset is used.


\subsection{A228180: the recurrence with its initial boundary}

\label{round7:r7entry:A228180}

\paragraph{Source and selected assertion.}

\href{https://oeis.org/A228180}{A228180}, consulted on 20 October 2026, records the conjecture attributed to R. J. Mathar, Aug 25 2013.

The source offset is $s=0$.  In the convention of \eqref{round7:r7:gf-operator}, the selected recurrence is $R_n=0$ for every integer $n\geq 4$, with $r=3$ and

\[\begin{aligned}
p_{0}(n)&=2 \left(n + 1\right)\\
p_{1}(n)&=5 - 13 n\\
p_{2}(n)&=23 n - 37\\
p_{3}(n)&=- 6 \left(2 n - 5\right)\\
\end{aligned}\]

\paragraph{Known inputs and scope.}

The recorded Catalan tree-boundary GF is an input. The selected recurrence requires n>=4; n=3 has residual 12. This boundary condition is part of the result.

\paragraph{Proof.}

With $C=(1-\sqrt{1-4x})/(2x)$, use the recorded series $(xC+2x^3C^4)/(1-x)$.

\[y^2=1 - 4 x,\qquad y(0)=1.\]

\[C=\frac{1-y}{2x},\qquad F=\frac{xC+2x^3C^4}{1-x}.\]

Exact rational differentiation and cancellation in \eqref{round7:r7:gf-rational-check} give

\[LF=P(x),\qquad P(x)=12 x^{3} - 9 x^{2} + 4 x.\]

Coefficient extraction in \eqref{round7:r7:gf-operator} therefore proves the displayed recurrence for all $n\geq 4$.

The first fully indexed value is $n=3$, where $R_{3}=12$.  This is exactly the coefficient of $x^{3}$ in $P$ and explains the stated starting index.


\subsection{A278472: the selected recurrence}

\label{round7:r7entry:A278472}

\paragraph{Source and selected assertion.}

\href{https://oeis.org/A278472}{A278472}, consulted on 20 October 2026, records the conjecture attributed to R. J. Mathar, Mar 06 2017.

The source offset is $s=0$.  In the convention of \eqref{round7:r7:gf-operator}, the selected recurrence is $R_n=0$ for every integer $n\geq 3$, with $r=3$ and

\[\begin{aligned}
p_{0}(n)&=\left(n + 2\right) \left(n^{2} - 4 n + 1\right)\\
p_{1}(n)&=- 2 \left(4 n^{3} - 9 n^{2} - 16 n + 6\right)\\
p_{2}(n)&=15 n^{3} - 38 n^{2} - 9 n + 14\\
p_{3}(n)&=2 \left(2 n - 1\right) \left(n^{2} - 2 n - 2\right)\\
\end{aligned}\]

\paragraph{Known inputs and scope.}

The Fibonacci-binomial sum and its GF are recorded inputs. The first marked recurrence, of order three, is selected. The second conjectured recurrence is not another count.

\paragraph{Proof.}

Use the credited generating function displayed below, with the formal branch $y(0)=1$.

\[y^2=1 - 4 x,\qquad y(0)=1.\]

\[F(x)=\frac{- 2 x - y + 1}{x^{2} \left(2 x y - 8 x + 2\right)}.\]

Exact rational differentiation and cancellation in \eqref{round7:r7:gf-rational-check} give

\[LF=P(x),\qquad P(x)=2.\]

Coefficient extraction in \eqref{round7:r7:gf-operator} therefore proves the displayed recurrence for all $n\geq 3$.

The stated lower bound also ensures $n\geq s+r=3$, so no term below the source offset is used.


\subsection{A372991: the selected recurrence}

\label{round7:r7entry:A372991}

\paragraph{Source and selected assertion.}

\href{https://oeis.org/A372991}{A372991}, consulted on 20 October 2026, records the conjecture attributed to R. J. Mathar, Jul 16 2024.

The source offset is $s=0$.  In the convention of \eqref{round7:r7:gf-operator}, the selected recurrence is $R_n=0$ for every integer $n\geq 3$, with $r=3$ and

\[\begin{aligned}
p_{0}(n)&=1\\
p_{1}(n)&=0\\
p_{2}(n)&=0\\
p_{3}(n)&=- 2 n \left(2 n - 1\right)\\
\end{aligned}\]

\paragraph{Known inputs and scope.}

The nonlinear factorial quotient is the definition. Its positivity permits cancellation at consecutive indices. The selected assertion is exactly Mathar's marked step-three linear recurrence.

\paragraph{Proof.}

The definition gives
$a_na_{n-1}a_{n-2}=(2n)!$ for $n\geq2$, with $a_0=a_1=1$.
All terms are positive by induction.  Divide the identity at $n$ by
the identity at $n-1$ (valid when $n\geq3$) to obtain
\[
 \frac{a_n}{a_{n-3}}=\frac{(2n)!}{(2n-2)!}=2n(2n-1).
\]
This is exactly the asserted recurrence.  In particular the cancellation
does not assume integrality or divide by a potentially zero term.


\subsection{A383702: the selected recurrence}

\label{round7:r7entry:A383702}

\paragraph{Source and selected assertion.}

\href{https://oeis.org/A383702}{A383702}, consulted on 20 October 2026, records the conjecture attributed to R. J. Mathar, May 07 2025.

The source offset is $s=0$.  In the convention of \eqref{round7:r7:gf-operator}, the selected recurrence is $R_n=0$ for every integer $n\geq 3$, with $r=3$ and

\[\begin{aligned}
p_{0}(n)&=1\\
p_{1}(n)&=- 3 \left(4 n - 5\right)\\
p_{2}(n)&=48 n^{2} - 168 n + 151\\
p_{3}(n)&=- \left(4 n - 9\right)^{3}\\
\end{aligned}\]

\paragraph{Known inputs and scope.}

The coefficient -of-product definition is an input. Israel's March 12, 2026 proofs on neighboring A383700/A383701 are acknowledged and those IDs are excluded. The present step 4, constant 3 coefficient problem is distinct; its own exact polynomial identity is supplied, without claiming novelty of the general coefficient -elimination method.

\paragraph{Proof.}

Let $P_n(z)=\prod_{h=1}^n(z+4h-1)$, so the definition is
$a_n=[z^2]P_n(z)$.  With the displayed $p_j$, direct multiplication gives
the polynomial identity in the two indeterminates $n,z$,
\[
 \sum_{j=0}^{3}p_j(n)
          \prod_{h=j}^{2}(z+4(n-h)-1)=z^3.
\]
Multiplication by $P_{n-3}(z)$ yields
$\sum_{j=0}^3p_j(n)P_{n-j}(z)=z^3P_{n-3}(z)$ for $n\geq3$.
Taking the coefficient of $z^2$ proves the recurrence, since the
right side has no such coefficient.  The empty product is one.
This argument needs no guessed initial values and no numerical fit.


\subsection{A386830: the selected recurrence}

\label{round7:r7entry:A386830}

\paragraph{Source and selected assertion.}

\href{https://oeis.org/A386830}{A386830}, consulted on 20 October 2026, records the conjecture attributed to R. J. Mathar, Aug 19 2025.

The source offset is $s=0$.  In the convention of \eqref{round7:r7:gf-operator}, the selected recurrence is $R_n=0$ for every integer $n\geq 2$, with $r=2$ and

\[\begin{aligned}
p_{0}(n)&=8 n \left(2 n - 3\right)\\
p_{1}(n)&=- 6 \left(108 n^{2} - 207 n + 80\right)\\
p_{2}(n)&=405 \left(3 n - 4\right) \left(3 n - 2\right)\\
\end{aligned}\]

\paragraph{Known inputs and scope.}

Manyama's finite binomial sum and coefficient -extraction formula are existing inputs. The diagonal-series identity converts that definition to a rational parametrization. Only Mathar's marked order-two recurrence is selected.

\paragraph{Proof.}

The defining binomial sum is the coefficient
\[
 a_n=[u^n]\frac{(1+3u)^{3n+1}}{(1-2u)^{2n+1}}.
\]
Expansion of the numerator and denominator gives the defining summand
\[
 3^k2^{n-k}\binom{3n+1}k\binom{3n-k}{n-k}.
\]
Apply \eqref{round7:r7:gf-diagonal} with
$B(t)=(1+3t)/(1-2t)$ and
$\phi(t)=(1+3t)^3/(1-2t)^2$.
This proves the rational parametrization used below.

\[x=\xi(t)=\frac{t \left(2 t - 1\right)^{2}}{\left(3 t + 1\right)^{3}},\qquad t(0)=0,\quad \xi'(0)=1.\]

\[\phi(t)=\frac{(1+3t)^3}{(1-2t)^2},\qquad F=\frac{(1+3t)/(1-2t)}{1-t\phi'(t)/\phi(t)}.\]

Use $D_0=1$, $R_0=F(\xi(t))$, and $M=\xi/\xi'$ in \eqref{round7:r7:gf-rational-check}.  All expressions are rational in $t$.

Exact rational differentiation and cancellation in \eqref{round7:r7:gf-rational-check} give

\[LF=P(x),\qquad P(x)=- 30 x.\]

Coefficient extraction in \eqref{round7:r7:gf-operator} therefore proves the displayed recurrence for all $n\geq 2$.

The stated lower bound also ensures $n\geq s+r=2$, so no term below the source offset is used.

\section{Forty-five empirical array recurrence packages}
\label{round7:sec:r7arrays}
This section proves forty-five entry-level recurrence assertions marked
\emph{Empirical} on the primary OEIS pages consulted on 20 October 2026.
Nineteen source-stated empirical or conjectural rational generating
functions are proved in the same packages. The original arrays and
empirical recurrences are credited to R.~H.~Hardin; the additional
generating-function conjectures are credited on their pages to Colin
Barker. The counting unit is one selected A-number, even when several
formulas on its page follow from the same proof. The forty-five entries
are absent from the 608 A-numbers selected in rounds I--VI.

\subsection{The local definitions}
An array has $n+1$ rows, width $w\ge2$, and entries in
$\{0,\ldots,q-1\}$. For a block
$B=\left(\begin{smallmatrix}a&b\\c&d\end{smallmatrix}\right)$ let
\[
 \Delta(B)=(b-a,d-b,c-d,a-c).
\]
All differences here are signed integers. The entry-specific parameters
below select one of the following predicates.
\begin{description}
\item[Cell values (C).] $\#\{a,b,c,d\}\in S$.
\item[Distinct edge differences (E).] The four entries of $\Delta(B)$
have a number of distinct values belonging to $S$.
\item[Stress (D).] $|a+d-b-c|=s$.
\item[Absolute edge sum (U).] $\sum_{e\in\Delta(B)}|e|=s$.
\item[Squared edge sum (V).] $\sum_{e\in\Delta(B)}e^2=s$.
\item[Exactly one orientation (J).] Exactly one of
$\max\Delta(B)\le1$ and $\min\Delta(B)\ge-1$ holds.
\item[At least one orientation (K).] At least one of the same two
inequalities holds, allowing both to hold.
\end{description}
The sums in U and V include all four edges with multiplicity.
When the particular source additionally prohibits equal adjacent
elements, every entry of $\Delta(B)$ is required to be nonzero.
This extra condition is not imposed on sources that allow equality.
Every horizontal and vertical edge of an array with at least two rows
and columns lies in a $2\times2$ block. Thus this local nonzero condition
is exactly the source's global adjacency prohibition, including boundary
edges; diagonal neighbors are not covered by that prohibition.

\subsection{Transfers derived from the array definition}
For C, let $\gamma(r)$ rename a row's colors in first-occurrence order.
For the numeric predicates E, D, U, V, J and K, let $\gamma(r)$ be the
lexicographically least of $r$, its horizontal reversal, its simultaneous
complement $q-1-r$, and the reversal of that complement. Let
$\mathcal O_{q,w}$ be the sorted set of representatives and define
\begin{align}
 m(r)&=\#\{v\in\{0,\ldots,q-1\}^w:\gamma(v)=r\},
 \label{round7:eq:r7arr-mass}\\
 T(r,s)&=\#\left\{v\in\{0,\ldots,q-1\}^w:
 \begin{array}{l}\gamma(v)=s,\\
 P(r_i,r_{i+1},v_i,v_{i+1})\quad(1\le i<w)\end{array}\right\}.
 \label{round7:eq:r7arr-T}
\end{align}
In C, simultaneous color permutations preserve the local predicate.
For all numeric predicates, reversal and complementation permute or
negate the edge differences and preserve absolute stress. They also
interchange the two orientation inequalities, preserving both exclusive
and inclusive disjunction. Consequently each actual row in an orbit has
the same number of successors in every target orbit. Arbitrary color
permutations are not used for numeric predicates.

It follows by adding rows that the exact labeled count is
\begin{equation}\label{round7:eq:r7arr-labeled}
 L_q(n)=m^{\mathsf T}T^n\mathbf1\qquad(n\ge1).
\end{equation}
The exponent is $n$ because an array has $n+1$ rows. All orbit masses are
retained, including fixed representatives and rows with no successors.
For a source with scalar normalization $\lambda$, the required count is
$a(n)=\lambda L_q(n)$, with exactly the stated scalar. No interpretation
of that scalar as a free symmetry quotient is needed.

The C entry normalizes global color names by row-major
first-occurrence order and requires all four colors. The general
counting formulas below also record the at-most-color variant used in
the earlier rounds. Let $U_q^{=}$ and $U_q^{\le}$ count normalized
arrays using exactly $q$ and at most $q$ colors, respectively. Since a
surjective labeling has $q!$ distinct color renamings, inclusion--exclusion
gives
\begin{align}
 U_q^{=}(n)&=\sum_{j=1}^{q}\frac{(-1)^{q-j}}{j!(q-j)!}L_j(n),
 \label{round7:eq:r7arr-exact}\\
 U_q^{\le}(n)&=\sum_{j=1}^{q}
 \left(\frac1{j!}\sum_{h=0}^{q-j}\frac{(-1)^h}{h!}\right)L_j(n).
 \label{round7:eq:r7arr-atmost}
\end{align}
The second formula sums the exact-color formula over the number of
colors actually used. The empty alphabet contributes zero. All colors
are required only when the entry explicitly says so.

\subsection{An exact finite certificate proves every index}
For the source signature $c=(c_1,\ldots,c_d)$ put
\[
 R(z)=z^d-\sum_{i=1}^{d}c_i z^{d-i},\qquad
 Q(x)=1-\sum_{i=1}^{d}c_i x^i.
\]
For each component with nonzero weight in the preceding counting formula,
the complete certificate establishes the integer vector identity
\begin{equation}\label{round7:eq:r7arr-cert}
 \boxed{TR(T)\mathbf1=\mathbf0.}
\end{equation}
Multiplying by $m^{\mathsf T}T^{n-d-1}$ for $n\ge d+1$, and then summing
the weighted components, proves
\begin{equation}\label{round7:eq:r7arr-recurrence}
 \boxed{a(n)=\sum_{i=1}^{d}c_i a(n-i)\qquad(n\ge d+1).}
\end{equation}
Every source term on the right has index at least 1. The auxiliary
transfer value at $n=0$ is not asserted to be an OEIS term and can be
nonintegral for a scalar-normalized definition.

The 48 distinct complete matrices have 8,784 nonzero weighted edges;
the largest has 1,056 states, and the largest recurrence order is 50.
All 890 source-prefix values agree with the defining transfers. These
comparisons check interpretation and indexing. The certificates contain
all component weights, matrices, vectors
$v_k=T^k\mathbf1$ for $0\le k\le d+1$, and residual coordinates.
The computation is exactly
\begin{Verbatim}[fontsize=\small]
v[0] = [1 for every state]
for k = 1,...,d+1:
    v[k][r] = sum(T[r,s]*v[k-1][s] for every state s)
for every state r:
    assert v[d+1][r] == sum(c[i]*v[d+1-i][r] for i=1,...,d)
\end{Verbatim}
The verifier regenerates all states and transitions directly from
\eqref{round7:eq:r7arr-T} and checks every stored vector entry with exact integer
and rational arithmetic. Thus this is a finite computer-assisted proof:
the vector identity implies every later index. Initial-value agreement
alone is not used as the proof. Complete expanded data are supplied in
\texttt{r7\_arrays\_matrices.json} and
\texttt{r7\_arrays\_certificates.json}; the definitions and coefficient
signatures in this manuscript determine the calculation without those
files.

For $F(x)=\sum_{n\ge1}a(n)x^n$, the proved recurrence gives $QF=P$,
where $\deg P\le d$ and
\begin{equation}\label{round7:eq:r7arr-numerator}
 [x^k]P=\sum_{\substack{0\le i\le\min(k,d)\\k-i\ge1}}
 q_i a(k-i),\qquad q_0=1,\quad q_i=-c_i\ (i\ge1).
\end{equation}
The initial values are computed from the defining transfers. Exact
polynomial cross-multiplication identifies this $P/Q$ with all nineteen
selected source generating functions. Their numerators are displayed
below and their original expressions are in the source audit. Separately
stated radical closed forms are not extra selected assertions.

\subsection{A hand-checkable example: A233637}
For squared edge sum $2$, $q=3$, and $w=2$, the representatives are
$00,01,02,11$. Direct enumeration gives
\[
 T=\begin{pmatrix}0&2&0&1\\1&1&0&1\\0&0&0&0\\2&4&0&0\end{pmatrix},
 \qquad m=\begin{pmatrix}2\\4\\2\\1\end{pmatrix}.
\]
The row $02$ cannot extend to a legal block and has no outgoing edges.
The first four vectors $T^k\mathbf1$ are
\[
 (1,1,1,1)^{\mathsf T},\quad(3,3,0,6)^{\mathsf T},\quad
 (12,12,0,18)^{\mathsf T},\quad(42,42,0,72)^{\mathsf T}.
\]
The last is twice the preceding vector plus six times the second one.
This is $T(T^2-2T-6I)\mathbf1=0$, proving
$a(n)=2a(n-1)+6a(n-2)$ for $n\ge3$. Since $a(1)=24$ and $a(2)=90$,
\[
 F(x)=\frac{x(24+42x)}{1-2x-6x^2}.
\]
This proves the empirical recurrence and conjectural rational generating
function on \url{https://oeis.org/A233637}.

\subsection{Source status, provenance and the selection unit}
The status used here is the explicit word \emph{Empirical} attached to
each recurrence, not OEIS's editorial label ``approved.'' The consulted
primary pages record no proof or proof link for these selected assertions.
Representative searches for the exact A-numbers and related families did
not identify a prior proof among the inspected results; many search hits
were irrelevant. This limited audit does not establish worldwide priority.

All forty-five IDs are disjoint from the prior 608 selected project IDs
and from the retained 97 distinct A-numbers in the 100 costly claims in
Appendix A of \emph{OEIS Open}. That check does not cover the entire
repository. Five-value diagnostics under shifts $-1,0,1$ and rational
rescaling found no aliases against 169 prior array certificates or among
these forty-five entries; this is not a complete mathematical alias test.
The transfer framework adapts rounds IV--VI. Two previously certified,
uncounted Round VI reserves and one previously sourced, uncounted
candidate have freshly checked sources here; their earlier research is
credited in the ledger. Two further valid certificates, A210104 and
A210170, are held outside this selection to keep the full dossier at
100 entries; this is count curation, not a failed proof. The additional
certified reserve A235306 is also uncounted. The freshly sourced
A234653, A235218 and A235284 reserves are included in the final selection. No OEIS editing or submission
has been performed.

\subsection{Individual definitions and full recurrence signatures}
Every coefficient list below gives $c_1,\ldots,c_d$ in order, including
zeros, in \eqref{round7:eq:r7arr-recurrence}. Component dimensions include every
nonzero weighted component. Initial values start at $n=1$.
\Needspace{7\baselineskip}
\subsubsection{A209550}
\label{round7:r7entry:A209550}
Arrays have $n+1$ rows and width $w=6$ on the alphabet $0,\ldots,3$. Predicate E requires $S=\{2\}$ distinct signed clockwise edge differences; the labeled count is multiplied by $\lambda=1/4$.
The recurrence holds for every $n\ge 38$, with order $d=37$ and component dimensions $(1056)$ (total $D=1056$). Source: \url{https://oeis.org/A209550}.
The first three values are $(603,16582,425803)$.
\begin{Verbatim}[fontsize=\scriptsize,breaklines=true]
52, -856, 4014, 34244, -455004, 1232892, 6401520, -47507560, 64609279, 283015278,
-991731422, -25138511, 4334710815, -4388934859, -8391654500, 15874524375, 6183586137,
-26768058473, 3068141187, 25533122171, -9123208293, -14680808558, 7130958556,
5229264986, -2809641978, -1165852380, 598144917, 161163696, -66836588, -13174080,
3588829, 586993, -75654, -11622, 320, 76, 2
\end{Verbatim}
\Needspace{7\baselineskip}
\subsubsection{A210179}
\label{round7:r7entry:A210179}
Arrays have $n+1$ rows and width $w=5$ on the alphabet $0,\ldots,3$. Predicate C requires $S=\{3,4\}$ distinct values per block. The array uses all $4$ colors and is renamed in row-major first-occurrence order.
The recurrence holds for every $n\ge 51$, with order $d=50$ and component dimensions $(1,16,41,51)$ (total $D=109$). Source: \url{https://oeis.org/A210179}.
The first three values are $(9042,2563795,677791782)$.
\begin{Verbatim}[fontsize=\scriptsize,breaklines=true]
305, -8964, -592720, 25109411, 152644028, -18195830806, 181226187451, 3092095684717,
-66471049453238, 90264180231796, 6772584563156660, -49066540883042923,
-192090370791348766, 3446057586143788976, -5299468411446658785,
-101206410270018640248, 453991188789119119647, 1192283901355404923002,
-11763947474565853612062, 4976777711797107269820, 154357702742341519669776,
-320357552305372941467372, -1038371855595680313674168, 4129078215906298103085296,
2119123146244235654868816, -27118380193817405050087984, 17929710091623807641222208,
97507644778196665361489152, -149297639238887453328571904,
-169191295081970510706676864, 490220164048331643495809024, 23379223105512261258727168,
-836167077385995836324062720, 406266067090180580354471936,
749819009543217881638707200, -676693879122818542026817536,
-310443123972805588169031680, 498203947536501588227915776, 9744748192081361437720576,
-187420996792950565377146880, 36325389074681359144845312, 34537447300013200634806272,
-11215019873068960788250624, -2580406693969511919386624, 1132682738571576929681408,
60903702041543212793856, -38872759178345280700416, -482390924757181661184,
345714358079352471552, -5401496006984466432
\end{Verbatim}
\Needspace{7\baselineskip}
\subsubsection{A233637}
\label{round7:r7entry:A233637}
Arrays have $n+1$ rows and width $w=2$ on the alphabet $0,\ldots,2$. Predicate V has target $s=2$, allowing equal adjacent values.
The recurrence holds for every $n\ge 3$, with order $d=2$ and component dimensions $(4)$ (total $D=4$). Source: \url{https://oeis.org/A233637}.
The first three values are $(24,90,324)$.
\[(c_1,\ldots,c_d)=(2,6).\]
The source rational generating function is also proved: $F=P/Q$, with $Q=1-\sum_{i=1}^{d}c_i x^i$ and
\[P(x)=24x+42x^{2}.\]
\Needspace{7\baselineskip}
\subsubsection{A233638}
\label{round7:r7entry:A233638}
Arrays have $n+1$ rows and width $w=3$ on the alphabet $0,\ldots,2$. Predicate V has target $s=2$, allowing equal adjacent values.
The recurrence holds for every $n\ge 8$, with order $d=7$ and component dimensions $(10)$ (total $D=10$). Source: \url{https://oeis.org/A233638}.
The first three values are $(90,534,3004)$.
\[(c_1,\ldots,c_d)=(6,9,-67,26,81,-48,2).\]
The source rational generating function is also proved: $F=P/Q$, with $Q=1-\sum_{i=1}^{d}c_i x^i$ and
\[P(x)=90x-6x^{2}-1010x^{3}+624x^{4}+1238x^{5}-810x^{6}+34x^{7}.\]
\Needspace{7\baselineskip}
\subsubsection{A233639}
\label{round7:r7entry:A233639}
Arrays have $n+1$ rows and width $w=4$ on the alphabet $0,\ldots,2$. Predicate V has target $s=2$, allowing equal adjacent values.
The recurrence holds for every $n\ge 14$, with order $d=13$ and component dimensions $(25)$ (total $D=25$). Source: \url{https://oeis.org/A233639}.
The first three values are $(324,3004,26192)$.
\begin{Verbatim}[fontsize=\scriptsize,breaklines=true]
2, 76, -5, -1165, 564, 5703, -4723, -9217, 10950, 1714, -5684, 1880, -96
\end{Verbatim}
\Needspace{7\baselineskip}
\subsubsection{A233640}
\label{round7:r7entry:A233640}
Arrays have $n+1$ rows and width $w=5$ on the alphabet $0,\ldots,2$. Predicate V has target $s=2$, allowing equal adjacent values.
The recurrence holds for every $n\ge 32$, with order $d=31$ and component dimensions $(70)$ (total $D=70$). Source: \url{https://oeis.org/A233640}.
The first three values are $(1188,17424,238570)$.
\begin{Verbatim}[fontsize=\scriptsize,breaklines=true]
18, 104, -2900, 1703, 142773, -323741, -3138407, 9628188, 36089489, -132426888,
-233184244, 1009363199, 867273484, -4570942635, -1819188109, 12605037643, 1987026933,
-21096877188, -968267399, 20856565656, 317665620, -11488323206, -375223522,
3165049833, 183421296, -361542217, -20354018, 12534108, 274216, -114944, -832
\end{Verbatim}
\Needspace{7\baselineskip}
\subsubsection{A233684}
\label{round7:r7entry:A233684}
Arrays have $n+1$ rows and width $w=2$ on the alphabet $0,\ldots,3$. Predicate V has target $s=10$, with every horizontal and vertical neighboring pair unequal.
The recurrence holds for every $n\ge 4$, with order $d=3$ and component dimensions $(6)$ (total $D=6$). Source: \url{https://oeis.org/A233684}.
The first three values are $(32,108,356)$.
\[(c_1,\ldots,c_d)=(2,5,-2).\]
The source rational generating function is also proved: $F=P/Q$, with $Q=1-\sum_{i=1}^{d}c_i x^i$ and
\[P(x)=32x+44x^{2}-20x^{3}.\]
\Needspace{7\baselineskip}
\subsubsection{A233685}
\label{round7:r7entry:A233685}
Arrays have $n+1$ rows and width $w=3$ on the alphabet $0,\ldots,3$. Predicate V has target $s=10$, with every horizontal and vertical neighboring pair unequal.
The recurrence holds for every $n\ge 10$, with order $d=9$ and component dimensions $(20)$ (total $D=20$). Source: \url{https://oeis.org/A233685}.
The first three values are $(108,548,2628)$.
\begin{Verbatim}[fontsize=\scriptsize,breaklines=true]
1, 28, -25, -103, 97, 68, -47, -12, 2
\end{Verbatim}
\Needspace{7\baselineskip}
\subsubsection{A233686}
\label{round7:r7entry:A233686}
Arrays have $n+1$ rows and width $w=4$ on the alphabet $0,\ldots,3$. Predicate V has target $s=10$, with every horizontal and vertical neighboring pair unequal.
The recurrence holds for every $n\ge 16$, with order $d=15$ and component dimensions $(72)$ (total $D=72$). Source: \url{https://oeis.org/A233686}.
The first three values are $(356,2628,19248)$.
\begin{Verbatim}[fontsize=\scriptsize,breaklines=true]
7, 30, -198, -190, 1708, 65, -6015, 1916, 9134, -3808, -5930, 2155, 1461, -312, -112
\end{Verbatim}
\Needspace{7\baselineskip}
\subsubsection{A233710}
\label{round7:r7entry:A233710}
Arrays have $n+1$ rows and width $w=2$ on the alphabet $0,\ldots,4$. Predicate V has target $s=14$, allowing equal adjacent values.
The recurrence holds for every $n\ge 9$, with order $d=8$ and component dimensions $(9)$ (total $D=9$). Source: \url{https://oeis.org/A233710}.
The first three values are $(96,444,1992)$.
\begin{Verbatim}[fontsize=\scriptsize,breaklines=true]
0, 24, 14, -129, -68, 228, 64, -128
\end{Verbatim}
The source rational generating function is also proved: $F=P/Q$, with $Q=1-\sum_{i=1}^{d}c_i x^i$ and
\[P(x)=96x+444x^{2}-312x^{3}-2900x^{4}-384x^{5}+5296x^{6}+832x^{7}-2944x^{8}.\]
\Needspace{7\baselineskip}
\subsubsection{A233711}
\label{round7:r7entry:A233711}
Arrays have $n+1$ rows and width $w=3$ on the alphabet $0,\ldots,4$. Predicate V has target $s=14$, allowing equal adjacent values.
The recurrence holds for every $n\ge 28$, with order $d=27$ and component dimensions $(39)$ (total $D=39$). Source: \url{https://oeis.org/A233711}.
The first three values are $(444,2624,16004)$.
\begin{Verbatim}[fontsize=\scriptsize,breaklines=true]
8, 56, -565, -766, 15139, -7664, -200684, 302993, 1413766, -3297757, -5115754,
18205660, 6468818, -56612688, 15267568, 100195210, -68399612, -95039845, 102899681,
36863396, -73128319, 5984214, 22126722, -7453432, -1251600, 860624, -95168
\end{Verbatim}
\Needspace{7\baselineskip}
\subsubsection{A233749}
\label{round7:r7entry:A233749}
Arrays have $n+1$ rows and width $w=2$ on the alphabet $0,\ldots,5$. Predicate U has target $s=10$, allowing equal adjacent values.
The recurrence holds for every $n\ge 10$, with order $d=9$ and component dimensions $(12)$ (total $D=12$). Source: \url{https://oeis.org/A233749}.
The first three values are $(292,3244,33372)$.
\begin{Verbatim}[fontsize=\scriptsize,breaklines=true]
8, 50, -191, -670, 1425, 2761, -3160, -644, 358
\end{Verbatim}
\Needspace{7\baselineskip}
\subsubsection{A233750}
\label{round7:r7entry:A233750}
Arrays have $n+1$ rows and width $w=3$ on the alphabet $0,\ldots,5$. Predicate U has target $s=10$, allowing equal adjacent values.
The recurrence holds for every $n\ge 45$, with order $d=44$ and component dimensions $(63)$ (total $D=63$). Source: \url{https://oeis.org/A233750}.
The first three values are $(3244,91116,2317964)$.
\begin{Verbatim}[fontsize=\scriptsize,breaklines=true]
28, 258, -8153, -21818, 884110, 641196, -49138214, 7881936, 1607832444, -1043204114,
-33437651311, 31731911114, 463391143344, -530799098345, -4407915878767, 5632602673516,
29289730685819, -40018456626533, -137180728733978, 194997456311814, 454206116121400,
-657831115817189, -1062348157268958, 1537508763715626, 1751866868878969,
-2476743191321207, -2036444198056430, 2723651246669803, 1674200095938075,
-2016791049187732, -976960955082435, 985471325349785, 401735680540938,
-307547802798483, -112638329429263, 58076876665437, 20060500861421, -6092017479043,
-2023405058657, 313961613786, 99780981302, -5783024558, -1748862224, -15322920
\end{Verbatim}
\Needspace{7\baselineskip}
\subsubsection{A233885}
\label{round7:r7entry:A233885}
Arrays have $n+1$ rows and width $w=2$ on the alphabet $0,\ldots,6$. Predicate V has target $s=30$, with every horizontal and vertical neighboring pair unequal.
The recurrence holds for every $n\ge 11$, with order $d=10$ and component dimensions $(16)$ (total $D=16$). Source: \url{https://oeis.org/A233885}.
The first three values are $(112,384,1264)$.
\begin{Verbatim}[fontsize=\scriptsize,breaklines=true]
0, 18, 4, -95, -30, 165, 36, -105, -8, 20
\end{Verbatim}
The source rational generating function is also proved: $F=P/Q$, with $Q=1-\sum_{i=1}^{d}c_i x^i$ and
\[\begin{aligned}P(x)&=112x+384x^{2}-752x^{3}-3068x^{4}+624x^{5}+5828x^{6}+224x^{7}-3760x^{8}-88x^{9}\\ &\quad +720x^{10}.\end{aligned}\]
\Needspace{7\baselineskip}
\subsubsection{A233886}
\label{round7:r7entry:A233886}
Arrays have $n+1$ rows and width $w=3$ on the alphabet $0,\ldots,6$. Predicate V has target $s=30$, with every horizontal and vertical neighboring pair unequal.
The recurrence holds for every $n\ge 33$, with order $d=32$ and component dimensions $(100)$ (total $D=100$). Source: \url{https://oeis.org/A233886}.
The first three values are $(384,1960,9376)$.
\begin{Verbatim}[fontsize=\scriptsize,breaklines=true]
4, 72, -278, -2093, 7846, 31943, -118256, -280293, 1059474, 1464690, -5968574,
-4541666, 21819710, 7753185, -52635660, -4645134, 84280382, -7066733, -89291380,
16609381, 61880902, -14399047, -27456604, 6388152, 7484866, -1438992, -1148164,
135572, 79124, -1776, -736, 32
\end{Verbatim}
\Needspace{7\baselineskip}
\subsubsection{A233910}
\label{round7:r7entry:A233910}
Arrays have $n+1$ rows and width $w=2$ on the alphabet $0,\ldots,7$. Predicate V has target $s=30$, with every horizontal and vertical neighboring pair unequal.
The recurrence holds for every $n\ge 12$, with order $d=11$ and component dimensions $(20)$ (total $D=20$). Source: \url{https://oeis.org/A233910}.
The first three values are $(160,640,2440)$.
\begin{Verbatim}[fontsize=\scriptsize,breaklines=true]
0, 24, 6, -177, -52, 516, 80, -646, 42, 284, -80
\end{Verbatim}
\Needspace{7\baselineskip}
\subsubsection{A233911}
\label{round7:r7entry:A233911}
Arrays have $n+1$ rows and width $w=3$ on the alphabet $0,\ldots,7$. Predicate V has target $s=30$, with every horizontal and vertical neighboring pair unequal.
The recurrence holds for every $n\ge 39$, with order $d=38$ and component dimensions $(144)$ (total $D=144$). Source: \url{https://oeis.org/A233911}.
The first three values are $(640,3944,22632)$.
\begin{Verbatim}[fontsize=\scriptsize,breaklines=true]
4, 122, -461, -6376, 22883, 188152, -648596, -3478827, 11715412, 42308852, -142371381,
-345931578, 1197824103, 1908669954, -7078856586, -6995027392, 29551303874,
16143386244, -87052258170, -19433628990, 179619345422, -2128109482, -255775086626,
46189589396, 245226941694, -73044737386, -152167442028, 56827325500, 57294922304,
-23672423108, -11706249412, 4940867468, 1040756704, -439680112, -20118176, 11155520,
-250624, -14592
\end{Verbatim}
\Needspace{7\baselineskip}
\subsubsection{A233951}
\label{round7:r7entry:A233951}
Arrays have $n+1$ rows and width $w=2$ on the alphabet $0,\ldots,5$. Predicate V has target $s=10$, with every horizontal and vertical neighboring pair unequal.
The recurrence holds for every $n\ge 5$, with order $d=4$ and component dimensions $(12)$ (total $D=12$). Source: \url{https://oeis.org/A233951}.
The first three values are $(80,380,1776)$.
\[(c_1,\ldots,c_d)=(4,5,-7,-6).\]
The source rational generating function is also proved: $F=P/Q$, with $Q=1-\sum_{i=1}^{d}c_i x^i$ and
\[P(x)=80x+60x^{2}-144x^{3}-108x^{4}.\]
\Needspace{7\baselineskip}
\subsubsection{A233952}
\label{round7:r7entry:A233952}
Arrays have $n+1$ rows and width $w=3$ on the alphabet $0,\ldots,5$. Predicate V has target $s=10$, with every horizontal and vertical neighboring pair unequal.
The recurrence holds for every $n\ge 20$, with order $d=19$ and component dimensions $(63)$ (total $D=63$). Source: \url{https://oeis.org/A233952}.
The first three values are $(380,2820,20616)$.
\begin{Verbatim}[fontsize=\scriptsize,breaklines=true]
3, 76, -161, -1932, 3277, 22377, -32880, -128727, 173771, 368053, -472908, -488012,
600604, 270148, -316676, -46496, 44736, 6656, -256
\end{Verbatim}
\Needspace{7\baselineskip}
\subsubsection{A234146}
\label{round7:r7entry:A234146}
Arrays have $n+1$ rows and width $w=2$ on the alphabet $0,\ldots,5$. Predicate U has target $s=10$, with every horizontal and vertical neighboring pair unequal.
The recurrence holds for every $n\ge 9$, with order $d=8$ and component dimensions $(12)$ (total $D=12$). Source: \url{https://oeis.org/A234146}.
The first three values are $(184,1264,8752)$.
\begin{Verbatim}[fontsize=\scriptsize,breaklines=true]
6, 25, -105, -192, 494, 478, -322, -204
\end{Verbatim}
The source rational generating function is also proved: $F=P/Q$, with $Q=1-\sum_{i=1}^{d}c_i x^i$ and
\[P(x)=184x+160x^{2}-3432x^{3}-2788x^{4}+16796x^{5}+12460x^{6}-10692x^{7}-6120x^{8}.\]
\Needspace{7\baselineskip}
\subsubsection{A234147}
\label{round7:r7entry:A234147}
Arrays have $n+1$ rows and width $w=3$ on the alphabet $0,\ldots,5$. Predicate U has target $s=10$, with every horizontal and vertical neighboring pair unequal.
The recurrence holds for every $n\ge 38$, with order $d=37$ and component dimensions $(63)$ (total $D=63$). Source: \url{https://oeis.org/A234147}.
The first three values are $(1264,13936,158812)$.
\begin{Verbatim}[fontsize=\scriptsize,breaklines=true]
10, 188, -1705, -13231, 115068, 493976, -4229448, -11190946, 96334439, 164224925,
-1453766839, -1614673985, 15099553644, 10806854528, -110319519409, -49482794946,
573293794815, 155510560962, -2125916423322, -344648421788, 5607869207405,
602950452380, -10426809489560, -1017179390188, 13440589196562, 1660271297890,
-11680287015490, -2067440715994, 6524300789888, 1631686761052, -2149016060048,
-727542691768, 347865584016, 160185863728, -13689852000, -13226096640, -1457602560
\end{Verbatim}
\Needspace{7\baselineskip}
\subsubsection{A234651}
\label{round7:r7entry:A234651}
Arrays have $n+1$ rows and width $w=2$ on the alphabet $0,\ldots,5$. Predicate D has stress $s=4$.
The recurrence holds for every $n\ge 6$, with order $d=5$ and component dimensions $(12)$ (total $D=12$). Source: \url{https://oeis.org/A234651}.
The first three values are $(160,728,3264)$.
\[(c_1,\ldots,c_d)=(4,44,-176,-480,1920).\]
The source rational generating function is also proved: $F=P/Q$, with $Q=1-\sum_{i=1}^{d}c_i x^i$ and
\[P(x)=160x+88x^{2}-6688x^{3}-1920x^{4}+69120x^{5}.\]
\Needspace{7\baselineskip}
\subsubsection{A234652}
\label{round7:r7entry:A234652}
Arrays have $n+1$ rows and width $w=3$ on the alphabet $0,\ldots,5$. Predicate D has stress $s=4$.
The recurrence holds for every $n\ge 16$, with order $d=15$ and component dimensions $(63)$ (total $D=63$). Source: \url{https://oeis.org/A234652}.
The first three values are $(728,2648,9616)$.
\begin{Verbatim}[fontsize=\scriptsize,breaklines=true]
9, 59, -741, -680, 23970, -24980, -382080, 836736, 3002976, -9941184, -8688384,
53360640, -15759360, -107827200, 99532800
\end{Verbatim}
\Needspace{7\baselineskip}
\subsubsection{A234653}
\label{round7:r7entry:A234653}
Arrays have $n+1$ rows and width $w=4$ on the alphabet $0,\ldots,5$. Predicate D has stress $s=4$.
The recurrence holds for every $n\ge 22$, with order $d=21$ and component dimensions $(342)$ (total $D=342$). Source: \url{https://oeis.org/A234653}.
The first three values are $(3264,9616,28496)$.
\begin{Verbatim}[fontsize=\scriptsize,breaklines=true]
8, 83, -804, -2479, 34112, 22457, -798036, 453896, 11262032, -15161188, -98135856,
191644944, 513355968, -1327868352, -1430555904, 5245827840, 1148359680, -11006668800,
3284582400, 9455616000, -5971968000
\end{Verbatim}
\Needspace{7\baselineskip}
\subsubsection{A234731}
\label{round7:r7entry:A234731}
Arrays have $n+1$ rows and width $w=2$ on the alphabet $0,\ldots,7$. Predicate D has stress $s=8$.
The recurrence holds for every $n\ge 8$, with order $d=7$ and component dimensions $(20)$ (total $D=20$). Source: \url{https://oeis.org/A234731}.
The first three values are $(168,672,2184)$.
\[(c_1,\ldots,c_d)=(4,34,-136,-369,1476,1260,-5040).\]
The source rational generating function is also proved: $F=P/Q$, with $Q=1-\sum_{i=1}^{d}c_i x^i$ and
\[P(x)=168x-6216x^{3}+74160x^{5}-282240x^{7}.\]
\Needspace{7\baselineskip}
\subsubsection{A234732}
\label{round7:r7entry:A234732}
Arrays have $n+1$ rows and width $w=3$ on the alphabet $0,\ldots,7$. Predicate D has stress $s=8$.
The recurrence holds for every $n\ge 24$, with order $d=23$ and component dimensions $(144)$ (total $D=144$). Source: \url{https://oeis.org/A234732}.
The first three values are $(672,2184,6048)$.
\begin{Verbatim}[fontsize=\scriptsize,breaklines=true]
8, 53, -564, -888, 17184, -2382, -296784, 308859, 3196608, -5284083, -22169196,
48357478, 97763896, -273091508, -253071696, 976703832, 271829664, -2148744960,
304750080, 2640988800, -1133222400, -1378944000, 870912000
\end{Verbatim}
\Needspace{7\baselineskip}
\subsubsection{A234733}
\label{round7:r7entry:A234733}
Arrays have $n+1$ rows and width $w=4$ on the alphabet $0,\ldots,7$. Predicate D has stress $s=8$.
The recurrence holds for every $n\ge 24$, with order $d=23$ and component dimensions $(1056)$ (total $D=1056$). Source: \url{https://oeis.org/A234733}.
The first three values are $(2184,6048,14760)$.
\begin{Verbatim}[fontsize=\scriptsize,breaklines=true]
8, 53, -564, -888, 17184, -2382, -296784, 308859, 3196608, -5284083, -22169196,
48357478, 97763896, -273091508, -253071696, 976703832, 271829664, -2148744960,
304750080, 2640988800, -1133222400, -1378944000, 870912000
\end{Verbatim}
\Needspace{7\baselineskip}
\subsubsection{A234753}
\label{round7:r7entry:A234753}
Arrays have $n+1$ rows and width $w=2$ on the alphabet $0,\ldots,3$. Predicate K permits at least one orientation whose four edge jumps are all at most $+1$, allowing both.
The recurrence holds for every $n\ge 7$, with order $d=6$ and component dimensions $(6)$ (total $D=6$). Source: \url{https://oeis.org/A234753}.
The first three values are $(110,1014,8968)$.
\[(c_1,\ldots,c_d)=(8,10,-16,-8,4,1).\]
The source rational generating function is also proved: $F=P/Q$, with $Q=1-\sum_{i=1}^{d}c_i x^i$ and
\[P(x)=110x+134x^{2}-244x^{3}-114x^{4}+66x^{5}+16x^{6}.\]
\Needspace{7\baselineskip}
\subsubsection{A234754}
\label{round7:r7entry:A234754}
Arrays have $n+1$ rows and width $w=3$ on the alphabet $0,\ldots,3$. Predicate K permits at least one orientation whose four edge jumps are all at most $+1$, allowing both.
The recurrence holds for every $n\ge 18$, with order $d=17$ and component dimensions $(20)$ (total $D=20$). Source: \url{https://oeis.org/A234754}.
The first three values are $(1014,26642,646036)$.
\begin{Verbatim}[fontsize=\scriptsize,breaklines=true]
27, -26, -884, 1597, 7345, -14758, -23881, 52625, 29984, -81577, -7015, 51495, -7850,
-9139, 1340, 518, -36
\end{Verbatim}
\Needspace{7\baselineskip}
\subsubsection{A234770}
\label{round7:r7entry:A234770}
Arrays have $n+1$ rows and width $w=2$ on the alphabet $0,\ldots,4$. Predicate J permits exactly one orientation whose four edge jumps are all at most $+1$.
The recurrence holds for every $n\ge 7$, with order $d=6$ and component dimensions $(9)$ (total $D=9$). Source: \url{https://oeis.org/A234770}.
The first three values are $(88,454,2056)$.
\[(c_1,\ldots,c_d)=(1,21,-1,-67,-10,46).\]
The source rational generating function is also proved: $F=P/Q$, with $Q=1-\sum_{i=1}^{d}c_i x^i$ and
\[P(x)=88x+366x^{2}-246x^{3}-1426x^{4}-98x^{5}+1026x^{6}.\]
\Needspace{7\baselineskip}
\subsubsection{A234771}
\label{round7:r7entry:A234771}
Arrays have $n+1$ rows and width $w=3$ on the alphabet $0,\ldots,4$. Predicate J permits exactly one orientation whose four edge jumps are all at most $+1$.
The recurrence holds for every $n\ge 24$, with order $d=23$ and component dimensions $(39)$ (total $D=39$). Source: \url{https://oeis.org/A234771}.
The first three values are $(454,3980,28840)$.
\begin{Verbatim}[fontsize=\scriptsize,breaklines=true]
9, 60, -657, -510, 14451, -9878, -143349, 195967, 711924, -1328829, -1738223, 4362370,
1713737, -7406892, 154943, 6781020, -1542017, -3326149, 1109694, 807890, -310656,
-75048, 30144
\end{Verbatim}
\Needspace{7\baselineskip}
\subsubsection{A235091}
\label{round7:r7entry:A235091}
Arrays have $n+1$ rows and width $w=2$ on the alphabet $0,\ldots,5$. Predicate D has stress $s=4$, with every horizontal and vertical neighboring pair unequal.
The recurrence holds for every $n\ge 6$, with order $d=5$ and component dimensions $(12)$ (total $D=12$). Source: \url{https://oeis.org/A235091}.
The first three values are $(80,256,808)$.
\[(c_1,\ldots,c_d)=(3,13,-39,-26,78).\]
The source rational generating function is also proved: $F=P/Q$, with $Q=1-\sum_{i=1}^{d}c_i x^i$ and
\[P(x)=80x+16x^{2}-1000x^{3}-52x^{4}+2004x^{5}.\]
\Needspace{7\baselineskip}
\subsubsection{A235092}
\label{round7:r7entry:A235092}
Arrays have $n+1$ rows and width $w=3$ on the alphabet $0,\ldots,5$. Predicate D has stress $s=4$, with every horizontal and vertical neighboring pair unequal.
The recurrence holds for every $n\ge 17$, with order $d=16$ and component dimensions $(63)$ (total $D=63$). Source: \url{https://oeis.org/A235092}.
The first three values are $(256,656,1728)$.
\begin{Verbatim}[fontsize=\scriptsize,breaklines=true]
6, 20, -178, -45, 2028, -1568, -11316, 15017, 32278, -57324, -42818, 104201, 14456,
-82056, 10920, 16380
\end{Verbatim}
\Needspace{7\baselineskip}
\subsubsection{A235093}
\label{round7:r7entry:A235093}
Arrays have $n+1$ rows and width $w=4$ on the alphabet $0,\ldots,5$. Predicate D has stress $s=4$, with every horizontal and vertical neighboring pair unequal.
The recurrence holds for every $n\ge 23$, with order $d=22$ and component dimensions $(342)$ (total $D=342$). Source: \url{https://oeis.org/A235093}.
The first three values are $(808,1728,3832)$.
\begin{Verbatim}[fontsize=\scriptsize,breaklines=true]
7, 20, -239, 36, 3285, -4086, -23439, 45882, 91959, -251226, -184209, 783240, 99087,
-1431666, 277611, 1499767, -513754, -855386, 293540, 256932, -54600, -32760
\end{Verbatim}
\Needspace{7\baselineskip}
\subsubsection{A235216}
\label{round7:r7entry:A235216}
Arrays have $n+1$ rows and width $w=2$ on the alphabet $0,\ldots,2$. Predicate K permits at least one orientation whose four edge jumps are all at most $+1$, allowing both.
The recurrence holds for every $n\ge 5$, with order $d=4$ and component dimensions $(4)$ (total $D=4$). Source: \url{https://oeis.org/A235216}.
The first three values are $(59,427,3037)$.
\[(c_1,\ldots,c_d)=(6,9,-6,-2).\]
The source rational generating function is also proved: $F=P/Q$, with $Q=1-\sum_{i=1}^{d}c_i x^i$ and
\[P(x)=59x+73x^{2}-56x^{3}-18x^{4}.\]
\Needspace{7\baselineskip}
\subsubsection{A235217}
\label{round7:r7entry:A235217}
Arrays have $n+1$ rows and width $w=3$ on the alphabet $0,\ldots,2$. Predicate K permits at least one orientation whose four edge jumps are all at most $+1$, allowing both.
The recurrence holds for every $n\ge 9$, with order $d=8$ and component dimensions $(10)$ (total $D=10$). Source: \url{https://oeis.org/A235217}.
The first three values are $(427,8029,145669)$.
\begin{Verbatim}[fontsize=\scriptsize,breaklines=true]
16, 52, -188, -10, 188, -29, -40, 10
\end{Verbatim}
\Needspace{7\baselineskip}
\subsubsection{A235218}
\label{round7:r7entry:A235218}
Arrays have $n+1$ rows and width $w=4$ on the alphabet $0,\ldots,2$. Predicate K permits at least one orientation whose four edge jumps are all at most $+1$, allowing both.
The recurrence holds for every $n\ge 22$, with order $d=21$ and component dimensions $(25)$ (total $D=25$). Source: \url{https://oeis.org/A235218}.
The first three values are $(3037,145669,6630699)$.
\begin{Verbatim}[fontsize=\scriptsize,breaklines=true]
39, 410, -3737, -13178, 103275, 89722, -961878, -181366, 4083283, -48872, -8784678,
444362, 9715586, -378256, -5306393, 93253, 1319247, 1378, -124156, -3168, 2456
\end{Verbatim}
\Needspace{7\baselineskip}
\subsubsection{A235282}
\label{round7:r7entry:A235282}
Arrays have $n+1$ rows and width $w=2$ on the alphabet $0,\ldots,3$. Predicate D has stress $s=4$.
The recurrence holds for every $n\ge 4$, with order $d=3$ and component dimensions $(6)$ (total $D=6$). Source: \url{https://oeis.org/A235282}.
The first three values are $(20,40,68)$.
\[(c_1,\ldots,c_d)=(2,3,-6).\]
The source rational generating function is also proved: $F=P/Q$, with $Q=1-\sum_{i=1}^{d}c_i x^i$ and
\[P(x)=20x-72x^{3}.\]
\Needspace{7\baselineskip}
\subsubsection{A235283}
\label{round7:r7entry:A235283}
Arrays have $n+1$ rows and width $w=3$ on the alphabet $0,\ldots,3$. Predicate D has stress $s=4$.
The recurrence holds for every $n\ge 7$, with order $d=6$ and component dimensions $(20)$ (total $D=20$). Source: \url{https://oeis.org/A235283}.
The first three values are $(40,68,104)$.
\[(c_1,\ldots,c_d)=(3,3,-15,4,18,-12).\]
The source rational generating function is also proved: $F=P/Q$, with $Q=1-\sum_{i=1}^{d}c_i x^i$ and
\[P(x)=40x-52x^{2}-220x^{3}+272x^{4}+288x^{5}-336x^{6}.\]
\Needspace{7\baselineskip}
\subsubsection{A235284}
\label{round7:r7entry:A235284}
Arrays have $n+1$ rows and width $w=4$ on the alphabet $0,\ldots,3$. Predicate D has stress $s=4$.
The recurrence holds for every $n\ge 7$, with order $d=6$ and component dimensions $(72)$ (total $D=72$). Source: \url{https://oeis.org/A235284}.
The first three values are $(68,104,148)$.
\[(c_1,\ldots,c_d)=(3,3,-15,4,18,-12).\]
\Needspace{7\baselineskip}
\subsubsection{A235291}
\label{round7:r7entry:A235291}
Arrays have $n+1$ rows and width $w=2$ on the alphabet $0,\ldots,3$. Predicate D has stress $s=3$.
The recurrence holds for every $n\ge 5$, with order $d=4$ and component dimensions $(6)$ (total $D=6$). Source: \url{https://oeis.org/A235291}.
The first three values are $(40,108,272)$.
\[(c_1,\ldots,c_d)=(0,14,0,-48).\]
The source rational generating function is also proved: $F=P/Q$, with $Q=1-\sum_{i=1}^{d}c_i x^i$ and
\[P(x)=40x+108x^{2}-288x^{3}-768x^{4}.\]
\Needspace{7\baselineskip}
\subsubsection{A235292}
\label{round7:r7entry:A235292}
Arrays have $n+1$ rows and width $w=3$ on the alphabet $0,\ldots,3$. Predicate D has stress $s=3$.
The recurrence holds for every $n\ge 11$, with order $d=10$ and component dimensions $(20)$ (total $D=20$). Source: \url{https://oeis.org/A235292}.
The first three values are $(108,232,476)$.
\begin{Verbatim}[fontsize=\scriptsize,breaklines=true]
3, 17, -57, -86, 372, 76, -972, 360, 864, -576
\end{Verbatim}
The source rational generating function is also proved: $F=P/Q$, with $Q=1-\sum_{i=1}^{d}c_i x^i$ and
\[\begin{aligned}P(x)&=108x-92x^{2}-2056x^{3}+1916x^{4}+13580x^{5}-13472x^{6}-36408x^{7}+37296x^{8}\\ &\quad +33408x^{9}-34560x^{10}.\end{aligned}\]
\Needspace{7\baselineskip}
\subsubsection{A235293}
\label{round7:r7entry:A235293}
Arrays have $n+1$ rows and width $w=4$ on the alphabet $0,\ldots,3$. Predicate D has stress $s=3$.
The recurrence holds for every $n\ge 11$, with order $d=10$ and component dimensions $(72)$ (total $D=72$). Source: \url{https://oeis.org/A235293}.
The first three values are $(272,476,832)$.
\begin{Verbatim}[fontsize=\scriptsize,breaklines=true]
3, 17, -57, -86, 372, 76, -972, 360, 864, -576
\end{Verbatim}
The source rational generating function is also proved: $F=P/Q$, with $Q=1-\sum_{i=1}^{d}c_i x^i$ and
\[\begin{aligned}P(x)&=272x-340x^{2}-5220x^{3}+6652x^{4}+34708x^{5}-44608x^{6}-93528x^{7}+119472x^{8}\\ &\quad +86400x^{9}-108288x^{10}.\end{aligned}\]
\Needspace{7\baselineskip}
\subsubsection{A235294}
\label{round7:r7entry:A235294}
Arrays have $n+1$ rows and width $w=5$ on the alphabet $0,\ldots,3$. Predicate D has stress $s=3$.
The recurrence holds for every $n\ge 11$, with order $d=10$ and component dimensions $(272)$ (total $D=272$). Source: \url{https://oeis.org/A235294}.
The first three values are $(744,1132,1736)$.
\begin{Verbatim}[fontsize=\scriptsize,breaklines=true]
3, 17, -57, -86, 372, 76, -972, 360, 864, -576
\end{Verbatim}
The source rational generating function is also proved: $F=P/Q$, with $Q=1-\sum_{i=1}^{d}c_i x^i$ and
\[\begin{aligned}P(x)&=744x-1100x^{2}-14308x^{3}+21140x^{4}+95228x^{5}-139616x^{6}-256584x^{7}\\ &\quad +369360x^{8}+237312x^{9}-331776x^{10}.\end{aligned}\]
\Needspace{7\baselineskip}
\subsubsection{A235314}
\label{round7:r7entry:A235314}
Arrays have $n+1$ rows and width $w=4$ on the alphabet $0,\ldots,6$. Predicate D has stress $s=7$.
The recurrence holds for every $n\ge 19$, with order $d=18$ and component dimensions $(625)$ (total $D=625$). Source: \url{https://oeis.org/A235314}.
The first three values are $(1120,2744,5968)$.
\begin{Verbatim}[fontsize=\scriptsize,breaklines=true]
4, 47, -200, -901, 4204, 8957, -48440, -47194, 334096, 101948, -1410080, 170136,
3549696, -1446912, -4861440, 2954880, 2764800, -2073600
\end{Verbatim}


\section{Verification, provenance and reuse}
\label{round7:r7:verification}
The accompanying archive contains the complete manuscript, numbered
ledger, defining data, exact certificates, verifiers, source audits and
cross-review reports. With Python, SymPy and mpmath installed, run
\begin{Verbatim}[fontsize=\small]
python3 r7_run_all.py
\end{Verbatim}
from its extracted directory. No network access is needed.
The runner checks the distinct-A-ID count, all previous selections,
the retained prior-work baseline, immutable review bindings, and each
verification program. It records the actual versions and results in
\texttt{oeis\_hunt\_round7\_verification\_summary.json}.

The array certificates use complete integer matrices derived from the
definitions. An identity $TR(T)\mathbf1=0$ implies the recurrence for
all later indices after multiplication by an arbitrary power of $T$.
The generating-function certificates prove exact differential identities
whose residuals are explicitly displayed polynomials; coefficients beyond
the polynomial degree vanish for every index. These are finite exact
certificates for infinite statements. Finite-prefix and numerical
comparisons check interpretation and implementation but are not used as
proofs of universal claims.

The arithmetic and formal-series parts give complete mathematical
arguments. Their verification programs check exact identities and
representative computations; they are not claimed to constitute
proof-assistant formalizations. Separate review implementations and
source-definition checks provide additional error detection.

\subsection{Source and priority scope}
Each selected theorem links to its OEIS source. The JSON records preserve
the exact target, offset, correction, attribution, retrieval reference and
the known formula used as an input. The audit distinguishes the source's
conjectural wording from its editorial approval status. Source wording
alone does not prove that a mathematical result was unknown elsewhere.

The combined audit compares the previous project ledgers and the retained
OEIS Open Appendix A baseline. The generating-function audit additionally
screens the relevant Fischer (2020) and Ekhad--Yang (2017) families.
Other already proved or equivalent claims found during this round are
excluded and explained in the component audits. The bibliography and
links are not represented as an exhaustive historical search.
In particular, the still-marked recurrences on A065942 and A226302 are
exact consequences of already recorded shorter recurrences and were
removed. A359643 was conservatively excluded after the official
Alkahest guide identified its own earlier recurrence result; that
project's complete result catalogue was not audited.

All previously selected IDs remain excluded, including the five Round IV
entries whose priority claim was withdrawn in Round V:
A176605, A176607, A176610, A176611 and A176612. Their valid formulas
were already implied by Georg Fischer's 2020 general derivation.
The current result count is a count of precisely scoped, source-marked
assertions resolved in this packet, not a cumulative worldwide-discovery
claim.

\subsection{Standalone manuscript and reproducibility}
The delivered \texttt{.tex} file contains all prose, formulas and
individual recurrence signatures directly; it has no external
\texttt{\string\input} or \texttt{\string\include} dependency.
Compile it twice with pdfLaTeX and the standard packages in its preamble.
\texttt{r7\_build\_dossier.py} reconstructs the manuscript and ledger
from the retained component records. The README gives the certificate
regeneration commands and dependencies. SHA-256 digests in the archive
manifest bind the delivered files; review reports separately bind the
proof and implementation versions they inspected.

\subsection{Selected common references}
\begin{enumerate}
\item The OEIS Foundation, \emph{The On-Line Encyclopedia of Integer
Sequences}, \url{https://oeis.org/}. Individual entry links accompany
the statements throughout the dossier.
\item NIST Digital Library of Mathematical Functions,
\href{https://dlmf.nist.gov/19.20}{\S19.20},
\href{https://dlmf.nist.gov/20.7}{\S20.7} and
\href{https://dlmf.nist.gov/20.9}{\S20.9}: classical elliptic special
values and theta identities used in the offset correction.
\item \emph{OEIS Open}, retained Appendix A comparison,
\url{https://arxiv.org/html/2608.11941v1}.
The retained scope is 100 costliest claims on 97 distinct A-IDs among
153 combined-model resolutions, not the complete repository.
\item Georg Fischer, 2020 family derivation linked at
\url{https://oeis.org/A176604/a176604_2.txt}; used for exclusion and
priority auditing, not as a newly claimed result.
\item S. B. Ekhad and D. Yang, 2017 recurrence family and supporting
output at \url{https://sites.math.rutgers.edu/~zeilberg/tokhniot/oMathar1maple12.txt};
used for exclusion and credit.
\item Alkahest, official user guide,
\url{https://alkahest-cas.github.io/alkahest/print.html},
novelty-filtering discussion identifying A359643 as an earlier result;
used as a conservative exclusion lead.
\end{enumerate}
% ===== END ORIGINAL ROUND 7 =====

% ===== BEGIN ORIGINAL ROUND 8 =====
\clearpage
\setcounter{section}{0}\setcounter{equation}{0}\setcounter{theorem}{0}
\setcounter{table}{0}\setcounter{figure}{0}
\normalsize\renewcommand{\arraystretch}{1}\renewcommand{\baselinestretch}{1}\selectfont
\part[Round 8: Denominators, exponent classification, arithmetic, and recurrence certificates]{Denominators, exponent classification, arithmetic, and recurrence certificates}
\label{combined:round8}
\rhead{\small Round 8 / KEIO}
\begingroup\small
\section*{Round overview}
\addcontentsline{toc}{section}{Round 8 overview and scope}
This eighth-round dossier resolves 100 selected assertion packages on
100 distinct OEIS entries, disjoint from the 708 A-numbers selected in
Rounds I--VII. It proves 98 selected claims, corrects and proves one
modulo-four support rule, and refutes one empirical classification by
an explicit counterexample and an infinite omitted family. The latter
package also proves the lower bound in the same source comment.

The four parts contain 13 packages on formal series, denominators and
congruences, 12 on arithmetic identities and a refutation, 26 on
generating-function recurrences, and 49 on constrained arrays. They
use coefficient-ring arguments, valuation and divisor identities,
exact differential equations, and finite transfer operators.

\subsection*{What the count means}
Each catalogue row selects one precisely scoped assertion package on
one distinct A-number. All selected assertions are explicitly marked
conjectural or empirical in the consulted sources. A recurrence and
its generating function on the same entry count once. Several entries
share a general lemma or method; in particular the 49 array packages
are distinct fixed-parameter counting problems certified within one
transfer framework. The count does not assert 100 unrelated methods.

The source wording is evidence of an unresolved label in the consulted
entry. It does not certify worldwide historical priority. Established
formulas, classical theorems and earlier project methods are credited
inputs. Every selected scope is recorded in the numbered catalogue,
the proofs and the machine-readable ledger.

\subsection*{Correction, refutation and auxiliary source issues}
For \oeis{A380680}, the proposed square support modulo 4 is false:
$a(1)=4$ and $a(2)=6$. The correct positive-index support is
$n=2j^2$, $j\ge1$, where the residue is 2; elsewhere it is 0.

For \oeis{A345993}, the proposed classification of $a(n)=3$ omits
$n=375$, where the least admissible $m$ is 125. More generally,
$n=3p^e$ works whenever $p\equiv2\pmod3$ is prime and $e$ is odd.
The statement $a(n)\ge3$ for $n\ge3$ is proved. A complete replacement
classification is not asserted.

All 26 selected generating-function recurrence coefficient lists are
proved unchanged, with explicit valid starting ranges. Auxiliary input
issues are disclosed: \oeis{A308815}'s title requires the alternating
sign, an older displayed GF for \oeis{A113409} conflicts with its
defining sum, and older unmarked formulas for \oeis{A174169} require
care. The proofs derive or identify the consistent inputs explicitly;
these repairs add no selected counts.

\subsection*{Prior work and uncounted material}
The audit excludes \oeis{A268149}, already covered by a linked paper's
general formula, and \oeis{A101601}, whose recurrence is an immediate
consequence of the recorded hypergeometric closed form. Other stale
labels and already proved claims found during the search are documented
in the component audits. The previous 708 is an exclusion ledger, not
a certified discovery total. The five Round IV priority withdrawals
disclosed in Round V remain in force and remain excluded.

The integrality argument for \oeis{A300617} is supplied as an uncounted
supplement because its companion already states integrality without a
conjecture qualifier. The stronger triangular-divisibility assertion
is not proved here. The companion numerator consequence for
\oeis{A079269}, the duplicate formulations \oeis{A174490} and
\oeis{A177431}, and the certified reserve \oeis{A235252} add no counts.

\subsection*{Evidence and limits}
Primary entries, relevant linked materials and targeted searches were
consulted on 20 October 2026; retrieval may use cached pages. The
retained OEIS Open comparison covers 100 costliest Appendix A claims
on 97 distinct A-numbers, with no selected overlap. It does not cover
that project's entire repository. The additional known-family and
linked-paper checks are bounded, not an exhaustive priority search.

No OEIS edit or submission was performed. The results here are
mathematical resolutions of selected assertions, separate from OEIS
editorial acceptance. Internal cross-reviews were performed by separate
agents in the same AI-assisted workflow; they are not external peer
review or proof-assistant formalizations.
\endgroup

\clearpage
\section{Catalogue of the 100 selected entries}
\label{round8:sec:r8catalogue}
\textbf{P}: selected claim proved.
\textbf{C}: support rule corrected and proved.
\textbf{R}: empirical classification refuted, with its lower bound proved.
Page references lead to the full argument or recurrence signature.

\small
\begin{longtable}{@{}r p{21mm} p{103mm} c r@{}}
\toprule No.&Entry&Selected assertion package&Result&Page\\
\midrule\endfirsthead
\toprule No.&Entry&Selected assertion package&Result&Page\\
\midrule\endhead
\midrule\multicolumn{5}{r}{\emph{Continued on the next page}}\\\endfoot
\bottomrule\endlastfoot

\multicolumn{5}{@{}l}{\textbf{Formal series, denominators and congruences}}\\[2pt]

1&\oeis{A381670}&Dyadic denominators in the compositional equation&P&\pageref{round8:r8entry:A381670}\\

2&\oeis{A352117}&The convergent central-binomial moment formula&P&\pageref{round8:r8entry:A352117}\\

3&\oeis{A174491}&The exact power-of-two denominator&P&\pageref{round8:r8entry:A174491}\\

4&\oeis{A079278}&The denominator recurrence after reduction&P&\pageref{round8:r8entry:A079278}\\

5&\oeis{A177430}&The integral exponential transform: all real parameters&P&\pageref{round8:r8entry:A177430}\\

6&\oeis{A383377}&Parity of the absolute-value self-referential series&P&\pageref{round8:r8entry:A383377}\\

7&\oeis{A382122}&Powers-of-two parity support for every real solution&P&\pageref{round8:r8entry:A382122}\\

8&\oeis{A365095}&The full two-parameter square divisibility&P&\pageref{round8:r8entry:A365095}\\

9&\oeis{A396102}&All positive coefficients are 1 modulo 3&P&\pageref{round8:r8entry:A396102}\\

10&\oeis{A378581}&Square support modulo 4 in the theta composition&P&\pageref{round8:r8entry:A378581}\\

11&\oeis{A380680}&Correct the support to twice a square, modulo 4&C&\pageref{round8:r8entry:A380680}\\

12&\oeis{A379764}&Signed square support modulo 4&P&\pageref{round8:r8entry:A379764}\\

13&\oeis{A380060}&Square support modulo 4 in the bilateral equation&P&\pageref{round8:r8entry:A380060}\\

\multicolumn{5}{@{}l}{\textbf{Arithmetic identities and one refutation}}\\[2pt]

14&\oeis{A023887}&The divisor-sum ratio for every positive exponent&P&\pageref{round8:r8entry:A023887}\\

15&\oeis{A092287}&The rectangular gcd-product valuation formula&P&\pageref{round8:r8entry:A092287}\\

16&\oeis{A119689}&Evenness and the extreme-prime-divisor criterion&P&\pageref{round8:r8entry:A119689}\\

17&\oeis{A143991}&Equality of reduced integral and series numerators&P&\pageref{round8:r8entry:A143991}\\

18&\oeis{A179504}&The doubled-argument divisor-power identity&P&\pageref{round8:r8entry:A179504}\\

19&\oeis{A193679}&The logarithmic Dirichlet series and convolution&P&\pageref{round8:r8entry:A193679}\\

20&\oeis{A226787}&The empirical inverse-nine rational generating function&P&\pageref{round8:r8entry:A226787}\\

21&\oeis{A326691}&The arithmetic-derivative gcd formula&P&\pageref{round8:r8entry:A326691}\\

22&\oeis{A345993}&Refute the value-three classification; prove its lower bound&R&\pageref{round8:r8entry:A345993}\\

23&\oeis{A367865}&The radical-average divisor identity&P&\pageref{round8:r8entry:A367865}\\

24&\oeis{A373561}&The fourfold gcd-indicator sum&P&\pageref{round8:r8entry:A373561}\\

25&\oeis{A383044}&The seven-ninths totient family&P&\pageref{round8:r8entry:A383044}\\

\multicolumn{5}{@{}l}{\textbf{Generating-function recurrence certificates}}\\[2pt]

26&\oeis{A014533}&Conjectured polynomial recurrence of order 2&P&\pageref{round8:r8entry:A014533}\\

27&\oeis{A024719}&Conjectured polynomial recurrence of order 2&P&\pageref{round8:r8entry:A024719}\\

28&\oeis{A026126}&Conjectured polynomial recurrence of order 4&P&\pageref{round8:r8entry:A026126}\\

29&\oeis{A098519}&Conjectured polynomial recurrence of order 2&P&\pageref{round8:r8entry:A098519}\\

30&\oeis{A103138}&Conjectured polynomial recurrence of order 3&P&\pageref{round8:r8entry:A103138}\\

31&\oeis{A112702}&Conjectured polynomial recurrence of order 2&P&\pageref{round8:r8entry:A112702}\\

32&\oeis{A113409}&Conjectured polynomial recurrence of order 4&P&\pageref{round8:r8entry:A113409}\\

33&\oeis{A114851}&Conjectured polynomial recurrence of order 6&P&\pageref{round8:r8entry:A114851}\\

34&\oeis{A162972}&Conjectured polynomial recurrence of order 7&P&\pageref{round8:r8entry:A162972}\\

35&\oeis{A168505}&Conjectured polynomial recurrence of order 4&P&\pageref{round8:r8entry:A168505}\\

36&\oeis{A174169}&Conjectured polynomial recurrence of order 4&P&\pageref{round8:r8entry:A174169}\\

37&\oeis{A182892}&Conjectured polynomial recurrence of order 8&P&\pageref{round8:r8entry:A182892}\\

38&\oeis{A220898}&Conjectured polynomial recurrence of order 3&P&\pageref{round8:r8entry:A220898}\\

39&\oeis{A225692}&Conjectured polynomial recurrence of order 3&P&\pageref{round8:r8entry:A225692}\\

40&\oeis{A240722}&Conjectured polynomial recurrence of order 2&P&\pageref{round8:r8entry:A240722}\\

41&\oeis{A246606}&Conjectured polynomial recurrence of order 2&P&\pageref{round8:r8entry:A246606}\\

42&\oeis{A268402}&Conjectured polynomial recurrence of order 4&P&\pageref{round8:r8entry:A268402}\\

43&\oeis{A278618}&Conjectured polynomial recurrence of order 3&P&\pageref{round8:r8entry:A278618}\\

44&\oeis{A302180}&Conjectured polynomial recurrence of order 3&P&\pageref{round8:r8entry:A302180}\\

45&\oeis{A302181}&Conjectured polynomial recurrence of order 2&P&\pageref{round8:r8entry:A302181}\\

46&\oeis{A302184}&Conjectured polynomial recurrence of order 3&P&\pageref{round8:r8entry:A302184}\\

47&\oeis{A302186}&Conjectured polynomial recurrence of order 4&P&\pageref{round8:r8entry:A302186}\\

48&\oeis{A308815}&Conjectured polynomial recurrence of order 3&P&\pageref{round8:r8entry:A308815}\\

49&\oeis{A344502}&Conjectured polynomial recurrence of order 4&P&\pageref{round8:r8entry:A344502}\\

50&\oeis{A379463}&Conjectured polynomial recurrence of order 5&P&\pageref{round8:r8entry:A379463}\\

51&\oeis{A379464}&Conjectured polynomial recurrence of order 7&P&\pageref{round8:r8entry:A379464}\\

\multicolumn{5}{@{}l}{\textbf{Finite-state array recurrence certificates}}\\[2pt]

52&\oeis{A233646}&Empirical array recurrence of order 9&P&\pageref{round8:r8entry:A233646}\\

53&\oeis{A233647}&Empirical array recurrence of order 38&P&\pageref{round8:r8entry:A233647}\\

54&\oeis{A233675}&Empirical array recurrence of order 4&P&\pageref{round8:r8entry:A233675}\\

55&\oeis{A233676}&Empirical array recurrence of order 10&P&\pageref{round8:r8entry:A233676}\\

56&\oeis{A233677}&Empirical array recurrence of order 17&P&\pageref{round8:r8entry:A233677}\\

57&\oeis{A233687}&Empirical array recurrence of order 41&P&\pageref{round8:r8entry:A233687}\\

58&\oeis{A233718}&Empirical array recurrence of order 16&P&\pageref{round8:r8entry:A233718}\\

59&\oeis{A233726}&Empirical array recurrence of order 4&P&\pageref{round8:r8entry:A233726}\\

60&\oeis{A233727}&Empirical array recurrence of order 12&P&\pageref{round8:r8entry:A233727}\\

61&\oeis{A233728}&Empirical array recurrence of order 28&P&\pageref{round8:r8entry:A233728}\\

62&\oeis{A233942}&Empirical array recurrence of order 4&P&\pageref{round8:r8entry:A233942}\\

63&\oeis{A233943}&Empirical array recurrence of order 10&P&\pageref{round8:r8entry:A233943}\\

64&\oeis{A233944}&Empirical array recurrence of order 25&P&\pageref{round8:r8entry:A233944}\\

65&\oeis{A234241}&Empirical array recurrence of order 8&P&\pageref{round8:r8entry:A234241}\\

66&\oeis{A234259}&Empirical array recurrence of order 3&P&\pageref{round8:r8entry:A234259}\\

67&\oeis{A234260}&Empirical array recurrence of order 6&P&\pageref{round8:r8entry:A234260}\\

68&\oeis{A234261}&Empirical array recurrence of order 6&P&\pageref{round8:r8entry:A234261}\\

69&\oeis{A234326}&Empirical array recurrence of order 4&P&\pageref{round8:r8entry:A234326}\\

70&\oeis{A234327}&Empirical array recurrence of order 11&P&\pageref{round8:r8entry:A234327}\\

71&\oeis{A234328}&Empirical array recurrence of order 48&P&\pageref{round8:r8entry:A234328}\\

72&\oeis{A234338}&Empirical array recurrence of order 13&P&\pageref{round8:r8entry:A234338}\\

73&\oeis{A234444}&Empirical array recurrence of order 4&P&\pageref{round8:r8entry:A234444}\\

74&\oeis{A234445}&Empirical array recurrence of order 17&P&\pageref{round8:r8entry:A234445}\\

75&\oeis{A234452}&Empirical array recurrence of order 6&P&\pageref{round8:r8entry:A234452}\\

76&\oeis{A234453}&Empirical array recurrence of order 26&P&\pageref{round8:r8entry:A234453}\\

77&\oeis{A234483}&Empirical array recurrence of order 3&P&\pageref{round8:r8entry:A234483}\\

78&\oeis{A234484}&Empirical array recurrence of order 5&P&\pageref{round8:r8entry:A234484}\\

79&\oeis{A234491}&Empirical array recurrence of order 5&P&\pageref{round8:r8entry:A234491}\\

80&\oeis{A234492}&Empirical array recurrence of order 23&P&\pageref{round8:r8entry:A234492}\\

81&\oeis{A234674}&Empirical array recurrence of order 4&P&\pageref{round8:r8entry:A234674}\\

82&\oeis{A234675}&Empirical array recurrence of order 18&P&\pageref{round8:r8entry:A234675}\\

83&\oeis{A234884}&Empirical array recurrence of order 27&P&\pageref{round8:r8entry:A234884}\\

84&\oeis{A234885}&Empirical array recurrence of order 46&P&\pageref{round8:r8entry:A234885}\\

85&\oeis{A234891}&Empirical array recurrence of order 10&P&\pageref{round8:r8entry:A234891}\\

86&\oeis{A234905}&Empirical array recurrence of order 12&P&\pageref{round8:r8entry:A234905}\\

87&\oeis{A234984}&Empirical array recurrence of order 8&P&\pageref{round8:r8entry:A234984}\\

88&\oeis{A234985}&Empirical array recurrence of order 28&P&\pageref{round8:r8entry:A234985}\\

89&\oeis{A235019}&Empirical array recurrence of order 6&P&\pageref{round8:r8entry:A235019}\\

90&\oeis{A235020}&Empirical array recurrence of order 21&P&\pageref{round8:r8entry:A235020}\\

91&\oeis{A235064}&Empirical array recurrence of order 10&P&\pageref{round8:r8entry:A235064}\\

92&\oeis{A235065}&Empirical array recurrence of order 40&P&\pageref{round8:r8entry:A235065}\\

93&\oeis{A235100}&Empirical array recurrence of order 6&P&\pageref{round8:r8entry:A235100}\\

94&\oeis{A235101}&Empirical array recurrence of order 19&P&\pageref{round8:r8entry:A235101}\\

95&\oeis{A235180}&Empirical array recurrence of order 18&P&\pageref{round8:r8entry:A235180}\\

96&\oeis{A235181}&Empirical array recurrence of order 26&P&\pageref{round8:r8entry:A235181}\\

97&\oeis{A235191}&Empirical array recurrence of order 6&P&\pageref{round8:r8entry:A235191}\\

98&\oeis{A235192}&Empirical array recurrence of order 23&P&\pageref{round8:r8entry:A235192}\\

99&\oeis{A235193}&Empirical array recurrence of order 39&P&\pageref{round8:r8entry:A235193}\\

100&\oeis{A235251}&Empirical array recurrence of order 10&P&\pageref{round8:r8entry:A235251}\\

\end{longtable}
\normalsize
\clearpage

\section{Formal series, divisibility, and corrected congruences}
\label{round8:r8sec:root}

This section selects one assertion package for each of thirteen distinct
A-numbers. The source statements are paraphrased, with their defining equations
retained. The proofs concern all indices in the stated ranges. Calculations in
\texttt{r8\_root\_verify.py} are supplementary checks, except where an exact
symbolic identity is explicitly identified. Classical inversion, Euler products,
and the six exponentials theorem are credited inputs, not new results.

\subsection{A381670: dyadic denominators under self-composition}
\label{round8:r8entry:A381670}
\textbf{Source.} \oeis{A381670}, Thomas Scheuerle, March 3, 2025, conjectures
that the reduced denominators of the coefficients of
\[
 A(x)=x+\tfrac12x^2-\tfrac1{16}x^4+\cdots,
 \qquad A(A(x))=x(1+A(x)),
\]
are powers of two. We include $1=2^0$ and use denominator $1$ for a zero
coefficient.

\begin{proposition}
There is a unique solution $A\in x+x^2\mathbb Q[[x]]$, and it belongs to
$x+x^2\mathbb Z[1/2][[x]]$. Consequently every term of A381670 is a power of two.
\end{proposition}
\begin{proof}
Write $A=x+\sum_{n\ge2}c_nx^n$. For $n\ge2$, the coefficient of $c_n$ in
$[x^n]A(A(x))$ is $2$: one occurrence comes from the outer linear term, and
one from $c_nA(x)^n$. All other contributions form an integer polynomial
$P_n(c_2,\ldots,c_{n-1})$. Thus
\[
 2c_n=c_{n-1}-P_n(c_2,\ldots,c_{n-1}),\qquad c_1=1.
\]
This determines the coefficients uniquely and inductively keeps them in
$\mathbb Z[1/2]$. Reduction of a dyadic rational cannot introduce an odd
prime in its denominator.
\end{proof}

\subsection{A352117: a convergent central-binomial moment formula}
\label{round8:r8entry:A352117}
\textbf{Source.} \oeis{A352117} is defined by the exponential generating
function $F(t)=(2-e^{2t})^{-1/2}$. The infinite-sum formula below is labelled
conjectural by Velin Yanev and Vaclav Kotesovec, July 21, 2026.
The separate Stieltjes-fraction assertion is not selected here.

\begin{proposition}
For every integer $n>0$,
\[
 a(n)=\sum_{k\ge1}2^{n-3k-1/2}\binom{2k}{k}k^n.
\]
The series converges absolutely.
\end{proposition}
\begin{proof}
The binomial expansion, initially for $|e^{2t}|<2$, gives
\[
 (2-e^{2t})^{-1/2}
 =\frac1{\sqrt2}\sum_{k\ge0}\binom{2k}{k}\frac{e^{2kt}}{8^k}.
\]
For any $0<\epsilon<\log(2)/2$ and $|t|\le\epsilon$, the $n$th
derivative of the $k$th summand has modulus at most a constant times
$k^n(e^{2\epsilon}/2)^k$, using $\binom{2k}{k}\le4^k$.
This is summable. Termwise differentiation is therefore valid near zero.
Evaluating $F^{(n)}(0)$ proves the formula; the $k=0$ summand vanishes when
$n>0$. The same estimate at $t=0$ proves absolute convergence. For $n=0$,
one must retain $k=0$ rather than apply the displayed $n>0$ formula.
\end{proof}

\subsection{A174491: the exact denominator, including odd-prime cancellation}
\label{round8:r8entry:A174491}
\textbf{Source.} \oeis{A174491}, Paul D. Hanna, March 25, 2010, defines the
denominator of $g_n=[x^n]G(x)$, where
\[
 G(x)=\exp\left(\sum_{j\ge1}\frac{x^j}{j2^{j^2}}\right).
\]
It conjectures the formula $2^{n^2+v_2(n)}$ for $n\ge1$.
The identical denominator assertion on \oeis{A174490} is not counted again.

\begin{lemma}\label{round8:r8lem:localized-euler}
Put $b_n=2^{-n^2}$ and
$c_n=n^{-1}\sum_{d\mid n}\mu(d)b_{n/d}$. Then $c_n\in\mathbb Z[1/2]$, and
\[
 G(x)=\prod_{n\ge1}(1-x^n)^{-c_n}\in\mathbb Z[1/2][[x]].
\]
\end{lemma}
\begin{proof}
Fix an odd prime $p$ and let $n=p^su$ with $s\ge1$ and $p\nmid u$.
For $p\nmid e$, Euler's theorem for units modulo $p^s$ implies
\[
 b_{p^se}-b_{p^{s-1}e}\in p^s\mathbb Z_{(p)}.
\]
Indeed, the difference of the exponents of $2$ is
$(p^2-1)p^{2s-2}e^2$, divisible by $p^{s-1}(p-1)$.
Pairing the divisors $d$ and $pd$ in the Mobius sum yields
\[
 \sum_{d\mid n}\mu(d)b_{n/d}
 =\sum_{d\mid u}\mu(d)
       (b_{p^su/d}-b_{p^{s-1}u/d})\in p^s\mathbb Z_{(p)}.
\]
Division by $n$ thus leaves no $p$ in the denominator. This holds for every
odd $p$, proving $c_n\in\mathbb Z[1/2]$.

For $r\in\mathbb Z[1/2]$, all binomial coefficients $\binom rj$ also lie
in that ring. To see this at an odd prime $p$, choose an integer $M$ congruent
to $r$ modulo $p^N$ with $N>v_p(j!)$; the falling products
show $\binom rj\in\mathbb Z_{(p)}$. Each factor of the product is therefore
dyadic coefficientwise. Finally, the logarithm of the product has coefficient
$n^{-1}\sum_{d\mid n}dc_d=b_n/n$ by Mobius inversion. All products and
logarithms are formal and coefficientwise finite.
\end{proof}

\begin{proposition}
The reduced denominator of $g_0$ is $1$. For every $n\ge1$, that of $g_n$
is exactly $2^{n^2+v_2(n)}$.
\end{proposition}
\begin{proof}
The exponential formula expands $g_n$ as
\[
 \sum_{\substack{m_1,m_2,\ldots\ge0\\\sum j m_j=n}}
 \prod_{j\ge1}\frac{1}{m_j!\,j^{m_j}\,2^{j^2m_j}}.
\]
The term $m_n=1$ has 2-adic valuation $-n^2-v_2(n)$.
For every other partition of $n\ge2$, convexity gives
$\sum j^2m_j\le(n-1)^2+1=n^2-2n+2$.
Write $r=\sum m_j$ and let $t$ be the number of nonzero $m_j$.
Since $v_2(j)\le j-1$ and $v_2(m!)\le m-1$ for $m\ge1$,
\[
 \sum m_jv_2(j)+\sum v_2(m_j!)\le(n-r)+(r-t)\le n-1.
\]
Thus every other term has denominator 2-adic exponent at most
$n^2-n+1<n^2+v_2(n)$. A sum with exactly one term of least 2-adic valuation
has that same valuation. Lemma~\ref{round8:r8lem:localized-euler} rules out all
odd denominator primes. The cases $n=0,1$ are immediate.
\end{proof}

\subsection{A079278: eliminate the numerators without losing reduction}
\label{round8:r8entry:A079278}
\textbf{Source.} \oeis{A079278} gives $b_1=1$,
$b_{n+1}=b_n+b_n/(1+b_n)$, and defines $q_n$ as the reduced denominator of
$b_n$. Leroy Quet's recurrence is explicitly conjectural.

\begin{proposition}
For $n\ge2$,
\[
 q_{n+1}=q_n^2+\frac{q_n^3}{q_{n-1}^2}-q_nq_{n-1}^2,
 \qquad q_1=1,\quad q_2=2.
\]
\end{proposition}
\begin{proof}
Write $b_n=p_n/q_n$ with positive coprime integers. Direct substitution gives
\[
 p_{n+1}=p_n(p_n+2q_n),\qquad q_{n+1}=q_n(p_n+q_n).
\]
These are already reduced: each of $p_n$ and $p_n+2q_n$ is coprime to
each of $q_n$ and $p_n+q_n$. In particular no implicit cancellation changes
the denominator recurrence. Now
\[
 p_{n-1}=q_n/q_{n-1}-q_{n-1},\qquad
 p_n=p_{n-1}^2+2p_{n-1}q_{n-1}
     =q_n^2/q_{n-1}^2-q_{n-1}^2.
\]
Substitution into $q_{n+1}=q_n(p_n+q_n)$ gives the result.
The quotients appearing here are integers by the first-order denominator
recurrence. The associated numerator recurrence on \oeis{A079269} is an
uncounted corollary, obtained by the same elimination.
\end{proof}

\subsection{A177430: which powers remain logarithmic derivatives?}
\label{round8:r8entry:A177430}
\textbf{Source.} \oeis{A177430}, Paul D. Hanna, September 6, 2010, is the
least strictly increasing sequence of squares $a_n$ for which
$\exp(\sum a_nx^n/n)$ has integer coefficients. Its selected conjecture says
that replacing $a_n$ by $a_n^{m/2}$ preserves this property precisely for
nonnegative even integers $m$. The duplicate statement at \oeis{A177431}
is uncounted. We prove the assertion even if $m$ is allowed to be real,
using positive real powers. The integer-parameter proof is elementary.

\begin{lemma}\label{round8:r8lem:ghost}
For a sequence $u_n$ of integers,
\[
 \exp\left(\sum_{n\ge1}u_nx^n/n\right)\in1+x\mathbb Z[[x]]
 \quad\Longleftrightarrow\quad
 u_n\equiv u_{n/p}\pmod{p^{v_p(n)}}
 \quad(p\mid n).
\]
Sequences satisfying these congruences are closed under nonnegative
integer powers taken term by term.
\end{lemma}
\begin{proof}
Every series in $1+x\mathbb Z[[x]]$ has a unique formal factorization
$\prod_{d\ge1}(1-x^d)^{-e_d}$ with $e_d\in\mathbb Z$: determine $e_d$
successively from the coefficient of $x^d$. Its logarithmic derivative has
$u_n=\sum_{d\mid n}de_d$. Subtracting the expression for $u_{n/p}$ proves
the congruences. Conversely $ne_n=\sum_{d\mid n}\mu(d)u_{n/d}$; pairing
divisors with and without $p$ shows this numerator divisible by every
$p^{v_p(n)}$, hence by $n$. Integer powers preserve the congruences because
$s^k-t^k$ is divisible by $s-t$; $k=0$ gives the constant sequence $1$.
\end{proof}

The defining greedy sequence exists at every step. If $n$ is fixed and its
earlier square roots are $r_j$, the Chinese remainder theorem supplies a
positive integer $r$ with $r\equiv r_{n/p}\pmod{p^{v_p(n)}}$ for every
prime $p\mid n$. Taking $r$ arbitrarily large makes $r^2$ an admissible
next square. Lemma~\ref{round8:r8lem:ghost} gives consistency. Its first four
terms are $1,9,16,25$: at $n=2,3,4$ the new square must be respectively
$1\pmod2$, $1\pmod3$, and $9\pmod4$, as well as exceed the previous one.

\begin{proposition}
For real $m$,
\[
 \exp\left(\sum_{n\ge1}a_n^{m/2}x^n/n\right)\in\mathbb Z[[x]]
 \quad\Longleftrightarrow\quad m\in\{0,2,4,6,\ldots\}.
\]
\end{proposition}
\begin{proof}
For $m=2k\ge0$, the assertion follows from
Lemma~\ref{round8:r8lem:ghost}. Conversely an integral unit series has an integral
logarithmic derivative $xF'/F$, so integrality forces $3^m,4^m,5^m$ to be
positive integers. If $m<0$, already $0<3^m<1$ is impossible. If $m$ is
rational, unique prime factorization applied to $3^m$ implies that $m$ is
a nonnegative integer.

For completeness, irrational real $m$ is excluded by the classical six
exponentials theorem. The numbers $1,m$ are linearly independent over
$\mathbb Q$, as are $\log3,\log4,\log5$. That theorem says that at least
one of the six exponentials of their pairwise products is transcendental.
Here those exponentials would be $3,4,5,3^m,4^m,5^m$, all integers, a
contradiction. We use the established theorem in the form stated and proved
in Samit Dasgupta, \emph{Ranks of matrices of logarithms of algebraic
numbers I}, Theorem 4.2 and Section 5.1 (2023),
\url{https://sites.math.duke.edu/~dasgupta/papers/Baker.pdf}.

Finally, for an odd integer $m$, the necessary ghost congruence at $n=4$
would require $5^m\equiv3^m\pmod4$. Its two sides are $1$ and $-1$,
respectively. This excludes every odd $m$ and completes the classification.
\end{proof}

\subsection{A383377: absolute values disappear modulo two}
\label{round8:r8entry:A383377}
\textbf{Source.} \oeis{A383377}, Paul D. Hanna, May 15, 2025, defines
\[
 A(x)=\sum_{k\ge0}x^k\operatorname{abs}(A(x)^{-k}),\qquad A(0)=1,
\]
where $\operatorname{abs}$ acts on each integer coefficient. It conjectures
that $a_n$ is even for $n>1$.

\begin{proposition}
The equation defines a unique integral series, and $A(x)\equiv1+x\pmod2$.
\end{proposition}
\begin{proof}
For positive degree $n$, the term $k=0$ contributes nothing, and each
term $k\ge1$ on the right uses coefficients of $A$ only through degree
$n-k<n$. This proves existence, uniqueness and integrality by recursion.
For an integer $c$, $|c|\equiv c\pmod2$. Reducing the equation in
$\mathbb F_2[[x]]$ therefore gives
\[
 A=\sum_{k\ge0}(x/A)^k=\frac{A}{A-x}.
\]
Both $A$ and $A-x$ are units. Cancellation yields $A-x=1$, which is
exactly the asserted parity pattern.
\end{proof}

\subsection{A382122: a Catalan target forces powers-of-two parity}
\label{round8:r8entry:A382122}
\textbf{Source.} \oeis{A382122}, Paul D. Hanna, March 16, 2025, uses
\[
 \sum_{k\ge0}x^k\operatorname{abs}(A(x)^{-k})=C(x),
 \qquad C(x)=1+xC(x)^2,
\]
and conjectures that its positive-index odd terms occur precisely at powers
of two. The conclusion below holds for every real-coefficient solution with
$A(0)=1$; it needs no uniqueness convention for the absolute-value branches.

\begin{proposition}
Every such solution is integral and satisfies
\[
 A(x)\equiv1+\sum_{j\ge0}x^{2^j}\pmod2.
\]
In particular the conjectured parity of A382122 holds.
\end{proposition}
\begin{proof}
Assume $a_1,\ldots,a_{n-1}$ are integers. In degree $n+1$ the summand
$x\operatorname{abs}(A^{-1})$ contributes $|I-a_n|$ for an integer $I$,
and all other terms use earlier coefficients only. The defining equation
therefore makes $|I-a_n|$ an integer, so $a_n$ is an integer too. This
starts with $|a_1|=1$ and proves integrality for every solution.

Now reduce modulo two to obtain $A/(A-x)=C$. Thus
$A(C-1)=Cx$. Using $C-1=xC^2$ and cancelling $x$ gives
$A=C^{-1}=1+xC$ in characteristic two. The unique Catalan series there
satisfies $C=1+xC^2$, hence by successive squaring
$C=\sum_{j\ge0}x^{2^j-1}$. Substitution proves the assertion. This argument
does not assert that the original absolute-value equation has only one
real or positive branch; the parity is identical on every branch.
\end{proof}

\subsection{A365095: a square-divisibility theorem for every integer parameter}
\label{round8:r8entry:A365095}
\textbf{Source.} \oeis{A365095}, Paul D. Hanna, September 3, 2023, is
defined by $A(0)=1$ and
\[
 [x^{n-1}]\bigl(A^{-1}+(n-1)xA\bigr)^n=0\qquad(n>1).
\]
Its selected conjecture replaces $n-1$ by $kn-1$ and asks for divisibility
by $n^2$ for every integer $k$. Other observations about related series on
the page are outside this selected assertion.

\begin{proposition}
The defining series is integral, and for $n\ge1$, $k\in\mathbb Z$,
\[
 n^2\ \bigm|\ [x^{n-1}]
 \frac{(1+(kn-1)xA(x)^2)^n}{A(x)^n}.
\]
\end{proposition}
\begin{proof}
For any $H\in\mathbb Z[[x]]$, differentiation gives
\[
 (n-1)[x^{n-1}]H^n=n[x^{n-2}]H'H^{n-1}.
\]
Since $n$ and $n-1$ are coprime, $n\mid[x^{n-1}]H^n$.
Suppose the coefficients of $A$ below degree $n-1$ are already integral,
and temporarily put $a_{n-1}=0$. Then
$H=A^{-1}+(n-1)xA$ is integral. Restoring $a_{n-1}$ changes the coefficient
under consideration by $-na_{n-1}$. The preceding divisibility determines
an integral $a_{n-1}$, proving existence and uniqueness inductively.

For $n>1$, put $H=A^{-1}+(n-1)xA$. The target expression is
$[x^{n-1}](H+n(k-1)xA)^n$. Subtract its known zero value $[x^{n-1}]H^n$.
The term with one copy of $n(k-1)xA$ has the factor
$\binom n1 n=n^2$, and each term with $j\ge2$ such copies has the factor
$n^j$, hence $n^2$. All remaining coefficients are integers. For $n=1$
the coefficient is $1$, divisible by $1$ as required.
\end{proof}

\subsection{A396102: a triple-composition congruence}
\label{round8:r8entry:A396102}
\textbf{Source.} \oeis{A396102}, Paul D. Hanna, June 6, 2026, has
$A=x+O(x^2)$ and
\[
 A^{\langle3\rangle}(x)=(1+x)A^{\langle2\rangle}(x),
\]
where angle brackets denote composition, not multiplication. It conjectures
$a_n\equiv1\pmod3$ for $n\ge1$.

\begin{proposition}
The solution is integral and $A(x)\equiv x/(1-x)\pmod3$.
\end{proposition}
\begin{proof}
For $A=x+\sum_{j\ge2}a_jx^j$, the coefficient of $a_n$ in the degree-$n$
coefficient of the $r$th iterate is $r$. Thus the defining difference has
coefficient $3-2=1$ on $a_n$; the extra factor $x$ uses only earlier
coefficients. This gives a unique solution by integer polynomial
recurrences over any commutative coefficient ring.

Let $B=x/(1-x)$. Direct composition gives
$B^{\langle r\rangle}=x/(1-rx)$. In characteristic three,
\[
 B^{\langle3\rangle}=x
     =(1+x)\frac{x}{1-2x}
     =(1+x)B^{\langle2\rangle}.
\]
Uniqueness over $\mathbb F_3$ now gives $A=B$ there.
The related but different iterate congruence in A196523 already has an
explicit proof by Amol Bhave dated October 13, 2026, and is excluded.
\end{proof}

\subsection{A378581: composition preserves the theta support modulo four}
\label{round8:r8entry:A378581}
\textbf{Source.} \oeis{A378581}, Paul D. Hanna, January 8, 2025, defines
\[
 A(xA(x))=1+2\sum_{j\ge1}x^{j^2},\qquad A(0)=1,
\]
and conjectures $a_n\equiv2\pmod4$ at positive squares and $0$ elsewhere.

\begin{proposition}
The defining equation has a unique integral solution and
$A(x)\equiv1+2\sum_{j\ge1}x^{j^2}\pmod4$.
\end{proposition}
\begin{proof}
In degree $n$, the coefficient $a_n$ occurs once, in
$a_nx^nA(x)^n$. Every other contribution uses $a_1,\ldots,a_{n-1}$.
This proves integral existence and uniqueness over every coefficient ring.
Modulo two the unique solution is $A=1$, since the right side is $1$.
Write $A=1+2B$ with $B\in x\mathbb Z[[x]]$. Since $xA\equiv x\pmod2$,
\[
 A(xA)=1+2B(xA)\equiv1+2B(x)=A(x)\pmod4.
\]
The defining equation gives the stated coefficientwise congruence.
The transform differs from A378580, selected in an earlier round; no
second count is assigned to that earlier A-number.
\end{proof}

\subsection{A380680: the support is twice a square, not a square}
\label{round8:r8entry:A380680}
\textbf{Source.} \oeis{A380680}, Paul D. Hanna, February 2025, defines
\[
 A(x)=\sum_{j\in\mathbb Z}x^{j^2}
      \left(\frac{2-x^j}{1-2x^j}\right)^j
\]
and labels the square-index residue pattern conjectural. Its own data
already have $a_1=4$ and $a_2=6$, contradicting that pattern.

\begin{proposition}
For every $n>0$,
\[
 a_n\equiv
 \begin{cases}2\pmod4,&n=2j^2\text{ for an integer }j\ge1,\\
 0\pmod4,&\text{otherwise}.
 \end{cases}
\]
\end{proposition}
\begin{proof}
Replacing $j$ by $-j$ leaves the summand invariant. Pairing the two signs
therefore gives the well-defined integral formal series
\[
 A(x)=1+2\sum_{j\ge1}x^{j^2}
       \left(\frac{2-x^j}{1-2x^j}\right)^j.
\]
For positive $j$, the factor in parentheses is $x^j$ modulo two, so its
$j$th power is $x^{j^2}$ modulo two. Hence
\[
 A(x)\equiv1+2\sum_{j\ge1}x^{2j^2}\pmod4.
\]
The positive exponents $2j^2$ are distinct. This proves the corrected
classification and refutes the original one. We count this once as a
corrected assertion, not once for the counterexample and again for its
replacement.
\end{proof}

\subsection{A379764: a bilateral series with an integral compositional inverse}
\label{round8:r8entry:A379764}
\textbf{Source.} \oeis{A379764}, Paul D. Hanna, January 13, 2025, has
$A=x+O(x^2)$ and $F(4A)=2\theta_4(x)$, where
\[
 F(z)=\sum_{j\in\mathbb Z}(z^j-z)^{j+1},\qquad
 \theta_4(x)=1+2\sum_{m\ge1}(-1)^m x^{m^2}.
\]
The two selected conjectures are the square support of odd coefficients
and its refinement $a_n\equiv(-1)^{n-1}\pmod4$ on squares, zero otherwise.

\begin{proposition}
Both conjectures hold for every positive index.
\end{proposition}
\begin{proof}
The terms $j=-1,0,1$ of $F$ are respectively $1,1-z,0$.
For $j\ge2$, the term is an integral polynomial of order $j+1$.
For $j=-r\le-2$ it is
\[
 z^{r(r-1)}(1-z^{r+1})^{-(r-1)},
\]
an integral power series of order $r(r-1)$. Thus the bilateral sum is
coefficientwise finite and
$F(z)=2-z+\sum_{j\ge2}f_jz^j$ with $f_j\in\mathbb Z$.
Dividing $F(4A)=2\theta_4(x)$ by $4$ after subtracting $2$ gives
\[
 A(x)=\sum_{m\ge1}(-1)^{m-1}x^{m^2}
       +\sum_{j\ge2}4^{j-1}f_jA(x)^j.
\]
This is a monic triangular equation, so it defines an integral solution.
All terms in its second sum vanish modulo four. At a square $n=m^2$,
the remaining coefficient is $(-1)^{m-1}=(-1)^{n-1}$, and at a
nonsquare it is zero. The parity claim follows as well.
\end{proof}

\subsection{A380060: the implicit bilateral theta-type series}
\label{round8:r8entry:A380060}
\textbf{Source.} \oeis{A380060}, Paul D. Hanna, February 6, 2025, defines
\[
 A(x)=\sum_{j\in\mathbb Z}x^{j^2}
 \left(\frac{A(x)-x^j}{1-x^jA(x)}\right)^j,\qquad A(0)=1.
\]
It conjectures that $a_n$ is $2$ modulo four at positive squares and $0$
otherwise. Unlike A380680, the numerator here contains the unknown series
$A(x)$ rather than the constant $2$.

\begin{proposition}
The equation has a unique integral formal solution, and
\[
 A(x)\equiv1+2\sum_{j\ge1}x^{j^2}\pmod4.
\]
\end{proposition}
\begin{proof}
Pair the terms with indices $j$ and $-j$. For $j>0$ the two rational
factors are reciprocals before the exponent is applied, so the summands
are equal. The equation becomes
\[
 A(x)=1+2\sum_{j\ge1}x^{j^2}
 \left(\frac{A(x)-x^j}{1-x^jA(x)}\right)^j.
\]
All denominators are units in $\mathbb Z[[x]]$, and the factor $x^{j^2}$
means the degree-$n$ coefficient on the right uses only earlier
coefficients of $A$. This proves unique integral existence. The same
equation first gives $A\equiv1\pmod2$. Consequently each quotient in
parentheses is $(1-x^j)/(1-x^j)=1$ modulo two. Multiplication by $2$
then proves the claimed reduction modulo four at every index.
\end{proof}

\subsection{Uncounted supplement: integrality for A300617}
\label{round8:r8uncounted:A300617}
\textbf{Source and scope.} \oeis{A300617}, Paul D. Hanna, March 10, 2018,
defines $A=\sum_{n\ge1}a_nx^n$ by
\[
 [x^n]e^{nA(x)}=n^2[x^{n-1}]e^{nA(x)}\qquad(n\ge1),
\]
and explicitly conjectures integrality. We give a proof here, but exclude
this A-number from the thirteen selected packages: its companion
\oeis{A300616} states the same integrality without a conjecture qualifier,
so the source status is inconsistent. No prior proof was located in the
consulted entries. This conservative exclusion does not depend on a claim
that the result was previously proved.
The stronger, separate claim $n(n+1)/2\mid a_n$ is \emph{not} settled
or counted by this package.

\begin{proposition}
The defining series is the unique solution in $x\mathbb Q[[x]]$ of
\[
 A(1-xA')^3=x+x^3A'',
\]
and all its coefficients are integers.
\end{proposition}
\begin{proof}
Let $W(t)$ be the inverse of $t=xe^{-A(x)}$, so
$W=t e^{A(W)}$. Lagrange inversion gives, for $n\ge1$,
\[
 [t^n]W=\frac1n[x^{n-1}]e^{nA(x)},\qquad
 [t^n]A(W)=\frac1n[x^{n-1}]A'e^{nA}
             =\frac1n[x^n]e^{nA}.
\]
The source equations are consequently equivalent to
$A(W)=\Theta_t^2W$, where $\Theta_t=t\,d/dt$. Since
$\Theta_t=\frac{x}{1-xA'}\frac{d}{dx}$ after setting $x=W(t)$,
\[
 \Theta_t W=\frac{x}{1-xA'},\qquad
 \Theta_t^2 W=\frac{x(1+x^2A'')}{(1-xA')^3}.
\]
This proves the differential equation, and every step is reversible for
formal series with $A(0)=0$.

In the equivalent form
\[
 A=x+x^3A''+3xAA'-3x^2A(A')^2+x^3A(A')^3,
\]
the coefficient of degree $n$ on the right uses only
$a_1,\ldots,a_{n-1}$, with integer coefficients. Thus $a_1=1$ and the
remaining coefficients are uniquely determined integers by induction.
This establishes precisely the stated integrality conjecture.
\end{proof}

\section{Arithmetic identities, exact counts, and one refutation}

The twelve packages in this section concern twelve previously uncounted
OEIS identifiers.  Each selected statement is explicitly conjectural or
empirical in the inspected primary entry.  Known formulas used below are
credited as inputs; the section does not claim that its elementary methods
or every consequence of those formulas are historically new.  In
particular, the square case of the gcd-product formula, and the
$m=2$ case of the divisor-sum ratio, already have published proofs.
Neither known special case is counted separately here.

\subsection{The divisor-sum ratio for arbitrary positive exponent: A023887}
\label{round8:r8entry:A023887}

Write $\sigma_m(n)=\sum_{d\mid n}d^m$, $\sigma=\sigma_1$, and
$R=\operatorname{rad}(n)=\prod_{p\mid n}p$.
The \href{https://oeis.org/A023887}{primary entry} labels Yanev's
August 24, 2017 identity conjectural for positive $m,n$.
The usual prime-power formula for $\sigma_m$ is already recorded there.
Fried's \href{https://oeis.org/A001157/a001157.pdf}{2025 paper},
Theorem 3, proves the special case $m=2$; that case is credited.

\begin{theorem}
For all integers $m,n\geq1$,
\[
 \sigma_m(n)=\frac{\sigma(n^mR^{m-1})}{\sigma(R^{m-1})}.
\]
\end{theorem}
\begin{proof}
If $n=\prod_p p^{e_p}$, where the displayed exponents are positive, then
$n^mR^{m-1}=\prod_p p^{m(e_p+1)-1}$.  Consequently the quotient on
the right is
\[
 \prod_{p\mid n}
 \frac{(p^{m(e_p+1)}-1)/(p-1)}{(p^m-1)/(p-1)}
 =\prod_{p\mid n}(1+p^m+\cdots+p^{me_p})
 =\sum_{d\mid n}d^m.
\]
For $n=1$, all products are empty and both sides are $1$.
This proves the unrestricted positive-integer parameter statement.
\end{proof}

\subsection{The rectangular gcd product: A092287}
\label{round8:r8entry:A092287}

The \href{https://oeis.org/A092287}{entry} records a proved formula
for the square product, and separately retains Kaydalov's March 11,
2019 conjecture for a rectangle.  Ma's
\href{https://arxiv.org/html/2509.13354v1}{2025 paper}, Section 2.2,
proves the square and equal-side higher-dimensional versions by threshold
counting.  We explicitly credit that argument and apply it to the two
independent side lengths in the selected statement.

\begin{theorem}
For nonnegative integers $n,m$ and every prime $p$, define
\[
 F(n,m)=\prod_{j=1}^{n}\prod_{k=1}^{m}\gcd(j,k).
\]
Then
\[
 v_p(F(n,m))=\sum_{r\geq1}
 \left\lfloor\frac n{p^r}\right\rfloor
 \left\lfloor\frac m{p^r}\right\rfloor.
\]
\end{theorem}
\begin{proof}
For positive $j,k$,
\[
 v_p(\gcd(j,k))
 =\min(v_p(j),v_p(k))
 =\sum_{r\geq1}\mathbf1_{p^r\mid j}\mathbf1_{p^r\mid k}.
\]
Sum this equality over the rectangle.  The number of pairs divisible
coordinatewise by $p^r$ is exactly
$\lfloor n/p^r\rfloor\lfloor m/p^r\rfloor$.
All sums are finite: the summands vanish once $p^r>\min(n,m)$.
When a side length is zero, the product is $1$ and both sides are zero.
\end{proof}

\subsection{Evenness of the extreme-prime-divisor sequence: A119689}
\label{round8:r8entry:A119689}

The members of \href{https://oeis.org/A119689}{A119689} are integers
$n\geq2$ whose smallest and largest prime divisors have prime sum.
Dale's September 22, 2016 evenness statement remains explicitly empirical.

\begin{theorem}
Every member is even.  More precisely, $n$ is a member if and only if
$n$ is even and its largest prime divisor $P$ satisfies that $P+2$ is
prime.
\end{theorem}
\begin{proof}
Let $p$ and $P$ be the smallest and largest prime divisors of $n$.
Their sum is at least $4$, so a prime value of $p+P$ must be odd.
Exactly one of $p,P$ is therefore the even prime $2$.  Since $p\leq P$,
this forces $p=2$, whence $n$ is even and $p+P=P+2$.
Conversely, for even $n$ the smallest prime divisor is $2$, so primality
of $P+2$ is precisely the defining condition.  This also handles powers
of two: they fail because $2+2=4$.
\end{proof}

\subsection{Equality of the integral and generating-function numerators: A143991}
\label{round8:r8entry:A143991}

Let
\[
 c_n=[x^n]\exp(1-\sqrt{1-x}),\qquad
 J_n=\frac e\pi\int_{-\infty}^{\infty}
             \frac{\cos t}{(1+t^2)^n}\,dt\quad(n\geq1).
\]
The numerator sequence of $c_n$ is
\href{https://oeis.org/A143991}{A143991}.  Spezia's March 12, 2025
conjecture says that $c_n$ and $J_n$ have the same reduced numerator.
The companion \href{https://oeis.org/A381845}{A381845} actually lists
the denominators of $J_n$; its cross-reference also marks the numerator
identification conjectural.  The Bessel numbers used in the proof are
the classical sequence \href{https://oeis.org/A001515}{A001515}.

\begin{theorem}
For every integer $n\geq1$,
\[
 J_n=2n c_n,
 \qquad 2n\mid\operatorname{den}(c_n).
\]
In particular $\operatorname{num}(J_n)=\operatorname{num}(c_n)$.
\end{theorem}
\begin{proof}
Put
\[
 B_m=\sum_{k=0}^{m}\frac{(m+k)!}{(m-k)!\,k!\,2^k}
 \quad(m\geq0).
\]
Every summand is an integer, since it equals
$\binom{m+k}{2k}(2k-1)!!$, with $(-1)!!=1$.
These are the Bessel polynomial values $y_m(1)$; their standard
recurrence is
\[
 B_0=1,\quad B_1=2,\qquad
 B_m=(2m-1)B_{m-1}+B_{m-2}\quad(m\geq2).
\]
For completeness, define the polynomial
\[
y_m(z)=\sum_{k=0}^{m}\frac{(m+k)!z^k}{(m-k)!\,k!\,2^k}.
\]
The recurrence follows by comparing coefficients:
the coefficient of $z^k$ in $(2m-1)z y_{m-1}(z)+y_{m-2}(z)$ equals
the coefficient in $y_m(z)$ after putting the two terms over the common
factorial denominator.  The boundary coefficients $k=0,m-1,m$ obey
the same identity directly.

First we identify the two rational numbers.  Set
$w=1-\sqrt{1-x}$, so $w=x/(2-w)$.  Lagrange inversion gives, for $n\geq1$,
\begin{align*}
 c_n
 &=\frac1n[u^{n-1}]e^u(2-u)^{-n}\\
 &=\frac1n\sum_{k=0}^{n-1}
       \frac{\binom{n+k-1}{k}}{2^{n+k}(n-1-k)!}
 =\frac{B_{n-1}}{2^n n!}.
\end{align*}
To evaluate the integral, integrate $e^{iz}/(1+z^2)^n$ over an upper
semicircle of radius $R>1$.  On the arc its modulus is at most
$(R^2-1)^{-n}$, so the arc integral is $O(R^{1-2n})$ and tends to zero.
The only enclosed pole is $z=i$.  Writing $z=i+u$, its residue is
\[
 e^{-1}[u^{n-1}]e^{iu}(2i+u)^{-n}
 =-i e^{-1}\sum_{k=0}^{n-1}
       \frac{\binom{n+k-1}{k}}{2^{n+k}(n-1-k)!}.
\]
The residue theorem, and oddness of the imaginary part on the real
axis, yield
\[
 J_n=2\sum_{k=0}^{n-1}
       \frac{\binom{n+k-1}{k}}{2^{n+k}(n-1-k)!}
 =\frac{B_{n-1}}{2^{n-1}(n-1)!}=2n c_n.
\]
Absolute convergence on the real line holds for every $n\geq1$.

It remains to control cancellation; the identity $J_n=2n c_n$ alone
would not prove equality of numerators.  We prove two elementary
congruence facts about $B_m$.

For an odd prime $p$, the finite factorial sum shows
\[
 B_{p-1}\equiv B_p\equiv1\pmod p.
\]
For $B_{p-1}$ every summand with $1\leq k\leq p-1$ has a factor $p$
in its numerator and none in its denominator.  For $B_p$, the same is
true when $1\leq k<p$.  Its final summand has valuation
$v_p((2p)!/p!)=1$, so it too vanishes modulo $p$.
Define $B_{-1}=1$; the recurrence is then valid also at $m=1$.
The coefficient $2m-1$ has period $p$ modulo $p$, and the state
$(B_{p-1},B_p)$ equals $(B_{-1},B_0)$ modulo $p$.  Induction in the
recurrence therefore proves
\[
 B_{m+p}\equiv B_m\pmod p\qquad(m\geq-1).
\]
In particular, if $p\mid n$, then $B_{n-1}\equiv1\pmod p$.

Modulo $4$, the recurrence has coefficient period $2$.  Its first
twelve values, indexed $0$ through $11$, are
\[
 1,2,3,1,2,3,3,2,1,3,2,1\pmod4.
\]
The next two values are again $1,2$.  Thus the initial pair and the
coefficient phase recur after twelve steps, which proves periodicity
for all subsequent indices by induction.  None of these twelve
residues is zero; hence $v_2(B_m)\leq1$ for every $m$.

Let $v=\operatorname{den}(c_n)$.  For every odd prime $p\mid n$,
the preceding congruence shows that cancellation removes no factor of
$p$ from $2^n n!$, so $v_p(v)=v_p(n!)\geq v_p(n)$.
For $n\geq2$, the modulo-$4$ result gives
\[
 v_2(v)\geq n+v_2(n!)-1
       \geq 1+v_2(n)=v_2(2n).
\]
For $n=1$, $c_1=1/2$ gives the same divisibility directly.
Thus $2n\mid v$.  If $c_n=u/v$ in lowest terms, then
$J_n=u/(v/(2n))$ is also in lowest terms, since its denominator
divides $v$.  Its numerator is therefore exactly $u$.
\end{proof}

\subsection{The doubled-argument divisor-power formula: A179504}
\label{round8:r8entry:A179504}

For $n>0$, \href{https://oeis.org/A179504}{A179504} is
$\sigma_n(2n)$.  Yanev's July 8, 2017 formula is still labelled
conjectural.  The diagonal divisor sum $\sigma_n(n)$ is A023887.

\begin{theorem}
For every integer $n\geq1$, with $e=v_2(n)$,
\[
 \sigma_n(2n)=\sigma_n(n)
 \left(2^n+\frac{2^n-1}{2^{n(e+1)}-1}\right).
\]
\end{theorem}
\begin{proof}
Write $n=2^e u$ with $u$ odd, and put $q=2^n$ and
$S_e=1+q+\cdots+q^e$.  Multiplicativity in the argument, with the
exponent $n$ fixed, gives
$\sigma_n(n)=\sigma_n(u)S_e$ and
$\sigma_n(2n)=\sigma_n(u)S_{e+1}$.
Since $S_{e+1}=qS_e+1$ and $S_e=(q^{e+1}-1)/(q-1)$, their ratio is
$q+(q-1)/(q^{e+1}-1)$, as claimed.  The sequence's separately specified
$n=0$ term is outside the source conjecture.
\end{proof}

\subsection{The logarithmic Dirichlet series of A193679}
\label{round8:r8entry:A193679}

The known cyclotomic-discriminant formula in
\href{https://oeis.org/A193679}{A193679} is
\[
 a(1)=1,\qquad a(n)=\prod_{p\mid n}p^{\phi(n)/(p-1)}.
\]
Irwin's July 14, 2018 conjecture concerns the Dirichlet series of its
logarithm and the equivalent von Mangoldt convolution.
We use the displayed classical product as an input, not as a new
discriminant formula.

\begin{theorem}
For every $n\geq1$,
\[
 \log a(n)=\sum_{d\mid n}\Lambda(d)\phi(n/d).
\]
For $\Re s>2$, both sides below converge absolutely and
\[
 \sum_{n\geq1}\frac{\log a(n)}{n^s}
 =-\frac{\zeta(s-1)\zeta'(s)}{\zeta(s)^2}.
\]
\end{theorem}
\begin{proof}
Fix a prime $p^e\Vert n$ and write $n=p^e u$, $(p,u)=1$.
The coefficient of $\log p$ in the asserted convolution is
\[
 \sum_{j=1}^{e}\phi(n/p^j)
 =\phi(u)\sum_{h=0}^{e-1}\phi(p^h)
 =\phi(u)p^{e-1}
 =\frac{\phi(n)}{p-1}.
\]
The middle identity follows by telescoping
$\phi(p^h)=p^h-p^{h-1}$ for $h\geq1$, together with $\phi(1)=1$.
These are exactly the coefficients of $\log p$ in the known product
for $a(n)$.  For $n=1$ the convolution is zero, as is $\log a(1)$.

For $\sigma=\Re s>2$, the series
$\sum\phi(m)m^{-s}$ and $\sum\Lambda(d)d^{-s}$ are absolutely
convergent, since $\phi(m)\leq m$ and $\Lambda(d)\leq\log d$.
The convolution may therefore be summed and factored absolutely.
The classical Euler-product identities give
\[
 \sum_{m\geq1}\frac{\phi(m)}{m^s}=\frac{\zeta(s-1)}{\zeta(s)},
 \qquad
 \sum_{d\geq1}\frac{\Lambda(d)}{d^s}=-\frac{\zeta'(s)}{\zeta(s)}.
\]
For the first identity one may use $\phi=\mu*\mathrm{id}$; the
second follows by logarithmically differentiating the absolutely
convergent Euler product.  Their product proves the stated formula
on the specified half-plane.
\end{proof}

\subsection{The empirical inverse-nine generating function: A226787}
\label{round8:r8entry:A226787}

Let $a(1)=0$, let $a(n)=0$ when $3\mid n$, and otherwise let $a(n)$
be the least nonnegative inverse of $9$ modulo $n$.
Barker's June 20, 2013 rational generating function in
\href{https://oeis.org/A226787/internal}{A226787} remains empirical.

\begin{theorem}
With $A(x)=\sum_{n\geq1}a(n)x^n$,
\[
 A(x)=\frac{x^2+x^4+4x^5+4x^7+x^8+9x^{10}
 +3x^{11}+x^{13}+3x^{14}+x^{16}-x^{19}}{(1-x^9)^2}.
\]
This is exactly the posted rational function after expanding its numerator.
\end{theorem}
\begin{proof}
For $n=9q+r$ with $r\in\{1,2,4,5,7,8\}$, choose $t_r$ with
$1\leq t_r\leq8$ and $rt_r\equiv-1\pmod9$.
The complete table is
\[
\begin{array}{c|rrrrrr}
r&1&2&4&5&7&8\\ \hline
t_r&8&4&2&7&5&1\\
c_r=(rt_r+1)/9&1&1&1&4&4&1
\end{array}
\]
For $n>1$, the integer $(t_r n+1)/9=t_rq+c_r$ lies strictly between
$0$ and $n$, and multiplying it by $9$ gives $1$ modulo $n$.
It is therefore $a(n)$.  The formula at $q=0,r=1$ instead gives $1$,
whereas the prescribed value $a(1)$ is zero; this requires subtracting $x$.
All three residue classes divisible by $3$ contribute zero.  Hence
\[
 A(x)=-x+\sum_{r\in\{1,2,4,5,7,8\}}x^r
 \left(\frac{c_r}{1-x^9}+\frac{t_rx^9}{(1-x^9)^2}\right).
\]
Putting this expression over the common denominator and using the
six displayed table entries gives the stated numerator identically.
The calculation is an identity of formal power series, not a fit to
a finite prefix.
\end{proof}

\subsection{The arithmetic-derivative gcd conjecture: A326691}
\label{round8:r8entry:A326691}

In \href{https://oeis.org/A326691}{A326691}, define
\[
 a(n)=\frac{n}{\operatorname{den}\left(\sum_{p\mid n}\frac1p-\frac1n\right)},
\]
where the denominator is positive and reduced.  Karttunen's March 15,
2021 conjecture identifies this with a gcd involving the arithmetic
derivative $D(n)=\sum_{p^e\Vert n}e n/p$, with $D(1)=0$.

\begin{theorem}
For every $n\geq1$,
\[
 a(n)=\gcd\bigl(n,1+(n-1)D(n)\bigr)=\gcd(n,D(n)-1).
\]
\end{theorem}
\begin{proof}
Set $S(n)=\sum_{p\mid n}n/p$.  Reduction of $(S(n)-1)/n$ gives
$a(n)=\gcd(n,S(n)-1)$, including $n=1$.
Fix $p^e\Vert n$.  If $e\geq2$, every summand of both $S(n)$ and
$D(n)$ is divisible by $p$.  Neither $S(n)-1$ nor $D(n)-1$ is then
divisible by $p$, so their gcds with $n$ both have $p$-valuation zero.
If $e=1$, the summand corresponding to $p$ has coefficient $1$ in
both sums, while all summands belonging to other primes are divisible
by $p$.  Therefore $S(n)-1\equiv D(n)-1\pmod p$.
Since the valuation in either gcd is at most $e=1$, the two valuations
again agree.  Comparing every prime divisor of $n$ proves the second
expression for $a(n)$.  Finally
$1+(n-1)D(n)\equiv1-D(n)\pmod n$, which gives the first expression.
\end{proof}

\subsection{A lower bound survives, but the proposed value-three classification fails: A345993}
\label{round8:r8entry:A345993}

Let $m(n)$ be the least positive integer $m$ with $n\mid m(m+1)$,
and put $a(n)=\gcd(n,m(n)+1)$, as in
\href{https://oeis.org/A345993}{A345993}.
The empirical comment proposes both $a(n)\geq3$ for $n\geq3$ and a
classification of $a(n)=3$ using only $n=3\cdot2^{\text{odd}}$ and
$n=3p$ with primes $p\equiv-1\pmod6$.

\begin{theorem}
For all $n\geq3$, $a(n)\geq3$.
However, the proposed classification is false: $m(375)=125$ and
$a(375)=3$.
More generally, if $p\equiv2\pmod3$ is prime and $e$ is any positive
odd integer, then
\[
 m(3p^e)=p^e,\qquad a(3p^e)=3.
\]
Thus $3\cdot5^{2r+1}$, $r\geq1$, is an infinite family omitted by
the proposed classification.
\end{theorem}
\begin{proof}
For $n\geq2$, the value $m=n-1$ is admissible, so $m(n)<n$.
Write $m=m(n)$ and $g=\gcd(n,m+1)$.  Since consecutive integers are
coprime, every prime power dividing $n$ divides exactly one of $m,m+1$;
in particular $n/g$ divides $m$.
If $g=1$, this would force $n\mid m$, impossible for $0<m<n$.
If $g=2$, then $m$ is a positive multiple of $n/2$ smaller than $n$,
so $m=n/2$.  Because $m+1$ is even, $m$ is odd.  For $n\geq3$ this
odd $m$ is at least $3$, and $m-1$ is a smaller positive admissible
integer: $2m=n$ divides $(m-1)m$ because $m-1$ is even.
This contradicts minimality.  Hence $g\geq3$.

For the family, put $q=p^e$.  Then $q\equiv2\pmod3$.
A root of $u(u+1)\equiv0\pmod q$ is either $u\equiv0$ or
$u\equiv-1\pmod q$, since $q$ is a prime power.
Combining either possibility with $u\equiv0$ or $-1\pmod3$ gives
exactly four roots modulo $3q$, namely
\[
 0,\qquad q,\qquad 2q-1,\qquad3q-1.
\]
They are distinct, and the least positive one is $q$.  Consequently
$m(3q)=q$ and $a(3q)=\gcd(3q,q+1)=3$.
Taking $p=5,e=3$ proves the explicit counterexample.  The integer
$125$ is neither a power of $2$ nor a prime, so $375$ belongs to
neither proposed family.  This resolves the two parts as a true lower
bound and a false classification; it does not assert a complete
replacement classification of all value-three inputs.
\end{proof}

\subsection{The radical average identity: A367865}
\label{round8:r8entry:A367865}

The definition of \href{https://oeis.org/A367865}{A367865} is
$a(n)=\sum_{d\mid n}d\phi(d)\mu(d)^2$.
Yanev's September 14, 2025 conjecture expresses it as a radical average.

\begin{theorem}
For every integer $n\geq1$,
\[
 a(n)=\frac{\operatorname{rad}(n)}{n}
       \sum_{k=1}^{n}\frac{\operatorname{rad}(nk)}{\operatorname{rad}(k)}.
\]
\end{theorem}
\begin{proof}
Let $R=\operatorname{rad}(n)$.  Prime by prime,
\[
 \frac{\operatorname{rad}(nk)}{\operatorname{rad}(k)}
 =\frac R{\gcd(R,k)}.
\]
This is periodic in $k$ with period $R$, and $R\mid n$.
The right side of the proposed identity therefore equals
\[
 \sum_{k=1}^{R}\frac R{\gcd(R,k)}.
\]
For a divisor $d\mid R$, exactly $\phi(d)$ of these $k$ satisfy
$\gcd(R,k)=R/d$.  Grouping them gives $\sum_{d\mid R}d\phi(d)$.
The divisors of $R$ are exactly the squarefree divisors of $n$, so
this is $\sum_{d\mid n}d\phi(d)\mu(d)^2=a(n)$.
The same argument includes $n=1$.
\end{proof}

\subsection{The fourfold gcd-indicator sum: A373561}
\label{round8:r8entry:A373561}

The polynomial sequence \href{https://oeis.org/A373561}{A373561}
is $a(n)=n^3(n+1)(2n+1)/6$ for $n\geq0$.
The entry conjectures the following weighted fourfold expression,
where $f(x,y,z)=x^2+y^2-z^2$.

\begin{theorem}
For every integer $n\geq0$,
\[
 \sum_{k=1}^{n}\sum_{z=1}^{n}\sum_{y=1}^{n}\sum_{x=1}^{n}
 \mathbf1_{\gcd(f(x,y,z),n)=k}\,f(x,y,z)
 =\frac{n^3(n+1)(2n+1)}6.
\]
\end{theorem}
\begin{proof}
For $n\geq1$ and every integer $f$, including zero and negative values,
$\gcd(f,n)$ is one positive integer in $\{1,\ldots,n\}$.
Thus the sum of its $n$ equality indicators is exactly $1$.
Interchanging the finite sums removes the outer indicator sum and leaves
\[
 \sum_{x,y,z=1}^{n}(x^2+y^2-z^2)
 =n^2\left(\sum_{j=1}^{n}j^2+\sum_{j=1}^{n}j^2
                  -\sum_{j=1}^{n}j^2\right)
 =\frac{n^3(n+1)(2n+1)}6.
\]
For $n=0$, the sums are empty and the polynomial is zero.
\end{proof}

\subsection{The empirical seven-ninths totient family: A383044}
\label{round8:r8entry:A383044}

The selected April 25, 2025 empirical statement in
\href{https://oeis.org/A383044}{A383044} concerns the ratio $7/9$,
not the entry's main ratio-one classification.  It proposes the family
$m=18\,\oeis{A003586}(n)$, where A003586 consists of $2^u3^v$ with
$u,v\geq0$.
Steinerberger's linked
\href{https://arxiv.org/html/2504.08023v1}{2025 paper} studies the
ratio-one equation and its six known families; that analysis is
credited and is outside this selected assertion.

\begin{theorem}
For all nonnegative integers $u,v$, if $m=18\cdot2^u3^v$, then
\[
 \phi(m)+\phi(m+\phi(m))=\frac79m.
\]
In particular this ratio has infinitely many positive integer solutions.
\end{theorem}
\begin{proof}
Here $m=2^{u+1}3^{v+2}$, so the ordinary prime-power formula gives
$\phi(m)=m(1-1/2)(1-1/3)=m/3$.
Consequently
$m+\phi(m)=4m/3=2^{u+3}3^{v+1}$ still has exactly the prime divisors
$2,3$.  Its totient is therefore $(4m/3)/3=4m/9$.
Adding gives $7m/9$.  Taking $v=0$ and letting $u$ increase provides
infinitely many distinct solutions.
\end{proof}

\section{Generating functions and exact recurrence certificates}
\label{round8:r8:gf}
This section proves 26 distinct source-marked polynomial recurrences.  Each
selected target is explicitly labelled a conjecture in its consulted OEIS entry.
The generating functions, finite sums and combinatorial definitions already in
those entries are credited inputs.  Auxiliary source formulas also need correction:
the missing alternating sign in the title of A308815, an incorrect generating
function printed for A113409, and older substitution and finite-sum expressions
for A174169.  These corrections do not add any counted items.

The selection excludes the 708 previously selected IDs, the inspected Fischer
2020 and Ekhad--Yang 2017 proof families, and the limited 97-ID Appendix A
baseline from OEIS Open.  Same-entry established recurrences and linked proof
papers were checked for stale labels.  In particular, A101601 and A268149 are
uncounted because of previously established hypergeometric evaluations;
A278992 is uncounted because its linked paper proves a recurrence system for
the same counts.  The audit is bounded, and does not assert worldwide priority,
official OEIS closure, or novelty of any credited input.  A220898 and the
previously treated A220899 are related by a finite difference, not a shift or
scalar multiple; their selected recurrence assertions are distinct.

\subsection*{Exact coefficient identities}
Set $a_n=0$ below the actual source offset $s$, and let
$F(x)=\sum_{n\geq s}a_nx^n$.  In every ordinary recurrence below the convention is
\begin{equation}
 R_n=\sum_{j=0}^r p_j(n)a_{n-j},\qquad
 L=\sum_{j=0}^r x^j p_j(\theta+j),\qquad
 \theta=x\frac{d}{dx}.
 \label{round8:r8:gf-operator}
\end{equation}
Termwise differentiation gives $[x^n]LF=R_n$.  Thus an exact identity
$LF=P(x)$ proves the recurrence for $n>\deg P$.  We also require $n\geq s+r$
so all displayed sequence indices belong to their source domain.  No division
by $p_0(n)$ is required, including at its integer zeros.

All the algebraic certificates below are finite rational identities.  For
$x=\xi(t)$, put $M=\xi/\xi'$.  Reduce the displayed seed modulo $y^2=D(t)$
to $F=R_0(t)+R_1(t)y$, where $y(0)=1$.  Set $D_0=1$, $D_1=D$, and define
\begin{equation}
 H_{i,0}=R_i,\quad
 H_{i,m+1}=M\left(H_{i,m}'+\frac{D_i'}{2D_i}H_{i,m}\right),\quad
 c_m=\sum_{j=0}^r \xi^j[z^m]p_j(z+j).
 \label{round8:r8:gf-rational}
\end{equation}
The certificate is exactly
$\sum_m c_mH_{0,m}=P(\xi)$ and $\sum_m c_mH_{1,m}=0$.
Each entry supplies every polynomial $p_j$, the seed, its radicand and
parameter, and $P$.  Formula \eqref{round8:r8:gf-rational} therefore specifies a finite,
independently reproducible proof using rational differentiation and polynomial
cancellation.  Its two residual numerators are zero identically, not merely on
a finite set of coefficients.  The accompanying verifier repeats these exact
operations and checks the source transcription.  For an ordinary unparametrized
series use $t=x$, $\xi=x$ and $M=x$.

If $t=x\phi(t)$, $\phi(0)=1$, formal inversion gives a unique $t(0)=0$.
For any formal series $B$ we use the diagonal form of Lagrange inversion,
\begin{equation}
 \sum_{n\geq0}x^n[u^n]B(u)\phi(u)^n
 =\frac{B(t)}{1-t\phi'(t)/\phi(t)}.
 \label{round8:r8:gf-diagonal}
\end{equation}
Indeed, substituting $x=u/\phi(u)$ in the coefficient residue on the right
cancels its denominator by the Jacobian and leaves
$\operatorname{Res}_u B(u)\phi(u)^n u^{-n-1}du$.  The formula includes $n=0$.
We also use $[x^n]H(t)=n^{-1}[u^{n-1}]H'(u)\phi(u)^n$ for $n\geq1$,
which follows from the same substitution and formal integration by parts.

For an exponential or even exponential series
$F(z)=\sum a_nz^{kn}/(kn)!$, where $k=1$ or $2$, set
\begin{equation}
 q_j(n)=p_j(n)\prod_{h=kj}^{kr-1}(kn-h),\qquad
 \Theta=\frac zk\frac{d}{dz},\qquad
 \mathcal L=\sum_{j=0}^r z^{kj}q_j(\Theta+j).
 \label{round8:r8:gf-egf}
\end{equation}
Then $[z^{kn}]\mathcal LF=R_n/(kn-kr)!$ for $n\geq r$.
This follows by cancellation of the factorials term by term.
A differential basis $b_i$ with $b_j'=\sum_i A_{ji}b_i$ is handled by
\[
 (R_i)_i\longmapsto
 \left(\frac zk\left(R_i'+\sum_jR_jA_{ji}\right)\right)_i.
\]
Repeated use of this rational matrix rule, with the polynomials $q_j$, is
again a finite exact certificate.

\subsection*{Walk components for the Bessel certificates}
Define the formal series
\[
 I_0(2z)=\sum_{m\geq0}\frac{z^{2m}}{(m!)^2},\qquad
 I_1(2z)=\sum_{m\geq0}\frac{z^{2m+1}}{m!(m+1)!}.
\]
Differentiating these sums gives $I_0'=2I_1$ and
$I_1'=2I_0-I_1/z$.  In one dimension, excursions constrained to stay
nonnegative have EGF $I_1(2z)/z$, unrestricted bridges have EGF $I_0(2z)$,
and nonnegative meanders have EGF $I_0(2z)+I_1(2z)$.  The coefficient
formulas are respectively $C_m$ at length $2m$, $\binom{2m}{m}$ at
length $2m$, and $\binom{j}{\lfloor j/2\rfloor}$ at length $j$.
The last formula follows from reflection of paths first reaching $-1$;
the first is the same reflection argument with final height zero.
Interleaving independent components corresponds to multiplication of EGFs,
since choosing their positions gives the multinomial coefficient.
One or two free directions contribute $e^z$ or $e^{2z}$.
These are the components used in Dershowitz's \emph{Touchard's Drunkard}
(2017), Sections 4--5 and Table 3.

For a fixed degree $d$ use the basis
$b_j=e^{cz}I_0(2z)^{d-j}I_1(2z)^j$, $0\leq j\leq d$.  Its derivative is
\begin{equation}
 b_j'=(c-j/z)b_j+2(d-j)b_{j+1}+2j b_{j-1},
 \label{round8:r8:gf-bessel}
\end{equation}
with out-of-range basis terms omitted.  This explicitly supplies every
matrix entry used by the four walk certificates below.



\subsection{A014533: the selected recurrence}

\label{round8:r8entry:A014533}

\paragraph{Source and assertion.}

\href{https://oeis.org/A014533}{A014533}, consulted on 20 October 2026, explicitly marks the recurrence as a conjecture attributed to R. J. Mathar, Feb 25 2015.

The source offset is $1$.  With $R_n=\sum_{j=0}^{2}p_j(n)a_{n-j}$, the exact source coefficients are:

\[\begin{aligned}
p_{0}(n)&=- \left(n - 1\right) \left(n + 7\right)\\
p_{1}(n)&=\left(n + 3\right) \left(2 n + 5\right)\\
p_{2}(n)&=3 \left(n + 2\right) \left(n + 3\right)\\
\end{aligned}\]

We prove $R_n=0$ for every integer $n\geq 3$.

\paragraph{Definition and generating-function derivation.}

The fourth trinomial column is
$a_n=[u^{n-1}](1+u+u^2)^{n+3}$ for $n\geq1$.
Equivalently it counts unrestricted Motzkin walks of length $n+3$ ending
at height four.  The bridge GF is $1/y$, $y=\sqrt{1-2x-3x^2}$, and the
GF for endpoint $h\geq0$ is $(xM)^h/y$, obtained by solving the walk
kernel recurrence with the branch $xM=O(x)$.
At $h=4$, shifting length down by three gives
$F=xM^4/y$, with no artificial constant term.
This agrees with the credited source GF and fixes the offset explicitly.

\paragraph{Exact certificate.}

\[F(x)=\frac{\left(- x - y + 1\right)^{4}}{16 x^{7} y}.\]

Here $y^2=- 3 x^{2} - 2 x + 1$, with $y(0)=1$.

Substitution in \eqref{round8:r8:gf-rational}, followed by exact cancellation, gives

\[LF=P(x)=0.\]

The two rational residuals are $P(\xi)$ and $0$.  Their coefficient interpretation proves the recurrence for all $n\geq 3$; this is at least the source offset plus the recurrence order.

\paragraph{Independent consistency check.} The companion validation computes the defining finite sum, combinatorial dynamic program, or nonlinear grammar through index 47, without using the conjectured recurrence.  All listed source coefficients match.  These checks audit the transcription and branch; the exact identities above supply the proof.


\subsection{A024719: the selected recurrence}

\label{round8:r8entry:A024719}

\paragraph{Source and assertion.}

\href{https://oeis.org/A024719}{A024719}, consulted on 20 October 2026, explicitly marks the recurrence as a conjecture attributed to R. J. Mathar, Sep 29 2012.

The source offset is $0$.  With $R_n=\sum_{j=0}^{2}p_j(n)a_{n-j}$, the exact source coefficients are:

\[\begin{aligned}
p_{0}(n)&=2 n \left(2 n - 1\right)\\
p_{1}(n)&=- 31 n^{2} + 29 n - 6\\
p_{2}(n)&=3 \left(3 n - 2\right) \left(3 n - 1\right)\\
\end{aligned}\]

We prove $R_n=0$ for every integer $n\geq 2$.

\paragraph{Definition and generating-function derivation.}

The source defines
\[
 a_n=\frac{2+\sum_{k=0}^n\binom{3k}{k}}{3}.
\]
The identity $\binom{3n}{n}=[u^n](1-u)^{-2n-1}$ lets us apply
the diagonal formula with $\phi=(1-t)^{-2}$, giving
\[
 \sum_{n\geq0}\binom{3n}{n}x^n=\frac1{1-3t},
 \qquad x=t(1-t)^2.
\]
Partial summation therefore gives
\[
 F=\frac{2+1/(1-3t)}{3(1-x)}
  =\frac{1-2t}{(1-3t)(1-x)},
\]
which is van Hoeij's November 9, 2011 parameterization.
Greubel's later first-difference formula is the same defining partial
sum in inhomogeneous form; no novelty is claimed for it.

\paragraph{Exact certificate.}

\[F(\xi(t))=\frac{1 - 2 t}{\left(3 t - 1\right) \left(t^{3} - 2 t^{2} + t - 1\right)}.\]

\[\xi(t)=t \left(1 - t\right)^{2},\qquad \phi(t)=\frac{1}{\left(1 - t\right)^{2}},\qquad y^2=1.\]

The formal parameter satisfies $t=x\phi(t)$; every square root uses its branch with constant term one.

Substitution in \eqref{round8:r8:gf-rational}, followed by exact cancellation, gives

\[LF=P(x)=- 4 x.\]

The two rational residuals are $P(\xi)$ and $0$.  Their coefficient interpretation proves the recurrence for all $n\geq 2$; this is at least the source offset plus the recurrence order.

\paragraph{Independent consistency check.} The companion validation computes the defining finite sum, combinatorial dynamic program, or nonlinear grammar through index 47, without using the conjectured recurrence.  All listed source coefficients match.  These checks audit the transcription and branch; the exact identities above supply the proof.


\subsection{A026126: the selected recurrence}

\label{round8:r8entry:A026126}

\paragraph{Source and assertion.}

\href{https://oeis.org/A026126}{A026126}, consulted on 20 October 2026, explicitly marks the recurrence as a conjecture attributed to R. J. Mathar, Jun 10 2013.

The source offset is $4$.  With $R_n=\sum_{j=0}^{4}p_j(n)a_{n-j}$, the exact source coefficients are:

\[\begin{aligned}
p_{0}(n)&=- \left(n - 4\right) \left(n + 8\right)\\
p_{1}(n)&=4 n^{2} + 5 n - 73\\
p_{2}(n)&=- 2 n^{2} + 13 n + 38\\
p_{3}(n)&=- \left(n - 3\right) \left(4 n + 5\right)\\
p_{4}(n)&=3 \left(n - 4\right) \left(n - 3\right)\\
\end{aligned}\]

We prove $R_n=0$ for every integer $n\geq 8$.

\paragraph{Definition and generating-function derivation.}

Let $M=1+xM+x^2M^2$ be the Motzkin GF.  A nonnegative Motzkin
walk from $0$ to height $h$ has GF $x^hM^{h+1}$, by cutting at the last
upward crossing of each level.  Walks from $1$ to $5$ whose minimum is
at least $1$ contribute $x^4M^5$.  The other walks first reach zero,
with first-passage GF $xM$, then continue with GF $x^5M^6$.
Thus their total GF is $x^4M^5(1+x^2M^2)=x^4(1-x)M^6$.
Forbidding a horizontal first step multiplies by $1-x$, proving
$F=x^4(1-x)^2M^6$ with the actual offset four.

\paragraph{Exact certificate.}

\[F(x)=\frac{\left(1 - x\right)^{2} \left(- x - y + 1\right)^{6}}{64 x^{8}}.\]

Here $y^2=- 3 x^{2} - 2 x + 1$, with $y(0)=1$.

Substitution in \eqref{round8:r8:gf-rational}, followed by exact cancellation, gives

\[LF=P(x)=0.\]

The two rational residuals are $P(\xi)$ and $0$.  Their coefficient interpretation proves the recurrence for all $n\geq 8$; this is at least the source offset plus the recurrence order.

\paragraph{Independent consistency check.} The companion validation computes the defining finite sum, combinatorial dynamic program, or nonlinear grammar through index 47, without using the conjectured recurrence.  All listed source coefficients match.  These checks audit the transcription and branch; the exact identities above supply the proof.


\subsection{A098519: the selected recurrence}

\label{round8:r8entry:A098519}

\paragraph{Source and assertion.}

\href{https://oeis.org/A098519}{A098519}, consulted on 20 October 2026, explicitly marks the recurrence as a conjecture attributed to R. J. Mathar, Nov 23 2011.

The source offset is $0$.  With $R_n=\sum_{j=0}^{2}p_j(n)a_{n-j}$, the exact source coefficients are:

\[\begin{aligned}
p_{0}(n)&=\left(n - 1\right) \left(n + 1\right)\\
p_{1}(n)&=- n \left(2 n - 1\right)\\
p_{2}(n)&=- 11 n \left(n - 1\right)\\
\end{aligned}\]

We prove $R_n=0$ for every integer $n\geq 2$.

\paragraph{Definition and generating-function derivation.}

The defining EGF is $e^zI_1(2\sqrt3z)/\sqrt3$.
Expanding its two factors gives
\[
 a_n=\sum_{k\geq0}\binom nk\binom{n-k}{k+1}3^k.
\]
The untransformed sequence has OGF
$2x/(1-12x^2+\sqrt{1-12x^2})$.
The binomial-transform rule $F(x)=B(x/(1-x))/(1-x)$, proved by summing
$\binom nj$, gives the source radical expression below.
No property of special functions beyond their displayed power series is
needed for this ordinary generating-function certificate.

\paragraph{Exact certificate.}

\[F(x)=\frac{2 x}{- 11 x^{2} - 2 x + y \left(1 - x\right) + 1}.\]

Here $y^2=- 11 x^{2} - 2 x + 1$, with $y(0)=1$.

Substitution in \eqref{round8:r8:gf-rational}, followed by exact cancellation, gives

\[LF=P(x)=0.\]

The two rational residuals are $P(\xi)$ and $0$.  Their coefficient interpretation proves the recurrence for all $n\geq 2$; this is at least the source offset plus the recurrence order.

\paragraph{Independent consistency check.} The companion validation computes the defining finite sum, combinatorial dynamic program, or nonlinear grammar through index 47, without using the conjectured recurrence.  All listed source coefficients match.  These checks audit the transcription and branch; the exact identities above supply the proof.


\subsection{A103138: the selected recurrence}

\label{round8:r8entry:A103138}

\paragraph{Source and assertion.}

\href{https://oeis.org/A103138}{A103138}, consulted on 20 October 2026, explicitly marks the recurrence as a conjecture attributed to R. J. Mathar, Dec 14 2011.

The source offset is $0$.  With $R_n=\sum_{j=0}^{3}p_j(n)a_{n-j}$, the exact source coefficients are:

\[\begin{aligned}
p_{0}(n)&=2 n\\
p_{1}(n)&=13 n - 20\\
p_{2}(n)&=8 n - 27\\
p_{3}(n)&=n - 5\\
\end{aligned}\]

We prove $R_n=0$ for every integer $n\geq 3$.

\paragraph{Definition and generating-function derivation.}

The source defining GF is
$F=(1-xS(-x))xS(-x)$, where $S=1+xS+xS^2$ is the large-Schroeder GF.
Put $t=xS(-x)$.  Its quadratic equation gives
$x=t(1+t)/(1-t)$, and $F=t-t^2$.
This also proves Hanna's August 6, 2013 composition identity used in
the certificate.  The parameter has $\xi'(0)=1$, selecting the
formal inverse uniquely.

\paragraph{Exact certificate.}

\[F(\xi(t))=- t^{2} + t.\]

\[\xi(t)=\frac{t \left(t + 1\right)}{1 - t},\qquad \phi(t)=\frac{1 - t}{t + 1},\qquad y^2=1.\]

The formal parameter satisfies $t=x\phi(t)$; every square root uses its branch with constant term one.

Substitution in \eqref{round8:r8:gf-rational}, followed by exact cancellation, gives

\[LF=P(x)=- 6 x^{2} + 2 x.\]

The two rational residuals are $P(\xi)$ and $0$.  Their coefficient interpretation proves the recurrence for all $n\geq 3$; this is at least the source offset plus the recurrence order.

\paragraph{Independent consistency check.} The companion validation computes the defining finite sum, combinatorial dynamic program, or nonlinear grammar through index 47, without using the conjectured recurrence.  All listed source coefficients match.  These checks audit the transcription and branch; the exact identities above supply the proof.


\subsection{A112702: the selected recurrence}

\label{round8:r8entry:A112702}

\paragraph{Source and assertion.}

\href{https://oeis.org/A112702}{A112702}, consulted on 20 October 2026, explicitly marks the recurrence as a conjecture attributed to R. J. Mathar, Jun 08 2016.

The source offset is $0$.  With $R_n=\sum_{j=0}^{2}p_j(n)a_{n-j}$, the exact source coefficients are:

\[\begin{aligned}
p_{0}(n)&=n + 1\\
p_{1}(n)&=- 3 \left(11 n - 5\right)\\
p_{2}(n)&=16 \left(2 n - 1\right)\\
\end{aligned}\]

We prove $R_n=0$ for every integer $n\geq 2$.

\paragraph{Definition and generating-function derivation.}

The definition is $a_n=\sum_{k=0}^n8^k C_k$.
Summing first over $n$ gives $F=C(8x)/(1-x)$, where
$C(x)=(1-\sqrt{1-4x})/(2x)$.  Thus the generating function used below
follows directly from the definition.

\paragraph{Exact certificate.}

\[F(x)=\frac{1 - y}{16 x \left(1 - x\right)}.\]

Here $y^2=1 - 32 x$, with $y(0)=1$.

Substitution in \eqref{round8:r8:gf-rational}, followed by exact cancellation, gives

\[LF=P(x)=1.\]

The two rational residuals are $P(\xi)$ and $0$.  Their coefficient interpretation proves the recurrence for all $n\geq 2$; this is at least the source offset plus the recurrence order.

\paragraph{Independent consistency check.} The companion validation computes the defining finite sum, combinatorial dynamic program, or nonlinear grammar through index 47, without using the conjectured recurrence.  All listed source coefficients match.  These checks audit the transcription and branch; the exact identities above supply the proof.


\subsection{A113409: the selected recurrence}

\label{round8:r8entry:A113409}

\paragraph{Source and assertion.}

\href{https://oeis.org/A113409}{A113409}, consulted on 20 October 2026, explicitly marks the recurrence as a conjecture attributed to R. J. Mathar, Nov 07 2012.

The source offset is $0$.  With $R_n=\sum_{j=0}^{4}p_j(n)a_{n-j}$, the exact source coefficients are:

\[\begin{aligned}
p_{0}(n)&=n + 2\\
p_{1}(n)&=- 2 \left(n + 1\right)\\
p_{2}(n)&=n - 4\\
p_{3}(n)&=2\\
p_{4}(n)&=- 4 \left(n - 2\right)\\
\end{aligned}\]

We prove $R_n=0$ for every integer $n\geq 4$.

\paragraph{Definition and generating-function derivation.}

The defining finite sum is
\[
 a_n=\sum_k\binom{n-k}{k}\binom{k}{\lfloor k/2\rfloor}.
\]
Put $B(z)=\sum_k\binom{k}{\lfloor k/2\rfloor}z^k$.
Separating even and odd $k$ gives
\[
 B(z)=\frac{1+2z-\sqrt{1-4z^2}}{2z\sqrt{1-4z^2}}.
\]
Since $\sum_{n\geq2k}\binom{n-k}{k}x^n=x^{2k}/(1-x)^{k+1}$,
we obtain $F=B(x^2/(1-x))/(1-x)$, which simplifies to the seed below.
The source's older expression
$(1-xC(x^2))/(1-x^2-x^4C(x^4))$ has coefficient $-1$ at $x$,
instead of the correct $a_1=1$.  It is explicitly rejected as a premise.
The recurrence itself is proved without changing any coefficient.

\paragraph{Exact certificate.}

\[F(x)=\frac{2 x^{2} - x - y + 1}{2 x^{2} y}.\]

Here $y^2=- 4 x^{4} + x^{2} - 2 x + 1$, with $y(0)=1$.

Substitution in \eqref{round8:r8:gf-rational}, followed by exact cancellation, gives

\[LF=P(x)=2 - x.\]

The two rational residuals are $P(\xi)$ and $0$.  Their coefficient interpretation proves the recurrence for all $n\geq 4$; this is at least the source offset plus the recurrence order.

\paragraph{Independent consistency check.} The companion validation computes the defining finite sum, combinatorial dynamic program, or nonlinear grammar through index 47, without using the conjectured recurrence.  All listed source coefficients match.  These checks audit the transcription and branch; the exact identities above supply the proof.


\subsection{A114851: the selected recurrence}

\label{round8:r8entry:A114851}

\paragraph{Source and assertion.}

\href{https://oeis.org/A114851}{A114851}, consulted on 20 October 2026, explicitly marks the recurrence as a conjecture attributed to R. J. Mathar, Mar 04 2015.

The source offset is $0$.  With $R_n=\sum_{j=0}^{6}p_j(n)a_{n-j}$, the exact source coefficients are:

\[\begin{aligned}
p_{0}(n)&=n + 2\\
p_{1}(n)&=- 2 \left(n + 1\right)\\
p_{2}(n)&=2 - n\\
p_{3}(n)&=4 \left(n - 2\right)\\
p_{4}(n)&=18 - 5 n\\
p_{5}(n)&=2 \left(n - 4\right)\\
p_{6}(n)&=n - 6\\
\end{aligned}\]

We prove $R_n=0$ for every integer $n\geq 6$.

\paragraph{Definition and generating-function derivation.}

A binary lambda term is a de Bruijn index of length at least two,
an abstraction of cost two followed by a term, or an application of cost
two followed by two ordered terms.  Thus
\[
 F=\frac{x^2}{1-x}+x^2F+x^2F^2,
 \qquad a_0=a_1=0.
\]
This is the nonlinear recurrence proved by Grygiel--Lescanne,
\href{https://arxiv.org/abs/1401.0379}{arXiv:1401.0379}, Section 5.
Solving the quadratic with zero constant term gives the seed below.
Its radicand is
$(1-x)^2(1-x^2)^2-4x^4(1-x)$.
The selected finite-order linear recurrence is certified separately.

\paragraph{Exact certificate.}

\[F(x)=\frac{- y + \left(1 - x\right) \left(1 - x^{2}\right)}{2 x^{2} \left(1 - x\right)}.\]

Here $y^2=- 4 x^{4} \left(1 - x\right) + \left(1 - x\right)^{2} \left(1 - x^{2}\right)^{2}$, with $y(0)=1$.

Substitution in \eqref{round8:r8:gf-rational}, followed by exact cancellation, gives

\[LF=P(x)=- x^{5} - 3 x^{3} + 4 x^{2}.\]

The two rational residuals are $P(\xi)$ and $0$.  Their coefficient interpretation proves the recurrence for all $n\geq 6$; this is at least the source offset plus the recurrence order.

\paragraph{Independent consistency check.} The companion validation computes the defining finite sum, combinatorial dynamic program, or nonlinear grammar through index 47, without using the conjectured recurrence.  All listed source coefficients match.  These checks audit the transcription and branch; the exact identities above supply the proof.


\subsection{A162972: the selected recurrence}

\label{round8:r8entry:A162972}

\paragraph{Source and assertion.}

\href{https://oeis.org/A162972}{A162972}, consulted on 20 October 2026, explicitly marks the recurrence as a conjecture attributed to R. J. Mathar, Jul 26 2022.

The source offset is $1$.  With $R_n=\sum_{j=0}^{7}p_j(n)a_{n-j}$, the exact source coefficients are:

\[\begin{aligned}
p_{0}(n)&=1\\
p_{1}(n)&=8 - 3 n\\
p_{2}(n)&=\left(n - 3\right) \left(3 n - 14\right)\\
p_{3}(n)&=- n^{3} + 19 n^{2} - 98 n + 156\\
p_{4}(n)&=- 4 n^{3} + 56 n^{2} - 258 n + 395\\
p_{5}(n)&=- 6 n^{3} + 84 n^{2} - 395 n + 624\\
p_{6}(n)&=- \left(n - 5\right) \left(4 n^{2} - 41 n + 106\right)\\
p_{7}(n)&=- \left(n - 6\right)^{2} \left(n - 5\right)\\
\end{aligned}\]

We prove $R_n=0$ for every integer $n\geq 8$.

\paragraph{Definition and generating-function derivation.}

Mark each cycle by $u$.  The EGF for permutations is
$(1-z)^{-u}$, and that for derangements is
$(e^{-z}/(1-z))^u$, by the exponential formula for disjoint cycles.
Differentiating their difference at $u=1$ counts cycles in the
non-derangement permutations and gives
\[
 F(z)=\frac{ze^{-z}+(e^{-z}-1)\log(1-z)}{1-z}.
\]
For $\ell=\log(1-z)$ and $e=e^{-z}$, the basis is
$(1,\ell,e,e\ell)$, whose derivatives are
$(0,-1/(1-z),-e,-e/(1-z)-e\ell)$.
The seed coefficients are
$(0,-1/(1-z),z/(1-z),1/(1-z))$.
Apply \eqref{round8:r8:gf-egf} with $k=1$ and $r=7$ using these derivatives.
All four rational residuals vanish identically.  This proves the
recurrence from $n=8$, where every source index is at least one.

\paragraph{Exact certificate.} The seed and basis derivatives just given, together with the displayed $p_j$ and \eqref{round8:r8:gf-egf}, give an identically zero rational vector.  The exact matrix and its zero residual are also recorded in the machine-readable certificate.  Thus every required coefficient vanishes, proving the all-index assertion.

\paragraph{Independent consistency check.} The companion validation computes the defining finite sum, combinatorial dynamic program, or nonlinear grammar through index 47, without using the conjectured recurrence.  All listed source coefficients match.  These checks audit the transcription and branch; the exact identities above supply the proof.


\subsection{A168505: the selected recurrence}

\label{round8:r8entry:A168505}

\paragraph{Source and assertion.}

\href{https://oeis.org/A168505}{A168505}, consulted on 20 October 2026, explicitly marks the recurrence as a conjecture attributed to R. J. Mathar, Feb 10 2015.

The source offset is $0$.  With $R_n=\sum_{j=0}^{4}p_j(n)a_{n-j}$, the exact source coefficients are:

\[\begin{aligned}
p_{0}(n)&=n + 1\\
p_{1}(n)&=- 3 n\\
p_{2}(n)&=3 \left(n - 1\right)\\
p_{3}(n)&=3 \left(n - 4\right)\\
p_{4}(n)&=- 2 \left(2 n - 7\right)\\
\end{aligned}\]

We prove $R_n=0$ for every integer $n\geq 4$.

\paragraph{Definition and generating-function derivation.}

The defining continued fraction repeats the three signs $-,+,-$.
Writing one full period gives
\[
 F=\frac1{1-\displaystyle\frac{x}{1+\displaystyle\frac{x}{1-xF}}}.
\]
Clearing denominators yields
$x(1-x)F^2-(1+x)F+(1+x)=0$.
The branch $F(0)=1$ is exactly the radical below.  This also derives
Somos's January 20, 2017 identity $(F-1)/F^2=x(1-x)/(1+x)$.
That known nonlinear identity is an input; it is not the selected linear
polynomial recurrence.

\paragraph{Exact certificate.}

\[F(x)=\frac{x - y + 1}{2 x \left(1 - x\right)}.\]

Here $y^2=4 x^{3} + x^{2} - 2 x + 1$, with $y(0)=1$.

Substitution in \eqref{round8:r8:gf-rational}, followed by exact cancellation, gives

\[LF=P(x)=- x^{3} - 3 x^{2} - x + 1.\]

The two rational residuals are $P(\xi)$ and $0$.  Their coefficient interpretation proves the recurrence for all $n\geq 4$; this is at least the source offset plus the recurrence order.

\paragraph{Independent consistency check.} The companion validation computes the defining finite sum, combinatorial dynamic program, or nonlinear grammar through index 47, without using the conjectured recurrence.  All listed source coefficients match.  These checks audit the transcription and branch; the exact identities above supply the proof.


\subsection{A174169: the selected recurrence}

\label{round8:r8entry:A174169}

\paragraph{Source and assertion.}

\href{https://oeis.org/A174169}{A174169}, consulted on 20 October 2026, explicitly marks the recurrence as a conjecture attributed to R. J. Mathar, Sep 30 2012.

The source offset is $0$.  With $R_n=\sum_{j=0}^{4}p_j(n)a_{n-j}$, the exact source coefficients are:

\[\begin{aligned}
p_{0}(n)&=n + 2\\
p_{1}(n)&=- 2 n - 1\\
p_{2}(n)&=3 \left(n - 1\right)\\
p_{3}(n)&=- 3 \left(2 n - 5\right)\\
p_{4}(n)&=9 \left(n - 4\right)\\
\end{aligned}\]

We prove $R_n=0$ for every integer $n\geq 4$.

\paragraph{Definition and generating-function derivation.}

Manyama recorded the correct quadratic identity on
October 17, 2025:
\[
 (1-x+3x^2)F=1+x^2F^2,\qquad a_0=a_1=1.
\]
It also follows from the Chebyshev substitution
\[
 F=\frac1{1+3x^2}M\left(\frac{x}{1+3x^2}\right).
\]
The branch with constant term one gives the displayed radical.
For clarity, the older printed substitution denominator $1+3x$ is
incorrect, and its older finite sum omits a factor:
the corrected sum is
\[
 a_n=\sum_k(-3)^k\binom{n-k}{k}M_{n-2k}.
\]
Neither erroneous expression is a premise, and these auxiliary fixes
are not additional counted assertions.

\paragraph{Exact certificate.}

\[F(x)=\frac{3 x^{2} - x - y + 1}{2 x^{2}}.\]

Here $y^2=9 x^{4} - 6 x^{3} + 3 x^{2} - 2 x + 1$, with $y(0)=1$.

Substitution in \eqref{round8:r8:gf-rational}, followed by exact cancellation, gives

\[LF=P(x)=2 - 6 x^{2}.\]

The two rational residuals are $P(\xi)$ and $0$.  Their coefficient interpretation proves the recurrence for all $n\geq 4$; this is at least the source offset plus the recurrence order.

\paragraph{Independent consistency check.} The companion validation computes the defining finite sum, combinatorial dynamic program, or nonlinear grammar through index 47, without using the conjectured recurrence.  All listed source coefficients match.  These checks audit the transcription and branch; the exact identities above supply the proof.


\subsection{A182892: the selected recurrence}

\label{round8:r8entry:A182892}

\paragraph{Source and assertion.}

\href{https://oeis.org/A182892}{A182892}, consulted on 20 October 2026, explicitly marks the recurrence as a conjecture attributed to R. J. Mathar, Jun 14 2016.

The source offset is $0$.  With $R_n=\sum_{j=0}^{8}p_j(n)a_{n-j}$, the exact source coefficients are:

\[\begin{aligned}
p_{0}(n)&=n\\
p_{1}(n)&=n - 2\\
p_{2}(n)&=- 2 \left(9 n - 16\right)\\
p_{3}(n)&=5 \left(2 n - 5\right)\\
p_{4}(n)&=10 n - 33\\
p_{5}(n)&=2 \left(26 n - 109\right)\\
p_{6}(n)&=13 n - 37\\
p_{7}(n)&=13 n - 63\\
p_{8}(n)&=- 10 \left(n - 7\right)\\
\end{aligned}\]

We prove $R_n=0$ for every integer $n\geq 8$.

\paragraph{Definition and generating-function derivation.}

The source permits heights of either sign.  An unrestricted weighted
bridge with horizontal weights $x,x^2$, upward weight $x^2$ and downward
weight $x$ has GF
\[
 B(x)=\sum_{m\geq0}\binom{2m}{m}
       \frac{x^{3m}}{(1-x-x^2)^{2m+1}}
 =\frac1{\sqrt{(1-x-x^2)^2-4x^3}}.
\]
Decompose a bridge into its successive returns to height zero.  Removing
the weight-two horizontal component at zero increases the inverse GF by
$x^2$, so the required GF is $F=1/(B^{-1}+x^2)$.
The radicand factors as $(1+x+x^2)(1-3x+x^2)$.
This proves the GF recorded by Deutsch, independently of the background
composition/fence-poset papers linked in the entry.

\paragraph{Exact certificate.}

\[F(x)=\frac{1}{x^{2} + y}.\]

Here $y^2=\left(x^{2} - 3 x + 1\right) \left(x^{2} + x + 1\right)$, with $y(0)=1$.

Substitution in \eqref{round8:r8:gf-rational}, followed by exact cancellation, gives

\[LF=P(x)=5 x^{6} + 6 x^{5} + 16 x^{4} - 7 x^{3} - 2 x^{2}.\]

The two rational residuals are $P(\xi)$ and $0$.  Their coefficient interpretation proves the recurrence for all $n\geq 8$; this is at least the source offset plus the recurrence order.

\paragraph{Independent consistency check.} The companion validation computes the defining finite sum, combinatorial dynamic program, or nonlinear grammar through index 47, without using the conjectured recurrence.  All listed source coefficients match.  These checks audit the transcription and branch; the exact identities above supply the proof.


\subsection{A220898: the selected recurrence}

\label{round8:r8entry:A220898}

\paragraph{Source and assertion.}

\href{https://oeis.org/A220898}{A220898}, consulted on 20 October 2026, explicitly marks the recurrence as a conjecture attributed to R. J. Mathar, May 30 2014.

The source offset is $0$.  With $R_n=\sum_{j=0}^{3}p_j(n)a_{n-j}$, the exact source coefficients are:

\[\begin{aligned}
p_{0}(n)&=2 \left(n + 1\right) \left(2 n + 1\right)\\
p_{1}(n)&=- 15 n^{2} + 19 n - 8\\
p_{2}(n)&=- 2 \left(n - 2\right) \left(21 n - 22\right)\\
p_{3}(n)&=- 23 \left(n - 3\right) \left(n - 2\right)\\
\end{aligned}\]

We prove $R_n=0$ for every integer $n\geq 3$.

\paragraph{Definition and generating-function derivation.}

Kitaev, Salimov, Severs and Ulfarsson (2012), Theorem 1,
prove that the primitive-map GF is $F(x)=G(x/(1+x))$.
The established map series has
$[z^n]G=2(3n)!/((2n+1)!(n+1)!)$, including constant term two.
Its parameterization is
$z=t(1-t)^2$, $G=(2-3t)/(1-t)^2$.
For completeness Lagrange inversion gives, for $n\geq1$,
\[
 [z^n]G=\frac1n\left(\binom{3n+1}{n-1}
                         -3\binom{3n}{n-2}\right)
 =\frac{2(3n)!}{(2n+1)!(n+1)!}.
\]
Thus with $u=t(1-t)^2$, setting $x=u/(1-u)$ gives the seed below.
The known map GF and composition are credited premises; the selected
2014 recurrence is the assertion certified here.

\paragraph{Exact certificate.}

\[F(\xi(t))=\frac{2 - 3 t}{\left(1 - t\right)^{2}}.\]

\[\xi(t)=\frac{t \left(1 - t\right)^{2}}{- t \left(1 - t\right)^{2} + 1},\qquad \phi(t)=\frac{- t \left(1 - t\right)^{2} + 1}{\left(1 - t\right)^{2}},\qquad y^2=1.\]

The formal parameter satisfies $t=x\phi(t)$; every square root uses its branch with constant term one.

Substitution in \eqref{round8:r8:gf-rational}, followed by exact cancellation, gives

\[LF=P(x)=4 x + 4.\]

The two rational residuals are $P(\xi)$ and $0$.  Their coefficient interpretation proves the recurrence for all $n\geq 3$; this is at least the source offset plus the recurrence order.

\paragraph{Independent consistency check.} The companion validation computes the defining finite sum, combinatorial dynamic program, or nonlinear grammar through index 47, without using the conjectured recurrence.  All listed source coefficients match.  These checks audit the transcription and branch; the exact identities above supply the proof.


\subsection{A225692: the selected recurrence}

\label{round8:r8entry:A225692}

\paragraph{Source and assertion.}

\href{https://oeis.org/A225692}{A225692}, consulted on 20 October 2026, explicitly marks the recurrence as a conjecture attributed to R. J. Mathar, Jul 09 2013.

The source offset is $2$.  With $R_n=\sum_{j=0}^{3}p_j(n)a_{n-j}$, the exact source coefficients are:

\[\begin{aligned}
p_{0}(n)&=- 10 \left(n - 4\right) \left(n + 1\right)\\
p_{1}(n)&=51 n^{2} - 196 n + 80\\
p_{2}(n)&=- 45 n^{2} + 134 n - 99\\
p_{3}(n)&=2 \left(n - 3\right) \left(2 n - 5\right)\\
\end{aligned}\]

We prove $R_n=0$ for every integer $n\geq 5$.

\paragraph{Definition and generating-function derivation.}

In the source's Dyck-pattern order, containment means a scattered
subword.  A Dyck path of semilength $n$ contains $U^{n-1}D^{n-1}$
exactly when its first $n$ steps have at most one down step.
Indeed a cut separating the selected up and down steps must have at least
$n-1$ ups and at most one down, and can be taken at the midpoint;
conversely that midpoint condition supplies the selected subword.
There is one path with no down step in its first half.  For exactly one
down step, the nonnegative first half has $n-1$ choices, and by reversal
the second half has $n-1$ choices.  Thus
\[
 a_n=C_n-1-(n-1)^2\qquad(n\geq2).
\]
This also follows by subtracting the first two terms from the
ballot-square sum in Bernini--Ferrari--Pinzani--West (2013), Section 3.3.
Writing $G=\sum_{n\geq0}(n^2-2n+2)x^n$, the actual offset-two series
is $F=C-G+1$; the extra $1$ removes the spurious constant coefficient.
The exact operator identity below includes that boundary adjustment.

\paragraph{Exact certificate.}

\[F(x)=\frac{2 x}{\left(1 - x\right)^{2}} - \frac{x \left(x + 1\right)}{\left(1 - x\right)^{3}} + 1 - \frac{2}{1 - x} + \frac{1 - y}{2 x}.\]

Here $y^2=1 - 4 x$, with $y(0)=1$.

Substitution in \eqref{round8:r8:gf-rational}, followed by exact cancellation, gives

\[LF=P(x)=0.\]

The two rational residuals are $P(\xi)$ and $0$.  Their coefficient interpretation proves the recurrence for all $n\geq 5$; this is at least the source offset plus the recurrence order.

\paragraph{Independent consistency check.} The companion validation computes the defining finite sum, combinatorial dynamic program, or nonlinear grammar through index 47, without using the conjectured recurrence.  All listed source coefficients match.  These checks audit the transcription and branch; the exact identities above supply the proof.


\subsection{A240722: the selected recurrence}

\label{round8:r8entry:A240722}

\paragraph{Source and assertion.}

\href{https://oeis.org/A240722}{A240722}, consulted on 20 October 2026, explicitly marks the recurrence as a conjecture attributed to R. J. Mathar, Jun 14 2016.

The source offset is $0$.  With $R_n=\sum_{j=0}^{2}p_j(n)a_{n-j}$, the exact source coefficients are:

\[\begin{aligned}
p_{0}(n)&=- \left(n + 1\right) \left(2 n + 1\right) \left(3 n - 1\right)\\
p_{1}(n)&=2 \left(21 n^{3} + 14 n^{2} - 2 n - 1\right)\\
p_{2}(n)&=8 n \left(2 n - 1\right) \left(3 n + 2\right)\\
\end{aligned}\]

We prove $R_n=0$ for every integer $n\geq 2$.

\paragraph{Definition and generating-function derivation.}

The source definition is $F=B'/B$, where
\[
 B(x)=\frac{-1-\sqrt{1-8x}+\sqrt{2+2\sqrt{1-8x}+8x}}{4x}.
\]
Set $x=t(1-2t)$ and $y^2=1-4t^2$, $y(0)=1$.
Then $\sqrt{1-8x}=1-4t$ and
\[
 B=\frac{-1+2t+y}{2t(1-2t)},\qquad
 F=\frac1{1-4t}\frac{dB/dt}{B},\qquad
 y'=-4t/y.
\]
These exact substitutions specify the seed used in the rational
certificate.  The removable singularities select $B(0)=1$.
The proof uses the original offset zero, not the shifted convention
of the source's later congruence comments.

\paragraph{Exact certificate.}

Use the $B,F,x,y$ expressions immediately above, differentiating with $y\prime=-4t/y$.

\[\xi(t)=t \left(1 - 2 t\right),\qquad \phi(t)=\frac{1}{1 - 2 t},\qquad y^2=1 - 4 t^{2}.\]

The formal parameter satisfies $t=x\phi(t)$; every square root uses its branch with constant term one.

Substitution in \eqref{round8:r8:gf-rational}, followed by exact cancellation, gives

\[LF=P(x)=1 - 20 x.\]

The two rational residuals are $P(\xi)$ and $0$.  Their coefficient interpretation proves the recurrence for all $n\geq 2$; this is at least the source offset plus the recurrence order.

\paragraph{Independent consistency check.} The companion validation computes the defining finite sum, combinatorial dynamic program, or nonlinear grammar through index 47, without using the conjectured recurrence.  All listed source coefficients match.  These checks audit the transcription and branch; the exact identities above supply the proof.


\subsection{A246606: the selected recurrence}

\label{round8:r8entry:A246606}

\paragraph{Source and assertion.}

\href{https://oeis.org/A246606}{A246606}, consulted on 20 October 2026, explicitly marks the recurrence as a conjecture attributed to R. J. Mathar, Jul 01 2022.

The source offset is $1$.  With $R_n=\sum_{j=0}^{2}p_j(n)a_{n-j}$, the exact source coefficients are:

\[\begin{aligned}
p_{0}(n)&=3 - 2 n\\
p_{1}(n)&=4 \left(n - 1\right) \left(2 n^{2} - 4 n + 1\right)\\
p_{2}(n)&=\left(n - 2\right) \left(n - 1\right) \left(2 n - 1\right)\\
\end{aligned}\]

We prove $R_n=0$ for every integer $n\geq 3$.

\paragraph{Definition and generating-function derivation.}

The terminating hypergeometric expression recorded by Luschny
on November 4, 2018 expands to
\[
 a_n=\sum_{j=0}^{n-1}(-1)^j\binom{n-1}{j}(2n-1-j)!
 =\int_0^\infty e^{-t}t^n(t-1)^{n-1}\,dt.
\]
The integral equality follows by expanding $(t-1)^{n-1}$ and integrating
monomials against $e^{-t}$.  For $n\geq3$, put
$g=e^{-t}t^{n-2}(t-1)^{n-3}$ and
$R=p_0t^2(t-1)^2+p_1t(t-1)+p_2$.
Direct polynomial multiplication gives
\[
 R=t(t-1)Q'+(-t^2+(2n-2)t-(n-1))Q,
\]
where
\[
 Q=(2n-3)t^2+(4n^2-10n+6)t-2n^2+5n-2.
\]
Thus the recurrence residual is the integral of
$\frac{d}{dt}\{e^{-t}t^{n-1}(t-1)^{n-2}Q\}$.
The bracket vanishes at zero because $n-1\geq2$, and at infinity because
of $e^{-t}$.  There is no singularity at $t=1$, since the exponents are
nonnegative integers.  The integral is therefore zero for every $n\geq3$.

\paragraph{Independent consistency check.} The companion validation computes the defining finite sum, combinatorial dynamic program, or nonlinear grammar through index 47, without using the conjectured recurrence.  All listed source coefficients match.  These checks audit the transcription and branch; the exact identities above supply the proof.


\subsection{A268402: the selected recurrence}

\label{round8:r8entry:A268402}

\paragraph{Source and assertion.}

\href{https://oeis.org/A268402}{A268402}, consulted on 20 October 2026, explicitly marks the recurrence as a conjecture attributed to R. J. Mathar, Oct 07 2016.

The source offset is $5$.  With $R_n=\sum_{j=0}^{4}p_j(n)a_{n-j}$, the exact source coefficients are:

\[\begin{aligned}
p_{0}(n)&=- \left(n - 5\right) \left(n - 3\right) \left(55 n - 618\right)\\
p_{1}(n)&=- 55 n^{3} - 840 n^{2} + 6969 n - 5742\\
p_{2}(n)&=3135 n^{3} - 37532 n^{2} + 138815 n - 163614\\
p_{3}(n)&=- 7645 n^{3} + 93072 n^{2} - 343985 n + 391386\\
p_{4}(n)&=- 18 \left(n - 3\right) \left(2 n - 5\right) \left(55 n - 126\right)\\
\end{aligned}\]

We prove $R_n=0$ for every integer $n\geq 9$.

\paragraph{Definition and generating-function derivation.}

Cut a north-east path at its returns to the diagonal.  Each nonempty
piece is a prime excursion, on either side, with GF $E=xC(x)$.
A right bounce occurs precisely between two consecutive pieces on the
right.  If $u$ marks right bounces and $U,V$ count nonempty words ending
on the left or right, then
$U=E(1+U+V)$ and $V=E(1+U+uV)$.
Consequently the total GF is
\[
 H(u,x)=\frac{1+(1-u)E}{1-E-E^2-uE(1-E)}.
\]
For $r\geq1$ its coefficient of $u^r$ is
$E^{r+1}(1-E)^{r-1}/(1-E-E^2)^{r+1}$.
Taking $r=4$ gives the required GF and, after substituting
$E=(1-\sqrt{1-4x})/2$, simplifies to the source expression below.
This also gives an independent derivation of the specialization of
Pan--Remmel's proved bivariate GF (2016, Section 3.2).

\paragraph{Exact certificate.}

\[F(x)=- \frac{x^{3} \left(y - 1\right)^{5} \left(2 x + y - 1\right)}{2 \left(x \left(y - 5\right) - y + 1\right)^{5}}.\]

Here $y^2=1 - 4 x$, with $y(0)=1$.

Substitution in \eqref{round8:r8:gf-rational}, followed by exact cancellation, gives

\[LF=P(x)=0.\]

The two rational residuals are $P(\xi)$ and $0$.  Their coefficient interpretation proves the recurrence for all $n\geq 9$; this is at least the source offset plus the recurrence order.

\paragraph{Independent consistency check.} The companion validation computes the defining finite sum, combinatorial dynamic program, or nonlinear grammar through index 47, without using the conjectured recurrence.  All listed source coefficients match.  These checks audit the transcription and branch; the exact identities above supply the proof.


\subsection{A278618: the selected recurrence}

\label{round8:r8entry:A278618}

\paragraph{Source and assertion.}

\href{https://oeis.org/A278618}{A278618}, consulted on 20 October 2026, explicitly marks the recurrence as a conjecture attributed to R. J. Mathar, Dec 02 2016.

The source offset is $0$.  With $R_n=\sum_{j=0}^{3}p_j(n)a_{n-j}$, the exact source coefficients are:

\[\begin{aligned}
p_{0}(n)&=2 n \left(n + 1\right) \left(2 n - 1\right) \left(5 n - 8\right)\\
p_{1}(n)&=- n \left(115 n^{3} - 344 n^{2} + 299 n - 82\right)\\
p_{2}(n)&=- 4 \left(2 n - 1\right) \left(5 n^{3} + 27 n^{2} - 74 n + 30\right)\\
p_{3}(n)&=36 \left(n - 2\right) \left(2 n - 3\right) \left(2 n - 1\right) \left(5 n - 3\right)\\
\end{aligned}\]

We prove $R_n=0$ for every integer $n\geq 3$.

\paragraph{Definition and generating-function derivation.}

Set $B(u)=1+u^2/(1-u)=(1-u+u^2)/(1-u)$.
Expanding its $(2n+1)$-st power gives
\[
 [u^n]B(u)^{2n+1}
 =\sum_j\binom{2n+1}{j}\binom{n-j-1}{n-2j}=a_n.
\]
The $j=0$ term is interpreted as $1$ for $n=0$ and $0$ otherwise.
Formula \eqref{round8:r8:gf-diagonal}, with $\phi=B^2$, gives the rational
parameterization below, without using the longer nested radical printed
in the source.

\paragraph{Exact certificate.}

\[F(\xi(t))=\frac{t^{2} - t + 1}{\left(1 - t\right) \left(- \frac{t \left(1 - t\right)^{2} \left(\frac{\left(4 t - 2\right) \left(t^{2} - t + 1\right)}{\left(1 - t\right)^{2}} + \frac{2 \left(t^{2} - t + 1\right)^{2}}{\left(1 - t\right)^{3}}\right)}{\left(t^{2} - t + 1\right)^{2}} + 1\right)}.\]

\[\xi(t)=\frac{t \left(1 - t\right)^{2}}{\left(t^{2} - t + 1\right)^{2}},\qquad \phi(t)=\frac{\left(t^{2} - t + 1\right)^{2}}{\left(1 - t\right)^{2}},\qquad y^2=1.\]

The formal parameter satisfies $t=x\phi(t)$; every square root uses its branch with constant term one.

Substitution in \eqref{round8:r8:gf-rational}, followed by exact cancellation, gives

\[LF=P(x)=12 x.\]

The two rational residuals are $P(\xi)$ and $0$.  Their coefficient interpretation proves the recurrence for all $n\geq 3$; this is at least the source offset plus the recurrence order.

\paragraph{Independent consistency check.} The companion validation computes the defining finite sum, combinatorial dynamic program, or nonlinear grammar through index 47, without using the conjectured recurrence.  All listed source coefficients match.  These checks audit the transcription and branch; the exact identities above supply the proof.


\subsection{A302180: the selected recurrence}

\label{round8:r8entry:A302180}

\paragraph{Source and assertion.}

\href{https://oeis.org/A302180}{A302180}, consulted on 20 October 2026, explicitly marks the recurrence as a conjecture attributed to R. J. Mathar, Oct 29 2025.

The source offset is $0$.  With $R_n=\sum_{j=0}^{3}p_j(n)a_{n-j}$, the exact source coefficients are:

\[\begin{aligned}
p_{0}(n)&=\left(n + 2\right) \left(n + 4\right)\\
p_{1}(n)&=- 3 \left(n^{2} + 3 n + 1\right)\\
p_{2}(n)&=- 13 \left(n - 1\right) \left(n + 1\right)\\
p_{3}(n)&=15 \left(n - 2\right) \left(n - 1\right)\\
\end{aligned}\]

We prove $R_n=0$ for every integer $n\geq 3$.

\paragraph{Definition and generating-function derivation.}

These are Dershowitz's walks of type $aad$: two independent
nonnegative excursions and one free direction.  Choosing the positions
of their steps proves
$F(z)=e^z(I_1(2z)/z)^2$.
Use the degree-two basis \eqref{round8:r8:gf-bessel} with $c=1$, seed
$(0,0,z^{-2})$, and $k=1$ in \eqref{round8:r8:gf-egf}.
Exact application of $\mathcal L$ gives the zero vector $(0,0,0)$.
All apparent Laurent poles cancel because $I_1(2z)=z+O(z^3)$.

\paragraph{Exact certificate.} The seed and basis derivatives just given, together with the displayed $p_j$ and \eqref{round8:r8:gf-egf}, give an identically zero rational vector.  The exact matrix and its zero residual are also recorded in the machine-readable certificate.  Thus every required coefficient vanishes, proving the all-index assertion.

\paragraph{Independent consistency check.} The companion validation computes the defining finite sum, combinatorial dynamic program, or nonlinear grammar through index 47, without using the conjectured recurrence.  All listed source coefficients match.  These checks audit the transcription and branch; the exact identities above supply the proof.


\subsection{A302181: the selected recurrence}

\label{round8:r8entry:A302181}

\paragraph{Source and assertion.}

\href{https://oeis.org/A302181}{A302181}, consulted on 20 October 2026, explicitly marks the recurrence as a conjecture attributed to R. J. Mathar, Oct 29 2025.

The source offset is $0$.  With $R_n=\sum_{j=0}^{2}p_j(n)a_{n-j}$, the exact source coefficients are:

\[\begin{aligned}
p_{0}(n)&=\left(n + 1\right)^{3}\\
p_{1}(n)&=- 2 \left(2 n - 1\right) \left(10 n^{2} + 10 n + 3\right)\\
p_{2}(n)&=36 n \left(2 n - 3\right) \left(2 n - 1\right)\\
\end{aligned}\]

We prove $R_n=0$ for every integer $n\geq 2$.

\paragraph{Definition and generating-function derivation.}

These are the even-length walks of type $abb$: a nonnegative
excursion and two unrestricted bridges.  Hence
\[
 \sum_{n\geq0}\frac{a_nz^{2n}}{(2n)!}
 =I_0(2z)^2I_1(2z)/z.
\]
The even-index convolution recorded by Tebni on November 6, 2024 is
another form of this decomposition.  Use $d=3$, $c=0$, seed
$(0,z^{-1},0,0)$, and crucially $k=2$ in \eqref{round8:r8:gf-egf}.
The basis derivative remains with respect to $z$, while the Euler
operator is $(z/2)d/dz$.  The resulting four rational residuals are
identically zero, proving the asserted recurrence in the index $n$.

\paragraph{Exact certificate.} The seed and basis derivatives just given, together with the displayed $p_j$ and \eqref{round8:r8:gf-egf}, give an identically zero rational vector.  The exact matrix and its zero residual are also recorded in the machine-readable certificate.  Thus every required coefficient vanishes, proving the all-index assertion.

\paragraph{Independent consistency check.} The companion validation computes the defining finite sum, combinatorial dynamic program, or nonlinear grammar through index 47, without using the conjectured recurrence.  All listed source coefficients match.  These checks audit the transcription and branch; the exact identities above supply the proof.


\subsection{A302184: the selected recurrence}

\label{round8:r8entry:A302184}

\paragraph{Source and assertion.}

\href{https://oeis.org/A302184}{A302184}, consulted on 20 October 2026, explicitly marks the recurrence as a conjecture attributed to R. J. Mathar, Oct 29 2025.

The source offset is $0$.  With $R_n=\sum_{j=0}^{3}p_j(n)a_{n-j}$, the exact source coefficients are:

\[\begin{aligned}
p_{0}(n)&=\left(n + 2\right)^{2}\\
p_{1}(n)&=- 2 \left(3 n^{2} + 5 n + 1\right)\\
p_{2}(n)&=- 4 \left(n - 1\right) \left(n + 3\right)\\
p_{3}(n)&=24 \left(n - 2\right) \left(n - 1\right)\\
\end{aligned}\]

We prove $R_n=0$ for every integer $n\geq 3$.

\paragraph{Definition and generating-function derivation.}

These are walks of type $abe$: one nonnegative excursion,
one unrestricted bridge and two free directions.  The multinomial
interleaving gives
$F(z)=e^{2z}I_0(2z)I_1(2z)/z$.
Equivalently this is Tebni's November 30, 2024 source convolution of
the aerated Catalan numbers with $\binom{2m}{m}$; the latter has EGF
$e^{2z}I_0(2z)$.
Use $d=2$, $c=2$, seed $(0,z^{-1},0)$ and $k=1$.
Formula \eqref{round8:r8:gf-bessel} gives $\mathcal LF=(0,0,0)$ exactly.

\paragraph{Exact certificate.} The seed and basis derivatives just given, together with the displayed $p_j$ and \eqref{round8:r8:gf-egf}, give an identically zero rational vector.  The exact matrix and its zero residual are also recorded in the machine-readable certificate.  Thus every required coefficient vanishes, proving the all-index assertion.

\paragraph{Independent consistency check.} The companion validation computes the defining finite sum, combinatorial dynamic program, or nonlinear grammar through index 47, without using the conjectured recurrence.  All listed source coefficients match.  These checks audit the transcription and branch; the exact identities above supply the proof.


\subsection{A302186: the selected recurrence}

\label{round8:r8entry:A302186}

\paragraph{Source and assertion.}

\href{https://oeis.org/A302186}{A302186}, consulted on 20 October 2026, explicitly marks the recurrence as a conjecture attributed to R. J. Mathar, Oct 29 2025.

The source offset is $0$.  With $R_n=\sum_{j=0}^{4}p_j(n)a_{n-j}$, the exact source coefficients are:

\[\begin{aligned}
p_{0}(n)&=\left(n + 2\right) \left(n + 3\right)\\
p_{1}(n)&=- 4 \left(3 n^{2} + 8 n + 3\right)\\
p_{2}(n)&=8 \left(4 n^{2} + 3 n - 3\right)\\
p_{3}(n)&=48 \left(n - 2\right)^{2}\\
p_{4}(n)&=- 144 \left(n - 3\right) \left(n - 2\right)\\
\end{aligned}\]

We prove $R_n=0$ for every integer $n\geq 4$.

\paragraph{Definition and generating-function derivation.}

These are walks of type $ace$: a nonnegative excursion,
a nonnegative meander and two free directions.  Interleaving gives
\[
 F(z)=e^{2z}I_1(2z)(I_0(2z)+I_1(2z))/z.
\]
This also derives the source double-binomial sum from the component
counts.  Use $d=2$, $c=2$, seed $(0,z^{-1},z^{-1})$ and $k=1$.
The rational matrix calculation in \eqref{round8:r8:gf-bessel} gives the zero
vector after applying $\mathcal L$.

\paragraph{Exact certificate.} The seed and basis derivatives just given, together with the displayed $p_j$ and \eqref{round8:r8:gf-egf}, give an identically zero rational vector.  The exact matrix and its zero residual are also recorded in the machine-readable certificate.  Thus every required coefficient vanishes, proving the all-index assertion.

\paragraph{Independent consistency check.} The companion validation computes the defining finite sum, combinatorial dynamic program, or nonlinear grammar through index 47, without using the conjectured recurrence.  All listed source coefficients match.  These checks audit the transcription and branch; the exact identities above supply the proof.


\subsection{A308815: the selected recurrence}

\label{round8:r8entry:A308815}

\paragraph{Source and assertion.}

\href{https://oeis.org/A308815}{A308815}, consulted on 20 October 2026, explicitly marks the recurrence as a conjecture attributed to R. J. Mathar, Jan 27 2020.

The source offset is $0$.  With $R_n=\sum_{j=0}^{3}p_j(n)a_{n-j}$, the exact source coefficients are:

\[\begin{aligned}
p_{0}(n)&=n \left(2 n + 1\right)\\
p_{1}(n)&=- 32 \left(n - 1\right) \left(9 n - 4\right)\\
p_{2}(n)&=1024 \left(3 n^{2} - 23 n + 29\right)\\
p_{3}(n)&=65536 \left(n - 2\right) \left(2 n - 5\right)\\
\end{aligned}\]

We prove $R_n=0$ for every integer $n\geq 3$.

\paragraph{Definition and generating-function derivation.}

Farkas, Moschetti, Naranjo and Pirola, Theorems 1.1--1.2 in
\href{https://arxiv.org/abs/1906.10406}{arXiv:1906.10406}, prove the alternating
sum and generating series.  The intended sequence is
\[
 a_n=16^n\sum_{i=0}^n(-2)^i\binom ni C_{2n-i}.
\]
The OEIS title prints $2^i$ instead of $(-2)^i$.  At $n=1$ the printed
title gives $64$, while the data, PARI program and cited paper give $0$.
We explicitly correct that title; the selected recurrence is unchanged.
Squaring and rationalizing the paper's displayed nested-radical GF gives
\[
 F(x)=\sqrt{\frac{2}{1+64x+\sqrt{1-128x}}}.
\]
With $x=t(1-32t)$ the inner square root is $1-64t$ and hence
$F=1/\sqrt{1-1024t^2}$.  Constant terms select the asserted branches.

\paragraph{Exact certificate.}

\[F(\xi(t))=\frac{1}{y}.\]

\[\xi(t)=t \left(1 - 32 t\right),\qquad \phi(t)=\frac{1}{1 - 32 t},\qquad y^2=1 - 1024 t^{2}.\]

The formal parameter satisfies $t=x\phi(t)$; every square root uses its branch with constant term one.

Substitution in \eqref{round8:r8:gf-rational}, followed by exact cancellation, gives

\[LF=P(x)=0.\]

The two rational residuals are $P(\xi)$ and $0$.  Their coefficient interpretation proves the recurrence for all $n\geq 3$; this is at least the source offset plus the recurrence order.

\paragraph{Independent consistency check.} The companion validation computes the defining finite sum, combinatorial dynamic program, or nonlinear grammar through index 47, without using the conjectured recurrence.  All listed source coefficients match.  These checks audit the transcription and branch; the exact identities above supply the proof.


\subsection{A344502: the selected recurrence}

\label{round8:r8entry:A344502}

\paragraph{Source and assertion.}

\href{https://oeis.org/A344502}{A344502}, consulted on 20 October 2026, explicitly marks the recurrence as a conjecture attributed to R. J. Mathar, Nov 02 2021.

The source offset is $0$.  With $R_n=\sum_{j=0}^{4}p_j(n)a_{n-j}$, the exact source coefficients are:

\[\begin{aligned}
p_{0}(n)&=- 4 \left(n + 1\right) \left(2 n + 1\right) \left(3035 n - 11997\right)\\
p_{1}(n)&=2 \left(153546 n^{3} - 490325 n^{2} - 62942 n + 71982\right)\\
p_{2}(n)&=- 2 \left(452090 n^{3} - 1745622 n^{2} + 1405285 n + 226200\right)\\
p_{3}(n)&=- 2 \left(n - 2\right) \left(83741 n^{2} - 200458 n - 19650\right)\\
p_{4}(n)&=- 3 \left(n - 3\right) \left(n - 2\right) \left(46423 n + 6938\right)\\
\end{aligned}\]

We prove $R_n=0$ for every integer $n\geq 4$.

\paragraph{Definition and generating-function derivation.}

Let $T(n,k)$ count nonnegative Motzkin prefixes of length $n$
ending at height $k$.  Reflection at the first visit to $-1$ gives
\[
 T(n,k)=[u^{n-k}](1-u^2)(1+u+u^2)^n.
\]
Expanding by the number $j$ of down steps gives
\[
 T(n,k)=\sum_{j\geq0}
 \frac{n!(k+1)}{j!(j+k+1)!(n-k-2j)!}.
\]
The terminating hypergeometric factor in the source, multiplied by
$\binom nk$, equals this sum, since
$4^j((k-n)/2)_j((k-n+1)/2)_j=(k-n)_{2j}$.
Therefore the source definition is $\sum_k\binom nkT(n,k)$, and the
binomial theorem proves
\[
 a_n=[u^n](1-u^2)\big((1+u)(1+u+u^2)\big)^n.
\]
Apply \eqref{round8:r8:gf-diagonal} with $B=1-t^2$ and
$\phi=(1+t)(1+t+t^2)$ to obtain the certificate seed.

\paragraph{Exact certificate.}

\[F(\xi(t))=\frac{1 - t^{2}}{- \frac{t \left(t^{2} + t + \left(t + 1\right) \left(2 t + 1\right) + 1\right)}{\left(t + 1\right) \left(t^{2} + t + 1\right)} + 1}.\]

\[\xi(t)=\frac{t}{\left(t + 1\right) \left(t^{2} + t + 1\right)},\qquad \phi(t)=\left(t + 1\right) \left(t^{2} + t + 1\right),\qquad y^2=1.\]

The formal parameter satisfies $t=x\phi(t)$; every square root uses its branch with constant term one.

Substitution in \eqref{round8:r8:gf-rational}, followed by exact cancellation, gives

\[LF=P(x)=47988 - 225302 x.\]

The two rational residuals are $P(\xi)$ and $0$.  Their coefficient interpretation proves the recurrence for all $n\geq 4$; this is at least the source offset plus the recurrence order.

\paragraph{Independent consistency check.} The companion validation computes the defining finite sum, combinatorial dynamic program, or nonlinear grammar through index 47, without using the conjectured recurrence.  All listed source coefficients match.  These checks audit the transcription and branch; the exact identities above supply the proof.


\subsection{A379463: the selected recurrence}

\label{round8:r8entry:A379463}

\paragraph{Source and assertion.}

\href{https://oeis.org/A379463}{A379463}, consulted on 20 October 2026, explicitly marks the recurrence as a conjecture attributed to R. J. Mathar, Jan 29 2025.

The source offset is $0$.  With $R_n=\sum_{j=0}^{5}p_j(n)a_{n-j}$, the exact source coefficients are:

\[\begin{aligned}
p_{0}(n)&=3 n \left(n - 3\right) \left(3 n + 4\right) \left(3 n + 8\right)\\
p_{1}(n)&=- 3 \left(45 n^{4} - 54 n^{3} - 192 n^{2} + 27 n + 20\right)\\
p_{2}(n)&=9 \left(n - 1\right) \left(30 n^{3} - 72 n^{2} - 7 n + 20\right)\\
p_{3}(n)&=- 3 \left(n - 2\right) \left(n - 1\right) \left(90 n^{2} - 234 n + 95\right)\\
p_{4}(n)&=- \left(n - 3\right) \left(n - 2\right) \left(n - 1\right) \left(121 n + 499\right)\\
p_{5}(n)&=229 \left(n - 4\right) \left(n - 3\right) \left(n - 2\right) \left(n - 1\right)\\
\end{aligned}\]

We prove $R_n=0$ for every integer $n\geq 5$.

\paragraph{Definition and generating-function derivation.}

First derive the finite sum from the path definition.  For
$r=2$ or $3$, remove the horizontal steps from a path with steps
$1,0,-3$, ending at zero and staying at least $-(r-1)$.
Its remaining $4k$ steps consist of $3k$ upward and $k$ downward steps.
Reverse their order, negate them, and start at height $r-1$.
This gives a nonnegative path with steps $-1,3$, ending at height $r-1$.
Append $r$ downward steps.  The resulting Lukasiewicz word encodes an
ordered forest of $r$ rooted four-ary trees with $k$ internal vertices:
each internal vertex contributes a step $3$ and each leaf a step $-1$.
Conversely, every nontrivial such forest ends with at least four leaf
steps, because its last internal vertex has four leaf children.
Deleting the final $r$ steps and reversing the transformation recovers
the original path.  The forest of $r$ isolated leaves gives the empty path.
The horizontal steps can be inserted in $\binom n{4k}$ ways.

For this entry $r=2$, so the defining path count has the finite sum
\[
 a_n=\sum_{k=0}^{\lfloor n/4\rfloor}
 \binom n{4k}\frac{2}{4k+2}\binom{4k+2}{k}.
\]
This is the same as Scheuerle's January 7, 2025 program, since
$\binom{4k+2}{k}/(4k+2)=\binom{4k+1}{k}/(3k+2)$.
Let $C=1+zC^4$.  Lagrange inversion gives
$[z^k]C^r=r\binom{4k+r}{k}/(4k+r)$.
The binomial insertion identity
$\sum_{n\geq m}\binom nm x^n=x^m/(1-x)^{m+1}$ therefore yields
$F=C((x/(1-x))^4)^2/(1-x)$.  Put $t=xC/(1-x)$.
Then $C=1+t^4$, $x=t/(1+t+t^4)$ and
$F=(1+t^4)(1+t+t^4)$, as used in the certificate.

\paragraph{Exact certificate.}

\[F(\xi(t))=\left(t^{4} + 1\right) \left(t^{4} + t + 1\right).\]

\[\xi(t)=\frac{t}{t^{4} + t + 1},\qquad \phi(t)=t^{4} + t + 1,\qquad y^2=1.\]

The formal parameter satisfies $t=x\phi(t)$; every square root uses its branch with constant term one.

Substitution in \eqref{round8:r8:gf-rational}, followed by exact cancellation, gives

\[LF=P(x)=0.\]

The two rational residuals are $P(\xi)$ and $0$.  Their coefficient interpretation proves the recurrence for all $n\geq 5$; this is at least the source offset plus the recurrence order.

\paragraph{Independent consistency check.} The companion validation computes the defining finite sum, combinatorial dynamic program, or nonlinear grammar through index 47, without using the conjectured recurrence.  All listed source coefficients match.  These checks audit the transcription and branch; the exact identities above supply the proof.


\subsection{A379464: the selected recurrence}

\label{round8:r8entry:A379464}

\paragraph{Source and assertion.}

\href{https://oeis.org/A379464}{A379464}, consulted on 20 October 2026, explicitly marks the recurrence as a conjecture attributed to R. J. Mathar, Jan 29 2025.

The source offset is $0$.  With $R_n=\sum_{j=0}^{7}p_j(n)a_{n-j}$, the exact source coefficients are:

\[\begin{aligned}
p_{0}(n)&=3 \left(n + 4\right) \left(3 n + 4\right) \left(3 n + 8\right)\\
p_{1}(n)&=- 3 \left(63 n^{3} + 297 n^{2} + 349 n + 60\right)\\
p_{2}(n)&=3 \left(189 n^{3} + 270 n^{2} - 229 n - 140\right)\\
p_{3}(n)&=- 15 \left(63 n^{3} - 117 n^{2} - 44 n + 64\right)\\
p_{4}(n)&=689 n^{3} - 5372 n^{2} + 6946 n - 1288\\
p_{5}(n)&=\left(n - 4\right) \left(201 n^{2} + 2767 n - 3011\right)\\
p_{6}(n)&=- \left(n - 5\right) \left(n - 4\right) \left(579 n + 257\right)\\
p_{7}(n)&=229 \left(n - 6\right) \left(n - 5\right) \left(n - 4\right)\\
\end{aligned}\]

We prove $R_n=0$ for every integer $n\geq 7$.

\paragraph{Definition and generating-function derivation.}

Apply the path-to-forest bijection just proved for A379463
with $r=3$: the defining lower boundary is $-2$, so deleting horizontal
steps, reversing and negating, shifting the starting height to two,
and appending three downward steps gives a forest of three four-ary trees.
Lagrange inversion and horizontal-step insertion give
\[
 a_n=\sum_{k=0}^{\lfloor n/4\rfloor}
 \binom n{4k}\frac{3}{4k+3}\binom{4k+3}{k}.
\]
For $C=1+zC^4$, Lagrange inversion gives the second factor as
$[z^k]C^3$.  Summing the binomial insertion factor gives
$F=C((x/(1-x))^4)^3/(1-x)$.  With $t=xC/(1-x)$ one obtains
$x=t/(1+t+t^4)$ and $F=(1+t^4)^2(1+t+t^4)$.
This differs from A379463 in its defining lower boundary and forest power;
it is not a shift or scalar multiple.

\paragraph{Exact certificate.}

\[F(\xi(t))=\left(t^{4} + 1\right)^{2} \left(t^{4} + t + 1\right).\]

\[\xi(t)=\frac{t}{t^{4} + t + 1},\qquad \phi(t)=t^{4} + t + 1,\qquad y^2=1.\]

The formal parameter satisfies $t=x\phi(t)$; every square root uses its branch with constant term one.

Substitution in \eqref{round8:r8:gf-rational}, followed by exact cancellation, gives

\[LF=P(x)=- 384 x^{3} + 1152 x^{2} - 1152 x + 384.\]

The two rational residuals are $P(\xi)$ and $0$.  Their coefficient interpretation proves the recurrence for all $n\geq 7$; this is at least the source offset plus the recurrence order.

\paragraph{Independent consistency check.} The companion validation computes the defining finite sum, combinatorial dynamic program, or nonlinear grammar through index 47, without using the conjectured recurrence.  All listed source coefficients match.  These checks audit the transcription and branch; the exact identities above supply the proof.

\section{Forty-nine empirical array recurrence packages}
\label{round8:sec:r8arrays}
This section proves forty-nine entry-level recurrence assertions that
are explicitly marked \emph{Empirical} on the primary OEIS pages
consulted on 20 October 2026. Seventeen source-stated empirical or
conjectural rational generating functions are proved in the same
packages. The arrays and empirical recurrences are credited to
R.~H.~Hardin; the supplementary generating functions are credited on
their pages to Colin Barker. One selected A-number contributes one
package, regardless of the number of formulas proved on that page.
All forty-nine A-numbers are absent from the 708 entries selected in
Rounds I--VII. These are freshly sourced and certified candidates;
none of the earlier rounds' uncounted reserves is promoted here.

\subsection{The local definitions and the counting convention}
An array has $n+1$ rows, width $w\ge2$, and entries in
$\{0,\ldots,q-1\}$. For a block
$B=\left(\begin{smallmatrix}a&b\\c&d\end{smallmatrix}\right)$,
write its four signed clockwise edge differences as
\[
 \Delta(B)=(b-a,d-b,c-d,a-c).
\]
The entry-specific parameters select one of three predicates:
\begin{description}
\item[Stress (D).] $|a+d-b-c|=s$.
\item[Absolute edge sum (U).] $\sum_{e\in\Delta(B)}|e|=s$.
\item[Squared edge sum (V).] $\sum_{e\in\Delta(B)}e^2=s$.
\end{description}
Both edge sums include all four edges with multiplicity. In entries
that additionally prohibit equal adjacent values, every component of
$\Delta(B)$ must be nonzero. Adjacency here is horizontal or vertical,
not diagonal. Every edge of an array with at least two rows and two
columns belongs to a $2\times2$ block, so this local test enforces the
source condition even on boundary edges.

All selected definitions count labeled arrays literally. Colors may be
absent; there is no requirement to use every alphabet symbol, no
first-occurrence renaming, and no division by an orbit size. The row
orbits below compress a counting computation, while retaining the full
labeled count through their multiplicities.

\subsection{Transfers obtained directly from the array definition}
Let $\gamma(r)$ be the lexicographically least of a row $r$, its
horizontal reversal, its simultaneous complement $q-1-r$, and the
reversal of that complement. Let $\mathcal O_{q,w}$ be the sorted set
of these representatives. Define the mass vector and transition matrix
by
\begin{align}
 m(r)&=\#\{u\in\{0,\ldots,q-1\}^w:\gamma(u)=r\},
 \label{round8:eq:r8arr-mass}\\
 T(r,t)&=\#\left\{u\in\{0,\ldots,q-1\}^w:
 \begin{array}{l}\gamma(u)=t,\\
 P(r_i,r_{i+1},u_i,u_{i+1})\quad(1\le i<w)
 \end{array}\right\}.
 \label{round8:eq:r8arr-T}
\end{align}
Reversal and complementation preserve each of D, U and V, as well as
the prohibition on zero edge differences. Applying either operation
simultaneously to a pair of rows preserves compatibility and permutes
each target orbit. Thus every row in a given orbit has the same number
of successors in every target orbit. Arbitrary permutations of the
numeric alphabet are not used.

Adding rows therefore gives the exact count
\begin{equation}\label{round8:eq:r8arr-count}
 a(n)=m^{\mathsf T}T^n\mathbf1\qquad(n\ge1).
\end{equation}
The exponent is $n$ because there are $n+1$ rows. In particular,
$\sum_r m(r)=q^w$; fixed rows, shorter orbits, and rows with no
successors retain their proper multiplicities. The auxiliary transfer
value at $n=0$ is not asserted to be a term of the source sequence,
whose offset is 1.

\subsection{A finite exact certificate proves every index}
For the source coefficient signature $c=(c_1,\ldots,c_d)$, put
\[
 R(z)=z^d-\sum_{i=1}^{d}c_i z^{d-i},\qquad
 Q(x)=1-\sum_{i=1}^{d}c_i x^i.
\]
The complete certificate for each selected definition establishes the
integer vector identity
\begin{equation}\label{round8:eq:r8arr-cert}
 \boxed{TR(T)\mathbf1=\mathbf0.}
\end{equation}
For $n\ge d+1$, multiply this identity by
$m^{\mathsf T}T^{n-d-1}$ and use \eqref{round8:eq:r8arr-count}. This proves
\begin{equation}\label{round8:eq:r8arr-rec}
 \boxed{a(n)=\sum_{i=1}^{d}c_i a(n-i)\qquad(n\ge d+1).}
\end{equation}
Every index on the right is at least 1. The proof does not assume an
empirical generating function or infer a recurrence from a numerical
prefix.

The forty-nine complete matrices contain 3,067 states in total and
5,879 nonzero weighted edges. The largest matrix has 625 states; the
largest recurrence order is 48. Every matrix, every orbit mass, and
every coordinate of $v_k=T^k\mathbf1$ for $0\le k\le d+1$ is
retained in the certificate data. The exact calculation is
\begin{Verbatim}[fontsize=\small]
v[0] = [1 for every state]
for k = 1,...,d+1:
    v[k][r] = sum(T[r,s]*v[k-1][s] for every state s)
for every state r:
    assert v[d+1][r] == sum(c[i]*v[d+1-i][r] for i=1,...,d)
\end{Verbatim}
The verifier rebuilds all states and transitions from the predicates,
then recomputes and compares every stored vector coordinate using
integer arithmetic. This is a finite computer-assisted proof of an
all-index assertion. The parameters and full coefficient lists below
specify the calculation independently of the data files; the expanded
matrices and vectors are supplied in
\texttt{r8\_arrays\_matrices.json} and
\texttt{r8\_arrays\_certificates.json}.

All 954 displayed source-prefix values agree with the defining
transfers. An additional program, which does not import the primary
engine, counts two-row and three-row arrays using unquotiented
column-by-column dynamic programming; all 98 values agree. These
checks validate the interpretation and offset. They are supplementary
to the exact vector identities that establish every later index.

For $F(x)=\sum_{n\ge1}a(n)x^n$, the proved recurrence implies
$QF=P$, where $\deg P\le d$ and
\begin{equation}\label{round8:eq:r8arr-numerator}
 [x^k]P=\sum_{\substack{0\le i\le\min(k,d)\\k-i\ge1}}
 q_i a(k-i),\qquad q_0=1,\quad q_i=-c_i\ (i\ge1).
\end{equation}
The initial values come from the defining transfer. Polynomial
cross-multiplication identifies this rational function exactly with all
seventeen selected source generating functions. Their numerators are
displayed below. Other formulas on a source page, such as radical
closed forms, are not additional selected assertions.

\subsection{A hand-checkable example: A234259}
For $q=3$, $w=2$, and stress $s=2$, the row representatives are
$00,01,02,11$. Direct enumeration from \eqref{round8:eq:r8arr-T} gives
\[
 T=\begin{pmatrix}
 0&0&2&0\\0&2&0&0\\2&0&0&1\\0&0&2&0
 \end{pmatrix},\qquad
 m=\begin{pmatrix}2\\4\\2\\1\end{pmatrix}.
\]
The vectors $v_k=T^k\mathbf1$, for $k=1,2,3,4$, are
\[
 (2,2,3,2)^{\mathsf T},\quad
 (6,4,6,6)^{\mathsf T},\quad
 (12,8,18,12)^{\mathsf T},\quad
 (36,16,36,36)^{\mathsf T}.
\]
They satisfy $v_4=2v_3+6v_2-12v_1$. Consequently
\[
 a(n)=2a(n-1)+6a(n-2)-12a(n-3)\qquad(n\ge4).
\]
The defining counts are $a(1)=20$, $a(2)=46$, $a(3)=104$, so
\[
 F(x)=\frac{20x+6x^2-108x^3}{1-2x-6x^2+12x^3}
     =\frac{2x(10+3x-54x^2)}{(1-2x)(1-6x^2)}.
\]
This proves the empirical recurrence and conjectural rational
generating function on \url{https://oeis.org/A234259}.

\subsection{Source status, prior work, and the selection unit}
Each selected primary page and each of the 25 associated parent-table
pages was retrieved and inspected for this round. All selected
recurrences are explicitly empirical; no proof link or established
encompassing recurrence was identified on those consulted pages.
Retrieval dates and the search tool's reported cache ages are recorded
separately in the source audit. A retrieval on the audit date does not
guarantee that an indexed page reflects every live edit.

Exact-ID and family searches did not identify a prior proof of a
selected assertion among the inspected results. Several searches
returned irrelevant matches, so this is a bounded audit, not a claim
of worldwide priority. The inspected README of the additional
\url{https://github.com/Eppie/oeis-proofs} lead lists no selected array
entry; its entire repository was not audited.

The IDs are also disjoint from the retained 97 distinct A-numbers in
the 100 costly claims of Appendix A of \emph{OEIS Open}. This comparison
does not cover every result in that project's repository. The method
here adapts the geometric-orbit transfers used in Rounds IV--VII,
with a new definition-derived matrix certificate for every selected
entry.

A separate exact comparison with 236 previously certified array
sequences from Rounds II--VII and between these forty-nine sequences
tests twenty values for each shift from $-3$ through $3$, allowing an
arbitrary rational scale. Every tested proportional equality has a
mismatch within the first three compared values. This rules out those
specific shift-and-scale aliases; unbounded shifts, affine changes,
and other transformations are not covered. A shared recurrence
polynomial alone does not make two different counting sequences the
same selected claim.

One further fresh certificate, A235252, is
retained as an uncounted reserve. It is not a failed proof and does not
contribute to the forty-nine selected packages. No OEIS editing,
submission, official closure, or worldwide-first claim is made.

\subsection{Individual definitions and complete coefficient signatures}
Every list below gives $c_1,\ldots,c_d$ in order, including zeros,
for \eqref{round8:eq:r8arr-rec}. Initial values start at $n=1$.
\Needspace{7\baselineskip}
\subsubsection{A233646}
\label{round8:r8entry:A233646}
Arrays have $n+1$ rows, width $w=2$, and labeled alphabet $0,\ldots,5$. Predicate V has target $s=14$, allowing equal adjacent values.
The recurrence holds for every $n\ge 10$, with order $d=9$ and a complete $D=12$ state matrix. The first three source values are $(144,772,3984)$.
\begin{Verbatim}[fontsize=\scriptsize,breaklines=true]
2, 31, -36, -261, 206, 714, -356, -528, 256
\end{Verbatim}
The defining transfer satisfies $TR(T)\mathbf1=0$ in every coordinate, so \eqref{round8:eq:r8arr-rec} follows. Source: \url{https://oeis.org/A233646}; parent column 1 of \url{https://oeis.org/A233653}. Both consulted pages mark the relevant recurrence empirical.
\Needspace{7\baselineskip}
\subsubsection{A233647}
\label{round8:r8entry:A233647}
Arrays have $n+1$ rows, width $w=3$, and labeled alphabet $0,\ldots,5$. Predicate V has target $s=14$, allowing equal adjacent values.
The recurrence holds for every $n\ge 39$, with order $d=38$ and a complete $D=63$ state matrix. The first three source values are $(772,5232,36784)$.
\begin{Verbatim}[fontsize=\scriptsize,breaklines=true]
13, 60, -1403, 915, 60037, -161542, -1305952, 5581957, 14886303, -99457303, -65691409,
1070532941, -439804539, -7350827624, 8431876793, 32422650909, -59093935245,
-87614113660, 242031773814, 116489014001, -628983746855, 50730140495, 1046367995789,
-482503791934, -1075303683128, 846072582698, 610370785910, -745472433880,
-117896421864, 344109162632, -46130175152, -71419513632, 22452177920, 3530443776,
-1777822976, 30745088, 22653952, -933888
\end{Verbatim}
The defining transfer satisfies $TR(T)\mathbf1=0$ in every coordinate, so \eqref{round8:eq:r8arr-rec} follows. Source: \url{https://oeis.org/A233647}; parent column 2 of \url{https://oeis.org/A233653}. Both consulted pages mark the relevant recurrence empirical.
\Needspace{7\baselineskip}
\subsubsection{A233675}
\label{round8:r8entry:A233675}
Arrays have $n+1$ rows, width $w=2$, and labeled alphabet $0,\ldots,3$. Predicate V has target $s=6$, allowing equal adjacent values.
The recurrence holds for every $n\ge 5$, with order $d=4$ and a complete $D=6$ state matrix. The first three source values are $(48,184,648)$.
\[(c_1,\ldots,c_d)=(1,12,-4,-16).\]
The defining transfer satisfies $TR(T)\mathbf1=0$ in every coordinate, so \eqref{round8:eq:r8arr-rec} follows. Source: \url{https://oeis.org/A233675}; parent column 1 of \url{https://oeis.org/A233682}. Both consulted pages mark the relevant recurrence empirical.
The source rational generating function is also proved: $F=P/Q$, with $Q=1-\sum_{i=1}^{d}c_i x^i$ and
\[P(x)=48x+136x^{2}-112x^{3}-224x^{4}.\]
\Needspace{7\baselineskip}
\subsubsection{A233676}
\label{round8:r8entry:A233676}
Arrays have $n+1$ rows, width $w=3$, and labeled alphabet $0,\ldots,3$. Predicate V has target $s=6$, allowing equal adjacent values.
The recurrence holds for every $n\ge 11$, with order $d=10$ and a complete $D=20$ state matrix. The first three source values are $(184,928,4448)$.
\begin{Verbatim}[fontsize=\scriptsize,breaklines=true]
7, 11, -157, 202, 554, -1224, -112, 1584, -928, 64
\end{Verbatim}
The defining transfer satisfies $TR(T)\mathbf1=0$ in every coordinate, so \eqref{round8:eq:r8arr-rec} follows. Source: \url{https://oeis.org/A233676}; parent column 2 of \url{https://oeis.org/A233682}. Both consulted pages mark the relevant recurrence empirical.
\Needspace{7\baselineskip}
\subsubsection{A233677}
\label{round8:r8entry:A233677}
Arrays have $n+1$ rows, width $w=4$, and labeled alphabet $0,\ldots,3$. Predicate V has target $s=6$, allowing equal adjacent values.
The recurrence holds for every $n\ge 18$, with order $d=17$ and a complete $D=72$ state matrix. The first three source values are $(648,4448,28544)$.
\begin{Verbatim}[fontsize=\scriptsize,breaklines=true]
5, 70, -332, -1384, 6780, 11520, -61568, -45216, 280096, 84864, -653696, -81664,
762880, 41984, -397312, -8192, 65536
\end{Verbatim}
The defining transfer satisfies $TR(T)\mathbf1=0$ in every coordinate, so \eqref{round8:eq:r8arr-rec} follows. Source: \url{https://oeis.org/A233677}; parent column 3 of \url{https://oeis.org/A233682}. Both consulted pages mark the relevant recurrence empirical.
\Needspace{7\baselineskip}
\subsubsection{A233687}
\label{round8:r8entry:A233687}
Arrays have $n+1$ rows, width $w=5$, and labeled alphabet $0,\ldots,3$. Predicate V has target $s=10$, with every horizontal and vertical neighboring pair unequal.
The recurrence holds for every $n\ge 42$, with order $d=41$ and a complete $D=272$ state matrix. The first three source values are $(1188,13324,145512)$.
\begin{Verbatim}[fontsize=\scriptsize,breaklines=true]
4, 222, -895, -14806, 61290, 479263, -2067766, -8897249, 40799029, 102512949,
-513857547, -760171900, 4350919810, 3628256006, -25577829290, -10454924632,
106480049675, 12586615989, -317308345201, 27167936641, 679221981450, -157201990004,
-1040817140533, 351215509546, 1129423930543, -469396649218, -850891317254,
403441486513, 431368323585, -223830914820, -140453027219, 77851050483, 27476109266,
-16042803856, -2950519700, 1797075893, 152255054, -93497323, -3429989, 1518688, 74280
\end{Verbatim}
The defining transfer satisfies $TR(T)\mathbf1=0$ in every coordinate, so \eqref{round8:eq:r8arr-rec} follows. Source: \url{https://oeis.org/A233687}; parent column 4 of \url{https://oeis.org/A233691}. Both consulted pages mark the relevant recurrence empirical.
\Needspace{7\baselineskip}
\subsubsection{A233718}
\label{round8:r8entry:A233718}
Arrays have $n+1$ rows, width $w=2$, and labeled alphabet $0,\ldots,7$. Predicate V has target $s=26$, allowing equal adjacent values.
The recurrence holds for every $n\ge 17$, with order $d=16$ and a complete $D=20$ state matrix. The first three source values are $(280,1820,11016)$.
\begin{Verbatim}[fontsize=\scriptsize,breaklines=true]
2, 58, -58, -1227, 461, 12296, 722, -62278, -22804, 155159, 90798, -162636, -121112,
35968, 28416, -1152
\end{Verbatim}
The defining transfer satisfies $TR(T)\mathbf1=0$ in every coordinate, so \eqref{round8:eq:r8arr-rec} follows. Source: \url{https://oeis.org/A233718}; parent column 1 of \url{https://oeis.org/A233724}. Both consulted pages mark the relevant recurrence empirical.
\Needspace{7\baselineskip}
\subsubsection{A233726}
\label{round8:r8entry:A233726}
Arrays have $n+1$ rows, width $w=2$, and labeled alphabet $0,\ldots,3$. Predicate U has target $s=6$, with every horizontal and vertical neighboring pair unequal.
The recurrence holds for every $n\ge 5$, with order $d=4$ and a complete $D=6$ state matrix. The first three source values are $(40,148,556)$.
\[(c_1,\ldots,c_d)=(3,5,-7,-2).\]
The defining transfer satisfies $TR(T)\mathbf1=0$ in every coordinate, so \eqref{round8:eq:r8arr-rec} follows. Source: \url{https://oeis.org/A233726}; parent column 1 of \url{https://oeis.org/A233733}. Both consulted pages mark the relevant recurrence empirical.
The source rational generating function is also proved: $F=P/Q$, with $Q=1-\sum_{i=1}^{d}c_i x^i$ and
\[P(x)=40x+28x^{2}-88x^{3}-24x^{4}.\]
\Needspace{7\baselineskip}
\subsubsection{A233727}
\label{round8:r8entry:A233727}
Arrays have $n+1$ rows, width $w=3$, and labeled alphabet $0,\ldots,3$. Predicate U has target $s=6$, with every horizontal and vertical neighboring pair unequal.
The recurrence holds for every $n\ge 13$, with order $d=12$ and a complete $D=20$ state matrix. The first three source values are $(148,756,3972)$.
\begin{Verbatim}[fontsize=\scriptsize,breaklines=true]
3, 30, -71, -192, 415, 410, -847, -269, 544, 38, -42, -4
\end{Verbatim}
The defining transfer satisfies $TR(T)\mathbf1=0$ in every coordinate, so \eqref{round8:eq:r8arr-rec} follows. Source: \url{https://oeis.org/A233727}; parent column 2 of \url{https://oeis.org/A233733}. Both consulted pages mark the relevant recurrence empirical.
\Needspace{7\baselineskip}
\subsubsection{A233728}
\label{round8:r8entry:A233728}
Arrays have $n+1$ rows, width $w=4$, and labeled alphabet $0,\ldots,3$. Predicate U has target $s=6$, with every horizontal and vertical neighboring pair unequal.
The recurrence holds for every $n\ge 29$, with order $d=28$ and a complete $D=72$ state matrix. The first three source values are $(556,3972,29816)$.
\begin{Verbatim}[fontsize=\scriptsize,breaklines=true]
9, 46, -508, -481, 10680, -3945, -112548, 115263, 657688, -958275, -2192200, 4058993,
4067658, -9776097, -3651898, 13759514, 407041, -11221093, 1862871, 5154086, -1437504,
-1269561, 444011, 147392, -60418, -4972, 2810, -123
\end{Verbatim}
The defining transfer satisfies $TR(T)\mathbf1=0$ in every coordinate, so \eqref{round8:eq:r8arr-rec} follows. Source: \url{https://oeis.org/A233728}; parent column 3 of \url{https://oeis.org/A233733}. Both consulted pages mark the relevant recurrence empirical.
\Needspace{7\baselineskip}
\subsubsection{A233942}
\label{round8:r8entry:A233942}
Arrays have $n+1$ rows, width $w=2$, and labeled alphabet $0,\ldots,2$. Predicate U has target $s=4$, allowing equal adjacent values.
The recurrence holds for every $n\ge 5$, with order $d=4$ and a complete $D=4$ state matrix. The first three source values are $(44,226,1150)$.
\[(c_1,\ldots,c_d)=(5,3,-13,2).\]
The defining transfer satisfies $TR(T)\mathbf1=0$ in every coordinate, so \eqref{round8:eq:r8arr-rec} follows. Source: \url{https://oeis.org/A233942}; parent column 1 of \url{https://oeis.org/A233949}. Both consulted pages mark the relevant recurrence empirical.
The source rational generating function is also proved: $F=P/Q$, with $Q=1-\sum_{i=1}^{d}c_i x^i$ and
\[P(x)=44x+6x^{2}-112x^{3}+18x^{4}.\]
\Needspace{7\baselineskip}
\subsubsection{A233943}
\label{round8:r8entry:A233943}
Arrays have $n+1$ rows, width $w=3$, and labeled alphabet $0,\ldots,2$. Predicate U has target $s=4$, allowing equal adjacent values.
The recurrence holds for every $n\ge 11$, with order $d=10$ and a complete $D=10$ state matrix. The first three source values are $(226,2108,19484)$.
\begin{Verbatim}[fontsize=\scriptsize,breaklines=true]
10, 9, -154, 24, 569, -305, -372, 163, 31, -6
\end{Verbatim}
The defining transfer satisfies $TR(T)\mathbf1=0$ in every coordinate, so \eqref{round8:eq:r8arr-rec} follows. Source: \url{https://oeis.org/A233943}; parent column 2 of \url{https://oeis.org/A233949}. Both consulted pages mark the relevant recurrence empirical.
\Needspace{7\baselineskip}
\subsubsection{A233944}
\label{round8:r8entry:A233944}
Arrays have $n+1$ rows, width $w=4$, and labeled alphabet $0,\ldots,2$. Predicate U has target $s=4$, allowing equal adjacent values.
The recurrence holds for every $n\ge 26$, with order $d=25$ and a complete $D=25$ state matrix. The first three source values are $(1150,19484,328508)$.
\begin{Verbatim}[fontsize=\scriptsize,breaklines=true]
22, 20, -2402, 8315, 57619, -297894, -394397, 3586512, -630788, -18290282, 13131067,
44278092, -40653252, -53889205, 50657365, 32681066, -27199613, -9122229, 6034975,
846023, -477739, -19115, 11259, -26, -40
\end{Verbatim}
The defining transfer satisfies $TR(T)\mathbf1=0$ in every coordinate, so \eqref{round8:eq:r8arr-rec} follows. Source: \url{https://oeis.org/A233944}; parent column 3 of \url{https://oeis.org/A233949}. Both consulted pages mark the relevant recurrence empirical.
\Needspace{7\baselineskip}
\subsubsection{A234241}
\label{round8:r8entry:A234241}
Arrays have $n+1$ rows, width $w=2$, and labeled alphabet $0,\ldots,7$. Predicate D has target $s=2$, allowing equal adjacent values.
The recurrence holds for every $n\ge 9$, with order $d=8$ and a complete $D=20$ state matrix. The first three source values are $(630,6706,72616)$.
\begin{Verbatim}[fontsize=\scriptsize,breaklines=true]
7, 181, -1022, -7657, 21504, 54144, -96768, -80640
\end{Verbatim}
The defining transfer satisfies $TR(T)\mathbf1=0$ in every coordinate, so \eqref{round8:eq:r8arr-rec} follows. Source: \url{https://oeis.org/A234241}; parent column 1 of \url{https://oeis.org/A234245}. Both consulted pages mark the relevant recurrence empirical.
\Needspace{7\baselineskip}
\subsubsection{A234259}
\label{round8:r8entry:A234259}
Arrays have $n+1$ rows, width $w=2$, and labeled alphabet $0,\ldots,2$. Predicate D has target $s=2$, allowing equal adjacent values.
The recurrence holds for every $n\ge 4$, with order $d=3$ and a complete $D=4$ state matrix. The first three source values are $(20,46,104)$.
\[(c_1,\ldots,c_d)=(2,6,-12).\]
The defining transfer satisfies $TR(T)\mathbf1=0$ in every coordinate, so \eqref{round8:eq:r8arr-rec} follows. Source: \url{https://oeis.org/A234259}; parent column 1 of \url{https://oeis.org/A234266}. Both consulted pages mark the relevant recurrence empirical.
The source rational generating function is also proved: $F=P/Q$, with $Q=1-\sum_{i=1}^{d}c_i x^i$ and
\[P(x)=20x+6x^{2}-108x^{3}.\]
\Needspace{7\baselineskip}
\subsubsection{A234260}
\label{round8:r8entry:A234260}
Arrays have $n+1$ rows, width $w=3$, and labeled alphabet $0,\ldots,2$. Predicate D has target $s=2$, allowing equal adjacent values.
The recurrence holds for every $n\ge 7$, with order $d=6$ and a complete $D=10$ state matrix. The first three source values are $(46,88,170)$.
\[(c_1,\ldots,c_d)=(3,6,-24,4,36,-24).\]
The defining transfer satisfies $TR(T)\mathbf1=0$ in every coordinate, so \eqref{round8:eq:r8arr-rec} follows. Source: \url{https://oeis.org/A234260}; parent column 2 of \url{https://oeis.org/A234266}. Both consulted pages mark the relevant recurrence empirical.
The source rational generating function is also proved: $F=P/Q$, with $Q=1-\sum_{i=1}^{d}c_i x^i$ and
\[P(x)=46x-50x^{2}-370x^{3}+424x^{4}+588x^{5}-648x^{6}.\]
\Needspace{7\baselineskip}
\subsubsection{A234261}
\label{round8:r8entry:A234261}
Arrays have $n+1$ rows, width $w=4$, and labeled alphabet $0,\ldots,2$. Predicate D has target $s=2$, allowing equal adjacent values.
The recurrence holds for every $n\ge 7$, with order $d=6$ and a complete $D=25$ state matrix. The first three source values are $(104,170,292)$.
\[(c_1,\ldots,c_d)=(3,6,-24,4,36,-24).\]
The defining transfer satisfies $TR(T)\mathbf1=0$ in every coordinate, so \eqref{round8:eq:r8arr-rec} follows. Source: \url{https://oeis.org/A234261}; parent column 3 of \url{https://oeis.org/A234266}. Both consulted pages mark the relevant recurrence empirical.
The source rational generating function is also proved: $F=P/Q$, with $Q=1-\sum_{i=1}^{d}c_i x^i$ and
\[P(x)=104x-142x^{2}-842x^{3}+1160x^{4}+1332x^{5}-1752x^{6}.\]
\Needspace{7\baselineskip}
\subsubsection{A234326}
\label{round8:r8entry:A234326}
Arrays have $n+1$ rows, width $w=2$, and labeled alphabet $0,\ldots,3$. Predicate D has target $s=1$, allowing equal adjacent values.
The recurrence holds for every $n\ge 5$, with order $d=4$ and a complete $D=6$ state matrix. The first three source values are $(80,432,2368)$.
\[(c_1,\ldots,c_d)=(0,32,0,-48).\]
The defining transfer satisfies $TR(T)\mathbf1=0$ in every coordinate, so \eqref{round8:eq:r8arr-rec} follows. Source: \url{https://oeis.org/A234326}; parent column 1 of \url{https://oeis.org/A234333}. Both consulted pages mark the relevant recurrence empirical.
The source rational generating function is also proved: $F=P/Q$, with $Q=1-\sum_{i=1}^{d}c_i x^i$ and
\[P(x)=80x+432x^{2}-192x^{3}-768x^{4}.\]
\Needspace{7\baselineskip}
\subsubsection{A234327}
\label{round8:r8entry:A234327}
Arrays have $n+1$ rows, width $w=3$, and labeled alphabet $0,\ldots,3$. Predicate D has target $s=1$, allowing equal adjacent values.
The recurrence holds for every $n\ge 12$, with order $d=11$ and a complete $D=20$ state matrix. The first three source values are $(432,3296,25728)$.
\begin{Verbatim}[fontsize=\scriptsize,breaklines=true]
10, 59, -768, 608, 7580, -11628, -18088, 32880, 3232, -17600, 4224
\end{Verbatim}
The defining transfer satisfies $TR(T)\mathbf1=0$ in every coordinate, so \eqref{round8:eq:r8arr-rec} follows. Source: \url{https://oeis.org/A234327}; parent column 2 of \url{https://oeis.org/A234333}. Both consulted pages mark the relevant recurrence empirical.
\Needspace{7\baselineskip}
\subsubsection{A234328}
\label{round8:r8entry:A234328}
Arrays have $n+1$ rows, width $w=4$, and labeled alphabet $0,\ldots,3$. Predicate D has target $s=1$, allowing equal adjacent values.
The recurrence holds for every $n\ge 49$, with order $d=48$ and a complete $D=72$ state matrix. The first three source values are $(2368,25728,287888)$.
\begin{Verbatim}[fontsize=\scriptsize,breaklines=true]
0, 507, 0, -102295, 0, 10646830, 0, -635915377, 0, 23832852746, 0, -596280829182, 0,
10371633339873, 0, -128744093807520, 0, 1159671655063068, 0, -7660667577596808, 0,
37368059037902308, 0, -135219186772494408, 0, 364034788822163472, 0,
-730004249849579520, 0, 1089413122161376064, 0, -1205481757762071168, 0,
981742220067000064, 0, -581000394777277440, 0, 244895745327540224, 0,
-71316845470162944, 0, 13706205520330752, 0, -1618885433491456, 0, 103796813856768, 0,
-2721872609280
\end{Verbatim}
The defining transfer satisfies $TR(T)\mathbf1=0$ in every coordinate, so \eqref{round8:eq:r8arr-rec} follows. Source: \url{https://oeis.org/A234328}; parent column 3 of \url{https://oeis.org/A234333}. Both consulted pages mark the relevant recurrence empirical.
\Needspace{7\baselineskip}
\subsubsection{A234338}
\label{round8:r8entry:A234338}
Arrays have $n+1$ rows, width $w=2$, and labeled alphabet $0,\ldots,7$. Predicate U has target $s=14$, allowing equal adjacent values.
The recurrence holds for every $n\ge 14$, with order $d=13$ and a complete $D=20$ state matrix. The first three source values are $(604,9076,120838)$.
\begin{Verbatim}[fontsize=\scriptsize,breaklines=true]
11, 100, -657, -3362, 13855, 47579, -118877, -254743, 334913, 167982, -130132, -15912,
10008
\end{Verbatim}
The defining transfer satisfies $TR(T)\mathbf1=0$ in every coordinate, so \eqref{round8:eq:r8arr-rec} follows. Source: \url{https://oeis.org/A234338}; parent column 1 of \url{https://oeis.org/A234343}. Both consulted pages mark the relevant recurrence empirical.
\Needspace{7\baselineskip}
\subsubsection{A234444}
\label{round8:r8entry:A234444}
Arrays have $n+1$ rows, width $w=2$, and labeled alphabet $0,\ldots,4$. Predicate D has target $s=1$, allowing equal adjacent values.
The recurrence holds for every $n\ge 5$, with order $d=4$ and a complete $D=9$ state matrix. The first three source values are $(160,1120,8064)$.
\[(c_1,\ldots,c_d)=(0,60,0,-344).\]
The defining transfer satisfies $TR(T)\mathbf1=0$ in every coordinate, so \eqref{round8:eq:r8arr-rec} follows. Source: \url{https://oeis.org/A234444}; parent column 1 of \url{https://oeis.org/A234450}. Both consulted pages mark the relevant recurrence empirical.
The source rational generating function is also proved: $F=P/Q$, with $Q=1-\sum_{i=1}^{d}c_i x^i$ and
\[P(x)=160x+1120x^{2}-1536x^{3}-8480x^{4}.\]
\Needspace{7\baselineskip}
\subsubsection{A234445}
\label{round8:r8entry:A234445}
Arrays have $n+1$ rows, width $w=3$, and labeled alphabet $0,\ldots,4$. Predicate D has target $s=1$, allowing equal adjacent values.
The recurrence holds for every $n\ge 18$, with order $d=17$ and a complete $D=39$ state matrix. The first three source values are $(1120,11574,125080)$.
\begin{Verbatim}[fontsize=\scriptsize,breaklines=true]
16, 122, -2854, 3035, 114162, -324660, -1581792, 6105836, 6998392, -41646592, 4188416,
100180352, -55355328, -54337536, 19902720, 7110912, -1815552
\end{Verbatim}
The defining transfer satisfies $TR(T)\mathbf1=0$ in every coordinate, so \eqref{round8:eq:r8arr-rec} follows. Source: \url{https://oeis.org/A234445}; parent column 2 of \url{https://oeis.org/A234450}. Both consulted pages mark the relevant recurrence empirical.
\Needspace{7\baselineskip}
\subsubsection{A234452}
\label{round8:r8entry:A234452}
Arrays have $n+1$ rows, width $w=2$, and labeled alphabet $0,\ldots,5$. Predicate D has target $s=1$, allowing equal adjacent values.
The recurrence holds for every $n\ge 7$, with order $d=6$ and a complete $D=12$ state matrix. The first three source values are $(280,2400,21344)$.
\[(c_1,\ldots,c_d)=(0,100,0,-1384,0,1440).\]
The defining transfer satisfies $TR(T)\mathbf1=0$ in every coordinate, so \eqref{round8:eq:r8arr-rec} follows. Source: \url{https://oeis.org/A234452}; parent column 1 of \url{https://oeis.org/A234458}. Both consulted pages mark the relevant recurrence empirical.
The source rational generating function is also proved: $F=P/Q$, with $Q=1-\sum_{i=1}^{d}c_i x^i$ and
\[P(x)=280x+2400x^{2}-6656x^{3}-47040x^{4}+8640x^{5}+51840x^{6}.\]
\Needspace{7\baselineskip}
\subsubsection{A234453}
\label{round8:r8entry:A234453}
Arrays have $n+1$ rows, width $w=3$, and labeled alphabet $0,\ldots,5$. Predicate D has target $s=1$, allowing equal adjacent values.
The recurrence holds for every $n\ge 27$, with order $d=26$ and a complete $D=63$ state matrix. The first three source values are $(2400,31112,427956)$.
\begin{Verbatim}[fontsize=\scriptsize,breaklines=true]
24, 202, -8375, 18366, 832456, -4372543, -32845569, 252119346, 475222218, -6386106484,
1228195800, 78155524488, -93005666672, -459032184864, 814396694400, 1253941668864,
-2968493257344, -1338994735872, 5126543725056, -147546633728, -4053307347968,
1013644265472, 1200673271808, -333045153792, -113446158336, 22613852160
\end{Verbatim}
The defining transfer satisfies $TR(T)\mathbf1=0$ in every coordinate, so \eqref{round8:eq:r8arr-rec} follows. Source: \url{https://oeis.org/A234453}; parent column 2 of \url{https://oeis.org/A234458}. Both consulted pages mark the relevant recurrence empirical.
\Needspace{7\baselineskip}
\subsubsection{A234483}
\label{round8:r8entry:A234483}
Arrays have $n+1$ rows, width $w=2$, and labeled alphabet $0,\ldots,3$. Predicate D has target $s=2$, allowing equal adjacent values.
The recurrence holds for every $n\ge 4$, with order $d=3$ and a complete $D=6$ state matrix. The first three source values are $(62,242,944)$.
\[(c_1,\ldots,c_d)=(7,-9,-12).\]
The defining transfer satisfies $TR(T)\mathbf1=0$ in every coordinate, so \eqref{round8:eq:r8arr-rec} follows. Source: \url{https://oeis.org/A234483}; parent column 1 of \url{https://oeis.org/A234490}. Both consulted pages mark the relevant recurrence empirical.
The source rational generating function is also proved: $F=P/Q$, with $Q=1-\sum_{i=1}^{d}c_i x^i$ and
\[P(x)=62x-192x^{2}-192x^{3}.\]
\Needspace{7\baselineskip}
\subsubsection{A234484}
\label{round8:r8entry:A234484}
Arrays have $n+1$ rows, width $w=3$, and labeled alphabet $0,\ldots,3$. Predicate D has target $s=2$, allowing equal adjacent values.
The recurrence holds for every $n\ge 6$, with order $d=5$ and a complete $D=20$ state matrix. The first three source values are $(242,926,3540)$.
\[(c_1,\ldots,c_d)=(10,-28,1,54,24).\]
The defining transfer satisfies $TR(T)\mathbf1=0$ in every coordinate, so \eqref{round8:eq:r8arr-rec} follows. Source: \url{https://oeis.org/A234484}; parent column 2 of \url{https://oeis.org/A234490}. Both consulted pages mark the relevant recurrence empirical.
The source rational generating function is also proved: $F=P/Q$, with $Q=1-\sum_{i=1}^{d}c_i x^i$ and
\[P(x)=242x-1494x^{2}+1056x^{3}+3840x^{4}+1536x^{5}.\]
\Needspace{7\baselineskip}
\subsubsection{A234491}
\label{round8:r8entry:A234491}
Arrays have $n+1$ rows, width $w=2$, and labeled alphabet $0,\ldots,5$. Predicate D has target $s=2$, allowing equal adjacent values.
The recurrence holds for every $n\ge 6$, with order $d=5$ and a complete $D=12$ state matrix. The first three source values are $(250,1834,13464)$.
\[(c_1,\ldots,c_d)=(5,74,-295,-1008,840).\]
The defining transfer satisfies $TR(T)\mathbf1=0$ in every coordinate, so \eqref{round8:eq:r8arr-rec} follows. Source: \url{https://oeis.org/A234491}; parent column 1 of \url{https://oeis.org/A234497}. Both consulted pages mark the relevant recurrence empirical.
The source rational generating function is also proved: $F=P/Q$, with $Q=1-\sum_{i=1}^{d}c_i x^i$ and
\[P(x)=250x+584x^{2}-14206x^{3}-30144x^{4}+29520x^{5}.\]
\Needspace{7\baselineskip}
\subsubsection{A234492}
\label{round8:r8entry:A234492}
Arrays have $n+1$ rows, width $w=3$, and labeled alphabet $0,\ldots,5$. Predicate D has target $s=2$, allowing equal adjacent values.
The recurrence holds for every $n\ge 24$, with order $d=23$ and a complete $D=63$ state matrix. The first three source values are $(1834,17022,158288)$.
\begin{Verbatim}[fontsize=\scriptsize,breaklines=true]
27, -7, -5642, 38876, 328643, -3991917, -1314423, 134764409, -274870060, -1754310602,
6188057089, 6383741530, -45869779722, 19646271368, 128555834088, -144208395792,
-115200081792, 233207698368, -26459811072, -102846620160, 35034937344, 13514784768,
-5918883840
\end{Verbatim}
The defining transfer satisfies $TR(T)\mathbf1=0$ in every coordinate, so \eqref{round8:eq:r8arr-rec} follows. Source: \url{https://oeis.org/A234492}; parent column 2 of \url{https://oeis.org/A234497}. Both consulted pages mark the relevant recurrence empirical.
\Needspace{7\baselineskip}
\subsubsection{A234674}
\label{round8:r8entry:A234674}
Arrays have $n+1$ rows, width $w=2$, and labeled alphabet $0,\ldots,5$. Predicate D has target $s=5$, allowing equal adjacent values.
The recurrence holds for every $n\ge 5$, with order $d=4$ and a complete $D=12$ state matrix. The first three source values are $(112,404,1264)$.
\[(c_1,\ldots,c_d)=(0,22,0,-120).\]
The defining transfer satisfies $TR(T)\mathbf1=0$ in every coordinate, so \eqref{round8:eq:r8arr-rec} follows. Source: \url{https://oeis.org/A234674}; parent column 1 of \url{https://oeis.org/A234681}. Both consulted pages mark the relevant recurrence empirical.
The source rational generating function is also proved: $F=P/Q$, with $Q=1-\sum_{i=1}^{d}c_i x^i$ and
\[P(x)=112x+404x^{2}-1200x^{3}-4320x^{4}.\]
\Needspace{7\baselineskip}
\subsubsection{A234675}
\label{round8:r8entry:A234675}
Arrays have $n+1$ rows, width $w=3$, and labeled alphabet $0,\ldots,5$. Predicate D has target $s=5$, allowing equal adjacent values.
The recurrence holds for every $n\ge 19$, with order $d=18$ and a complete $D=63$ state matrix. The first three source values are $(404,1152,2924)$.
\begin{Verbatim}[fontsize=\scriptsize,breaklines=true]
4, 47, -200, -901, 4204, 8957, -48440, -47194, 334096, 101948, -1410080, 170136,
3549696, -1446912, -4861440, 2954880, 2764800, -2073600
\end{Verbatim}
The defining transfer satisfies $TR(T)\mathbf1=0$ in every coordinate, so \eqref{round8:eq:r8arr-rec} follows. Source: \url{https://oeis.org/A234675}; parent column 2 of \url{https://oeis.org/A234681}. Both consulted pages mark the relevant recurrence empirical.
\Needspace{7\baselineskip}
\subsubsection{A234884}
\label{round8:r8entry:A234884}
Arrays have $n+1$ rows, width $w=3$, and labeled alphabet $0,\ldots,5$. Predicate D has target $s=1$, with every horizontal and vertical neighboring pair unequal.
The recurrence holds for every $n\ge 28$, with order $d=27$ and a complete $D=63$ state matrix. The first three source values are $(272,980,3164)$.
\begin{Verbatim}[fontsize=\scriptsize,breaklines=true]
2, 46, -91, -852, 1670, 8364, -16318, -48531, 94856, 177211, -349287, -422181, 844527,
666796, -1362305, -697540, 1464198, 475379, -1029412, -204338, 455683, 52461, -118912,
-7142, 16187, 376, -846
\end{Verbatim}
The defining transfer satisfies $TR(T)\mathbf1=0$ in every coordinate, so \eqref{round8:eq:r8arr-rec} follows. Source: \url{https://oeis.org/A234884}; parent column 2 of \url{https://oeis.org/A234890}. Both consulted pages mark the relevant recurrence empirical.
\Needspace{7\baselineskip}
\subsubsection{A234885}
\label{round8:r8entry:A234885}
Arrays have $n+1$ rows, width $w=4$, and labeled alphabet $0,\ldots,5$. Predicate D has target $s=1$, with every horizontal and vertical neighboring pair unequal.
The recurrence holds for every $n\ge 47$, with order $d=46$ and a complete $D=342$ state matrix. The first three source values are $(856,3164,10640)$.
\begin{Verbatim}[fontsize=\scriptsize,breaklines=true]
2, 86, -171, -3254, 6420, 72208, -140996, -1058026, 2037844, 10934873, -20686144,
-83167418, 153691479, 479339096, -859425200, -2137842471, 3687468872, 7489913552,
-12299566373, -20830586705, 32156843576, 46298748518, -66176437253, -82513440614,
107256565124, 117915942463, -136506202014, -134651989808, 135526272729, 121980787038,
-103852594014, -86649823814, 60484275256, 47465104086, -26205941074, -19592874823,
8200629244, 5905922924, -1777329681, -1244167978, 250186878, 171453636, -20378088,
-13746096, 723168, 482112
\end{Verbatim}
The defining transfer satisfies $TR(T)\mathbf1=0$ in every coordinate, so \eqref{round8:eq:r8arr-rec} follows. Source: \url{https://oeis.org/A234885}; parent column 3 of \url{https://oeis.org/A234890}. Both consulted pages mark the relevant recurrence empirical.
\Needspace{7\baselineskip}
\subsubsection{A234891}
\label{round8:r8entry:A234891}
Arrays have $n+1$ rows, width $w=2$, and labeled alphabet $0,\ldots,6$. Predicate D has target $s=1$, with every horizontal and vertical neighboring pair unequal.
The recurrence holds for every $n\ge 11$, with order $d=10$ and a complete $D=16$ state matrix. The first three source values are $(160,760,3456)$.
\begin{Verbatim}[fontsize=\scriptsize,breaklines=true]
0, 30, 0, -194, 0, 428, 0, -288, 0, 48
\end{Verbatim}
The defining transfer satisfies $TR(T)\mathbf1=0$ in every coordinate, so \eqref{round8:eq:r8arr-rec} follows. Source: \url{https://oeis.org/A234891}; parent column 1 of \url{https://oeis.org/A234898}. Both consulted pages mark the relevant recurrence empirical.
\Needspace{7\baselineskip}
\subsubsection{A234905}
\label{round8:r8entry:A234905}
Arrays have $n+1$ rows, width $w=2$, and labeled alphabet $0,\ldots,7$. Predicate D has target $s=1$, with every horizontal and vertical neighboring pair unequal.
The recurrence holds for every $n\ge 13$, with order $d=12$ and a complete $D=20$ state matrix. The first three source values are $(280,1720,10344)$.
\begin{Verbatim}[fontsize=\scriptsize,breaklines=true]
0, 55, 0, -725, 0, 3663, 0, -7525, 0, 6335, 0, -1805
\end{Verbatim}
The defining transfer satisfies $TR(T)\mathbf1=0$ in every coordinate, so \eqref{round8:eq:r8arr-rec} follows. Source: \url{https://oeis.org/A234905}; parent column 1 of \url{https://oeis.org/A234911}. Both consulted pages mark the relevant recurrence empirical.
\Needspace{7\baselineskip}
\subsubsection{A234984}
\label{round8:r8entry:A234984}
Arrays have $n+1$ rows, width $w=2$, and labeled alphabet $0,\ldots,5$. Predicate D has target $s=2$, with every horizontal and vertical neighboring pair unequal.
The recurrence holds for every $n\ge 9$, with order $d=8$ and a complete $D=12$ state matrix. The first three source values are $(90,370,1640)$.
\begin{Verbatim}[fontsize=\scriptsize,breaklines=true]
5, 6, -33, -24, 56, 34, -20, -4
\end{Verbatim}
The defining transfer satisfies $TR(T)\mathbf1=0$ in every coordinate, so \eqref{round8:eq:r8arr-rec} follows. Source: \url{https://oeis.org/A234984}; parent column 1 of \url{https://oeis.org/A234991}. Both consulted pages mark the relevant recurrence empirical.
The source rational generating function is also proved: $F=P/Q$, with $Q=1-\sum_{i=1}^{d}c_i x^i$ and
\[P(x)=90x-80x^{2}-750x^{3}+104x^{4}+1832x^{5}+692x^{6}-648x^{7}-120x^{8}.\]
\Needspace{7\baselineskip}
\subsubsection{A234985}
\label{round8:r8entry:A234985}
Arrays have $n+1$ rows, width $w=3$, and labeled alphabet $0,\ldots,5$. Predicate D has target $s=2$, with every horizontal and vertical neighboring pair unequal.
The recurrence holds for every $n\ge 29$, with order $d=28$ and a complete $D=63$ state matrix. The first three source values are $(370,1490,6908)$.
\begin{Verbatim}[fontsize=\scriptsize,breaklines=true]
10, 9, -317, 216, 4392, -4481, -34461, 37067, 168161, -177688, -531245, 547375,
1099678, -1121930, -1478131, 1530027, 1252813, -1355278, -636426, 748324, 180286,
-245436, -25096, 44796, 1272, -4080, 0, 144
\end{Verbatim}
The defining transfer satisfies $TR(T)\mathbf1=0$ in every coordinate, so \eqref{round8:eq:r8arr-rec} follows. Source: \url{https://oeis.org/A234985}; parent column 2 of \url{https://oeis.org/A234991}. Both consulted pages mark the relevant recurrence empirical.
\Needspace{7\baselineskip}
\subsubsection{A235019}
\label{round8:r8entry:A235019}
Arrays have $n+1$ rows, width $w=2$, and labeled alphabet $0,\ldots,5$. Predicate D has target $s=3$, with every horizontal and vertical neighboring pair unequal.
The recurrence holds for every $n\ge 7$, with order $d=6$ and a complete $D=12$ state matrix. The first three source values are $(88,332,1248)$.
\[(c_1,\ldots,c_d)=(0,18,0,-56,0,32).\]
The defining transfer satisfies $TR(T)\mathbf1=0$ in every coordinate, so \eqref{round8:eq:r8arr-rec} follows. Source: \url{https://oeis.org/A235019}; parent column 1 of \url{https://oeis.org/A235026}. Both consulted pages mark the relevant recurrence empirical.
The source rational generating function is also proved: $F=P/Q$, with $Q=1-\sum_{i=1}^{d}c_i x^i$ and
\[P(x)=88x+332x^{2}-336x^{3}-1272x^{4}+192x^{5}+768x^{6}.\]
\Needspace{7\baselineskip}
\subsubsection{A235020}
\label{round8:r8entry:A235020}
Arrays have $n+1$ rows, width $w=3$, and labeled alphabet $0,\ldots,5$. Predicate D has target $s=3$, with every horizontal and vertical neighboring pair unequal.
The recurrence holds for every $n\ge 22$, with order $d=21$ and a complete $D=63$ state matrix. The first three source values are $(332,1200,4312)$.
\begin{Verbatim}[fontsize=\scriptsize,breaklines=true]
5, 40, -213, -616, 3529, 4898, -29657, -24043, 140862, 80302, -394226, -183852,
650340, 268328, -614608, -222048, 318656, 91776, -82432, -14336, 8192
\end{Verbatim}
The defining transfer satisfies $TR(T)\mathbf1=0$ in every coordinate, so \eqref{round8:eq:r8arr-rec} follows. Source: \url{https://oeis.org/A235020}; parent column 2 of \url{https://oeis.org/A235026}. Both consulted pages mark the relevant recurrence empirical.
\Needspace{7\baselineskip}
\subsubsection{A235064}
\label{round8:r8entry:A235064}
Arrays have $n+1$ rows, width $w=2$, and labeled alphabet $0,\ldots,6$. Predicate D has target $s=3$, with every horizontal and vertical neighboring pair unequal.
The recurrence holds for every $n\ge 11$, with order $d=10$ and a complete $D=16$ state matrix. The first three source values are $(168,848,4352)$.
\begin{Verbatim}[fontsize=\scriptsize,breaklines=true]
0, 35, 0, -235, 0, 543, 0, -458, 0, 96
\end{Verbatim}
The defining transfer satisfies $TR(T)\mathbf1=0$ in every coordinate, so \eqref{round8:eq:r8arr-rec} follows. Source: \url{https://oeis.org/A235064}; parent column 1 of \url{https://oeis.org/A235071}. Both consulted pages mark the relevant recurrence empirical.
\Needspace{7\baselineskip}
\subsubsection{A235065}
\label{round8:r8entry:A235065}
Arrays have $n+1$ rows, width $w=3$, and labeled alphabet $0,\ldots,6$. Predicate D has target $s=3$, with every horizontal and vertical neighboring pair unequal.
The recurrence holds for every $n\ge 41$, with order $d=40$ and a complete $D=100$ state matrix. The first three source values are $(848,4348,22936)$.
\begin{Verbatim}[fontsize=\scriptsize,breaklines=true]
7, 114, -826, -5390, 40509, 141982, -1101107, -2376591, 18759386, 27110607,
-214422274, -218647229, 1714423248, 1266295692, -9847282701, -5264745392, 41312388336,
15449236953, -127755960707, -30541686755, 292034903938, 35307790177, -491841781764,
-7382786774, 604462209267, -48701691090, -533018275889, 88708374022, 328737830185,
-79723218859, -136723211260, 42223796030, 36429154260, -13346749528, -5768901376,
2417530544, 474735456, -225680256, -14561280, 8110080
\end{Verbatim}
The defining transfer satisfies $TR(T)\mathbf1=0$ in every coordinate, so \eqref{round8:eq:r8arr-rec} follows. Source: \url{https://oeis.org/A235065}; parent column 2 of \url{https://oeis.org/A235071}. Both consulted pages mark the relevant recurrence empirical.
\Needspace{7\baselineskip}
\subsubsection{A235100}
\label{round8:r8entry:A235100}
Arrays have $n+1$ rows, width $w=2$, and labeled alphabet $0,\ldots,6$. Predicate D has target $s=4$, with every horizontal and vertical neighboring pair unequal.
The recurrence holds for every $n\ge 7$, with order $d=6$ and a complete $D=16$ state matrix. The first three source values are $(154,670,2900)$.
\[(c_1,\ldots,c_d)=(7,-2,-47,3,84,36).\]
The defining transfer satisfies $TR(T)\mathbf1=0$ in every coordinate, so \eqref{round8:eq:r8arr-rec} follows. Source: \url{https://oeis.org/A235100}; parent column 1 of \url{https://oeis.org/A235107}. Both consulted pages mark the relevant recurrence empirical.
The source rational generating function is also proved: $F=P/Q$, with $Q=1-\sum_{i=1}^{d}c_i x^i$ and
\[P(x)=154x-408x^{2}-1482x^{3}+856x^{4}+3312x^{5}+1296x^{6}.\]
\Needspace{7\baselineskip}
\subsubsection{A235101}
\label{round8:r8entry:A235101}
Arrays have $n+1$ rows, width $w=3$, and labeled alphabet $0,\ldots,6$. Predicate D has target $s=4$, with every horizontal and vertical neighboring pair unequal.
The recurrence holds for every $n\ge 20$, with order $d=19$ and a complete $D=100$ state matrix. The first three source values are $(670,2498,9328)$.
\begin{Verbatim}[fontsize=\scriptsize,breaklines=true]
12, -14, -322, 949, 3384, -12901, -21740, 86236, 107744, -322761, -412538, 641088,
1018740, -478292, -1306432, -247248, 584736, 372096, 65664
\end{Verbatim}
The defining transfer satisfies $TR(T)\mathbf1=0$ in every coordinate, so \eqref{round8:eq:r8arr-rec} follows. Source: \url{https://oeis.org/A235101}; parent column 2 of \url{https://oeis.org/A235107}. Both consulted pages mark the relevant recurrence empirical.
\Needspace{7\baselineskip}
\subsubsection{A235180}
\label{round8:r8entry:A235180}
Arrays have $n+1$ rows, width $w=3$, and labeled alphabet $0,\ldots,5$. Predicate D has target $s=5$, with every horizontal and vertical neighboring pair unequal.
The recurrence holds for every $n\ge 19$, with order $d=18$ and a complete $D=63$ state matrix. The first three source values are $(212,560,1400)$.
\begin{Verbatim}[fontsize=\scriptsize,breaklines=true]
3, 29, -88, -327, 1003, 1816, -5577, -5295, 15566, 8705, -20575, -9813, 12552, 5962,
-3436, -1676, 336, 168
\end{Verbatim}
The defining transfer satisfies $TR(T)\mathbf1=0$ in every coordinate, so \eqref{round8:eq:r8arr-rec} follows. Source: \url{https://oeis.org/A235180}; parent column 2 of \url{https://oeis.org/A235186}. Both consulted pages mark the relevant recurrence empirical.
\Needspace{7\baselineskip}
\subsubsection{A235181}
\label{round8:r8entry:A235181}
Arrays have $n+1$ rows, width $w=4$, and labeled alphabet $0,\ldots,5$. Predicate D has target $s=5$, with every horizontal and vertical neighboring pair unequal.
The recurrence holds for every $n\ge 27$, with order $d=26$ and a complete $D=342$ state matrix. The first three source values are $(576,1400,3240)$.
\begin{Verbatim}[fontsize=\scriptsize,breaklines=true]
3, 44, -133, -842, 2563, 9221, -28202, -64059, 195733, 296446, -893437, -937956,
2714749, 2064161, -5442684, -3205526, 7001040, 3522128, -5502080, -2622832, 2452848,
1173648, -555264, -271584, 48384, 24192
\end{Verbatim}
The defining transfer satisfies $TR(T)\mathbf1=0$ in every coordinate, so \eqref{round8:eq:r8arr-rec} follows. Source: \url{https://oeis.org/A235181}; parent column 3 of \url{https://oeis.org/A235186}. Both consulted pages mark the relevant recurrence empirical.
\Needspace{7\baselineskip}
\subsubsection{A235191}
\label{round8:r8entry:A235191}
Arrays have $n+1$ rows, width $w=2$, and labeled alphabet $0,\ldots,6$. Predicate D has target $s=5$, with every horizontal and vertical neighboring pair unequal.
The recurrence holds for every $n\ge 7$, with order $d=6$ and a complete $D=16$ state matrix. The first three source values are $(136,512,1880)$.
\[(c_1,\ldots,c_d)=(0,30,0,-257,0,468).\]
The defining transfer satisfies $TR(T)\mathbf1=0$ in every coordinate, so \eqref{round8:eq:r8arr-rec} follows. Source: \url{https://oeis.org/A235191}; parent column 1 of \url{https://oeis.org/A235198}. Both consulted pages mark the relevant recurrence empirical.
The source rational generating function is also proved: $F=P/Q$, with $Q=1-\sum_{i=1}^{d}c_i x^i$ and
\[P(x)=136x+512x^{2}-2200x^{3}-8288x^{4}+4592x^{5}+17504x^{6}.\]
\Needspace{7\baselineskip}
\subsubsection{A235192}
\label{round8:r8entry:A235192}
Arrays have $n+1$ rows, width $w=3$, and labeled alphabet $0,\ldots,6$. Predicate D has target $s=5$, with every horizontal and vertical neighboring pair unequal.
The recurrence holds for every $n\ge 24$, with order $d=23$ and a complete $D=100$ state matrix. The first three source values are $(512,1544,4724)$.
\begin{Verbatim}[fontsize=\scriptsize,breaklines=true]
6, 52, -367, -995, 9350, 7034, -128869, 26637, 1045428, -801048, -5098437, 5754163,
14714982, -20381230, -23985995, 38063780, 20215810, -35199080, -8246780, 13170048,
2801112, -1572480, -393120
\end{Verbatim}
The defining transfer satisfies $TR(T)\mathbf1=0$ in every coordinate, so \eqref{round8:eq:r8arr-rec} follows. Source: \url{https://oeis.org/A235192}; parent column 2 of \url{https://oeis.org/A235198}. Both consulted pages mark the relevant recurrence empirical.
\Needspace{7\baselineskip}
\subsubsection{A235193}
\label{round8:r8entry:A235193}
Arrays have $n+1$ rows, width $w=4$, and labeled alphabet $0,\ldots,6$. Predicate D has target $s=5$, with every horizontal and vertical neighboring pair unequal.
The recurrence holds for every $n\ge 40$, with order $d=39$ and a complete $D=625$ state matrix. The first three source values are $(1880,4724,12176)$.
\begin{Verbatim}[fontsize=\scriptsize,breaklines=true]
9, 57, -731, -878, 26304, -15900, -552844, 899723, 7521105, -18125633, -69122561,
220983779, 431291855, -1831340243, -1746606119, 10803108780, 3716404250, -46315168788,
2550194940, 145490417695, -46954531755, -335197745077, 166429786799, 564880422518,
-334400168840, -694120308016, 424547635156, 620744392928, -343613520080,
-401608244880, 170947437904, 182786646240, -46865973504, -54609655680, 4665398976,
9369063936, 551978496, -679311360, -113218560
\end{Verbatim}
The defining transfer satisfies $TR(T)\mathbf1=0$ in every coordinate, so \eqref{round8:eq:r8arr-rec} follows. Source: \url{https://oeis.org/A235193}; parent column 3 of \url{https://oeis.org/A235198}. Both consulted pages mark the relevant recurrence empirical.
\Needspace{7\baselineskip}
\subsubsection{A235251}
\label{round8:r8entry:A235251}
Arrays have $n+1$ rows, width $w=2$, and labeled alphabet $0,\ldots,7$. Predicate D has target $s=7$, with every horizontal and vertical neighboring pair unequal.
The recurrence holds for every $n\ge 11$, with order $d=10$ and a complete $D=20$ state matrix. The first three source values are $(184,744,2672)$.
\begin{Verbatim}[fontsize=\scriptsize,breaklines=true]
0, 43, 0, -617, 0, 3073, 0, -1762, 0, 264
\end{Verbatim}
The defining transfer satisfies $TR(T)\mathbf1=0$ in every coordinate, so \eqref{round8:eq:r8arr-rec} follows. Source: \url{https://oeis.org/A235251}; parent column 1 of \url{https://oeis.org/A235258}. Both consulted pages mark the relevant recurrence empirical.


\section{Verification, provenance and reuse}
\label{round8:r8:verification}
The accompanying archive contains the standalone manuscript, numbered
ledger, source definitions, complete exact certificates, verifiers,
compact source audits and independent internal review reports.
With Python, SymPy and mpmath installed, run
\begin{Verbatim}[fontsize=\small]
python3 r8_run_all.py
\end{Verbatim}
from its extracted directory. No network access is needed. The runner
checks the 100 distinct selections, all 708 previous selections,
the retained prior-work baseline, the immutable review bindings and
every included verification program. It records actual versions and
results in \texttt{oeis\_hunt\_round8\_verification\_summary.json}.

\subsection{Why the finite certificates prove all indices}
For each array problem, the complete transfer matrix is derived from
the definition. The identity $TR(T)\mathbf1=0$ implies every later
recurrence instance after multiplication by an arbitrary power of $T$.
An independently written implementation exhaustively constructs the
uncompressed labeled-row relation: all 10,978 states and 24,364
directed edges across the 49 problems. It checks the annihilator on
every full state and verifies every row of the proposed orbit quotient.
Thus the symmetry reduction has a separate complete check.

The generating-function certificates establish exact differential
identities $\mathcal L F=P$, with every nonzero boundary polynomial
displayed. Coefficients beyond its degree vanish for all indices.
Radical-field, rational-parameter and differential-module calculations
are exact. They do not infer an infinite recurrence from a fitted
prefix. Independent coefficient calculations additionally check the
intended sequence and offset.

The arithmetic and formal-series parts supply written universal
proofs. Their finite rational, modular, divisor and numerical checks
are supplementary. In particular, numerical integration is not the
proof of reduced-numerator equality in \oeis{A143991}: the proof also
controls every possible cancellation prime. The real-parameter result
in \oeis{A177430} explicitly invokes the classical six-exponentials
theorem for irrational parameters.

\subsection{Source and priority scope}
Each selected proof links to its primary entry. The compact JSON records
retain its precise target, uncertainty marker, offset, attribution,
retrieval reference, relevant known input and bounded prior-work notes.
Source hashes identify consulted snapshots without redistributing full
raw web pages. An entry's editorial approved status is not interpreted
as a mathematical proof of every formula on that page.

The audit compares all prior project selections and the retained OEIS
Open Appendix A list. It also compares the Fischer (2020) and
Ekhad--Yang (2017) proof-family lists. Relevant linked arithmetic and
generating-function papers were inspected as described in the component
audits. Finding a known input does not itself establish priority for
the selected consequence; no worldwide-first-proof claim is made.

All previously selected IDs remain excluded, including A176605,
A176607, A176610, A176611 and A176612, whose Round IV priority claims
were withdrawn in Round V because Fischer's 2020 general derivation
already covered them. The current and cumulative selection counts are
not presented as counts of verified new discoveries.

\subsection{Standalone manuscript and reproducibility}
The delivered \texttt{.tex} contains all prose, formulas and individual
recurrence signatures directly, with no external
\texttt{\string\input} or \texttt{\string\include} fragments.
Compile with pdfLaTeX until cross-references stabilize.
\texttt{r8\_build\_dossier.py} reconstructs the manuscript, ledger and
combined audit. The README gives all offline certificate-regeneration
commands and dependencies. An isolated regeneration check compares
the regenerated immutable outputs with the delivered versions.

The archive manifest records SHA-256 digests. The internal review
reports separately bind the exact proof, code and certificate versions
they inspected. The layout review binds the visually inspected PDF.

\subsection{Selected common references}
\begin{enumerate}
\item The OEIS Foundation, \emph{The On-Line Encyclopedia of Integer
Sequences}, \url{https://oeis.org/}. Individual primary-entry links
accompany every selected assertion.
\item S. Dasgupta, \emph{Ranks of matrices of logarithms of algebraic
numbers I: the theorems of Baker and Waldschmidt--Masser},
3 March 2023, Theorem 4.2 and Section 5.1,
\url{https://sites.math.duke.edu/~dasgupta/papers/Baker.pdf}.
Used for the six-exponentials theorem, not claimed as a new result.
\item \emph{OEIS Open}, retained Appendix A comparison,
\url{https://arxiv.org/html/2608.11941v1}. The retained scope is 100
costliest claims on 97 distinct A-IDs among 153 combined-model
resolutions, not the complete repository.
\item Georg Fischer, 2020 family derivation linked at
\url{https://oeis.org/A176604/a176604_2.txt}; used for exclusion and credit.
\item S. B. Ekhad and D. Yang, 2017 recurrence family and supporting
output at
\url{https://sites.math.rutgers.edu/~zeilberg/tokhniot/oMathar1maple12.txt};
used for the bounded known-family comparison.
\end{enumerate}
% ===== END ORIGINAL ROUND 8 =====

% ===== BEGIN ORIGINAL ROUND 9 =====
\clearpage
\setcounter{section}{0}\setcounter{equation}{0}\setcounter{theorem}{0}
\setcounter{table}{0}\setcounter{figure}{0}
\normalsize\renewcommand{\arraystretch}{1}\renewcommand{\baselinestretch}{1}\selectfont
\part[Round 9: Gauss congruences, matrix symmetry, valuations, and exact algebra]{Gauss congruences, matrix symmetry, valuations, and exact algebra}
\label{combined:round9}
\rhead{\small Round 9 / KEIO}
\begingroup\small
\section*{Round overview}
\addcontentsline{toc}{section}{Round 9 overview and scope}
This ninth-round dossier resolves 100 precisely selected assertion packages
on 100 distinct OEIS entries, disjoint from the 808 A-numbers selected in
Rounds I--VIII. It proves 96 selected assertions and gives explicit
corrections, with proofs over the full stated domains, for four others.
The four parts contain ten formal-series, congruence and matrix packages,
nine arithmetic packages, nineteen generating-function packages, and
sixty-two constrained-array recurrence packages.

The methods include formal inversion in finite characteristic, integer
exponential series, cyclic actions, an operator symmetry, elementary
divisor and valuation arguments, exact generating-function identities,
and finite transfer operators derived directly from the definitions.

\subsection*{What the count means}
One catalogue row selects one assertion package on one distinct A-number.
All selected statements are explicitly conjectural or empirical in the
consulted primary entries. A recurrence and its generating function count
once, and duplicate formulations of a general assertion are not promoted
to separate results. The 62 array entries are distinct fixed-parameter
counting problems within a shared transfer framework. Thus the count is
not a claim of 100 independent techniques or 100 worldwide-first proofs.

The word ``proved'' refers to the exact target and domain stated in each
package. Other conjectures on the same OEIS page may remain untouched.
Established formulas, prior project methods and linked results are
credited as inputs, including cases where the new package is an elementary
consequence or an explicit repair of an old formula.

\subsection*{The four corrections}
\begin{description}
\item[A202943.] The parity-product sum must start at $j=0$, not $j>0$.
The omitted constant already contradicts $a_0=1$. The corrected full
product and integrality of the underlying cube root are proved.
\item[A048896.] The 2026 factorial sum ending at $n+1$ has reduced
numerator $13$ when $n=1$, whereas $a(1)=1$. Its upper limit must be
$n+2$; the repaired numerator identity is proved for every $n\ge0$.
\item[A001784.] The printed recurrence predates an offset change.
Its coefficient index is shifted by one, giving the correct recurrence
for $n\ge2$. The previously recorded closed form is credited.
\item[A052736.] The stated identity fails at $n=1$. Its corrected domain
is $n\ge2$, established from the defining exponential generating function.
\end{description}
No correction is counted twice. Additional notation and convention issues
are disclosed within the relevant proofs; in particular A180654's
integrality theorem covers both the printed extra-$x^m$ convention and the
program/example convention. The recurrences of A025256 and A025260 start
at $n=6$, as their exact boundary polynomials show.

\subsection*{Source and prior-work audit}
The primary entries, relevant linked material and targeted searches were
consulted on 20 October 2026. Cached retrievals can predate that date;
source wording does not certify every live edit or worldwide priority.
The earlier 808 is an exclusion ledger, not a cumulative discovery count.
The five Round IV priority withdrawals disclosed in Round V remain in
force and remain excluded. The retained OEIS Open comparison covers
100 Appendix A claims on 97 distinct A-numbers, with zero overlap; it
does not cover that project's complete repository.

Several initially promising targets were removed after finding prior
proofs, including A354766 and A278085 in an October 2026 paper, the linked
A192069/A192248/A192249 family, and source entries that already contain
proofs or immediate duplicate recurrences. A234202 has a fresh exact
certificate but remains an uncounted reserve. Previously certified
reserves from earlier rounds are not counted as new work here.

All 62 selected array primary pages were retrieved. Of the 24 parent
tables for the selected and reserve set, 21 were retrievable. The
unretrieved parents A233883, A234114 and A234175 are disclosed in the
affected records. The child definitions and certificates do not rely on
those missing pages, but their historical proof status is not asserted to
have been audited.

No OEIS submission or edit was performed. These are mathematical
resolutions of selected assertions, separate from editorial acceptance.
Independent internal reviews were performed by separate agents in the
same AI-assisted workflow. They are not external peer review or
proof-assistant formalization.
\endgroup

\clearpage
\section{Catalogue of the 100 selected entries}
\label{round9:sec:r9catalogue}
\textbf{P}: selected assertion proved over its stated domain.
\textbf{C}: source formula or domain corrected and proved.
Page references lead to the complete argument or exact certificate.

\small
\begin{longtable}{@{}r p{21mm} p{103mm} c r@{}}
\toprule No.&Entry&Selected assertion package&Result&Page\\
\midrule\endfirsthead
\toprule No.&Entry&Selected assertion package&Result&Page\\
\midrule\endhead
\midrule\multicolumn{5}{r}{\emph{Continued on the next page}}\\\endfoot
\bottomrule\endlastfoot

\multicolumn{5}{@{}l}{\textbf{Formal series, congruences and matrix symmetry}}\\[2pt]

1&\oeis{A158119}&Powers-of-two parity from the continued fraction&P&\pageref{round9:r9entry:A158119}\\

2&\oeis{A180654}&Three exponential-integrality assertions&P&\pageref{round9:r9entry:A180654}\\

3&\oeis{A384270}&Dyadic parity of the composition equation&P&\pageref{round9:r9entry:A384270}\\

4&\oeis{A376226}&The complete ternary residue rule&P&\pageref{round9:r9entry:A376226}\\

5&\oeis{A397588}&Dyadic parity of the differential equation&P&\pageref{round9:r9entry:A397588}\\

6&\oeis{A355414}&The prime-index logarithmic congruence&P&\pageref{round9:r9entry:A355414}\\

7&\oeis{A032312}&All prime-power Gauss congruences&P&\pageref{round9:r9entry:A032312}\\

8&\oeis{A167138}&Integrality under every positive pointwise power&P&\pageref{round9:r9entry:A167138}\\

9&\oeis{A091351}&Symmetry between matrix-power parameters&P&\pageref{round9:r9entry:A091351}\\

10&\oeis{A202943}&Corrected parity-product lower limit&C&\pageref{round9:r9entry:A202943}\\

\multicolumn{5}{@{}l}{\textbf{Arithmetic identities and a factorial-sum correction}}\\[2pt]

11&\oeis{A026741}&The general-prime gcd and lcm identities&P&\pageref{round9:r9entry:A026741}\\

12&\oeis{A040003}&A gcd formula for the partial quotients of sqrt(6)&P&\pageref{round9:r9entry:A040003}\\

13&\oeis{A048896}&Corrected upper limit in the factorial sum&C&\pageref{round9:r9entry:A048896}\\

14&\oeis{A084068}&The square-triangular Pell gcd&P&\pageref{round9:r9entry:A084068}\\

15&\oeis{A096617}&The reduced-numerator recurrence for n H(n)&P&\pageref{round9:r9entry:A096617}\\

16&\oeis{A329354}&The summatory asymptotic and its constant&P&\pageref{round9:r9entry:A329354}\\

17&\oeis{A363409}&The selected exact 7-adic valuation&P&\pageref{round9:r9entry:A363409}\\

18&\oeis{A363411}&The selected exact 5-adic valuation&P&\pageref{round9:r9entry:A363411}\\

19&\oeis{A363412}&The selected 3-adic assertions&P&\pageref{round9:r9entry:A363412}\\

\multicolumn{5}{@{}l}{\textbf{Generating functions and polynomial recurrences}}\\[2pt]

20&\oeis{A000483}&Conjectured polynomial recurrence of order 3&P&\pageref{round9:r9entry:A000483}\\

21&\oeis{A001784}&Correction of the recurrence after the offset change&C&\pageref{round9:r9entry:A001784}\\

22&\oeis{A025256}&Conjectured polynomial recurrence of order 4&P&\pageref{round9:r9entry:A025256}\\

23&\oeis{A025260}&Conjectured polynomial recurrence of order 4&P&\pageref{round9:r9entry:A025260}\\

24&\oeis{A026023}&Conjectured polynomial recurrence of order 3&P&\pageref{round9:r9entry:A026023}\\

25&\oeis{A026027}&Conjectured polynomial recurrence of order 3&P&\pageref{round9:r9entry:A026027}\\

26&\oeis{A026570}&Conjectured polynomial recurrence of order 3&P&\pageref{round9:r9entry:A026570}\\

27&\oeis{A026572}&Conjectured polynomial recurrence of order 5&P&\pageref{round9:r9entry:A026572}\\

28&\oeis{A026582}&Conjectured polynomial recurrence of order 6&P&\pageref{round9:r9entry:A026582}\\

29&\oeis{A026732}&Conjectured polynomial recurrence of order 8&P&\pageref{round9:r9entry:A026732}\\

30&\oeis{A026765}&Conjectured polynomial recurrence of order 6&P&\pageref{round9:r9entry:A026765}\\

31&\oeis{A052736}&Correction of the boundary range in the identity&C&\pageref{round9:r9entry:A052736}\\

32&\oeis{A071715}&Conjectured polynomial recurrence of order 3&P&\pageref{round9:r9entry:A071715}\\

33&\oeis{A098520}&Conjectured polynomial recurrence of order 2&P&\pageref{round9:r9entry:A098520}\\

34&\oeis{A123071}&Conjectured polynomial recurrence of order 7&P&\pageref{round9:r9entry:A123071}\\

35&\oeis{A143013}&Conjectured polynomial recurrence of order 3&P&\pageref{round9:r9entry:A143013}\\

36&\oeis{A372018}&The identity for the odd-index subsequence&P&\pageref{round9:r9entry:A372018}\\

37&\oeis{A212341}&The empirical coefficient generating function&P&\pageref{round9:r9entry:A212341}\\

38&\oeis{A110668}&The empirical generating function of partial sums&P&\pageref{round9:r9entry:A110668}\\

\multicolumn{5}{@{}l}{\textbf{Finite-state array recurrence certificates}}\\[2pt]

39&\oeis{A233787}&Empirical array recurrence of order 34&P&\pageref{round9:r9entry:A233787}\\

40&\oeis{A233877}&Empirical array recurrence of order 6&P&\pageref{round9:r9entry:A233877}\\

41&\oeis{A233878}&Empirical array recurrence of order 12&P&\pageref{round9:r9entry:A233878}\\

42&\oeis{A233879}&Empirical array recurrence of order 29&P&\pageref{round9:r9entry:A233879}\\

43&\oeis{A233880}&Empirical array recurrence of order 52&P&\pageref{round9:r9entry:A233880}\\

44&\oeis{A233960}&Empirical array recurrence of order 5&P&\pageref{round9:r9entry:A233960}\\

45&\oeis{A233961}&Empirical array recurrence of order 14&P&\pageref{round9:r9entry:A233961}\\

46&\oeis{A233962}&Empirical array recurrence of order 43&P&\pageref{round9:r9entry:A233962}\\

47&\oeis{A234029}&Empirical array recurrence of order 4&P&\pageref{round9:r9entry:A234029}\\

48&\oeis{A234030}&Empirical array recurrence of order 9&P&\pageref{round9:r9entry:A234030}\\

49&\oeis{A234031}&Empirical array recurrence of order 16&P&\pageref{round9:r9entry:A234031}\\

50&\oeis{A234032}&Empirical array recurrence of order 44&P&\pageref{round9:r9entry:A234032}\\

51&\oeis{A234107}&Empirical array recurrence of order 6&P&\pageref{round9:r9entry:A234107}\\

52&\oeis{A234108}&Empirical array recurrence of order 18&P&\pageref{round9:r9entry:A234108}\\

53&\oeis{A234109}&Empirical array recurrence of order 40&P&\pageref{round9:r9entry:A234109}\\

54&\oeis{A234117}&Empirical array recurrence of order 5&P&\pageref{round9:r9entry:A234117}\\

55&\oeis{A234118}&Empirical array recurrence of order 11&P&\pageref{round9:r9entry:A234118}\\

56&\oeis{A234119}&Empirical array recurrence of order 17&P&\pageref{round9:r9entry:A234119}\\

57&\oeis{A234124}&Empirical array recurrence of order 6&P&\pageref{round9:r9entry:A234124}\\

58&\oeis{A234125}&Empirical array recurrence of order 14&P&\pageref{round9:r9entry:A234125}\\

59&\oeis{A234126}&Empirical array recurrence of order 27&P&\pageref{round9:r9entry:A234126}\\

60&\oeis{A234133}&Empirical array recurrence of order 3&P&\pageref{round9:r9entry:A234133}\\

61&\oeis{A234135}&Empirical array recurrence of order 6&P&\pageref{round9:r9entry:A234135}\\

62&\oeis{A234154}&Empirical array recurrence of order 3&P&\pageref{round9:r9entry:A234154}\\

63&\oeis{A234155}&Empirical array recurrence of order 6&P&\pageref{round9:r9entry:A234155}\\

64&\oeis{A234156}&Empirical array recurrence of order 11&P&\pageref{round9:r9entry:A234156}\\

65&\oeis{A234157}&Empirical array recurrence of order 26&P&\pageref{round9:r9entry:A234157}\\

66&\oeis{A234158}&Empirical array recurrence of order 53&P&\pageref{round9:r9entry:A234158}\\

67&\oeis{A234170}&Empirical array recurrence of order 6&P&\pageref{round9:r9entry:A234170}\\

68&\oeis{A234171}&Empirical array recurrence of order 35&P&\pageref{round9:r9entry:A234171}\\

69&\oeis{A234177}&Empirical array recurrence of order 5&P&\pageref{round9:r9entry:A234177}\\

70&\oeis{A234178}&Empirical array recurrence of order 11&P&\pageref{round9:r9entry:A234178}\\

71&\oeis{A234179}&Empirical array recurrence of order 22&P&\pageref{round9:r9entry:A234179}\\

72&\oeis{A234201}&Empirical array recurrence of order 10&P&\pageref{round9:r9entry:A234201}\\

73&\oeis{A234210}&Empirical array recurrence of order 4&P&\pageref{round9:r9entry:A234210}\\

74&\oeis{A234211}&Empirical array recurrence of order 15&P&\pageref{round9:r9entry:A234211}\\

75&\oeis{A234212}&Empirical array recurrence of order 56&P&\pageref{round9:r9entry:A234212}\\

76&\oeis{A234659}&Empirical array recurrence of order 5&P&\pageref{round9:r9entry:A234659}\\

77&\oeis{A234660}&Empirical array recurrence of order 16&P&\pageref{round9:r9entry:A234660}\\

78&\oeis{A234661}&Empirical array recurrence of order 39&P&\pageref{round9:r9entry:A234661}\\

79&\oeis{A234667}&Empirical array recurrence of order 7&P&\pageref{round9:r9entry:A234667}\\

80&\oeis{A234668}&Empirical array recurrence of order 31&P&\pageref{round9:r9entry:A234668}\\

81&\oeis{A234683}&Empirical array recurrence of order 6&P&\pageref{round9:r9entry:A234683}\\

82&\oeis{A234684}&Empirical array recurrence of order 24&P&\pageref{round9:r9entry:A234684}\\

83&\oeis{A234697}&Empirical array recurrence of order 5&P&\pageref{round9:r9entry:A234697}\\

84&\oeis{A234698}&Empirical array recurrence of order 15&P&\pageref{round9:r9entry:A234698}\\

85&\oeis{A234699}&Empirical array recurrence of order 31&P&\pageref{round9:r9entry:A234699}\\

86&\oeis{A234706}&Empirical array recurrence of order 31&P&\pageref{round9:r9entry:A234706}\\

87&\oeis{A234707}&Empirical array recurrence of order 36&P&\pageref{round9:r9entry:A234707}\\

88&\oeis{A234721}&Empirical array recurrence of order 8&P&\pageref{round9:r9entry:A234721}\\

89&\oeis{A234722}&Empirical array recurrence of order 30&P&\pageref{round9:r9entry:A234722}\\

90&\oeis{A234816}&Empirical array recurrence of order 7&P&\pageref{round9:r9entry:A234816}\\

91&\oeis{A234817}&Empirical array recurrence of order 23&P&\pageref{round9:r9entry:A234817}\\

92&\oeis{A234818}&Empirical array recurrence of order 23&P&\pageref{round9:r9entry:A234818}\\

93&\oeis{A234992}&Empirical array recurrence of order 10&P&\pageref{round9:r9entry:A234992}\\

94&\oeis{A234993}&Empirical array recurrence of order 49&P&\pageref{round9:r9entry:A234993}\\

95&\oeis{A235169}&Empirical array recurrence of order 9&P&\pageref{round9:r9entry:A235169}\\

96&\oeis{A235170}&Empirical array recurrence of order 46&P&\pageref{round9:r9entry:A235170}\\

97&\oeis{A235205}&Empirical array recurrence of order 6&P&\pageref{round9:r9entry:A235205}\\

98&\oeis{A235206}&Empirical array recurrence of order 29&P&\pageref{round9:r9entry:A235206}\\

99&\oeis{A235233}&Empirical array recurrence of order 28&P&\pageref{round9:r9entry:A235233}\\

100&\oeis{A235234}&Empirical array recurrence of order 40&P&\pageref{round9:r9entry:A235234}\\

\end{longtable}
\normalsize
\clearpage

\section{Formal series, congruences, and matrix symmetry}
\label{round9:r9sec:root}

The ten packages in this section concern precisely identified conjectures
on ten primary OEIS entries. Nine are proved as stated; the parity formula
at A202943 requires its summation to include the zero term. Known definitions,
classical identities, and closely related earlier results are credited below.
Finite calculations in \texttt{r9\_root\_verify.py} check transcription and
boundary cases; the arguments in this section prove the universal statements.

\subsection{A158119: parity from a continued fraction}
\label{round9:r9entry:A158119}
\textbf{Source.} \oeis{A158119} conjectures that $a(n)$ is odd exactly when
$n=0$ or $n$ is a power of two. Its ordinary generating function, with
$A(0)=1$, is characterized by Hanna's displayed continued fraction
\[
 1=A(x)-\cfrac{x}{A(x)-\cfrac{4x}{A(x)-\cfrac{9x}{A(x)-\cdots}}}.
\]
The continued-fraction characterization is a known input from the entry.

\begin{proposition}
The integral solution satisfies
\[
 A(x)\equiv1+\sum_{j\ge0}x^{2^j}\pmod2.
\]
\end{proposition}
\begin{proof}
All denominators of the fraction are units. Truncation at depth $N$ fixes
coefficients through a degree tending to infinity with $N$, because each
additional level is multiplied by $x$. The fraction therefore has a precise
coefficientwise meaning. In degree $n\ge1$, the unknown coefficient of $A$
occurs with coefficient $1$ in the initial $A$; the rest uses only earlier
coefficients. This proves existence and uniqueness over $\mathbb Z$.

Modulo two the numerator $4x$ vanishes, so
$1=A-x/A$, or $A^2+A=x$. Set $B=A-1$. In characteristic two this is
$B+B(x^2)=x$. Comparing coefficients gives $[x]B=1$, all other odd-index
coefficients zero, and $[x^{2n}]B=[x^n]B$. Iteration gives exactly the stated
support, including the constant coefficient of $A$.
\end{proof}
\textbf{Prior-work scope.} The linked Bozejko--Hasebe paper concerns free
infinite divisibility; its inspected text supplied no parity theorem for
this sequence. The related entries A157308, A157310, A155585, and A157309
describe the transforms underlying the continued fraction. They are not
additional selected packages.

\subsection{A180654: integrality of the logarithmic construction}
\label{round9:r9entry:A180654}
\textbf{Source.} \oeis{A180654} states three integrality conjectures: the
binomial-square case, every fixed positive integer power, and the case in
which the power equals the summation index. Its program and expanded example
use the following series:
\[
 P_{m,p}(x)=\sum_{k=0}^m\binom{m}{k}^{p}x^k,\qquad
 F(x)=1+\sum_{m\ge1}\frac1{m!}
       \left(\frac{\log P_{m,p_m}(x)}m\right)^m.
\]
Here each $p_m$ is a positive integer. The source title and two displayed
generalizations contain an extra factor $x^m$ absent from the program and
coefficient list. The theorem covers both conventions; the numerical sequence
is checked against the program convention.

\begin{proposition}
For every choice of positive integers $(p_m)_{m\ge1}$, all coefficients
$n![x^n]F$ are integers. The same assertion holds after multiplying the $m$th
summand by $x^m$. Thus all three source integrality claims hold under either
of the two displayed conventions.
\end{proposition}
\begin{proof}
For $1\le k\le m$,
\[
 \frac{k}{m}\binom{m}{k}^{p_m}
 =\binom{m-1}{k-1}\binom{m}{k}^{p_m-1}\in\mathbb Z.
\]
Since $P_{m,p_m}(0)=1$, this shows
\[
 \left(\frac{\log P_{m,p_m}}m\right)'
   =\frac{P'_{m,p_m}/m}{P_{m,p_m}}\in\mathbb Z[[x]].
\]
Consequently $H_m=\log(P_{m,p_m})/m$ has zero constant term and integral
exponential coefficients: if $H_m'=\sum b_jx^j$, then
$r![x^r]H_m=(r-1)!b_{r-1}$.

For any such $H=\sum_{r\ge1}h_rx^r/r!$ with $h_r\in\mathbb Z$,
\[
 n![x^n]\frac{H^m}{m!}
 =\sum_{\substack{\sum j_r=m\\\sum rj_r=n}}
    \frac{n!}{\prod_r(r!)^{j_r}j_r!}\prod_r h_r^{j_r}\in\mathbb Z.
\]
The prefactor counts set partitions with prescribed block sizes. Each
$H_m^m$ starts in degree $m$, so the sum is locally finite. Finally,
$n![x^n](x^mG)=n!/(n-m)!\cdot (n-m)![x^{n-m}]G$, preserving integrality.
Taking $p_m=2$, a fixed $p$, or $p_m=m$ gives the three cases. They form one
package at this A-number, not three separate counts.
\end{proof}

\subsection{A384270: a composition quotient modulo two}
\label{round9:r9entry:A384270}
\textbf{Source.} \oeis{A384270}, Hanna, July 10, 2025, defines
\[
 A=x+O(x^2),\qquad
 A\,A(x^3+3xA^3)=A(x^4+4xA^4),
\]
and conjectures odd coefficients precisely at powers of two.

\begin{proposition}
There is a unique integral solution and
$A\equiv\sum_{j\ge0}x^{2^j}\pmod2$.
\end{proposition}
\begin{proof}
For $n\ge2$, the coefficient $a_n$ appears only in the leading product
$A x^3$ when the defining equation is read in degree $n+3$; every other
occurrence has larger degree. Its coefficient is $1$. This gives a monic
integer recursion, also unique over every commutative coefficient ring.

Over $\mathbb F_2$, the equation becomes
$A\,A(x^3+xA^3)=A(x^4)=A^4$. Let
$S=\sum_{j\ge0}x^{2^j}$, so $x=S+S^2$. A direct identity gives
\[
 x^3+xS^3=S^3+S^6.
\]
The series $S(y)$ is the compositional inverse of $y+y^2$; hence
$S(x^3+xS^3)=S^3$. Thus $S$ satisfies the reduced equation. The monic
uniqueness just proved forces $A=S$ modulo two.
\end{proof}

\subsection{A376226: the full ternary support rule}
\label{round9:r9entry:A376226}
\textbf{Source.} \oeis{A376226}, Hanna, October 4, 2024, defines
\[
 A=x+O(x^2),\qquad
 A\,A(x^2+4xA^2)=A(x^3+6xA^3).
\]
Its two conjectures describe the coefficients congruent to $1$ and $2$
modulo three.

\begin{proposition}
For every $n\ge1$,
\[
 a_n\equiv
 \begin{cases}
  1,&n=3^i\quad(i\ge0),\\
  2,&n=(3^i+3^j)/2\quad(0\le i<j),\\
  0,&\text{otherwise}
 \end{cases}\pmod3.
\]
\end{proposition}
\begin{proof}
In degree $n+2$ the new coefficient $a_n$ occurs only in $Ax^2$, with
coefficient $1$. All other contributions involve earlier coefficients.
This establishes integral existence and uniqueness, also modulo three.
The reduced equation, after cancellation of the nonzero series $A$, is
\[
 A(x^2+xA^2)=A^2.
\]
Let $B$ be the inverse of $u+u^2+u^3$. For $x=B+B^2+B^3$ one has
\[
 x^2+xB^2=B^2+B^4+B^6.
\]
The right side is $u+u^2+u^3$ evaluated at $u=B^2$, so $B$ satisfies the
reduced equation. Uniqueness gives $A=B$ modulo three.

To read its coefficients, put $U(t)=\sum_{i\ge0}t^{3^i}$. Then
$U-U^3=t$, and
\[
 t^2=(U-U^3)^2=U^2+U^4+U^6.
\]
Thus $B(t^2)=U(t)^2$. Its diagonal terms are $t^{2\cdot3^i}$ and its
off-diagonal terms are $2t^{3^i+3^j}$. Base-three expansions show that these
exponents are distinct and that no diagonal exponent is off-diagonal.
Dividing all exponents by two proves the three cases.
\end{proof}

\subsection{A397588: parity of a differential recursion}
\label{round9:r9entry:A397588}
\textbf{Source.} \oeis{A397588}, Hanna, July 3, 2026, has
$A=x+(xA^2)'$, $A\in x\mathbb Z[[x]]$. The selected conjecture says that
$a_n$ is odd exactly at powers of two, for $n\ge1$.

\begin{proposition}
$A\equiv\sum_{j\ge0}x^{2^j}\pmod2$.
\end{proposition}
\begin{proof}
The defining equation gives $a_1=1$ and
\[
 a_n=(n+1)\sum_{k=1}^{n-1}a_ka_{n-k}\qquad(n\ge2),
\]
which establishes integrality and uniqueness. In characteristic two,
$(xA^2)'=A^2+2xAA'=A^2$, so $A=x+A^2$.
Coefficient comparison gives the required support as in A158119.
\end{proof}

\subsection{A355414: the prime-index logarithmic congruence}
\label{round9:r9entry:A355414}
\textbf{Source.} \oeis{A355414}, Melika Tebni, July 1, 2022, defines
\[
 \sum_{n\ge0}a_n\frac{x^n}{n!}
 =\frac{\log((1-x)/(1-2x))}{(1-x)^5}
\]
and conjectures $a_p\equiv-1\pmod p$ for every prime $p$.

\begin{proposition}
More generally, if $H\in\mathbb Z[[x]]$, $H(0)=1$, and
$b_n=n![x^n]H(x)\log((1-x)/(1-qx))$ with integer $q$, then
$b_n\in\mathbb Z$ and
\[
 b_p\equiv1-q\pmod p\qquad(p\text{ prime}).
\]
\end{proposition}
\begin{proof}
Write $H=\sum h_jx^j$. For $n\ge1$,
\[
 b_n=\sum_{k=1}^n \frac{n!}{k}h_{n-k}(q^k-1).
\]
Each $n!/k$ is an integer. At $n=p$ every term with $k<p$ is divisible
by $p$. The remaining term is $(p-1)!(q^p-1)$, congruent to
$-(q-1)$ by Wilson's theorem and $q^p\equiv q\pmod p$.
Taking $H=(1-x)^{-5}$ and $q=2$ proves the selected assertion, including
$p=2$.
\end{proof}
\textbf{Credit and counting.} Fried's 2025 note and the 2026 compilation,
Section 3.9, Theorem 52, already prove the power-three case A355372.
The proof above is a routine extension of that type of argument, not an
independent new method. A355372 is excluded; the power-four consequence
A355407 is not counted separately. Only the still-marked power-five
assertion is selected here.

\subsection{A032312: all Gauss congruences by a cyclic action}
\label{round9:r9entry:A032312}
\textbf{Source.} \oeis{A032312} has exponential generating function
$\prod_{r\ge1}(1+x^r/r!)^2$. Peter Bala, September 5, 2022, conjectures its
Gauss congruences and the equivalent integrality of its logarithmic
exponential transform.

\begin{proposition}
For every prime $p$ and all positive integers $n,k$,
\[
 a(np^k)\equiv a(np^{k-1})\pmod{p^k}.
\]
Furthermore $\exp(\sum_{m\ge1}a(m)x^m/m)\in\mathbb Z[[x]]$.
The congruence proof holds for any fixed nonnegative number of colors.
\end{proposition}
\begin{proof}
The exponential product counts set partitions whose blocks have one of two
colors, with at most one block of each given size and color. Let a cyclic
group $C_{p^k}$ act freely on $np^k$ labels, as $n$ disjoint regular orbits,
and hence on these partitions. An orbit of partitions whose size is less
than $p^k$ is fixed by the unique subgroup $H$ of order $p$; conversely,
a partition not fixed by $H$ has orbit size exactly $p^k$.

If a partition is fixed by $H$, each block is fixed individually: the action
preserves its size and color, and a second block with the same pair is
forbidden. Each block is therefore a union of $H$-orbits of labels.
Quotienting by $H$ gives exactly the same kind of colored partition on
$np^{k-1}$ labels, with block sizes divided by $p$. This is a bijection.
All other partitions occur in orbits of size $p^k$, proving the congruence.

For completeness, put
\[
 c_m=\frac1m\sum_{d\mid m}\mu(d)a(m/d).
\]
If $m=p^su$, $p\nmid u$, pair divisors $d$ and $pd$ to see that the
numerator equals
$\sum_{d\mid u}\mu(d)(a(p^su/d)-a(p^{s-1}u/d))$, divisible by $p^s$.
Doing this for every prime factor gives $c_m\in\mathbb Z$. Formal logarithms
then give
\[
 \exp\left(\sum_{m\ge1}\frac{a(m)}m x^m\right)
 =\prod_{m\ge1}(1-x^m)^{-c_m}\in\mathbb Z[[x]].
\]
All products and sums are locally finite coefficientwise.
\end{proof}
\textbf{Prior-work scope.} A007837 is the one-color case, and its entry
credits Jovovic with a September 26, 2026 proof of the Gauss congruences.
That case and its companion A168268 are excluded. The selected two-color
case is proved directly above; the common Euler-product criterion is
classical, not a new theorem.

\subsection{A167138: closure under powers after logarithmic substitution}
\label{round9:r9entry:A167138}
\textbf{Source.} \oeis{A167138} asks whether integrality of
$\sum_{n\ge0}L(n)\log(1+x)^n/n!$ implies integrality of the same expression
with $L(n)$ replaced by $L(n)^m$, for every positive integer $m$.
The same general conjecture is repeated at A167139, A167141, and A167532;
these repetitions contribute no additional count.

\begin{proposition}
The asserted closure holds. In fact, the class is closed under pointwise
products of any finite number of sequences.
\end{proposition}
\begin{proof}
Define $s_j(n)=j!\,\genfrac{\{}{\}}{0pt}{}{n}{j}$, the number of surjections from an $n$-set
to a fixed $j$-set, with $s_0(0)=1$ and $s_0(n)=0$ for $n>0$.
If the original integral series is $F(x)=\sum c_jx^j$, substitution
$x=e^t-1$ shows
\[
 L(n)=\sum_{j=0}^n c_j s_j(n),\qquad c_j\in\mathbb Z.
\]
For nonnegative $j_1,\ldots,j_m$, let $C(j_1,\ldots,j_m;r)$ be the number
of $r$-element subsets of $[j_1]\times\cdots\times[j_m]$ whose projection
onto each factor is surjective. A tuple of surjections has such a joint
image; after choosing that image there are $s_r(n)$ surjections onto it.
Hence
\[
 \prod_{i=1}^m s_{j_i}(n)
 =\sum_{r\ge0}C(j_1,\ldots,j_m;r)s_r(n).
\]
This includes empty sets with the stated conventions. For a fixed $r$ a
nonzero coefficient forces every $j_i\le r$, except that the $r=0$ case
forces all $j_i=0$. Thus collecting coefficients in a product of the
expansions of $L(n)$ is finite at each $r$ and yields integral coefficients
$d_r$. Finally
\[
 \sum_{n\ge0}s_r(n)\frac{\log(1+x)^n}{n!}=x^r.
\]
The transformed product is $\sum d_rx^r\in\mathbb Z[[x]]$.
For the entry itself, $c_j$ is the partition number $p(j)$ and $m=2$.
\end{proof}

\subsection{A091351: symmetry of two matrix-power parameters}
\label{round9:r9entry:A091351}
\textbf{Source.} \oeis{A091351} is the lower triangular matrix $T$ in which
$T(n,k)$ equals the row sum in row $n-k$ of $T^k$. Hanna, July 16, 2026,
conjectures
\[
 (T^p)_{n+q-1,q-1}=(T^q)_{n+p-1,p-1}
 \quad(n\ge0,\ p,q\ge1).
\]
We prove it from a formal operator identity and also justify that identity
from an equivalent recursive construction of this matrix.

\begin{proposition}
The displayed symmetry holds for every stated $n,p,q$.
\end{proposition}
\begin{proof}
Let $S$ be the down-shift matrix, $S_{i,j}=1$ if $i=j+1$ and $0$ otherwise,
and let $e_0=(1,0,\ldots)^t$. First construct a lower triangular matrix $T$
by $T(n,0)=1$ and
\[
 TS=ST+(ST)^2. \tag{\the\value{section}.M1}\label{round9:r9eq:matrix-shift}
\]
This determines column $k+1$ in row $n$ using only rows smaller than $n$,
so there is a unique integral matrix with diagonal entries $1$.
Put $D=ST$ and $f(z)=z+z^2$, and let $g$ be the compositional inverse of
$f$. Since $D$ is strictly lower triangular, every power series in $D$ is
well-defined entrywise. The shift relation gives
\[
 TD=f(D)T,\qquad DS=Sf(D),\qquad SD=g(D)S.
\]
The last identity follows by applying $g$ to the middle intertwining
identity. Thus, for any formal $h$,
\[
 T h(D)=h(f(D))T,\qquad S h(D)=h(g(D))S.
\]
Since $Te_0$ is the all-ones vector, $Te_0=(1+D)e_0$.
In particular, $De_0=S(1+D)e_0=(1+g(D))Se_0$. As
$D=g(D)(1+g(D))$, cancellation of the unit $1+g(D)$ gives
$Se_0=g(D)e_0$. It follows that
\[
 Se_0=g(D)e_0,\quad
 S^r e_0=\prod_{i=1}^r g^{\circ i}(D)e_0,\quad
 T^r e_0=\frac{f^{\circ r}(D)}D e_0.
\]
The quotient in the last formula is a power series with constant term $1$;
no inverse of the shift matrix is being assumed. These identities follow
by induction and $f(z)/z=1+z$.

Define $H_1(z)=1$ and
$H_r(z)=\prod_{i=0}^{r-2}g^{\circ i}(z)$ for $r\ge2$. Applying the preceding
intertwining identities and cancelling equal formal factors yields
\[
 S^{p-1}T^pS^{q-1}e_0
  =(1+D)H_p(D)H_q(D)e_0. \tag{\the\value{section}.M2}\label{round9:r9eq:matrix-sym}
\]
One direct check of the cancellation starts from
\[
 T^pS^{q-1}e_0=
 \frac{f^{\circ p}(D)}D
   \prod_{i=1}^{q-1} f^{\circ(p-i)}(D)e_0,
\]
where negative iterates denote iterates of $g$. Apply $S^{p-1}$, which
substitutes $g^{\circ(p-1)}(D)$ and contributes
$\prod_{i=1}^{p-1}g^{\circ i}(D)$. The cancelled factors leave exactly
\eqref{round9:r9eq:matrix-sym}. Each quotient after cancellation is an ordinary
power series, so this calculation is valid in the strictly triangular
matrix algebra.

Putting $q=1$ in \eqref{round9:r9eq:matrix-sym} proves
$S^{p-1}T^p e_0=T S^{p-1}e_0$. Therefore column $k$ of $T$, shifted up by
$k$, equals $T^{k+1}e_0=T^k\mathbf1$. This is precisely the source
definition. That definition also determines rows successively: for $k>0$,
the relevant power uses only rows through $n-k<n$. Hence the matrix just
constructed is the source matrix.

Finally the right side of \eqref{round9:r9eq:matrix-sym} is symmetric in $p,q$.
Comparing coordinate $n+p+q-2$ proves the desired equality.
\end{proof}

\subsection{A202943: a corrected parity product and integral cube root}
\label{round9:r9entry:A202943}
\textbf{Source.} \oeis{A202943}, Hanna, December 26, 2011, defines the
cube root with constant term $1$ by
\[
 A(x)^3=R(x):=\sum_{n\ge0}\binom{n+2}{2}2^{n(n-1)}x^n.
\]
The conjectured parity product uses $\sum_{j>0}x^{32M(j)}$, where
$M=$A000695 has offset $0$ and $M(0)=0$. That lower limit omits the
constant term, contradicting $a_0=1$ (and $a_1=1$). The source's expanded
example already includes the missing $1$. We make the correction explicit.

\begin{proposition}
$A\in\mathbb Z[[x]]$, and its corrected parity formula is
\[
 A(x)\equiv(1+x)(1+x^2)(1+x^8)
              \sum_{j\ge0}x^{32M(j)}\pmod2.
\]
In addition, $a_n\equiv1\pmod3$ for all $n\ge0$; this last assertion is
already unmarked in the entry and is not an extra selected claim.
\end{proposition}
\begin{proof}
First, $R(x)\equiv(1-x)^{-3}\pmod9$. Indeed, if $n\equiv0$ or $1\pmod3$,
then $n(n-1)$ is divisible by six and $2^{n(n-1)}\equiv1\pmod9$. If
$n\equiv2\pmod3$, the binomial factor is divisible by three and
$2^{n(n-1)}-1$ is also divisible by three. Thus
\[
 (1-x)^3R(x)=1+9K(x),\qquad K\in x\mathbb Z[[x]].
\]
For every $j\ge1$, the coefficient
\[
 9^j\binom{1/3}{j}
   =\frac{3^j\prod_{i=0}^{j-1}(1-3i)}{j!}
\]
is an integer divisible by three. For a prime $\ell\ne3$, among $j$
successive terms of the numerator progression at least
$\lfloor j/\ell^r\rfloor$ are divisible by $\ell^r$, for every $r$.
This supplies the $\ell$-valuation of $j!$. At three the valuation is
$j-v_3(j!)\ge1$. Hence
$A=(1-x)^{-1}(1+9K)^{1/3}$ is integral and congruent to $(1-x)^{-1}$
modulo three.

Modulo two the defining equation is $A^3=1+x$. Its constant-term-one root
is unique, since the new coefficient at each positive degree has multiplier
$3=1$ in $\mathbb F_2$. Let
\[
 B=(1+x)\prod_{i\ge0}(1+x^{2\cdot4^i}).
\]
Then, coefficientwise,
\[
 B^3=(1+x)^3\prod_{r\ge1}(1+x^{2^r})
      =\frac{(1+x)^3}{1+x^2}=1+x.
\]
Thus $A=B$. Separating the factors with exponents $2$ and $8$ leaves
$\prod_{i\ge0}(1+x^{32\cdot4^i})$, the generating function for all
finite sums of distinct powers of four scaled by $32$, including the empty
sum. This is the corrected formula.
\end{proof}

\section{Arithmetic identities, valuations, and an exact correction}
\label{round9:r9sec:arithmetic}

The nine packages below concern nine identifiers absent from the 808
previously selected identifiers. Eight selected assertions are proved;
the ninth, a 2026 factorial-sum assertion on A048896, is refuted as written
and repaired for every index. The primary entries were inspected on
October 20, 2026. Each package states its selected scope explicitly.
In particular, the valuation results do not purport to settle the wider
prime-classification or asymptotic conjectures on the same pages.
The accompanying verifier distinguishes exact finite proof certificates
from numerical checks of the universal arguments.

\subsection{A026741: the prime-parameter gcd and lcm formulas}
\label{round9:r9entry:A026741}
\textbf{Source.} \oeis{A026741}, Timothy Hopper, February 26, 2017,
labels the following pair of formulas ``Conjecture''. For a prime $p$,
put $f_p(n)=n/p$ if $p\mid n$, and $f_p(n)=n$ otherwise.
The original sequence is the case $p=2$.

\begin{proposition}
For every prime $p$ and positive integer $n$,
\[
 \operatorname{lcm}(n,p)=p f_p(n),\qquad
 \gcd(n,p)=\frac{n}{f_p(n)}.
\]
\end{proposition}
\begin{proof}
If $p\mid n$, then $\gcd(n,p)=p$, $\operatorname{lcm}(n,p)=n$,
and $f_p(n)=n/p$. If $p\nmid n$, primality gives $\gcd(n,p)=1$,
$\operatorname{lcm}(n,p)=np$, and $f_p(n)=n$. Substitution proves
both identities in both cases.
\end{proof}
\textbf{Scope and prior work.} These are elementary consequences of the
definitions. The $p=2$ gcd and lcm formulas already occur elsewhere in the
entry and are credited; the selected package is the stated general-prime
assertion, counted once. The linked $p=3,5,7$ entries
\oeis{A051176}, \oeis{A060791}, and \oeis{A106608} were inspected as
companions and are not separate results. At $n=0$ the lcm identity remains
valid, but the printed gcd quotient is $0/0$; its meaningful domain is
$n\ge1$. Fermat's little theorem also immediately verifies Hopper's
associated formula involving $n^{p-1}\bmod p$.

\subsection{A040003: a gcd expression for the continued fraction of sqrt(6)}
\label{round9:r9entry:A040003}
\textbf{Source.} Michel Marcus's September 3, 2020 formula on
\oeis{A040003} is marked ``empirical''. Let $a(j)$, $j\ge0$, be the
successive partial quotients of $\sqrt6$.

\begin{proposition}
For every integer $n\ge1$,
\[
 a(n-1)=\gcd(2^n,3^n+1).
\]
\end{proposition}
\begin{proof}
Writing $s=\sqrt6$, we have $2<s<3$ and
\[
 \frac1{s-2}=\frac{s+2}{2},\qquad
 \frac1{(s+2)/2-2}=s+2.
\]
The latter two complete quotients have integer parts $2$ and $4$ and
alternate. Thus $\sqrt6=[2;\overline{2,4}]$. The proposed gcd is $2$
when $n=1$. For positive even $n$, $3^n+1\equiv2\pmod8$, so the
gcd is $2$. For odd $n\ge3$, $3^n+1\equiv4\pmod8$, so the gcd
is $4$. These are exactly $a(n-1)$ in the three cases.
\end{proof}
\textbf{Prior-work scope.} The continued fraction is classical and is the
definition of the entry. The inspected companion pages \oeis{A040001}
and \oeis{A010694} describe the related periodic sequences; they supplied
no proof of this gcd expression. No claim of novelty for the continued
fraction itself is made.

\subsection{A048896: a factorial-sum counterexample and its corrected limit}
\label{round9:r9entry:A048896}
\textbf{Source.} The entry \oeis{A048896} defines
$a(n)=2^{s_2(n+1)-1}$ for $n\ge0$, where $s_2$ counts binary ones.
Stefano Spezia's August 20, 2026 assertion, introduced by ``It appears
that'', identifies $a(n)$ with the reduced numerator of
\[
 R_n=\sum_{k=1}^{n+1}
       \frac{k^2}{(n+2-k)!(n+2+k)!}.
\]
Here and below the numerator of a rational number means the numerator
after reduction to lowest terms with positive denominator.

\begin{proposition}
The displayed assertion is false: $R_1=13/240$, while $a(1)=1$.
Changing the upper summation limit to $n+2$ gives a valid identity for
every $n\ge0$:
\[
 a(n)=\operatorname{num}\left(
       \sum_{k=1}^{n+2}\frac{k^2}{(n+2-k)!(n+2+k)!}\right).
\]
\end{proposition}
\begin{proof}
Direct exact arithmetic gives
\[
 R_1=\frac1{2!4!}+\frac4{1!5!}
     =\frac1{48}+\frac1{30}=\frac{13}{240}.
\]
Since $\gcd(13,240)=1$, this is a counterexample to the literal formula.
For the correction put $m=n+2$. The binomial second-moment identity,
obtained by differentiating $(1+x)^{2m}$ twice and once and evaluating at
$x=1$, is
\[
 \sum_{j=0}^{2m}(j-m)^2\binom{2m}{j}=m2^{2m-1}.
\]
Pairing $j=m-k$ and $j=m+k$ therefore yields
\[
 \sum_{k=1}^{m}\frac{k^2}{(m-k)!(m+k)!}
   =\frac{m2^{2m-2}}{(2m)!}
   =\frac{4^{n+1}}{2(2n+3)!}.
\]
Legendre's identity $v_2(N!)=N-s_2(N)$ follows by summing
$\lfloor N/2^r\rfloor$ over $r\ge1$. The exponent of two remaining
in the reduced numerator of the last fraction is consequently
\[
 2n+1-v_2((2n+3)!)
   =s_2(2n+3)-2=s_2(n+1)-1\ge0.
\]
Its unreduced numerator is a power of two, so no odd prime can remain
in the numerator. This proves the corrected identity.
\end{proof}
\textbf{Prior-work scope.} The older Lenard parity-sum observation is
not the counted target. Its content already appears as an unmarked
central-entry identity on \oeis{A261363} (August 16, 2015): the middle
partial sum in an even row of Pascal's triangle modulo two is A048896.
The present package counts only the literal 2026 factorial-sum assertion
and its exact correction, once.

\subsection{A084068: the square-triangular Pell gcd}
\label{round9:r9entry:A084068}
\textbf{Source.} Joseph A. Stocke's July 28, 2025 conjecture on
\oeis{A084068} states
\[
 a(n)=\gcd(\mathrm{A001108}(n),\mathrm{A001109}(n)).
\]
The source defines $a(1)=1$, $a(2)=2$,
$a(2k)=2a(2k-1)-a(2k-2)$, and
$a(2k+1)=4a(2k)-a(2k-1)$.

\begin{proposition}
For every $n\ge1$, the displayed gcd equals $a(n)$.
\end{proposition}
\begin{proof}
Define positive integers $u_n,v_n$ by
\[
 (1+\sqrt2)^n=u_n+v_n\sqrt2,
 \qquad u_n^2-2v_n^2=(-1)^n.
\]
Their norm shows $\gcd(u_n,v_n)=1$. Also $u_n$ is odd, since
$u_{n+1}=u_n+2v_n$ and $u_0=1$.
The standard Pell formulas in \oeis{A001108} and \oeis{A001109} give
\[
 T_n:=\mathrm{A001108}(n)=\frac{u_n^2+2v_n^2-1}{2},\qquad
 S_n:=\mathrm{A001109}(n)=u_nv_n.
\]
These formulas also follow by squaring $u_n+v_n\sqrt2$; they have
$T_0=S_0=0$, $T_1=S_1=1$ and the familiar recurrences
$T_{n+2}=6T_{n+1}-T_n+2$ and $S_{n+2}=6S_{n+1}-S_n$.
For even $n$ the norm identity gives $T_n=2v_n^2$, whence
\[
 \gcd(T_n,S_n)=v_n\gcd(2v_n,u_n)=v_n.
\]
For odd $n$ it gives $T_n=u_n^2$, whence
$\gcd(T_n,S_n)=u_n$.
Finally set $c_n=v_n$ for even $n$ and $c_n=u_n$ for odd $n$.
The relations $u_{n+1}=u_n+2v_n$ and $v_{n+1}=u_n+v_n$ imply
\[
 v_{2k}=2u_{2k-1}-v_{2k-2},\qquad
 u_{2k+1}=4v_{2k}-u_{2k-1}.
\]
Together with $c_1=1,c_2=2$, these are precisely the defining
recurrences of $a$. Thus $c_n=a(n)$ and the gcd assertion follows.
\end{proof}
\textbf{Prior-work scope.} Mathar's October 15, 2021 even/odd Pell
description of $a(n)$ is already on the primary entry and is credited.
The companion square-triangular pages were inspected. This package
concerns Stocke's gcd formula only; the separate conjecture about first
occurrences in A348295 is not selected.

\subsection{A096617: the reduced recurrence for harmonic numerators}
\label{round9:r9entry:A096617}
\textbf{Source.} M. F. Hasler's July 4, 2019 observation on
\oeis{A096617} is introduced by ``It appears that'' and asks whether
the recurrence is easy to prove. The sequence is the reduced numerator
of $nH_n$, where $H_n=\sum_{j=1}^n1/j$.

\begin{proposition}
Set $b_1=f_0=1$. For $n\ge1$ define
\[
 g_n=\gcd((n+1)b_n,nf_{n-1}),\quad
 f_n=\frac{nf_{n-1}}{g_n},\quad
 b_{n+1}=\frac{(n+1)b_n}{g_n}+f_n.
\]
Then $b_n$ is the reduced numerator of $nH_n$ for every $n\ge1$.
This initialization gives $g_1=f_1=1$, as in the source.
\end{proposition}
\begin{proof}
We prove the stronger invariant
\[
 nH_n=\frac{b_n}{f_{n-1}},\qquad \gcd(b_n,f_{n-1})=1.
\]
It holds at $n=1$. If it holds at $n$, then
\[
 (n+1)H_{n+1}
 =\frac{(n+1)b_n+nf_{n-1}}{nf_{n-1}}.
\]
The gcd of this numerator and denominator is
$\gcd((n+1)b_n,nf_{n-1})=g_n$. Dividing by it gives exactly
$b_{n+1}/f_n$ in lowest terms. Induction proves the assertion.
\end{proof}
\textbf{Scope and prior work.} The auxiliary $f_0=1$ makes explicit
the first denominator that is implicit in the printed initialization.
The standard Stirling/harmonic identity and the linked A000254 recurrence
are known inputs motivating the source. The inspected pages
\oeis{A000254} and \oeis{A193758} do not supply this reduction invariant.

\subsection{A329354: the summatory asymptotic, including its constant}
\label{round9:r9entry:A329354}
\textbf{Source.} Vaclav Kotesovec's March 3, 2023 conjecture on
\oeis{A329354} concerns $a(n)=\sum_{d\mid n}d\omega(d)$, where
$\omega$ counts distinct prime divisors. Let $B_1$ denote the prime
Mertens constant, \oeis{A077761}.

\begin{proposition}
As real $x\to\infty$,
\[
 \sum_{n\le x}a(n)
 =\frac{\pi^2x^2}{12}(\log\log x+B_1)+o(x^2).
\]
In particular the asymptotic equivalence in the source holds.
\end{proposition}
\begin{proof}
We use the classical prime-harmonic estimate
\[
 \sum_{p\le y}\frac1p=\log\log y+B_1+O(1/\log y)
\]
and $\pi(y)=O(y/\log y)$. These are established inputs; see
Axler, \emph{New estimates for some functions defined over primes},
\href{https://math.colgate.edu/~integers/s52/s52.pdf}{Integers 18 (2018),
A52}, equation (5.1) and the prime-counting estimates in its introduction.
Put $W(y)=\sum_{d\le y}d\omega(d)$. Counting each distinct prime
divisor once gives
\[
 W(y)=\sum_{p\le y}p\sum_{m\le y/p}m
 =\frac{y^2}{2}\sum_{p\le y}\frac1p+O(y\pi(y))
 =\frac{y^2}{2}(\log\log y+B_1)+o(y^2).
\]
The floor in the triangular sum introduces $O(y)$ per prime because
$y/p\ge1$. This argument also gives the uniform bound
$W(y)=O(y^2\log\log(3y))$ for $y\ge1$.

Changing the order of the finite divisor sum yields
\[
 \sum_{n\le x}a(n)=\sum_{m\le x}W(x/m).
\]
Let $M=\lfloor(\log x)^2\rfloor$. For $m\le M$, the error in the
formula for $W(x/m)$ is uniformly $o(x^2/m^2)$ because $x/M\to\infty$.
Moreover,
\[
 \log\log(x/m)=\log\log x+O(\log M/\log x)
 \quad(1\le m\le M).
\]
Since $\sum_{m\le M}m^{-2}=\zeta(2)+O(M^{-1})$, the contribution
from $m\le M$ is
$\tfrac12\zeta(2)x^2(\log\log x+B_1)+o(x^2)$.
The nonnegative tail is at most
\[
 O\left(x^2\log\log(3x)\sum_{m>M}m^{-2}\right)
 =O\left(\frac{x^2\log\log(3x)}{M}\right)=o(x^2).
\]
Finally $\zeta(2)=\pi^2/6$ proves the result.
\end{proof}
\textbf{Prior-work scope.} The entry already gives its divisor and prime
convolutions and its Dirichlet generating function. Those are credited,
not counted as discoveries. The inspected companions A328260, A329375,
and A323599 did not provide the selected summatory formula. Our stated
$o(x^2)$ remainder is stronger than the displayed asymptotic equivalence.
Numerical checks of finite sums are transcription checks, not evidence
used to prove an asymptotic limit.

\subsection*{A finite-block lemma for the next two packages}
Write
\[
 P_D(n)=\prod_{k=1}^n(1+ks)=A_D(n)+B_D(n)s,
 \qquad s^2=-D,\quad P_D(0)=1.
\]
All products are taken in the commutative ring
$\mathbb Z[s]/(s^2+D)$. Their coefficients are generated exactly by
\[
 A_D(n)=A_D(n-1)-DnB_D(n-1),\qquad
 B_D(n)=B_D(n-1)+nA_D(n-1).
\]
If $p$ is an odd prime and $-D$ is a quadratic nonresidue modulo $p$,
then
\begin{equation}\label{round9:r9arith:block}
 P_D(n+p)\equiv2P_D(n)\pmod p,\qquad
 P_D(n+p^2)\equiv2^pP_D(n)\pmod {p^2}.
\end{equation}
Indeed, over $\mathbb F_p$,
$\prod_{a=0}^{p-1}(1+as)=1-s^{p-1}=2$.
For a full block modulo $p^2$, group the factors as $k=a+pt$,
$1\le a\le p$, $0\le t<p$. Expansion gives
\[
 \prod_{t=0}^{p-1}(1+(a+pt)s)\equiv(1+as)^p\pmod {p^2}:
\]
every term of order at least two in $pt$ is divisible by $p^2$, and
the linear coefficient is divisible by
$p\sum_{t=0}^{p-1}t=p^2(p-1)/2$. Thus the full block is congruent to
$C^p$, where $C=\prod_{a=1}^p(1+as)\equiv2\pmod p$, and hence to
$2^p\pmod {p^2}$. Any $p$ or $p^2$ consecutive factors have the same
residue multiset, proving both shift formulas for every $n\ge0$.

\subsection{A363409: the complete 7-adic valuation rule}
\label{round9:r9entry:A363409}
\textbf{Source.} The entry \oeis{A363409} (Peter Bala, June 1, 2023)
defines $a(n)=A_2(n)$ and says ``It appears that'' before the rule
$v_7(a(7n+3))=1$, with valuation zero at all other indices.
The assertion has no separate signature or date within the entry.

\begin{proposition}
For every $n\ge0$,
\[
 v_7(A_2(n))=
 \begin{cases}1,&n\equiv3\pmod7,\\0,&n\not\equiv3\pmod7.\end{cases}
\]
\end{proposition}
\begin{proof}
Here $-2$ is a nonresidue modulo $7$, so
\eqref{round9:r9arith:block} gives the multipliers $2$ modulo $7$ and
$2^7\equiv30\pmod {49}$. The exact coefficient recurrence in the
lemma gives the following small certificates:
\[
 \begin{array}{c|rrrrrrr}
 r&0&1&2&3&4&5&6\\\hline
 A_2(r)\bmod7&1&1&4&0&6&3&2\\
 \bigl(A_2(7r+3)\bmod49\bigr)/7&4&6&4&5&2&5&4
 \end{array}
\]
In the second row of values, the residue modulo $49$ is chosen in
$\{0,\ldots,48\}$ before division. The first shift formula and first
row show that divisibility by $7$ occurs exactly at indices $3\pmod7$.
Write any such index uniquely as $49q+7r+3$ with $0\le r<7$.
The second shift formula multiplies its residue modulo $49$ by the
unit $30^q$. Every entry in the second row is nonzero modulo $7$,
so none of these values is divisible by $49$. This proves the exact
valuation at every index.
\end{proof}
\textbf{Audit and verification.} The displayed 14 residues are generated
in integer arithmetic by \texttt{r9\_arith\_verify.py}; together with
the proved shifts they are a finite certificate for the universal result.
The linked \href{https://www.tulane.edu/~vhm/papers_html/xn-final.pdf}
{Moll paper (2012)} studies $D=1$. Its statements and examples do not
give the selected $D=2$, $p=7$ rule. Other prime classifications on the
OEIS entry remain outside this package.

\subsection{A363411: the complete 5-adic valuation rule}
\label{round9:r9entry:A363411}
\textbf{Source.} \oeis{A363411} (Peter Bala, June 1, 2023) defines
$a(n)=A_3(n)$ and explicitly conjectures
$v_5(a(n+2))=\mathrm{A079998}(n)$, reporting checks through 5000.
The assertion has no separate signature or date within the entry.
The companion \oeis{A079998} is the characteristic function of
multiples of five, with offset zero.

\begin{proposition}
For every $n\ge0$,
\[
 v_5(A_3(n))=
 \begin{cases}1,&n\equiv2\pmod5,\\0,&n\not\equiv2\pmod5.\end{cases}
\]
Consequently the source's shifted identity holds for every $n\ge0$.
\end{proposition}
\begin{proof}
Here $-3$ is a nonresidue modulo $5$, so
\eqref{round9:r9arith:block} gives multipliers $2$ modulo $5$ and
$2^5\equiv7\pmod {25}$. Exact multiplication gives
\[
 \begin{array}{c|rrrrr}
 r&0&1&2&3&4\\\hline
 A_3(r)\bmod5&1&1&0&3&2\\
 \bigl(A_3(5r+2)\bmod25\bigr)/5&4&4&2&1&4
 \end{array}
\]
The first shift and first row give exactly the residue class $2\pmod5$
for divisibility by $5$. Each such index is $25q+5r+2$ with
$0\le r<5$. The second shift multiplies the corresponding entry
modulo $25$ by the unit $7^q$; all five quotients in the table are
nonzero modulo $5$. Thus the valuation is exactly one there.
\end{proof}
\textbf{Audit and verification.} The ten table entries and full-block
multiplier are reproduced exactly in the verifier. This closes the
unbounded shifted valuation assertion, rather than merely extending its
5000-term check. The linked Moll paper is credited as background for
$D=1$; the wider $D=3$ prime-type conjectures are not selected.

\subsection{A363412: a stronger exact 3-adic valuation formula}
\label{round9:r9entry:A363412}
\textbf{Source.} \oeis{A363412} (Peter Bala, June 1, 2023) defines
$a(n)=B_3(n)$ and labels two statements ``Conjectures'':
$3\nmid a(3n+1)$, and $v_3(a(3n))=v_3(a(3n-1))$.
The second statement is meaningful for $n\ge1$; the first includes
$n=0$. These assertions have no separate signature or date in the entry.

\begin{proposition}
For every $n\ge1$,
\[
 v_3(B_3(n))=v_3\left(\frac{n(n+1)}2\right).
\]
In particular $B_3(n)\ne0$ for $n\ge1$, and both selected conjectures
hold on their stated meaningful domains.
\end{proposition}
\begin{proof}
Let $e_r=e_r(1,2,\ldots,n)$ denote an elementary symmetric function
and put $T_n=n(n+1)/2$. We use the established Hultman-number formula
\[
 H(N,k)=\frac{c(N+2,k)}{\binom{N+2}{2}}
 \quad\text{if }N-k\text{ is odd},
\]
where $H(N,k)$ counts permutations and $c$ is an unsigned Stirling
number of the first kind. See B\'ona and Flynn,
\href{https://arxiv.org/pdf/0811.0740}{arXiv:0811.0740}, Section 3,
the final displayed formula; \oeis{A164652} also records this known
formula. Since
\[
 e_{2j+1}=c(n+1,n-2j),\qquad
 (n-1)-(n-2j)=2j-1,
\]
it follows that $e_{2j+1}/T_n$ is an integer whenever $2j+1\le n$.
Also $e_1/T_n=1$.

Extracting the coefficient of $s$ in $P_3(n)$ gives
\[
 \frac{B_3(n)}{T_n}
 =\sum_{j\ge0}(-3)^j\frac{e_{2j+1}}{T_n}
 \equiv1\pmod3.
\]
The finite sum on the right is an integer not divisible by three.
This proves the valuation formula and nonvanishing.
At $n=3m+1$, neither $n$ nor $n+1$ is divisible by three.
For $m\ge1$ the two triangular numbers at $3m$ and $3m-1$ have
the same 3-adic valuation $1+v_3(m)$. The two source conjectures
follow.
\end{proof}
\textbf{Prior-work scope.} Hultman integrality is a credited classical
input, not a new theorem of this report. The linked Moll paper concerns
the related $D=1$ product. No assertion is made here about the source's
global Type 2 or Type 3 prime classification. The empty-Type-1 statement
is also not counted: its literal version includes the zero term $a(0)=0$.

\section{Generating functions and exact algebra: 19 selected packages}
\label{round9:sec:r9gf}
These nineteen packages concern primary OEIS assertions carrying an explicit
conjectural or empirical marker on 20 October 2026.  Seventeen packages prove
the selected assertion with its index range stated precisely.  A001784 and
A052736 instead correct an index or boundary error, give a counterexample
to the unqualified printed version, and prove the corrected statement.
Existing generating functions, formulas and combinatorial interpretations
are credited as inputs.  One A-number contributes one package, even if both
its generating function and recurrence are established.  No claim of first
publication, independent research novelty, or OEIS acceptance is made.

\subsection{A finite exact certificate for every recurrence}
Put $\theta=x\,d/dx$ and extend $a_n$ by zero below its stated offset.
If $F(x)=\sum a_nx^n$ and $p_0,\ldots,p_r$ are polynomials, then
\[
 L=\sum_{j=0}^r x^jp_j(\theta+j),\qquad
 [x^n]LF=\sum_{j=0}^r p_j(n)a_{n-j}.
\]
Consequently an exact identity $LF=P(x)$, with $P$ polynomial, proves
the recurrence for every $n>\deg P$ at which the sequence indices exist.
This is an algebraic proof for all indices, not a finite numerical test.

For an exponential generating function $F=\sum a_nx^n/n!$, replace
$p_j(n)$ by $q_j(n)=p_j(n)(n-j)_{r-j}$, where $(v)_h$ is the falling
factorial.  For $n\ge r$ the coefficient of the resulting operator is
$\sum_jp_j(n)a_{n-j}/(n-r)!$.  The same polynomial-residual argument
therefore applies.

For a radical $F=R_0+R_1\sqrt D$, use
\[
 \theta(R_0+R_1\sqrt D)=xR_0'
 +x\left(R_1'+\frac{D'}{2D}R_1\right)\sqrt D.
\]
For a rational parameter $x=X(t)$, replace $x\,d/dx$ by
$X(t)/X'(t)\,d/dt$.  Every table below gives the actual $p_j$, the actual
series and branch, and the exact residual.  All coefficients lie in a
rational-function field; the scripts clear denominators and verify zero
polynomials.  A displayed radical is the branch with constant term $+1$.
Removable negative powers in the rational components do not introduce
negative coefficients in the complete formal series.

The independently calculated prefixes are checks on source indexing and
transcription.  They are not substitutes for the identities above.
\subsection{A000483}
\label{round9:r9entry:A000483}
\noindent\textbf{Source and selected claim.} \href{https://oeis.org/A000483}{A000483}, checked 20 October 2026. Conjecture. The selected polynomial recurrence is
\[\sum_{j=0}^{3}p_j(n)a_{n-j}=0\qquad(n\ge 9).\]
\begin{align*}
p_{0}(n)&=\left(n - 4\right) \left(n - 2\right),\\
p_{1}(n)&=- \left(n - 1\right) \left(3 n^{2} - 21 n + 35\right),\\
p_{2}(n)&=\left(n - 2\right) \left(n - 1\right) \left(3 n^{2} - 24 n + 47\right),\\
p_{3}(n)&=- \left(n - 5\right) \left(n - 4\right) \left(n - 3\right) \left(n - 2\right) \left(n - 1\right),\\
\end{align*}
\noindent\textbf{Offset and initial data.} Offset $6$.
\[\left(a_{6},\ldots,a_{13}\right)=
\left(15,210,2380,26432,303660,3678840,47324376,647536032\right).\]
\noindent\textbf{Proof.} 
The defining objects are permutations with exactly three cycles, each of
length at least two.  A labelled cycle of length $j$ contributes $x^j/j$;
the unordered set of three cycles therefore has exponential generating
function
\[
 F(x)=\frac{1}{6}\left(\sum_{j\ge2}\frac{x^j}{j}\right)^3
 =-\frac{(x+\log(1-x))^3}{6}.
\]
This argument derives the series from the objects, without using either
conjectural recurrence.  Put $\ell=x+\log(1-x)$, so
$\ell'=-x/(1-x)$.  In the basis $1,\ell,\ell^2,\ell^3$ the coefficient
vector of $F$ is $(0,0,0,-1/6)$.  For a vector $R$ the Euler derivative has
component $x(R_i'+(i+1)\ell'R_{i+1})$, with $R_4=0$.  Applying the
exponential version of the certificate lemma above gives the displayed
exact residual, and hence the recurrence over the stated full range.

\noindent\textbf{Exact certificate.} Using $q_j(n)=p_j(n)(n-j)_{3-j}$, $LF=20 x^{6}$; all three positive log-power residuals are zero.
\noindent\textbf{Verification.} The standalone verifier reproduces the exact identity, compares 20 published terms with an independent definition calculation, and checks all applicable indices through $100$.\par\medskip
\subsection{A001784}
\label{round9:r9entry:A001784}
\noindent\textbf{Source and selected claim.} \href{https://oeis.org/A001784}{A001784}, checked 20 October 2026. Conjecture. The printed first-order recurrence requires the index correction proved below.
\[\sum_{j=0}^{1}p_j(n)a_{n-j}=0\qquad(n\ge 2).\]
\begin{align*}
p_{0}(n)&=\left(n - 1\right) \left(20 n^{2} + 35 n + 12\right),\\
p_{1}(n)&=- \left(n + 1\right) \left(2 n + 3\right) \left(20 n^{2} + 75 n + 67\right),\\
\end{align*}
\noindent\textbf{Offset and initial data.} Offset $0$.
\[\left(a_{0},\ldots,a_{7}\right)=
\left(1,24,924,26432,705320,18858840,520059540,14980405440\right).\]
\noindent\textbf{Proof.} 
Here $a_0=1$ is an added initial convention, and for $n\ge1$ the sequence
counts permutations of $2n+3$ elements having $n$ cycles and no fixed points.
The source's first-order conjecture was entered in 2015, before the offset
was changed to zero in 2016.  With the current indexing its residual at
$n=3$ is
\[
 162a_3-6237a_2=162\cdot26432-6237\cdot924=-1481004.
\]
Thus that printed formula is false with the current offset.  The corrected
coefficients displayed here result from replacing its coefficient variable
$n$ by $n+1$ while leaving $a_n,a_{n-1}$ in place.

To prove the correction, subtract two from every cycle length.  The excess
is three, so the only possibilities are one 5-cycle; one 4-cycle and one
3-cycle; or three 3-cycles.  All other cycles have length two.  The standard
labelled cycle count gives
\[
 a_n=(2n+3)!\left\{
 \frac{1}{5\,2^{n-1}(n-1)!}
 +\frac{1}{12\,2^{n-2}(n-2)!}
 +\frac{1}{162\,2^{n-3}(n-3)!}\right\},
\]
where a term with a negative factorial is omitted.  Factoring the first
denominator and using the factors $n-1$ and $(n-1)(n-2)$ gives, for every
$n\ge1$,
\[
 a_n=\frac{(20n^2+75n+67)(2n+3)!}{405\,2^n(n-1)!}.
\]
This closed form is already credited to Kotesovec (2016) in the entry;
the present package is an index correction, not a claim of a new closed
form or a newly discovered recurrence.  Dividing successive expressions,
and clearing the nonzero denominators for $n\ge2$, proves the correction.

\noindent\textbf{Verification.} The standalone verifier reproduces the exact identity, compares 16 published terms with an independent definition calculation, and checks all applicable indices through $100$.\par\medskip
\subsection{A025256}
\label{round9:r9entry:A025256}
\noindent\textbf{Source and selected claim.} \href{https://oeis.org/A025256}{A025256}, checked 20 October 2026. Conjecture. The selected polynomial recurrence is
\[\sum_{j=0}^{4}p_j(n)a_{n-j}=0\qquad(n\ge 6).\]
\begin{align*}
p_{0}(n)&=n,\\
p_{1}(n)&=3 - 2 n,\\
p_{2}(n)&=3 - n,\\
p_{3}(n)&=- 5 \left(2 n - 9\right),\\
p_{4}(n)&=n - 6,\\
\end{align*}
\noindent\textbf{Offset and initial data.} Offset $1$.
\[\left(a_{1},\ldots,a_{8}\right)=
\left(1,1,3,3,6,18,42,105\right).\]
\noindent\textbf{Proof.} 
The source defines $a_1=1$, $a_2=1$, $a_3=3$ and
$a_n=\sum_{j=1}^{n-3}a_j a_{n-j}$ for $n\ge4$.
Set $F=\sum_{n\ge1}a_nx^n$.  The omitted last two terms in the
full convolution are $a_1a_{n-1}+a_2a_{n-2}$.  Summing over
$n\ge4$ and retaining the three initial coefficients gives
\[
 F^2-(1+1x+(1)x^2)F+1x+(1)x^2+3x^3=0.
\]
The branch with zero constant term is precisely the radical displayed
below.  Inserting it in the certificate lemma proves the selected
linear recurrence, independent of the nonlinear recurrence used to define
the terms.  The other quadratic branch has constant term one and is excluded.

\noindent\textbf{Exact certificate.} $F=R_0+R_1\sqrt D$, with\begin{align*}
D&=x^{4} - 10 x^{3} - x^{2} - 2 x + 1,\\
R_0&=\frac{x^{2} + x + 1}{2},\\
R_1&=- \frac{1}{2},\\
\end{align*}
Exact substitution gives $LF=P(x)=- 3 x^{5} + x^{4} + 6 x^{3} + x^{2} + x$. The degree-five boundary term is why the recurrence starts at $n=6$, not $n=5$. The coefficient of the radical in the residual is identically zero.
\noindent\textbf{Verification.} The standalone verifier reproduces the exact identity, compares 30 published terms with an independent definition calculation, and checks all applicable indices through $100$.\par\medskip
\subsection{A025260}
\label{round9:r9entry:A025260}
\noindent\textbf{Source and selected claim.} \href{https://oeis.org/A025260}{A025260}, checked 20 October 2026. Conjecture. The selected polynomial recurrence is
\[\sum_{j=0}^{4}p_j(n)a_{n-j}=0\qquad(n\ge 6).\]
\begin{align*}
p_{0}(n)&=n,\\
p_{1}(n)&=- 2 \left(2 n - 3\right),\\
p_{2}(n)&=8 \left(n - 3\right),\\
p_{3}(n)&=- 6 \left(2 n - 9\right),\\
p_{4}(n)&=4 \left(n - 6\right),\\
\end{align*}
\noindent\textbf{Offset and initial data.} Offset $1$.
\[\left(a_{1},\ldots,a_{8}\right)=
\left(2,-2,1,2,2,1,2,10\right).\]
\noindent\textbf{Proof.} 
The source defines $a_1=2$, $a_2=-2$, $a_3=1$ and
$a_n=\sum_{j=1}^{n-3}a_j a_{n-j}$ for $n\ge4$.
Set $F=\sum_{n\ge1}a_nx^n$.  The omitted last two terms in the
full convolution are $a_1a_{n-1}+a_2a_{n-2}$.  Summing over
$n\ge4$ and retaining the three initial coefficients gives
\[
 F^2-(1+2x+(-2)x^2)F+2x+(-2)x^2+1x^3=0.
\]
The branch with zero constant term is precisely the radical displayed
below.  Inserting it in the certificate lemma proves the selected
linear recurrence, independent of the nonlinear recurrence used to define
the terms.  The other quadratic branch has constant term one and is excluded.

\noindent\textbf{Exact certificate.} $F=R_0+R_1\sqrt D$, with\begin{align*}
D&=4 x^{4} - 12 x^{3} + 8 x^{2} - 4 x + 1,\\
R_0&=- \frac{2 x^{2} - 2 x - 1}{2},\\
R_1&=- \frac{1}{2},\\
\end{align*}
Exact substitution gives $LF=P(x)=2 x^{5} - 6 x^{4} + 15 x^{3} - 8 x^{2} + 2 x$. The degree-five boundary term is why the recurrence starts at $n=6$, not $n=5$. The coefficient of the radical in the residual is identically zero.
\noindent\textbf{Verification.} The standalone verifier reproduces the exact identity, compares 35 published terms with an independent definition calculation, and checks all applicable indices through $100$.\par\medskip
\subsection{A026023}
\label{round9:r9entry:A026023}
\noindent\textbf{Source and selected claim.} \href{https://oeis.org/A026023}{A026023}, checked 20 October 2026. Conjecture. The selected polynomial recurrence is
\[\sum_{j=0}^{3}p_j(n)a_{n-j}=0\qquad(n\ge 3).\]
\begin{align*}
p_{0}(n)&=\left(n - 1\right) \left(n + 4\right),\\
p_{1}(n)&=\left(n - 1\right) \left(n + 1\right),\\
p_{2}(n)&=- 2 \left(n + 1\right) \left(2 n + 1\right),\\
p_{3}(n)&=- 4 \left(n - 1\right) \left(n + 1\right),\\
\end{align*}
\noindent\textbf{Offset and initial data.} Offset $0$.
\[\left(a_{0},\ldots,a_{7}\right)=
\left(1,2,4,8,15,30,56,112\right).\]
\noindent\textbf{Proof.} 
For a nonnegative nearest-neighbor walk starting at height three, the
reflection principle gives
$T(n,k)=\binom nk-\binom n{k-4}$ when the terminal height is
$n+3-2k\ge0$.  The reflected forbidden walks start at $-5$, which explains
the difference of four in the binomial indices.  Summing in $k$ telescopes
to four neighboring binomial coefficients.  Pascal's identity yields
\[
 a_{2m}=\binom{2m+2}{m},\qquad a_{2m+1}=2\binom{2m+2}{m}.
\]
These bisections are known in the entry; the reflection argument also
provides their justification from the definition.  Writing
$C(z)=(1-\sqrt{1-4z})/(2z)$, the binomial generating function is
$C(z)^2/\sqrt{1-4z}$.  Consequently
$F(x)=(1+2x)C(x^2)^2/\sqrt{1-4x^2}$, which reduces to the displayed
radical coefficients.  The exact differential residual proves the
conjectured mixed-parity recurrence.

\noindent\textbf{Exact certificate.} $F=R_0+R_1\sqrt D$, with\begin{align*}
D&=- \left(2 x - 1\right) \left(2 x + 1\right),\\
R_0&=- \frac{2 x + 1}{2 x^{4}},\\
R_1&=\frac{2 x^{2} - 1}{2 x^{4} \left(2 x - 1\right)},\\
\end{align*}
Exact substitution gives $LF=P(x)=-4$. The coefficient of the radical in the residual is identically zero.
\noindent\textbf{Verification.} The standalone verifier reproduces the exact identity, compares 34 published terms with an independent definition calculation, and checks all applicable indices through $100$.\par\medskip
\subsection{A026027}
\label{round9:r9entry:A026027}
\noindent\textbf{Source and selected claim.} \href{https://oeis.org/A026027}{A026027}, checked 20 October 2026. Conjecture. The selected polynomial recurrence is
\[\sum_{j=0}^{3}p_j(n)a_{n-j}=0\qquad(n\ge 5).\]
\begin{align*}
p_{0}(n)&=\left(n + 5\right) \left(13 n + 118\right),\\
p_{1}(n)&=13 n^{2} - 969 n - 3490,\\
p_{2}(n)&=- 2 \left(203 n^{2} - 1080 n - 758\right),\\
p_{3}(n)&=4 \left(2 n - 5\right) \left(73 n - 182\right),\\
\end{align*}
\noindent\textbf{Offset and initial data.} Offset $2$.
\[\left(a_{2},\ldots,a_{9}\right)=
\left(1,5,21,84,329,1274,4900,18768\right).\]
\noindent\textbf{Proof.} 
Use the equivalent triangle definition in the source: the path has
$2n-1$ nearest-neighbor steps from height three to height six and remains
nonnegative.  The title's step-condition range ending at $n$ is a typo;
the triangle and the binomial formula specify all $2n-1$ steps.
Reflection gives
\[
 a_n=\binom{2n-1}{n-2}-\binom{2n-1}{n-6},\quad n\ge2,
\]
where an out-of-range binomial coefficient is zero.  For $q\ge0$,
$\sum_{m\ge0}\binom{2m+q}{m}x^m=C(x)^q/\sqrt{1-4x}$.
Apply this identity with $q=3,11$ to the two sums.  It follows that
\[
 F(x)=\frac{x^2 C(x)^3-x^6 C(x)^{11}}{\sqrt{1-4x}}
      =x^2\{C(x)^5+x^2C(x)^9\}.
\]
The source writes the latter expression without $x^2$, because that
display starts its exponent at the sequence offset; here exponents are
the actual indices.  Put $t=C(x)-1$, so $x=t/(1+t)^2$ and
$F=t^2(1+t)(1+t^2)$.  The rational parameter is locally invertible at
zero.  The parameter version of the certificate lemma now proves the
recurrence for every admissible index.

\noindent\textbf{Exact certificate.} $x=\frac{t}{\left(t + 1\right)^{2}}$, $F=t^{2} \left(t + 1\right) \left(t^{2} + 1\right)$.
Exact substitution gives $LF=P(x)=1008 x^{2}$. The coefficient of the radical in the residual is identically zero.
\noindent\textbf{Verification.} The standalone verifier reproduces the exact identity, compares 24 published terms with an independent definition calculation, and checks all applicable indices through $100$.\par\medskip
\subsection{A026570}
\label{round9:r9entry:A026570}
\noindent\textbf{Source and selected claim.} \href{https://oeis.org/A026570}{A026570}, checked 20 October 2026. Conjecture. The selected polynomial recurrence is
\[\sum_{j=0}^{3}p_j(n)a_{n-j}=0\qquad(n\ge 4).\]
\begin{align*}
p_{0}(n)&=n + 1,\\
p_{1}(n)&=- 2 n,\\
p_{2}(n)&=- 3 n - 1,\\
p_{3}(n)&=2 \left(2 n - 3\right),\\
\end{align*}
\noindent\textbf{Offset and initial data.} Offset $1$.
\[\left(a_{1},\ldots,a_{8}\right)=
\left(1,1,4,7,20,43,111,259\right).\]
\noindent\textbf{Proof.} 
The parent triangle A026568 counts walks on the integers starting at zero,
with steps $+1,-1$, and with a zero step allowed only at an even height.
Let $e_j,o_j$ be the generating functions for endpoint $2j,2j+1$.
Arrival at a vertex gives
\[
 (1-x)e_j=\delta_{j0}+x(o_{j-1}+o_j),\qquad
 o_j=x(e_j+e_{j+1}).
\]
Set $d=(1-x)(1-x-4x^2)$, $h=1-x-2x^2$, and
$r=(h-\sqrt d)/(2x^2)=O(x^2)$.  Eliminating the odd vertices and using
symmetry about zero gives
\[
 e_j=\frac{r^{|j|}}{\sqrt d},\qquad o_j=x(e_j+e_{j+1}).
\]
Indeed $x^2r^2-hr+x^2=0$, and the equation at zero reduces to
$(h-2x^2r)/\sqrt d=1$.  This is the unique formal walk solution:
the coefficient of degree $N$ depends only on paths of length at most
$N$, hence on finitely many vertices, and the arrival equations determine
it inductively.  Thus no analytic convergence or unproved recurrence is
being assumed.

For endpoint one, $F=o_0=x(1+r)/\sqrt d$
$=((1-x)/\sqrt d-1)/(2x)$.  This produces the displayed radical
certificate.  In particular it also proves the entry's conditional
identity with the difference of A026569, but that identity is not counted
as an additional package.

\noindent\textbf{Exact certificate.} $F=R_0+R_1\sqrt D$, with\begin{align*}
D&=\left(x - 1\right) \left(4 x^{2} + x - 1\right),\\
R_0&=- \frac{1}{2 x},\\
R_1&=- \frac{1}{2 x \left(4 x^{2} + x - 1\right)},\\
\end{align*}
Exact substitution gives $LF=P(x)=- x^{2} + 2 x$. The coefficient of the radical in the residual is identically zero.
\noindent\textbf{Verification.} The standalone verifier reproduces the exact identity, compares 29 published terms with an independent definition calculation, and checks all applicable indices through $100$.\par\medskip
\subsection{A026572}
\label{round9:r9entry:A026572}
\noindent\textbf{Source and selected claim.} \href{https://oeis.org/A026572}{A026572}, checked 20 October 2026. Conjecture. The selected polynomial recurrence is
\[\sum_{j=0}^{5}p_j(n)a_{n-j}=0\qquad(n\ge 8).\]
\begin{align*}
p_{0}(n)&=n + 3,\\
p_{1}(n)&=- 3 n - 5,\\
p_{2}(n)&=- 4 \left(n + 1\right),\\
p_{3}(n)&=13 n - 9,\\
p_{4}(n)&=5 n - 9,\\
p_{5}(n)&=- 6 \left(2 n - 7\right),\\
\end{align*}
\noindent\textbf{Offset and initial data.} Offset $3$.
\[\left(a_{3},\ldots,a_{10}\right)=
\left(1,2,8,19,56,140,376,953\right).\]
\noindent\textbf{Proof.} 
The parent triangle A026568 counts walks on the integers starting at zero,
with steps $+1,-1$, and with a zero step allowed only at an even height.
Let $e_j,o_j$ be the generating functions for endpoint $2j,2j+1$.
Arrival at a vertex gives
\[
 (1-x)e_j=\delta_{j0}+x(o_{j-1}+o_j),\qquad
 o_j=x(e_j+e_{j+1}).
\]
Set $d=(1-x)(1-x-4x^2)$, $h=1-x-2x^2$, and
$r=(h-\sqrt d)/(2x^2)=O(x^2)$.  Eliminating the odd vertices and using
symmetry about zero gives
\[
 e_j=\frac{r^{|j|}}{\sqrt d},\qquad o_j=x(e_j+e_{j+1}).
\]
Indeed $x^2r^2-hr+x^2=0$, and the equation at zero reduces to
$(h-2x^2r)/\sqrt d=1$.  This is the unique formal walk solution:
the coefficient of degree $N$ depends only on paths of length at most
$N$, hence on finitely many vertices, and the arrival equations determine
it inductively.  Thus no analytic convergence or unproved recurrence is
being assumed.

Endpoint three has $F=o_1=x(r+r^2)/\sqrt d$.  Reduction using
$r=(1-x-2x^2-\sqrt d)/(2x^2)$ gives
\[
 F=\frac{(1-x-3x^2)\sqrt{(1-x)/(1-x-4x^2)}-1+x+x^2}{2x^3},
\]
which is exactly the conjectured generating function (its apparent pole
is removable).  The separate conjectural recurrence follows from the
exact residual below.  These two linked assertions constitute one package.

\noindent\textbf{Exact certificate.} $F=R_0+R_1\sqrt D$, with\begin{align*}
D&=\left(x - 1\right) \left(4 x^{2} + x - 1\right),\\
R_0&=\frac{x^{2} + x - 1}{2 x^{3}},\\
R_1&=\frac{3 x^{2} + x - 1}{2 x^{3} \left(4 x^{2} + x - 1\right)},\\
\end{align*}
Exact substitution gives $LF=P(x)=- 3 x^{4} + 6 x^{3}$. The coefficient of the radical in the residual is identically zero.
\noindent\textbf{Verification.} The standalone verifier reproduces the exact identity, compares 27 published terms with an independent definition calculation, and checks all applicable indices through $100$.\par\medskip
\subsection{A026582}
\label{round9:r9entry:A026582}
\noindent\textbf{Source and selected claim.} \href{https://oeis.org/A026582}{A026582}, checked 20 October 2026. Conjecture. The selected polynomial recurrence is
\[\sum_{j=0}^{6}p_j(n)a_{n-j}=0\qquad(n\ge 6).\]
\begin{align*}
p_{0}(n)&=n,\\
p_{1}(n)&=- 2 \left(2 n - 1\right),\\
p_{2}(n)&=- 2 \left(n - 3\right),\\
p_{3}(n)&=20 \left(n - 2\right),\\
p_{4}(n)&=6 - 7 n,\\
p_{5}(n)&=- 2 \left(12 n - 41\right),\\
p_{6}(n)&=8 \left(2 n - 7\right),\\
\end{align*}
\noindent\textbf{Offset and initial data.} Offset $0$.
\[\left(a_{0},\ldots,a_{7}\right)=
\left(1,3,8,20,50,125,314,792\right).\]
\noindent\textbf{Proof.} 
The parent triangle A026568 counts walks on the integers starting at zero,
with steps $+1,-1$, and with a zero step allowed only at an even height.
Let $e_j,o_j$ be the generating functions for endpoint $2j,2j+1$.
Arrival at a vertex gives
\[
 (1-x)e_j=\delta_{j0}+x(o_{j-1}+o_j),\qquad
 o_j=x(e_j+e_{j+1}).
\]
Set $d=(1-x)(1-x-4x^2)$, $h=1-x-2x^2$, and
$r=(h-\sqrt d)/(2x^2)=O(x^2)$.  Eliminating the odd vertices and using
symmetry about zero gives
\[
 e_j=\frac{r^{|j|}}{\sqrt d},\qquad o_j=x(e_j+e_{j+1}).
\]
Indeed $x^2r^2-hr+x^2=0$, and the equation at zero reduces to
$(h-2x^2r)/\sqrt d=1$.  This is the unique formal walk solution:
the coefficient of degree $N$ depends only on paths of length at most
$N$, hence on finitely many vertices, and the arrival equations determine
it inductively.  Thus no analytic convergence or unproved recurrence is
being assumed.

The definition sums, through row $n$, the half-row with nonnegative
endpoint.  Reflection in height zero shows that its row sum is half the
sum of all walks plus the number ending at zero.  If $E,O$ sum over all
even and odd endpoints, then $(1-x)E=1+2xO$ and $O=2xE$.
Thus $E+O=(1+2x)/(1-x-4x^2)$ and $e_0=1/\sqrt d$.  Taking cumulative
sums gives
\[
 F=\frac{1}{2(1-x)}\left\{\frac{1+2x}{1-x-4x^2}
                      +\frac1{\sqrt d}\right\}.
\]
This derivation fixes both the half-row boundary and the initial term.
The displayed certificate proves the selected recurrence.

\noindent\textbf{Exact certificate.} $F=R_0+R_1\sqrt D$, with\begin{align*}
D&=\left(x - 1\right) \left(4 x^{2} + x - 1\right),\\
R_0&=\frac{2 x + 1}{2 \left(x - 1\right) \left(4 x^{2} + x - 1\right)},\\
R_1&=- \frac{1}{2 \left(x - 1\right)^{2} \left(4 x^{2} + x - 1\right)},\\
\end{align*}
Exact substitution gives $LF=P(x)=2 x^{4} + x$. The coefficient of the radical in the residual is identically zero.
\noindent\textbf{Verification.} The standalone verifier reproduces the exact identity, compares 28 published terms with an independent definition calculation, and checks all applicable indices through $100$.\par\medskip
\subsection{A026732}
\label{round9:r9entry:A026732}
\noindent\textbf{Source and selected claim.} \href{https://oeis.org/A026732}{A026732}, checked 20 October 2026. Conjecture. The selected polynomial recurrence is
\[\sum_{j=0}^{8}p_j(n)a_{n-j}=0\qquad(n\ge 8).\]
\begin{align*}
p_{0}(n)&=1 - n,\\
p_{1}(n)&=2,\\
p_{2}(n)&=3 \left(3 n - 7\right),\\
p_{3}(n)&=-10,\\
p_{4}(n)&=95 - 23 n,\\
p_{5}(n)&=6,\\
p_{6}(n)&=11 n - 95,\\
p_{7}(n)&=2,\\
p_{8}(n)&=4 \left(n - 7\right),\\
\end{align*}
\noindent\textbf{Offset and initial data.} Offset $0$.
\[\left(a_{0},\ldots,a_{7}\right)=
\left(1,2,4,9,18,40,80,176\right).\]
\noindent\textbf{Proof.} 
In A026725 the Pascal rule has an extra $T(2m-1,m-1)$ only at
position $(2m+1,m)$ for $m\ge1$.  Boundary entries remain one.
Put $b_m=T(2m+1,m)$ and $P_n(u)=\sum_kT(n,k)u^k$.
Propagation of each extra entry by Pascal's rule gives
\[
 P_n(u)=(1+u)^n+
 \sum_{\substack{j\ge1\\2j+1\le n}}
 b_{j-1}u^j(1+u)^{n-2j-1}.
\]
The corresponding row-polynomial identity also occurs in Fried's 2026
paper, equation (57), where it is used to study a different Riordan-array
claim; it is not claimed here as new.  Extracting $u^m$ in $P_{2m+1}$
gives
\[
 b_m=\binom{2m+1}{m}+\sum_{j=1}^m b_{j-1}\binom{2(m-j)}{m-j}.
\]
Hence $B(z)=\sum b_mz^m=C(z)/(\sqrt{1-4z}-z)$.
Putting $u=1$ in the row identity gives the row-sum generating function
\[
 F(x)=\frac{1+x^3B(x^2)}{1-2x}.
\]
This is the radical used below; the differential certificate proves the
selected row-sum recurrence, which is distinct from the diagonal and
Riordan-array assertions in the cited earlier work.

\noindent\textbf{Exact certificate.} $F=R_0+R_1\sqrt D$, with\begin{align*}
D&=- \left(2 x - 1\right) \left(2 x + 1\right),\\
R_0&=- \frac{x^{3} - 2 x^{2} + 3 x + 2}{2 \left(x^{4} + 4 x^{2} - 1\right)},\\
R_1&=- \frac{x \left(x - 1\right) \left(x + 1\right)}{2 \left(2 x - 1\right) \left(x^{4} + 4 x^{2} - 1\right)},\\
\end{align*}
Exact substitution gives $LF=P(x)=- 5 x^{6} + 18 x^{5} + 7 x^{4} - 8 x^{3} - 3 x^{2} + 2 x + 1$. The coefficient of the radical in the residual is identically zero.
\noindent\textbf{Verification.} The standalone verifier reproduces the exact identity, compares 31 published terms with an independent definition calculation, and checks all applicable indices through $100$.\par\medskip
\subsection{A026765}
\label{round9:r9entry:A026765}
\noindent\textbf{Source and selected claim.} \href{https://oeis.org/A026765}{A026765}, checked 20 October 2026. Conjecture. The selected polynomial recurrence is
\[\sum_{j=0}^{6}p_j(n)a_{n-j}=0\qquad(n\ge 6).\]
\begin{align*}
p_{0}(n)&=n + 1,\\
p_{1}(n)&=-2,\\
p_{2}(n)&=- 2 \left(5 n - 3\right),\\
p_{3}(n)&=12,\\
p_{4}(n)&=29 n - 71,\\
p_{5}(n)&=-10,\\
p_{6}(n)&=- 20 \left(n - 5\right),\\
\end{align*}
\noindent\textbf{Offset and initial data.} Offset $0$.
\[\left(a_{0},\ldots,a_{7}\right)=
\left(1,2,4,9,18,41,82,188\right).\]
\noindent\textbf{Proof.} 
The parent array A026758 adds $T(n-2,k-1)$ to the Pascal rule exactly
when $n$ is odd and $1\le k\le(n-1)/2$.  At odd row $2m+1$ use the
integer coordinate $j=m-k$.  Two successive ordinary Pascal steps move
$j$ by $-1,0,1$ with weights $1,2,1$; the extra term raises the middle
weight to three on $j\ge0$.  The initial odd row places one at each of
$j=0,-1$.  These statements also hold at the finite row boundaries after
extending the rows by zero, because the extra predecessor there is zero.

Let $U_j(z)$ sum these odd-row entries over $m\ge0$.  Write
\[
 r_+=\frac{1-3z-\sqrt{(1-z)(1-5z)}}{2z},\qquad
 r_-=\frac{1-2z-\sqrt{1-4z}}{2z}.
\]
Both roots have zero constant term.  The unique solution with finite
support at each coefficient is $U_j=A r_+^j$ for $j\ge0$ and
$U_{-1-j}=B r_-^j$ for $j\ge0$.  The two boundary equations are
$A/r_+-B=1/z$ and $B/r_--A=1/z$.  Therefore
\[
 A=\frac{r_+(1+r_-)}{z(1-r_+r_-)},\qquad
 B=\frac{r_-(1+r_+)}{z(1-r_+r_-)}.
\]
Their sum over all endpoints simplifies exactly to
\[
 H(z)=\frac{A}{1-r_+}+\frac{B}{1-r_-}
     =\frac{\sqrt{(1-z)/((1-4z)(1-5z))}-1}{2z}.
\]
The accompanying certificate reduces this equality modulo the two
quadratic radical equations, so it is an algebraic identity, not a fit.
Every positive even row has twice the preceding odd row sum, whence
$F(x)=1+(x+2x^2)H(x^2)$.  This is exactly the conjectured generating
function printed below.  Applying the differential certificate also proves
the conjectural recurrence.  These are one package; no asymptotic claim is
needed for its count.

\noindent\textbf{Exact certificate.} $F=R_0+R_1\sqrt D$, with\begin{align*}
D&=- \left(x - 1\right) \left(x + 1\right) \left(2 x - 1\right) \left(2 x + 1\right) \left(5 x^{2} - 1\right),\\
R_0&=- \frac{1}{2 x},\\
R_1&=\frac{1}{2 x \left(2 x - 1\right) \left(5 x^{2} - 1\right)},\\
\end{align*}
Exact substitution gives $LF=P(x)=5 x^{4} - 8 x^{3} - 6 x^{2} + 2 x + 1$. The coefficient of the radical in the residual is identically zero.
\noindent\textbf{Verification.} The standalone verifier reproduces the exact identity, compares 30 published terms with an independent definition calculation, and checks all applicable indices through $100$.\par\medskip
\subsection{A052736}
\label{round9:r9entry:A052736}
\noindent\textbf{Source and selected claim.} \href{https://oeis.org/A052736}{A052736}, checked 20 October 2026. Conjecture. The printed identity requires the lower-bound correction proved below.
\noindent\textbf{Offset and initial data.} Offset $0$.
\[\left(a_{0},\ldots,a_{7}\right)=
\left(0,0,2,24,384,8160,218880,7116480\right).\]
\noindent\textbf{Proof.} 
The source defines $a_n$ by its exponential generating function
\[
 F(x)=\frac{1-3x-\sqrt{1-6x+x^2}-x(1-x-\sqrt{1-6x+x^2})}{2}.
\]
Let $b_n$ denote A006319.  Its known ordinary generating function, credited
in that entry to Catalani (2003), is
\[
 B(x)=\frac{(1-x)^2-(1-x)\sqrt{1-6x+x^2}}{2x}.
\]
For example it follows from $B=C(x/(1-x)^2)$, the standard Catalan
composition for the royal-path sequence.  Direct subtraction gives the
exact identity $F(x)=x(B(x)-1)$.  Consequently
\[
 a_0=a_1=0,\qquad a_n=n!\,b_{n-1}\quad(n\ge2).
\]
The entry's conjecture omits this lower-bound qualification: at $n=1$,
$a_1=0$ but $1!b_0=1$.  The result is therefore recorded as a corrected
boundary statement, with a counterexample and a full-range proof, rather
than as a proof of the literally unqualified claim.  The existing
unmarked recurrence in A052736 is not counted.

\noindent\textbf{Verification.} The standalone verifier reproduces the exact identity, compares 20 published terms with an independent definition calculation, and checks all applicable indices through $100$.\par\medskip
\subsection{A071715}
\label{round9:r9entry:A071715}
\noindent\textbf{Source and selected claim.} \href{https://oeis.org/A071715}{A071715}, checked 20 October 2026. Conjecture. The selected polynomial recurrence is
\[\sum_{j=0}^{3}p_j(n)a_{n-j}=0\qquad(n\ge 3).\]
\begin{align*}
p_{0}(n)&=\left(n + 7\right) \left(891 n + 6824\right),\\
p_{1}(n)&=891 n^{2} - 38453 n - 250970,\\
p_{2}(n)&=- 4 \left(5040 n^{2} + 4271 n - 8095\right),\\
p_{3}(n)&=4 \left(2 n - 1\right) \left(1170 n + 3847\right),\\
\end{align*}
\noindent\textbf{Offset and initial data.} Offset $0$.
\[\left(a_{0},\ldots,a_{7}\right)=
\left(1,5,22,92,373,1482,5810,22576\right).\]
\noindent\textbf{Proof.} 
The defining generating function is $F=C(x)^4+xC(x)^8$, where
$C=1+xC^2$ and $C(0)=1$.  Put $t=C-1$.  Then
$x=t/(1+t)^2$ and $F=(1+t)^4+t(1+t)^6$.
This parameter change is invertible as a formal power series since
$dx/dt$ is one at zero.  Euler differentiation becomes
$\theta=t(1+t)/(1-t)\,d/dt$.  Applying the displayed conjectural
operator to this rational expression gives the exact polynomial residual
below and proves the recurrence for all indicated indices.  The entry's
2025 binomial formula is a known alternative description, not an additional
claim of this package.

\noindent\textbf{Exact certificate.} $x=\frac{t}{\left(t + 1\right)^{2}}$, $F=t \left(t + 1\right)^{6} + \left(t + 1\right)^{4}$.
Exact substitution gives $LF=P(x)=20068 x + 47768$. The coefficient of the radical in the residual is identically zero.
\noindent\textbf{Verification.} The standalone verifier reproduces the exact identity, compares 26 published terms with an independent definition calculation, and checks all applicable indices through $100$.\par\medskip
\subsection{A098520}
\label{round9:r9entry:A098520}
\noindent\textbf{Source and selected claim.} \href{https://oeis.org/A098520}{A098520}, checked 20 October 2026. Conjecture. The selected polynomial recurrence is
\[\sum_{j=0}^{2}p_j(n)a_{n-j}=0\qquad(n\ge 2).\]
\begin{align*}
p_{0}(n)&=\left(n - 1\right) \left(n + 1\right),\\
p_{1}(n)&=- n \left(2 n - 1\right),\\
p_{2}(n)&=- 15 n \left(n - 1\right),\\
\end{align*}
\noindent\textbf{Offset and initial data.} Offset $0$.
\[\left(a_{0},\ldots,a_{7}\right)=
\left(0,1,2,15,52,285,1206,6027\right).\]
\noindent\textbf{Proof.} 
Expand the defining exponential generating function
$e^x I_1(4x)/2$.  Since
$I_1(4x)/2=\sum_{k\ge0}4^k x^{2k+1}/(k!(k+1)!)$, its coefficients are
\[
 a_n=\sum_{2k+1\le n}\frac{n!\,4^k}{k!(k+1)!(n-2k-1)!}.
\]
Summing first in $n$ gives
\[
 F(x)=\frac{x}{(1-x)^2}
 \sum_{k\ge0}\binom{2k+1}{k}
       \left(\frac{4x^2}{(1-x)^2}\right)^k.
\]
Use $\sum_{k\ge0}\binom{2k+1}{k}z^k
=(1/\sqrt{1-4z}-1)/(2z)$.
After simplification,
$F=((1-x)/\sqrt{1-2x-15x^2}-1)/(8x)$.
The displayed radical differential certificate now proves the selected
recurrence.  This argument fixes $a_0=0$ and the offset without using the
conjectural recurrence to generate data.

\noindent\textbf{Exact certificate.} $F=R_0+R_1\sqrt D$, with\begin{align*}
D&=- \left(3 x + 1\right) \left(5 x - 1\right),\\
R_0&=- \frac{1}{8 x},\\
R_1&=\frac{x - 1}{8 x \left(3 x + 1\right) \left(5 x - 1\right)},\\
\end{align*}
Exact substitution gives $LF=P(x)=0$. The coefficient of the radical in the residual is identically zero.
\noindent\textbf{Verification.} The standalone verifier reproduces the exact identity, compares 27 published terms with an independent definition calculation, and checks all applicable indices through $100$.\par\medskip
\subsection{A123071}
\label{round9:r9entry:A123071}
\noindent\textbf{Source and selected claim.} \href{https://oeis.org/A123071}{A123071}, checked 20 October 2026. Conjecture. The selected polynomial recurrence is
\[\sum_{j=0}^{7}p_j(n)a_{n-j}=0\qquad(n\ge 7).\]
\begin{align*}
p_{0}(n)&=2,\\
p_{1}(n)&=1,\\
p_{2}(n)&=- 2 n,\\
p_{3}(n)&=- n - 10,\\
p_{4}(n)&=- 2 \left(n - 2\right) \left(n + 2\right),\\
p_{5}(n)&=- n^{2} - 2 n + 23,\\
p_{6}(n)&=2 \left(n - 5\right) \left(n^{2} - 7 n + 11\right),\\
p_{7}(n)&=\left(n - 6\right) \left(n - 5\right)^{2},\\
\end{align*}
\noindent\textbf{Offset and initial data.} Offset $0$.
\[\left(a_{0},\ldots,a_{7}\right)=
\left(1,2,4,12,36,120,400,1520\right).\]
\noindent\textbf{Proof.} 
Use the sequence specification $a_n=B(n)B(n+1)$ in A123071, with
$B(2m)=B(2m+1)=u_m$ and
\[
 u_0=1,\quad u_1=2,\quad
 u_m=2u_{m-1}+2(m-1)u_{m-2}\quad(m\ge2).
\]
This is the established symmetric-involution sequence A000898; the source
attributes the bishop interpretation to Robinson (1976), equation (24).
Thus $a_{2m}=u_m^2$ and $a_{2m+1}=u_m u_{m+1}$.
The elementary involution recurrence also follows by choosing whether a
distinguished item is a singleton (two colors) or belongs to a pair
(two colors and $m-1$ choices for its partner).

The proof of the selected order-seven recurrence is a finite polynomial
identity in two independent indeterminates.  For even $n=2m$, $m\ge4$,
set $u_{m-4}=X,u_{m-3}=Y$ and use the displayed second-order recurrence
to express $u_{m-2},u_{m-1},u_m$.  Substituting the eight products into
$\sum_{j=0}^7p_j(2m)a_{2m-j}$ gives the zero polynomial in $m,X,Y$.
For odd $n=2m+1$, $m\ge3$, start with $u_{m-3}=X,u_{m-2}=Y$, extend
through $u_{m+1}$, and the analogous residual is again identically zero.
The machine-readable certificate records every product and the zero
polynomial in each parity.  These two identities cover every $n\ge7$;
they require no induction based on the target recurrence and no division
by a possibly zero coefficient.

\noindent\textbf{Verification.} The standalone verifier reproduces the exact identity, compares 26 published terms with an independent definition calculation, and checks all applicable indices through $100$.\par\medskip
\subsection{A143013}
\label{round9:r9entry:A143013}
\noindent\textbf{Source and selected claim.} \href{https://oeis.org/A143013}{A143013}, checked 20 October 2026. Conjecture. The selected polynomial recurrence is
\[\sum_{j=0}^{3}p_j(n)a_{n-j}=0\qquad(n\ge 3).\]
\begin{align*}
p_{0}(n)&=n + 2,\\
p_{1}(n)&=- 2 n - 1,\\
p_{2}(n)&=- 3 \left(n - 1\right),\\
p_{3}(n)&=- 2 \left(2 n - 5\right),\\
\end{align*}
\noindent\textbf{Offset and initial data.} Offset $0$.
\[\left(a_{0},\ldots,a_{7}\right)=
\left(1,2,3,7,17,43,114,310\right).\]
\noindent\textbf{Proof.} 
In the source's precise word specification the special level step $F$
may be followed only by a down step or the end of the word.  At the
ground level a nonempty admissible Motzkin path is either the singleton
$F$, an ordinary level step followed by a path, or a first-return
$UADB$, with $A,B$ admissible paths.  Consequently its generating function
obeys $F(x)=1+x+xF(x)+x^2F(x)^2$.
The root with constant term one is
$(1-x-\sqrt{1-2x-3x^2-4x^3})/(2x^2)$.
The linked 2023 papers give this known algebraic generating function and
related enumerations; the selected claim here is the particular
polynomial recurrence still marked conjectural.  Its exact differential
residual is given below, completing its full-range proof.

\noindent\textbf{Exact certificate.} $F=R_0+R_1\sqrt D$, with\begin{align*}
D&=- 4 x^{3} - 3 x^{2} - 2 x + 1,\\
R_0&=- \frac{x - 1}{2 x^{2}},\\
R_1&=- \frac{1}{2 x^{2}},\\
\end{align*}
Exact substitution gives $LF=P(x)=- x^{2} + 3 x + 2$. The coefficient of the radical in the residual is identically zero.
\noindent\textbf{Verification.} The standalone verifier reproduces the exact identity, compares 28 published terms with an independent definition calculation, and checks all applicable indices through $100$.\par\medskip
\subsection{A372018}
\label{round9:r9entry:A372018}
\noindent\textbf{Source and selected claim.} \href{https://oeis.org/A372018}{A372018}, checked 20 October 2026. Conjecture. The selected generating-function or coefficient identity is stated and proved below.
\noindent\textbf{Offset and initial data.} Offset $0$.
\[\left(a_{0},\ldots,a_{7}\right)=
\left(1,2,4,10,30,98,336,1194\right).\]
\noindent\textbf{Proof.} 
Write $A(x)=\sum a_nx^n$ for the defining series.  Squaring its defining
square-root equation and using the branch $A(0)=1$ gives
\[
 xA^3-A^2+3xA+1=0.
\]
Separate parity by $A(x)=E(z)+xO(z)$, $z=x^2$.
The odd part of this equation is
$E(E^2+3zO^2-2O+3)=0$.  Since $E(0)=1$, it is a unit, and hence
$E^2=2O-3-3zO^2$.
The even part is
$E^2(3zO-1)+z^2O^3-zO^2+3zO+1=0$.
Eliminating $E^2$ yields
\[
 4z^2O^3-4zO^2+3zO+O-2=0.
\]
For $B=O/2$ this becomes $B(1-4zB)^2=1-3zB$, with $B(0)=1$.
If $W=zB$, then $W$ is the compositional inverse of
$u(1-4u)^2/(1-3u)$, exactly the definition of A371364.
The derivative of the last implicit equation with respect to $B$ is one
at $(z,B)=(0,1)$, so that formal solution is unique.  Therefore
\[
 a_{2n+1}=2\,\mathrm{A371364}(n)\qquad(n\ge0).
\]
The source prints empty parentheses after A371364; the intended index
$n$ is made explicit here.  Neither a separate unmarked recurrence in
A372018 nor the partner sequence is counted as another package.

\noindent\textbf{Verification.} The standalone verifier reproduces the exact identity, compares 27 published terms with an independent definition calculation, and checks all applicable indices through $100$.\par\medskip
\subsection{A212341}
\label{round9:r9entry:A212341}
\noindent\textbf{Source and selected claim.} \href{https://oeis.org/A212341}{A212341}, checked 20 October 2026. Empirical generating function. The selected generating-function or coefficient identity is stated and proved below.
\noindent\textbf{Offset and initial data.} Offset $5$.
\[\left(a_{5},\ldots,a_{12}\right)=
\left(14,56,188,603,1907,5615,15968,44464\right).\]
\noindent\textbf{Proof.} 
Let $Q(x,u)$ count 132-avoiding permutations, with $u$ marking entries
having at least four smaller entries to their left.  This is the source's
quadrant mesh statistic.  Split a permutation at its maximum.  All values
on the left exceed those on the right, so the smaller-to-left statistic in
either subpermutation is unchanged.  The maximum is marked precisely
when its left subtree has at least four entries.  Subtrees of sizes zero
through three have no marked entries and have Catalan counts
$P(x)=1+x+2x^2+5x^3$.  Therefore
\[
 Q=1+xP Q+xu(Q-P)Q.
\]
This is the bivariate equation of Kitaev--Remmel--Tiefenbruck (2012),
Theorem 9; it is credited as a known starting point.  Put
$Q_0=Q(x,0)=1/(1-xP)$ and $Q_1=[u]Q$.  Taking the coefficient of $u$
in the equation gives $(1-xP)Q_1=x(Q_0-P)Q_0$, and thus
\[
 Q_1=\frac{x(1-P+xP^2)}{(1-xP)^3}
 =\frac{x^5(14+14x+20x^2+25x^3)}{(1-x-x^2-2x^3-5x^4)^3}.
\]
This is the entry's empirical generating function, now proved for every
coefficient.  Section 4 of the cited paper gives the bivariate formula
and initial terms and explicitly leaves this coefficient extraction
without a simple formula; the present one-line extraction is the
selected assertion.  The initial zero coefficients through degree four
explain the source offset five.

\noindent\textbf{Verification.} The standalone verifier reproduces the exact identity, compares 26 published terms with an independent definition calculation, and checks all applicable indices through $100$.\par\medskip
\subsection{A110668}
\label{round9:r9entry:A110668}
\noindent\textbf{Source and selected claim.} \href{https://oeis.org/A110668}{A110668}, checked 20 October 2026. Conjecture. The selected generating-function or coefficient identity is stated and proved below.
\noindent\textbf{Offset and initial data.} Offset $0$.
\[\left(a_{0},\ldots,a_{7}\right)=
\left(0,1,3,3,-3,-15,-27,-32\right).\]
\noindent\textbf{Proof.} 
Let $u_n=a(0,n)$ denote the base row in A110665 and let
$U_m(x)=\sum_{n\ge0}a(m,n)x^n$.  The defining array recurrence and the
zero first column give $(1-x)U_m=U_{m-1}$, hence $U_3=U_0/(1-x)^3$.
The base-row generating function
$U_0=x(1-2x)/(1-x+x^2)^2$ is already credited to
Adams-Watters (2006) in A110665.  For completeness it can be recovered
from the diagonal condition alone: if $C=1+xC^2$ and $z=xC$, then
\[
 \sum_{n\ge0}a(n,n)x^n
 =\frac{U_0(z)}{C(x)\sqrt{1-4x}},\qquad x=z(1-z).
\]
Indeed the coefficient of $u_k$ on the diagonal is
$\binom{2n-k-1}{n-1}$ for $n\ge k\ge1$, whose generating function
is $(xC)^k/(C(x)\sqrt{1-4x})$.
Since the diagonal equals $n$, the left side is $x/(1-x)^2$.
Substituting $x=z(1-z)$, $C(x)=1/(1-z)$, and
$\sqrt{1-4x}=1-2z$ gives the stated
$U_0$.  Thus
\[
 U_3(x)=\frac{x(1-2x)}{(1-x+x^2)^2(1-x)^3}
       =\frac{x(-1+2x)}{(x^2-x+1)^2(x-1)^3},
\]
which proves the selected conjectural formula.  The base row and the
other rows are credited or uncounted; this package selects only A110668.

\noindent\textbf{Verification.} The standalone verifier reproduces the exact identity, compares 54 published terms with an independent definition calculation, and checks all applicable indices through $100$.\par\medskip

\section{Empirical array recurrences: 62 selected packages}
\label{round9:sec:r9arrays}
This section proves 62 entry-level recurrences explicitly marked
\emph{Empirical} on primary OEIS pages retrieved on 20 October 2026.
The local array definitions and recurrence conjectures are due to
R.~H. Hardin. In 20 of the same packages, a source rational
generating function, credited to Colin Barker, is also proved.
One selected A-number contributes one package. No additional count is
assigned to its generating function or to a common theorem.

\subsection{Four local predicates and labeled counting}
For a block $B=\left(\begin{smallmatrix}a&b\\c&d\end{smallmatrix}\right)$,
let
\[
 \mathcal D(B)=(b-a,d-b,c-d,a-c,d-a,c-b).
\]
These are all six pairwise differences among its four positions,
with signs chosen immaterially for squares or absolute values.
Each selected sequence counts arrays with $n+1$ rows, width $w\ge2$,
and entries in $\{0,\ldots,q-1\}$, subject to one of the following
requirements in every $2\times2$ block:
\begin{description}
\item[Stress (D).] $|a+d-b-c|=s$.
\item[Absolute six-difference sum (U).]
 $\sum_{u\in\mathcal D(B)}|u|=s$.
\item[Squared six-difference sum (V).]
 $\sum_{u\in\mathcal D(B)}u^2=s$.
\item[Bounded six-difference range (B).]
 $|u|\le s$ for every $u\in\mathcal D(B)$, equivalently
 $\max(a,b,c,d)-\min(a,b,c,d)\le s$.
\end{description}
An additional inequality condition is imposed only where the source
requires every horizontally or vertically adjacent pair to be unequal.
In a block this says $a\ne b$, $b\ne d$, $d\ne c$, and $c\ne a$.
Every edge belongs to a block when both array dimensions are at least
two, so this also treats boundary edges correctly.

All counts retain the given numeric labels. Symbols may be absent;
there is no requirement to use every symbol and no quotient by color
permutations or geometric symmetries. The row orbits used below are
only a compression of the exact labeled count.

\subsection{A transfer obtained directly from the definition}
For a row $r\in\{0,\ldots,q-1\}^w$, let $\gamma(r)$ be the
lexicographically least among $r$, its reversal, its simultaneous
complement $q-1-r$, and the reversal of that complement. Let
$\mathcal O_{q,w}$ be the sorted set of representatives. Define
\begin{align}
 m(r)&=\#\{u\in\{0,\ldots,q-1\}^w:\gamma(u)=r\},
 \label{round9:eq:r9arr-mass}\\
 T(r,t)&=\#\left\{u\in\{0,\ldots,q-1\}^w:
 \begin{array}{l}\gamma(u)=t,\\
 P(r_i,r_{i+1},u_i,u_{i+1})\quad(1\le i<w)
 \end{array}\right\}.
 \label{round9:eq:r9arr-T}
\end{align}
Here $P$ is the complete local predicate specified above, including the
edge inequalities when required. Reversal and complementation preserve
each predicate and permute every target orbit when applied to both
rows. Consequently every row in a source orbit has the same number of
successors in each target orbit. The weighted quotient is therefore
equitable, and adding rows gives
\begin{equation}\label{round9:eq:r9arr-count}
 a(n)=m^{\mathsf T}T^n\mathbf1\qquad(n\ge1).
\end{equation}
The exponent $n$ counts the joins between $n+1$ rows. The mass sum is
$\sum_r m(r)=q^w$; short orbits, fixed rows, and zero-successor rows
retain their actual multiplicities. The auxiliary value at $n=0$ is
not asserted to be an OEIS term; all selected sources have offset 1.

\subsection{An exact finite identity proves all later indices}
For the source signature $c=(c_1,\ldots,c_d)$, define
\[
 R(z)=z^d-\sum_{i=1}^{d}c_i z^{d-i},\qquad
 Q(x)=1-\sum_{i=1}^{d}c_i x^i.
\]
The full integer certificate for every selected entry verifies
\begin{equation}\label{round9:eq:r9arr-cert}
 \boxed{TR(T)\mathbf1=\mathbf0.}
\end{equation}
Multiplication by $m^{\mathsf T}T^{n-d-1}$ and
\eqref{round9:eq:r9arr-count} prove
\begin{equation}\label{round9:eq:r9arr-rec}
 \boxed{a(n)=\sum_{i=1}^{d}c_i a(n-i)\qquad(n\ge d+1).}
\end{equation}
Thus every index on the right belongs to the original offset-one
sequence. No recurrence is fitted to a numerical prefix, and no
conjectural generating function is assumed.

The 62 complete matrices have 10537 states and 30926 nonzero
weighted edges in total. The largest dimension is 1056; the largest
source recurrence order is 56. For every matrix the data retain
the representative rows, all orbit masses, every weighted transition,
and every coordinate of $v_k=T^k\mathbf1$ for $0\le k\le d+1$.
The decisive calculation is
\begin{Verbatim}[fontsize=\small]
v[0] = [1 for every state]
for k = 1,...,d+1:
    v[k][r] = sum(T[r,s]*v[k-1][s] for every state s)
for every state r:
    assert v[d+1][r] == sum(c[i]*v[d+1-i][r] for i=1,...,d)
\end{Verbatim}
The verifier rebuilds the matrix from the definition and recomputes all
coordinates with exact integer arithmetic. This is a finite
computer-assisted proof of an all-index assertion. The full parameter
and coefficient lists below specify the calculation without relying
on opaque numerical output. Expanded data are in
\texttt{r9\_arrays\_matrices.json} and
\texttt{r9\_arrays\_certificates.json}.

All 1112 displayed source terms match the defining transfers. A
separate program, which does not import this engine, uses labeled
column-by-column dynamic programming to obtain two-row and three-row
counts. All 124 values agree. These checks confirm interpretation
and offset; the vector identity supplies the universal proof.

For $F(x)=\sum_{n\ge1}a(n)x^n$, the recurrence gives $QF=P$ with
$\deg P\le d$, where
\begin{equation}\label{round9:eq:r9arr-numerator}
 [x^k]P=\sum_{\substack{0\le i\le\min(k,d)\\k-i\ge1}}
 q_i a(k-i),\quad q_0=1,\quad q_i=-c_i\ (i\ge1).
\end{equation}
Exact polynomial cross-multiplication proves agreement with all 20
source rational generating functions. Their numerators are printed in
the corresponding packages; other source formulas are not counted.
For every order-$d$ recurrence, the certificate file also retains all
initial values $a(1),\ldots,a(d)$, computed from the defining transfer;
the entry paragraphs print the first three values for orientation.

\subsection{A hand-checkable example: A234154}
For $q=4$, width $w=2$, and predicate B with $s=1$, the orbit
representatives are $00,01,02,03,11,12$. Direct enumeration gives
\[
 T=\begin{pmatrix}
 1&2&0&0&1&0\\1&2&0&0&1&0\\0&0&0&0&0&0\\
 0&0&0&0&0&0\\1&2&0&0&2&2\\0&0&0&0&2&2
 \end{pmatrix},\qquad m=(2,4,4,2,2,2)^{\mathsf T}.
\]
The successive column vectors are
\begin{align*}
 v_1&=(4,4,0,0,7,4)^{\mathsf T},&
 v_2&=(19,19,0,0,34,22)^{\mathsf T},\\
 v_3&=(91,91,0,0,169,112)^{\mathsf T},&
 v_4&=(442,442,0,0,835,562)^{\mathsf T}.
\end{align*}
Every coordinate satisfies $v_4=7v_3-9v_2-6v_1$. It follows that
\[
 a(n)=7a(n-1)-9a(n-2)-6a(n-3)\qquad(n\ge4).
\]
The first values $46,226,1108$ come from $m^{\mathsf T}v_i$, so
\[
 F(x)=\frac{46x-96x^2-60x^3}{1-7x+9x^2+6x^3}.
\]
This proves both empirical formulas on
\url{https://oeis.org/A234154}.

\subsection{Source status, prior work and audit limits}
All selected primary entries were retrieved for this round and
explicitly mark the selected recurrence empirical. No proof or
established encompassing recurrence was identified in their consulted
content. The 24 associated parent tables for the selected and reserve
packages were sought separately: 21 were retrieved and their relevant
statements are also empirical.
The contents of A233883, A234114, and A234175 could not be retrieved
after ordinary, internal-format and public text-endpoint attempts.
Their proof status is not claimed to have been audited. The affected
child entries remain directly sourced, and their mathematical
certificates are independent of those inaccessible parent pages.

Retrieval dates and reported cache ages are retained separately.
Some retrieved copies predate this round; no claim is made that they
reflect every live edit. Exact-ID and family searches did not identify
a prior proof among the inspected results, but several searches had
poor recall or irrelevant results. This bounded audit does not
establish worldwide priority.

The selected IDs are absent from all 808 earlier selected IDs and from
the retained 97 distinct IDs in the 100 costly claims of Appendix A of
\emph{OEIS Open}. The latter comparison is not a full audit of that
project's repository. Previously proved array reserves were also
excluded from new selection.

The alias screen compares each new sequence with 289 retained array
packages from Rounds II--VIII, including four earlier uncounted
reserves, and with earlier entries in this new set. It compares 20
exact values at each shift $-3,\ldots,3$, allowing an arbitrary
rational scale. Every proposed equality has a mismatch within the
first 3 compared values. This
disproves those specific shift-and-scale identities; unbounded shifts
and other transformations are not covered. Shared recurrence
polynomials alone do not identify two different counting sequences.

One further fresh certificate is retained as an uncounted reserve.
None was proved in an earlier round. No OEIS submission,
edit, or message was performed, and no official closure is claimed.

\subsection{Entry-level coefficients and certificates}
In each package below, the displayed coefficient signature is ordered
as $(c_1,\ldots,c_d)$ in \eqref{round9:eq:r9arr-rec}. Every zero coefficient is
retained. All arrays have offset 1.
\Needspace{7\baselineskip}
\subsubsection{A233787}
\label{round9:r9entry:A233787}
Arrays have $n+1$ rows, width $w=4$, and labeled alphabet $0,\ldots,3$. Predicate V has target $s=11$, allowing equal adjacent values.
The recurrence holds for every $n\ge 35$, with order $d=34$ and a complete $D=72$ state matrix. The first three source values are $(720,3932,20816)$.
\begin{Verbatim}[fontsize=\scriptsize,breaklines=true]
0, 139, 0, -7929, 0, 248319, 0, -4798442, 0, 60800038, 0, -523147257, 0, 3118918441,
0, -13016374051, 0, 38139962547, 0, -78218266028, 0, 111158046394, 0, -107387771932,
0, 68279402362, 0, -27119287096, 0, 6185757504, 0, -702040320, 0, 29030400
\end{Verbatim}
The defining transfer satisfies $TR(T)\mathbf1=0$ in every coordinate, so \eqref{round9:eq:r9arr-rec} follows. Source: \url{https://oeis.org/A233787}; parent column 3 of \url{https://oeis.org/A233792}. Both retrieved entries mark the relevant recurrence empirical.
\Needspace{7\baselineskip}
\subsubsection{A233877}
\label{round9:r9entry:A233877}
Arrays have $n+1$ rows, width $w=3$, and labeled alphabet $0,\ldots,2$. Predicate V has target $s=11$, allowing equal adjacent values.
The recurrence holds for every $n\ge 7$, with order $d=6$ and a complete $D=10$ state matrix. The first three source values are $(76,300,1224)$.
\[(c_1,\ldots,c_d)=(3,19,-60,8,36,-12).\]
The defining transfer satisfies $TR(T)\mathbf1=0$ in every coordinate, so \eqref{round9:eq:r9arr-rec} follows. Source: \url{https://oeis.org/A233877}; parent column 2 of \url{https://oeis.org/A233883}. The primary entry marks the recurrence empirical; the parent content was not retrieved, as disclosed in the source audit.
The source rational generating function is also proved: $F=P/Q$, with $Q=1-\sum_{i=1}^{d}c_i x^i$ and
\[P(x)=76x+72x^{2}-1120x^{3}+344x^{4}+688x^{5}-256x^{6}.\]
\Needspace{7\baselineskip}
\subsubsection{A233878}
\label{round9:r9entry:A233878}
Arrays have $n+1$ rows, width $w=4$, and labeled alphabet $0,\ldots,2$. Predicate V has target $s=11$, allowing equal adjacent values.
The recurrence holds for every $n\ge 13$, with order $d=12$ and a complete $D=25$ state matrix. The first three source values are $(240,1224,6200)$.
\begin{Verbatim}[fontsize=\scriptsize,breaklines=true]
0, 56, 0, -819, 0, 2791, 0, -3096, 0, 961, 0, -84
\end{Verbatim}
The defining transfer satisfies $TR(T)\mathbf1=0$ in every coordinate, so \eqref{round9:eq:r9arr-rec} follows. Source: \url{https://oeis.org/A233878}; parent column 3 of \url{https://oeis.org/A233883}. The primary entry marks the recurrence empirical; the parent content was not retrieved, as disclosed in the source audit.
\Needspace{7\baselineskip}
\subsubsection{A233879}
\label{round9:r9entry:A233879}
Arrays have $n+1$ rows, width $w=5$, and labeled alphabet $0,\ldots,2$. Predicate V has target $s=11$, allowing equal adjacent values.
The recurrence holds for every $n\ge 30$, with order $d=29$ and a complete $D=70$ state matrix. The first three source values are $(760,5156,33656)$.
\begin{Verbatim}[fontsize=\scriptsize,breaklines=true]
10, 111, -1347, -2952, 61962, -23589, -1303743, 2061738, 13676265, -32627870,
-74205472, 248050297, 184715702, -1058894761, 18879451, 2642739884, -1291495298,
-3771977707, 3244720225, 2774675096, -3693043205, -677649206, 2114258648, -285661000,
-570117698, 190279800, 48545460, -28512000, 2851200
\end{Verbatim}
The defining transfer satisfies $TR(T)\mathbf1=0$ in every coordinate, so \eqref{round9:eq:r9arr-rec} follows. Source: \url{https://oeis.org/A233879}; parent column 4 of \url{https://oeis.org/A233883}. The primary entry marks the recurrence empirical; the parent content was not retrieved, as disclosed in the source audit.
\Needspace{7\baselineskip}
\subsubsection{A233880}
\label{round9:r9entry:A233880}
Arrays have $n+1$ rows, width $w=6$, and labeled alphabet $0,\ldots,2$. Predicate V has target $s=11$, allowing equal adjacent values.
The recurrence holds for every $n\ge 53$, with order $d=52$ and a complete $D=196$ state matrix. The first three source values are $(2400,22020,178704)$.
\begin{Verbatim}[fontsize=\scriptsize,breaklines=true]
0, 390, 0, -63115, 0, 5633508, 0, -311287762, 0, 11395143071, 0, -290608508897, 0,
5364876782159, 0, -73770016903887, 0, 771455610322102, 0, -6226548780129951, 0,
39172372910038001, 0, -193237451888156699, 0, 749448322438960607, 0,
-2284924382785662062, 0, 5461859355511846682, 0, -10185644885765121071, 0,
14709447103939892100, 0, -16281042118217455735, 0, 13617119182197701100, 0,
-8438611741506693900, 0, 3769223995920232000, 0, -1166605374110880000, 0,
236321446856400000, 0, -28818965889000000, 0, 1865383020000000, 0, -48478500000000
\end{Verbatim}
The defining transfer satisfies $TR(T)\mathbf1=0$ in every coordinate, so \eqref{round9:eq:r9arr-rec} follows. Source: \url{https://oeis.org/A233880}; parent column 5 of \url{https://oeis.org/A233883}. The primary entry marks the recurrence empirical; the parent content was not retrieved, as disclosed in the source audit.
\Needspace{7\baselineskip}
\subsubsection{A233960}
\label{round9:r9entry:A233960}
Arrays have $n+1$ rows, width $w=2$, and labeled alphabet $0,\ldots,3$. Predicate U has target $s=10$, allowing equal adjacent values.
The recurrence holds for every $n\ge 6$, with order $d=5$ and a complete $D=6$ state matrix. The first three source values are $(48,188,728)$.
\[(c_1,\ldots,c_d)=(4,10,-36,-8,16).\]
The defining transfer satisfies $TR(T)\mathbf1=0$ in every coordinate, so \eqref{round9:eq:r9arr-rec} follows. Source: \url{https://oeis.org/A233960}; parent column 1 of \url{https://oeis.org/A233967}. Both retrieved entries mark the relevant recurrence empirical.
The source rational generating function is also proved: $F=P/Q$, with $Q=1-\sum_{i=1}^{d}c_i x^i$ and
\[P(x)=48x-4x^{2}-504x^{3}-48x^{4}+224x^{5}.\]
\Needspace{7\baselineskip}
\subsubsection{A233961}
\label{round9:r9entry:A233961}
Arrays have $n+1$ rows, width $w=3$, and labeled alphabet $0,\ldots,3$. Predicate U has target $s=10$, allowing equal adjacent values.
The recurrence holds for every $n\ge 15$, with order $d=14$ and a complete $D=20$ state matrix. The first three source values are $(188,864,3976)$.
\begin{Verbatim}[fontsize=\scriptsize,breaklines=true]
9, 6, -231, 402, 1164, -3026, -973, 5685, -1687, -2838, 1862, -188, -88, 16
\end{Verbatim}
The defining transfer satisfies $TR(T)\mathbf1=0$ in every coordinate, so \eqref{round9:eq:r9arr-rec} follows. Source: \url{https://oeis.org/A233961}; parent column 2 of \url{https://oeis.org/A233967}. Both retrieved entries mark the relevant recurrence empirical.
\Needspace{7\baselineskip}
\subsubsection{A233962}
\label{round9:r9entry:A233962}
Arrays have $n+1$ rows, width $w=4$, and labeled alphabet $0,\ldots,3$. Predicate U has target $s=10$, allowing equal adjacent values.
The recurrence holds for every $n\ge 44$, with order $d=43$ and a complete $D=72$ state matrix. The first three source values are $(728,3976,21272)$.
\begin{Verbatim}[fontsize=\scriptsize,breaklines=true]
15, 31, -1396, 2653, 51714, -174135, -1048642, 4568556, 13206552, -69675134,
-110040243, 692530682, 626223287, -4730984068, -2448897493, 22858856946, 6357801357,
-79394291518, -9325291758, 199923021296, -23907331, -365908349411, 34128121167,
485101885188, -82901695912, -461141546148, 108852864503, 308761092196, -88956808820,
-141933147555, 45953728822, 43469138372, -14674989574, -8649583694, 2801908488,
1098165952, -309043744, -85338496, 18684416, 3638272, -557056, -64000, 6144
\end{Verbatim}
The defining transfer satisfies $TR(T)\mathbf1=0$ in every coordinate, so \eqref{round9:eq:r9arr-rec} follows. Source: \url{https://oeis.org/A233962}; parent column 3 of \url{https://oeis.org/A233967}. Both retrieved entries mark the relevant recurrence empirical.
\Needspace{7\baselineskip}
\subsubsection{A234029}
\label{round9:r9entry:A234029}
Arrays have $n+1$ rows, width $w=2$, and labeled alphabet $0,\ldots,4$. Predicate B has target $s=1$, allowing equal adjacent values.
The recurrence holds for every $n\ge 5$, with order $d=4$ and a complete $D=9$ state matrix. The first three source values are $(61,307,1543)$.
\[(c_1,\ldots,c_d)=(7,-6,-18,-6).\]
The defining transfer satisfies $TR(T)\mathbf1=0$ in every coordinate, so \eqref{round9:eq:r9arr-rec} follows. Source: \url{https://oeis.org/A234029}; parent column 1 of \url{https://oeis.org/A234036}. Both retrieved entries mark the relevant recurrence empirical.
The source rational generating function is also proved: $F=P/Q$, with $Q=1-\sum_{i=1}^{d}c_i x^i$ and
\[P(x)=61x-120x^{2}-240x^{3}-78x^{4}.\]
\Needspace{7\baselineskip}
\subsubsection{A234030}
\label{round9:r9entry:A234030}
Arrays have $n+1$ rows, width $w=3$, and labeled alphabet $0,\ldots,4$. Predicate B has target $s=1$, allowing equal adjacent values.
The recurrence holds for every $n\ge 10$, with order $d=9$ and a complete $D=39$ state matrix. The first three source values are $(307,3059,30263)$.
\begin{Verbatim}[fontsize=\scriptsize,breaklines=true]
18, -79, -56, 459, 56, -597, -88, 90, -10
\end{Verbatim}
The defining transfer satisfies $TR(T)\mathbf1=0$ in every coordinate, so \eqref{round9:eq:r9arr-rec} follows. Source: \url{https://oeis.org/A234030}; parent column 2 of \url{https://oeis.org/A234036}. Both retrieved entries mark the relevant recurrence empirical.
\Needspace{7\baselineskip}
\subsubsection{A234031}
\label{round9:r9entry:A234031}
Arrays have $n+1$ rows, width $w=4$, and labeled alphabet $0,\ldots,4$. Predicate B has target $s=1$, allowing equal adjacent values.
The recurrence holds for every $n\ge 17$, with order $d=16$ and a complete $D=169$ state matrix. The first three source values are $(1543,30263,591711)$.
\begin{Verbatim}[fontsize=\scriptsize,breaklines=true]
41, -473, 378, 15090, -24456, -191256, 160111, 1124270, 180876, -2156866, -1359614,
738304, 501820, -106672, -46728, 7920
\end{Verbatim}
The defining transfer satisfies $TR(T)\mathbf1=0$ in every coordinate, so \eqref{round9:eq:r9arr-rec} follows. Source: \url{https://oeis.org/A234031}; parent column 3 of \url{https://oeis.org/A234036}. Both retrieved entries mark the relevant recurrence empirical.
\Needspace{7\baselineskip}
\subsubsection{A234032}
\label{round9:r9entry:A234032}
Arrays have $n+1$ rows, width $w=5$, and labeled alphabet $0,\ldots,4$. Predicate B has target $s=1$, allowing equal adjacent values.
The recurrence holds for every $n\ge 45$, with order $d=44$ and a complete $D=819$ state matrix. The first three source values are $(7783,300597,11637549)$.
\begin{Verbatim}[fontsize=\scriptsize,breaklines=true]
110, -4333, 67258, -49727, -9381974, 63529143, 464873722, -5105834204, -11512298940,
199030863273, 171749489662, -4640901058548, -2279932923269, 69427761078243,
37219734180920, -679717373578792, -517634627173445, 4315879589285745,
4575357081875100, -17189611279421840, -24043742578423549, 40454429041355663,
73834278266333876, -51489559946124003, -132873735034278651, 30290331332947967,
146882875991746168, -208125301638949, -104800983506559948, -9888619515464597,
50085499296310162, 5296975913473720, -16275620572050062, -948007832112512,
3507290470767932, -86886144072040, -454894672567616, 54652514222304, 27282841319424,
-5902990277920, -173722766592, 96048316160, -753748992, -473149440
\end{Verbatim}
The defining transfer satisfies $TR(T)\mathbf1=0$ in every coordinate, so \eqref{round9:eq:r9arr-rec} follows. Source: \url{https://oeis.org/A234032}; parent column 4 of \url{https://oeis.org/A234036}. Both retrieved entries mark the relevant recurrence empirical.
\Needspace{7\baselineskip}
\subsubsection{A234107}
\label{round9:r9entry:A234107}
Arrays have $n+1$ rows, width $w=2$, and labeled alphabet $0,\ldots,4$. Predicate U has target $s=12$, allowing equal adjacent values.
The recurrence holds for every $n\ge 7$, with order $d=6$ and a complete $D=9$ state matrix. The first three source values are $(56,186,604)$.
\[(c_1,\ldots,c_d)=(5,6,-50,16,80,-16).\]
The defining transfer satisfies $TR(T)\mathbf1=0$ in every coordinate, so \eqref{round9:eq:r9arr-rec} follows. Source: \url{https://oeis.org/A234107}; parent column 1 of \url{https://oeis.org/A234114}. The primary entry marks the recurrence empirical; the parent content was not retrieved, as disclosed in the source audit.
The source rational generating function is also proved: $F=P/Q$, with $Q=1-\sum_{i=1}^{d}c_i x^i$ and
\[P(x)=56x-94x^{2}-662x^{3}+728x^{4}+1424x^{5}-304x^{6}.\]
\Needspace{7\baselineskip}
\subsubsection{A234108}
\label{round9:r9entry:A234108}
Arrays have $n+1$ rows, width $w=3$, and labeled alphabet $0,\ldots,4$. Predicate U has target $s=12$, allowing equal adjacent values.
The recurrence holds for every $n\ge 19$, with order $d=18$ and a complete $D=39$ state matrix. The first three source values are $(186,608,2086)$.
\begin{Verbatim}[fontsize=\scriptsize,breaklines=true]
8, 6, -172, 144, 1438, -1805, -6160, 8467, 14464, -20123, -17822, 25464, 9216, -15948,
136, 3552, -928, 64
\end{Verbatim}
The defining transfer satisfies $TR(T)\mathbf1=0$ in every coordinate, so \eqref{round9:eq:r9arr-rec} follows. Source: \url{https://oeis.org/A234108}; parent column 2 of \url{https://oeis.org/A234114}. The primary entry marks the recurrence empirical; the parent content was not retrieved, as disclosed in the source audit.
\Needspace{7\baselineskip}
\subsubsection{A234109}
\label{round9:r9entry:A234109}
Arrays have $n+1$ rows, width $w=4$, and labeled alphabet $0,\ldots,4$. Predicate U has target $s=12$, allowing equal adjacent values.
The recurrence holds for every $n\ge 41$, with order $d=40$ and a complete $D=169$ state matrix. The first three source values are $(604,2086,7732)$.
\begin{Verbatim}[fontsize=\scriptsize,breaklines=true]
7, 60, -438, -1782, 12360, 35330, -205978, -512290, 2229224, 5503649, -16243635,
-43372518, 79850408, 248228062, -254990128, -1024802383, 454865089, 3032850159,
-63270053, -6374308278, -1896654228, 9366697820, 5199067584, -9374661980, -7431981604,
6110120008, 6536677712, -2372079792, -3650095360, 416839296, 1285380768, 31536704,
-277159232, -25366784, 34816896, 3406336, -2369024, -107520, 69632, -4096
\end{Verbatim}
The defining transfer satisfies $TR(T)\mathbf1=0$ in every coordinate, so \eqref{round9:eq:r9arr-rec} follows. Source: \url{https://oeis.org/A234109}; parent column 3 of \url{https://oeis.org/A234114}. The primary entry marks the recurrence empirical; the parent content was not retrieved, as disclosed in the source audit.
\Needspace{7\baselineskip}
\subsubsection{A234117}
\label{round9:r9entry:A234117}
Arrays have $n+1$ rows, width $w=4$, and labeled alphabet $0,\ldots,2$. Predicate B has target $s=1$, allowing equal adjacent values.
The recurrence holds for every $n\ge 6$, with order $d=5$ and a complete $D=25$ state matrix. The first three source values are $(673,12593,231713)$.
\[(c_1,\ldots,c_d)=(20,-10,-324,-277,144).\]
The defining transfer satisfies $TR(T)\mathbf1=0$ in every coordinate, so \eqref{round9:eq:r9arr-rec} follows. Source: \url{https://oeis.org/A234117}; parent column 3 of \url{https://oeis.org/A234122}. Both retrieved entries mark the relevant recurrence empirical.
The source rational generating function is also proved: $F=P/Q$, with $Q=1-\sum_{i=1}^{d}c_i x^i$ and
\[P(x)=673x-867x^{2}-13417x^{3}-11213x^{4}+5904x^{5}.\]
\Needspace{7\baselineskip}
\subsubsection{A234118}
\label{round9:r9entry:A234118}
Arrays have $n+1$ rows, width $w=5$, and labeled alphabet $0,\ldots,2$. Predicate B has target $s=1$, allowing equal adjacent values.
The recurrence holds for every $n\ge 12$, with order $d=11$ and a complete $D=70$ state matrix. The first three source values are $(3127,116801,4279065)$.
\begin{Verbatim}[fontsize=\scriptsize,breaklines=true]
50, -427, -2748, 22158, 50476, -220478, -236864, 237555, 152850, -74775, -7956
\end{Verbatim}
The defining transfer satisfies $TR(T)\mathbf1=0$ in every coordinate, so \eqref{round9:eq:r9arr-rec} follows. Source: \url{https://oeis.org/A234118}; parent column 4 of \url{https://oeis.org/A234122}. Both retrieved entries mark the relevant recurrence empirical.
\Needspace{7\baselineskip}
\subsubsection{A234119}
\label{round9:r9entry:A234119}
Arrays have $n+1$ rows, width $w=6$, and labeled alphabet $0,\ldots,2$. Predicate B has target $s=1$, allowing equal adjacent values.
The recurrence holds for every $n\ge 18$, with order $d=17$ and a complete $D=196$ state matrix. The first three source values are $(14527,1082977,79003521)$.
\begin{Verbatim}[fontsize=\scriptsize,breaklines=true]
104, -1893, -33284, 535522, 4829920, -36672922, -308181352, 266551299, 4516883912,
2174714575, -20675388468, -11598332564, 40919536464, 9741869952, -32943946752,
6233518080, 2468413440
\end{Verbatim}
The defining transfer satisfies $TR(T)\mathbf1=0$ in every coordinate, so \eqref{round9:eq:r9arr-rec} follows. Source: \url{https://oeis.org/A234119}; parent column 5 of \url{https://oeis.org/A234122}. Both retrieved entries mark the relevant recurrence empirical.
\Needspace{7\baselineskip}
\subsubsection{A234124}
\label{round9:r9entry:A234124}
Arrays have $n+1$ rows, width $w=2$, and labeled alphabet $0,\ldots,6$. Predicate B has target $s=1$, allowing equal adjacent values.
The recurrence holds for every $n\ge 7$, with order $d=6$ and a complete $D=16$ state matrix. The first three source values are $(91,469,2413)$.
\[(c_1,\ldots,c_d)=(10,-24,-18,57,54,12).\]
The defining transfer satisfies $TR(T)\mathbf1=0$ in every coordinate, so \eqref{round9:eq:r9arr-rec} follows. Source: \url{https://oeis.org/A234124}; parent column 1 of \url{https://oeis.org/A234131}. Both retrieved entries mark the relevant recurrence empirical.
The source rational generating function is also proved: $F=P/Q$, with $Q=1-\sum_{i=1}^{d}c_i x^i$ and
\[P(x)=91x-441x^{2}-93x^{3}+1221x^{4}+1038x^{5}+228x^{6}.\]
\Needspace{7\baselineskip}
\subsubsection{A234125}
\label{round9:r9entry:A234125}
Arrays have $n+1$ rows, width $w=3$, and labeled alphabet $0,\ldots,6$. Predicate B has target $s=1$, allowing equal adjacent values.
The recurrence holds for every $n\ge 15$, with order $d=14$ and a complete $D=100$ state matrix. The first three source values are $(469,4757,47933)$.
\begin{Verbatim}[fontsize=\scriptsize,breaklines=true]
26, -218, 474, 1603, -5034, -4660, 14440, 7897, -13910, -6694, 2654, 581, -250, 20
\end{Verbatim}
The defining transfer satisfies $TR(T)\mathbf1=0$ in every coordinate, so \eqref{round9:eq:r9arr-rec} follows. Source: \url{https://oeis.org/A234125}; parent column 2 of \url{https://oeis.org/A234131}. Both retrieved entries mark the relevant recurrence empirical.
\Needspace{7\baselineskip}
\subsubsection{A234126}
\label{round9:r9entry:A234126}
Arrays have $n+1$ rows, width $w=4$, and labeled alphabet $0,\ldots,6$. Predicate B has target $s=1$, allowing equal adjacent values.
The recurrence holds for every $n\ge 28$, with order $d=27$ and a complete $D=625$ state matrix. The first three source values are $(2413,47933,951709)$.
\begin{Verbatim}[fontsize=\scriptsize,breaklines=true]
62, -1377, 11706, 3232, -539522, 1362867, 10121999, -31760118, -116725179, 307724737,
923620665, -1359253203, -4570974318, 1827227267, 11806964282, 3166332198,
-13354128233, -8963914586, 4581307330, 4739020728, -551989326, -1042026878, 26127344,
107962116, -4177212, -4538160, 475200
\end{Verbatim}
The defining transfer satisfies $TR(T)\mathbf1=0$ in every coordinate, so \eqref{round9:eq:r9arr-rec} follows. Source: \url{https://oeis.org/A234126}; parent column 3 of \url{https://oeis.org/A234131}. Both retrieved entries mark the relevant recurrence empirical.
\Needspace{7\baselineskip}
\subsubsection{A234133}
\label{round9:r9entry:A234133}
Arrays have $n+1$ rows, width $w=2$, and labeled alphabet $0,\ldots,3$. Predicate U has target $s=9$, allowing equal adjacent values.
The recurrence holds for every $n\ge 4$, with order $d=3$ and a complete $D=6$ state matrix. The first three source values are $(32,80,192)$.
\[(c_1,\ldots,c_d)=(2,8,-16).\]
The defining transfer satisfies $TR(T)\mathbf1=0$ in every coordinate, so \eqref{round9:eq:r9arr-rec} follows. Source: \url{https://oeis.org/A234133}; parent column 1 of \url{https://oeis.org/A234140}. Both retrieved entries mark the relevant recurrence empirical.
The source rational generating function is also proved: $F=P/Q$, with $Q=1-\sum_{i=1}^{d}c_i x^i$ and
\[P(x)=32x+16x^{2}-224x^{3}.\]
\Needspace{7\baselineskip}
\subsubsection{A234135}
\label{round9:r9entry:A234135}
Arrays have $n+1$ rows, width $w=4$, and labeled alphabet $0,\ldots,3$. Predicate U has target $s=9$, allowing equal adjacent values.
The recurrence holds for every $n\ge 7$, with order $d=6$ and a complete $D=72$ state matrix. The first three source values are $(192,296,488)$.
\[(c_1,\ldots,c_d)=(3,8,-30,4,48,-32).\]
The defining transfer satisfies $TR(T)\mathbf1=0$ in every coordinate, so \eqref{round9:eq:r9arr-rec} follows. Source: \url{https://oeis.org/A234135}; parent column 3 of \url{https://oeis.org/A234140}. Both retrieved entries mark the relevant recurrence empirical.
The source rational generating function is also proved: $F=P/Q$, with $Q=1-\sum_{i=1}^{d}c_i x^i$ and
\[P(x)=192x-280x^{2}-1936x^{3}+2896x^{4}+3296x^{5}-4672x^{6}.\]
\Needspace{7\baselineskip}
\subsubsection{A234154}
\label{round9:r9entry:A234154}
Arrays have $n+1$ rows, width $w=2$, and labeled alphabet $0,\ldots,3$. Predicate B has target $s=1$, allowing equal adjacent values.
The recurrence holds for every $n\ge 4$, with order $d=3$ and a complete $D=6$ state matrix. The first three source values are $(46,226,1108)$.
\[(c_1,\ldots,c_d)=(7,-9,-6).\]
The defining transfer satisfies $TR(T)\mathbf1=0$ in every coordinate, so \eqref{round9:eq:r9arr-rec} follows. Source: \url{https://oeis.org/A234154}; parent column 1 of \url{https://oeis.org/A234161}. Both retrieved entries mark the relevant recurrence empirical.
The source rational generating function is also proved: $F=P/Q$, with $Q=1-\sum_{i=1}^{d}c_i x^i$ and
\[P(x)=46x-96x^{2}-60x^{3}.\]
\Needspace{7\baselineskip}
\subsubsection{A234155}
\label{round9:r9entry:A234155}
Arrays have $n+1$ rows, width $w=3$, and labeled alphabet $0,\ldots,3$. Predicate B has target $s=1$, allowing equal adjacent values.
The recurrence holds for every $n\ge 7$, with order $d=6$ and a complete $D=20$ state matrix. The first three source values are $(226,2210,21428)$.
\[(c_1,\ldots,c_d)=(16,-59,-31,118,35,-10).\]
The defining transfer satisfies $TR(T)\mathbf1=0$ in every coordinate, so \eqref{round9:eq:r9arr-rec} follows. Source: \url{https://oeis.org/A234155}; parent column 2 of \url{https://oeis.org/A234161}. Both retrieved entries mark the relevant recurrence empirical.
The source rational generating function is also proved: $F=P/Q$, with $Q=1-\sum_{i=1}^{d}c_i x^i$ and
\[P(x)=226x-1406x^{2}-598x^{3}+3078x^{4}+900x^{5}-260x^{6}.\]
\Needspace{7\baselineskip}
\subsubsection{A234156}
\label{round9:r9entry:A234156}
Arrays have $n+1$ rows, width $w=4$, and labeled alphabet $0,\ldots,3$. Predicate B has target $s=1$, allowing equal adjacent values.
The recurrence holds for every $n\ge 12$, with order $d=11$ and a complete $D=72$ state matrix. The first three source values are $(1108,21428,411712)$.
\begin{Verbatim}[fontsize=\scriptsize,breaklines=true]
38, -412, 740, 5548, -11466, -28213, 27544, 43860, -11608, -8172, 1584
\end{Verbatim}
The defining transfer satisfies $TR(T)\mathbf1=0$ in every coordinate, so \eqref{round9:eq:r9arr-rec} follows. Source: \url{https://oeis.org/A234156}; parent column 3 of \url{https://oeis.org/A234161}. Both retrieved entries mark the relevant recurrence empirical.
\Needspace{7\baselineskip}
\subsubsection{A234157}
\label{round9:r9entry:A234157}
Arrays have $n+1$ rows, width $w=5$, and labeled alphabet $0,\ldots,3$. Predicate B has target $s=1$, allowing equal adjacent values.
The recurrence holds for every $n\ge 27$, with order $d=26$ and a complete $D=272$ state matrix. The first three source values are $(5446,208530,7952532)$.
\begin{Verbatim}[fontsize=\scriptsize,breaklines=true]
90, -2653, 24943, 89825, -2218609, 1432019, 71490923, -78048525, -1164227315,
761044775, 9924372783, -667076603, -41948446347, -15809876567, 78900501505,
41327430666, -75139994803, -37071907734, 40119015622, 12686351262, -12371945864,
-754232304, 1729058752, -288379744, -6642240, 2957184
\end{Verbatim}
The defining transfer satisfies $TR(T)\mathbf1=0$ in every coordinate, so \eqref{round9:eq:r9arr-rec} follows. Source: \url{https://oeis.org/A234157}; parent column 4 of \url{https://oeis.org/A234161}. Both retrieved entries mark the relevant recurrence empirical.
\Needspace{7\baselineskip}
\subsubsection{A234158}
\label{round9:r9entry:A234158}
Arrays have $n+1$ rows, width $w=6$, and labeled alphabet $0,\ldots,3$. Predicate B has target $s=1$, allowing equal adjacent values.
The recurrence holds for every $n\ge 54$, with order $d=53$ and a complete $D=1056$ state matrix. The first three source values are $(26794,2031286,153793284)$.
\begin{Verbatim}[fontsize=\scriptsize,breaklines=true]
217, -17184, 575274, -4536464, -187521807, 3793789039, 16856009949, -838186556820,
566323425072, 100894927353562, -209284767431095, -7820748769614354, 15947814680267338,
413533369489268493, -552537406964179485, -15031120049686806579, 6646617078976384635,
370659665896453216859, 135453696144913493680, -6128560315932176342756,
-6096479285302210528097, 67946244106778846892950, 98850271203664956510736,
-512865758968691586413028, -912755411333156100506867, 2714978138747552717986492,
5387366562109720554218542, -10494361350111126990456413, -21466605718448227486566037,
30983177388450723856730013, 59383999283149719935113021, -72421333273384167090658391,
-114767457860090624328317490, 135071826569372046710954792,
151628599152427140815306228, -195506583085447410842602076,
-126909336149740745701919608, 208023789140373924579568064, 50774598029282885589249600,
-151206017630995324948864864, 11766570188661225224139904, 67386656389693313023099968,
-23587940089762795376954112, -14522140434485866540641792, 10175560813561814524537344,
535272730559164256256, -1562960541170454941048832, 425775164053979683454976,
9006299347542687350784, -19411131518479690629120, 2198238267084808126464,
101315221269019361280, -21080039086463385600
\end{Verbatim}
The defining transfer satisfies $TR(T)\mathbf1=0$ in every coordinate, so \eqref{round9:eq:r9arr-rec} follows. Source: \url{https://oeis.org/A234158}; parent column 5 of \url{https://oeis.org/A234161}. Both retrieved entries mark the relevant recurrence empirical.
The empirical coefficient list is in the primary entry's linked attachment: \url{https://oeis.org/A234158/a234158.txt}.
\Needspace{7\baselineskip}
\subsubsection{A234170}
\label{round9:r9entry:A234170}
Arrays have $n+1$ rows, width $w=2$, and labeled alphabet $0,\ldots,6$. Predicate D has target $s=1$, allowing equal adjacent values.
The recurrence holds for every $n\ge 7$, with order $d=6$ and a complete $D=16$ state matrix. The first three source values are $(448,4536,47936)$.
\[(c_1,\ldots,c_d)=(0,154,0,-4144,0,16176).\]
The defining transfer satisfies $TR(T)\mathbf1=0$ in every coordinate, so \eqref{round9:eq:r9arr-rec} follows. Source: \url{https://oeis.org/A234170}; parent column 1 of \url{https://oeis.org/A234175}. The primary entry marks the recurrence empirical; the parent content was not retrieved, as disclosed in the source audit.
The source rational generating function is also proved: $F=P/Q$, with $Q=1-\sum_{i=1}^{d}c_i x^i$ and
\[P(x)=448x+4536x^{2}-21056x^{3}-181440x^{4}+105216x^{5}+787584x^{6}.\]
\Needspace{7\baselineskip}
\subsubsection{A234171}
\label{round9:r9entry:A234171}
Arrays have $n+1$ rows, width $w=3$, and labeled alphabet $0,\ldots,6$. Predicate D has target $s=1$, allowing equal adjacent values.
The recurrence holds for every $n\ge 36$, with order $d=35$ and a complete $D=100$ state matrix. The first three source values are $(4536,70486,1173940)$.
\begin{Verbatim}[fontsize=\scriptsize,breaklines=true]
33, 323, -20479, 66827, 4196961, -32408967, -381895309, 4395837462, 15113145790,
-293548976068, -51289577932, 10922776315912, -16766805352904, -235952101172000,
626750889441184, 2967198544937088, -10881957631245888, -20625630149777664,
106077230575696640, 62012982701743616, -603993107147691008, 80220923857108992,
1983943732124129280, -1120944653090611200, -3571475321182953472, 3062535903549603840,
3148640202567516160, -3550878140723822592, -1004211355174305792, 1711073233311105024,
-31999222443147264, -285978658647048192, 29150317966786560, 14870074192035840,
-1882287872409600
\end{Verbatim}
The defining transfer satisfies $TR(T)\mathbf1=0$ in every coordinate, so \eqref{round9:eq:r9arr-rec} follows. Source: \url{https://oeis.org/A234171}; parent column 2 of \url{https://oeis.org/A234175}. The primary entry marks the recurrence empirical; the parent content was not retrieved, as disclosed in the source audit.
\Needspace{7\baselineskip}
\subsubsection{A234177}
\label{round9:r9entry:A234177}
Arrays have $n+1$ rows, width $w=2$, and labeled alphabet $0,\ldots,5$. Predicate B has target $s=1$, allowing equal adjacent values.
The recurrence holds for every $n\ge 6$, with order $d=5$ and a complete $D=12$ state matrix. The first three source values are $(76,388,1978)$.
\[(c_1,\ldots,c_d)=(10,-27,3,36,12).\]
The defining transfer satisfies $TR(T)\mathbf1=0$ in every coordinate, so \eqref{round9:eq:r9arr-rec} follows. Source: \url{https://oeis.org/A234177}; parent column 1 of \url{https://oeis.org/A234184}. Both retrieved entries mark the relevant recurrence empirical.
The source rational generating function is also proved: $F=P/Q$, with $Q=1-\sum_{i=1}^{d}c_i x^i$ and
\[P(x)=76x-372x^{2}+150x^{3}+588x^{4}+192x^{5}.\]
\Needspace{7\baselineskip}
\subsubsection{A234178}
\label{round9:r9entry:A234178}
Arrays have $n+1$ rows, width $w=3$, and labeled alphabet $0,\ldots,5$. Predicate B has target $s=1$, allowing equal adjacent values.
The recurrence holds for every $n\ge 12$, with order $d=11$ and a complete $D=63$ state matrix. The first three source values are $(388,3908,39098)$.
\begin{Verbatim}[fontsize=\scriptsize,breaklines=true]
24, -182, 349, 938, -2021, -1914, 2375, 1747, -219, -140, 20
\end{Verbatim}
The defining transfer satisfies $TR(T)\mathbf1=0$ in every coordinate, so \eqref{round9:eq:r9arr-rec} follows. Source: \url{https://oeis.org/A234178}; parent column 2 of \url{https://oeis.org/A234184}. Both retrieved entries mark the relevant recurrence empirical.
\Needspace{7\baselineskip}
\subsubsection{A234179}
\label{round9:r9entry:A234179}
Arrays have $n+1$ rows, width $w=4$, and labeled alphabet $0,\ldots,5$. Predicate B has target $s=1$, allowing equal adjacent values.
The recurrence holds for every $n\ge 23$, with order $d=22$ and a complete $D=342$ state matrix. The first three source values are $(1978,39098,771710)$.
\begin{Verbatim}[fontsize=\scriptsize,breaklines=true]
59, -1253, 10658, -12393, -283822, 938714, 2905662, -12071951, -18482610, 67209819,
84336539, -168519304, -218405166, 167608180, 253438774, -36105806, -96411410, 1835064,
14870944, -529572, -839520, 95040
\end{Verbatim}
The defining transfer satisfies $TR(T)\mathbf1=0$ in every coordinate, so \eqref{round9:eq:r9arr-rec} follows. Source: \url{https://oeis.org/A234179}; parent column 3 of \url{https://oeis.org/A234184}. Both retrieved entries mark the relevant recurrence empirical.
\Needspace{7\baselineskip}
\subsubsection{A234201}
\label{round9:r9entry:A234201}
Arrays have $n+1$ rows, width $w=2$, and labeled alphabet $0,\ldots,5$. Predicate U has target $s=15$, allowing equal adjacent values.
The recurrence holds for every $n\ge 11$, with order $d=10$ and a complete $D=12$ state matrix. The first three source values are $(104,440,1792)$.
\begin{Verbatim}[fontsize=\scriptsize,breaklines=true]
0, 36, 0, -396, 0, 1392, 0, -1664, 0, 512
\end{Verbatim}
The defining transfer satisfies $TR(T)\mathbf1=0$ in every coordinate, so \eqref{round9:eq:r9arr-rec} follows. Source: \url{https://oeis.org/A234201}; parent column 1 of \url{https://oeis.org/A234208}. Both retrieved entries mark the relevant recurrence empirical.
\Needspace{7\baselineskip}
\subsubsection{A234210}
\label{round9:r9entry:A234210}
Arrays have $n+1$ rows, width $w=2$, and labeled alphabet $0,\ldots,4$. Predicate V has target $s=11$, allowing equal adjacent values.
The recurrence holds for every $n\ge 5$, with order $d=4$ and a complete $D=9$ state matrix. The first three source values are $(72,308,1248)$.
\[(c_1,\ldots,c_d)=(0,24,0,-112).\]
The defining transfer satisfies $TR(T)\mathbf1=0$ in every coordinate, so \eqref{round9:eq:r9arr-rec} follows. Source: \url{https://oeis.org/A234210}; parent column 1 of \url{https://oeis.org/A234217}. Both retrieved entries mark the relevant recurrence empirical.
The source rational generating function is also proved: $F=P/Q$, with $Q=1-\sum_{i=1}^{d}c_i x^i$ and
\[P(x)=72x+308x^{2}-480x^{3}-2048x^{4}.\]
\Needspace{7\baselineskip}
\subsubsection{A234211}
\label{round9:r9entry:A234211}
Arrays have $n+1$ rows, width $w=3$, and labeled alphabet $0,\ldots,4$. Predicate V has target $s=11$, allowing equal adjacent values.
The recurrence holds for every $n\ge 16$, with order $d=15$ and a complete $D=39$ state matrix. The first three source values are $(308,1540,7348)$.
\begin{Verbatim}[fontsize=\scriptsize,breaklines=true]
6, 51, -354, -656, 6790, -457, -52786, 47381, 158504, -209462, -148184, 260800, 12192,
-94128, 21312
\end{Verbatim}
The defining transfer satisfies $TR(T)\mathbf1=0$ in every coordinate, so \eqref{round9:eq:r9arr-rec} follows. Source: \url{https://oeis.org/A234211}; parent column 2 of \url{https://oeis.org/A234217}. Both retrieved entries mark the relevant recurrence empirical.
\Needspace{7\baselineskip}
\subsubsection{A234212}
\label{round9:r9entry:A234212}
Arrays have $n+1$ rows, width $w=4$, and labeled alphabet $0,\ldots,4$. Predicate V has target $s=11$, allowing equal adjacent values.
The recurrence holds for every $n\ge 57$, with order $d=56$ and a complete $D=169$ state matrix. The first three source values are $(1248,7348,40384)$.
\begin{Verbatim}[fontsize=\scriptsize,breaklines=true]
0, 223, 0, -22136, 0, 1303842, 0, -51198565, 0, 1426660738, 0, -29308622249, 0,
455180602960, 0, -5436994724661, 0, 50553279723332, 0, -368988227986457, 0,
2126351180635715, 0, -9709003146923883, 0, 35189269788128342, 0, -101259728165773203,
0, 231060390001397360, 0, -417017561764832798, 0, 592925182387254959, 0,
-660503855304387479, 0, 572258984227035715, 0, -381879096938313788, 0,
193753758860748742, 0, -73452183198353516, 0, 20317503519116744, 0, -3966504433028624,
0, 520682605459552, 0, -42615909301440, 0, 1905181236480, 0, -34505662464
\end{Verbatim}
The defining transfer satisfies $TR(T)\mathbf1=0$ in every coordinate, so \eqref{round9:eq:r9arr-rec} follows. Source: \url{https://oeis.org/A234212}; parent column 3 of \url{https://oeis.org/A234217}. Both retrieved entries mark the relevant recurrence empirical.
\Needspace{7\baselineskip}
\subsubsection{A234659}
\label{round9:r9entry:A234659}
Arrays have $n+1$ rows, width $w=2$, and labeled alphabet $0,\ldots,6$. Predicate D has target $s=4$, allowing equal adjacent values.
The recurrence holds for every $n\ge 6$, with order $d=5$ and a complete $D=16$ state matrix. The first three source values are $(298,1820,11112)$.
\[(c_1,\ldots,c_d)=(11,17,-492,1050,1260).\]
The defining transfer satisfies $TR(T)\mathbf1=0$ in every coordinate, so \eqref{round9:eq:r9arr-rec} follows. Source: \url{https://oeis.org/A234659}; parent column 1 of \url{https://oeis.org/A234665}. Both retrieved entries mark the relevant recurrence empirical.
The source rational generating function is also proved: $F=P/Q$, with $Q=1-\sum_{i=1}^{d}c_i x^i$ and
\[P(x)=298x-1458x^{2}-13974x^{3}+61530x^{4}+61740x^{5}.\]
\Needspace{7\baselineskip}
\subsubsection{A234660}
\label{round9:r9entry:A234660}
Arrays have $n+1$ rows, width $w=3$, and labeled alphabet $0,\ldots,6$. Predicate D has target $s=4$, allowing equal adjacent values.
The recurrence holds for every $n\ge 17$, with order $d=16$ and a complete $D=100$ state matrix. The first three source values are $(1820,9774,52646)$.
\begin{Verbatim}[fontsize=\scriptsize,breaklines=true]
24, -116, -1573, 17655, -7512, -599532, 2201681, 4899218, -40315144, 38576320,
130316460, -132305400, -263116800, 8496000, 134568000, 45360000
\end{Verbatim}
The defining transfer satisfies $TR(T)\mathbf1=0$ in every coordinate, so \eqref{round9:eq:r9arr-rec} follows. Source: \url{https://oeis.org/A234660}; parent column 2 of \url{https://oeis.org/A234665}. Both retrieved entries mark the relevant recurrence empirical.
\Needspace{7\baselineskip}
\subsubsection{A234661}
\label{round9:r9entry:A234661}
Arrays have $n+1$ rows, width $w=4$, and labeled alphabet $0,\ldots,6$. Predicate D has target $s=4$, allowing equal adjacent values.
The recurrence holds for every $n\ge 40$, with order $d=39$ and a complete $D=625$ state matrix. The first three source values are $(11112,52646,250792)$.
\begin{Verbatim}[fontsize=\scriptsize,breaklines=true]
37, -377, -2618, 77447, -294799, -4757619, 45540630, 61891376, -2437351502,
6395028748, 63315544498, -379499690085, -574635388877, 9755368316951, -11625591036256,
-133055653485659, 418266455931349, 823106610584361, -5467891684835430,
1031251784123218, 36193728208088096, -44705964178611952, -124471844677347304,
241456322415418264, 272081034415319536, -671531537886950352, -480780534963511680,
1128701330288740800, 770576980895712000, -1120320571251686400, -920761444521216000,
573508008695808000, 665567689121280000, -74650786306560000, -247547977113600000,
-49525866547200000, 32441969664000000, 15456946176000000, 1881169920000000
\end{Verbatim}
The defining transfer satisfies $TR(T)\mathbf1=0$ in every coordinate, so \eqref{round9:eq:r9arr-rec} follows. Source: \url{https://oeis.org/A234661}; parent column 3 of \url{https://oeis.org/A234665}. Both retrieved entries mark the relevant recurrence empirical.
\Needspace{7\baselineskip}
\subsubsection{A234667}
\label{round9:r9entry:A234667}
Arrays have $n+1$ rows, width $w=2$, and labeled alphabet $0,\ldots,7$. Predicate D has target $s=4$, allowing equal adjacent values.
The recurrence holds for every $n\ge 8$, with order $d=7$ and a complete $D=20$ state matrix. The first three source values are $(492,3816,29568)$.
\[(c_1,\ldots,c_d)=(14,25,-950,2091,7326,-3780,-10080).\]
The defining transfer satisfies $TR(T)\mathbf1=0$ in every coordinate, so \eqref{round9:eq:r9arr-rec} follows. Source: \url{https://oeis.org/A234667}; parent column 1 of \url{https://oeis.org/A234672}. Both retrieved entries mark the relevant recurrence empirical.
The source rational generating function is also proved: $F=P/Q$, with $Q=1-\sum_{i=1}^{d}c_i x^i$ and
\[P(x)=492x-3072x^{2}-36156x^{3}+187320x^{4}+425832x^{5}-322560x^{6}-645120x^{7}.\]
\Needspace{7\baselineskip}
\subsubsection{A234668}
\label{round9:r9entry:A234668}
Arrays have $n+1$ rows, width $w=3$, and labeled alphabet $0,\ldots,7$. Predicate D has target $s=4$, allowing equal adjacent values.
The recurrence holds for every $n\ge 32$, with order $d=31$ and a complete $D=144$ state matrix. The first three source values are $(3816,28860,217968)$.
\begin{Verbatim}[fontsize=\scriptsize,breaklines=true]
35, -234, -5393, 83088, 72599, -7727281, 30601466, 288437560, -2223662740,
-2827663592, 61030210280, -68134000167, -723492961431, 1379135043174, 4956504550981,
-10336340466744, -22249573436775, 42072453243081, 68185197096618, -103181037099516,
-142550751258288, 157651414286112, 199205608411008, -149554834911360,
-179827747799040, 84774205497600, 99173800742400, -26026116480000, -29990476800000,
3292047360000, 3762339840000
\end{Verbatim}
The defining transfer satisfies $TR(T)\mathbf1=0$ in every coordinate, so \eqref{round9:eq:r9arr-rec} follows. Source: \url{https://oeis.org/A234668}; parent column 2 of \url{https://oeis.org/A234672}. Both retrieved entries mark the relevant recurrence empirical.
\Needspace{7\baselineskip}
\subsubsection{A234683}
\label{round9:r9entry:A234683}
Arrays have $n+1$ rows, width $w=2$, and labeled alphabet $0,\ldots,6$. Predicate D has target $s=5$, allowing equal adjacent values.
The recurrence holds for every $n\ge 7$, with order $d=6$ and a complete $D=16$ state matrix. The first three source values are $(232,1148,5472)$.
\[(c_1,\ldots,c_d)=(0,72,0,-1712,0,13440).\]
The defining transfer satisfies $TR(T)\mathbf1=0$ in every coordinate, so \eqref{round9:eq:r9arr-rec} follows. Source: \url{https://oeis.org/A234683}; parent column 1 of \url{https://oeis.org/A234690}. Both retrieved entries mark the relevant recurrence empirical.
The source rational generating function is also proved: $F=P/Q$, with $Q=1-\sum_{i=1}^{d}c_i x^i$ and
\[P(x)=232x+1148x^{2}-11232x^{3}-55312x^{4}+134400x^{5}+658560x^{6}.\]
\Needspace{7\baselineskip}
\subsubsection{A234684}
\label{round9:r9entry:A234684}
Arrays have $n+1$ rows, width $w=3$, and labeled alphabet $0,\ldots,6$. Predicate D has target $s=5$, allowing equal adjacent values.
The recurrence holds for every $n\ge 25$, with order $d=24$ and a complete $D=100$ state matrix. The first three source values are $(1148,4416,16580)$.
\begin{Verbatim}[fontsize=\scriptsize,breaklines=true]
7, 136, -1036, -7837, 67291, 245614, -2526790, -4355332, 60808804, 36086728,
-982312744, 135510432, 10839179904, -6945594624, -81450996480, 81465827328,
407151166464, -514409232384, -1284949370880, 1874710241280, 2296420761600,
-3685898649600, -1755758592000, 3009871872000
\end{Verbatim}
The defining transfer satisfies $TR(T)\mathbf1=0$ in every coordinate, so \eqref{round9:eq:r9arr-rec} follows. Source: \url{https://oeis.org/A234684}; parent column 2 of \url{https://oeis.org/A234690}. Both retrieved entries mark the relevant recurrence empirical.
\Needspace{7\baselineskip}
\subsubsection{A234697}
\label{round9:r9entry:A234697}
Arrays have $n+1$ rows, width $w=2$, and labeled alphabet $0,\ldots,7$. Predicate D has target $s=5$, allowing equal adjacent values.
The recurrence holds for every $n\ge 6$, with order $d=5$ and a complete $D=20$ state matrix. The first three source values are $(408,2628,16848)$.
\[(c_1,\ldots,c_d)=(6,90,-540,-2016,12096).\]
The defining transfer satisfies $TR(T)\mathbf1=0$ in every coordinate, so \eqref{round9:eq:r9arr-rec} follows. Source: \url{https://oeis.org/A234697}; parent column 1 of \url{https://oeis.org/A234702}. Both retrieved entries mark the relevant recurrence empirical.
The source rational generating function is also proved: $F=P/Q$, with $Q=1-\sum_{i=1}^{d}c_i x^i$ and
\[P(x)=408x+180x^{2}-35640x^{3}-8064x^{4}+774144x^{5}.\]
\Needspace{7\baselineskip}
\subsubsection{A234698}
\label{round9:r9entry:A234698}
Arrays have $n+1$ rows, width $w=3$, and labeled alphabet $0,\ldots,7$. Predicate D has target $s=5$, allowing equal adjacent values.
The recurrence holds for every $n\ge 16$, with order $d=15$ and a complete $D=144$ state matrix. The first three source values are $(2628,13952,74308)$.
\begin{Verbatim}[fontsize=\scriptsize,breaklines=true]
15, 125, -2865, -1490, 219360, -507740, -8437740, 34905816, 161734320, -1006299360,
-1062547200, 13887417600, -8926848000, -75188736000, 121927680000
\end{Verbatim}
The defining transfer satisfies $TR(T)\mathbf1=0$ in every coordinate, so \eqref{round9:eq:r9arr-rec} follows. Source: \url{https://oeis.org/A234698}; parent column 2 of \url{https://oeis.org/A234702}. Both retrieved entries mark the relevant recurrence empirical.
\Needspace{7\baselineskip}
\subsubsection{A234699}
\label{round9:r9entry:A234699}
Arrays have $n+1$ rows, width $w=4$, and labeled alphabet $0,\ldots,7$. Predicate D has target $s=5$, allowing equal adjacent values.
The recurrence holds for every $n\ge 32$, with order $d=31$ and a complete $D=1056$ state matrix. The first three source values are $(16848,74308,329984)$.
\begin{Verbatim}[fontsize=\scriptsize,breaklines=true]
20, 141, -5340, 5451, 616740, -2678805, -39706740, 278440182, 1491646680,
-16092380232, -27470360160, 601403119168, -194332088960, -15257818574160,
26729605928640, 263906583339744, -801813357648000, -2982678610973184,
13948745101747200, 18588639920403456, -156794289434542080, -1681574850293760,
1140373919202508800, -1022866437252096000, -5015897891205120000, 8536532135362560000,
10656426280550400000, -31534442652303360000, 815374290124800000, 46191901645209600000,
-31856483893248000000
\end{Verbatim}
The defining transfer satisfies $TR(T)\mathbf1=0$ in every coordinate, so \eqref{round9:eq:r9arr-rec} follows. Source: \url{https://oeis.org/A234699}; parent column 3 of \url{https://oeis.org/A234702}. Both retrieved entries mark the relevant recurrence empirical.
\Needspace{7\baselineskip}
\subsubsection{A234706}
\label{round9:r9entry:A234706}
Arrays have $n+1$ rows, width $w=3$, and labeled alphabet $0,\ldots,7$. Predicate D has target $s=6$, allowing equal adjacent values.
The recurrence holds for every $n\ge 32$, with order $d=31$ and a complete $D=144$ state matrix. The first three source values are $(1730,7154,28170)$.
\begin{Verbatim}[fontsize=\scriptsize,breaklines=true]
12, 159, -2412, -9452, 217248, 161262, -11578680, 10718427, 405901116, -824307633,
-9828766476, 28650887914, 167599827384, -632763934044, -2003077953360, 9664372451624,
16138030431072, -105143721192576, -76625123046912, 819575849887488, 75093856109568,
-4528143100320768, 1633754973450240, 17236678723430400, -11938042414080000,
-42782336704512000, 40275310510080000, 62010467205120000, -70124998164480000,
-39609913835520000, 50565847449600000
\end{Verbatim}
The defining transfer satisfies $TR(T)\mathbf1=0$ in every coordinate, so \eqref{round9:eq:r9arr-rec} follows. Source: \url{https://oeis.org/A234706}; parent column 2 of \url{https://oeis.org/A234712}. Both retrieved entries mark the relevant recurrence empirical.
\Needspace{7\baselineskip}
\subsubsection{A234707}
\label{round9:r9entry:A234707}
Arrays have $n+1$ rows, width $w=4$, and labeled alphabet $0,\ldots,7$. Predicate D has target $s=6$, allowing equal adjacent values.
The recurrence holds for every $n\ge 37$, with order $d=36$ and a complete $D=1056$ state matrix. The first three source values are $(8738,28170,87698)$.
\begin{Verbatim}[fontsize=\scriptsize,breaklines=true]
13, 163, -2779, -9420, 268200, 60770, -15439130, 22995763, 589369361, -1588530069,
-15656100243, 58787313370, 294085396250, -1450484076640, -3845621898440,
25550700910752, 32289252593664, -330371143423376, -113428722066352, 3163397610252320,
-1019517140684000, -22338341568842880, 18872130722296320, 114348342295028736,
-148610586143628288, -409381709070163968, 722386331480346624, 952125730341396480,
-2277950012653363200, -1180887662923776000, 4529957329207296000, 69777190748160000,
-5142605210910720000, 1663475920404480000, 2524921315983360000, -1415843728588800000
\end{Verbatim}
The defining transfer satisfies $TR(T)\mathbf1=0$ in every coordinate, so \eqref{round9:eq:r9arr-rec} follows. Source: \url{https://oeis.org/A234707}; parent column 3 of \url{https://oeis.org/A234712}. Both retrieved entries mark the relevant recurrence empirical.
\Needspace{7\baselineskip}
\subsubsection{A234721}
\label{round9:r9entry:A234721}
Arrays have $n+1$ rows, width $w=2$, and labeled alphabet $0,\ldots,7$. Predicate D has target $s=7$, allowing equal adjacent values.
The recurrence holds for every $n\ge 9$, with order $d=8$ and a complete $D=20$ state matrix. The first three source values are $(240,1096,4192)$.
\begin{Verbatim}[fontsize=\scriptsize,breaklines=true]
0, 68, 0, -1724, 0, 19312, 0, -80640
\end{Verbatim}
The defining transfer satisfies $TR(T)\mathbf1=0$ in every coordinate, so \eqref{round9:eq:r9arr-rec} follows. Source: \url{https://oeis.org/A234721}; parent column 1 of \url{https://oeis.org/A234728}. Both retrieved entries mark the relevant recurrence empirical.
The source rational generating function is also proved: $F=P/Q$, with $Q=1-\sum_{i=1}^{d}c_i x^i$ and
\[\begin{aligned}P(x)&=240x+1096x^{2}-12128x^{3}-55408x^{4}+203200x^{5}+928672x^{6}-1128960x^{7}\\ &\quad -5160960x^{8}.\end{aligned}\]
\Needspace{7\baselineskip}
\subsubsection{A234722}
\label{round9:r9entry:A234722}
Arrays have $n+1$ rows, width $w=3$, and labeled alphabet $0,\ldots,7$. Predicate D has target $s=7$, allowing equal adjacent values.
The recurrence holds for every $n\ge 31$, with order $d=30$ and a complete $D=144$ state matrix. The first three source values are $(1096,3976,12464)$.
\begin{Verbatim}[fontsize=\scriptsize,breaklines=true]
4, 137, -560, -8456, 35504, 310702, -1349320, -7561533, 34294092, 128206281,
-615707400, -1548288698, 8040276992, 13333179376, -77453548480, -79946374816,
552146144704, 309222145232, -2893327015040, -567224779296, 10948880162304,
-974085452928, -28950298675200, 9103597862400, 50436504576000, -24908850432000,
-51674112000000, 32903055360000, 23410114560000, -17557585920000
\end{Verbatim}
The defining transfer satisfies $TR(T)\mathbf1=0$ in every coordinate, so \eqref{round9:eq:r9arr-rec} follows. Source: \url{https://oeis.org/A234722}; parent column 2 of \url{https://oeis.org/A234728}. Both retrieved entries mark the relevant recurrence empirical.
\Needspace{7\baselineskip}
\subsubsection{A234816}
\label{round9:r9entry:A234816}
Arrays have $n+1$ rows, width $w=2$, and labeled alphabet $0,\ldots,6$. Predicate D has target $s=6$, allowing equal adjacent values.
The recurrence holds for every $n\ge 8$, with order $d=7$ and a complete $D=16$ state matrix. The first three source values are $(168,686,2380)$.
\[(c_1,\ldots,c_d)=(4,41,-164,-558,2232,2520,-10080).\]
The defining transfer satisfies $TR(T)\mathbf1=0$ in every coordinate, so \eqref{round9:eq:r9arr-rec} follows. Source: \url{https://oeis.org/A234816}; parent column 1 of \url{https://oeis.org/A234823}. Both retrieved entries mark the relevant recurrence empirical.
The source rational generating function is also proved: $F=P/Q$, with $Q=1-\sum_{i=1}^{d}c_i x^i$ and
\[P(x)=168x+14x^{2}-7252x^{3}-378x^{4}+103896x^{5}+2520x^{6}-493920x^{7}.\]
\Needspace{7\baselineskip}
\subsubsection{A234817}
\label{round9:r9entry:A234817}
Arrays have $n+1$ rows, width $w=3$, and labeled alphabet $0,\ldots,6$. Predicate D has target $s=6$, allowing equal adjacent values.
The recurrence holds for every $n\ge 24$, with order $d=23$ and a complete $D=100$ state matrix. The first three source values are $(686,2220,6282)$.
\begin{Verbatim}[fontsize=\scriptsize,breaklines=true]
8, 60, -620, -1210, 20740, 1580, -392180, 353267, 4606324, -7135240, -34677440,
72387960, 165390560, -443380480, -464032320, 1696321584, 566607168, -3948347520,
465384960, 5084985600, -2142028800, -2757888000, 1741824000
\end{Verbatim}
The defining transfer satisfies $TR(T)\mathbf1=0$ in every coordinate, so \eqref{round9:eq:r9arr-rec} follows. Source: \url{https://oeis.org/A234817}; parent column 2 of \url{https://oeis.org/A234823}. Both retrieved entries mark the relevant recurrence empirical.
\Needspace{7\baselineskip}
\subsubsection{A234818}
\label{round9:r9entry:A234818}
Arrays have $n+1$ rows, width $w=4$, and labeled alphabet $0,\ldots,6$. Predicate D has target $s=6$, allowing equal adjacent values.
The recurrence holds for every $n\ge 24$, with order $d=23$ and a complete $D=625$ state matrix. The first three source values are $(2380,6282,15224)$.
\begin{Verbatim}[fontsize=\scriptsize,breaklines=true]
8, 60, -620, -1210, 20740, 1580, -392180, 353267, 4606324, -7135240, -34677440,
72387960, 165390560, -443380480, -464032320, 1696321584, 566607168, -3948347520,
465384960, 5084985600, -2142028800, -2757888000, 1741824000
\end{Verbatim}
The defining transfer satisfies $TR(T)\mathbf1=0$ in every coordinate, so \eqref{round9:eq:r9arr-rec} follows. Source: \url{https://oeis.org/A234818}; parent column 3 of \url{https://oeis.org/A234823}. Both retrieved entries mark the relevant recurrence empirical.
\Needspace{7\baselineskip}
\subsubsection{A234992}
\label{round9:r9entry:A234992}
Arrays have $n+1$ rows, width $w=2$, and labeled alphabet $0,\ldots,6$. Predicate D has target $s=2$, with every horizontal and vertical neighboring pair unequal.
The recurrence holds for every $n\ge 11$, with order $d=10$ and a complete $D=16$ state matrix. The first three source values are $(172,916,5280)$.
\begin{Verbatim}[fontsize=\scriptsize,breaklines=true]
5, 21, -70, -144, 326, 381, -549, -351, 198, 54
\end{Verbatim}
The defining transfer satisfies $TR(T)\mathbf1=0$ in every coordinate, so \eqref{round9:eq:r9arr-rec} follows. Source: \url{https://oeis.org/A234992}; parent column 1 of \url{https://oeis.org/A234998}. Both retrieved entries mark the relevant recurrence empirical.
\Needspace{7\baselineskip}
\subsubsection{A234993}
\label{round9:r9entry:A234993}
Arrays have $n+1$ rows, width $w=3$, and labeled alphabet $0,\ldots,6$. Predicate D has target $s=2$, with every horizontal and vertical neighboring pair unequal.
The recurrence holds for every $n\ge 50$, with order $d=49$ and a complete $D=100$ state matrix. The first three source values are $(916,5112,32256)$.
\begin{Verbatim}[fontsize=\scriptsize,breaklines=true]
6, 145, -752, -9334, 41421, 350517, -1354277, -8617439, 29691561, 147908106,
-465456923, -1844354978, 5418592827, 17142962010, -47956757782, -120746606447,
327447020322, 650926343809, -1740121457481, -2698809487907, 7230222191327,
8610347154836, -23519481347436, -21060770608040, 59812533502435, 39165695095186,
-118451096631369, -54588838411578, 181491721699373, 55729099282996, -213160956817772,
-40145457720166, 189498178497244, 19155239362192, -125401972381171, -5433880826151,
60455696914684, 890128451025, -20656751347661, -330921215331, 4825908656200,
227383374666, -729157775008, -75605102388, 62860038344, 11433294720, -2022154176,
-653147712, -43972992
\end{Verbatim}
The defining transfer satisfies $TR(T)\mathbf1=0$ in every coordinate, so \eqref{round9:eq:r9arr-rec} follows. Source: \url{https://oeis.org/A234993}; parent column 2 of \url{https://oeis.org/A234998}. Both retrieved entries mark the relevant recurrence empirical.
\Needspace{7\baselineskip}
\subsubsection{A235169}
\label{round9:r9entry:A235169}
Arrays have $n+1$ rows, width $w=2$, and labeled alphabet $0,\ldots,7$. Predicate D has target $s=4$, with every horizontal and vertical neighboring pair unequal.
The recurrence holds for every $n\ge 10$, with order $d=9$ and a complete $D=20$ state matrix. The first three source values are $(268,1528,8688)$.
\begin{Verbatim}[fontsize=\scriptsize,breaklines=true]
6, 37, -234, -141, 1290, 108, -2340, 0, 1296
\end{Verbatim}
The defining transfer satisfies $TR(T)\mathbf1=0$ in every coordinate, so \eqref{round9:eq:r9arr-rec} follows. Source: \url{https://oeis.org/A235169}; parent column 1 of \url{https://oeis.org/A235175}. Both retrieved entries mark the relevant recurrence empirical.
\Needspace{7\baselineskip}
\subsubsection{A235170}
\label{round9:r9entry:A235170}
Arrays have $n+1$ rows, width $w=3$, and labeled alphabet $0,\ldots,7$. Predicate D has target $s=4$, with every horizontal and vertical neighboring pair unequal.
The recurrence holds for every $n\ge 47$, with order $d=46$ and a complete $D=144$ state matrix. The first three source values are $(1528,8380,45776)$.
\begin{Verbatim}[fontsize=\scriptsize,breaklines=true]
19, 34, -2719, 8988, 156665, -916999, -4550486, 40037121, 64400627, -1012400300,
-148677233, 16600041004, -9873660689, -189875213441, 187284339366, 1593590554819,
-1874014494735, -10158650544956, 12545534952885, 50324143130662, -60649685406985,
-196489095492359, 219272447037552, 609171221202828, -603073067708424,
-1502594545151892, 1270051970766120, 2940961682306376, -2045313082168560,
-4535601595588512, 2497426659701280, 5448277567918080, -2274890615189568,
-5012380431150912, 1506311626822272, 3449065504346496, -697075003450368,
-1716839475671808, 212360144961024, 588874155660288, -38772262462464,
-129032521869312, 3632324124672, 15786025500672, -130606940160, -783641640960
\end{Verbatim}
The defining transfer satisfies $TR(T)\mathbf1=0$ in every coordinate, so \eqref{round9:eq:r9arr-rec} follows. Source: \url{https://oeis.org/A235170}; parent column 2 of \url{https://oeis.org/A235175}. Both retrieved entries mark the relevant recurrence empirical.
\Needspace{7\baselineskip}
\subsubsection{A235205}
\label{round9:r9entry:A235205}
Arrays have $n+1$ rows, width $w=2$, and labeled alphabet $0,\ldots,7$. Predicate D has target $s=5$, with every horizontal and vertical neighboring pair unequal.
The recurrence holds for every $n\ge 7$, with order $d=6$ and a complete $D=20$ state matrix. The first three source values are $(240,1172,5696)$.
\[(c_1,\ldots,c_d)=(3,35,-81,-264,144,384).\]
The defining transfer satisfies $TR(T)\mathbf1=0$ in every coordinate, so \eqref{round9:eq:r9arr-rec} follows. Source: \url{https://oeis.org/A235205}; parent column 1 of \url{https://oeis.org/A235212}. Both retrieved entries mark the relevant recurrence empirical.
The source rational generating function is also proved: $F=P/Q$, with $Q=1-\sum_{i=1}^{d}c_i x^i$ and
\[P(x)=240x+452x^{2}-6220x^{3}-10856x^{4}+11016x^{5}+19136x^{6}.\]
\Needspace{7\baselineskip}
\subsubsection{A235206}
\label{round9:r9entry:A235206}
Arrays have $n+1$ rows, width $w=3$, and labeled alphabet $0,\ldots,7$. Predicate D has target $s=5$, with every horizontal and vertical neighboring pair unequal.
The recurrence holds for every $n\ge 30$, with order $d=29$ and a complete $D=144$ state matrix. The first three source values are $(1172,4616,18444)$.
\begin{Verbatim}[fontsize=\scriptsize,breaklines=true]
15, 23, -1373, 3547, 50847, -236177, -957405, 6739749, 8784137, -112122375, -8734283,
1205207433, -682793243, -8845512699, 8015164257, 45997463512, -47860413548,
-173398534628, 175140255652, 476432770592, -402722060848, -936013467040, 551524974720,
1245278007168, -366681768960, -999476834304, 6506053632, 362586931200, 89712230400
\end{Verbatim}
The defining transfer satisfies $TR(T)\mathbf1=0$ in every coordinate, so \eqref{round9:eq:r9arr-rec} follows. Source: \url{https://oeis.org/A235206}; parent column 2 of \url{https://oeis.org/A235212}. Both retrieved entries mark the relevant recurrence empirical.
\Needspace{7\baselineskip}
\subsubsection{A235233}
\label{round9:r9entry:A235233}
Arrays have $n+1$ rows, width $w=3$, and labeled alphabet $0,\ldots,6$. Predicate D has target $s=6$, with every horizontal and vertical neighboring pair unequal.
The recurrence holds for every $n\ge 29$, with order $d=28$ and a complete $D=100$ state matrix. The first three source values are $(420,1288,3688)$.
\begin{Verbatim}[fontsize=\scriptsize,breaklines=true]
6, 47, -331, -846, 7735, 6531, -99876, -4062, 780931, -294777, -3829232, 2187642,
11922251, -7106223, -23840076, 11482195, 31200479, -9197246, -26295345, 2967600,
13895070, 468804, -4403844, -634896, 759672, 171504, -54432, -15552
\end{Verbatim}
The defining transfer satisfies $TR(T)\mathbf1=0$ in every coordinate, so \eqref{round9:eq:r9arr-rec} follows. Source: \url{https://oeis.org/A235233}; parent column 2 of \url{https://oeis.org/A235239}. Both retrieved entries mark the relevant recurrence empirical.
\Needspace{7\baselineskip}
\subsubsection{A235234}
\label{round9:r9entry:A235234}
Arrays have $n+1$ rows, width $w=4$, and labeled alphabet $0,\ldots,6$. Predicate D has target $s=6$, with every horizontal and vertical neighboring pair unequal.
The recurrence holds for every $n\ge 41$, with order $d=40$ and a complete $D=625$ state matrix. The first three source values are $(1328,3688,9656)$.
\begin{Verbatim}[fontsize=\scriptsize,breaklines=true]
6, 75, -499, -2477, 18893, 46844, -431641, -544644, 6650508, 3692164, -73169776,
-7965360, 593911464, -108754278, -3624939268, 1348847209, 16812874388, -8151710835,
-59533275997, 31896084725, 160970283883, -85813326662, -331183400635, 159993289668,
514651400122, -202610873932, -596529887736, 165384664128, 505501523616, -76023552288,
-303948460656, 8914601088, 124404317856, 9463220928, -32567277312, -5301282816,
4869517824, 1128950784, -313528320, -89579520
\end{Verbatim}
The defining transfer satisfies $TR(T)\mathbf1=0$ in every coordinate, so \eqref{round9:eq:r9arr-rec} follows. Source: \url{https://oeis.org/A235234}; parent column 3 of \url{https://oeis.org/A235239}. Both retrieved entries mark the relevant recurrence empirical.


\section{Verification, provenance and reuse}
\label{round9:r9:verification}
The accompanying archive contains this standalone manuscript, the
numbered ledger, precise source statements, complete exact certificates,
verifiers, compact source audits and independent internal review reports.
With the Python dependencies in \texttt{requirements.txt} installed, run
\begin{Verbatim}[fontsize=\small]
python3 r9_run_all.py
\end{Verbatim}
from the extracted directory. No network access is needed. The runner
checks all 100 selections, the 808-ID exclusion ledger, the retained
prior-work baseline, review bindings and the included verification
programs. Actual versions and results are recorded in
\texttt{oeis\_hunt\_round9\_verification\_summary.json}.

\subsection{Universal proofs and supplementary finite checks}
The array definitions give complete finite transition matrices. The
identity $TR(T)\mathbf1=0$ proves the recurrence at every subsequent
index by multiplication by arbitrary nonnegative powers of $T$.
The 62 compressed matrices have 10,537 states and 30,926 nonzero
weighted edges. A separately written verifier, importing none of the
author's transfer engine, constructs all compatible pairs of labelled
rows. It checks the annihilator on all 39,370 uncompressed
states and every one of their 139,886 directed edges, then verifies
every row of the proposed equitable quotient. Source generating
functions in 20 of these packages also follow by exact polynomial
cross-multiplication, without an additional count.

Generating-function certificates establish exact differential or
rational identities. Every inhomogeneous boundary term is retained;
vanishing beyond its degree proves the stated recurrence for every
index in the specified range. The seed series is derived from the
definition or explicitly credited established input, independently of
the conjectural target. This avoids assuming the recurrence being proved.

The arithmetic and formal-series sections supply written universal
arguments. Complete residue tables are finite proof certificates where
the proof reduces all indices to those classes. Initial-term comparisons,
additional modular tests and numerical asymptotic checks are labelled
supplementary; a long finite prefix is never used as an infinite proof.

\subsection{Counting and priority limits}
The compact source records retain each selected target, its uncertainty
marker, attribution, offset, date, known inputs and audit notes. Source
hashes identify consulted snapshots without distributing entire raw web
pages. An entry's editorial ``approved'' status is not a mathematical
proof of every conjecture on that entry.

The audit excludes all prior project selections, including A176605,
A176607, A176610, A176611 and A176612, whose Round IV priority claims
were withdrawn in Round V. The present and cumulative counts refer to
selection packages. They are not certified totals of new discoveries.
Exact-ID and family searches, linked-paper inspections and bounded alias
screens are recorded with their actual coverage. Unretrieved material
and transformations outside the alias screen remain explicit limits.

\subsection{Standalone manuscript and reproducibility}
The delivered \texttt{.tex} includes every proof and recurrence signature
directly, without external \texttt{\string\input} or
\texttt{\string\include} fragments. Compile with pdfLaTeX until
cross-references stabilize. \texttt{r9\_build\_dossier.py} reconstructs
the manuscript, ledger and combined audit. The README lists all offline
regeneration commands. An isolated regeneration check compares the
reconstructed immutable mathematical outputs byte for byte.

The archive manifest records SHA-256 digests. Internal review reports
bind the precise mathematical files reviewed. The manuscript consistency
check compares printed coefficient lists, definitions, indices and
catalogue statuses against the exact records. The layout review binds
the PDF whose rendered pages were inspected.

\subsection{Common references}
\begin{enumerate}
\item The OEIS Foundation, \emph{The On-Line Encyclopedia of Integer
Sequences}, \url{https://oeis.org/}. Individual links and attributions
accompany each selected statement; established input results are credited
in the relevant proofs.
\item \emph{OEIS Open}, retained Appendix A comparison,
\url{https://arxiv.org/html/2608.11941v1}. The comparison covers the
100 costliest Appendix A claims on 97 distinct A-numbers, not the entire
project repository.
\item Daniel Axler, \emph{New estimates for some functions defined over
primes}, Integers 18 (2018), A52,
\url{https://math.colgate.edu/~integers/s52/s52.pdf}. Used as an established
input for the prime-harmonic estimate in A329354.
\item Mikl\'os B\'ona and Ryan Flynn, \emph{The average number of block
interchanges needed to sort a permutation and a recent result of Stanley},
\url{https://arxiv.org/abs/0811.0740}. The Hultman-number identity is a
credited input to A363412.
\end{enumerate}
% ===== END ORIGINAL ROUND 9 =====

% ===== BEGIN ORIGINAL ROUND 10 =====
\clearpage
\setcounter{section}{0}\setcounter{equation}{0}\setcounter{theorem}{0}
\setcounter{table}{0}\setcounter{figure}{0}
\normalsize\renewcommand{\arraystretch}{1}\renewcommand{\baselinestretch}{1}\selectfont
\part[Round 10: Universal periods, dimension formulas, tatami polynomials, and word arrays]{Universal periods, dimension formulas, tatami polynomials, and word arrays}
\label{combined:round10}
\rhead{\small Round 10 / KEIO}
\begingroup\small
\section*{Round overview and scope}
\addcontentsline{toc}{section}{Round 10 overview and scope}
This tenth-round dossier gives mathematical resolutions for 100 precisely
specified assertion packages on 100 distinct OEIS entries, excluding the 908
A-numbers selected in Rounds I--IX. It proves 97 selected assertions or
explicit questions and supplies corrected statements with full proofs for
three others. The four parts contain ten formal-series congruence packages,
five arithmetic packages, twenty generating-function or recurrence packages,
and sixty-five constrained-word recurrence packages.

\subsection*{Counting unit and source status}
One catalogue row corresponds to one selected package on one A-number.
A recurrence and its generating function count once. The 65 word entries
belong to twelve fixed-pattern families and share an automaton theorem;
thirteen modular-form dimension packages specialize established general
dimension formulas. These are not 100 independent proof techniques.

The selected source wording is conjectural or empirical in the consulted
primary material, except for A243046, which explicitly asks for a direct
formula. The exact result and domain are stated for each entry. A parity
proof at A337758 or A338328 does not resolve its separate modulo-four
claim. A339100 is treated only for the selected three-adic valuation.
The no-three-in-line comparison mentioned at A326478 is not asserted.
Known definitions, theorems and related formulas are credited as inputs.

\subsection*{Corrections and boundary conventions}
\begin{description}
\item[A339100.] The empirical three-adic valuation has its parity reversed.
The correct valuation is one at positive even indices and zero at odd indices.
\item[A094287 and A094288.] The printed conjectural generating functions
omit the constant coefficient added when the entries' offsets were changed
to include the empty path. Adding one restores the complete present-index
generating functions. Each correction is counted once; the underlying known
bounded-path theory is credited.
\end{description}
Within the A326478 package, the auxiliary comparison with A103517 also
requires the index $n-1$ in place of $n+1$. Its selected empirical even-index
generating function is true as stated, so this auxiliary repair is not
counted as another correction package.

\subsection*{Audit and priority limits}
Primary entries, relevant linked material and targeted searches were
consulted on 20 October 2026. Retrieval can be cached, so these source
records do not certify the latest live edit state. The retained OEIS Open
comparison covers 100 Appendix A claims on 97 unique A-numbers, with no
overlap; it does not cover that project's entire repository.

Known proofs found during the audit were excluded. Examples include the
general A272462/A272468 integrality assertion, A195986 after a recent
explicit proof was found, nine fractional-factorial candidates resolved
in an updated public repository, and the A242468--A242471 triple-equality
avoidance family covered by an earlier generating-function theorem.
Earlier priority withdrawals remain in force; 908 is a selection ledger,
not a cumulative total of worldwide-new discoveries. Fresh reserve
certificates are kept outside the 100 selected entries.

All selected word entries and their parent definitions were retrieved.
An unsuccessful retrieval of root companion A373313 is disclosed in its
audit; that companion supplies no needed proof input. Independent reviews
are internal checks by separate agents within one AI-assisted workflow,
not external peer review or proof-assistant formalization.

No OEIS submission or edit was performed. The dossier proves or repairs
the selected assertions; editorial acceptance is a separate event.
\endgroup

\clearpage
\section{Catalogue of the 100 selected entries}
\label{round10:r10:catalogue}
\textbf{P}: the selected assertion or question is proved over its stated domain.
\textbf{C}: the source formula is corrected and the corrected form is proved.
Page references lead to the complete argument or exact finite certificate.

\small
\begin{longtable}{@{}r p{21mm} p{103mm} c r@{}}
\toprule No.&Entry&Selected assertion package&Result&Page\\
\midrule\endfirsthead
\toprule No.&Entry&Selected assertion package&Result&Page\\
\midrule\endhead
\midrule\multicolumn{5}{r}{\emph{Continued on the next page}}\\\endfoot
\bottomrule\endlastfoot

\multicolumn{5}{@{}l}{\textbf{Formal series and universal congruences}}\\[2pt]

1&\oeis{A373312}&Binary-minus-one parity&P&\pageref{round10:r10entry:A373312}\\

2&\oeis{A325286}&Ternary digit residues&P&\pageref{round10:r10entry:A325286}\\

3&\oeis{A320949}&Divisibility by 3 of rationally summed moments&P&\pageref{round10:r10entry:A320949}\\

4&\oeis{A205572}&Full powers-of-two congruence&P&\pageref{round10:r10entry:A205572}\\

5&\oeis{A227051}&Last decimal digit 3&P&\pageref{round10:r10entry:A227051}\\

6&\oeis{A292177}&Prime-index vanishing in a bilateral limit&P&\pageref{round10:r10entry:A292177}\\

7&\oeis{A382124}&Cube-root Frobenius congruence&P&\pageref{round10:r10entry:A382124}\\

8&\oeis{A338328}&Parity of a diagonal coefficient definition&P&\pageref{round10:r10entry:A338328}\\

9&\oeis{A337758}&Parity with even diagonal exponents&P&\pageref{round10:r10entry:A337758}\\

10&\oeis{A230218}&Square-exponent parity&P&\pageref{round10:r10entry:A230218}\\

\multicolumn{5}{@{}l}{\textbf{Arithmetic, a direct formula and a valuation correction}}\\[2pt]

11&\oeis{A165886}&Consecutive normalized Hirzebruch quotient&P&\pageref{round10:r10entry:A165886}\\

12&\oeis{A243046}&Direct formula for the simultaneous integrality count&P&\pageref{round10:r10entry:A243046}\\

13&\oeis{A308833}&Tetrahedral divisibility and odd composites&P&\pageref{round10:r10entry:A308833}\\

14&\oeis{A326478}&Even-index Bernoulli-denominator generating function&P&\pageref{round10:r10entry:A326478}\\

15&\oeis{A339100}&Corrected three-adic valuation&C&\pageref{round10:r10entry:A339100}\\

\multicolumn{5}{@{}l}{\textbf{Generating functions, dimensions and polynomial recurrences}}\\[2pt]

16&\oeis{A063081}&Conjectural modular-form dimension formulas&P&\pageref{round10:r10entry:A063081}\\

17&\oeis{A063087}&Conjectural modular-form dimension formulas&P&\pageref{round10:r10entry:A063087}\\

18&\oeis{A063098}&Conjectural modular-form dimension formulas&P&\pageref{round10:r10entry:A063098}\\

19&\oeis{A063105}&Conjectural modular-form dimension formulas&P&\pageref{round10:r10entry:A063105}\\

20&\oeis{A063110}&Conjectural modular-form dimension formulas&P&\pageref{round10:r10entry:A063110}\\

21&\oeis{A063116}&Conjectural modular-form dimension formulas&P&\pageref{round10:r10entry:A063116}\\

22&\oeis{A063118}&Conjectural modular-form dimension formulas&P&\pageref{round10:r10entry:A063118}\\

23&\oeis{A063126}&Conjectural modular-form dimension formulas&P&\pageref{round10:r10entry:A063126}\\

24&\oeis{A063140}&Conjectural modular-form dimension formulas&P&\pageref{round10:r10entry:A063140}\\

25&\oeis{A063147}&Conjectural modular-form dimension formulas&P&\pageref{round10:r10entry:A063147}\\

26&\oeis{A063148}&Conjectural modular-form dimension formulas&P&\pageref{round10:r10entry:A063148}\\

27&\oeis{A063168}&Conjectural modular-form dimension formulas&P&\pageref{round10:r10entry:A063168}\\

28&\oeis{A063247}&Conjectural modular-form dimension formulas&P&\pageref{round10:r10entry:A063247}\\

29&\oeis{A006502}&The second Carlitz polynomial diagonal&P&\pageref{round10:r10entry:A006502}\\

30&\oeis{A259709}&The first Carlitz polynomial diagonal&P&\pageref{round10:r10entry:A259709}\\

31&\oeis{A174015}&Polynomial recurrence from the signed Catalan fraction&P&\pageref{round10:r10entry:A174015}\\

32&\oeis{A226302}&Tatami-polynomial specialization at minus one&P&\pageref{round10:r10entry:A226302}\\

33&\oeis{A273905}&Polynomial recurrence for symmetric bargraphs&P&\pageref{round10:r10entry:A273905}\\

34&\oeis{A094287}&Corrected constant term for bounded Motzkin paths&C&\pageref{round10:r10entry:A094287}\\

35&\oeis{A094288}&Corrected constant term for bounded Motzkin paths&C&\pageref{round10:r10entry:A094288}\\

\multicolumn{5}{@{}l}{\textbf{Consecutive-pattern word recurrences}}\\[2pt]

36&\oeis{A212824}&Empirical word recurrence of order 4&P&\pageref{round10:r10entry:A212824}\\

37&\oeis{A212825}&Empirical word recurrence of order 6&P&\pageref{round10:r10entry:A212825}\\

38&\oeis{A212826}&Empirical word recurrence of order 8&P&\pageref{round10:r10entry:A212826}\\

39&\oeis{A212827}&Empirical word recurrence of order 10&P&\pageref{round10:r10entry:A212827}\\

40&\oeis{A242543}&Empirical word recurrence of order 4&P&\pageref{round10:r10entry:A242543}\\

41&\oeis{A242544}&Empirical word recurrence of order 8&P&\pageref{round10:r10entry:A242544}\\

42&\oeis{A242546}&Empirical word recurrence of order 16&P&\pageref{round10:r10entry:A242546}\\

43&\oeis{A242547}&Empirical word recurrence of order 20&P&\pageref{round10:r10entry:A242547}\\

44&\oeis{A242548}&Empirical word recurrence of order 24&P&\pageref{round10:r10entry:A242548}\\

45&\oeis{A242578}&Empirical word recurrence of order 4&P&\pageref{round10:r10entry:A242578}\\

46&\oeis{A242579}&Empirical word recurrence of order 7&P&\pageref{round10:r10entry:A242579}\\

47&\oeis{A242580}&Empirical word recurrence of order 10&P&\pageref{round10:r10entry:A242580}\\

48&\oeis{A242581}&Empirical word recurrence of order 13&P&\pageref{round10:r10entry:A242581}\\

49&\oeis{A242582}&Empirical word recurrence of order 16&P&\pageref{round10:r10entry:A242582}\\

50&\oeis{A243027}&Empirical word recurrence of order 3&P&\pageref{round10:r10entry:A243027}\\

51&\oeis{A243028}&Empirical word recurrence of order 6&P&\pageref{round10:r10entry:A243028}\\

52&\oeis{A243029}&Empirical word recurrence of order 9&P&\pageref{round10:r10entry:A243029}\\

53&\oeis{A243038}&Empirical word recurrence of order 4&P&\pageref{round10:r10entry:A243038}\\

54&\oeis{A243039}&Empirical word recurrence of order 7&P&\pageref{round10:r10entry:A243039}\\

55&\oeis{A243040}&Empirical word recurrence of order 10&P&\pageref{round10:r10entry:A243040}\\

56&\oeis{A243041}&Empirical word recurrence of order 13&P&\pageref{round10:r10entry:A243041}\\

57&\oeis{A243042}&Empirical word recurrence of order 16&P&\pageref{round10:r10entry:A243042}\\

58&\oeis{A243043}&Empirical word recurrence of order 19&P&\pageref{round10:r10entry:A243043}\\

59&\oeis{A243383}&Empirical word recurrence of order 4&P&\pageref{round10:r10entry:A243383}\\

60&\oeis{A243384}&Empirical word recurrence of order 7&P&\pageref{round10:r10entry:A243384}\\

61&\oeis{A243385}&Empirical word recurrence of order 10&P&\pageref{round10:r10entry:A243385}\\

62&\oeis{A243386}&Empirical word recurrence of order 13&P&\pageref{round10:r10entry:A243386}\\

63&\oeis{A243387}&Empirical word recurrence of order 16&P&\pageref{round10:r10entry:A243387}\\

64&\oeis{A243388}&Empirical word recurrence of order 19&P&\pageref{round10:r10entry:A243388}\\

65&\oeis{A243513}&Empirical word recurrence of order 3&P&\pageref{round10:r10entry:A243513}\\

66&\oeis{A243514}&Empirical word recurrence of order 4&P&\pageref{round10:r10entry:A243514}\\

67&\oeis{A243515}&Empirical word recurrence of order 5&P&\pageref{round10:r10entry:A243515}\\

68&\oeis{A243516}&Empirical word recurrence of order 6&P&\pageref{round10:r10entry:A243516}\\

69&\oeis{A243517}&Empirical word recurrence of order 7&P&\pageref{round10:r10entry:A243517}\\

70&\oeis{A243518}&Empirical word recurrence of order 8&P&\pageref{round10:r10entry:A243518}\\

71&\oeis{A243600}&Empirical word recurrence of order 5&P&\pageref{round10:r10entry:A243600}\\

72&\oeis{A243601}&Empirical word recurrence of order 7&P&\pageref{round10:r10entry:A243601}\\

73&\oeis{A243602}&Empirical word recurrence of order 9&P&\pageref{round10:r10entry:A243602}\\

74&\oeis{A243603}&Empirical word recurrence of order 11&P&\pageref{round10:r10entry:A243603}\\

75&\oeis{A243604}&Empirical word recurrence of order 13&P&\pageref{round10:r10entry:A243604}\\

76&\oeis{A243605}&Empirical word recurrence of order 15&P&\pageref{round10:r10entry:A243605}\\

77&\oeis{A243606}&Empirical word recurrence of order 17&P&\pageref{round10:r10entry:A243606}\\

78&\oeis{A243635}&Empirical word recurrence of order 6&P&\pageref{round10:r10entry:A243635}\\

79&\oeis{A243636}&Empirical word recurrence of order 8&P&\pageref{round10:r10entry:A243636}\\

80&\oeis{A243637}&Empirical word recurrence of order 10&P&\pageref{round10:r10entry:A243637}\\

81&\oeis{A243638}&Empirical word recurrence of order 12&P&\pageref{round10:r10entry:A243638}\\

82&\oeis{A243640}&Empirical word recurrence of order 16&P&\pageref{round10:r10entry:A243640}\\

83&\oeis{A243724}&Empirical word recurrence of order 7&P&\pageref{round10:r10entry:A243724}\\

84&\oeis{A243727}&Empirical word recurrence of order 13&P&\pageref{round10:r10entry:A243727}\\

85&\oeis{A244177}&Empirical word recurrence of order 6&P&\pageref{round10:r10entry:A244177}\\

86&\oeis{A244178}&Empirical word recurrence of order 12&P&\pageref{round10:r10entry:A244178}\\

87&\oeis{A244179}&Empirical word recurrence of order 18&P&\pageref{round10:r10entry:A244179}\\

88&\oeis{A244180}&Empirical word recurrence of order 24&P&\pageref{round10:r10entry:A244180}\\

89&\oeis{A244181}&Empirical word recurrence of order 30&P&\pageref{round10:r10entry:A244181}\\

90&\oeis{A244182}&Empirical word recurrence of order 36&P&\pageref{round10:r10entry:A244182}\\

91&\oeis{A244183}&Empirical word recurrence of order 42&P&\pageref{round10:r10entry:A244183}\\

92&\oeis{A244184}&Empirical word recurrence of order 48&P&\pageref{round10:r10entry:A244184}\\

93&\oeis{A244694}&Empirical word recurrence of order 7&P&\pageref{round10:r10entry:A244694}\\

94&\oeis{A244695}&Empirical word recurrence of order 13&P&\pageref{round10:r10entry:A244695}\\

95&\oeis{A244696}&Empirical word recurrence of order 19&P&\pageref{round10:r10entry:A244696}\\

96&\oeis{A244697}&Empirical word recurrence of order 25&P&\pageref{round10:r10entry:A244697}\\

97&\oeis{A244698}&Empirical word recurrence of order 31&P&\pageref{round10:r10entry:A244698}\\

98&\oeis{A244699}&Empirical word recurrence of order 37&P&\pageref{round10:r10entry:A244699}\\

99&\oeis{A244700}&Empirical word recurrence of order 43&P&\pageref{round10:r10entry:A244700}\\

100&\oeis{A244701}&Empirical word recurrence of order 49&P&\pageref{round10:r10entry:A244701}\\

\end{longtable}
\normalsize
\clearpage

\section{Formal series and universal congruences}
\label{round10:r10sec:root}

This section proves ten precisely selected assertions marked conjectural
in the consulted primary entries. Definitions and unmarked companion formulas
are credited as existing inputs. The proofs are over the full stated domains;
the coefficient checks in \texttt{r10\_root\_verify.py} are auxiliary checks of
the definitions and their transcription. In particular, a proof of a parity
clause is not a claim to settle every other conjecture on the same entry.

\subsection{A373312: parity after inversion of a binary product}
\label{round10:r10entry:A373312}
\textbf{Source.} \oeis{A373312}, Hanna, June 25, 2024, defines
\[
 A=x+O(x^2),\qquad A(x)^2=A\left(\frac{xA(x)}{(1-A(x))^2}\right),
\]
and conjectures that $a_n$ is odd precisely when $n=2^j-1$, $j\ge1$.
The product and reversion characterizations are also displayed in the entry.

\begin{proposition}
The defining equation has a unique solution in $x+x^2\mathbb Z[[x]]$, and
\[
 A(x)\equiv\sum_{j\ge1}x^{2^j-1}\pmod2.
\]
\end{proposition}
\begin{proof}
Let $R$ be the compositional inverse of $A$. Substitution of $t=A(x)$ gives
\[
 R(t^2)=\frac{tR(t)}{(1-t)^2}.
\]
Writing $R(t)=tP(t)$, with $P(0)=1$, yields
$P(t)=(1-t)^2P(t^2)$. Iterating this identity coefficientwise proves
\[
 R(t)=t\prod_{j\ge0}(1-t^{2^j})^2.
\]
Conversely, this product satisfies the equation for $R$ and its inverse
satisfies the original equation for $A$. All coefficients of $R$ are integral
and its linear coefficient is one, so coefficient recursion proves integral
existence and uniqueness of its inverse.

Over $\mathbb F_2$, uniqueness of binary expansions gives
$\prod_{j\ge0}(1+t^{2^j})=1/(1+t)$. Hence
$R(t)=t/(1+t)^2$ in this field. Its inverse obeys
$A=x(1+A)^2=x+xA^2$. Starting with its zero constant term and comparing
successive coefficients gives exactly $A=\sum_{j\ge1}x^{2^j-1}$.
\end{proof}
The companion entries \oeis{A106407} and \oeis{A372530} describe the binary
product and a related inverse. They are known constructions, not additional
selected packages. Direct retrieval of the further companion A373313 failed;
no assertion about that entry is needed here.

\subsection{A320949 and A325286: defining rationally summed moments}

The two source sums below are not locally finite as ordinary formal sums in
$x$: their summands already have nonzero constant terms. Their precise meaning
is obtained by summing each polynomial moment as a rational function first.
The following elementary construction also proves that reduction modulo three
is legitimate.

For an integer-valued polynomial $b(r)$ of degree at most $d$ on nonnegative
integers, Newton expansion gives
\[
 b(r)=\sum_{j=0}^d c_j\binom rj,\qquad c_j=\Delta^j b(0)\in\mathbb Z.
\]
Consequently, its ordinary moment series is the rational function
\begin{equation}\label{round10:r10eq:newton}
 \sum_{r\ge0}b(r)t^r
 =\sum_{j=0}^d c_j\frac{t^j}{(1-t)^{j+1}}.
\end{equation}
Both identities follow by finite differences and
$\sum_{r\ge0}\binom rj t^r=t^j/(1-t)^{j+1}$.

\begin{lemma}[Odd-period reflection]\label{round10:r10lem:reflection}
For $k\ge1$, let $M_k(t)$ be the rational function associated by
\eqref{round10:r10eq:newton} with $b(r)=\binom{r^2}{k}$. Its reduction over
$\mathbb F_3$ is defined at $t=-1$, and $M_k(-1)=0$ there.
\end{lemma}
\begin{proof}
Choose $L=3^h>k$. The coefficient of $z^k$ in $(1+z)^v$ depends, modulo
three, only on $v\bmod L$, because $(1+z)^L=1+z^L$ in characteristic three.
Thus $b(r)$ has period $L$ modulo three and $b(L-r)=b(r)$ for $1\le r<L$.
As rational functions over $\mathbb F_3$,
\[
 M_k(t)=\frac{\sum_{r=0}^{L-1}b(r)t^r}{1-t^L}.
\]
At $t=-1$, the denominator is $2\ne0$. The term $r=0$ vanishes for $k>0$;
the terms $r$ and $L-r$ cancel because $L$ is odd. This proves the lemma.
The evaluation is also allowed in \eqref{round10:r10eq:newton}, whose denominators
are powers of $1-t$.
\end{proof}

\subsubsection{A320949: every positive coefficient is a multiple of three}
\label{round10:r10entry:A320949}
\textbf{Source.} \oeis{A320949}, Hanna, November 13, 2018, defines $A(0)=1$
by
\[
 \sum_{r\ge0}\frac{(1+x)^{r^2}}{(1+A(x))^{r+1}}=1
\]
and marks $a_n\equiv0\pmod3$, $n>0$, as a conjecture.

\begin{proposition}
The rational-moment interpretation of the defining equation has a unique
integral solution, and $A(x)\equiv1\pmod3$.
\end{proposition}
\begin{proof}
Let $c_{k,j}=\Delta^j\binom{r^2}{k}|_{r=0}$, with $0\le j\le2k$.
Put $t=(1+A)^{-1}$ in \eqref{round10:r10eq:newton}, and multiply by $t$.
The defining equation becomes
\[
 1=\sum_{k\ge0}x^k\sum_{j=0}^{2k}c_{k,j}A^{-j-1},
 \qquad
 A=1+\sum_{k\ge1}x^k\sum_{j=0}^{2k}c_{k,j}A^{-j}.
\]
The second expression is now locally finite. In degree $n>0$ its right side
depends only on coefficients of $A$ below degree $n$, because every summand
has a factor $x^k$, $k\ge1$. This is a monic recursion over $\mathbb Z$,
and over any commutative coefficient ring, with constant term one.

In $\mathbb F_3$, try $A=1$, for which $t=1/2=-1$. Every positive-degree
moment is zero by Lemma~\ref{round10:r10lem:reflection}. The constant moment is
$t/(1-t)=1$ at this value. Thus the constant series $1$ satisfies the
reduced defining equation. Uniqueness of the monic recursion proves the claim.
\end{proof}

\subsubsection{A325286: the complete ternary-digit rule}
\label{round10:r10entry:A325286}
\textbf{Source.} \oeis{A325286}, Hanna, April 18, 2019, defines $A(0)=1$ by
\[
 \frac12\sum_{r\ge0}\frac{(1+x)^{r(r-1)/2}}{(2A(x))^r}=1.
\]
Its conjecture assigns residue one to indices containing only zeroes and
ones in base three, and residue zero to every other index.

\begin{proposition}
There is a unique integral solution and
\[
 A(x)\equiv\prod_{j\ge0}(1+x^{3^j})\pmod3.
\]
\end{proposition}
\begin{proof}
Set $c_{k,j}=\Delta^j\binom{r(r-1)/2}{k}|_{r=0}$ for $0\le j\le2k$.
Taking $t=(2A)^{-1}$ in \eqref{round10:r10eq:newton} gives
\[
 1=A\sum_{k\ge0}x^k\sum_{j=0}^{2k}
          \frac{c_{k,j}}{(2A-1)^{j+1}},
 \qquad
 A-1=A\sum_{k\ge1}x^k\sum_{j=0}^{2k}
          \frac{c_{k,j}}{(2A-1)^j}.
\]
Since $2A-1$ has constant term one, the second equation gives a monic integral
recursion exactly as before, including uniqueness after reduction modulo three.

Work in $\mathbb F_3[[u]]$ and substitute $x=2u+u^2$ and
$A=(1+u)^{-1}$. The change of variable has invertible linear coefficient.
The original rational-moment expression becomes
\[
 \frac12\sum_{r\ge0}\frac{(1+u)^{r^2}}{2^r}.
\]
This substitution is legitimate coefficientwise: before reduction every
coefficient is a polynomial in $r$ times $2^{-r}$, summed by
\eqref{round10:r10eq:newton}. Its positive-degree coefficients are zero modulo three
by Lemma~\ref{round10:r10lem:reflection}, and its constant term is one.
Uniqueness therefore gives $A(x)^2(1+x)=1$ over $\mathbb F_3$.

Let $D(x)=\prod_{j\ge0}(1+x^{3^j})$. Frobenius gives
$D=(1+x)D^3$, so $D^2(1+x)=1$. A square root with constant term one is
unique here, since two such roots have sum with invertible constant term two.
Hence $A=D$. Unique ternary expansion of the selected exponents proves both
residue cases, including $a_0=1$.
\end{proof}
The related moment equation at \oeis{A303290} is an existing construction;
no additional claim for that entry is counted.

\subsection{A205572: the full family of powers-of-two periods}
\label{round10:r10entry:A205572}
\textbf{Source.} \oeis{A205572}, Hanna, January 29, 2012, defines
$F(x)=1/(\cos x-\sinh x)=\sum a_nx^n/n!$ and conjectures the congruence
below for every indicated pair of indices.

\begin{proposition}
For all $r\ge1$ and $k\ge0$,
\[
 a_{k+2^r}\equiv a_k\pmod{2^r}.
\]
This proves that $2^r$ is a period; it does not assert that it is the least
period.
\end{proposition}
\begin{proof}
Use the graded integral ring
\[
 R=\mathbb Z[C,S,H,K,(C-H)^{-1}],\qquad
 DC=-S,\quad DS=C,\quad DH=K,\quad DK=H.
\]
Each of $C,S,H,K$ has degree one and $(C-H)^{-1}$ has degree minus one.
The derivation $D$ preserves degree and extends to the localization by the
quotient rule. Modulo two, $D^2$ is a derivation: the middle term in the
second-derivative product formula vanishes. On the four generators it equals
the Euler derivation $E=C\partial_C+S\partial_S+H\partial_H+K\partial_K$.
These derivations therefore agree on the localization as well. On its
homogeneous degree-minus-one component $V$, $E$ is multiplication by $-1$,
which is the identity modulo two. Thus $D^2-I$ maps $V$ into $2V$.

The ring is torsion-free and the kernel of reduction modulo two is exactly
$2R$. Consequently $B=(D^2-I)/2$ is a well-defined integral linear
endomorphism of $V$. For $r\ge1$,
\[
 D^{2^r}=(I+2B)^{2^{r-1}}\equiv I\pmod{2^r}
 \quad\hbox{on }V.
\]
Indeed, for $1\le j\le2^{r-1}$,
\[
 v_2\left(2^j\binom{2^{r-1}}j\right)
 \ge j+r-1-v_2(j)\ge r.
\]
The first inequality follows from
$j\binom Nj=N\binom{N-1}{j-1}$.

Take $f=(C-H)^{-1}\in V$, apply $D^k$, and evaluate at
$(C,S,H,K)=(1,0,0,1)$. This evaluation is integral because the localized
denominator is one. Under $(C,S,H,K)=(\cos x,\sin x,\sinh x,\cosh x)$,
the derivation is ordinary differentiation, so the evaluated values are
precisely $a_{k+2^r}$ and $a_k$. This proves the assertion for every $k,r$.
\end{proof}

\subsection{A227051: every positive term ends in three}
\label{round10:r10entry:A227051}
\textbf{Source.} \oeis{A227051}, Hanna, June 29, 2013, defines
\[
 F(x)=\frac1{\cos x-\sin x\sinh x}
     =\sum_{m\ge0}b_m\frac{x^m}{m!},\qquad a_n=b_{2n},
\]
and conjectures $a_n\equiv3\pmod{10}$ for $n>0$.

\begin{proposition}
For every $n\ge1$, $a_n\equiv3\pmod{10}$.
\end{proposition}
\begin{proof}
The denominator has integral exponential coefficients and constant term one.
Its reciprocal has integral exponential coefficients by the recursion
$b_m=-\sum_{j=1}^m\binom mj g_jb_{m-j}$, where $g_j$ are the denominator's
exponential coefficients.

First work modulo five in
$\mathbb F_5[C,S,H,K,(C-SH)^{-1}]$ with the preceding derivation $D$.
In characteristic five, $D^5$ is a derivation. On each of $C,S,H,K$,
$D^5=D$ because $D^4$ is the identity on these generators. The derivations
therefore agree on the localized ring. For $f=(C-SH)^{-1}$ this gives
$D^{m+4}f=D^mf$ whenever $m\ge1$. Evaluation at $(1,0,0,1)$ implies
$b_{m+4}\equiv b_m\pmod5$ for $m\ge1$. The initial values
$b_2=3$ and $b_4=53$ give $b_{2n}\equiv3\pmod5$ for all $n>0$.

For the modulo-two part, use the integral Hurwitz-series ring: the product
of exponential coefficient sequences is
$(uv)_m=\sum_{j=0}^m\binom mj u_jv_{m-j}$. Modulo two, the square of every
zero-constant series is zero. Terms with $j\ne m-j$ cancel in pairs, and
the middle coefficient $\binom{2j}{j}$ is even for $j>0$.
Since $\sin x\equiv\sinh x$ in this ring, their product is zero.
Also $(\cos x)^2=1$ modulo two. Thus $F\equiv\cos x$ in the Hurwitz ring,
and every even-index exponential coefficient is one modulo two.
Combining residues modulo two and five proves the claim.
\end{proof}
The different square-root EGF at \oeis{A185071} has its own unmarked
congruence. That statement is not used to infer this result and is not counted.

\subsection{A292177: a justified bilateral limit and prime-index zeroes}
\label{round10:r10entry:A292177}
\textbf{Source.} \oeis{A292177}, Hanna, September 10, 2017, defines
\[
 S_K(x)=\sum_{n\in\mathbb Z}x^{n-K}
       \left(1-x^n+\frac{n(n+1)}6x^{n+K}\right)^n,
 \qquad A(x)=\lim_{K\to\infty}S_K(x),
\]
and conjectures $a_p=0$ for odd primes $p$.

\begin{proposition}
The limit exists coefficientwise in $\mathbb Q[[x]]$ and equals the
integral series
\begin{align*}
 A(x)={}&\sum_{d\ge1}\frac{d^2(d+1)}6x^{2d}(1-x^d)^{d-1}\\
 &+\sum_{d\ge2}\frac{(-1)^dd^2(d-1)}6
                  \frac{x^{d(d-1)}}{(1-x^d)^{d+1}}.
\end{align*}
In particular, $a_p=0$ for every odd prime $p$.
\end{proposition}
\begin{proof}
First establish the bilateral base identity
\[
 \sum_{n\in\mathbb Z}x^n(1-x^n)^n=0.
\]
For $n=-m<0$, rewrite the summand as
$(-1)^mx^{m^2-m}(1-x^m)^{-m}$. These summands are now locally finite;
their constant term at $m=1$ cancels the $n=0$ term. For a positive degree
$N=dr$, the contribution of the positive index $d$ is
$(-1)^{r-1}\binom d{r-1}$, and that of the negative index $-r$ is
$(-1)^r\binom d{r-1}$. Summing over divisors proves cancellation in every
degree, with binomial coefficients outside their usual range interpreted
as zero.

Expand $S_K$ in powers $j$ of the perturbation
$n(n+1)x^{n+K}/6$. The $j=0$ part is $x^{-K}$ times the preceding zero
series. The $j=1$ part is independent of $K$ and equals
\[
 \sum_{n\in\mathbb Z}\frac{n^2(n+1)}6
              x^{2n}(1-x^n)^{n-1}.
\]
For $n\ge1$, a nonzero $j\ge2$ term has valuation at least
$(j-1)K+(j+1)n\ge K+3$. For $n=-m<0$, rewriting its negative powers gives
valuation at least $m^2-m+(j-1)K$. The case $m=1$ has zero perturbation;
for $m\ge2$ and $j\ge2$ this is at least $K+2$. For each fixed $K\ge1$
the rewritings define a locally finite Laurent series, and the bounds are
uniform in the indices. The higher perturbations therefore disappear
coefficientwise as $K$ tends to infinity. Rewriting the negative indices
in the linear term gives the displayed two unilateral sums.

For completeness their coefficients also have the paired-divisor formula
\begin{equation}\label{round10:r10eq:bilateraldivisor}
 a_N=\sum_{dr=N}(-1)^r
       \frac{r(r-1)(d+r)}6\binom{d+1}{r}\quad(N\ge1).
\end{equation}
This follows by pairing the positive index $d$ with the negative index
$-r$ in the two unilateral sums. Every summand is integral. Divisibility
by two follows from $r(r-1)$. For divisibility by three, only
$r\equiv2\pmod3$ needs attention. If $d\equiv1\pmod3$, then $d+r$ is
a multiple of three. Otherwise the units ternary digit of $d+1$ is zero
or one, less than the units digit two of $r$, so
$\binom{d+1}{r}\equiv0\pmod3$. This last fact follows directly by
expanding $(1+z)^{d+1}$ using $(1+z)^3=1+z^3$ modulo three.

Finally, a coefficient in either unilateral sum can occur only at a
multiple of $d$. For an odd prime degree $p$, the only possible indices
are $d=1,p$. The positive $d=1$ summand is just $x^2/3$ and the negative
$d=1$ summand is absent. The $d=p$ summands start at $2p$ and $p(p-1)$,
both greater than $p$. Hence $a_p=0$.
\end{proof}
The base bilateral identity and several different weighted identities also
appear at \oeis{A291937}; they are known context. Here the perturbation limit
and the exact selected prime-index assertion are proved explicitly.

\subsection{A382124: integrality before Frobenius}
\label{round10:r10entry:A382124}
\textbf{Source.} \oeis{A382124}, Hanna, April 6, 2025, defines
\[
 L_n=\frac{\sigma(n)\sigma(2n)}3,\qquad
 A(x)=\exp\left(\sum_{n\ge1}\frac{L_n}{n}x^n\right).
\]
The companion \oeis{A382125} is $B(x)=A(x)^3$. The selected conjecture is
$[x^n]A\equiv[x^{3n}]B\pmod3$ for all $n\ge0$.

\begin{proposition}
The coefficients of $A$ are integers and $B(x)\equiv A(x^3)\pmod3$.
\end{proposition}
\begin{proof}
The prime-power formulas for $L_n$ are known at \oeis{A382123}; they also
follow immediately from multiplicativity of the divisor sum. For odd primes
$p$ and $r\ge1$,
\[
 L_{p^r}=\sigma(p^r)^2,
 \qquad L_{2^r}=C(r)=\frac{(2^{r+1}-1)(2^{r+2}-1)}3,
 \qquad L_1=C(0)=1.
\]
Define the Euler exponents
$e_n=n^{-1}\sum_{d\mid n}\mu(d)L_{n/d}$. They are multiplicative. On odd
prime powers,
\[
 e_{p^r}=\frac{\sigma(p^r)^2-\sigma(p^{r-1})^2}{p^r}
         =\sigma(p^r)+\sigma(p^{r-1})\in\mathbb Z.
\]
For powers of two,
$C(r)-C(r-1)=2^r(2^{r+1}-1)$, so $e_{2^r}=2^{r+1}-1$; also $e_1=1$.
Thus every Euler exponent is a nonnegative integer. Taking formal logarithms
and using $L_n=\sum_{d\mid n}d e_d$ proves
\[
 A(x)=\prod_{n\ge1}(1-x^n)^{-e_n}\in\mathbb Z[[x]].
\]
Frobenius may now be applied in $\mathbb F_3[[x]]$ to give
$A(x)^3=A(x^3)$. Extracting degree $3n$ proves the selected assertion,
including $n=0$.
\end{proof}
The zero residues of the other two index classes in A382125 follow at once
but are uncounted consequences. Its duplicate statement of the same cube
congruence does not supply a second package.

\subsection{A338328, A337758 and A230218: integral diagonal recursions}

The same standard cyclic-action argument justifies the divisions needed in
three different source definitions. The entries and selected assertions are
distinct; shared lemmas do not turn them into three independent new methods.

\begin{lemma}[Coprime diagonal division]\label{round10:r10lem:necklace}
Let $n,q\ge1$ with $\gcd(n,q)=1$, and let
$F=1+\sum_{i\ge1}f_ix^i$. Then
\[
 \frac1q[x^n]F(x)^q
\]
is a polynomial with integer coefficients in $f_1,\ldots,f_n$. Its coefficient
of $f_n$ is one and it is affine in $f_n$.
\end{lemma}
\begin{proof}
Expand $F^q$ using ordered $q$-tuples of nonnegative indices with total
weight $n$. The cyclic group of order $q$ acts by rotation. If a tuple were
a repetition of $d>1$ identical blocks, then $d$ would divide both $q$ and
$n$, a contradiction. Every orbit therefore has $q$ elements, all contributing
the same monomial. Dividing orbit sums by $q$ proves integrality as a
polynomial identity, so it remains valid for arbitrary signed integer
specializations. The terms involving $f_n$ have exactly one nonzero index
and together give $qf_n$ before division.
\end{proof}

Write
\[
 S(x)=\sum_{j\ge0}x^{2^j}\in\mathbb F_2[[x]],\qquad S+S^2=x.
\]
We will use the coefficient identity
\begin{equation}\label{round10:r10eq:residue}
 [x^n]H(S(x))=[t^n]H(t)(1+t)^{-n-1}.
\end{equation}
Indeed, substitute $x=t+t^2$ in the formal residue giving the left side.
The derivative is one in characteristic two, and
$(t+t^2)^{-n-1}=t^{-n-1}(1+t)^{-n-1}$. This is a valid change of formal
variable because its linear coefficient is one.

\subsubsection{A338328: odd coefficients exactly at powers of two}
\label{round10:r10entry:A338328}
\textbf{Source.} \oeis{A338328}, Hanna, October 24, 2020, defines $A(0)=0$ by
\[
 [x^n](1+nx-A(x))^{2n-1}=0\qquad(n>0).
\]
We select only its conjectural parity statement. Its separate modulo-four
assertion and its asymptotic formula are outside this package.

\begin{proposition}
The defining recursion is integral and $A(x)\equiv S(x)\pmod2$.
\end{proposition}
\begin{proof}
Since $\gcd(n,2n-1)=1$, Lemma~\ref{round10:r10lem:necklace} makes the divided
coefficient equation monic with coefficient $-1$ in the new coefficient
$a_n$. It gives a unique integral solution, and a unique solution of the
divided equations modulo two.

Test $A=S$. If $n$ is odd, the reduced base is
$1+x+S=1+S^2$, which contains only even degrees; hence its odd-degree
coefficient $n$ is zero. If $n$ is even, $n\ge2$, the reduced coefficient is
\[
 [x^n](1+S)^{2n-1}=[t^n](1+t)^{n-2}=0
\]
by \eqref{round10:r10eq:residue}. The exponent $2n-1$ is odd, so the divided
equation also vanishes modulo two. Monic uniqueness proves $A=S$.
\end{proof}

\subsubsection{A337758: handling the even exponents correctly}
\label{round10:r10entry:A337758}
\textbf{Source.} \oeis{A337758}, Hanna, October 24, 2020, defines $A(0)=0$ by
\[
 [x^n](1+nx+nx^2-A(x))^{n+1}=0\qquad(n>0).
\]
Only the parity conjecture is selected; the separate modulo-four clause
is not asserted here.

\begin{proposition}
The coefficients are integral and $A(x)\equiv S(x)\pmod2$.
\end{proposition}
\begin{proof}
Here $\gcd(n,n+1)=1$, so Lemma~\ref{round10:r10lem:necklace} again gives a monic
integer recursion. Let $S_0=\sum_{j\ge0}x^{2^j}$ denote an integral lift
of $S$, and put $F=1+nx+nx^2-S_0$ and $q=n+1$. It suffices to show that
the integer $q^{-1}[x^n]F^q$ is even.

If $n$ is even, $q$ is odd and the reduced base is $1+S$. Then
\[
 [x^n](1+S)^{n+1}=[t^n]1=0.
\]
If $n$ is odd, division by $q$ directly modulo two would be invalid.
Instead use the exact derivative identity
\[
 n\left(\frac1q[x^n]F^q\right)=[x^{n-1}]F'F^{q-1}.
\]
The quantity in parentheses is integral by the lemma. Also
$F'=n+2nx-S_0'\equiv1-1=0\pmod2$. As $n$ is odd, the quantity in
parentheses must be even. Thus $S$ satisfies all reduced divided equations,
and monic uniqueness completes the proof.
\end{proof}

\subsubsection{A230218: a missing binary digit in a square}
\label{round10:r10entry:A230218}
\textbf{Source.} \oeis{A230218}, Hanna, October 11, 2013, has
$a_0=a_1=1$ and
\[
 [x^n]A(x)^{n^2+n+1}=0\qquad(n>1).
\]
The selected conjecture says that, for $n>0$, $a_n$ is odd precisely at
powers of two.

\begin{proposition}
The recursion is integral and $A(x)\equiv1+S(x)\pmod2$.
\end{proposition}
\begin{proof}
For $q=n^2+n+1$ one has $\gcd(n,q)=1$ and $q$ is odd.
Lemma~\ref{round10:r10lem:necklace} gives integral monic uniqueness from the specified
initial coefficients. By \eqref{round10:r10eq:residue}, the candidate $1+S$ gives
\[
 [x^n](1+S)^{n^2+n+1}=[t^n](1+t)^{n^2}=\binom{n^2}{n}\pmod2.
\]
This binomial coefficient is even for every $n>1$. If $n$ is even and its
least nonzero binary digit is at position $r\ge1$, then $2^{2r}\mid n^2$,
so that digit is missing from $n^2$. If $n>1$ is odd, let $r\ge1$ be its
first nonzero binary digit above the units digit. Then
$n\equiv1+2^r\pmod{2^{r+1}}$ but $n^2\equiv1\pmod{2^{r+1}}$, so again
digit $r$ is missing from the square. The binary factorization of
$(1+t)^{n^2}$ in characteristic two shows that its coefficient of $t^n$
is zero whenever such a digit is missing. Thus all defining equations hold
for $1+S$, and uniqueness proves the result.
\end{proof}
The inspected companion entries A185072, A229041, A229044, A171791 and
A292877 concern known reversion relations or other parity patterns. None
is counted as a further result of this package.

\subsection*{Scope of the source and prior-work audit}
All ten selected A-numbers are disjoint from the previous 908 selected
A-numbers, the retained uncounted earlier reserves, and the 97 unique
A-numbers in the inspected OEIS Open Appendix A. The audit includes the
primary entries, the named companions above and targeted searches. It is a
bounded source audit: cached retrievals and unsuccessful companion retrievals
cannot certify worldwide historical priority or the latest live edit state.
Known tools used here include finite differences, formal inversion, cyclic
group actions, formal residues, Euler products, and elementary congruences.
The count records ten source-marked assertion packages with written proofs,
not ten new general methods or an official change to OEIS.

\section{Arithmetic identities, a divisor formula, and a valuation correction}
\label{round10:r10sec:arithmetic}

These five packages concern five distinct entries absent from the previous
908 selected identifiers and from the supplied 97-identifier OEIS Open
Appendix A baseline. Four selected assertions or questions are resolved
affirmatively; one empirical valuation clause is corrected. The primary
entries and relevant companion records were inspected on October 20, 2026.
The source audit is bounded: it does not establish historical priority.
Every theorem below has a proof valid on its stated infinite domain.
The finite computations in the accompanying standard-library verifier are
independent checks of definitions and boundary cases.

\subsection{A165886: the quotient of consecutive normalized Hirzebruch numbers}
\label{round10:r10entry:A165886}
\textbf{Source and selected assertion.} The comment on \oeis{A165886}
marks the quotient identity with \oeis{A141459} and \oeis{A006519}
as a conjecture. The entry was submitted by Paul Curtz on September 29,
2009; the comment is not separately dated. R.~J. Mathar records a finite
check through 2000 on July 4, 2011. Define
\[
 H(n)=\prod_{p\ {\rm prime}}p^{\lfloor n/(p-1)\rfloor},\qquad
 F(n)=\frac{H(n)}{2^{s_2(n)}}\quad(n\ge0),
\]
where $s_2(n)$ counts the ones in the binary expansion of $n$.
The product is finite, since primes greater than $n+1$ have exponent zero.
These are the definitions of \oeis{A091137} and \oeis{A165641};
\oeis{A165886} is $a(n)=F(n+1)/F(n)$.

\begin{proposition}
For every integer $n\ge0$,
\[
 a(n)=2^{v_2(n+1)}
       \prod_{\substack{p\ {\rm odd\ prime}\\p-1\mid n+1}}p
     =\mathrm{A006519}(n+1)\,\mathrm{A141459}(n+1).
\]
\end{proposition}
\begin{proof}
For any positive integer $d$,
\[
 \left\lfloor\frac{n+1}{d}\right\rfloor-
 \left\lfloor\frac n d\right\rfloor
 =\begin{cases}1,&d\mid n+1,\\0,&d\nmid n+1.\end{cases}
\]
Consequently
\[
 \frac{H(n+1)}{H(n)}
  =2\prod_{\substack{p\ {\rm odd\ prime}\\p-1\mid n+1}}p.
\]
Put $t=v_2(n+1)$. Adding one to $n$ erases exactly $t$ trailing
binary ones and inserts one binary one, including the case $n=0$.
Thus $s_2(n+1)-s_2(n)=1-t$. Dividing the preceding product by
$2^{s_2(n+1)-s_2(n)}$ gives the assertion. The companion definitions
are precisely the odd-prime product and the greatest power of two
dividing the argument.
\end{proof}
\textbf{Prior-work scope.} The Hirzebruch prime product, its consecutive
quotient, and Gould's binary-digit formula are established inputs,
explicitly credited to \oeis{A091137}, \oeis{A001316},
\oeis{A141459}, and \oeis{A006519}. The selected result is their
combination in the still-marked assertion on \oeis{A165886}. A search
also located the same formula in the conjecture-list mirror and in
Sequence Machine; neither inspected result supplied a proof. No novelty
claim is made for the elementary floor or binary-carry identities.

\subsection{A243046: a finite formula using paired divisors}
\label{round10:r10entry:A243046}
\textbf{Source and selected question.} \oeis{A243046} counts positive
integers $k\ne n$ for which both $kn/(k+n)$ and $kn/(k-n)$ are
integers. Antti Karttunen's February 18, 2023 comment explicitly asks
for a direct formula. The following formula uses only the positive
divisors of $n^2$; it does not search over the original $k$ interval.

\begin{proposition}
Let $n\ge1$, let $D_n=\{d>0:d\mid n^2\}$, and let $[P]$ be
one or zero according as the proposition $P$ is true or false. Then
\[
 \boxed{\displaystyle
 a(n)=\sum_{d\in D_n}[d+2n\in D_n]
       +\sum_{\substack{d\in D_n\\d<n}}[2n-d\in D_n].}
 \tag{A243046}
\]
The full set of admissible values of $k$ is the disjoint union
\[
 \{n+d:d\in D_n,\ d+2n\in D_n\}
 \ \dot\cup\
 \{n-d:d\in D_n,\ d<n,\ 2n-d\in D_n\}.
\]
In particular, for a prime $p$ and integer $e\ge1$,
\[
 a(p^e)=
 \begin{cases}
 1,&p=3,\\
 1,&p=2\text{ and }e\ge2,\\
 0,&\text{otherwise}.
 \end{cases}
\]
\end{proposition}
\begin{proof}
For $k>0$ and $k\ne n$, the exact identities
\[
 \frac{kn}{k+n}=n-\frac{n^2}{k+n},\qquad
 \frac{kn}{k-n}=n+\frac{n^2}{k-n}
\]
show that admissibility is equivalent to $(k+n)\mid n^2$ and
$(k-n)\mid n^2$, with the second divisor allowed to be negative.
If $k>n$, put $d=k-n>0$. Then the two positive divisors are $d$
and $d+2n$, and $k=n+d$. This is a bijection with the first set.
If $0<k<n$, put $d=n-k$. Then $0<d<n$ and the two positive
divisors are $d$ and $2n-d$, with $k=n-d$. This is a bijection
with the second set. The ranges $k>n$ and $0<k<n$ are disjoint,
which proves both the formula and the stated solution set.

For the prime-power corollary, distinct divisors of $p^{2e}$ have the
form $p^b<p^c$. A pair contributing to the first sum must satisfy
\[
 p^b(p^{c-b}-1)=2p^e.
\]
If $p$ is odd, comparison of $p$-adic valuations gives $b=e$;
then $p^{c-b}=3$, forcing $p=3$ and $c=e+1$. This pair exists
for every $e\ge1$. If $p=2$, the odd factor forces $b=e+1$
and $c=b+1=e+2$, which lies at most $2e$ exactly when $e\ge2$.
A pair contributing to the second sum would satisfy
$p^b(p^{c-b}+1)=2p^e$ with $b<c$. For odd $p$ this would give
$b=e$ and $p^{c-b}+1=2$, which is impossible. For $p=2$, its
odd factor $2^{c-b}+1>1$ cannot divide a power of two. Thus the
second sum has no contribution for a prime power.
\end{proof}
\textbf{Effective range and scope.} If
$n=\prod_i p_i^{e_i}$, the set in the formula is explicitly
$D_n=\{\prod_i p_i^{f_i}:0\le f_i\le2e_i\}$, of size
$\prod_i(2e_i+1)=\tau(n^2)$. Once these divisors have been generated,
the answer requires fewer than $2\tau(n^2)$ set-membership checks.
Any admissible $k$ also satisfies $k\le n^2-n$, because $k+n$
is a positive divisor of $n^2$. This proves finiteness even for the
boundary $n=1$, when no $k$ exists. The inspected companions
\oeis{A063647} and \oeis{A146564} count either individual
integrality condition, while \oeis{A243017} and \oeis{A243045}
record the minimum and maximum. Umberto Cerruti's linked note
\emph{Percorsi tra i numeri}, Theorems 1--2 and Corollary 4,
proves the signed-divisor parametrization for the single minus
condition and the associated companion counts. Those results are
credited inputs, with the elementary equivalences reproved above.
The simultaneous divisor-pair formula is the answer selected here;
the companion counts are not additional packages.

\subsection{A308833: tetrahedral divisibility is governed by an odd composite}
\label{round10:r10entry:A308833}
\textbf{Source and selected assertion.} \oeis{A308833}, submitted by
Lorenzo Sauras Altuzarra on June 28, 2019, consists of the positive
integers $r$ for which $r(r+1)(r+2)/6$ divides $r!$. Its conjectural
comment, not separately dated, compares consecutive odd/even indices
with the compositeness of the following odd integer.

\begin{proposition}
For every odd integer $r>1$, the following statements are equivalent:
\[
 \frac{r(r+1)(r+2)}6\mid r!,\qquad
 \frac{(r+1)(r+2)(r+3)}6\mid(r+1)!,\qquad
 r+2\text{ is composite}.
\]
\end{proposition}
\begin{proof}
Write $m=r+2$, so $m$ is odd and $m\ge5$. If $m$ is prime,
it is coprime to 6. Each of the two tetrahedral numbers contains
a factor $m$, whereas neither $(m-2)!$ nor $(m-1)!$ does.
Both divisibility assertions are therefore false.

Suppose instead that $m$ is composite; then $m\ge9$. We first
prove $m\mid(m-3)!$. Choose a factorization $m=uv$ with
$3\le u\le v$. If $u<v$, both distinct factors lie at most
$m-3$ and their product occurs in the factorial. If $u=v$,
the distinct numbers $u$ and $2u$ lie at most $u^2-3=m-3$
for $u\ge3$, and their product is $2m$. This proves the claim
in both cases.

The even integer $m-1\ge8$ is the product of the distinct numbers
$2$ and $(m-1)/2$, both at most $m-3$. Hence $m-1$ also divides
$(m-3)!$. Since $\gcd(m,m-1)=1$, we obtain
\[
 \frac{(m-2)!}{(m-2)(m-1)m/6}
   =\frac{6(m-3)!}{m(m-1)}\in\mathbb Z.
\]
Likewise $m+1\ge10$ is the product of the distinct integers $2$
and $(m+1)/2$, both at most $m-2$. Together with
$m\mid(m-3)!$ and $\gcd(m,m+1)=1$, this gives
\[
 \frac{(m-1)!}{(m-1)m(m+1)/6}
   =\frac{6(m-2)!}{m(m+1)}\in\mathbb Z.
\]
Thus both divisibilities hold exactly in the composite case.
\end{proof}
\textbf{Prior-work scope.} The companion \oeis{A153238} describes
the associated odd composites. The inspected source and companion
provide definitions and computational tables, but no proof of this
selected three-way equivalence. The factorial arguments above are
elementary; neither that technique nor the companion classification
is claimed as new mathematics.

\subsection{A326478: the even Bernoulli-denominator subsequence}
\label{round10:r10entry:A326478}
\textbf{Source and selected assertion.} \oeis{A326478}, submitted by
Peter Luschny on July 16, 2019, defines
\[
 a(m)=m\frac{\operatorname{den}(mB_{m-1})}
                {\operatorname{den}(B_{m-1})}\quad(m\ge1),
\]
where $B_j$ is the $j$th Bernoulli number, $B_1=-1/2$, and
$\operatorname{den}(0)=1$. The even-index generating-function
comment is explicitly marked empirical; that comment is not separately
dated. Only this even-index assertion is selected.

\begin{proposition}
For every integer $n\ge1$,
\[
 a(2n)=\begin{cases}1,&n=1,\\2n,&n\ge2,\end{cases}
 \qquad
 \sum_{n\ge1}a(2n)x^n=x\left(\frac{2}{(1-x)^2}-1\right).
\]
With the currently stated offset of \oeis{A103517}, the correct
companion relation is
\[
 a(2n)=\mathrm{A103517}(n-1).
\]
\end{proposition}
\begin{proof}
The defining Bernoulli series satisfies
\[
 \frac{z}{e^z-1}+\frac z2
 =\frac z2\frac{e^z+1}{e^z-1},
\]
an even formal power series, since substitution of $-z$ leaves it
unchanged. Consequently $B_j=0$ for every odd $j>1$.
For $n\ge2$, both $B_{2n-1}$ and $2nB_{2n-1}$ are zero;
their reduced denominators are 1, so $a(2n)=2n$.
For $n=1$, the denominators of $B_1=-1/2$ and $2B_1=-1$
are 2 and 1, respectively, so $a(2)=1$. Therefore
\[
 \sum_{n\ge1}a(2n)x^n
   =\frac{2x}{(1-x)^2}-x.
\]
The source \oeis{A103517} has offset zero and generating function
$2/(1-x)^2-1$, so shifting its coefficients by one yields the
companion relation.
\end{proof}
\textbf{Exact scope and an offset correction.} The same source sentence
prints the comparison with $\mathrm{A103517}(n+1)$. This index is
incorrect with its present offset: at $n=1$, $a(2)=1$ but
$\mathrm{A103517}(2)=6$. The correction to $n-1$ is included in
this one package and is not counted separately. The further comparison
with \oeis{A272651}, a maximum in the no-three-in-line problem,
does not follow from a Bernoulli calculation and is not asserted here.
The odd-fixed-point question on \oeis{A326478} is also outside the
selected scope. The Bernoulli vanishing fact and the companion rational
generating function are classical inputs, not novelty claims.

\subsection{A339100: correcting the three-adic valuation clause}
\label{round10:r10entry:A339100}
\textbf{Source and selected clause.} \oeis{A339100}, submitted by
Luc Rousseau on November 23, 2020, uses the triangle \oeis{A214406}
with
\[
 T(0,0)=1,\qquad
 T(n,k)=(4n-2k-1)T(n-1,k-1)+(2k+1)T(n-1,k),
\]
and $T(n,k)=0$ outside $0\le k\le n$. It defines
\[
 a(n)=\gcd_{0\le k\le n}
       \big((2n-k)T(n,k)+(k+1)T(n,k+1)\big),\quad n\ge1.
\]
The empirical prime-valuation comment includes the clause
$v_3(a(n))=n\bmod2$. This particular clause is reversed.

\begin{proposition}
The literal three-adic clause is false already at $n=1$, where
$a(1)=\gcd(3,1)=1$. Its corrected version, valid for every $n\ge1$, is
\[
 \boxed{\displaystyle v_3(a(n))=
 \begin{cases}1,&n\text{ even},\\0,&n\text{ odd}.\end{cases}}
\]
\end{proposition}
\begin{proof}
Put $P_n(y)=\sum_{k=0}^nT(n,k)y^k$ and
\[
 R_n(x)=\frac{P_n(x^2)}{(1-x^2)^{2n}},\qquad
 L(f)=\frac{d}{dx}\left(\frac{xf}{1-x^2}\right).
\]
Summing the triangular recurrence gives
\[
 P_{n+1}(y)=2y(1-y)P_n'(y)+(1+(4n+1)y)P_n(y).
\]
Differentiation then shows $R_0=1$ and $R_{n+1}=L(R_n)$.
If $Q_n(y)$ denotes the polynomial whose coefficients occur in the
defining gcd, a further exact differentiation gives
\[
 R_n'(x)=\frac{2x}{(1-x^2)^{2n+1}}Q_n(x^2),\qquad
 Q_n(y)=(1-y)P_n'(y)+2nP_n(y).
\]
Reduction modulo 3 is legitimate in the rational-function field
$\mathbb F_3(x)$, since all displayed denominators remain nonzero.
In that field $3\mid a(n)$ is equivalent to $R_n'=0$.

The first two iterates are
\[
 R_1=\frac{1+x^2}{(1-x^2)^2},\qquad
 R_2=\frac{1+8x^2+3x^4}{(1-x^2)^4}
       \equiv(1-x^2)^{-3}=:H\pmod3.
\]
Here $H'=0$ in characteristic 3. Hence
$L(H^m f)=H^mL(f)$, and induction gives
\[
 R_{2m}=H^m,\qquad R_{2m+1}=H^mR_1
       \quad\text{in }\mathbb F_3(x),\quad m\ge0.
\]
Thus every even iterate has zero derivative, whereas
\[
 R_{2m+1}'=H^mR_1',\qquad
 R_1'=\frac{2x(3+x^2)}{(1-x^2)^3}
       \equiv\frac{2x^3}{(1-x^2)^3}\ne0\pmod3.
\]
It follows that 3 divides the defining gcd exactly when $n$ is even.

It remains to show that the even gcd has only one factor of 3.
The triangular recurrence gives $T(n,0)=1$. Setting $c_n=T(n,1)$,
it also gives $c_0=0$ and $c_n=3c_{n-1}+4n-3$; induction solves this as
\[
 c_n=\frac{3^{n+1}-4n-3}{2}.
\]
Therefore the constant coefficient of $Q_n$ is
\[
 Q_n(0)=2n+c_n=\frac{3(3^n-1)}2,
\]
whose three-adic valuation is exactly 1 for every $n\ge1$.
As $a(n)$ divides this coefficient, its valuation is at most 1.
Combined with the preceding divisibility criterion, this proves the
corrected formula.
\end{proof}
\textbf{Scope and prior work.} This is one correction package for the
three-adic clause alone. It proves neither the two-adic clause nor the
claims for primes at least 5 in the same empirical paragraph.
The recurrence and rational-function construction are credited to
\oeis{A214406} and \oeis{A339100}. The related quotient sequence
\oeis{A165886} is not identified with this gcd sequence:
$a(20)=3300$, whereas $\mathrm{A165886}(19)=660$. The source's
own displayed factorization table agrees with the corrected parity.
The bounded search and inspected companion notes supplied no prior
proof of the selected correction.

\section{Generating functions, dimension formulas, and exact recurrence certificates}
\label{round10:r10sec:gf}
This section contains twenty selected A-number packages: eighteen proofs of
source-marked assertions and two corrections to generating functions after
an offset change. The thirteen modular-form packages are direct applications
of classical dimension theorems. The continued-fraction and bargraph generating
functions are also credited to earlier work. A226302 includes an independent
proof of the generating-function conjecture in its defining paper.
One A-number contributes one package, even when several displayed claims are
addressed. These counts do not assert twenty independent research discoveries,
first publication, or acceptance of an edit by OEIS.

\subsection{Exact coefficient and operator certificates}
For $F(x)=\sum a_nx^n$, extend $a_n$ by zero below the stated offset and set
$\theta=x\,d/dx$. For polynomial coefficients $p_j$,
\begin{equation}\label{round10:r10gf:eq:operator}
 L=\sum_{j=0}^r x^jp_j(\theta+j),\qquad
 [x^n]LF=\sum_{j=0}^r p_j(n)a_{n-j}.
\end{equation}
This follows by applying $(\theta+j)$ to $x^{n-j}$. Therefore an exact
identity $LF=P$, with $P$ polynomial, proves the recurrence at every
$n>\deg P$ for which all source indices exist. Constant $p_j$ reduce this
certificate to multiplication by $Q(x)=\sum_jp_jx^j$.

For a radical, always take $\sqrt{D(0)}=1$. In the basis $1,\sqrt D$,
\[
 \theta(R_0+R_1\sqrt D)
 =xR'_0+x\left(R'_1+\frac{D'}{2D}R_1\right)\sqrt D.
\]
The displayed residuals below are exact rational identities, checked after
clearing denominators. Each proof first obtains the series from its independent
definition. The accompanying coefficient comparisons through index 100, and
comparisons with the published source terms, are supplementary checks of
indexing and transcription; they do not replace these identities.

\subsection{A classical input for thirteen dimension assertions}
For this subsection let $a_N(n)=\dim S_{2n}(\Gamma_0(N))$ for $n\ge1$, and
put $a_N(0)=0$ only as a generating-function convention. We use the classical
dimension theorem in William Stein, \emph{Modular Forms, a Computational
Approach}, Propositions 6.1 and 6.4
(\href{https://wstein.org/books/modform/modform/dimension_formulas.html}{Chapter 6}).
Define
\[
 \mu=N\prod_{p\mid N}(1+1/p),\quad
 c=\sum_{d\mid N}\varphi(\gcd(d,N/d)),\quad
 g=1+\frac\mu{12}-\frac{e_2}{4}-\frac{e_3}{3}-\frac c2,
\]
where $e_2$ and $e_3$ count roots modulo $N$ of $u^2+1$ and $u^2+u+1$,
respectively. The theorem gives
\begin{equation}\label{round10:r10gf:eq:dim}
 a_N(n)=(2n-1)(g-1)+(n-1)c
       +e_2\lfloor n/2\rfloor+e_3\lfloor2n/3\rfloor+[n=1].
\end{equation}
For newforms it gives
\begin{equation}\label{round10:r10gf:eq:new}
 a_N^{\rm new}(n)=\sum_{M\mid N}\overline\mu(N/M)a_M(n),
\end{equation}
where $\overline\mu$ is multiplicative with values $-2,1,0$ on
$p,p^2,p^e$ ($e\ge3$). These theorems are existing results.

We now sum \eqref{round10:r10gf:eq:dim}. Splitting indices by their residues modulo
2 or 3 gives
\[
 \sum_{n\ge1}\lfloor n/2\rfloor x^n=\frac{x^2}{(1-x)(1-x^2)},\qquad
 \sum_{n\ge1}\lfloor2n/3\rfloor x^n
 =\frac{x^2(1+x)}{(1-x)(1-x^3)}.
\]
Together with $\sum nx^n=x/(1-x)^2$, these prove
\begin{equation}\label{round10:r10gf:eq:dimGF}
 G_N(x)=x+\frac{(g-1)x(1+x)+cx^2}{(1-x)^2}
 +\frac{e_2x^2}{(1-x)(1-x^2)}
 +\frac{e_3x^2(1+x)}{(1-x)(1-x^3)}.
\end{equation}
The root counts and divisor sums are finite, exact computations. Substitution
in this identity proves each rational function below. Multiplication by its
stated $Q$ proves the recurrence on its full range. None of these series is
defined by the recurrence being proved. Closed forms follow from
\eqref{round10:r10gf:eq:dim}; the verifier checks the six symbolic identities
$n=6m+r$, $r=2,\ldots,7$, $m\ge0$, covering all $n\ge2$.

\subsection{A063081: cusp forms of level 13}
\label{round10:r10entry:A063081}

\noindent\textbf{Source and scope.} \href{https://oeis.org/A063081}{A063081}, checked 20 October 2026. Conjecture; Chai Wah Wu; Bill McEachen (2016-05-27; 2025-06-16).

\noindent\textbf{Indexing.} Offset $1$, with
\[ (a_{1},\ldots,a_{8})=(0,3,5,7,9,13,13,17). \]

The defining weight is $2n$. The selected ordinary generating function is

\[ G(x)=\frac{x^{2} \left(x^{4} + x^{3} + 4 x^{2} + 5 x + 3\right)}{x^{5} - x^{3} - x^{2} + 1}. \]

The selected recurrence, in coefficient-vector notation, is
\[ \sum_{j=0}^{5}p_j(n)a_{n-j}=0\qquad(n\ge 7). \]

\[ (p_0,\ldots,p_5)=(1,0,-1,-1,0,1). \]

The further selected assertion is $a_n=a_{n-6}+14$ for $n\ge8$.

\noindent\textbf{Proof.} This is a specialization of the credited dimension theorem.

The exact invariants are
\[ (\mu,e_2,e_3,c,g)=(14,2,2,2,0). \]

Substituting these values in \eqref{round10:r10gf:eq:dimGF} gives the stated $G$.

With $Q(x)=\sum_jp_jx^j$, direct rational simplification gives
\[ Q(x)G(x)=x^{6} + x^{5} + 4 x^{4} + 5 x^{3} + 3 x^{2}. \]

The residual has degree below $7$, so coefficient extraction proves the selected recurrence for every stated index.

In \eqref{round10:r10gf:eq:dim}, increasing $n\ge2$ by 6 increases the dimension by 14, proving the additional assertion.

\subsection{A063087: cusp forms of level 19}
\label{round10:r10entry:A063087}

\noindent\textbf{Source and scope.} \href{https://oeis.org/A063087}{A063087}, checked 20 October 2026. Conjecture; Wesley Ivan Hurt; Colin Barker (2016-12-27; 2019-06-09).

\noindent\textbf{Indexing.} Offset $1$, with
\[ (a_{1},\ldots,a_{8})=(1,4,8,10,14,18,20,24). \]

The defining weight is $2n$. The selected ordinary generating function is

\[ G(x)=\frac{x \left(x^{4} + x^{3} + 4 x^{2} + 3 x + 1\right)}{\left(1 - x\right)^{2} \left(x^{2} + x + 1\right)}. \]

The selected recurrence, in coefficient-vector notation, is
\[ \sum_{j=0}^{4}p_j(n)a_{n-j}=0\qquad(n\ge 6). \]

\[ (p_0,\ldots,p_4)=(1,-1,0,-1,1). \]

The selected closed form is $a_n=4 n - 2 \left\lfloor{\frac{n}{3} - \frac{1}{3}}\right\rfloor - 4$ for $n\ge2$.

\noindent\textbf{Proof.} This is a specialization of the credited dimension theorem.

The exact invariants are
\[ (\mu,e_2,e_3,c,g)=(20,0,2,2,1). \]

Substituting these values in \eqref{round10:r10gf:eq:dimGF} gives the stated $G$.

With $Q(x)=\sum_jp_jx^j$, direct rational simplification gives
\[ Q(x)G(x)=x^{5} + x^{4} + 4 x^{3} + 3 x^{2} + x. \]

The residual has degree below $6$, so coefficient extraction proves the selected recurrence for every stated index.

The closed form follows by substituting the same invariants in \eqref{round10:r10gf:eq:dim}; its exception at $n=1$ is retained in the displayed initial data.

\subsection{A063098: cusp forms of level 30}
\label{round10:r10entry:A063098}

\noindent\textbf{Source and scope.} \href{https://oeis.org/A063098}{A063098}, checked 20 October 2026. Conjecture; Colin Barker (2019-06-09).

\noindent\textbf{Indexing.} Offset $1$, with
\[ (a_{1},\ldots,a_{8})=(3,14,26,38,50,62,74,86). \]

The defining weight is $2n$. The selected ordinary generating function is

\[ G(x)=\frac{x \left(x^{2} + 8 x + 3\right)}{\left(1 - x\right)^{2}}. \]

The selected recurrence, in coefficient-vector notation, is
\[ \sum_{j=0}^{2}p_j(n)a_{n-j}=0\qquad(n\ge 4). \]

\[ (p_0,\ldots,p_2)=(1,-2,1). \]

The selected closed form is $a_n=12 n - 10$ for $n\ge2$.

\noindent\textbf{Proof.} This is a specialization of the credited dimension theorem.

The exact invariants are
\[ (\mu,e_2,e_3,c,g)=(72,0,0,8,3). \]

Substituting these values in \eqref{round10:r10gf:eq:dimGF} gives the stated $G$.

With $Q(x)=\sum_jp_jx^j$, direct rational simplification gives
\[ Q(x)G(x)=x^{3} + 8 x^{2} + 3 x. \]

The residual has degree below $4$, so coefficient extraction proves the selected recurrence for every stated index.

The closed form follows by substituting the same invariants in \eqref{round10:r10gf:eq:dim}; its exception at $n=1$ is retained in the displayed initial data.

\subsection{A063105: cusp forms of level 37}
\label{round10:r10entry:A063105}

\noindent\textbf{Source and scope.} \href{https://oeis.org/A063105}{A063105}, checked 20 October 2026. Conjecture; Colin Barker (2012-09-27).

\noindent\textbf{Indexing.} Offset $1$, with
\[ (a_{1},\ldots,a_{8})=(2,9,15,21,27,35,39,47). \]

The defining weight is $2n$. The selected ordinary generating function is

\[ G(x)=\frac{x \left(x^{5} + 3 x^{4} + 10 x^{3} + 13 x^{2} + 9 x + 2\right)}{\left(1 - x\right)^{2} \left(x + 1\right) \left(x^{2} + x + 1\right)}. \]

The selected recurrence, in coefficient-vector notation, is
\[ \sum_{j=0}^{5}p_j(n)a_{n-j}=0\qquad(n\ge 7). \]

\[ (p_0,\ldots,p_5)=(1,0,-1,-1,0,1). \]

\noindent\textbf{Proof.} This is a specialization of the credited dimension theorem.

The exact invariants are
\[ (\mu,e_2,e_3,c,g)=(38,2,2,2,2). \]

Substituting these values in \eqref{round10:r10gf:eq:dimGF} gives the stated $G$.

With $Q(x)=\sum_jp_jx^j$, direct rational simplification gives
\[ Q(x)G(x)=x^{6} + 3 x^{5} + 10 x^{4} + 13 x^{3} + 9 x^{2} + 2 x. \]

The residual has degree below $7$, so coefficient extraction proves the selected recurrence for every stated index.

\subsection{A063110: cusp forms of level 42}
\label{round10:r10entry:A063110}

\noindent\textbf{Source and scope.} \href{https://oeis.org/A063110}{A063110}, checked 20 October 2026. Conjecture; Colin Barker (2012-09-23).

\noindent\textbf{Indexing.} Offset $1$, with
\[ (a_{1},\ldots,a_{8})=(5,20,36,52,68,84,100,116). \]

The defining weight is $2n$. The selected ordinary generating function is

\[ G(x)=\frac{x \left(x^{2} + 10 x + 5\right)}{\left(1 - x\right)^{2}}. \]

The selected recurrence, in coefficient-vector notation, is
\[ \sum_{j=0}^{2}p_j(n)a_{n-j}=0\qquad(n\ge 4). \]

\[ (p_0,\ldots,p_2)=(1,-2,1). \]

The selected closed form is $a_n=16 n - 12$ for $n\ge2$.

\noindent\textbf{Proof.} This is a specialization of the credited dimension theorem.

The exact invariants are
\[ (\mu,e_2,e_3,c,g)=(96,0,0,8,5). \]

Substituting these values in \eqref{round10:r10gf:eq:dimGF} gives the stated $G$.

With $Q(x)=\sum_jp_jx^j$, direct rational simplification gives
\[ Q(x)G(x)=x^{3} + 10 x^{2} + 5 x. \]

The residual has degree below $4$, so coefficient extraction proves the selected recurrence for every stated index.

The closed form follows by substituting the same invariants in \eqref{round10:r10gf:eq:dim}; its exception at $n=1$ is retained in the displayed initial data.

\subsection{A063116: cusp forms of level 48}
\label{round10:r10entry:A063116}

\noindent\textbf{Source and scope.} \href{https://oeis.org/A063116}{A063116}, checked 20 October 2026. Conjecture; Colin Barker (2019-06-09).

\noindent\textbf{Indexing.} Offset $1$, with
\[ (a_{1},\ldots,a_{8})=(3,18,34,50,66,82,98,114). \]

The defining weight is $2n$. The selected ordinary generating function is

\[ G(x)=\frac{x \left(x^{2} + 12 x + 3\right)}{\left(1 - x\right)^{2}}. \]

The selected recurrence, in coefficient-vector notation, is
\[ \sum_{j=0}^{2}p_j(n)a_{n-j}=0\qquad(n\ge 4). \]

\[ (p_0,\ldots,p_2)=(1,-2,1). \]

The selected closed form is $a_n=16 n - 14$ for $n\ge2$.

\noindent\textbf{Proof.} This is a specialization of the credited dimension theorem.

The exact invariants are
\[ (\mu,e_2,e_3,c,g)=(96,0,0,12,3). \]

Substituting these values in \eqref{round10:r10gf:eq:dimGF} gives the stated $G$.

With $Q(x)=\sum_jp_jx^j$, direct rational simplification gives
\[ Q(x)G(x)=x^{3} + 12 x^{2} + 3 x. \]

The residual has degree below $4$, so coefficient extraction proves the selected recurrence for every stated index.

The closed form follows by substituting the same invariants in \eqref{round10:r10gf:eq:dim}; its exception at $n=1$ is retained in the displayed initial data.

\subsection{A063118: cusp forms of level 50}
\label{round10:r10entry:A063118}

\noindent\textbf{Source and scope.} \href{https://oeis.org/A063118}{A063118}, checked 20 October 2026. Conjecture; Klaus Brockhaus; Bruno Berselli; Vincenzo Librandi (2011-03-10; 2011-03-24).

\noindent\textbf{Indexing.} Offset $1$, with
\[ (a_{1},\ldots,a_{8})=(2,17,31,47,61,77,91,107). \]

The defining weight is $2n$. The selected ordinary generating function is

\[ G(x)=\frac{x \left(x^{3} + 12 x^{2} + 15 x + 2\right)}{\left(x - 1\right)^{2} \left(x + 1\right)}. \]

The selected recurrence, in coefficient-vector notation, is
\[ \sum_{j=0}^{3}p_j(n)a_{n-j}=0\qquad(n\ge 5). \]

\[ (p_0,\ldots,p_3)=(1,-1,-1,1). \]

The selected closed form is $a_n=\frac{\left(-1\right)^{n}}{2} + 15 n - \frac{27}{2}$ for $n\ge2$.

Also selected are $a_n=a_{n-2}+30$ and $a_n=2a_{n-1}-a_{n-2}+2(-1)^n$ for $n\ge4$.

\noindent\textbf{Proof.} This is a specialization of the credited dimension theorem.

The exact invariants are
\[ (\mu,e_2,e_3,c,g)=(90,2,0,12,2). \]

Substituting these values in \eqref{round10:r10gf:eq:dimGF} gives the stated $G$.

With $Q(x)=\sum_jp_jx^j$, direct rational simplification gives
\[ Q(x)G(x)=x^{4} + 12 x^{3} + 15 x^{2} + 2 x. \]

The residual has degree below $5$, so coefficient extraction proves the selected recurrence for every stated index.

The closed form follows by substituting the same invariants in \eqref{round10:r10gf:eq:dim}; its exception at $n=1$ is retained in the displayed initial data.

The two further identities follow by subtracting the closed form at indices differing by 2, and by taking its second difference, respectively.

\subsection{A063126: cusp forms of level 58}
\label{round10:r10entry:A063126}

\noindent\textbf{Source and scope.} \href{https://oeis.org/A063126}{A063126}, checked 20 October 2026. Conjecture; Colin Barker (2013-02-07).

\noindent\textbf{Indexing.} Offset $1$, with
\[ (a_{1},\ldots,a_{8})=(6,21,35,51,65,81,95,111). \]

The defining weight is $2n$. The selected ordinary generating function is

\[ G(x)=\frac{x \left(x + 2\right) \left(x^{2} + 6 x + 3\right)}{\left(x - 1\right)^{2} \left(x + 1\right)}. \]

The selected recurrence, in coefficient-vector notation, is
\[ \sum_{j=0}^{3}p_j(n)a_{n-j}=0\qquad(n\ge 5). \]

\[ (p_0,\ldots,p_3)=(1,-1,-1,1). \]

The selected closed form is $a_n=\frac{\left(-1\right)^{n}}{2} + 15 n - \frac{19}{2}$ for $n\ge2$.

\noindent\textbf{Proof.} This is a specialization of the credited dimension theorem.

The exact invariants are
\[ (\mu,e_2,e_3,c,g)=(90,2,0,4,6). \]

Substituting these values in \eqref{round10:r10gf:eq:dimGF} gives the stated $G$.

With $Q(x)=\sum_jp_jx^j$, direct rational simplification gives
\[ Q(x)G(x)=x^{4} + 8 x^{3} + 15 x^{2} + 6 x. \]

The residual has degree below $5$, so coefficient extraction proves the selected recurrence for every stated index.

The closed form follows by substituting the same invariants in \eqref{round10:r10gf:eq:dim}; its exception at $n=1$ is retained in the displayed initial data.

\subsection{A063140: cusp forms of level 72}
\label{round10:r10entry:A063140}

\noindent\textbf{Source and scope.} \href{https://oeis.org/A063140}{A063140}, checked 20 October 2026. Conjecture; Colin Barker; Stefano Spezia (2012-09-23; 2025-03-02).

\noindent\textbf{Indexing.} Offset $1$, with
\[ (a_{1},\ldots,a_{8})=(5,28,52,76,100,124,148,172). \]

The defining weight is $2n$. The selected ordinary generating function is

\[ G(x)=\frac{x \left(x^{2} + 18 x + 5\right)}{\left(1 - x\right)^{2}}. \]

The selected recurrence, in coefficient-vector notation, is
\[ \sum_{j=0}^{2}p_j(n)a_{n-j}=0\qquad(n\ge 4). \]

\[ (p_0,\ldots,p_2)=(1,-2,1). \]

The selected closed form is $a_n=24 n - 20$ for $n\ge2$.

The selected exponential generating function is $20+x+4e^x(6x-5)$.

\noindent\textbf{Proof.} This is a specialization of the credited dimension theorem.

The exact invariants are
\[ (\mu,e_2,e_3,c,g)=(144,0,0,16,5). \]

Substituting these values in \eqref{round10:r10gf:eq:dimGF} gives the stated $G$.

With $Q(x)=\sum_jp_jx^j$, direct rational simplification gives
\[ Q(x)G(x)=x^{3} + 18 x^{2} + 5 x. \]

The residual has degree below $4$, so coefficient extraction proves the selected recurrence for every stated index.

The closed form follows by substituting the same invariants in \eqref{round10:r10gf:eq:dim}; its exception at $n=1$ is retained in the displayed initial data.

Since $a_0=0$, $a_1=5$ and $a_n=24n-20$ for $n\ge2$, summing $(24n-20)x^n/n!$ and correcting its first two coefficients gives $24xe^x-20e^x+20+x$, exactly the selected exponential generating function.

\subsection{A063147: cusp forms of level 79}
\label{round10:r10entry:A063147}

\noindent\textbf{Source and scope.} \href{https://oeis.org/A063147}{A063147}, checked 20 October 2026. Conjecture; Colin Barker (2019-06-10).

\noindent\textbf{Indexing.} Offset $1$, with
\[ (a_{1},\ldots,a_{8})=(6,19,33,45,59,73,85,99). \]

The defining weight is $2n$. The selected ordinary generating function is

\[ G(x)=\frac{x \left(x^{4} + 6 x^{3} + 14 x^{2} + 13 x + 6\right)}{\left(1 - x\right)^{2} \left(x^{2} + x + 1\right)}. \]

The selected recurrence, in coefficient-vector notation, is
\[ \sum_{j=0}^{4}p_j(n)a_{n-j}=0\qquad(n\ge 6). \]

\[ (p_0,\ldots,p_4)=(1,-1,0,-1,1). \]

\noindent\textbf{Proof.} This is a specialization of the credited dimension theorem.

The exact invariants are
\[ (\mu,e_2,e_3,c,g)=(80,0,2,2,6). \]

Substituting these values in \eqref{round10:r10gf:eq:dimGF} gives the stated $G$.

With $Q(x)=\sum_jp_jx^j$, direct rational simplification gives
\[ Q(x)G(x)=x^{5} + 6 x^{4} + 14 x^{3} + 13 x^{2} + 6 x. \]

The residual has degree below $6$, so coefficient extraction proves the selected recurrence for every stated index.

The nearby newform sequence \href{https://oeis.org/A063233}{A063233} is not the same sequence: their values at $n=6$ are 73 and 71.

\subsection{A063148: cusp forms of level 80}
\label{round10:r10entry:A063148}

\noindent\textbf{Source and scope.} \href{https://oeis.org/A063148}{A063148}, checked 20 October 2026. Conjecture; R. J. Mathar; Colin Barker (2015-07-21; 2019-06-08).

\noindent\textbf{Indexing.} Offset $1$, with
\[ (a_{1},\ldots,a_{8})=(7,30,54,78,102,126,150,174). \]

The defining weight is $2n$. The selected ordinary generating function is

\[ G(x)=\frac{x \left(x^{2} + 16 x + 7\right)}{\left(1 - x\right)^{2}}. \]

The selected recurrence, in coefficient-vector notation, is
\[ \sum_{j=0}^{2}p_j(n)a_{n-j}=0\qquad(n\ge 4). \]

\[ (p_0,\ldots,p_2)=(1,-2,1). \]

The selected closed form is $a_n=24 n - 18$ for $n\ge2$.

\noindent\textbf{Proof.} This is a specialization of the credited dimension theorem.

The exact invariants are
\[ (\mu,e_2,e_3,c,g)=(144,0,0,12,7). \]

Substituting these values in \eqref{round10:r10gf:eq:dimGF} gives the stated $G$.

With $Q(x)=\sum_jp_jx^j$, direct rational simplification gives
\[ Q(x)G(x)=x^{3} + 16 x^{2} + 7 x. \]

The residual has degree below $4$, so coefficient extraction proves the selected recurrence for every stated index.

The closed form follows by substituting the same invariants in \eqref{round10:r10gf:eq:dim}; its exception at $n=1$ is retained in the displayed initial data.

The level-60 companion \href{https://oeis.org/A063128}{A063128} has the same invariants, so it is equal termwise. It is credited as a companion and contributes no package.

\subsection{A063168: cusp forms of level 100}
\label{round10:r10entry:A063168}

\noindent\textbf{Source and scope.} \href{https://oeis.org/A063168}{A063168}, checked 20 October 2026. Conjecture; Chai Wah Wu; Bill McEachen (2016-06-15; 2025-06-16).

\noindent\textbf{Indexing.} Offset $1$, with
\[ (a_{1},\ldots,a_{8})=(7,36,66,96,126,156,186,216). \]

The defining weight is $2n$. The selected ordinary generating function is

\[ G(x)=\frac{x \left(x^{2} + 22 x + 7\right)}{\left(x - 1\right)^{2}}. \]

The selected recurrence, in coefficient-vector notation, is
\[ \sum_{j=0}^{2}p_j(n)a_{n-j}=0\qquad(n\ge 4). \]

\[ (p_0,\ldots,p_2)=(1,-2,1). \]

The further selected assertion is $a_n=a_{n-1}+30$ for $n\ge3$.

\noindent\textbf{Proof.} This is a specialization of the credited dimension theorem.

The exact invariants are
\[ (\mu,e_2,e_3,c,g)=(180,0,0,18,7). \]

Substituting these values in \eqref{round10:r10gf:eq:dimGF} gives the stated $G$.

With $Q(x)=\sum_jp_jx^j$, direct rational simplification gives
\[ Q(x)G(x)=x^{3} + 22 x^{2} + 7 x. \]

The residual has degree below $4$, so coefficient extraction proves the selected recurrence for every stated index.

For $n\ge2$, \eqref{round10:r10gf:eq:dim} reduces to $30n-24$, proving the further first-difference assertion.

\subsection{A063247: newforms of level 100}
\label{round10:r10entry:A063247}

\noindent\textbf{Source and scope.} \href{https://oeis.org/A063247}{A063247}, checked 20 October 2026. Conjecture; Colin Barker (2019-06-08).

\noindent\textbf{Indexing.} Offset $1$, with
\[ (a_{1},\ldots,a_{8})=(1,5,8,11,14,18,20,24). \]

The defining weight is $2n$. The selected ordinary generating function is

\[ G(x)=\frac{x \left(x^{4} + 5 x^{3} + 7 x^{2} + 5 x + 1\right)}{\left(1 - x\right)^{2} \left(x + 1\right) \left(x^{2} + x + 1\right)}. \]

The selected recurrence, in coefficient-vector notation, is
\[ \sum_{j=0}^{5}p_j(n)a_{n-j}=0\qquad(n\ge 6). \]

\[ (p_0,\ldots,p_5)=(1,0,-1,-1,0,1). \]

\noindent\textbf{Proof.} This is a specialization of the credited dimension theorem.

For \eqref{round10:r10gf:eq:new}, the divisor data are as follows; $b=\overline\mu(100/M)$.

\[\begin{array}{r|rrrrrr}M&\mu&e_2&e_3&c&g&b\\\hline

1&1&1&1&1&0&1\\

2&3&1&0&2&0&-2\\

4&6&0&0&3&0&1\\

5&6&2&0&2&0&-2\\

10&18&2&0&4&0&4\\

20&36&0&0&6&1&-2\\

25&30&2&0&6&0&1\\

50&90&2&0&12&2&-2\\

100&180&0&0&18&7&1\\

\end{array}\]

The sum $\sum_{M\mid100}bG_M$ from \eqref{round10:r10gf:eq:dimGF} simplifies exactly to the stated $G$.

With $Q(x)=\sum_jp_jx^j$, direct rational simplification gives
\[ Q(x)G(x)=x^{5} + 5 x^{4} + 7 x^{3} + 5 x^{2} + x. \]

The residual has degree below $6$, so coefficient extraction proves the selected recurrence for every stated index.

\subsection{Two diagonals of Carlitz's polynomial triangle}
The independent definition is the polynomial recurrence in
\href{https://oeis.org/A259708}{A259708}, from L. Carlitz,
\emph{Some Polynomials Related to Fibonacci and Eulerian Numbers},
Fibonacci Quarterly 16 (1978), 216--226,
\href{https://www.fq.math.ca/Scanned/16-3/carlitz1.pdf}{equation (1.6)}:
\[
 r_0(z)=1,\quad r_1(z)=z,\qquad
 r_{m+1}(z)=(m+1)zr_m(z)+z(1-z)r'_m(z)+(1-z)^2r_{m-1}(z).
\]
The initial value $r_1=z$ corrects the printed $r_1=1$ typo, as explicitly
recorded by Eric Schmidt in A259708; the displayed polynomial table also has
$r_1=z$. Let $T(m,k)=[z^k]r_m(z)$, zero outside $0\le k\le m$.
Coefficient extraction gives
\begin{align*}
 T(m,k)={}&(m-k+1)T(m-1,k-1)+kT(m-1,k)\\
 &+T(m-2,k)-2T(m-2,k-1)+T(m-2,k-2).
\end{align*}
Set $D_j(x)=\sum_{k\ge0}T(k+j,k)x^k$, with $D_{-1}=D_{-2}=0$.
Multiplying this coefficient identity by $x^k$ and summing proves
\begin{equation}\label{round10:r10gf:eq:carlitz}
 (1-(j+1)x-x^2)D_j=xD'_{j-1}-2xD_{j-1}+D_{j-2}+[j=0].
\end{equation}
Indeed, the five terms give $(j+1)xD_j$, $xD'_{j-1}$,
$D_{j-2}$, $-2xD_{j-1}$ and $x^2D_j$; the exceptional cell is $T(0,0)=1$.
Thus, entirely from the defining triangle,
\[
 D_0=\frac1{1-x-x^2},\quad
 D_1=\frac{xD'_0-2xD_0}{1-2x-x^2},\quad
 D_2=\frac{xD'_1-2xD_1+D_0}{1-3x-x^2}.
\]

\subsection{A006502: a Carlitz diagonal}
\label{round10:r10entry:A006502}

\noindent\textbf{Source and scope.} \href{https://oeis.org/A006502}{A006502}, checked 20 October 2026. Conjecture; R. J. Mathar (2019-06-18).

\noindent\textbf{Indexing.} Offset $0$, with
\[ (a_{0},\ldots,a_{7})=(1,3,14,60,279,1251,5329,21532). \]

The independent sequence is $a_n=T(n+2,n)$ for $n\ge0$.

The source-marked rational function is $G=P/Q$, where

\begin{align*}P={}&- 3 x^{8} - 20 x^{7} - 39 x^{6} - 15 x^{5} + 5 x^{4} - 10 x^{3} + 18 x^{2} - 7 x + 1,\\Q={}&\left(x^{2} + x - 1\right)^{3} \left(x^{2} + 2 x - 1\right)^{2} \left(x^{2} + 3 x - 1\right).\end{align*}

\noindent\textbf{Proof.} Equation \eqref{round10:r10gf:eq:carlitz} gives $G=D_2$; substitution and rational simplification give exactly the stated numerator and denominator. This proves the generating function for all coefficients.

The two diagonals have different reduced denominators, so they are not scalar multiples or index shifts of each other.

\subsection{A259709: a Carlitz diagonal}
\label{round10:r10entry:A259709}

\noindent\textbf{Source and scope.} \href{https://oeis.org/A259709}{A259709}, checked 20 October 2026. Empirical; Eric M. Schmidt; R. J. Mathar (2015-07-10; 2015-07-15).

\noindent\textbf{Indexing.} Offset $1$, with
\[ (a_{1},\ldots,a_{8})=(0,-1,0,4,22,78,240,679). \]

The independent sequence is $a_n=T(n,n-1)$ for $n\ge1$.

The source-marked rational function is $G=P/Q$, where

\begin{align*}P={}&2 x^{4} + 4 x^{3} - x^{2},\\Q={}&- \left(x^{2} + x - 1\right)^{2} \left(x^{2} + 2 x - 1\right).\end{align*}

The selected recurrence, in coefficient-vector notation, is
\[ \sum_{j=0}^{6}p_j(n)a_{n-j}=0\qquad(n\ge 7). \]

\begin{align*}
&p_{0}(n)=1,\qquad p_{1}(n)=-4,\qquad p_{2}(n)=2,\\
&p_{3}(n)=6,\qquad p_{4}(n)=-2,\qquad p_{5}(n)=-4,\\
&p_{6}(n)=-1.
\end{align*}

\noindent\textbf{Proof.} Equation \eqref{round10:r10gf:eq:carlitz} gives $G=xD_1$; substitution and rational simplification give exactly the stated numerator and denominator. This proves the generating function for all coefficients.

For the recurrence polynomial $Q_0=\sum_jp_jx^j$, the exact identity is
\[ Q_0(x)G(x)=2 x^{4} + 4 x^{3} - x^{2}.\]

Its residual has degree 4, and all source indices exist for $n\ge7$, proving the selected recurrence.

The two diagonals have different reduced denominators, so they are not scalar multiples or index shifts of each other.

\subsection{A174015: a period-three continued fraction}
\label{round10:r10entry:A174015}

\noindent\textbf{Source and scope.} \href{https://oeis.org/A174015}{A174015}, checked 20 October 2026. Conjecture; R. J. Mathar (2012-11-13).

\noindent\textbf{Indexing.} Offset $0$, with
\[ (a_{0},\ldots,a_{7})=(1,1,2,3,4,4,2,-3). \]

The selected recurrence, in coefficient-vector notation, is
\[ \sum_{j=0}^{4}p_j(n)a_{n-j}=0\qquad(n\ge 4). \]

\begin{align*}
&p_{0}(n)=n + 1,\qquad p_{1}(n)=- 3 n,\qquad p_{2}(n)=3 n - 3,\\
&p_{3}(n)=3 n - 12,\qquad p_{4}(n)=14 - 4 n.
\end{align*}

\noindent\textbf{Independent construction and prior credit.}
A174015 is the first column of \href{https://oeis.org/A174014}{A174014}.
The defining Del\'eham transform has period-three data
$r=(1,1,-1,\ldots)$, $s=(1,0,0,\ldots)$.
The constant-column continued fraction is the case $y=0$ of
\href{https://oeis.org/A084938}{A084938}. Its three tails obey
\[
 C_0=\frac1{1-xC_1},\qquad C_1=\frac1{1-xC_2},\qquad
 C_2=\frac1{1+xC_0}.
\]
The generating function was already treated by Mathar,
\emph{Properties of Del\'eham's Delta Transformation} (2015),
\href{https://oeis.org/A084938/a084938.pdf}{Theorem 4 and Examples 17, 20}.
That generating function is an existing input, not a new claim here.

\noindent\textbf{Proof.}
Eliminating $C_1,C_2$ and setting $F=C_0$ gives
\[
 x(1-x)F^2+(1-3x)F+x-1=0.
\]
The derivative with respect to $F$ at $(x,F)=(0,1)$ is 1, so this has a
unique formal solution with constant term 1. Solving the quadratic gives
\[
 F=\frac{3x-1+\sqrt{1-2x+x^2+4x^3}}{2x(1-x)}.
\]
Apply \eqref{round10:r10gf:eq:operator} with the displayed $p_j$. In the basis
$1,\sqrt{1-2x+x^2+4x^3}$, the exact result is
\[
 LF=1-x+3x^2-3x^3+0\sqrt{1-2x+x^2+4x^3}.
\]
The radical coefficient cancels identically and the remaining polynomial
has degree 3. Consequently the recurrence holds for all $n\ge4$.

\subsection{A226302: tatami polynomials at minus one}
\label{round10:r10entry:A226302}

\noindent\textbf{Source and scope.} \href{https://oeis.org/A226302}{A226302}, checked 20 October 2026. Conjecture; R. J. Mathar (2013-11-06).

\noindent\textbf{Indexing.} Offset $2$, with
\[ (a_{2},\ldots,a_{9})=(1,-1,2,-4,6,-14,20,-48). \]

The selected recurrence, in coefficient-vector notation, is
\[ \sum_{j=0}^{5}p_j(n)a_{n-j}=0\qquad(n\ge 7). \]

\begin{align*}
&p_{0}(n)=2 - n,\qquad p_{1}(n)=2 - n,\qquad p_{2}(n)=6 n - 22,\\
&p_{3}(n)=6 n - 28,\qquad p_{4}(n)=36 - 8 n,\qquad p_{5}(n)=48 - 8 n.
\end{align*}

\noindent\textbf{Defining data and selected scope.}
The sequence is $a_n=P_n(-1)$, $n\ge2$, for the polynomials of
A. Erickson and F. Ruskey, \emph{Enumerating Maximal Tatami Mat Coverings of Square Grids
with v Vertical Dominoes},
\href{https://arxiv.org/abs/1304.0070}{arXiv:1304.0070}.
Their equations (1)--(2) define $P_n=V_n/D_n$ using $S_m(z)=\prod_{j=1}^m(1+z^j)$:
\begin{align*}
 V_n(z)&=2\sum_{i=1}^{\lfloor(n-1)/2\rfloor}
 S_{n-i-2}(z)S_{i-1}(z)z^{n-i-1}
 +S_{\lfloor(n-2)/2\rfloor}(z)^2,\\
 D_n(z)&=\prod_{j\ge1}S_{\lfloor(n-2)/2^j\rfloor}(z).
\end{align*}
The product is finite because $S_0=1$. The paper proves the divisibility
needed to make $P_n$ a polynomial. Its Conjecture 2(e) states the generating
function of $P_{n+2}(-1)$. We prove it below from the defining products;
we do not assume the later unmarked generating-function entry as a theorem.
It and the OEIS recurrence form one package.

\noindent\textbf{Proof by leading terms at $z=-1$.}
Put $z=-1+h$ and
\[
 v_m=\lceil m/2\rceil,\qquad L_m=\frac{m!}{\lfloor m/2\rfloor!}.
\]
Each odd $j$ contributes $jh+O(h^2)$ and each even $j$ contributes
$2+O(h)$, hence
\begin{equation}\label{round10:r10gf:eq:lead}
 S_m(-1+h)=L_mh^{v_m}+O(h^{v_m+1}).
\end{equation}
For $q=\lfloor(n-2)/2\rfloor$, the denominator has valuation $q$ and leading
coefficient $q!$. Both statements follow inductively from
$\lceil q/2\rceil+\lfloor q/2\rfloor=q$ and
$L_q\lfloor q/2\rfloor!=q!$.

First let $n=2m+2$, $m\ge1$. Every summand in $V_n$ has valuation at
least $m$, and its coefficient at $h^m$ is
$2(-1)^{i+1}L_{2m-i}L_{i-1}$. The terms $i=2r$ and $i=2r+1$ cancel,
because $L_{2s}=2L_{2s-1}$ for $s\ge1$. If $m$ is even, the last
unpaired negative term cancels the central square's coefficient
$L_m^2=2L_mL_{m-1}$. If $m$ is odd, that square has larger valuation.
Only the $i=1$ term remains. Thus, also checking $m=0$ directly,
\begin{equation}\label{round10:r10gf:eq:tatamieven}
 P_{2m+2}(-1)=\frac{2L_{2m-1}}{m!}=\binom{2m}{m}\quad(m\ge1),
 \qquad P_2(-1)=1.
\end{equation}

Next let $n=2m+3$. Only odd $i=2r+1$ contribute at order $m$;
each contributes $-2L_{2(m-r)}L_{2r}$. The central square, when $m$ is
even, compensates exactly for the duplicated middle term. With
$u_r=L_{2r}=(2r)!/r!$ and $T_m=\sum_{r=0}^mu_ru_{m-r}$, it follows that
\[
 P_{2m+3}(-1)=-\frac{T_m}{m!}.
\]
The identity $u_r=2(2r-1)u_{r-1}$ and symmetry of this convolution give,
for $m\ge1$,
\begin{align*}
 T_m-u_m
 &=2\sum_{r=0}^{m-1}(2r+1)u_ru_{m-1-r}
 =2mT_{m-1}.
\end{align*}
Therefore $b_m=T_m/m!$ satisfies $b_0=1$ and
$b_m=2b_{m-1}+\binom{2m}{m}$. Induction proves
\begin{equation}\label{round10:r10gf:eq:tatamiodd}
 P_{2m+3}(-1)=-\sum_{j=0}^m2^{m-j}\binom{2j}{j}\qquad(m\ge0).
\end{equation}
Using $\sum_{m\ge0}\binom{2m}{m}t^m=(1-4t)^{-1/2}$ now proves
\[
 F(x)=\sum_{n\ge2}P_n(-1)x^n
 =\frac{x^2}{\sqrt{1-4x^2}}
   \left(1-\frac{x}{1-2x^2}\right).
\]
After dividing by $x^2$, this is precisely the paper's Conjecture 2(e).
Finally write
\[
 F=\frac{x^2(1-x-2x^2)}{(1-2x^2)(1-4x^2)}\sqrt{1-4x^2}.
\]
Substitution in \eqref{round10:r10gf:eq:operator} gives $LF=0$ exactly in both
quadratic-field components. All source indices exist for $n\ge7$, so the
selected OEIS recurrence follows for its entire stated range.

\subsection{A273905: symmetric bargraphs}
\label{round10:r10entry:A273905}

\noindent\textbf{Source and scope.} \href{https://oeis.org/A273905}{A273905}, checked 20 October 2026. Conjecture; R. J. Mathar (2022-07-22).

\noindent\textbf{Indexing.} Offset $2$, with
\[ (a_{2},\ldots,a_{9})=(1,2,3,5,9,15,27,46). \]

The selected recurrence, in coefficient-vector notation, is
\[ \sum_{j=0}^{8}p_j(n)a_{n-j}=0\qquad(n\ge 10). \]

\begin{align*}
&p_{0}(n)=n + 1,\qquad p_{1}(n)=-2,\qquad p_{2}(n)=6 - 4 n,\\
&p_{3}(n)=0,\qquad p_{4}(n)=2 n - 12,\qquad p_{5}(n)=6,\\
&p_{6}(n)=-2,\qquad p_{7}(n)=0,\qquad p_{8}(n)=n - 9.
\end{align*}

\noindent\textbf{Independent definition and prior credit.}
The objects are bargraphs invariant under vertical reflection, counted by semiperimeter.
Deutsch and Elizalde, \emph{Statistics on Bargraphs Viewed as Cornerless
Motzkin Paths}, \href{https://arxiv.org/abs/1609.00088}{Section 4.1},
prove the generating function. Heinz's 2016 OEIS program already contains
the recurrence; its exact proof, rather than discovery priority, is at issue.

Their decomposition can be summarized as follows. Let $M(u,v)$ count
cornerless Motzkin paths by horizontal and up steps. First return gives
$M=(1+v(M-1))(1+uM)$. Prefixes ending at height $h$ have series
$P_h=M(v/K)^h$, where $K=1-u-uv(M-1)$, because
$P_h=vP_{h-1}+(u+uv(M-1))P_h$.
Every nonempty symmetric path is uniquely $AHA'$ or $AHHA'$, with $A'$
reflected. If $A$ ends at height $h$, its reflection supplies $h$ fewer
up steps. The bargraph bijection adds one up step. Thus
\begin{equation}\label{round10:r10gf:eq:bargraph}
 F=x(x+x^2)M(x^2,x^2)\frac K{K-x},\qquad
 K=1-x^2-x^4(M(x^2,x^2)-1).
\end{equation}
Solving the quadratic with constant term 1 gives
\[
 M(x^2,x^2)=\frac{(1-x^2)^2-\sqrt D}{2x^4},\qquad
 D=(1-x^2)(1-x-x^2-x^3)(1+x-x^2+x^3).
\]
Substitution and exact reduction modulo $(\sqrt D)^2-D$ yield
\[
 F=\frac{(1+x)(x^4+2x^3+2x^2-1+\sqrt D)}
         {2x(1-x-x^2-x^3)}.
\]
The verifier independently enumerates palindromic positive column heights.

\noindent\textbf{Recurrence proof.}
The operator \eqref{round10:r10gf:eq:operator} with the displayed $p_j$ gives
\[
 LF=3x^2+6x^3+x^4-4x^5-x^6+2x^7+x^8+0\sqrt D.
\]
The residual has degree 8. The zero-extended identity starts at $n=9$;
for the source sequence of offset 2, every listed term exists at $n\ge10$.
This proves the selected recurrence on that range.

\subsection{Two current-offset corrections for bounded Motzkin walks}
The classical finite-path enumeration is credited in advance: Felsner and
Heldt, \emph{Lattice Path Enumeration and Toeplitz Matrices},
Journal of Integer Sequences 18 (2015),
\href{https://cs.uwaterloo.ca/journals/JIS/VOL18/Felsner/felsner2.pdf}{Fact 10 and Example 17},
give the general solution and explicitly list both A-numbers below.
Accordingly these packages correct a boundary mismatch in the currently
printed generating functions, not an unsolved enumeration theorem.

Let $G_m$ count closed walks beginning at the bottom of a path with $m$
states, with a loop at each state. Decomposing at the first return gives
$G_1=1/(1-x)$ and $G_m=1/(1-x-x^2G_{m-1})$. Equivalently, with
\[
 D_0=1,\quad D_1=1-x,\quad
 D_m=(1-x)D_{m-1}-x^2D_{m-2},
\]
induction proves $G_m=D_{m-1}/D_m$. All denominators have constant term 1,
so these are identities of formal power series with $G_m(0)=1$.

\subsection{A094287: a boundary correction after adding the empty walk}
\label{round10:r10entry:A094287}

\noindent\textbf{Source and scope.} \href{https://oeis.org/A094287}{A094287}, checked 20 October 2026. Conjecture; R. J. Mathar (2011-12-20).

\noindent\textbf{Indexing.} Offset $0$, with
\[ (a_{0},\ldots,a_{7})=(1,1,2,4,9,21,51,127). \]

The definition counts walks on states $1,\ldots,6$, starting and ending at 1, with increments in $\{-1,0,1\}$.

The source-marked 2011 generating function currently printed in the entry is

\[ H(x)=- \frac{x \left(x^{5} + 2 x^{4} - 5 x^{3} - 2 x^{2} + 4 x - 1\right)}{\left(x^{3} - x^{2} - 2 x + 1\right) \left(x^{3} + 3 x^{2} - 4 x + 1\right)}. \]

The selected recurrence, in coefficient-vector notation, is
\[ \sum_{j=0}^{6}p_j(n)a_{n-j}=0\qquad(n\ge 6). \]

\begin{align*}
&p_{0}(n)=1,\qquad p_{1}(n)=-6,\qquad p_{2}(n)=10,\\
&p_{3}(n)=0,\qquad p_{4}(n)=-9,\qquad p_{5}(n)=2,\\
&p_{6}(n)=1.
\end{align*}

\noindent\textbf{Correction and proof.} The empty walk gives $a_0=1$, whereas $H(0)=0$. The term $a_0=1$ was prepended in 2023, after the 2011 statement. Thus this is a mismatch with the current offset; the old positive-index coefficients are unaffected.

The continuant calculation gives $G=P/Q$, where
\begin{align*}P={}&- 4 x^{4} + 2 x^{3} + 6 x^{2} - 5 x + 1,\\Q={}&x^{6} + 2 x^{5} - 9 x^{4} + 10 x^{2} - 6 x + 1.\end{align*}

Exact rational simplification yields $G=1+H$. This proves the corrected generating function at every index and supplies the explicit counterexample to the unadjusted constant term. Since $Q=\sum_jp_jx^j$ and $\deg P<6$, the selected recurrence also follows for all $n\ge6$.

\subsection{A094288: a boundary correction after adding the empty walk}
\label{round10:r10entry:A094288}

\noindent\textbf{Source and scope.} \href{https://oeis.org/A094288}{A094288}, checked 20 October 2026. Conjecture; R. J. Mathar (2011-12-20).

\noindent\textbf{Indexing.} Offset $0$, with
\[ (a_{0},\ldots,a_{7})=(1,1,2,4,9,21,51,127). \]

The definition counts walks on states $1,\ldots,7$, starting and ending at 1, with increments in $\{-1,0,1\}$.

The source-marked 2011 generating function currently printed in the entry is

\[ H(x)=- \frac{x \left(x^{6} - 2 x^{5} - 7 x^{4} + 6 x^{3} + 5 x^{2} - 5 x + 1\right)}{\left(x - 1\right) \left(x^{2} + 2 x - 1\right) \left(x^{4} - 4 x^{3} - 2 x^{2} + 4 x - 1\right)}. \]

The selected recurrence, in coefficient-vector notation, is
\[ \sum_{j=0}^{7}p_j(n)a_{n-j}=0\qquad(n\ge 7). \]

\begin{align*}
&p_{0}(n)=1,\qquad p_{1}(n)=-7,\qquad p_{2}(n)=15,\\
&p_{3}(n)=-5,\qquad p_{4}(n)=-15,\qquad p_{5}(n)=9,\\
&p_{6}(n)=3,\qquad p_{7}(n)=-1.
\end{align*}

\noindent\textbf{Correction and proof.} The empty walk gives $a_0=1$, whereas $H(0)=0$. The term $a_0=1$ was prepended in 2023, after the 2011 statement. Thus this is a mismatch with the current offset; the old positive-index coefficients are unaffected.

The continuant calculation gives $G=P/Q$, where
\begin{align*}P={}&x^{6} + 2 x^{5} - 9 x^{4} + 10 x^{2} - 6 x + 1,\\Q={}&- x^{7} + 3 x^{6} + 9 x^{5} - 15 x^{4} - 5 x^{3} + 15 x^{2} - 7 x + 1.\end{align*}

Exact rational simplification yields $G=1+H$. This proves the corrected generating function at every index and supplies the explicit counterexample to the unadjusted constant term. Since $Q=\sum_jp_jx^j$ and $\deg P<7$, the selected recurrence also follows for all $n\ge7$.

\subsection{Scope of the audit and reproducibility}
The selected IDs were compared against the parent's 908 previously selected
IDs, its prior uncounted reserve and known-result exclusions, and the
97-ID OEIS Open appendix baseline. The portable exclusion copy and the
per-entry source statements accompany this fragment. Freshly retrieved OEIS
pages and the defining references listed above delimit the prior-result audit;
no exhaustive claim about all unpublished or unindexed proofs is made.
The modular dimension formulas, the A174015 and A273905 generating functions,
and the bounded-walk general theorem are expressly credited as known results.
The two blocked candidates A099553 and A308620 were not selected.

The dimension packages are related specializations. In particular, for
$n\ge1$,
\[
 a_{063110}(n)=a_{063116}(n)+2,\qquad
 a_{063148}(n)=a_{063140}(n)+2.
\]
Also $a_{063140}(n)=2a_{063098}(n)$ and
$2a_{063148}(n)=3a_{063110}(n)$ for $n\ge2$, with respective differences
$-1$ at $n=1$. These exact identities are disclosed in the companion audit;
the count is of source-assertion packages, not unrelated mathematical methods.
The exact companion A063128 is not counted.

Run \texttt{python r10\_gf\_verify.py} for the exact certificates, or
\texttt{python r10\_gf\_build.py} to regenerate this fragment and those
certificates. Python 3 and SymPy are sufficient; the scripts use no network
and import no prior-round module. The JSON certificates contain the actual
operator residuals, rational numerators and denominators, all dimension
parameters, source-prefix comparisons, and independent coefficient checks
through index 100. Their universal identities, together with the mathematical
definition-to-series arguments above, prove the stated full ranges.

\section{Restricted-growth word arrays: exact recurrence certificates}
\label{round10:sec:r10arrays}

This section proves \textbf{65 distinct primary-entry assertion packages} and retains
3 further freshly certified packages outside that count.  In each selected
OEIS snapshot the recurrence is explicitly marked \emph{Empirical}.  The selection
contains twelve closely related local-avoidance families.  It is a collection of
specific source assertions proved by a common finite-state argument, not a claim
to 65 unrelated major open problems.  The 31 listed source generating
functions that are also verified add no further count.

The primary entries and all 12 named parent tables were retrieved for this
round on October 20, 2026.  Cache ages are retained in the source audit; the retrieval
date does not establish that the snapshots include every later live edit.  The
primary recurrence of A212827 is empirical, although its column is absent from
the parent table's displayed empirical formula list.  The primary entries list
numeric b-files rather than linked proofs; those b-files were not expanded, and
the 1452 printed selected-entry terms were compared with the definition.
The retained source authorship and formula attributions distinguish R. H. Hardin's
entries from later generating-function contributions.

\subsection{Definitions and the twelve local predicates}

A word $w=w_1\cdots w_\ell$ over $\{0,\ldots,q-1\}$ is a
\emph{restricted-growth word} if $w_1=0$ and, for every $j>1$,
\[
  0\leq w_j\leq\min\bigl(q-1,1+\max(w_1,\ldots,w_{j-1})\bigr).
\]
Thus new labels appear in the order $0,1,\ldots$; a word need not use every
available label.  This is the precise meaning of the entries' introduction-order
condition.  Existing labels remain distinct and are never divided out by a
color-permutation factor.

For any finite word $v$, let $\operatorname{can}(v)$ relabel its distinct letters
by $0,1,\ldots$ in their order of first appearance.  For example,
$\operatorname{can}(7,2,7,2)=0101$ and
$\operatorname{can}(7,7,7,7)=0000$.  The forbidden patterns below concern
\emph{equality}, not the numerical order of distinct letters.  Consequently a
forbidden $aba$ with $a\ne b$ excludes both numerical orders of its two letters.
Forbidding $abc$ with three unequal letters excludes all six numerical orders.
Only consecutive windows are tested.

Each entry specifies an alphabet size $q$, a memory $K$, a length shift $s$,
and a set $\mathcal B$ of canonical words of length $K+1$.
Its source term $a(n)$ counts the permitted words of length $n+s$ for $n\geq1$.
The complete translation of the local conditions is as follows.
\begin{center}
\small
\begin{tabular}{lll}
\hline
Source local condition & $\mathcal B$ & $(K,s)$\\
\hline
Repeated adjacent pair, allowing a constant pair & $\{0000,0101\}$ & $(3,0)$\\
$aabb$, possibly $a=b$ & $\{0000,0011\}$ & $(3,3)$\\
$aabb$, $a\ne b$ & $\{0011\}$ & $(3,3)$\\
$abba$, possibly $a=b$ & $\{0000,0110\}$ & $(3,3)$\\
$abab$, $a\ne b$ & $\{0101\}$ & $(3,3)$\\
$abba$, $a\ne b$ & $\{0110\}$ & $(3,3)$\\
$aba$ or $abb$, $a\ne b$ & $\{010,011\}$ & $(2,2)$\\
$aba$, $a\ne b$ & $\{010\}$ & $(2,2)$\\
Three pairwise unequal letters & $\{012\}$ & $(2,2)$\\
Three equal or three pairwise unequal letters & $\{000,012\}$ & $(2,2)$\\
$ababa$ or $abbba$, possibly $a=b$ & $\{00000,01010,01110\}$ & $(4,4)$\\
$ababa$ or $abbba$, $a\ne b$ & $\{01010,01110\}$ & $(4,4)$\\
\hline
\end{tabular}
\end{center}
In the first row, two consecutive adjacent pairs are equal exactly when the
four-letter window has its first and third letters equal and its second and
fourth letters equal.  The word length is $n$ in that family; shifting it to
$n+3$ would change both the source sequence and its recurrence threshold.

\subsection{The complete state space and weighted transitions}

Let $\mathcal P_{K,q}$ be all canonical restricted-growth words of length $K$
using at most $q$ letters.  Define
\[
 \mathcal S=\{(m,p):1\leq m\leq q,\ p\in\mathcal P_{K,q},\ r(p)\leq m\},
 \qquad r(p)=1+\max p.
\]
The state records the number $m$ of labels introduced anywhere in the prefix
and the equality pattern $p$ of its final $K$ letters.  The state order is
increasing $m$, then lexicographic $p$.  All these states are retained, including
states unreachable from the initial words.  In particular the matrix dimension is
\[
 D=\sum_{r=1}^{\min(K,q)}\left\{\begin{matrix}K\\r\end{matrix}\right\}(q-r+1),
\]
where the braces count partitions of a $K$-element set into $r$ blocks.

Let $T$ be the row-to-column transition matrix.  From $(m,p)$, put $r=r(p)$.
Every legal append belongs to exactly one of the following three cases.
\begin{enumerate}
\item For each letter $c\in\{0,\ldots,r-1\}$ already present in the suffix,
  if $\operatorname{can}(pc)\notin\mathcal B$, add weight one to
  $(m,\operatorname{can}(p_2\cdots p_Kc))$.
\item There are $m-r$ previously introduced letters absent from the suffix.
  They all have the same equality pattern when appended, represented by $c=r$.
  If $\operatorname{can}(pr)\notin\mathcal B$, add weight $m-r$ to
  $(m,\operatorname{can}(p_2\cdots p_Kr))$.
\item If $m<q$, the single next new label may be introduced.  With the same
  test $\operatorname{can}(pr)\notin\mathcal B$, add weight one to
  $(m+1,\operatorname{can}(p_2\cdots p_Kr))$.
\end{enumerate}
Weights with the same destination are added; zero weights are omitted.
These cases exhaust the actual labels available to an introduced-order word.
The forbidden-window test and destination depend only on the stated equality
pattern, so two actual prefixes with the same state have exactly these same
numbers of extensions to each destination.  This proves that the weighted
compression preserves the number of words, rather than assuming a symmetry
quotient for the counted objects.

Let $u$ be the row vector with $u_{(m,p)}=1$ when $m=r(p)$ and zero otherwise.
Every length-$K$ restricted-growth word is one of these $p$'s, and cannot yet
contain a forbidden window of length $K+1$.  Hence $u$ is the exact initial
vector.  With $\mathbf1$ the all-ones column, induction on the number of appended
letters gives
\begin{equation}
 \label{round10:eq:r10arr-transfer}
 a(n)=uT^{\,n+s-K}\mathbf1\qquad(n+s\geq K).
\end{equation}
When $n+s<K$, the finitely many short words are counted directly.  Only
A212824--A212827 use this exception, and their first two source terms are $1,2$.

\subsection{A finite vector identity proves every later recurrence term}

For an entry let the displayed source coefficients be $c_1,\ldots,c_d$ and let
\[
 b=d+1+K-s,\qquad e=b+s-K=d+1,\qquad
 v_j=T^j\mathbf1\quad(0\leq j\leq e).
\]
The exact integer certificate retains every coordinate of every $v_j$ and of
\begin{equation}
 \label{round10:eq:r10arr-vector}
 z=v_e-\sum_{i=1}^d c_i v_{e-i}.
\end{equation}
Let $W$ be the set of states reachable by nonzero transitions from the support
of $u$.  It contains that support and is closed under transitions by definition;
the verifiers also check both properties directly, and independently reconstruct
the whole reachable set.  Every certificate has $z_t=0$ for every $t\in W$.
In fact, in all 68 retained packages the stronger statement $z=0$ on
the full state space holds: all 2188 residual coordinates are zero.
No unreachable coordinate has been removed from the certificate.

For $n\geq b$ all indices $n-i$ in the recurrence are source indices at least
one, all powers in \eqref{round10:eq:r10arr-transfer} are nonnegative, and
\[
\begin{split}
 a(n)-\sum_{i=1}^d c_i a(n-i)
 &=uT^{\,n-b}\left(T^e\mathbf1-\sum_{i=1}^d c_iT^{e-i}\mathbf1\right)\\
 &=uT^{\,n-b}z=0.
\end{split}
\]
The last equality already follows from $z|_W=0$ and the closure of $W$.
We have therefore proved, universally,
\begin{equation}
 \label{round10:eq:r10arr-rec}
 a(n)=\sum_{i=1}^d c_i a(n-i)\qquad(n\geq b).
\end{equation}
For length $n+K$ entries this means $b=d+1$.  For A212824--A212827,
$s=0,K=3$ and therefore $b=d+4$, namely $8,10,12,14$ respectively,
exactly the explicit thresholds in their source assertions.

This proof does not infer a recurrence from a long list of scalar values.
The source supplies the candidate coefficients, the definition supplies $T,u$,
and \eqref{round10:eq:r10arr-vector} is checked as an exact vector equality.
The finite identity is sufficient for all subsequent $n$ by the displayed
algebra.  No assumption about minimal recurrence order is needed.

For reproducibility, the central identity check is the following finite integer
calculation.  The full verifier additionally reconstructs $T,u$ from the definition,
checks the source coefficients and threshold, checks the reachability set, and
checks every stored vector against this calculation.
\begin{Verbatim}[fontsize=\scriptsize,breaklines=true]
# rows[t] contains every (destination, positive integer weight).
# initial[t] is u_t. c contains the d source coefficients.
v = [[1] * len(rows)]
for j in range(d + 1):
    v.append([sum(weight * v[-1][dest] for dest, weight in row)
              for row in rows])
z = [v[d+1][t] - sum(c[i-1] * v[d+1-i][t]
                    for i in range(1, d+1))
     for t in range(len(rows))]
W = {t for t, value in enumerate(initial) if value}
pending = list(W)
while pending:
    t = pending.pop()
    for dest, weight in rows[t]:
        if dest not in W:
            W.add(dest)
            pending.append(dest)
assert all(dest in W for t in W for dest, weight in rows[t])
assert all(z[t] == 0 for t in W)
# All coordinates were retained; the present 68 cases also satisfy:
assert z == [0] * len(rows)
\end{Verbatim}

\subsection{Generating functions and an explicit small certificate}

Put $F(x)=\sum_{n\geq1}a(n)x^n$ and $Q(x)=1-\sum_{i=1}^d c_i x^i$.
For this formal series let $A_0=0$ and $A_j=a(j)$ for $j\geq1$; $A_0$ is
a series convention, not an assertion about an auxiliary transfer value at
source index zero.  If $Q(x)=\sum_{i=0}^d q_i x^i$, define
\[
 P(x)=\sum_{j=0}^{b-1}\left(\sum_{i=0}^{\min(j,d)}q_iA_{j-i}\right)x^j.
\]
Equation~\eqref{round10:eq:r10arr-rec} makes every coefficient of $QF$ of degree at
least $b$ zero; the lower coefficients are exactly those displayed in $P$.
Thus $QF=P$ as formal power series.  Since $Q(0)=1$, division is legitimate,
and $F=P/Q$.  This proves each recorded source generating function by exact
polynomial cross-multiplication, not by agreement through finitely many terms.

For example A243513 has $q=5$, $K=s=2$, and
$\mathcal B=\{010,011\}$.  The full matrix has nine states.  Its reachable
states, in order, are
\[
 (1,00),\ (2,01),\ (3,01),\ (4,01),\ (5,01).
\]
Their closed transition block and initial row are
\[
 T_W=\begin{pmatrix}
 1&1&0&0&0\\0&0&1&0&0\\0&0&1&1&0\\
 0&0&0&2&1\\0&0&0&0&3
 \end{pmatrix},\qquad u_W=(1,1,0,0,0).
\]
The four successive vectors, restricted here to $W$ for a readable worked
example, are
\[
\begin{array}{c|rrrrr}
j&\multicolumn{5}{c}{(T_W^j\mathbf1)^{\mathsf T}}\\\hline
1&2&1&2&3&3\\
2&3&2&5&9&9\\
3&5&5&14&27&27\\
4&10&14&41&81&81
\end{array}
\]
Every column satisfies $v_4-5v_3+7v_2-3v_1=0$.
Therefore $a(n)=5a(n-1)-7a(n-2)+3a(n-3)$ for all $n\geq4$,
with initial terms $3,5,10$, and
\[
 F(x)=\frac{3x-10x^2+6x^3}{1-5x+7x^2-3x^3}.
\]
The machine certificate also retains all four other state coordinates and checks
them; the small display is not a deletion from the full proof data.
The explicitly empirical source and its parent are
\url{https://oeis.org/A243513} and \url{https://oeis.org/A243519}.

\subsection{Independent checks, exclusions, and scope}

The 65 selected matrices have 2098 full states and
6816 nonzero weighted edges.  They retain 55400 Krylov-vector
coordinates, with maximum dimension 128 and maximum recurrence order
49.  The independent implementation enumerates every actual suffix over
the introduced labels, tests forbidden windows by pairwise equality, and checks
that every actual transition aggregates to the stated compressed row.  It imports
neither the compressed engine nor its verifier.  For the selected set it checked
131159 labelled states, 1052986 legal labelled transitions, all
2098 compressed rows, and 406 direct short-word/source comparisons.
It independently checks the exact vector identity and closed reachable set as well.
This is an internal independent implementation check, not external peer review or
proof-assistant formalization.

No selected ID occurs among the retained 908 earlier selected IDs, five earlier
uncounted reserves, 22 retained known-result exclusions, or the 97 distinct IDs
of the OEIS Open Appendix A baseline.  The latter audit concerns that appendix's
100 costly claims, not every result in its repository.  Identity and arbitrary
rational-scale relations were screened at shifts $-3,\ldots,3$ against all 352
retained earlier array packages, including their available reserves.  Each of the
167552 prior comparisons failed within 20 terms.  Each of the 15946
within-round comparisons failed within 60 terms.  These mismatches exclude those
particular relations; they do not exclude larger shifts or other transforms.

Four initially considered three-equal-letter avoidance entries,
A242468--A242471, were removed after a direct prior result was located.
Mansour and Munagi, \emph{Enumeration of Partitions by Long Rises, Levels,
and Descents}, J. Integer Sequences 12 (2009), Article 09.1.8,
Theorem 2.1 and Section 2.1, give for exactly $k$ blocks
\[
 L_3(x,0;k)=\frac{(x+x^2)^k}{\prod_{j=0}^{k-1}(1-j(x+x^2))}.
\]
Summing this expression over $1\leq k\leq q$ and removing the length-one
and length-two terms gives exactly the source series
$(\sum_{k=1}^qL_3(x,0;k)-x-2x^2)/x^2$.
Those four direct specializations have count zero.
The primary paper is
\url{https://cs.uwaterloo.ca/journals/JIS/VOL12/Munagi/munagi15.pdf}.

The bounded literature review also read Mansour, Shattuck and Yan,
\emph{Counting Subwords in a Partition of a Set}, Electron. J. Combin.
17 (2010), R19, at
\url{https://www.combinatorics.org/ojs/index.php/eljc/article/download/v17i1r19/pdf/},
and Mansour and Shattuck, \emph{Counting Peaks and Valleys in a Partition of
a Set}, J. Integer Sequences 13 (2010), Article 10.6.8, at
\url{https://emis.de/ft/3989}.
The former treats individual numerical-order subwords, so its separate $121$
and $212$ assertions are not themselves a joint equality-pattern avoidance
formula.  The latter's inequalities for peaks and valleys differ from the
selected equal-endpoint conditions.  No direct proof of a retained selected
assertion was identified in the inspected portions.  A potentially relevant
2020 IEEE paper by Sabri and Tran on restricted-growth functions avoiding
vincular patterns was found, but its full text could not be retrieved; that
is an explicit coverage gap.  The compact literature audit preserves the
primary URLs, search scope, and this limitation.

These facts establish the proofs and the source-marked status of the selected
assertions in the retrieved material.  They do not establish worldwide priority
or the absence of every earlier proof.  No OEIS edit, message, or submission
was performed, and no official closure is claimed.  Other formulas in the same
entries and unselected columns of parent tables are not silently counted.

\subsection{Entry-by-entry assertions and certificates}

Each coefficient list below is in the order $(c_1,\ldots,c_d)$ used in
\eqref{round10:eq:r10arr-rec}.  Each package supplies its full $D$-state matrix,
initial vector, all $e+1=d+2$ full vectors, and every residual coordinate in
the portable JSON certificate files.  The transition definition above and
the per-entry parameters completely specify the finite calculation; the
listed assertion follows from \eqref{round10:eq:r10arr-vector} and the universal
argument, not from the three displayed initial terms.
\Needspace{7\baselineskip}
\subsubsection{A212824}
\label{round10:r10entry:A212824}
Parameters: $q=4$, $K=3$, $s=0$, and $\mathcal B=\{0000,0101\}$. The source counts words of length $n$ for $n\geq1$.  Its empirical recurrence has order $d=4$ and is proved for every $n\geq 8$, with coefficients
\[(c_1,\ldots,c_d)=(4,1,-6,-3).\]
The complete matrix has $D=15$ states, with $|W|=15$ reachable states. All $90$ coordinates of $v_0,\ldots,v_{5}$ are retained. The residual in \eqref{round10:eq:r10arr-vector} is zero in all $15$ coordinates; the initial support lies in the verified closed set $W$. Thus \eqref{round10:eq:r10arr-rec} follows for every asserted index. The first three source values are $(1,2,5)$.
Source: \url{https://oeis.org/A212824}; column 3 of \url{https://oeis.org/A212829}. The primary recurrence and the relevant parent formula block are explicitly empirical.
The source rational generating function is also proved: $F=P/Q$, where $Q=1-\sum_{i=1}^d c_i x^i$ and
\[P(x)=x-2x^{2}-4x^{3}-3x^{4}+x^{5}+3x^{6}+x^{7}.\]
\Needspace{7\baselineskip}
\subsubsection{A212825}
\label{round10:r10entry:A212825}
Parameters: $q=5$, $K=3$, $s=0$, and $\mathcal B=\{0000,0101\}$. The source counts words of length $n$ for $n\geq1$.  Its empirical recurrence has order $d=6$ and is proved for every $n\geq 10$, with coefficients
\[(c_1,\ldots,c_d)=(7,-7,-20,10,24,8).\]
The complete matrix has $D=20$ states, with $|W|=20$ reachable states. All $160$ coordinates of $v_0,\ldots,v_{7}$ are retained. The residual in \eqref{round10:eq:r10arr-vector} is zero in all $20$ coordinates; the initial support lies in the verified closed set $W$. Thus \eqref{round10:eq:r10arr-rec} follows for every asserted index. The first three source values are $(1,2,5)$.
Source: \url{https://oeis.org/A212825}; column 4 of \url{https://oeis.org/A212829}. The primary recurrence and the relevant parent formula block are explicitly empirical.
The source rational generating function is also proved: $F=P/Q$, where $Q=1-\sum_{i=1}^d c_i x^i$ and
\[P(x)=x-5x^{2}-2x^{3}+12x^{4}+18x^{5}+6x^{6}-12x^{7}-12x^{8}-3x^{9}.\]
\Needspace{7\baselineskip}
\subsubsection{A212826}
\label{round10:r10entry:A212826}
Parameters: $q=6$, $K=3$, $s=0$, and $\mathcal B=\{0000,0101\}$. The source counts words of length $n$ for $n\geq1$.  Its empirical recurrence has order $d=8$ and is proved for every $n\geq 12$, with coefficients
\begin{Verbatim}[fontsize=\scriptsize,breaklines=true]
11, -30, -21, 112, 63, -119, -120, -30
\end{Verbatim}
The complete matrix has $D=25$ states, with $|W|=25$ reachable states. All $250$ coordinates of $v_0,\ldots,v_{9}$ are retained. The residual in \eqref{round10:eq:r10arr-vector} is zero in all $25$ coordinates; the initial support lies in the verified closed set $W$. Thus \eqref{round10:eq:r10arr-rec} follows for every asserted index. The first three source values are $(1,2,5)$.
Source: \url{https://oeis.org/A212826}; column 5 of \url{https://oeis.org/A212829}. The primary recurrence and the relevant parent formula block are explicitly empirical.
The source rational generating function is also proved: $F=P/Q$, where $Q=1-\sum_{i=1}^d c_i x^i$ and
\[\begin{aligned}P(x)&=x-9x^{2}+13x^{3}+39x^{4}-19x^{5}-108x^{6}-106x^{7}+6x^{8}+84x^{9}\\ &\quad +55x^{10}+11x^{11}.\end{aligned}\]
\Needspace{7\baselineskip}
\subsubsection{A212827}
\label{round10:r10entry:A212827}
Parameters: $q=7$, $K=3$, $s=0$, and $\mathcal B=\{0000,0101\}$. The source counts words of length $n$ for $n\geq1$.  Its empirical recurrence has order $d=10$ and is proved for every $n\geq 14$, with coefficients
\begin{Verbatim}[fontsize=\scriptsize,breaklines=true]
16, -79, 70, 361, -372, -964, 144, 1116, 720, 144
\end{Verbatim}
The complete matrix has $D=30$ states, with $|W|=30$ reachable states. All $360$ coordinates of $v_0,\ldots,v_{11}$ are retained. The residual in \eqref{round10:eq:r10arr-vector} is zero in all $30$ coordinates; the initial support lies in the verified closed set $W$. Thus \eqref{round10:eq:r10arr-rec} follows for every asserted index. The first three source values are $(1,2,5)$.
Source: \url{https://oeis.org/A212827}; column 6 of \url{https://oeis.org/A212829}. The primary recurrence is explicitly empirical. The parent definition was read, but its displayed empirical formulas do not include this column.
\Needspace{7\baselineskip}
\subsubsection{A242543}
\label{round10:r10entry:A242543}
Parameters: $q=3$, $K=3$, $s=3$, and $\mathcal B=\{0000,0011\}$. The source counts words of length $n+3$ for $n\geq1$.  Its empirical recurrence has order $d=4$ and is proved for every $n\geq 5$, with coefficients
\[(c_1,\ldots,c_d)=(2,0,4,4).\]
The complete matrix has $D=10$ states, with $|W|=10$ reachable states. All $60$ coordinates of $v_0,\ldots,v_{5}$ are retained. The residual in \eqref{round10:eq:r10arr-vector} is zero in all $10$ coordinates; the initial support lies in the verified closed set $W$. Thus \eqref{round10:eq:r10arr-rec} follows for every asserted index. The first three source values are $(12,32,88)$.
Source: \url{https://oeis.org/A242543}; column 2 of \url{https://oeis.org/A242549}. The primary recurrence and the relevant parent formula block are explicitly empirical.
The source rational generating function is also proved: $F=P/Q$, where $Q=1-\sum_{i=1}^d c_i x^i$ and
\[P(x)=12x+8x^{2}+24x^{3}+18x^{4}.\]
\Needspace{7\baselineskip}
\subsubsection{A242544}
\label{round10:r10entry:A242544}
Parameters: $q=4$, $K=3$, $s=3$, and $\mathcal B=\{0000,0011\}$. The source counts words of length $n+3$ for $n\geq1$.  Its empirical recurrence has order $d=8$ and is proved for every $n\geq 9$, with coefficients
\begin{Verbatim}[fontsize=\scriptsize,breaklines=true]
4, -3, 10, -2, -12, -9, -18, -9
\end{Verbatim}
The complete matrix has $D=15$ states, with $|W|=15$ reachable states. All $150$ coordinates of $v_0,\ldots,v_{9}$ are retained. The residual in \eqref{round10:eq:r10arr-vector} is zero in all $15$ coordinates; the initial support lies in the verified closed set $W$. Thus \eqref{round10:eq:r10arr-rec} follows for every asserted index. The first three source values are $(13,42,150)$.
Source: \url{https://oeis.org/A242544}; column 3 of \url{https://oeis.org/A242549}. The primary recurrence and the relevant parent formula block are explicitly empirical.
The source rational generating function is also proved: $F=P/Q$, where $Q=1-\sum_{i=1}^d c_i x^i$ and
\[P(x)=13x-10x^{2}+21x^{3}-50x^{4}-88x^{5}-78x^{6}-99x^{7}-42x^{8}.\]
\Needspace{7\baselineskip}
\subsubsection{A242546}
\label{round10:r10entry:A242546}
Parameters: $q=6$, $K=3$, $s=3$, and $\mathcal B=\{0000,0011\}$. The source counts words of length $n+3$ for $n\geq1$.  Its empirical recurrence has order $d=16$ and is proved for every $n\geq 17$, with coefficients
\begin{Verbatim}[fontsize=\scriptsize,breaklines=true]
11, -41, 100, -259, 283, -178, 383, 1393, 312, 723, -1707, -5129, -5430, -5400, -3600,
-900
\end{Verbatim}
The complete matrix has $D=25$ states, with $|W|=25$ reachable states. All $450$ coordinates of $v_0,\ldots,v_{17}$ are retained. The residual in \eqref{round10:eq:r10arr-vector} is zero in all $25$ coordinates; the initial support lies in the verified closed set $W$. Thus \eqref{round10:eq:r10arr-rec} follows for every asserted index. The first three source values are $(13,43,166)$.
Source: \url{https://oeis.org/A242546}; column 5 of \url{https://oeis.org/A242549}. The primary recurrence and the relevant parent formula block are explicitly empirical.
\Needspace{7\baselineskip}
\subsubsection{A242547}
\label{round10:r10entry:A242547}
Parameters: $q=7$, $K=3$, $s=3$, and $\mathcal B=\{0000,0011\}$. The source counts words of length $n+3$ for $n\geq1$.  Its empirical recurrence has order $d=20$ and is proved for every $n\geq 21$, with coefficients
\begin{Verbatim}[fontsize=\scriptsize,breaklines=true]
16, -95, 326, -998, 2268, -2953, 5002, -877, -8932, -5460, -46536, -36728, -4608,
43488, 180288, 260208, 254016, 207360, 103680, 20736
\end{Verbatim}
The complete matrix has $D=30$ states, with $|W|=30$ reachable states. All $660$ coordinates of $v_0,\ldots,v_{21}$ are retained. The residual in \eqref{round10:eq:r10arr-vector} is zero in all $30$ coordinates; the initial support lies in the verified closed set $W$. Thus \eqref{round10:eq:r10arr-rec} follows for every asserted index. The first three source values are $(13,43,166)$.
Source: \url{https://oeis.org/A242547}; column 6 of \url{https://oeis.org/A242549}. The primary recurrence and the relevant parent formula block are explicitly empirical.
\Needspace{7\baselineskip}
\subsubsection{A242548}
\label{round10:r10entry:A242548}
Parameters: $q=8$, $K=3$, $s=3$, and $\mathcal B=\{0000,0011\}$. The source counts words of length $n+3$ for $n\geq1$.  Its empirical recurrence has order $d=24$ and is proved for every $n\geq 25$, with coefficients
\begin{Verbatim}[fontsize=\scriptsize,breaklines=true]
22, -190, 924, -3465, 10962, -24402, 45828, -74016, 19016, -23090, -159064, 571211,
704098, 1846562, 2774132, 616507, -2599138, -8906720, -16117180, -18096476, -15823920,
-10584000, -4233600, -705600
\end{Verbatim}
The complete matrix has $D=35$ states, with $|W|=35$ reachable states. All $910$ coordinates of $v_0,\ldots,v_{25}$ are retained. The residual in \eqref{round10:eq:r10arr-vector} is zero in all $35$ coordinates; the initial support lies in the verified closed set $W$. Thus \eqref{round10:eq:r10arr-rec} follows for every asserted index. The first three source values are $(13,43,166)$.
Source: \url{https://oeis.org/A242548}; column 7 of \url{https://oeis.org/A242549}. The primary recurrence and the relevant parent formula block are explicitly empirical.
\Needspace{7\baselineskip}
\subsubsection{A242578}
\label{round10:r10entry:A242578}
Parameters: $q=3$, $K=3$, $s=3$, and $\mathcal B=\{0011\}$. The source counts words of length $n+3$ for $n\geq1$.  Its empirical recurrence has order $d=4$ and is proved for every $n\geq 5$, with coefficients
\[(c_1,\ldots,c_d)=(4,-5,6,-4).\]
The complete matrix has $D=10$ states, with $|W|=10$ reachable states. All $60$ coordinates of $v_0,\ldots,v_{5}$ are retained. The residual in \eqref{round10:eq:r10arr-vector} is zero in all $10$ coordinates; the initial support lies in the verified closed set $W$. Thus \eqref{round10:eq:r10arr-rec} follows for every asserted index. The first three source values are $(13,35,97)$.
Source: \url{https://oeis.org/A242578}; column 2 of \url{https://oeis.org/A242584}. The primary recurrence and the relevant parent formula block are explicitly empirical.
The source rational generating function is also proved: $F=P/Q$, where $Q=1-\sum_{i=1}^d c_i x^i$ and
\[P(x)=13x-17x^{2}+22x^{3}-20x^{4}.\]
\Needspace{7\baselineskip}
\subsubsection{A242579}
\label{round10:r10entry:A242579}
Parameters: $q=4$, $K=3$, $s=3$, and $\mathcal B=\{0011\}$. The source counts words of length $n+3$ for $n\geq1$.  Its empirical recurrence has order $d=7$ and is proved for every $n\geq 8$, with coefficients
\begin{Verbatim}[fontsize=\scriptsize,breaklines=true]
7, -18, 32, -45, 37, -21, 9
\end{Verbatim}
The complete matrix has $D=15$ states, with $|W|=15$ reachable states. All $135$ coordinates of $v_0,\ldots,v_{8}$ are retained. The residual in \eqref{round10:eq:r10arr-vector} is zero in all $15$ coordinates; the initial support lies in the verified closed set $W$. Thus \eqref{round10:eq:r10arr-rec} follows for every asserted index. The first three source values are $(14,45,159)$.
Source: \url{https://oeis.org/A242579}; column 3 of \url{https://oeis.org/A242584}. The primary recurrence and the relevant parent formula block are explicitly empirical.
The source rational generating function is also proved: $F=P/Q$, where $Q=1-\sum_{i=1}^d c_i x^i$ and
\[P(x)=14x-53x^{2}+96x^{3}-164x^{4}+152x^{5}-87x^{6}+45x^{7}.\]
\Needspace{7\baselineskip}
\subsubsection{A242580}
\label{round10:r10entry:A242580}
Parameters: $q=5$, $K=3$, $s=3$, and $\mathcal B=\{0011\}$. The source counts words of length $n+3$ for $n\geq1$.  Its empirical recurrence has order $d=10$ and is proved for every $n\geq 11$, with coefficients
\begin{Verbatim}[fontsize=\scriptsize,breaklines=true]
11, -48, 131, -282, 452, -541, 530, -364, 176, -64
\end{Verbatim}
The complete matrix has $D=20$ states, with $|W|=20$ reachable states. All $240$ coordinates of $v_0,\ldots,v_{11}$ are retained. The residual in \eqref{round10:eq:r10arr-vector} is zero in all $20$ coordinates; the initial support lies in the verified closed set $W$. Thus \eqref{round10:eq:r10arr-rec} follows for every asserted index. The first three source values are $(14,46,174)$.
Source: \url{https://oeis.org/A242580}; column 4 of \url{https://oeis.org/A242584}. The primary recurrence and the relevant parent formula block are explicitly empirical.
\Needspace{7\baselineskip}
\subsubsection{A242581}
\label{round10:r10entry:A242581}
Parameters: $q=6$, $K=3$, $s=3$, and $\mathcal B=\{0011\}$. The source counts words of length $n+3$ for $n\geq1$.  Its empirical recurrence has order $d=13$ and is proved for every $n\geq 14$, with coefficients
\begin{Verbatim}[fontsize=\scriptsize,breaklines=true]
16, -106, 425, -1284, 3083, -5823, 9091, -11623, 11784, -9883, 6151, -2730, 900
\end{Verbatim}
The complete matrix has $D=25$ states, with $|W|=25$ reachable states. All $375$ coordinates of $v_0,\ldots,v_{14}$ are retained. The residual in \eqref{round10:eq:r10arr-vector} is zero in all $25$ coordinates; the initial support lies in the verified closed set $W$. Thus \eqref{round10:eq:r10arr-rec} follows for every asserted index. The first three source values are $(14,46,175)$.
Source: \url{https://oeis.org/A242581}; column 5 of \url{https://oeis.org/A242584}. The primary recurrence and the relevant parent formula block are explicitly empirical.
\Needspace{7\baselineskip}
\subsubsection{A242582}
\label{round10:r10entry:A242582}
Parameters: $q=7$, $K=3$, $s=3$, and $\mathcal B=\{0011\}$. The source counts words of length $n+3$ for $n\geq1$.  Its empirical recurrence has order $d=16$ and is proved for every $n\geq 17$, with coefficients
\begin{Verbatim}[fontsize=\scriptsize,breaklines=true]
22, -206, 1156, -4714, 15358, -40643, 89400, -167115, 262308, -349485, 393316,
-361252, 276048, -160848, 67392, -20736
\end{Verbatim}
The complete matrix has $D=30$ states, with $|W|=30$ reachable states. All $540$ coordinates of $v_0,\ldots,v_{17}$ are retained. The residual in \eqref{round10:eq:r10arr-vector} is zero in all $30$ coordinates; the initial support lies in the verified closed set $W$. Thus \eqref{round10:eq:r10arr-rec} follows for every asserted index. The first three source values are $(14,46,175)$.
Source: \url{https://oeis.org/A242582}; column 6 of \url{https://oeis.org/A242584}. The primary recurrence and the relevant parent formula block are explicitly empirical.
\Needspace{7\baselineskip}
\subsubsection{A243027}
\label{round10:r10entry:A243027}
Parameters: $q=3$, $K=3$, $s=3$, and $\mathcal B=\{0000,0110\}$. The source counts words of length $n+3$ for $n\geq1$.  Its empirical recurrence has order $d=3$ and is proved for every $n\geq 4$, with coefficients
\[(c_1,\ldots,c_d)=(2,1,2).\]
The complete matrix has $D=10$ states, with $|W|=10$ reachable states. All $50$ coordinates of $v_0,\ldots,v_{4}$ are retained. The residual in \eqref{round10:eq:r10arr-vector} is zero in all $10$ coordinates; the initial support lies in the verified closed set $W$. Thus \eqref{round10:eq:r10arr-rec} follows for every asserted index. The first three source values are $(12,32,85)$.
Source: \url{https://oeis.org/A243027}; column 2 of \url{https://oeis.org/A243033}. The primary recurrence and the relevant parent formula block are explicitly empirical.
The source rational generating function is also proved: $F=P/Q$, where $Q=1-\sum_{i=1}^d c_i x^i$ and
\[P(x)=12x+8x^{2}+9x^{3}.\]
\Needspace{7\baselineskip}
\subsubsection{A243028}
\label{round10:r10entry:A243028}
Parameters: $q=4$, $K=3$, $s=3$, and $\mathcal B=\{0000,0110\}$. The source counts words of length $n+3$ for $n\geq1$.  Its empirical recurrence has order $d=6$ and is proved for every $n\geq 7$, with coefficients
\[(c_1,\ldots,c_d)=(4,-1,2,-6,-2,-3).\]
The complete matrix has $D=15$ states, with $|W|=15$ reachable states. All $120$ coordinates of $v_0,\ldots,v_{7}$ are retained. The residual in \eqref{round10:eq:r10arr-vector} is zero in all $15$ coordinates; the initial support lies in the verified closed set $W$. Thus \eqref{round10:eq:r10arr-rec} follows for every asserted index. The first three source values are $(13,42,147)$.
Source: \url{https://oeis.org/A243028}; column 3 of \url{https://oeis.org/A243033}. The primary recurrence and the relevant parent formula block are explicitly empirical.
The source rational generating function is also proved: $F=P/Q$, where $Q=1-\sum_{i=1}^d c_i x^i$ and
\[P(x)=13x-10x^{2}-8x^{3}-36x^{4}-16x^{5}-14x^{6}.\]
\Needspace{7\baselineskip}
\subsubsection{A243029}
\label{round10:r10entry:A243029}
Parameters: $q=5$, $K=3$, $s=3$, and $\mathcal B=\{0000,0110\}$. The source counts words of length $n+3$ for $n\geq1$.  Its empirical recurrence has order $d=9$ and is proved for every $n\geq 10$, with coefficients
\begin{Verbatim}[fontsize=\scriptsize,breaklines=true]
7, -10, 1, -21, 13, 6, 27, 10, 8
\end{Verbatim}
The complete matrix has $D=20$ states, with $|W|=20$ reachable states. All $220$ coordinates of $v_0,\ldots,v_{10}$ are retained. The residual in \eqref{round10:eq:r10arr-vector} is zero in all $20$ coordinates; the initial support lies in the verified closed set $W$. Thus \eqref{round10:eq:r10arr-rec} follows for every asserted index. The first three source values are $(13,43,162)$.
Source: \url{https://oeis.org/A243029}; column 4 of \url{https://oeis.org/A243033}. The primary recurrence and the relevant parent formula block are explicitly empirical.
The source rational generating function is also proved: $F=P/Q$, where $Q=1-\sum_{i=1}^d c_i x^i$ and
\[P(x)=13x-48x^{2}-9x^{3}-45x^{4}+108x^{5}+98x^{6}+154x^{7}+63x^{8}+37x^{9}.\]
\Needspace{7\baselineskip}
\subsubsection{A243038}
\label{round10:r10entry:A243038}
Parameters: $q=3$, $K=3$, $s=3$, and $\mathcal B=\{0101\}$. The source counts words of length $n+3$ for $n\geq1$.  Its empirical recurrence has order $d=4$ and is proved for every $n\geq 5$, with coefficients
\[(c_1,\ldots,c_d)=(3,0,-1,-1).\]
The complete matrix has $D=10$ states, with $|W|=10$ reachable states. All $60$ coordinates of $v_0,\ldots,v_{5}$ are retained. The residual in \eqref{round10:eq:r10arr-vector} is zero in all $10$ coordinates; the initial support lies in the verified closed set $W$. Thus \eqref{round10:eq:r10arr-rec} follows for every asserted index. The first three source values are $(13,36,101)$.
Source: \url{https://oeis.org/A243038}; column 2 of \url{https://oeis.org/A243044}. The primary recurrence and the relevant parent formula block are explicitly empirical.
The source rational generating function is also proved: $F=P/Q$, where $Q=1-\sum_{i=1}^d c_i x^i$ and
\[P(x)=13x-3x^{2}-7x^{3}-5x^{4}.\]
\Needspace{7\baselineskip}
\subsubsection{A243039}
\label{round10:r10entry:A243039}
Parameters: $q=4$, $K=3$, $s=3$, and $\mathcal B=\{0101\}$. The source counts words of length $n+3$ for $n\geq1$.  Its empirical recurrence has order $d=7$ and is proved for every $n\geq 8$, with coefficients
\begin{Verbatim}[fontsize=\scriptsize,breaklines=true]
5, -3, -5, -3, 3, 3, 1
\end{Verbatim}
The complete matrix has $D=15$ states, with $|W|=15$ reachable states. All $135$ coordinates of $v_0,\ldots,v_{8}$ are retained. The residual in \eqref{round10:eq:r10arr-vector} is zero in all $15$ coordinates; the initial support lies in the verified closed set $W$. Thus \eqref{round10:eq:r10arr-rec} follows for every asserted index. The first three source values are $(14,46,163)$.
Source: \url{https://oeis.org/A243039}; column 3 of \url{https://oeis.org/A243044}. The primary recurrence and the relevant parent formula block are explicitly empirical.
The source rational generating function is also proved: $F=P/Q$, where $Q=1-\sum_{i=1}^d c_i x^i$ and
\[P(x)=14x-24x^{2}-25x^{3}-5x^{4}+22x^{5}+17x^{6}+5x^{7}.\]
\Needspace{7\baselineskip}
\subsubsection{A243040}
\label{round10:r10entry:A243040}
Parameters: $q=5$, $K=3$, $s=3$, and $\mathcal B=\{0101\}$. The source counts words of length $n+3$ for $n\geq1$.  Its empirical recurrence has order $d=10$ and is proved for every $n\geq 11$, with coefficients
\begin{Verbatim}[fontsize=\scriptsize,breaklines=true]
8, -14, -10, 13, 28, 9, -12, -14, -6, -1
\end{Verbatim}
The complete matrix has $D=20$ states, with $|W|=20$ reachable states. All $240$ coordinates of $v_0,\ldots,v_{11}$ are retained. The residual in \eqref{round10:eq:r10arr-vector} is zero in all $20$ coordinates; the initial support lies in the verified closed set $W$. Thus \eqref{round10:eq:r10arr-rec} follows for every asserted index. The first three source values are $(14,47,178)$.
Source: \url{https://oeis.org/A243040}; column 4 of \url{https://oeis.org/A243044}. The primary recurrence and the relevant parent formula block are explicitly empirical.
\Needspace{7\baselineskip}
\subsubsection{A243041}
\label{round10:r10entry:A243041}
Parameters: $q=6$, $K=3$, $s=3$, and $\mathcal B=\{0101\}$. The source counts words of length $n+3$ for $n\geq1$.  Its empirical recurrence has order $d=13$ and is proved for every $n\geq 14$, with coefficients
\begin{Verbatim}[fontsize=\scriptsize,breaklines=true]
12, -41, 13, 96, 33, -134, -166, -33, 82, 87, 41, 10, 1
\end{Verbatim}
The complete matrix has $D=25$ states, with $|W|=25$ reachable states. All $375$ coordinates of $v_0,\ldots,v_{14}$ are retained. The residual in \eqref{round10:eq:r10arr-vector} is zero in all $25$ coordinates; the initial support lies in the verified closed set $W$. Thus \eqref{round10:eq:r10arr-rec} follows for every asserted index. The first three source values are $(14,47,179)$.
Source: \url{https://oeis.org/A243041}; column 5 of \url{https://oeis.org/A243044}. The primary recurrence and the relevant parent formula block are explicitly empirical.
\Needspace{7\baselineskip}
\subsubsection{A243042}
\label{round10:r10entry:A243042}
Parameters: $q=7$, $K=3$, $s=3$, and $\mathcal B=\{0101\}$. The source counts words of length $n+3$ for $n\geq1$.  Its empirical recurrence has order $d=16$ and is proved for every $n\geq 17$, with coefficients
\begin{Verbatim}[fontsize=\scriptsize,breaklines=true]
17, -95, 154, 222, -448, -773, 213, 1374, 1181, 27, -736, -690, -334, -95, -15, -1
\end{Verbatim}
The complete matrix has $D=30$ states, with $|W|=30$ reachable states. All $540$ coordinates of $v_0,\ldots,v_{17}$ are retained. The residual in \eqref{round10:eq:r10arr-vector} is zero in all $30$ coordinates; the initial support lies in the verified closed set $W$. Thus \eqref{round10:eq:r10arr-rec} follows for every asserted index. The first three source values are $(14,47,179)$.
Source: \url{https://oeis.org/A243042}; column 6 of \url{https://oeis.org/A243044}. The primary recurrence and the relevant parent formula block are explicitly empirical.
\Needspace{7\baselineskip}
\subsubsection{A243043}
\label{round10:r10entry:A243043}
Parameters: $q=8$, $K=3$, $s=3$, and $\mathcal B=\{0101\}$. The source counts words of length $n+3$ for $n\geq1$.  Its empirical recurrence has order $d=19$ and is proved for every $n\geq 20$, with coefficients
\begin{Verbatim}[fontsize=\scriptsize,breaklines=true]
23, -190, 614, -135, -2559, 369, 7121, 5143, -7257, -15107, -9219, 2219, 7955, 6635,
3195, 982, 190, 21, 1
\end{Verbatim}
The complete matrix has $D=35$ states, with $|W|=35$ reachable states. All $735$ coordinates of $v_0,\ldots,v_{20}$ are retained. The residual in \eqref{round10:eq:r10arr-vector} is zero in all $35$ coordinates; the initial support lies in the verified closed set $W$. Thus \eqref{round10:eq:r10arr-rec} follows for every asserted index. The first three source values are $(14,47,179)$.
Source: \url{https://oeis.org/A243043}; column 7 of \url{https://oeis.org/A243044}. The primary recurrence and the relevant parent formula block are explicitly empirical.
\Needspace{7\baselineskip}
\subsubsection{A243383}
\label{round10:r10entry:A243383}
Parameters: $q=3$, $K=3$, $s=3$, and $\mathcal B=\{0110\}$. The source counts words of length $n+3$ for $n\geq1$.  Its empirical recurrence has order $d=4$ and is proved for every $n\geq 5$, with coefficients
\[(c_1,\ldots,c_d)=(4,-4,2,-1).\]
The complete matrix has $D=10$ states, with $|W|=10$ reachable states. All $60$ coordinates of $v_0,\ldots,v_{5}$ are retained. The residual in \eqref{round10:eq:r10arr-vector} is zero in all $10$ coordinates; the initial support lies in the verified closed set $W$. Thus \eqref{round10:eq:r10arr-rec} follows for every asserted index. The first three source values are $(13,35,96)$.
Source: \url{https://oeis.org/A243383}; column 2 of \url{https://oeis.org/A243389}. The primary recurrence and the relevant parent formula block are explicitly empirical.
The source rational generating function is also proved: $F=P/Q$, where $Q=1-\sum_{i=1}^d c_i x^i$ and
\[P(x)=13x-17x^{2}+8x^{3}-5x^{4}.\]
\Needspace{7\baselineskip}
\subsubsection{A243384}
\label{round10:r10entry:A243384}
Parameters: $q=4$, $K=3$, $s=3$, and $\mathcal B=\{0110\}$. The source counts words of length $n+3$ for $n\geq1$.  Its empirical recurrence has order $d=7$ and is proved for every $n\geq 8$, with coefficients
\begin{Verbatim}[fontsize=\scriptsize,breaklines=true]
7, -16, 18, -15, 9, -3, 1
\end{Verbatim}
The complete matrix has $D=15$ states, with $|W|=15$ reachable states. All $135$ coordinates of $v_0,\ldots,v_{8}$ are retained. The residual in \eqref{round10:eq:r10arr-vector} is zero in all $15$ coordinates; the initial support lies in the verified closed set $W$. Thus \eqref{round10:eq:r10arr-rec} follows for every asserted index. The first three source values are $(14,45,158)$.
Source: \url{https://oeis.org/A243384}; column 3 of \url{https://oeis.org/A243389}. The primary recurrence and the relevant parent formula block are explicitly empirical.
The source rational generating function is also proved: $F=P/Q$, where $Q=1-\sum_{i=1}^d c_i x^i$ and
\[P(x)=14x-53x^{2}+67x^{3}-59x^{4}+40x^{5}-13x^{6}+5x^{7}.\]
\Needspace{7\baselineskip}
\subsubsection{A243385}
\label{round10:r10entry:A243385}
Parameters: $q=5$, $K=3$, $s=3$, and $\mathcal B=\{0110\}$. The source counts words of length $n+3$ for $n\geq1$.  Its empirical recurrence has order $d=10$ and is proved for every $n\geq 11$, with coefficients
\begin{Verbatim}[fontsize=\scriptsize,breaklines=true]
11, -44, 87, -107, 101, -71, 37, -16, 4, -1
\end{Verbatim}
The complete matrix has $D=20$ states, with $|W|=20$ reachable states. All $240$ coordinates of $v_0,\ldots,v_{11}$ are retained. The residual in \eqref{round10:eq:r10arr-vector} is zero in all $20$ coordinates; the initial support lies in the verified closed set $W$. Thus \eqref{round10:eq:r10arr-rec} follows for every asserted index. The first three source values are $(14,46,173)$.
Source: \url{https://oeis.org/A243385}; column 4 of \url{https://oeis.org/A243389}. The primary recurrence and the relevant parent formula block are explicitly empirical.
\Needspace{7\baselineskip}
\subsubsection{A243386}
\label{round10:r10entry:A243386}
Parameters: $q=6$, $K=3$, $s=3$, and $\mathcal B=\{0110\}$. The source counts words of length $n+3$ for $n\geq1$.  Its empirical recurrence has order $d=13$ and is proved for every $n\geq 14$, with coefficients
\begin{Verbatim}[fontsize=\scriptsize,breaklines=true]
16, -99, 313, -584, 752, -757, 592, -370, 191, -74, 25, -5, 1
\end{Verbatim}
The complete matrix has $D=25$ states, with $|W|=25$ reachable states. All $375$ coordinates of $v_0,\ldots,v_{14}$ are retained. The residual in \eqref{round10:eq:r10arr-vector} is zero in all $25$ coordinates; the initial support lies in the verified closed set $W$. Thus \eqref{round10:eq:r10arr-rec} follows for every asserted index. The first three source values are $(14,46,174)$.
Source: \url{https://oeis.org/A243386}; column 5 of \url{https://oeis.org/A243389}. The primary recurrence and the relevant parent formula block are explicitly empirical.
\Needspace{7\baselineskip}
\subsubsection{A243387}
\label{round10:r10entry:A243387}
Parameters: $q=7$, $K=3$, $s=3$, and $\mathcal B=\{0110\}$. The source counts words of length $n+3$ for $n\geq1$.  Its empirical recurrence has order $d=16$ and is proved for every $n\geq 17$, with coefficients
\begin{Verbatim}[fontsize=\scriptsize,breaklines=true]
22, -195, 914, -2539, 4594, -6074, 6382, -5368, 3726, -2167, 1026, -419, 130, -36, 6,
-1
\end{Verbatim}
The complete matrix has $D=30$ states, with $|W|=30$ reachable states. All $540$ coordinates of $v_0,\ldots,v_{17}$ are retained. The residual in \eqref{round10:eq:r10arr-vector} is zero in all $30$ coordinates; the initial support lies in the verified closed set $W$. Thus \eqref{round10:eq:r10arr-rec} follows for every asserted index. The first three source values are $(14,46,174)$.
Source: \url{https://oeis.org/A243387}; column 6 of \url{https://oeis.org/A243389}. The primary recurrence and the relevant parent formula block are explicitly empirical.
\Needspace{7\baselineskip}
\subsubsection{A243388}
\label{round10:r10entry:A243388}
Parameters: $q=8$, $K=3$, $s=3$, and $\mathcal B=\{0110\}$. The source counts words of length $n+3$ for $n\geq1$.  Its empirical recurrence has order $d=19$ and is proved for every $n\geq 20$, with coefficients
\begin{Verbatim}[fontsize=\scriptsize,breaklines=true]
29, -349, 2287, -9065, 23201, -41175, 55373, -60041, 53229, -39663, 25113, -13424,
6232, -2386, 806, -209, 49, -7, 1
\end{Verbatim}
The complete matrix has $D=35$ states, with $|W|=35$ reachable states. All $735$ coordinates of $v_0,\ldots,v_{20}$ are retained. The residual in \eqref{round10:eq:r10arr-vector} is zero in all $35$ coordinates; the initial support lies in the verified closed set $W$. Thus \eqref{round10:eq:r10arr-rec} follows for every asserted index. The first three source values are $(14,46,174)$.
Source: \url{https://oeis.org/A243388}; column 7 of \url{https://oeis.org/A243389}. The primary recurrence and the relevant parent formula block are explicitly empirical.
\Needspace{7\baselineskip}
\subsubsection{A243513}
\label{round10:r10entry:A243513}
Parameters: $q=5$, $K=2$, $s=2$, and $\mathcal B=\{010,011\}$. The source counts words of length $n+2$ for $n\geq1$.  Its empirical recurrence has order $d=3$ and is proved for every $n\geq 4$, with coefficients
\[(c_1,\ldots,c_d)=(5,-7,3).\]
The complete matrix has $D=9$ states, with $|W|=5$ reachable states. All $45$ coordinates of $v_0,\ldots,v_{4}$ are retained. The residual in \eqref{round10:eq:r10arr-vector} is zero in all $9$ coordinates; the initial support lies in the verified closed set $W$. Thus \eqref{round10:eq:r10arr-rec} follows for every asserted index. The first three source values are $(3,5,10)$.
Source: \url{https://oeis.org/A243513}; column 4 of \url{https://oeis.org/A243519}. The primary recurrence and the relevant parent formula block are explicitly empirical.
The source rational generating function is also proved: $F=P/Q$, where $Q=1-\sum_{i=1}^d c_i x^i$ and
\[P(x)=3x-10x^{2}+6x^{3}.\]
\Needspace{7\baselineskip}
\subsubsection{A243514}
\label{round10:r10entry:A243514}
Parameters: $q=6$, $K=2$, $s=2$, and $\mathcal B=\{010,011\}$. The source counts words of length $n+2$ for $n\geq1$.  Its empirical recurrence has order $d=4$ and is proved for every $n\geq 5$, with coefficients
\[(c_1,\ldots,c_d)=(8,-21,22,-8).\]
The complete matrix has $D=11$ states, with $|W|=6$ reachable states. All $66$ coordinates of $v_0,\ldots,v_{5}$ are retained. The residual in \eqref{round10:eq:r10arr-vector} is zero in all $11$ coordinates; the initial support lies in the verified closed set $W$. Thus \eqref{round10:eq:r10arr-rec} follows for every asserted index. The first three source values are $(3,5,10)$.
Source: \url{https://oeis.org/A243514}; column 5 of \url{https://oeis.org/A243519}. The primary recurrence and the relevant parent formula block are explicitly empirical.
The source rational generating function is also proved: $F=P/Q$, where $Q=1-\sum_{i=1}^d c_i x^i$ and
\[P(x)=3x-19x^{2}+33x^{3}-16x^{4}.\]
\Needspace{7\baselineskip}
\subsubsection{A243515}
\label{round10:r10entry:A243515}
Parameters: $q=7$, $K=2$, $s=2$, and $\mathcal B=\{010,011\}$. The source counts words of length $n+2$ for $n\geq1$.  Its empirical recurrence has order $d=5$ and is proved for every $n\geq 6$, with coefficients
\[(c_1,\ldots,c_d)=(12,-52,102,-91,30).\]
The complete matrix has $D=13$ states, with $|W|=7$ reachable states. All $91$ coordinates of $v_0,\ldots,v_{6}$ are retained. The residual in \eqref{round10:eq:r10arr-vector} is zero in all $13$ coordinates; the initial support lies in the verified closed set $W$. Thus \eqref{round10:eq:r10arr-rec} follows for every asserted index. The first three source values are $(3,5,10)$.
Source: \url{https://oeis.org/A243515}; column 6 of \url{https://oeis.org/A243519}. The primary recurrence and the relevant parent formula block are explicitly empirical.
The source rational generating function is also proved: $F=P/Q$, where $Q=1-\sum_{i=1}^d c_i x^i$ and
\[P(x)=3x-31x^{2}+106x^{3}-141x^{4}+60x^{5}.\]
\Needspace{7\baselineskip}
\subsubsection{A243516}
\label{round10:r10entry:A243516}
Parameters: $q=8$, $K=2$, $s=2$, and $\mathcal B=\{010,011\}$. The source counts words of length $n+2$ for $n\geq1$.  Its empirical recurrence has order $d=6$ and is proved for every $n\geq 7$, with coefficients
\[(c_1,\ldots,c_d)=(17,-111,355,-584,468,-144).\]
The complete matrix has $D=15$ states, with $|W|=8$ reachable states. All $120$ coordinates of $v_0,\ldots,v_{7}$ are retained. The residual in \eqref{round10:eq:r10arr-vector} is zero in all $15$ coordinates; the initial support lies in the verified closed set $W$. Thus \eqref{round10:eq:r10arr-rec} follows for every asserted index. The first three source values are $(3,5,10)$.
Source: \url{https://oeis.org/A243516}; column 7 of \url{https://oeis.org/A243519}. The primary recurrence and the relevant parent formula block are explicitly empirical.
The source rational generating function is also proved: $F=P/Q$, where $Q=1-\sum_{i=1}^d c_i x^i$ and
\[P(x)=3x-46x^{2}+258x^{3}-655x^{4}+739x^{5}-288x^{6}.\]
\Needspace{7\baselineskip}
\subsubsection{A243517}
\label{round10:r10entry:A243517}
Parameters: $q=9$, $K=2$, $s=2$, and $\mathcal B=\{010,011\}$. The source counts words of length $n+2$ for $n\geq1$.  Its empirical recurrence has order $d=7$ and is proved for every $n\geq 8$, with coefficients
\begin{Verbatim}[fontsize=\scriptsize,breaklines=true]
23, -212, 1010, -2669, 3887, -2878, 840
\end{Verbatim}
The complete matrix has $D=17$ states, with $|W|=9$ reachable states. All $153$ coordinates of $v_0,\ldots,v_{8}$ are retained. The residual in \eqref{round10:eq:r10arr-vector} is zero in all $17$ coordinates; the initial support lies in the verified closed set $W$. Thus \eqref{round10:eq:r10arr-rec} follows for every asserted index. The first three source values are $(3,5,10)$.
Source: \url{https://oeis.org/A243517}; column 8 of \url{https://oeis.org/A243519}. The primary recurrence and the relevant parent formula block are explicitly empirical.
The source rational generating function is also proved: $F=P/Q$, where $Q=1-\sum_{i=1}^d c_i x^i$ and
\[P(x)=3x-64x^{2}+531x^{3}-2175x^{4}+4579x^{5}-4607x^{6}+1680x^{7}.\]
\Needspace{7\baselineskip}
\subsubsection{A243518}
\label{round10:r10entry:A243518}
Parameters: $q=10$, $K=2$, $s=2$, and $\mathcal B=\{010,011\}$. The source counts words of length $n+2$ for $n\geq1$.  Its empirical recurrence has order $d=8$ and is proved for every $n\geq 9$, with coefficients
\begin{Verbatim}[fontsize=\scriptsize,breaklines=true]
30, -372, 2478, -9639, 22260, -29588, 20592, -5760
\end{Verbatim}
The complete matrix has $D=19$ states, with $|W|=10$ reachable states. All $190$ coordinates of $v_0,\ldots,v_{9}$ are retained. The residual in \eqref{round10:eq:r10arr-vector} is zero in all $19$ coordinates; the initial support lies in the verified closed set $W$. Thus \eqref{round10:eq:r10arr-rec} follows for every asserted index. The first three source values are $(3,5,10)$.
Source: \url{https://oeis.org/A243518}; column 9 of \url{https://oeis.org/A243519}. The primary recurrence and the relevant parent formula block are explicitly empirical.
The source rational generating function is also proved: $F=P/Q$, where $Q=1-\sum_{i=1}^d c_i x^i$ and
\[P(x)=3x-85x^{2}+976x^{3}-5849x^{4}+19574x^{5}-36095x^{6}+33305x^{7}-11520x^{8}.\]
\Needspace{7\baselineskip}
\subsubsection{A243600}
\label{round10:r10entry:A243600}
Parameters: $q=4$, $K=2$, $s=2$, and $\mathcal B=\{010\}$. The source counts words of length $n+2$ for $n\geq1$.  Its empirical recurrence has order $d=5$ and is proved for every $n\geq 6$, with coefficients
\[(c_1,\ldots,c_d)=(5,-5,-3,3,1).\]
The complete matrix has $D=7$ states, with $|W|=7$ reachable states. All $49$ coordinates of $v_0,\ldots,v_{6}$ are retained. The residual in \eqref{round10:eq:r10arr-vector} is zero in all $7$ coordinates; the initial support lies in the verified closed set $W$. Thus \eqref{round10:eq:r10arr-rec} follows for every asserted index. The first three source values are $(4,10,28)$.
Source: \url{https://oeis.org/A243600}; column 3 of \url{https://oeis.org/A243607}. The primary recurrence and the relevant parent formula block are explicitly empirical.
The source rational generating function is also proved: $F=P/Q$, where $Q=1-\sum_{i=1}^d c_i x^i$ and
\[P(x)=4x-10x^{2}-2x^{3}+7x^{4}+2x^{5}.\]
\Needspace{7\baselineskip}
\subsubsection{A243601}
\label{round10:r10entry:A243601}
Parameters: $q=5$, $K=2$, $s=2$, and $\mathcal B=\{010\}$. The source counts words of length $n+2$ for $n\geq1$.  Its empirical recurrence has order $d=7$ and is proved for every $n\geq 8$, with coefficients
\begin{Verbatim}[fontsize=\scriptsize,breaklines=true]
8, -18, 5, 17, -4, -6, -1
\end{Verbatim}
The complete matrix has $D=9$ states, with $|W|=9$ reachable states. All $81$ coordinates of $v_0,\ldots,v_{8}$ are retained. The residual in \eqref{round10:eq:r10arr-vector} is zero in all $9$ coordinates; the initial support lies in the verified closed set $W$. Thus \eqref{round10:eq:r10arr-rec} follows for every asserted index. The first three source values are $(4,10,29)$.
Source: \url{https://oeis.org/A243601}; column 4 of \url{https://oeis.org/A243607}. The primary recurrence and the relevant parent formula block are explicitly empirical.
The source rational generating function is also proved: $F=P/Q$, where $Q=1-\sum_{i=1}^d c_i x^i$ and
\[P(x)=4x-22x^{2}+21x^{3}+24x^{4}-15x^{5}-13x^{6}-2x^{7}.\]
\Needspace{7\baselineskip}
\subsubsection{A243602}
\label{round10:r10entry:A243602}
Parameters: $q=6$, $K=2$, $s=2$, and $\mathcal B=\{010\}$. The source counts words of length $n+2$ for $n\geq1$.  Its empirical recurrence has order $d=9$ and is proved for every $n\geq 10$, with coefficients
\begin{Verbatim}[fontsize=\scriptsize,breaklines=true]
12, -48, 65, 18, -74, -9, 26, 10, 1
\end{Verbatim}
The complete matrix has $D=11$ states, with $|W|=11$ reachable states. All $121$ coordinates of $v_0,\ldots,v_{10}$ are retained. The residual in \eqref{round10:eq:r10arr-vector} is zero in all $11$ coordinates; the initial support lies in the verified closed set $W$. Thus \eqref{round10:eq:r10arr-rec} follows for every asserted index. The first three source values are $(4,10,29)$.
Source: \url{https://oeis.org/A243602}; column 5 of \url{https://oeis.org/A243607}. The primary recurrence and the relevant parent formula block are explicitly empirical.
The source rational generating function is also proved: $F=P/Q$, where $Q=1-\sum_{i=1}^d c_i x^i$ and
\[P(x)=4x-38x^{2}+101x^{3}-31x^{4}-129x^{5}+18x^{6}+63x^{7}+21x^{8}+2x^{9}.\]
\Needspace{7\baselineskip}
\subsubsection{A243603}
\label{round10:r10entry:A243603}
Parameters: $q=7$, $K=2$, $s=2$, and $\mathcal B=\{010\}$. The source counts words of length $n+2$ for $n\geq1$.  Its empirical recurrence has order $d=11$ and is proved for every $n\geq 12$, with coefficients
\begin{Verbatim}[fontsize=\scriptsize,breaklines=true]
17, -106, 286, -245, -231, 329, 147, -106, -74, -15, -1
\end{Verbatim}
The complete matrix has $D=13$ states, with $|W|=13$ reachable states. All $169$ coordinates of $v_0,\ldots,v_{12}$ are retained. The residual in \eqref{round10:eq:r10arr-vector} is zero in all $13$ coordinates; the initial support lies in the verified closed set $W$. Thus \eqref{round10:eq:r10arr-rec} follows for every asserted index. The first three source values are $(4,10,29)$.
Source: \url{https://oeis.org/A243603}; column 6 of \url{https://oeis.org/A243607}. The primary recurrence and the relevant parent formula block are explicitly empirical.
\Needspace{7\baselineskip}
\subsubsection{A243604}
\label{round10:r10entry:A243604}
Parameters: $q=8$, $K=2$, $s=2$, and $\mathcal B=\{010\}$. The source counts words of length $n+2$ for $n\geq1$.  Its empirical recurrence has order $d=13$ and is proved for every $n\geq 14$, with coefficients
\begin{Verbatim}[fontsize=\scriptsize,breaklines=true]
23, -206, 894, -1814, 902, 1940, -1536, -1306, 394, 526, 162, 21, 1
\end{Verbatim}
The complete matrix has $D=15$ states, with $|W|=15$ reachable states. All $225$ coordinates of $v_0,\ldots,v_{14}$ are retained. The residual in \eqref{round10:eq:r10arr-vector} is zero in all $15$ coordinates; the initial support lies in the verified closed set $W$. Thus \eqref{round10:eq:r10arr-rec} follows for every asserted index. The first three source values are $(4,10,29)$.
Source: \url{https://oeis.org/A243604}; column 7 of \url{https://oeis.org/A243607}. The primary recurrence and the relevant parent formula block are explicitly empirical.
\Needspace{7\baselineskip}
\subsubsection{A243605}
\label{round10:r10entry:A243605}
Parameters: $q=9$, $K=2$, $s=2$, and $\mathcal B=\{010\}$. The source counts words of length $n+2$ for $n\geq1$.  Its empirical recurrence has order $d=15$ and is proved for every $n\geq 16$, with coefficients
\begin{Verbatim}[fontsize=\scriptsize,breaklines=true]
30, -365, 2297, -7771, 12461, -2366, -15807, 7237, 10936, -841, -3849, -1625, -307,
-28, -1
\end{Verbatim}
The complete matrix has $D=17$ states, with $|W|=17$ reachable states. All $289$ coordinates of $v_0,\ldots,v_{16}$ are retained. The residual in \eqref{round10:eq:r10arr-vector} is zero in all $17$ coordinates; the initial support lies in the verified closed set $W$. Thus \eqref{round10:eq:r10arr-rec} follows for every asserted index. The first three source values are $(4,10,29)$.
Source: \url{https://oeis.org/A243605}; column 8 of \url{https://oeis.org/A243607}. The primary recurrence and the relevant parent formula block are explicitly empirical.
\Needspace{7\baselineskip}
\subsubsection{A243606}
\label{round10:r10entry:A243606}
Parameters: $q=10$, $K=2$, $s=2$, and $\mathcal B=\{010\}$. The source counts words of length $n+2$ for $n\geq1$.  Its empirical recurrence has order $d=17$ and is proved for every $n\geq 18$, with coefficients
\begin{Verbatim}[fontsize=\scriptsize,breaklines=true]
38, -603, 5165, -25592, 71533, -92714, -10048, 134330, -29886, -94396, -8366, 29547,
16394, 4033, 529, 36, 1
\end{Verbatim}
The complete matrix has $D=19$ states, with $|W|=19$ reachable states. All $361$ coordinates of $v_0,\ldots,v_{18}$ are retained. The residual in \eqref{round10:eq:r10arr-vector} is zero in all $19$ coordinates; the initial support lies in the verified closed set $W$. Thus \eqref{round10:eq:r10arr-rec} follows for every asserted index. The first three source values are $(4,10,29)$.
Source: \url{https://oeis.org/A243606}; column 9 of \url{https://oeis.org/A243607}. The primary recurrence and the relevant parent formula block are explicitly empirical.
\Needspace{7\baselineskip}
\subsubsection{A243635}
\label{round10:r10entry:A243635}
Parameters: $q=5$, $K=2$, $s=2$, and $\mathcal B=\{012\}$. The source counts words of length $n+2$ for $n\geq1$.  Its empirical recurrence has order $d=6$ and is proved for every $n\geq 7$, with coefficients
\[(c_1,\ldots,c_d)=(7,-14,0,21,-7,-6).\]
The complete matrix has $D=9$ states, with $|W|=9$ reachable states. All $72$ coordinates of $v_0,\ldots,v_{7}$ are retained. The residual in \eqref{round10:eq:r10arr-vector} is zero in all $9$ coordinates; the initial support lies in the verified closed set $W$. Thus \eqref{round10:eq:r10arr-rec} follows for every asserted index. The first three source values are $(4,9,21)$.
Source: \url{https://oeis.org/A243635}; column 4 of \url{https://oeis.org/A243641}. The primary recurrence and the relevant parent formula block are explicitly empirical.
The source rational generating function is also proved: $F=P/Q$, where $Q=1-\sum_{i=1}^d c_i x^i$ and
\[P(x)=4x-19x^{2}+14x^{3}+30x^{4}-20x^{5}-12x^{6}.\]
\Needspace{7\baselineskip}
\subsubsection{A243636}
\label{round10:r10entry:A243636}
Parameters: $q=6$, $K=2$, $s=2$, and $\mathcal B=\{012\}$. The source counts words of length $n+2$ for $n\geq1$.  Its empirical recurrence has order $d=8$ and is proved for every $n\geq 9$, with coefficients
\begin{Verbatim}[fontsize=\scriptsize,breaklines=true]
9, -25, 7, 64, -54, -48, 32, 16
\end{Verbatim}
The complete matrix has $D=11$ states, with $|W|=11$ reachable states. All $110$ coordinates of $v_0,\ldots,v_{9}$ are retained. The residual in \eqref{round10:eq:r10arr-vector} is zero in all $11$ coordinates; the initial support lies in the verified closed set $W$. Thus \eqref{round10:eq:r10arr-rec} follows for every asserted index. The first three source values are $(4,9,21)$.
Source: \url{https://oeis.org/A243636}; column 5 of \url{https://oeis.org/A243641}. The primary recurrence and the relevant parent formula block are explicitly empirical.
The source rational generating function is also proved: $F=P/Q$, where $Q=1-\sum_{i=1}^d c_i x^i$ and
\[P(x)=4x-27x^{2}+40x^{3}+59x^{4}-126x^{5}-51x^{6}+80x^{7}+32x^{8}.\]
\Needspace{7\baselineskip}
\subsubsection{A243637}
\label{round10:r10entry:A243637}
Parameters: $q=7$, $K=2$, $s=2$, and $\mathcal B=\{012\}$. The source counts words of length $n+2$ for $n\geq1$.  Its empirical recurrence has order $d=10$ and is proved for every $n\geq 11$, with coefficients
\begin{Verbatim}[fontsize=\scriptsize,breaklines=true]
11, -39, 21, 151, -217, -181, 339, 130, -154, -60
\end{Verbatim}
The complete matrix has $D=13$ states, with $|W|=13$ reachable states. All $156$ coordinates of $v_0,\ldots,v_{11}$ are retained. The residual in \eqref{round10:eq:r10arr-vector} is zero in all $13$ coordinates; the initial support lies in the verified closed set $W$. Thus \eqref{round10:eq:r10arr-rec} follows for every asserted index. The first three source values are $(4,9,21)$.
Source: \url{https://oeis.org/A243637}; column 6 of \url{https://oeis.org/A243641}. The primary recurrence and the relevant parent formula block are explicitly empirical.
The source rational generating function is also proved: $F=P/Q$, where $Q=1-\sum_{i=1}^d c_i x^i$ and
\[\begin{aligned}P(x)&=4x-35x^{2}+78x^{3}+87x^{4}-408x^{5}-16x^{6}+668x^{7}+57x^{8}-368x^{9}\\ &\quad -120x^{10}.\end{aligned}\]
\Needspace{7\baselineskip}
\subsubsection{A243638}
\label{round10:r10entry:A243638}
Parameters: $q=8$, $K=2$, $s=2$, and $\mathcal B=\{012\}$. The source counts words of length $n+2$ for $n\geq1$.  Its empirical recurrence has order $d=12$ and is proved for every $n\geq 13$, with coefficients
\begin{Verbatim}[fontsize=\scriptsize,breaklines=true]
13, -56, 44, 305, -633, -476, 1772, 308, -2060, -368, 864, 288
\end{Verbatim}
The complete matrix has $D=15$ states, with $|W|=15$ reachable states. All $210$ coordinates of $v_0,\ldots,v_{13}$ are retained. The residual in \eqref{round10:eq:r10arr-vector} is zero in all $15$ coordinates; the initial support lies in the verified closed set $W$. Thus \eqref{round10:eq:r10arr-rec} follows for every asserted index. The first three source values are $(4,9,21)$.
Source: \url{https://oeis.org/A243638}; column 7 of \url{https://oeis.org/A243641}. The primary recurrence and the relevant parent formula block are explicitly empirical.
The source rational generating function is also proved: $F=P/Q$, where $Q=1-\sum_{i=1}^d c_i x^i$ and
\[\begin{aligned}P(x)&=4x-43x^{2}+128x^{3}+106x^{4}-976x^{5}+392x^{6}+2696x^{7}-1239x^{8}\\ &\quad -3714x^{9}+363x^{10}+2016x^{11}+576x^{12}.\end{aligned}\]
\Needspace{7\baselineskip}
\subsubsection{A243640}
\label{round10:r10entry:A243640}
Parameters: $q=10$, $K=2$, $s=2$, and $\mathcal B=\{012\}$. The source counts words of length $n+2$ for $n\geq1$.  Its empirical recurrence has order $d=16$ and is proved for every $n\geq 17$, with coefficients
\begin{Verbatim}[fontsize=\scriptsize,breaklines=true]
17, -99, 125, 931, -3185, -1785, 19063, -6754, -56936, 32132, 99764, -41096, -100896,
5152, 42048, 11520
\end{Verbatim}
The complete matrix has $D=19$ states, with $|W|=19$ reachable states. All $342$ coordinates of $v_0,\ldots,v_{17}$ are retained. The residual in \eqref{round10:eq:r10arr-vector} is zero in all $19$ coordinates; the initial support lies in the verified closed set $W$. Thus \eqref{round10:eq:r10arr-rec} follows for every asserted index. The first three source values are $(4,9,21)$.
Source: \url{https://oeis.org/A243640}; column 9 of \url{https://oeis.org/A243641}. The primary recurrence and the relevant parent formula block are explicitly empirical.
\Needspace{7\baselineskip}
\subsubsection{A243724}
\label{round10:r10entry:A243724}
Parameters: $q=6$, $K=2$, $s=2$, and $\mathcal B=\{000,012\}$. The source counts words of length $n+2$ for $n\geq1$.  Its empirical recurrence has order $d=7$ and is proved for every $n\geq 8$, with coefficients
\begin{Verbatim}[fontsize=\scriptsize,breaklines=true]
5, 0, -29, 20, 51, -31, -30
\end{Verbatim}
The complete matrix has $D=11$ states, with $|W|=11$ reachable states. All $99$ coordinates of $v_0,\ldots,v_{8}$ are retained. The residual in \eqref{round10:eq:r10arr-vector} is zero in all $11$ coordinates; the initial support lies in the verified closed set $W$. Thus \eqref{round10:eq:r10arr-rec} follows for every asserted index. The first three source values are $(3,6,12)$.
Source: \url{https://oeis.org/A243724}; column 5 of \url{https://oeis.org/A243729}. The primary recurrence and the relevant parent formula block are explicitly empirical.
The source rational generating function is also proved: $F=P/Q$, where $Q=1-\sum_{i=1}^d c_i x^i$ and
\[P(x)=3x-9x^{2}-18x^{3}+52x^{4}+42x^{5}-74x^{6}-49x^{7}.\]
\Needspace{7\baselineskip}
\subsubsection{A243727}
\label{round10:r10entry:A243727}
Parameters: $q=9$, $K=2$, $s=2$, and $\mathcal B=\{000,012\}$. The source counts words of length $n+2$ for $n\geq1$.  Its empirical recurrence has order $d=13$ and is proved for every $n\geq 14$, with coefficients
\begin{Verbatim}[fontsize=\scriptsize,breaklines=true]
8, 0, -139, 196, 960, -1827, -3456, 7076, 7239, -12860, -9148, 9072, 5760
\end{Verbatim}
The complete matrix has $D=17$ states, with $|W|=17$ reachable states. All $255$ coordinates of $v_0,\ldots,v_{14}$ are retained. The residual in \eqref{round10:eq:r10arr-vector} is zero in all $17$ coordinates; the initial support lies in the verified closed set $W$. Thus \eqref{round10:eq:r10arr-rec} follows for every asserted index. The first three source values are $(3,6,12)$.
Source: \url{https://oeis.org/A243727}; column 8 of \url{https://oeis.org/A243729}. The primary recurrence and the relevant parent formula block are explicitly empirical.
The source rational generating function is also proved: $F=P/Q$, where $Q=1-\sum_{i=1}^d c_i x^i$ and
\[\begin{aligned}P(x)&=3x-18x^{2}-36x^{3}+346x^{4}+99x^{5}-2696x^{6}+175x^{7}+10799x^{8}\\ &\quad -195x^{9}-22430x^{10}-4582x^{11}+19296x^{12}+9401x^{13}.\end{aligned}\]
\Needspace{7\baselineskip}
\subsubsection{A244177}
\label{round10:r10entry:A244177}
Parameters: $q=3$, $K=4$, $s=4$, and $\mathcal B=\{00000,01010,01110\}$. The source counts words of length $n+4$ for $n\geq1$.  Its empirical recurrence has order $d=6$ and is proved for every $n\geq 7$, with coefficients
\[(c_1,\ldots,c_d)=(1,3,4,5,3,2).\]
The complete matrix has $D=23$ states, with $|W|=23$ reachable states. All $184$ coordinates of $v_0,\ldots,v_{7}$ are retained. The residual in \eqref{round10:eq:r10arr-vector} is zero in all $23$ coordinates; the initial support lies in the verified closed set $W$. Thus \eqref{round10:eq:r10arr-rec} follows for every asserted index. The first three source values are $(38,108,306)$.
Source: \url{https://oeis.org/A244177}; column 2 of \url{https://oeis.org/A244185}. The primary recurrence and the relevant parent formula block are explicitly empirical.
The source rational generating function is also proved: $F=P/Q$, where $Q=1-\sum_{i=1}^d c_i x^i$ and
\[P(x)=38x+70x^{2}+84x^{3}+85x^{4}+50x^{5}+27x^{6}.\]
\Needspace{7\baselineskip}
\subsubsection{A244178}
\label{round10:r10entry:A244178}
Parameters: $q=4$, $K=4$, $s=4$, and $\mathcal B=\{00000,01010,01110\}$. The source counts words of length $n+4$ for $n\geq1$.  Its empirical recurrence has order $d=12$ and is proved for every $n\geq 13$, with coefficients
\begin{Verbatim}[fontsize=\scriptsize,breaklines=true]
2, 6, 6, 3, -10, -23, -34, -36, -28, -19, -8, -3
\end{Verbatim}
The complete matrix has $D=38$ states, with $|W|=38$ reachable states. All $532$ coordinates of $v_0,\ldots,v_{13}$ are retained. The residual in \eqref{round10:eq:r10arr-vector} is zero in all $38$ coordinates; the initial support lies in the verified closed set $W$. Thus \eqref{round10:eq:r10arr-rec} follows for every asserted index. The first three source values are $(48,173,650)$.
Source: \url{https://oeis.org/A244178}; column 3 of \url{https://oeis.org/A244185}. The primary recurrence and the relevant parent formula block are explicitly empirical.
\Needspace{7\baselineskip}
\subsubsection{A244179}
\label{round10:r10entry:A244179}
Parameters: $q=5$, $K=4$, $s=4$, and $\mathcal B=\{00000,01010,01110\}$. The source counts words of length $n+4$ for $n\geq1$.  Its empirical recurrence has order $d=18$ and is proved for every $n\geq 19$, with coefficients
\begin{Verbatim}[fontsize=\scriptsize,breaklines=true]
4, 8, -5, -36, -86, -106, -52, 102, 337, 559, 701, 701, 582, 403, 227, 105, 34, 8
\end{Verbatim}
The complete matrix has $D=53$ states, with $|W|=53$ reachable states. All $1060$ coordinates of $v_0,\ldots,v_{19}$ are retained. The residual in \eqref{round10:eq:r10arr-vector} is zero in all $53$ coordinates; the initial support lies in the verified closed set $W$. Thus \eqref{round10:eq:r10arr-rec} follows for every asserted index. The first three source values are $(49,188,790)$.
Source: \url{https://oeis.org/A244179}; column 4 of \url{https://oeis.org/A244185}. The primary recurrence and the relevant parent formula block are explicitly empirical.
\Needspace{7\baselineskip}
\subsubsection{A244180}
\label{round10:r10entry:A244180}
Parameters: $q=6$, $K=4$, $s=4$, and $\mathcal B=\{00000,01010,01110\}$. The source counts words of length $n+4$ for $n\geq1$.  Its empirical recurrence has order $d=24$ and is proved for every $n\geq 25$, with coefficients
\begin{Verbatim}[fontsize=\scriptsize,breaklines=true]
7, 4, -44, -104, -72, 290, 1095, 2087, 2522, 1293, -2295, -8084, -14734, -20253,
-22955, -22128, -18438, -13301, -8283, -4402, -1946, -694, -179, -30
\end{Verbatim}
The complete matrix has $D=68$ states, with $|W|=68$ reachable states. All $1768$ coordinates of $v_0,\ldots,v_{25}$ are retained. The residual in \eqref{round10:eq:r10arr-vector} is zero in all $68$ coordinates; the initial support lies in the verified closed set $W$. Thus \eqref{round10:eq:r10arr-rec} follows for every asserted index. The first three source values are $(49,189,811)$.
Source: \url{https://oeis.org/A244180}; column 5 of \url{https://oeis.org/A244185}. The primary recurrence and the relevant parent formula block are explicitly empirical.
\Needspace{7\baselineskip}
\subsubsection{A244181}
\label{round10:r10entry:A244181}
Parameters: $q=7$, $K=4$, $s=4$, and $\mathcal B=\{00000,01010,01110\}$. The source counts words of length $n+4$ for $n\geq1$.  Its empirical recurrence has order $d=30$ and is proved for every $n\geq 31$, with coefficients
\begin{Verbatim}[fontsize=\scriptsize,breaklines=true]
11, -14, -109, -50, 561, 1905, 2485, -1765, -16073, -42234, -71345, -80685, -38274,
82534, 284947, 538989, 786421, 962575, 1022055, 956487, 794930, 587882, 386232,
224067, 113614, 49515, 18041, 5264, 1116, 144
\end{Verbatim}
The complete matrix has $D=83$ states, with $|W|=83$ reachable states. All $2656$ coordinates of $v_0,\ldots,v_{31}$ are retained. The residual in \eqref{round10:eq:r10arr-vector} is zero in all $83$ coordinates; the initial support lies in the verified closed set $W$. Thus \eqref{round10:eq:r10arr-rec} follows for every asserted index. The first three source values are $(49,189,812)$.
Source: \url{https://oeis.org/A244181}; column 6 of \url{https://oeis.org/A244185}. The primary recurrence and the relevant parent formula block are explicitly empirical.
\Needspace{7\baselineskip}
\subsubsection{A244182}
\label{round10:r10entry:A244182}
Parameters: $q=8$, $K=4$, $s=4$, and $\mathcal B=\{00000,01010,01110\}$. The source counts words of length $n+4$ for $n\geq1$.  Its empirical recurrence has order $d=36$ and is proved for every $n\geq 37$, with coefficients
\begin{Verbatim}[fontsize=\scriptsize,breaklines=true]
16, -57, -146, 471, 2030, 1606, -9926, -40696, -72180, -23242, 262788, 948320,
2028394, 3106669, 3270400, 1237472, -4174278, -13394281, -25599412, -38689072,
-49852288, -56526936, -57307380, -52390492, -43372994, -32559403, -22139372,
-13590504, -7488422, -3670971, -1580476, -586225, -181844, -44951, -8042, -840
\end{Verbatim}
The complete matrix has $D=98$ states, with $|W|=98$ reachable states. All $3724$ coordinates of $v_0,\ldots,v_{37}$ are retained. The residual in \eqref{round10:eq:r10arr-vector} is zero in all $98$ coordinates; the initial support lies in the verified closed set $W$. Thus \eqref{round10:eq:r10arr-rec} follows for every asserted index. The first three source values are $(49,189,812)$.
Source: \url{https://oeis.org/A244182}; column 7 of \url{https://oeis.org/A244185}. The primary recurrence and the relevant parent formula block are explicitly empirical.
\Needspace{7\baselineskip}
\subsubsection{A244183}
\label{round10:r10entry:A244183}
Parameters: $q=9$, $K=4$, $s=4$, and $\mathcal B=\{00000,01010,01110\}$. The source counts words of length $n+4$ for $n\geq1$.  Its empirical recurrence has order $d=42$ and is proved for every $n\geq 43$, with coefficients
\begin{Verbatim}[fontsize=\scriptsize,breaklines=true]
22, -139, 1, 1784, 1931, -12600, -50646, -48753, 215387, 1038438, 2236793, 1983939,
-4361986, -23908864, -62445686, -116557476, -164083468, -158681803, -35333656,
270032838, 791193237, 1504649448, 2319758240, 3094945131, 3677150092, 3949611368,
3868169748, 3471359234, 2861776683, 2168940385, 1510165759, 963934762, 562027256,
297750070, 142267724, 60680383, 22773619, 7366519, 1990304, 426740, 65808, 5760
\end{Verbatim}
The complete matrix has $D=113$ states, with $|W|=113$ reachable states. All $4972$ coordinates of $v_0,\ldots,v_{43}$ are retained. The residual in \eqref{round10:eq:r10arr-vector} is zero in all $113$ coordinates; the initial support lies in the verified closed set $W$. Thus \eqref{round10:eq:r10arr-rec} follows for every asserted index. The first three source values are $(49,189,812)$.
Source: \url{https://oeis.org/A244183}; column 8 of \url{https://oeis.org/A244185}. The primary recurrence and the relevant parent formula block are explicitly empirical.
\Needspace{7\baselineskip}
\subsubsection{A244184}
\label{round10:r10entry:A244184}
Parameters: $q=10$, $K=4$, $s=4$, and $\mathcal B=\{00000,01010,01110\}$. The source counts words of length $n+4$ for $n\geq1$.  Its empirical recurrence has order $d=48$ and is proved for every $n\geq 49$, with coefficients
\begin{Verbatim}[fontsize=\scriptsize,breaklines=true]
29, -277, 655, 3400, -7802, -49606, -30159, 403108, 1503590, 1635899, -5652490,
-31339466, -77219594, -98298345, 54151570, 669953408, 2120355686, 4618747926,
7785441378, 10142532705, 8820633444, -184894715, -21013605262, -56342602223,
-105733609417, -164683705023, -224995549990, -276580511002, -310166430329,
-319929413204, -305057789374, -269700758984, -221428750153, -168899208188,
-119626976395, -78560246496, -47719787532, -26718125970, -13723716584, -6426330933,
-2720546834, -1029660492, -343102672, -98490588, -23573519, -4460436, -603108, -45360
\end{Verbatim}
The complete matrix has $D=128$ states, with $|W|=128$ reachable states. All $6400$ coordinates of $v_0,\ldots,v_{49}$ are retained. The residual in \eqref{round10:eq:r10arr-vector} is zero in all $128$ coordinates; the initial support lies in the verified closed set $W$. Thus \eqref{round10:eq:r10arr-rec} follows for every asserted index. The first three source values are $(49,189,812)$.
Source: \url{https://oeis.org/A244184}; column 9 of \url{https://oeis.org/A244185}. The primary recurrence and the relevant parent formula block are explicitly empirical.
\Needspace{7\baselineskip}
\subsubsection{A244694}
\label{round10:r10entry:A244694}
Parameters: $q=3$, $K=4$, $s=4$, and $\mathcal B=\{01010,01110\}$. The source counts words of length $n+4$ for $n\geq1$.  Its empirical recurrence has order $d=7$ and is proved for every $n\geq 8$, with coefficients
\begin{Verbatim}[fontsize=\scriptsize,breaklines=true]
3, 0, -1, 0, -1, 1, -1
\end{Verbatim}
The complete matrix has $D=23$ states, with $|W|=23$ reachable states. All $207$ coordinates of $v_0,\ldots,v_{8}$ are retained. The residual in \eqref{round10:eq:r10arr-vector} is zero in all $23$ coordinates; the initial support lies in the verified closed set $W$. Thus \eqref{round10:eq:r10arr-rec} follows for every asserted index. The first three source values are $(39,111,317)$.
Source: \url{https://oeis.org/A244694}; column 2 of \url{https://oeis.org/A244702}. The primary recurrence and the relevant parent formula block are explicitly empirical.
The source rational generating function is also proved: $F=P/Q$, where $Q=1-\sum_{i=1}^d c_i x^i$ and
\[P(x)=39x-6x^{2}-16x^{3}-4x^{4}-11x^{5}+9x^{6}-14x^{7}.\]
\Needspace{7\baselineskip}
\subsubsection{A244695}
\label{round10:r10entry:A244695}
Parameters: $q=4$, $K=4$, $s=4$, and $\mathcal B=\{01010,01110\}$. The source counts words of length $n+4$ for $n\geq1$.  Its empirical recurrence has order $d=13$ and is proved for every $n\geq 14$, with coefficients
\begin{Verbatim}[fontsize=\scriptsize,breaklines=true]
5, -3, -5, 1, -3, 4, -4, 1, 3, 0, 2, -1, 1
\end{Verbatim}
The complete matrix has $D=38$ states, with $|W|=38$ reachable states. All $570$ coordinates of $v_0,\ldots,v_{14}$ are retained. The residual in \eqref{round10:eq:r10arr-vector} is zero in all $38$ coordinates; the initial support lies in the verified closed set $W$. Thus \eqref{round10:eq:r10arr-rec} follows for every asserted index. The first three source values are $(49,176,661)$.
Source: \url{https://oeis.org/A244695}; column 3 of \url{https://oeis.org/A244702}. The primary recurrence and the relevant parent formula block are explicitly empirical.
\Needspace{7\baselineskip}
\subsubsection{A244696}
\label{round10:r10entry:A244696}
Parameters: $q=5$, $K=4$, $s=4$, and $\mathcal B=\{01010,01110\}$. The source counts words of length $n+4$ for $n\geq1$.  Its empirical recurrence has order $d=19$ and is proved for every $n\geq 20$, with coefficients
\begin{Verbatim}[fontsize=\scriptsize,breaklines=true]
8, -14, -9, 15, 7, 22, -9, 9, 13, -9, 0, -17, 0, -5, -6, -1, -3, 1, -1
\end{Verbatim}
The complete matrix has $D=53$ states, with $|W|=53$ reachable states. All $1113$ coordinates of $v_0,\ldots,v_{20}$ are retained. The residual in \eqref{round10:eq:r10arr-vector} is zero in all $53$ coordinates; the initial support lies in the verified closed set $W$. Thus \eqref{round10:eq:r10arr-rec} follows for every asserted index. The first three source values are $(50,191,801)$.
Source: \url{https://oeis.org/A244696}; column 4 of \url{https://oeis.org/A244702}. The primary recurrence and the relevant parent formula block are explicitly empirical.
\Needspace{7\baselineskip}
\subsubsection{A244697}
\label{round10:r10entry:A244697}
Parameters: $q=6$, $K=4$, $s=4$, and $\mathcal B=\{01010,01110\}$. The source counts words of length $n+4$ for $n\geq1$.  Its empirical recurrence has order $d=25$ and is proved for every $n\geq 26$, with coefficients
\begin{Verbatim}[fontsize=\scriptsize,breaklines=true]
12, -41, 16, 81, 15, -19, -158, -58, -45, -106, -38, -53, 85, 57, 46, 67, 37, 56, 15,
15, 10, 3, 4, -1, 1
\end{Verbatim}
The complete matrix has $D=68$ states, with $|W|=68$ reachable states. All $1836$ coordinates of $v_0,\ldots,v_{26}$ are retained. The residual in \eqref{round10:eq:r10arr-vector} is zero in all $68$ coordinates; the initial support lies in the verified closed set $W$. Thus \eqref{round10:eq:r10arr-rec} follows for every asserted index. The first three source values are $(50,192,822)$.
Source: \url{https://oeis.org/A244697}; column 5 of \url{https://oeis.org/A244702}. The primary recurrence and the relevant parent formula block are explicitly empirical.
\Needspace{7\baselineskip}
\subsubsection{A244698}
\label{round10:r10entry:A244698}
Parameters: $q=7$, $K=4$, $s=4$, and $\mathcal B=\{01010,01110\}$. The source counts words of length $n+4$ for $n\geq1$.  Its empirical recurrence has order $d=31$ and is proved for every $n\geq 32$, with coefficients
\begin{Verbatim}[fontsize=\scriptsize,breaklines=true]
17, -95, 160, 159, -307, -498, -480, 692, 1026, 917, 1051, 886, 1139, 217, -344, -582,
-767, -618, -822, -592, -433, -300, -169, -146, -50, -34, -15, -6, -5, 1, -1
\end{Verbatim}
The complete matrix has $D=83$ states, with $|W|=83$ reachable states. All $2739$ coordinates of $v_0,\ldots,v_{32}$ are retained. The residual in \eqref{round10:eq:r10arr-vector} is zero in all $83$ coordinates; the initial support lies in the verified closed set $W$. Thus \eqref{round10:eq:r10arr-rec} follows for every asserted index. The first three source values are $(50,192,823)$.
Source: \url{https://oeis.org/A244698}; column 6 of \url{https://oeis.org/A244702}. The primary recurrence and the relevant parent formula block are explicitly empirical.
\Needspace{7\baselineskip}
\subsubsection{A244699}
\label{round10:r10entry:A244699}
Parameters: $q=8$, $K=4$, $s=4$, and $\mathcal B=\{01010,01110\}$. The source counts words of length $n+4$ for $n\geq1$.  Its empirical recurrence has order $d=37$ and is proved for every $n\geq 38$, with coefficients
\begin{Verbatim}[fontsize=\scriptsize,breaklines=true]
23, -190, 624, -299, -1735, -321, 3257, 7830, 2362, -6411, -12773, -15698, -15096,
-18309, -14657, -7203, 1319, 7754, 9212, 11865, 11915, 10867, 8583, 5912, 4520, 2834,
1796, 975, 505, 321, 117, 65, 21, 10, 6, -1, 1
\end{Verbatim}
The complete matrix has $D=98$ states, with $|W|=98$ reachable states. All $3822$ coordinates of $v_0,\ldots,v_{38}$ are retained. The residual in \eqref{round10:eq:r10arr-vector} is zero in all $98$ coordinates; the initial support lies in the verified closed set $W$. Thus \eqref{round10:eq:r10arr-rec} follows for every asserted index. The first three source values are $(50,192,823)$.
Source: \url{https://oeis.org/A244699}; column 7 of \url{https://oeis.org/A244702}. The primary recurrence and the relevant parent formula block are explicitly empirical.
\Needspace{7\baselineskip}
\subsubsection{A244700}
\label{round10:r10entry:A244700}
Parameters: $q=9$, $K=4$, $s=4$, and $\mathcal B=\{01010,01110\}$. The source counts words of length $n+4$ for $n\geq1$.  Its empirical recurrence has order $d=43$ and is proved for every $n\geq 44$, with coefficients
\begin{Verbatim}[fontsize=\scriptsize,breaklines=true]
30, -343, 1785, -3418, -2904, 10172, 18515, -1513, -70620, -94811, -32588, 94041,
211477, 265558, 316407, 329408, 286665, 168596, 19643, -85855, -164056, -206321,
-227539, -216968, -179819, -141979, -104021, -74269, -48384, -29216, -17666, -9756,
-5441, -2576, -1217, -624, -231, -111, -28, -15, -7, 1, -1
\end{Verbatim}
The complete matrix has $D=113$ states, with $|W|=113$ reachable states. All $5085$ coordinates of $v_0,\ldots,v_{44}$ are retained. The residual in \eqref{round10:eq:r10arr-vector} is zero in all $113$ coordinates; the initial support lies in the verified closed set $W$. Thus \eqref{round10:eq:r10arr-rec} follows for every asserted index. The first three source values are $(50,192,823)$.
Source: \url{https://oeis.org/A244700}; column 8 of \url{https://oeis.org/A244702}. The primary recurrence and the relevant parent formula block are explicitly empirical.
\Needspace{7\baselineskip}
\subsubsection{A244701}
\label{round10:r10entry:A244701}
Parameters: $q=10$, $K=4$, $s=4$, and $\mathcal B=\{01010,01110\}$. The source counts words of length $n+4$ for $n\geq1$.  Its empirical recurrence has order $d=49$ and is proved for every $n\geq 50$, with coefficients
\begin{Verbatim}[fontsize=\scriptsize,breaklines=true]
38, -574, 4276, -15029, 12263, 48204, -16847, -203774, -274320, 337901, 1264190,
1574718, 434569, -1814318, -3975030, -5698180, -6798034, -7216705, -6361941, -4141269,
-1415098, 1117645, 2963339, 4238015, 4880440, 4838799, 4315267, 3525637, 2742853,
2014392, 1390851, 910679, 565426, 343188, 195749, 106036, 55321, 27440, 13678, 5894,
2555, 1107, 410, 175, 36, 21, 8, -1, 1
\end{Verbatim}
The complete matrix has $D=128$ states, with $|W|=128$ reachable states. All $6528$ coordinates of $v_0,\ldots,v_{50}$ are retained. The residual in \eqref{round10:eq:r10arr-vector} is zero in all $128$ coordinates; the initial support lies in the verified closed set $W$. Thus \eqref{round10:eq:r10arr-rec} follows for every asserted index. The first three source values are $(50,192,823)$.
Source: \url{https://oeis.org/A244701}; column 9 of \url{https://oeis.org/A244702}. The primary recurrence and the relevant parent formula block are explicitly empirical.


\section{Verification, provenance and reuse}
\label{round10:r10:verification}
The archive contains this standalone manuscript, the numbered ledger,
compact primary-source assertions, complete exact certificates, portable
programs and internal review reports. After installing the dependency
listed in \texttt{requirements.txt}, run
\begin{Verbatim}[fontsize=\small]
python3 r10_run_all.py
\end{Verbatim}
from the extracted directory. The checks require no network access and
no previous-round program. Actual versions, counts and outcomes are
recorded in the verification summary.

\subsection{Why the finite certificates prove infinite assertions}
For the word families, the state records the number of introduced letters
and the equality pattern of the last few letters. Old letters absent from
that suffix retain their exact multiplicity. This gives a finite transition
matrix directly from the definition, before any empirical recurrence is
consulted. The exact vector
\[
 z=T^b\mathbf1-\sum_{i=1}^d c_iT^{b-i}\mathbf1
\]
vanishes on a transition-closed set containing the initial support.
Multiplying by the initial row and every subsequent power of $T$ proves
every recurrence instance in the stated domain. It is not a fitted-prefix
argument. Boundary shifts between word length and OEIS index are explicit.

The 65 selected compressed matrices have 2098 states. An independent
implementation, which imports none of the compressed engine, keeps the
actual letter names. It checks all 131159 states and 1052986 transitions
in these full alphabets, including every row of the proposed compression.
Direct enumeration of short words supplies 406 additional source checks.
The main word verifier checks 1452 published terms and proves 31 source
rational generating functions without increasing the count.

For the other generating functions, exact rational identities, radical
operator identities, classical dimension formulas or defining polynomial
recursions supply all-index arguments. Every inhomogeneous boundary
polynomial is retained. The tatami specialization is derived at $z=-1$
from the defining polynomial quotient, rather than assumed from the
conjectural generating function being proved.

The arithmetic and formal-series sections give written universal proofs.
Numerical checks verify transcription, source coefficients and boundary
cases. In particular, checking many periods is not the proof of A205572;
the integral derivation argument proves all moduli and all indices at once.
Reduction of non-locally-finite moment sums is performed only after finite
Newton expansions have made the sums into rational functions.

\subsection{Reproducibility and integrity}
The delivered \texttt{.tex} embeds all proof sections and entry signatures;
it has no external \texttt{\string\input} or \texttt{\string\include}
dependencies. Compile with pdfLaTeX until references stabilize.
\texttt{r10\_build\_dossier.py} rebuilds the manuscript and catalogue.
The README lists the retained inputs and reconstruction commands.
An isolated regeneration check compares immutable mathematical outputs.

The archive manifest gives SHA-256 digests. Review reports bind the exact
proof, source, code and certificate files inspected. The consistency check
matches catalogue counts, proof labels, source offsets and retained
mathematical inputs. The layout report identifies the PDF whose rendered
pages were inspected. Timing fields and logs may change on a rerun;
the mathematical assertions and exact certificates do not depend on timing.

\subsection{Common references and limits}
Individual OEIS links, original contributors and specialized references are
given with the relevant statements. The general source is the OEIS
Foundation's \emph{On-Line Encyclopedia of Integer Sequences},
\url{https://oeis.org/}. The retained OEIS Open Appendix A comparison is
\url{https://arxiv.org/html/2608.11941v1}; its scope is explicitly limited
to the 97 unique identifiers in that appendix.

The component audits document named literature, known results removed,
duplicate formulations and bounded alias screens. Such a screen cannot
exclude every possible transformation between sequences. The resulting
count is 100 precisely specified source-assertion packages with complete
proofs or corrections. It is not a certificate of 100 worldwide-first
discoveries, and it is not an official closure of 100 OEIS entries.
% ===== END ORIGINAL ROUND 10 =====

\clearpage
\rhead{\small Identifier locator}
\section*{OEIS identifier locator}
\addcontentsline{toc}{section}{OEIS identifier locator}
This locator includes all identifiers mentioned in the original manuscript
bodies, including selected targets, supporting sequences, earlier work,
uncounted supplements, and excluded candidates. Its size is not a
selection count. The linked numbers identify the rounds containing the
identifier; each round's catalogue gives the selected assertions and
links to their proofs.
\begingroup\small
\setlength{\tabcolsep}{4pt}
\renewcommand{\arraystretch}{1.06}
\begin{longtable}{@{}l p{.12\textwidth} l p{.12\textwidth} l p{.12\textwidth} l p{.12\textwidth}@{}}
\toprule
OEIS & Rounds & OEIS & Rounds & OEIS & Rounds & OEIS & Rounds\\
\midrule\endfirsthead
\toprule OEIS & Rounds & OEIS & Rounds & OEIS & Rounds & OEIS & Rounds\\
\midrule\endhead
\bottomrule\endfoot
\oeis{A000139} & \hyperref[combined:round7]{7} & \oeis{A121860} & \hyperref[combined:round3]{3} & \oeis{A210302} & \hyperref[combined:round4]{4} & \oeis{A243041} & \hyperref[combined:round10]{10}\\
\oeis{A000159} & \hyperref[combined:round2]{2} & \oeis{A123071} & \hyperref[combined:round9]{9} & \oeis{A210303} & \hyperref[combined:round4]{4} & \oeis{A243042} & \hyperref[combined:round10]{10}\\
\oeis{A000181} & \hyperref[combined:round2]{2} & \oeis{A126089} & \hyperref[combined:round2]{2} & \oeis{A210304} & \hyperref[combined:round4]{4} & \oeis{A243043} & \hyperref[combined:round10]{10}\\
\oeis{A000185} & \hyperref[combined:round2]{2} & \oeis{A126156} & \hyperref[combined:round4]{4} & \oeis{A210305} & \hyperref[combined:round5]{5} & \oeis{A243044} & \hyperref[combined:round10]{10}\\
\oeis{A000254} & \hyperref[combined:round9]{9} & \oeis{A127473} & \hyperref[combined:round2]{2} & \oeis{A210313} & \hyperref[combined:round4]{4} & \oeis{A243045} & \hyperref[combined:round10]{10}\\
\oeis{A000364} & \hyperref[combined:round3]{3} & \oeis{A127751} & \hyperref[combined:round1]{1} & \oeis{A210314} & \hyperref[combined:round4]{4} & \oeis{A243046} & \hyperref[combined:round10]{10}\\
\oeis{A000483} & \hyperref[combined:round9]{9} & \oeis{A127752} & \hyperref[combined:round1]{1} & \oeis{A210315} & \hyperref[combined:round4]{4} & \oeis{A243255} & \hyperref[combined:round6]{6}\\
\oeis{A000594} & \hyperref[combined:round2]{2} & \oeis{A127905} & \hyperref[combined:round1]{1} & \oeis{A210316} & \hyperref[combined:round4]{4} & \oeis{A243383} & \hyperref[combined:round10]{10}\\
\oeis{A000657} & \hyperref[combined:round4]{4} & \oeis{A128057} & \hyperref[combined:round4]{4} & \oeis{A210328} & \hyperref[combined:round5]{5} & \oeis{A243384} & \hyperref[combined:round10]{10}\\
\oeis{A000670} & \hyperref[combined:round3]{3} & \oeis{A129366} & \hyperref[combined:round5]{5} & \oeis{A210329} & \hyperref[combined:round5]{5} & \oeis{A243385} & \hyperref[combined:round10]{10}\\
\oeis{A000695} & \hyperref[combined:round9]{9} & \oeis{A129454} & \hyperref[combined:round5]{5} & \oeis{A210330} & \hyperref[combined:round5]{5} & \oeis{A243386} & \hyperref[combined:round10]{10}\\
\oeis{A000781} & \hyperref[combined:round6]{6} & \oeis{A129455} & \hyperref[combined:round5]{5} & \oeis{A210348} & \hyperref[combined:round5]{5} & \oeis{A243387} & \hyperref[combined:round10]{10}\\
\oeis{A000898} & \hyperref[combined:round9]{9} & \oeis{A130540} & \hyperref[combined:round2]{2} & \oeis{A210349} & \hyperref[combined:round5]{5} & \oeis{A243388} & \hyperref[combined:round10]{10}\\
\oeis{A000986} & \hyperref[combined:round4]{4} & \oeis{A131418} & \hyperref[combined:round3]{3} & \oeis{A210350} & \hyperref[combined:round5]{5} & \oeis{A243389} & \hyperref[combined:round10]{10}\\
\oeis{A001108} & \hyperref[combined:round9]{9} & \oeis{A132393} & \hyperref[combined:round2]{2} & \oeis{A210383} & \hyperref[combined:round5]{5} & \oeis{A243513} & \hyperref[combined:round10]{10}\\
\oeis{A001109} & \hyperref[combined:round9]{9} & \oeis{A133158} & \hyperref[combined:round2]{2} & \oeis{A210384} & \hyperref[combined:round5]{5} & \oeis{A243514} & \hyperref[combined:round10]{10}\\
\oeis{A001157} & \hyperref[combined:round8]{8} & \oeis{A133205} & \hyperref[combined:round7]{7} & \oeis{A210386} & \hyperref[combined:round5]{5} & \oeis{A243515} & \hyperref[combined:round10]{10}\\
\oeis{A001222} & \hyperref[combined:round2]{2} & \oeis{A133907} & \hyperref[combined:round6]{6} & \oeis{A210397} & \hyperref[combined:round4]{4} & \oeis{A243516} & \hyperref[combined:round10]{10}\\
\oeis{A001277} & \hyperref[combined:round1]{1} & \oeis{A133910} & \hyperref[combined:round2]{2} & \oeis{A210398} & \hyperref[combined:round4]{4} & \oeis{A243517} & \hyperref[combined:round10]{10}\\
\oeis{A001278} & \hyperref[combined:round1]{1} & \oeis{A134389} & \hyperref[combined:round7]{7} & \oeis{A210400} & \hyperref[combined:round4]{4} & \oeis{A243518} & \hyperref[combined:round10]{10}\\
\oeis{A001279} & \hyperref[combined:round1]{1} & \oeis{A134416} & \hyperref[combined:round5]{5} & \oeis{A210406} & \hyperref[combined:round6]{6} & \oeis{A243519} & \hyperref[combined:round10]{10}\\
\oeis{A001280} & \hyperref[combined:round1]{1} & \oeis{A134424} & \hyperref[combined:round5]{5} & \oeis{A210407} & \hyperref[combined:round6]{6} & \oeis{A243600} & \hyperref[combined:round10]{10}\\
\oeis{A001316} & \hyperref[combined:round10]{10} & \oeis{A135925} & \hyperref[combined:round2]{2} & \oeis{A210408} & \hyperref[combined:round6]{6} & \oeis{A243601} & \hyperref[combined:round10]{10}\\
\oeis{A001497} & \hyperref[combined:round2]{2} & \oeis{A136097} & \hyperref[combined:round1]{1} & \oeis{A210409} & \hyperref[combined:round6]{6} & \oeis{A243602} & \hyperref[combined:round10]{10}\\
\oeis{A001515} & \hyperref[combined:round8]{8} & \oeis{A136123} & \hyperref[combined:round1]{1} & \oeis{A210496} & \hyperref[combined:round4]{4} & \oeis{A243603} & \hyperref[combined:round10]{10}\\
\oeis{A001711} & \hyperref[combined:round2]{2} & \oeis{A140110} & \hyperref[combined:round7]{7} & \oeis{A210826} & \hyperref[combined:round6]{6} & \oeis{A243604} & \hyperref[combined:round10]{10}\\
\oeis{A001712} & \hyperref[combined:round2]{2} & \oeis{A141117} & \hyperref[combined:round3]{3} & \oeis{A212341} & \hyperref[combined:round9]{9} & \oeis{A243605} & \hyperref[combined:round10]{10}\\
\oeis{A001722} & \hyperref[combined:round1]{1} & \oeis{A141351} & \hyperref[combined:round6]{6} & \oeis{A212824} & \hyperref[combined:round10]{10} & \oeis{A243606} & \hyperref[combined:round10]{10}\\
\oeis{A001723} & \hyperref[combined:round2]{2} & \oeis{A141353} & \hyperref[combined:round2]{2} & \oeis{A212825} & \hyperref[combined:round10]{10} & \oeis{A243607} & \hyperref[combined:round10]{10}\\
\oeis{A001724} & \hyperref[combined:round2]{2} & \oeis{A141459} & \hyperref[combined:round10]{10} & \oeis{A212826} & \hyperref[combined:round10]{10} & \oeis{A243635} & \hyperref[combined:round10]{10}\\
\oeis{A001784} & \hyperref[combined:round9]{9} & \oeis{A142243} & \hyperref[combined:round1]{1} & \oeis{A212827} & \hyperref[combined:round10]{10} & \oeis{A243636} & \hyperref[combined:round10]{10}\\
\oeis{A001840} & \hyperref[combined:round2]{2} & \oeis{A143013} & \hyperref[combined:round9]{9} & \oeis{A212829} & \hyperref[combined:round10]{10} & \oeis{A243637} & \hyperref[combined:round10]{10}\\
\oeis{A001969} & \hyperref[combined:round1]{1} & \oeis{A143020} & \hyperref[combined:round7]{7} & \oeis{A214406} & \hyperref[combined:round10]{10} & \oeis{A243638} & \hyperref[combined:round10]{10}\\
\oeis{A002026} & \hyperref[combined:round4]{4} & \oeis{A143991} & \hyperref[combined:round8]{8} & \oeis{A216971} & \hyperref[combined:round1]{1} & \oeis{A243640} & \hyperref[combined:round10]{10}\\
\oeis{A002050} & \hyperref[combined:round3]{3} & \oeis{A144575} & \hyperref[combined:round2]{2} & \oeis{A217333} & \hyperref[combined:round6]{6} & \oeis{A243641} & \hyperref[combined:round10]{10}\\
\oeis{A002293} & \hyperref[combined:round2]{2} & \oeis{A146564} & \hyperref[combined:round10]{10} & \oeis{A218045} & \hyperref[combined:round2]{2} & \oeis{A243642} & \hyperref[combined:round4]{4}\\
\oeis{A002512} & \hyperref[combined:round5]{5} & \oeis{A146955} & \hyperref[combined:round1]{1} & \oeis{A218185} & \hyperref[combined:round5]{5} & \oeis{A243724} & \hyperref[combined:round10]{10}\\
\oeis{A002605} & \hyperref[combined:round7]{7} & \oeis{A152920} & \hyperref[combined:round2]{2} & \oeis{A220119} & \hyperref[combined:round1]{1} & \oeis{A243727} & \hyperref[combined:round10]{10}\\
\oeis{A002618} & \hyperref[combined:round2]{2} & \oeis{A153197} & \hyperref[combined:round2]{2} & \oeis{A220433} & \hyperref[combined:round1]{1} & \oeis{A243729} & \hyperref[combined:round10]{10}\\
\oeis{A002620} & \hyperref[combined:round2]{2} & \oeis{A153238} & \hyperref[combined:round10]{10} & \oeis{A220898} & \hyperref[combined:round7]{7}, \hyperref[combined:round8]{8} & \oeis{A243763} & \hyperref[combined:round5]{5}\\
\oeis{A003586} & \hyperref[combined:round8]{8} & \oeis{A155051} & \hyperref[combined:round7]{7} & \oeis{A220899} & \hyperref[combined:round7]{7}, \hyperref[combined:round8]{8} & \oeis{A243764} & \hyperref[combined:round5]{5}\\
\oeis{A003961} & \hyperref[combined:round4]{4} & \oeis{A155585} & \hyperref[combined:round9]{9} & \oeis{A224511} & \hyperref[combined:round1]{1} & \oeis{A243816} & \hyperref[combined:round4]{4}\\
\oeis{A004123} & \hyperref[combined:round3]{3} & \oeis{A155828} & \hyperref[combined:round1]{1} & \oeis{A224515} & \hyperref[combined:round1]{1} & \oeis{A244177} & \hyperref[combined:round10]{10}\\
\oeis{A004395} & \hyperref[combined:round4]{4} & \oeis{A156833} & \hyperref[combined:round1]{1} & \oeis{A225692} & \hyperref[combined:round8]{8} & \oeis{A244178} & \hyperref[combined:round10]{10}\\
\oeis{A004737} & \hyperref[combined:round2]{2} & \oeis{A156834} & \hyperref[combined:round1]{1} & \oeis{A226302} & \hyperref[combined:round7]{7}, \hyperref[combined:round10]{10} & \oeis{A244179} & \hyperref[combined:round10]{10}\\
\oeis{A005046} & \hyperref[combined:round3]{3} & \oeis{A157003} & \hyperref[combined:round6]{6} & \oeis{A226602} & \hyperref[combined:round5]{5} & \oeis{A244180} & \hyperref[combined:round10]{10}\\
\oeis{A005329} & \hyperref[combined:round6]{6} & \oeis{A157021} & \hyperref[combined:round2]{2} & \oeis{A226787} & \hyperref[combined:round8]{8} & \oeis{A244181} & \hyperref[combined:round10]{10}\\
\oeis{A005351} & \hyperref[combined:round1]{1} & \oeis{A157125} & \hyperref[combined:round2]{2} & \oeis{A227012} & \hyperref[combined:round1]{1} & \oeis{A244182} & \hyperref[combined:round10]{10}\\
\oeis{A005352} & \hyperref[combined:round1]{1} & \oeis{A157308} & \hyperref[combined:round9]{9} & \oeis{A227013} & \hyperref[combined:round1]{1} & \oeis{A244183} & \hyperref[combined:round10]{10}\\
\oeis{A005558} & \hyperref[combined:round3]{3} & \oeis{A157309} & \hyperref[combined:round9]{9} & \oeis{A227014} & \hyperref[combined:round1]{1} & \oeis{A244184} & \hyperref[combined:round10]{10}\\
\oeis{A005559} & \hyperref[combined:round3]{3} & \oeis{A157310} & \hyperref[combined:round9]{9} & \oeis{A227015} & \hyperref[combined:round1]{1} & \oeis{A244185} & \hyperref[combined:round10]{10}\\
\oeis{A005560} & \hyperref[combined:round3]{3} & \oeis{A157418} & \hyperref[combined:round7]{7} & \oeis{A227016} & \hyperref[combined:round1]{1} & \oeis{A244694} & \hyperref[combined:round10]{10}\\
\oeis{A005561} & \hyperref[combined:round3]{3} & \oeis{A158119} & \hyperref[combined:round9]{9} & \oeis{A227017} & \hyperref[combined:round1]{1} & \oeis{A244695} & \hyperref[combined:round10]{10}\\
\oeis{A005562} & \hyperref[combined:round3]{3} & \oeis{A158826} & \hyperref[combined:round2]{2} & \oeis{A227018} & \hyperref[combined:round2]{2} & \oeis{A244696} & \hyperref[combined:round10]{10}\\
\oeis{A005573} & \hyperref[combined:round2]{2} & \oeis{A159772} & \hyperref[combined:round6]{6} & \oeis{A227051} & \hyperref[combined:round10]{10} & \oeis{A244697} & \hyperref[combined:round10]{10}\\
\oeis{A005773} & \hyperref[combined:round2]{2} & \oeis{A160823} & \hyperref[combined:round6]{6} & \oeis{A228178} & \hyperref[combined:round6]{6} & \oeis{A244698} & \hyperref[combined:round10]{10}\\
\oeis{A006068} & \hyperref[combined:round1]{1} & \oeis{A161789} & \hyperref[combined:round5]{5} & \oeis{A228180} & \hyperref[combined:round7]{7} & \oeis{A244699} & \hyperref[combined:round10]{10}\\
\oeis{A006154} & \hyperref[combined:round3]{3} & \oeis{A162477} & \hyperref[combined:round4]{4} & \oeis{A228884} & \hyperref[combined:round4]{4} & \oeis{A244700} & \hyperref[combined:round10]{10}\\
\oeis{A006319} & \hyperref[combined:round9]{9} & \oeis{A162481} & \hyperref[combined:round4]{4} & \oeis{A228885} & \hyperref[combined:round4]{4} & \oeis{A244701} & \hyperref[combined:round10]{10}\\
\oeis{A006343} & \hyperref[combined:round2]{2} & \oeis{A162482} & \hyperref[combined:round5]{5} & \oeis{A229041} & \hyperref[combined:round10]{10} & \oeis{A244702} & \hyperref[combined:round10]{10}\\
\oeis{A006502} & \hyperref[combined:round10]{10} & \oeis{A162533} & \hyperref[combined:round6]{6} & \oeis{A229044} & \hyperref[combined:round10]{10} & \oeis{A244884} & \hyperref[combined:round1]{1}\\
\oeis{A006519} & \hyperref[combined:round10]{10} & \oeis{A162972} & \hyperref[combined:round8]{8} & \oeis{A230089} & \hyperref[combined:round3]{3} & \oeis{A244886} & \hyperref[combined:round1]{1}\\
\oeis{A006531} & \hyperref[combined:round3]{3} & \oeis{A163774} & \hyperref[combined:round2]{2} & \oeis{A230218} & \hyperref[combined:round10]{10} & \oeis{A246606} & \hyperref[combined:round8]{8}\\
\oeis{A006789} & \hyperref[combined:round2]{2} & \oeis{A163824} & \hyperref[combined:round7]{7} & \oeis{A233637} & \hyperref[combined:round7]{7} & \oeis{A247102} & \hyperref[combined:round1]{1}\\
\oeis{A007837} & \hyperref[combined:round9]{9} & \oeis{A164586} & \hyperref[combined:round5]{5} & \oeis{A233638} & \hyperref[combined:round7]{7} & \oeis{A247173} & \hyperref[combined:round1]{1}\\
\oeis{A008292} & \hyperref[combined:round1]{1} & \oeis{A164652} & \hyperref[combined:round5]{5}, \hyperref[combined:round9]{9} & \oeis{A233639} & \hyperref[combined:round7]{7} & \oeis{A247471} & \hyperref[combined:round2]{2}\\
\oeis{A008364} & \hyperref[combined:round1]{1} & \oeis{A165201} & \hyperref[combined:round7]{7} & \oeis{A233640} & \hyperref[combined:round7]{7} & \oeis{A250917} & \hyperref[combined:round2]{2}\\
\oeis{A008365} & \hyperref[combined:round1]{1} & \oeis{A165641} & \hyperref[combined:round10]{10} & \oeis{A233646} & \hyperref[combined:round8]{8} & \oeis{A254314} & \hyperref[combined:round6]{6}\\
\oeis{A008486} & \hyperref[combined:round1]{1} & \oeis{A165886} & \hyperref[combined:round10]{10} & \oeis{A233647} & \hyperref[combined:round8]{8} & \oeis{A257104} & \hyperref[combined:round1]{1}\\
\oeis{A008578} & \hyperref[combined:round1]{1} & \oeis{A166228} & \hyperref[combined:round4]{4} & \oeis{A233653} & \hyperref[combined:round8]{8} & \oeis{A257388} & \hyperref[combined:round6]{6}\\
\oeis{A008930} & \hyperref[combined:round6]{6} & \oeis{A166306} & \hyperref[combined:round2]{2} & \oeis{A233675} & \hyperref[combined:round8]{8} & \oeis{A258409} & \hyperref[combined:round4]{4}\\
\oeis{A010694} & \hyperref[combined:round9]{9} & \oeis{A166756} & \hyperref[combined:round2]{2} & \oeis{A233676} & \hyperref[combined:round8]{8} & \oeis{A258899} & \hyperref[combined:round3]{3}\\
\oeis{A010786} & \hyperref[combined:round5]{5} & \oeis{A166757} & \hyperref[combined:round2]{2} & \oeis{A233677} & \hyperref[combined:round8]{8} & \oeis{A258903} & \hyperref[combined:round3]{3}\\
\oeis{A014533} & \hyperref[combined:round8]{8} & \oeis{A166758} & \hyperref[combined:round2]{2} & \oeis{A233682} & \hyperref[combined:round8]{8} & \oeis{A259708} & \hyperref[combined:round10]{10}\\
\oeis{A014865} & \hyperref[combined:round3]{3} & \oeis{A166759} & \hyperref[combined:round2]{2} & \oeis{A233684} & \hyperref[combined:round7]{7} & \oeis{A259709} & \hyperref[combined:round10]{10}\\
\oeis{A014866} & \hyperref[combined:round3]{3} & \oeis{A166760} & \hyperref[combined:round2]{2} & \oeis{A233685} & \hyperref[combined:round7]{7} & \oeis{A261363} & \hyperref[combined:round9]{9}\\
\oeis{A014868} & \hyperref[combined:round3]{3} & \oeis{A166761} & \hyperref[combined:round2]{2} & \oeis{A233686} & \hyperref[combined:round7]{7} & \oeis{A262735} & \hyperref[combined:round6]{6}\\
\oeis{A016789} & \hyperref[combined:round1]{1} & \oeis{A166762} & \hyperref[combined:round2]{2} & \oeis{A233687} & \hyperref[combined:round8]{8} & \oeis{A262768} & \hyperref[combined:round5]{5}\\
\oeis{A021009} & \hyperref[combined:round2]{2} & \oeis{A166776} & \hyperref[combined:round2]{2} & \oeis{A233691} & \hyperref[combined:round8]{8} & \oeis{A265649} & \hyperref[combined:round2]{2}\\
\oeis{A023887} & \hyperref[combined:round8]{8} & \oeis{A166777} & \hyperref[combined:round2]{2} & \oeis{A233710} & \hyperref[combined:round7]{7} & \oeis{A266489} & \hyperref[combined:round7]{7}\\
\oeis{A024172} & \hyperref[combined:round2]{2} & \oeis{A166778} & \hyperref[combined:round2]{2} & \oeis{A233711} & \hyperref[combined:round7]{7} & \oeis{A268149} & \hyperref[combined:round8]{8}\\
\oeis{A024180} & \hyperref[combined:round2]{2} & \oeis{A166779} & \hyperref[combined:round2]{2} & \oeis{A233718} & \hyperref[combined:round8]{8} & \oeis{A268293} & \hyperref[combined:round7]{7}\\
\oeis{A024219} & \hyperref[combined:round1]{1} & \oeis{A166780} & \hyperref[combined:round2]{2} & \oeis{A233724} & \hyperref[combined:round8]{8} & \oeis{A268402} & \hyperref[combined:round8]{8}\\
\oeis{A024220} & \hyperref[combined:round1]{1} & \oeis{A166781} & \hyperref[combined:round2]{2} & \oeis{A233726} & \hyperref[combined:round8]{8} & \oeis{A268600} & \hyperref[combined:round1]{1}\\
\oeis{A024385} & \hyperref[combined:round1]{1} & \oeis{A166782} & \hyperref[combined:round2]{2} & \oeis{A233727} & \hyperref[combined:round8]{8} & \oeis{A268601} & \hyperref[combined:round1]{1}\\
\oeis{A024388} & \hyperref[combined:round1]{1} & \oeis{A166783} & \hyperref[combined:round2]{2} & \oeis{A233728} & \hyperref[combined:round8]{8} & \oeis{A269027} & \hyperref[combined:round1]{1}\\
\oeis{A024389} & \hyperref[combined:round1]{1} & \oeis{A166786} & \hyperref[combined:round2]{2} & \oeis{A233733} & \hyperref[combined:round8]{8} & \oeis{A270490} & \hyperref[combined:round5]{5}\\
\oeis{A024390} & \hyperref[combined:round1]{1} & \oeis{A166796} & \hyperref[combined:round2]{2} & \oeis{A233749} & \hyperref[combined:round7]{7} & \oeis{A270530} & \hyperref[combined:round5]{5}\\
\oeis{A024399} & \hyperref[combined:round1]{1} & \oeis{A166797} & \hyperref[combined:round2]{2} & \oeis{A233750} & \hyperref[combined:round7]{7} & \oeis{A272261} & \hyperref[combined:round2]{2}\\
\oeis{A024401} & \hyperref[combined:round1]{1} & \oeis{A166798} & \hyperref[combined:round2]{2} & \oeis{A233787} & \hyperref[combined:round9]{9} & \oeis{A272462} & \hyperref[combined:round10]{10}\\
\oeis{A024403} & \hyperref[combined:round1]{1} & \oeis{A166799} & \hyperref[combined:round2]{2} & \oeis{A233792} & \hyperref[combined:round9]{9} & \oeis{A272468} & \hyperref[combined:round10]{10}\\
\oeis{A024719} & \hyperref[combined:round8]{8} & \oeis{A166800} & \hyperref[combined:round2]{2} & \oeis{A233877} & \hyperref[combined:round9]{9} & \oeis{A272651} & \hyperref[combined:round10]{10}\\
\oeis{A025060} & \hyperref[combined:round4]{4} & \oeis{A166801} & \hyperref[combined:round2]{2} & \oeis{A233878} & \hyperref[combined:round9]{9} & \oeis{A273228} & \hyperref[combined:round6]{6}\\
\oeis{A025256} & \hyperref[combined:round9]{9} & \oeis{A166802} & \hyperref[combined:round2]{2} & \oeis{A233879} & \hyperref[combined:round9]{9} & \oeis{A273351} & \hyperref[combined:round5]{5}\\
\oeis{A025258} & \hyperref[combined:round2]{2} & \oeis{A167137} & \hyperref[combined:round3]{3} & \oeis{A233880} & \hyperref[combined:round9]{9} & \oeis{A273905} & \hyperref[combined:round10]{10}\\
\oeis{A025259} & \hyperref[combined:round5]{5} & \oeis{A167138} & \hyperref[combined:round9]{9} & \oeis{A233883} & \hyperref[combined:round9]{9} & \oeis{A276646} & \hyperref[combined:round5]{5}\\
\oeis{A025260} & \hyperref[combined:round9]{9} & \oeis{A167139} & \hyperref[combined:round9]{9} & \oeis{A233885} & \hyperref[combined:round7]{7} & \oeis{A278023} & \hyperref[combined:round6]{6}\\
\oeis{A025268} & \hyperref[combined:round6]{6} & \oeis{A167141} & \hyperref[combined:round9]{9} & \oeis{A233886} & \hyperref[combined:round7]{7} & \oeis{A278083} & \hyperref[combined:round5]{5}\\
\oeis{A025269} & \hyperref[combined:round6]{6} & \oeis{A167532} & \hyperref[combined:round9]{9} & \oeis{A233910} & \hyperref[combined:round7]{7} & \oeis{A278085} & \hyperref[combined:round9]{9}\\
\oeis{A025270} & \hyperref[combined:round6]{6} & \oeis{A168268} & \hyperref[combined:round9]{9} & \oeis{A233911} & \hyperref[combined:round7]{7} & \oeis{A278472} & \hyperref[combined:round7]{7}\\
\oeis{A025272} & \hyperref[combined:round5]{5} & \oeis{A168452} & \hyperref[combined:round2]{2} & \oeis{A233921} & \hyperref[combined:round6]{6} & \oeis{A278618} & \hyperref[combined:round8]{8}\\
\oeis{A025273} & \hyperref[combined:round2]{2}, \hyperref[combined:round6]{6} & \oeis{A168494} & \hyperref[combined:round7]{7} & \oeis{A233942} & \hyperref[combined:round8]{8} & \oeis{A278992} & \hyperref[combined:round8]{8}\\
\oeis{A025275} & \hyperref[combined:round6]{6} & \oeis{A168503} & \hyperref[combined:round7]{7} & \oeis{A233943} & \hyperref[combined:round8]{8} & \oeis{A279013} & \hyperref[combined:round5]{5}\\
\oeis{A026023} & \hyperref[combined:round9]{9} & \oeis{A168505} & \hyperref[combined:round8]{8} & \oeis{A233944} & \hyperref[combined:round8]{8} & \oeis{A279014} & \hyperref[combined:round5]{5}\\
\oeis{A026026} & \hyperref[combined:round7]{7} & \oeis{A171791} & \hyperref[combined:round10]{10} & \oeis{A233949} & \hyperref[combined:round8]{8} & \oeis{A280098} & \hyperref[combined:round5]{5}\\
\oeis{A026027} & \hyperref[combined:round9]{9} & \oeis{A171847} & \hyperref[combined:round4]{4} & \oeis{A233951} & \hyperref[combined:round7]{7} & \oeis{A284005} & \hyperref[combined:round1]{1}\\
\oeis{A026030} & \hyperref[combined:round6]{6} & \oeis{A171849} & \hyperref[combined:round4]{4} & \oeis{A233952} & \hyperref[combined:round7]{7} & \oeis{A285195} & \hyperref[combined:round4]{4}\\
\oeis{A026126} & \hyperref[combined:round8]{8} & \oeis{A171851} & \hyperref[combined:round4]{4} & \oeis{A233960} & \hyperref[combined:round9]{9} & \oeis{A286779} & \hyperref[combined:round5]{5}\\
\oeis{A026268} & \hyperref[combined:round7]{7} & \oeis{A173693} & \hyperref[combined:round1]{1} & \oeis{A233961} & \hyperref[combined:round9]{9} & \oeis{A288417} & \hyperref[combined:round2]{2}\\
\oeis{A026269} & \hyperref[combined:round7]{7} & \oeis{A173993} & \hyperref[combined:round2]{2} & \oeis{A233962} & \hyperref[combined:round9]{9} & \oeis{A290815} & \hyperref[combined:round2]{2}\\
\oeis{A026270} & \hyperref[combined:round7]{7} & \oeis{A174013} & \hyperref[combined:round5]{5} & \oeis{A233967} & \hyperref[combined:round9]{9} & \oeis{A291937} & \hyperref[combined:round10]{10}\\
\oeis{A026568} & \hyperref[combined:round9]{9} & \oeis{A174014} & \hyperref[combined:round7]{7}, \hyperref[combined:round10]{10} & \oeis{A233982} & \hyperref[combined:round6]{6} & \oeis{A292177} & \hyperref[combined:round10]{10}\\
\oeis{A026569} & \hyperref[combined:round9]{9} & \oeis{A174015} & \hyperref[combined:round10]{10} & \oeis{A233983} & \hyperref[combined:round6]{6} & \oeis{A292394} & \hyperref[combined:round7]{7}\\
\oeis{A026570} & \hyperref[combined:round9]{9} & \oeis{A174016} & \hyperref[combined:round7]{7} & \oeis{A233984} & \hyperref[combined:round6]{6} & \oeis{A292395} & \hyperref[combined:round7]{7}\\
\oeis{A026571} & \hyperref[combined:round6]{6} & \oeis{A174107} & \hyperref[combined:round6]{6} & \oeis{A234029} & \hyperref[combined:round9]{9} & \oeis{A292461} & \hyperref[combined:round2]{2}\\
\oeis{A026572} & \hyperref[combined:round9]{9} & \oeis{A174169} & \hyperref[combined:round8]{8} & \oeis{A234030} & \hyperref[combined:round9]{9} & \oeis{A292877} & \hyperref[combined:round10]{10}\\
\oeis{A026582} & \hyperref[combined:round9]{9} & \oeis{A174193} & \hyperref[combined:round1]{1} & \oeis{A234031} & \hyperref[combined:round9]{9} & \oeis{A295012} & \hyperref[combined:round5]{5}\\
\oeis{A026641} & \hyperref[combined:round2]{2} & \oeis{A174490} & \hyperref[combined:round8]{8} & \oeis{A234032} & \hyperref[combined:round9]{9} & \oeis{A299406} & \hyperref[combined:round6]{6}\\
\oeis{A026725} & \hyperref[combined:round9]{9} & \oeis{A174491} & \hyperref[combined:round8]{8} & \oeis{A234036} & \hyperref[combined:round9]{9} & \oeis{A300616} & \hyperref[combined:round8]{8}\\
\oeis{A026732} & \hyperref[combined:round9]{9} & \oeis{A174808} & \hyperref[combined:round6]{6} & \oeis{A234107} & \hyperref[combined:round9]{9} & \oeis{A300617} & \hyperref[combined:round8]{8}\\
\oeis{A026741} & \hyperref[combined:round9]{9} & \oeis{A175505} & \hyperref[combined:round7]{7} & \oeis{A234108} & \hyperref[combined:round9]{9} & \oeis{A300732} & \hyperref[combined:round7]{7}\\
\oeis{A026743} & \hyperref[combined:round6]{6} & \oeis{A176604} & \hyperref[combined:round5]{5}, \hyperref[combined:round6]{6}, \hyperref[combined:round7]{7}, \hyperref[combined:round8]{8} & \oeis{A234109} & \hyperref[combined:round9]{9} & \oeis{A300734} & \hyperref[combined:round7]{7}\\
\oeis{A026758} & \hyperref[combined:round9]{9} & \oeis{A176605} & \hyperref[combined:round4]{4}, \hyperref[combined:round5]{5}, \hyperref[combined:round7]{7}, \hyperref[combined:round8]{8}, \hyperref[combined:round9]{9} & \oeis{A234114} & \hyperref[combined:round9]{9} & \oeis{A300984} & \hyperref[combined:round5]{5}\\
\oeis{A026765} & \hyperref[combined:round9]{9} & \oeis{A176607} & \hyperref[combined:round4]{4}, \hyperref[combined:round5]{5}, \hyperref[combined:round7]{7}, \hyperref[combined:round8]{8}, \hyperref[combined:round9]{9} & \oeis{A234117} & \hyperref[combined:round9]{9} & \oeis{A301921} & \hyperref[combined:round3]{3}\\
\oeis{A027023} & \hyperref[combined:round1]{1}, \hyperref[combined:round2]{2} & \oeis{A176610} & \hyperref[combined:round4]{4}, \hyperref[combined:round5]{5}, \hyperref[combined:round7]{7}, \hyperref[combined:round8]{8}, \hyperref[combined:round9]{9} & \oeis{A234118} & \hyperref[combined:round9]{9} & \oeis{A302180} & \hyperref[combined:round8]{8}\\
\oeis{A027027} & \hyperref[combined:round2]{2} & \oeis{A176611} & \hyperref[combined:round4]{4}, \hyperref[combined:round5]{5}, \hyperref[combined:round7]{7}, \hyperref[combined:round8]{8}, \hyperref[combined:round9]{9} & \oeis{A234119} & \hyperref[combined:round9]{9} & \oeis{A302181} & \hyperref[combined:round8]{8}\\
\oeis{A027028} & \hyperref[combined:round2]{2} & \oeis{A176612} & \hyperref[combined:round4]{4}, \hyperref[combined:round5]{5}, \hyperref[combined:round7]{7}, \hyperref[combined:round8]{8}, \hyperref[combined:round9]{9} & \oeis{A234122} & \hyperref[combined:round9]{9} & \oeis{A302184} & \hyperref[combined:round8]{8}\\
\oeis{A027050} & \hyperref[combined:round1]{1} & \oeis{A176615} & \hyperref[combined:round4]{4} & \oeis{A234124} & \hyperref[combined:round9]{9} & \oeis{A302186} & \hyperref[combined:round8]{8}\\
\oeis{A027051} & \hyperref[combined:round2]{2} & \oeis{A177258} & \hyperref[combined:round1]{1} & \oeis{A234125} & \hyperref[combined:round9]{9} & \oeis{A303290} & \hyperref[combined:round10]{10}\\
\oeis{A027052} & \hyperref[combined:round1]{1}, \hyperref[combined:round2]{2} & \oeis{A177259} & \hyperref[combined:round1]{1} & \oeis{A234126} & \hyperref[combined:round9]{9} & \oeis{A305550} & \hyperref[combined:round3]{3}\\
\oeis{A027056} & \hyperref[combined:round2]{2} & \oeis{A177260} & \hyperref[combined:round1]{1} & \oeis{A234131} & \hyperref[combined:round9]{9} & \oeis{A306082} & \hyperref[combined:round3]{3}\\
\oeis{A027058} & \hyperref[combined:round1]{1} & \oeis{A177430} & \hyperref[combined:round8]{8} & \oeis{A234133} & \hyperref[combined:round9]{9} & \oeis{A308620} & \hyperref[combined:round10]{10}\\
\oeis{A027059} & \hyperref[combined:round1]{1} & \oeis{A177431} & \hyperref[combined:round8]{8} & \oeis{A234135} & \hyperref[combined:round9]{9} & \oeis{A308726} & \hyperref[combined:round4]{4}\\
\oeis{A027082} & \hyperref[combined:round1]{1}, \hyperref[combined:round2]{2} & \oeis{A177510} & \hyperref[combined:round6]{6} & \oeis{A234140} & \hyperref[combined:round9]{9} & \oeis{A308815} & \hyperref[combined:round8]{8}\\
\oeis{A027087} & \hyperref[combined:round2]{2} & \oeis{A178113} & \hyperref[combined:round2]{2} & \oeis{A234146} & \hyperref[combined:round7]{7} & \oeis{A308833} & \hyperref[combined:round10]{10}\\
\oeis{A027088} & \hyperref[combined:round1]{1} & \oeis{A179504} & \hyperref[combined:round8]{8} & \oeis{A234147} & \hyperref[combined:round7]{7} & \oeis{A309263} & \hyperref[combined:round7]{7}\\
\oeis{A027089} & \hyperref[combined:round1]{1} & \oeis{A179648} & \hyperref[combined:round4]{4} & \oeis{A234154} & \hyperref[combined:round9]{9} & \oeis{A320192} & \hyperref[combined:round4]{4}\\
\oeis{A027113} & \hyperref[combined:round2]{2} & \oeis{A179873} & \hyperref[combined:round7]{7} & \oeis{A234155} & \hyperref[combined:round9]{9} & \oeis{A320352} & \hyperref[combined:round3]{3}\\
\oeis{A027116} & \hyperref[combined:round2]{2} & \oeis{A179929} & \hyperref[combined:round5]{5} & \oeis{A234156} & \hyperref[combined:round9]{9} & \oeis{A320949} & \hyperref[combined:round10]{10}\\
\oeis{A027117} & \hyperref[combined:round2]{2} & \oeis{A180195} & \hyperref[combined:round1]{1} & \oeis{A234157} & \hyperref[combined:round9]{9} & \oeis{A322327} & \hyperref[combined:round2]{2}\\
\oeis{A027170} & \hyperref[combined:round1]{1} & \oeis{A180654} & \hyperref[combined:round9]{9} & \oeis{A234158} & \hyperref[combined:round9]{9} & \oeis{A323599} & \hyperref[combined:round9]{9}\\
\oeis{A027175} & \hyperref[combined:round1]{1} & \oeis{A182879} & \hyperref[combined:round5]{5} & \oeis{A234161} & \hyperref[combined:round9]{9} & \oeis{A323689} & \hyperref[combined:round1]{1}\\
\oeis{A027182} & \hyperref[combined:round1]{1} & \oeis{A182892} & \hyperref[combined:round8]{8} & \oeis{A234162} & \hyperref[combined:round6]{6} & \oeis{A324300} & \hyperref[combined:round1]{1}\\
\oeis{A027907} & \hyperref[combined:round6]{6} & \oeis{A182894} & \hyperref[combined:round5]{5} & \oeis{A234163} & \hyperref[combined:round6]{6} & \oeis{A325286} & \hyperref[combined:round10]{10}\\
\oeis{A030204} & \hyperref[combined:round5]{5} & \oeis{A182899} & \hyperref[combined:round4]{4} & \oeis{A234170} & \hyperref[combined:round9]{9} & \oeis{A326478} & \hyperref[combined:round10]{10}\\
\oeis{A030212} & \hyperref[combined:round7]{7} & \oeis{A182904} & \hyperref[combined:round5]{5} & \oeis{A234171} & \hyperref[combined:round9]{9} & \oeis{A326691} & \hyperref[combined:round8]{8}\\
\oeis{A032115} & \hyperref[combined:round2]{2} & \oeis{A183586} & \hyperref[combined:round6]{6} & \oeis{A234175} & \hyperref[combined:round9]{9} & \oeis{A327096} & \hyperref[combined:round2]{2}\\
\oeis{A032123} & \hyperref[combined:round2]{2}, \hyperref[combined:round4]{4} & \oeis{A183587} & \hyperref[combined:round6]{6} & \oeis{A234177} & \hyperref[combined:round9]{9} & \oeis{A328260} & \hyperref[combined:round9]{9}\\
\oeis{A032312} & \hyperref[combined:round9]{9} & \oeis{A183589} & \hyperref[combined:round6]{6} & \oeis{A234178} & \hyperref[combined:round9]{9} & \oeis{A329354} & \hyperref[combined:round9]{9}\\
\oeis{A033955} & \hyperref[combined:round2]{2} & \oeis{A183596} & \hyperref[combined:round6]{6} & \oeis{A234179} & \hyperref[combined:round9]{9} & \oeis{A329369} & \hyperref[combined:round1]{1}\\
\oeis{A034261} & \hyperref[combined:round5]{5} & \oeis{A184018} & \hyperref[combined:round7]{7} & \oeis{A234184} & \hyperref[combined:round9]{9} & \oeis{A329375} & \hyperref[combined:round9]{9}\\
\oeis{A034267} & \hyperref[combined:round5]{5} & \oeis{A184120} & \hyperref[combined:round4]{4} & \oeis{A234201} & \hyperref[combined:round9]{9} & \oeis{A329649} & \hyperref[combined:round1]{1}\\
\oeis{A034444} & \hyperref[combined:round5]{5} & \oeis{A184534} & \hyperref[combined:round2]{2} & \oeis{A234202} & \hyperref[combined:round9]{9} & \oeis{A329697} & \hyperref[combined:round1]{1}\\
\oeis{A035263} & \hyperref[combined:round1]{1} & \oeis{A185010} & \hyperref[combined:round4]{4} & \oeis{A234208} & \hyperref[combined:round9]{9} & \oeis{A330358} & \hyperref[combined:round1]{1}\\
\oeis{A037223} & \hyperref[combined:round2]{2} & \oeis{A185020} & \hyperref[combined:round7]{7} & \oeis{A234210} & \hyperref[combined:round9]{9} & \oeis{A334109} & \hyperref[combined:round1]{1}\\
\oeis{A040001} & \hyperref[combined:round9]{9} & \oeis{A185071} & \hyperref[combined:round10]{10} & \oeis{A234211} & \hyperref[combined:round9]{9} & \oeis{A334110} & \hyperref[combined:round1]{1}\\
\oeis{A040003} & \hyperref[combined:round9]{9} & \oeis{A185072} & \hyperref[combined:round10]{10} & \oeis{A234212} & \hyperref[combined:round9]{9} & \oeis{A334592} & \hyperref[combined:round1]{1}\\
\oeis{A045406} & \hyperref[combined:round4]{4} & \oeis{A186185} & \hyperref[combined:round1]{1} & \oeis{A234217} & \hyperref[combined:round9]{9} & \oeis{A334595} & \hyperref[combined:round2]{2}\\
\oeis{A045742} & \hyperref[combined:round1]{1} & \oeis{A186335} & \hyperref[combined:round2]{2} & \oeis{A234220} & \hyperref[combined:round6]{6} & \oeis{A336838} & \hyperref[combined:round4]{4}\\
\oeis{A045743} & \hyperref[combined:round4]{4} & \oeis{A186338} & \hyperref[combined:round4]{4} & \oeis{A234221} & \hyperref[combined:round6]{6} & \oeis{A337500} & \hyperref[combined:round1]{1}\\
\oeis{A048651} & \hyperref[combined:round6]{6} & \oeis{A186776} & \hyperref[combined:round5]{5} & \oeis{A234222} & \hyperref[combined:round6]{6} & \oeis{A337758} & \hyperref[combined:round10]{10}\\
\oeis{A048702} & \hyperref[combined:round3]{3} & \oeis{A186940} & \hyperref[combined:round4]{4} & \oeis{A234241} & \hyperref[combined:round8]{8} & \oeis{A338328} & \hyperref[combined:round10]{10}\\
\oeis{A048868} & \hyperref[combined:round3]{3} & \oeis{A187011} & \hyperref[combined:round5]{5} & \oeis{A234245} & \hyperref[combined:round8]{8} & \oeis{A338632} & \hyperref[combined:round5]{5}\\
\oeis{A048896} & \hyperref[combined:round9]{9} & \oeis{A187250} & \hyperref[combined:round1]{1} & \oeis{A234259} & \hyperref[combined:round8]{8} & \oeis{A338636} & \hyperref[combined:round5]{5}\\
\oeis{A048993} & \hyperref[combined:round2]{2} & \oeis{A187252} & \hyperref[combined:round1]{1} & \oeis{A234260} & \hyperref[combined:round8]{8} & \oeis{A339100} & \hyperref[combined:round10]{10}\\
\oeis{A048994} & \hyperref[combined:round2]{2} & \oeis{A187256} & \hyperref[combined:round4]{4} & \oeis{A234261} & \hyperref[combined:round8]{8} & \oeis{A339384} & \hyperref[combined:round5]{5}\\
\oeis{A049803} & \hyperref[combined:round1]{1} & \oeis{A188570} & \hyperref[combined:round3]{3} & \oeis{A234266} & \hyperref[combined:round8]{8} & \oeis{A339387} & \hyperref[combined:round5]{5}\\
\oeis{A049804} & \hyperref[combined:round1]{1} & \oeis{A188571} & \hyperref[combined:round3]{3} & \oeis{A234269} & \hyperref[combined:round1]{1} & \oeis{A342912} & \hyperref[combined:round2]{2}\\
\oeis{A050168} & \hyperref[combined:round4]{4} & \oeis{A188572} & \hyperref[combined:round3]{3} & \oeis{A234307} & \hyperref[combined:round2]{2} & \oeis{A344502} & \hyperref[combined:round8]{8}\\
\oeis{A051176} & \hyperref[combined:round9]{9} & \oeis{A188573} & \hyperref[combined:round3]{3} & \oeis{A234326} & \hyperref[combined:round8]{8} & \oeis{A345651} & \hyperref[combined:round1]{1}\\
\oeis{A051283} & \hyperref[combined:round2]{2} & \oeis{A189176} & \hyperref[combined:round5]{5} & \oeis{A234327} & \hyperref[combined:round8]{8} & \oeis{A345993} & \hyperref[combined:round8]{8}\\
\oeis{A051699} & \hyperref[combined:round5]{5} & \oeis{A190160} & \hyperref[combined:round2]{2} & \oeis{A234328} & \hyperref[combined:round8]{8} & \oeis{A346074} & \hyperref[combined:round1]{1}\\
\oeis{A052174} & \hyperref[combined:round3]{3} & \oeis{A190166} & \hyperref[combined:round5]{5}, \hyperref[combined:round6]{6} & \oeis{A234333} & \hyperref[combined:round8]{8} & \oeis{A346418} & \hyperref[combined:round2]{2}\\
\oeis{A052183} & \hyperref[combined:round6]{6} & \oeis{A190169} & \hyperref[combined:round6]{6} & \oeis{A234338} & \hyperref[combined:round8]{8} & \oeis{A347293} & \hyperref[combined:round2]{2}\\
\oeis{A052731} & \hyperref[combined:round4]{4} & \oeis{A190363} & \hyperref[combined:round4]{4} & \oeis{A234343} & \hyperref[combined:round8]{8} & \oeis{A348295} & \hyperref[combined:round3]{3}, \hyperref[combined:round9]{9}\\
\oeis{A052736} & \hyperref[combined:round9]{9} & \oeis{A190725} & \hyperref[combined:round7]{7} & \oeis{A234444} & \hyperref[combined:round8]{8} & \oeis{A348299} & \hyperref[combined:round3]{3}\\
\oeis{A053532} & \hyperref[combined:round2]{2} & \oeis{A190736} & \hyperref[combined:round7]{7} & \oeis{A234445} & \hyperref[combined:round8]{8} & \oeis{A348410} & \hyperref[combined:round2]{2}\\
\oeis{A053644} & \hyperref[combined:round1]{1} & \oeis{A190788} & \hyperref[combined:round7]{7} & \oeis{A234450} & \hyperref[combined:round8]{8} & \oeis{A348713} & \hyperref[combined:round6]{6}\\
\oeis{A053645} & \hyperref[combined:round1]{1} & \oeis{A191107} & \hyperref[combined:round5]{5} & \oeis{A234452} & \hyperref[combined:round8]{8} & \oeis{A349802} & \hyperref[combined:round2]{2}\\
\oeis{A054120} & \hyperref[combined:round1]{1} & \oeis{A191319} & \hyperref[combined:round7]{7} & \oeis{A234453} & \hyperref[combined:round8]{8} & \oeis{A349890} & \hyperref[combined:round6]{6}\\
\oeis{A054523} & \hyperref[combined:round7]{7} & \oeis{A191398} & \hyperref[combined:round7]{7} & \oeis{A234458} & \hyperref[combined:round8]{8} & \oeis{A350512} & \hyperref[combined:round1]{1}\\
\oeis{A057660} & \hyperref[combined:round2]{2} & \oeis{A191579} & \hyperref[combined:round2]{2} & \oeis{A234483} & \hyperref[combined:round8]{8} & \oeis{A350557} & \hyperref[combined:round1]{1}\\
\oeis{A058087} & \hyperref[combined:round2]{2} & \oeis{A191792} & \hyperref[combined:round4]{4} & \oeis{A234484} & \hyperref[combined:round8]{8} & \oeis{A350900} & \hyperref[combined:round2]{2}\\
\oeis{A059259} & \hyperref[combined:round2]{2} & \oeis{A192069} & \hyperref[combined:round9]{9} & \oeis{A234490} & \hyperref[combined:round8]{8} & \oeis{A351933} & \hyperref[combined:round1]{1}\\
\oeis{A059279} & \hyperref[combined:round7]{7} & \oeis{A192248} & \hyperref[combined:round9]{9} & \oeis{A234491} & \hyperref[combined:round8]{8} & \oeis{A352117} & \hyperref[combined:round8]{8}\\
\oeis{A059297} & \hyperref[combined:round2]{2} & \oeis{A192249} & \hyperref[combined:round9]{9} & \oeis{A234492} & \hyperref[combined:round8]{8} & \oeis{A354242} & \hyperref[combined:round3]{3}\\
\oeis{A059837} & \hyperref[combined:round1]{1} & \oeis{A193679} & \hyperref[combined:round8]{8} & \oeis{A234497} & \hyperref[combined:round8]{8} & \oeis{A354253} & \hyperref[combined:round3]{3}\\
\oeis{A060043} & \hyperref[combined:round7]{7} & \oeis{A193758} & \hyperref[combined:round9]{9} & \oeis{A234557} & \hyperref[combined:round6]{6} & \oeis{A354766} & \hyperref[combined:round9]{9}\\
\oeis{A060791} & \hyperref[combined:round9]{9} & \oeis{A193844} & \hyperref[combined:round2]{2} & \oeis{A234558} & \hyperref[combined:round6]{6} & \oeis{A354933} & \hyperref[combined:round2]{2}\\
\oeis{A061002} & \hyperref[combined:round6]{6} & \oeis{A195194} & \hyperref[combined:round3]{3} & \oeis{A234559} & \hyperref[combined:round6]{6} & \oeis{A354995} & \hyperref[combined:round2]{2}\\
\oeis{A062803} & \hyperref[combined:round7]{7} & \oeis{A195986} & \hyperref[combined:round10]{10} & \oeis{A234651} & \hyperref[combined:round7]{7} & \oeis{A355282} & \hyperref[combined:round1]{1}\\
\oeis{A063081} & \hyperref[combined:round10]{10} & \oeis{A196523} & \hyperref[combined:round8]{8} & \oeis{A234652} & \hyperref[combined:round7]{7} & \oeis{A355372} & \hyperref[combined:round9]{9}\\
\oeis{A063087} & \hyperref[combined:round10]{10} & \oeis{A199716} & \hyperref[combined:round3]{3} & \oeis{A234653} & \hyperref[combined:round7]{7} & \oeis{A355401} & \hyperref[combined:round1]{1}\\
\oeis{A063098} & \hyperref[combined:round10]{10} & \oeis{A202943} & \hyperref[combined:round9]{9} & \oeis{A234659} & \hyperref[combined:round9]{9} & \oeis{A355407} & \hyperref[combined:round9]{9}\\
\oeis{A063105} & \hyperref[combined:round10]{10} & \oeis{A203150} & \hyperref[combined:round1]{1} & \oeis{A234660} & \hyperref[combined:round9]{9} & \oeis{A355414} & \hyperref[combined:round9]{9}\\
\oeis{A063110} & \hyperref[combined:round10]{10} & \oeis{A203157} & \hyperref[combined:round1]{1} & \oeis{A234661} & \hyperref[combined:round9]{9} & \oeis{A356459} & \hyperref[combined:round7]{7}\\
\oeis{A063116} & \hyperref[combined:round10]{10} & \oeis{A203159} & \hyperref[combined:round1]{1} & \oeis{A234665} & \hyperref[combined:round9]{9} & \oeis{A359643} & \hyperref[combined:round7]{7}\\
\oeis{A063118} & \hyperref[combined:round10]{10} & \oeis{A203160} & \hyperref[combined:round1]{1} & \oeis{A234667} & \hyperref[combined:round9]{9} & \oeis{A359914} & \hyperref[combined:round4]{4}\\
\oeis{A063126} & \hyperref[combined:round10]{10} & \oeis{A203163} & \hyperref[combined:round1]{1} & \oeis{A234668} & \hyperref[combined:round9]{9} & \oeis{A361552} & \hyperref[combined:round4]{4}\\
\oeis{A063128} & \hyperref[combined:round10]{10} & \oeis{A203164} & \hyperref[combined:round1]{1} & \oeis{A234672} & \hyperref[combined:round9]{9} & \oeis{A362560} & \hyperref[combined:round1]{1}\\
\oeis{A063140} & \hyperref[combined:round10]{10} & \oeis{A203165} & \hyperref[combined:round1]{1} & \oeis{A234674} & \hyperref[combined:round8]{8} & \oeis{A363409} & \hyperref[combined:round9]{9}\\
\oeis{A063147} & \hyperref[combined:round10]{10} & \oeis{A203166} & \hyperref[combined:round1]{1} & \oeis{A234675} & \hyperref[combined:round8]{8} & \oeis{A363411} & \hyperref[combined:round9]{9}\\
\oeis{A063148} & \hyperref[combined:round10]{10} & \oeis{A203167} & \hyperref[combined:round1]{1} & \oeis{A234681} & \hyperref[combined:round8]{8} & \oeis{A363412} & \hyperref[combined:round9]{9}\\
\oeis{A063168} & \hyperref[combined:round10]{10} & \oeis{A203231} & \hyperref[combined:round1]{1} & \oeis{A234683} & \hyperref[combined:round9]{9} & \oeis{A365095} & \hyperref[combined:round8]{8}\\
\oeis{A063233} & \hyperref[combined:round10]{10} & \oeis{A203233} & \hyperref[combined:round1]{1} & \oeis{A234684} & \hyperref[combined:round9]{9} & \oeis{A366561} & \hyperref[combined:round7]{7}\\
\oeis{A063247} & \hyperref[combined:round10]{10} & \oeis{A203234} & \hyperref[combined:round1]{1} & \oeis{A234690} & \hyperref[combined:round9]{9} & \oeis{A367865} & \hyperref[combined:round8]{8}\\
\oeis{A063647} & \hyperref[combined:round10]{10} & \oeis{A203235} & \hyperref[combined:round1]{1} & \oeis{A234697} & \hyperref[combined:round9]{9} & \oeis{A368628} & \hyperref[combined:round3]{3}\\
\oeis{A063946} & \hyperref[combined:round1]{1} & \oeis{A204243} & \hyperref[combined:round6]{6} & \oeis{A234698} & \hyperref[combined:round9]{9} & \oeis{A368630} & \hyperref[combined:round4]{4}\\
\oeis{A063956} & \hyperref[combined:round6]{6} & \oeis{A205355} & \hyperref[combined:round3]{3} & \oeis{A234699} & \hyperref[combined:round9]{9} & \oeis{A368631} & \hyperref[combined:round4]{4}\\
\oeis{A064618} & \hyperref[combined:round4]{4} & \oeis{A205356} & \hyperref[combined:round3]{3} & \oeis{A234702} & \hyperref[combined:round9]{9} & \oeis{A368633} & \hyperref[combined:round2]{2}\\
\oeis{A065919} & \hyperref[combined:round2]{2} & \oeis{A205357} & \hyperref[combined:round3]{3} & \oeis{A234706} & \hyperref[combined:round9]{9} & \oeis{A368634} & \hyperref[combined:round3]{3}\\
\oeis{A065942} & \hyperref[combined:round7]{7} & \oeis{A205358} & \hyperref[combined:round3]{3} & \oeis{A234707} & \hyperref[combined:round9]{9} & \oeis{A369083} & \hyperref[combined:round4]{4}\\
\oeis{A068991} & \hyperref[combined:round6]{6} & \oeis{A205572} & \hyperref[combined:round10]{10} & \oeis{A234712} & \hyperref[combined:round9]{9} & \oeis{A371364} & \hyperref[combined:round9]{9}\\
\oeis{A069208} & \hyperref[combined:round6]{6} & \oeis{A206377} & \hyperref[combined:round3]{3} & \oeis{A234721} & \hyperref[combined:round9]{9} & \oeis{A372018} & \hyperref[combined:round9]{9}\\
\oeis{A071264} & \hyperref[combined:round1]{1} & \oeis{A206378} & \hyperref[combined:round3]{3} & \oeis{A234722} & \hyperref[combined:round9]{9} & \oeis{A372530} & \hyperref[combined:round10]{10}\\
\oeis{A071715} & \hyperref[combined:round9]{9} & \oeis{A206379} & \hyperref[combined:round3]{3} & \oeis{A234728} & \hyperref[combined:round9]{9} & \oeis{A372991} & \hyperref[combined:round7]{7}\\
\oeis{A072100} & \hyperref[combined:round4]{4} & \oeis{A206380} & \hyperref[combined:round3]{3} & \oeis{A234731} & \hyperref[combined:round7]{7} & \oeis{A373312} & \hyperref[combined:round10]{10}\\
\oeis{A075002} & \hyperref[combined:round1]{1} & \oeis{A208355} & \hyperref[combined:round6]{6} & \oeis{A234732} & \hyperref[combined:round7]{7} & \oeis{A373313} & \hyperref[combined:round10]{10}\\
\oeis{A075045} & \hyperref[combined:round2]{2} & \oeis{A209505} & \hyperref[combined:round5]{5} & \oeis{A234733} & \hyperref[combined:round7]{7} & \oeis{A373561} & \hyperref[combined:round8]{8}\\
\oeis{A077761} & \hyperref[combined:round9]{9} & \oeis{A209506} & \hyperref[combined:round5]{5} & \oeis{A234753} & \hyperref[combined:round7]{7} & \oeis{A375838} & \hyperref[combined:round7]{7}\\
\oeis{A078121} & \hyperref[combined:round1]{1} & \oeis{A209507} & \hyperref[combined:round5]{5} & \oeis{A234754} & \hyperref[combined:round7]{7} & \oeis{A376182} & \hyperref[combined:round2]{2}\\
\oeis{A078122} & \hyperref[combined:round1]{1} & \oeis{A209508} & \hyperref[combined:round5]{5} & \oeis{A234762} & \hyperref[combined:round6]{6} & \oeis{A376226} & \hyperref[combined:round9]{9}\\
\oeis{A078995} & \hyperref[combined:round2]{2} & \oeis{A209509} & \hyperref[combined:round5]{5} & \oeis{A234763} & \hyperref[combined:round6]{6} & \oeis{A376336} & \hyperref[combined:round2]{2}\\
\oeis{A078999} & \hyperref[combined:round2]{2} & \oeis{A209530} & \hyperref[combined:round5]{5} & \oeis{A234764} & \hyperref[combined:round6]{6} & \oeis{A376484} & \hyperref[combined:round6]{6}\\
\oeis{A079269} & \hyperref[combined:round8]{8} & \oeis{A209531} & \hyperref[combined:round5]{5} & \oeis{A234770} & \hyperref[combined:round7]{7} & \oeis{A376499} & \hyperref[combined:round6]{6}\\
\oeis{A079278} & \hyperref[combined:round8]{8} & \oeis{A209546} & \hyperref[combined:round5]{5} & \oeis{A234771} & \hyperref[combined:round7]{7} & \oeis{A377109} & \hyperref[combined:round3]{3}\\
\oeis{A079998} & \hyperref[combined:round9]{9} & \oeis{A209547} & \hyperref[combined:round5]{5} & \oeis{A234816} & \hyperref[combined:round9]{9} & \oeis{A377110} & \hyperref[combined:round3]{3}\\
\oeis{A080243} & \hyperref[combined:round6]{6} & \oeis{A209548} & \hyperref[combined:round5]{5} & \oeis{A234817} & \hyperref[combined:round9]{9} & \oeis{A377111} & \hyperref[combined:round3]{3}\\
\oeis{A080253} & \hyperref[combined:round3]{3} & \oeis{A209549} & \hyperref[combined:round5]{5} & \oeis{A234818} & \hyperref[combined:round9]{9} & \oeis{A377112} & \hyperref[combined:round3]{3}\\
\oeis{A080894} & \hyperref[combined:round2]{2} & \oeis{A209550} & \hyperref[combined:round7]{7} & \oeis{A234823} & \hyperref[combined:round9]{9} & \oeis{A377113} & \hyperref[combined:round3]{3}\\
\oeis{A081199} & \hyperref[combined:round1]{1} & \oeis{A209709} & \hyperref[combined:round3]{3} & \oeis{A234875} & \hyperref[combined:round6]{6} & \oeis{A377114} & \hyperref[combined:round3]{3}\\
\oeis{A081200} & \hyperref[combined:round1]{1} & \oeis{A209710} & \hyperref[combined:round5]{5} & \oeis{A234876} & \hyperref[combined:round6]{6} & \oeis{A377115} & \hyperref[combined:round3]{3}\\
\oeis{A081201} & \hyperref[combined:round1]{1} & \oeis{A209711} & \hyperref[combined:round5]{5} & \oeis{A234877} & \hyperref[combined:round6]{6} & \oeis{A377116} & \hyperref[combined:round3]{3}\\
\oeis{A081577} & \hyperref[combined:round1]{1} & \oeis{A209722} & \hyperref[combined:round5]{5} & \oeis{A234884} & \hyperref[combined:round8]{8} & \oeis{A377117} & \hyperref[combined:round3]{3}\\
\oeis{A082397} & \hyperref[combined:round2]{2} & \oeis{A209729} & \hyperref[combined:round3]{3} & \oeis{A234885} & \hyperref[combined:round8]{8} & \oeis{A377118} & \hyperref[combined:round3]{3}\\
\oeis{A082953} & \hyperref[combined:round7]{7} & \oeis{A209807} & \hyperref[combined:round3]{3} & \oeis{A234890} & \hyperref[combined:round8]{8} & \oeis{A377119} & \hyperref[combined:round3]{3}\\
\oeis{A083365} & \hyperref[combined:round6]{6} & \oeis{A209808} & \hyperref[combined:round3]{3} & \oeis{A234891} & \hyperref[combined:round8]{8} & \oeis{A377278} & \hyperref[combined:round6]{6}\\
\oeis{A084068} & \hyperref[combined:round9]{9} & \oeis{A209809} & \hyperref[combined:round3]{3} & \oeis{A234898} & \hyperref[combined:round8]{8} & \oeis{A378277} & \hyperref[combined:round2]{2}\\
\oeis{A084938} & \hyperref[combined:round10]{10} & \oeis{A209810} & \hyperref[combined:round3]{3} & \oeis{A234905} & \hyperref[combined:round8]{8} & \oeis{A378575} & \hyperref[combined:round4]{4}\\
\oeis{A085361} & \hyperref[combined:round5]{5} & \oeis{A209811} & \hyperref[combined:round3]{3} & \oeis{A234911} & \hyperref[combined:round8]{8} & \oeis{A378578} & \hyperref[combined:round4]{4}\\
\oeis{A085765} & \hyperref[combined:round1]{1} & \oeis{A209822} & \hyperref[combined:round3]{3} & \oeis{A234984} & \hyperref[combined:round8]{8} & \oeis{A378580} & \hyperref[combined:round3]{3}, \hyperref[combined:round8]{8}\\
\oeis{A086450} & \hyperref[combined:round1]{1} & \oeis{A209823} & \hyperref[combined:round3]{3} & \oeis{A234985} & \hyperref[combined:round8]{8} & \oeis{A378581} & \hyperref[combined:round8]{8}\\
\oeis{A087726} & \hyperref[combined:round2]{2} & \oeis{A209824} & \hyperref[combined:round3]{3} & \oeis{A234991} & \hyperref[combined:round8]{8} & \oeis{A379463} & \hyperref[combined:round8]{8}\\
\oeis{A089324} & \hyperref[combined:round5]{5} & \oeis{A209825} & \hyperref[combined:round3]{3} & \oeis{A234992} & \hyperref[combined:round9]{9} & \oeis{A379464} & \hyperref[combined:round8]{8}\\
\oeis{A089354} & \hyperref[combined:round2]{2} & \oeis{A209826} & \hyperref[combined:round3]{3} & \oeis{A234993} & \hyperref[combined:round9]{9} & \oeis{A379764} & \hyperref[combined:round8]{8}\\
\oeis{A091137} & \hyperref[combined:round10]{10} & \oeis{A209841} & \hyperref[combined:round3]{3} & \oeis{A234998} & \hyperref[combined:round9]{9} & \oeis{A379907} & \hyperref[combined:round2]{2}\\
\oeis{A091351} & \hyperref[combined:round9]{9} & \oeis{A209842} & \hyperref[combined:round3]{3} & \oeis{A235011} & \hyperref[combined:round6]{6} & \oeis{A380056} & \hyperref[combined:round4]{4}\\
\oeis{A091539} & \hyperref[combined:round1]{1} & \oeis{A209843} & \hyperref[combined:round3]{3} & \oeis{A235012} & \hyperref[combined:round6]{6} & \oeis{A380060} & \hyperref[combined:round8]{8}\\
\oeis{A092278} & \hyperref[combined:round3]{3} & \oeis{A209844} & \hyperref[combined:round3]{3} & \oeis{A235019} & \hyperref[combined:round8]{8} & \oeis{A380191} & \hyperref[combined:round2]{2}\\
\oeis{A092287} & \hyperref[combined:round8]{8} & \oeis{A209845} & \hyperref[combined:round3]{3} & \oeis{A235020} & \hyperref[combined:round8]{8} & \oeis{A380680} & \hyperref[combined:round8]{8}\\
\oeis{A092843} & \hyperref[combined:round7]{7} & \oeis{A209851} & \hyperref[combined:round5]{5} & \oeis{A235026} & \hyperref[combined:round8]{8} & \oeis{A380709} & \hyperref[combined:round4]{4}\\
\oeis{A094287} & \hyperref[combined:round10]{10} & \oeis{A209852} & \hyperref[combined:round5]{5} & \oeis{A235064} & \hyperref[combined:round8]{8} & \oeis{A380711} & \hyperref[combined:round4]{4}\\
\oeis{A094288} & \hyperref[combined:round10]{10} & \oeis{A209853} & \hyperref[combined:round5]{5} & \oeis{A235065} & \hyperref[combined:round8]{8} & \oeis{A381353} & \hyperref[combined:round6]{6}\\
\oeis{A095981} & \hyperref[combined:round6]{6} & \oeis{A209890} & \hyperref[combined:round3]{3} & \oeis{A235071} & \hyperref[combined:round8]{8} & \oeis{A381355} & \hyperref[combined:round6]{6}\\
\oeis{A096471} & \hyperref[combined:round7]{7} & \oeis{A209891} & \hyperref[combined:round3]{3} & \oeis{A235080} & \hyperref[combined:round6]{6} & \oeis{A381670} & \hyperref[combined:round8]{8}\\
\oeis{A096617} & \hyperref[combined:round9]{9} & \oeis{A209892} & \hyperref[combined:round3]{3} & \oeis{A235081} & \hyperref[combined:round6]{6} & \oeis{A381845} & \hyperref[combined:round8]{8}\\
\oeis{A096920} & \hyperref[combined:round6]{6} & \oeis{A209893} & \hyperref[combined:round3]{3} & \oeis{A235091} & \hyperref[combined:round7]{7} & \oeis{A382122} & \hyperref[combined:round8]{8}\\
\oeis{A097331} & \hyperref[combined:round5]{5} & \oeis{A209894} & \hyperref[combined:round3]{3} & \oeis{A235092} & \hyperref[combined:round7]{7} & \oeis{A382123} & \hyperref[combined:round10]{10}\\
\oeis{A097332} & \hyperref[combined:round4]{4} & \oeis{A209906} & \hyperref[combined:round5]{5} & \oeis{A235093} & \hyperref[combined:round7]{7} & \oeis{A382124} & \hyperref[combined:round10]{10}\\
\oeis{A097724} & \hyperref[combined:round2]{2} & \oeis{A209907} & \hyperref[combined:round5]{5} & \oeis{A235100} & \hyperref[combined:round8]{8} & \oeis{A382125} & \hyperref[combined:round10]{10}\\
\oeis{A098275} & \hyperref[combined:round2]{2} & \oeis{A209908} & \hyperref[combined:round5]{5} & \oeis{A235101} & \hyperref[combined:round8]{8} & \oeis{A383044} & \hyperref[combined:round8]{8}\\
\oeis{A098519} & \hyperref[combined:round8]{8} & \oeis{A209953} & \hyperref[combined:round5]{5} & \oeis{A235107} & \hyperref[combined:round8]{8} & \oeis{A383165} & \hyperref[combined:round5]{5}\\
\oeis{A098520} & \hyperref[combined:round9]{9} & \oeis{A209955} & \hyperref[combined:round5]{5} & \oeis{A235169} & \hyperref[combined:round9]{9} & \oeis{A383166} & \hyperref[combined:round5]{5}\\
\oeis{A099038} & \hyperref[combined:round1]{1} & \oeis{A209956} & \hyperref[combined:round5]{5} & \oeis{A235170} & \hyperref[combined:round9]{9} & \oeis{A383377} & \hyperref[combined:round8]{8}\\
\oeis{A099553} & \hyperref[combined:round10]{10} & \oeis{A209962} & \hyperref[combined:round3]{3} & \oeis{A235175} & \hyperref[combined:round9]{9} & \oeis{A383700} & \hyperref[combined:round7]{7}\\
\oeis{A100095} & \hyperref[combined:round4]{4} & \oeis{A209963} & \hyperref[combined:round3]{3} & \oeis{A235180} & \hyperref[combined:round8]{8} & \oeis{A383701} & \hyperref[combined:round7]{7}\\
\oeis{A100311} & \hyperref[combined:round1]{1} & \oeis{A209964} & \hyperref[combined:round3]{3} & \oeis{A235181} & \hyperref[combined:round8]{8} & \oeis{A383702} & \hyperref[combined:round7]{7}\\
\oeis{A100631} & \hyperref[combined:round2]{2} & \oeis{A209965} & \hyperref[combined:round3]{3} & \oeis{A235186} & \hyperref[combined:round8]{8} & \oeis{A384270} & \hyperref[combined:round9]{9}\\
\oeis{A101601} & \hyperref[combined:round8]{8} & \oeis{A210005} & \hyperref[combined:round4]{4} & \oeis{A235191} & \hyperref[combined:round8]{8} & \oeis{A385001} & \hyperref[combined:round7]{7}\\
\oeis{A102086} & \hyperref[combined:round2]{2} & \oeis{A210006} & \hyperref[combined:round4]{4} & \oeis{A235192} & \hyperref[combined:round8]{8} & \oeis{A385252} & \hyperref[combined:round6]{6}\\
\oeis{A102318} & \hyperref[combined:round4]{4} & \oeis{A210007} & \hyperref[combined:round4]{4} & \oeis{A235193} & \hyperref[combined:round8]{8} & \oeis{A385590} & \hyperref[combined:round6]{6}\\
\oeis{A102319} & \hyperref[combined:round5]{5} & \oeis{A210008} & \hyperref[combined:round4]{4} & \oeis{A235198} & \hyperref[combined:round8]{8} & \oeis{A386830} & \hyperref[combined:round7]{7}\\
\oeis{A102880} & \hyperref[combined:round6]{6} & \oeis{A210045} & \hyperref[combined:round3]{3} & \oeis{A235205} & \hyperref[combined:round9]{9} & \oeis{A387084} & \hyperref[combined:round5]{5}\\
\oeis{A103138} & \hyperref[combined:round8]{8} & \oeis{A210046} & \hyperref[combined:round3]{3} & \oeis{A235206} & \hyperref[combined:round9]{9} & \oeis{A388390} & \hyperref[combined:round7]{7}\\
\oeis{A103517} & \hyperref[combined:round10]{10} & \oeis{A210047} & \hyperref[combined:round3]{3} & \oeis{A235212} & \hyperref[combined:round9]{9} & \oeis{A389425} & \hyperref[combined:round3]{3}\\
\oeis{A104624} & \hyperref[combined:round5]{5} & \oeis{A210048} & \hyperref[combined:round3]{3} & \oeis{A235216} & \hyperref[combined:round7]{7} & \oeis{A389536} & \hyperref[combined:round3]{3}\\
\oeis{A105696} & \hyperref[combined:round1]{1} & \oeis{A210054} & \hyperref[combined:round4]{4} & \oeis{A235217} & \hyperref[combined:round7]{7} & \oeis{A389537} & \hyperref[combined:round6]{6}\\
\oeis{A105794} & \hyperref[combined:round2]{2} & \oeis{A210055} & \hyperref[combined:round4]{4} & \oeis{A235218} & \hyperref[combined:round7]{7} & \oeis{A392211} & \hyperref[combined:round1]{1}\\
\oeis{A105864} & \hyperref[combined:round6]{6} & \oeis{A210056} & \hyperref[combined:round4]{4} & \oeis{A235233} & \hyperref[combined:round9]{9} & \oeis{A392976} & \hyperref[combined:round1]{1}\\
\oeis{A106391} & \hyperref[combined:round7]{7} & \oeis{A210057} & \hyperref[combined:round4]{4} & \oeis{A235234} & \hyperref[combined:round9]{9} & \oeis{A393170} & \hyperref[combined:round6]{6}\\
\oeis{A106407} & \hyperref[combined:round10]{10} & \oeis{A210063} & \hyperref[combined:round5]{5} & \oeis{A235239} & \hyperref[combined:round9]{9} & \oeis{A393171} & \hyperref[combined:round6]{6}\\
\oeis{A106608} & \hyperref[combined:round9]{9} & \oeis{A210070} & \hyperref[combined:round5]{5} & \oeis{A235251} & \hyperref[combined:round8]{8} & \oeis{A393172} & \hyperref[combined:round3]{3}, \hyperref[combined:round6]{6}\\
\oeis{A107026} & \hyperref[combined:round7]{7} & \oeis{A210071} & \hyperref[combined:round5]{5} & \oeis{A235252} & \hyperref[combined:round8]{8} & \oeis{A393173} & \hyperref[combined:round6]{6}\\
\oeis{A107587} & \hyperref[combined:round5]{5} & \oeis{A210072} & \hyperref[combined:round5]{5} & \oeis{A235258} & \hyperref[combined:round8]{8} & \oeis{A393856} & \hyperref[combined:round6]{6}\\
\oeis{A107635} & \hyperref[combined:round5]{5} & \oeis{A210079} & \hyperref[combined:round4]{4} & \oeis{A235271} & \hyperref[combined:round6]{6} & \oeis{A393857} & \hyperref[combined:round6]{6}\\
\oeis{A108296} & \hyperref[combined:round1]{1} & \oeis{A210080} & \hyperref[combined:round4]{4} & \oeis{A235272} & \hyperref[combined:round6]{6} & \oeis{A393858} & \hyperref[combined:round6]{6}\\
\oeis{A108600} & \hyperref[combined:round5]{5} & \oeis{A210081} & \hyperref[combined:round4]{4} & \oeis{A235273} & \hyperref[combined:round6]{6} & \oeis{A393859} & \hyperref[combined:round6]{6}\\
\oeis{A108625} & \hyperref[combined:round7]{7} & \oeis{A210082} & \hyperref[combined:round5]{5} & \oeis{A235282} & \hyperref[combined:round7]{7} & \oeis{A394399} & \hyperref[combined:round3]{3}\\
\oeis{A108629} & \hyperref[combined:round2]{2} & \oeis{A210088} & \hyperref[combined:round3]{3} & \oeis{A235283} & \hyperref[combined:round7]{7} & \oeis{A394454} & \hyperref[combined:round5]{5}\\
\oeis{A108630} & \hyperref[combined:round7]{7} & \oeis{A210089} & \hyperref[combined:round3]{3} & \oeis{A235284} & \hyperref[combined:round7]{7} & \oeis{A394578} & \hyperref[combined:round1]{1}\\
\oeis{A110198} & \hyperref[combined:round7]{7} & \oeis{A210091} & \hyperref[combined:round3]{3} & \oeis{A235291} & \hyperref[combined:round7]{7} & \oeis{A394757} & \hyperref[combined:round3]{3}\\
\oeis{A110616} & \hyperref[combined:round2]{2} & \oeis{A210092} & \hyperref[combined:round4]{4} & \oeis{A235292} & \hyperref[combined:round7]{7} & \oeis{A394815} & \hyperref[combined:round3]{3}\\
\oeis{A110665} & \hyperref[combined:round9]{9} & \oeis{A210100} & \hyperref[combined:round4]{4} & \oeis{A235293} & \hyperref[combined:round7]{7} & \oeis{A395454} & \hyperref[combined:round2]{2}\\
\oeis{A110668} & \hyperref[combined:round9]{9} & \oeis{A210103} & \hyperref[combined:round4]{4} & \oeis{A235294} & \hyperref[combined:round7]{7} & \oeis{A395833} & \hyperref[combined:round7]{7}\\
\oeis{A111205} & \hyperref[combined:round4]{4} & \oeis{A210104} & \hyperref[combined:round7]{7} & \oeis{A235303} & \hyperref[combined:round6]{6} & \oeis{A396099} & \hyperref[combined:round4]{4}\\
\oeis{A111942} & \hyperref[combined:round1]{1} & \oeis{A210118} & \hyperref[combined:round4]{4} & \oeis{A235304} & \hyperref[combined:round6]{6} & \oeis{A396102} & \hyperref[combined:round8]{8}\\
\oeis{A111967} & \hyperref[combined:round1]{1} & \oeis{A210119} & \hyperref[combined:round4]{4} & \oeis{A235305} & \hyperref[combined:round6]{6} & \oeis{A396722} & \hyperref[combined:round6]{6}\\
\oeis{A112302} & \hyperref[combined:round2]{2} & \oeis{A210128} & \hyperref[combined:round5]{5} & \oeis{A235306} & \hyperref[combined:round7]{7} & \oeis{A396724} & \hyperref[combined:round1]{1}\\
\oeis{A112526} & \hyperref[combined:round6]{6} & \oeis{A210129} & \hyperref[combined:round4]{4} & \oeis{A235312} & \hyperref[combined:round6]{6} & \oeis{A396793} & \hyperref[combined:round6]{6}\\
\oeis{A112702} & \hyperref[combined:round8]{8} & \oeis{A210130} & \hyperref[combined:round4]{4} & \oeis{A235313} & \hyperref[combined:round6]{6} & \oeis{A396794} & \hyperref[combined:round6]{6}\\
\oeis{A113249} & \hyperref[combined:round1]{1} & \oeis{A210136} & \hyperref[combined:round4]{4} & \oeis{A235314} & \hyperref[combined:round7]{7} & \oeis{A396795} & \hyperref[combined:round6]{6}\\
\oeis{A113250} & \hyperref[combined:round1]{1} & \oeis{A210137} & \hyperref[combined:round4]{4} & \oeis{A239201} & \hyperref[combined:round5]{5} & \oeis{A396797} & \hyperref[combined:round6]{6}\\
\oeis{A113256} & \hyperref[combined:round1]{1} & \oeis{A210138} & \hyperref[combined:round4]{4} & \oeis{A239230} & \hyperref[combined:round2]{2} & \oeis{A396798} & \hyperref[combined:round6]{6}\\
\oeis{A113409} & \hyperref[combined:round8]{8} & \oeis{A210139} & \hyperref[combined:round4]{4} & \oeis{A240722} & \hyperref[combined:round8]{8} & \oeis{A396803} & \hyperref[combined:round3]{3}\\
\oeis{A114464} & \hyperref[combined:round4]{4} & \oeis{A210149} & \hyperref[combined:round5]{5} & \oeis{A242468} & \hyperref[combined:round10]{10} & \oeis{A396804} & \hyperref[combined:round5]{5}\\
\oeis{A114584} & \hyperref[combined:round6]{6} & \oeis{A210150} & \hyperref[combined:round5]{5} & \oeis{A242471} & \hyperref[combined:round10]{10} & \oeis{A396805} & \hyperref[combined:round6]{6}\\
\oeis{A114851} & \hyperref[combined:round8]{8} & \oeis{A210151} & \hyperref[combined:round5]{5} & \oeis{A242543} & \hyperref[combined:round10]{10} & \oeis{A396806} & \hyperref[combined:round6]{6}\\
\oeis{A116184} & \hyperref[combined:round6]{6} & \oeis{A210158} & \hyperref[combined:round4]{4} & \oeis{A242544} & \hyperref[combined:round10]{10} & \oeis{A396807} & \hyperref[combined:round6]{6}\\
\oeis{A116388} & \hyperref[combined:round5]{5} & \oeis{A210159} & \hyperref[combined:round4]{4} & \oeis{A242546} & \hyperref[combined:round10]{10} & \oeis{A396808} & \hyperref[combined:round3]{3}\\
\oeis{A116391} & \hyperref[combined:round4]{4} & \oeis{A210160} & \hyperref[combined:round4]{4} & \oeis{A242547} & \hyperref[combined:round10]{10} & \oeis{A396843} & \hyperref[combined:round3]{3}\\
\oeis{A117436} & \hyperref[combined:round1]{1} & \oeis{A210161} & \hyperref[combined:round4]{4} & \oeis{A242548} & \hyperref[combined:round10]{10} & \oeis{A396844} & \hyperref[combined:round3]{3}\\
\oeis{A118093} & \hyperref[combined:round6]{6} & \oeis{A210167} & \hyperref[combined:round4]{4} & \oeis{A242549} & \hyperref[combined:round10]{10} & \oeis{A396845} & \hyperref[combined:round3]{3}\\
\oeis{A119689} & \hyperref[combined:round8]{8} & \oeis{A210168} & \hyperref[combined:round4]{4} & \oeis{A242578} & \hyperref[combined:round10]{10} & \oeis{A396846} & \hyperref[combined:round7]{7}\\
\oeis{A119975} & \hyperref[combined:round7]{7} & \oeis{A210169} & \hyperref[combined:round4]{4} & \oeis{A242579} & \hyperref[combined:round10]{10} & \oeis{A397048} & \hyperref[combined:round2]{2}\\
\oeis{A121320} & \hyperref[combined:round6]{6} & \oeis{A210170} & \hyperref[combined:round7]{7} & \oeis{A242580} & \hyperref[combined:round10]{10} & \oeis{A397241} & \hyperref[combined:round3]{3}, \hyperref[combined:round5]{5}\\
\oeis{A121553} & \hyperref[combined:round1]{1} & \oeis{A210176} & \hyperref[combined:round6]{6} & \oeis{A242581} & \hyperref[combined:round10]{10} & \oeis{A397242} & \hyperref[combined:round5]{5}, \hyperref[combined:round7]{7}\\
\oeis{A121580} & \hyperref[combined:round1]{1} & \oeis{A210177} & \hyperref[combined:round6]{6} & \oeis{A242582} & \hyperref[combined:round10]{10} & \oeis{A397244} & \hyperref[combined:round7]{7}\\
\oeis{A121582} & \hyperref[combined:round1]{1} & \oeis{A210178} & \hyperref[combined:round6]{6} & \oeis{A242584} & \hyperref[combined:round10]{10} & \oeis{A397245} & \hyperref[combined:round7]{7}\\
\oeis{A121584} & \hyperref[combined:round1]{1} & \oeis{A210179} & \hyperref[combined:round7]{7} & \oeis{A243017} & \hyperref[combined:round10]{10} & \oeis{A397345} & \hyperref[combined:round7]{7}\\
\oeis{A121639} & \hyperref[combined:round1]{1} & \oeis{A210260} & \hyperref[combined:round4]{4} & \oeis{A243022} & \hyperref[combined:round6]{6} & \oeis{A397346} & \hyperref[combined:round7]{7}\\
\oeis{A121696} & \hyperref[combined:round1]{1} & \oeis{A210261} & \hyperref[combined:round4]{4} & \oeis{A243027} & \hyperref[combined:round10]{10} & \oeis{A397347} & \hyperref[combined:round7]{7}\\
\oeis{A121697} & \hyperref[combined:round1]{1} & \oeis{A210262} & \hyperref[combined:round4]{4} & \oeis{A243028} & \hyperref[combined:round10]{10} & \oeis{A397348} & \hyperref[combined:round7]{7}\\
\oeis{A121745} & \hyperref[combined:round1]{1} & \oeis{A210263} & \hyperref[combined:round5]{5} & \oeis{A243029} & \hyperref[combined:round10]{10} & \oeis{A397349} & \hyperref[combined:round7]{7}\\
\oeis{A121747} & \hyperref[combined:round1]{1} & \oeis{A210269} & \hyperref[combined:round5]{5} & \oeis{A243033} & \hyperref[combined:round10]{10} & \oeis{A397356} & \hyperref[combined:round3]{3}\\
\oeis{A121751} & \hyperref[combined:round1]{1} & \oeis{A210270} & \hyperref[combined:round5]{5} & \oeis{A243038} & \hyperref[combined:round10]{10} & \oeis{A397588} & \hyperref[combined:round9]{9}\\
\oeis{A121752} & \hyperref[combined:round1]{1} & \oeis{A210271} & \hyperref[combined:round5]{5} & \oeis{A243039} & \hyperref[combined:round10]{10} &  & \\
\oeis{A121754} & \hyperref[combined:round4]{4} & \oeis{A210272} & \hyperref[combined:round5]{5} & \oeis{A243040} & \hyperref[combined:round10]{10} &  & \\
\end{longtable}\endgroup
\end{document}
